% ############################################### PREAMBLE AND CONTROL PANEL #

% ########################################################### DOCUMENT CLASS #
\documentclass[  %   draft
             % , demo  % black boxes instead of figures
              , 8pt, twoside, openright]{extbook}

% ########################################################## AND CONTROL PANEL #

% true - hmarginration 1:1, urlcolor - darkgreen, version - for searching
\newif\ifscreen 	   \screentrue
% true - no apendix, no author of grammer, no backmatter ISBN, 
\newif\ifPDF               \PDFtrue
% true - rotate pages of morph verb tables
\newif\ifPDFrotate         \PDFrotatetrue
% true - adds white searchable value !
\newif\ifPDFsearchvalue    \PDFsearchvaluetrue
\newif\ifmakecovers        \makecoverstrue
\newif\ifmakefrontmatter   \makefrontmattertrue
\newif\ifinputintroduction \inputintroductiontrue
\newif\ifinputletters      \inputletterstrue
\newif\ifinputphon         \inputphontrue
\newif\ifinputmorpho       \inputmorphotrue
\newif\ifmakebackmatter    \makebackmattertrue
\newif\ifshowcrop           \showcroptrue
% corrections
\newif\ifC                  \Ctrue

\newif\ifPRINTHOUSE 	   \PRINTHOUSEfalse
\newif\ifPDFSEARCH  	   \PDFSEARCHfalse
\newif\ifPDFPRINT 	   \PDFPRINTfalse
\newif\ifPDFtest	   \PDFtesttrue

%versions
%PRINTHOUSE
\ifPRINTHOUSE
\makecoversfalse
\screenfalse
\PDFfalse
\PDFrotatefalse
\PDFsearchvaluefalse
\showcroptrue
\Ctrue
\fi

%PDFSEARCH
\ifPDFSEARCH
\screentrue
\PDFtrue
\PDFrotatetrue
\PDFsearchvaluetrue
\showcropfalse
\Ctrue
\fi

%PDFPRINT
\ifPDFPRINT
\screenfalse
\PDFtrue
\PDFrotatefalse
\PDFsearchvaluefalse
\showcropfalse
\Ctrue
\fi

%PDFtest
\ifPDFtest
\makecoversfalse
\makefrontmatterfalse
\inputintroductionfalse
\inputletterstrue
\inputphonfalse
\inputmorphofalse
\makebackmatterfalse
\screentrue
\PDFtrue
\PDFrotatefalse
\PDFsearchvaluefalse
\showcropfalse
\Cfalse
\fi

% 1. PDF (for printing)  
% a) searchable value, b) no apendix, c) version= PDF d) no ISBN and no ISBN bar
% e) automatic rotation of landscape pages, f) no author of Icelandic Grammer
% 2. PDF (for viewing offline and searching) + screen (horizontal ratio 1:1) 
% 3. PRINTHOUSE
% a) no searchable value, b) apendix, c) version ISBN d) ISBN bar
% e) no rotation of landscape pages, f)urllinkcolor black
% f) A4 paper dimentions, C5 layout, shows cropping marks

% ########################################################## PAGE GEOMETRIES #

\usepackage%[showframe]
           {geometry}
           

\ifshowcrop
% Basic geometry of document
\geometry
  { headsep    =   \baselineskip
  , textwidth  = 42\baselineskip
  , textheight = 60\baselineskip
  , hmarginratio = \ifscreen 1:1\else 2:3\fi
  , vmarginratio = 2:3
  , bindingoffset = 0cm
  , onecolumn
  , a4paper 
  , layout=c5paper,
  , layouthoffset=\dimexpr(\paperwidth-\csname Gm@layoutwidth\endcsname)/2\relax
   , layoutvoffset=\dimexpr(\paperheight-\csname Gm@layoutheight\endcsname)/2\relax
  , showcrop
  }
\else
\geometry
  { headsep    =   \baselineskip
  , textwidth  = 42\baselineskip
  , textheight = 60\baselineskip
  , hmarginratio = \ifscreen 1:1\else 2:3\fi
  , vmarginratio = 2:3
  , bindingoffset = 0cm
  , onecolumn,
  , c5paper
  }
\fi

\ifshowcrop
% Geometry of the dictionary section
\newcommand\dictionarygeometry{\newgeometry
  { textwidth  = 42\baselineskip
  , textheight = 60\baselineskip
  , hmarginratio = \ifscreen 1:1\else 2:3\fi
  , vmarginratio = 3:2
  , bindingoffset = 0cm
  , twocolumn
  , a4paper 
  , layout=c5paper
  , layouthoffset=\dimexpr(\paperwidth-\csname Gm@layoutwidth\endcsname)/2\relax
   , layoutvoffset=\dimexpr(\paperheight-\csname Gm@layoutheight\endcsname)/2\relax
  , showcrop
  }}
   % Geometry of the cover pages
\newcommand\covergeometry{\newgeometry
  { textwidth  = 42\baselineskip
  , textheight = 60\baselineskip
  , hmarginratio = 1:1
  , vmarginratio = 2:3
  , bindingoffset = 0cm
  , onecolumn
  , layout=c5paper
 , layouthoffset=\dimexpr(\paperwidth-\csname Gm@layoutwidth\endcsname)/2\relax
   , layoutvoffset=\dimexpr(\paperheight-\csname Gm@layoutheight\endcsname)/2\relax
  , showcrop
  }} 
\else
% Geometry of the dictionary section
\newcommand\dictionarygeometry{\newgeometry
  { textwidth  = 42\baselineskip
  , textheight = 60\baselineskip
  , hmarginratio = \ifscreen 1:1\else 2:3\fi
  , vmarginratio = 3:2
  , bindingoffset = 0cm
  , twocolumn,
  , c5paper
  }}
   % Geometry of the cover pages
\newcommand\covergeometry{\newgeometry
  { textwidth  = 42\baselineskip
  , textheight = 60\baselineskip
  , hmarginratio = 1:1
  , vmarginratio = 2:3
  , bindingoffset = 0cm
  , onecolumn
  , c5paper
  }} 
\fi

% ################################################################# LANGUAGES #

\usepackage[icelandic, latin, czech]{babel}
\usepackage{csquotes}


% ################################################################## ENCODING #

\usepackage[utf8]{inputenc}

% Smashing uppercase letters with diacritics
% ------------------ THIS IS EVIL AND SHOULD NEVER BE DONE -------------------
% Don't smash anything if we are on the first run to produce the index
\IfFileExists{figures.idx}{
  % Czech uppercase letters with diacritics:
  \DeclareUnicodeCharacter{00C1}{\strut\smash{\' A}} % Á
  \DeclareUnicodeCharacter{010C}{\strut\smash{\v C}} % Č
  \DeclareUnicodeCharacter{010E}{\strut\smash{\v D}} % Ď
  \DeclareUnicodeCharacter{00C9}{\strut\smash{\' E}} % É
  \DeclareUnicodeCharacter{011A}{\strut\smash{\v E}} % Ě
  \DeclareUnicodeCharacter{00CD}{\strut\smash{\' I}} % Í
  \DeclareUnicodeCharacter{0147}{\strut\smash{\v N}} % Ň
  \DeclareUnicodeCharacter{00D3}{\strut\smash{\' O}} % Ó
  \DeclareUnicodeCharacter{0158}{\strut\smash{\v R}} % Ř
  \DeclareUnicodeCharacter{0160}{\strut\smash{\v S}} % Š
  \DeclareUnicodeCharacter{0164}{\strut\smash{\v T}} % Ť
  \DeclareUnicodeCharacter{00DA}{\strut\smash{\' U}} % Ú
  \DeclareUnicodeCharacter{016E}{\strut\smash{\r U}} % Ů
  \DeclareUnicodeCharacter{00DD}{\strut\smash{\' Y}} % Ý
  \DeclareUnicodeCharacter{017D}{\strut\smash{\v Z}} % Ž
  % Icelandic uppercase letters with diacritics: ÁÉÍÓÚÝ (ÐÞÆ) and
  \DeclareUnicodeCharacter{00D6}{\strut\smash{\" O}} % Ö
}{}
% ------------------ THAT WAS EVIL AND SHOULD NEVER BE DONE ------------------

% Non-breakable space
\DeclareUnicodeCharacter{00A0}{~}

% Non-breakable hyphen
\newcommand*{\nobreakhyphen}{\mbox{-}}
\DeclareUnicodeCharacter{2011}{\nobreakhyphen}


% ############################################################ FONTS, SYMBOLS #

\usepackage[T1]{fontenc}

\usepackage{tgpagella}
\usepackage[scaled=0.88]{helvet}      % relative scale of the two fonts

\def\phvfamily{\fontfamily{phv}\selectfont}

%\usepackage{stmaryrd}
% Let's keep this minimal instead of using the package:
\DeclareSymbolFont{stmry}{U}{stmry}{m}{n}
%\SetSymbolFont{stmry}{bold}{U}{stmry}{b}{n}
\DeclareMathSymbol\shortrightarrow\mathrel{stmry}{"01}
\DeclareMathSymbol\shortuparrow\mathrel{stmry}{"02}

% Methinks these are not necessary
%\usepackage{dingbat}
%\usepackage{manfnt}
%\usepackage{latexsym}

\usepackage{pifont}




% #################################################################### LAYOUT #

% \normalsize should be {8pt}{9.6pt}
\def\HUGE {\fontsize{23.887872pt}{3\baselineskip}\selectfont}
\def\Huge {\fontsize{19.906560pt}{3\baselineskip}\selectfont}
\def\huge {\fontsize{16.588800pt}{3\baselineskip}\selectfont}
\def\LARGE{\fontsize{13.824000pt}{2\baselineskip}\selectfont}
\def\Large{\fontsize{11.520000pt}{2\baselineskip}\selectfont}
\def\large{\fontsize{ 9.600000pt}{2\baselineskip}\selectfont}

% Two columns layout ruler
\setlength\columnsep    {2\baselineskip}
\setlength\columnseprule{0.4pt}

% Necessary for baseline alignment
\topskip=\baselineskip
\raggedbottom
\setlength\parskip{0pt} % it's better to avoid glue

% Temporarily suppress warnings
\hbadness=5000
\vbadness=5000

\setlength\emergencystretch{17pt}

% Allow smaller emergencystretch in several cases
\newenvironment{xtolerant}[2]{%
  \par
  \ifx\empty#1\empty\else\tolerance=#1\relax\fi
  \ifx\empty#2\empty\else\emergencystretch=#2\relax\fi
}{%
  \par
}

% ######################################################### ASSORTED PACKAGES #

\usepackage{tikz}
\usetikzlibrary{calc}

\usepackage{tipa}                      % IPA phonetics

\usepackage{fourier-orns}              % ornaments
\usepackage{amsmath}                   % non-breakable dash
\usepackage{hyphenat}                  % no hyphen in abbreviations

\usepackage{hanging}
\usepackage{fix2col}
\usepackage{fixltx2e}                  % subscript
\usepackage{pagecolor}
\usepackage{afterpage}                 % Used in coloring trick for covers
\usepackage{pdfpages}

\usepackage{multirow}
\usepackage{multicol}
\usepackage{rotating}

\usepackage{etoolbox} % http://ctan.org/pkg/etoolbox
%\usepackage{newunicodechar} 		% new definition of middle dot
\usepackage[ISBN=978-80-260-2385-2,SC0]{ean13isbn}

\usepackage{pgffor}

\usepackage{newunicodechar}

\newunicodechar{·}{\chejnikcdot}
\makeatletter
\newcommand{\chejnikcdot}{%
  \ifnum\lastpenalty=10042
    \kern-0.06667em\relax % dots should always be kerned
  \fi
  \textperiodcentered
  \@ifnextchar|{\kern\chejnikkern\relax}{\penalty10042 }%
}

\newcommand{\closeupdotbar}{%
  \renewcommand\chejnikkern{-0.06667em}%
}
\newcommand\chejnikkern{0pt} % default

% Middle dot + | little bit closer, nicer
%\newunicodechar{·}{\advancedcdot}
%\makeatletter
%\newcommand*\advancedcdot{\textperiodcentered
%  \@ifnextchar|{\hspace{-0.03667em}}{\checkifcdotnext}}
%\newcommand*\checkifcdotnext[2]
%  {\def\tmpa{·}\def\tmpb{#1#2}\ifx\tmpa\tmpb\hspace{-0.06667em}\fi#1#2}
%\makeatother

% ############################################################ BASELINE GRID #

\usepackage{eso-pic}
\usepackage{tikzpagenodes}

% Command to draw a baseline grid
\newcommand\drawbaselinegrid{%
  \begin{tikzpicture}[overlay,remember picture]
    \draw [red!30!white, ultra thin, dashed]
      (0,0) grid [ ystep  = \baselineskip, xstep = \textwidth
                 , shift  = (current page text area.north west)
                 , yshift = -\dp\strutbox
                 ] ++(\textwidth,-\textheight);
    \draw [red!30!white, thin]
      (0,0) grid [ step = \baselineskip, xstep = \textwidth
                 , shift=(current page text area.north)
                 ] ++(0.5\textwidth,-\textheight)
      (0,0) grid [ step = \baselineskip, xstep = \textwidth
                 , shift=(current page text area.north)
                 ] ++(-0.5\textwidth,-\textheight);
    \draw[black!10!white]
      (current page text area.north west)
        rectangle (current page text area.south east);
  \end{tikzpicture}}

% Draw a baseline grid on every page
% \AddToShipoutPicture{\drawbaselinegrid}


% ################################################################## METADATA #

\title{%
  \textbf{Islandsko-český studijní slovník}%
  \thanks{Tato kniha byla vytvořena v \LaTeX{}u pod Ubuntu 14.04.%
    Poděkování patří všem autorům, kteří publikují pod svobodnými licencemi.}}
\author{Aleš Chejn, Ján Zaťko, Jón Gíslason}
\date{říjen 2011}


% ############################################################## BIBLIOGRAPHY #

\usepackage[style=authoryear]{biblatex}
\addbibresource{slovnik.bib}


% ##################################################################### INDEX #

\usepackage{imakeidx}

% Index of authors of photographs
\makeindex
  [ name = figures
  , title = {Seznam autorů fotografií a ilustrací}
  , intoc
  , columnseprule
  , columnsep = \columnsep
  , noautomatic                       % We handle this in make.sh
  ]

% icelandic rules start
\usepackage{filecontents}
\def\mygroup#1{{\bfseries#1}}
\begin{filecontents*}{style.xdy}
;; style.xdy
(markup-letter-group :open-head "~n\mygroup{" :close-head "}")
\end{filecontents*}
% icelandic rules end


% ############################################################### PAGE STYLES #

\usepackage{fancyhdr}

% -------------------------------------- BASIC PAGE STYLE
\fancypagestyle{plain}{
  \fancyhf\relax                       % Clear header/footer
  \renewcommand \headrulewidth {0.0pt} % No header rule
  \renewcommand \footrulewidth {0.0pt} % No footer rule
  \fancyfoot[C]{\thepage}              % Page number in footer, centred
}

% -------------------------------------- INDEX PAGE STYLE  
%\fancypagestyle{indexstyle}{
%  \fancyhf\relax                       % Clear header/footer
%  \renewcommand \headrulewidth {0.4pt} % Thin header rule
%  \renewcommand \footrulewidth {0.0pt} % No footer rule
%  \fancyhead[C]{\thepage}              % Page number in footer, centred
%}

% -------------------------------------- DICTIONARY PAGE STYLE
\fancypagestyle{myheadings}{
  \fancyhf\relax                       % Clear header/footer
  \renewcommand \headrulewidth {0.4pt} % Thin header rule
  \fancyhead[CO,CE]{\thepage}          % Page number in header, centred
  % NOTE: the page numbers will be printed when the dictionary is ready
  \fancyhead[LE,LO]{\phvfamily\bfseries\rightmark}
  \fancyhead[RE,RO]{\phvfamily\bfseries\leftmark}
}

\pagestyle{plain}  

% ##################################################################### LISTS #

\usepackage{expdlist}
\usepackage{enumitem}

\setlist{nosep, topsep=\baselineskip}
\setlist[itemize]{leftmargin=*,label=\scalebox{.6}{\textbullet}}


% #################################################################### FLOATS #

\usepackage{caption}
\captionsetup{labelformat=empty}

% Alter some LaTeX defaults for better treatment of figures:
% See p.105 of "TeX Unbound" for suggested values. 
% See pp. 199-200 of Lamport's "LaTeX" book for details.
% ------------------------------------- Parameters for ALL pages:
\renewcommand \topfraction    {0.9}   %   max fraction of floats at top
\renewcommand \bottomfraction {0.8}   %   max fraction of floats at bottom
% ------------------------------------- Parameters for TEXT (non float) pages:
\setcounter{topnumber}    {1}
\setcounter{bottomnumber} {1}
\setcounter{totalnumber}  {1}         %   2 may work better
\setcounter{dbltopnumber} {2}         %   for 2-column pages
\renewcommand \dbltopfraction   {0.9} %   fit big float above 2-col. text
\renewcommand \textfraction     {0.07}%   allow minimal text w. figs
% ------------------------------------- Parameters for FLOAT (non text) pages:
\renewcommand \floatpagefraction{0.7} %   require fuller float pages
% N.B.: floatpagefraction MUST be less than topfraction !!
\renewcommand{\dblfloatpagefraction}{0.7} % require fuller float pages
% remember to use [htp] or [htpb] for placement

\makeatletter
\setlength{\@fptop}{0pt}
\setlength{\@fpbot}{0pt plus 1fil}
\makeatother

% This kills spacing around floats
\setlength \intextsep        {0pt}
\setlength \floatsep         {0pt}
\setlength \textfloatsep     {0pt}
\setlength \dbltextfloatsep  {0pt}
\setlength \dblfloatsep      {0pt}

% This kills spacing around captions
\setlength \abovecaptionskip {0pt}
\setlength \belowcaptionskip {0pt}


% #################################################################### TABLES #

\usepackage{booktabs}
\usepackage{tabularx}
\usepackage{longtable}
\usepackage{ltxtable}
\usepackage{lscape}

\ifPDFrotate
\else
\makeatletter			       
\renewcommand\PLS@Rotate[1]{}          % this command prevents pdflscape 
\makeatother                           % to rotate pages
\fi

%\LetLtxMacro\oldtabularx\tabularx
%\LetLtxMacro\oldendtabularx\endtabularx
%\renewenvironment{tabularx}[1]
%  {\oldtabularx{\columnwidth}{#2}}
%  {\oldendtabularx}

\newcommand\LTXfw[1]{\LTXtable{\columnwidth}{#1}}

%%% \LetLtxMacro\oldtabularx\tabularx
%%% \LetLtxMacro\oldendtabularx\endtabularx
%%% \newenvironment{longtabularx}[2]
%%%   {\oldtabularx{#1}{#2}}
%%%   {\oldendtabularx}

% ######################################################### czech babel vs cmidrule
% ###################################### http://tex.stackexchange.com/a/112001/8380
\begingroup
  \makeatletter
  \catcode`\-=\active
  \AtBeginDocument{
  \def\@@@cmidrule[#1-#2]#3#4{\global\@cmidla#1\relax
    \global\advance\@cmidla\m@ne
    \ifnum\@cmidla>0\global\let\@gtempa\@cmidrulea\else
    \global\let\@gtempa\@cmidruleb\fi
    \global\@cmidlb#2\relax
    \global\advance\@cmidlb-\@cmidla
    \global\@thisrulewidth=#3
    \@setrulekerning{#4}
    \ifnum\@lastruleclass=\z@\vskip \aboverulesep\fi
    \ifnum0=`{\fi}\@gtempa
    \noalign{\ifnum0=`}\fi\futurenonspacelet\@tempa\@xcmidrule}
  }
\endgroup

% macro that supress cmidrule at page brake
% http://tex.stackexchange.com/questions/110548/using-longtable-with-cmidrule-page-breaking-issues
\makeatletter
\def\LT@output{%
  \ifnum\outputpenalty <-\@Mi
    \ifnum\outputpenalty > -\LT@end@pen
      \LT@err{floats and marginpars not allowed in a longtable}\@ehc
    \else
      \setbox\z@\vbox{\unvbox\@cclv}%
      \ifdim \ht\LT@lastfoot>\ht\LT@foot
        \dimen@\pagegoal
        \advance\dimen@-\ht\LT@lastfoot
        \ifdim\dimen@<\ht\z@
          \setbox\@cclv\vbox{\unvbox\z@
\copy\LT@foot\vss}%
          \@makecol
          \@outputpage
          \setbox\z@\vbox{\box\LT@head}%
        \fi
      \fi
      \global\@colroom\@colht
      \global\vsize\@colht
      \vbox
        {\unvbox\z@
\box\ifvoid\LT@lastfoot\LT@foot\else\LT@lastfoot\fi}%
    \fi
  \else
    \setbox\@cclv\vbox{\unvbox\@cclv

\setbox0\lastbox
\ifdim\ht0=0.29999pt \unskip % hope that was a cline we threw away
\else
\nointerlineskip
\box0     % put it back, whatever it was
\fi

\copy\LT@foot\vss}%
    \@makecol
    \@outputpage
      \global\vsize\@colroom
    \copy\LT@head\nobreak
  \fi}

\makeatother

% These values ensure that for top, mid and bottom rules
% space above + rule width + space below = baselineskip
\setlength \abovetopsep    {0.43\baselineskip}
\setlength \heavyrulewidth {0.10\baselineskip}
\setlength \belowbottomsep {0.43\baselineskip}
\setlength \aboverulesep   {0.47\baselineskip}
\setlength \lightrulewidth {0.06\baselineskip}
\setlength \belowrulesep   {0.47\baselineskip}

%\usepackage{letltxmacro}

%%% \LetLtxMacro\oldtabular\tabular
%%% \LetLtxMacro\oldendtabular\endtabular
%%% \renewenvironment{tabular}[2][t]
%%%   {\oldtabular[#1]{#2}}
%%%   {\oldendtabular}

%%% \LetLtxMacro\oldtable\table
%%% \LetLtxMacro\oldendtable\endtable
%%% \renewenvironment{table}[1][h]
%%%   {\leavevmode\oldtable[#1]\vspace{\dp\strutbox}}
%%%   {\vspace{-\dp\strutbox}\oldendtable}
%%%   % This aligns a [t] aligned tab-envo to baseline grid


%%% % declention and conjugation tables
%%% \usepackage{floatrow}
%%% \DeclareFloatFont{footnotesize}{\footnotesize}
%%% % "scriptsize" is defined by floatrow, "tiny" not
%%% \floatsetup[table]{font=footnotesize}

% ################################################################### FIGURES #

\usepackage{adjustbox} % loads also graphicx

\graphicspath{%
  {/home/chejnik/Dokumenty/hvalur.org/images/biolib/full/}%
  {/home/chejnik/Dokumenty/hvalur.org/images/uploaded_files/}}

\newlength\dicfigheight
\newlength\dicfigskip

\newcommand\dicFigure[3]{ % {filename}{caption}{index entry}
  \setlength\fboxsep{0pt}\setlength\fboxrule{0.5pt}
  \def\dicfigbox{\fbox{%
    \adjincludegraphics [ max height = 0.8\columnwidth
                        , max width  = 0.8\columnwidth ] {#1}}}
  \setlength\abovecaptionskip {\dp\strutbox}
  \setlength\belowcaptionskip {\ht\strutbox} 
  \setlength\dicfigheight {\heightof\dicfigbox+\fboxrule}
  \setlength\dicfigskip
    {\baselineskip*\numexpr1+\dicfigheight/\baselineskip\relax-\dicfigheight}
  \begin{figure}[ht]
    \centering\vspace*{\dicfigskip}\dicfigbox\caption{#2}\index[figures]{#3}
  \end{figure}
}


% #################################################################### COLORS #

\usepackage{color}

% http://www.colorschemer.com/online.html
% main color - grammar color
% #66023C
\definecolor {darkgreen}          {rgb} {0.40, 0.01, 0.24}

\definecolor {royalazure}         {rgb} {0.00, 0.22, 0.66}
\definecolor {brown}              {rgb} {0.40, 0.01, 0.24}

% COLORS FOR THUMB INDEXES
\definecolor {babyblueeyes}       {rgb} {0.63, 0.79, 0.95}
\definecolor {unitednationsblue}  {rgb} {0.36, 0.57, 0.90}
\definecolor {blue(ryb)}          {rgb} {0.01, 0.28, 1.00}
\definecolor {darkblue}           {rgb} {0.00, 0.00, 0.55}
\definecolor {screamingreen}      {rgb} {0.46, 1.00, 0.44}
\definecolor {limegreen}          {rgb} {0.20, 0.80, 0.20}
\definecolor {islamicgreen}       {rgb} {0.00, 0.56, 0.00}
\definecolor {upforestgreen}      {rgb} {0.00, 0.27, 0.13}
\definecolor {icterine}           {rgb} {0.99, 0.97, 0.37}
\definecolor {orange(colorwheel)} {rgb} {1.00, 0.50, 0.00}
\definecolor {orange-red}         {rgb} {1.00, 0.27, 0.00}
\definecolor {oucrimsonred}       {rgb} {0.60, 0.00, 0.00}
\definecolor {cottoncandy}        {rgb} {1.00, 0.74, 0.85}
\definecolor {orchid}             {rgb} {0.85, 0.44, 0.84}
\definecolor {vividcerise}        {rgb} {0.85, 0.11, 0.51}
\definecolor {patriarch}          {rgb} {0.50, 0.00, 0.50}

\definecolor {title}              {RGB} {  16,   13,   32}
\definecolor {golden}             {RGB} { 241,  184,   45}
\definecolor {freqstar}           {RGB} {  91,    3,   99}
\definecolor {newdev}             {RGB} {   3,   11,   99}

\definecolor {color1}             {RGB} { 182,   86,    0}
\definecolor {color2}             {RGB} { 143,    9,    6}
\definecolor {color3}             {RGB} {   3,   23,  118}
\definecolor {color4}             {RGB} {   0,   82,  168}
\definecolor {color5}             {RGB} {   0,   85,  142}
\definecolor {color6}             {RGB} {   0,  115,  162}
\definecolor {color7}             {RGB} {  34,  146,  186}
\definecolor {color8}             {RGB} {  40,  159,  153}
\definecolor {color9}             {RGB} {   0,  125,  111}
\definecolor {color10}            {RGB} {   4,  107,   60}
\definecolor {color11}            {RGB} {  71,  134,   81}
\definecolor {color12}            {RGB} { 109,  134,   42}
\definecolor {color13}            {RGB} { 205,  194,   18}
\definecolor {color14}            {RGB} { 204,  162,   24}

%\definecolor {color1}             {RGB} { 239,  222,  205}
%\definecolor {color2}             {RGB} { 251,  206,  177}
%\definecolor {color3}             {RGB} { 251,  185,  142}
%\definecolor {color4}             {RGB} { 248,  166,  112}
%\definecolor {color5}             {RGB} { 216,  231,  246}
%\definecolor {color6}             {RGB} { 186,  215,  244}
%\definecolor {color7}             {RGB} { 152,  197,  243}
%\definecolor {color8}             {RGB} { 118,  181,  245}
%\definecolor {color9}             {RGB} { 181,  249,  185}
%\definecolor {color10}            {RGB} { 160,  248,  167}
%\definecolor {color11}            {RGB} { 128,  247,  136}
%\definecolor {color12}            {RGB} {  87,  248,   98}
%\definecolor {color13}            {RGB} { 250,  241,  138}
%\definecolor {color14}            {RGB} { 250,  240,  111}
%\definecolor {color15}            {RGB} { 244,  239,   85}
%\definecolor {color16}            {RGB} { 245,  239,   57}

\definecolor {lightgrey}          {RGB} { 105,  105,  105}
\definecolor {darkgreen_real}     {rgb} {0.00, 0.50, 0.00}


% ############################################################# THUMB INDEXES #

% Thumb indexes' colors
%\newcommand\BoxColor{%
%  \ifcase\theletternum darkgreen!30\or babyblueeyes\or unitednationsblue\or blue(ryb)\or screamingreen\or limegreen\or islamicgreen\or upforestgreen\or icterine\or orange(colorwheel)\or orange-red%
%  \or oucrimsonred\or cottoncandy\or orchid\or vividcerise\or patriarch\or babyblueeyes\or unitednationsblue\or blue(ryb)\or screamingreen\or limegreen\or islamicgreen\or upforestgreen\or icterine\or orange(colorwheel)\or orange-red%
%  \or oucrimsonred\or cottoncandy\or orchid\or vividcerise\or patriarch\else darkgreen!30\fi}
  
% Thumb indexes' colors
%\newcommand\BoxColor{%
%  \ifcase\theletternum darkgreen!30\or color1\or color2\or color3\or color4\or color5\or color6\or color7\or color8\or color9\or color10%
%  \or color11\or color12\or color13\or color14\or color1\or color2\or color3\or color4\or color5\or color6\or color7\or color8\or color9%
%  \or color10\or color11\or color12\or color13\or color14\or color1\or color2\or color3\or color4\or color5\or color6\else color7\or color8\or color9\or color10\or color11\or color12\fi}

%\def\BoxColor#1{darkgreen!\the\numexpr103-#1-#1-#1\relax!blue}

\newcommand\BoxColor[1]{%
\ifcase#1 darkgreen!30\or color13\or color14\or color1\or color2\or color3\or color4\or color5\or color6\or color7\or color8\or color9\or color10%
\or color11\or color12\or color13\or color14\or color1\or color2\or color3\or color4\or color5\or color6\or color7\or color8\or color9%
\or color10\or color11\or color12\or color13\or color14\or color1\or color2\or color3\or color4\or color5\or color6\else color7\or color8\or color9\or color10\or color11\or color12\fi}

\ifshowcrop
%%%%%%%%%%%%%%%%%%%%%%%%%%%%%%%%%%%%%%%%%%%%%%%%%%%%%%%%%%%%%%%%%%%%%%%%%%%%%%%%

\newcounter{letternum}
\newcounter{lettersum}                 \setcounter{lettersum}{34}
% these margins separate top thumb from top of page and last thumb from
% bottom of page. BEWARE, it's from top of the (cut out) page, not the folio!
\newlength{\thumbtopmargin}            \setlength{\thumbtopmargin}{1cm}
\newlength{\thumbbottommargin}         \setlength{\thumbbottommargin}{1cm}
\newlength{\thumbwidth}                \setlength{\thumbwidth}{2.20cm}
\newlength{\thumbheight}
\pgfmathsetlength{\thumbheight}%
  {(\csname Gm@layoutheight\endcsname-\thumbtopmargin-\thumbbottommargin)/\value{lettersum}}
% inner xsep of letter from box margin
\newlength{\thumbsep}                  \setlength{\thumbsep}{.30cm}
% overlap over cutting line
\newlength{\thumboverlap}              \setlength{\thumboverlap}{3cm}

\tikzset{
  thumb/.style={
    text=white,
    outer sep=0pt,
    inner xsep=\thumbsep,
    minimum height=\thumbheight,
    text width=\thumbwidth-2\thumbsep,
    font=\sffamily\bfseries,},
  thumb east/.style={thumb,
    xshift=\dimexpr(+\thumboverlap-\paperwidth+\csname Gm@layoutwidth\endcsname)/2\relax,
    align=left,anchor=north east,},
  thumb west/.style={thumb,
    xshift=\dimexpr(-\thumboverlap+\paperwidth-\csname Gm@layoutwidth\endcsname)/2\relax,
    align=right,anchor=north west,},
}

\def\thumbnew{}
\def\thumbold{}
\usepackage{everypage}
\AddEverypageHook{\if\relax\thumbnew\relax\xdef\thumbnew{\thumbold}\fi}

\def\ethumbs#1,#2\relax%
  {\if\relax#1\relax\else\eventhumb{#1}\fi%
   \if\relax#2\relax\else\ethumbs#2\relax\fi%
   \gdef\thumbnew{}%
   \gdef\thumbold{#1,}}

\def\othumbs#1,#2\relax%
  {\if\relax#1\relax\else\oddthumb{#1}\fi%
   \if\relax#2\relax\else\othumbs#2\relax\fi%
   \gdef\thumbold{#1,}%
   \gdef\thumbnew{}}

\newcommand{\drawthumb}[2]{%
  \begin{tikzpicture}[remember picture, overlay]
    \node [thumb #2, fill=\BoxColor{#1}]
       at ($(current page.north #2)-%
            (0,\csname Gm@layoutvoffset\endcsname+\thumbtopmargin-\thumbheight+%
               #1*\thumbheight)$) {\csname Let#1\endcsname};
   \end{tikzpicture}}



%%%%%%%%%%%%%%%%%%%%%%%%%%%%%%%%%%%%%%%%%%%%%%%%%%%%%%%%%%%%%%%%%%%%%%%%%%%%%%%%
\else
% THUMB INDEXES
% new counter to hold the current number of the letter to determine the vertical position
\newcounter{letternum}
% newcounter for the sum of all letters to get the right height of a box
\newcounter{lettersum}
\setcounter{lettersum}{34}
% some margin settings
\newlength{\thumbtopmargin}
\setlength{\thumbtopmargin}{\ifshowcrop 3cm\else 1.5cm\fi} %1cm
\newlength{\thumbbottommargin}
\setlength{\thumbbottommargin}{\ifshowcrop 6cm\else 2cm\fi} %2.5cm
% calculate the box height by dividing the page height
\newlength{\thumbheight}
\pgfmathsetlength{\thumbheight}{%
(\paperheight-\thumbtopmargin-\thumbbottommargin)%
/%
\value{lettersum}
}
% box width
\newlength{\thumbwidth}
\setlength{\thumbwidth}{\ifshowcrop 2cm\else 0.5cm\fi} %0.5cm
% style the boxes
\tikzset{
thumb/.style={
   text=white,
   minimum height=\thumbheight,
   text width=\thumbwidth,
   outer sep=0pt,
   font=\sffamily\bfseries,
 }
 }

\def\thumbnew{}
\def\thumbold{}
\usepackage{everypage}
\AddEverypageHook{\if\relax\thumbnew\relax\xdef\thumbnew{\thumbold}\fi}

\def\ethumbs#1,#2\relax{\if\relax#1\relax\else\eventhumb{#1}\fi%
                        \if\relax#2\relax\else\ethumbs#2\relax\fi%
                        \gdef\thumbnew{}%
                        \gdef\thumbold{#1,}%
}

\def\othumbs#1,#2\relax{\if\relax#1\relax\else\oddthumb{#1}\fi%
                        \if\relax#2\relax\else\othumbs#2\relax\fi%
                        \gdef\thumbold{#1,}%
                        \gdef\thumbnew{}%
}

\newcommand{\drawthumb}[2]{%
  % see pgfmanual.pdf for more information about this part
  \begin{tikzpicture}[remember picture, overlay]
    \node [thumb, fill=\BoxColor{#1}, text centered, anchor=north #2]
       at ($(current page.north #2)-%
%            (0,\thumbtopmargin+\value{letternum}*\thumbheight)%
            (0,\thumbtopmargin+#1*\thumbheight)$) {\csname Let#1\endcsname};
   \end{tikzpicture}}
\fi

\newcommand{\oddthumb} [1]{\drawthumb{#1}{west}}
\newcommand{\eventhumb}[1]{\drawthumb{#1}{east}}

\newcommand{\lettergroup}[1]%
  {\refstepcounter{letternum}%
   \expandafter\gdef\csname Let\theletternum\endcsname{#1}%
   \xdef\thumbnew{\theletternum,\thumbnew}%
   \fancyhead[LO]{\phvfamily\bfseries\rightmark%
     \expandafter\ethumbs\thumbnew\relax\relax\relax}%
   \fancyhead[RE]{\phvfamily\bfseries\leftmark%
     \expandafter\othumbs\thumbnew\relax\relax\relax}}
     
% ############################################################## LOCALIZATION #

% renames the index name
%\addto\captionsczech{%
%  \renewcommand\indexname{{Seznam autorů fotografií a ilustrací}}}

% renames the contents name
\addto\captionsczech{%
  \renewcommand\contentsname{Obsah}}

\renewcommand \contentsname {Obsah}
\renewcommand \chaptername  {Kapitola}


% ################################################################ SECTIONING #
\usepackage[explicit]{titlesec}

\setlength{\beforetitleunit}{\baselineskip}
\setlength{\aftertitleunit} {\baselineskip}

\titlespacing*{\chapter}   {0em}{*0}{*3} % Remembere there's a topskip too
\titlespacing*{\section}   {0em}{*1}{*2}
\titlespacing*{\subsection}{0em}{*1}{*2}

\newcommand\ghostdrop[2]{%
  \raisebox{-#1\baselineskip}[0pt][0pt]{#2}}

\titleformat {\chapter} {} {} {0pt} % {#1}
  {\ghostdrop{1.0}{\parbox[t]{\columnwidth}
     {\LARGE\bfseries\centering\MakeUppercase{#1}}}}

\def\longchapterskip{\vspace{2\baselineskip}}
\def\longsectionskip{\vspace{2\baselineskip}}

\titleformat {\section} {\bfseries} {} {0pt}
  {\ghostdrop{1.0}{\parbox[t]{\columnwidth}
     {\Large\bfseries\centering#1}}}

\titleformat {\subsection} {\bfseries} {} {0pt}
  {\ghostdrop{1.0}{\hspace{1em}\parbox[t]{\columnwidth-1em}
     {\large\bfseries#1}}}


\setlength\multicolsep{0pt}
%\BeforeBeginEnvironment{multicols}{\vspace{-\dp\strutbox}}


% ######################################################## TABLE OF CONTENTS #
\usepackage{titletoc}

\titlecontents {chapter} [0pt] {\addvspace\baselineskip\bfseries}
  {} {} {\titlerule*[2\baselineskip]{}\contentspage}

\titlecontents {section} [\baselineskip] {} {} {}
  {\titlerule*[1\baselineskip]{.}\contentspage}


% #################################################### ASSORTED CUSTOM MACROS #

\def\textIS#1{\foreignlanguage{icelandic}{#1}}
\def\textCS#1{\foreignlanguage{czech}{#1}}
\def\textLA#1{\foreignlanguage{latin}{#1}}

\usepackage{accsupp}

\makeatletter                                                                                                          
  \def\ACCSUPP@bdc{\special{pdf:code \ACCSUPP@span BDC}}                                                               
  \def\ACCSUPP@emc{\special{pdf:code EMC}}                                                                             
\makeatother

\usepackage{pgffor}

\newcommand\blspace[1][1]{\vspace{#1\baselineskip}}

\usepackage{xparse}


\newcommand \dicsymFrequent  {\raisebox{-.2ex}{\color{darkgreen}\ding{167}}}
\newcommand \dicsymSee       {$\shortrightarrow$}
\newcommand \dicsymCompare   {$\shortuparrow$}
\newcommand \dicsymIdiom     {$\diamondsuit$}
\newcommand \dicsymExampleIS {$\triangleright$}
\newcommand \dicsymExampleCS {\guilsinglright}
\newcommand \dicsymProverb   {{\color{newdev}\ding{96}}}
\newcommand \dicsymPhoto
  {\includegraphics[keepaspectratio,height=0.65em]{photo_icon}}

%\NewDocumentCommand\dicTerm{mo}
 % {{\phvfamily\bfseries\textIS{#1}}\IfValueT{#2}{, \phvfamily\bfseries\textIS{#2}}}

\NewDocumentCommand\dicEntry{o}
  {\par\hangpara{0.25\baselineskip}{1}%
   \IfValueT{#1}{\markboth{#1}{#1}}%
    \ifPDFsearchvalue
    \settowidth{\eyja}{#1}%
  \makebox[\eyja]{\color{white}#1}\hspace{-\eyja}%
  \fi
   \ignorespaces}
   
   \NewDocumentCommand\dicEntryGuide{o}
  {\par\hangpara{0.29\baselineskip}{1}%
   \IfValueT{#1}{\markboth{#1}{#1}}%
    \ifPDFsearchvalue
    \settowidth{\eyja}{#1}%
  \makebox[\eyja]{\color{white}#1}\hspace{-\eyja}%
  \fi
   \ignorespaces}
   
\newcommand\dicSubEntry
  {\\\hspace*{-\hangindent}}

\NewDocumentCommand\textAlt{mm}
  {\BeginAccSupp{method=pdfstringdef,ActualText={#2}}#1\EndAccSupp{}}

\newcounter{terms}

\NewDocumentCommand\dicStyleIntTerm{m}{\textIS{{\phvfamily\bfseries#1}}}
\NewDocumentCommand\dicStyleExtTerm{m}{\textIS{\textit{#1}}}

   \NewDocumentCommand\dicTerm{mo}
  {{\closeupdotbar \phvfamily\bfseries\textIS{#1}}\IfValueT{#2}{, \phvfamily\bfseries\textIS{#2}}}

  \newcommand\dicTermList[1]
  {\foreach \t [count=\n] in {#1}{\ifnum\n=1\else, \fi\dicTerm{\t}}}
  
\def\dicIPA#1{\textipa{[#1]}}

\NewDocumentCommand\dicPos{smo}
  {{\color{darkgreen}\textbf{\small#2}%
   \IfValueT{#3}{\nobreak\textsubscript{#3}}}}

%\NewDocumentCommand\dicFlx{smo}
 % {{\color{darkgreen}\footnotesize#1\IfValueT{#2}{\textsubscript{#2}}}}

  \NewDocumentCommand\dicFlx{smo}
  {{\color{darkgreen}\footnotesize#2%
   \IfValueT{#3}{\nobreak\textsubscript{#3}}}}
  
\newcommand\dicExampleCS[1]{\dicsymExampleCS~\textit{#1}}
\newcommand\dicExampleIS[1]{\dicsymExampleIS~\textIS{\textit{#1}}}

\NewDocumentCommand\dicLink{sm}
  {\IfBooleanF#1{\dicsymSee\nobreak\hspace{.08667em plus .03333em}}\nobreak\hspace{.03667em}\dicTerm{#2}}
\NewDocumentCommand\dicSynonym{sm}
  {(\IfBooleanF#1{\dicsymSee\nobreak\hspace{.08667em plus .03333em}}\nobreak\hspace{.03667em}\textIS{\textit{#2}})}
\NewDocumentCommand\dicAntonym{sm}
  {(\IfBooleanF#1{\dicsymCompare\nobreak\hspace{.08667em plus .03333em}}\nobreak\hspace{.03667em}\textIS{\textit{#2}})}
\NewDocumentCommand\dicCompare{sm}
  {(\IfBooleanF#1{\dicsymCompare\nobreak\hspace{.08667em plus .03333em}}\nobreak\hspace{.03667em}\textIS{\textit{#2}})}


\newcommand\dicPhraseIS[1]{\textIS{\textbf{#1}}}

\NewDocumentCommand\dicIdiom{mo}
  {\dicSubEntry{\color{newdev}\dicStyleIntTerm{#1} %
     \IfValueTF{#2}{+ \dicStyleIntTerm{#2} \markboth{#1 + #2}{#1 + #2}}{\markboth{#1}{#1}}}\dicsymIdiom}
     
\newcommand\dicProverb
  {\dicSubEntry\dicsymProverb}



\newcommand\dicLangCat[1]{{\footnotesize#1}}
\newcommand\dicFieldCat[1]{{\footnotesize#1}}


\newcommand\dicDirectTranslationCS[1]{#1}
\newcommand\dicIndirectTranslationCS[1]{{\footnotesize#1}}

\usepackage{placeins}

\newcommand*{\dicLetter}[2]{%
  \FloatBarrier\newpage\phantomsection\pdfbookmark[1]{#1}{#2}%
  \noindent\parbox[b][9\baselineskip][c]{\columnwidth}
    {\centering\HUGE\strut\smash{\MakeUppercase{#1}~\MakeLowercase{#1}}}%
      %\par\xdef\tpd{\the\prevdepth}
      %\par\prevdepth\tpd%
  \expandafter\lettergroup{\MakeLowercase{#1}}}

  
%tabulky
\newcommand{\specialcell}[2][l]{%
\begin{tabular}[#1]{@{}l@{}}#2\end{tabular}}


%\newcommand{\devisionguide}[1]{\hspace*{0.3cm}{{{{\foreignlanguage{icelandic}{\color{newdev}{\phvfamily{\textbf{#1}}}}}}}}}
%\newcommand{\devisionguideindent}[1]{\hspace*{-0.1cm}{{{{\foreignlanguage{icelandic}{\color{newdev}{\phvfamily{\textbf{#1}}}}}}}}}

%\newcommand\n{$n$\nobreakdash-\hspace{0pt}} %nonbreakable dash in grammar endings
%\def \nobreakseq {\nobreak \hskip 0pt \hbox} %nolinebreak in case like (-s, -)

% rule line
\newcommand{\HRule}{\rule{\linewidth}{0.1mm}}

% width of pdf found box
\newlength\eyja

%\newcommand\strutsmash[1]{\strut\smash{#1}}


% Dictionary symbols


% Command synonyms
%\let\freqstar\dicsymFrequent
%\let\photoicon\dicsymPhoto

\newcommand{\vstretch}[1]{\vspace*{\stretch{#1}}}

\newcommand*{\addthin}{\hskip0.16667em\relax}
\newcommand*{\addthinS}{\hskip0.06667em\relax}
\newcommand*{\addthinM}{\hskip0.10007em\relax}
\newcommand*{\addthinSS}{\hskip0.00007em\relax}
% #################################################################### COVERS #

\def\dictnameCZ{islandsko-český studijní slovník}
\def\dictnameIS{íslensk-tékknesk stúdentaorðabók}

\newcommand*{\frontcoverimages}{%
  ds_image_blaberjalyng_0_2.jpg,
  ds_image_lundi_0_1.jpg,
  ds_image_hestur_0_1.jpg,
  ds_image_hrafnafifa_0_2.jpg,
  ds_image_baldursbra_0_1.jpg,
  20948.jpg,
  ds_image_hvalur_0_1.jpg,
  ds_image_blodberg_0_1.jpg}

\newcommand*{\backcoverimages}{%
  ds_image_spoi_0_1.jpg,
  ds_image_mosalyng_0_1.jpg,
  ds_image_blaklukkulyng_0_1.jpg,
  20787.jpg,
  ds_image_ufsi_0_2.jpg,
  %ds_image_gullbra_0_1.jpg,
  ds_image_vatnsberi_0_1.jpg,
  ds_image_islandsfifill_0_2.jpg,
  21137.jpg}

\newcommand\makecoverwith[1]{
 \pagecolor{title}\afterpage\nopagecolor
 \begin{center}
 \vspace*{0.3cm}%{\fill}
 {\Huge\scshape\color{white}%
   \dictnameCZ\\\dictnameIS}
 \vspace*{1em}                        % I think this is necessary
 \begin{figure}[ht]\centering%
   \setlength\fboxsep{0pt}\setlength\fboxrule{0.5pt}%
   \foreach \file in #1{
     \fbox{\includegraphics[width=5.5cm]{\file}}}
 \end{figure}
 \vspace*{\fill}
 \end{center}
}


% ###################################################### HYPERREF & PDF INFO #

% unicode neccessary so that national characters in hypersetup appear ok
\usepackage[pdftex, unicode, hyperfootnotes=false]{hyperref}

% hyperlinks in black
\makeatletter
\let\Hy@linktoc\Hy@linktoc@none
\makeatother

\hypersetup
  { pdftitle={Islandsko-český studijní slovník}
  , pdfauthor={Aleš Chejn, Jón Gíslason, Ján Zaťko}
  , pdfsubject={Islandsko-český studijní slovník}
  , pdfkeywords={Islandsko-český studijní slovník%
               , islandština, čeština, slovník}
  , bookmarks=true
  , colorlinks=true
  , citecolor=black
  , urlcolor=\ifscreen darkgreen\else black\fi
  , linkcolor=black}


% ################################################################# MICROTYPE #

\usepackage[final]{microtype}

% ############################################################### HYPHENATION #

% this allows to include brackets inside hyphenation rules f.e (pro)-dis-ku-to-vat
\lccode`\(`\(
\lccode`\)`\)
\lccode`\[`\[
\lccode`\]`\]

\begin{hyphenrules}{icelandic}
\hyphenation{
  frum-semja
  burm-neskur
  ó-breytan-le-gur
  pla-tó-nskt
  skatta-fram-tal
  skrúf-gang-ur
  strand-rauð-viður
  prent-sver-ta
  grun-nur
  uppi-skrop-pa
  ger-vi-mennska
  sö-lu-stað-ur
  sý-na-gó-ga
  mark-mið
  trau-st-ur
  ís-kyggi-leg-ur
  brjó-sta-höld
  kné-krjú-pa
  nel-li-ka
  kaffi-kvörn
  margs-konar
  skrúf-gang-ur
  uppi-skrop-pa
  veð-rað-ur
  vest-ur-átt
  mis-mun-andi
  Bar-ba-do-si
  tung-u-mál
  au-ga-brún
  há-sker-pu-sjón-varp
  rúnn-stykki
  ger-vi-tungl
  yfir-borðs-menn-ska
  gæ-fu-mun-ur
  hlí-fi-skil-di
  á-vís-un
  ó-um-flýj-an-leg-ur
  san-mar-ín-sk-ur
  gæ-ða-stað-all
}
\end{hyphenrules}
\begin{hyphenrules}{czech}
\hyphenation{
  teo-rie
  mi-liardy
  (z)-od-po-věd-nost
  tan-tié-ma
  alian-ce
  pe-dia-tr(ička)
  pa-cienty
  zá-chra-nář(ka)
  dia-man-to-vá
  poe-zie
  ma-te-riál-ní
  egoi-s-ta
  egoist-ka
  nai-vi-ta
  ne-upřím-nost
  pě-ti-úhel-ník
  še-s-ti-úhel-ník
  za-šmo-dr-chat
  ar-cheo-lo-gic-ký
  ne-omlu-vi-tel-ný
  (vy)-pro-du-ko-vat
  dia-gno-s-ti-ko-vat
  teo-lož-ka
  ateist-ka
  oceá-nu
  im-pe-ria-li-s-ta
  (vy)-zdvih-nout
  ne-úměr-ný
  kon-ti-nuál-ní
  neu-t-rum
  ne-opráv-ně-ně
  vig-vam
  pro-zaic-ký
  geo-lož-ka
  seiz-mic-ká
  teo-rém
  olean-dr
  poe-ta
  (vy)-fo-to-gra-fo-vat
  (vy)-ven-ti-lo-vat
  pneu-mo-nie
  Lu-cem-bur-čan-(ka)
  ma-te-riály
  am-bi-cióz-ní
  Myan-ma(r)
  Myan-ma-řan-(ka)
  mol-dav-šti-na
  si-tua-cí
  ne-ukáz-ně-ný
  he-te-ro-trof-ní
  ne-omlu-vi-tel-ná
  ne-upřím-ný
  dvou-bliz-ný
  ne-úmy-sl-ný
  ne-ustá-lý-mi
  ne-ohle-du-pl-ný
  ne-osob-ní
  ne-oso-le-ný
  ne-uvě-ři-tel-ný
  ne-ohro-že-ný
  ne-oby-čej-ný
  ne-útul-ný
  ne-únav-ný
  ne-ože-nit
  po-ukáz-ka
  za-úto-čit
  aero-dy-na-mic-ký
  hru-bián-ství
  (pro)-di-s-ku-to-vat
  (vy)-ča-lou-nit
  (na)-oč-ko-vá-ní
  Bhú-tán-ka
  (za)-šně-ro-vat
  soud-ní
  re-vo-lu-cio-nář(ka)
  á-zer-báj-d-žán
  čtyř-hran-ný
  ze-m-dle-lý
  klou-zač-ka
  Vest-mann-ské
  tří-stran-ný
  (za)-bu-rá-ce-ní
  os-tud-ný
  pro-du-cent-(ka)
  chvá-s-tal-(ka)
  apar-theid
  asis-ten-ce
  asis-tent-(ka)
  (u)-či-nit
  (u)-za-mknout
  past-vi-na
  (za)-zá-lo-ho-vat
  Af-ri-čan-(ka)
  (za)-an-ga-žo-vá-ní
  se-rióz-ní
  opo-nent-(ka)
  pro-tiv-ník
  zvlášt-ní
  (po)-div-ný
  ho-mo-se-xuál-ní
  ko-alič-ní
  si-tua-ce
  si-tua-cí
  egoi-s-mus
  spe-ciál-ní
  spe-cia-li-zo-va-ný
  spe-cia-li-zo-va-né
  sobě-stač-ný
  sa-mo-zvo-le-ný
  sa-mo-čin-ně
 chvá-s-ta-vost
 klient-(ka)
 ku-rióz-ní
 (na)-šrou-bo-vat
 (za)-šrou-bo-vat
 krea-ti-vi-ta
 pneu-ma-ti-ka
 toa-le-ta
 so-cia-list-ka
 (vy)-dez-in-fi-ko-va-ný
 vše-uměl
 vše-uměl-(kyně)
 (pře)-škrt-nout
 (vy)-svlék-nout
 vy-učo-va-cí
 vy-učo-va-cích
 (vy)-chla-s-tat
 (vy)-te-to-vat
 pa-cient-(ka)
 ne-opod-s-tat-ně-ný
 de-se-ti-úhel-ník
 pe-rio-di-kum
 (za)-mum-lá-ní
 ne-smě-lá
 ne-dů-vě-ři-vá
 tee-na-ger-(ka)
 (z)-de-ge-ne-ro-vá-ní
 mo-ci-chti-vý
 bez-úhon-nost
 zne-uctít
 vo-dě-odol-ný
 we-bo-vý
 web-de-sign
 in-že-nýr-(ka)
 zá-pa-do-is-land-ský
 klient-(ka)
 za-uja-tý
 vy-uží-va-né
 vy-uží-vat
 (za)-re-gi-s-t-ro-vat
 (z)-kom-pri-mo-va-ný
 na-cio-na-li-s-mus
 ma-k-ro-eko-no-mic-ká
 bez-bo-lest-ný
 viol-ka
 troj-úhel-ník
 (za)-bu-rá-ce-ní
 ne-ústup-nost
 (vy)-cvi-če-ný
 zprá-va
 ge-ria-tr-(ička)
 osud-ný
 ná-pra-vy
 vstup-ní
 obec-ně
 re-pub-li-ka
 (po)-div-ný
 pa-slun-ce
 vliv-ný
 ba-ham-ský
 pe-dia-tr-(ička)
 ha-sič-(ka)
 vnuč-ka
 bro-s-kvoň
 lo-di-vod-(ka)
 mi-s-tr-(yně)
 in-že-nýr-(ka)
 dr-ži-tel-(ka)
 to-pič-(ka)
 pro-mi-nent-ní
 takt-ní
 frag-men-tár-ní
 po-zo-ru-hod-ný
 úspě-š-ný
 trans-fu-ze
 bur-žoa-zie
 pře-stro-je-ní
 oh-ni-vzdor-ný
 dis-tan-ční
 de-frau-da-ce
 fi-s-kál-ní
 zběh-lý
 kvap-ný
 střed-ní
 sa-mo-hlá-s-ka
 pře-vo-di-tel-ný
 zban-k-ro-to-va-ný
 zá-clo-na
 zná-mý
 zá-klad-ní
 na-uka
 ro-š-ťa-č-ka
 lí-bán-ky
 sou-du
 zkou-š-ka
 mlý-nek
 prá-šiv-ka
 kul-ma
 la-hod-ný
 zá-meč-ni-ce
 ne-vý-znam-ný
 ovzdu-ší
 mno-ho-úhel-ní-ko-vý
 žra-vý
 zcest-ný
 esen-ce
 (při)-nu-ce-ný
 skr-b-lic-ký
 skle-s-lý
 bez-bož-ný
 ne-opráv-ně-ný
 ošu-mě-lý
 pru-der-ní
 chvá-s-ta-vý
 sou-drž-ný
 van-dal-(ka)
 sa-bo-tér-(ka)
 (po)-div-ně
 mo-k-řad
 plyt-ký
 na-spo-du
 izo-la-ce
 le-toun
 za-han-be-ný
 ří-ze-ní
 kom-plex-ní
 pr-kén-ko
 ba-vič-(ka)
 plyt-ký
 ne-mo-tor-ně
 šéf-(ka)
 přes-čas
 zji-s-ti-tel-ný
 ta-ko-vý
 kňou-ral-(ka)
 se-me-s-tr
 po-pra-s-ka-ný
 han-ba
 vzpří-me-ný
 chys-ta-ný
 lik-vi-da-ce
 or-ga-ni-za-ce
 ri-s-kant-ní
 (za)-pís-kat
 (za)-hvíz-dat
 vrou-cí
  }
\end{hyphenrules}


% ############################################################# CZECH TYPOGRAPHY #
% https://math.feld.cvut.cz/tkadlec/ftp/text/textypo.pdf



% ################################################################# DOCUMENT #

\begin{document}
\pagenumbering{gobble}

% -------------------------------------------------------------- Front cover -

\ifmakecovers
\covergeometry
  \makecoverwith\frontcoverimages
\restoregeometry  
\fi


% ============================================================= FRONT MATTER =

\ifmakefrontmatter

% ------------------------------------------------------------------ Patitul -

\cleardoublepage

\begin{center}
  %\vspace*{2cm}{\fontsize{20}{20}\selectfont\scshape\dictnameCZ\\\dictnameIS}
  \vspace*{\stretch 1}
  {\Huge\scshape\dictnameCZ\\\dictnameIS}
  \vspace{\stretch 4}
\end{center}

% -------------------------------------------------------------------- Titul -

\cleardoublepage

\begin{center}
  %\vspace*{2cm}{\fontsize{24}{24}\selectfont\scshape\dictnameCZ\\\dictnameIS}
  \vspace*{\stretch 1}
  {\HUGE\scshape\dictnameCZ\\\dictnameIS}
  \vspace{\stretch 2}
\end{center}

% ----------------------------------------------------------------- Impresum -

\clearpage

\noindent\begin{minipage}[t][\textheight][t]{\textwidth}
Autoři \textit{Islandsko-českého studijního slovníku}:\\
  \textbf{Aleš Chejn, Ján Zaťko, Jón Gíslason}\\[\baselineskip]
\ifPDF
\else
Autor \textit{Islandské gramatiky}:\\
  \textbf{Vojtěch Kupča}\\[\baselineskip]
\fi
Poděkování:\\
  \textbf{Dorotě Nierychlewské-Chejn, Jiřině Hanzlové, Jiřině Chejnové,
          Simoně Petrásové,\\ Petru Mikešovi, Vojtěchu Kupčovi, 
          Renatě Peškové Emilsson, Amiru Mulahumici, \\ Kristýně Antonové, 
          Martině Kašparové, Lucii Brabcové, Samoriele, Adamu Kožouškovi}\\[\baselineskip]
Finanční příspěvek na výtisk exemplářů určených knihovnám\\ v České republice,
Slovenské republice a na Islandu poskytli:\\
  \textbf{Petr Mikeš, Dorota Nierychlewska-Chejn, Jiřina Chejnová, Otta Chejn}
  \ifPDF
  \\[39\baselineskip]
  \else
  \\[36\baselineskip]
  \fi
\textbf{Hvalur.org - 2016}\\
Licence: \textbf{GNU Free Documentation License}\\
Sazba: \textbf{Aleš Chejn} v prostředí  \textbf\LaTeX\\
  Zvláštní poděkování patří \textbf{Paolu al-Imk\=an\=\i\ Brasolin} a uživatelům  \textbf{\TeX\ - \LaTeX\ Stack Exchange}\\[\baselineskip]
\ifPDF
Verze: \textbf{PDF verze} \textit{Islandsko-českého studijního slovníku} vhodná pro \textbf{\ifscreen offline prohlížení a vyhledávání\else vlastní tisk\fi}\\ volně dostupná na \url{www.hvalur.org}\\[\baselineskip]
\else
\textbf\ISBN\\[\baselineskip]
\fi
\ifPDF
\else
Tisk: \textbf{Tiskárna Melmen, Pardubice}\\
\fi
Heslová část vysázena patkovým fontem \TeX\ Gyre Pagella,
hesla bezpatkovým fontem {\phvfamily Helvetica}.
\end{minipage}

% -------------------------------------------------------- Table of contents -

\cleardoublepage 

\pdfbookmark\contentsname{toc}    
\tableofcontents
\thispagestyle{empty}
\cleardoublepage 

\fi

\pagenumbering{arabic}
\setcounter{page}{3}
% ============================================================= INTRODUCTION =

\ifinputintroduction

% ------------------------------------- The study Icelandic-Czech dictionary -

\chapter{O Islandsko-českém studijním slovníku}
Cílem vytvoření\textit{ Islandsko-českého studijního slovníku} je zpřístupnění islandského jazyka všem zájemcům o~tento jazyk. 

\textit{Islandsko-český studijní slovník} je určen především českým a~slovenským studentům islandštiny, začátečníkům i~pokročilým, kteří zde naleznou ucelené fonetické, morfologické a~syntaktické informace, jakož
i~slovní zásobu potřebnou k~aktivnímu i~pasivnímu užití s~rozsahem vhodným ke studiu islandštiny na vysoké škole.
Slovník je rovněž určen všem zájemcům o~islandštinu a~Island, jimž slovník nabízí encyklopedické informace, úvod do fonetiky a~gramatiky a~velké množství fotografií a~ilustrací. 
Třebaže se jedná o~slovník primárně zaměřený na potřeby českých a~slovenských uživatelů, obsahuje také řadu informací, které usnadní práci islandskému uživateli (např. použití více jak 10 000 synonym a~antonym, překlady příkladů ap.).

Slovní zásoba původně vychází z~\textit{Concise Icelandic-English Dictionary} (\cite {ic_en}), značně však byla rozšířena díky přístupu k~webovému projektu \textit{ISLEX} (\cite {int1}). Velké množství hesel bylo
přidáno na podnět uživatelů online verze \textit{Islandsko-českého studijního slovníku}.

Slovník obsahuje 30~575 hesel s~46~725 významy nebo slovními spojeními. Z~celkového počtu 30~575 je 28~097 svébytných hesel obsahujících alespoň jeden význam nebo slovní spojení. 
Zbylá hesla (v~počtu 2~478) jsou varianty, dublety nebo nepravidelné tvary.

Slovník je primárně online elektronický slovník, volně přístupný na internetových stránkách \url{www.hvalur.org}  (\cite {int14}). 
Tištěná verze slovníku je publikována na vlastní náklady v~počtu 30 kusů, přičemž většina exemplářů je určena knihovnám a~institucím v~České republice, Slovenské republice a~na Islandu
a~několik exemplářů je určeno autorům nebo zájemcům z~řad veřejnosti. Dvě PDF verze (vyhledávání\,/\addthin prohlížení offline a~vlastní tisk) jsou dostupné zdarma na internetových stránkách slovníku (viz výše).

\textit{Islandsko-český studijní slovník} je otevřený svobodný projekt, jehož se v~průběhu tvorby zúčastnilo velké množství spoluautorů (viz Historie tvorby IČSS), kterým tímto jménem autorů velmi děkuji.
Děkuji také všem uživatelům online verze slovníku za trpělivost, dobrou vůli a~veškeré připomínky a~návrhy, jmenovitě především Petru Mikešovi, Adamu Kožouškovi a~Samoriele za neutuchající víru v~náš projekt.
Poděkování patří také našim rodinám a~blízkým, díky kterým jsme se mohli soustavně věnovat tvorbě slovníku.
Srdečné poděkování patří mé ženě Dorotce za podporu a~nadšení potřebné k~tvorbě slovníku.

\blspace[5]

{\centering Aleš Chejn, Brwinów, 24. prosince 2015\par}


% ------------ The history of the formation Icelandic-Czech dictionary study -

\chapter{Historie tvorby Islandsko-českého studijního slovníku}
\longchapterskip
Projekt \textit{Islandsko-českého studijního slovníku} (IČSS) vznikl v~roce 2001 na nepřímý podnět doc. PhDr. Heleny Kadečkové, CSc. z~Univerzity Karlovy,
kdy byl vypracován seznam slov. 
Tento seznam obsahoval islandská slova a~jejich překlady, později byly přidány základní morfologické informace o~slovním druhu a~deklinační\,/\addthin konjugační koncovky podle vzoru z~\textit{Concise Icelandic-English Dictionary} (\cite {ic_en}). 
Seznam v~této době obsahoval okolo 1 200 slov. Byl zveřejněn pomocí statických webových stránek a~později byl umístěn na webu Univerzity Karlovy. 

V~roce 2006 byl slovník rozšířen o~nových 1 500 heslových slov a~stávající hesla byla rozšířena o~nové významy, zvláště významy frázových sloves. 
Ke konci roku 2006  byl slovník převeden z programu Excel do databáze MySQL a~byly zahájeny práce na vytvoření aplikace \textit{Dictionary System} (DS) (\cite {int19}), která by podporovala tvorbu IČSS. 
Obsah slovníku byl zveřejněn pod svobodnou licencí na webových stránkách \url{www.hvalur.org} (\cite {int14}), kde se nachází dodnes. 

V~roce 2007 byla vytvořena první verze aplikace DS. V~této době se k~projektu přidala Renata Pešková Emilsson, která se nyní věnuje podpoře výuky češtiny na Islandu. 
Ta slovník obohatila rozsáhlým seznamem slov z~praktického života na Islandu. Dále se slovník rozšiřoval o~heslová slova podle \textit{Concise Icelandic-English Dictionary} (\cite {ic_en}). 
V průběhu následujících třech let bylo doplněno přibližně 15~000 slov. V této době byla hesla rozšířena o~synonyma, příklady, frekvenci, oborovou a~stylovou charakteristiku. K projektu se na Islandu připojila na krátkou dobu Kristýna Antonová, studentka islandštiny na Islandské univerzitě, která 
nás upozornila na nesrovnalosti v~uživatelském rozhraní aplikace DS.

Během léta 2009 byl vytvořen skript na generování skloňování a~časování heslových slov. Výsledky generování byly porovnány s~\textit{Beygingarlýsing íslensks nútímamáls} (\cite {DalvikVM}) a~během následujících let konzultovány
s~Jónem Gíslasonem. Kristín Bjarnadóttir, vedoucí projektu \textit{Beygingarlýsing íslensks nútímamáls}, zodpověděla řadu našich dotazů týkajících se deklinací a~konjugací.

V~roce 2009 začala spolupráce s~\textit{Biolib.cz} (\cite {int5}), jmenovitě s~Ondřejem Zichou. Obdrželi jsme  databázi s~náhledy obrázků, autory a~licencemi a~latinskými názvy botanických a~zoologických druhů.

V~zimním semestru 2009 jsme požádali s~Amirem Mulahumicem, bosenským studentem islandštiny na Islandské univerzitě, Studentský fond Islandské univerzity o~grant na tvorbu zvukové databáze islandských slov. 
Studentský fond návrh přijal. Peníze z~grantu jsme určili pro islandského rodilého mluvčího, Jóna Gíslasona, lektora islandského jazyka na Islandské univerzitě v~Reykjavíku, který nahrál více jak 22~000 islandských slov do elektronické formy. 
S~Dorotou Nierychlewskou-Chejn jsme připravili dva seznamy heslových slov (jeden pro nahrávajícího a~druhý pro zpracování nahrávek). Na podporu nahrávek jsme vytvořili webové stránky nazvané \textit{Databáze s~výslovností islandských slov} (\cite {int13}) a~tyto stránky jsme přeložili do 6~jazyků (Jón Gíslason do islandštiny, Ján Zaťko do slovenštiny, francouzštiny a~angličtiny, Amir Mulamuhic do srbochorvatštiny a~Dorota Nierychlewska-Chejn do polštiny).

V~roce 2010 jsme přidali do databáze pravidla islandské výslovnosti, která jsou uvedena v~knize \textit{Handbók um íslenskan framburð} (\cite {is2}). Z~těchto pravidel jsme automaticky vygenerovali fonetický zápis výslovnosti všech hesel ve slovníku. 
Jón Gíslason ověřil správnost fonetického zápisu u~všech heslových slov.

Na konci roku 2011 jsme spojili webové stránky slovníku se stránkami aplikace a~sjednotili jsme a~vylepšili vzhled webových stránek. Cennými připomínkami se na vytvoření aplikace podílela
RNDr. Ingrid Nagyová, PhD., vysokoškolská pedagožka Ostravské Univerzity. Ján Zaťko přeložil nové webové stránky, Průvodce po slovníku do slovenštiny, francouzštiny a~angličtiny a~celý manuál aplikace do angličtiny. 
Dorota Nierychlewska-Chejn přeložila stránky do polštiny, Jón Gíslason do islandštiny a~Marcos Helena do portugalštiny.

Od prosince 2011 do března 2014 jsme realizovali první čtení, ve kterém jsme doplnili chybějící informace, přidali nová hesla a~opravili pravopisné chyby. 
Obsah slovníku jsme porovnávali především s~\textit{Íslensk orðabók} (\mbox{\cite {is}}) a~\textit{Islandsko-českým slovníkem} (\cite {is7}). S~Vojtěchem Kupčou jsme konzultovali mnohé nesnáze při překladech.

V~letech 2013 až 2015 byl vytvořen morfologický klíč, který rozšiřuje klasický systém zobrazení morfologické informace v~islandských slovnících o~přesné zařazení daného hesla do určité deklinační\,/\addthin konjugační skupiny, jejíž
skloňování\,/\addthin časování je ukázáno na příkladu. V~roce 2014 jsme obdrželi roční grant od Ostravské univerzity v~Ostravě na tvorbu morfologického klíče. S~realizací grantu nám velmi pomohla doc. PhDr. Diana Svobodová, Ph.D. a~Ing. Zuzana Havrlantová.

Od března 2014 do srpna 2015 jsme realizovali druhé finální čtení, ve kterém jsme porovnávali informace s~webovým slovníkem \textit{ISLEX} (\cite {int1}). 
Na požádání jsme obdrželi od Þórdís Úlfarsdóttir (vedoucí projektu \textit{ISLEX}) seznam víceslovných příslovečných frází. Takto poskytnutý seznam nám usnadnil opravy slovníku. 
S~Þórdís Úlfarsdóttir jsme konzultovali různé gramatické jevy a~záležitosti týkající se deklinací\,/\addthin konjugací.
Martina Kašparová, studentka islandštiny na Islandské univerzitě, prošla celý slovník, přičemž se soustředila mimo jiné na stylistické a~pravopisné korektury překladů. 
Lucie Brabcová udělala jazykové a~typografické korektury celého slovníku.

Vojtěch Kupča vydal v~roce 2014 \textit{Islandskou gramatiku} (\cite {is77}) a~souhlasil s~připojením této publikace ke~knižní verzi IČSS.
V~roce 2015 jsme s~přispěním Paola al-Imk\=an\=\i\ Brasolin vysázeli slovník v~typografickém prostředí {\LaTeX}. V~únoru 2016 byl slovník vytištěn v~Tiskárně MELMEN, s.r.o v~Pardubicích a svázán 
v~Knihařství Vojtíšková Hrubá v~Pardubicích. 

\blspace[5]

{\centering Všem spoluautorům děkujeme za spolupráci.\par}


% -------------------------------------------------  Guide to the dictionary -

\chapter{Průvodce po slovníku}
\begin{multicols}{2}
\def\tableA{\noindent\begin{tabularx}{\columnwidth}{>{\bfseries}lll}
  \toprule
  \textit{přípona} & \textit{přípona s~koncovkou} & \textit{příklad} \\
  \midrule
  -ing-   & -ingur   & rugl··ingur \\
          & -ing     & hreyf··ing \\
          & -ingi    & erf··ingi \\
          & -ingar   & lækn··ingar \\
  -ling-  & -lingur  & vett··lingur \\
  -ning-  & -ningur  & flut··ningur \\
          & -ning    & set··ning \\
          & -ningar  & kos··ningar \\
  -að-    & -aður    & mark··aður \\
  -nað-   & -naður   & skil··naður \\
  -ar-    & -ari     & dóm··ari \\
  -skap-  & -skapur  & kveð··skapur \\
  -dóm-   & -dómur   & lær··dómur \\
  -leik-  & -leikur  & sann··leikur \\
          & -leiki   & mögu··leiki \\
  -lát-   & -látur   & laus··látur \\
  -átt-   & -átta    & víð··átta \\
  -sam-   & -samur   & starf··samur \\
  -sem-   & -semi    & hamingju··semi \\
  -nesk-  & -neskur  & tékk··neskur \\
          & -neska   & tékk··neska \\
  -leg-   & -legur   & blíð··legur \\
          & -lega    & blíð··lega \\
  -ung-   & -ungur   & gler··ungur \\
          & -unga    & snopp··unga \\
          & -ungi    & sveit··ungi \\
          & -ung     & nýj··ung \\
  -ótt-   & -óttur   & blett··óttur \\
  -ræn-   & -rænn    & suð··rænn \\
          & -ræna    & nor··ræna \\
  -ern-   & -erni    & fað··erni \\
  -ind-   & -indi    & vís··indi \\
  -and-   & -andi    & gef··andi \\
  -ug-    & -ugur    & göf··ugur \\
          & -ugt     & stöð··ugt \\
  -ust-   & -usta    & þjón··usta \\
  -ald-   & -aldur   & far··aldur \\
  -uð-    & -uður    & söfn··uður \\
  -úð-    & -úð      & dul··úð \\
  -versk- & -verskur & ung··verskur \\
          & -verska  & ung··verska \\
  -ist-   & -isti    & hass··isti \\
  -ism-   & -ismi    & húman··ismi \\
  -ang-   & -angur   & far··angur \\
  \bottomrule
  \end{tabularx}
  \captionof{table}{Seznam přípon označených ve slovníku}
  \label{table1}}

\def\tableB{\noindent\begin{tabularx}{\columnwidth}{>{\bfseries}lX}
  \toprule
  \textbf{symbol} & \textbf{vysvětlivka} \\
  \midrule
  m                   & označuje slovní druh \\
  \textsubscript{1}   & označuje třídu \\
  (-a, -ar)           & označuje deklinační / konjugační koncovky hesla \\
  \textsubscript{8}   & označuje podtřídu\\
  \bottomrule
  \end{tabularx}
  \captionof{table}{Morfologická informace}
  \label{table2}}

\def\tableC{\noindent\begin{tabularx}{\columnwidth}
    {>{\bfseries}l>{\footnotesize\itshape}lllll}
  \toprule
  \multirow{3}{*}{{{\textbf{f{\textsubscript{1}}} \Large{\textbf{7}}}}}
    & nom & kisa & kisan    & kisur & kisurnar \\
    & acc & kisu & kisuna   & kisur & kisurnar \\
    & dat & kisu & kisunni  & kisum & kisunum \\
    & gen & kisu & kisunnar & kisa  & kisanna \\
  \bottomrule
  \end{tabularx}
  \captionof{table}{Skloňování hesla \textit{kisa}}
  \label{table3}}

\def\tableD{\noindent\begin{tabularx}{\columnwidth}{>{\bfseries}lllll}
  \toprule
  \textit{zkratka} & \textit{slovo} & \textit{pád} & \textit{číslo}
    & \specialcell{\textit{předmět}\\\textit{podmět}} \\
  \midrule
  e-að  & eitthvað   & 1. & sg & neživotný \\
  e-ð   & eitthvað   & 4. & sg & neživotný \\
  e-ir  & einhverjir & 1. & pl & životný \\
  e-ja  & einhverja  & 4. & pl & životný \\
  e-jum & einhverjum & 3. & pl & životný \\
  e-m   & einhverjum & 3. & sg & životný \\
  e-n   & einhvern   & 4. & sg & životný \\
  e-r   & einhver    & 1. & sg & životný \\
  e-rra & einhverra  & 2. & pl & životný \\
  e-rs  & einhvers   & 2. & sg & životný \\
  e-s   & einhvers  & 2. & sg & neživotný \\
  e-u   & einhverju  & 3. & sg & neživotný \\
  \bottomrule
  \end{tabularx}
  \captionof{table}{Seznam použitých zkratek neurčitých zájmen}
  \label{table4}}

%#############################################################################

\section{Heslo}

\subsection*{Řazení slov}

Heslová slova jsou zobrazena tučně a~jsou seřazena podle islandské abecedy:

\blspace
{\bfseries\centering
  a~á b c d ð e é f g h i~í j k~l m n\\
  o~ó p q r s~t u~ú v~w x y ý z~þ æ ö\par}
\blspace

K~snadnější orientaci v~heslové části je použit u~každého písmene tzv. palcový index, který se nachází na hřbetě slovníku. Každé písmeno je v~palcovém indexu barevně odlišeno.

\subsection*{Homonyma}

Homonyma (heslová slova stejného reprezentativního tvaru) jsou označena povýšenými arabskými číslicemi uváděnými za heslovým slovem.

\blspace
  \dicEntry \dicTerm{vor\textsuperscript{1}} \dicPos{n} \dicFlx{(-s, -)}
  \dicEntry \dicTerm{vor\textsuperscript{2}} \dicPos{pron} \dicFlx{pers gen pl}
\blspace

Homonyma jsou seřazena podle českého ustáleného řazení slovních druhů. Poté následují předpony a~zkratky.
V~případě dvou podstatných jmen je pořadí mužský rod, ženský rod, střední rod. V~případě dvou podstatných jmen stejného rodu má přednost slovo se slabým skloňováním před slovem se silným skloňováním. V~případě dvou sloves stejného reprezentativního tvaru má přednost sloveso se slabým časováním před slovesem se silným časováním.

\subsection*{Varianty a~dublety}

U~některých heslových slov jsou v~záhlaví hesla uvedeny a) varianty (obměny)

\blspace
  \dicEntry \dicTerm{ærsla··fenginn}, \dicTerm{ærslafullur} \dicPos{adj}
\blspace

nebo b) dublety (dvojtvary).

\blspace
  \dicEntry \dicTerm{{ævin··týri}, \dicTerm{ævintýr}} \dicPos{n}
\blspace

Varianty a~dublety jsou méně frekventované než samotné heslo. Lze je vyhledat ve slovníku jako hesla obsahující fonetické a~morfologické informace s~odkazem na původní heslové slovo.

\blspace
  \dicEntry \dicTerm{ærsla··|fullur} \dicIPA{{a}{i}{\textsubring{r}}{s}{\textsubring{d}}{l}{a}{f}{\textscy}{\textsubring{d}}{l}{\textscy}{\textsubring{r}}} \dicPos{adj} \dicFlx{(comp -fyllri, sup -fyllstur)} \dicLink{ærslafenginn}
  \dicEntry \dicTerm{ævin··týr} \dicIPA{{a}{i}{\textlengthmark}{v}{\textsci}{n}{t\textsuperscript{h}}{i}{\textsubring{r}}} \dicPos{n} \dicFlx{(-s, -)} \dicLink{ævintýri}
\blspace

\subsection*{Vybrané tvary jako hesla}

Ve slovníku jsou uvedeny vybrané tvary heslových slov, u~kterých by mohl mít uživatel potíže s~identifikací. Jedná se o~a) genitiv jednotného čísla ({\textbf{sg}}) a~nominativ množného čísla ({\textbf{pl}}) podstatných jmen, b) ženský a~střední rod přídavných jmen a~2. a~3. stupeň rodu mužského přídavných jmen, c) oznamovací způsob ({\textbf{ind}}) přítomného času ({\textbf{praes}}) a~minulého času ({\textbf{pf}}) 1. osoby ({\textbf{pers}}) sg a~pl a~spojovací způsob ({\textbf{con}}) minulého času 1. osoby sg a~jiné nepravidelné tvary.

\blspace
  \dicEntry \dicTerm{öttum} \dicPos{v} \dicFlx{ind pf pl 1 pers} \dicLink{etja}
\blspace

\subsection*{Dělení heslových slov}

K~dělení heslových slov je použita střední tečka \textbf{'·'}, která odděluje části složeného slova, přípony a~předpony. Platí následující pravidla:

\renewcommand{\labelenumi}{\alph{enumi})}
\begin{enumerate}
\item jediné dělení heslového slova je označeno dvěma středními tečkami \textbf{'··'}. Může se jednat o~dělení částí složeného slova
\blspace
\dicEntry \dicTerm{ráð··hús}
\blspace

nebo o~dělení mezi kořenem a~předponou / příponou.

\blspace
\dicEntry \dicTerm{þjón··usta}
\blspace
\item V~případě, že se v~hesle vyskytuje více dělení, jsou použity dvě střední tečky \textbf{'··'} pro významově poslední dělení a~jedna střední tečka \textbf{'·'} pro zbylé / zbylá dělení.
\blspace
\dicEntry \dicTerm{við·bótar··líf·eyris·sparn·aður}
\blspace
\item V~hesle s~předponou je označena jedna nebo i~více předpon (každá zvlášť).
\blspace
\dicEntry \dicTerm{ó··að·finnan·legur}
\blspace
\item V~hesle s~příponou je označena jedna nebo i~více přípon (každá zvlášť).
\blspace
\dicEntry \dicTerm{létt·úð··ugur}
\blspace
\item V~hesle s~příponou je označena pouze přípona / přípony poslední části složeného slova. Důvodem pro neoznačení přípon jiných částí složených slov je lepší čitelnost hesla. Příklad označení pouze přípony poslední části složeniny:
\blspace
\dicEntry \dicTerm{flutninga··miðl·ari}
\end{enumerate}

Následuje seznam přípon, které jsou ve slovníku označeny:

\tableA

\section{Frekvence slov}

Symbolem \dicsymFrequent\ je označeno 2~825 hesel patřících mezi velmi frekventovaná islandská slova. Tento seznam vychází z~korpusu \textit{Íslenskur Orðasjóður} (\cite{int7}) a~knihy \textit{Íslensk orðtíðnibók} (\cite{is2}).

\blspace
  \dicEntry \dicTerm{koma}\textsuperscript{2} \dicsymFrequent
\blspace

\section{Fonetický zápis}

V~hranatých závorkách v~záhlaví hesla je uveden fonetický zápis výslovnosti hesla v~mezinárodní fonetické abecedě (IPA, angl. International Phonetic Alphabet).

\blspace
  \dicEntry \dicTerm{öldu··rót} \dicIPA{{\oe}{l}{\textsubring{d}}{\textscy}{r}{ou}{\textsubring{d}}}
\blspace

Znak '{\textlengthmark}' označuje dlouhou samohlásku nebo souhlásku.

\blspace
  \dicEntry \dicTerm{ör}\textsuperscript{2} \dicIPA{{\oe}{\textlengthmark}{\textsubring{r}}}
\blspace

\section{Slovní druhy, předpony a~zkratky}

V~záhlaví hesla je uvedena morfologická informace popisující a) slovní druh hesla, b) předponu nebo c) zkratku. Pokud se heslo řadí do více slovních druhů, jsou jednotlivé zkratky pro slovní druhy odděleny lomítkem.

\blspace
  \dicEntry \dicTerm{af} \dicPos{prep / adv}
\blspace

Informace o~slovním druhu, předponě či zkratce je vysázena barevně a~tučně.

\subsection*{Podstatná jména}

Podstatná jména jsou uvedena v~1. pádě (\textbf{nom}) jednotného čísla (\textbf{sg}) nebo v~1. pádě množného čísla (\textbf{pl}), pokud se heslo v~jednotném čísle nevyskytuje. Podstatná jména jsou označena následovně: \textbf{m} (l. masculinum) - podstatné jméno rodu mužského,
\textbf{f} (l. femininum) - podstatné jméno rodu ženského, \textbf{n} (l. neutrum) - podstatné jméno rodu středního, \textbf{subs} (l. substantivum) - podstatné jméno bez jednoznačného mluvnického rodu.

\blspace
  \dicEntry \dicTerm{hestur} \dicPos{m}
  \dicEntry \dicTerm{Hanoí} \dicPos{subs}
\blspace

\subsection*{Přídavná jména}

Přídavná jména jsou uvedena v~1. pádě (\textbf{nom}) rodu mužského (\textbf{m}) jednotného čísla (\textbf{sg}) 1. stupně (\textbf{pos}) (l. positivus) a~jsou označena zkratkou \textbf{adj} (l. adjectivum).

\blspace
  \dicEntry \dicTerm{fallegur} \dicPos{adj}
\blspace

V~případě, že u~přídavného jména dochází v~rodě ženském (\textbf{f}) nebo středním (\textbf{n}) k~přehlásce anebo je tvar druhého stupně (\textbf{comp} (l. comparativus)) či třetího stupně (\textbf{sup} (l. superlativus)) nepravidelný, pak jsou tyto tvary uvedeny v~nominativu jako heslová slova s~odkazem na hlavní heslový tvar

\blspace
  \dicEntry \dicTerm{yngri} \dicPos{adj} \dicFlx{comp m} \dicLink{ungur}
\blspace

nebo jako svébytná hesla s~významy.

\blspace
  \dicEntry \dicTerm{fremri} \dicPos{adj} \dicFlx{comp (sup fremstur)}
\blspace

\subsection*{Zájmena}

Zájmena jsou označena zkratkou \textbf{pron} (l. pronomen) a~jsou tříděna na zájmena ukazovací \textbf{dem} (angl. demonstrative), přivlastňovací \textbf{poss} (angl. possessive), osobní \textbf{pers} (angl. personal), neurčitá \textbf{indef} (angl. indefinite), zvratná \textbf{refl} (angl. reflexive) a~vztažná \textbf{rel} (angl. relative).

\blspace
  \dicEntry \dicTerm{þessi} \dicPos{pron} \dicFlx{dem}
\blspace

U~zájmen přivlastňovacích je v~3. osobě označen rod.

\blspace
  \dicEntry \dicTerm{hans} \dicPos{pron} \dicFlx{poss m}
\blspace

U~zájmen osobních je označen pád a~číslo.

\blspace
  \dicEntry \dicTerm{mig} \dicPos{pron} \dicFlx{pers acc sg}
\blspace

\subsection*{Číslovky}

Číslovky jsou označeny zkratkou \textbf{num} (l. numerale). V~případě vybraných číslovek (1 až 4) jsou ve slovníku uvedeny jako heslová slova také tvary ženského i~středního rodu. V~takovém případě jsou označeny zkratkou \textbf{m} pro mužský rod, zkratkou \textbf{f} pro ženský rod a~zkratkou \textbf{n} pro střední rod. Zkratkou \textbf{ord} (angl. ordinal number) jsou označeny číslovky řadové.

\blspace
  \dicEntry \dicTerm{tut··tugu} \dicPos{num}
  \dicEntry \dicTerm{tveir} \dicPos{num} \dicFlx{pl m}
  \dicEntry \dicTerm{fyrsti} \dicPos{num} \dicFlx{ord}
\blspace

\subsection*{Slovesa}

Slovesa jsou uvedena v~infinitivu a~jsou označena zkratkou \textbf{v} (l. verbum).

\blspace
  \dicEntry \dicTerm{fara} \dicPos{v}
\blspace

Slovesa ve středním rodě jsou označena jako \textbf{refl} (angl. reflexive).

\blspace
  \dicEntry \dicTerm{nálgast} \dicPos{v} \dicFlx{refl}
\blspace

Slovesa, která jsou vždy neosobní, jsou označena \textbf{impers} (angl. impersonal).

\blspace
  \dicEntry \dicTerm{svima} \dicPos{v} \dicFlx{impers}
\blspace

U~sloves je dále uvedeno, s~jakým pádem se pojí. Zkratka \textbf{acc} označuje 4. pád, \textbf{dat} 3. pád, \textbf{gen} 2. pád a~\textbf{nom} 1. pád. V~případě, že se sloveso pojí s~různými pády, je použito lomítko.

\blspace
  \dicEntry \dicTerm{klóra} \dicPos{v} \dicFlx{acc / dat}
\blspace

V~případě, že se sloveso pojí s~více pády, je použito znaménko plus.

\blspace
  \dicEntry \dicTerm{gefa} \dicPos{v} \dicFlx{dat + acc}
\blspace

\subsection*{Příslovce}

Příslovce jsou označena zkratkou \textbf{adv} (l. adverbium). Pokud chybí informace o~stupni, jedná se o~1. stupeň stupňování (\textbf{pos}). 2. stupeň je vždy označen zkratkou \textbf{comp} a~3. stupeň zkratkou \textbf{sup}.

\blspace
  \dicEntry \dicTerm{þar} \dicPos{adv}
  \dicEntry \dicTerm{ofar} \dicPos{adv} \dicFlx{comp (pos uppi, sup efst)}
  \dicEntry \dicTerm{síðast} \dicPos{adv} \dicFlx{sup (comp síðar)}
\blspace

\subsection*{Předložky}

Předložky jsou označeny zkratkou \textbf{prep} (l. praepositio). Pokud se předložka pojí výhradně s~jedním pádem, jsou použity zkratky pro 4. pád (\textbf{acc}), 3. pád (\textbf{dat}), 2. pád (\textbf{gen}).

\blspace
  \dicEntry \dicTerm{af} \dicPos{prep / adv} \dicFlx{dat}
\blspace

V~případě, že se předložka pojí s~různými pády, je použito lomítko.

\blspace
  \dicEntry \dicTerm{fyrir} \dicPos{prep / adv} \dicFlx{acc / dat}
\blspace

\subsection*{Spojky}

Spojky jsou označeny zkratkou \textbf{conj} (l. conjunctio).

\subsection*{Částice}

Částice jsou označeny zkratkou \textbf{part} (l. particula).

\subsection*{Citoslovce}

Citoslovce jsou označeny zkratkou \textbf{inter} (l. interjectio).

\subsection*{Předpony}

Předpony jsou označeny zkratkou \textbf{predp}.

\subsection*{Zkratky}

Zkratky jsou označeny zkratkou \textbf{zkr}.

\section{Skloňování, časování a~stupňování}

Každé heslo obsahuje zápis morfologické informace, která přesně informuje, jak se heslo skloňuje / časuje nebo stupňuje.

V~slovníku jsou použity dva odlišné systémy zápisu morfologické informace, které se vzájemně doplňují:

\begin{enumerate}
\item tradiční systém - do závorky za slovním druhem se zapisují deklinační / konjugační koncovky. Počet koncovek závisí na slovním druhu (viz dále).
\item numerický systém (vychází ze systému Helgi Haraldssona (\cite {is6})) - číslice odkazující na morfologický klíč, kde jsou podrobně vypsány deklinační / konjugační třídy.
Další informace o~numerickém systému lze nalézt na straně \pageref{sec:dec}.
\end{enumerate}

Deklinační / konjugační tvary nebo koncovky podstatných jmen, přídavných jmen, zájmen, číslovek, sloves a~příslovcí jsou uvedeny v~závorkách za slovním druhem.

Hnízdovací šev \textbf{'|'} označuje místo ve slově, kde se deklinační koncovka připojuje ke slovu.

\blspace
  \dicEntry \dicTerm{heils|a{\textsuperscript{1}}} \dicPos{f} \dicFlx{(-u)}
\blspace

Pokud není hnízdovací šev \textbf{'|'} uveden, tak se koncovka připojuje k~samotnému heslu.

\blspace
  \dicEntry \dicTerm{mynd} \dicPos{f} \dicFlx{(-ar, -ir)}
\blspace

Pokud dochází v~ohýbání k~přehlásce nebo je ohýbání nepravidelné, pak jsou uvedeny celé tvary hesla.

\blspace
  \dicEntry \dicTerm{maður} \dicPos{m} \dicFlx{(manns, menn)}
\blspace

V~případě, že heslové slovo má více tvarů, jsou varianty koncovky odděleny lomítkem.

\blspace
  \dicEntry \dicTerm{beð|ur} \dicPos{m} \dicFlx{(-s / -jar, -ir)}
\blspace

\subsection*{Podstatná jména}

U~podstatných jmen je uveden tvar hesla pro 2.~pád jednotného čísla a~1.~pád množného čísla.

\blspace
  \dicEntry \dicTerm{hval|ur} \dicPos{m} \dicFlx{(-s, -ir)}
\blspace

kde '\textit{hvals}' je 2.~pád jednotného čísla a~'\textit{hvalir}' 1.~pád množného čísla.
Pokud je uveden pouze jeden tvar, jedná se o~2.~pád jednotného čísla a~znamená to rovněž, že heslové slovo se neohýbá v~množném čísle.

\blspace
  \dicEntry \dicTerm{heils|a{\textsuperscript{1}}} \dicPos{f} \dicFlx{(-u)}
\blspace

Pokud je použita zkratka pro množné číslo (\textbf{pl}), znamená to, že se heslo ohýbá pouze v~množném čísle.

\blspace
  \dicEntry \dicTerm{afar··kostir} \dicPos{m} \dicFlx{pl}
\blspace

Pokud je použita zkratka \textbf{indecl}, znamená to, že podstatné jméno je nesklonné.

\blspace
  \dicEntry \dicTerm{október} \dicPos{m} \dicFlx{indecl}
\blspace

\subsection*{Přídavná jména}

U~přídavných jmen jsou uvedeny v~závorkách tvary, které poukazují na nepravidelnosti ve skloňování.

Pokud přídavné jméno v~ženském rodě neztrácí koncové '\textit{-r}', je uveden tvar hesla v~ženském rodě.

\blspace
  \dicEntry \dicTerm{hýr} \dicPos{adj} \dicFlx{(f -)}
\blspace

Dále jsou uvedeny tvary ženského rodu, ve~kterých dochází k~přehlásce:

\blspace
  \dicEntry \dicTerm{reglu··|samur} \dicPos{adj} \dicFlx{(f -söm)}
\blspace

Pokud je přídavné jméno nesklonné, je uvedena zkratka \textbf{indecl}.

\blspace
  \dicEntry \dicTerm{tví··mála} \dicPos{adj} \dicFlx{indecl}
\blspace

Pokud má přídavné jméno nepravidelné tvary 2.~a~3.~stupně stupňování, jsou tyto tvary uvedeny v~závorce spolu se zkratkou pro 2.~stupeň (\textbf{comp}) a~3.~stupeň (\textbf{sup}).

\blspace
  \dicEntry \dicTerm{þunnur} \dicPos{adj} \dicFlx{(comp þynnri, sup þynnstur)}
\blspace

V~případě, že je uveden pouze jeden ze dvou zbývajících stupňů, znamená to, že heslo se stupňuje pouze v~uvedených stupních.

\blspace
  \dicEntry \dicTerm{syðri} \dicPos{adj} \dicFlx{comp (sup syðstur)}
\blspace

\subsection*{Slovesa}

Islandská slovesa se dělí na slabá, silná a~nepravidelná. Vzhledem k~tomuto dělení se liší počet slovních tvarů uvedených v~závorce za zkratkou \textbf{v}. Princip uvádění slovních tvarů spočívá v~tom, že uživatel
je schopen (více méně) z~těchto tvarů odvodit časování slovesa. V~tomto slovníku jsme
záměrně umístili některá slabá slovesa mezi silná (místo dvou slovesných tvarů je uvedeno pět tvarů) kvůli tomu, aby byl uživatel přesněji informován o~časování slovesa.

Slabá slovesa, která tvoří 1.~osobu jednotného čísla minulého času oznamovacího způsobu koncovkou '\textit{-aði}', mají v~závorce uvedenu jen jednu koncovku, jmenovitě '\textit{(-aði)}'.

\blspace
  \dicEntry \dicTerm{ætl|a} \dicPos{v} \dicFlx{(-aði)}
\blspace

U~zbývajících skupin slabých sloves se v~závorce nachází koncovka 1.~osoby jednotného čísla minulého času oznamovacího způsobu a~supinum.

\blspace
  \dicEntry \dicTerm{kenn|a} \dicPos{v} \dicFlx{(-di, -t)}
\blspace

Silná nebo nepravidelná slovesa mají v~závorce vždy pět tvarů, jmenovitě 1.~osobu jednotného čísla přítomného času oznamovacího způsobu, 1.~osobu jednotného čísla minulého času oznamovacího způsobu, 1.~osobu množného čísla minulého času oznamovacího způsobu, 1.~osobu jednotného čísla minulého času spojovacího způsobu a~supinum.

\blspace
  \dicEntry \dicTerm{grípa} \dicPos{v} \dicFlx{(gríp, greip, gripum, gripi, gripið)}
\blspace

Silná nebo nepravidelná slovesa vyskytující se pouze v~neosobním užití mají v~závorce 4 tvary, jmenovitě 3.~osobu jednotného čísla přítomného času oznamovacího způsobu, 3.~osobu jednotného čísla minulého času oznamovacího způsobu, 3.~osobu jednotného čísla minulého času spojovacího způsobu a~supinum.

\blspace
  \dicEntry \dicTerm{kross··|bregða} \dicPos{v} \dicFlx{(-bregður, -brá, -brygði, -brugðið)}
\blspace

U~některých významů v~heslovém slově se vyskytují gramatické informace, které popisují chování daného významu.

Zkratka \textbf{impers} označuje, že se daný význam slovesa vyskytuje v~neosobním užití. Pomocí slovního spojení je vyjádřeno užití slovesa.

\blspace
  \dicEntry \dicTerm{batn|a} \dicPos{v} \dicFlx{(-aði)} \dicTerm{e-m batnar} \dicFlx{impers} (kdo) se uzdravuje, (kdo) se zotavuje, (komu) je lépe
\blspace

Zkratka \textbf{refl} označuje mediopasivní tvary slovesa.

\blspace
  \dicEntry \dicTerm{að··lag|a} \dicPos{v} \dicFlx{(-aði)} \dicIdiom{aðlagast} {\dicPhraseIS{aðlagast e-u}} \dicFlx{refl} {přizpůsobit se (čemu)}
\blspace

\subsection*{Příslovce}

U~příslovcí se v~záhlaví nachází tvary stupňování. Tyto tvary jsou uvedeny v~závorce spolu se zkratkou pro 2.~stupeň (\textbf{comp}) a~3.~stupeň (\textbf{sup}).

\blspace
  \dicEntry \dicTerm{vel\textsuperscript{2}} \dicPos{adv} \dicFlx{(comp betur, sup best)}
\blspace

Nepravidelné tvary 2.~a~3.~stupně jsou uvedeny buď jako hesla odkazující na základní tvar v~1.~stupni (\textbf{pos}), anebo jako svébytná hesla. V~obou případech jsou uvedeny v~závorce zbývající stupně.

\blspace
  \dicEntry \dicTerm{síst} \dicPos{adv} \dicFlx{sup (pos varla, comp síður)}
\blspace

Pokud některý ze stupňů uveden není, znamená to, že pro dané slovo neexistuje.

\blspace
  \dicEntry \dicTerm{síðar} \dicPos{adv} \dicFlx{comp (sup síðast)}
\blspace

\section{Morfologická informace} \label{sec:dec}

V~záhlaví hesla se kromě deklinačních / konjugačních koncovek nachází také morfologické numerické informace, které odkazují na morfologický klíč (viz strana \pageref{sec:morpho}), kde jsou podrobně vypsány kompletní deklinační / konjugační tabulky.

\blspace
  \dicEntry \dicTerm{ap|i} \dicPos{m}[1] \dicFlx{(-a, -ar)}[8]
\blspace

Morfologická informace je zapsána podle následujícího schématu: \textbf{m{\textsubscript{1}} (-a, -ar){\textsubscript{8}}}

\blspace
\tableB
\blspace

Deklinační / konjugační třídy a~jejich počet nejsou v~islandštině ustálené a~slouží pouze pro rozlišení odlišných skloňování / časování.

Názvy podtříd slouží pro rozlišení odlišných skloňování / časování.

\subsection*{Práce s~morfologickým klíčem}

Pro vyhledání kompletních deklinačních / konjugačních tabulek slouží následující postup, který demonstrujeme na ukázce:

\blspace
  \dicEntry \dicTerm{kis|a} \dicPos{f}[1] \dicFlx{(-u, -ur)}[7]
\blspace

Nejdříve určíme slovní druh (v~případě podstatných jmen také rod), v~tomto případě podstatné jméno rodu ženského, a~navštívíme Morfologický klíč (str. \pageref{sec:morpho}) pro podstatná jména rodu ženského (str. \pageref{sec:morpho_f}).
Následně dohledáme třídu (v~tomto případě číslo {\textsubscript{\textbf{1}}}) a~následně dohledáme dodatečné numerické informace za závorkou (v~tomto případě číslo {\textsubscript{\textbf{7}}}).

\tableC

\section{Definice}

\subsection*{Základní tvar}

Základním tvarem definice je, že jednomu islandskému slovu odpovídá jeden nebo více českých významů.

\blspace
  \dicEntry \dicTerm{landa··fræð|i} \dicPos{f} \dicFlx{(-i)} \textCS{zeměpis, geografie}
\blspace

V~některých případech je v~české části použita závorka. V~prvním případě je závorka užívaná ke~zkrácení zápisu dvou slov ('\textit{pomoc}', '\textit{výpomoc}').

\blspace
  \dicEntry \dicTerm{að··stoð} \dicPos{f} \dicFlx{(-ar)} \textCS{(vý)pomoc, asistence}
\blspace

V~druhém případě se v~závorce nachází upřesňující informace, která zabraňuje dvojznačnosti.

\blspace
  \dicEntry \dicTerm{fíkj|a} \dicPos{f} \dicFlx{(-u, -ur)} \textCS{fík} {\footnotesize {\textCS{(plod)}}}
\blspace

V~mnoha případech je závorka v~české části použita pro příklad užití českého slova a~tím vymezení jeho významu.

\blspace
  \dicEntry \dicTerm{að·gengi··legur} \dicPos{adj} {přístupný, dostupný (informace ap.)}
\blspace

Lomítko \textbf{'/'} slouží v~české části významu k~oddělení více významů pojících se s~jedním nebo více slovy a~je použito v~české části následujícím způsobem:

\blspace
  \dicEntry \dicTerm{aðal··|gata} \dicPos{f} \dicFlx{(-götu, -götur)} \textCS{hlavní třída / ulice}
\blspace

Význam přečteme jako '\textit{hlavní třída, hlavní ulice}'.

\subsection*{Opis heslového slova}

V~případech, kdy islandské slovo nemá ekvivalent v~českém jazyce, jsou významy islandských slov opsány delším opisem a~zobrazeny menším písmem. Příkladem jsou slova z~lidových pověstí, kulinářské speciality nebo botanické pojmy.

\blspace
  \dicEntry \dicTerm{til··ber|i} \dicPos{m} \dicFlx{(-a, -ar)} {\footnotesize{pov.}} {\footnotesize{stvoření, které dojí krávy a~ovce hospodářů}}
\blspace

\subsection*{Slovní spojení}

V~definici hesla jsou kromě samotných významů uváděna také slovní spojení, která jsou sázena tučně. Po slovním spojení následuje překlad slovního spojení.

\blspace
  \dicEntry \dicTerm{lær|a} \dicPos{v} \dicFlx{(-ði, -t)} \textbf{læra e-ð utanbókar} {učit se (co) zpaměti, memorovat (co)}
\blspace

Závorky jsou použity na tu část slovního spojení, kterou je možné vypustit.

\blspace
  \dicEntry \dicTerm{ár\textsuperscript{1}} \dicPos{f} \dicFlx{(-ar, -ar)} \textbf{taka (of) djúpt í árinni} {\footnotesize{přen.}} {přehnat (co v~tvrzení)}
\blspace

Příklad můžeme přečíst jako '\textit{taka djúpt í árinni}' nebo '\textit{taka of djúpt í árinni}'.

Mezi slovní spojení v~tomto slovníku řadíme a) tvary heslového slova, b) příkladová spojení, c) ustálená slovní spojení a~d) rčení nebo přísloví.

\begin{enumerate}
\item Tvary heslového slova\\
Do této skupiny řadíme tvar množného čísla hesla (zkratka \textbf{pl}).
\blspace
\dicEntry \dicTerm{bót} \dicPos{f} \dicFlx{(-ar, bætur)} \textbf{bætur} \dicFlx{pl} {dávky, přídavky, příspěvky}\\

Dále do této skupiny řadíme tvar slovesa v~středním slovesném rodě (zkratka \textbf{refl}).\\

\dicEntry \dicTerm{finna} \dicPos{v} \dicFlx{(finn, fann, fundum, fyndi, fundið)} \dicIdiom{finnast} {\dicPhraseIS{finnast}} \dicFlx{refl} {být nalezen, najít se, nalézt se, nacházet se}
\blspace
\item Příkladová spojení\\
Příkladová spojení ilustrují komunikační chování heslových slov v~příslušných významech a~v~různých syntagmatech (např. sloveso s~určitou předložkou ap.). Jsou označena tučně modrou barvou a~odsazená pro snadnou orientaci jako svébytné heslo.
\blspace
\dicEntry \dicTerm{koma\textsuperscript{2}} \dicPos{v} \dicFlx{(kem, kom, komum, kæmi, komið)} \dicIdiom{koma}[að] {\dicPhraseIS{koma að e-m}} {přiblížit se ke (komu), přijít ke (komu)}
\blspace
\item Ustálená slovní spojení\\
Tento druh slovních spojení je charakterizován většinou přeneseným významem. Tato spojení jsou uváděna za posledním významem nebo příkladovým spojením a~jsou označená tučně modrou barvou a~odsazená pro snadnou orientaci jako svébytné heslo.
\blspace
\dicEntry \dicTerm{kast} \dicPos{n} \dicFlx{(-s, köst)} \dicIdiom{kast} {\dicPhraseIS{þegar til kastanna kemur}} {\footnotesize{přen.}} {když jde do tuhého}
\blspace
\item Rčení a~přísloví\\
Rčení a~přísloví se nacházejí vždy na konci definice a~jsou odlišená symbolem \dicsymProverb\ a~odsazená pro snadnou orientaci.
\blspace
\dicEntry \dicTerm{æfing} \dicPos{f} \dicFlx{(-ar, -ar)} \dicProverb{} \dicPhraseIS{Æfing skapar meistara.} {\footnotesize{přís.}} \textCS{Cvičení dělá mistra.}
\end{enumerate}

\section{Vícevýznamovost}

\subsection*{Vícevýznamová hesla}

Značné množství hesel má více než jeden význam. Významy jsou odlišeny arabskými číslicemi a~v~některých případech i~islandskými synonymy.

\blspace
  \dicEntry \dicTerm{land} \dicPos{n} \dicFlx{(-s, lönd)} {\textbf{1.}} \dicSynonym{þurrlendi} {souš, pevnina, země} {\textbf{2.}} \dicSynonym{árbakki} {břeh} {\textbf{3.}} \dicSynonym{ríki} {země, stát, vlast} {\textbf{4.}} \dicSynonym{landareign} {země, půda, pozemek}
\blspace

\subsection*{Vícevýznamová slovní spojení}

V~případě, že slovní spojení uvnitř definice má více významů, jsou jednotlivé významy odlišeny písmeny {\textbf{a., b., c.}} atd.

\blspace
  \dicEntry \dicTerm{drag} {\dicPos{n} \dicFlx{(-s, drög)} {\textbf{1.}} {mokřina, podmoklý terén;} {\textbf{2.}} \textbf{drög} {\footnotesize{pl}} {\textbf{a.}} \dicSynonym{uppsprettur} {prameny (řeky ap.)} {\textbf{b.}} \dicSynonym{undirbúningur} {náčrt, návrh.}}
\blspace

\section{Řazení významů a~slovních spojení v~definici}
\longsectionskip

V~definici jsou nejdříve seřazeny významy slova a~odlišeny arabskými číslicemi.

\subsection*{V rámci významu}

Slovní spojení jsou uváděna v~rámci jednotlivého významu, pokud se s~ním významově vážou.

\blspace
  \dicEntry \dicTerm{tím|i} \dicPos{m} \dicFlx{(-a, -ar)} {\textbf{1.}} \dicSynonym{tíð} {čas, doba}; \textbf{í þann tíma} \dicFlx{adv} \dicSynonym{þá\textsuperscript{2}} {pak} {\textbf{2.}} \dicSynonym{rétt stund} {(správný) čas / moment}
\blspace

\subsection*{Za posledním očíslovaným významem}

Pokud se slovní spojení neváže k~žádnému významu, je řazeno za posledním očíslovaným významem. Ustálená slovní spojení jsou uvedena symbolem \dicsymIdiom.

\blspace
  \dicEntry \dicTerm{finna} \dicPos{v} \dicFlx{(finn, fann, fundum, fyndi, fundið)} {\textbf{3.}} \dicSynonym{skynja} {cítit, vnímat} \dicIdiom{finna}[að] {\dicPhraseIS{finna að e-u}} {kritizovat (co), nacházet chyby na (čem)}
\blspace

\subsection*{Způsob řazení slovních spojení}

Slovní spojení, která se nachází za posledním očíslovaným významem, jsou řazena podle předložky / příslovce / zájmena, se kterou se pojí. Pro snadnou orientaci v~rozsáhlých definicích je uvedeno heslové slovo nebo tvar heslového slova plus předložka / příslovce / zájmeno a~barevně odlišeno modrou barvou.

\blspace
  \dicEntry \dicTerm{koma\textsuperscript{2}} \dicPos{v} \dicFlx{(kem, kom, komum, kæmi, komið)} \dicIdiom{koma}[út] {\dicPhraseIS{koma út}} {vyjít (kniha ap.), objevit se (na trhu ap.)}
\blspace

Poté následují slovní spojení v~slovesném středním rodě řazená opět podle předložky / příslovce / zájmena, se kterou se pojí:

\blspace
  \dicEntry \dicTerm{koma\textsuperscript{2}} \dicPos{v} \dicFlx{(kem, kom, komum, kæmi, komið)} \dicIdiom{komast} {\dicPhraseIS{komast}} \dicFlx{refl} {dostat se, dostávat se}
\blspace

Poté následují ustálená slovní spojení, která nebylo možné zařadit do předchozích položek:

\blspace
  \dicEntry \dicTerm{koma\textsuperscript{2}} \dicPos{v} \dicFlx{(kem, kom, komum, kæmi, komið)} \dicIdiom{koma} {\dicPhraseIS{(komdu) sæll (og blessaður)}} {ahoj, nazdar, dobrý den}
\blspace

Nakonec jsou uvedena přísloví a~rčení označená symbolem \dicsymProverb\:

\blspace
  \dicEntry \dicTerm{eik} \dicPos{n} \dicFlx{(-ar / -ur, -ur)} \dicProverb{} \dicPhraseIS{Ekki / eigi fellur eik við fyrsta högg.} {\footnotesize{přís.}} {Jednou ranou dub nespadne.}
\blspace

\section{Synonyma a~antonyma / hesla ke~srovnání}
\longsectionskip

\subsection*{Synonyma}

Synonyma nebo synonymní slovní spojení jsou uvedeny před opisem významu hesla. Jsou umístěna v~závorkách. Islandská synonyma slouží a) islandskému uživateli k~rozlišení českých významů b) českému uživateli jako dodatečná informace o~významu slova.

\blspace
  \dicEntry \dicTerm{staða} \dicPos{f} \dicFlx{(stöðu, stöður)} {\textbf{1.}} \dicSynonym{ástand} {situace, stav, poměry} {\textbf{2.}} \dicSynonym{embætti} {post, (pracovní) místo / pozice (ve firmě ap.)}
\blspace

V~případě, že synonymum lze nalézt jako heslo ve slovníku, je použita šipka '\dicsymSee\ '.

\blspace
  \dicEntry \dicTerm{bygging} {\dicPos{f} \dicFlx{(-ar, -ar)} {\textbf{1.}} \dicSynonym{hús} {budova, stavba, stavení} {\textbf{2.}} \textit{(það að byggja)} {(vý)stavba} }
\blspace

\subsection*{Antonyma / Hesla ke~srovnání}

Antonyma či hesla ke~srovnání jsou uvedena za opisem významu heslového slova a~jsou označena šipkou '$\shortuparrow$'.

\blspace
  \dicEntry \dicTerm{kaldur} \dicPos{adj} \dicFlx{(f köld)} {\textbf{1.}} {studený, chladný} \textit{($\shortuparrow$ heitur)}
\blspace

\section{Oborové a~stylové charakteristiky}

K~upřesnění významu heslového slova slouží oborové a~stylové charakteristiky.

\subsection*{Oborové charakteristiky}

Oborové charakteristiky zařazují islandské slovo do určitého oboru a~naznačují použití daného slova v~jazyce

\blspace
  \dicEntry \dicTerm{berg} \dicPos{n} \dicFlx{(-s, -)} {\textbf{1.}} {\footnotesize{geol.}} {hornina, kámen} {\textbf{2.}} \dicSynonym{klöpp\textsuperscript{1}} {skalní stěna, skalisko}
\blspace

a~vymezují význam českého slova.

\blspace
  \dicEntry \dicTerm{mús} \dicPos{f} \dicFlx{(-ar, mýs)} {\textbf{1.}} {\footnotesize{zool.}} {myš} (\textit{l. Mus}) {\textbf{2.}} {\footnotesize{poč.}} {myš}
\blspace

\subsection*{Stylové charakteristiky}

Stylové charakteristiky jsou uváděny v~případě, kdy význam heslového slova je možné použít pouze při určité příležitosti nebo je význam jazykově zabarven.

\blspace
  \dicEntry \dicTerm{gaur} \dicPos{m} \dicFlx{(-s, -ar)} {\textbf{1.}} {\footnotesize{hovor.}} {chlápek, maník, člověk}
\blspace

\section{Syntax}

Syntaktické informace popisují:

\begin{enumerate}
\item zdali je sloveso tranzitivní nebo intranzitivní a~s~jakým(i) pádem (pády) se pojí, pokud je tranzitivní,
\item zdali je podmět / předmět životný či neživotný,
\item zdali je sloveso osobní či neosobní a~v~jakém pádu je podmět, pokud je sloveso neosobní,
\item zdali je podmět gramatický,
\item zdali je sloveso použito ve středním rodě,
\item nejčastější příklady, které se pojí se slovesem,
\item se kterými příslovci nebo předložkami se sloveso pojí.
\end{enumerate}

Syntaktické informace se ve slovníku vyskytují ve třech formách: a) informace uvedené zkratkou, b) teoretické příklady a~c) praktické příklady.

\subsection*{Syntaktické informace uvedené zkratkou}

Zde máme na mysli informace o~pádě (\textbf{nom, acc, dat, gen}), slovesu v~středním slovesném rodě (\textbf{refl}), neosobním slovesu (\textbf{impers}), čísle (\textbf{sg, pl}), rozkazovacím způsobu (\textbf{imper}) a~spojovacím způsobu (\textbf{con}).

\subsection*{Teoretické příklady}

Teoretickým příkladem nazýváme příklad, kdy podmět nebo předmět je zastoupen zájmenem neurčitým v~pádě, se~kterým se slovo pojí a~sloveso je zpravidla v~infinitivu. V~případě, že podmět není v~1.~pádě (např. u~sloves neosobních), je sloveso uvedeno v~3.~osobě jednotného čísla přítomného času.
Ve slovních spojeních používáme zkratky pro tvary neurčitého zájmena \textbf{einhver} (někdo) a~\textbf{eitthvað} (něco). Tyto zkratky označují a) pád, b) číslo a~c) životnost nebo neživotnost podmětu / předmětu.\\
Následuje seznam zkratek neurčitých zájmen použitých ve slovníku: \\

\tableD
\blspace
Příklad použití zkratek s~neurčitými zájmeny:

\blspace
  \dicEntry \dicTerm{leit|a} \dicPos{v} \dicFlx{(-aði)} \textbf{leita að e-u / e-m} {hledat (co / koho)} \dicExampleIS{leita að lyklunum} \dicExampleCS{hledat klíče}
\blspace

Příklad tedy přečteme jako '\textit{leita að einhverju}' (\textit{hledat co}) a~'\textit{leita að einhverjum}' (\textit{hledat koho}). V~tomto případě se sloveso pojí s~3.~pádem a~předmět může být životný i~neživotný.

\subsection*{Praktické příklady}

Ve slovníku se vyskytuje velké množství praktických příkladů a~jejich překladů. Praktický příklad se liší od příkladu teoretického tím, že místo zájmen neurčitých jsou uvedena frekventovaná slova z~běžného projevu. Praktické příklady slouží k~ilustraci použití
heslového slova.\\
Praktickým příkladem může být a) (složené) slovo, b) slovní spojení nebo c) věta.

\begin{enumerate}
\item Slovo jako praktický příklad:
\blspace
\dicEntry \dicTerm{hett|a} \dicPos{f} \dicFlx{(-u, -ur)} \textbf{1.} {kapuce, kápě} \dicExampleIS{hettuúlpa} \dicExampleCS{bunda s~kapucí}
\blspace
\item Slovní spojení jako praktický příklad:
\blspace
\dicEntry \dicTerm{tím|i} \dicPos{m} \dicFlx{(-a, -ar)} \textbf{3.} \dicSynonym{klukkustund} {hodina} {\footnotesize{(šedesát minut)}} \dicExampleIS{tveggja tíma gangur} \dicExampleCS{dvouhodinový pochod}
\blspace
\item Věta jako praktický příklad:
\blspace
\dicEntry \dicTerm{mæt|a} \dicPos{v} \dicFlx{(-ti, -t)} \dicIdiom{mætast} \dicPhraseIS{mætast} \dicSynonym*{ná saman} \dicFlx{refl} {setkat se, potkat se} \dicExampleIS{Þau mættust á gangstéttinni.} \dicExampleCS{Potkali se na chodníku.}
\end{enumerate}

\section{Fotografie a~ilustrace}

Ve slovníku je publikováno 630 fotografií rostlin, živočichů a~fotografií, která dokumentují místopisně heslová slova (např. kulinářské speciality).
Fotografie jsou publikovány pod veřejnými licencemi. Jména autorů a~názvy licencí lze nalézt v~Seznamu autorů fotografií a~ilustrací na straně \pageref{sec:photo}. U~hesla a~konkrétně u~daného významu, ke kterému se fotografie vztahuje, je použit symbol fotoaparátu \dicsymPhoto. Fotografie je umístěna buď pod heslem, anebo v~horní části stránky, kde se heslo nachází, anebo na stránce (stránkách) následující(ch).

%\par\begin{center}\setlength\fboxsep{0pt}\setlength\fboxrule{0.5pt}
%\fbox{\includegraphics[width=6cm]{ds_image_posthus_0_2.jpg}}\end{center}
%\par\begin{center}\footnotesize {Pósthús}\end{center}

\end{multicols}

% ---------------------------------------- List of abbreviations and symbols -

\chapter{Seznam použitých zkratek a symbolů}
\begin{multicols}{2}
\begin{description}[leftmargin=!,labelwidth=1.5cm]
%[\breaklabel\setleftmargin{40pt}\setlabelstyle{\bfseries}
%\addtolength{\itemsep}{-0.5\baselineskip}]
\item[{abb}] {zkrácený}
\item[{acc}] {4. pád, akuzativ}
\item[{adj}] {přídavné jméno}
\item[{adv}] {příslovce}
\item[{akt}] {rod činný, aktivum}
\item[{anat.}] {anatomie}
\item[{angl.}] {anglicky}
\item[{astro.}] {astronomie, výzkum vesmíru}
\item[{ap.}] {a podobně}
\item[{básn.}] {básnický, poetický výraz}
\item[{biol.}] {biologie}
\item[{bot.}] {botanika}

\item[{comp}] {2. stupeň stupňování, komparativ}
\item[{con}] {spojovací způsob, konjunktiv}
\item[{conj}] {spojka}
\item[{dat}] {3. pád, dativ}
\item[{def}] {určitý tvar}
\item[{dem}] {zájmeno ukazovací}
\item[{dět.}] {dětsky}

\item[{e-að}] {co (sg nom neživotný) (\textit{eitthvað})}
\item[{e-ð}] {co (sg acc neživotný) (\textit{eitthvað})}
\item[{e-ir}] {kdo (pl nom životný) (\textit{einhverjir})}
\item[{e-ja}] {koho (pl acc životný) (\textit{einhverja})}
\item[{e-jum}] {komu (pl dat životný) (\textit{einhverjum})}
\item[{e-m}] {komu (sg dat životný) (\textit{einhverjum})}
\item[{e-n}] {koho (sg acc životný) (\textit{einhvern})}
\item[{e-r}] {kdo (sg nom životný) (\textit{einhver})}
\item[{e-rra}] {koho (pl gen životný) (\textit{einhverra})}
\item[{e-rs}] {koho (sg gen životný) (\textit{einhvers})}
\item[{e-s}] {čeho (sg gen neživotný) (\textit{einhvers})}
\item[{e-u}] {čemu (sg dat neživotný) (\textit{einhverju})}

\item[{ekon.}] {ekonomika, obchod}
\item[{elek.}] {elektřina}
\item[{f}] {rod ženský}
\item[{filos.}] {filosofie, logika}
\item[{form.}] {formálně}
\item[{fyz.}] {fyzika}
\item[{gen}] {2. pád, genitiv}
\item[{geog.}] {zeměpis, geografie}
\item[{geol.}] {geologie}
\item[{han.}] {hanlivě, pejorativně}
\item[{hist.}] {historie}
\item[{hovor.}] {hovorově}
\item[{hrub.}] {hrubě}
\item[{hud.}] {hudba}
\item[{imper}] {rozkazovací způsob, imperativ}
\item[{impers}] {sloveso neosobní}
\item[{ind}] {oznamovací způsob, indikativ}
\item[{indecl}] {nesklonný}
\item[{indef}] {neurčitý tvar, (pron +) zájmeno neurčité}
\item[{inf}] {infinitiv}
\item[{int}] {tázací zájmeno}
\item[{inter}] {citoslovce}
\item[{chem.}] {chemie}
\item[{jaz.}] {jazykověda}
\item[{kulin.}] {kulinářství, vaření}
\item[{l.}] {latinsky}
\item[{let.}] {letectví}
\item[{lit.}] {literatura, vydavatelství}
\item[{m}] {rod mužský}
\item[{mat.}] {matematika}
\item[{med. }] {lékařství}
\item[{meteo.}] {meteorologie}
\item[{myt.}] {mytologie}
\item[{n}] {rod střední}
\item[{náb.}] {náboženství}
\item[{nám.}] {námořnictví, rybolov}
\item[{nom}] {1. pád, nominativ}
\item[{num}] {číslovka}
\item[{ord}] {řadová číslovka}
\item[{p}] {osoba}
\item[{part}] {částice}
\item[{pers}] {(v +) osoba, (pron +) zájmeno osobní}
\item[{pf}] {minulý čas}
\item[{pl}] {množné číslo, plurál}
\item[{poč.}] {informatika}
\item[{pol.}] {politika, politologie}
\item[{poss}] {zájmeno přivlastňovací}
\item[{pos}] {1. stupeň stupňování, pozitiv}
\item[{pov.}] {lidové pověsti, folkloristika}
\item[{pp}] {příčestí minulé}
\item[{praes}] {přítomný čas}
\item[{práv.}] {právnictví, soudnictví}
\item[{predp}] {předpona}
\item[{prep}] {předložka}
\item[{presp}] {příčestí přítomné}
\item[{pron}] {zájmeno}
\item[{prop}] {vlastní jméno, proprium}
\item[{přen.}] {přeneseně}
\item[{přís.}] {přísloví}
\item[{psych.}] {psychologie}
\item[{refl}] {(v +) médium, (pron +) zvratné zájmeno}
\item[{rel}] {vztažné zájmeno}
\item[{sg}] {jednotné číslo, singulár}
\item[{slang.}] {slang}
\item[{sport.}] {sport}
\item[{stav.}] {stavebnictví, architektura}
\item[{subs}] {podstatné jméno (bez rodu)}
\item[{sup}] {3. stupeň stupňování, superlativ}
\item[{supin}] {supinum}
\item[{škol.}] {školství}
\item[{techn.}] {technika, mechanika}
\item[{v}] {sloveso}
\item[{voj.}] {vojenství}
\item[{zast.}] {zastarale}
\item[{zkr}] {zkratka}
\item[{zool.}] {zoologie}

\item[{;}] {středník uprostřed definice odděluje věty; středník na konci definice odděluje významy}

\item[{(...)}] {v kulatých závorkách jsou uvedeny doplňkové informace}
\item[{[...]}] {v hranatých závorkách je uveden fonetický zápis výslovnosti}\footnote{Vysvětlení fonetických symbolů je uvedeno v kapitole Fonetika \ref{sec:phon_phonems} na straně \pageref{sec:phon_phonems}}
\item[{|}] {naznačuje místo, kde se deklinační koncovka připojuje ke slovnímu základu}
\item[{·}] {vedlejší dělení složeného slova}
\item[{··}] {hlavní dělení složeného slova}
\item[{+}] {plus}
\item[{/}] {lomítko označuje více možností}
\item[{\dicsymSee}] {šipka odkazuje na heslo ve slovníku}
\item[{\dicsymCompare}] {šipka odkazuje na heslo ve slovníku ke srovnání}
\item[{\dicsymExampleIS}] {uvozuje islandský příklad}
\item[{\dicsymExampleCS}] {uvozuje český překlad příkladu}
\item[{\dicsymIdiom}] {uvozuje slovní spojení}
\item[{\dicsymFrequent}] {heslo patří mezi frekventované výrazy}

\item[{\dicsymPhoto}] {fotografie, ilustrace}
\item[{\dicsymProverb}] {přísloví, rčení}

\end{description}
\end{multicols}

\chapter{Skrá yfir skammstafanir og tákn}
\begin{multicols}{2}
\begin{description}[leftmargin=!,labelwidth=1.5cm]
%[\breaklabel\setleftmargin{40pt}\setlabelstyle{\bfseries}
%\addtolength{\itemsep}{-0.5\baselineskip}]
\item[{abb}] {stýfður}
\item[{acc}] {þolfall}
\item[{adj}] {lýsingarorð}
\item[{adv}] {atviksorð}
\item[{akt}] {germynd}
\item[{anat.}] {líffærafræði}
\item[{angl.}] {enska}
\item[{astro.}] {geimvísindi, stjörnufræði}
\item[{ap.}] {e.þ.h.}
\item[{básn.}] {ljóðrænn}
\item[{biol.}] {líffræði}
\item[{bot.}] {grasafræði}

\item[{comp}] {miðstig}
\item[{con}] {viðtengingarháttur}
\item[{conj}] {samtenging}
\item[{dat}] {þágufall}
\item[{def}] {með greini}
\item[{dem}] {ábendingarfornafn}
\item[{dět.}] {barnamál}

\item[{e-að}] {(\textit{eitthvað})}
\item[{e-ð}] {(\textit{eitthvað})}
\item[{e-ir}] {(\textit{einhverjir})}
\item[{e-ja}] {(\textit{einhverja})}
\item[{e-jum}] {(\textit{einhverjum})}
\item[{e-m}] {(\textit{einhverjum})}
\item[{e-n}] {(\textit{einhvern})}
\item[{e-r}] { (\textit{einhver})}
\item[{e-rra}] {(\textit{einhverra})}
\item[{e-rs}] {(\textit{einhvers})}
\item[{e-s}] {(\textit{einhvers})}
\item[{e-u}] { (\textit{einhverju})}

\item[{ekon.}] {hagfræði, viðskipti}
\item[{elek.}] {rafmagn}
\item[{f}] {kvenkyn}
\item[{filos.}] {heimspeki, rökfræði}
\item[{form.}] {formlegt mál}
\item[{fyz.}] {eðlisfræði}
\item[{gen}] {eignarfall}
\item[{geog.}] {landafræði}
\item[{geol.}] {jarðfræði}
\item[{han.}] {niðrandi}
\item[{hist.}] {sagnfræði}
\item[{hovor.}] {talmál}
\item[{hrub.}] {gróft}
\item[{hud.}] {tónlist}
\item[{imper}] {boðháttur}
\item[{impers}] {ópersónuleg sögn}
\item[{ind}] {framsöguháttur}
\item[{indecl}] {óbeygjanlegur}
\item[{indef}] {óákveðinn, (pron +) óákveðið fornafn}
\item[{inf}] {nafnháttur}
\item[{int}] {spurnigarfornafn}
\item[{inter}] {upphrópun}
\item[{chem.}] {efnafræði}
\item[{jaz.}] {málfræði, málvísindi}
\item[{kulin.}] {eldhús}
\item[{l.}] {latína}
\item[{let.}] {flug, flugmál}
\item[{lit.}] {skáldskapur}
\item[{m}] {nafnorð, karlkyn}
\item[{mat.}] {stærðfræði}
\item[{med. }] {læknisfræði}
\item[{meteo.}] {veðurfræði}
\item[{myt.}] {goðafræði}
\item[{n}] {nafnorð, hvorugkyn}
\item[{náb.}] {trúfræði}
\item[{nám.}] {sjómennska}
\item[{nom}] {nefnifall}
\item[{num}] {töluorð}
\item[{ord}] {raðtala}
\item[{p}] {persóna}
\item[{part}] {ögn}
\item[{pers}] {(v +) persóna, (pron +) persónufornafn}
\item[{pf}] {þátíð}
\item[{pl}] {fleirtala}
\item[{poč.}] {tölvufræði, tölvur}
\item[{pol.}] {stjórnmál, stjórnmálafræði}
\item[{poss}] {eignarfornafn}
\item[{pos}] {frumstig}
\item[{pov.}] {þjóðtrú}
\item[{pp}] {lýsingarháttur þátíðar}
\item[{praes}] {nútíð}
\item[{práv.}] {lögfræði, dómsmál}
\item[{predp}] {forskeyti}
\item[{prep}] {forsetning}
\item[{presp}] {lýsingarháttur nútíðar}
\item[{pron}] {fornafn}
\item[{prop}] {sérnafn}
\item[{přen.}] {afleidd merking}
\item[{přís.}] {málsháttur}
\item[{psych.}] {sálfræði}
\item[{refl}] {(v +) miðmynd, (pron +) afturbeygt fornafn}
\item[{rel}] {tilvísunarfornafn}
\item[{sg}] {eintala}
\item[{slang.}] {slangur}
\item[{sport.}] {íþróttir}
\item[{stav.}] {húsasmíði, byggingarlist}
\item[{subs}] {nafnorð (án kyns)}
\item[{sup}] {efstastig}
\item[{supin}] {sagnbót}
\item[{škol.}] {skólamál}
\item[{techn.}] {tækni, verkfræði}
\item[{v}] {sagnorð}
\item[{voj.}] {her}
\item[{zast.}] {úrelt mál}
\item[{zkr}] {skammstöfun}
\item[{zool.}] {dýrafræði}

\item[{;}] {semíkomma í miðri skýringu aðgreinir setningu; semíkomma í enda skýringar aðgreinir mismunandi merkingarsvið}

\item[{(...)}] {innan sviga eru viðbótarupplýsingar}
\item[{[...]}] {innan hornklofa er hljóðritaður framburður}
\item[{|}] {merkir stað, þar sem beygingarending bætist við stofn}
\item[{·}] {aukasamskeyti samsetts orðs}
\item[{··}] {aðalsamskeyti samsetts orðs}
\item[{+}] {plús}
\item[{/}] {skástrik táknar fleiri möguleika}
\item[{\dicsymSee}] {ör vísar til uppflettiorðs í orðabókinni}
\item[{\dicsymCompare}] {ör vísar til uppflettiorðs í bókinni til samanburðar}
\item[{\dicsymExampleIS}] {kemur á undan dæmi á íslensku}
\item[{\dicsymExampleCS}] {kemur á undan tékkneskri þýðingu á dæminu}
\item[{\dicsymIdiom}] {orð með háa tíðni}

\item[{\dicsymPhoto}] {ljósmynd}
\item[{\dicsymProverb}] {málsháttur}

\end{description}
\end{multicols}

% ------------------------------------------------- Headword part and alphabet

\chapter{Heslová část}
{\renewcommand{\arraystretch}{2}%
%\topskip0pt
\vstretch{1}
{\setlength{\tabcolsep}{1.5em}
\begin{table}[htpb]
\centering

\begin{tabular}{ccccccc} \hline
A a &  & G g &  & O o &  & V v \\
Á á  &  & H h &  & Ó ó &  & W w \\
B b &  & I i &  & P p &  & X x \\
C c &  & Í í &  & Q q &  & Y y \\
D d &  & J j &  & R r &  & Ý ý \\
Ð ð &  & K k &  & S s &  & Z z \\
E e &  & L l &  & T t &  & Þ þ \\
É é  &  & M m &  & U u &  & Æ æ \\
F f &  & N n &  & Ú ú &  & Ö ö \\
\hline
\end{tabular}

%\caption{\textbf{Islandská abeceda}}
\label{my-islandska-abeceda}
\end{table}
}
\vstretch{1}
{\renewcommand{\arraystretch}{1}%

\clearpage

\fi

% =========================================================== THE DICTIONARY =

\ifinputletters

\cleardoublepage

\dictionarygeometry
\pagestyle{myheadings}

\newcommand*{\alphabetF}{%
  a,aa,b,c,d,e,ee,f,g,h,i,ii,j,k,l,m,
  n,o,oo,p,r,s1,s2, t,u,uu,v,w,x,y,yy,th,ae,oe}

\ifPDFtest
\dicEntry[slétta] \dicTerm{slétt|a\smash{\textsuperscript{1}}} \dicIPA{{s}{\textsubring{d}}{l}{j}{\textepsilon}{h}{\textsubring{d}}{a}} \dicPos{f}[1] \dicFlx{(‑u, ‑ur)}[19] \dicSynonym{flöt} \dicDirectTranslationCS{rovina, planina}
\dicEntry[slétta] \dicTerm{slétt|a\smash{\textsuperscript{2}}} \dicIPA{{s}{\textsubring{d}}{l}{j}{\textepsilon}{h}{\textsubring{d}}{a}} \dicPos{v}[3] \dicFlx{(‑aði)\,/\addthin (‑i, ‑)}[2] \dicFlx{acc} \dicSynonym*{gera sléttan} \dicDirectTranslationCS{uhladit, uhlazovat, zahladit, zahlazovat, srovnat, srovnávat} \dicExampleIS{slétta lóðina} \dicExampleCS{srovnat pozemek}
\dicEntry[sléttbakur] \dicTerm{slétt··bak|ur} \dicIPA{{s}{\textsubring{d}}{l}{j}{\textepsilon}{h}{\textsubring{d}}{\textsubring{b}}{a}{\r{g}}{\textscy}{\textsubring{r}}} \dicPos{m}[6] \dicFlx{(‑s, ‑ar)}[25] \dicFieldCat{zool.} \dicDirectTranslationCS{velryba černá} \textit{(l.~{\textLA{Eubalaena glacialis}})}  \dicsymPhoto\ 
\dicFigure{ds_image_slettbakur_0_1.jpg}{Sléttbakur}{Sléttbakur - United States Federal Government, PD}
\dicEntry[sléttlendi] \dicTerm{slétt··lendi} \dicIPA{{s}{\textsubring{d}}{l}{j}{\textepsilon}{h}{\textsubring{d}}{l}{\textepsilon}{n}{\textsubring{d}}{\textsci}} \dicPos{n}[2] \dicFlx{(‑s)}[20] \dicDirectTranslationCS{rovina, planina}
\dicEntry[sléttprjón] \dicTerm{slétt··prjón} \dicIPA{{s}{\textsubring{d}}{l}{j}{\textepsilon}{h}{\textsubring{d}}{p\smash{\textsuperscript{h}}}{r}{j}{ou}{\textsubring{n}}} \dicPos{n}[2] \dicFlx{(‑s)}[2] \dicDirectTranslationCS{pletení hladce}
\dicEntry[sléttubönd] \dicTerm{sléttu··bönd} \dicIPA{{s}{\textsubring{d}}{l}{j}{\textepsilon}{h}{\textsubring{d}}{\textscy}{\textsubring{b}}{\oe}{n}{\textsubring{d}}} \dicPos{n}[2] \dicFlx{pl}[9] \dicFieldCat{lit.} \dicIndirectTranslationCS{čtyřřádkový verš, který je možné číst popředu i~pozpátku, aniž by se rým a~aliterace změnili}
\dicEntry[sléttur] \dicTerm{sléttur} \dicsymFrequent\  \dicIPA{{s}{\textsubring{d}}{l}{j}{\textepsilon}{h}{\textsubring{d}}{\textscy}{\textsubring{r}}} \dicPos{adj}[1]\dicFlx{}[-10] \textbf{1.} \dicSynonym{flatur} \dicDirectTranslationCS{rovný, plochý} \dicExampleIS{slétt land} \dicExampleCS{rovná země}  \textbf{2.} \dicFieldCat{mat.} \dicDirectTranslationCS{sudý} \dicExampleIS{slétt tala} \dicExampleCS{sudé číslo};  \dicIdiom{sléttur}{ \dicPhraseIS{e‑að er slétt og fellt}} \dicDirectTranslationCS{(co) je v~pořádku (na první pohled)};  \dicPhraseIS{réttur og sléttur} \dicFlx{adj} \dicDirectTranslationCS{obyčejný, zvyklý};  \dicPhraseIS{segja farir sínar ekki sléttar} \dicLangCat{přen.} \dicDirectTranslationCS{nemít na růžích ustláno}
\dicEntry[sléttuúlfur] \dicTerm{sléttu··úlf|ur} \dicIPA{{s}{\textsubring{d}}{l}{j}{\textepsilon}{h}{\textsubring{d}}{\textscy}{u}{l}{v}{\textscy}{\textsubring{r}}} \dicPos{m}[6] \dicFlx{(‑s, ‑ar)}[4] \dicFieldCat{zool.} \dicDirectTranslationCS{kojot, kojot prérijní} \textit{(l.~{\textLA{Canis latrans}})}  \dicsymPhoto\ 
\dicFigure{ds_image_slettuulfur_0_1.jpg}{Sléttuúlfur}{Sléttuúlfur - Justin Johnsen, CC BY 2.0}
\dicEntry[slifsi] \dicTerm{slifsi}\dicTerm{, slipsi} \dicIPA{{s}\-{\textsubring{d}}\-{l}\-{\textsci}\-{f}\-{s}\-{\textsci}\-} \dicPos{n}[2] \dicFlx{(‑s, ‑)}[14] \textbf{1.} \dicIndirectTranslationCS{hedvábný šátek nošený k~islandskému národnímu kroji}  \textbf{2.} \dicSynonym{hálsbindi} \dicDirectTranslationCS{kravata, vázanka}
\dicEntry[slig] \dicTerm{slig} \dicIPA{{s}{\textsubring{d}}{l}{\textsci}{\textlengthmark}{x}} \dicPos{n}[2] \dicFlx{(‑s)}[2] \dicDirectTranslationCS{uštvání, strhání};  \dicPhraseIS{ríða e‑u á slig} \dicDirectTranslationCS{strhat (co), uštvat (co)}
\dicEntry[sliga] \dicTerm{slig|a} \dicIPA{{s}{\textsubring{d}}{l}{\textsci}{\textlengthmark}{\textbabygamma}{a}} \dicPos{v}[1] \dicFlx{(‑aði)}[1] \dicFlx{acc} \dicDirectTranslationCS{prohnout, prohýbat, zatížit, zatěžovat} \dicExampleIS{Snjórinn sligaði þakið.} \dicExampleCS{Sníh prohnul střechu.};  \dicIdiom{sligast}{ \dicPhraseIS{sligast}} \dicFlx{refl} \dicDirectTranslationCS{prohnout se, prohýbat se}
\dicEntry[slikja] \dicTerm{slikj|a}\dicTerm{, slíkja} \dicIPA{{s}\-{\textsubring{d}}\-{l}\-{\textsci}\-{\textlengthmark}\-{\r{\textObardotlessj}}\-{a}\-} \dicPos{f}[1] \dicFlx{(‑u, ‑ur)}[7] \dicDirectTranslationCS{zabarvení, odstín, tón} \dicExampleIS{slikja á flík} \dicExampleCS{zabarvení prádla}
\dicEntry[slikk] \dicTerm{slikk} \dicIPA{{s}{\textsubring{d}}{l}{\textsci}{h}{\r{g}}} \dicPos{n}[2] \dicFlx{(‑s)}[2] \dicPhraseIS{fyrir slikk} \dicFlx{adv} \dicDirectTranslationCS{za hubičku\,/\addthin babku} \dicExampleIS{fá e‑ð fyrir slikk} \dicExampleCS{dostat (co) za babku}
\dicEntry[slinkur] \dicTerm{slink|ur} \dicIPA{{s}{\textsubring{d}}{l}{i}{\r{\ng}}{\r{g}}{\textscy}{\textsubring{r}}} \dicPos{m}[9] \dicFlx{(‑s, ‑ir)}[8] \dicSynonym{kippur} \dicDirectTranslationCS{cuknutí, trhnutí, škubnutí} \dicExampleIS{Það kom slinkur á bílinn í beygjunum.} \dicExampleCS{Auto sebou cuklo v~zatáčce.}
\dicEntry[slippur] \dicTerm{slipp|ur} \dicIPA{{s}{\textsubring{d}}{l}{\textsci}{h}{\textsubring{b}}{\textscy}{\textsubring{r}}} \dicPos{m}[9] \dicFlx{(‑s, ‑ir)}[8] \dicDirectTranslationCS{loděnice}
\dicEntry[slipsi] \dicTerm{slipsi} \dicIPA{{s}{\textsubring{d}}{l}{\textsci}{\textsubring{b}}{s}{\textsci}} \dicPos{n}[2] \dicFlx{(‑s, ‑)}[14] \dicLink{slifsi}
\dicEntry[slit] \dicTerm{slit} \dicIPA{{s}{\textsubring{d}}{l}{\textsci}{\textlengthmark}{\textsubring{d}}} \dicPos{n}[2] \dicFlx{(‑s, ‑)}[5] \textbf{1.} \dicSynonym{eyðing} \dicDirectTranslationCS{opotřebení, opotřebování} \dicExampleIS{slit á fötum} \dicExampleCS{opotřebení šatů}  \textbf{2.} \dicPhraseIS{slit} \dicFlx{pl} \dicSynonym{endir} \dicDirectTranslationCS{konec, ukončení}  \textbf{3.} \dicPhraseIS{slit} \dicFlx{pl} \dicFieldCat{práv.} \dicDirectTranslationCS{likvidace, zrušení}
\dicEntry[sliti] \dicTerm{sliti} \dicIPA{{s}{\textsubring{d}}{l}{\textsci}{\textlengthmark}{\textsubring{d}}{\textsci}} \dicPos{v} \dicFlx{con pf sg 1 pers} \dicLink{slíta}
\dicEntry[slitið] \dicTerm{slitið} \dicIPA{{s}{\textsubring{d}}{l}{\textsci}{\textlengthmark}{\textsubring{d}}{\textsci}{\texttheta}} \dicPos{v} \dicFlx{supin} \dicLink{slíta}
\dicEntry[slitinn] \dicTerm{slitinn} \dicsymFrequent\  \dicIPA{{s}{\textsubring{d}}{l}{\textsci}{\textlengthmark}{\textsubring{d}}{\textsci}{\textsubring{n}}} \dicPos{adj}[6]\dicFlx{}[-2] \textbf{1.} \dicSynonym*{farinn í sundur} \dicDirectTranslationCS{natržený, přetržený}  \textbf{2.} \dicSynonym*{mikið notaður} \dicDirectTranslationCS{opotřebovaný, opotřebený, obnošený} \dicExampleIS{slitnir skór} \dicExampleCS{opotřebované boty}  \textbf{3.} \dicSynonym{þreklítill} \dicDirectTranslationCS{vysílený, vyžilý}
\dicEntry[slitlag] \dicTerm{slit··|lag} \dicIPA{{s}{\textsubring{d}}{l}{\textsci}{\textlengthmark}{\textsubring{d}}{l}{a}{x}} \dicPos{n}[2] \dicFlx{(‑lags, ‑lög)}[8] \dicDirectTranslationCS{ochranná vrstva} \dicIndirectTranslationCS{(zvláště při stavbě silnic)}
\dicEntry[slitna] \dicTerm{slitn|a} \dicIPA{{s}{\textsubring{d}}{l}{\textsci}{h}{\textsubring{d}}{n}{a}} \dicPos{v}[1] \dicFlx{(‑aði)}[75] \textbf{1.} \dicSynonym{bresta} \dicDirectTranslationCS{roztrhnout se, přetrhnout se, trhat se} \dicExampleIS{Bandið slitnar strax.} \dicExampleCS{Provázek se ihned přetrhne.}  \textbf{2.} \dicSynonym*{eyðast} \dicDirectTranslationCS{opotřebovat se, opotřebit se} \dicExampleIS{Vélin slitnaði.} \dicExampleCS{Stroj se opotřeboval.}  \textbf{3.} \dicSynonym{rofna} \dicDirectTranslationCS{přerušit se, být přerušen} \dicExampleIS{Sambandið slitnaði.} \dicExampleCS{Spojení se přerušilo.};  \dicIdiom{slitna}[upp]{ \dicPhraseIS{e‑að slitnar upp}} \dicDirectTranslationCS{(co) se uvolňuje (od kotviště) (člun ap.)};  \dicIdiom{slitna}[upp úr]{ \dicPhraseIS{það slitnar upp úr e‑u}} \dicFlx{impers} \dicDirectTranslationCS{schází z~(čeho) (ze smlouvy ap.)} \dicExampleIS{Það slitnaði upp úr samningaviðræðum.} \dicExampleCS{Sešlo ze smluvních rozhovorů.}
\dicEntry[slitróttur] \dicTerm{slitr··óttur} \dicIPA{{s}{\textsubring{d}}{l}{\textsci}{\textlengthmark}{\textsubring{d}}{r}{ou}{h}{\textsubring{d}}{\textscy}{\textsubring{r}}} \dicPos{adj}[1]\dicFlx{}[-10] \dicDirectTranslationCS{nesouvislý, přerušovaný, útržkovitý, sporadický} \dicExampleIS{slitróttur vegur} \dicExampleCS{nesouvislá cesta}
\dicEntry[slitum] \dicTerm{slitum} \dicIPA{{s}{\textsubring{d}}{l}{\textsci}{\textlengthmark}{\textsubring{d}}{\textscy}{\textsubring{m}}} \dicPos{v} \dicFlx{ind pf pl 1 pers} \dicLink{slíta}
\dicEntry[slitur] \dicTerm{slitur} \dicIPA{{s}{\textsubring{d}}{l}{\textsci}{\textlengthmark}{\textsubring{d}}{\textscy}{\textsubring{r}}} \dicPos{n}[2] \dicFlx{(‑s, ‑)}[25] \dicDirectTranslationCS{útržek, fragment} \dicExampleIS{slitur af\,/\addthin úr bókinni} \dicExampleCS{fragment z~knihy}
\dicEntry[slíður] \dicTerm{slíður} \dicIPA{{s}{\textsubring{d}}{l}{i}{\textlengthmark}{ð}{\textscy}{\textsubring{r}}} \dicPos{n}[2] \dicFlx{(‑s, ‑)}[25] \dicDirectTranslationCS{pochva, pouzdro (na zbraň ap.)} \dicExampleIS{draga sverð úr slíðrum} \dicExampleCS{vytáhnout meč z~pochvy}
\dicEntry[slíkja] \dicTerm{slíkj|a} \dicIPA{{s}{\textsubring{d}}{l}{i}{\textlengthmark}{\r{\textObardotlessj}}{a}} \dicPos{f}[1] \dicFlx{(‑u, ‑ur)}[7] \dicLink{slikja}
\dicEntry[slíkur] \dicTerm{slíkur} \dicsymFrequent\  \dicIPA{{s}{\textsubring{d}}{l}{i}{\textlengthmark}{\r{g}}{\textscy}{\textsubring{r}}} \dicPos{pron}[13] \dicFlx{dem} \dicSynonym{þvílíkur} \dicDirectTranslationCS{takový} \dicExampleIS{slíkur maður} \dicExampleCS{takový člověk}
\dicEntry[slím] \dicTerm{slím} \dicIPA{{s}{\textsubring{d}}{l}{i}{\textlengthmark}{\textsubring{m}}} \dicPos{n}[2] \dicFlx{(‑s)}[2] \dicDirectTranslationCS{hlen, sliz} \dicExampleIS{slím í hálsi} \dicExampleCS{hlen v~krku}
\dicEntry[slímhúð] \dicTerm{slím··húð} \dicIPA{{s}{\textsubring{d}}{l}{i}{\textlengthmark}{\textsubring{m}}{h}{u}{\texttheta}} \dicPos{f}[7] \dicFlx{(‑ar)}[3] \dicFieldCat{anat.} \dicDirectTranslationCS{sliznice}
\dicEntry[slímhúðarpest] \dicTerm{slím·húðar··pest} \dicIPA{{s}{\textsubring{d}}{l}{i}{\textlengthmark}{\textsubring{m}}{h}{u}{ð}{a}{\textsubring{r}}{p\smash{\textsuperscript{h}}}{\textepsilon}{s}{\textsubring{d}}} \dicPos{f}[7] \dicFlx{(‑ar, ‑ir)}[1] \dicFieldCat{med. } \dicDirectTranslationCS{bovinní virová diarrhea} \dicIndirectTranslationCS{(infekční virové onemocnění skotu)}
\dicEntry[slímkenndur] \dicTerm{slím··kenndur} \dicIPA{{s}{\textsubring{d}}{l}{i}{m}{c\smash{\textsuperscript{h}}}{\textepsilon}{n}{\textsubring{d}}{\textscy}{\textsubring{r}}} \dicPos{adj}[2]\dicFlx{}[-14] \dicDirectTranslationCS{hlenovitý, hleno\-vý, slizovitý}
\dicEntry[slípa] \dicTerm{slíp|a} \dicIPA{{s}{\textsubring{d}}{l}{i}{\textlengthmark}{\textsubring{b}}{a}} \dicPos{v}[1] \dicFlx{(‑aði)}[1] \dicFlx{acc} \textbf{1.} \dicSynonym{fægja} \dicDirectTranslationCS{(vy)leštit, (vy)cídit} \dicExampleIS{slípa gólf} \dicExampleCS{vyleštit podlahu}  \textbf{2.} \dicDirectTranslationCS{(vy)pilovat, (vy)tříbit} \dicExampleIS{slípa málfarið} \dicExampleCS{vytříbit výrazivo};  \dicIdiom{slípast}{ \dicPhraseIS{slípast}} \dicFlx{refl} {\textbf{a.}} \dicDirectTranslationCS{(vy)leštit se, (vy)cídit se};  {\textbf{b.}} \dicDirectTranslationCS{(vy)pilovat se, (vy)tříbit se}
\dicEntry[slít] \dicTerm{slít} \dicIPA{{s}{\textsubring{d}}{l}{i}{\textlengthmark}{\textsubring{d}}} \dicPos{v} \dicFlx{ind praes sg 1 pers} \dicLink{slíta}
\dicEntry[slíta] \dicTerm{slíta} \dicsymFrequent\  \dicIPA{{s}{\textsubring{d}}{l}{i}{\textlengthmark}{\textsubring{d}}{a}} \dicPos{v}[6] \dicFlx{(slít, sleit, slitum, sliti, slitið)}[68] \dicFlx{acc\,/\addthin dat} \textbf{1.} \dicFlx{acc} \dicSynonym{rífa} \dicDirectTranslationCS{roztrhnout, roztrhat, přetrhnout, přetrhat} \dicExampleIS{slíta band} \dicExampleCS{přetrhnout provázek};  \dicPhraseIS{slíta sig lausan} \dicDirectTranslationCS{uvolnit se, vytrhnout se (kůň ap.)}  \textbf{2.} \dicFlx{dat} \dicSynonym*{eyða með notkun} \dicDirectTranslationCS{opotřebit, opotřebovat, obnosit} \dicExampleIS{slíta buxunum} \dicExampleCS{opotřebit kalhoty}  \textbf{3.} \dicFlx{dat} \dicSynonym{enda\smash{\textsuperscript{1}}} \dicDirectTranslationCS{(u)končit, skončit, zakončit} \dicExampleIS{slíta sambandi við e‑n} \dicExampleCS{ukončit s~(kým) kontakt};  \dicPhraseIS{slíta talinu} \dicDirectTranslationCS{ukončit rozhovor};  \dicIdiom{slíta}[frá]{ \dicPhraseIS{slíta sig frá e‑u}} \dicDirectTranslationCS{odtrhnout se od (čeho), odpoutat se od (čeho)} \dicExampleIS{eiga erfitt með að slíta sig frá bókinni} \dicExampleCS{mít potíže s~odtržením se od knížky};  \dicIdiom{slíta}[upp]{ \dicPhraseIS{slíta upp e‑ð}} \dicDirectTranslationCS{vytrhnout (co), trhat (co) (květinu ap.)};  \dicIdiom{slíta}[út]{ \dicPhraseIS{slíta e‑u út}} \dicDirectTranslationCS{obnosit (co)}; { \dicPhraseIS{slíta sér út}} \dicDirectTranslationCS{namoct se, vysílit se}
\dicEntry[slítandi] \dicTerm{slít··andi} \dicIPA{{s}{\textsubring{d}}{l}{i}{\textlengthmark}{\textsubring{d}}{a}{n}{\textsubring{d}}{\textsci}} \dicPos{adj}[13] \dicFlx{indecl}[1] \dicDirectTranslationCS{vyčerpávající, vysilující, namáhavý} \dicExampleIS{slítandi vinna} \dicExampleCS{vyčerpávající práce};  \dicPhraseIS{það er slítandi að (gera e‑ð)} \dicDirectTranslationCS{je vyčerpávající (dělat (co))}
\dicEntry[sljákka] \dicTerm{sljákk|a} \dicIPA{{s}{\textsubring{d}}{l}{j}{au}{h}{\r{g}}{a}} \dicPos{v}[1] \dicFlx{(‑aði)}[88] \dicFlx{impers} \dicPhraseIS{það sljákkar í e‑m} \dicDirectTranslationCS{(kdo) se uklidňuje} \dicExampleIS{Það sljákkaði í honum.} \dicExampleCS{Uklidnil se.}
\dicEntry[sljóleiki] \dicTerm{sljó··leik|i} \dicIPA{{s}{\textsubring{d}}{l}{j}{ou}{\textlengthmark}{l}{ei}{\r{\textObardotlessj}}{\textsci}} \dicPos{m}[1] \dicFlx{(‑a)}[3] \dicSynonym{deyfð} \dicDirectTranslationCS{netečnost, otupělost, apatie}
\dicEntry[sljór] \dicTerm{sljór} \dicIPA{{s}{\textsubring{d}}{l}{j}{ou}{\textlengthmark}{\textsubring{r}}} \dicPos{adj}[4]\dicFlx{}[-1] \textbf{1.} \dicDirectTranslationCS{netečný, apatický}  \textbf{2.} \dicSynonym{bitlaus} \dicDirectTranslationCS{tupý, ztupený};  \dicPhraseIS{sljótt horn} \dicFieldCat{mat.} \dicDirectTranslationCS{tupý úhel}
\dicEntry[sljóvga] \dicTerm{sljóvg|a} \dicIPA{{s}{\textsubring{d}}{l}{j}{ou}{v}{\textbabygamma}{a}} \dicPos{v}[1] \dicFlx{(‑aði)}[1] \dicFlx{acc} \dicSynonym{slæva} \dicDirectTranslationCS{otupit, znecitlivit (smysly ap.)};  \dicIdiom{sljóvgast}{ \dicPhraseIS{sljóvgast}} \dicFlx{refl} \dicDirectTranslationCS{otupit se, znecitlivět, stát se otupělým\,/\addthin znecitlivělým}
\dicEntry[slokkna] \dicTerm{slokkn|a} \dicIPA{{s}{\textsubring{d}}{l}{\textopeno}{h}{\r{g}}{n}{a}} \dicPos{v}[1] \dicFlx{(‑aði)}[7] \dicDirectTranslationCS{zhasnout, vyhasnout, pohasnout (oheň ap.)};  \dicPhraseIS{það slokknar á e‑u} \dicFlx{impers} \dicDirectTranslationCS{(co) zhasíná, (co) vyhasíná}
\dicEntry[slokra] \dicTerm{slokr|a} \dicIPA{{s}{\textsubring{d}}{l}{\textopeno}{\textlengthmark}{\r{g}}{r}{a}} \dicPos{v}[1] \dicFlx{(‑aði)}[1] \dicFlx{acc} \dicDirectTranslationCS{srkat, usrkávat}
\dicEntry[sloppið] \dicTerm{sloppið} \dicIPA{{s}{\textsubring{d}}{l}{\textopeno}{h}{\textsubring{b}}{\textsci}{\texttheta}} \dicPos{v} \dicFlx{supin} \dicLink{sleppa\smash{\textsuperscript{2}}}
\dicEntry[sloppur] \dicTerm{slopp|ur} \dicsymFrequent\  \dicIPA{{s}{\textsubring{d}}{l}{\textopeno}{h}{\textsubring{b}}{\textscy}{\textsubring{r}}} \dicPos{m}[6] \dicFlx{(‑s, ‑ar)}[24] \dicDirectTranslationCS{župan, plášť (pracovní ap.)} \dicExampleIS{vinnusloppur} \dicExampleCS{pracovní plášť}
\dicEntry[slor] \dicTerm{slor} \dicIPA{{s}{\textsubring{d}}{l}{\textopeno}{\textlengthmark}{\textsubring{r}}} \dicPos{n}[2] \dicFlx{(‑s)}[2] \textbf{1.} \dicSynonym*{slím á fiski} \dicDirectTranslationCS{rybí sliz}  \textbf{2.} \dicSynonym{slóg} \dicDirectTranslationCS{vnitřnosti\,/\addthin droby ryby};  \dicIdiom{slor}{ \dicPhraseIS{e‑að er ekkert slor}} \dicDirectTranslationCS{(co) není tak špatné, (co) je fantastické}
\dicEntry[sló] \dicTerm{sló} \dicIPA{{s}{\textsubring{d}}{l}{ou}{\textlengthmark}} \dicPos{v} \dicFlx{ind pf sg 1 pers} \dicLink{slá\smash{\textsuperscript{3}}}
\dicEntry[slóð] \dicTerm{slóð} \dicsymFrequent\  \dicIPA{{s}{\textsubring{d}}{l}{ou}{\textlengthmark}{\texttheta}} \dicPos{f}[7] \dicFlx{(‑ar, ‑ir)}[1] \textbf{1.} \dicSynonym{ferill} \dicDirectTranslationCS{pěšina, stezka, cesta} \dicExampleIS{fjárslóð} \dicExampleCS{ovčí pěšina}  \textbf{2.} \dicPhraseIS{slóðir} \dicFlx{pl} \dicSynonym{svæði} \dicDirectTranslationCS{oblast, lokalita}  \textbf{3.} \dicFieldCat{poč.} \dicDirectTranslationCS{adresa, URL};  \dicIdiom{slóð}{ \dicPhraseIS{fara ótroðnar slóðir}} \dicLangCat{přen.} \dicDirectTranslationCS{jít nevyšlapanými stezkami}
\dicEntry[slóðalegur] \dicTerm{slóða··legur} \dicIPA{{s}{\textsubring{d}}{l}{ou}{\textlengthmark}{ð}{a}{l}{\textepsilon}{\textbabygamma}{\textscy}{\textsubring{r}}} \dicPos{adj}[1]\dicFlx{}[-8] \dicDirectTranslationCS{nedbalý, šlendriánský}
\dicEntry[slóðaskapur] \dicTerm{slóða··skap|ur} \dicIPA{{s}{\textsubring{d}}{l}{ou}{\textlengthmark}{ð}{a}{s}{\r{g}}{a}{\textsubring{b}}{\textscy}{\textsubring{r}}} \dicPos{m}[10] \dicFlx{(‑ar)}[15] \dicDirectTranslationCS{nedbalost, šlendriánství}
\dicEntry[slóði] \dicTerm{slóð|i} \dicIPA{{s}{\textsubring{d}}{l}{ou}{\textlengthmark}{ð}{\textsci}} \dicPos{m}[1] \dicFlx{(‑a, ‑ar)}[1] \textbf{1.} \dicSynonym{stígur} \dicDirectTranslationCS{stezka, pěšina, cesta}  \textbf{2.} \dicSynonym{ónytjungur} \dicDirectTranslationCS{šlendrián, povaleč(ka), darmošlap(ka)}  \textbf{3.} \dicSynonym*{hældrag} \dicDirectTranslationCS{vlečka} \dicIndirectTranslationCS{(prodloužená zadní část ženských šatů vlekoucí se po zemi)}
\dicEntry[slóg] \dicTerm{slóg} \dicIPA{{s}{\textsubring{d}}{l}{ou}{\textlengthmark}{x}} \dicPos{n}[2] \dicFlx{(‑s)}[2] \dicSynonym{slor} \dicDirectTranslationCS{vnitřnosti\,/\addthin droby ryby}
\dicEntry[slógum] \dicTerm{slógum} \dicIPA{{s}{\textsubring{d}}{l}{ou}{\textlengthmark}{\textscy}{\textsubring{m}}} \dicPos{v} \dicFlx{ind pf pl 1 pers} \dicLink{slá\smash{\textsuperscript{3}}}
\dicEntry[slór] \dicTerm{slór} \dicIPA{{s}{\textsubring{d}}{l}{ou}{\textlengthmark}{\textsubring{r}}} \dicPos{n}[2] \dicFlx{(‑s)}[2] \textbf{1.} \dicDirectTranslationCS{ledabylost, nedbalost, šlendriánství}  \textbf{2.} \dicSynonym{seinlæti} \dicDirectTranslationCS{zdlouhavost, liknavost}
\dicEntry[slóra] \dicTerm{slór|a} \dicIPA{{s}{\textsubring{d}}{l}{ou}{\textlengthmark}{r}{a}} \dicPos{v}[1] \dicFlx{(‑aði)}[1] \dicDirectTranslationCS{zahálet, flákat se} \dicExampleIS{slóra vinnulaus helming dagsins} \dicExampleCS{zahálet bez práce půlku dne}
\dicEntry[slóttugur] \dicTerm{slótt··ugur} \dicIPA{{s}{\textsubring{d}}{l}{ou}{h}{\textsubring{d}}{\textscy}{\textbabygamma}{\textscy}{\textsubring{r}}} \dicPos{adj}[1]\dicFlx{}[-8] \dicSynonym{slægur} \dicDirectTranslationCS{potměšilý, záludný}
\dicEntry[Slóvaki] \dicTerm{Slóvak|i} \dicIPA{{s}{\textsubring{d}}{l}{ou}{\textlengthmark}{v}{a}{\r{\textObardotlessj}}{\textsci}} \dicPos{m}[1] \dicFlx{(‑a, ‑ar)}[8] \dicDirectTranslationCS{Slovák, Slovenka}
\dicEntry[Slóvakía] \dicTerm{Slóvakí|a} \dicIPA{{s}{\textsubring{d}}{l}{ou}{\textlengthmark}{v}{a}{\r{\textObardotlessj}}{i}{j}{a}} \dicPos{f}[1] \dicFlx{(‑u)}[6] \dicFieldCat{geog.} \dicDirectTranslationCS{Slovensko}
\dicEntry[slóvakíska] \dicTerm{slóvakísk|a} \dicIPA{{s}{\textsubring{d}}{l}{ou}{\textlengthmark}{v}{a}{\r{\textObardotlessj}}{i}{s}{\r{g}}{a}} \dicPos{f}[1] \dicFlx{(‑u)}[5] \dicDirectTranslationCS{slovenština}
\dicEntry[slóvakískur] \dicTerm{slóvakískur}\dicTerm{, slóvaskur} \dicIPA{{s}\-{\textsubring{d}}\-{l}\-{ou}\-{\textlengthmark}\-{v}\-{a}\-{\r{\textObardotlessj}}\-{i}\-{s}\-{\r{g}}\-{\textscy}\-{\textsubring{r}}\-} \dicPos{adj}[1]\dicFlx{}[-6] \dicDirectTranslationCS{slovenský}
\dicEntry[slóvaskur] \dicTerm{slóvaskur} \dicIPA{{s}{\textsubring{d}}{l}{ou}{\textlengthmark}{v}{a}{s}{\r{g}}{\textscy}{\textsubring{r}}} \dicPos{adj}[1] \dicFlx{(f slóvösk)}[3] \dicLink{slóvakískur}
\dicEntry[Slóveni] \dicTerm{Slóven|i} \dicIPA{{s}{\textsubring{d}}{l}{ou}{\textlengthmark}{v}{\textepsilon}{n}{\textsci}} \dicPos{m}[1] \dicFlx{(‑a, ‑ar)}[1] \dicDirectTranslationCS{Slovinec, Slovinka}
\dicEntry[Slóvenía] \dicTerm{Slóvení|a} \dicIPA{{s}{\textsubring{d}}{l}{ou}{\textlengthmark}{v}{\textepsilon}{n}{i}{j}{a}} \dicPos{f}[1] \dicFlx{(‑u)}[6] \dicFieldCat{geog.} \dicDirectTranslationCS{Slovinsko}
\dicEntry[slóvenska] \dicTerm{slóvensk|a} \dicIPA{{s}{\textsubring{d}}{l}{ou}{\textlengthmark}{v}{\textepsilon}{n}{s}{\r{g}}{a}} \dicPos{f}[1] \dicFlx{(‑u)}[5] \dicDirectTranslationCS{slovinština}
\dicEntry[slóvenskur] \dicTerm{slóvenskur} \dicIPA{{s}{\textsubring{d}}{l}{ou}{\textlengthmark}{v}{\textepsilon}{n}{s}{\r{g}}{\textscy}{\textsubring{r}}} \dicPos{adj}[1]\dicFlx{}[-6] \dicDirectTranslationCS{slovinský}
\dicEntry[slumpa] \dicTerm{slump|a} \dicIPA{{s}{\textsubring{d}}{l}{\textscy}{\textsubring{m}}{\textsubring{b}}{a}} \dicPos{v}[1] \dicFlx{(‑aði)}[1] \dicSynonym{meta} \dicDirectTranslationCS{odhadnout, odhadovat};  \dicPhraseIS{slumpa á e‑ð} \dicDirectTranslationCS{odhadnout (co)} \dicExampleIS{slumpa á stærðina} \dicExampleCS{odhadnout velikost}
\dicEntry[slunginn] \dicTerm{slunginn} \dicIPA{{s}{\textsubring{d}}{l}{u}{\textltailn}{\r{\textObardotlessj}}{\textsci}{\textsubring{n}}} \dicPos{adj}[6]\dicFlx{}[-2] \dicSynonym{kænn} \dicDirectTranslationCS{mazaný, šikovný} \dicExampleIS{slunginn lögfræðingur} \dicExampleCS{mazaný právník}
\dicEntry[sluppum] \dicTerm{sluppum} \dicIPA{{s}{\textsubring{d}}{l}{\textscy}{h}{\textsubring{b}}{\textscy}{\textsubring{m}}} \dicPos{v} \dicFlx{ind pf pl 1 pers} \dicLink{sleppa\smash{\textsuperscript{2}}}
\dicEntry[slúðra] \dicTerm{slúðr|a} \dicIPA{{s}{\textsubring{d}}{l}{u}{ð}{r}{a}} \dicPos{v}[1] \dicFlx{(‑aði)}[1] \dicSynonym*{kjafta frá} \dicDirectTranslationCS{klevetit, pomlouvat, klepat} \dicExampleIS{slúðra um nágrannana} \dicExampleCS{pomlouvat sousedy}
\dicEntry[slúður] \dicTerm{slúður} \dicIPA{{s}{\textsubring{d}}{l}{u}{\textlengthmark}{ð}{\textscy}{\textsubring{r}}} \dicPos{n}[2] \dicFlx{(‑s)}[28] \dicSynonym{söguburður} \dicDirectTranslationCS{kleveta, pomluva, klep} \dicExampleIS{hafa heyrt slúður um hann} \dicExampleCS{slyšet o~něm klep}
\dicEntry[slúppa] \dicTerm{slúpp|a} \dicIPA{{s}{\textsubring{d}}{l}{u}{h}{\textsubring{b}}{a}} \dicPos{f}[1] \dicFlx{(‑u, ‑ur)}[19] \dicFieldCat{nám.} \dicDirectTranslationCS{šalupa}
\dicEntry[slúta] \dicTerm{slút|a} \dicIPA{{s}{\textsubring{d}}{l}{u}{\textlengthmark}{\textsubring{d}}{a}} \dicPos{v}[2] \dicFlx{(‑ti, ‑t)}[62] \dicSynonym{hanga} \dicDirectTranslationCS{čnít, přečnívat, vyčnívat} \dicExampleIS{Kletturinn slútir yfir opið.} \dicExampleCS{Skála přečnívá nad otvorem v~zemi.}
\dicEntry[slydda] \dicTerm{slydd|a} \dicIPA{{s}{\textsubring{d}}{l}{\textsci}{\textsubring{d}}{\textlengthmark}{a}} \dicPos{f}[1] \dicFlx{(‑u, ‑ur)}[19] \dicDirectTranslationCS{plískanice, déšť se sněhem}
\dicEntry[slyngur] \dicTerm{slyngur} \dicIPA{{s}{\textsubring{d}}{l}{i}{\ng}{\r{g}}{\textscy}{\textsubring{r}}} \dicPos{adj}[1]\dicFlx{}[-1] \dicSynonym{leikinn} \dicDirectTranslationCS{šikovný, zručný, dovedný} \dicExampleIS{slyngur veiðimaður} \dicExampleCS{zručný rybář}
\dicEntry[slyppi] \dicTerm{slyppi} \dicIPA{{s}{\textsubring{d}}{l}{\textsci}{h}{\textsubring{b}}{\textsci}} \dicPos{v} \dicFlx{con pf sg 1 pers} \dicLink{sleppa\smash{\textsuperscript{2}}}
\dicEntry[slyppur] \dicTerm{slyppur} \dicIPA{{s}{\textsubring{d}}{l}{\textsci}{h}{\textsubring{b}}{\textscy}{\textsubring{r}}} \dicPos{adj}[1]\dicFlx{}[-1] \dicSynonym{tómhentur} \dicDirectTranslationCS{(jsoucí) s~prázdnýma rukama} \dicIndirectTranslationCS{(jsoucí bez zisku nebo něčeho žádaného)}
\dicEntry[slys] \dicTerm{slys} \dicsymFrequent\  \dicIPA{{s}{\textsubring{d}}{l}{\textsci}{\textlengthmark}{s}} \dicPos{n}[2] \dicFlx{(‑s, ‑)}[5] \dicDirectTranslationCS{nehoda, havárie};  \dicPhraseIS{lenda í slysi} \dicDirectTranslationCS{být účastníkem nehody, mít nehodu, havarovat};  \dicPhraseIS{verða fyrir slysi} \dicDirectTranslationCS{mít nehodu}
\dicEntry[slysabætur] \dicTerm{slysa··bætur} \dicIPA{{s}{\textsubring{d}}{l}{\textsci}{\textlengthmark}{s}{a}{\textsubring{b}}{a}{i}{\textsubring{d}}{\textscy}{\textsubring{r}}} \dicPos{f}[8] \dicFlx{pl}[6] \dicDirectTranslationCS{odškodné (při nehodě)}
\dicEntry[slysadeild] \dicTerm{slysa··deild} \dicIPA{{s}{\textsubring{d}}{l}{\textsci}{\textlengthmark}{s}{a}{\textsubring{d}}{ei}{l}{\textsubring{d}}} \dicPos{f}[7] \dicFlx{(‑ar, ‑ir)}[1] \dicDirectTranslationCS{pohotovost, pohotovostní lékařská služba}
\dicEntry[slysalegur] \dicTerm{slysa··legur} \dicIPA{{s}{\textsubring{d}}{l}{\textsci}{\textlengthmark}{s}{a}{l}{\textepsilon}{\textbabygamma}{\textscy}{\textsubring{r}}} \dicPos{adj}[1]\dicFlx{}[-8] \dicDirectTranslationCS{nešťastný, smolný} \dicExampleIS{Það var slysalegt að týna vegabréfinu.} \dicExampleCS{Byla to smůla ztratit pas.}
\dicEntry[slysatrygging] \dicTerm{slysa··trygg·ing} \dicIPA{{s}{\textsubring{d}}{l}{\textsci}{\textlengthmark}{s}{a}{t\smash{\textsuperscript{h}}}{r}{\textsci}{\r{\textObardotlessj}}{i}{\ng}{\r{g}}} \dicPos{f}[4] \dicFlx{(‑ar, ‑ar)}[5] \dicDirectTranslationCS{úrazové pojištění}
\dicEntry[slysatryggja] \dicTerm{slysa··trygg|ja} \dicIPA{{s}{\textsubring{d}}{l}{\textsci}{\textlengthmark}{s}{a}{t\smash{\textsuperscript{h}}}{r}{\textsci}{\r{\textObardotlessj}}{a}} \dicPos{v}[2] \dicFlx{(‑ði, ‑t)}[92] \dicFlx{acc} \dicDirectTranslationCS{pojistit proti nehodě\,/\addthin úrazu}
\dicEntry[slysavarðstofa] \dicTerm{slysa··varð·stof|a} \dicIPA{{s}{\textsubring{d}}{l}{\textsci}{\textlengthmark}{s}{a}{v}{a}{r}{ð}{s}{\textsubring{d}}{\textopeno}{v}{a}} \dicPos{f}[1] \dicFlx{(‑u, ‑ur)}[7] \dicDirectTranslationCS{ambulance} \dicIndirectTranslationCS{(nemocniční ordinace pro docházející pacienty)}
\dicEntry[slysfarir] \dicTerm{slys··farir} \dicIPA{{s}{\textsubring{d}}{l}{\textsci}{\textlengthmark}{s}{f}{a}{r}{\textsci}{\textsubring{r}}} \dicPos{f} \dicFlx{pl nom} \dicLink{slysför}
\dicEntry[slysför] \dicTerm{slys··|för} \dicIPA{{s}{\textsubring{d}}{l}{\textsci}{\textlengthmark}{s}{f}{\oe}{\textsubring{r}}} \dicPos{f}[7] \dicFlx{(‑farar, ‑farir)}[16] \dicPhraseIS{slysfarir} \dicFlx{pl} \dicSynonym{slys} \dicDirectTranslationCS{(vážný) úraz,  nehoda};  \dicPhraseIS{látast af slysförum} \dicFlx{refl} \dicDirectTranslationCS{zemřít při nehodě}
\dicEntry[slysni] \dicTerm{slysn|i} \dicIPA{{s}{\textsubring{d}}{l}{\textsci}{s}{\textsubring{d}}{n}{\textsci}} \dicPos{f}[3] \dicFlx{(‑i)}[3] \dicSynonym{óheppni} \dicDirectTranslationCS{nešťastná náhoda};  \dicPhraseIS{fyrir slysni} \dicFlx{adv} \dicDirectTranslationCS{nešťastnou náhodou} \dicExampleIS{Þetta gerðist fyrir slysni.} \dicExampleCS{Stalo se to nešťastnou náhodou.}
\dicEntry[slæ] \dicTerm{slæ} \dicIPA{{s}{\textsubring{d}}{l}{a}{i}{\textlengthmark}} \dicPos{v} \dicFlx{ind praes sg 1 pers} \dicLink{slá\smash{\textsuperscript{3}}}
\dicEntry[slæða] \dicTerm{slæð|a\smash{\textsuperscript{1}}} \dicIPA{{s}{\textsubring{d}}{l}{a}{i}{\textlengthmark}{ð}{a}} \dicPos{f}[1] \dicFlx{(‑u, ‑ur)}[19] \dicSynonym{klútur} \dicDirectTranslationCS{šátek, šál(a)}
\dicEntry[slæða] \dicTerm{slæ|ða\smash{\textsuperscript{2}}} \dicIPA{{s}{\textsubring{d}}{l}{a}{i}{\textlengthmark}{ð}{a}} \dicPos{v}[2] \dicFlx{(‑ddi, ‑tt)}[167] \dicFlx{acc} \dicDirectTranslationCS{(pro)zkoumat dno};  \dicIdiom{slæðast}{ \dicPhraseIS{slæðast}} \dicFlx{refl} \dicDirectTranslationCS{zavléci se, vloudit se (chyba do tisku ap.)}
\dicEntry[slægð] \dicTerm{slægð} \dicIPA{{s}{\textsubring{d}}{l}{a}{i}{\textbabygamma}{\texttheta}} \dicPos{f}[4] \dicFlx{(‑ar)}[3] \dicDirectTranslationCS{prohnanost, mazanost, vychytralost}
\dicEntry[slægi] \dicTerm{slægi} \dicIPA{{s}{\textsubring{d}}{l}{a}{i}{j}{\textlengthmark}{\textsci}} \dicPos{v} \dicFlx{con pf sg 1 pers} \dicLink{slá\smash{\textsuperscript{3}}}
\dicEntry[slægja] \dicTerm{slægj|a\smash{\textsuperscript{1}}} \dicIPA{{s}{\textsubring{d}}{l}{a}{i}{j}{\textlengthmark}{a}} \dicPos{f}[1] \dicFlx{(‑u, ‑ur)}[25] \dicDirectTranslationCS{čerstvě posekaná tráva}
\dicEntry[slægja] \dicTerm{slæg|ja\smash{\textsuperscript{2}}} \dicIPA{{s}{\textsubring{d}}{l}{a}{i}{j}{\textlengthmark}{a}} \dicPos{v}[2] \dicFlx{(‑ði, ‑t)}[89] \dicFlx{acc} \dicDirectTranslationCS{(vy)kuchat (rybu)} \dicExampleIS{slægja fiskinn} \dicExampleCS{vykuchat rybu};  \dicIdiom{slægjast}[eftir]{ \dicPhraseIS{slægjast eftir e‑u}} \dicFlx{refl} \dicDirectTranslationCS{pátrat po (čem), pídit se po (čem)} \dicExampleIS{slægja eftir merkilegum bókum} \dicExampleCS{pátrat po zajímavých knihách}
\dicEntry[slægur] \dicTerm{slægur} \dicIPA{{s}{\textsubring{d}}{l}{a}{i}{\textlengthmark}{\textbabygamma}{\textscy}{\textsubring{r}}} \dicPos{adj}[1]\dicFlx{}[-1] \dicSynonym{slóttugur} \dicDirectTranslationCS{prohnaný, mazaný, vykutálený}
\dicEntry[slæmur] \dicTerm{slæmur} \dicsymFrequent\  \dicIPA{{s}{\textsubring{d}}{l}{a}{i}{\textlengthmark}{m}{\textscy}{\textsubring{r}}} \dicPos{adj}[11] \dicFlx{(comp verri, sup verstur)}[4] \textbf{1.} \dicSynonym{illur} \dicDirectTranslationCS{špatný, zlý, hanebný} \dicExampleIS{slæmar fréttir} \dicExampleCS{špatné zprávy}  \textbf{2.} \dicSynonym{lélegur} \dicDirectTranslationCS{špatný, neschopný} \dicExampleIS{slæmur kennari} \dicExampleCS{špatný učitel}  \textbf{3.} \dicSynonym{lasinn} \dicDirectTranslationCS{slabý, churavý} \dicExampleIS{Sjúklingurinn er slæmur í dag.} \dicExampleCS{Pacient se cítí dneska slabě.};  \dicPhraseIS{vera slæmur af e‑u} \dicDirectTranslationCS{trpět (čím) (rýmou ap.)};  \dicPhraseIS{vera slæmur í e‑u} \dicDirectTranslationCS{mít bolesti (čeho) (zad ap.)}
\dicEntry[slæpast] \dicTerm{slæp|ast} \dicIPA{{s}{\textsubring{d}}{l}{a}{i}{\textlengthmark}{\textsubring{b}}{a}{s}{\textsubring{d}}} \dicPos{v}[2] \dicFlx{(‑tist, ‑st)}[193] \dicFlx{refl} \dicSynonym{slóra} \dicDirectTranslationCS{flákat se, flinkat se, poflakovat se} \dicExampleIS{slæpast niðri í bæ} \dicExampleCS{flákat se v~centru města}
\dicEntry[slæpingi] \dicTerm{slæp··ing|i} \dicIPA{{s}{\textsubring{d}}{l}{a}{i}{\textlengthmark}{\textsubring{b}}{i}{\textltailn}{\r{\textObardotlessj}}{\textsci}} \dicPos{m}[1] \dicFlx{(‑ja, ‑jar)}[14] \dicSynonym*{iðjuleysingi} \dicDirectTranslationCS{flákač(ka), povaleč(ka)}
\dicEntry[slættir] \dicTerm{slættir} \dicIPA{{s}{\textsubring{d}}{l}{a}{i}{h}{\textsubring{d}}{\textsci}{\textsubring{r}}} \dicPos{m} \dicFlx{pl nom} \dicLink{sláttur}
\dicEntry[slæva] \dicTerm{slæv|a} \dicIPA{{s}{\textsubring{d}}{l}{a}{i}{\textlengthmark}{v}{a}} \dicPos{v}[2] \dicFlx{(‑ði, ‑t)}[99] \dicFlx{acc} \dicSynonym{deyfa} \dicDirectTranslationCS{otupit, ztlumit (pozornost ap.)} \dicExampleIS{slæva athyglina} \dicExampleCS{otupit pozornost};  \dicIdiom{slævast}{ \dicPhraseIS{slævast}} \dicFlx{refl} \dicDirectTranslationCS{otupit se, ztlumit se}
\dicEntry[slög] \dicTerm{slög} \dicIPA{{s}{\textsubring{d}}{l}{\oe}{\textlengthmark}{x}} \dicPos{n} \dicFlx{pl nom} \dicLink{slag}
\dicEntry[slökkva] \dicTerm{slökk|va} \dicsymFrequent\  \dicIPA{{s}{\textsubring{d}}{l}{\oe}{h}{\r{g}}{v}{a}} \dicPos{v}[4] \dicFlx{(‑ti, ‑t)}[68] \dicFlx{acc} \textbf{1.} \dicDirectTranslationCS{uhasit, zhasnout, zhasit} \dicExampleIS{slökkva ljósin} \dicExampleCS{zhasnout světlo};  \dicPhraseIS{slökkva á e‑u} \dicDirectTranslationCS{zhasnout (co) (lampu ap.)};  \dicPhraseIS{slökkva í e‑u} \dicDirectTranslationCS{uhasit (co), zhasit (co)} \dicExampleIS{slökkva í sígarettunni} \dicExampleCS{uhasit cigaretu}  \textbf{2.} \dicPhraseIS{slökkva á e‑u} \dicDirectTranslationCS{vypnout (co), vypínat (co) (televizi ap.)} \dicExampleIS{slökkva á tölvunni} \dicExampleCS{vypnout počítač};  \dicPhraseIS{það er slökkt á e‑u} \dicFlx{impers} \dicDirectTranslationCS{(co) je vypnuté}  \textbf{3.} \dicDirectTranslationCS{(u)hasit, (u)tišit} \dicExampleIS{slökkva þorsta} \dicExampleCS{uhasit žízeň}  \textbf{4.} \dicPhraseIS{slökkva kalk} \dicDirectTranslationCS{hasit vápno}
\dicEntry[slökkvari] \dicTerm{slökkv··ar|i} \dicIPA{{s}{\textsubring{d}}{l}{\oe}{h}{\r{g}}{v}{a}{r}{\textsci}} \dicPos{m}[1] \dicFlx{(‑a, ‑ar)}[13] \dicDirectTranslationCS{vypínač}
\dicEntry[slökkvibifreið] \dicTerm{slökkvi··bif·reið} \dicIPA{{s}{\textsubring{d}}{l}{\oe}{h}{\r{g}}{v}{\textsci}{\textsubring{b}}{\textsci}{v}{r}{ei}{\texttheta}} \dicPos{f}[4] \dicFlx{(‑ar, ‑ar\,/\addthin ‑ir)}[30] \dicDirectTranslationCS{hasičský vůz}
\dicEntry[slökkvibíll] \dicTerm{slökkvi··bíl|l} \dicIPA{{s}{\textsubring{d}}{l}{\oe}{h}{\r{g}}{v}{\textsci}{\textsubring{b}}{i}{\textsubring{d}}{\textsubring{l}}} \dicPos{m}[6] \dicFlx{(‑s, ‑ar)}[48] \dicDirectTranslationCS{hasičský vůz}
\dicEntry[slökkvilið] \dicTerm{slökkvi··lið} \dicIPA{{s}{\textsubring{d}}{l}{\oe}{h}{\r{g}}{v}{\textsci}{l}{\textsci}{\texttheta}} \dicPos{n}[2] \dicFlx{(‑s, ‑)}[5] \dicDirectTranslationCS{hasiči, hasičský\,/\addthin požární sbor}
\dicEntry[slökkvistöð] \dicTerm{slökkvi··stöð} \dicIPA{{s}{\textsubring{d}}{l}{\oe}{h}{\r{g}}{v}{\textsci}{s}{\textsubring{d}}{\oe}{\texttheta}} \dicPos{f}[6] \dicFlx{(‑var, ‑var)}[1] \dicDirectTranslationCS{hasičská\,/\addthin požární zbrojnice}
\dicEntry[slökkvitæki] \dicTerm{slökkvi··tæki} \dicIPA{{s}{\textsubring{d}}{l}{\oe}{h}{\r{g}}{v}{\textsci}{t\smash{\textsuperscript{h}}}{a}{i}{\r{\textObardotlessj}}{\textsci}} \dicPos{n}[2] \dicFlx{(‑s, ‑)}[16] \dicDirectTranslationCS{hasicí přístroj}
\dicEntry[slökun] \dicTerm{slök|un} \dicIPA{{s}{\textsubring{d}}{l}{\oe}{\textlengthmark}{\r{g}}{\textscy}{\textsubring{n}}} \dicPos{f}[7] \dicFlx{(‑unar)}[12] \textbf{1.} \dicDirectTranslationCS{uvolnění, relaxace} \dicExampleIS{taugaslökun} \dicExampleCS{uvolnění nervů}  \textbf{2.} \dicDirectTranslationCS{détente, uvolnění mezinárodního napětí}
\dicEntry[slökunaræfing] \dicTerm{slökunar··æf·ing} \dicIPA{{s}{\textsubring{d}}{l}{\oe}{\textlengthmark}{\r{g}}{\textscy}{n}{a}{r}{a}{i}{f}{i}{\ng}{\r{g}}} \dicPos{f}[4] \dicFlx{(‑ar, ‑ar)}[5] \dicDirectTranslationCS{relaxační cvičení}
\dicEntry[slöngu] \dicTerm{slöngu} \dicIPA{{s}{\textsubring{d}}{l}{\oe i}{\ng}{\r{g}}{\textscy}} \dicPos{f} \dicFlx{sg gen} \dicLink{slanga}
\dicEntry[slöngva] \dicTerm{slöngv|a\smash{\textsuperscript{1}}} \dicIPA{{s}{\textsubring{d}}{l}{\oe i}{\ng}{\r{g}}{v}{a}} \dicPos{f}[1] \dicFlx{(‑u, ‑ur)}[7] \dicDirectTranslationCS{prak}
\dicEntry[slöngva] \dicTerm{slöngv|a\smash{\textsuperscript{2}}} \dicIPA{{s}{\textsubring{d}}{l}{\oe i}{\ng}{\r{g}}{v}{a}} \dicPos{v}[1] \dicFlx{(‑aði)}[1] \dicFlx{dat} \dicSynonym{kasta} \dicDirectTranslationCS{(vy)mrštit, vrhnout}
\dicEntry[slöngvivaður] \dicTerm{slöngvi··vað|ur} \dicIPA{{s}{\textsubring{d}}{l}{\oe i}{\ng}{\r{g}}{v}{\textsci}{v}{a}{ð}{\textscy}{\textsubring{r}}} \dicPos{m}[10] \dicFlx{(‑ar\,/\addthin ‑s, ‑ir)}[38] \dicDirectTranslationCS{laso}
\dicEntry[slör] \dicTerm{slör} \dicIPA{{s}{\textsubring{d}}{l}{\oe}{\textlengthmark}{\textsubring{r}}} \dicPos{n}[2] \dicFlx{(‑s, ‑)}[5] \dicDirectTranslationCS{závoj, rouška}
\dicEntry[slörk] \dicTerm{slörk} \dicIPA{{s}{\textsubring{d}}{l}{\oe}{\textsubring{r}}{\r{g}}} \dicPos{n} \dicFlx{pl nom} \dicLink{slark}
\dicEntry[sm] \dicTerm{sm} \dicPos{zkr} \dicPhraseIS{sentimetri} \dicDirectTranslationCS{centimetr}
\dicEntry[smakka] \dicTerm{smakk|a} \dicIPA{{s}{m}{a}{h}{\r{g}}{a}} \dicPos{v}[1] \dicFlx{(‑aði)}[13] \dicFlx{acc} \textbf{1.} \dicSynonym*{bragða á} \dicDirectTranslationCS{(o)chutnat, okusit, okoštovat} \dicExampleIS{smakka vínið} \dicExampleCS{ochutnat víno};  \dicPhraseIS{smakka á e‑u} \dicDirectTranslationCS{ochutnat (co)}  \textbf{2.} \dicSynonym*{hafa bragð} \dicDirectTranslationCS{mít chuť, chutnat} \dicExampleIS{e‑að smakkar vel} \dicExampleCS{(co) chutná dobře};  \dicIdiom{smakkast}{ \dicPhraseIS{smakkast}} \dicFlx{refl} \dicDirectTranslationCS{chutnat, šmakovat}
\dicEntry[smala] \dicTerm{smal|a} \dicIPA{{s}{m}{a}{\textlengthmark}{l}{a}} \dicPos{v}[1] \dicFlx{(‑aði)}[13] \dicFlx{dat} \textbf{1.} \dicDirectTranslationCS{shánět (ovce ap.)} \dicExampleIS{smala fénu} \dicExampleCS{shánět ovce}  \textbf{2.} \dicPhraseIS{smala saman e‑m} \dicDirectTranslationCS{shromáždit (koho), shromažďovat (koho)} \dicExampleIS{smala saman áskrifendum að blaðinu} \dicExampleCS{shromažďovat předplatitele novin}
\dicEntry[smalahundur] \dicTerm{smala··hund|ur} \dicIPA{{s}{m}{a}{\textlengthmark}{l}{a}{h}{\textscy}{n}{\textsubring{d}}{\textscy}{\textsubring{r}}} \dicPos{m}[6] \dicFlx{(‑s, ‑ar)}[4] \dicDirectTranslationCS{ovčácký pes}
\begin{xtolerant}{}{1pt}
\dicEntry[smalamaður] \dicTerm{smala··|maður}\addthinS\textls[-5]{\dicIPA{{s}{m}{a}{\textlengthmark}{l}{a}{m}{a}{ð}{\textscy}{\textsubring{r}}}}\addthinS\dicPos{m}[13]\addthinS\dicFlx{(‑manns, ‑menn)}[2]\addthinS\dicSynonym{smali} \dicDirectTranslationCS{ovčák, ovčačka, pastýř(ka)}
\end{xtolerant}
\dicEntry[smalamennska] \dicTerm{smala··mennsk|a} \dicIPA{{s}{m}{a}{\textlengthmark}{l}{a}{m}{\textepsilon}{n}{s}{\r{g}}{a}} \dicPos{f}[1] \dicFlx{(‑u)}[5] \dicDirectTranslationCS{shánění (ovcí ap.)}
\dicEntry[smali] \dicTerm{smal|i} \dicIPA{{s}{m}{a}{\textlengthmark}{l}{\textsci}} \dicPos{m}[1] \dicFlx{(‑a, ‑ar)}[8] \dicSynonym{smalamaður} \dicDirectTranslationCS{ovčák, ovčačka, pastýř(ka)}
\dicEntry[small] \dicTerm{small} \dicIPA{{s}{m}{a}{\textsubring{d}}{\textsubring{l}}} \dicPos{v} \dicFlx{ind pf sg 1 pers} \dicLink{smella\smash{\textsuperscript{3}}}
\dicEntry[smaragður] \dicTerm{smaragð|ur} \dicIPA{{s}{m}{a}{\textlengthmark}{r}{a}{\textbabygamma}{ð}{\textscy}{\textsubring{r}}} \dicPos{m}[6] \dicFlx{(‑s, ‑ar)}[6] \dicFieldCat{geol.} \dicDirectTranslationCS{smaragd}
\dicEntry[smart] \textls[15]{\dicTerm{smart} \dicIPA{{s}{m}{a}{\textsubring{r}}{\textsubring{d}}} \dicPos{adj}[13] \dicFlx{indecl}[1] \dicSynonym{sniðugur} \dicDirectTranslationCS{elegantní, vkusný}}
\dicEntry[smaug] \dicTerm{smaug} \dicIPA{{s}{m}{\oe i}{\textlengthmark}{x}} \dicPos{v} \dicFlx{ind pf sg 1 pers} \dicLink{smjúga}
\dicEntry[smá] \dicTerm{smá} \dicsymFrequent\  \dicIPA{{s}{m}{au}{\textlengthmark}} \dicPos{adv} \dicSynonym{smávegis\smash{\textsuperscript{1}}} \dicDirectTranslationCS{trochu, málo, nepatrně} \dicExampleIS{Hún hló smá.} \dicExampleCS{Trochu se zasmála.}
\dicEntry[smáatriði] \dicTerm{smá··at·riði} \dicsymFrequent\  \dicIPA{{s}{m}{au}{\textlengthmark}{a}{\textsubring{d}}{r}{\textsci}{ð}{\textsci}} \dicPos{n}[2] \dicFlx{(‑s, ‑)}[14] \dicDirectTranslationCS{(po)drobnost, detail, maličkost} \dicExampleIS{Hvert smáatriði skiptir máli.} \dicExampleCS{Na každé maličkosti záleží.};  \dicPhraseIS{í smáatriðum} \dicFlx{adv} \dicDirectTranslationCS{detailně, podrobně}
\dicEntry[smáauglýsing] \dicTerm{smá··aug·lýs·ing} \dicIPA{{s}{m}{au}{\textlengthmark}{\oe i}{\textbabygamma}{l}{i}{s}{i}{\ng}{\r{g}}} \dicPos{f}[4] \dicFlx{(‑ar, ‑ar)}[5] \dicDirectTranslationCS{(malý) inzerát}
\dicEntry[smábarn] \dicTerm{smá··|barn} \dicIPA{{s}{m}{au}{\textlengthmark}{\textsubring{b}}{a}{r}{\textsubring{d}}{\textsubring{n}}} \dicPos{n}[2] \dicFlx{(‑barns, ‑börn)}[8] \dicDirectTranslationCS{děťátko, nemluvně, miminko}
\dicEntry[smáborgari] \dicTerm{smá··borg·ar|i} \dicIPA{{s}{m}{au}{\textlengthmark}{\textsubring{b}}{\textopeno}{r}{\r{g}}{a}{r}{\textsci}} \dicPos{m}[1] \dicFlx{(‑a, ‑ar)}[13] \dicDirectTranslationCS{maloměšťák, maloměšťačka}
\dicEntry[smábær] \dicTerm{smá··bæ|r} \dicIPA{{s}{m}{au}{\textlengthmark}{\textsubring{b}}{a}{i}{\textsubring{r}}} \dicPos{m}[10] \dicFlx{(‑jar, ‑ir)}[25] \dicDirectTranslationCS{maloměsto, městečko, malé\,/\addthin menší město}
\dicEntry[smáfiskur] \dicTerm{smá··fisk|ur} \dicIPA{{s}{m}{au}{\textlengthmark}{f}{\textsci}{s}{\r{g}}{\textscy}{\textsubring{r}}} \dicPos{m}[6] \dicFlx{(‑s, ‑ar)}[10] \dicDirectTranslationCS{rybka, rybička}
\dicEntry[smáfugl] \dicTerm{smá··fugl} \dicIPA{{s}{m}{au}{\textlengthmark}{f}{\textscy}{\r{g}}{\textsubring{l}}} \dicPos{m}[4] \dicFlx{(‑s, ‑ar)}[3] \dicDirectTranslationCS{ptáček, malý pták}
\dicEntry[smáger] \dicTerm{smá··ger} \dicIPA{{s}{m}{au}{\textlengthmark}{\r{\textObardotlessj}}{\textepsilon}{\textsubring{r}}} \dicPos{adj}[5] \dicFlx{(f ‑)}[8] \textbf{1.} \dicDirectTranslationCS{jemný, drobný, delikátní} \dicExampleIS{smáger kona} \dicExampleCS{drobná žena}  \textbf{2.} \dicDirectTranslationCS{mikroskopický}
\dicEntry[smágerður] \dicTerm{smá··gerður} \dicIPA{{s}{m}{au}{\textlengthmark}{\r{\textObardotlessj}}{\textepsilon}{r}{ð}{\textscy}{\textsubring{r}}} \dicPos{adj}[2]\dicFlx{}[-1] \textbf{1.} \dicSynonym{fíngerður} \dicDirectTranslationCS{jemný, delikátní}  \textbf{2.} \dicSynonym{smár} \dicDirectTranslationCS{jemný, drobný, malý}
\dicEntry[smáiðnaður] \dicTerm{smá··iðn·að|ur} \dicIPA{{s}{m}{au}{\textlengthmark}{\textsci}{ð}{n}{a}{ð}{\textscy}{\textsubring{r}}} \dicPos{m}[10] \dicFlx{(‑ar)}[9] \dicDirectTranslationCS{domácký průmysl, živnost}
\dicEntry[smákaka] \dicTerm{smá··|kaka} \dicIPA{{s}{m}{au}{\textlengthmark}{k\smash{\textsuperscript{h}}}{a}{\r{g}}{a}} \dicPos{f}[1] \dicFlx{(‑köku, ‑kökur)}[20] \dicFieldCat{kulin.} \dicDirectTranslationCS{sušenka, koláček}
\begin{xtolerant}{}{1pt}
\dicEntry[smákaupmaður] \dicTerm{smá··kaup·|maður}\addthinSS\textls[-15]{\dicIPA{{s}{m}{au}{\textlengthmark}{k\smash{\textsuperscript{h}}}{\oe i}{\textsubring{b}}{m}{a}{ð}{\textscy}{\textsubring{r}}}} \dicPos{m}[1\hspace{-0.10667em}3]\addthinSS\dicFlx{(‑manns,\addthinSS‑menn)}[2] \dicDirectTranslationCS{maloobchodník, maloobchodnice, kupec, hokynář(ka)}
\end{xtolerant}
\begin{xtolerant}{}{1pt}
\dicEntry[smálegur] \dicTerm{smá··legur}\addthinS\dicIPA{{s}{m}{au}{\textlengthmark}{l}{\textepsilon}{\textbabygamma}{\textscy}{\textsubring{r}}}\addthinS\dicPos{adj}[1]\dicFlx{}[-8]\addthinS\dicDirectTranslationCS{drobný, nepatrný, titěrný}
\end{xtolerant}
\dicEntry[smálest] \dicTerm{smá··lest} \dicIPA{{s}{m}{au}{\textlengthmark}{l}{\textepsilon}{s}{\textsubring{d}}} \dicPos{f}[7] \dicFlx{(‑ar, ‑ir)}[1] \dicDirectTranslationCS{(metrická) tuna} \dicIndirectTranslationCS{(jednotka objemu k~měření kubického objemu lodi)} \dicAntonym{tonn}
\dicEntry[smáletur] \dicTerm{smá··letur} \dicIPA{{s}{m}{au}{\textlengthmark}{l}{\textepsilon}{\textsubring{d}}{\textscy}{\textsubring{r}}} \dicPos{n}[2] \dicFlx{(‑s, ‑)}[25] \dicDirectTranslationCS{malé\,/\addthin drobné písmo}
\dicEntry[smám] \dicTerm{smám} \dicsymFrequent\  \dicIPA{{s}{m}{au}{\textlengthmark}{\textsubring{m}}} \dicPos{adv} \dicPhraseIS{smám saman} \dicFlx{adv} \dicDirectTranslationCS{postupně, krok za krokem}
\begin{xtolerant}{}{6pt}
\dicEntry[smámunasamur] \dicTerm{smá·muna··|samur}\addthinS\textls[-5]{\dicIPA{{s}{m}{au}{\textlengthmark}{m}{\textscy}{n}{a}{s}{a}{m}{\textscy}{\textsubring{r}}}}\addthinS\dicPos{adj}[1]\addthinS\dicFlx{(f ‑söm)}[2] \dicDirectTranslationCS{pun\-tič\-kářský, pedantský}
\end{xtolerant}
\dicEntry[smámunasemi] \dicTerm{smá·muna··sem|i} \dicIPA{{s}{m}{au}{\textlengthmark}{m}{\textscy}{n}{a}{s}{\textepsilon}{m}{\textsci}} \dicPos{f}[2] \dicFlx{(‑i)}[2] \dicDirectTranslationCS{puntičkářství, pedantství, perfekcionismus}
\dicEntry[smámunir] \dicTerm{smá··munir} \dicsymFrequent\  \dicIPA{{s}{m}{au}{\textlengthmark}{m}{\textscy}{n}{\textsci}{\textsubring{r}}} \dicPos{m}[10] \dicFlx{pl}[1] \dicDirectTranslationCS{maličkosti, drobnosti, detaily};  \dicPhraseIS{kippa sér ekki upp við smámuni} \dicDirectTranslationCS{nenervovat se kvůli maličkostem}
\dicEntry[smámæli] \dicTerm{smá··mæli} \dicIPA{{s}{m}{au}{\textlengthmark}{m}{a}{i}{l}{\textsci}} \dicPos{n}[2] \dicFlx{(‑s, ‑)}[14] \dicDirectTranslationCS{šišlání}
\dicEntry[smámæltur] \dicTerm{smá··mæltur} \dicIPA{{s}{m}{au}{\textlengthmark}{m}{a}{i}{\textsubring{l}}{\textsubring{d}}{\textscy}{\textsubring{r}}} \dicPos{adj}[1]\dicFlx{}[-10] \dicDirectTranslationCS{šišlavý, šišlající}
\dicEntry[smán] \dicTerm{smán} \dicIPA{{s}{m}{au}{\textlengthmark}{\textsubring{n}}} \dicPos{f}[7] \dicFlx{(‑ar, ‑ir)}[1] \dicSynonym{skömm\smash{\textsuperscript{1}}} \dicDirectTranslationCS{ostuda, hanba, potupa}
\dicEntry[smána] \dicTerm{smán|a} \dicIPA{{s}{m}{au}{\textlengthmark}{n}{a}} \dicPos{v}[1] \dicFlx{(‑aði)}[1] \dicFlx{acc} \dicDirectTranslationCS{pokořit, pokořovat, ponížit, ponižovat, (po)tupit, (po)hanět} \dicExampleIS{smána minningu drottningarinnar} \dicExampleCS{potupit památku královny}
\dicEntry[smánarlega] \dicTerm{smánar··lega} \dicIPA{{s}{m}{au}{\textlengthmark}{n}{a}{r}{l}{\textepsilon}{\textbabygamma}{a}} \dicPos{adv} \dicDirectTranslationCS{potupně, ostudně}
\begin{xtolerant}{}{1pt}
\dicEntry[smánarlegur] \dicTerm{smánar··legur}\addthinS\dicIPA{{s}{m}{au}{\textlengthmark}{n}{a}{r}{l}{\textepsilon}{\textbabygamma}{\textscy}{\textsubring{r}}}\addthinS\dicPos{adj}[1]\dicFlx{}[-8]\addthinS\dicDirectTranslationCS{potupný, ostudný, ponižující} \dicExampleIS{smánarleg laun} \dicExampleCS{potupný plat}
\end{xtolerant}
\dicEntry[smánaryrði] \dicTerm{smánar··yrði} \dicIPA{{s}{m}{au}{\textlengthmark}{n}{a}{r}{\textsci}{r}{ð}{\textsci}} \dicPos{n}[2] \dicFlx{(‑s, ‑)}[14] \dicSynonym{skammaryrði} \dicDirectTranslationCS{urážka, potupa}
\dicEntry[smáorð] \dicTerm{smá··orð} \dicIPA{{s}{m}{au}{\textlengthmark}{\textopeno}{r}{\texttheta}} \dicPos{n}[2] \dicFlx{(‑s, ‑)}[5] \textbf{1.} \dicDirectTranslationCS{slovíčko}  \textbf{2.} \dicFieldCat{jaz.} \dicIndirectTranslationCS{slovo ze skupiny slov, která se ani neskloňují ani nečasují}
\dicEntry[smápeningar] \dicTerm{smá··pen·ingar} \dicIPA{{s}{m}{au}{\textlengthmark}{p\smash{\textsuperscript{h}}}{\textepsilon}{n}{i}{\ng}{\r{g}}{a}{\textsubring{r}}} \dicPos{m}[6] \dicFlx{pl}[2] \dicDirectTranslationCS{drobné}
\dicEntry[smár] \dicTerm{smár} \dicsymFrequent\  \dicIPA{{s}{m}{au}{\textlengthmark}{\textsubring{r}}} \dicPos{adj}[10] \dicFlx{(comp smærri, sup smæstur)}[14] \textbf{1.} \dicSynonym{lítill} \dicDirectTranslationCS{malý} \dicExampleIS{Loðnan er smár fiskur.} \dicExampleCS{Huňáček je malá ryba.}  \textbf{2.} \dicSynonym{fíngerður} \dicDirectTranslationCS{jemný, drobný}
\dicEntry[smári] \dicTerm{smár|i} \dicIPA{{s}{m}{au}{\textlengthmark}{r}{\textsci}} \dicPos{m}[1] \dicFlx{(‑a, ‑ar)}[1] \textbf{1.} \dicFieldCat{bot.} \dicDirectTranslationCS{jetel} \textit{(l.~{\textLA{Trifolium}})}  \dicsymPhoto\   \textbf{2.} \dicFieldCat{elek.} \dicSynonym{transistor} \dicDirectTranslationCS{tranzistor}
\dicFigure{ds_image_smari_0_2.jpg}{Smári}{Smári - Alvesgaspar, CC BY-SA 3.0}
\dicEntry[smáræði] \dicTerm{smá··ræði} \dicIPA{{s}{m}{au}{\textlengthmark}{r}{a}{i}{ð}{\textsci}} \dicPos{n}[2] \dicFlx{(‑s)}[20] \dicSynonym{smámunir} \dicDirectTranslationCS{maličkost, drobnost}
\dicEntry[smásaga] \dicTerm{smá··|saga} \dicIPA{{s}{m}{au}{\textlengthmark}{s}{a}{\textbabygamma}{a}} \dicPos{f}[1] \dicFlx{(‑sögu, ‑sögur)}[14] \dicFieldCat{lit.} \dicDirectTranslationCS{povídka, novela}
\dicEntry[smásagnahöfundur] \dicTerm{smá·sagna··höf·und|ur} \dicIPA{{s}\-{m}\-{au}\-{\textlengthmark}\-{s}\-{a}\-{\r{g}}\-{n}\-{a}\-{h}\-{\oe}\-{v}\-{\textscy}\-{n}\-{\textsubring{d}}\-{\textscy}\-{\textsubring{r}}\-} \dicPos{m}[6] \dicFlx{(‑ar, ‑ar)}[56] \dicDirectTranslationCS{povídkář(ka), autor(ka) povídek, novelista, novelistka}
\dicEntry[smásala] \textls[15]{\dicTerm{smá··|sala} \dicIPA{{s}{m}{au}{\textlengthmark}{s}{a}{l}{a}} \dicPos{f}[1] \dicFlx{(‑sölu, ‑sölur)}[20] \dicDirectTranslationCS{maloobchod}} \dicAntonym{heildsala}
\dicEntry[smásali] \dicTerm{smá··sal|i} \dicIPA{{s}{m}{au}{\textlengthmark}{s}{a}{l}{\textsci}} \dicPos{m}[1] \dicFlx{(‑a, ‑ar)}[8] \dicDirectTranslationCS{maloobchodník, maloobchodnice, kupec}
\dicEntry[smásálarlegur] \dicTerm{smá·sálar··legur} \dicIPA{{s}{m}{au}{\textlengthmark}{s}{au}{l}{a}{r}{l}{\textepsilon}{\textbabygamma}{\textscy}{\textsubring{r}}} \dicPos{adj}[1]\dicFlx{}[-8] \dicDirectTranslationCS{malicherný, hnidopišský}
\dicEntry[smásjá] \dicTerm{smá··sjá} \dicIPA{{s}{m}{au}{\textlengthmark}{s}{j}{au}} \dicPos{f}[4] \dicFlx{(‑r, ‑r)}[18] \dicDirectTranslationCS{mikroskop};  \dicPhraseIS{vera undir smásjá(nni)} \dicLangCat{přen.} \dicDirectTranslationCS{být pod lupou} \dicIndirectTranslationCS{(pečlivě pozorován)}
\dicEntry[smáskammtalækningar] \dicTerm{smá·skammta··lækn·ingar} \dicIPA{{s}\-{m}\-{au}\-{\textlengthmark}\-{s}\-{\r{g}}\-{a}\-{\textsubring{m}}\-{\textsubring{d}}\-{a}\-{l}\-{a}\-{i}\-{h}\-{\r{g}}\-{n}\-{i}\-{\ng}\-{\r{g}}\-{a}\-{\textsubring{r}}\-} \dicPos{f}[4] \dicFlx{pl}[6] \dicFieldCat{med.} \dicDirectTranslationCS{homeopatie}
\dicEntry[smáskammtalæknir] \textls[5]{\dicTerm{smá·skammta··lækn|ir} \dicIPA{{s}{m}{au}{\textlengthmark}{s}{\r{g}}{a}{\textsubring{m}}{\textsubring{d}}{a}{l}{a}{i}{h}{\r{g}}{n}{\textsci}{\textsubring{r}}}} \dicPos{m}[7] \dicFlx{(‑is, ‑ar)}[1] \dicDirectTranslationCS{homeopat(ka)}
\dicEntry[smásmíði] \dicTerm{smá··smíð|i\smash{\textsuperscript{1}}} \dicIPA{{s}{m}{au}{\textlengthmark}{s}{m}{i}{ð}{\textsci}} \dicPos{f}[2] \dicFlx{(‑i, ‑ar)}[1] \dicLink{smásmíði\smash{\textsuperscript{2}}}
\dicEntry[smásmíði] \dicTerm{smá··smíði\smash{\textsuperscript{2}}}\dicTerm{, smásmíði\smash{\textsuperscript{1}}} \dicIPA{{s}\-{m}\-{au}\-{\textlengthmark}\-{s}\-{m}\-{i}\-{ð}\-{\textsci}\-} \dicPos{n}[2] \dicFlx{(‑s, ‑)}[14] \dicDirectTranslationCS{drobnost, drobná\,/\addthin malá věc} \dicIndirectTranslationCS{(používané zvláště v~záporu)};  \dicPhraseIS{e‑að er ekkert\,/\addthin engin smásmíði} \dicDirectTranslationCS{(co) není zrovna malé} \dicExampleIS{Turninn er ekkert smásmíði.} \dicExampleCS{Věž není zrovna malá.}
\begin{xtolerant}{}{1pt}
\dicEntry[smásmuglegur] \dicTerm{smá·smug··legur}\addthinS\textls[-10]{\dicIPA{{s}{m}{au}{\textlengthmark}{s}{m}{\textscy}{\textbabygamma}{l}{\textepsilon}{\textbabygamma}{\textscy}{\textsubring{r}}}} \dicPos{adj}[1]\dicFlx{}[-8]\addthinS\dicLink{smásmugulegur}
\end{xtolerant}
\dicEntry[smásmugulegur] \textls[5]{\dicTerm{smá·smugu··legur}\dicTerm{, smásmuglegur}\addthin\dicIPA{{s}\-{m}\-{au}\-{\textlengthmark}\-{s}\-{m}\-{\textscy}\-{\textbabygamma}\-{\textscy}\-{l}\-{\textepsilon}\-{\textbabygamma}\-{\textscy}\-{\textsubring{r}}}} \textls[15]{\dicPos{adj}[1]\dicFlx{}[-8] \dicSynonym{smámunasamur} \dicDirectTranslationCS{puntičkářský, pedantský}} \dicExampleIS{smásmugulegar spurningar} \dicExampleCS{puntičkářské otázky}
\dicEntry[smásmygli] \dicTerm{smá··smygl|i} \dicIPA{{s}{m}{au}{\textlengthmark}{s}{m}{\textsci}{\r{g}}{l}{\textsci}} \dicPos{f}[3] \dicFlx{(‑i)}[3] \dicDirectTranslationCS{pedantství, puntičkářství, malichernost}
\dicEntry[smástirni] \dicTerm{smá··stirni} \dicIPA{{s}{m}{au}{\textlengthmark}{s}{\textsubring{d}}{\textsci}{r}{\textsubring{d}}{n}{\textsci}} \dicPos{n}[2] \dicFlx{(‑s, ‑)}[14] \textbf{1.} \dicFieldCat{astro.} \dicDirectTranslationCS{asteroid(a)}  \textbf{2.} \dicLangCat{přen.} \dicDirectTranslationCS{vycházející hvězda, hvězdička (mladá herečka ap.)}
\dicEntry[smástreymi] \dicTerm{smá··streymi} \dicIPA{{s}{m}{au}{\textlengthmark}{s}{\textsubring{d}}{r}{ei}{m}{\textsci}} \dicPos{n}[2] \dicFlx{(‑s, ‑)}[14] \dicDirectTranslationCS{malý příliv}
\dicEntry[smástund] \dicTerm{smá··stund} \dicsymFrequent\  \dicIPA{{s}{m}{au}{\textlengthmark}{s}{\textsubring{d}}{\textscy}{n}{\textsubring{d}}} \dicPos{f}[7] \dicFlx{(‑ar)}[3] \dicDirectTranslationCS{okamžik, chvilka, chvilenka} \dicExampleIS{eftir smástund} \dicExampleCS{za chvilku}
\dicEntry[smásær] \dicTerm{smá··sær} \dicIPA{{s}{m}{au}{\textlengthmark}{s}{a}{i}{\textsubring{r}}} \dicPos{adj}[4]\dicFlx{}[-7] \dicDirectTranslationCS{mikroskopický}
\dicEntry[smásöluverð] \dicTerm{smá·sölu··verð} \dicIPA{{s}{m}{au}{\textlengthmark}{s}{\oe}{l}{\textscy}{v}{\textepsilon}{r}{\texttheta}} \dicPos{n}[2] \dicFlx{(‑s)}[2] \dicDirectTranslationCS{maloobchodní cena}
\dicEntry[smásöluverslun] \dicTerm{smá·sölu··versl|un} \dicIPA{{s}{m}{au}{\textlengthmark}{s}{\oe}{l}{\textscy}{v}{\textepsilon}{\textsubring{r}}{s}{\textsubring{d}}{l}{\textscy}{\textsubring{n}}} \dicPos{f}[7] \dicFlx{(‑unar, ‑anir)}[8] \dicDirectTranslationCS{maloobchod}
\dicEntry[smátt] \dicTerm{smátt} \dicsymFrequent\  \dicIPA{{s}{m}{au}{h}{\textsubring{d}}} \dicPos{adv} \dicFlx{(comp smærra, sup smæst)} \dicSynonym{lítt} \dicDirectTranslationCS{drob\-ně, nepatrně};  \dicPhraseIS{smátt og smátt} \dicFlx{adv} \dicDirectTranslationCS{postupně, krůček po krůčku}
\dicEntry[smátækur] \dicTerm{smá··tækur} \dicIPA{{s}{m}{au}{\textlengthmark}{t\smash{\textsuperscript{h}}}{a}{i}{\r{g}}{\textscy}{\textsubring{r}}} \dicPos{adj}[1]\dicFlx{}[-1] \dicSynonym{naumur} \dicDirectTranslationCS{lakomý, chamtivý, skoupý}
\dicEntry[smávaxinn] \dicTerm{smá··vaxinn} \dicIPA{{s}{m}{au}{\textlengthmark}{v}{a}{x}{s}{\textsci}{\textsubring{n}}} \dicPos{adj}[6]\dicFlx{}[-3] \dicDirectTranslationCS{malý (vzrůstem)} \dicAntonym{hávaxinn}
\dicEntry[smávegis] \dicTerm{smá··vegis\smash{\textsuperscript{1}}} \dicIPA{{s}{m}{au}{\textlengthmark}{v}{ei}{\textsci}{s}} \dicPos{adj}[13] \dicFlx{indecl}[1] \dicSynonym{lítill} \dicDirectTranslationCS{malý, nepatrný} \dicExampleIS{engin smávegis ferð} \dicExampleCS{pořádná cesta}
\dicEntry[smávegis] \dicTerm{smá··vegis\smash{\textsuperscript{2}}} \dicIPA{{s}{m}{au}{\textlengthmark}{v}{ei}{\textsci}{s}} \dicPos{adv} \dicDirectTranslationCS{trochu, nepatrně}
\dicEntry[smávægilegur] \dicTerm{smá·vægi··legur} \dicIPA{{s}{m}{au}{\textlengthmark}{v}{a}{i}{j}{\textsci}{l}{\textepsilon}{\textbabygamma}{\textscy}{\textsubring{r}}} \dicPos{adj}[1]\dicFlx{}[-8] \dicDirectTranslationCS{bezvýznamný, nevýznamný, banální, nepatrný} \dicExampleIS{smávægileg bilun} \dicExampleCS{banální porucha}
\dicEntry[smáþarmar] \dicTerm{smá··þarmar} \dicIPA{{s}{m}{au}{\textlengthmark}{\texttheta}{a}{r}{m}{a}{\textsubring{r}}} \dicPos{m}[6] \dicFlx{pl}[1] \dicFieldCat{anat.} \dicDirectTranslationCS{tenké střevo}
\dicEntry[smáþjóð] \dicTerm{smá··þjóð} \dicIPA{{s}{m}{au}{\textlengthmark}{\texttheta}{j}{ou}{\texttheta}} \dicPos{f}[7] \dicFlx{(‑ar, ‑ir)}[1] \dicDirectTranslationCS{malý národ}
\dicEntry[smáþjófnaður] \dicTerm{smá··þjóf·nað|ur} \dicIPA{{s}{m}{au}{\textlengthmark}{\texttheta}{j}{ou}{\textsubring{b}}{n}{a}{ð}{\textscy}{\textsubring{r}}} \dicPos{m}[10] \dicFlx{(‑ar, ‑ir)}[8] \dicDirectTranslationCS{drobná krádež}
\dicEntry[smáþjófur] \dicTerm{smá··þjóf|ur} \dicIPA{{s}{m}{au}{\textlengthmark}{\texttheta}{j}{ou}{v}{\textscy}{\textsubring{r}}} \dicPos{m}[6] \dicFlx{(‑s, ‑ar)}[22] \dicDirectTranslationCS{drobný zloděj, drobná zlodějka}
\dicEntry[smáögn] \dicTerm{smá··|ögn} \dicIPA{{s}{m}{au}{\textlengthmark}{\oe}{\r{g}}{\textsubring{n}}} \dicPos{f}[7] \dicFlx{(‑agnar, ‑agnir)}[16] \dicDirectTranslationCS{malé množství, malá porce, troška} \dicExampleIS{smáögn af sykri} \dicExampleCS{troška cukru}
\dicEntry[smeðjulegur] \dicTerm{smeðju··legur} \dicIPA{{s}{m}{\textepsilon}{ð}{j}{\textscy}{l}{\textepsilon}{\textbabygamma}{\textscy}{\textsubring{r}}} \dicPos{adj}[1]\dicFlx{}[-8] \dicSynonym{fleðulegur} \dicDirectTranslationCS{podlézavý, servilní}
\dicEntry[smekklaus] \dicTerm{smekk··laus} \dicIPA{{s}{m}{\textepsilon}{h}{\r{g}}{l}{\oe i}{s}} \dicPos{adj}[5]\dicFlx{}[-1] \dicDirectTranslationCS{nevkusný, nevhodný} \dicExampleIS{smekklaus klæðnaður} \dicExampleCS{nevkusné oblečení}
\dicEntry[smekklás] \dicTerm{smekk··lás} \dicIPA{{s}{m}{\textepsilon}{h}{\r{g}}{l}{au}{s}} \dicPos{m}[4] \dicFlx{(‑s, ‑ar)}[25] \dicDirectTranslationCS{západka}
\dicEntry[smekklegur] \dicTerm{smekk··legur} \dicIPA{{s}{m}{\textepsilon}{h}{\r{g}}{l}{\textepsilon}{\textbabygamma}{\textscy}{\textsubring{r}}} \dicPos{adj}[1]\dicFlx{}[-8] \dicDirectTranslationCS{vkusný, vytříbený} \dicExampleIS{smekkleg húsgögn} \dicExampleCS{vkusný nábytek}
\dicEntry[smekkleysi] \dicTerm{smekk··leysi} \dicIPA{{s}{m}{\textepsilon}{h}{\r{g}}{l}{ei}{s}{\textsci}} \dicPos{n}[2] \dicFlx{(‑s)}[20] \dicDirectTranslationCS{nevkusnost, nevkus}
\dicEntry[smekkur] \dicTerm{smekk|ur} \dicIPA{{s}{m}{\textepsilon}{h}{\r{g}}{\textscy}{\textsubring{r}}} \dicPos{m}[9] \dicFlx{(‑s, ‑ir)}[23] \textbf{1.} \dicSynonym*{persónulegt mat} \dicDirectTranslationCS{vkus, vkusnost} \dicExampleIS{góður smekkur} \dicExampleCS{dobrý vkus}  \textbf{2.} \dicDirectTranslationCS{bryndáček}
\dicEntry[smekkvís] \dicTerm{smekk··vís} \dicIPA{{s}{m}{\textepsilon}{h}{\r{g}}{v}{i}{s}} \dicPos{adj}[5]\dicFlx{}[-1] \textbf{1.} \dicDirectTranslationCS{(jsoucí) s~dobrým vkusem}  \textbf{2.} \dicDirectTranslationCS{vkusný, vytříbený, estetický}
\dicEntry[smekkvísi] \dicTerm{smekk··vís|i} \dicIPA{{s}{m}{\textepsilon}{h}{\r{g}}{v}{i}{s}{\textsci}} \dicPos{f}[3] \dicFlx{(‑i)}[3] \dicDirectTranslationCS{vkus, vytříbenost}
\dicEntry[smella] \dicTerm{smell|a\smash{\textsuperscript{1}}} \dicIPA{{s}{m}{\textepsilon}{\textsubring{d}}{l}{a}} \dicPos{f}[1] \dicFlx{(‑u, ‑ur)}[19] \dicDirectTranslationCS{patentka, patentní stiskací knoflík}
\dicEntry[smella] \dicTerm{smell|a\smash{\textsuperscript{2}}} \dicIPA{{s}{m}{\textepsilon}{\textsubring{d}}{l}{a}} \dicPos{v}[2] \dicFlx{(‑ti, ‑t)}[63] \dicFlx{dat} \textbf{1.} \dicDirectTranslationCS{cvaknout, cvakat, klapnout, klapat, lusknout, luskat} \dicExampleIS{smella fingrunum} \dicExampleCS{lusknout prsty}  \textbf{2.} \dicSynonym*{loka e‑u með smellu} \dicDirectTranslationCS{zacvaknout\,/\addthin zapnout (na patentky) (sukni ap.)} \dicExampleIS{smella blússunni} \dicExampleCS{zapnout blůzu}  \textbf{3.} \dicFieldCat{poč.} \dicDirectTranslationCS{kliknout, klikat, ťuknout, ťukat (myší ap.)} \dicExampleIS{smella músinni} \dicExampleCS{kliknout myší};  \dicIdiom{smella}[af]{ \dicPhraseIS{smella af}} \dicDirectTranslationCS{zmáčknout\,/\addthin stisknout spoušť (u~fotoaparátu ap.)};  \dicIdiom{smella}[sér]{ \dicPhraseIS{smella sér e‑s staðar}} \dicDirectTranslationCS{vypravit se (kam), udělat si výlet (kam), podniknout výlet (kam)}
\dicEntry[smella] \dicTerm{smella\smash{\textsuperscript{3}}} \dicIPA{{s}{m}{\textepsilon}{\textsubring{d}}{l}{a}} \dicPos{v}[6] \dicFlx{(smell, small, smullum, smylli, smollið)}[28] \textbf{1.} \dicSynonym{skella\smash{\textsuperscript{2}}} \dicDirectTranslationCS{bouchnout, bouchat, bušit, tlouct} \dicExampleIS{Hurðin small í lás.} \dicExampleCS{Dveře se zabouchly.};  \dicPhraseIS{það smellur í e‑u} \dicFlx{impers} \dicDirectTranslationCS{(co) bouchá, (co) práská} \dicExampleIS{Það small í svipunni.} \dicExampleCS{Bič zapráskal.}  \textbf{2.} \dicDirectTranslationCS{zapadnout, zapadat, přilehnout, přiléhat};  \dicPhraseIS{smella saman} \dicDirectTranslationCS{pasovat, sedět} \dicExampleIS{Endarnir smella saman.} \dicExampleCS{Konce do sebe zapadají.}
\dicEntry[smellinn] \dicTerm{smellinn} \dicIPA{{s}{m}{\textepsilon}{\textsubring{d}}{l}{\textsci}{\textsubring{n}}} \dicPos{adj}[6]\dicFlx{}[-2] \dicSynonym{hnyttinn} \dicDirectTranslationCS{vtipný, trefný, lapidární} \dicExampleIS{smellið svar} \dicExampleCS{vtipná odpověď}
\dicEntry[smellur] \dicTerm{smell|ur} \dicIPA{{s}{m}{\textepsilon}{\textsubring{d}}{l}{\textscy}{\textsubring{r}}} \dicPos{m}[9] \dicFlx{(‑s, ‑ir)}[4] \textbf{1.} \dicSynonym*{smellandi hljóð} \dicDirectTranslationCS{cvaknutí, klapnutí, lusknutí}  \textbf{2.} \dicSynonym*{vinsælt lag} \dicDirectTranslationCS{hit, úspěšná píseň}  \textbf{3.} \dicFieldCat{poč.} \dicDirectTranslationCS{ťuknutí, klepnutí, kliknutí (myší ap.)}
\dicEntry[smergill] \dicTerm{smerg|ill} \dicIPA{{s}{m}{\textepsilon}{r}{\r{\textObardotlessj}}{\textsci}{\textsubring{d}}{\textsubring{l}}} \dicPos{m}[6] \dicFlx{(‑ils, ‑lar)}[35] \dicFieldCat{geol.} \dicDirectTranslationCS{smirek}
\dicEntry[smetti] \dicTerm{smetti} \dicIPA{{s}{m}{\textepsilon}{h}{\textsubring{d}}{\textsci}} \dicPos{n}[2] \dicFlx{(‑s, ‑)}[14] \dicLangCat{han.} \dicDirectTranslationCS{držka, ksicht}
\dicEntry[smeygja] \dicTerm{smeyg|ja} \dicsymFrequent\  \dicIPA{{s}{m}{ei}{j}{\textlengthmark}{a}} \dicPos{v}[2] \dicFlx{(‑ði, ‑t)}[89] \dicFlx{dat} \textbf{1.} \dicDirectTranslationCS{vklouznout, vsunout} \dicIndirectTranslationCS{(rychle a~snadno)};  \dicPhraseIS{smeygja sér í e‑ð} \dicDirectTranslationCS{nasadit si (co), navléct si (co)} \dicExampleIS{smeygja sér í buxurnar} \dicExampleCS{navléct si kalhoty};  \dicPhraseIS{smeygja sér úr e‑u} \dicDirectTranslationCS{svléknout si (co)} \dicExampleIS{smeygja sér úr sokkunum} \dicExampleCS{svléknout si ponožky}  \textbf{2.} \dicDirectTranslationCS{prosmýknout se (dveřmi ap.)}
\dicEntry[smeykur] \dicTerm{smeykur} \dicIPA{{s}{m}{ei}{\textlengthmark}{\r{g}}{\textscy}{\textsubring{r}}} \dicPos{adj}[1]\dicFlx{}[-1] \dicSynonym{banginn} \dicDirectTranslationCS{mající obavy, ustaraný, nervózní}
\dicEntry[smiðja] \dicTerm{smiðj|a} \dicIPA{{s}{m}{\textsci}{ð}{j}{a}} \dicPos{f}[1] \dicFlx{(‑u, ‑ur)}[7] \textbf{1.} \dicSynonym{járnsmiðja} \dicDirectTranslationCS{kovárna}  \textbf{2.} \dicSynonym{verkstæði} \dicDirectTranslationCS{dílna};  \dicPhraseIS{leita í smiðju til e‑rs} \dicLangCat{přen.} \dicDirectTranslationCS{vyhledat (odbornou) radu u~(koho)}
\dicEntry[smiður] \dicTerm{smið|ur} \dicIPA{{s}{m}{\textsci}{\textlengthmark}{ð}{\textscy}{\textsubring{r}}} \dicPos{m}[9] \dicFlx{(‑s, ‑ir)}[11] \textbf{1.} \dicDirectTranslationCS{řemeslník, řemeslnice, dělník, dělnice} \dicExampleIS{góður smiður} \dicExampleCS{dobrý řemeslník}  \textbf{2.} \dicSynonym{trésmiður} \dicDirectTranslationCS{tesař(ka)}  \textbf{3.} \dicSynonym{járnsmiður} \dicDirectTranslationCS{kovář(ka)};  \dicPhraseIS{hengja bakara fyrir smið} \dicLangCat{přen.} \dicDirectTranslationCS{pověsit pekaře místo tesaře} \dicIndirectTranslationCS{(potrestat člověka za něco, co udělal někdo jiný)}
\dicEntry[smit] \dicTerm{smit} \dicIPA{{s}{m}{\textsci}{\textlengthmark}{\textsubring{d}}} \dicPos{n}[2] \dicFlx{(‑s)}[2] \dicFieldCat{med.} \dicDirectTranslationCS{nákaza, nakažení, infikování} \dicExampleIS{verða fyrir smiti} \dicExampleCS{nakazit se}
\dicEntry[smita] \dicTerm{smit|a} \dicIPA{{s}{m}{\textsci}{\textlengthmark}{\textsubring{d}}{a}} \dicPos{v}[1] \dicFlx{(‑aði)}[1] \dicFlx{acc} \textbf{1.} \dicDirectTranslationCS{nakazit, infikovat (rýmou ap.)};  \dicPhraseIS{smita e‑n af e‑u} \dicDirectTranslationCS{nakazit (koho čím)} \dicExampleIS{smita hana af inflúensu} \dicExampleCS{nakazit ji chřipkou}  \textbf{2.} \dicLangCat{přen.} \dicDirectTranslationCS{nakazit, infikovat (optimismem ap.)};  \dicPhraseIS{smita e‑n af e‑u} \dicLangCat{přen.} \dicDirectTranslationCS{nakazit (koho čím)};  \dicPhraseIS{e‑að smitar út frá sér} \dicLangCat{přen.} \dicDirectTranslationCS{(co) působí nakažlivě};  \dicIdiom{smitast}{ \dicPhraseIS{smitast}} \dicFlx{refl} {\textbf{a.}} \dicDirectTranslationCS{nakazit se, infikovat se} \dicExampleIS{smitast af e‑u} \dicExampleCS{nakazit se (čím)};  {\textbf{b.}} \dicLangCat{přen.} \dicDirectTranslationCS{nakazit se, infikovat se (radostí ap.)}
\dicEntry[smitaður] \dicTerm{smit··|aður} \dicIPA{{s}{m}{\textsci}{\textlengthmark}{\textsubring{d}}{a}{ð}{\textscy}{\textsubring{r}}} \dicPos{adj}[3] \dicFlx{(f ‑uð)}[3] \dicDirectTranslationCS{nakažený, infikovaný} \dicExampleIS{smitaður af e‑u} \dicExampleCS{nakažený (čím)}
\dicEntry[smitandi] \dicTerm{smit··andi} \dicIPA{{s}{m}{\textsci}{\textlengthmark}{\textsubring{d}}{a}{n}{\textsubring{d}}{\textsci}} \dicPos{adj}[13] \dicFlx{indecl}[1] \textbf{1.} \dicDirectTranslationCS{nakažlivý, infekční, přenosný} \dicExampleIS{smitandi sjúkdómur} \dicExampleCS{nakažlivé onemocnění}  \textbf{2.} \dicDirectTranslationCS{nakažlivý (smích ap.)} \dicExampleIS{Hlátur hennar var smitandi.} \dicExampleCS{Její smích byl nakažlivý.}
\dicEntry[smitberi] \dicTerm{smit··ber|i} \dicIPA{{s}{m}{\textsci}{\textlengthmark}{\textsubring{d}}{\textsubring{b}}{\textepsilon}{r}{\textsci}} \dicPos{m}[1] \dicFlx{(‑a, ‑ar)}[1] \dicDirectTranslationCS{přenašeč(ka), bacilonosič(ka)}
\dicEntry[smithætta] \dicTerm{smit··hætt|a} \dicIPA{{s}{m}{\textsci}{\textlengthmark}{\textsubring{d}}{h}{a}{i}{\textsubring{d}}{a}} \dicPos{f}[1] \dicFlx{(‑u)}[5] \dicFieldCat{med.} \dicDirectTranslationCS{nebezpečí\,/\addthin riziko nakažení}
\dicEntry[smitnæmi] \dicTerm{smit··næmi} \dicIPA{{s}{m}{\textsci}{\textlengthmark}{\textsubring{d}}{n}{a}{i}{m}{\textsci}} \dicPos{n}[2] \dicFlx{(‑s)}[20] \dicDirectTranslationCS{nakažlivost, infekčnost, přenosnost}
\dicEntry[smitnæmur] \dicTerm{smit··næmur} \dicIPA{{s}{m}{\textsci}{\textlengthmark}{\textsubring{d}}{n}{a}{i}{m}{\textscy}{\textsubring{r}}} \dicPos{adj}[1]\dicFlx{}[-1] \dicFieldCat{med. } \dicDirectTranslationCS{nakažlivý, infekční, přenosný}
\dicEntry[smitsjúkdómur] \dicTerm{smit··sjúk·dóm|ur} \dicIPA{{s}{m}{\textsci}{\textlengthmark}{\textsubring{d}}{s}{j}{u}{\r{g}}{\textsubring{d}}{ou}{m}{\textscy}{\textsubring{r}}} \dicPos{m}[6] \dicFlx{(‑s, ‑ar)}[10] \dicDirectTranslationCS{nakažlivá nemoc, infekční onemocnění}
\dicEntry[smitun] \dicTerm{smit|un} \dicIPA{{s}{m}{\textsci}{\textlengthmark}{\textsubring{d}}{\textscy}{\textsubring{n}}} \dicPos{f}[7] \dicFlx{(‑unar)}[9] \dicFieldCat{med.} \dicDirectTranslationCS{nákaza, infekce} \dicExampleIS{koma í veg fyrir smitun} \dicExampleCS{zabránit infekci}
\dicEntry[smíð] \dicTerm{smíð} \dicsymFrequent\  \dicIPA{{s}{m}{i}{\textlengthmark}{\texttheta}} \dicPos{f}[4] \dicFlx{(‑ar, ‑ar)}[1] \dicDirectTranslationCS{(řemeslný) výrobek} \dicExampleIS{fagrar smíðar} \dicExampleCS{krásné výrobky}
\dicEntry[smíða] \dicTerm{smíð|a} \dicsymFrequent\  \dicIPA{{s}{m}{i}{\textlengthmark}{ð}{a}} \dicPos{v}[1] \dicFlx{(‑aði)}[1] \dicFlx{acc} \dicDirectTranslationCS{vyrobit, vyrábět, zhotovit, zhotovovat, postavit, stavět} \dicExampleIS{smíða stól úr tré} \dicExampleCS{vyrobit židli ze dřeva}
\dicEntry[smíðatól] \dicTerm{smíða··tól} \dicIPA{{s}{m}{i}{\textlengthmark}{ð}{a}{t\smash{\textsuperscript{h}}}{ou}{\textsubring{l}}} \dicPos{n}[2] \dicFlx{(‑s, ‑)}[5] \dicDirectTranslationCS{(pracovní) nástroj\,/\addthin nářadí}
\dicEntry[smíði] \dicTerm{smíð|i\smash{\textsuperscript{1}}} \dicIPA{{s}{m}{i}{\textlengthmark}{ð}{\textsci}} \dicPos{f}[3] \dicFlx{(‑i)}[3] \dicLink{smíði\smash{\textsuperscript{2}}}
\dicEntry[smíði] \dicTerm{smíði\smash{\textsuperscript{2}}}\dicTerm{, smíði\smash{\textsuperscript{1}}} \dicIPA{{s}\-{m}\-{i}\-{\textlengthmark}\-{ð}\-{\textsci}\-} \dicPos{n}[2] \dicFlx{(‑s, ‑)}[14] \textbf{1.} \dicSynonym*{smíðavinna} \dicDirectTranslationCS{vyrábění, zhotovení, stavění} \dicExampleIS{húsasmíði} \dicExampleCS{stavění domu}  \textbf{2.} \dicDirectTranslationCS{práce se dřevem} \dicIndirectTranslationCS{(školní předmět)}  \textbf{3.} \dicSynonym{smíð} \dicDirectTranslationCS{(řemeslný) výrobek\,/\addthin produkt}
\dicEntry[smíðisgripur] \dicTerm{smíðis··grip|ur} \dicIPA{{s}{m}{i}{\textlengthmark}{ð}{\textsci}{s}{\r{g}}{r}{\textsci}{\textsubring{b}}{\textscy}{\textsubring{r}}} \dicPos{m}[9] \dicFlx{(‑s, ‑ir)}[8] \dicDirectTranslationCS{(řemeslný) výrobek\,/\addthin výtvor}
\dicEntry[smjaðra] \textls[15]{\dicTerm{smjaðr|a} \dicIPA{{s}{m}{j}{a}{ð}{r}{a}} \dicPos{v}[1] \dicFlx{(‑aði)}[13] \dicDirectTranslationCS{lichotit, pochlebovat};}  \dicPhraseIS{smjað\-ra fyrir e‑m} \dicDirectTranslationCS{lichotit (komu)}
\dicEntry[smjaður] \dicTerm{smjaður} \dicIPA{{s}{m}{j}{a}{\textlengthmark}{ð}{\textscy}{\textsubring{r}}} \dicPos{n}[2] \dicFlx{(‑s)}[28] \dicDirectTranslationCS{lichocení, pochlebování} \dicExampleIS{vera með smjaður} \dicExampleCS{pochlebovat}
\dicEntry[smjatt] \dicTerm{smjatt} \dicIPA{{s}{m}{j}{a}{h}{\textsubring{d}}} \dicPos{n}[2] \dicFlx{(‑s)}[2] \dicDirectTranslationCS{mlasknutí, (za)mlaskání}
\dicEntry[smjatta] \dicTerm{smjatt|a} \dicIPA{{s}{m}{j}{a}{h}{\textsubring{d}}{a}} \dicPos{v}[1] \dicFlx{(‑aði)}[13] \dicSynonym{kjamsa} \dicDirectTranslationCS{mlasknout, (za)mlaskat} \dicExampleIS{smjatta á matnum} \dicExampleCS{mlaskat u~jídla}
\dicEntry[smjúga] \dicTerm{smjúga} \dicIPA{{s}{m}{j}{u}{\textlengthmark}{a}} \dicPos{v}[6] \dicFlx{(smýg, smaug, smugum, smygi, smogið)}[52] \dicDirectTranslationCS{vniknout, vnikat, proniknout, pronikat} \dicExampleIS{smjúga inn í holuna} \dicExampleCS{vniknout do díry};  \dicPhraseIS{smjúga gegnum e‑ð} \dicDirectTranslationCS{proniknout naskrz (čím)}
\dicEntry[smjör] \dicTerm{smjör} \dicIPA{{s}{m}{j}{\oe}{\textlengthmark}{\textsubring{r}}} \dicPos{n}[2] \dicFlx{(‑s)}[2] \dicDirectTranslationCS{máslo};  \dicPhraseIS{áfram með smjörið} \dicDirectTranslationCS{dál, pokračuj (ve čtení ap.)};  \dicPhraseIS{e‑s staðar drýpur smjör af hverju strái} \dicLangCat{přen.} \dicDirectTranslationCS{(kde) to oplývá mlékem a~strdím};  \dicPhraseIS{verða eins og bráðið smjör} \dicLangCat{přen.} \dicDirectTranslationCS{rozplývat se\,/\addthin být jako máslo}
\dicEntry[smjördeig] \dicTerm{smjör··deig} \dicIPA{{s}{m}{j}{\oe}{r}{\textsubring{d}}{ei}{x}} \dicPos{n}[2] \dicFlx{(‑s)}[2] \dicDirectTranslationCS{listové těsto}
\dicEntry[smjörgras] \dicTerm{smjör··|gras} \dicIPA{{s}{m}{j}{\oe}{r}{\r{g}}{r}{a}{s}} \dicPos{n}[2] \dicFlx{(‑grass, ‑grös)}[8] \dicFieldCat{bot.} \dicDirectTranslationCS{lepnice alpská} \textit{(l.~{\textLA{Bartsia alpina}})}  \dicsymPhoto\ 
\dicFigure{32490.jpg}{Smjörgras}{Smjörgras - Blanka Brandová, Biolib, Copyright/CC-BY-NC}
\dicEntry[smjörkúpa] \dicTerm{smjör··kúp|a} \dicIPA{{s}{m}{j}{\oe}{\textsubring{r}}{k\smash{\textsuperscript{h}}}{u}{\textsubring{b}}{a}} \dicPos{f}[1] \dicFlx{(‑u, ‑ur)}[19] \dicDirectTranslationCS{máselník, máslenka}
\dicEntry[smjörlíki] \dicTerm{smjör··líki} \dicIPA{{s}{m}{j}{\oe}{r}{l}{i}{\r{\textObardotlessj}}{\textsci}} \dicPos{n}[2] \dicFlx{(‑s, ‑)}[16] \dicDirectTranslationCS{margarín}
\dicEntry[smjörpappír] \dicTerm{smjör··pappír} \dicIPA{{s}{m}{j}{\oe}{\textsubring{r}}{p\smash{\textsuperscript{h}}}{a}{h}{\textsubring{b}}{i}{\textsubring{r}}} \dicPos{m}[4] \dicFlx{(‑s)}[18] \dicDirectTranslationCS{pečicí papír, papír na pečení, svačinový papír}
\dicEntry[smjörstykki] \dicTerm{smjör··stykki} \dicIPA{{s}{m}{j}{\oe}{\textsubring{r}}{s}{\textsubring{d}}{\textsci}{h}{\r{\textObardotlessj}}{\textsci}} \dicPos{n}[2] \dicFlx{(‑s, ‑)}[16] \dicDirectTranslationCS{balení másla}
\dicEntry[smogið] \dicTerm{smogið} \dicIPA{{s}{m}{\textopeno i}{j}{\textlengthmark}{\textsci}{\texttheta}} \dicPos{v} \dicFlx{supin} \dicLink{smjúga}
\dicEntry[smokka] \dicTerm{smokk|a} \dicIPA{{s}{m}{\textopeno}{h}{\r{g}}{a}} \dicPos{v}[1] \dicFlx{(‑aði)}[1] \dicLangCat{zast.} \dicSynonym{smeygja} \dicDirectTranslationCS{nasoukat, nasunout, vsunout} \dicExampleIS{smokka lykkjunni upp á staurinn} \dicExampleCS{nasunout poutko na sloupek}
\dicEntry[smokkfiskur] \dicTerm{smokk··fisk|ur} \dicIPA{{s}{m}{\textopeno}{h}{\r{g}}{f}{\textsci}{s}{\r{g}}{\textscy}{\textsubring{r}}} \dicPos{m}[6] \dicFlx{(‑s, ‑ar)}[10] \dicFieldCat{zool.} \dicDirectTranslationCS{krakatice, oliheň, sépie} \textit{(l.~{\textLA{Teuthida}})}
\dicEntry[smokkur] \dicTerm{smokk|ur} \dicIPA{{s}{m}{\textopeno}{h}{\r{g}}{\textscy}{\textsubring{r}}} \dicPos{m}[6] \dicFlx{(‑s, ‑ar)}[22] \textbf{1.} \dicSynonym{getnaðarvörn} \dicDirectTranslationCS{kondom, prezervativ}  \textbf{2.} \dicFieldCat{zool.} \dicDirectTranslationCS{hlavonožec} \textit{(l.~{\textLA{Cephalopoda}})}
\dicEntry[smokra] \dicTerm{smokr|a} \dicIPA{{s}{m}{\textopeno}{\textlengthmark}{\r{g}}{r}{a}} \dicPos{v}[1] \dicFlx{(‑aði)}[1] \dicFlx{acc} \dicSynonym{smokka} \dicDirectTranslationCS{vsunout, nasunout, nasoukat} \dicExampleIS{smokra sér í buxurnar} \dicExampleCS{nasoukat se do kalhot}
\dicEntry[smollið] \dicTerm{smollið} \dicIPA{{s}{m}{\textopeno}{\textsubring{d}}{l}{\textsci}{\texttheta}} \dicPos{v} \dicFlx{supin} \dicLink{smella\smash{\textsuperscript{3}}}
\dicEntry[smotterí] \dicTerm{smotterí}\dicTerm{, smottirí} \dicIPA{{s}\-{m}\-{\textopeno}\-{h}\-{\textsubring{d}}\-{\textepsilon}\-{r}\-{i}\-} \dicPos{n}[2] \dicFlx{(‑s)}[2] \dicDirectTranslationCS{maličkost, drobnost, trivialita, triviálnost}
\dicEntry[smottirí] \dicTerm{smottirí} \dicIPA{{s}{m}{\textopeno}{h}{\textsubring{d}}{\textsci}{r}{i}} \dicPos{n}[2] \dicFlx{(‑s)}[2] \dicLink{smotterí}
\dicEntry[smók] \dicTerm{smók} \dicIPA{{s}{m}{ou}{\textlengthmark}{\r{g}}} \dicPos{m}[4] \dicFlx{(‑s, ‑ar)}[10] \dicLink{smókur}
\dicEntry[smóking] \dicTerm{smóking} \dicIPA{{s}{m}{ou}{\textlengthmark}{\r{\textObardotlessj}}{i}{\ng}{\r{g}}} \dicPos{m}[4] \dicFlx{(‑s, ‑ar)}[14] \dicDirectTranslationCS{smoking, smokink}
\dicEntry[smókur] \dicTerm{smók|ur}\dicTerm{, smók} \dicIPA{{s}\-{m}\-{ou}\-{\textlengthmark}\-{\r{g}}\-{\textscy}\-{\textsubring{r}}\-} \dicPos{m}[6] \dicFlx{(‑s, ‑ar)}[24] \dicLangCat{hovor.} \dicSynonym{sígaretta} \dicDirectTranslationCS{pokouření, cigareta};  \dicPhraseIS{fá sér smók} \dicDirectTranslationCS{pokouřit si, zakouřit si, zapálit si}
\dicEntry[smuga] \dicTerm{smug|a} \dicIPA{{s}{m}{\textscy}{\textlengthmark}{\textbabygamma}{a}} \dicPos{f}[1] \dicFlx{(‑u, ‑ur)}[19] \textbf{1.} \dicSynonym{glufa} \dicDirectTranslationCS{skulina, štěrbina} \dicExampleIS{Það er ofurlítil smuga á glugganum.} \dicExampleCS{V~okně je maličká skulina.}  \textbf{2.} \dicLangCat{přen.} \dicSynonym{viðlit} \dicDirectTranslationCS{mezera, skulina (v~zákoně ap.)}
\dicEntry[smugum] \dicTerm{smugum} \dicIPA{{s}{m}{\textscy}{\textlengthmark}{\textbabygamma}{\textscy}{\textsubring{m}}} \dicPos{v} \dicFlx{ind pf pl 1 pers} \dicLink{smjúga}
\dicEntry[smullum] \dicTerm{smullum} \dicIPA{{s}{m}{\textscy}{\textsubring{d}}{l}{\textscy}{\textsubring{m}}} \dicPos{v} \dicFlx{ind pf pl 1 pers} \dicLink{smella\smash{\textsuperscript{3}}}
\dicEntry[smurbrauð] \dicTerm{smur··brauð} \dicIPA{{s}{m}{\textscy}{r}{\textsubring{b}}{r}{\oe i}{\texttheta}} \dicPos{n}[2] \dicFlx{(‑s)}[2] \dicDirectTranslationCS{obložený chléb\,/\addthin chlebíček}
\dicEntry[smurbrauðsstofa] \dicTerm{smur·brauðs··stof|a} \dicIPA{{s}{m}{\textscy}{r}{\textsubring{b}}{r}{\oe i}{ð}{s}{\textsubring{d}}{\textopeno}{v}{a}} \dicPos{f}[1] \dicFlx{(‑u, ‑ur)}[7] \dicDirectTranslationCS{lahůdkářství, bufet, bistro}
\dicEntry[smurði] \dicTerm{smurði} \dicIPA{{s}{m}{\textscy}{r}{ð}{\textsci}} \dicPos{v} \dicFlx{ind pf sg 1 pers} \dicLink{smyrja}
\dicEntry[smurðum] \dicTerm{smurðum} \dicIPA{{s}{m}{\textscy}{r}{ð}{\textscy}{\textsubring{m}}} \dicPos{v} \dicFlx{ind pf pl 1 pers} \dicLink{smyrja}
\dicEntry[smurning] \dicTerm{smur··ning} \dicIPA{{s}{m}{\textscy}{r}{\textsubring{d}}{n}{i}{\ng}{\r{g}}} \dicPos{f}[4] \dicFlx{(‑ar, ‑ar)}[5] \textbf{1.} \dicDirectTranslationCS{(na)mazání, promazání} \dicExampleIS{fara með bílinn í smurningu} \dicExampleCS{zajet s~autem na promazání}  \textbf{2.} \dicSynonym{smurolía} \dicDirectTranslationCS{mazivo, olej}  \textbf{3.} \dicFieldCat{náb.} \dicDirectTranslationCS{pomazání};  \dicPhraseIS{hinsta\,/\addthin síðasta smurning} \dicFieldCat{náb.} \dicDirectTranslationCS{poslední pomazání}
\dicEntry[smurningsolía] \dicTerm{smurnings··olí|a} \dicIPA{{s}{m}{\textscy}{r}{\textsubring{d}}{n}{i}{\ng}{s}{\textopeno}{l}{i}{j}{a}} \dicPos{f}[1] \dicFlx{(‑u, ‑ur)}[7] \dicDirectTranslationCS{mazací\,/\addthin motorový olej}
\dicEntry[smurolía] \dicTerm{smur··olí|a} \dicIPA{{s}{m}{\textscy}{\textlengthmark}{r}{\textopeno}{l}{i}{j}{a}} \dicPos{f}[1] \dicFlx{(‑u, ‑ur)}[7] \dicSynonym{smurningsolía} \dicDirectTranslationCS{mazací\,/\addthin motorový olej}
\dicEntry[smurostur] \dicTerm{smur··ost|ur} \dicIPA{{s}{m}{\textscy}{\textlengthmark}{r}{\textopeno}{s}{\textsubring{d}}{\textscy}{\textsubring{r}}} \dicPos{m}[6] \dicFlx{(‑s, ‑ar)}[4] \dicDirectTranslationCS{tavený sýr}
\dicEntry[smurstöð] \dicTerm{smur··stöð} \dicIPA{{s}{m}{\textscy}{\textsubring{r}}{s}{\textsubring{d}}{\oe}{\texttheta}} \dicPos{f}[6] \dicFlx{(‑var, ‑var)}[1] \dicDirectTranslationCS{servis, autoopravna}
\dicEntry[smurt] \dicTerm{smurt} \dicIPA{{s}{m}{\textscy}{\textsubring{r}}{\textsubring{d}}} \dicPos{v} \dicFlx{supin} \dicLink{smyrja}
\dicEntry[smygi] \dicTerm{smygi} \dicIPA{{s}{m}{\textsci}{j}{\textlengthmark}{\textsci}} \dicPos{v} \dicFlx{con pf sg 1 pers} \dicLink{smjúga}
\dicEntry[smygildi] \dicTerm{smy··gildi} \dicIPA{{s}{m}{\textsci}{\textlengthmark}{\r{\textObardotlessj}}{\textsci}{l}{\textsubring{d}}{\textsci}} \dicPos{n}[2] \dicFlx{(‑s, ‑)}[14] \dicFieldCat{poč.} \dicSynonym{vafrakaka} \dicDirectTranslationCS{cookie}
\dicEntry[smygl] \dicTerm{smygl} \dicIPA{{s}{m}{\textsci}{\r{g}}{\textsubring{l}}} \dicPos{n}[2] \dicFlx{(‑s)}[2] \textbf{1.} \dicDirectTranslationCS{pašování}  \textbf{2.} \dicDirectTranslationCS{pašované zboží, kontraband}
\dicEntry[smygla] \dicTerm{smygl|a} \dicIPA{{s}{m}{\textsci}{\r{g}}{l}{a}} \dicPos{v}[1] \dicFlx{(‑aði)}[1] \dicFlx{dat} \dicDirectTranslationCS{pašovat} \dicExampleIS{smygla fíkniefnum til landsins} \dicExampleCS{pašovat do země narkotika}
\dicEntry[smyglari] \dicTerm{smygl··ar|i} \dicIPA{{s}{m}{\textsci}{\r{g}}{l}{a}{r}{\textsci}} \dicPos{m}[1] \dicFlx{(‑a, ‑ar)}[13] \dicDirectTranslationCS{pašerák, pašeračka}
\dicEntry[smyglvarningur] \dicTerm{smygl··varn·ing|ur} \dicIPA{{s}{m}{\textsci}{\r{g}}{l}{v}{a}{r}{\textsubring{d}}{n}{i}{\ng}{\r{g}}{\textscy}{\textsubring{r}}} \dicPos{m}[6] \dicFlx{(‑s)}[9] \dicDirectTranslationCS{pašované zboží, kontraband}
\dicEntry[smylli] \dicTerm{smylli} \dicIPA{{s}{m}{\textsci}{\textsubring{d}}{l}{\textsci}} \dicPos{v} \dicFlx{con pf sg 1 pers} \dicLink{smella\smash{\textsuperscript{3}}}
\dicEntry[smyr] \dicTerm{smyr} \dicIPA{{s}{m}{\textsci}{\textlengthmark}{\textsubring{r}}} \dicPos{v} \dicFlx{ind praes sg 1 pers} \dicLink{smyrja}
\dicEntry[smyrði] \dicTerm{smyrði} \dicIPA{{s}{m}{\textsci}{r}{ð}{\textsci}} \dicPos{v} \dicFlx{con pf sg 1 pers} \dicLink{smyrja}
\dicEntry[smyrðlingur] \dicTerm{smyrð··ling|ur}\dicTerm{, smyrlingur} \dicIPA{{s}\-{m}\-{\textsci}\-{r}\-{l}\-{i}\-{\ng}\-{\r{g}}\-{\textscy}\-{\textsubring{r}}\-} \dicPos{m}[6] \dicFlx{(‑s, ‑ar)}[8] \dicDirectTranslationCS{mumie}
\dicEntry[smyrill] \dicTerm{smyr|ill} \dicIPA{{s}{m}{\textsci}{\textlengthmark}{r}{\textsci}{\textsubring{d}}{\textsubring{l}}} \dicPos{m}[6] \dicFlx{(‑ils, ‑lar)}[35] \dicFieldCat{zool.} \dicDirectTranslationCS{dřemlík, dřemlík tundrový} \textit{(l.~{\textLA{Falco columbarius}})}  \dicsymPhoto\ 
\dicFigure{ds_image_smyrill_0_2.jpg}{Smyrill}{Smyrill - John Harrison, CC BY-SA 3.0}
\dicEntry[smyrja] \dicTerm{smyrja} \dicsymFrequent\  \dicIPA{{s}{m}{\textsci}{r}{j}{a}} \dicPos{v}[4] \dicFlx{(smyr, smurði, smurðum, smyrði, smurt)}[7] \dicFlx{acc\,/\addthin dat} \textbf{1.} \dicFlx{acc\,/\addthin dat} \dicDirectTranslationCS{(na)mazat (chléb ap.)} \dicExampleIS{smyrja brauðið} \dicExampleCS{namazat chléb};  \dicPhraseIS{smyrja e‑u á e‑ð} \dicDirectTranslationCS{namazat (co) na (co)} \dicExampleIS{smyrja smjöri á brauð} \dicExampleCS{namazat máslo na chleba}  \textbf{2.} \dicFlx{acc} \dicDirectTranslationCS{(na)mazat, promazat} \dicExampleIS{smyrja vélina} \dicExampleCS{promazat motor}
\dicEntry[smyrlingur] \dicTerm{smyr··ling|ur} \dicIPA{{s}{m}{\textsci}{r}{\textsubring{d}}{l}{i}{\ng}{\r{g}}{\textscy}{\textsubring{r}}} \dicPos{m}[6] \dicFlx{(‑s, ‑ar)}[8] \dicLink{smyrðlingur}
\dicEntry[smyrsl] \dicTerm{smyrsl}\dicTerm{, smyrsli} \dicIPA{{s}\-{m}\-{\textsci}\-{\textsubring{r}}\-{s}\-{\textsubring{d}}\-{\textsubring{l}}\-} \dicPos{n}[2] \dicFlx{(‑s, ‑)}[5] \dicDirectTranslationCS{mast, krém}
\dicEntry[smyrsli] \dicTerm{smyrsli} \dicIPA{{s}{m}{\textsci}{\textsubring{r}}{s}{\textsubring{d}}{l}{\textsci}} \dicPos{n}[2] \dicFlx{(‑s, ‑)}[14] \dicLink{smyrsl}
\dicEntry[smýg] \dicTerm{smýg} \dicIPA{{s}{m}{i}{\textlengthmark}{x}} \dicPos{v} \dicFlx{ind praes sg 1 pers} \dicLink{smjúga}
\dicEntry[smæð] \dicTerm{smæð} \dicIPA{{s}{m}{a}{i}{\textlengthmark}{\texttheta}} \dicPos{f}[7] \dicFlx{(‑ar, ‑ir)}[1] \textbf{1.} \dicSynonym*{það að vera smár} \dicDirectTranslationCS{malost, malá velikost}  \textbf{2.} \dicSynonym*{lítilvægi} \dicDirectTranslationCS{nedůležitost, nepodstatnost} \dicExampleIS{finna til smæðar sinnar frammi fyrir náttúruöflunum} \dicExampleCS{cítit svoji nepodstatnost ve srovnání se silami přírody}
\dicEntry[smækka] \dicTerm{smækk|a} \dicIPA{{s}{m}{a}{i}{h}{\r{g}}{a}} \dicPos{v}[1] \dicFlx{(‑aði)}[34] \textbf{1.} \dicSynonym{minnka} \dicDirectTranslationCS{(z)menšit se, zmenšovat se, ubýt, ubývat} \dicExampleIS{Fiskurinn er alltaf að smækka.} \dicExampleCS{Ryb stále ubývá.}  \textbf{2.} \dicFlx{acc} \dicDirectTranslationCS{zmenšit, zmenšovat, (z)redukovat} \dicExampleIS{smækka myndina um helming} \dicExampleCS{zmenšit obrázek o~polovinu}
\dicEntry[smækkaður] \dicTerm{smækk··|aður} \dicIPA{{s}{m}{a}{i}{h}{\r{g}}{a}{ð}{\textscy}{\textsubring{r}}} \dicPos{adj}[3] \dicFlx{(f ‑uð)}[4] \dicDirectTranslationCS{zmenšený}
\dicEntry[smælingi] \textls[15]{\dicTerm{smæl··ing|i} \dicIPA{{s}{m}{a}{i}{\textlengthmark}{l}{i}{\textltailn}{\r{\textObardotlessj}}{\textsci}} \dicPos{m}[1] \dicFlx{(‑ja, ‑jar)}[14] \dicSynonym{fátæklingur}} \dicDirectTranslationCS{chudák, chudačka, chuďas(ka)}
\dicEntry[smælki] \dicTerm{smælki} \dicIPA{{s}{m}{a}{i}{\textsubring{l}}{\r{\textObardotlessj}}{\textsci}} \dicPos{n}[2] \dicFlx{(‑s)}[20] \textbf{1.} \dicDirectTranslationCS{drobnůstky, drobnosti}  \textbf{2.} \dicDirectTranslationCS{malé brambory}
\dicEntry[smærra] \dicTerm{smærra} \dicIPA{{s}{m}{a}{i}{r}{\textlengthmark}{a}} \dicPos{adv} \dicFlx{comp} \dicLink{smátt}
\dicEntry[smæst] \dicTerm{smæst} \dicIPA{{s}{m}{a}{i}{s}{\textsubring{d}}} \dicPos{adv} \dicFlx{sup} \dicLink{smátt}
\dicEntry[smölun] \dicTerm{smölun} \dicIPA{{s}{m}{\oe}{\textlengthmark}{l}{\textscy}{\textsubring{n}}} \dicPos{f}[7] \dicFlx{(‑ar, smalanir)}[16] \textbf{1.} \dicSynonym*{smölun kinda} \dicDirectTranslationCS{shánění (ovcí)}  \textbf{2.} \dicSynonym*{söfnun atkvæða} \dicDirectTranslationCS{sbírání (volebních) hlasů}
\dicEntry[snafs] \dicTerm{snafs} \dicIPA{{s}{\textsubring{d}}{n}{a}{f}{s}} \dicPos{m}[4] \dicFlx{(‑, ‑ar)}[36] \dicDirectTranslationCS{panák, štamprle} \dicExampleIS{snafs af brennivíni} \dicExampleCS{panák pálenky}
\dicEntry[snaggaralegur] \dicTerm{snaggara··legur} \dicIPA{{s}{\textsubring{d}}{n}{a}{\r{g}}{\textlengthmark}{a}{r}{a}{l}{\textepsilon}{\textbabygamma}{\textscy}{\textsubring{r}}} \dicPos{adj}[1]\dicFlx{}[-8] \dicDirectTranslationCS{živý, energický} \dicExampleIS{snaggaralegar hreyfingar} \dicExampleCS{energické pohyby}
\dicEntry[snagi] \dicTerm{snag|i} \dicIPA{{s}{\textsubring{d}}{n}{a}{i}{j}{\textlengthmark}{\textsci}} \dicPos{m}[1] \dicFlx{(‑a, ‑ar)}[8] \dicDirectTranslationCS{háček, věšák (na oděv ap.)}
\dicEntry[snakk] \dicTerm{snakk} \dicIPA{{s}{\textsubring{d}}{n}{a}{h}{\r{g}}} \dicPos{n}[2] \dicFlx{(‑s)}[2] \textbf{1.} \dicSynonym{tal} \dicDirectTranslationCS{hovor, řeč, povídání}  \textbf{2.} \dicSynonym{nasl} \dicDirectTranslationCS{svačina, lehké jídlo}
\dicEntry[snakkur] \dicTerm{snakk|ur} \dicIPA{{s}{\textsubring{d}}{n}{a}{h}{\r{g}}{\textscy}{\textsubring{r}}} \dicPos{m}[6] \dicFlx{(‑s, ‑ar)}[25] \dicFieldCat{pov.} \dicIndirectTranslationCS{stvoření, které krade vlnu a~seno hospodářům}
\dicEntry[snapa] \dicTerm{snap|a} \dicIPA{{s}{\textsubring{d}}{n}{a}{\textlengthmark}{\textsubring{b}}{a}} \dicPos{v}[1] \dicFlx{(‑aði)}[13] \dicFlx{(dat +) acc} \dicSynonym{sníkja} \dicDirectTranslationCS{(vy)žadonit, (vy)škemrat} \dicExampleIS{snapa sér vinnu} \dicExampleCS{vyžadonit si práci}
\dicEntry[snar] \dicTerm{snar} \dicIPA{{s}{\textsubring{d}}{n}{a}{\textlengthmark}{\textsubring{r}}} \dicPos{adj}[5] \dicFlx{(f snör)}[9] \textbf{1.} \dicSynonym{fljótur} \dicDirectTranslationCS{rychlý, čilý, hbitý};  \dicPhraseIS{hafa snör handtök} \dicDirectTranslationCS{prudce reagovat};  \dicPhraseIS{vera snar í snúningum} \dicDirectTranslationCS{být hbitý}  \textbf{2.} \dicDirectTranslationCS{důležitý, vážný};  \dicPhraseIS{e‑að er snar þáttur í e‑u} \dicDirectTranslationCS{(co) je důležitá část (čeho)}  \textbf{3.} \dicLangCat{hovor.} \dicSynonym{snarvitlaus} \dicDirectTranslationCS{šílený, (jsoucí) mimo} \dicExampleIS{Ertu alveg snar?} \dicExampleCS{Jsi úplně mimo?}
\dicEntry[snara] \dicTerm{snara\smash{\textsuperscript{1}}} \dicIPA{{s}{\textsubring{d}}{n}{a}{\textlengthmark}{r}{a}} \dicPos{f}[1] \dicFlx{(snöru, snörur)}[8] \textbf{1.} \dicSynonym*{snærislykkja} \dicDirectTranslationCS{smyčka, oko, klička}  \textbf{2.} \dicSynonym{slöngvivaður} \dicDirectTranslationCS{laso};  \dicIdiom{snara}{ \dicPhraseIS{ganga í snöruna}} \dicLangCat{přen.} \dicDirectTranslationCS{padnout\,/\addthin vlézt do pasti};  \dicPhraseIS{leggja snörur fyrir e‑n} \dicLangCat{přen.} \dicDirectTranslationCS{nalíčit\,/\addthin nastrojit (komu) past};  \dicPhraseIS{vera fastur í snörunni} \dicLangCat{přen.} \dicDirectTranslationCS{být v~pasti}
\dicEntry[snara] \dicTerm{snar|a\smash{\textsuperscript{2}}} \dicsymFrequent\  \dicIPA{{s}{\textsubring{d}}{n}{a}{\textlengthmark}{r}{a}} \dicPos{v}[1] \dicFlx{(‑aði)}[13] \dicFlx{acc\,/\addthin dat} \textbf{1.} \dicFlx{acc} \dicSynonym*{veiða í snöru} \dicDirectTranslationCS{chytit\,/\addthin chytat do smyčky\,/\addthin oka\,/\addthin lasa} \dicExampleIS{snara dýrið} \dicExampleCS{chytit zvíře do lasa}  \textbf{2.} \dicFlx{dat} \dicSynonym{þýða} \dicDirectTranslationCS{přeložit, překládat (knihu ap.)} \dicExampleIS{snara bókinni úr ensku á íslensku} \dicExampleCS{přeložit knížku z~angličtiny do islandštiny};  \dicIdiom{snara}[í]{ \dicPhraseIS{snara sér í e‑ð}} \dicDirectTranslationCS{hodit na sebe (co), rychle se ustrojit do (čeho)} \dicExampleIS{snara sér í buxurnar} \dicExampleCS{hodit na sebe kalhoty};  \dicIdiom{snara}[út]{ \dicPhraseIS{snara út e‑u}} \dicDirectTranslationCS{(na místě) zaplatit (co), vysolit (co) (peníze ap.)} \dicExampleIS{snara út fimm milljónum} \dicExampleCS{vysolit pět milionů};  \dicIdiom{snarast}{ \dicPhraseIS{snarast}} \dicFlx{refl} \dicDirectTranslationCS{vyrazit, vyletět}
\dicEntry[snarbrattur] \dicTerm{snar··|brattur} \dicIPA{{s}{\textsubring{d}}{n}{a}{r}{\textsubring{b}}{r}{a}{h}{\textsubring{d}}{\textscy}{\textsubring{r}}} \dicPos{adj}[1] \dicFlx{(f ‑brött)}[12] \dicDirectTranslationCS{strmý, přík\-rý, srázný}
\dicEntry[snark] \dicTerm{snark} \dicIPA{{s}{\textsubring{d}}{n}{a}{\textsubring{r}}{\r{g}}} \dicPos{n}[2] \dicFlx{(‑s)}[2] \dicDirectTranslationCS{(za)praskání (v~ohni ap.)}
\dicEntry[snarka] \dicTerm{snark|a} \dicIPA{{s}{\textsubring{d}}{n}{a}{\textsubring{r}}{\r{g}}{a}} \dicPos{v}[1] \dicFlx{(‑aði)}[8] \dicPhraseIS{það snarkar í e‑u} \dicFlx{impers} \dicDirectTranslationCS{praská v~(čem) (v~ohni ap.)} \dicExampleIS{Það snarkar í eldinum.} \dicExampleCS{V~ohni praská.}
\dicEntry[snarl] \dicTerm{snarl} \dicIPA{{s}{\textsubring{d}}{n}{a}{r}{\textsubring{d}}{\textsubring{l}}} \dicPos{n}[2] \dicFlx{(‑s)}[2] \dicDirectTranslationCS{svačina, lehké jídlo}
\dicEntry[snarpur] \dicTerm{snarpur} \dicIPA{{s}{\textsubring{d}}{n}{a}{\textsubring{r}}{\textsubring{b}}{\textscy}{\textsubring{r}}} \dicPos{adj}[1] \dicFlx{(f snörp)}[2] \textbf{1.} \dicSynonym{kraftmikill} \dicDirectTranslationCS{mohutný, silný} \dicExampleIS{snarpur jarðskjálfti} \dicExampleCS{mohutné zemětřesení}  \textbf{2.} \dicSynonym{beittur} \dicDirectTranslationCS{ostrý (nůž ap.)} \dicExampleIS{snarpur hnífur} \dicExampleCS{ostrý nůž}  \textbf{3.} \dicSynonym{hrjúfur} \dicDirectTranslationCS{hrubý, drsný (kůže ap.)} \dicExampleIS{snörp húð} \dicExampleCS{hrubá kůže}  \textbf{4.} \dicSynonym{bitur} \dicDirectTranslationCS{řezavý, ostrý (vítr ap.)} \dicExampleIS{snarpur kuldi} \dicExampleCS{řezavý chlad}  \textbf{5.} \dicSynonym{ákafur} \dicDirectTranslationCS{intenzivní, razantní} \dicExampleIS{snarpar umræður} \dicExampleCS{intenzivní diskuze}
\dicEntry[snarráður] \dicTerm{snar··ráður} \dicIPA{{s}{\textsubring{d}}{n}{a}{r}{\textlengthmark}{au}{ð}{\textscy}{\textsubring{r}}} \dicPos{adj}[2]\dicFlx{}[-6] \dicDirectTranslationCS{duchapřítomný, pohotový}
\dicEntry[snarrótarpuntur] \dicTerm{snar·rótar··punt|ur} \dicIPA{{s}{\textsubring{d}}{n}{a}{r}{\textlengthmark}{ou}{\textsubring{d}}{a}{\textsubring{r}}{p\smash{\textsuperscript{h}}}{\textscy}{\textsubring{n}}{\textsubring{d}}{\textscy}{\textsubring{r}}} \dicPos{m}[6] \dicFlx{(‑s)}[7] \dicFieldCat{bot.} \dicDirectTranslationCS{metlice trsnatá} \textit{(l.~{\textLA{Deschampsia cespitosa}})}  \dicsymPhoto\ 
\dicFigure{ds_image_snarrotarpuntur_0_2.jpg}{Snarrótarpuntur}{Snarrótarpuntur - Christian Fischer, CC BY-SA 3.0}
\dicEntry[snarræði] \dicTerm{snar··ræði} \dicIPA{{s}{\textsubring{d}}{n}{a}{r}{\textlengthmark}{a}{i}{ð}{\textsci}} \dicPos{n}[2] \dicFlx{(‑s)}[20] \dicDirectTranslationCS{duchapřítomnost, pohotovost} \dicExampleIS{sýna mikið snarræði} \dicExampleCS{prokázat velkou duchapřítomnost}
\dicEntry[snarsneri] \dicTerm{snar··sneri} \dicIPA{{s}{\textsubring{d}}{n}{a}{\textsubring{r}}{s}{\textsubring{d}}{n}{\textepsilon}{r}{\textsci}} \dicPos{v} \dicFlx{ind\,/\addthin con pf sg 1 pers} \dicLink{snarsnúa}
\dicEntry[snarsnerum] \dicTerm{snar··snerum} \dicIPA{{s}{\textsubring{d}}{n}{a}{\textsubring{r}}{s}{\textsubring{d}}{n}{\textepsilon}{r}{\textscy}{\textsubring{m}}} \dicPos{v} \dicFlx{ind pf pl 1 pers} \dicLink{snarsnúa}
\dicEntry[snarsnéri] \dicTerm{snar··snéri} \dicIPA{{s}{\textsubring{d}}{n}{a}{\textsubring{r}}{s}{\textsubring{d}}{n}{j}{\textepsilon}{r}{\textsci}} \dicPos{v} \dicFlx{ind\,/\addthin con pf sg 1 pers} \dicLink{snarsnúa}
\dicEntry[snarsnérum] \dicTerm{snar··snérum} \dicIPA{{s}{\textsubring{d}}{n}{a}{\textsubring{r}}{s}{\textsubring{d}}{n}{j}{\textepsilon}{r}{\textscy}{\textsubring{m}}} \dicPos{v} \dicFlx{ind pf pl 1 pers} \dicLink{snarsnúa}
\dicEntry[snarsnúa] \dicTerm{snar··|snúa} \dicIPA{{s}{\textsubring{d}}{n}{a}{\textsubring{r}}{s}{\textsubring{d}}{n}{u}{a}} \dicPos{v}[5] \dicFlx{(‑sný, ‑sneri\,/\addthin ‑snéri, ‑snerum\,/\addthin ‑snérum, ‑sneri\,/\addthin ‑snéri, ‑snúið)}[5] \dicDirectTranslationCS{rychle otočit\,/\addthin obrátit};  \dicIdiom{snarsnúast}{ \dicPhraseIS{snarsnúast}} \dicFlx{refl} \dicDirectTranslationCS{rychle se otočit\,/\addthin obrátit}
\dicEntry[snarsnúið] \dicTerm{snar··snúið} \dicIPA{{s}{\textsubring{d}}{n}{a}{\textsubring{r}}{s}{\textsubring{d}}{n}{u}{\textsci}{\texttheta}} \dicPos{v} \dicFlx{supin} \dicLink{snarsnúa}
\dicEntry[snarsný] \dicTerm{snar··sný} \dicIPA{{s}{\textsubring{d}}{n}{a}{\textsubring{r}}{s}{\textsubring{d}}{n}{i}} \dicPos{v} \dicFlx{ind praes sg 1 pers} \dicLink{snarsnúa}
\dicEntry[snarstansa] \dicTerm{snar··stans|a} \dicIPA{{s}{\textsubring{d}}{n}{a}{\textsubring{r}}{s}{\textsubring{d}}{a}{n}{s}{a}} \dicPos{v}[1] \dicFlx{(‑aði)}[50] \dicDirectTranslationCS{prudce\,/\addthin rychle zastavit}
\dicEntry[snartúlkun] \dicTerm{snar··túlk|un} \dicIPA{{s}{\textsubring{d}}{n}{a}{\textsubring{r}}{t\smash{\textsuperscript{h}}}{u}{\textsubring{l}}{\r{g}}{\textscy}{\textsubring{n}}} \dicPos{f}[7] \dicFlx{(‑unar, ‑anir)}[8] \dicDirectTranslationCS{simultánní tlumočení}
\dicEntry[snarvitlaus] \dicTerm{snar··vit·laus} \dicIPA{{s}{\textsubring{d}}{n}{a}{r}{v}{\textsci}{h}{\textsubring{d}}{l}{\oe i}{s}} \dicPos{adj}[5]\dicFlx{}[-1] \textbf{1.} \dicLangCat{přen.} \dicDirectTranslationCS{padlý na hlavu}  \textbf{2.} \dicDirectTranslationCS{bláznivý, šílený} \dicIndirectTranslationCS{(o~počasí)}
\dicEntry[snasar] \dicTerm{snasar} \dicIPA{{s}{\textsubring{d}}{n}{a}{\textlengthmark}{s}{a}{\textsubring{r}}} \dicPos{f} \dicFlx{sg gen} \dicLink{snös}
\dicEntry[snatt] \dicTerm{snatt} \dicIPA{{s}{\textsubring{d}}{n}{a}{h}{\textsubring{d}}} \dicPos{n}[2] \dicFlx{(‑s)}[2] \dicDirectTranslationCS{zařizování, běhání (po úřadech ap.)}
\dicEntry[snatur] \dicTerm{snatur} \dicIPA{{s}{\textsubring{d}}{n}{a}{\textlengthmark}{\textsubring{d}}{\textscy}{\textsubring{r}}} \dicPos{n}[2] \dicFlx{(‑s)}[28] \dicPhraseIS{í snatri} \dicFlx{adv} \dicSynonym{fljótt} \dicDirectTranslationCS{honem, ve spěchu}
\dicEntry[snauður] \dicTerm{snauður} \dicIPA{{s}{\textsubring{d}}{n}{\oe i}{\textlengthmark}{ð}{\textscy}{\textsubring{r}}} \dicPos{adj}[2]\dicFlx{}[-6] \textbf{1.} \dicSynonym{fátækur} \dicDirectTranslationCS{nuzný, chudý, (jsoucí) bez prostředků}  \textbf{2.} \dicPhraseIS{snauður af e‑u} \dicDirectTranslationCS{prostý (čeho), postrádající (co)} \dicExampleIS{Fæðan er snauð af bætiefnum.} \dicExampleCS{Jídlo je prosté přídavných látek.}
\dicEntry[snauta] \dicTerm{snaut|a} \dicIPA{{s}{\textsubring{d}}{n}{\oe i}{\textlengthmark}{\textsubring{d}}{a}} \dicPos{v}[1] \dicFlx{(‑aði)}[1] \dicSynonym{snáfa} \dicDirectTranslationCS{vytratit se, klidit se, zmizet};  \dicPhraseIS{snauta burt} \dicDirectTranslationCS{ztratit se, klidit se (z~očí ap.)};  \dicPhraseIS{snautaðu burt!} \dicFlx{imper} \dicDirectTranslationCS{zmiz!, kliď se mi z~očí!}
\dicEntry[snautlegur] \dicTerm{snaut··legur} \dicIPA{{s}{\textsubring{d}}{n}{\oe i}{\textlengthmark}{\textsubring{d}}{l}{\textepsilon}{\textbabygamma}{\textscy}{\textsubring{r}}} \dicPos{adj}[1]\dicFlx{}[-8] \dicSynonym{sviplaus} \dicDirectTranslationCS{jednotvárný, mdlý, fádní}
\dicEntry[snáði] \dicTerm{snáð|i} \dicIPA{{s}{\textsubring{d}}{n}{au}{\textlengthmark}{ð}{\textsci}} \dicPos{m}[1] \dicFlx{(‑a, ‑ar)}[1] \dicSynonym{drengur} \dicDirectTranslationCS{kluk, výrostek}
\dicEntry[snáfa] \dicTerm{snáf|a} \dicIPA{{s}{\textsubring{d}}{n}{au}{\textlengthmark}{v}{a}} \dicPos{v}[1] \dicFlx{(‑aði)}[1] \dicSynonym{snauta} \dicDirectTranslationCS{klidit se, ztratit se (z~očí ap.)};  \dicPhraseIS{snáfaðu burt!} \dicFlx{imper} \dicDirectTranslationCS{kliď se odtud!}
\dicEntry[snákur] \dicTerm{snák|ur} \dicIPA{{s}{\textsubring{d}}{n}{au}{\textlengthmark}{\r{g}}{\textscy}{\textsubring{r}}} \dicPos{m}[6] \dicFlx{(‑s, ‑ar)}[22] \dicSynonym{slanga} \dicDirectTranslationCS{had (užovka ap.)} \dicIndirectTranslationCS{(spíše nejedovatý)}
\dicEntry[snefilefni] \dicTerm{snefil··efni} \dicIPA{{s}{\textsubring{d}}{n}{\textepsilon}{\textlengthmark}{v}{\textsci}{l}{\textepsilon}{\textsubring{b}}{n}{\textsci}} \dicPos{n}[2] \dicFlx{(‑s, ‑)}[14] \dicFieldCat{chem.} \dicDirectTranslationCS{stopový prvek}
\dicEntry[sneggra] \dicTerm{sneggra} \dicIPA{{s}{\textsubring{d}}{n}{\textepsilon}{\r{g}}{r}{a}} \dicPos{adv} \dicFlx{comp} \dicLink{snöggt}
\dicEntry[sneggst] \dicTerm{sneggst} \dicIPA{{s}{\textsubring{d}}{n}{\textepsilon}{x}{s}{\textsubring{d}}} \dicPos{adv} \dicFlx{sup} \dicLink{snöggt}
\dicEntry[sneið] \dicTerm{sneið\smash{\textsuperscript{1}}} \dicIPA{{s}{\textsubring{d}}{n}{ei}{\textlengthmark}{\texttheta}} \dicPos{f}[4] \dicFlx{(‑ar, ‑ar)}[1] \textbf{1.} \dicSynonym{flaga} \dicDirectTranslationCS{plátek, krajíc} \dicExampleIS{sneið af brauðinu} \dicExampleCS{krajíc chleba}  \textbf{2.} \dicSynonym{aðdróttun} \dicDirectTranslationCS{narážka, aluze};  \dicPhraseIS{senda e‑m sneið} \dicDirectTranslationCS{mít narážky na (koho)};  \dicPhraseIS{skilja sneiðina} \dicDirectTranslationCS{rozumět narážkám}  \textbf{3.} \dicFieldCat{mat.} \dicDirectTranslationCS{úseč}
\dicEntry[sneið] \dicTerm{sneið\smash{\textsuperscript{2}}} \dicIPA{{s}{\textsubring{d}}{n}{ei}{\textlengthmark}{\texttheta}} \dicPos{v} \dicFlx{ind pf sg 1 pers} \dicLink{sníða}
\dicEntry[sneiða] \dicTerm{snei|ða} \dicIPA{{s}{\textsubring{d}}{n}{ei}{\textlengthmark}{ð}{a}} \dicPos{v}[2] \dicFlx{(‑ddi, ‑tt)}[167] \dicFlx{acc} \textbf{1.} \dicSynonym*{skera í sneiðar} \dicDirectTranslationCS{(na)krájet na plátky} \dicExampleIS{sneiða brauðið} \dicExampleCS{nakrájet chleba na plátky}  \textbf{2.} \dicPhraseIS{sneiða að e‑m} \dicDirectTranslationCS{mít narážky na (koho)}  \textbf{3.} \dicPhraseIS{sneiða hjá e‑u} \dicDirectTranslationCS{vyvarovat se (čeho), vystříhat se (čeho), vyhnout se (čemu)} \dicExampleIS{sneiða hjá  þessari spurningu} \dicExampleCS{vyhnout se této otázce}  \textbf{4.} \dicSynonym*{ganga á snið} \dicDirectTranslationCS{jít šikmo (do svahu ap.)} \dicExampleIS{sneiða brekkuna} \dicExampleCS{jít šikmo po svahu}
\dicEntry[sneisafullur] \textls[15]{\dicTerm{sneisa··fullur} \dicIPA{{s}{\textsubring{d}}{n}{ei}{\textlengthmark}{s}{a}{f}{\textscy}{\textsubring{d}}{l}{\textscy}{\textsubring{r}}} \dicPos{adj}[10]\dicFlx{}[-8] \dicSynonym{barmafullur}} \dicDirectTranslationCS{přeplněný, plný po okraj}
\dicEntry[snekkja] \dicTerm{snekkj|a} \dicIPA{{s}{\textsubring{d}}{n}{\textepsilon}{h}{\r{\textObardotlessj}}{a}} \dicPos{f}[1] \dicFlx{(‑u, ‑ur)}[7] \textbf{1.} \dicFieldCat{nám.} \dicSynonym{lystibátur} \dicDirectTranslationCS{jachta} \dicExampleIS{skemmtisnekkja} \dicExampleCS{zábavní jachta}  \textbf{2.} \dicFieldCat{nám., hist.} \dicSynonym{herskip} \dicDirectTranslationCS{válečná loď}
\dicEntry[snemma] \dicTerm{snemma} \dicsymFrequent\  \dicIPA{{s}{\textsubring{d}}{n}{\textepsilon}{m}{\textlengthmark}{a}} \dicPos{adv} \dicFlx{(comp fyrr, sup fyrst)} \textbf{1.} \dicSynonym{árla} \dicDirectTranslationCS{brzy, brzo, časně} \dicExampleIS{vera snemma á fótum} \dicExampleCS{být brzo na nohou}  \textbf{2.} \dicSynonym*{eftir stuttan tíma} \dicDirectTranslationCS{brzy, brzo, zakrátko, záhy}
\dicEntry[snepill] \dicTerm{snep|ill} \dicIPA{{s}{\textsubring{d}}{n}{\textepsilon}{\textlengthmark}{\textsubring{b}}{\textsci}{\textsubring{d}}{\textsubring{l}}} \dicPos{m}[6] \dicFlx{(‑ils, ‑lar)}[35] \textbf{1.} \dicSynonym*{pappírssnifsi} \dicDirectTranslationCS{kousek\,/\addthin útržek papíru, papírek}  \textbf{2.} \dicDirectTranslationCS{bulvár, (bulvární) plátek}  \textbf{3.} \dicFieldCat{anat.} \dicDirectTranslationCS{lalok}
\dicEntry[sneri] \dicTerm{sneri} \dicIPA{{s}{\textsubring{d}}{n}{\textepsilon}{\textlengthmark}{r}{\textsci}} \dicPos{v} \dicFlx{ind\,/\addthin con pf sg 1 pers} \dicLink{snúa}
\dicEntry[snerill] \dicTerm{sner|ill} \dicIPA{{s}{\textsubring{d}}{n}{\textepsilon}{\textlengthmark}{r}{\textsci}{\textsubring{d}}{\textsubring{l}}} \dicPos{m}[6] \dicFlx{(‑ils, ‑lar)}[35] \dicDirectTranslationCS{kulatá klika} \dicExampleIS{snúa snerlinum} \dicExampleCS{otočit klikou}
\dicEntry[snerpa] \dicTerm{snerp|a\smash{\textsuperscript{1}}} \dicIPA{{s}{\textsubring{d}}{n}{\textepsilon}{\textsubring{r}}{\textsubring{b}}{a}} \dicPos{f}[1] \dicFlx{(‑u, ‑ur)}[19] \textbf{1.} \dicSynonym{röskleiki} \dicDirectTranslationCS{čilost, hbitost}  \textbf{2.} \dicSynonym*{hröð viðbrögð} \dicDirectTranslationCS{rychlost a~síla}
\dicEntry[snerpa] \dicTerm{snerp|a\smash{\textsuperscript{2}}} \dicIPA{{s}{\textsubring{d}}{n}{\textepsilon}{\textsubring{r}}{\textsubring{b}}{a}} \dicPos{v}[2] \dicFlx{(‑ti, ‑t)}[74] \dicPhraseIS{snerpa á e‑u} \dicDirectTranslationCS{ohřát (co), zahřát (co) (kafe ap.)}
\dicEntry[snerta] \dicTerm{snert|a} \dicsymFrequent\  \dicIPA{{s}{\textsubring{d}}{n}{\textepsilon}{\textsubring{r}}{\textsubring{d}}{a}} \dicPos{v}[5] \dicFlx{(‑i, ‑)\,/\addthin (snerti, snart, snurtum, snerti, snortið)}[1] \dicFlx{acc} \textbf{1.} \dicSynonym*{koma við} \dicDirectTranslationCS{dotknout se, dotýkat se (ruky ap.)} \dicExampleIS{snerta hönd hennar} \dicExampleCS{dotýkat se její ruky}  \textbf{2.} \dicSynonym*{hafa áhrif á} \dicDirectTranslationCS{dotknout se, dotýkat se, ublížit, ubližovat} \dicExampleIS{Það snerti mig djúpt.} \dicExampleCS{Hluboce se mě to dotklo.}  \textbf{3.} \dicSynonym{varða\smash{\textsuperscript{2}}} \dicDirectTranslationCS{(do)týkat se};  \dicPhraseIS{hvað þetta snertir} \dicDirectTranslationCS{co se toho týče, pokud jde o~to};  \dicPhraseIS{hvað mig snertir} \dicDirectTranslationCS{co se mě týče, pokud jde o~mě} \dicExampleIS{Þetta er í lagi hvað mig snertir.} \dicExampleCS{Co se mě týče, je to v~pořádku.};  \dicIdiom{snertast}{ \dicPhraseIS{snertast}} \dicFlx{refl} \dicDirectTranslationCS{dotknout\,/\addthin dotýkat se (navzájem)}
\dicEntry[snertiflötur] \dicTerm{snerti··|flötur} \dicIPA{{s}{\textsubring{d}}{n}{\textepsilon}{\textsubring{r}}{\textsubring{d}}{\textsci}{f}{l}{\oe}{\textsubring{d}}{\textscy}{\textsubring{r}}} \dicPos{m}[11] \dicFlx{(‑flatar, ‑fletir)}[5] \textbf{1.} \dicDirectTranslationCS{styčná plocha}  \textbf{2.} \dicFieldCat{poč.} \dicDirectTranslationCS{dotyková plocha, touchpad}
\dicEntry[snertill] \dicTerm{snert|ill} \dicIPA{{s}{\textsubring{d}}{n}{\textepsilon}{\textsubring{r}}{\textsubring{d}}{\textsci}{\textsubring{d}}{\textsubring{l}}} \dicPos{m}[6] \dicFlx{(‑ils, ‑lar)}[35] \dicFieldCat{mat.} \dicDirectTranslationCS{tečna}
\dicEntry[snerting] \dicTerm{snert··ing} \dicsymFrequent\  \dicIPA{{s}{\textsubring{d}}{n}{\textepsilon}{\textsubring{r}}{\textsubring{d}}{i}{\ng}{\r{g}}} \dicPos{f}[4] \dicFlx{(‑ar, ‑ar)}[5] \dicDirectTranslationCS{dotek, dotyk, dotknutí, kontakt};  \dicPhraseIS{komast í snertingu við e‑ð} \dicFlx{refl} \dicDirectTranslationCS{přijít do kontaktu s~(čím)};  \dicPhraseIS{vera í snertingu við e‑n} \dicDirectTranslationCS{být v~kontaktu s~(kým)}
\dicEntry[snertiskjár] \dicTerm{snerti··skjá|r} \dicIPA{{s}{\textsubring{d}}{n}{\textepsilon}{\textsubring{r}}{\textsubring{d}}{\textsci}{s}{\r{\textObardotlessj}}{au}{\textsubring{r}}} \dicPos{m}[9] \dicFlx{(‑s, ‑ir)}[12] \dicFieldCat{poč.} \dicDirectTranslationCS{dotyková obrazovka}
\dicEntry[snertiskyn] \dicTerm{snerti··skyn} \dicIPA{{s}{\textsubring{d}}{n}{\textepsilon}{\textsubring{r}}{\textsubring{d}}{\textsci}{s}{\r{\textObardotlessj}}{\textsci}{\textsubring{n}}} \dicPos{n}[2] \dicFlx{(‑s)}[2] \dicDirectTranslationCS{hmat} \dicIndirectTranslationCS{(smysl k~vnímání dotyku)}
\dicEntry[snertur] \dicTerm{snert|ur} \dicIPA{{s}{\textsubring{d}}{n}{\textepsilon}{\textsubring{r}}{\textsubring{d}}{\textscy}{\textsubring{r}}} \dicPos{m}[9] \dicFlx{(‑s, ‑ir)}[7] \textbf{1.} \dicSynonym{keimur} \dicDirectTranslationCS{náznak, nádech} \dicExampleIS{snertur af hæðni í orðum hans} \dicExampleCS{náznak sarkasmu v~jeho slovech}  \textbf{2.} \dicSynonym{aðkenning} \dicDirectTranslationCS{náběh (nachlazení ap.)} \dicExampleIS{snertur af kvefi} \dicExampleCS{náběh rýmy}
\dicEntry[sneyða] \dicTerm{sney|ða} \dicIPA{{s}{\textsubring{d}}{n}{ei}{\textlengthmark}{ð}{a}} \dicPos{v}[2] \dicFlx{(‑ddi, ‑tt)}[167] \dicFlx{acc + dat} \dicDirectTranslationCS{zbavit, vzít, odejmout};  \dicPhraseIS{sneyða e‑n e‑u} \dicDirectTranslationCS{odejmout (komu co)};  \dicIdiom{sneyðast}{ \dicPhraseIS{það sneyðist um e‑ð}} \dicFlx{refl impers} \dicDirectTranslationCS{(co) dochází, (čeho) ubývá (zásob ap.)} \dicExampleIS{Það fer að sneyðast um birgðirnar.} \dicExampleCS{Začínají docházet zásoby.}
\dicEntry[sneypa] \dicTerm{sneyp|a\smash{\textsuperscript{1}}} \dicIPA{{s}{\textsubring{d}}{n}{ei}{\textlengthmark}{\textsubring{b}}{a}} \dicPos{f}[1] \dicFlx{(‑u, ‑ur)}[19] \dicSynonym{skömm\smash{\textsuperscript{1}}} \dicDirectTranslationCS{hanba, os\-tu\-da, potupa};  \dicPhraseIS{gera e‑m þá sneypu að (gera e‑ð)} \dicDirectTranslationCS{zahanbit (koho) (uděláním (čeho))}
\dicEntry[sneypa] \dicTerm{sneyp|a\smash{\textsuperscript{2}}} \dicIPA{{s}{\textsubring{d}}{n}{ei}{\textlengthmark}{\textsubring{b}}{a}} \dicPos{v}[2] \dicFlx{(‑ti, ‑t)}[63] \dicFlx{acc} \dicSynonym{svívirða\smash{\textsuperscript{2}}} \dicDirectTranslationCS{zahanbit, potupit};  \dicIdiom{sneypast}{ \dicPhraseIS{sneypast}} \dicFlx{refl} \dicSynonym*{skammast sín} \dicDirectTranslationCS{stydět se, cítit se zahanbený}
\dicEntry[sneypulegur] \dicTerm{sneypu··legur} \dicIPA{{s}{\textsubring{d}}{n}{ei}{\textlengthmark}{\textsubring{b}}{\textscy}{l}{\textepsilon}{\textbabygamma}{\textscy}{\textsubring{r}}} \dicPos{adj}[1]\dicFlx{}[-8] \dicSynonym{skammarlegur} \dicDirectTranslationCS{zahanbený, potupený} \dicExampleIS{vera sneypulegur á svipinn} \dicExampleCS{mít zahanbený výraz}
\dicEntry[snéri] \dicTerm{snéri} \dicIPA{{s}{\textsubring{d}}{n}{j}{\textepsilon}{\textlengthmark}{r}{\textsci}} \dicPos{v} \dicFlx{ind\,/\addthin con pf sg 1 pers} \dicLink{snúa}
\dicEntry[snérum] \dicTerm{snérum} \dicIPA{{s}{\textsubring{d}}{n}{j}{\textepsilon}{\textlengthmark}{r}{\textscy}{\textsubring{m}}} \dicPos{v} \dicFlx{ind pf pl 1 pers} \dicLink{snúa}
\dicEntry[snið] \dicTerm{snið} \dicsymFrequent\  \dicIPA{{s}{\textsubring{d}}{n}{\textsci}{\textlengthmark}{\texttheta}} \dicPos{n}[2] \dicFlx{(‑s, ‑)}[5] \textbf{1.} \dicSynonym{form} \dicDirectTranslationCS{tvar, forma, podoba}  \textbf{2.} \dicSynonym*{fatasnið} \dicDirectTranslationCS{střih (kalhot ap.)} \dicExampleIS{Kjóllinn er fallegur í sniðinu.} \dicExampleCS{Šaty mají hezký střih.}  \textbf{3.} \dicFieldCat{poč.} \dicDirectTranslationCS{formát};  \dicIdiom{snið}{ \dicPhraseIS{e‑að er stórt\,/\addthin smátt í sniðum}} \dicDirectTranslationCS{(co) je velké\,/\addthin malé (firma ap.)};  \dicPhraseIS{fara\,/\addthin ganga á snið við e‑ð} \dicDirectTranslationCS{obejít\,/\addthin obcházet (co) (zákon ap.)}
\dicEntry[sniðganga] \dicTerm{snið··|ganga} \dicIPA{{s}{\textsubring{d}}{n}{\textsci}{ð}{\r{g}}{au}{\ng}{\r{g}}{a}} \dicPos{v}[5] \dicFlx{(‑geng, ‑gekk, ‑gengum, ‑gengi, ‑gengið)}[10] \dicFlx{acc} \dicDirectTranslationCS{obejít, obcházet, vyhnout se, vyhýbat se, bojkotovat} \dicExampleIS{sniðganga lögin} \dicExampleCS{obcházet zákon}
\dicEntry[sniðgekk] \dicTerm{snið··gekk} \dicIPA{{s}{\textsubring{d}}{n}{\textsci}{ð}{\r{\textObardotlessj}}{\textepsilon}{h}{\r{g}}} \dicPos{v} \dicFlx{ind pf sg 1 pers} \dicLink{sniðganga}
\dicEntry[sniðgeng] \dicTerm{snið··geng} \dicIPA{{s}{\textsubring{d}}{n}{\textsci}{ð}{\r{\textObardotlessj}}{ei}{\ng}{\r{g}}} \dicPos{v} \dicFlx{ind praes sg 1 pers} \dicLink{sniðganga}
\dicEntry[sniðgengi] \dicTerm{snið··gengi} \dicIPA{{s}{\textsubring{d}}{n}{\textsci}{ð}{\r{\textObardotlessj}}{ei}{\textltailn}{\r{\textObardotlessj}}{\textsci}} \dicPos{v} \dicFlx{con pf sg 1 pers} \dicLink{sniðganga}
\dicEntry[sniðgengið] \dicTerm{snið··gengið} \dicIPA{{s}{\textsubring{d}}{n}{\textsci}{ð}{\r{\textObardotlessj}}{ei}{\textltailn}{\r{\textObardotlessj}}{\textsci}{\texttheta}} \dicPos{v} \dicFlx{supin} \dicLink{sniðganga}
\dicEntry[sniðgengum] \dicTerm{snið··gengum} \dicIPA{{s}{\textsubring{d}}{n}{\textsci}{ð}{\r{\textObardotlessj}}{ei}{\ng}{\r{g}}{\textscy}{\textsubring{m}}} \dicPos{v} \dicFlx{ind pf pl 1 pers} \dicLink{sniðganga}
\dicEntry[sniði] \dicTerm{sniði} \dicIPA{{s}{\textsubring{d}}{n}{\textsci}{\textlengthmark}{ð}{\textsci}} \dicPos{v} \dicFlx{con pf sg 1 pers} \dicLink{sníða}
\dicEntry[sniðið] \dicTerm{sniðið} \dicIPA{{s}{\textsubring{d}}{n}{\textsci}{\textlengthmark}{ð}{\textsci}{\texttheta}} \dicPos{v} \dicFlx{supin} \dicLink{sníða}
\dicEntry[sniðugur] \dicTerm{snið··ugur} \dicsymFrequent\  \dicIPA{{s}{\textsubring{d}}{n}{\textsci}{\textlengthmark}{ð}{\textscy}{\textbabygamma}{\textscy}{\textsubring{r}}} \dicPos{adj}[1]\dicFlx{}[-8] \textbf{1.} \dicSynonym*{hugkvæmnislegur} \dicDirectTranslationCS{chytrý, geniální} \dicExampleIS{Hann kemur með sniðugar hugmyndir.} \dicExampleCS{Má geniální nápady.}  \textbf{2.} \dicSynonym{snjall} \dicDirectTranslationCS{praktický, vynalézavý}  \textbf{3.} \dicSynonym{fyndinn} \dicDirectTranslationCS{legrační, vtipný} \dicExampleIS{sniðug saga} \dicExampleCS{vtipný příběh}
\dicEntry[sniðum] \dicTerm{sniðum} \dicIPA{{s}{\textsubring{d}}{n}{\textsci}{\textlengthmark}{ð}{\textscy}{\textsubring{m}}} \dicPos{v} \dicFlx{ind pf pl 1 pers} \dicLink{sníða}
\dicEntry[sniffa] \dicTerm{sniff|a} \dicIPA{{s}{\textsubring{d}}{n}{\textsci}{f}{\textlengthmark}{a}} \dicPos{v}[1] \dicFlx{(‑aði)}[1] \dicFlx{acc} \dicLangCat{hovor.} \dicDirectTranslationCS{čichat (toluen ap.)}
\dicEntry[snifsi] \dicTerm{snifsi}\dicTerm{, snyfsi} \dicIPA{{s}\-{\textsubring{d}}\-{n}\-{\textsci}\-{f}\-{s}\-{\textsci}\-} \dicPos{n}[2] \dicFlx{(‑s, ‑)}[14] \dicDirectTranslationCS{útržek, fragment} \dicExampleIS{snifsi af handriti} \dicExampleCS{útržek rukopisu}
\dicEntry[snigill] \dicTerm{snig|ill} \dicIPA{{s}{\textsubring{d}}{n}{i}{j}{\textlengthmark}{\textsci}{\textsubring{d}}{\textsubring{l}}} \dicPos{m}[6] \dicFlx{(‑ils, ‑lar)}[35] \textbf{1.} \dicFieldCat{zool.} \dicDirectTranslationCS{plž, hlemýžď, šnek} \textit{(l.~{\textLA{Gastropoda}})}  \dicsymPhoto\   \textbf{2.} \dicDirectTranslationCS{šnek (dopravníku)}  \textbf{3.} \dicSynonym*{spíralbor} \dicDirectTranslationCS{vrták, nebozez}  \textbf{4.} \dicFieldCat{anat.} \dicDirectTranslationCS{hlemýžď} \dicIndirectTranslationCS{(část vnitřního ucha)}
\dicFigure{ds_image_snigill_0_1.jpg}{Snigill}{Snigill - Hakan Svensson, CC BY-SA 3.0}
\dicEntry[snikkari] \dicTerm{snikk··ar|i} \dicIPA{{s}{\textsubring{d}}{n}{\textsci}{h}{\r{g}}{a}{r}{\textsci}} \dicPos{m}[1] \dicFlx{(‑a, ‑ar)}[13] \dicLangCat{zast.} \dicSynonym{trésmiður} \dicDirectTranslationCS{truhlář(ka), stolař(ka)}
\dicEntry[snilld] \dicTerm{snilld} \dicIPA{{s}{\textsubring{d}}{n}{\textsci}{l}{\textsubring{d}}} \dicPos{f}[7] \dicFlx{(‑ar, ‑ir)}[1] \dicSynonym{snilligáfa} \dicDirectTranslationCS{genialita, mistrovství, bravurnost} \dicExampleIS{tefla skákina af hreinustu snilld} \dicExampleCS{hrát šachy s~absolutním mistrovstvím}
\dicEntry[snilldarlegur] \dicTerm{snilldar··legur} \dicIPA{{s}{\textsubring{d}}{n}{\textsci}{l}{\textsubring{d}}{a}{r}{l}{\textepsilon}{\textbabygamma}{\textscy}{\textsubring{r}}} \dicPos{adj}[1]\dicFlx{}[-8] \dicSynonym{ágætur} \dicDirectTranslationCS{geniální, mistrovský, mistrný, bravurní}
\dicEntry[snilldarverk] \dicTerm{snilldar··verk} \dicIPA{{s}{\textsubring{d}}{n}{\textsci}{l}{\textsubring{d}}{a}{r}{v}{\textepsilon}{\textsubring{r}}{\r{g}}} \dicPos{n}[2] \dicFlx{(‑s, ‑)}[5] \dicDirectTranslationCS{mistrovské dílo}
\dicEntry[snilligáfa] \dicTerm{snilli··gáf|a} \dicIPA{{s}{\textsubring{d}}{n}{\textsci}{\textsubring{d}}{l}{\textsci}{\r{g}}{au}{v}{a}} \dicPos{f}[1] \dicFlx{(‑u)}[5] \dicDirectTranslationCS{genialita, geniální nadání}
\dicEntry[snillingur] \dicTerm{snill··ing|ur} \dicIPA{{s}{\textsubring{d}}{n}{\textsci}{\textsubring{d}}{l}{i}{\ng}{\r{g}}{\textscy}{\textsubring{r}}} \dicPos{m}[6] \dicFlx{(‑s, ‑ar)}[8] \dicDirectTranslationCS{génius, virtuos(ka), mistr(yně)} \dicExampleIS{vera snillingur í fótbolta} \dicExampleCS{být fotbalový génius}
\dicEntry[snitta] \dicTerm{snitt|a\smash{\textsuperscript{1}}} \dicIPA{{s}{\textsubring{d}}{n}{\textsci}{h}{\textsubring{d}}{a}} \dicPos{f}[1] \dicFlx{(‑u, ‑ur)}[7] \dicFieldCat{kulin.} \dicDirectTranslationCS{jednohubka, chlebíček}
\dicEntry[snitta] \dicTerm{snitt|a\smash{\textsuperscript{2}}} \dicIPA{{s}{\textsubring{d}}{n}{\textsci}{h}{\textsubring{d}}{a}} \dicPos{v}[1] \dicFlx{(‑aði)}[1] \dicFlx{acc} \dicDirectTranslationCS{(vy)řezat závit}
\dicEntry[snitti] \dicTerm{snitti} \dicIPA{{s}{\textsubring{d}}{n}{\textsci}{h}{\textsubring{d}}{\textsci}} \dicPos{n}[2] \dicFlx{(‑s, ‑)}[14] \textbf{1.} \dicSynonym{ögn\smash{\textsuperscript{1}}} \dicDirectTranslationCS{zrnko, troška}  \textbf{2.} \dicFieldCat{techn.} \dicDirectTranslationCS{závitník}
\dicEntry[sníð] \dicTerm{sníð} \dicIPA{{s}{\textsubring{d}}{n}{i}{\textlengthmark}{\texttheta}} \dicPos{v} \dicFlx{ind praes sg 1 pers} \dicLink{sníða}
\dicEntry[sníða] \dicTerm{sníða} \dicIPA{{s}{\textsubring{d}}{n}{i}{\textlengthmark}{ð}{a}} \dicPos{v}[6] \dicFlx{(sníð, sneið, sniðum, sniði, sniðið)}[65] \dicFlx{acc} \textbf{1.} \dicSynonym{skera} \dicDirectTranslationCS{(na)stříhat} \dicExampleIS{sníða e‑ð af e‑u} \dicExampleCS{odstřihnout (co) od (čeho)}  \textbf{2.} \dicDirectTranslationCS{stříhat, zastřihnout, přistřihnout (šaty ap.)} \dicExampleIS{sníða buxurnar} \dicExampleCS{stříhat šaty}  \textbf{3.} \dicDirectTranslationCS{(u)tvořit, (z)formovat}  \textbf{4.} \dicFieldCat{poč.} \dicDirectTranslationCS{(z)formátovat (text ap.)}
\dicEntry[sníkill] \dicTerm{sník|ill} \dicIPA{{s}{\textsubring{d}}{n}{i}{\textlengthmark}{\r{\textObardotlessj}}{\textsci}{\textsubring{d}}{\textsubring{l}}} \dicPos{m}[6] \dicFlx{(‑ils, ‑lar)}[35] \dicFieldCat{biol.} \dicDirectTranslationCS{parazit}
\dicEntry[sníkja] \dicTerm{sník|ja} \dicIPA{{s}{\textsubring{d}}{n}{i}{\textlengthmark}{\r{\textObardotlessj}}{a}} \dicPos{v}[2] \dicFlx{(‑ti, ‑t)}[33] \dicFlx{acc} \dicSynonym{betla} \dicDirectTranslationCS{(vy)loudit, (vy)škemrat};  \dicPhraseIS{sníkja e‑ð af\,/\addthin hjá e‑m} \dicDirectTranslationCS{vyloudit (co) od\,/\addthin u~(koho)}
\dicEntry[sníkjudýr] \dicTerm{sníkju··dýr} \dicIPA{{s}{\textsubring{d}}{n}{i}{\textlengthmark}{\r{\textObardotlessj}}{\textscy}{\textsubring{d}}{i}{\textsubring{r}}} \dicPos{n}[2] \dicFlx{(‑s, ‑)}[5] \textbf{1.} \dicFieldCat{zool.} \dicDirectTranslationCS{cizopasník, parazit}  \textbf{2.} \dicDirectTranslationCS{příživník, příživnice, vyžírka}
\dicEntry[sníkjudýrafræði] \dicTerm{sníkju·dýra··fræð|i} \dicIPA{{s}{\textsubring{d}}{n}{i}{\textlengthmark}{\r{\textObardotlessj}}{\textscy}{\textsubring{d}}{i}{r}{a}{f}{r}{a}{i}{ð}{\textsci}} \dicPos{f}[3] \dicFlx{(‑i)}[3] \dicDirectTranslationCS{parazitologie}
\dicEntry[sníkjudýrafræðingur] \dicTerm{sníkju·dýra·fræð··ing|ur} \dicIPA{{s}\-{\textsubring{d}}\-{n}\-{i}\-{\textlengthmark}\-{\r{\textObardotlessj}}\-{\textscy}\-{\textsubring{d}}\-{i}\-{r}\-{a}\-{f}\-{r}\-{a}\-{i}\-{ð}\-{i}\-{\ng}\-{\r{g}}\-{\textscy}\-{\textsubring{r}}\-} \dicPos{m}[6] \dicFlx{(‑s, ‑ar)}[8] \dicDirectTranslationCS{parazitolog, parazitoložka}
\dicEntry[sníkjujurt] \dicTerm{sníkju··jurt} \dicIPA{{s}{\textsubring{d}}{n}{i}{\textlengthmark}{\r{\textObardotlessj}}{\textscy}{j}{\textscy}{\textsubring{r}}{\textsubring{d}}} \dicPos{f}[7] \dicFlx{(‑ar, ‑ir)}[1] \dicDirectTranslationCS{parazitní rostlina}
\dicEntry[sníkjur] \dicTerm{sníkjur} \dicIPA{{s}{\textsubring{d}}{n}{i}{\textlengthmark}{\r{\textObardotlessj}}{\textscy}{\textsubring{r}}} \dicPos{f}[1] \dicFlx{pl}[12] \dicDirectTranslationCS{příživnictví, parazitismus}
\dicEntry[sníkjuvespa] \dicTerm{sníkju··vesp|a} \dicIPA{{s}{\textsubring{d}}{n}{i}{\textlengthmark}{\r{\textObardotlessj}}{\textscy}{v}{\textepsilon}{s}{\textsubring{b}}{a}} \dicPos{f}[1] \dicFlx{(‑u, ‑ur)}[19] \dicFieldCat{zool.} \dicDirectTranslationCS{lumek} \textit{(l.~{\textLA{Ichneumonidae}})}  \dicsymPhoto\ 
\dicFigure{ds_image_snikjuvespa_0_1.jpg}{Sníkjuvespa}{Sníkjuvespa - Hectonichus, CC BY-SA 3.0}
\dicEntry[snípur] \dicTerm{sníp|ur} \dicIPA{{s}{\textsubring{d}}{n}{i}{\textlengthmark}{\textsubring{b}}{\textscy}{\textsubring{r}}} \dicPos{m}[9] \dicFlx{(‑s, ‑ir)}[8] \dicFieldCat{anat.} \dicDirectTranslationCS{poštěváček, klitoris}
\dicEntry[snjall] \dicTerm{snjall} \dicIPA{{s}{\textsubring{d}}{n}{j}{a}{\textsubring{d}}{\textsubring{l}}} \dicPos{adj}[5] \dicFlx{(f snjöll)}[11] \dicDirectTranslationCS{geniální, mistrovský, brilantní, bravurní, vynikající} \dicExampleIS{snjöll lausn} \dicExampleCS{geniální řešení}
\dicEntry[snjallræði] \dicTerm{snjall··ræði} \dicIPA{{s}{\textsubring{d}}{n}{j}{a}{\textsubring{d}}{l}{r}{a}{i}{ð}{\textsci}} \dicPos{n}[2] \dicFlx{(‑s)}[20] \dicDirectTranslationCS{vynikající\,/\addthin geniální nápad}
\dicEntry[snjallsími] \dicTerm{snjall··sím|i} \dicIPA{{s}{\textsubring{d}}{n}{j}{a}{\textsubring{d}}{\textsubring{l}}{s}{i}{m}{\textsci}} \dicPos{m}[1] \dicFlx{(‑a, ‑ar)}[1] \dicDirectTranslationCS{smartphone, chytrý telefon}
\dicEntry[snjáður] \dicTerm{snjáður} \dicIPA{{s}{\textsubring{d}}{n}{j}{au}{\textlengthmark}{ð}{\textscy}{\textsubring{r}}} \dicPos{adj}[2]\dicFlx{}[-12] \dicDirectTranslationCS{ošoupaný, ošuntělý, obnošený} \dicExampleIS{Skyrtan er orðin snjáð og upplituð.} \dicExampleCS{Košile je už ošuntělá a~vybledlá.}
\dicEntry[snjáldur] \dicTerm{snjáldur} \dicIPA{{s}{\textsubring{d}}{n}{j}{au}{l}{\textsubring{d}}{\textscy}{\textsubring{r}}} \dicPos{n}[2] \dicFlx{(‑s, ‑)}[25] \dicSynonym{trýni} \dicDirectTranslationCS{čumák, čenich, rypák}
\dicEntry[snjóa] \dicTerm{snjó|a} \dicIPA{{s}{\textsubring{d}}{n}{j}{ou}{\textlengthmark}{a}} \dicPos{v}[1] \dicFlx{(‑aði)}[70] \dicFlx{impers} \dicPhraseIS{það snjóar} \dicDirectTranslationCS{sněží, chumelí, padá sníh} \dicExampleIS{Það snjóar í fjöll.} \dicExampleCS{Na horách sněží.}
\dicEntry[snjóavetur] \dicTerm{snjóa··vet|ur} \dicIPA{{s}{\textsubring{d}}{n}{j}{ou}{\textlengthmark}{a}{v}{\textepsilon}{\textsubring{d}}{\textscy}{\textsubring{r}}} \dicPos{m}[13] \dicFlx{(‑rar, ‑ur)}[6] \dicDirectTranslationCS{zima se sněhem, sněžná zima}
\dicEntry[snjóbíll] \dicTerm{snjó··bíl|l} \dicIPA{{s}{\textsubring{d}}{n}{j}{ou}{\textlengthmark}{\textsubring{b}}{i}{\textsubring{d}}{\textsubring{l}}} \dicPos{m}[6] \dicFlx{(‑s, ‑ar)}[48] \dicDirectTranslationCS{(sněžná) rolba}
\dicEntry[snjóblinda] \dicTerm{snjó··blind|a} \dicIPA{{s}{\textsubring{d}}{n}{j}{ou}{\textlengthmark}{\textsubring{b}}{l}{\textsci}{n}{\textsubring{d}}{a}} \dicPos{f}[1] \dicFlx{(‑u)}[5] \dicDirectTranslationCS{sněžná slepota}
\dicEntry[snjóblindur] \dicTerm{snjó··blindur} \dicIPA{{s}{\textsubring{d}}{n}{j}{ou}{\textlengthmark}{\textsubring{b}}{l}{\textsci}{n}{\textsubring{d}}{\textscy}{\textsubring{r}}} \dicPos{adj}[2]\dicFlx{}[-17] \dicDirectTranslationCS{oslepený sněhem, mající sněžnou slepotu}
\dicEntry[snjóbolti] \dicTerm{snjó··bolt|i} \dicIPA{{s}{\textsubring{d}}{n}{j}{ou}{\textlengthmark}{\textsubring{b}}{\textopeno}{\textsubring{l}}{\textsubring{d}}{\textsci}} \dicPos{m}[1] \dicFlx{(‑a, ‑ar)}[1] \dicDirectTranslationCS{(sněhová) koule}
\dicEntry[snjóbretti] \dicTerm{snjó··bretti} \dicIPA{{s}{\textsubring{d}}{n}{j}{ou}{\textlengthmark}{\textsubring{b}}{r}{\textepsilon}{h}{\textsubring{d}}{\textsci}} \dicPos{n}[2] \dicFlx{(‑s, ‑)}[14] \dicFieldCat{sport.} \dicDirectTranslationCS{snowboard}
\dicEntry[snjóbræðsla] \dicTerm{snjó··bræðsl|a} \dicIPA{{s}{\textsubring{d}}{n}{j}{ou}{\textlengthmark}{\textsubring{b}}{r}{a}{i}{ð}{s}{\textsubring{d}}{l}{a}} \dicPos{f}[1] \dicFlx{(‑u, ‑ur)}[13] \dicDirectTranslationCS{odstraňování sněhu} \dicIndirectTranslationCS{(pomocí tepla\,/\addthin vytápění)}
\dicEntry[snjóbylur] \dicTerm{snjó··byl|ur} \dicIPA{{s}{\textsubring{d}}{n}{j}{ou}{\textlengthmark}{\textsubring{b}}{\textsci}{l}{\textscy}{\textsubring{r}}} \dicPos{m}[9] \dicFlx{(‑s\,/\addthin ‑jar, ‑jir)}[32] \dicSynonym{bylur} \dicDirectTranslationCS{sněhová vichřice}
\dicEntry[snjódekk] \dicTerm{snjó··dekk} \dicIPA{{s}{\textsubring{d}}{n}{j}{ou}{\textlengthmark}{\textsubring{d}}{\textepsilon}{h}{\r{g}}} \dicPos{n}[2] \dicFlx{(‑s, ‑)}[13] \dicDirectTranslationCS{zimní\,/\addthin sněžná pneumatika}
\dicEntry[snjóflóð] \dicTerm{snjó··flóð} \dicIPA{{s}{\textsubring{d}}{n}{j}{ou}{\textlengthmark}{f}{l}{ou}{\texttheta}} \dicPos{n}[2] \dicFlx{(‑s, ‑)}[5] \dicDirectTranslationCS{(sněhová) lavina}
\dicEntry[snjóflóðahætta] \dicTerm{snjó·flóða··hætt|a} \dicIPA{{s}{\textsubring{d}}{n}{j}{ou}{\textlengthmark}{f}{l}{ou}{ð}{a}{h}{a}{i}{h}{\textsubring{d}}{a}} \dicPos{f}[1] \dicFlx{(‑u)}[5] \dicDirectTranslationCS{lavino\-vé nebezpečí, nebezpečí lavin}
\dicEntry[snjóflyksa] \dicTerm{snjó··flyks|a} \dicIPA{{s}{\textsubring{d}}{n}{j}{ou}{\textlengthmark}{f}{l}{\textsci}{\r{g}}{s}{a}} \dicPos{f}[1] \dicFlx{(‑u, ‑ur)}[13] \dicDirectTranslationCS{sněhová vločka}
\dicEntry[snjógleraugu] \dicTerm{snjó··gler·augu} \dicIPA{{s}{\textsubring{d}}{n}{j}{ou}{\textlengthmark}{\r{g}}{l}{\textepsilon}{r}{\oe i}{\textbabygamma}{\textscy}} \dicPos{n}[1] \dicFlx{pl}[3] \dicDirectTranslationCS{sněhové brýle}
\dicEntry[snjógæs] \dicTerm{snjó··gæs} \dicIPA{{s}{\textsubring{d}}{n}{j}{ou}{\textlengthmark}{\r{\textObardotlessj}}{a}{i}{s}} \dicPos{f}[7] \dicFlx{(‑ar, ‑ir)}[1] \dicFieldCat{zool.} \dicDirectTranslationCS{husa sněžní} \textit{(l.~{\textLA{Chen caerulescens}})}  \dicsymPhoto\ 
\dicFigure{ds_image_snjogaes_0_1.jpg}{Snjógæs}{Snjógæs - Walter Siegmund, CC BY-SA 3.0}
\dicEntry[snjóhús] \dicTerm{snjó··hús} \dicIPA{{s}{\textsubring{d}}{n}{j}{ou}{\textlengthmark}{h}{u}{s}} \dicPos{n}[2] \dicFlx{(‑s, ‑)}[5] \dicDirectTranslationCS{iglú}
\dicEntry[snjóhvítur] \dicTerm{snjó··hvítur} \dicIPA{{s}{\textsubring{d}}{n}{j}{ou}{\textlengthmark}{k\smash{\textsuperscript{h}}}{v}{i}{\textsubring{d}}{\textscy}{\textsubring{r}}} \dicPos{adj}[1]\dicFlx{}[-6] \dicDirectTranslationCS{sněhobílý, sněhově bílý}
\dicEntry[snjókarl] \dicTerm{snjó··karl} \dicIPA{{s}{\textsubring{d}}{n}{j}{ou}{\textlengthmark}{k\smash{\textsuperscript{h}}}{a}{r}{\textsubring{d}}{\textsubring{l}}} \dicPos{m}[4] \dicFlx{(‑s, ‑ar)}[4] \dicDirectTranslationCS{sněhulák}
\dicEntry[snjókast] \dicTerm{snjó··|kast} \dicIPA{{s}{\textsubring{d}}{n}{j}{ou}{\textlengthmark}{k\smash{\textsuperscript{h}}}{a}{s}{\textsubring{d}}} \dicPos{n}[2] \dicFlx{(‑kasts, ‑köst)}[8] \dicDirectTranslationCS{koulování, koulovačka}
\dicEntry[snjókeðja] \dicTerm{snjó··keðj|a} \dicIPA{{s}{\textsubring{d}}{n}{j}{ou}{\textlengthmark}{c\smash{\textsuperscript{h}}}{\textepsilon}{ð}{j}{a}} \dicPos{f}[1] \dicFlx{(‑u, ‑ur)}[7] \dicDirectTranslationCS{sněhový řetěz}
\dicEntry[snjókerling] \dicTerm{snjó··kerl·ing} \dicIPA{{s}{\textsubring{d}}{n}{j}{ou}{\textlengthmark}{c\smash{\textsuperscript{h}}}{\textepsilon}{r}{\textsubring{d}}{l}{i}{\ng}{\r{g}}} \dicPos{f}[4] \dicFlx{(‑ar, ‑ar)}[5] \dicDirectTranslationCS{sněhulačka} \dicIndirectTranslationCS{(panna vyrobená ze sněhu)}
\dicEntry[snjókoma] \dicTerm{snjó··kom|a} \dicIPA{{s}{\textsubring{d}}{n}{j}{ou}{\textlengthmark}{k\smash{\textsuperscript{h}}}{\textopeno}{m}{a}} \dicPos{f}[1] \dicFlx{(‑u)}[5] \dicDirectTranslationCS{sněžení, chumelení}
\dicEntry[snjókorn] \dicTerm{snjó··korn} \dicIPA{{s}{\textsubring{d}}{n}{j}{ou}{\textlengthmark}{k\smash{\textsuperscript{h}}}{\textopeno}{r}{\textsubring{d}}{\textsubring{n}}} \dicPos{n}[2] \dicFlx{(‑s, ‑)}[5] \dicDirectTranslationCS{sněhová vločka}
\dicEntry[snjólaus] \dicTerm{snjó··laus} \dicIPA{{s}{\textsubring{d}}{n}{j}{ou}{\textlengthmark}{l}{\oe i}{s}} \dicPos{adj}[5]\dicFlx{}[-1] \dicDirectTranslationCS{(jsoucí) bez sněhu, nezasněžený}
\dicEntry[snjólína] \dicTerm{snjó··lín|a} \dicIPA{{s}{\textsubring{d}}{n}{j}{ou}{\textlengthmark}{l}{i}{n}{a}} \dicPos{f}[1] \dicFlx{(‑u, ‑ur)}[7] \dicFieldCat{meteo.} \dicSynonym{snælína} \dicDirectTranslationCS{sněžná čára}
\dicEntry[snjóplógur] \dicTerm{snjó··plóg|ur} \dicIPA{{s}{\textsubring{d}}{n}{j}{ou}{\textlengthmark}{p\smash{\textsuperscript{h}}}{l}{ou}{\textscy}{\textsubring{r}}} \dicPos{m}[6] \dicFlx{(‑s, ‑ar)}[22] \dicDirectTranslationCS{sněžný pluh}
\dicEntry[snjór] \dicTerm{snjó|r} \dicsymFrequent\  \dicIPA{{s}{\textsubring{d}}{n}{j}{ou}{\textlengthmark}{\textsubring{r}}} \dicPos{m}[6] \dicFlx{(‑s\,/\addthin ‑var\,/\addthin ‑ar, ‑ar\,/\addthin ‑var)}[78] \dicDirectTranslationCS{sníh} \dicExampleIS{Snjórinn fellur til jarðar.} \dicExampleCS{Sníh padá na zem.}
\dicEntry[snjóruðningur] \dicTerm{snjó··ruð·ning|ur} \dicIPA{{s}{\textsubring{d}}{n}{j}{ou}{\textlengthmark}{r}{\textscy}{ð}{\textsubring{n}}{i}{\ng}{\r{g}}{\textscy}{\textsubring{r}}} \dicPos{m}[6] \dicFlx{(‑s, ‑ar)}[8] \dicDirectTranslationCS{odklízení\,/\addthin odhrnování sněhu}
\dicEntry[snjóskafl] \dicTerm{snjó··skafl} \dicIPA{{s}{\textsubring{d}}{n}{j}{ou}{\textlengthmark}{s}{\r{g}}{a}{\textsubring{b}}{\textsubring{l}}} \dicPos{m}[4] \dicFlx{(‑s, ‑ar)}[4] \dicDirectTranslationCS{(sněhová) závěj}
\dicEntry[snjósleði] \dicTerm{snjó··sleð|i} \dicIPA{{s}{\textsubring{d}}{n}{j}{ou}{\textlengthmark}{s}{\textsubring{d}}{l}{\textepsilon}{ð}{\textsci}} \dicPos{m}[1] \dicFlx{(‑a, ‑ar)}[1] \dicDirectTranslationCS{sněžný skútr}
\dicEntry[snjótittlingur] \dicTerm{snjó··titt·ling|ur} \dicIPA{{s}{\textsubring{d}}{n}{j}{ou}{\textlengthmark}{t\smash{\textsuperscript{h}}}{\textsci}{h}{\textsubring{d}}{l}{i}{\ng}{\r{g}}{\textscy}{\textsubring{r}}} \dicPos{m}[6] \dicFlx{(‑s, ‑ar)}[8] \dicFieldCat{zool.} \dicDirectTranslationCS{sněhule severní} \dicIndirectTranslationCS{(označení ptáka v~zimním období)} \textit{(l.~{\textLA{Plectrophenax nivalis}})}  \dicsymPhoto\  \dicAntonym{sólskríkja}
\dicFigure{ds_image_snjotittlingur_0_1.jpg}{Snjótittlingur}{Snjótittlingur - U.S. Fish and Wildlife Service, PD}
\dicEntry[snjóugur] \dicTerm{snjó··ugur} \dicIPA{{s}{\textsubring{d}}{n}{j}{ou}{\textlengthmark}{\textscy}{\textbabygamma}{\textscy}{\textsubring{r}}} \dicPos{adj}[1]\dicFlx{}[-8] \dicDirectTranslationCS{zasněžený, bílý} \dicExampleIS{snjóug fjöll} \dicExampleCS{zasněžené hory}
\dicEntry[snjóþota] \dicTerm{snjó··þot|a} \dicIPA{{s}{\textsubring{d}}{n}{j}{ou}{\textlengthmark}{\texttheta}{\textopeno}{\textsubring{d}}{a}} \dicPos{f}[1] \dicFlx{(‑u, ‑ur)}[19] \dicDirectTranslationCS{bob} \dicIndirectTranslationCS{(nízké plastové řiditelné sáňky)}
\dicEntry[snjóþungur] \dicTerm{snjó··|þungur} \dicIPA{{s}{\textsubring{d}}{n}{j}{ou}{\textlengthmark}{\texttheta}{u}{\ng}{\r{g}}{\textscy}{\textsubring{r}}} \dicPos{adj}[10] \dicFlx{(comp ‑þyngri, sup ‑þyng\-st\-ur)}[7] \dicPhraseIS{snjóþungur staður} \dicDirectTranslationCS{místo, kde hodně sněží};  \dicPhraseIS{snjóþungur vetur} \dicDirectTranslationCS{zima s~velkým množstvím sněhu}
\dicEntry[snobb] \dicTerm{snobb} \dicIPA{{s}{\textsubring{d}}{n}{\textopeno}{\textsubring{b}}{\textlengthmark}} \dicPos{n}[2] \dicFlx{(‑s)}[2] \textbf{1.} \dicDirectTranslationCS{snobství}  \textbf{2.} \dicSynonym*{snobbari} \dicDirectTranslationCS{snob(ka)}
\dicEntry[snobba] \dicTerm{snobb|a} \dicIPA{{s}{\textsubring{d}}{n}{\textopeno}{\textsubring{b}}{\textlengthmark}{a}} \dicPos{v}[1] \dicFlx{(‑aði)}[1] \dicPhraseIS{snobba fyrir e‑m} \dicDirectTranslationCS{podlézat (komu), pochlebovat (komu)} \dicExampleIS{snobba fyrir fína fólkinu} \dicExampleCS{podlézat bohatým lidem}
\dicEntry[snobbaður] \dicTerm{snobb··|aður} \dicIPA{{s}{\textsubring{d}}{n}{\textopeno}{\textsubring{b}}{\textlengthmark}{a}{ð}{\textscy}{\textsubring{r}}} \dicPos{adj}[3] \dicFlx{(f ‑uð)}[3] \dicDirectTranslationCS{snobský}
\dicEntry[snoða] \dicTerm{snoð|a} \dicIPA{{s}{\textsubring{d}}{n}{\textopeno}{\textlengthmark}{ð}{a}} \dicPos{v}[1] \dicFlx{(‑aði)}[1] \dicFlx{acc} \dicDirectTranslationCS{(vy)holit, oholit, přistřihnout\,/\addthin stříhat nakrátko}
\dicEntry[snoðir] \dicTerm{snoðir} \dicIPA{{s}{\textsubring{d}}{n}{\textopeno}{\textlengthmark}{ð}{\textsci}{\textsubring{r}}} \dicPos{f}[12] \dicFlx{pl}[10] \dicSynonym{ávæningur} \dicDirectTranslationCS{náznak, stopa};  \dicPhraseIS{komast á snoðir um e‑ð} \dicFlx{refl} \dicDirectTranslationCS{dozvědět se o~(čem)} \dicExampleIS{komast á snoðir um ólöglega starfsemi} \dicExampleCS{dozvědět se o~nelegální činnosti}
\dicEntry[snoðklippa] \dicTerm{snoð··klipp|a} \dicIPA{{s}{\textsubring{d}}{n}{\textopeno}{ð}{k\smash{\textsuperscript{h}}}{l}{\textsci}{h}{\textsubring{b}}{a}} \dicPos{v}[2] \dicFlx{(‑ti, ‑t)}[69] \dicFlx{acc} \dicDirectTranslationCS{přistřihnout\,/\addthin stříhat nakrátko (vlasy ap.)}
\dicEntry[snoppa] \dicTerm{snopp|a} \dicIPA{{s}{\textsubring{d}}{n}{\textopeno}{h}{\textsubring{b}}{a}} \dicPos{f}[1] \dicFlx{(‑u, ‑ur)}[19] \dicSynonym{trýni} \dicDirectTranslationCS{čenich, čumák}
\begin{xtolerant}{}{1pt}
\dicEntry[snoppufríður] \dicTerm{snoppu··fríður}\addthinS\dicIPA{{s}{\textsubring{d}}{n}{\textopeno}{h}{\textsubring{b}}{\textscy}{f}{r}{i}{ð}{\textscy}{\textsubring{r}}}\addthinS\dicPos{adj}[2]\dicFlx{}[-6]\addthinS\dicDirectTranslationCS{pohledný\,/\addthinSS líbez\-ný v~obličeji}
\end{xtolerant}
\dicEntry[snoppunga] \dicTerm{snopp··ung|a} \dicIPA{{s}{\textsubring{d}}{n}{\textopeno}{h}{\textsubring{b}}{u}{\ng}{\r{g}}{a}} \dicPos{v}[1] \dicFlx{(‑aði)}[1] \dicFlx{acc} \dicDirectTranslationCS{(z)fackovat, dát políček}
\dicEntry[snoppungur] \dicTerm{snopp··ung|ur} \dicIPA{{s}{\textsubring{d}}{n}{\textopeno}{h}{\textsubring{b}}{u}{\ng}{\r{g}}{\textscy}{\textsubring{r}}} \dicPos{m}[6] \dicFlx{(‑s, ‑ar)}[8] \dicSynonym{löðrungur} \dicDirectTranslationCS{facka, políček, pohlavek} \dicExampleIS{gefa honum gildan snoppung} \dicExampleCS{dát mu pořádnou facku}
\dicEntry[snortinn] \dicTerm{snortinn} \dicIPA{{s}{\textsubring{d}}{n}{\textopeno}{\textsubring{r}}{\textsubring{d}}{\textsci}{\textsubring{n}}} \dicPos{adj}[6]\dicFlx{}[-2] \dicDirectTranslationCS{dojatý, pohnutý};  \dicPhraseIS{vera djúpt snortinn af e‑u} \dicDirectTranslationCS{být hluboce dojatý (čím)}
\dicEntry[snotur] \dicTerm{snotur} \dicIPA{{s}{\textsubring{d}}{n}{\textopeno}{\textlengthmark}{\textsubring{d}}{\textscy}{\textsubring{r}}} \dicPos{adj}[9] \dicFlx{(f ‑)}[1] \dicSynonym{laglegur} \dicDirectTranslationCS{půvabný, pohledný, líbezný, sličný} \dicExampleIS{snotur stúlka} \dicExampleCS{půvabná dívka}
\dicEntry[snóker] \dicTerm{snóker} \dicIPA{{s}{\textsubring{d}}{n}{ou}{\textlengthmark}{\r{\textObardotlessj}}{\textepsilon}{\textsubring{r}}} \dicPos{m}[4] \dicFlx{(‑s)}[18] \dicFieldCat{sport.} \dicDirectTranslationCS{snooker} \dicIndirectTranslationCS{(kulečníková hra)}
\dicEntry[snót] \dicTerm{snót} \dicIPA{{s}{\textsubring{d}}{n}{ou}{\textlengthmark}{\textsubring{d}}} \dicPos{f}[7] \dicFlx{(‑ar, ‑ir)}[1] \dicLangCat{básn.} \dicSynonym{kona} \dicDirectTranslationCS{dívka, žena}
\dicEntry[snubbóttur] \dicTerm{snubb··óttur} \dicIPA{{s}{\textsubring{d}}{n}{\textscy}{\textsubring{b}}{\textlengthmark}{ou}{h}{\textsubring{d}}{\textscy}{\textsubring{r}}} \dicPos{adj}[1]\dicFlx{}[-10] \textbf{1.} \dicSynonym{endasleppur} \dicDirectTranslationCS{náhle\,/\addthin nenadále ukončený} \dicExampleIS{snubbótt ræða} \dicExampleCS{náhle ukončená řeč}  \textbf{2.} \dicSynonym{önugur} \dicDirectTranslationCS{strohý, příkrý (v~řeči ap.)} \dicExampleIS{snubbóttur í tali} \dicExampleCS{příkrý v~řeči}
\dicEntry[snudda] \dicTerm{snudd|a\smash{\textsuperscript{1}}} \dicIPA{{s}{\textsubring{d}}{n}{\textscy}{\textsubring{d}}{\textlengthmark}{a}} \dicPos{f}[1] \dicFlx{(‑u, ‑ur)}[19] \dicLangCat{dět.} \dicDirectTranslationCS{dudlík, šidítko}
\dicEntry[snudda] \dicTerm{snudd|a\smash{\textsuperscript{2}}} \dicIPA{{s}{\textsubring{d}}{n}{\textscy}{\textsubring{d}}{\textlengthmark}{a}} \dicPos{v}[1] \dicFlx{(‑aði)}[1] \dicDirectTranslationCS{potulovat se, slídit kolem}
\dicEntry[snuð] \dicTerm{snuð} \dicIPA{{s}{\textsubring{d}}{n}{\textscy}{\textlengthmark}{\texttheta}} \dicPos{n}[2] \dicFlx{(‑s, ‑)}[5] \dicSynonym{tútta} \dicDirectTranslationCS{dudlík, šidítko}
\dicEntry[snuða] \dicTerm{snuð|a} \dicIPA{{s}{\textsubring{d}}{n}{\textscy}{\textlengthmark}{ð}{a}} \dicPos{v}[1] \dicFlx{(‑aði)}[1] \dicFlx{acc} \dicDirectTranslationCS{ošidit, podvést, ošulit};  \dicPhraseIS{snuða e‑n um e‑ð} \dicDirectTranslationCS{ošidit (koho) o~(co)} \dicExampleIS{Strákurinn snuðaði mig um hundrað krónur.} \dicExampleCS{Ten kluk mě ošulil o~sto korun.}
\dicEntry[snuðra] \dicTerm{snuðr|a} \dicIPA{{s}{\textsubring{d}}{n}{\textscy}{ð}{r}{a}} \dicPos{v}[1] \dicFlx{(‑aði)}[1] \dicSynonym{hnusa} \dicDirectTranslationCS{čmuchat, čenichat, slídit} \dicExampleIS{Hundurinn snuðraði allt í kringum húsið.} \dicExampleCS{Pes čenichal všude kolem domu.};  \dicPhraseIS{snuðra e‑ð uppi} \dicSynonym{finna} \dicDirectTranslationCS{vyčenichat (co), vyslídit (co)} \dicExampleIS{snuðra uppi ný sönnunargögn} \dicExampleCS{vyslídit nové důkazy}
\dicEntry[snuðrari] \dicTerm{snuðr··ar|i} \dicIPA{{s}{\textsubring{d}}{n}{\textscy}{ð}{r}{a}{r}{\textsci}} \dicPos{m}[1] \dicFlx{(‑a, ‑ar)}[13] \dicDirectTranslationCS{slídič(ka), slídil(ka), špeh}
\dicEntry[snuður] \dicTerm{snuður} \dicIPA{{s}{\textsubring{d}}{n}{\textscy}{\textlengthmark}{ð}{\textscy}{\textsubring{r}}} \dicPos{n}[2] \dicFlx{(‑s)}[28] \dicDirectTranslationCS{slídění, čmuchání}
\dicEntry[snupra] \dicTerm{snupr|a\smash{\textsuperscript{1}}} \dicIPA{{s}{\textsubring{d}}{n}{\textscy}{\textlengthmark}{\textsubring{b}}{r}{a}} \dicPos{f}[1] \dicFlx{(‑u, ‑ur)}[7] \dicDirectTranslationCS{(po)kárání, důtka}
\dicEntry[snupra] \dicTerm{snupr|a\smash{\textsuperscript{2}}} \dicIPA{{s}{\textsubring{d}}{n}{\textscy}{\textlengthmark}{\textsubring{b}}{r}{a}} \dicPos{v}[1] \dicFlx{(‑aði)}[1] \dicFlx{acc} \dicSynonym{ávíta} \dicDirectTranslationCS{(po)kárat} \dicExampleIS{snupra e‑n fyrir e‑ð} \dicExampleCS{pokárat (koho) za (co)}
\dicEntry[snurða] \dicTerm{snurð|a} \dicIPA{{s}{\textsubring{d}}{n}{\textscy}{r}{ð}{a}} \dicPos{f}[1] \dicFlx{(‑u, ‑ur)}[19] \textbf{1.} \dicSynonym{hnökri} \dicDirectTranslationCS{uzlík, zašmodrchání (na niti ap.)}  \textbf{2.} \dicLangCat{přen.} \dicDirectTranslationCS{zádrhel, nesnáz};  \dicPhraseIS{það er hlaupin snurða á þráðinn} \dicLangCat{přen.} \dicDirectTranslationCS{vzniká zádrhel}
\dicEntry[snurðóttur] \dicTerm{snurð··óttur} \dicIPA{{s}{\textsubring{d}}{n}{\textscy}{r}{ð}{ou}{h}{\textsubring{d}}{\textscy}{\textsubring{r}}} \dicPos{adj}[1]\dicFlx{}[-10] \textbf{1.} \dicDirectTranslationCS{zašmodrchaný, zauzlený (nit ap.)}  \textbf{2.} \dicLangCat{přen.} \dicSynonym{torveldur} \dicDirectTranslationCS{plný zádrhelů, zamotaný}
\dicEntry[snurðulaus] \dicTerm{snurðu··laus} \dicIPA{{s}{\textsubring{d}}{n}{\textscy}{r}{ð}{\textscy}{l}{\oe i}{s}} \dicPos{adj}[5]\dicFlx{}[-1] \textbf{1.} \dicDirectTranslationCS{nezašmodrchaný, (jsoucí) bez uzlíků (nit ap.)}  \textbf{2.} \dicLangCat{přen.} \dicDirectTranslationCS{hladký, (jsoucí) bez zádrhelů (průběh ap.)} \dicExampleIS{snurðulaust samstarf} \dicExampleCS{hladká spolupráce}
\dicEntry[snurfusa] \dicTerm{snur··fus|a} \dicIPA{{s}{\textsubring{d}}{n}{\textscy}{\textsubring{r}}{f}{\textscy}{s}{a}} \dicPos{v}[1] \dicFlx{(‑aði)}[1] \dicFlx{acc} \dicSynonym{snyrta} \dicDirectTranslationCS{zkrášlit, vylepšit (vzhled ap.)}
\dicEntry[snurpinót] \dicTerm{snurpi··|nót} \dicIPA{{s}{\textsubring{d}}{n}{\textscy}{\textsubring{r}}{\textsubring{b}}{\textsci}{n}{ou}{\textsubring{d}}} \dicPos{f}[7] \dicFlx{(‑nótar, ‑nótir\,/\addthin ‑nætur)}[27] \dicFieldCat{nám.} \dicDirectTranslationCS{vakový nevod}
\dicEntry[snúa] \dicTerm{snúa} \dicsymFrequent\  \dicIPA{{s}{\textsubring{d}}{n}{u}{\textlengthmark}{a}} \dicPos{v}[5] \dicFlx{(sný, sneri\,/\addthin snéri, snerum\,/\addthin snérum, sneri\,/\addthin snéri, snúið)}[4] \dicFlx{dat\,/\addthin acc} \textbf{1.} \dicFlx{dat} \dicSynonym*{breyta um stefnu e‑s} \dicDirectTranslationCS{(o)točit, obrátit, obracet} \dicExampleIS{snúa lyklinum} \dicExampleCS{otočit klíčem};  \dicPhraseIS{snúa sér í hring} \dicDirectTranslationCS{otočit se kolem dokola}  \textbf{2.} \dicFlx{dat} \dicSynonym{velta\smash{\textsuperscript{3}}} \dicDirectTranslationCS{obrátit, obracet, převrátit, převracet} \dicExampleIS{snúa heyinu} \dicExampleCS{obrátit seno};  \dicPhraseIS{snúa sér} \dicDirectTranslationCS{převracet se, obracet se (ve spánku ap.)}  \textbf{3.} \dicSynonym{horfa} \dicDirectTranslationCS{ležet, sedět, stát, směřovat, být obrácený (směrem k~něčemu)} \dicExampleIS{Dyrnar snúa í norður.} \dicExampleCS{Dveře leží na sever.}  \textbf{4.} \dicDirectTranslationCS{vrátit se, vracet se, obrátit se (zpět)} \dicExampleIS{snúa heim} \dicExampleCS{vracet se domů}  \textbf{5.} \dicFlx{dat} \dicDirectTranslationCS{přeložit, překládat (knihu ap.)} \dicExampleIS{snúa kvæðinu úr íslensku á ensku} \dicExampleCS{přeložit báseň z~islandštiny do angličtiny}  \textbf{6.} \dicFlx{dat} \dicDirectTranslationCS{přesvědčit, přesvědčovat, (z)měnit (čí) názor} \dicExampleIS{snúa honum í málinu} \dicExampleCS{přesvědčit ho, aby změnil názor}  \textbf{7.} \dicFlx{dat} \dicDirectTranslationCS{svinout, svíjet, (s)motat} \dicExampleIS{snúa reipi} \dicExampleCS{svinout provaz}  \textbf{8.} \dicPhraseIS{snúa sig} \dicFlx{acc} \dicDirectTranslationCS{vyvrtnout si, vymknout si} \dicExampleIS{snúa sig um ökklann} \dicExampleCS{vymknout si kotník};  \dicIdiom{snúa}[að]{ \dicPhraseIS{snúa sér að e‑u}} \dicDirectTranslationCS{pustit se do (čeho), dát se do (čeho)} \dicExampleIS{snúa sér að verkefninu} \dicExampleCS{pustit se do zadání};  \dicIdiom{snúa}[aftur]{ \dicPhraseIS{snúa aftur}} \dicDirectTranslationCS{vrátit se zpět};  \dicIdiom{snúa}[á]{ \dicPhraseIS{snúa á e‑n}} \dicDirectTranslationCS{oklamat (koho), napálit (koho), obelstít (koho)} \dicExampleIS{snúa illilega á hana í viðskiptum} \dicExampleCS{šeredně ji obelstít při obchodu}; { \dicPhraseIS{snúa e‑u á e‑ð}} \dicDirectTranslationCS{přeložit (co) do (čeho)} \dicIndirectTranslationCS{(vyjádřit v~jiném jazyce)};  \dicIdiom{snúa}[frá]{ \dicPhraseIS{snúa frá e‑u}} \dicDirectTranslationCS{být obrácený směrem od (čeho)} \dicExampleIS{Glugginn snýr frá götunni.} \dicExampleCS{Okno je obrácené směrem od ulice.};  \dicIdiom{snúa}[niður]{ \dicPhraseIS{snúa e‑n niður}} \dicDirectTranslationCS{přemoct (koho), překonat (koho)};  \dicIdiom{snúa}[saman]{ \dicPhraseIS{snúa e‑ð saman}} \dicDirectTranslationCS{splést\,/\addthin splétat (co) (provázky ap.)};  \dicIdiom{snúa}[til]{ \dicPhraseIS{snúa til e‑rs}} \dicDirectTranslationCS{obrátit se na (koho), oslovit (koho)};  \dicIdiom{snúa}[undan]{ \dicPhraseIS{snúa sér undan}} \dicDirectTranslationCS{otočit se, odvrátit se} \dicIndirectTranslationCS{(uhnout pohledem)};  \dicIdiom{snúa}[upp á]{ \dicPhraseIS{snúa upp á e‑ð}} \dicDirectTranslationCS{vykroutit (co) (ruku ap.)}; { \dicPhraseIS{snúa upp á sig}} \dicDirectTranslationCS{urazit se, urážet se};  \dicIdiom{snúa}[upp í]{ \dicPhraseIS{snúa e‑u upp í grín\,/\addthin spaug}} \dicDirectTranslationCS{obrátit (co) v~žert};  \dicIdiom{snúa}[út úr]{ \dicPhraseIS{snúa sig út úr e‑u}} \dicLangCat{přen.} \dicDirectTranslationCS{vykroutit se z~(čeho)} \dicExampleIS{snúa sig út úr vandanum} \dicExampleCS{vykroutit se z~potíží}; { \dicPhraseIS{snúa út úr fyrir e‑m}} \dicDirectTranslationCS{překroutit\,/\addthin zkomolit (čí) slova};  \dicIdiom{snúa}[við]{ \dicPhraseIS{snúa baki(nu) við e‑m}} \dicLangCat{přen.} \dicDirectTranslationCS{otočit se ke (komu) zády}; { \dicPhraseIS{snúa e‑u við}} \dicDirectTranslationCS{obrátit\,/\addthin obracet (co), otočit\,/\addthin otáčet (co)} \dicExampleIS{snúa plötunni við} \dicExampleCS{otočit desku}; { \dicPhraseIS{snúa sér við}} \dicDirectTranslationCS{obrátit se, obracet se, otočit se, otáčet se}; { \dicPhraseIS{snúa við}} \dicDirectTranslationCS{vrátit se, vracet se}; { \dicPhraseIS{snúa við blaðinu}} \dicLangCat{přen.} \dicDirectTranslationCS{obrátit list, nabrat (lepší) směr};  \dicIdiom{snúast}{ \dicPhraseIS{snúast}} \dicFlx{refl} {\textbf{a.}} \dicDirectTranslationCS{točit se, otáčet se, obracet se} \dicExampleIS{snúast í hringi} \dicExampleCS{točit se dokola};  {\textbf{b.}} \dicDirectTranslationCS{obrátit se, přestoupit, konvertovat (na víru ap.)} \dicExampleIS{snúast til kristindóms} \dicExampleCS{obrátit se na křesťanství}; { \dicPhraseIS{e‑m snýst hugur}} \dicFlx{refl impers} \dicDirectTranslationCS{(kdo) mění názor};  \dicIdiom{snúast}[fyrir]{ \dicPhraseIS{snúast fyrir e‑n}} \dicFlx{refl} \dicDirectTranslationCS{obstarat\,/\addthin obstarávat pochůzky pro (koho)};  \dicIdiom{snúast}[gegn]{ \dicPhraseIS{snúast gegn e‑u}} \dicFlx{refl} \dicDirectTranslationCS{vystoupit proti (čemu), obrátit se proti (čemu)} \dicExampleIS{snúast gegn stefnu flokksins} \dicExampleCS{vystoupit proti programu strany};  \dicIdiom{snúast}[í]{ \dicPhraseIS{snúast í e‑ð}} \dicFlx{refl} \dicDirectTranslationCS{změnit se v~(co), obrátit se v~(co)} \dicExampleIS{Ást hennar hefur snúist í hatur.} \dicExampleCS{Její láska se obrátila v~nenávist.};  \dicIdiom{snúast}[í kringum]{ \dicPhraseIS{snúast í kringum e‑n}} \dicFlx{refl} \dicDirectTranslationCS{točit\,/\addthin otáčet se kolem (koho) (číšník kolem hostů ap.)}; { \dicPhraseIS{e‑að snýst í kringum e‑n}} \dicFlx{refl} \dicDirectTranslationCS{(co) se točí kolem (koho) (život matky kolem dítěte ap.)};  \dicIdiom{snúast}[til]{ \dicPhraseIS{snúast til e‑s}} \dicFlx{refl} \dicDirectTranslationCS{uchýlit se k~(čemu), přistoupit k~(čemu)};  \dicIdiom{snúast}[um]{ \dicPhraseIS{snúast um e‑ð}} \dicFlx{refl} \dicDirectTranslationCS{být o~(čem), pojednávat o~(čem)};  \dicIdiom{snúast}[upp]{ \dicPhraseIS{snúast upp í e‑ð}} \dicFlx{refl} \dicDirectTranslationCS{změnit se v~(co), obrátit se v~(co)};  \dicIdiom{snúast}[við]{ \dicPhraseIS{snúast við e‑u}} \dicFlx{refl} \dicDirectTranslationCS{(za)reagovat na (co)} \dicExampleIS{snúast við erfiðleikunum} \dicExampleCS{reagovat na potíže}
\dicEntry[snúðugur] \dicTerm{snúð··ugur} \dicIPA{{s}{\textsubring{d}}{n}{u}{\textlengthmark}{ð}{\textscy}{\textbabygamma}{\textscy}{\textsubring{r}}} \dicPos{adj}[1]\dicFlx{}[-8] \dicSynonym{þóttafullur} \dicDirectTranslationCS{povýšený, namyšlený, nadutý}
\dicEntry[snúður] \dicTerm{snúð|ur} \dicIPA{{s}{\textsubring{d}}{n}{u}{\textlengthmark}{ð}{\textscy}{\textsubring{r}}} \dicPos{m}[6] \dicFlx{(‑s\,/\addthin ‑ar, ‑ar)}[66] \textbf{1.} \dicSynonym{snúningur} \dicDirectTranslationCS{zauzlení, uzel} \dicExampleIS{snúður á bandi} \dicExampleCS{zauzlení na provázku}  \textbf{2.} \dicFieldCat{kulin.} \dicSynonym{bolla} \dicDirectTranslationCS{(sladká) houska, brioška}  \textbf{3.} \dicSynonym{reigingur} \dicDirectTranslationCS{povýšenost, namyšlenost, nadutost};  \dicIdiom{snúður}{ \dicPhraseIS{fá\,/\addthin hafa e‑ð fyrir snúð sinn}} \dicLangCat{přen.} \dicDirectTranslationCS{být (čím) náležitě odměněn};  \dicPhraseIS{setja á sig snúð} \dicLangCat{přen.} \dicDirectTranslationCS{nafouknout se, urazit se}
\dicEntry[snúið] \dicTerm{snúið} \dicIPA{{s}{\textsubring{d}}{n}{u}{\textlengthmark}{\textsci}{\texttheta}} \dicPos{v} \dicFlx{supin} \dicLink{snúa}
\dicEntry[snúinn] \dicTerm{snúinn} \dicIPA{{s}{\textsubring{d}}{n}{u}{\textlengthmark}{\textsci}{\textsubring{n}}} \dicPos{adj}[6]\dicFlx{}[-2] \textbf{1.} \dicDirectTranslationCS{zatočený, zkroucený}  \textbf{2.} \dicSynonym{erfiður} \dicDirectTranslationCS{těžký, složitý}  \textbf{3.} \dicSynonym{óvingjarnlegur} \dicDirectTranslationCS{nepříjemný, protivný, ošklivý};  \dicPhraseIS{vera snúinn við e‑n} \dicDirectTranslationCS{být na (koho) protivný}
\dicEntry[snúningalipur] \dicTerm{snúninga··lipur} \dicIPA{{s}{\textsubring{d}}{n}{u}{\textlengthmark}{n}{i}{\ng}{\r{g}}{a}{l}{\textsci}{\textsubring{b}}{\textscy}{\textsubring{r}}} \dicPos{adj}[9] \dicFlx{(f ‑)}[1] \dicDirectTranslationCS{nápomocný, úslužný}
\dicEntry[snúningshlið] \dicTerm{snúnings··hlið} \dicIPA{{s}{\textsubring{d}}{n}{u}{\textlengthmark}{n}{i}{\ng}{s}{\textsubring{l}}{\textsci}{\texttheta}} \dicPos{n}[2] \dicFlx{(‑s, ‑)}[5] \dicDirectTranslationCS{turniket, otáčivé dveře}
\dicEntry[snúningshraði] \dicTerm{snúnings··hrað|i} \dicIPA{{s}{\textsubring{d}}{n}{u}{\textlengthmark}{n}{i}{\ng}{s}{\textsubring{r}}{a}{ð}{\textsci}} \dicPos{m}[1] \dicFlx{(‑a)}[3] \dicDirectTranslationCS{otáčky, rychlost otáčení}
\dicEntry[snúningur] \dicTerm{snún··ing|ur} \dicIPA{{s}{\textsubring{d}}{n}{u}{\textlengthmark}{n}{i}{\ng}{\r{g}}{\textscy}{\textsubring{r}}} \dicPos{m}[6] \dicFlx{(‑s, ‑ar)}[8] \textbf{1.} \dicSynonym*{það að snúa} \dicDirectTranslationCS{otáčení, točení, rotace} \dicExampleIS{snúningur hjóls} \dicExampleCS{točení kola}  \textbf{2.} \dicSynonym{dans} \dicDirectTranslationCS{tanec} \dicExampleIS{taka snúning á dansgólfinu} \dicExampleCS{zatancovat si na tanečním parketě}  \textbf{3.} \dicSynonym{viðvik} \dicDirectTranslationCS{drobná\,/\addthin pomocná práce} \dicIndirectTranslationCS{(zvláště v~pl)} \dicExampleIS{hafa dreng til snúninga} \dicExampleCS{mít chlapce na drobné práce};  \dicIdiom{snúningur}{ \dicPhraseIS{e‑að er á síðasta snúningi}} \dicLangCat{přen.} \dicDirectTranslationCS{(co) mele z~posledního (auto ap.)}; { \dicPhraseIS{standast e‑m ekki snúning}} \dicFlx{refl} \dicDirectTranslationCS{nemoct se srovnávat s~(kým)};  \dicPhraseIS{vera á síðasta snúningi með e‑u} \dicDirectTranslationCS{být v~posledním momentě (čeho), končit (co) (přípravy ap.)}
\dicEntry[snúra] \dicTerm{snúr|a} \dicIPA{{s}{\textsubring{d}}{n}{u}{\textlengthmark}{r}{a}} \dicPos{f}[1] \dicFlx{(‑u, ‑ur)}[7] \textbf{1.} \dicSynonym*{snúið band} \dicDirectTranslationCS{provaz, provázek}  \textbf{2.} \dicSynonym{þvottasnúra} \dicDirectTranslationCS{šňůra, špagát} \dicExampleIS{hengja þvottinn upp á snúru} \dicExampleCS{pověsit prádlo na šňůru}  \textbf{3.} \dicSynonym*{rafmagnssnúra} \dicDirectTranslationCS{šňůra, kabel}
\dicEntry[snússa] \dicTerm{snúss|a} \dicIPA{{s}{\textsubring{d}}{n}{u}{s}{\textlengthmark}{a}} \dicPos{v}[1] \dicFlx{(‑aði)}[44] \dicPhraseIS{snússa sig} \dicDirectTranslationCS{šňupnout si}
\dicEntry[snyfsi] \dicTerm{snyfsi} \dicIPA{{s}{\textsubring{d}}{n}{\textsci}{f}{s}{\textsci}} \dicPos{n}[2] \dicFlx{(‑s, ‑)}[14] \dicLink{snifsi}
\dicEntry[snyrta] \dicTerm{snyrt|a} \dicIPA{{s}{\textsubring{d}}{n}{\textsci}{\textsubring{r}}{\textsubring{d}}{a}} \dicPos{v}[2] \dicFlx{(‑i, ‑)}[16] \dicFlx{acc} \textbf{1.} \dicSynonym*{laga til} \dicDirectTranslationCS{poklidit, poklízet, uklidit, uklízet} \dicExampleIS{snyrta í herberginu} \dicExampleCS{poklidit v~pokoji}  \textbf{2.} \dicSynonym{fegra} \dicDirectTranslationCS{zkrášlit, zkrášlovat, (na)zdobit} \dicExampleIS{snyrta sig} \dicExampleCS{zkrášlit se}
\dicEntry[snyrtiborð] \dicTerm{snyrti··borð} \dicIPA{{s}{\textsubring{d}}{n}{\textsci}{\textsubring{r}}{\textsubring{d}}{\textsci}{\textsubring{b}}{\textopeno}{r}{\texttheta}} \dicPos{n}[2] \dicFlx{(‑s, ‑)}[5] \dicDirectTranslationCS{toaletní stolek}
\dicEntry[snyrtifræðingur] \dicTerm{snyrti·fræð··ing|ur} \dicIPA{{s}{\textsubring{d}}{n}{\textsci}{\textsubring{r}}{\textsubring{d}}{\textsci}{f}{r}{a}{i}{ð}{i}{\ng}{\r{g}}{\textscy}{\textsubring{r}}} \dicPos{m}[6] \dicFlx{(‑s, ‑ar)}[8] \dicDirectTranslationCS{kosmetolog, kosmetoložka}
\dicEntry[snyrtiherbergi] \dicTerm{snyrti··her·bergi} \dicIPA{{s}{\textsubring{d}}{n}{\textsci}{\textsubring{r}}{\textsubring{d}}{\textsci}{h}{\textepsilon}{r}{\textsubring{b}}{\textepsilon}{r}{\r{\textObardotlessj}}{\textsci}} \dicPos{n}[2] \dicFlx{(‑s, ‑)}[16] \dicSynonym{bað\smash{\textsuperscript{1}}} \dicDirectTranslationCS{toaleta, koupelna, umývárna}
\dicEntry[snyrtilegur] \dicTerm{snyrti··legur} \dicsymFrequent\  \dicIPA{{s}{\textsubring{d}}{n}{\textsci}{\textsubring{r}}{\textsubring{d}}{\textsci}{l}{\textepsilon}{\textbabygamma}{\textscy}{\textsubring{r}}} \dicPos{adj}[1]\dicFlx{}[-8] \dicDirectTranslationCS{upravený, dobře udržovaný} \dicExampleIS{snyrtilegt yfirskegg} \dicExampleCS{upravený knírek}
\dicEntry[snyrtimenni] \dicTerm{snyrti··menni} \dicIPA{{s}{\textsubring{d}}{n}{\textsci}{\textsubring{r}}{\textsubring{d}}{\textsci}{m}{\textepsilon}{n}{\textsci}} \dicPos{n}[2] \dicFlx{(‑s, ‑)}[14] \dicDirectTranslationCS{elegantní\,/\addthin uhlazený člověk, elegán(ka), fešák, fešanda}
\dicEntry[snyrtimennska] \dicTerm{snyrti··mennsk|a} \dicIPA{{s}{\textsubring{d}}{n}{\textsci}{\textsubring{r}}{\textsubring{d}}{\textsci}{m}{\textepsilon}{n}{s}{\r{g}}{a}} \dicPos{f}[1] \dicFlx{(‑u)}[5] \dicDirectTranslationCS{elegance, uhlazenost}
\dicEntry[snyrting] \dicTerm{snyrt··ing} \dicIPA{{s}{\textsubring{d}}{n}{\textsci}{\textsubring{r}}{\textsubring{d}}{i}{\ng}{\r{g}}} \dicPos{f}[4] \dicFlx{(‑ar, ‑ar)}[5] \textbf{1.} \dicSynonym{fegrunaraðgerð} \dicDirectTranslationCS{kosmetika} \dicIndirectTranslationCS{(odborné ošetření pleti)}  \textbf{2.} \dicSynonym*{snyrting fisks} \dicDirectTranslationCS{(finální) úprava\,/\addthin zpracování}  \textbf{3.} \dicSynonym{salerni} \dicDirectTranslationCS{toaleta, umývárna}
\dicEntry[snyrtistofa] \dicTerm{snyrti··stof|a} \dicIPA{{s}{\textsubring{d}}{n}{\textsci}{\textsubring{r}}{\textsubring{d}}{\textsci}{s}{\textsubring{d}}{\textopeno}{v}{a}} \dicPos{f}[1] \dicFlx{(‑u, ‑ur)}[7] \dicDirectTranslationCS{kosmetický salón}
\dicEntry[snyrtivörur] \dicTerm{snyrti··vörur} \dicIPA{{s}{\textsubring{d}}{n}{\textsci}{\textsubring{r}}{\textsubring{d}}{\textsci}{v}{\oe}{r}{\textscy}{\textsubring{r}}} \dicPos{f}[1] \dicFlx{pl}[9] \dicDirectTranslationCS{kosmetika} \dicIndirectTranslationCS{(prostředky používané k~péči o~tělo)}
\dicEntry[sný] \dicTerm{sný} \dicIPA{{s}{\textsubring{d}}{n}{i}{\textlengthmark}} \dicPos{v} \dicFlx{ind praes sg 1 pers} \dicLink{snúa}
\dicEntry[snýta] \dicTerm{snýt|a} \dicIPA{{s}{\textsubring{d}}{n}{i}{\textlengthmark}{\textsubring{d}}{a}} \dicPos{v}[2] \dicFlx{(‑ti, ‑t)}[62] \dicFlx{dat} \dicDirectTranslationCS{(vy)smrkat};  \dicPhraseIS{snýta sér} \dicDirectTranslationCS{vysmrkat se} \dicExampleIS{snýta sér í vasaklútinn} \dicExampleCS{vysmrkat se do kapesníku}
\dicEntry[snýtibréf] \dicTerm{snýti··bréf}\dicTerm{, snýtubréf} \dicIPA{{s}\-{\textsubring{d}}\-{n}\-{i}\-{\textlengthmark}\-{\textsubring{d}}\-{\textsci}\-{\textsubring{b}}\-{r}\-{j}\-{\textepsilon}\-{f}\-} \dicPos{n}[2] \dicFlx{(‑s, ‑)}[5] \dicDirectTranslationCS{papírový kapesník}
\dicEntry[snýtubréf] \dicTerm{snýtu··bréf} \dicIPA{{s}{\textsubring{d}}{n}{i}{\textlengthmark}{\textsubring{d}}{\textscy}{\textsubring{b}}{r}{j}{\textepsilon}{f}} \dicPos{n}[2] \dicFlx{(‑s, ‑)}[5] \dicLink{snýtibréf}
\dicEntry[snýtur] \dicTerm{snýtur} \dicIPA{{s}{\textsubring{d}}{n}{i}{\textlengthmark}{\textsubring{d}}{\textscy}{\textsubring{r}}} \dicPos{f}[1] \dicFlx{pl}[12] \dicDirectTranslationCS{vysmrkání, vyčištění nosu}
\dicEntry[snæða] \dicTerm{snæ|ða} \dicIPA{{s}{\textsubring{d}}{n}{a}{i}{\textlengthmark}{ð}{a}} \dicPos{v}[2] \dicFlx{(‑ddi, ‑tt)}[174] \dicFlx{acc} \dicLangCat{form.} \dicSynonym{borða} \dicDirectTranslationCS{jíst} \dicExampleIS{snæða fisk} \dicExampleCS{jíst rybu};  \dicPhraseIS{snæða kvöldverð} \dicLangCat{form.} \dicDirectTranslationCS{večeřet}
\dicEntry[snæðingur] \dicTerm{snæð··ing|ur} \dicIPA{{s}{\textsubring{d}}{n}{a}{i}{\textlengthmark}{ð}{i}{\ng}{\r{g}}{\textscy}{\textsubring{r}}} \dicPos{m}[6] \dicFlx{(‑s)}[9] \dicDirectTranslationCS{jídlo, pokrm};  \dicPhraseIS{setjast\,/\addthin sitja að snæðingi} \dicFlx{refl} \dicDirectTranslationCS{jíst, stolovat, zasednout k~jídlu}
\dicEntry[snælda] \dicTerm{snæld|a} \dicIPA{{s}{\textsubring{d}}{n}{a}{i}{l}{\textsubring{d}}{a}} \dicPos{f}[1] \dicFlx{(‑u, ‑ur)}[19] \textbf{1.} \dicSynonym*{spunasnælda} \dicDirectTranslationCS{vřeteno}  \textbf{2.} \dicSynonym{kassetta} \dicDirectTranslationCS{kazeta} \dicIndirectTranslationCS{(výměnná jednotka s~magnetickou páskou)}
\dicEntry[snælína] \dicTerm{snæ··lín|a} \dicIPA{{s}{\textsubring{d}}{n}{a}{i}{\textlengthmark}{l}{i}{n}{a}} \dicPos{f}[1] \dicFlx{(‑u, ‑ur)}[7] \dicFieldCat{meteo.} \dicSynonym{snjólína} \dicDirectTranslationCS{sněž\-ná čára}
\dicEntry[snæljós] \dicTerm{snæ··ljós} \dicIPA{{s}{\textsubring{d}}{n}{a}{i}{\textlengthmark}{l}{j}{ou}{s}} \dicPos{n}[2] \dicFlx{(‑s, ‑)}[5] \dicSynonym*{ljósglampi} \dicDirectTranslationCS{blýskavice, záblesk};  \dicPhraseIS{vera eins og snæljós} \dicLangCat{přen.} \dicDirectTranslationCS{být rychlý jako blesk}
\dicEntry[snær] \dicTerm{snæ|r} \dicIPA{{s}{\textsubring{d}}{n}{a}{i}{\textlengthmark}{\textsubring{r}}} \dicPos{m}[8] \dicFlx{(‑s\,/\addthin ‑var)}[5] \dicSynonym{snjór} \dicDirectTranslationCS{sníh}
\dicEntry[snæri] \dicTerm{snæri} \dicIPA{{s}{\textsubring{d}}{n}{a}{i}{\textlengthmark}{r}{\textsci}} \dicPos{n}[2] \dicFlx{(‑s, ‑)}[14] \dicSynonym{band} \dicDirectTranslationCS{šňůra, provázek, motouz};  \dicPhraseIS{vera á snærum e‑s} \dicDirectTranslationCS{reprezentovat (co), být představitelem\,/\addthin reprezentantem (čeho)};  \dicPhraseIS{það hleypur á snærið hjá e‑m} \dicFlx{impers} \dicLangCat{přen.} \dicDirectTranslationCS{(kdo) má z~pekla štěstí}
\dicEntry[snæugla] \dicTerm{snæ··ugl|a} \dicIPA{{s}{\textsubring{d}}{n}{a}{i}{\textlengthmark}{\textscy}{\r{g}}{l}{a}} \dicPos{f}[1] \dicFlx{(‑u, ‑ur)}[19] \dicFieldCat{zool.} \dicDirectTranslationCS{sovice sněžní} \textit{(l.~{\textLA{Bubo scandiacus}})}  \dicsymPhoto\ 
\dicFigure{13499.jpg}{Snæugla}{Snæugla - Antušek Ivo, Biolib, Copyright/CC-BY-NC}
\dicEntry[snöggklipptur] \dicTerm{snögg··klipptur} \dicIPA{{s}{\textsubring{d}}{n}{\oe}{\r{g}}{k\smash{\textsuperscript{h}}}{l}{\textsci}{f}{\textsubring{d}}{\textscy}{\textsubring{r}}} \dicPos{adj}[1]\dicFlx{}[-10] \dicDirectTranslationCS{ostříhaný nakrátko}
\dicEntry[snöggklæddur] \dicTerm{snögg··klæddur} \dicIPA{{s}{\textsubring{d}}{n}{\oe}{\r{g}}{k\smash{\textsuperscript{h}}}{l}{a}{i}{\textsubring{d}}{\textscy}{\textsubring{r}}} \dicPos{adj}[2]\dicFlx{}[-18] \dicDirectTranslationCS{oblečený nalehko}
\dicEntry[snögglega] \dicTerm{snögg··lega} \dicsymFrequent\  \dicIPA{{s}{\textsubring{d}}{n}{\oe}{\r{g}}{l}{\textepsilon}{\textbabygamma}{a}} \dicPos{adv} \dicSynonym{skyndilega} \dicDirectTranslationCS{nenadále, nečekaně} \dicExampleIS{Fundinum lauk snögglega.} \dicExampleCS{Schůze se nečekaně skončila.}
\dicEntry[snöggsoðinn] \dicTerm{snögg··soðinn} \dicIPA{{s}{\textsubring{d}}{n}{\oe}{\r{g}}{s}{\textopeno}{ð}{\textsci}{\textsubring{n}}} \dicPos{adj}[6]\dicFlx{}[-2] \textbf{1.} \dicDirectTranslationCS{krátce vařený}  \textbf{2.} \dicSynonym*{illa undirbúinn} \dicDirectTranslationCS{krátce\,/\addthin nedostatečně připravený (projev ap.)}
\dicEntry[snöggt] \dicTerm{snöggt} \dicsymFrequent\  \dicIPA{{s}{\textsubring{d}}{n}{\oe}{x}{\textsubring{d}}} \dicPos{adv} \dicFlx{(comp sneggra, sup sneggst)} \dicSynonym{stuttlega} \dicDirectTranslationCS{krátce, (na)krátko} \dicExampleIS{klippa hárið snöggt} \dicExampleCS{ostříhat vlasy nakrátko}
\dicEntry[snöggur] \dicTerm{snöggur} \dicsymFrequent\  \dicIPA{{s}{\textsubring{d}}{n}{\oe}{\r{g}}{\textlengthmark}{\textscy}{\textsubring{r}}} \dicPos{adj}[1]\dicFlx{}[-8] \textbf{1.} \dicSynonym{skjótur} \dicDirectTranslationCS{rychlý, pohotový}  \textbf{2.} \dicSynonym{skyndilegur} \dicDirectTranslationCS{náhlý, nenadálý} \dicExampleIS{snögg ferð} \dicExampleCS{náhlá cesta}  \textbf{3.} \dicSynonym{stutthærður} \dicDirectTranslationCS{nakrátko ostříhaný};  \dicIdiom{snöggur}{ \dicPhraseIS{hitta snöggan blett á e‑m}} \dicLangCat{přen.} \dicDirectTranslationCS{zasáhnout (čí) slabé místo}
\dicEntry[snöggvast] \dicTerm{snöggvast} \dicIPA{{s}{\textsubring{d}}{n}{\oe}{\r{g}}{v}{a}{s}{\textsubring{d}}} \dicPos{adv} \dicFlx{sup} \textbf{1.} \dicSynonym*{eitt andartak} \dicDirectTranslationCS{na okamžik, na chvilku};  \dicPhraseIS{sem snöggvast} \dicFlx{adv} \dicDirectTranslationCS{na sekundu, na moment}  \textbf{2.} \dicSynonym{núna} \dicDirectTranslationCS{teď, nyní}
\dicEntry[snökt] \dicTerm{snökt} \dicIPA{{s}{\textsubring{d}}{n}{\oe}{x}{\textsubring{d}}} \dicPos{n}[2] \dicFlx{(‑s)}[2] \dicSynonym{grátur\smash{\textsuperscript{1}}} \dicDirectTranslationCS{štkaní, lkaní, vzlykání}
\dicEntry[snökta] \dicTerm{snökt|a} \dicIPA{{s}{\textsubring{d}}{n}{\oe}{x}{\textsubring{d}}{a}} \dicPos{v}[2] \dicFlx{(‑i, ‑)}[17] \dicDirectTranslationCS{štkát, lkát, vzlykat}
\dicEntry[snörl] \dicTerm{snörl} \dicIPA{{s}{\textsubring{d}}{n}{\oe}{r}{\textsubring{d}}{\textsubring{l}}} \dicPos{n}[2] \dicFlx{(‑s)}[2] \dicSynonym{korr} \dicDirectTranslationCS{chrchlání, chraptění}
\dicEntry[snörla] \dicTerm{snörl|a} \dicIPA{{s}{\textsubring{d}}{n}{\oe}{r}{\textsubring{d}}{l}{a}} \dicPos{v}[1] \dicFlx{(‑aði)}[48] \dicDirectTranslationCS{chrchlat, chraptět};  \dicIdiom{snörla}[í]{ \dicPhraseIS{það snörlar í e‑m}} \dicFlx{impers} \dicSynonym*{það korrar} \dicDirectTranslationCS{(kdo) chrchlá\,/\addthin chraptí}
\dicEntry[snöru] \dicTerm{snöru} \dicIPA{{s}{\textsubring{d}}{n}{\oe}{\textlengthmark}{r}{\textscy}} \dicPos{f} \dicFlx{sg gen} \dicLink{snara\smash{\textsuperscript{1}}}
\dicEntry[snös] \textls[15]{\dicTerm{snös} \dicIPA{{s}{\textsubring{d}}{n}{\oe}{\textlengthmark}{s}} \dicPos{f}[7] \dicFlx{(snasar, snasir)}[16] \dicDirectTranslationCS{okraj skály\,/\addthin útesu}} \dicExampleIS{standa frammi á snösinni} \dicExampleCS{stát na okraji skály}
\dicEntry[so.] \dicTerm{so.} \dicPos{zkr} \dicPhraseIS{sagnorð} \dicFieldCat{jaz.} \dicDirectTranslationCS{sloveso}
\dicEntry[soddan] \dicTerm{soddan}\dicTerm{, svoddan} \dicIPA{{s}\-{\textopeno}\-{\textsubring{d}}\-{\textlengthmark}\-{a}\-{\textsubring{n}}\-} \dicPos{adj}[13] \dicFlx{indecl}[3] \dicSynonym{slíkur} \dicDirectTranslationCS{takový}
\dicEntry[soð] \dicTerm{soð} \dicIPA{{s}{\textopeno}{\textlengthmark}{\texttheta}} \dicPos{n}[2] \dicFlx{(‑s)}[2] \dicDirectTranslationCS{vývar} \dicExampleIS{kjötsoð} \dicExampleCS{masový vývar};  \dicPhraseIS{fá sér í soðið} \dicDirectTranslationCS{koupit si rybu k~jídlu}
\dicEntry[soðið] \dicTerm{soðið} \dicIPA{{s}{\textopeno}{\textlengthmark}{ð}{\textsci}{\texttheta}} \dicPos{v} \dicFlx{supin} \dicLink{sjóða}
\dicEntry[soðinn] \dicTerm{soðinn} \dicIPA{{s}{\textopeno}{\textlengthmark}{ð}{\textsci}{\textsubring{n}}} \dicPos{adj}[6]\dicFlx{}[-6] \dicDirectTranslationCS{(u)vařený} \dicExampleIS{soðið egg} \dicExampleCS{vařené vajíčko}
\dicEntry[soðna] \dicTerm{soðn|a} \dicIPA{{s}{\textopeno}{ð}{n}{a}} \dicPos{v}[1] \dicFlx{(‑aði)}[54] \dicDirectTranslationCS{uvařit se, být uvařen} \dicExampleIS{Hrísgrjónin soðna á 20 mínútum.} \dicExampleCS{Rýže se uvaří za 20 minut.}
\dicEntry[soðning] \dicTerm{soð··ning} \dicIPA{{s}{\textopeno}{ð}{n}{i}{\ng}{\r{g}}} \dicPos{f}[4] \dicFlx{(‑ar, ‑ar)}[5] \dicSynonym*{soðinn fiskur} \dicDirectTranslationCS{vařená ryba}
\dicEntry[sofa] \dicTerm{sofa} \dicsymFrequent\  \dicIPA{{s}{\textopeno}{\textlengthmark}{v}{a}} \dicPos{v}[6] \dicFlx{(sef, svaf, sváfum, svæfi, sofið)}[3] \dicDirectTranslationCS{spát, vyspat se} \dicExampleIS{sofa undir berum himni} \dicExampleCS{spát pod širákem};  \dicPhraseIS{fara að sofa} \dicDirectTranslationCS{jít spát} \dicExampleIS{fara seint að sofa} \dicExampleCS{jít pozdě spát};  \dicPhraseIS{sofa laust\,/\addthin fast} \dicDirectTranslationCS{spát lehce\,/\addthin tvrdě};  \dicIdiom{sofa}[hjá]{ \dicPhraseIS{sofa hjá e‑m}} \dicDirectTranslationCS{spát s~(kým), milovat se s~(kým)};  \dicIdiom{sofa}[saman]{ \dicPhraseIS{sofa saman}} \dicDirectTranslationCS{spát spolu, milovat se};  \dicIdiom{sofa}[úr]{ \dicPhraseIS{sofa úr sér e‑u}} \dicDirectTranslationCS{vyspat se z~(čeho)};  \dicIdiom{sofa}[út]{ \dicPhraseIS{sofa út}} \dicDirectTranslationCS{vyspávat, prospat se};  \dicIdiom{sofa}[yfir]{ \dicPhraseIS{sofa yfir sig}} \dicDirectTranslationCS{zaspat}
\dicEntry[sofandaháttur] \dicTerm{sofanda··hátt|ur} \dicIPA{{s}{\textopeno}{\textlengthmark}{v}{a}{n}{\textsubring{d}}{a}{h}{au}{h}{\textsubring{d}}{\textscy}{\textsubring{r}}} \dicPos{m}[12] \dicFlx{(‑ar)}[8] \dicDirectTranslationCS{otupělost, netečnost}
\dicEntry[sofandi] \dicTerm{sof··andi} \dicsymFrequent\  \dicIPA{{s}{\textopeno}{\textlengthmark}{v}{a}{n}{\textsubring{d}}{\textsci}} \dicPos{adj}[13] \dicFlx{indecl}[1] \textbf{1.} \dicSynonym*{sem sefur} \dicDirectTranslationCS{spící, podřimující} \dicExampleIS{Ég var ekki sofandi.} \dicExampleCS{Nespal jsem.}  \textbf{2.} \dicSynonym*{athyglislaus} \dicDirectTranslationCS{nepozorný, nesoustředěný}
\dicEntry[sofið] \dicTerm{sofið} \dicIPA{{s}{\textopeno}{\textlengthmark}{v}{\textsci}{\texttheta}} \dicPos{v} \dicFlx{supin} \dicLink{sofa}
\dicEntry[sofinn] \dicTerm{sofinn} \dicIPA{{s}{\textopeno}{\textlengthmark}{v}{\textsci}{\textsubring{n}}} \dicPos{adj}[6]\dicFlx{}[-6] \dicDirectTranslationCS{vyspalý, prospaný};  \dicPhraseIS{vera vel\,/\addthin illa sofinn} \dicDirectTranslationCS{být dobře\,/\addthin špatně vyspalý};  \dicPhraseIS{vera vakinn og sofinn að (gera e‑ð)} \dicLangCat{přen.} \dicDirectTranslationCS{udělat první poslední pro (udělání (čeho))}
\dicEntry[Sofía] \dicTerm{Sofí|a} \dicIPA{{s}{\textopeno}{\textlengthmark}{f}{i}{j}{a}} \dicPos{f}[1] \dicFlx{(‑u)}[6] \dicLink{Sófía}
\dicEntry[sofna] \dicTerm{sofn|a} \dicsymFrequent\  \dicIPA{{s}{\textopeno}{\textsubring{b}}{n}{a}} \dicPos{v}[1] \dicFlx{(‑aði)}[1] \textbf{1.} \dicSynonym*{falla í svefn} \dicDirectTranslationCS{usnout, usínat} \dicExampleIS{sofna fast} \dicExampleCS{usnout tvrdě}  \textbf{2.} \dicSynonym{deyja} \dicDirectTranslationCS{usnout, odejít} \dicIndirectTranslationCS{(zemřít)}
\dicEntry[sofnaður] \dicTerm{sofn··|aður} \dicIPA{{s}{\textopeno}{\textsubring{b}}{n}{a}{ð}{\textscy}{\textsubring{r}}} \dicPos{adj}[3] \dicFlx{(f ‑uð)}[4] \dicDirectTranslationCS{usnuvší, spící}
\dicEntry[sog] \dicTerm{sog} \dicIPA{{s}{\textopeno}{\textlengthmark}{x}} \dicPos{n}[2] \dicFlx{(‑s, ‑)}[5] \textbf{1.} \dicSynonym*{það að sogast} \dicDirectTranslationCS{(na)sání, nasávání}  \textbf{2.} \dicSynonym{korr} \dicDirectTranslationCS{chrčení, chrchlání} \dicExampleIS{sog í lungum} \dicExampleCS{chrčení na plicích}  \textbf{3.} \dicSynonym*{öldusog} \dicDirectTranslationCS{příboj}
\dicEntry[soga] \dicTerm{sog|a} \dicIPA{{s}{\textopeno}{\textlengthmark}{\textbabygamma}{a}} \dicPos{v}[1] \dicFlx{(‑aði)}[1] \dicFlx{acc} \dicDirectTranslationCS{sát, nasávat, odsávat} \dicExampleIS{soga reykinn úr pípunni} \dicExampleCS{sát kouř z~dýmky};  \dicPhraseIS{soga í sig e‑ð} \dicDirectTranslationCS{vsát (co), vdechovat (co), inhalovat (co)};  \dicIdiom{sogast}{ \dicPhraseIS{sogast}} \dicFlx{refl} {\textbf{a.}} \dicDirectTranslationCS{být vysán\,/\addthin odsán};  {\textbf{b.}} \dicLangCat{přen.} \dicDirectTranslationCS{být vtažen\,/\addthin zatažen (do příběhu ap.)}
\dicEntry[sogdæla] \dicTerm{sog··dæl|a} \dicIPA{{s}{\textopeno}{\textbabygamma}{\textsubring{d}}{a}{i}{l}{a}} \dicPos{f}[1] \dicFlx{(‑u, ‑ur)}[19] \dicFieldCat{fyz.} \dicDirectTranslationCS{vývěva, vakuová pumpa}
\dicEntry[sogið] \dicTerm{sogið} \dicIPA{{s}{\textopeno i}{j}{\textlengthmark}{\textsci}{\texttheta}} \dicPos{v} \dicFlx{supin} \dicLink{sjúga}
\dicEntry[sogpípa] \dicTerm{sog··píp|a} \dicIPA{{s}{\textopeno}{x}{p\smash{\textsuperscript{h}}}{i}{\textsubring{b}}{a}} \dicPos{f}[1] \dicFlx{(‑u, ‑ur)}[19] \dicDirectTranslationCS{násoska, pipeta}
\dicEntry[sogrör] \dicTerm{sog··rör} \dicIPA{{s}{\textopeno}{\textbabygamma}{r}{\oe}{\textsubring{r}}} \dicPos{n}[2] \dicFlx{(‑s, ‑)}[5] \dicDirectTranslationCS{brčko} \dicIndirectTranslationCS{(dlouhá tenká trubička k~nasávání nápoje)}
\dicEntry[sogskál] \dicTerm{sog··skál} \dicIPA{{s}{\textopeno}{x}{s}{\r{g}}{au}{\textsubring{l}}} \dicPos{f}[4] \dicFlx{(‑ar, ‑ar\,/\addthin ‑ir)}[30] \textbf{1.} \dicDirectTranslationCS{přísavka} \dicIndirectTranslationCS{(pomůcka miskovitého tvaru, kterou lze připevnit na povrch přisátím)}  \textbf{2.} \dicFieldCat{biol.} \dicDirectTranslationCS{přísavka} \dicIndirectTranslationCS{(přísavný orgán na povrchu některých živočichů)}
\dicEntry[sogæð] \dicTerm{sog··æð} \dicIPA{{s}{\textopeno}{\textlengthmark}{\textbabygamma}{a}{i}{\texttheta}} \dicPos{f}[4] \dicFlx{(‑ar, ‑ar)}[1] \dicFieldCat{anat.} \dicDirectTranslationCS{lymfatická\,/\addthin mízní céva}
\dicEntry[sogæðavökvi] \dicTerm{sog·æða··vökv|i} \dicIPA{{s}{\textopeno}{\textlengthmark}{\textbabygamma}{a}{i}{ð}{a}{v}{\oe}{\r{g}}{v}{\textsci}} \dicPos{m}[1] \dicFlx{(‑a)}[3] \dicFieldCat{anat.} \dicDirectTranslationCS{lymfa, míza}
\dicEntry[sojaolía] \dicTerm{soja··olí|a} \dicIPA{{s}{\textopeno}{\textlengthmark}{j}{a}{\textopeno}{l}{i}{j}{a}} \dicPos{f}[1] \dicFlx{(‑u)}[5] \dicDirectTranslationCS{sojový olej}
\dicEntry[sojasósa] \dicTerm{soja··sós|a} \dicIPA{{s}{\textopeno}{\textlengthmark}{j}{a}{s}{ou}{s}{a}} \dicPos{f}[1] \dicFlx{(‑u, ‑ur)}[7] \dicDirectTranslationCS{sojová omáčka}
\dicEntry[sokkaband] \dicTerm{sokka··|band} \dicIPA{{s}{\textopeno}{h}{\r{g}}{a}{\textsubring{b}}{a}{n}{\textsubring{d}}} \dicPos{n}[2] \dicFlx{(‑bands, ‑bönd)}[8] \dicDirectTranslationCS{podvazek}
\dicEntry[sokkabandabelti] \dicTerm{sokka·banda··belti} \dicIPA{{s}{\textopeno}{h}{\r{g}}{a}{\textsubring{b}}{a}{n}{\textsubring{d}}{a}{\textsubring{b}}{\textepsilon}{\textsubring{l}}{\textsubring{d}}{\textsci}} \dicPos{n}[2] \dicFlx{(‑s, ‑)}[14] \dicDirectTranslationCS{podvazkový pás}
\dicEntry[sokkabandsár] \dicTerm{sokka·bands··ár} \dicIPA{{s}{\textopeno}{h}{\r{g}}{a}{\textsubring{b}}{a}{n}{\textsubring{d}}{s}{au}{\textsubring{r}}} \dicPos{n}[2] \dicFlx{(‑s, ‑)}[5] \dicDirectTranslationCS{mladá léta, mládí} \dicExampleIS{á sínum sokkabandsárum} \dicExampleCS{v~mládí}
\dicEntry[sokkabuxur] \dicTerm{sokka··buxur} \dicIPA{{s}{\textopeno}{h}{\r{g}}{a}{\textsubring{b}}{\textscy}{x}{s}{\textscy}{\textsubring{r}}} \dicPos{f}[12] \dicFlx{pl}[7] \dicDirectTranslationCS{punčochy, punčocháče}
\dicEntry[sokkaleistur] \dicTerm{sokka··leist|ur} \dicIPA{{s}{\textopeno}{h}{\r{g}}{a}{l}{ei}{s}{\textsubring{d}}{\textscy}{\textsubring{r}}} \dicPos{m}[6] \dicFlx{(‑s, ‑ar)}[4] \dicDirectTranslationCS{podkolenka, punčocha};  \dicPhraseIS{á sokkaleistunum} \dicFlx{adv} \dicDirectTranslationCS{jen v~ponožkách\,/\addthin punčochách, bez bot}
\dicEntry[sokkið] \dicTerm{sokkið} \dicIPA{{s}{\textopeno}{h}{\r{\textObardotlessj}}{\textsci}{\texttheta}} \dicPos{v} \dicFlx{supin} \dicLink{sökkva\smash{\textsuperscript{2}}}
\dicEntry[sokkinn] \dicTerm{sokkinn} \dicIPA{{s}{\textopeno}{h}{\r{\textObardotlessj}}{\textsci}{\textsubring{n}}} \dicPos{adj}[6]\dicFlx{}[-6] \dicDirectTranslationCS{potopený, ponořený} \dicExampleIS{sokkið skip} \dicExampleCS{potopená loď};  \dicPhraseIS{vera djúpt sokkinn í e‑u} \dicLangCat{přen.} \dicDirectTranslationCS{být hluboce ponořený v~(čem) (práci ap.)}
\dicEntry[sokkur] \dicTerm{sokk|ur} \dicsymFrequent\  \dicIPA{{s}{\textopeno}{h}{\r{g}}{\textscy}{\textsubring{r}}} \dicPos{m}[6] \dicFlx{(‑s, ‑ar)}[15] \dicDirectTranslationCS{ponožka} \dicExampleIS{einir sokkar} \dicExampleCS{jeden pár ponožek}
\dicEntry[soldán] \dicTerm{soldán} \dicIPA{{s}{\textopeno}{l}{\textsubring{d}}{au}{\textsubring{n}}} \dicPos{m}[4] \dicFlx{(‑s, ‑ar)}[10] \dicDirectTranslationCS{sultán}
\dicEntry[soldið] \dicTerm{soldið} \dicIPA{{s}{\textopeno}{l}{\textsubring{d}}{\textsci}{\texttheta}} \dicPos{adv} \dicSynonym{svolítið} \dicDirectTranslationCS{trochu, trošku}
\dicEntry[soldill] \dicTerm{soldill} \dicIPA{{s}{\textopeno}{l}{\textsubring{d}}{\textsci}{\textsubring{d}}{\textsubring{l}}} \dicPos{adj}[8]\dicFlx{}[-3] \dicSynonym{svolítill} \dicDirectTranslationCS{malý, maličký}
\dicEntry[sollið] \dicTerm{sollið} \dicIPA{{s}{\textopeno}{\textsubring{d}}{l}{\textsci}{\texttheta}} \dicPos{v} \dicFlx{supin} \dicLink{svella\smash{\textsuperscript{2}}}
\dicEntry[sollur] \dicTerm{soll|ur} \dicIPA{{s}{\textopeno}{\textsubring{d}}{l}{\textscy}{\textsubring{r}}} \dicPos{m}[6] \dicFlx{(‑s)}[7] \dicSynonym{svall\smash{\textsuperscript{1}}} \dicDirectTranslationCS{flám, flámování, hýření}
\dicEntry[soltið] \dicTerm{soltið} \dicIPA{{s}{\textopeno}{\textsubring{l}}{\textsubring{d}}{\textsci}{\texttheta}} \dicPos{v} \dicFlx{supin} \dicLink{svelta\smash{\textsuperscript{2}}}
\dicEntry[soltinn] \dicTerm{soltinn} \dicIPA{{s}{\textopeno}{\textsubring{l}}{\textsubring{d}}{\textsci}{\textsubring{n}}} \dicPos{adj}[6]\dicFlx{}[-2] \dicSynonym{svangur} \dicDirectTranslationCS{hladový, vyhladovělý}
\dicEntry[sonar] \dicTerm{sonar} \dicIPA{{s}{\textopeno}{\textlengthmark}{n}{a}{\textsubring{r}}} \dicPos{m} \dicFlx{sg gen} \dicLink{sonur}
\dicEntry[sonardóttir] \dicTerm{sonar··|dóttir} \dicIPA{{s}{\textopeno}{\textlengthmark}{n}{a}{r}{\textsubring{d}}{ou}{h}{\textsubring{d}}{\textsci}{\textsubring{r}}} \dicPos{f}[11] \dicFlx{(‑dóttur, ‑dætur)}[2] \dicDirectTranslationCS{vnučka} \dicIndirectTranslationCS{(dcera syna)} \dicAntonym{dótturdóttir}
\dicEntry[sonarsonur] \dicTerm{sonar··|sonur} \dicIPA{{s}{\textopeno}{\textlengthmark}{n}{a}{\textsubring{r}}{s}{\textopeno}{n}{\textscy}{\textsubring{r}}} \dicPos{m}[12] \dicFlx{(‑sonar, ‑synir)}[9] \dicDirectTranslationCS{vnuk} \dicIndirectTranslationCS{(syn syna)} \dicAntonym{dóttursonur}
\dicEntry[sonnetta] \dicTerm{sonnett|a} \dicIPA{{s}{\textopeno}{n}{\textlengthmark}{\textepsilon}{h}{\textsubring{d}}{a}} \dicPos{f}[1] \dicFlx{(‑u, ‑ur)}[7] \dicFieldCat{lit.} \dicDirectTranslationCS{sonet}
\dicEntry[sonur] \dicTerm{sonur} \dicsymFrequent\  \dicIPA{{s}{\textopeno}{\textlengthmark}{n}{\textscy}{\textsubring{r}}} \dicPos{m}[12] \dicFlx{(sonar, synir)}[9] \dicDirectTranslationCS{syn} \dicExampleIS{Hann átti tvo syni.} \dicExampleCS{Měl dva syny.}
\dicEntry[sopi] \dicTerm{sop|i} \dicIPA{{s}{\textopeno}{\textlengthmark}{\textsubring{b}}{\textsci}} \dicPos{m}[1] \dicFlx{(‑a, ‑ar)}[1] \dicDirectTranslationCS{doušek, lok, hlt} \dicExampleIS{fá sér sopa úr glasinu} \dicExampleCS{loknout si ze skleničky};  \dicPhraseIS{e‑m þykir sopinn góður} \dicFlx{impers} \dicDirectTranslationCS{(kdo) se rád napije (alkoholu)}
\dicEntry[sopið] \dicTerm{sopið} \dicIPA{{s}{\textopeno}{\textlengthmark}{\textsubring{b}}{\textsci}{\texttheta}} \dicPos{v} \dicFlx{supin} \dicLink{súpa\smash{\textsuperscript{2}}}
\dicEntry[sorfið] \dicTerm{sorfið} \dicIPA{{s}{\textopeno}{r}{v}{\textsci}{\texttheta}} \dicPos{v} \dicFlx{supin} \dicLink{sverfa}
\dicEntry[sorg] \dicTerm{sorg} \dicsymFrequent\  \dicIPA{{s}{\textopeno}{r}{\r{g}}} \dicPos{f}[7] \dicFlx{(‑ar, ‑ir)}[1] \dicSynonym{hryggð} \dicDirectTranslationCS{smutek, zármutek, žal} \dicExampleIS{djúp sorg} \dicExampleCS{hluboký smutek}
\dicEntry[sorgarbúningur] \dicTerm{sorgar··bú·ning|ur} \dicIPA{{s}{\textopeno}{r}{\r{g}}{a}{r}{\textsubring{b}}{u}{n}{i}{\ng}{\r{g}}{\textscy}{\textsubring{r}}} \dicPos{m}[6] \dicFlx{(‑s, ‑ar)}[8] \dicDirectTranslationCS{smuteční šaty\,/\addthin oblek}
\dicEntry[sorgarfrétt] \dicTerm{sorgar··frétt} \dicIPA{{s}{\textopeno}{r}{\r{g}}{a}{\textsubring{r}}{f}{r}{j}{\textepsilon}{h}{\textsubring{d}}} \dicPos{f}[7] \dicFlx{(‑ar, ‑ir)}[1] \dicDirectTranslationCS{smutná zpráva}
\dicEntry[sorgarlag] \dicTerm{sorgar··|lag} \dicIPA{{s}{\textopeno}{r}{\r{g}}{a}{r}{l}{a}{x}} \dicPos{n}[2] \dicFlx{(‑lags, ‑lög)}[8] \dicDirectTranslationCS{smuteční píseň}
\dicEntry[sorgarsaga] \dicTerm{sorgar··|saga} \dicIPA{{s}{\textopeno}{r}{\r{g}}{a}{\textsubring{r}}{s}{a}{\textbabygamma}{a}} \dicPos{f}[1] \dicFlx{(‑sögu, ‑sögur)}[14] \dicDirectTranslationCS{smutný\,/\addthin truchlivý příběh}
\dicEntry[sorgbitinn] \dicTerm{sorg··bitinn} \dicIPA{{s}{\textopeno}{r}{\r{g}}{\textsubring{b}}{\textsci}{\textsubring{d}}{\textsci}{\textsubring{n}}} \dicPos{adj}[6]\dicFlx{}[-2] \dicDirectTranslationCS{smutný, zarmoucený, sklíčený, zkroušený}
\dicEntry[sorglegur] \dicTerm{sorg··legur} \dicIPA{{s}{\textopeno}{r}{\r{g}}{l}{\textepsilon}{\textbabygamma}{\textscy}{\textsubring{r}}} \dicPos{adj}[1]\dicFlx{}[-8] \dicDirectTranslationCS{smutný, žalostný, žalostivý, truchlivý}
\begin{xtolerant}{}{1pt}
\dicEntry[sorgmæddur] \dicTerm{sorg··mæddur}\addthinS\dicIPA{{s}{\textopeno}{r}{\r{g}}{m}{a}{i}{\textsubring{d}}{\textscy}{\textsubring{r}}}\addthinS\dicPos{adj}[2]\dicFlx{}[-18]\addthinS\dicSynonym{hryggur\smash{\textsuperscript{2}}} \dicDirectTranslationCS{smutný, zarmoucený, sklíčený}
\end{xtolerant}
\dicEntry[sori] \dicTerm{sor|i} \dicIPA{{s}{\textopeno}{\textlengthmark}{r}{\textsci}} \dicPos{m}[1] \dicFlx{(‑a)}[3] \textbf{1.} \dicSynonym{viðbjóður} \dicDirectTranslationCS{hnus, nechutnost, šlichta}  \textbf{2.} \dicSynonym{grugg} \dicDirectTranslationCS{kal, sedlina}
\dicEntry[sorp] \dicTerm{sorp} \dicIPA{{s}{\textopeno}{\textsubring{r}}{\textsubring{b}}} \dicPos{n}[2] \dicFlx{(‑s)}[2] \dicSynonym{úrgangur} \dicDirectTranslationCS{smetí, odpad, odpadky} \dicExampleIS{fleygja sorpi í gáminn} \dicExampleCS{vyhodit odpadky do kontejneru}
\dicEntry[sorpbíll] \dicTerm{sorp··bíl|l} \dicIPA{{s}{\textopeno}{\textsubring{r}}{\textsubring{b}}{i}{\textsubring{d}}{\textsubring{l}}} \dicPos{m}[6] \dicFlx{(‑s, ‑ar)}[48] \dicDirectTranslationCS{popelářský vůz}
\dicEntry[sorpblað] \dicTerm{sorp··|blað} \dicIPA{{s}{\textopeno}{\textsubring{r}}{\textsubring{b}}{l}{a}{\texttheta}} \dicPos{n}[2] \dicFlx{(‑blaðs, ‑blöð)}[8] \dicDirectTranslationCS{bulvár, bulvární plátek}
\dicEntry[sorpeyðing] \dicTerm{sorp··eyð·ing} \dicIPA{{s}{\textopeno}{\textsubring{r}}{\textsubring{b}}{ei}{ð}{i}{\ng}{\r{g}}} \dicPos{f}[4] \dicFlx{(‑ar)}[7] \dicDirectTranslationCS{zneškodňování\,/\addthin likvidace odpadu}
\dicEntry[sorpflokkun] \dicTerm{sorp··flokk|un} \dicIPA{{s}{\textopeno}{\textsubring{r}}{\textsubring{b}}{f}{l}{\textopeno}{h}{\r{g}}{\textscy}{\textsubring{n}}} \dicPos{f}[7] \dicFlx{(‑unar)}[9] \dicDirectTranslationCS{třídění odpadu}
\dicEntry[sorphaugur] \dicTerm{sorp··haug|ur} \dicIPA{{s}{\textopeno}{\textsubring{r}}{\textsubring{b}}{h}{\oe i}{\textbabygamma}{\textscy}{\textsubring{r}}} \dicPos{m}[6] \dicFlx{(‑s, ‑ar)}[22] \dicDirectTranslationCS{smetiště, skládka (smetí)}
\dicEntry[sorphreinsun] \dicTerm{sorp··hreins|un} \dicIPA{{s}{\textopeno}{\textsubring{r}}{\textsubring{b}}{\textsubring{r}}{ei}{n}{s}{\textscy}{\textsubring{n}}} \dicPos{f}[7] \dicFlx{(‑unar)}[9] \dicDirectTranslationCS{čištění\,/\addthin odklízení odpadu}
\dicEntry[sorppoki] \dicTerm{sorp··pok|i} \dicIPA{{s}{\textopeno}{\textsubring{r}}{\textsubring{b}}{p\smash{\textsuperscript{h}}}{\textopeno}{\r{\textObardotlessj}}{\textsci}} \dicPos{m}[1] \dicFlx{(‑a, ‑ar)}[1] \dicDirectTranslationCS{pytel na odpadky\,/\addthin smetí}
\dicEntry[sorprenna] \dicTerm{sorp··renn|a} \dicIPA{{s}{\textopeno}{\textsubring{r}}{\textsubring{b}}{r}{\textepsilon}{n}{a}} \dicPos{f}[1] \dicFlx{(‑u, ‑ur)}[7] \dicDirectTranslationCS{šachta na odpadky (v~činžovním domě ap.)}
\dicEntry[sorprit] \dicTerm{sorp··rit} \dicIPA{{s}{\textopeno}{\textsubring{r}}{\textsubring{b}}{r}{\textsci}{\textsubring{d}}} \dicPos{n}[2] \dicFlx{(‑s, ‑)}[5] \dicDirectTranslationCS{braková literatura}
\dicEntry[sorptunna] \dicTerm{sorp··tunn|a} \dicIPA{{s}{\textopeno}{\textsubring{r}}{\textsubring{b}}{t\smash{\textsuperscript{h}}}{\textscy}{n}{a}} \dicPos{f}[1] \dicFlx{(‑u, ‑ur)}[7] \dicDirectTranslationCS{popelnice}
\dicEntry[sort] \dicTerm{sort} \dicIPA{{s}{\textopeno}{\textsubring{r}}{\textsubring{d}}} \dicPos{f}[7] \dicFlx{(‑ar, ‑ir)}[1] \dicSynonym{tegund} \dicDirectTranslationCS{druh, typ, odrůda, sorta}
\dicEntry[sortera] \dicTerm{sorter|a}\dicTerm{, sortéra} \dicIPA{{s}\-{\textopeno}\-{\textsubring{r}}\-{\textsubring{d}}\-{\textepsilon}\-{r}\-{a}\-} \dicPos{v}[1] \dicFlx{(‑aði)}[1] \dicFlx{acc} \dicSynonym*{skilja sundur} \dicDirectTranslationCS{(roz)třídit, (roz)dělit, (roz)členit}
\dicEntry[sortéra] \dicTerm{sortér|a} \dicIPA{{s}{\textopeno}{\textsubring{r}}{\textsubring{d}}{j}{\textepsilon}{r}{a}} \dicPos{v}[1] \dicFlx{(‑aði)}[1] \dicLink{sortera}
\dicEntry[sorti] \dicTerm{sort|i} \dicIPA{{s}{\textopeno}{\textsubring{r}}{\textsubring{d}}{\textsci}} \dicPos{m}[1] \dicFlx{(‑a)}[3] \dicSynonym{myrkur\smash{\textsuperscript{1}}} \dicDirectTranslationCS{temno, temnota, černo}
\dicEntry[sortna] \textls[15]{\dicTerm{sortn|a} \dicIPA{{s}{\textopeno}{\textsubring{r}}{\textsubring{d}}{n}{a}} \dicPos{v}[1] \dicFlx{(‑aði)}[39] \dicSynonym*{verða svartur} \dicDirectTranslationCS{(z)černat, (z)tmavnout};}  \dicIdiom{sortna}[fyrir]{ \dicPhraseIS{e‑m sortnar fyrir augum}} \dicFlx{impers} \dicDirectTranslationCS{(komu) se dělá černo před očima, (komu) se zatmívá před očima}
\dicEntry[sortulyng] \dicTerm{sortu··lyng} \dicIPA{{s}{\textopeno}{\textsubring{r}}{\textsubring{d}}{\textscy}{l}{i}{\ng}{\r{g}}} \dicPos{n}[2] \dicFlx{(‑s)}[2] \dicFieldCat{bot.} \dicDirectTranslationCS{medvědice lékařská} \textit{(l.~{\textLA{Arctostaphylos uva-ursi}})}  \dicsymPhoto\ 
\dicFigure{38554.jpg}{Sortulyng}{Sortulyng - Vendula Šimková, Biolib, Copyright/CC-BY-NC}
\dicEntry[sosum] \dicTerm{so··sum} \dicsymFrequent\  \dicIPA{{s}{\textopeno}{\textlengthmark}{s}{\textscy}{\textsubring{m}}} \dicPos{adv} \dicSynonym*{svo sem} \dicDirectTranslationCS{vlastně, docela} \dicExampleIS{Það er sosum allt í lagi.} \dicExampleCS{Je to vlastně v~pořádku.}
\dicEntry[sovéskur] \dicTerm{sovéskur} \dicIPA{{s}{\textopeno}{v}{j}{\textepsilon}{s}{\r{g}}{\textscy}{\textsubring{r}}} \dicPos{adj}[1]\dicFlx{}[-1] \dicDirectTranslationCS{sovětský}
\dicEntry[Sovétmaður] \dicTerm{Sovét··|maður} \dicIPA{{s}{\textopeno}{v}{j}{\textepsilon}{\textsubring{d}}{m}{a}{ð}{\textscy}{\textsubring{r}}} \dicPos{m}[13] \dicFlx{(‑manns, ‑menn)}[2] \dicDirectTranslationCS{Sovět(ka)}
\dicEntry[Sovétríkin] \dicTerm{Sovét··ríkin} \dicIPA{{s}{\textopeno}{v}{j}{\textepsilon}{\textsubring{d}}{r}{i}{\r{\textObardotlessj}}{\textsci}{\textsubring{n}}} \dicPos{n}[2] \dicFlx{pl}[18] \dicFlx{def} \dicFieldCat{geog. hist.} \dicDirectTranslationCS{Sovětský svaz}
\dicEntry[sóa] \dicTerm{só|a} \dicIPA{{s}{ou}{\textlengthmark}{a}} \dicPos{v}[1] \dicFlx{(‑aði)}[1] \dicFlx{dat} \dicSynonym{glutra} \dicDirectTranslationCS{(pro)marnit, (pro)mrhat} \dicExampleIS{sóa orku í e‑ð} \dicExampleCS{promarnit energii na (čem)}
\dicEntry[sódavatn] \dicTerm{sóda··vatn} \dicIPA{{s}{ou}{\textlengthmark}{\textsubring{d}}{a}{v}{a}{h}{\textsubring{d}}{\textsubring{n}}} \dicPos{n}[2] \dicFlx{(‑s)}[2] \dicDirectTranslationCS{sodovka, sodová voda, soda}
\dicEntry[sódi] \dicTerm{sód|i} \dicIPA{{s}{ou}{\textlengthmark}{\textsubring{d}}{\textsci}} \dicPos{m}[1] \dicFlx{(‑a)}[3] \dicFieldCat{chem.} \dicDirectTranslationCS{soda}
\dicEntry[sóða] \dicTerm{sóð|a} \dicIPA{{s}{ou}{\textlengthmark}{ð}{a}} \dicPos{v}[1] \dicFlx{(‑aði)}[1] \dicFlx{acc} \dicPhraseIS{sóða e‑ð út} \dicDirectTranslationCS{zašpinit (co), znečistit (co)} \dicExampleIS{sóða út herbergið} \dicExampleCS{znečistit pokoj}
\dicEntry[sóðalegur] \dicTerm{sóða··legur} \dicIPA{{s}{ou}{\textlengthmark}{ð}{a}{l}{\textepsilon}{\textbabygamma}{\textscy}{\textsubring{r}}} \dicPos{adj}[1]\dicFlx{}[-8] \dicSynonym{óhreinn} \dicDirectTranslationCS{špinavý, ušpiněný, zašpiněný, znečištěný, ušmudlaný} \dicExampleIS{sóðalegt heimili} \dicExampleCS{špinavá domácnost}
\dicEntry[sóðaskapur] \dicTerm{sóða··skap|ur} \dicIPA{{s}{ou}{\textlengthmark}{ð}{a}{s}{\r{g}}{a}{\textsubring{b}}{\textscy}{\textsubring{r}}} \dicPos{m}[10] \dicFlx{(‑ar)}[15] \dicDirectTranslationCS{špinavost, ušpiněnost, zašpiněnost, ušmudlanost}
\dicEntry[sóði] \dicTerm{sóð|i} \dicIPA{{s}{ou}{\textlengthmark}{ð}{\textsci}} \dicPos{m}[1] \dicFlx{(‑a, ‑ar)}[1] \dicDirectTranslationCS{čuně, prase, špindíra, šmudla, nepořádník, nepořádnice} \dicExampleIS{vera voðalegur sóði} \dicExampleCS{být hrozné čuně}
\dicEntry[sófaborð] \dicTerm{sófa··borð} \dicIPA{{s}{ou}{\textlengthmark}{f}{a}{\textsubring{b}}{\textopeno}{r}{\texttheta}} \dicPos{n}[2] \dicFlx{(‑s, ‑)}[5] \dicDirectTranslationCS{konferenční stolek}
\dicEntry[sófasett] \dicTerm{sófa··sett} \dicIPA{{s}{ou}{\textlengthmark}{f}{a}{s}{\textepsilon}{h}{\textsubring{d}}} \dicPos{n}[2] \dicFlx{(‑s, ‑)}[5] \dicDirectTranslationCS{sedací souprava}
\dicEntry[sófi] \dicTerm{sóf|i} \dicIPA{{s}{ou}{\textlengthmark}{f}{\textsci}} \dicPos{m}[1] \dicFlx{(‑a, ‑ar)}[1] \dicDirectTranslationCS{pohovka, gauč, sofa}
\dicEntry[Sófía] \dicTerm{Sófí|a}\dicTerm{, Sofía} \dicIPA{{s}\-{ou}\-{\textlengthmark}\-{f}\-{i}\-{j}\-{a}\-} \dicPos{f}[1] \dicFlx{(‑u)}[6] \dicFieldCat{geog.} \dicDirectTranslationCS{Sofie} \dicIndirectTranslationCS{(hlavní město Bulharska)}
\dicEntry[sókn] \dicTerm{sókn} \dicsymFrequent\  \dicIPA{{s}{ou}{h}{\r{g}}{\textsubring{n}}} \dicPos{f}[7] \dicFlx{(‑ar, ‑ir)}[1] \textbf{1.} \dicSynonym{árás} \dicDirectTranslationCS{útok, ofenziva}  \textbf{2.} \dicSynonym{ásókn} \dicDirectTranslationCS{využívání (přírodních) zdrojů}  \textbf{3.} \dicFieldCat{práv.} \dicDirectTranslationCS{trestní stíhání\,/\addthin řízení}  \textbf{4.} \dicSynonym{kirkjusókn} \dicDirectTranslationCS{farnost}  \textbf{5.} \dicFieldCat{sport.} \dicDirectTranslationCS{útok} \dicExampleIS{Sóknin var ekkert sérstök.} \dicExampleCS{Útok nebyl nijak výjimečný.} \dicAntonym{vörn};  \dicPhraseIS{vera í sókn} \dicDirectTranslationCS{být v~útoku, (za)útočit}
\dicEntry[sóknaraðili] \dicTerm{sóknar··aðil|i} \dicIPA{{s}{ou}{h}{\r{g}}{n}{a}{r}{a}{ð}{\textsci}{l}{\textsci}} \dicPos{m}[1] \dicFlx{(‑a\,/\addthin ‑ja, ‑ar\,/\addthin ‑jar)}[15] \dicFieldCat{práv.} \dicDirectTranslationCS{navrhovatel(ka), žalobce, žalobkyně}
\dicEntry[sóknarbarn] \dicTerm{sóknar··|barn} \dicIPA{{s}{ou}{h}{\r{g}}{n}{a}{r}{\textsubring{b}}{a}{r}{\textsubring{d}}{\textsubring{n}}} \dicPos{n}[2] \dicFlx{(‑barns, ‑börn)}[8] \dicDirectTranslationCS{farník, farnice}
\dicEntry[sóknarkirkja] \dicTerm{sóknar··kirkj|a} \dicIPA{{s}{ou}{h}{\r{g}}{n}{a}{\textsubring{r}}{c\smash{\textsuperscript{h}}}{\textsci}{\textsubring{r}}{\r{\textObardotlessj}}{a}} \dicPos{f}[1] \dicFlx{(‑u, ‑ur)}[17] \dicFieldCat{náb.} \dicDirectTranslationCS{farní kostel}
\dicEntry[sóknarmaður] \dicTerm{sóknar··|maður} \dicIPA{{s}{ou}{h}{\r{g}}{n}{a}{r}{m}{a}{ð}{\textscy}{\textsubring{r}}} \dicPos{m}[13] \dicFlx{(‑manns, ‑menn)}[2] \textbf{1.} \dicSynonym{sóknarbarn} \dicDirectTranslationCS{farník, farnice}  \textbf{2.} \dicFieldCat{sport.} \dicDirectTranslationCS{útočník, útočnice}
\dicEntry[sóknarmark] \dicTerm{sóknar··|mark} \dicIPA{{s}{ou}{h}{\r{g}}{n}{a}{r}{m}{a}{\textsubring{r}}{\r{g}}} \dicPos{n}[2] \dicFlx{(‑marks, ‑mörk)}[8] \dicDirectTranslationCS{časové omezení rybolovu} \dicIndirectTranslationCS{(regulace rybolovu)}
\dicEntry[sóknarnefnd] \dicTerm{sóknar··nefnd} \dicIPA{{s}{ou}{h}{\r{g}}{n}{a}{r}{n}{\textepsilon}{m}{\textsubring{d}}} \dicPos{f}[7] \dicFlx{(‑ar, ‑ir)}[1] \dicDirectTranslationCS{farní rada}
\dicEntry[sóknarprestur] \dicTerm{sóknar··prest|ur} \dicIPA{{s}{ou}{h}{\r{g}}{n}{a}{\textsubring{r}}{p\smash{\textsuperscript{h}}}{r}{\textepsilon}{s}{\textsubring{d}}{\textscy}{\textsubring{r}}} \dicPos{m}[6] \dicFlx{(‑s, ‑ar)}[4] \dicDirectTranslationCS{farář(ka)}
\dicEntry[sól] \dicTerm{sól} \dicsymFrequent\  \dicIPA{{s}{ou}{\textlengthmark}{\textsubring{l}}} \dicPos{f}[7] \dicFlx{(‑ar, ‑ir)}[7] \dicDirectTranslationCS{slunce} \dicExampleIS{Sólin lýsir fjöllin.} \dicExampleCS{Slunce osvětluje hory.};  \dicPhraseIS{sólin kemur upp} \dicDirectTranslationCS{slunce vychází};  \dicPhraseIS{sólin sest} \dicFlx{refl} \dicDirectTranslationCS{slunce zapadá};  \dicPhraseIS{það sést ekki til sólar} \dicFlx{refl impers} \dicDirectTranslationCS{je zataženo, slunce se schovává za mraky};  \dicIdiom{sól}{ \dicPhraseIS{eins og dögg fyrir sólu}} \dicFlx{adv} \dicLangCat{přen.} \dicDirectTranslationCS{jako mávnutím kouzelného proutku}; { \dicPhraseIS{sjá ekki sólina fyrir e‑m}} \dicLangCat{přen.} \dicDirectTranslationCS{mít oči jen pro (koho)}
\dicEntry[sóla] \dicTerm{sól|a\smash{\textsuperscript{1}}} \dicIPA{{s}{ou}{\textlengthmark}{l}{a}} \dicPos{v}[1] \dicFlx{(‑aði)}[1] \dicFlx{acc} \textbf{1.} \dicDirectTranslationCS{podrazit (boty ap.)}  \textbf{2.} \dicDirectTranslationCS{protektorovat (pneumatiku ap.)}
\dicEntry[sóla] \dicTerm{sól|a\smash{\textsuperscript{2}}} \dicIPA{{s}{ou}{\textlengthmark}{l}{a}} \dicPos{v}[1] \dicFlx{(‑aði)}[1] \dicFlx{acc} \dicPhraseIS{sóla sig} \dicDirectTranslationCS{slunit se, vyhřívat se na slunci}
\dicEntry[sólarár] \dicTerm{sólar··ár} \dicIPA{{s}{ou}{\textlengthmark}{l}{a}{r}{au}{\textsubring{r}}} \dicPos{n}[2] \dicFlx{(‑s, ‑)}[5] \dicFieldCat{astro.} \dicDirectTranslationCS{sluneční\,/\addthin solární rok}
\dicEntry[sólargeisli] \dicTerm{sólar··geisl|i} \dicIPA{{s}{ou}{\textlengthmark}{l}{a}{r}{\r{\textObardotlessj}}{ei}{s}{\textsubring{d}}{l}{\textsci}} \dicPos{m}[1] \dicFlx{(‑a, ‑ar)}[1] \textbf{1.} \dicDirectTranslationCS{sluneční paprsek}  \textbf{2.} \dicLangCat{přen.} \dicDirectTranslationCS{sluníčko} \dicIndirectTranslationCS{(usměvavý pohodový člověk)}
\dicEntry[sólarhringur] \dicTerm{sólar··hring|ur} \dicsymFrequent\  \dicIPA{{s}{ou}{\textlengthmark}{l}{a}{\textsubring{r}}{i}{\ng}{\r{g}}{\textscy}{\textsubring{r}}} \dicPos{m}[6] \dicFlx{(‑s, ‑ar)}[15] \dicDirectTranslationCS{den} \dicIndirectTranslationCS{(doba 24 hodin)} \dicExampleIS{Skipið var þrjá sólarhringa á leiðinni.} \dicExampleCS{Loď byla tři dny na cestě.}
\dicEntry[sólarlag] \dicTerm{sólar··|lag} \dicIPA{{s}{ou}{\textlengthmark}{l}{a}{r}{l}{a}{x}} \dicPos{n}[2] \dicFlx{(‑lags, ‑lög)}[8] \dicSynonym{sólsetur} \dicDirectTranslationCS{západ slunce} \dicExampleIS{fagurt sólarlag} \dicExampleCS{krásný západ slunce}
\dicEntry[sólarlampi] \dicTerm{sólar··lamp|i} \dicIPA{{s}{ou}{\textlengthmark}{l}{a}{r}{l}{a}{\textsubring{m}}{\textsubring{b}}{\textsci}} \dicPos{m}[1] \dicFlx{(‑a, ‑ar)}[8] \dicDirectTranslationCS{solární lampa}
\dicEntry[sólarland] \dicTerm{sólar··|land} \dicIPA{{s}{ou}{\textlengthmark}{l}{a}{r}{l}{a}{n}{\textsubring{d}}} \dicPos{n}[2] \dicFlx{(‑lands, ‑lönd)}[8] \dicDirectTranslationCS{teplá krajina} \dicIndirectTranslationCS{(země s~teplým klimatem, zvláště v~pl)}
\dicEntry[sólarljós] \dicTerm{sólar··ljós} \dicIPA{{s}{ou}{\textlengthmark}{l}{a}{r}{l}{j}{ou}{s}} \dicPos{n}[2] \dicFlx{(‑s, ‑)}[5] \dicDirectTranslationCS{sluneční světlo}
\dicEntry[sólarmegin] \dicTerm{sólar··megin} \dicIPA{{s}{ou}{\textlengthmark}{l}{a}{r}{m}{ei}{\textsci}{\textsubring{n}}} \dicPos{adv} \dicDirectTranslationCS{na sluneční straně};  \dicPhraseIS{vera sólarmegin í lífinu} \dicLangCat{přen.} \dicDirectTranslationCS{být šťastný, nestrádat}
\dicEntry[sólarolía] \dicTerm{sólar··olí|a}\dicTerm{, sólolía} \dicIPA{{s}\-{ou}\-{\textlengthmark}\-{l}\-{a}\-{r}\-{\textopeno}\-{l}\-{i}\-{j}\-{a}\-} \dicPos{f}[1] \dicFlx{(‑u, ‑ur)}[7] \dicDirectTranslationCS{opalovací krém}
\dicEntry[sólarorka] \dicTerm{sólar··ork|a} \dicIPA{{s}{ou}{\textlengthmark}{l}{a}{r}{\textopeno}{\textsubring{r}}{\r{g}}{a}} \dicPos{f}[1] \dicFlx{(‑u)}[5] \dicDirectTranslationCS{sluneční energie}
\begin{xtolerant}{}{1pt}
\dicEntry[sólarströnd] \dicTerm{sólar··|strönd}\addthinS\dicIPA{{s}{ou}{\textlengthmark}{l}{a}{\textsubring{r}}{s}{\textsubring{d}}{r}{\oe}{n}{\textsubring{d}}}\addthinS\dicPos{f}[8]\addthinS\dicFlx{(‑strandar, ‑strendur\,/\addthin ‑strand\-ir)}[9] \dicDirectTranslationCS{pláž} \dicIndirectTranslationCS{(v~jižních zemích)}
\end{xtolerant}
\dicEntry[sólaruppkoma] \textls[10]{\dicTerm{sólar··upp·kom|a} \dicIPA{{s}{ou}{\textlengthmark}{l}{a}{r}{\textscy}{h}{\textsubring{b}}{k\smash{\textsuperscript{h}}}{\textopeno}{m}{a}} \dicPos{f}[1] \dicFlx{(‑u)}[5] \dicDirectTranslationCS{východ slunce}}
\begin{xtolerant}{}{1pt}
\dicEntry[sólarupprás] \dicTerm{sólar··upp·rás}\addthinS\dicIPA{{s}{ou}{\textlengthmark}{l}{a}{r}{\textscy}{h}{\textsubring{b}}{r}{au}{s}}\addthinS\dicPos{f}[7]\addthinS\dicFlx{(‑ar, ‑ir)}[1] \dicDirectTranslationCS{východ slunce}
\end{xtolerant}
\dicEntry[sólarvörn] \dicTerm{sólar··|vörn} \dicIPA{{s}{ou}{\textlengthmark}{l}{a}{r}{v}{\oe}{r}{\textsubring{d}}{\textsubring{n}}} \dicPos{f}[7] \dicFlx{(‑varnar, ‑varnir)}[16] \dicDirectTranslationCS{krém na opalování, opalovací krém}
\dicEntry[sólbað] \dicTerm{sól··|bað} \dicIPA{{s}{ou}{l}{\textsubring{b}}{a}{\texttheta}} \dicPos{n}[2] \dicFlx{(‑baðs, ‑böð)}[8] \dicDirectTranslationCS{opalování, slunění} \dicExampleIS{fara í sólbað} \dicExampleCS{opalovat se};  \dicPhraseIS{liggja í sólbaði} \dicDirectTranslationCS{opalovat se, slunit se}
\dicEntry[sólbaðsstofa] \dicTerm{sól·baðs··stof|a} \dicIPA{{s}{ou}{l}{\textsubring{b}}{a}{ð}{s}{\textsubring{d}}{\textopeno}{v}{a}} \dicPos{f}[1] \dicFlx{(‑u, ‑ur)}[7] \dicDirectTranslationCS{solárium}
\dicEntry[sólbekkur] \dicTerm{sól··bekk|ur} \dicIPA{{s}{ou}{l}{\textsubring{b}}{\textepsilon}{h}{\r{g}}{\textscy}{\textsubring{r}}} \dicPos{m}[9] \dicFlx{(‑jar\,/\addthin ‑s, ‑ir)}[26] \textbf{1.} \dicDirectTranslationCS{lehátko (na opalování)}  \textbf{2.} \dicSynonym{gluggakista} \dicDirectTranslationCS{okenní římsa, parapet}
\dicEntry[sólber] \dicTerm{sól··ber} \dicIPA{{s}{ou}{l}{\textsubring{b}}{\textepsilon}{\textsubring{r}}} \dicPos{n}[2] \dicFlx{(‑s, ‑)}[13] \dicDirectTranslationCS{černý rybíz} \dicIndirectTranslationCS{(bobule)}
\dicEntry[sólberjarunni] \dicTerm{sól·berja··runn|i} \dicIPA{{s}{ou}{l}{\textsubring{b}}{\textepsilon}{r}{j}{a}{r}{\textscy}{n}{\textsci}} \dicPos{m}[1] \dicFlx{(‑a, ‑ar)}[1] \dicFieldCat{bot.} \dicDirectTranslationCS{černý rybíz} \textit{(l.~{\textLA{Ribes nigrum}})}  \dicsymPhoto\ 
\dicFigure{ds_image_solberjarunni_0_1.jpg}{Sólberjarunni}{Sólberjarunni - Saxo, PD}
\dicEntry[sólblettur] \textls[15]{\dicTerm{sól··blett|ur} \dicIPA{{s}{ou}{l}{\textsubring{b}}{l}{\textepsilon}{h}{\textsubring{d}}{\textscy}{\textsubring{r}}} \dicPos{m}[9] \dicFlx{(‑s, ‑ir)}[4] \dicFieldCat{astro.} \dicDirectTranslationCS{sluneční skvrna}}
\dicEntry[sólblóm] \dicTerm{sól··blóm} \dicIPA{{s}{ou}{l}{\textsubring{b}}{l}{ou}{\textsubring{m}}} \dicPos{n}[2] \dicFlx{(‑s, ‑)}[5] \dicFieldCat{bot.} \dicDirectTranslationCS{slunečnice, slunečnice roční} \textit{(l.~{\textLA{Helianthus annuus}})}  \dicsymPhoto\ 
\dicFigure{ds_image_solblom_0_1.jpg}{Sólblóm}{Sólblóm - Bruce Fritz, PD}
\dicEntry[sólblómafræ] \dicTerm{sól·blóma··fræ} \dicIPA{{s}{ou}{l}{\textsubring{b}}{l}{ou}{m}{a}{f}{r}{a}{i}} \dicPos{n}[2] \dicFlx{(‑s, ‑)}[13] \dicDirectTranslationCS{slunečnicové semínko}
\dicEntry[sólblómaolía] \dicTerm{sól·blóma··olí|a} \dicIPA{{s}{ou}{l}{\textsubring{b}}{l}{ou}{m}{a}{\textopeno}{l}{i}{j}{a}} \dicPos{f}[1] \dicFlx{(‑u)}[5] \dicDirectTranslationCS{slunečnicový olej}
\dicEntry[sólbrenna] \dicTerm{sól··|brenna} \dicIPA{{s}{ou}{l}{\textsubring{b}}{r}{\textepsilon}{n}{a}} \dicPos{v}[6] \dicFlx{(‑brenn, ‑brann, ‑brunnum, ‑brynni, ‑brunnið)}[24] \dicDirectTranslationCS{spálit se (na slunci)}
\dicEntry[sólbrenndur] \dicTerm{sól··brenndur} \dicIPA{{s}{ou}{l}{\textsubring{b}}{r}{\textepsilon}{n}{\textsubring{d}}{\textscy}{\textsubring{r}}} \dicPos{adj}[2]\dicFlx{}[-14] \dicDirectTranslationCS{spálený (sluncem), opálený}
\dicEntry[sólbruni] \dicTerm{sól··brun|i} \dicIPA{{s}{ou}{l}{\textsubring{b}}{r}{\textscy}{n}{\textsci}} \dicPos{m}[1] \dicFlx{(‑a)}[3] \dicDirectTranslationCS{spálení (sluncem)}
\dicEntry[sólbrunninn] \dicTerm{sól··brunninn} \dicIPA{{s}{ou}{l}{\textsubring{b}}{r}{\textscy}{n}{\textsci}{\textsubring{n}}} \dicPos{adj}[6]\dicFlx{}[-15] \dicDirectTranslationCS{spálený (sluncem)}
\dicEntry[sólbrúnka] \dicTerm{sól··brúnk|a} \dicIPA{{s}{ou}{l}{\textsubring{b}}{r}{u}{\r{\ng}}{\r{g}}{a}} \dicPos{f}[1] \dicFlx{(‑u, ‑ur)}[19] \dicDirectTranslationCS{opálení}
\dicEntry[sólbrúnn] \dicTerm{sól··brúnn} \dicIPA{{s}{ou}{l}{\textsubring{b}}{r}{u}{\textsubring{d}}{\textsubring{n}}} \dicPos{adj}[7]\dicFlx{}[-1] \dicDirectTranslationCS{opálený}
\dicEntry[sólböð] \dicTerm{sól··böð} \dicIPA{{s}{ou}{l}{\textsubring{b}}{\oe}{\texttheta}} \dicPos{n} \dicFlx{pl nom} \dicLink{sólbað}
\dicEntry[sóldögg] \dicTerm{sól··|dögg} \dicIPA{{s}{ou}{l}{\textsubring{d}}{\oe}{\r{g}}} \dicPos{f}[7] \dicFlx{(‑daggar, ‑daggir)}[16] \dicFieldCat{bot.} \dicDirectTranslationCS{rosnatka okrouhlolistá} \textit{(l.~{\textLA{Drosera rotundifolia}})}  \dicsymPhoto\ 
\dicFigure{657.jpg}{Sóldögg}{Sóldögg - Zicha Ondřej, Biolib, Copyright/CC-BY-NC}
\dicEntry[sóley] \dicTerm{sól··ey} \dicIPA{{s}{ou}{\textlengthmark}{l}{ei}} \dicPos{f}[4] \dicFlx{(‑jar, ‑jar)}[16] \dicFieldCat{bot.} \dicDirectTranslationCS{pryskyřník, pryskyřník prudký} \textit{(l.~{\textLA{Ranunculus acris}})}
\dicEntry[sólfar] \dicTerm{sól··far} \dicIPA{{s}{ou}{l}{f}{a}{\textsubring{r}}} \dicPos{n}[2] \dicFlx{(‑s)}[2] \dicSynonym{sólskin} \dicDirectTranslationCS{sluneční světlo} \dicExampleIS{mikið sólfar} \dicExampleCS{hodně slunečního světla}
\dicEntry[sólginn] \dicTerm{sólginn} \dicIPA{{s}{ou}{l}{\r{\textObardotlessj}}{\textsci}{\textsubring{n}}} \dicPos{adj}[6]\dicFlx{}[-2] \dicSynonym{fíkinn} \dicDirectTranslationCS{dychtivý, žádostivý, lačný};  \dicPhraseIS{sólginn í e‑ð} \dicDirectTranslationCS{dychtivý (čeho)} \dicExampleIS{sólginn í fréttir} \dicExampleCS{dychtivý zpráv}
\dicEntry[sólgleraugu] \dicTerm{sól··gler·augu} \dicIPA{{s}{ou}{l}{\r{g}}{l}{\textepsilon}{r}{\oe i}{\textbabygamma}{\textscy}} \dicPos{n}[1] \dicFlx{pl}[3] \dicDirectTranslationCS{sluneční brýle}
\dicEntry[sólhattur] \dicTerm{sól··hatt|ur} \dicIPA{{s}{ou}{\textlengthmark}{\textsubring{l}}{h}{a}{h}{\textsubring{d}}{\textscy}{\textsubring{r}}} \dicPos{m}[6] \dicFlx{(‑s, ‑ar)}[5] \textbf{1.} \dicSynonym{hattur} \dicDirectTranslationCS{(sluneční) klobouk}  \textbf{2.} \dicFieldCat{bot.} \dicDirectTranslationCS{třapatka, třapatka nachová} \textit{(l.~{\textLA{Echinacea purpurea}})}  \dicsymPhoto\ 
\dicFigure{ds_image_solhattur_0_1.jpg}{Sólhattur}{Sólhattur - Hedwig Storch, CC BY-SA 3.0}
\dicEntry[sólhlíf] \dicTerm{sól··hlíf} \dicIPA{{s}{ou}{l}{\textsubring{l}}{i}{f}} \dicPos{f}[4] \dicFlx{(‑ar, ‑ar)}[1] \dicDirectTranslationCS{slunečník}
\dicEntry[sólhvörf] \dicTerm{sól··hvörf} \dicIPA{{s}{ou}{l}{k\smash{\textsuperscript{h}}}{v}{\oe}{r}{f}} \dicPos{n}[2] \dicFlx{pl}[9] \dicDirectTranslationCS{slunovrat}
\dicEntry[sóli] \dicTerm{sól|i} \dicIPA{{s}{ou}{\textlengthmark}{l}{\textsci}} \dicPos{m}[1] \dicFlx{(‑a, ‑ar)}[1] \dicDirectTranslationCS{podrážka, podešev}
\dicEntry[sólkerfi] \dicTerm{sól··kerfi} \dicIPA{{s}{ou}{l}{c\smash{\textsuperscript{h}}}{\textepsilon}{r}{v}{\textsci}} \dicPos{n}[2] \dicFlx{(‑s, ‑)}[14] \dicFieldCat{astro.} \dicDirectTranslationCS{sluneční soustava}
\dicEntry[sólmyrkvi] \dicTerm{sól··myrkv|i} \dicIPA{{s}{ou}{l}{m}{\textsci}{\textsubring{r}}{\r{g}}{v}{\textsci}} \dicPos{m}[1] \dicFlx{(‑a, ‑ar)}[1] \dicFieldCat{astro.} \dicDirectTranslationCS{zatmění slunce}
\dicEntry[sólnánd] \dicTerm{sól··nánd} \dicIPA{{s}{ou}{l}{n}{au}{n}{\textsubring{d}}} \dicPos{f}[7] \dicFlx{(‑ar)}[3] \dicFieldCat{astro.} \dicDirectTranslationCS{přísluní, perihélium}
\dicEntry[sólolía] \dicTerm{sól··olí|a} \dicIPA{{s}{ou}{\textlengthmark}{l}{\textopeno}{l}{i}{j}{a}} \dicPos{f}[1] \dicFlx{(‑u, ‑ur)}[7] \dicLink{sólarolía}
\dicEntry[sóló] \dicTerm{sóló\smash{\textsuperscript{1}}} \dicIPA{{s}{ou}{\textlengthmark}{l}{ou}} \dicPos{f}[7] \dicFlx{(‑ar, ‑ir)}[1] \dicLink{sóló\smash{\textsuperscript{2}}}
\dicEntry[sóló] \dicTerm{sóló\smash{\textsuperscript{2}}}\dicTerm{, sóló} \dicIPA{{s}\-{ou}\-{\textlengthmark}\-{l}\-{ou}\-} \dicPos{n}[2] \dicFlx{(‑s, ‑)}[5] \textbf{1.} \dicFieldCat{hud.} \dicDirectTranslationCS{sólo} \dicIndirectTranslationCS{(skladba určená jedinému interpretovi)}  \textbf{2.} \dicFieldCat{sport.} \dicDirectTranslationCS{sólo, sólová hra}
\dicEntry[sólríkur] \dicTerm{sól··ríkur} \dicIPA{{s}{ou}{l}{r}{i}{\r{g}}{\textscy}{\textsubring{r}}} \dicPos{adj}[1]\dicFlx{}[-1] \dicDirectTranslationCS{slunečný, slunný}
\dicEntry[sólsetur] \dicTerm{sól··setur} \dicIPA{{s}{ou}{l}{s}{\textepsilon}{\textsubring{d}}{\textscy}{\textsubring{r}}} \dicPos{n}[2] \dicFlx{(‑s, ‑)}[25] \dicDirectTranslationCS{západ slunce}
\dicEntry[sólskin] \dicTerm{sól··skin} \dicsymFrequent\  \dicIPA{{s}{ou}{l}{s}{\r{\textObardotlessj}}{\textsci}{\textsubring{n}}} \dicPos{n}[2] \dicFlx{(‑s)}[2] \dicDirectTranslationCS{sluneční světlo\,/\addthin záře\,/\addthin svit} \dicExampleIS{Það er sólskin og hiti.} \dicExampleCS{Je slunečno a~teplo.}
\dicEntry[sólskinsskap] \dicTerm{sól·skins··|skap} \dicIPA{{s}{ou}{l}{s}{\r{\textObardotlessj}}{\textsci}{n}{s}{\r{g}}{a}{\textsubring{b}}} \dicPos{n}[2] \dicFlx{(‑skaps, ‑sköp)}[8] \dicDirectTranslationCS{výbor\-ná\,/\addthin dobrá nálada} \dicExampleIS{vera í sólskinsskapi} \dicExampleCS{mít výbor\-nou náladu}
\dicEntry[sólskríkja] \dicTerm{sól··skríkj|a} \dicIPA{{s}{ou}{l}{s}{\r{g}}{r}{i}{\r{\textObardotlessj}}{a}} \dicPos{f}[1] \dicFlx{(‑u, ‑ur)}[7] \dicFieldCat{zool.} \dicDirectTranslationCS{sněhule severní} \dicIndirectTranslationCS{(označení ptáka v~letním období)} \textit{(l.~{\textLA{Plectrophenax nivalis}})}  \dicsymPhoto\  \dicAntonym{snjótittlingur}
\dicFigure{ds_image_solskrikja_0_1.jpg}{Sólskríkja}{Sólskríkja - U.S. Fish and Wildlife Service, PD}
\dicEntry[sólstingur] \dicTerm{sól··sting|ur} \dicIPA{{s}{ou}{l}{s}{\textsubring{d}}{i}{\ng}{\r{g}}{\textscy}{\textsubring{r}}} \dicPos{m}[9] \dicFlx{(‑s, ‑ir)}[15] \dicDirectTranslationCS{sluneční úžeh, úpal}
\dicEntry[sólstóll] \dicTerm{sól··stól|l} \dicIPA{{s}{ou}{l}{s}{\textsubring{d}}{ou}{\textsubring{d}}{\textsubring{l}}} \dicPos{m}[6] \dicFlx{(‑s, ‑ar)}[48] \dicDirectTranslationCS{lehátko (na opalování ap.)}
\dicEntry[sóltjald] \dicTerm{sól··|tjald} \dicIPA{{s}{ou}{l}{t\smash{\textsuperscript{h}}}{j}{a}{l}{\textsubring{d}}} \dicPos{n}[2] \dicFlx{(‑tjalds, ‑tjöld)}[8] \dicDirectTranslationCS{(zahradní) slunečník, markýza} \dicIndirectTranslationCS{(skládací stříška)}
\dicEntry[sólunda] \dicTerm{sólund|a} \dicIPA{{s}{ou}{\textlengthmark}{l}{\textscy}{n}{\textsubring{d}}{a}} \dicPos{v}[1] \dicFlx{(‑aði)}[1] \dicFlx{dat} \dicSynonym{sóa} \dicDirectTranslationCS{(pro)marnit, (pro)mrhat} \dicExampleIS{sólunda peningum} \dicExampleCS{promarnit peníze}
\dicEntry[sólúr] \dicTerm{sól··úr} \dicIPA{{s}{ou}{\textlengthmark}{l}{u}{\textsubring{r}}} \dicPos{n}[2] \dicFlx{(‑s, ‑)}[5] \dicDirectTranslationCS{sluneční hodiny}
\dicEntry[sómafólk] \dicTerm{sóma··fólk} \dicIPA{{s}{ou}{\textlengthmark}{m}{a}{f}{ou}{\textsubring{l}}{\r{g}}} \dicPos{n}[2] \dicFlx{(‑s)}[2] \dicDirectTranslationCS{slušní\,/\addthin vážení lidé}
\dicEntry[sómakennd] \dicTerm{sóma··kennd} \dicIPA{{s}{ou}{\textlengthmark}{m}{a}{c\smash{\textsuperscript{h}}}{\textepsilon}{n}{\textsubring{d}}} \dicPos{f}[7] \dicFlx{(‑ar)}[3] \dicDirectTranslationCS{slušnost, váženost, počestnost}
\dicEntry[sómakær] \dicTerm{sóma··kær} \dicIPA{{s}{ou}{\textlengthmark}{m}{a}{c\smash{\textsuperscript{h}}}{a}{i}{\textsubring{r}}} \dicPos{adj}[5] \dicFlx{(f ‑)}[8] \dicDirectTranslationCS{slušný, vážený, počestný}
\dicEntry[Sómali] \dicTerm{Sómal|i} \dicIPA{{s}{ou}{\textlengthmark}{m}{a}{l}{\textsci}} \dicPos{m}[1] \dicFlx{(‑a, ‑ar)}[8] \dicDirectTranslationCS{Somálec, Somálka}
\dicEntry[Sómalía] \dicTerm{Sómalí|a} \dicIPA{{s}{ou}{\textlengthmark}{m}{a}{l}{i}{j}{a}} \dicPos{f}[1] \dicFlx{(‑u)}[6] \dicFieldCat{geog.} \dicDirectTranslationCS{Somálsko}
\dicEntry[sómalískur] \dicTerm{sómalískur} \dicIPA{{s}{ou}{\textlengthmark}{m}{a}{l}{i}{s}{\r{g}}{\textscy}{\textsubring{r}}} \dicPos{adj}[1]\dicFlx{}[-6] \dicDirectTranslationCS{somálský}
\dicEntry[sómalska] \dicTerm{sómalska} \dicIPA{{s}{ou}{\textlengthmark}{m}{a}{l}{s}{\r{g}}{a}} \dicPos{f}[1] \dicFlx{(sómölsku)}[2] \dicDirectTranslationCS{somálština}
\begin{xtolerant}{}{1pt}
\dicEntry[sómamaður] \dicTerm{sóma··|maður}\addthinS\dicIPA{{s}{ou}{\textlengthmark}{m}{a}{m}{a}{ð}{\textscy}{\textsubring{r}}}\addthinS\dicPos{m}[13]\addthinS\dicFlx{(‑manns, ‑menn)}[2]\addthinS\dicDirectTranslationCS{slušný\,/\addthin vážený člověk}
\end{xtolerant}
\dicEntry[sómasamlegur] \dicTerm{sóma·sam··legur} \dicIPA{{s}{ou}{\textlengthmark}{m}{a}{s}{a}{m}{l}{\textepsilon}{\textbabygamma}{\textscy}{\textsubring{r}}} \dicPos{adj}[1]\dicFlx{}[-8] \dicSynonym{sæmilegur} \dicDirectTranslationCS{slušný, úctyhodný, reprezentativní}
\dicEntry[sómatilfinning] \dicTerm{sóma··til·finn·ing} \dicIPA{{s}{ou}{\textlengthmark}{m}{a}{t\smash{\textsuperscript{h}}}{\textsci}{l}{f}{\textsci}{n}{i}{\ng}{\r{g}}} \dicPos{f}[4] \dicFlx{(‑ar, ‑ar)}[5] \dicDirectTranslationCS{pocit slušnosti, smysl pro čest}
\dicEntry[sómi] \dicTerm{sóm|i} \dicIPA{{s}{ou}{\textlengthmark}{m}{\textsci}} \dicPos{m}[1] \dicFlx{(‑a)}[3] \dicSynonym{heiður\smash{\textsuperscript{1}}} \dicDirectTranslationCS{čest, úcta} \dicExampleIS{veita e‑m mikinn sóma með e‑u} \dicExampleCS{projevit (komu) velkou čest (čím)};  \dicPhraseIS{e‑að er e‑m til sóma} \dicDirectTranslationCS{(co) je (komu) ke cti};  \dicPhraseIS{sjá sóma sinn í því að (gera e‑ð)} \dicDirectTranslationCS{udělat čest svému jménu a~(udělat (co))};  \dicPhraseIS{það er allt í sómanum} \dicLangCat{přen.} \dicDirectTranslationCS{všechno v~nejlepším pořádku}
\dicEntry[sónar] \dicTerm{sónar} \dicIPA{{s}{ou}{\textlengthmark}{n}{a}{\textsubring{r}}} \dicPos{m}[4] \dicFlx{(‑s)}[18] \textbf{1.} \dicFieldCat{med. } \dicSynonym{ómskoðun} \dicDirectTranslationCS{ultrazvukové\,/\addthin sonografické vyšetření, sonografie, ultrazvuk}  \textbf{2.} \dicFieldCat{med. } \dicSynonym{ómsjá} \dicDirectTranslationCS{sonograf}
\dicEntry[sónarskoðun] \dicTerm{sónar··skoð|un} \dicIPA{{s}{ou}{\textlengthmark}{n}{a}{\textsubring{r}}{s}{\r{g}}{\textopeno}{ð}{\textscy}{\textsubring{n}}} \dicPos{f}[7] \dicFlx{(‑unar)}[9] \dicFieldCat{med. } \dicDirectTranslationCS{ultrazvukové\,/\addthin sonografické vyšetření}
\dicEntry[sónartæki] \dicTerm{sónar··tæki} \dicIPA{{s}{ou}{\textlengthmark}{n}{a}{\textsubring{r}}{t\smash{\textsuperscript{h}}}{a}{i}{\r{\textObardotlessj}}{\textsci}} \dicPos{n}[2] \dicFlx{(‑s, ‑)}[16] \dicFieldCat{med. } \dicDirectTranslationCS{sonograf}
\dicEntry[sónata] \dicTerm{sónata} \dicIPA{{s}{ou}{\textlengthmark}{n}{a}{\textsubring{d}}{a}} \dicPos{f}[1] \dicFlx{(sónötu, sónötur)}[8] \dicFieldCat{hud.} \dicDirectTranslationCS{sonáta}
\dicEntry[sónn] \dicTerm{són|n} \dicIPA{{s}{ou}{\textsubring{d}}{\textsubring{n}}} \dicPos{m}[6] \dicFlx{(‑s, ‑ar)}[42] \dicDirectTranslationCS{oznamovací tón} \dicIndirectTranslationCS{(v~telefonu)}
\dicEntry[sónötu] \dicTerm{sónötu} \dicIPA{{s}{ou}{\textlengthmark}{n}{\oe}{\textsubring{d}}{\textscy}} \dicPos{f} \dicFlx{sg gen} \dicLink{sónata}
\dicEntry[sópa] \dicTerm{sóp|a} \dicIPA{{s}{ou}{\textlengthmark}{\textsubring{b}}{a}} \dicPos{v}[1] \dicFlx{(‑aði)}[7] \dicFlx{acc\,/\addthin dat} \dicSynonym*{hreinsa með sópi} \dicDirectTranslationCS{(za)mést, zametat} \dicExampleIS{sópa gólfið\,/\addthin gólfinu} \dicExampleCS{zamést podlahu};  \dicIdiom{sópa}[að]{ \dicPhraseIS{það sópar að e‑m}} \dicFlx{impers} \dicDirectTranslationCS{(kdo) je okouzlen, (komu) imponuje};  \dicIdiom{sópa}[saman]{ \dicPhraseIS{sópa e‑u saman}} \dicDirectTranslationCS{shromáždit (co), posbírat (co) (účtenky ap.)}
\dicEntry[sópran] \dicTerm{sópran} \dicIPA{{s}{ou}{\textlengthmark}{\textsubring{b}}{r}{a}{\textsubring{n}}} \dicPos{m}[4] \dicFlx{(‑s, ‑ar)}[11] \dicFieldCat{hud.} \dicDirectTranslationCS{soprán}
\dicEntry[sópun] \dicTerm{sóp|un} \dicIPA{{s}{ou}{\textlengthmark}{\textsubring{b}}{\textscy}{\textsubring{n}}} \dicPos{f}[7] \dicFlx{(‑unar)}[9] \dicDirectTranslationCS{zametání}
\dicEntry[sópur] \dicTerm{sóp|ur} \dicIPA{{s}{ou}{\textlengthmark}{\textsubring{b}}{\textscy}{\textsubring{r}}} \dicPos{m}[6] \dicFlx{(‑s, ‑ar)}[15] \dicSynonym{kústur} \dicDirectTranslationCS{koště, smeták, metla}
\dicEntry[sór] \dicTerm{sór} \dicIPA{{s}{ou}{\textlengthmark}{\textsubring{r}}} \dicPos{v} \dicFlx{ind pf sg 1 pers} \dicLink{sverja}
\dicEntry[sóri] \dicTerm{sór|i} \dicIPA{{s}{ou}{\textlengthmark}{r}{\textsci}} \dicPos{m}[1] \dicFlx{(‑a)}[3] \dicFieldCat{med. } \dicDirectTranslationCS{lupénka, psoriáza}
\dicEntry[sórum] \dicTerm{sórum} \dicIPA{{s}{ou}{\textlengthmark}{r}{\textscy}{\textsubring{m}}} \dicPos{v} \dicFlx{ind pf pl 1 pers} \dicLink{sverja}
\dicEntry[sósa] \dicTerm{sós|a} \dicIPA{{s}{ou}{\textlengthmark}{s}{a}} \dicPos{f}[1] \dicFlx{(‑u, ‑ur)}[7] \dicDirectTranslationCS{omáčka, sos}
\dicEntry[sósíaldemókrati] \dicTerm{sósíal··demókrat|i} \dicIPA{{s}{ou}{\textlengthmark}{s}{i}{j}{a}{l}{\textsubring{d}}{\textepsilon}{m}{ou}{\r{g}}{r}{a}{\textsubring{d}}{\textsci}} \dicPos{m}[1] \dicFlx{(‑a, ‑ar)}[8] \dicDirectTranslationCS{sociální demokrat(ka)}
\dicEntry[sósíalismi] \dicTerm{sósíal··ism|i} \dicIPA{{s}{ou}{\textlengthmark}{s}{i}{j}{a}{l}{\textsci}{s}{m}{\textsci}} \dicPos{m}[1] \dicFlx{(‑a)}[3] \dicFieldCat{pol.} \dicDirectTranslationCS{socialismus}
\dicEntry[sósíalistaflokkur] \dicTerm{sósíalista··flokk|ur} \dicIPA{{s}{ou}{\textlengthmark}{s}{i}{j}{a}{l}{\textsci}{s}{\textsubring{d}}{a}{f}{l}{\textopeno}{h}{\r{g}}{\textscy}{\textsubring{r}}} \dicPos{m}[6] \dicFlx{(‑s, ‑ar)}[8] \dicFieldCat{pol.} \dicDirectTranslationCS{socialistická strana}
\dicEntry[sósíalisti] \dicTerm{sósíal··ist|i} \dicIPA{{s}{ou}{\textlengthmark}{s}{i}{j}{a}{l}{\textsci}{s}{\textsubring{d}}{\textsci}} \dicPos{m}[1] \dicFlx{(‑a, ‑ar)}[1] \dicFieldCat{pol.} \dicDirectTranslationCS{socialista, socialistka}
\dicEntry[sósíalískur] \dicTerm{sósíalískur} \dicIPA{{s}{ou}{\textlengthmark}{s}{i}{j}{a}{l}{i}{s}{\r{g}}{\textscy}{\textsubring{r}}} \dicPos{adj}[1]\dicFlx{}[-1] \dicDirectTranslationCS{socialistický}
\dicEntry[sósukanna] \dicTerm{sósu··|kanna} \dicIPA{{s}{ou}{\textlengthmark}{s}{\textscy}{k\smash{\textsuperscript{h}}}{a}{n}{a}} \dicPos{f}[1] \dicFlx{(‑könnu, ‑könnur)}[8] \dicDirectTranslationCS{omáčník}
\dicEntry[sót] \dicTerm{sót} \dicIPA{{s}{ou}{\textlengthmark}{\textsubring{d}}} \dicPos{n}[2] \dicFlx{(‑s)}[2] \dicDirectTranslationCS{saze, mour} \dicExampleIS{sót í reykháfi} \dicExampleCS{saze v~komíně}
\dicEntry[sótari] \dicTerm{sót··ar|i} \dicIPA{{s}{ou}{\textlengthmark}{\textsubring{d}}{a}{r}{\textsci}} \dicPos{m}[1] \dicFlx{(‑a, ‑ar)}[13] \dicDirectTranslationCS{kominík, kominice}
\dicEntry[sótrauður] \dicTerm{sót··rauður} \dicIPA{{s}{ou}{\textlengthmark}{\textsubring{d}}{r}{\oe i}{ð}{\textscy}{\textsubring{r}}} \dicPos{adj}[2]\dicFlx{}[-6] \dicDirectTranslationCS{rudý, (temně) červený} \dicExampleIS{sótrauður af bræði} \dicExampleCS{rudý vztekem}
\dicEntry[sótroðna] \dicTerm{sót··roðn|a} \dicIPA{{s}{ou}{\textlengthmark}{\textsubring{d}}{r}{\textopeno}{ð}{n}{a}} \dicPos{v}[1] \dicFlx{(‑aði)}[44] \dicDirectTranslationCS{zrudnout, zčervenat (vzteky ap.)} \dicExampleIS{sótroðna af reiði} \dicExampleCS{zrudnout vzteky}
\dicEntry[sótt] \dicTerm{sótt\smash{\textsuperscript{1}}} \dicIPA{{s}{ou}{h}{\textsubring{d}}} \dicPos{f}[7] \dicFlx{(‑ar, ‑ir)}[1] \textbf{1.} \dicLangCat{zast.} \dicSynonym{sjúkleiki} \dicDirectTranslationCS{onemocnění, nemoc, nákaza} \dicExampleIS{taka sótt} \dicExampleCS{onemocnět}  \textbf{2.} \dicSynonym{jóðsótt} \dicDirectTranslationCS{porodní bolest}
\dicEntry[sótt] \dicTerm{sótt\smash{\textsuperscript{2}}} \dicIPA{{s}{ou}{h}{\textsubring{d}}} \dicPos{v} \dicFlx{supin} \dicLink{sækja}
\dicEntry[sóttarsæng] \dicTerm{sóttar··sæng} \dicIPA{{s}{ou}{h}{\textsubring{d}}{a}{\textsubring{r}}{s}{a}{i}{\ng}{\r{g}}} \dicPos{f}[10] \dicFlx{(‑ur)}[2] \dicDirectTranslationCS{lůžko nemocného, nemocenské lůžko}
\dicEntry[sótthiti] \dicTerm{sótt··hit|i} \dicIPA{{s}{ou}{h}{\textsubring{d}}{h}{\textsci}{\textsubring{d}}{\textsci}} \dicPos{m}[1] \dicFlx{(‑a)}[3] \dicSynonym{hiti} \dicDirectTranslationCS{horečka, teplota}
\dicEntry[sótthreinsa] \dicTerm{sótt··hreins|a} \dicIPA{{s}{ou}{h}{\textsubring{d}}{\textsubring{r}}{ei}{n}{s}{a}} \dicPos{v}[1] \dicFlx{(‑aði)}[1] \dicFlx{acc} \dicDirectTranslationCS{(vy)dezinfikovat}
\dicEntry[sótthreinsaður] \dicTerm{sótt··hreins·|aður} \dicIPA{{s}{ou}{h}{\textsubring{d}}{\textsubring{r}}{ei}{n}{s}{a}{ð}{\textscy}{\textsubring{r}}} \dicPos{adj}[3] \dicFlx{(f ‑uð)}[4] \dicDirectTranslationCS{(vy)dezinfikovaný}
\dicEntry[sótthreinsun] \dicTerm{sótt··hreins|un} \dicIPA{{s}{ou}{h}{\textsubring{d}}{\textsubring{r}}{ei}{n}{s}{\textscy}{\textsubring{n}}} \dicPos{f}[7] \dicFlx{(‑unar)}[9] \dicDirectTranslationCS{dezinfekce, dezinfikování, sterilizace}
\dicEntry[sótthreinsunarefni] \textls[10]{\dicTerm{sótt·hreinsunar··efni} \dicIPA{{s}\-{ou}\-{h}\-{\textsubring{d}}\-{\textsubring{r}}\-{ei}\-{n}\-{s}\-{\textscy}\-{n}\-{a}\-{r}\-{\textepsilon}\-{\textsubring{b}}\-{n}\-{\textsci}} \dicPos{n}[2] \dicFlx{(‑s, ‑)}[14]} \dicDirectTranslationCS{dezinfekční látka\,/\addthin prostředek}
\dicEntry[sótti] \dicTerm{sótti} \dicIPA{{s}{ou}{h}{\textsubring{d}}{\textsci}} \dicPos{v} \dicFlx{ind pf sg 1 pers} \dicLink{sækja}
\dicEntry[sóttkví] \dicTerm{sótt··kví} \dicIPA{{s}{ou}{h}{\textsubring{d}}{k\smash{\textsuperscript{h}}}{v}{i}} \dicPos{f}[4] \dicFlx{(‑ar)}[7] \dicDirectTranslationCS{karanténa};  \dicPhraseIS{setja e‑n í sóttkví} \dicDirectTranslationCS{umístit (koho) do karantény}
\dicEntry[sóttnæmi] \dicTerm{sótt··næmi} \dicIPA{{s}{ou}{h}{\textsubring{d}}{n}{a}{i}{m}{\textsci}} \dicPos{n}[2] \dicFlx{(‑s)}[20] \dicSynonym{smitnæmi} \dicDirectTranslationCS{nakažlivost, infekčnost}
\dicEntry[sóttnæmur] \dicTerm{sótt··næmur} \dicIPA{{s}{ou}{h}{\textsubring{d}}{n}{a}{i}{m}{\textscy}{\textsubring{r}}} \dicPos{adj}[1]\dicFlx{}[-1] \dicSynonym{smitnæmur} \dicDirectTranslationCS{nakažlivý, infekční}
\dicEntry[sóttum] \dicTerm{sóttum} \dicIPA{{s}{ou}{h}{\textsubring{d}}{\textscy}{\textsubring{m}}} \dicPos{v} \dicFlx{ind pf pl 1 pers} \dicLink{sækja}
\dicEntry[sóttverjandi] \dicTerm{sótt··verj·andi} \dicIPA{{s}{ou}{h}{\textsubring{d}}{v}{\textepsilon}{r}{j}{a}{n}{\textsubring{d}}{\textsci}} \dicPos{adj}[13] \dicFlx{indecl}[1] \dicDirectTranslationCS{antiseptický, dezinfekční}
\dicEntry[sóttvörn] \dicTerm{sótt··|vörn} \dicIPA{{s}{ou}{h}{\textsubring{d}}{v}{\oe}{r}{\textsubring{d}}{\textsubring{n}}} \dicPos{f}[7] \dicFlx{(‑varnar, ‑varnir)}[16] \dicDirectTranslationCS{antisepse, prevence infekce\,/\addthin nákazy} \dicIndirectTranslationCS{(zvláště v~pl)}
\dicEntry[sótugur] \dicTerm{sót··ugur} \dicIPA{{s}{ou}{\textlengthmark}{\textsubring{d}}{\textscy}{\textbabygamma}{\textscy}{\textsubring{r}}} \dicPos{adj}[1]\dicFlx{}[-8] \dicDirectTranslationCS{umazaný\,/\addthin černý od sazí}
\dicEntry[sótvondur] \dicTerm{sót··vondur} \dicIPA{{s}{ou}{\textlengthmark}{\textsubring{d}}{v}{\textopeno}{n}{\textsubring{d}}{\textscy}{\textsubring{r}}} \dicPos{adj}[11]\dicFlx{}[-6] \dicSynonym{bálreiður} \dicDirectTranslationCS{zuřivý, vzteklý, rozzuřený}
\dicEntry[sóun] \dicTerm{só|un} \dicIPA{{s}{ou}{\textlengthmark}{\textscy}{\textsubring{n}}} \dicPos{f}[7] \dicFlx{(‑unar)}[9] \dicDirectTranslationCS{plýtvání, mrhání} \dicExampleIS{sóun á fé} \dicExampleCS{plýtvání penězi}
\dicEntry[spað] \dicTerm{spað} \dicIPA{{s}{\textsubring{b}}{a}{\textlengthmark}{\texttheta}} \dicPos{n}[2] \dicFlx{(‑s)}[2] \dicPhraseIS{brytja e‑n í spað} \dicLangCat{přen.} \dicDirectTranslationCS{nenechat na (kom) nit suchou}
\dicEntry[spaði] \dicTerm{spað|i} \dicIPA{{s}{\textsubring{b}}{a}{\textlengthmark}{ð}{\textsci}} \dicPos{m}[1] \dicFlx{(‑a, ‑ar)}[8] \textbf{1.} \dicSynonym{skófla\smash{\textsuperscript{1}}} \dicDirectTranslationCS{rýč}  \textbf{2.} \dicSynonym{skeið\smash{\textsuperscript{1}}} \dicDirectTranslationCS{obracečka} \dicExampleIS{fiskspaði} \dicExampleCS{obracečka na ryby}  \textbf{3.} \dicSynonym*{litur í spilum} \dicDirectTranslationCS{piky, dýky} \dicIndirectTranslationCS{(barva v~kartách)}  \textbf{4.} \dicFieldCat{sport.} \dicDirectTranslationCS{raketa, pálka}  \textbf{5.} \dicDirectTranslationCS{lopata, lopatka (vrtule ap.)};  \dicIdiom{spaði}{ \dicPhraseIS{halda (vel) á spöðunum}} \dicLangCat{přen.} \dicDirectTranslationCS{vynasnažit se, postarat se, přičinit se}; { \dicPhraseIS{taka í spaðann á e‑m}} \dicLangCat{hovor.} \dicDirectTranslationCS{podat si s~(kým) ruku}
\dicEntry[spaghettí] \dicTerm{spaghettí} \dicIPA{{s}{\textsubring{b}}{a}{\textlengthmark}{\textbabygamma}{\textepsilon}{h}{\textsubring{d}}{i}} \dicPos{n}[2] \dicFlx{(‑s)}[31] \dicDirectTranslationCS{špagety}
\dicEntry[spaklegur] \dicTerm{spak··legur} \dicIPA{{s}{\textsubring{b}}{a}{\textlengthmark}{\r{g}}{l}{\textepsilon}{\textbabygamma}{\textscy}{\textsubring{r}}} \dicPos{adj}[1]\dicFlx{}[-8] \dicSynonym{hygginn} \dicDirectTranslationCS{moudrý, rozumný, inteligentní}
\dicEntry[spakmæli] \dicTerm{spak··mæli} \dicIPA{{s}{\textsubring{b}}{a}{\textlengthmark}{\r{g}}{m}{a}{i}{l}{\textsci}} \dicPos{n}[2] \dicFlx{(‑s, ‑)}[14] \dicDirectTranslationCS{moudro, aforismus}
\dicEntry[spakur] \dicTerm{spakur} \dicIPA{{s}{\textsubring{b}}{a}{\textlengthmark}{\r{g}}{\textscy}{\textsubring{r}}} \dicPos{adj}[1] \dicFlx{(f spök)}[2] \textbf{1.} \dicSynonym{kyrrlátur} \dicDirectTranslationCS{klidný, krotký} \dicExampleIS{spakur hestur} \dicExampleCS{krotký kůň}  \textbf{2.} \dicSynonym{vitur} \dicDirectTranslationCS{moudrý, prozíravý}
\dicEntry[spalar] \dicTerm{spalar} \dicIPA{{s}{\textsubring{b}}{a}{\textlengthmark}{l}{a}{\textsubring{r}}} \dicPos{m} \dicFlx{sg gen} \dicLink{spölur}
\dicEntry[span] \dicTerm{span} \dicIPA{{s}{\textsubring{b}}{a}{\textlengthmark}{\textsubring{n}}} \dicPos{n}[2] \dicFlx{(‑s)}[2] \textbf{1.} \dicSynonym*{hröð ferð} \dicDirectTranslationCS{shon, spěch};  \dicPhraseIS{vera á spani} \dicDirectTranslationCS{být ve shonu, spěchat};  \dicPhraseIS{það er span á e‑m} \dicFlx{impers} \dicDirectTranslationCS{(kdo) pospíchá}  \textbf{2.} \dicFieldCat{elek.} \dicDirectTranslationCS{indukce}
\dicEntry[spana] \dicTerm{span|a} \dicIPA{{s}{\textsubring{b}}{a}{\textlengthmark}{n}{a}} \dicPos{v}[1] \dicFlx{(‑aði)}[13] \dicFlx{acc} \textbf{1.} \dicSynonym{æsa} \dicDirectTranslationCS{podnítit, podněcovat, (vy)provokovat};  \dicPhraseIS{spana e‑n upp í e‑ð} \dicDirectTranslationCS{podněcovat (koho) k~(čemu), vyprovokovat (koho) k~(čemu)} \dicExampleIS{spana hann upp í að hrekkja kennarann} \dicExampleCS{podněcovat ho ke zlobení učitele}  \textbf{2.} \dicSynonym*{flýta sér} \dicDirectTranslationCS{hnát se, letět, uhánět}  \textbf{3.} \dicFieldCat{elek.} \dicDirectTranslationCS{indukovat}
\dicEntry[spandera] \dicTerm{spander|a}\dicTerm{, spandéra} \dicIPA{{s}\-{\textsubring{b}}\-{a}\-{n}\-{\textsubring{d}}\-{\textepsilon}\-{r}\-{a}\-} \dicPos{v}[1] \dicFlx{(‑aði)}[1] \dicFlx{dat} \dicLangCat{hovor.} \dicSynonym{sól\-unda} \dicDirectTranslationCS{utratit, promrhat, rozfofrovat (peníze ap.)} \dicExampleIS{span\-dera stórfé} \dicExampleCS{utratit těžké peníze}
\dicEntry[spandéra] \dicTerm{spandér|a} \dicIPA{{s}{\textsubring{b}}{a}{n}{\textsubring{d}}{j}{\textepsilon}{r}{a}} \dicPos{v}[1] \dicFlx{(‑aði)}[1] \dicLink{spandera}
\dicEntry[spangar] \dicTerm{spangar} \dicIPA{{s}{\textsubring{b}}{au}{\ng}{\r{g}}{a}{\textsubring{r}}{}} \dicPos{f} \dicFlx{sg gen} \dicLink{spöng}
\dicEntry[spangól] \dicTerm{span··gól} \dicIPA{{s}{\textsubring{b}}{a}{n}{\r{g}}{ou}{\textsubring{l}}} \dicPos{n}[2] \dicFlx{(‑s)}[2] \dicDirectTranslationCS{(za)vytí (psa ap.)}
\dicEntry[spangóla] \dicTerm{span··gól|a} \dicIPA{{s}{\textsubring{b}}{a}{n}{\r{g}}{ou}{l}{a}} \dicPos{v}[1] \dicFlx{(‑aði)}[1] \dicDirectTranslationCS{(za)výt (o~psovi ap.)}
\dicEntry[spann] \dicTerm{spann} \dicIPA{{s}{\textsubring{b}}{a}{n}{\textlengthmark}} \dicPos{v} \dicFlx{ind pf sg 1 pers} \dicLink{spinna}
\dicEntry[spanna] \dicTerm{spann|a} \dicIPA{{s}{\textsubring{b}}{a}{n}{\textlengthmark}{a}} \dicPos{v}[1] \dicFlx{(‑aði)}[13] \dicDirectTranslationCS{zahrnovat, pokrývat (časový úsek ap.)} \dicExampleIS{Bókin spannar hálfa öld.} \dicExampleCS{Knížka pokrývá polovinu století.}
\dicEntry[spannar] \dicTerm{spannar} \dicIPA{{s}{\textsubring{b}}{a}{n}{\textlengthmark}{a}{\textsubring{r}}} \dicPos{f} \dicFlx{sg gen} \dicLink{spönn}
\dicEntry[spanskgræna] \dicTerm{spansk··græn|a} \dicIPA{{s}{\textsubring{b}}{a}{n}{s}{\r{g}}{r}{a}{i}{n}{a}} \dicPos{f}[1] \dicFlx{(‑u)}[5] \dicDirectTranslationCS{patina, měděnka}
\dicEntry[spar] \dicTerm{spar} \dicIPA{{s}{\textsubring{b}}{a}{\textlengthmark}{\textsubring{r}}} \dicPos{adj}[5] \dicFlx{(f spör)}[9] \dicSynonym{sparsamur} \dicDirectTranslationCS{šetrný, spořivý, hospodárný};  \dicPhraseIS{vera spar á e‑ð} \dicDirectTranslationCS{šetřit (čím) (jídlem, pochvalami ap.)}
\dicEntry[spara] \dicTerm{spar|a} \dicsymFrequent\  \dicIPA{{s}{\textsubring{b}}{a}{\textlengthmark}{r}{a}} \dicPos{v}[1] \dicFlx{(‑aði)}[13] \dicFlx{dat + acc} \textbf{1.} \dicSynonym*{eyða litlu} \dicDirectTranslationCS{(u)šetřit, (u)spořit (peníze ap.)} \dicExampleIS{spara 1 200 krónur á kaupinu} \dicExampleCS{ušetřit 1 200 korun za nákup}  \textbf{2.} \dicDirectTranslationCS{(u)šetřit, uchovat (šaty na svatbu ap.)} \dicExampleIS{spara kjólinn} \dicExampleCS{šetřit šaty}  \textbf{3.} \dicSynonym{hlífa} \dicDirectTranslationCS{ušetřit, zachovat (život ap.)} \dicExampleIS{spara líf e‑rs} \dicExampleCS{ušetřit (čí) život}  \textbf{4.} \dicPhraseIS{spara sér e‑ð} \dicDirectTranslationCS{ušetřit si (co) (zbytečnou cestu ap.)};  \dicIdiom{sparast}{ \dicPhraseIS{sparast}} \dicFlx{refl} \dicDirectTranslationCS{ušetřit se, uspořit se}
\dicEntry[sparð] \dicTerm{sparð} \dicIPA{{s}{\textsubring{b}}{a}{r}{\texttheta}} \dicPos{n}[2] \dicFlx{(‑s, spörð)}[8] \dicDirectTranslationCS{ovčí bobek\,/\addthin trus}
\dicEntry[sparðatíningur] \dicTerm{sparða··tín·ing|ur} \dicIPA{{s}{\textsubring{b}}{a}{r}{ð}{a}{t\smash{\textsuperscript{h}}}{i}{n}{i}{\ng}{\r{g}}{\textscy}{\textsubring{r}}} \dicPos{m}[6] \dicFlx{(‑s)}[9] \dicDirectTranslationCS{hnidopišství, nimrání se}
\dicEntry[spari] \dicTerm{spari\smash{\textsuperscript{1}}} \dicIPA{{s}{\textsubring{b}}{a}{\textlengthmark}{r}{\textsci}} \dicPos{adv} \dicDirectTranslationCS{svátečně, ke sváteční příležitosti}
\dicEntry[spari] \dicTerm{spari-\smash{\textsuperscript{2}}} \dicIPA{{s}{\textsubring{b}}{a}{\textlengthmark}{r}{\textsci}} \dicPos{predp} \dicDirectTranslationCS{sváteční, slavnostní} \dicAntonym{hversdags\smash{\textsuperscript{2}}}
\dicEntry[sparibaukur] \dicTerm{spari··bauk|ur} \dicIPA{{s}{\textsubring{b}}{a}{\textlengthmark}{r}{\textsci}{\textsubring{b}}{\oe i}{\r{g}}{\textscy}{\textsubring{r}}} \dicPos{m}[6] \dicFlx{(‑s, ‑ar)}[22] \dicDirectTranslationCS{pokladnička, kasička, prasátko}
\dicEntry[sparibúinn] \dicTerm{spari··búinn} \dicIPA{{s}{\textsubring{b}}{a}{\textlengthmark}{r}{\textsci}{\textsubring{b}}{u}{\textsci}{\textsubring{n}}} \dicPos{adj}[6]\dicFlx{}[-6] \dicDirectTranslationCS{svátečně oblečený, vystrojený}
\dicEntry[sparifé] \dicTerm{spari··|fé} \dicIPA{{s}{\textsubring{b}}{a}{\textlengthmark}{r}{\textsci}{f}{j}{\textepsilon}} \dicPos{n}[3] \dicFlx{(‑fjár)}[5] \dicDirectTranslationCS{úspory, našetřené peníze}
\dicEntry[spariföt] \dicTerm{spari··föt} \dicIPA{{s}{\textsubring{b}}{a}{\textlengthmark}{r}{\textsci}{f}{\oe}{\textsubring{d}}} \dicPos{n}[2] \dicFlx{pl}[9] \dicDirectTranslationCS{sváteční\,/\addthin slavnostní oblečení}
\dicEntry[spariklæddur] \dicTerm{spari··klæddur} \dicIPA{{s}{\textsubring{b}}{a}{\textlengthmark}{r}{\textsci}{k\smash{\textsuperscript{h}}}{l}{a}{i}{\textsubring{d}}{\textscy}{\textsubring{r}}} \dicPos{adj}[2]\dicFlx{}[-18] \dicDirectTranslationCS{svátečně\,/\addthin slavnostně oblečený}
\dicEntry[spariklæðnaður] \dicTerm{spari··klæð·nað|ur} \dicIPA{{s}{\textsubring{b}}{a}{\textlengthmark}{r}{\textsci}{k\smash{\textsuperscript{h}}}{l}{a}{i}{ð}{n}{a}{ð}{\textscy}{\textsubring{r}}} \dicPos{m}[10] \dicFlx{(‑ar)}[9] \dicDirectTranslationCS{sváteční\,/\addthin slavnostní oblečení}
\begin{xtolerant}{}{1pt}
\dicEntry[sparimerki] \dicTerm{spari··merki}\addthinS\dicIPA{{s}{\textsubring{b}}{a}{\textlengthmark}{r}{\textsci}{m}{\textepsilon}{\textsubring{r}}{\r{\textObardotlessj}}{\textsci}}\addthinS\dicPos{n}[2] \dicFlx{(‑s, ‑)}[16] \dicFieldCat{hist.} \dicDirectTranslationCS{spořitelní známka}
\end{xtolerant}
\dicEntry[sparireikningur] \dicTerm{spari··reikn·ing|ur} \dicIPA{{s}{\textsubring{b}}{a}{\textlengthmark}{r}{\textsci}{r}{ei}{h}{\r{g}}{n}{i}{\ng}{\r{g}}{\textscy}{\textsubring{r}}} \dicPos{m}[6] \dicFlx{(‑s, ‑ar)}[8] \dicFieldCat{ekon.} \dicDirectTranslationCS{spořicí\,/\addthin spořitelní účet}
\dicEntry[sparisjóðsbók] \dicTerm{spari·sjóðs··|bók} \dicIPA{{s}{\textsubring{b}}{a}{\textlengthmark}{r}{\textsci}{s}{j}{ou}{ð}{s}{\textsubring{b}}{ou}{\r{g}}} \dicPos{f}[8] \dicFlx{(‑bókar, ‑bækur)}[5] \dicFieldCat{ekon.} \dicDirectTranslationCS{spořitelní knížka}
\dicEntry[sparisjóðsreikningur] \dicTerm{spari·sjóðs··reikn·ing|ur} \dicIPA{{s}\-{\textsubring{b}}\-{a}\-{\textlengthmark}\-{r}\-{\textsci}\-{s}\-{j}\-{ou}\-{ð}\-{s}\-{r}\-{ei}\-{h}\-{\r{g}}\-{n}\-{i}\-{\ng}\-{\r{g}}\-{\textscy}\-{\textsubring{r}}\-} \dicPos{m}[6] \dicFlx{(‑s, ‑ar)}[8] \dicFieldCat{ekon.} \dicDirectTranslationCS{spořicí\,/\addthin spořitelní účet}
\dicEntry[sparisjóður] \dicTerm{spari··sjóð|ur} \dicIPA{{s}{\textsubring{b}}{a}{\textlengthmark}{r}{\textsci}{s}{j}{ou}{ð}{\textscy}{\textsubring{r}}} \dicPos{m}[9] \dicFlx{(‑s, ‑ir)}[6] \dicFieldCat{ekon.} \dicDirectTranslationCS{spořitelna}
\dicEntry[spariskírteini] \dicTerm{spari··skír·teini} \dicIPA{{s}{\textsubring{b}}{a}{\textlengthmark}{r}{\textsci}{s}{\r{\textObardotlessj}}{i}{\textsubring{r}}{\textsubring{d}}{ei}{n}{\textsci}} \dicPos{n}[2] \dicFlx{(‑s, ‑)}[14] \dicFieldCat{ekon.} \dicDirectTranslationCS{depozitní certifikát};  \dicPhraseIS{spariskírteini ríkissjóðs} \dicFieldCat{ekon.} \dicDirectTranslationCS{státní dluhopis}
\dicEntry[spark] \dicTerm{spark} \dicIPA{{s}{\textsubring{b}}{a}{\textsubring{r}}{\r{g}}} \dicPos{n}[2] \dicFlx{(‑s, spörk)}[8] \dicDirectTranslationCS{kop, kopnutí};  \dicPhraseIS{fá spark} \dicLangCat{přen.} \dicDirectTranslationCS{dostat padáka, být vykopnutý (z~práce ap.)};  \dicPhraseIS{gefa e‑m spark í rassinn} \dicDirectTranslationCS{kopnout (koho) do zadku}
\dicEntry[sparka] \dicTerm{spark|a} \dicsymFrequent\  \dicIPA{{s}{\textsubring{b}}{a}{\textsubring{r}}{\r{g}}{a}} \dicPos{v}[1] \dicFlx{(‑aði)}[13] \dicFlx{dat} \textbf{1.} \dicDirectTranslationCS{(na)kopnout, kopat} \dicExampleIS{sparka boltanum í mark} \dicExampleCS{kopnout míč do brány};  \dicPhraseIS{sparka í e‑n} \dicDirectTranslationCS{kopnout (koho)}  \textbf{2.} \dicDirectTranslationCS{vyhodit, vykopnout, vylít (z~práce ap.)}
\dicEntry[sparnaðarráðstöfun] \dicTerm{sparnaðar··ráð·|stöfun} \dicIPA{{s}\-{\textsubring{b}}\-{a}\-{r}\-{\textsubring{d}}\-{n}\-{a}\-{ð}\-{a}\-{r}\-{au}\-{ð}\-{s}\-{\textsubring{d}}\-{\oe}\-{v}\-{\textscy}\-{\textsubring{n}}\-} \dicPos{f}[7] \dicFlx{(‑stöf\-unar, ‑stafanir)}[11] \dicDirectTranslationCS{úsporné opatření}
\dicEntry[sparnaður] \dicTerm{spar··nað|ur} \dicIPA{{s}{\textsubring{b}}{a}{r}{\textsubring{d}}{n}{a}{ð}{\textscy}{\textsubring{r}}} \dicPos{m}[10] \dicFlx{(‑ar)}[9] \textbf{1.} \dicSynonym{sparifé} \dicDirectTranslationCS{úspory}  \textbf{2.} \dicSynonym{sparsemi} \dicDirectTranslationCS{šetrnost, spořivost}
\dicEntry[sparneytinn] \dicTerm{spar··neytinn} \dicIPA{{s}{\textsubring{b}}{a}{r}{n}{ei}{\textsubring{d}}{\textsci}{\textsubring{n}}} \dicPos{adj}[6]\dicFlx{}[-2] \dicDirectTranslationCS{úsporný, hospodárný, ekonomický} \dicExampleIS{sparneytinn bíll} \dicExampleCS{úsporné auto}
\dicEntry[sparneytni] \dicTerm{spar··neytn|i} \dicIPA{{s}{\textsubring{b}}{a}{r}{n}{ei}{h}{\textsubring{d}}{n}{\textsci}} \dicPos{f}[3] \dicFlx{(‑i)}[3] \dicDirectTranslationCS{úspornost, hospodárnost, ekonomičnost}
\dicEntry[sparrhaukur] \dicTerm{sparr··hauk|ur} \dicIPA{{s}{\textsubring{b}}{a}{r}{\textlengthmark}{h}{\oe i}{\r{g}}{\textscy}{\textsubring{r}}} \dicPos{m}[6] \dicFlx{(‑s, ‑ar)}[22] \dicFieldCat{zool.} \dicDirectTranslationCS{krahujec obecný} \textit{(l.~{\textLA{Accipiter nisus}})}  \dicsymPhoto\ 
\dicFigure{ds_image_sparrhaukur_0_2.jpg}{Sparrhaukur}{Sparrhaukur - Meneer Zjeroen, CC BY 2.0}
\dicEntry[sparsamur] \dicTerm{spar··|samur} \dicIPA{{s}{\textsubring{b}}{a}{\textsubring{r}}{s}{a}{m}{\textscy}{\textsubring{r}}} \dicPos{adj}[1] \dicFlx{(f ‑söm)}[2] \dicSynonym{nýtinn} \dicDirectTranslationCS{šetrný, spořivý, hospodárný}
\dicEntry[sparsemi] \dicTerm{spar··sem|i} \dicIPA{{s}{\textsubring{b}}{a}{\textsubring{r}}{s}{\textepsilon}{m}{\textsci}} \dicPos{f}[2] \dicFlx{(‑i)}[2] \dicDirectTranslationCS{šetrnost, spořivost, hospodárnost}
\dicEntry[sparsl] \dicTerm{sparsl}\dicTerm{, spartl} \dicIPA{{s}\-{\textsubring{b}}\-{a}\-{\textsubring{r}}\-{s}\-{\textsubring{d}}\-{\textsubring{l}}\-} \dicPos{n}[2] \dicFlx{(‑s)}[2] \dicDirectTranslationCS{tmel, kyt}
\dicEntry[sparsla] \dicTerm{sparsl|a}\dicTerm{, spartla} \dicIPA{{s}\-{\textsubring{b}}\-{a}\-{\textsubring{r}}\-{s}\-{\textsubring{d}}\-{l}\-{a}\-} \dicPos{v}[1] \dicFlx{(‑aði)}[13] \dicFlx{acc} \dicDirectTranslationCS{(za)tmelit, (za)kytovat}
\dicEntry[spartl] \dicTerm{spartl} \dicIPA{{s}{\textsubring{b}}{a}{\textsubring{r}}{\textsubring{d}}{\textsubring{l}}} \dicPos{n}[2] \dicFlx{(‑s)}[2] \dicLink{sparsl}
\dicEntry[spartla] \dicTerm{spartl|a} \dicIPA{{s}{\textsubring{b}}{a}{\textsubring{r}}{\textsubring{d}}{l}{a}} \dicPos{v}[1] \dicFlx{(‑aði)}[13] \dicLink{sparsla}
\dicEntry[spaug] \dicTerm{spaug} \dicIPA{{s}{\textsubring{b}}{\oe i}{\textlengthmark}{x}} \dicPos{n}[2] \dicFlx{(‑s)}[2] \dicDirectTranslationCS{žert, vtip, šprým, anekdota};  \dicPhraseIS{það er ekkert spaug} \dicLangCat{přen.} \dicDirectTranslationCS{není to žádná legrace (nalézt dobrou knížku ap.)}
\dicEntry[spauga] \dicTerm{spaug|a} \dicIPA{{s}{\textsubring{b}}{\oe i}{\textlengthmark}{\textbabygamma}{a}} \dicPos{v}[1] \dicFlx{(‑aði)}[1] \dicDirectTranslationCS{žertovat, vtipkovat, šprýmovat} \dicExampleIS{spauga um e‑ð við e‑n} \dicExampleCS{žertovat o~(čem) s~(kým)}
\dicEntry[spaugari] \dicTerm{spaug··ar|i} \dicIPA{{s}{\textsubring{b}}{\oe i}{\textlengthmark}{\textbabygamma}{a}{r}{\textsci}} \dicPos{m}[1] \dicFlx{(‑a, ‑ar)}[13] \dicDirectTranslationCS{vtipálek, srandista, srandistka, humorista, humoristka, šprýmař(ka)}
\dicEntry[spaugilegur] \dicTerm{spaugi··legur} \dicIPA{{s}{\textsubring{b}}{\oe i}{j}{\textlengthmark}{\textsci}{l}{\textepsilon}{\textbabygamma}{\textscy}{\textsubring{r}}} \dicPos{adj}[1]\dicFlx{}[-8] \dicSynonym{kátlegur} \dicDirectTranslationCS{žertovný, vtipný, legrační, komický} \dicExampleIS{spaugileg saga} \dicExampleCS{žertovný příběh}
\dicEntry[spaugsamur] \dicTerm{spaug··|samur} \dicIPA{{s}{\textsubring{b}}{\oe i}{x}{s}{a}{m}{\textscy}{\textsubring{r}}} \dicPos{adj}[1] \dicFlx{(f ‑söm)}[2] \dicSynonym{gamansamur} \dicDirectTranslationCS{legrační, žertovný, komický, humorný}
\dicEntry[spá] \dicTerm{spá\smash{\textsuperscript{1}}} \dicIPA{{s}{\textsubring{b}}{au}{\textlengthmark}} \dicPos{f}[4] \dicFlx{(‑r, ‑r)}[18] \textbf{1.} \dicSynonym{spádómur} \dicDirectTranslationCS{věštba, proroctví} \dicExampleIS{treysta ekki spánni} \dicExampleCS{nevěřit proroctví}  \textbf{2.} \dicDirectTranslationCS{předpověď, prognóza (počasí ap.)} \dicExampleIS{veðurspá} \dicExampleCS{předpověď počasí}
\dicEntry[spá] \dicTerm{spá\smash{\textsuperscript{2}}} \dicIPA{{s}{\textsubring{b}}{au}{\textlengthmark}} \dicPos{v}[2] \dicFlx{(‑ði, ‑ð)}[116] \dicFlx{(dat +) dat} \textbf{1.} \dicDirectTranslationCS{věštit, prorokovat};  \dicPhraseIS{spá fyrir e‑m} \dicDirectTranslationCS{věštit (komu) budoucnost};  \dicPhraseIS{spá í lófa} \dicDirectTranslationCS{číst\,/\addthin věštit z~ruky};  \dicPhraseIS{spá í spil} \dicDirectTranslationCS{věštit z~karet}  \textbf{2.} \dicDirectTranslationCS{předpovědět, předvídat, prognózovat} \dicExampleIS{spá sólskini} \dicExampleCS{předvídat slunečné počasí};  \dicPhraseIS{spá fyrir um e‑ð} \dicDirectTranslationCS{předpovědět (co), vypočítat (co)}  \textbf{3.} \dicPhraseIS{spá í e‑ð} \dicDirectTranslationCS{přemýšlet o~(čem), uvažovat o~(čem)}
\dicEntry[spádómur] \dicTerm{spá··dóm|ur} \dicIPA{{s}{\textsubring{b}}{au}{\textlengthmark}{\textsubring{d}}{ou}{m}{\textscy}{\textsubring{r}}} \dicPos{m}[6] \dicFlx{(‑s, ‑ar)}[10] \dicDirectTranslationCS{věštba, proroctví} \dicExampleIS{spádómur um framtíðina} \dicExampleCS{věštba o~budoucnosti}
\begin{xtolerant}{}{1pt}
\dicEntry[spákaupmaður] \dicTerm{spá··kaup·\hspace{-0.06667em}|maður}\addthinS\dicIPA{{s}{\textsubring{b}}{au}{\textlengthmark}{k\smash{\textsuperscript{h}}}{\oe i}{\textsubring{b}}{m}{a}{ð}{\textscy}{\textsubring{r}}}\addthinS\dicPos{m}[13]\addthinS\dicFlx{(‑manns, ‑menn)}[2] \dicFieldCat{ekon.} \dicDirectTranslationCS{spekulant(ka)}
\end{xtolerant}
\dicEntry[spákaupmennska] \dicTerm{spá··kaup·mennsk|a} \dicIPA{{s}{\textsubring{b}}{au}{\textlengthmark}{k\smash{\textsuperscript{h}}}{\oe i}{\textsubring{b}}{m}{\textepsilon}{n}{s}{\r{g}}{a}} \dicPos{f}[1] \dicFlx{(‑u)}[5] \dicFieldCat{ekon.} \dicDirectTranslationCS{spekulace, spekulování}
\dicEntry[spákona] \dicTerm{spá··kon|a} \dicIPA{{s}{\textsubring{b}}{au}{\textlengthmark}{k\smash{\textsuperscript{h}}}{\textopeno}{n}{a}} \dicPos{f}[1] \dicFlx{(‑u, ‑ur)}[27] \dicDirectTranslationCS{věštkyně, prorokyně, jasnovidka}
\dicEntry[spámaður] \dicTerm{spá··|maður} \dicIPA{{s}{\textsubring{b}}{au}{\textlengthmark}{m}{a}{ð}{\textscy}{\textsubring{r}}} \dicPos{m}[13] \dicFlx{(‑manns, ‑menn)}[2] \dicDirectTranslationCS{věštec, prorok, jasnovidec};  \dicProverb\  \dicPhraseIS{Enginn er spámaður í sínu föðurlandi.} \dicLangCat{přís.} \dicDirectTranslationCS{Nikdo není doma prorokem.}
\dicEntry[Spánn] \dicTerm{Spán|n} \dicsymFrequent\  \dicIPA{{s}{\textsubring{b}}{au}{\textsubring{d}}{\textsubring{n}}} \dicPos{m}[10] \dicFlx{(‑ar)}[46] \dicFieldCat{geog.} \dicDirectTranslationCS{Španělsko} \dicExampleIS{fljúga til Spánar} \dicExampleCS{letět do Španělska}
\dicEntry[spánnýr] \dicTerm{spán··nýr} \dicIPA{{s}{\textsubring{b}}{au}{n}{\textlengthmark}{i}{\textsubring{r}}} \dicPos{adj}[4]\dicFlx{}[-8] \dicDirectTranslationCS{úplně\,/\addthin zbrusu nový}
\dicEntry[spánskur] \dicTerm{spánskur} \dicIPA{{s}{\textsubring{b}}{au}{n}{s}{\r{g}}{\textscy}{\textsubring{r}}} \dicPos{adj}[1]\dicFlx{}[-1] \dicSynonym{spænskur} \dicDirectTranslationCS{španělský};  \dicPhraseIS{e‑að kemur e‑m spánskt fyrir sjónir} \dicLangCat{přen.} \dicDirectTranslationCS{(co) je pro (koho) španělská vesnice}
\dicEntry[Spánverji] \dicTerm{Spán··verj|i} \dicIPA{{s}{\textsubring{b}}{au}{n}{v}{\textepsilon}{r}{j}{\textsci}} \dicPos{m}[1] \dicFlx{(‑a, ‑ar)}[1] \dicDirectTranslationCS{Španěl(ka)}
\dicEntry[spássera] \dicTerm{spásser|a}\dicTerm{, spásséra} \dicIPA{{s}\-{\textsubring{b}}\-{au}\-{s}\-{\textlengthmark}\-{\textepsilon}\-{r}\-{a}\-} \dicPos{v}[1] \dicFlx{(‑aði)}[1] \dicSynonym*{labba um} \dicDirectTranslationCS{procházet se, špacírovat (si)} \dicExampleIS{spássera um göturnar} \dicExampleCS{procházet se po ulicích}
\dicEntry[spássertúr] \dicTerm{spásser··túr} \dicIPA{{s}{\textsubring{b}}{au}{s}{\textlengthmark}{\textepsilon}{\textsubring{r}}{t\smash{\textsuperscript{h}}}{u}{\textsubring{r}}} \dicPos{m}[4] \dicFlx{(‑s, ‑ar)}[14] \dicLink{spássitúr}
\dicEntry[spásséra] \dicTerm{spássér|a} \dicIPA{{s}{\textsubring{b}}{au}{s}{\textlengthmark}{j}{\textepsilon}{r}{a}} \dicPos{v}[1] \dicFlx{(‑aði)}[1] \dicLink{spássera}
\dicEntry[spássértúr] \dicTerm{spássér··túr} \dicIPA{{s}{\textsubring{b}}{au}{s}{j}{\textepsilon}{\textsubring{r}}{t\smash{\textsuperscript{h}}}{u}{\textsubring{r}}} \dicPos{m}[4] \dicFlx{(‑s, ‑ar)}[14] \dicLink{spássitúr}
\dicEntry[spássitúr] \dicTerm{spássi··túr}\dicTerm{, spássertúr}\dicTerm{,  spássértúr} \dicIPA{{s}\-{\textsubring{b}}\-{au}\-{s}\-{\textlengthmark}\-{\textsci}\-{t\smash{\textsuperscript{h}}}\-{u}\-{\textsubring{r}}\-} \dicPos{m}[4] \dicFlx{(‑s, ‑ar)}[14] \dicLangCat{hovor.} \dicSynonym{gönguferð} \dicDirectTranslationCS{procházka, špacír}
\dicEntry[spássía] \dicTerm{spássí|a} \dicIPA{{s}{\textsubring{b}}{au}{s}{\textlengthmark}{i}{j}{a}} \dicPos{f}[1] \dicFlx{(‑u, ‑ur)}[7] \dicDirectTranslationCS{okraj papíru\,/\addthin stránky} \dicExampleIS{skrifa athugasemdir á spássíuna} \dicExampleCS{psát si poznámky na okraj papíru}
\dicEntry[spássíugrein] \dicTerm{spássíu··grein} \dicIPA{{s}{\textsubring{b}}{au}{s}{\textlengthmark}{i}{j}{\textscy}{\r{g}}{r}{ei}{\textsubring{n}}} \dicPos{f}[4] \dicFlx{(‑ar, ‑ar)}[1] \dicDirectTranslationCS{poznámka na okraji stránky}
\dicEntry[spegill] \dicTerm{speg|ill} \dicsymFrequent\  \dicIPA{{s}{\textsubring{b}}{ei}{\textsci}{\textsubring{d}}{\textsubring{l}}} \dicPos{m}[6] \dicFlx{(‑ils, ‑lar)}[35] \dicDirectTranslationCS{zrcadlo} \dicExampleIS{sjá sig í spegli} \dicExampleCS{vidět se v~zrcadle};  \dicProverb\  \dicPhraseIS{Augun eru spegill sálarinnar.} \dicLangCat{přís.} \dicDirectTranslationCS{Oči jsou zrcadlem duše.}
\dicEntry[spegilmynd] \dicTerm{spegil··mynd} \dicIPA{{s}{\textsubring{b}}{ei}{\textsci}{l}{m}{\textsci}{n}{\textsubring{d}}} \dicPos{f}[7] \dicFlx{(‑ar, ‑ir)}[1] \textbf{1.} \dicDirectTranslationCS{(zrcadlový) odraz\,/\addthin obraz, zrcadlení}  \textbf{2.} \dicDirectTranslationCS{věrný obraz, podobizna}
\dicEntry[spegilmyndavél] \textls[15]{\dicTerm{spegil··mynda·vél} \dicIPA{{s}{\textsubring{b}}{ei}{\textsci}{l}{m}{\textsci}{n}{\textsubring{d}}{a}{v}{j}{\textepsilon}{\textsubring{l}}} \dicPos{f}[4] \dicFlx{(‑ar, ‑ar)}[1] \dicDirectTranslationCS{zrcadlovka}}
\dicEntry[spegilsléttur] \dicTerm{spegil··sléttur} \dicIPA{{s}{\textsubring{b}}{ei}{\textsci}{\textsubring{l}}{s}{\textsubring{d}}{l}{j}{\textepsilon}{h}{\textsubring{d}}{\textscy}{\textsubring{r}}} \dicPos{adj}[1]\dicFlx{}[-10] \dicDirectTranslationCS{hladký jako zrcadlo (hladina ap.)}
\dicEntry[spegla] \dicTerm{spegl|a} \dicIPA{{s}{\textsubring{b}}{\textepsilon}{\r{g}}{l}{a}} \dicPos{v}[1] \dicFlx{(‑aði)}[1] \dicFlx{acc} \textbf{1.} \dicDirectTranslationCS{prohlédnout\,/\addthin prohlížet v~zrcadle};  \dicPhraseIS{spegla sig} \dicDirectTranslationCS{prohlédnout\,/\addthin prohlížet se v~zrcadle}  \textbf{2.} \dicDirectTranslationCS{zrcadlit, odrážet (v~zrcadle ap.)}  \textbf{3.} \dicFieldCat{med.} \dicDirectTranslationCS{vyšetřit pomocí endoskopu} \dicExampleIS{spegla magann} \dicExampleCS{vyšetřit břicho pomocí endoskopu};  \dicIdiom{speglast}{ \dicPhraseIS{speglast}} \dicFlx{refl} \dicDirectTranslationCS{zrcadlit se, odrážet se (v~zrcadle ap.)} \dicExampleIS{Húsin speglast í vatninu.} \dicExampleCS{Domy se odráží ve vodě.}
\dicEntry[speglun] \dicTerm{spegl|un} \dicIPA{{s}{\textsubring{b}}{\textepsilon}{\r{g}}{l}{\textscy}{\textsubring{n}}} \dicPos{f}[7] \dicFlx{(‑unar, ‑anir)}[8] \textbf{1.} \dicSynonym*{endurvarp} \dicDirectTranslationCS{zrcadlení, (zrcadlový) odraz\,/\addthin obraz}  \textbf{2.} \dicFieldCat{med. } \dicDirectTranslationCS{endoskopie}
\dicEntry[speki] \dicTerm{spek|i} \dicIPA{{s}{\textsubring{b}}{\textepsilon}{\textlengthmark}{\r{\textObardotlessj}}{\textsci}} \dicPos{f}[3] \dicFlx{(‑i)}[3] \dicSynonym{viska\smash{\textsuperscript{1}}} \dicDirectTranslationCS{moudro, moudrost, hlubokomyslnost} \dicExampleIS{tala speki} \dicExampleCS{říkat moudrosti}
\dicEntry[spekingslegur] \dicTerm{spekings··legur} \dicIPA{{s}{\textsubring{b}}{\textepsilon}{\textlengthmark}{\r{\textObardotlessj}}{i}{\ng}{s}{l}{\textepsilon}{\textbabygamma}{\textscy}{\textsubring{r}}} \dicPos{adj}[1]\dicFlx{}[-8] \dicDirectTranslationCS{moudrý, hlubokomyslný} \dicExampleIS{spekingslegur svipur} \dicExampleCS{moudrý výraz}
\dicEntry[spekingur] \dicTerm{spek··ing|ur} \dicIPA{{s}{\textsubring{b}}{\textepsilon}{\textlengthmark}{\r{\textObardotlessj}}{i}{\ng}{\r{g}}{\textscy}{\textsubring{r}}} \dicPos{m}[6] \dicFlx{(‑s, ‑ar)}[8] \textbf{1.} \dicSynonym*{vitur maður} \dicDirectTranslationCS{učenec, učenkyně, vzdělanec, vzdělankyně}  \textbf{2.} \dicSynonym{hugsuður} \dicDirectTranslationCS{myslitel(ka), mudrc, filosof(ka)}
\dicEntry[spekt] \dicTerm{spekt} \dicIPA{{s}{\textsubring{b}}{\textepsilon}{x}{\textsubring{d}}} \dicPos{f}[4] \dicFlx{(‑ar)}[3] \dicSynonym{rósemi} \dicDirectTranslationCS{(po)klid, klidnost, pokoj} \dicExampleIS{friður og spekt} \dicExampleCS{klid a~mír}
\dicEntry[spekúlant] \dicTerm{spekúlant} \dicIPA{{s}{\textsubring{b}}{\textepsilon}{\textlengthmark}{\r{g}}{u}{l}{a}{\textsubring{n}}{\textsubring{d}}} \dicPos{m}[4] \dicFlx{(‑s, ‑ar)}[11] \dicSynonym{braskari} \dicDirectTranslationCS{spekulant(ka), kšeftař(ka)}
\dicEntry[spekúlera] \dicTerm{spekúler|a} \dicIPA{{s}{\textsubring{b}}{\textepsilon}{\textlengthmark}{\r{g}}{u}{l}{\textepsilon}{r}{a}} \dicPos{v}[1] \dicFlx{(‑aði)}[1] \textbf{1.} \dicDirectTranslationCS{hloubat, přemýšlet, spekulovat};  \dicPhraseIS{spekúlera í e‑u} \dicDirectTranslationCS{hloubat o~(čem), spekulovat o~(čem)}  \textbf{2.} \dicFieldCat{ekon.} \dicSynonym{braska} \dicDirectTranslationCS{spekulovat (s~akciemi ap.)}
\dicEntry[speldi] \dicTerm{speldi} \dicIPA{{s}{\textsubring{b}}{\textepsilon}{l}{\textsubring{d}}{\textsci}} \dicPos{n}[2] \dicFlx{(‑s, ‑)}[14] \textbf{1.} \dicDirectTranslationCS{klapka, záklopka}  \textbf{2.} \dicFieldCat{anat.} \dicDirectTranslationCS{příklopka hrtanová} \textit{(l.~{\textLA{epiglottis}})}
\dicEntry[spelir] \dicTerm{spelir} \dicIPA{{s}{\textsubring{b}}{\textepsilon}{\textlengthmark}{l}{\textsci}{\textsubring{r}}} \dicPos{m} \dicFlx{pl nom} \dicLink{spölur}
\dicEntry[spelka] \dicTerm{spelk|a\smash{\textsuperscript{1}}} \dicIPA{{s}{\textsubring{b}}{\textepsilon}{\textsubring{l}}{\r{g}}{a}} \dicPos{f}[1] \dicFlx{(‑u, ‑ur)}[13] \dicDirectTranslationCS{dlaha, výztuha}
\dicEntry[spelka] \dicTerm{spelk|a\smash{\textsuperscript{2}}} \dicIPA{{s}{\textsubring{b}}{\textepsilon}{\textsubring{l}}{\r{g}}{a}} \dicPos{v}[1] \dicFlx{(‑aði)}[1] \dicFlx{acc} \dicDirectTranslationCS{dát do dlahy, vyztužit} \dicExampleIS{spelka fótinn} \dicExampleCS{dát nohu do dlahy}
\dicEntry[spellvirki] \dicTerm{spell··virk|i\smash{\textsuperscript{1}}} \dicIPA{{s}{\textsubring{b}}{\textepsilon}{\textsubring{d}}{\textsubring{l}}{v}{\textsci}{\textsubring{r}}{\r{\textObardotlessj}}{\textsci}} \dicPos{m}[1] \dicFlx{(‑ja, ‑jar)}[14] \dicSynonym{afbrotamaður} \dicDirectTranslationCS{vandal(ka), sabotér(ka)}
\dicEntry[spellvirki] \dicTerm{spell··virki\smash{\textsuperscript{2}}} \dicIPA{{s}{\textsubring{b}}{\textepsilon}{\textsubring{d}}{\textsubring{l}}{v}{\textsci}{\textsubring{r}}{\r{\textObardotlessj}}{\textsci}} \dicPos{n}[2] \dicFlx{(‑s, ‑)}[16] \dicDirectTranslationCS{vandalství, vandalismus}
\dicEntry[spendýr] \dicTerm{spen··dýr} \dicIPA{{s}{\textsubring{b}}{\textepsilon}{n}{\textsubring{d}}{i}{\textsubring{r}}} \dicPos{n}[2] \dicFlx{(‑s, ‑)}[5] \dicFieldCat{zool.} \dicDirectTranslationCS{savec} \textit{(l.~{\textLA{Mammalia}})}
\begin{xtolerant}{}{1pt}
\dicEntry[spengilegur] \dicTerm{spengi··legur}\addthinS\dicIPA{{s}{\textsubring{b}}{ei}{\textltailn}{\r{\textObardotlessj}}{\textsci}{l}{\textepsilon}{\textbabygamma}{\textscy}{\textsubring{r}}}\addthinS\dicPos{adj}[1]\dicFlx{}[-8]\addthinS\dicSynonym{grannvaxinn} \dicDirectTranslationCS{štíhlý, útlý}
\end{xtolerant}
\dicEntry[spengur] \dicTerm{spengur} \dicIPA{{s}{\textsubring{b}}{ei}{\ng}{\r{g}}{\textscy}{\textsubring{r}}} \dicPos{f} \dicFlx{pl nom} \dicLink{spöng}
\dicEntry[speni] \dicTerm{spen|i} \dicIPA{{s}{\textsubring{b}}{\textepsilon}{\textlengthmark}{n}{\textsci}} \dicPos{m}[1] \dicFlx{(‑a, ‑ar)}[1] \dicDirectTranslationCS{struk, cecík};  \dicPhraseIS{komast á spena(nn)} \dicFlx{refl} {\textbf{a.}} \dicDirectTranslationCS{nalézt cecík} \dicIndirectTranslationCS{(o~novorozeném zvířeti)};  {\textbf{b.}} \dicLangCat{přen.} \dicDirectTranslationCS{dostat se ke korytu}
\dicEntry[spenna] \dicTerm{spenn|a\smash{\textsuperscript{1}}} \dicsymFrequent\  \dicIPA{{s}{\textsubring{b}}{\textepsilon}{n}{\textlengthmark}{a}} \dicPos{f}[1] \dicFlx{(‑u, ‑ur)}[7] \textbf{1.} \dicSynonym*{þan} \dicDirectTranslationCS{(na)pnutí} \dicExampleIS{spenna á streng} \dicExampleCS{pnutí provazu}  \textbf{2.} \dicSynonym{taugaspenna} \dicDirectTranslationCS{napětí, nervozita, neklid};  \dicPhraseIS{slaka á spennunni} \dicDirectTranslationCS{uvolnit se, uklidnit se}  \textbf{3.} \dicFieldCat{elek.} \dicSynonym*{rafspenna} \dicDirectTranslationCS{napětí}  \textbf{4.} \dicSynonym{sylgja} \dicDirectTranslationCS{spona}  \textbf{5.} \dicSynonym{hárspenna} \dicDirectTranslationCS{sponka, vlásenka (do vlasů ap.)}
\dicEntry[spenna] \dicTerm{spenn|a\smash{\textsuperscript{2}}} \dicIPA{{s}{\textsubring{b}}{\textepsilon}{n}{\textlengthmark}{a}} \dicPos{v}[2] \dicFlx{(‑ti, ‑t)}[63] \dicFlx{acc} \dicDirectTranslationCS{napnout, napínat, natáhnout, natahovat, (na)šponovat} \dicExampleIS{spenna boga} \dicExampleCS{napnout luk};  \dicPhraseIS{spenna greipar} \dicDirectTranslationCS{sepnout ruce};  \dicIdiom{spenna}[upp]{ \dicPhraseIS{spenna e‑ð upp}} \dicDirectTranslationCS{násilím otevřít (co), vylomit (co), vylámat (co), vypáčit (co), vyrazit (co)} \dicExampleIS{spenna upp hurð} \dicExampleCS{vyrazit dveře};  \dicIdiom{spennast}{ \dicPhraseIS{spennast}} \dicFlx{refl} \dicDirectTranslationCS{napnout se, natáhnout se}
\dicEntry[spennandi] \dicTerm{spenn··andi} \dicsymFrequent\  \dicIPA{{s}{\textsubring{b}}{\textepsilon}{n}{\textlengthmark}{a}{n}{\textsubring{d}}{\textsci}} \dicPos{adj}[13] \dicFlx{indecl}[1] \dicDirectTranslationCS{vzrušující, napínavý, poutavý} \dicExampleIS{spennandi bækur} \dicExampleCS{poutavé knížky}
\dicEntry[spenningur] \dicTerm{spenn··ing|ur} \dicIPA{{s}{\textsubring{b}}{\textepsilon}{n}{\textlengthmark}{i}{\ng}{\r{g}}{\textscy}{\textsubring{r}}} \dicPos{m}[6] \dicFlx{(‑s)}[9] \dicDirectTranslationCS{napětí, vzrušení} \dicExampleIS{finna fyrir spenningi} \dicExampleCS{cítit napětí}
\dicEntry[spennistöð] \dicTerm{spenni··stöð} \dicIPA{{s}{\textsubring{b}}{\textepsilon}{n}{\textlengthmark}{\textsci}{s}{\textsubring{d}}{\oe}{\texttheta}} \dicPos{f}[6] \dicFlx{(‑var, ‑var)}[1] \dicFieldCat{elek.} \dicDirectTranslationCS{transformátorová stanice, trafostanice}
\dicEntry[spennitreyja] \dicTerm{spenni··treyj|a} \dicIPA{{s}{\textsubring{b}}{\textepsilon}{n}{\textlengthmark}{\textsci}{t\smash{\textsuperscript{h}}}{r}{ei}{j}{a}} \dicPos{f}[1] \dicFlx{(‑u, ‑ur)}[7] \dicDirectTranslationCS{svěrací kazajka}
\dicEntry[spenntur] \dicTerm{spenntur} \dicsymFrequent\  \dicIPA{{s}{\textsubring{b}}{\textepsilon}{\textsubring{n}}{\textsubring{d}}{\textscy}{\textsubring{r}}} \dicPos{adj}[1]\dicFlx{}[-10] \dicDirectTranslationCS{napjatý, napnutý, vzrušený} \dicExampleIS{Ég var mjög spenntur fyrir að sjá handritið.} \dicExampleCS{Byl jsem moc napjatý, že uvidím rukopis.}
\dicEntry[spennubók] \dicTerm{spennu··|bók} \dicIPA{{s}{\textsubring{b}}{\textepsilon}{n}{\textlengthmark}{\textscy}{\textsubring{b}}{ou}{\r{g}}} \dicPos{f}[8] \dicFlx{(‑bókar, ‑bækur)}[5] \dicDirectTranslationCS{napínavá kniha, thriller} \dicIndirectTranslationCS{(kniha)}
\dicEntry[spennubreytir] \dicTerm{spennu··breyt|ir} \dicIPA{{s}{\textsubring{b}}{\textepsilon}{n}{\textlengthmark}{\textscy}{\textsubring{b}}{r}{ei}{\textsubring{d}}{\textsci}{\textsubring{r}}} \dicPos{m}[7] \dicFlx{(‑is, ‑ar)}[1] \dicFieldCat{elek.} \dicDirectTranslationCS{transformátor}
\dicEntry[spennufall] \dicTerm{spennu··fall} \dicIPA{{s}{\textsubring{b}}{\textepsilon}{n}{\textlengthmark}{\textscy}{f}{a}{\textsubring{d}}{\textsubring{l}}} \dicPos{n}[2] \dicFlx{(‑s)}[2] \dicDirectTranslationCS{uvolnění napětí} \dicExampleIS{Íbúar eru í spennufalli.} \dicExampleCS{U~obyvatel dochází k~uvolnění napětí.}
\dicEntry[spennumynd] \dicTerm{spennu··mynd} \dicIPA{{s}{\textsubring{b}}{\textepsilon}{n}{\textlengthmark}{\textscy}{m}{\textsci}{n}{\textsubring{d}}} \dicPos{f}[7] \dicFlx{(‑ar, ‑ir)}[1] \dicDirectTranslationCS{napínavý film, thriller} \dicIndirectTranslationCS{(film)}
\dicEntry[sperðill] \dicTerm{sperð|ill} \dicIPA{{s}{\textsubring{b}}{\textepsilon}{r}{ð}{\textsci}{\textsubring{d}}{\textsubring{l}}} \dicPos{m}[6] \dicFlx{(‑ils, ‑lar)}[35] \dicFieldCat{kulin.} \dicSynonym{bjúga} \dicDirectTranslationCS{uzenka}
\dicEntry[spergilkál] \dicTerm{spergil··kál} \dicIPA{{s}{\textsubring{b}}{\textepsilon}{r}{\r{\textObardotlessj}}{\textsci}{\textsubring{l}}{k\smash{\textsuperscript{h}}}{au}{\textsubring{l}}} \dicPos{n}[2] \dicFlx{(‑s)}[2] \dicFieldCat{bot.} \dicDirectTranslationCS{brokolice} \textit{(l.~{\textLA{Brassica oleracea var. italica}})}  \dicsymPhoto\ 
\dicFigure{7997.jpg}{Spergilkál}{Spergilkál - Deml Miroslav, Biolib, Copyright/CC-BY}
\dicEntry[spergill] \dicTerm{sperg|ill} \dicIPA{{s}{\textsubring{b}}{\textepsilon}{r}{\r{\textObardotlessj}}{\textsci}{\textsubring{d}}{\textsubring{l}}} \dicPos{m}[6] \dicFlx{(‑ils, ‑lar)}[35] \dicFieldCat{bot.} \dicDirectTranslationCS{chřest, asparágus} \textit{(l.~{\textLA{Asparagus officinalis}})}
\dicEntry[sperra] \dicTerm{sperr|a\smash{\textsuperscript{1}}} \dicIPA{{s}{\textsubring{b}}{\textepsilon}{r}{\textlengthmark}{a}} \dicPos{f}[1] \dicFlx{(‑u, ‑ur)}[7] \dicDirectTranslationCS{krokev}
\dicEntry[sperra] \dicTerm{sperr|a\smash{\textsuperscript{2}}} \dicIPA{{s}{\textsubring{b}}{\textepsilon}{r}{\textlengthmark}{a}} \dicPos{v}[2] \dicFlx{(‑ti, ‑t)}[63] \dicFlx{acc} \textbf{1.} \dicSynonym*{rétta fram} \dicDirectTranslationCS{vzpřímit, vzpřimovat, vztyčit, vztyčovat} \dicExampleIS{Kötturinn sperrti eyrun.} \dicExampleCS{Kočka nastražila uši.}  \textbf{2.} \dicPhraseIS{sperra sig} {\textbf{a.}} \dicSynonym*{reigja sig} \dicDirectTranslationCS{vzpřímit se, vztyčit se};  {\textbf{b.}} \dicSynonym*{herða sig} \dicDirectTranslationCS{naparovat se, vypínat se, kasat se};  \dicIdiom{sperra}[upp]{ \dicPhraseIS{sperra upp e‑ð}} \dicDirectTranslationCS{otevřít (co) dokořán} \dicExampleIS{sperra upp hurðina} \dicExampleCS{otevřít dveře dokořán};  \dicIdiom{sperrast}{ \dicPhraseIS{sperrast við að (gera e‑ð)}} \dicFlx{refl} \dicSynonym{strita} \dicDirectTranslationCS{dřít se s~(děláním (čeho)), mořit se s~(děláním (čeho))}
\dicEntry[sperrileggur] \dicTerm{sperri··legg|ur} \dicIPA{{s}{\textsubring{b}}{\textepsilon}{r}{\textlengthmark}{\textsci}{l}{\textepsilon}{\r{g}}{\textscy}{\textsubring{r}}} \dicPos{m}[9] \dicFlx{(‑jar\,/\addthin ‑s, ‑ir)}[26] \textbf{1.} \dicFieldCat{anat.} \dicDirectTranslationCS{lýtková kost, fibula}  \textbf{2.} \dicDirectTranslationCS{kecal, mluvka, žvanil}
\dicEntry[sperringur] \dicTerm{sperr··ing|ur} \dicIPA{{s}{\textsubring{b}}{\textepsilon}{r}{\textlengthmark}{i}{\ng}{\r{g}}{\textscy}{\textsubring{r}}} \dicPos{m}[6] \dicFlx{(‑s)}[9] \dicSynonym{rembingur} \dicDirectTranslationCS{naparování, chvástání, kasání se}
\dicEntry[sperrtur] \dicTerm{sperrtur} \dicIPA{{s}{\textsubring{b}}{\textepsilon}{\textsubring{r}}{\textsubring{d}}{\textscy}{\textsubring{r}}} \dicPos{adj}[1]\dicFlx{}[-10] \textbf{1.} \dicSynonym{montinn} \dicDirectTranslationCS{nadutý, namyšlený, povýšený}  \textbf{2.} \dicSynonym{uppréttur} \dicDirectTranslationCS{vzpřímený, vztyčený}
\dicEntry[spes] \dicTerm{spes} \dicIPA{{s}{\textsubring{b}}{\textepsilon}{\textlengthmark}{s}} \dicPos{adj}[13] \dicFlx{indecl}[1] \dicLangCat{hovor.} \dicSynonym{sérstakur} \dicDirectTranslationCS{speciální, výjimečný}
\dicEntry[spé] \dicTerm{spé} \dicIPA{{s}{\textsubring{b}}{j}{\textepsilon}{\textlengthmark}} \dicPos{n}[2] \dicFlx{(‑s)}[2] \dicSynonym{háð\smash{\textsuperscript{1}}} \dicDirectTranslationCS{úšklebky, posměch} \dicExampleIS{fara með spé} \dicExampleCS{posmívat se}
\dicEntry[spéhræddur] \dicTerm{spé··hræddur} \dicIPA{{s}{\textsubring{b}}{j}{\textepsilon}{\textlengthmark}{\textsubring{r}}{a}{i}{\textsubring{d}}{\textscy}{\textsubring{r}}} \dicPos{adj}[2]\dicFlx{}[-18] \dicDirectTranslationCS{nedůtklivý, bojící se posměchu}
\dicEntry[spékoppur] \dicTerm{spé··kopp|ur} \dicIPA{{s}{\textsubring{b}}{j}{\textepsilon}{\textlengthmark}{k\smash{\textsuperscript{h}}}{\textopeno}{h}{\textsubring{b}}{\textscy}{\textsubring{r}}} \dicPos{m}[6] \dicFlx{(‑s, ‑ar)}[24] \dicDirectTranslationCS{dolíček, důlek} \dicIndirectTranslationCS{(ve tváři od smíchu)}
\dicEntry[spéspegill] \dicTerm{spé··speg|ill} \dicIPA{{s}{\textsubring{b}}{j}{\textepsilon}{\textlengthmark}{s}{\textsubring{b}}{ei}{\textsci}{\textsubring{d}}{\textsubring{l}}} \dicPos{m}[6] \dicFlx{(‑ils, ‑lar)}[35] \textbf{1.} \dicDirectTranslationCS{zkreslující zrcadlo}  \textbf{2.} \dicLangCat{přen.} \dicDirectTranslationCS{křivé zrcadlo}
\dicEntry[spik] \dicTerm{spik} \dicIPA{{s}{\textsubring{b}}{\textsci}{\textlengthmark}{\r{g}}} \dicPos{n}[2] \dicFlx{(‑s)}[2] \dicSynonym{fita\smash{\textsuperscript{1}}} \dicDirectTranslationCS{špek, sádlo};  \dicPhraseIS{hlaupa af sér spikið} \dicLangCat{přen.} \dicDirectTranslationCS{shodit špeky, zhubnout};  \dicPhraseIS{hlaupa í spik} \dicDirectTranslationCS{zakládat na špek\,/\addthin sádlo, tloustnout};  \dicPhraseIS{róa í spikinu} \dicLangCat{přen.} \dicDirectTranslationCS{být širší než vyšší}
\dicEntry[spikfeitur] \dicTerm{spik··feitur} \dicIPA{{s}{\textsubring{b}}{\textsci}{\textlengthmark}{\r{g}}{f}{ei}{\textsubring{d}}{\textscy}{\textsubring{r}}} \dicPos{adj}[1]\dicFlx{}[-6] \textbf{1.} \dicDirectTranslationCS{obézní, otylý}  \textbf{2.} \dicDirectTranslationCS{tlustý (maso ap.)}
\dicEntry[spil] \dicTerm{spil} \dicsymFrequent\  \dicIPA{{s}{\textsubring{b}}{\textsci}{\textlengthmark}{\textsubring{l}}} \dicPos{n}[2] \dicFlx{(‑s, ‑)}[5] \textbf{1.} \dicDirectTranslationCS{(hrací) karta} \dicExampleIS{pakki af spilum} \dicExampleCS{balíček karet};  \dicPhraseIS{spá í spil} \dicDirectTranslationCS{věštit z~karet}  \textbf{2.} \dicSynonym{leikur} \dicDirectTranslationCS{hra (karetní ap.)}  \textbf{3.} \dicDirectTranslationCS{hra, hraní (na piano ap.)} \dicExampleIS{harmónikuspil} \dicExampleCS{hra na harmoniku}  \textbf{4.} \dicDirectTranslationCS{naviják};  \dicIdiom{spil}{ \dicPhraseIS{halda (vel) á spilunum}} \dicDirectTranslationCS{vynasnažit se, přičinit se};  \dicPhraseIS{leggja spilin á borðið} \dicLangCat{přen.} \dicDirectTranslationCS{vyložit karty na stůl}
\dicEntry[spila] \dicTerm{spil|a} \dicsymFrequent\  \dicIPA{{s}{\textsubring{b}}{\textsci}{\textlengthmark}{l}{a}} \dicPos{v}[1] \dicFlx{(‑aði)}[1] \textbf{1.} \dicDirectTranslationCS{hrát} \dicIndirectTranslationCS{(provozovat jako sport)} \dicExampleIS{spila handbolta} \dicExampleCS{hrát házenou}  \textbf{2.} \dicDirectTranslationCS{(za)hrát} \dicIndirectTranslationCS{(provozovat hudbu)} \dicExampleIS{spila á píanó} \dicExampleCS{hrát na piano}  \textbf{3.} \dicDirectTranslationCS{(za)hrát} \dicIndirectTranslationCS{(zaměstnávat se společenskou hrou)} \dicExampleIS{spila á spil} \dicExampleCS{hrát karty}  \textbf{4.} \dicDirectTranslationCS{hrát, sázet (loterii ap.)} \dicExampleIS{spila í happdrætti} \dicExampleCS{hrát loterii};  \dicIdiom{spila}[burt]{ \dicPhraseIS{spila e‑u burt}} \dicDirectTranslationCS{prohrát (co)} \dicExampleIS{spila burt eignum sínum} \dicExampleCS{prohrát jmění};  \dicIdiom{spila}[með]{ \dicPhraseIS{spila með e‑n}} \dicLangCat{přen.} \dicDirectTranslationCS{hrát si s~(kým), manipulovat (s) (kým)};  \dicIdiom{spila}[um]{ \dicPhraseIS{spila um e‑ð}} \dicDirectTranslationCS{hrát o~(co)};  \dicIdiom{spila}[úr]{ \dicPhraseIS{hafa úr litlu að spila}} \dicDirectTranslationCS{mít málo (prostředků)} \dicIndirectTranslationCS{(být chudý)};  \dicIdiom{spila}[út]{ \dicPhraseIS{vera að spila út}} \dicLangCat{hovor.} \dicDirectTranslationCS{hrabat, blbnout} \dicExampleIS{Hún er að spila út.} \dicExampleCS{Hrabe jí.}
\dicEntry[spilaborð] \dicTerm{spila··borð} \dicIPA{{s}{\textsubring{b}}{\textsci}{\textlengthmark}{l}{a}{\textsubring{b}}{\textopeno}{r}{\texttheta}} \dicPos{n}[2] \dicFlx{(‑s, ‑)}[5] \dicDirectTranslationCS{hrací stůl\,/\addthin stolek}
\dicEntry[spilaborg] \dicTerm{spila··borg} \dicIPA{{s}{\textsubring{b}}{\textsci}{\textlengthmark}{l}{a}{\textsubring{b}}{\textopeno}{r}{\r{g}}} \dicPos{f}[7] \dicFlx{(‑ar, ‑ir)}[1] \textbf{1.} \dicDirectTranslationCS{domeček z~karet} \dicIndirectTranslationCS{(dětská hra)}  \textbf{2.} \dicLangCat{přen.} \dicDirectTranslationCS{domeček z~karet (něco, co se snadno zhroutí ap.)}
\dicEntry[spiladós] \dicTerm{spila··dós} \dicIPA{{s}{\textsubring{b}}{\textsci}{\textlengthmark}{l}{a}{\textsubring{d}}{ou}{s}} \dicPos{f}[7] \dicFlx{(‑ar, ‑ir)}[1] \dicDirectTranslationCS{hrací skříňka}
\dicEntry[spilafíkill] \dicTerm{spila··fík|ill} \dicIPA{{s}{\textsubring{b}}{\textsci}{\textlengthmark}{l}{a}{f}{i}{\r{\textObardotlessj}}{\textsci}{\textsubring{d}}{\textsubring{l}}} \dicPos{m}[6] \dicFlx{(‑ils, ‑lar)}[35] \dicDirectTranslationCS{gambler(ka), závislý\,/\addthin závislá na hazardních hrách}
\dicEntry[spilafíkn] \dicTerm{spila··fíkn} \dicIPA{{s}{\textsubring{b}}{\textsci}{\textlengthmark}{l}{a}{f}{i}{h}{\r{g}}{\textsubring{n}}} \dicPos{f}[4] \dicFlx{(‑ar)}[3] \dicDirectTranslationCS{hráčská závislost, závislost na hazardních hrách, gambling}
\dicEntry[spilagaldur] \dicTerm{spila··gald|ur} \dicIPA{{s}{\textsubring{b}}{\textsci}{\textlengthmark}{l}{a}{\r{g}}{a}{l}{\textsubring{d}}{\textscy}{\textsubring{r}}} \dicPos{m}[5] \dicFlx{(‑urs, ‑rar)}[3] \dicDirectTranslationCS{karetní trik\,/\addthin kouzlo}
\dicEntry[spilakassi] \dicTerm{spila··kass|i} \dicIPA{{s}{\textsubring{b}}{\textsci}{\textlengthmark}{l}{a}{k\smash{\textsuperscript{h}}}{a}{s}{\textsci}} \dicPos{m}[1] \dicFlx{(‑a, ‑ar)}[8] \dicDirectTranslationCS{výherní\,/\addthin hrací automat}
\dicEntry[spilapeningur] \dicTerm{spila··pen·ing|ur} \dicIPA{{s}{\textsubring{b}}{\textsci}{\textlengthmark}{l}{a}{p\smash{\textsuperscript{h}}}{\textepsilon}{n}{i}{\ng}{\r{g}}{\textscy}{\textsubring{r}}} \dicPos{m}[6] \dicFlx{(‑s, ‑ar)}[8] \dicDirectTranslationCS{(hrací) žeton}
\dicEntry[spilari] \dicTerm{spil··ar|i} \dicIPA{{s}{\textsubring{b}}{\textsci}{\textlengthmark}{l}{a}{r}{\textsci}} \dicPos{m}[1] \dicFlx{(‑a, ‑ar)}[13] \dicDirectTranslationCS{přehrávač} \dicExampleIS{geislaspilari} \dicExampleCS{CD přehrávač}
\dicEntry[spilasalur] \dicTerm{spila··sal|ur} \dicIPA{{s}{\textsubring{b}}{\textsci}{\textlengthmark}{l}{a}{s}{a}{l}{\textscy}{\textsubring{r}}} \dicPos{m}[10] \dicFlx{(‑ar, ‑ir)}[14] \dicDirectTranslationCS{herna, kasino}
\dicEntry[spilaskuld] \dicTerm{spila··skuld} \dicIPA{{s}{\textsubring{b}}{\textsci}{\textlengthmark}{l}{a}{s}{\r{g}}{\textscy}{l}{\textsubring{d}}} \dicPos{f}[7] \dicFlx{(‑ar, ‑ir)}[1] \dicDirectTranslationCS{dluh z~hazardu}
\dicEntry[spilastokkur] \dicTerm{spila··stokk|ur} \dicIPA{{s}{\textsubring{b}}{\textsci}{\textlengthmark}{l}{a}{s}{\textsubring{d}}{\textopeno}{h}{\r{g}}{\textscy}{\textsubring{r}}} \dicPos{m}[6] \dicFlx{(‑s, ‑ar)}[22] \dicDirectTranslationCS{balíček karet}
\dicEntry[spilavíti] \dicTerm{spila··víti} \dicIPA{{s}{\textsubring{b}}{\textsci}{\textlengthmark}{l}{a}{v}{i}{\textsubring{d}}{\textsci}} \dicPos{n}[2] \dicFlx{(‑s, ‑)}[14] \dicDirectTranslationCS{herna, kasino}
\dicEntry[spilda] \dicTerm{spild|a}\dicTerm{, spílda} \dicIPA{{s}\-{\textsubring{b}}\-{\textsci}\-{l}\-{\textsubring{d}}\-{a}\-} \dicPos{f}[1] \dicFlx{(‑u, ‑ur)}[13] \dicDirectTranslationCS{kus\,/\addthin pás\,/\addthin lán země} \dicExampleIS{hraunspilda} \dicExampleCS{lávový pás země}
\dicEntry[spilla] \dicTerm{spill|a} \dicsymFrequent\  \dicIPA{{s}{\textsubring{b}}{\textsci}{\textsubring{d}}{l}{a}} \dicPos{v}[2] \dicFlx{(‑ti, ‑t)}[63] \dicFlx{dat} \textbf{1.} \dicSynonym{skemma\smash{\textsuperscript{2}}} \dicDirectTranslationCS{(z)kazit, pokazit} \dicExampleIS{spilla matnum} \dicExampleCS{zkazit jídlo}  \textbf{2.} \dicSynonym{siðspilla} \dicDirectTranslationCS{(z)korumpovat};  \dicIdiom{spilla}[fyrir]{ \dicPhraseIS{spilla fyrir e‑u}} \dicDirectTranslationCS{zkazit (co), překazit (co)};  \dicIdiom{spillast}{ \dicPhraseIS{spillast}} \dicFlx{refl} \dicDirectTranslationCS{zkazit se, pokazit se}
\dicEntry[spilling] \dicTerm{spill··ing} \dicIPA{{s}{\textsubring{b}}{\textsci}{\textsubring{d}}{l}{i}{\ng}{\r{g}}} \dicPos{f}[4] \dicFlx{(‑ar)}[7] \dicDirectTranslationCS{korupce, úplatkářství} \dicExampleIS{koma í veg fyrir spillingu} \dicExampleCS{zabránit korupci}
\dicEntry[spillingarbæli] \dicTerm{spillingar··bæli} \dicIPA{{s}{\textsubring{b}}{\textsci}{\textsubring{d}}{l}{i}{\ng}{\r{g}}{a}{r}{\textsubring{b}}{a}{i}{l}{\textsci}} \dicPos{n}[2] \dicFlx{(‑s, ‑)}[14] \dicDirectTranslationCS{doupě neřesti}
\dicEntry[spilltur] \dicTerm{spilltur} \dicIPA{{s}{\textsubring{b}}{\textsci}{\textsubring{l}}{\textsubring{d}}{\textscy}{\textsubring{r}}} \dicPos{adj}[1]\dicFlx{}[-10] \textbf{1.} \dicSynonym{siðlaus} \dicDirectTranslationCS{zkorumpovaný, prodejný, úplatný} \dicExampleIS{spilltur embættismaður} \dicExampleCS{zkorumpovaný úředník}  \textbf{2.} \dicSynonym{skemmdur} \dicDirectTranslationCS{zkažený (jídlo ap.)} \dicExampleIS{spilltur matur} \dicExampleCS{zkažené jídlo}
\dicEntry[spinn] \dicTerm{spinn} \dicIPA{{s}{\textsubring{b}}{\textsci}{n}{\textlengthmark}} \dicPos{v} \dicFlx{ind praes sg 1 pers} \dicLink{spinna}
\dicEntry[spinna] \dicTerm{spinna} \dicIPA{{s}{\textsubring{b}}{\textsci}{n}{\textlengthmark}{a}} \dicPos{v}[6] \dicFlx{(spinn, spann, spunnum, spynni, spunnið)}[22] \dicFlx{acc} \textbf{1.} \dicDirectTranslationCS{(u)příst, spřádat (vlákno ap.)} \dicExampleIS{spinna ull} \dicExampleCS{příst vlnu}  \textbf{2.} \dicPhraseIS{spinna upp e‑ð} \dicDirectTranslationCS{vymyslet (si) (co), vybásnit (si) (co)} \dicExampleIS{spinna upp sögu} \dicExampleCS{vymyslet si historku}  \textbf{3.} \dicFieldCat{let.} \dicDirectTranslationCS{(u)dělat vývrtku};  \dicIdiom{spinnast}{ \dicPhraseIS{spinnast}} \dicFlx{refl} \dicDirectTranslationCS{rozvíjet se, odvíjet se (debata ap.)}
\dicEntry[spígspora] \dicTerm{spíg··spor|a} \dicIPA{{s}{\textsubring{b}}{i}{x}{s}{\textsubring{b}}{\textopeno}{r}{a}} \dicPos{v}[1] \dicFlx{(‑aði)}[44] \dicDirectTranslationCS{chodit sem a~tam, vykračovat si}
\dicEntry[spílda] \dicTerm{spíld|a} \dicIPA{{s}{\textsubring{b}}{i}{l}{\textsubring{d}}{a}} \dicPos{f}[1] \dicFlx{(‑u, ‑ur)}[13] \dicLink{spilda}
\dicEntry[spínat] \dicTerm{spínat} \dicIPA{{s}{\textsubring{b}}{i}{\textlengthmark}{n}{a}{\textsubring{d}}} \dicPos{n}[2] \dicFlx{(‑s)}[2] \dicFieldCat{bot.} \dicDirectTranslationCS{špenát} \textit{(l.~{\textLA{Spinacia oleracea}})}  \dicsymPhoto\ 
\dicFigure{ds_image_spinat_0_1.jpg}{Spínat}{Spínat - Wikimedia Commons, CC BY-SA 3.0}
\dicEntry[spíra] \dicTerm{spír|a\smash{\textsuperscript{1}}} \dicIPA{{s}{\textsubring{b}}{i}{\textlengthmark}{r}{a}} \dicPos{f}[1] \dicFlx{(‑u, ‑ur)}[7] \textbf{1.} \dicDirectTranslationCS{špice, špička (kostelní věže ap.)} \dicExampleIS{turnspíra} \dicExampleCS{špička věže}  \textbf{2.} \dicFieldCat{bot.} \dicDirectTranslationCS{klíček}  \textbf{3.} \dicLangCat{han.} \dicDirectTranslationCS{všeuměl(kyně), všeználek} \dicIndirectTranslationCS{(zvláště ve spojeních)} \dicExampleIS{listaspíra} \dicExampleCS{umělecký všeuměl}
\dicEntry[spíra] \dicTerm{spír|a\smash{\textsuperscript{2}}} \dicIPA{{s}{\textsubring{b}}{i}{\textlengthmark}{r}{a}} \dicPos{v}[1] \dicFlx{(‑aði)}[1] \dicDirectTranslationCS{(na)klíčit, vzklíčit}
\dicEntry[spíraður] \dicTerm{spír··|aður} \dicIPA{{s}{\textsubring{b}}{i}{\textlengthmark}{r}{a}{ð}{\textscy}{\textsubring{r}}} \dicPos{adj}[3] \dicFlx{(f ‑uð)}[4] \dicDirectTranslationCS{vyklíčený, naklíčený}
\dicEntry[spírall] \dicTerm{spíral|l} \dicIPA{{s}{\textsubring{b}}{i}{\textlengthmark}{r}{a}{\textsubring{d}}{\textsubring{l}}} \dicPos{m}[6] \dicFlx{(‑s, ‑ar)}[49] \dicDirectTranslationCS{spirála}
\dicEntry[spíri] \dicTerm{spír|i} \dicIPA{{s}{\textsubring{b}}{i}{\textlengthmark}{r}{\textsci}} \dicPos{m}[1] \dicFlx{(‑a)}[3] \dicLangCat{hovor.} \dicDirectTranslationCS{pálenka, kořalka, chlast}
\dicEntry[spíritismi] \dicTerm{spírit··ism|i} \dicIPA{{s}{\textsubring{b}}{i}{\textlengthmark}{r}{\textsci}{\textsubring{d}}{\textsci}{s}{m}{\textsci}} \dicPos{m}[1] \dicFlx{(‑a)}[3] \dicDirectTranslationCS{spiritismus}
\dicEntry[spíritisti] \dicTerm{spírit··ist|i} \dicIPA{{s}{\textsubring{b}}{i}{\textlengthmark}{r}{\textsci}{\textsubring{d}}{\textsci}{s}{\textsubring{d}}{\textsci}} \dicPos{m}[1] \dicFlx{(‑a, ‑ar)}[1] \dicDirectTranslationCS{spiritista, spiritistka}
\dicEntry[spíritus] \dicTerm{spíritus} \dicIPA{{s}{\textsubring{b}}{i}{\textlengthmark}{r}{\textsci}{\textsubring{d}}{\textscy}{s}} \dicPos{m}[4] \dicFlx{(‑s)}[18] \dicFieldCat{chem.} \dicDirectTranslationCS{líh, čistý alkohol}
\dicEntry[spírun] \dicTerm{spír|un} \dicIPA{{s}{\textsubring{b}}{i}{\textlengthmark}{r}{\textscy}{\textsubring{n}}} \dicPos{f}[7] \dicFlx{(‑unar)}[9] \dicDirectTranslationCS{(na)klíčení}
\dicEntry[spítali] \dicTerm{spítal|i} \dicIPA{{s}{\textsubring{b}}{i}{\textlengthmark}{\textsubring{d}}{a}{l}{\textsci}} \dicPos{m}[1] \dicFlx{(‑a, ‑ar)}[8] \dicDirectTranslationCS{nemocnice, špitál}
\dicEntry[spítt] \dicTerm{spítt} \dicIPA{{s}{\textsubring{b}}{i}{h}{\textsubring{d}}} \dicPos{n}[2] \dicFlx{(‑s, ‑)}[5] \textbf{1.} \dicLangCat{hovor.} \dicSynonym{hraði} \dicDirectTranslationCS{rychlost} \dicExampleIS{spíttbátur} \dicExampleCS{rychlý člun}  \textbf{2.} \dicLangCat{slang.} \dicSynonym{amfetamín} \dicDirectTranslationCS{speed} \dicIndirectTranslationCS{(droga)}
\dicEntry[spjald] \dicTerm{spjald} \dicsymFrequent\  \dicIPA{{s}{\textsubring{b}}{j}{a}{l}{\textsubring{d}}} \dicPos{n}[2] \dicFlx{(‑s, spjöld)}[8] \textbf{1.} \dicSynonym*{pappaspjald} \dicDirectTranslationCS{kartička, lístek} \dicExampleIS{prenta orðin á spjöld} \dicExampleCS{vytisknout si slova na kartičky}  \textbf{2.} \dicSynonym*{hörð bókarkápa} \dicDirectTranslationCS{obal\,/\addthin přebal knihy}  \textbf{3.} \dicFieldCat{sport.} \dicDirectTranslationCS{karta};  \dicPhraseIS{gult spjald} \dicFieldCat{sport.} \dicDirectTranslationCS{žlutá karta};  \dicPhraseIS{rautt spjald} \dicFieldCat{sport.} \dicDirectTranslationCS{červená karta};  \dicIdiom{spjald}{ \dicPhraseIS{spjaldanna á milli}} \dicFlx{adv} \dicDirectTranslationCS{na jeden zátah, od začátku do konce}
\dicEntry[spjaldbein] \dicTerm{spjald··bein} \dicIPA{{s}{\textsubring{b}}{j}{a}{l}{\textsubring{d}}{\textsubring{b}}{ei}{\textsubring{n}}} \dicPos{n}[2] \dicFlx{(‑s, ‑)}[5] \dicFieldCat{anat.} \dicDirectTranslationCS{křížová kost}
\dicEntry[spjaldskrá] \dicTerm{spjald··skrá} \dicIPA{{s}{\textsubring{b}}{j}{a}{l}{\textsubring{d}}{s}{\r{g}}{r}{au}} \dicPos{f}[4] \dicFlx{(‑r\,/\addthin ‑ar, ‑r)}[21] \dicDirectTranslationCS{kartotéka, lístkovnice} \dicExampleIS{bókaspjaldskrá} \dicExampleCS{knižní kartotéka}
\dicEntry[spjaldtölva] \dicTerm{spjald··tölv|a} \dicIPA{{s}{\textsubring{b}}{j}{a}{l}{\textsubring{d}}{t\smash{\textsuperscript{h}}}{\oe}{l}{v}{a}} \dicPos{f}[1] \dicFlx{(‑u, ‑ur)}[7] \dicFieldCat{poč.} \dicDirectTranslationCS{tablet} \dicIndirectTranslationCS{(počítač)}
\dicEntry[spjall] \dicTerm{spjall} \dicIPA{{s}{\textsubring{b}}{j}{a}{\textsubring{d}}{\textsubring{l}}} \dicPos{n}[2] \dicFlx{(‑s, spjöll)}[8] \textbf{1.} \dicSynonym{skraf} \dicDirectTranslationCS{(po)povídání, pokec} \dicExampleIS{spjall um daginn og veginn} \dicExampleCS{popovídání o~všem možném}  \textbf{2.} \dicPhraseIS{spjöll} \dicFlx{pl} \dicDirectTranslationCS{spoušť, zničení, vandalismus}
\dicEntry[spjalla] \dicTerm{spjall|a} \dicsymFrequent\  \dicIPA{{s}{\textsubring{b}}{j}{a}{\textsubring{d}}{l}{a}} \dicPos{v}[1] \dicFlx{(‑aði)}[13] \dicSynonym*{tala saman} \dicDirectTranslationCS{(po)povídat si} \dicExampleIS{spjalla við e‑n um e‑ð} \dicExampleCS{popovídat si s~(kým) o~(čem)}
\dicEntry[spjallborð] \dicTerm{spjall··borð} \dicIPA{{s}{\textsubring{b}}{j}{a}{\textsubring{d}}{l}{\textsubring{b}}{\textopeno}{r}{\texttheta}} \dicPos{n}[2] \dicFlx{(‑s, ‑)}[5] \dicFieldCat{poč.} \dicDirectTranslationCS{(diskuzní) fórum}
\dicEntry[spjallrás] \dicTerm{spjall··rás} \dicIPA{{s}{\textsubring{b}}{j}{a}{\textsubring{d}}{l}{r}{au}{s}} \dicPos{f}[7] \dicFlx{(‑ar, ‑ir)}[1] \dicFieldCat{poč.} \dicDirectTranslationCS{chat, diskuzní kanál}
\dicEntry[spjallþáttur] \dicTerm{spjall··|þáttur} \dicIPA{{s}{\textsubring{b}}{j}{a}{\textsubring{d}}{\textsubring{l}}{\texttheta}{au}{h}{\textsubring{d}}{\textscy}{\textsubring{r}}} \dicPos{m}[12] \dicFlx{(‑þáttar, ‑þættir)}[7] \dicDirectTranslationCS{diskuzní pořad\,/\addthin relace, talk show}
\dicEntry[spjallþráður] \dicTerm{spjall··|þráður} \dicIPA{{s}{\textsubring{b}}{j}{a}{\textsubring{d}}{\textsubring{l}}{\texttheta}{r}{au}{ð}{\textscy}{\textsubring{r}}} \dicPos{m}[12] \dicFlx{(‑þráðar, ‑þræðir)}[7] \dicFieldCat{poč.} \dicDirectTranslationCS{příspěvek diskuze}
\dicEntry[spjara] \dicTerm{spjar|a} \dicIPA{{s}{\textsubring{b}}{j}{a}{\textlengthmark}{r}{a}} \dicPos{v}[1] \dicFlx{(‑aði)}[13] \dicPhraseIS{spjara sig} \dicSynonym*{sjá um sig} \dicDirectTranslationCS{zvládnout, poradit si, vědět si rady} \dicExampleIS{Hún spjarar sig einhvern veginn.} \dicExampleCS{Nějak to zvládne.}
\dicEntry[spjarar] \dicTerm{spjarar} \dicIPA{{s}{\textsubring{b}}{j}{a}{\textlengthmark}{r}{a}{\textsubring{r}}} \dicPos{f} \dicFlx{sg gen} \dicLink{spjör}
\dicEntry[spjátrungslegur] \dicTerm{spjátr·ungs··legur} \dicIPA{{s}{\textsubring{b}}{j}{au}{\textlengthmark}{\textsubring{d}}{r}{u}{\ng}{s}{l}{\textepsilon}{\textbabygamma}{\textscy}{\textsubring{r}}} \dicPos{adj}[1]\dicFlx{}[-8] \dicDirectTranslationCS{fintivý, parádivý, frajerský} \dicExampleIS{vera spjátrungslegur í klæðaburði} \dicExampleCS{být fintivý v~oblékání}
\dicEntry[spjátrungur] \dicTerm{spjátr··ung|ur} \dicIPA{{s}{\textsubring{b}}{j}{au}{\textlengthmark}{\textsubring{d}}{r}{u}{\ng}{\r{g}}{\textscy}{\textsubring{r}}} \dicPos{m}[6] \dicFlx{(‑s, ‑ar)}[8] \dicDirectTranslationCS{frajer, floutek, hejsek, metrosexuál}
\dicEntry[spjó] \dicTerm{spjó} \dicIPA{{s}{\textsubring{b}}{j}{ou}{\textlengthmark}} \dicPos{v} \dicFlx{ind pf sg 1 pers} \dicLink{spýja\smash{\textsuperscript{2}}}
\dicEntry[spjót] \dicTerm{spjót} \dicIPA{{s}{\textsubring{b}}{j}{ou}{\textlengthmark}{\textsubring{d}}} \dicPos{n}[2] \dicFlx{(‑s, ‑)}[5] \textbf{1.} \dicDirectTranslationCS{oštěp, kopí}  \textbf{2.} \dicFieldCat{sport.} \dicDirectTranslationCS{oštěp};  \dicIdiom{spjót}{ \dicPhraseIS{beina spjótum sínum að e‑m}} \dicLangCat{přen.} \dicDirectTranslationCS{mít namířeno na (koho), zaměřit se (kriticky) na (koho)}; { \dicPhraseIS{hafa öll spjót úti til að (gera e‑ð)}} \dicLangCat{přen.} \dicDirectTranslationCS{dělat všechno možné pro (udělání (čeho)) (k~dosáhnutí určité věci ap.)}; { \dicPhraseIS{það standa á e‑m öll spjót}} \dicFlx{impers} \dicLangCat{přen.} \dicDirectTranslationCS{(kdo) je pod velkým tlakem}
\dicEntry[spjótkast] \dicTerm{spjót··kast} \dicIPA{{s}{\textsubring{b}}{j}{ou}{\textlengthmark}{\textsubring{d}}{k\smash{\textsuperscript{h}}}{a}{s}{\textsubring{d}}} \dicPos{n}[2] \dicFlx{(‑s)}[2] \dicFieldCat{sport.} \dicDirectTranslationCS{hod oštěpem}
\dicEntry[spjöld] \dicTerm{spjöld} \dicIPA{{s}{\textsubring{b}}{j}{\oe}{l}{\textsubring{d}}} \dicPos{n} \dicFlx{pl nom} \dicLink{spjald}
\dicEntry[spjöll] \dicTerm{spjöll} \dicIPA{{s}{\textsubring{b}}{j}{\oe}{\textsubring{d}}{\textsubring{l}}} \dicPos{n}[2] \dicFlx{pl}[9] \dicLink{spjall}
\dicEntry[spjör] \dicTerm{spjör} \dicIPA{{s}{\textsubring{b}}{j}{\oe}{\textlengthmark}{\textsubring{r}}} \dicPos{f}[7] \dicFlx{(spjarar, spjarir)}[16] \dicSynonym{fat} \dicDirectTranslationCS{(kus) oblečení, oděv};  \dicPhraseIS{fara úr hverri spjör} \dicSynonym*{berhátta} \dicDirectTranslationCS{sundat si všechno, svléknout se do naha};  \dicPhraseIS{spyrja e‑n spjörunum úr} \dicDirectTranslationCS{vyptat se (koho), vyzpovídat (koho)}
\dicEntry[splitt] \dicTerm{splitt} \dicIPA{{s}{\textsubring{b}}{l}{\textsci}{h}{\textsubring{d}}} \dicPos{n}[2] \dicFlx{(‑s, ‑)}[5] \dicLink{splitti}
\dicEntry[splitti] \dicTerm{splitti}\dicTerm{, splitt} \dicIPA{{s}\-{\textsubring{b}}\-{l}\-{\textsci}\-{h}\-{\textsubring{d}}\-{\textsci}\-} \dicPos{n}[2] \dicFlx{(‑s, ‑)}[14] \dicDirectTranslationCS{závlačka}
\dicEntry[splundra] \dicTerm{splundr|a} \dicIPA{{s}{\textsubring{b}}{l}{\textscy}{n}{\textsubring{d}}{r}{a}} \dicPos{v}[1] \dicFlx{(‑aði)}[1] \dicFlx{dat} \textbf{1.} \dicSynonym{sóa} \dicDirectTranslationCS{(pro)marnit, (pro)mrhat}  \textbf{2.} \dicSynonym*{sprengja í smátt} \dicDirectTranslationCS{roztříštit, rozbít na kusy};  \dicIdiom{splundrast}{ \dicPhraseIS{splundrast}} \dicFlx{refl} \dicDirectTranslationCS{roztříštit se, rozbít se na kusy}
\dicEntry[splunkunýr] \dicTerm{splunku··nýr} \dicIPA{{s}{\textsubring{b}}{l}{u}{\r{\ng}}{\r{g}}{\textscy}{n}{i}{\textsubring{r}}} \dicPos{adj}[4]\dicFlx{}[-8] \dicSynonym{glænýr} \dicDirectTranslationCS{zbrusu\,/\addthin zcela nový}
\dicEntry[splæsa] \dicTerm{splæs|a} \dicIPA{{s}{\textsubring{b}}{l}{a}{i}{\textlengthmark}{s}{a}} \dicPos{v}[2] \dicFlx{(‑ti, ‑t)}[63] \dicFlx{acc\,/\addthin dat} \textbf{1.} \dicFlx{acc} \dicDirectTranslationCS{splést, splétat} \dicIndirectTranslationCS{(pletením spojit)} \dicExampleIS{splæsa kaðal} \dicExampleCS{splést provaz}  \textbf{2.} \dicFlx{dat} \dicLangCat{hovor.} \dicDirectTranslationCS{zacvaknout, zacvakat, zaplatit (za všechny v~restauraci ap.)} \dicExampleIS{Komdu með, ég splæsi!} \dicExampleCS{Připoj se, já to platím!};  \dicPhraseIS{splæsa e‑u á e‑n} \dicLangCat{hovor.} \dicDirectTranslationCS{zacvaknout\,/\addthin zacvakat (co) za (koho)} \dicExampleIS{Hann splæsti á mig öðru glasi.} \dicExampleCS{Zacvaknul za mě druhou skleničku.};  \dicPhraseIS{splæsa í e‑ð} \dicLangCat{hovor.} \dicDirectTranslationCS{vypláznout za (co) peníze (za něco bezcenného ap.)}
\dicEntry[spor] \dicTerm{spor} \dicsymFrequent\  \dicIPA{{s}{\textsubring{b}}{\textopeno}{\textlengthmark}{\textsubring{r}}} \dicPos{n}[2] \dicFlx{(‑s, ‑)}[5] \textbf{1.} \dicSynonym{fótspor} \dicDirectTranslationCS{stopa, šlápota, šlépěj} \dicIndirectTranslationCS{(otisk chodidla)} \dicExampleIS{spor á gólfi} \dicExampleCS{stopa na podlaze};  \dicPhraseIS{komast varla\,/\addthin ekki úr sporunum} \dicFlx{refl} \dicLangCat{přen.} \dicDirectTranslationCS{nehnout se z~místa (s~psaním práce ap.)};  \dicPhraseIS{standa í sporum e‑rs} \dicLangCat{přen.} \dicDirectTranslationCS{být ve stejné situaci jako (kdo)};  \dicPhraseIS{standa í sömu sporum} \dicLangCat{přen.} \dicDirectTranslationCS{přešlapovat na místě, stagnovat}  \textbf{2.} \dicSynonym{skref} \dicDirectTranslationCS{krok} \dicExampleIS{stíga fáein spor} \dicExampleCS{popojít několik kroků};  \dicPhraseIS{þetta er spor í rétta átt} \dicLangCat{přen.} \dicDirectTranslationCS{to je krok správným směrem}  \textbf{3.} \dicSynonym{dansspor} \dicDirectTranslationCS{(taneční) krok};  \dicPhraseIS{ná sporinu} \dicDirectTranslationCS{chytit krok (v~tanci ap.)}  \textbf{4.} \dicDirectTranslationCS{krok, chůze} \dicExampleIS{léttur í spori} \dicExampleCS{mající lehký krok};  \dicPhraseIS{greikka sporið} \dicDirectTranslationCS{přidat do kroku};  \dicPhraseIS{spretta úr spori} \dicDirectTranslationCS{rozběhnout se}  \textbf{5.} \dicDirectTranslationCS{krok} \dicIndirectTranslationCS{(způsob chůze koně)}  \textbf{6.} \dicSynonym*{saumspor} \dicDirectTranslationCS{steh}  \textbf{7.} \dicSynonym*{brautarteinar} \dicDirectTranslationCS{dráha, koleje} \dicExampleIS{járnbrautarspor} \dicExampleCS{železniční koleje};  \dicPhraseIS{fara út af sporinu} \dicDirectTranslationCS{rozhodit, vykolejit}
\dicEntry[spora] \dicTerm{spor|a} \dicIPA{{s}{\textsubring{b}}{\textopeno}{\textlengthmark}{r}{a}} \dicPos{v}[1] \dicFlx{(‑aði)}[1] \dicFlx{acc} \dicDirectTranslationCS{(po)šlapat, zanechat stopu} \dicExampleIS{spora gólfið} \dicExampleCS{pošlapat podlahu};  \dicIdiom{sporast}{ \dicPhraseIS{sporast}} \dicFlx{refl} \dicDirectTranslationCS{pošlapat se}
\dicEntry[sporbaugur] \dicTerm{spor··baug|ur} \dicIPA{{s}{\textsubring{b}}{\textopeno}{r}{\textsubring{b}}{\oe i}{\textbabygamma}{\textscy}{\textsubring{r}}} \dicPos{m}[6] \dicFlx{(‑s, ‑ar)}[24] \dicFieldCat{mat.} \dicDirectTranslationCS{elipsa}
\dicEntry[sporbraut] \dicTerm{spor··braut} \dicIPA{{s}{\textsubring{b}}{\textopeno}{r}{\textsubring{b}}{r}{\oe i}{\textsubring{d}}} \dicPos{f}[7] \dicFlx{(‑ar, ‑ir)}[1] \textbf{1.} \dicFieldCat{astro.} \dicDirectTranslationCS{oběžná dráha, orbita};  \dicPhraseIS{sporbraut jarðarinnar} \dicFieldCat{astro.} \dicDirectTranslationCS{oběžná dráha Země}  \textbf{2.} \dicDirectTranslationCS{trať, dráha}
\dicEntry[sporðdreki] \dicTerm{sporð··drek|i} \dicIPA{{s}{\textsubring{b}}{\textopeno}{r}{ð}{\textsubring{d}}{r}{\textepsilon}{\r{\textObardotlessj}}{\textsci}} \dicPos{m}[1] \dicFlx{(‑a, ‑ar)}[1] \textbf{1.} \dicFieldCat{zool.} \dicDirectTranslationCS{štír, škorpion} \textit{(l.~{\textLA{Scorpiones}})}  \dicsymPhoto\   \textbf{2.} \dicDirectTranslationCS{Štír} \dicIndirectTranslationCS{(znamení zvěrokruhu)} \textit{(l.~{\textLA{Scorpius}})}
\dicFigure{107211.jpg}{Sporðdreki}{Sporðdreki - Ševčík Jan, Biolib, Copyright/PD}
\dicEntry[sporður] \dicTerm{sporð|ur} \dicIPA{{s}{\textsubring{b}}{\textopeno}{r}{ð}{\textscy}{\textsubring{r}}} \dicPos{m}[6] \dicFlx{(‑s, ‑ar)}[4] \dicDirectTranslationCS{ocas (ryby, hada ap.)};  \dicIdiom{sporður}{ \dicPhraseIS{standa e‑m á sporði}} \dicDirectTranslationCS{mít na (koho), vyrovnat se (komu)};  \dicPhraseIS{þekkja\,/\addthin vita hvorki haus\,/\addthin höfuð né sporð á e‑u} \dicLangCat{přen.} \dicDirectTranslationCS{nemít o~(čem) ponětí\,/\addthin páru}
\dicEntry[sporgöngumaður] \textls[10]{\dicTerm{spor·göngu··|maður} \dicIPA{{s}{\textsubring{b}}{\textopeno}{r}{\r{g}}{\oe i}{\ng}{\r{g}}{\textscy}{m}{a}{ð}{\textscy}{\textsubring{r}}} \dicPos{m}[13] \dicFlx{(‑manns, ‑menn)}[2]} \textbf{1.} \dicSynonym{eftirmaður} \dicDirectTranslationCS{následník, následnice, nástupce, nástupkyně}  \textbf{2.} \dicSynonym{fylgismaður} \dicDirectTranslationCS{následovník, následovnice, stoupenec, stoupenkyně}
\dicEntry[sporhundur] \dicTerm{spor··hund|ur} \dicIPA{{s}{\textsubring{b}}{\textopeno}{\textlengthmark}{\textsubring{r}}{h}{\textscy}{n}{\textsubring{d}}{\textscy}{\textsubring{r}}} \dicPos{m}[6] \dicFlx{(‑s, ‑ar)}[4] \dicDirectTranslationCS{stopařský pes}
\dicEntry[spori] \dicTerm{spor|i} \dicIPA{{s}{\textsubring{b}}{\textopeno}{\textlengthmark}{r}{\textsci}} \dicPos{m}[1] \dicFlx{(‑a, ‑ar)}[1] \textbf{1.} \dicDirectTranslationCS{ostruha} \dicIndirectTranslationCS{(kovový bodec na opatku jezdecké boty sloužící k~pobízení koně)}  \textbf{2.} \dicFieldCat{biol.} \dicDirectTranslationCS{ostruha} \dicIndirectTranslationCS{(zvláště v~pl)}
\dicEntry[sporjárn] \dicTerm{spor··járn} \dicIPA{{s}{\textsubring{b}}{\textopeno}{r}{j}{au}{r}{\textsubring{d}}{\textsubring{n}}} \dicPos{n}[2] \dicFlx{(‑s, ‑)}[5] \dicDirectTranslationCS{dláto}
\dicEntry[sporna] \dicTerm{sporn|a} \dicIPA{{s}{\textsubring{b}}{\textopeno}{r}{\textsubring{d}}{n}{a}} \dicPos{v}[1] \dicFlx{(‑aði)}[1] \dicPhraseIS{sporna við\,/\addthin gegn e‑u} \dicDirectTranslationCS{čelit (čemu), vystoupit proti (čemu)} \dicExampleIS{sporna gegn glæpum} \dicExampleCS{čelit zločinům}
\dicEntry[sport] \dicTerm{sport} \dicIPA{{s}{\textsubring{b}}{\textopeno}{\textsubring{r}}{\textsubring{d}}} \dicPos{n}[2] \dicFlx{(‑s)}[2] \textbf{1.} \dicSynonym*{íþróttir} \dicDirectTranslationCS{sport} \dicExampleIS{stunda sport} \dicExampleCS{sportovat}  \textbf{2.} \dicSynonym{skemmtun} \dicDirectTranslationCS{zábava, legrace} \dicExampleIS{eiga hesta upp á sport} \dicExampleCS{mít koně pro zábavu}
\dicEntry[sportbíll] \dicTerm{sport··bíl|l} \dicIPA{{s}{\textsubring{b}}{\textopeno}{\textsubring{r}}{\textsubring{d}}{\textsubring{b}}{i}{\textsubring{d}}{\textsubring{l}}} \dicPos{m}[6] \dicFlx{(‑s, ‑ar)}[48] \dicDirectTranslationCS{sportovní auto}
\dicEntry[sportfatnaður] \dicTerm{sport··fat·nað|ur} \dicIPA{{s}{\textsubring{b}}{\textopeno}{\textsubring{r}}{\textsubring{d}}{f}{a}{\textsubring{d}}{n}{a}{ð}{\textscy}{\textsubring{r}}} \dicPos{m}[10] \dicFlx{(‑ar)}[9] \dicDirectTranslationCS{sportovní oblečení}
\dicEntry[sportittlingur] \dicTerm{spor··titt·ling|ur} \dicIPA{{s}{\textsubring{b}}{\textopeno}{\textsubring{r}}{t\smash{\textsuperscript{h}}}{\textsci}{h}{\textsubring{d}}{l}{i}{\ng}{\r{g}}{\textscy}{\textsubring{r}}} \dicPos{m}[6] \dicFlx{(‑s, ‑ar)}[8] \dicFieldCat{zool.} \dicDirectTranslationCS{strnad severní} \textit{(l.~{\textLA{Calcarius lapponicus}})}  \dicsymPhoto\ 
\dicFigure{ds_image_sportittlingur_0_1.jpg}{Sportittlingur}{Sportittlingur - qmnonic at Flickr, CC BY 2.0}
\dicEntry[sportvöruverslun] \dicTerm{sport·vöru··versl|un} \dicIPA{{s}{\textsubring{b}}{\textopeno}{\textsubring{r}}{\textsubring{d}}{v}{\oe}{r}{\textscy}{v}{\textepsilon}{\textsubring{r}}{s}{\textsubring{d}}{l}{\textscy}{\textsubring{n}}} \dicPos{f}[7] \dicFlx{(‑unar, ‑anir)}[8] \dicDirectTranslationCS{obchod se sportovním zbožím}
\dicEntry[sporvagn] \dicTerm{spor··vagn} \dicIPA{{s}{\textsubring{b}}{\textopeno}{r}{v}{a}{\r{g}}{\textsubring{n}}} \dicPos{m}[4] \dicFlx{(‑s, ‑ar)}[4] \dicDirectTranslationCS{tramvaj}
\dicEntry[sporöskjulaga] \dicTerm{spor·öskju··laga} \dicIPA{{s}{\textsubring{b}}{\textopeno}{\textlengthmark}{r}{\oe}{s}{\r{\textObardotlessj}}{\textscy}{l}{a}{\textbabygamma}{a}} \dicPos{adj}[13] \dicFlx{indecl}[1] \dicDirectTranslationCS{oválný, elipsovitý}
\dicEntry[sposkur] \dicTerm{sposkur} \dicIPA{{s}{\textsubring{b}}{\textopeno}{s}{\r{g}}{\textscy}{\textsubring{r}}} \dicPos{adj}[1]\dicFlx{}[-1] \dicSynonym{hæðinn} \dicDirectTranslationCS{posměvačný, sarkastický, ironický}
\dicEntry[spott] \dicTerm{spott} \dicIPA{{s}{\textsubring{b}}{\textopeno}{h}{\textsubring{d}}} \dicPos{n}[2] \dicFlx{(‑s)}[2] \dicSynonym{háð\smash{\textsuperscript{1}}} \dicDirectTranslationCS{posměšek, úšklebek} \dicExampleIS{hafa e‑n að spotti} \dicExampleCS{posmívat se (komu)}
\dicEntry[spotta] \dicTerm{spott|a} \dicIPA{{s}{\textsubring{b}}{\textopeno}{h}{\textsubring{d}}{a}} \dicPos{v}[1] \dicFlx{(‑aði)}[1] \dicFlx{acc} \dicDirectTranslationCS{posmívat se, vysmívat se, zesměšňovat} \dicExampleIS{spotta e‑n} \dicExampleCS{vysmívat se (komu)}
\dicEntry[spotti] \dicTerm{spott|i} \dicIPA{{s}{\textsubring{b}}{\textopeno}{h}{\textsubring{d}}{\textsci}} \dicPos{m}[1] \dicFlx{(‑a, ‑ar)}[1] \textbf{1.} \dicSynonym*{stuttur þráður} \dicDirectTranslationCS{krátký provázek}  \textbf{2.} \dicSynonym*{vegarspotti} \dicDirectTranslationCS{kousek, krátká vzdálenost}
\dicEntry[spói] \dicTerm{spó|i} \dicIPA{{s}{\textsubring{b}}{ou}{\textlengthmark}{\textsci}} \dicPos{m}[1] \dicFlx{(‑a, ‑ar)}[1] \dicFieldCat{zool.} \dicDirectTranslationCS{koliha, koliha malá} \textit{(l.~{\textLA{Numenius phaeopus}})}  \dicsymPhoto\ 
\dicFigure{ds_image_spoi_0_1.jpg}{Spói}{Spói - Andreas Trepte, CC BY-SA 2.5}
\dicEntry[spóka] \dicTerm{spók|a} \dicIPA{{s}{\textsubring{b}}{ou}{\textlengthmark}{\r{g}}{a}} \dicPos{v}[1] \dicFlx{(‑aði)}[1] \dicFlx{acc} \dicPhraseIS{spóka sig} \dicDirectTranslationCS{promenádovat se, korzovat, procházet se} \dicExampleIS{spóka sig í góða veðrinu} \dicExampleCS{promenádovat se v~dobrém počasí}
\dicEntry[spóla] \dicTerm{spól|a\smash{\textsuperscript{1}}} \dicIPA{{s}{\textsubring{b}}{ou}{\textlengthmark}{l}{a}} \dicPos{f}[1] \dicFlx{(‑u, ‑ur)}[13] \textbf{1.} \dicSynonym{kefli} \dicDirectTranslationCS{cívka, špulka, kotouč (nití ap.)}  \textbf{2.} \dicSynonym{kassetta} \dicDirectTranslationCS{kazeta, videokazeta, audiokazeta} \dicExampleIS{leigja spólu} \dicExampleCS{půjčit si kazetu}  \textbf{3.} \dicSynonym*{kvikmyndafilma} \dicDirectTranslationCS{(filmový) kotouč}
\dicEntry[spóla] \dicTerm{spól|a\smash{\textsuperscript{2}}} \dicIPA{{s}{\textsubring{b}}{ou}{\textlengthmark}{l}{a}} \dicPos{v}[1] \dicFlx{(‑aði)}[1] \textbf{1.} \dicFlx{acc} \dicSynonym*{vinda upp} \dicDirectTranslationCS{převinout, převíjet, přetočit, přetáčet (kazetu ap.)};  \dicPhraseIS{spóla áfram} \dicDirectTranslationCS{převinout dopředu};  \dicPhraseIS{spóla til baka} \dicDirectTranslationCS{převinout dozadu}  \textbf{2.} \dicDirectTranslationCS{protáčet se (kola u~auta ap.)}
\dicEntry[spónlagning] \dicTerm{spón··lag·ning} \dicIPA{{s}{\textsubring{b}}{ou}{n}{l}{a}{\r{g}}{n}{i}{\ng}{\r{g}}} \dicPos{f}[4] \dicFlx{(‑ar)}[7] \dicDirectTranslationCS{dýhování}
\dicEntry[spónleggja] \dicTerm{spón··|leggja} \dicIPA{{s}{\textsubring{b}}{ou}{n}{l}{\textepsilon}{\r{\textObardotlessj}}{a}} \dicPos{v}[4] \dicFlx{(‑legg, ‑lagði, ‑lögðum, ‑legði, ‑lagt)}[19] \dicFlx{acc} \dicDirectTranslationCS{dýhovat}
\dicEntry[spónn] \dicTerm{spónn} \dicIPA{{s}{\textsubring{b}}{ou}{\textsubring{d}}{\textsubring{n}}} \dicPos{m}[6] \dicFlx{(spóns\,/\addthin spónar, spænir\,/\addthin spónar)}[77] \textbf{1.} \dicSynonym*{hefilspónn} \dicDirectTranslationCS{hoblina}  \textbf{2.} \dicSynonym*{viðarþynna} \dicDirectTranslationCS{dýha} \dicExampleIS{spónn úr rauðaviði} \dicExampleCS{dýha z~červeného dřeva}  \textbf{3.} \dicSynonym{matskeið} \dicDirectTranslationCS{lžíce} \dicIndirectTranslationCS{(ze dřeva nebo rohu)};  \dicIdiom{spónn}{ \dicPhraseIS{missa spón úr aski sínum}} \dicLangCat{přen.} \dicDirectTranslationCS{přijít o~výhody (o~výnosné místo, pozici ap.)}
\dicEntry[spraka] \dicTerm{spraka} \dicIPA{{s}{\textsubring{b}}{r}{a}{\textlengthmark}{\r{g}}{a}} \dicPos{f}[1] \dicFlx{(spröku, sprökur)}[8] \dicFieldCat{zool.} \dicSynonym{lúða} \dicDirectTranslationCS{halibut, platýs obecný} \textit{(l.~{\textLA{Hippoglossus vulgaris}})}
\dicEntry[sprakk] \dicTerm{sprakk} \dicIPA{{s}{\textsubring{b}}{r}{a}{h}{\r{g}}} \dicPos{v} \dicFlx{ind pf sg 1 pers} \dicLink{springa}
\dicEntry[spranga] \dicTerm{sprang|a} \dicIPA{{s}{\textsubring{b}}{r}{au}{\ng}{\r{g}}{a}} \dicPos{v}[1] \dicFlx{(‑aði)}[13] \dicDirectTranslationCS{kráčet, vykračovat si} \dicExampleIS{spranga um göturnar} \dicExampleCS{kráčet po ulicích}
\dicEntry[spratt] \dicTerm{spratt} \dicIPA{{s}{\textsubring{b}}{r}{a}{h}{\textsubring{d}}} \dicPos{v} \dicFlx{ind pf sg 1 pers} \dicLink{spretta\smash{\textsuperscript{2}}}
\dicEntry[sprauta] \dicTerm{spraut|a\smash{\textsuperscript{1}}} \dicIPA{{s}{\textsubring{b}}{r}{\oe i}{\textlengthmark}{\textsubring{d}}{a}} \dicPos{f}[1] \dicFlx{(‑u, ‑ur)}[13] \textbf{1.} \dicDirectTranslationCS{(injekční) stříkačka}  \textbf{2.} \dicSynonym*{lyfjagjöf} \dicDirectTranslationCS{injekce, vstříknutí (léku)}
\dicEntry[sprauta] \dicTerm{spraut|a\smash{\textsuperscript{2}}} \dicIPA{{s}{\textsubring{b}}{r}{\oe i}{\textlengthmark}{\textsubring{d}}{a}} \dicPos{v}[1] \dicFlx{(‑aði)}[1] \dicFlx{dat\,/\addthin acc} \textbf{1.} \dicFlx{dat} \dicDirectTranslationCS{(po)stříkat, nastříkat, rozstřikovat} \dicExampleIS{sprauta rjóma á kökuna} \dicExampleCS{nastříkat šlehačku na dort}  \textbf{2.} \dicFlx{acc} \dicDirectTranslationCS{dát\,/\addthin vstříknout injekci, očkovat} \dicExampleIS{sprauta sjúklinginn með pensilíni} \dicExampleCS{vstříknout pacientovi injekci s~penicilínem};  \dicPhraseIS{sprauta sig} {\textbf{a.}} \dicDirectTranslationCS{vstříknout si injekci};  {\textbf{b.}} \dicDirectTranslationCS{vstříknout si dávku (heroinu ap.)};  \dicIdiom{sprautast}{ \dicPhraseIS{sprautast}} \dicFlx{refl} \dicDirectTranslationCS{(vy)stříknout, stříkat}
\dicEntry[sprek] \dicTerm{sprek} \dicIPA{{s}{\textsubring{b}}{r}{\textepsilon}{\textlengthmark}{\r{g}}} \dicPos{n}[2] \dicFlx{(‑s, ‑)}[5] \dicDirectTranslationCS{roští, chrastí, klestí}
\dicEntry[sprell] \dicTerm{sprell} \dicIPA{{s}{\textsubring{b}}{r}{\textepsilon}{l}{\textlengthmark}} \dicPos{n}[2] \dicFlx{(‑s)}[2] \dicSynonym{ærsl} \dicDirectTranslationCS{rošťárna, vylomenina}
\dicEntry[sprella] \dicTerm{sprell|a} \dicIPA{{s}{\textsubring{b}}{r}{\textepsilon}{l}{\textlengthmark}{a}} \dicPos{v}[1] \dicFlx{(‑aði)}[1] \dicDirectTranslationCS{šaškovat, dovádět, blbnout}
\dicEntry[sprellikarl] \dicTerm{sprelli··karl} \dicIPA{{s}{\textsubring{b}}{r}{\textepsilon}{l}{\textlengthmark}{\textsci}{k\smash{\textsuperscript{h}}}{a}{r}{\textsubring{d}}{\textsubring{l}}} \dicPos{m}[4] \dicFlx{(‑s, ‑ar)}[4] \textbf{1.} \dicDirectTranslationCS{tahací panák} \dicIndirectTranslationCS{(hračka)}  \textbf{2.} \dicSynonym*{æringi} \dicDirectTranslationCS{rošťák, rošťačka}
\dicEntry[sprelllifandi] \dicTerm{sprell··lif·andi} \dicIPA{{s}{\textsubring{b}}{r}{\textepsilon}{l}{\textlengthmark}{\textsci}{v}{a}{n}{\textsubring{d}}{\textsci}} \dicPos{adj}[13] \dicFlx{indecl}[1] \dicDirectTranslationCS{(úplně) živý, mající se čile k~světu} \dicExampleIS{Kötturinn kom inn með sprelllifandi fugl.} \dicExampleCS{Kočka přišla s~úplně živým ptáčkem.}
\dicEntry[sprengfullur] \dicTerm{spreng··fullur} \dicIPA{{s}{\textsubring{b}}{r}{ei}{\ng}{\r{g}}{f}{\textscy}{\textsubring{d}}{l}{\textscy}{\textsubring{r}}} \dicPos{adj}[10]\dicFlx{}[-8] \dicSynonym{troðfullur} \dicDirectTranslationCS{nar\-vaný, přecpaný (autobus ap.)}
\dicEntry[sprengidagur] \dicTerm{sprengi··dag|ur} \dicIPA{{s}{\textsubring{b}}{r}{ei}{\textltailn}{\r{\textObardotlessj}}{\textsci}{\textsubring{d}}{a}{\textbabygamma}{\textscy}{\textsubring{r}}} \dicPos{m}[6] \dicFlx{(‑s, ‑ar)}[62] \dicDirectTranslationCS{masopustní úterý}
\dicEntry[sprengiefni] \dicTerm{sprengi··efni} \dicIPA{{s}{\textsubring{b}}{r}{ei}{\textltailn}{\r{\textObardotlessj}}{\textsci}{\textepsilon}{\textsubring{b}}{n}{\textsci}} \dicPos{n}[2] \dicFlx{(‑s, ‑)}[14] \dicDirectTranslationCS{výbušnina, třaskavina}
\dicEntry[sprengigígur] \dicTerm{sprengi··gíg|ur} \dicIPA{{s}{\textsubring{b}}{r}{ei}{\textltailn}{\r{\textObardotlessj}}{\textsci}{\r{\textObardotlessj}}{i}{\textbabygamma}{\textscy}{\textsubring{r}}} \dicPos{m}[6] \dicFlx{(‑s, ‑ar\,/\addthin ‑ir)}[70] \dicFieldCat{geol.} \dicDirectTranslationCS{maar, explozivní kráter}
\dicEntry[sprengihreyfill] \dicTerm{sprengi··hreyf|ill} \dicIPA{{s}{\textsubring{b}}{r}{ei}{\textltailn}{\r{\textObardotlessj}}{\textsci}{\textsubring{r}}{ei}{v}{\textsci}{\textsubring{d}}{\textsubring{l}}} \dicPos{m}[6] \dicFlx{(‑ils, ‑lar)}[35] \dicDirectTranslationCS{spalovací motor}
\dicEntry[sprengihætta] \dicTerm{sprengi··hætt|a} \dicIPA{{s}{\textsubring{b}}{r}{ei}{\textltailn}{\r{\textObardotlessj}}{\textsci}{h}{a}{i}{h}{\textsubring{d}}{a}} \dicPos{f}[1] \dicFlx{(‑u)}[5] \dicDirectTranslationCS{nebezpečí výbuchu}
\dicEntry[sprengikúla] \dicTerm{sprengi··kúl|a} \dicIPA{{s}{\textsubring{b}}{r}{ei}{\textltailn}{\r{\textObardotlessj}}{\textsci}{k\smash{\textsuperscript{h}}}{u}{l}{a}} \dicPos{f}[1] \dicFlx{(‑u, ‑ur)}[13] \dicDirectTranslationCS{(dělová) koule, (dělostřelecký) granát}
\dicEntry[sprenging] \dicTerm{spreng··ing} \dicsymFrequent\  \dicIPA{{s}{\textsubring{b}}{r}{ei}{\textltailn}{\r{\textObardotlessj}}{i}{\ng}{\r{g}}} \dicPos{f}[4] \dicFlx{(‑ar, ‑ar)}[5] \textbf{1.} \dicSynonym*{það að e‑að springur} \dicDirectTranslationCS{výbuch, exploze} \dicExampleIS{öflug sprenging} \dicExampleCS{mohutný výbuch}  \textbf{2.} \dicLangCat{přen.} \dicSynonym{aukning} \dicDirectTranslationCS{exploze, boom, (náhlý) nárůst}
\dicEntry[sprengja] \dicTerm{sprengj|a\smash{\textsuperscript{1}}} \dicIPA{{s}{\textsubring{b}}{r}{ei}{\textltailn}{\r{\textObardotlessj}}{a}} \dicPos{f}[1] \dicFlx{(‑u, ‑ur)}[25] \dicDirectTranslationCS{bomba, puma, mina}
\dicEntry[sprengja] \dicTerm{spreng|ja\smash{\textsuperscript{2}}} \dicsymFrequent\  \dicIPA{{s}{\textsubring{b}}{r}{ei}{\textltailn}{\r{\textObardotlessj}}{a}} \dicPos{v}[2] \dicFlx{(‑di, ‑t)}[137] \dicFlx{acc} \textbf{1.} \dicDirectTranslationCS{odpálit, odpalovat, detonovat, vyhodit do vzduchu\,/\addthin povětří} \dicExampleIS{sprengja tímasprengju} \dicExampleCS{odpálit časovanou bombu};  \dicPhraseIS{sprengja e‑ð upp} \dicDirectTranslationCS{vyhodit (co) do vzduchu\,/\addthin povětří (bombou ap.)} \dicExampleIS{sprengja upp hurð} \dicExampleCS{vyhodit dveře do povětří}  \textbf{2.} \dicDirectTranslationCS{bouchnout, prásknout (balónek ap.)};  \dicIdiom{sprengja}[utan af]{ \dicPhraseIS{sprengja e‑ð utan af sér}} \dicDirectTranslationCS{vyrůst z~(čeho) (kvůli nadváze ap.)}
\dicEntry[sprengjuárás] \dicTerm{sprengju··á·rás} \dicIPA{{s}{\textsubring{b}}{r}{ei}{\textltailn}{\r{\textObardotlessj}}{\textscy}{au}{r}{au}{s}} \dicPos{f}[7] \dicFlx{(‑ar, ‑ir)}[1] \dicDirectTranslationCS{bombový\,/\addthin pumový útok}
\dicEntry[sprengjuflugvél] \dicTerm{sprengju··flug·vél} \dicIPA{{s}{\textsubring{b}}{r}{ei}{\textltailn}{\r{\textObardotlessj}}{\textscy}{f}{l}{\textscy}{\textbabygamma}{v}{j}{\textepsilon}{\textsubring{l}}} \dicPos{f}[4] \dicFlx{(‑ar, ‑ar)}[1] \dicFieldCat{voj.} \dicDirectTranslationCS{bombardér}
\begin{xtolerant}{}{1pt}
\dicEntry[sprenglærður] \dicTerm{spreng··lærður}\addthinS\dicIPA{{s}{\textsubring{b}}{r}{ei}{\ng}{\r{g}}{l}{a}{i}{r}{ð}{\textscy}{\textsubring{r}}}\addthinS\dicPos{adj}[2]\dicFlx{}[-4]\addthinS\dicDirectTranslationCS{erudovaný, učený, vzdělaný}
\end{xtolerant}
\dicEntry[sprengmóður] \dicTerm{spreng··móður} \dicIPA{{s}{\textsubring{b}}{r}{ei}{\ng}{\r{g}}{m}{ou}{ð}{\textscy}{\textsubring{r}}} \dicPos{adj}[2]\dicFlx{}[-10] \dicDirectTranslationCS{udýchaný, popadající dech}
\dicEntry[sprengur] \dicTerm{spreng|ur} \dicIPA{{s}{\textsubring{b}}{r}{ei}{\ng}{\r{g}}{\textscy}{\textsubring{r}}} \dicPos{m}[9] \dicFlx{(‑s, ‑ir)}[15] \dicPhraseIS{vera kominn í spreng} \dicLangCat{přen.} \dicDirectTranslationCS{potřebovat nutně na malou};  \dicPhraseIS{í einum spreng} \dicFlx{adv} \dicDirectTranslationCS{najednou, na jeden zátah} \dicExampleIS{hlaupa heim í einum spreng} \dicExampleCS{doběhnout domů na jeden zátah}
\dicEntry[sprett] \dicTerm{sprett} \dicIPA{{s}{\textsubring{b}}{r}{\textepsilon}{h}{\textsubring{d}}} \dicPos{v} \dicFlx{ind praes sg 1 pers} \dicLink{spretta\smash{\textsuperscript{2}}}
\dicEntry[spretta] \dicTerm{sprett|a\smash{\textsuperscript{1}}} \dicIPA{{s}{\textsubring{b}}{r}{\textepsilon}{h}{\textsubring{d}}{a}} \dicPos{v}[2] \dicFlx{(‑i, ‑)}[4] \dicIdiom{spretta}[af]{ \dicPhraseIS{spretta af hestinum}} \dicDirectTranslationCS{odsedlat koně};  \dicIdiom{spretta}[upp]{ \dicPhraseIS{spretta upp e‑u}} \dicDirectTranslationCS{(od)párat (co), rozpárat (co)} \dicExampleIS{spretta upp saumnum} \dicExampleCS{rozpárat šev};  \dicIdiom{spretta}[úr]{ \dicPhraseIS{spretta úr spori}} \dicDirectTranslationCS{vystartovat, rozběhnout se}
\dicEntry[spretta] \dicTerm{spretta\smash{\textsuperscript{2}}} \dicsymFrequent\  \dicIPA{{s}{\textsubring{b}}{r}{\textepsilon}{h}{\textsubring{d}}{a}} \dicPos{v}[6] \dicFlx{(sprett, spratt, spruttum, sprytti, sprottið)}[29] \textbf{1.} \dicSynonym{vaxa} \dicDirectTranslationCS{(vy)rašit, pučet, vyrážet} \dicExampleIS{Grasið sprettur vel í hlýindunum.} \dicExampleCS{Tráva dobře raší v~teplém počasí.};  \dicPhraseIS{e‑að er sprottið úr sér} \dicDirectTranslationCS{(co) je přerostlé (tráva ap.)}  \textbf{2.} \dicDirectTranslationCS{vyskočit, vyrazit};  \dicPhraseIS{spretta á fætur} \dicDirectTranslationCS{vyskočit na nohy}  \textbf{3.} \dicDirectTranslationCS{pocházet, mít svůj původ};  \dicIdiom{spretta}[fram]{ \dicPhraseIS{spretta fram}} \dicDirectTranslationCS{vynořit se, vyrazit (pot ap.)} \dicExampleIS{Svitinn spratt fram á enni hans.} \dicExampleCS{Na čele mu vyrazil pot.};  \dicIdiom{spretta}[upp]{ \dicPhraseIS{spretta upp}} \dicDirectTranslationCS{vynořit se, zjevit se, objevit se}
\dicEntry[spretthlaup] \dicTerm{sprett··hlaup} \dicIPA{{s}{\textsubring{b}}{r}{\textepsilon}{h}{\textsubring{d}}{\textsubring{l}}{\oe i}{\textsubring{b}}} \dicPos{n}[2] \dicFlx{(‑s, ‑)}[5] \dicFieldCat{sport.} \dicDirectTranslationCS{sprint}
\dicEntry[spretthlaupari] \dicTerm{sprett··hlaup·ar|i} \dicIPA{{s}{\textsubring{b}}{r}{\textepsilon}{h}{\textsubring{d}}{\textsubring{l}}{\oe i}{\textsubring{b}}{a}{r}{\textsci}} \dicPos{m}[1] \dicFlx{(‑a, ‑ar)}[13] \dicDirectTranslationCS{sprinter(ka)}
\dicEntry[sprettigluggi] \dicTerm{spretti··glugg|i} \dicIPA{{s}{\textsubring{b}}{r}{\textepsilon}{h}{\textsubring{d}}{\textsci}{\r{g}}{l}{\textscy}{\r{\textObardotlessj}}{\textsci}} \dicPos{m}[1] \dicFlx{(‑a, ‑ar)}[1] \dicFieldCat{poč.} \dicDirectTranslationCS{vyskakovací okno}
\dicEntry[sprettur] \dicTerm{sprett|ur} \dicIPA{{s}{\textsubring{b}}{r}{\textepsilon}{h}{\textsubring{d}}{\textscy}{\textsubring{r}}} \dicPos{m}[9] \dicFlx{(‑s, ‑ir)}[4] \dicDirectTranslationCS{sprint, spurt};  \dicPhraseIS{taka á sprett} \dicDirectTranslationCS{sprintovat, spurtovat}
\dicEntry[spreyta] \dicTerm{spreyt|a} \dicIPA{{s}{\textsubring{b}}{r}{ei}{\textlengthmark}{\textsubring{d}}{a}} \dicPos{v}[2] \dicFlx{(‑ti, ‑t)}[62] \dicPhraseIS{spreyta sig á e‑u} \dicDirectTranslationCS{vyzkoušet si (co), zkusit si (co), změřit se s~(čím)}
\dicEntry[sprikl] \dicTerm{sprikl} \dicIPA{{s}{\textsubring{b}}{r}{\textsci}{h}{\r{g}}{\textsubring{l}}} \dicPos{n}[2] \dicFlx{(‑s)}[2] \textbf{1.} \dicSynonym*{það að sprikla} \dicDirectTranslationCS{svíjení, zmítání, kopání}  \textbf{2.} \dicLangCat{hovor.} \dicSynonym{leikfimi} \dicDirectTranslationCS{skákání, cvičení (v~tělocvičně na hudbu ap.)}
\dicEntry[sprikla] \dicTerm{sprikl|a} \dicIPA{{s}{\textsubring{b}}{r}{\textsci}{h}{\r{g}}{l}{a}} \dicPos{v}[1] \dicFlx{(‑aði)}[1] \dicDirectTranslationCS{svíjet se, zmítat se, kopat sebou} \dicExampleIS{Fiskurinn spriklaði í netinu.} \dicExampleCS{Ryba se zmítala v~síti.}
\dicEntry[spring] \dicTerm{spring} \dicIPA{{s}{\textsubring{b}}{r}{\textsci}{\ng}{\r{g}}} \dicPos{v} \dicFlx{ind praes sg 1 pers} \dicLink{springa}
\dicEntry[springa] \dicTerm{springa} \dicsymFrequent\  \dicIPA{{s}{\textsubring{b}}{r}{i}{\ng}{\r{g}}{a}} \dicPos{v}[6] \dicFlx{(spring, sprakk, sprungum, spryngi, sprungið)}[41] \textbf{1.} \dicSynonym{rifna} \dicDirectTranslationCS{prasknout, puknout, roztrhnout se} \dicExampleIS{Pokinn sprakk.} \dicExampleCS{Sáček se roztrhl.};  \dicPhraseIS{það springur á e‑u} \dicFlx{impers} \dicDirectTranslationCS{(co) píchá (pneumatika auta ap.)} \dicExampleIS{Það sprakk á bílnum.} \dicExampleCS{Auto píchlo.}  \textbf{2.} \dicDirectTranslationCS{(vy)bouchnout, explodovat};  \dicIdiom{springa}[úr]{ \dicPhraseIS{springa úr hlátri}} \dicDirectTranslationCS{pukat smíchy};  \dicIdiom{springa}[út]{ \dicPhraseIS{springa út}} \dicDirectTranslationCS{rozvinout se, rozvíjet se,  vykvést, vykvétat, rozkvést, rozkvétat (květ ap.)} \dicExampleIS{Blómið er sprungið út.} \dicExampleCS{Květina rozkvetla.}
\dicEntry[springdýna] \dicTerm{spring··dýn|a} \dicIPA{{s}{\textsubring{b}}{r}{i}{\ng}{\r{g}}{\textsubring{d}}{i}{n}{a}} \dicPos{f}[1] \dicFlx{(‑u, ‑ur)}[7] \dicDirectTranslationCS{pružinová matrace}
\dicEntry[spritt] \dicTerm{spritt} \dicIPA{{s}{\textsubring{b}}{r}{\textsci}{h}{\textsubring{d}}} \dicPos{n}[2] \dicFlx{(‑s)}[2] \dicDirectTranslationCS{líh, alkohol}
\dicEntry[sproti] \dicTerm{sprot|i} \dicIPA{{s}{\textsubring{b}}{r}{\textopeno}{\textlengthmark}{\textsubring{d}}{\textsci}} \dicPos{m}[1] \dicFlx{(‑a, ‑ar)}[1] \textbf{1.} \dicSynonym{angi} \dicDirectTranslationCS{výhonek, odnož}  \textbf{2.} \dicSynonym{prik} \dicDirectTranslationCS{hůlka, tyčka}
\dicEntry[sprottið] \dicTerm{sprottið} \dicIPA{{s}{\textsubring{b}}{r}{\textopeno}{h}{\textsubring{d}}{\textsci}{\texttheta}} \dicPos{v} \dicFlx{supin} \dicLink{spretta\smash{\textsuperscript{2}}}
\dicEntry[sprottinn] \dicTerm{sprottinn} \dicIPA{{s}{\textsubring{b}}{r}{\textopeno}{h}{\textsubring{d}}{\textsci}{\textsubring{n}}} \dicPos{adj}[6]\dicFlx{}[-2] \textbf{1.} \dicDirectTranslationCS{vyrašený, vzrostlý}  \textbf{2.} \dicSynonym*{upprunninn} \dicDirectTranslationCS{mající kořeny\,/\addthin původ}
\dicEntry[sprunga] \dicTerm{sprung|a} \dicIPA{{s}{\textsubring{b}}{r}{u}{\ng}{\r{g}}{a}} \dicPos{f}[1] \dicFlx{(‑u, ‑ur)}[13] \textbf{1.} \dicSynonym*{skora} \dicDirectTranslationCS{štěrbina, trhlina} \dicExampleIS{Sprunga myndast í vegginn.} \dicExampleCS{Na stěně se tvoří trhlina.}  \textbf{2.} \dicFieldCat{geol.} \dicSynonym{gjá} \dicDirectTranslationCS{puklina, trhlina}
\dicEntry[sprungið] \dicTerm{sprungið} \dicIPA{{s}{\textsubring{b}}{r}{\textscy}{\textltailn}{\r{\textObardotlessj}}{\textsci}{\texttheta}} \dicPos{v} \dicFlx{supin} \dicLink{springa}
\dicEntry[sprungum] \dicTerm{sprungum} \dicIPA{{s}{\textsubring{b}}{r}{\textscy}{\ng}{\r{g}}{\textscy}{\textsubring{m}}} \dicPos{v} \dicFlx{ind pf pl 1 pers} \dicLink{springa}
\dicEntry[spruttum] \dicTerm{spruttum} \dicIPA{{s}{\textsubring{b}}{r}{\textscy}{h}{\textsubring{d}}{\textscy}{\textsubring{m}}} \dicPos{v} \dicFlx{ind pf pl 1 pers} \dicLink{spretta\smash{\textsuperscript{2}}}
\dicEntry[sprútt] \dicTerm{sprútt} \dicIPA{{s}{\textsubring{b}}{r}{u}{h}{\textsubring{d}}} \dicPos{n}[2] \dicFlx{(‑s)}[2] \dicLangCat{hovor.} \dicDirectTranslationCS{chlast, pití}
\dicEntry[sprúttsali] \dicTerm{sprútt··sal|i} \dicIPA{{s}{\textsubring{b}}{r}{u}{h}{\textsubring{d}}{s}{a}{l}{\textsci}} \dicPos{m}[1] \dicFlx{(‑a, ‑ar)}[8] \dicDirectTranslationCS{nezákonný prodejce alkoholu}
\dicEntry[spryngi] \dicTerm{spryngi} \dicIPA{{s}{\textsubring{b}}{r}{\textsci}{\textltailn}{\r{\textObardotlessj}}{\textsci}} \dicPos{v} \dicFlx{con pf sg 1 pers} \dicLink{springa}
\dicEntry[sprytti] \dicTerm{sprytti} \dicIPA{{s}{\textsubring{b}}{r}{\textsci}{h}{\textsubring{d}}{\textsci}} \dicPos{v} \dicFlx{con pf sg 1 pers} \dicLink{spretta\smash{\textsuperscript{2}}}
\dicEntry[sprækur] \dicTerm{sprækur} \dicIPA{{s}{\textsubring{b}}{r}{a}{i}{\textlengthmark}{\r{g}}{\textscy}{\textsubring{r}}} \dicPos{adj}[1]\dicFlx{}[-1] \dicSynonym{hress} \dicDirectTranslationCS{živý, čiperný, čilý}
\dicEntry[spræna] \dicTerm{spræn|a\smash{\textsuperscript{1}}} \dicIPA{{s}{\textsubring{b}}{r}{a}{i}{\textlengthmark}{n}{a}} \dicPos{f}[1] \dicFlx{(‑u, ‑ur)}[7] \textbf{1.} \dicSynonym{lækur} \dicDirectTranslationCS{potok, potůček}  \textbf{2.} \dicSynonym*{hlandbuna} \dicDirectTranslationCS{proud (moči)}
\dicEntry[spræna] \dicTerm{spræn|a\smash{\textsuperscript{2}}} \dicIPA{{s}{\textsubring{b}}{r}{a}{i}{\textlengthmark}{n}{a}} \dicPos{v}[2] \dicFlx{(‑di, ‑t)}[149] \dicSynonym{buna\smash{\textsuperscript{2}}} \dicDirectTranslationCS{vytrysknout, tryskat, řinout se}
\dicEntry[spröku] \dicTerm{spröku} \dicIPA{{s}{\textsubring{b}}{r}{\oe}{\textlengthmark}{\r{g}}{\textscy}} \dicPos{f} \dicFlx{sg gen} \dicLink{spraka}
\dicEntry[spuni] \dicTerm{spun|i} \dicIPA{{s}{\textsubring{b}}{\textscy}{\textlengthmark}{n}{\textsci}} \dicPos{m}[1] \dicFlx{(‑a)}[3] \dicDirectTranslationCS{improvizace, improvizování};  \dicPhraseIS{vera stuttur í spuna} \dicDirectTranslationCS{být úsečný\,/\addthin strohý}
\dicEntry[spunnið] \dicTerm{spunnið} \dicIPA{{s}{\textsubring{b}}{\textscy}{n}{\textlengthmark}{\textsci}{\texttheta}} \dicPos{v} \dicFlx{supin} \dicLink{spinna}
\dicEntry[spunnum] \dicTerm{spunnum} \dicIPA{{s}{\textsubring{b}}{\textscy}{n}{\textlengthmark}{\textscy}{\textsubring{m}}} \dicPos{v} \dicFlx{ind pf pl 1 pers} \dicLink{spinna}
\dicEntry[spurði] \dicTerm{spurði} \dicIPA{{s}{\textsubring{b}}{\textscy}{r}{ð}{\textsci}} \dicPos{v} \dicFlx{ind pf sg 1 pers} \dicLink{spyrja}
\dicEntry[spurðum] \dicTerm{spurðum} \dicIPA{{s}{\textsubring{b}}{\textscy}{r}{ð}{\textscy}{\textsubring{m}}} \dicPos{v} \dicFlx{ind pf pl 1 pers} \dicLink{spyrja}
\dicEntry[spurn] \dicTerm{spurn} \dicsymFrequent\  \dicIPA{{s}{\textsubring{b}}{\textscy}{r}{\textsubring{d}}{\textsubring{n}}} \dicPos{f}[7] \dicFlx{(‑ar, ‑ir)}[1] \dicPhraseIS{hafa  spurnir af e‑u} \dicDirectTranslationCS{mít o~(čem) zprávy} \dicExampleIS{hafa engar spurnir af atburðum} \dicExampleCS{nemít žádné zprávy o~událostech};  \dicPhraseIS{í spurn} \dicFlx{adv} \dicDirectTranslationCS{tázavě, zkoumavě, pátravě};  \dicPhraseIS{mér er spurn hvort} \dicDirectTranslationCS{rád bych věděl, zda}
\dicEntry[spurnarform] \dicTerm{spurnar··form} \dicIPA{{s}{\textsubring{b}}{\textscy}{r}{\textsubring{d}}{n}{a}{\textsubring{r}}{f}{\textopeno}{r}{\textsubring{m}}} \dicPos{n}[2] \dicFlx{(‑s, ‑)}[5] \dicFieldCat{jaz.} \dicDirectTranslationCS{tázací forma}
\dicEntry[spurnarfornafn] \dicTerm{spurnar··for·|nafn} \dicIPA{{s}{\textsubring{b}}{\textscy}{r}{\textsubring{d}}{n}{a}{\textsubring{r}}{f}{\textopeno}{r}{n}{a}{\textsubring{b}}{\textsubring{n}}} \dicPos{n}[2] \dicFlx{(‑nafns, ‑nöfn)}[8] \dicFieldCat{jaz.} \dicDirectTranslationCS{zájmeno tázací}
\dicEntry[spurnarsetning] \dicTerm{spurnar··set·ning} \dicIPA{{s}{\textsubring{b}}{\textscy}{r}{\textsubring{d}}{n}{a}{\textsubring{r}}{s}{\textepsilon}{h}{\textsubring{d}}{n}{i}{\ng}{\r{g}}} \dicPos{f}[4] \dicFlx{(‑ar, ‑ar)}[5] \dicFieldCat{jaz.} \dicDirectTranslationCS{tázací věta}
\dicEntry[spurnarsvipur] \dicTerm{spurnar··svip|ur} \dicIPA{{s}{\textsubring{b}}{\textscy}{r}{\textsubring{d}}{n}{a}{\textsubring{r}}{s}{v}{\textsci}{\textsubring{b}}{\textscy}{\textsubring{r}}} \dicPos{m}[9] \dicFlx{(‑s, ‑ir)}[8] \dicDirectTranslationCS{tázavý výraz}
\dicEntry[spurnartenging] \textls[10]{\dicTerm{spurnar··teng·ing} \dicIPA{{s}{\textsubring{b}}{\textscy}{r}{\textsubring{d}}{n}{a}{\textsubring{r}}{t\smash{\textsuperscript{h}}}{ei}{\textltailn}{\r{\textObardotlessj}}{i}{\ng}{\r{g}}} \dicPos{f}[4] \dicFlx{(‑ar, ‑ar)}[5] \dicFieldCat{jaz.}} \dicDirectTranslationCS{spojka tázací}
\dicEntry[spurning] \dicTerm{spur··ning} \dicsymFrequent\  \dicIPA{{s}{\textsubring{b}}{\textscy}{r}{\textsubring{d}}{n}{i}{\ng}{\r{g}}} \dicPos{f}[4] \dicFlx{(‑ar, ‑ar)}[5] \textbf{1.} \dicDirectTranslationCS{otázka, dotaz} \textls[8]{\dicExampleIS{Nemendurnir spurðu kennarann margra spurninga.}} \dicExampleCS{Žáci se ptali učitele na mnoho otázek.};  \textls[15]{\dicPhraseIS{bera fram spurningu} \dicDirectTranslationCS{položit otázku};  \dicPhraseIS{leggja fyrir e‑n spurningu} \dicDirectTranslationCS{položit (komu) otázku}}  \textbf{2.} \dicSynonym{vafamál} \dicDirectTranslationCS{otázka, problém};  \dicPhraseIS{spurning um líf og dauða} \dicDirectTranslationCS{otázka života a~smrti};  \dicPhraseIS{það er engin spurning} \dicDirectTranslationCS{o~tom není pochyb}
\dicEntry[spurningakeppni] \dicTerm{spurninga··keppn|i} \dicIPA{{s}{\textsubring{b}}{\textscy}{r}{\textsubring{d}}{n}{i}{\ng}{\r{g}}{a}{c\smash{\textsuperscript{h}}}{\textepsilon}{h}{\textsubring{b}}{n}{\textsci}} \dicPos{f}[3] \dicFlx{(‑i, ‑ir)}[1] \dicDirectTranslationCS{(soutěžní) kvíz, znalostní\,/\addthin vědomostní soutěž}
\dicEntry[spurningalisti] \dicTerm{spurninga··list|i} \dicIPA{{s}{\textsubring{b}}{\textscy}{r}{\textsubring{d}}{n}{i}{\ng}{\r{g}}{a}{l}{\textsci}{s}{\textsubring{d}}{\textsci}} \dicPos{m}[1] \dicFlx{(‑a, ‑ar)}[1] \dicDirectTranslationCS{dotazník}
\dicEntry[spurningarmerki] \dicTerm{spurningar··merki} \dicIPA{{s}{\textsubring{b}}{\textscy}{r}{\textsubring{d}}{n}{i}{\ng}{\r{g}}{a}{r}{m}{\textepsilon}{\textsubring{r}}{\r{\textObardotlessj}}{\textsci}} \dicPos{n}[2] \dicFlx{(‑s, ‑)}[16] \dicFieldCat{jaz.} \dicDirectTranslationCS{otazník} \dicIndirectTranslationCS{(interpunkční znaménko)}
\begin{xtolerant}{}{1pt}
\dicEntry[spurningaþáttur] \dicTerm{spurninga··|þáttur}\addthinS\dicIPA{{s}{\textsubring{b}}{\textscy}{r}{\textsubring{d}}{n}{i}{\ng}{\r{g}}{a}{\texttheta}{au}{h}{\textsubring{d}}{\textscy}{\textsubring{r}}}\addthinS\dicPos{m}[12]\addthinS\dicFlx{(‑þáttar, ‑þætt\-ir)}[7] \dicDirectTranslationCS{(soutěžní) kvíz, kvízový pořad}
\end{xtolerant}
\dicEntry[spursmál] \dicTerm{spurs··mál} \dicIPA{{s}{\textsubring{b}}{\textscy}{\textsubring{r}}{s}{m}{au}{\textsubring{l}}} \dicPos{n}[2] \dicFlx{(‑s, ‑)}[5] \dicSynonym{úrlausnarefni} \dicDirectTranslationCS{otázka, problém, záležitost}
\dicEntry[spurt] \dicTerm{spurt} \dicIPA{{s}{\textsubring{b}}{\textscy}{\textsubring{r}}{\textsubring{d}}} \dicPos{v} \dicFlx{supin} \dicLink{spyrja}
\dicEntry[spurull] \dicTerm{spurull} \dicIPA{{s}{\textsubring{b}}{\textscy}{\textlengthmark}{r}{\textscy}{\textsubring{d}}{\textsubring{l}}} \dicPos{adj}[8]\dicFlx{}[-4] \dicSynonym{forvitinn} \dicDirectTranslationCS{zvídavý, všetečný} \dicExampleIS{vera spurull um einkahagi gestanna} \dicExampleCS{vyzvídat o~osobním životě hostů}
\dicEntry[spúa] \dicTerm{spú|a} \dicIPA{{s}{\textsubring{b}}{u}{\textlengthmark}{a}} \dicPos{v}[2] \dicFlx{(‑ði, ‑ð)}[114] \dicFlx{dat} \dicSynonym*{spýja} \dicDirectTranslationCS{chrlit, plivat, soptit (oheň ap.)} \dicExampleIS{Eldfjöllin spúa ösku yfir landið.} \dicExampleCS{Sopky chrlí po krajině popel.}
\dicEntry[spúð] \dicTerm{spúð} \dicIPA{{s}{\textsubring{b}}{u}{\textlengthmark}{\texttheta}} \dicPos{v} \dicFlx{supin} \dicLink{spýja\smash{\textsuperscript{2}}}
\dicEntry[spúði] \dicTerm{spúði} \dicIPA{{s}{\textsubring{b}}{u}{\textlengthmark}{ð}{\textsci}} \dicPos{v} \dicFlx{ind pf sg 1 pers} \dicLink{spýja\smash{\textsuperscript{2}}}
\dicEntry[spúðum] \dicTerm{spúðum} \dicIPA{{s}{\textsubring{b}}{u}{\textlengthmark}{ð}{\textscy}{\textsubring{m}}} \dicPos{v} \dicFlx{ind pf pl 1 pers} \dicLink{spýja\smash{\textsuperscript{2}}}
\dicEntry[spúnn] \dicTerm{spún|n} \dicIPA{{s}{\textsubring{b}}{u}{\textsubring{d}}{\textsubring{n}}} \dicPos{m}[6] \dicFlx{(‑s, ‑ar)}[42] \dicDirectTranslationCS{(rotační) třpytka}
\dicEntry[spynni] \dicTerm{spynni} \dicIPA{{s}{\textsubring{b}}{\textsci}{n}{\textlengthmark}{\textsci}} \dicPos{v} \dicFlx{con pf sg 1 pers} \dicLink{spinna}
\dicEntry[spyr] \dicTerm{spyr} \dicIPA{{s}{\textsubring{b}}{\textsci}{\textlengthmark}{\textsubring{r}}} \dicPos{v} \dicFlx{ind praes sg 1 pers} \dicLink{spyrja}
\dicEntry[spyrða] \dicTerm{spyrð|a\smash{\textsuperscript{1}}} \dicIPA{{s}{\textsubring{b}}{\textsci}{r}{ð}{a}} \dicPos{f}[1] \dicFlx{(‑u, ‑ur)}[19] \dicDirectTranslationCS{svazek ryb svázaných za ocas}
\dicEntry[spyrða] \dicTerm{spyr|ða\smash{\textsuperscript{2}}} \dicIPA{{s}{\textsubring{b}}{\textsci}{r}{ð}{a}} \dicPos{v}[2] \dicFlx{(‑ti, ‑t)}[47] \dicFlx{acc} \textbf{1.} \dicDirectTranslationCS{svázat ryby za ocas}  \textbf{2.} \dicPhraseIS{spyrða saman e‑ð} \dicDirectTranslationCS{spojit (co), sjednotit (co)}
\dicEntry[spyrði] \dicTerm{spyrði} \dicIPA{{s}{\textsubring{b}}{\textsci}{r}{ð}{\textsci}} \dicPos{v} \dicFlx{con pf sg 1 pers} \dicLink{spyrja}
\dicEntry[spyrill] \dicTerm{spyr|ill} \dicIPA{{s}{\textsubring{b}}{\textsci}{\textlengthmark}{r}{\textsci}{\textsubring{d}}{\textsubring{l}}} \dicPos{m}[6] \dicFlx{(‑ils, ‑lar)}[35] \dicDirectTranslationCS{tazatel(ka), dotazovatel(ka)}
\dicEntry[spyrja] \dicTerm{spyrja} \dicsymFrequent\  \dicIPA{{s}{\textsubring{b}}{\textsci}{r}{j}{a}} \dicPos{v}[4] \dicFlx{(spyr, spurði, spurðum, spyrði, spurt)}[8] \dicFlx{acc (+ gen)} \dicDirectTranslationCS{zeptat se, ptát se, (do)tázat se, dotazovat se} \dicExampleIS{Ég spurði ekki að því.} \dicExampleCS{Na to jsem se neptal.};  \dicPhraseIS{spyrja e‑n um e‑ð, spyrja e‑n að e‑u, spyrja e‑n e‑s} \dicDirectTranslationCS{ptát se (koho) na (co)};  \dicPhraseIS{spyrja e‑n að nafni} \dicDirectTranslationCS{zeptat se (koho) na jméno};  \dicPhraseIS{spyrja e‑n frétta} \dicDirectTranslationCS{ptát se (koho), co je nového};  \dicIdiom{spyrja}[til]{ \dicPhraseIS{spyrja til vegar}} \dicDirectTranslationCS{ptát se na cestu};  \dicIdiom{spyrja}[út úr]{ \dicPhraseIS{spyrja e‑n út úr}} \dicDirectTranslationCS{vyzpovídat (koho), vyptat se (koho)};  \dicIdiom{spyrjast}{ \dicPhraseIS{spyrjast}} \dicFlx{refl} \dicDirectTranslationCS{roznést se, rozkřiknout se (zpráva ap.)} \dicExampleIS{Þetta spurðist til landsins.} \dicExampleCS{Rozneslo se to po kraji.};  \dicIdiom{spyrjast}[fyrir]{ \dicPhraseIS{spyrjast fyrir um e‑ð}} \dicFlx{refl} \dicDirectTranslationCS{informovat se o~(čem)};  \dicIdiom{spyrjast}[út]{ \dicPhraseIS{spyrjast út}} \dicFlx{refl} \dicDirectTranslationCS{roznést se, rozkřiknout se (zpráva ap.)}
\dicEntry[spyrjandi] \dicTerm{spyrj··|andi\smash{\textsuperscript{1}}} \dicIPA{{s}{\textsubring{b}}{\textsci}{r}{j}{a}{n}{\textsubring{d}}{\textsci}} \dicPos{m}[2] \dicFlx{(‑anda, ‑endur)}[1] \dicDirectTranslationCS{tazatel(ka)}
\dicEntry[spyrjandi] \dicTerm{spyrj··andi\smash{\textsuperscript{2}}} \dicIPA{{s}{\textsubring{b}}{\textsci}{r}{j}{a}{n}{\textsubring{d}}{\textsci}} \dicPos{adj}[13] \dicFlx{indecl}[1] \dicDirectTranslationCS{tázavý, zkoumavý}
\dicEntry[spyrna] \dicTerm{spyrn|a\smash{\textsuperscript{1}}} \dicIPA{{s}{\textsubring{b}}{\textsci}{r}{\textsubring{d}}{n}{a}} \dicPos{f}[1] \dicFlx{(‑u, ‑ur)}[7] \dicSynonym{spark} \dicDirectTranslationCS{(vy)kopnutí}
\dicEntry[spyrna] \dicTerm{spyrn|a\smash{\textsuperscript{2}}} \dicIPA{{s}{\textsubring{b}}{\textsci}{r}{\textsubring{d}}{n}{a}} \dicPos{v}[2] \dicFlx{(‑ti, ‑t)}[63] \dicFlx{dat} \textbf{1.} \dicSynonym{sparka} \dicDirectTranslationCS{(vy)kopnout, kopat} \dicExampleIS{spyrna boltanum í mark} \dicExampleCS{kopnout míč do brány}  \textbf{2.} \dicSynonym*{veita viðnám} \dicDirectTranslationCS{stavět se, vzpírat se};  \dicPhraseIS{spyrna á móti e‑u} \dicDirectTranslationCS{stavět se proti (čemu)};  \dicPhraseIS{spyrna við fótum} \dicDirectTranslationCS{postavit se na odpor}
\dicEntry[spý] \dicTerm{spý} \dicIPA{{s}{\textsubring{b}}{i}{\textlengthmark}} \dicPos{v} \dicFlx{ind praes sg 1 pers} \dicLink{spýja\smash{\textsuperscript{2}}}
\dicEntry[spýði] \dicTerm{spýði} \dicIPA{{s}{\textsubring{b}}{i}{\textlengthmark}{ð}{\textsci}} \dicPos{v} \dicFlx{con pf sg 1 pers} \dicLink{spýja\smash{\textsuperscript{2}}}
\dicEntry[spýja] \dicTerm{spýj|a\smash{\textsuperscript{1}}} \dicIPA{{s}{\textsubring{b}}{i}{j}{\textlengthmark}{a}} \dicPos{f}[1] \dicFlx{(‑u, ‑ur)}[7] \textbf{1.} \dicSynonym{æla\smash{\textsuperscript{1}}} \dicDirectTranslationCS{zvratek, vydávenina} \dicExampleIS{liggja á gólfinu í spýju sinni} \dicExampleCS{ležet na podlaze ve svých zvratcích}  \textbf{2.} \dicDirectTranslationCS{sprška (sněhová ap.)}
\dicEntry[spýja] \dicTerm{spýja\smash{\textsuperscript{2}}} \dicIPA{{s}{\textsubring{b}}{i}{j}{\textlengthmark}{a}} \dicPos{v}[5] \dicFlx{(spý, spúði\,/\addthin spjó, spúðum, spýði, spúð)}[17] \dicFlx{dat} \textbf{1.} \dicSynonym*{kasta upp} \dicDirectTranslationCS{(vy)zvracet, (vy)dávit}  \textbf{2.} \dicSynonym{spúa} \dicDirectTranslationCS{chrlit, soptit} \dicExampleIS{Hekla spúði eldi.} \dicExampleCS{Hekla chrlila oheň.}
\dicEntry[spýta] \dicTerm{spýt|a\smash{\textsuperscript{1}}} \dicsymFrequent\  \dicIPA{{s}{\textsubring{b}}{i}{\textlengthmark}{\textsubring{d}}{a}} \dicPos{f}[1] \dicFlx{(‑u, ‑ur)}[13] \textbf{1.} \dicSynonym{fjöl\smash{\textsuperscript{1}}} \dicDirectTranslationCS{deska, prkno, fošna} \dicExampleIS{saga spýtuna í sundur} \dicExampleCS{rozřezat prkno}  \textbf{2.} \dicSynonym{eldspýta} \dicDirectTranslationCS{sirka, zápalka};  \dicIdiom{spýta}{ \dicPhraseIS{e‑að hangir á spýtunni}} \dicLangCat{přen.} \dicDirectTranslationCS{(co) se za tím skrývá, (co) za tím je}; { \dicPhraseIS{á eigin spýtur}} \dicFlx{adv} \dicDirectTranslationCS{na vlastní pěst}
\dicEntry[spýta] \dicTerm{spýt|a\smash{\textsuperscript{2}}} \dicsymFrequent\  \dicIPA{{s}{\textsubring{b}}{i}{\textlengthmark}{\textsubring{d}}{a}} \dicPos{v}[2] \dicFlx{(‑ti, ‑t)}[59] \dicFlx{dat} \textbf{1.} \dicDirectTranslationCS{(od)plivnout (si), vyplivnout, plivat} \dicExampleIS{spýta matnum út úr sér} \dicExampleCS{vyplivnout jídlo};  \dicPhraseIS{spýta í lófana} \dicLangCat{přen.} \dicDirectTranslationCS{plivnout si do dlaní}  \textbf{2.} \dicDirectTranslationCS{napnout kůži (na dřevěný rám ap.)};  \dicIdiom{spýtast}{ \dicPhraseIS{spýtast}} \dicFlx{refl} \dicDirectTranslationCS{stříkat, tryskat} \dicExampleIS{Vatnið spýttist um allt.} \dicExampleCS{Voda stříkala všude kolem.}
\dicEntry[spægipylsa] \dicTerm{spægi··pyls|a} \dicIPA{{s}{\textsubring{b}}{a}{i}{j}{\textlengthmark}{\textsci}{p\smash{\textsuperscript{h}}}{\textsci}{l}{s}{a}} \dicPos{f}[1] \dicFlx{(‑u, ‑ur)}[13] \dicFieldCat{kulin.} \dicDirectTranslationCS{klobása}
\dicEntry[spæja] \dicTerm{spæj|a} \dicIPA{{s}{\textsubring{b}}{a}{i}{j}{\textlengthmark}{a}} \dicPos{v}[1] \dicFlx{(‑aði)}[1] \dicDirectTranslationCS{získávat tajné informace} \dicIndirectTranslationCS{(o~špionovi)}
\dicEntry[spæjari] \dicTerm{spæj··ar|i} \dicIPA{{s}{\textsubring{b}}{a}{i}{j}{\textlengthmark}{a}{r}{\textsci}} \dicPos{m}[1] \dicFlx{(‑a, ‑ar)}[13] \textbf{1.} \dicSynonym{njósnari} \dicDirectTranslationCS{špion(ka), špeh, agent(ka)}  \textbf{2.} \dicSynonym{leynilögreglumaður} \dicDirectTranslationCS{detektiv} \dicExampleIS{einkaspæjari} \dicExampleCS{soukromý detektiv}
\dicEntry[spæla] \dicTerm{spæl|a} \dicIPA{{s}{\textsubring{b}}{a}{i}{\textlengthmark}{l}{a}} \dicPos{v}[2] \dicFlx{(‑di, ‑t)}[141] \dicFlx{acc} \textbf{1.} \dicDirectTranslationCS{(u)smažit (vajíčka)} \dicExampleIS{spæla egg} \dicExampleCS{smažit vajíčko}  \textbf{2.} \dicSynonym{hrella} \dicDirectTranslationCS{zlobit, štvát} \dicExampleIS{spæla e‑n} \dicExampleCS{zlobit (koho)}
\dicEntry[spældur] \dicTerm{spældur} \dicIPA{{s}{\textsubring{b}}{a}{i}{l}{\textsubring{d}}{\textscy}{\textsubring{r}}} \dicPos{adj}[2]\dicFlx{}[-14] \textbf{1.} \dicDirectTranslationCS{usmažený (vajíčko ap.)}  \textbf{2.} \dicSynonym{vonsvikinn} \dicDirectTranslationCS{zdecimovaný, zničený, zklamaný}
\dicEntry[spælegg] \dicTerm{spæl··egg} \dicIPA{{s}{\textsubring{b}}{a}{i}{\textlengthmark}{l}{\textepsilon}{\r{g}}} \dicPos{n}[2] \dicFlx{(‑s, ‑)}[13] \dicDirectTranslationCS{smažené vejce\,/\addthin vajíčko}
\dicEntry[spæling] \dicTerm{spæl··ing} \dicIPA{{s}{\textsubring{b}}{a}{i}{\textlengthmark}{l}{i}{\ng}{\r{g}}} \dicPos{f}[4] \dicFlx{(‑ar)}[7] \dicDirectTranslationCS{zklamání, rozčarování}
\dicEntry[spæna] \dicTerm{spæn|a} \dicIPA{{s}{\textsubring{b}}{a}{i}{\textlengthmark}{n}{a}} \dicPos{v}[2] \dicFlx{(‑di, ‑t)}[141] \dicFlx{acc} \textbf{1.} \dicPhraseIS{spæna í sig e‑ð} \dicDirectTranslationCS{házet do sebe (co), hltat (co) (jídlo ap.)} \dicExampleIS{Hann spændi í sig grautinn.} \dicExampleCS{Zhltnul kaši.}  \textbf{2.} \dicPhraseIS{spæna upp e‑ð} \dicDirectTranslationCS{rozjezdit (co), rozhrabat (co), rozrýt (co) (zem, povrch ap.)}  \textbf{3.} \dicDirectTranslationCS{odfičet, odsvištět} \dicIndirectTranslationCS{(rychle odjet autem)};  \dicIdiom{spænast}[upp]{ \dicPhraseIS{spænast upp}} \dicFlx{refl} \dicSynonym*{eyðast} \dicDirectTranslationCS{opotřebovat se}
\dicEntry[spænir] \dicTerm{spænir} \dicIPA{{s}{\textsubring{b}}{a}{i}{\textlengthmark}{n}{\textsci}{\textsubring{r}}} \dicPos{m} \dicFlx{pl nom} \dicLink{spónn}
\dicEntry[spænska] \dicTerm{spænsk|a} \dicIPA{{s}{\textsubring{b}}{a}{i}{n}{s}{\r{g}}{a}} \dicPos{f}[1] \dicFlx{(‑u)}[5] \dicDirectTranslationCS{španělština}
\begin{xtolerant}{}{1pt}
\dicEntry[spænskumælandi] \dicTerm{spænsku··mæl·andi}\addthinS\dicIPA{{s}{\textsubring{b}}{a}{i}{n}{s}{\r{g}}{\textscy}{m}{a}{i}{l}{a}{n}{\textsubring{d}}{\textsci}} \dicPos{adj}[13] \dicFlx{indecl}[1] \dicDirectTranslationCS{špa\-nělsky mluvící, hispanofonní}
\end{xtolerant}
\dicEntry[spænskur] \dicTerm{spænskur} \dicIPA{{s}{\textsubring{b}}{a}{i}{n}{s}{\r{g}}{\textscy}{\textsubring{r}}} \dicPos{adj}[1]\dicFlx{}[-1] \dicDirectTranslationCS{španělský}
\dicEntry[spæta] \dicTerm{spæt|a} \dicIPA{{s}{\textsubring{b}}{a}{i}{\textlengthmark}{\textsubring{d}}{a}} \dicPos{f}[1] \dicFlx{(‑u, ‑ur)}[7] \dicFieldCat{zool.} \dicDirectTranslationCS{datel} \textit{(l.~{\textLA{Picidae}})}  \dicsymPhoto\ 
\dicFigure{113029.jpg}{Spæta}{Spæta - Došlý Martin, Biolib, Copyright/PD}
\dicEntry[spölkorn] \dicTerm{spöl··korn} \dicIPA{{s}{\textsubring{b}}{\oe}{l}{k\smash{\textsuperscript{h}}}{\textopeno}{r}{\textsubring{d}}{\textsubring{n}}} \dicPos{n}[2] \dicFlx{(‑s)}[2] \dicDirectTranslationCS{kousek} \dicIndirectTranslationCS{(o~vzdálenosti)}
\dicEntry[spölur] \dicTerm{spölur} \dicsymFrequent\  \dicIPA{{s}{\textsubring{b}}{\oe}{\textlengthmark}{l}{\textscy}{\textsubring{r}}} \dicPos{m}[11] \dicFlx{(spalar, spelir)}[7] \dicSynonym*{stutt leið} \dicDirectTranslationCS{krátká vzdálenost, kousek} \dicExampleIS{Við gengum dálítinn spöl.} \dicExampleCS{Půjdeme malý kousek.}
\dicEntry[spöng] \dicTerm{spöng} \dicIPA{{s}{\textsubring{b}}{\oe i}{\ng}{\r{g}}} \dicPos{f}[8] \dicFlx{(spangar, spengur\,/\addthin spangir)}[9] \textbf{1.} \dicDirectTranslationCS{spona, sponka} \dicExampleIS{hárspöng} \dicExampleCS{sponka do vlasů}  \textbf{2.} \dicSynonym*{höfuðdjásn} \dicDirectTranslationCS{diadém}  \textbf{3.} \dicSynonym*{ísspöng} \dicDirectTranslationCS{ledový plát, ledová plotna}  \textbf{4.} \dicPhraseIS{spangir} \dicFlx{pl} \dicDirectTranslationCS{(zubní) rovnátka}
\dicEntry[spönn] \dicTerm{spönn} \dicIPA{{s}{\textsubring{b}}{\oe}{n}{\textlengthmark}} \dicPos{f}[7] \dicFlx{(spannar, spannir)}[16] \dicFieldCat{hist.} \dicDirectTranslationCS{píď};  \dicPhraseIS{hreyfa sig ekki spönn frá rassi} \dicLangCat{přen.} \dicDirectTranslationCS{nehnout ani prstem};  \dicPhraseIS{komast ekki spönn frá rassi} \dicFlx{refl} \dicLangCat{přen.} \dicDirectTranslationCS{nehnout se ani o~píď (s~úkolem ap.)}
\dicEntry[spör] \dicTerm{spör\smash{\textsuperscript{1}}} \dicIPA{{s}{\textsubring{b}}{\oe}{\textlengthmark}{\textsubring{r}}} \dicPos{m}[8] \dicFlx{(‑s, ‑var)}[1] \dicFieldCat{zool.} \dicDirectTranslationCS{vrabec} \textit{(l.~{\textLA{Passer}})}
\dicEntry[spör] \dicTerm{spör\smash{\textsuperscript{2}}} \dicIPA{{s}{\textsubring{b}}{\oe}{\textlengthmark}{\textsubring{r}}} \dicPos{adj}[4]\dicFlx{}[-5] \dicFlx{f sg nom pos} \dicLink{spar}
\dicEntry[spörð] \dicTerm{spörð} \dicIPA{{s}{\textsubring{b}}{\oe}{r}{\texttheta}} \dicPos{n} \dicFlx{pl nom} \dicLink{sparð}
\dicEntry[spörfugl] \dicTerm{spör··fugl} \dicIPA{{s}{\textsubring{b}}{\oe}{\textsubring{r}}{f}{\textscy}{\r{g}}{\textsubring{l}}} \dicPos{m}[4] \dicFlx{(‑s, ‑ar)}[3] \dicPhraseIS{spörfuglar} \dicFlx{pl} \dicFieldCat{zool.} \dicDirectTranslationCS{pěvci} \textit{(l.~{\textLA{Passeriformes}})}
\dicEntry[spörk] \dicTerm{spörk} \dicIPA{{s}{\textsubring{b}}{\oe}{\textsubring{r}}{\r{g}}} \dicPos{n} \dicFlx{pl nom} \dicLink{spark}
\dicEntry[sr.] \dicTerm{sr.} \dicPos{zkr} \dicPhraseIS{séra} \dicDirectTranslationCS{duchovní, pastor(ka)}
\dicEntry[Srí Lanka] \dicTerm{Srí Lanka} \dicIPA{{s}{r}{i}{\textlengthmark}{l}{au}{\r{\ng}}{\r{g}}{a}} \dicPos{n}[4] \dicFlx{indecl}[2] \dicFieldCat{geog.} \dicDirectTranslationCS{Srí Lanka, Šrí Lanka}
\dicEntry[Srí-Lankamaður] \dicTerm{Srí-Lanka··|maður} \dicIPA{{s}{r}{i}{\textlengthmark}{l}{au}{\r{\ng}}{\r{g}}{a}{m}{a}{ð}{\textscy}{\textsubring{r}}} \dicPos{m}[13] \dicFlx{(‑manns, ‑menn)}[2] \dicDirectTranslationCS{Srílančan(ka)}
\dicEntry[srílanskur] \dicTerm{srílanskur} \dicIPA{{s}{r}{i}{\textlengthmark}{l}{a}{n}{s}{\r{g}}{\textscy}{\textsubring{r}}} \dicPos{adj}[1] \dicFlx{(f srílönsk)}[3] \dicDirectTranslationCS{srílanský}
\dicEntry[st.] \dicTerm{st.} \dicPos{zkr} \dicPhraseIS{samtenging} \dicFieldCat{jaz.} \dicDirectTranslationCS{spojka}
\dicEntry[st.f.] \dicTerm{st.f.} \dicPos{zkr} \dicPhraseIS{staðinn fyrir} \dicDirectTranslationCS{(na) místo (čeho)}
\dicEntry[staddur] \dicTerm{staddur} \dicsymFrequent\  \dicIPA{{s}{\textsubring{d}}{a}{\textsubring{d}}{\textlengthmark}{\textscy}{\textsubring{r}}} \dicPos{adj}[2] \dicFlx{(f stödd)}[20] \dicDirectTranslationCS{situovaný, umístěný, nacházející se} \dicExampleIS{Ég var stödd uppi á fjalli þegar það fór að rigna.} \dicExampleCS{Nacházela jsem se na kopci, když začalo pršet.};  \dicPhraseIS{að svo stöddu} \dicFlx{adv} \dicDirectTranslationCS{momentálně, v~této chvíli};  \dicPhraseIS{vera í vanda staddur} \dicDirectTranslationCS{mít potíže\,/\addthin problémy};  \dicPhraseIS{vera í háska\,/\addthin hættu staddur} \dicDirectTranslationCS{být v~nebezpečí};  \dicPhraseIS{vera í nauðum staddur} \dicDirectTranslationCS{mít nesnáze};  \dicPhraseIS{vera vel\,/\addthin illa staddur} \dicDirectTranslationCS{být dobře\,/\addthin špatně situovaný}
\dicEntry[staða] \dicTerm{staða} \dicsymFrequent\  \dicIPA{{s}{\textsubring{d}}{a}{\textlengthmark}{ð}{a}} \dicPos{f}[1] \dicFlx{(stöðu, stöður)}[20] \textbf{1.} \dicSynonym{ástand} \dicDirectTranslationCS{situace, stav, poměry}  \textbf{2.} \dicSynonym{embætti} \dicDirectTranslationCS{post, (pracovní) místo\,/\addthin pozice (ve firmě ap.)} \dicExampleIS{laus staða} \dicExampleCS{volná pracovní pozice}  \textbf{3.} \dicSynonym*{standandi stelling} \dicDirectTranslationCS{postoj, postavení, poloha} \dicExampleIS{taka sér stöðu við gluggann} \dicExampleCS{postavit se u~okna}  \textbf{4.} \dicSynonym{staðsetning} \dicDirectTranslationCS{poloha, položení, umístění} \dicExampleIS{staða landsins} \dicExampleCS{položení země}  \textbf{5.} \dicFieldCat{sport.} \dicDirectTranslationCS{stav, skóre}  \textbf{6.} \dicDirectTranslationCS{stav, zůstatek (účtu ap.)}
\dicEntry[staðaldur] \dicTerm{stað··aldur} \dicIPA{{s}{\textsubring{d}}{a}{\textlengthmark}{ð}{a}{l}{\textsubring{d}}{\textscy}{\textsubring{r}}} \dicPos{n}[2] \dicFlx{(‑s)}[28] \dicPhraseIS{að staðaldri} \dicFlx{adv} \dicDirectTranslationCS{trvale, permanentně} \dicExampleIS{vinna þar að staðaldri} \dicExampleCS{pracovat tam trvale}
\dicEntry[staðalfrávik] \dicTerm{staðal··frá·vik} \dicIPA{{s}{\textsubring{d}}{a}{\textlengthmark}{ð}{a}{l}{f}{r}{au}{v}{\textsci}{\r{g}}} \dicPos{n}[2] \dicFlx{(‑s, ‑)}[5] \dicFieldCat{mat.} \dicDirectTranslationCS{směrodat\-ná\,/\addthin standardní odchylka}
\dicEntry[staðall] \dicTerm{stað|all} \dicIPA{{s}{\textsubring{d}}{a}{\textlengthmark}{ð}{a}{\textsubring{d}}{\textsubring{l}}} \dicPos{m}[6] \dicFlx{(‑als, ‑lar)}[28] \dicDirectTranslationCS{standard, norma} \dicExampleIS{gæðastaðall} \dicExampleCS{standard kvality}
\dicEntry[staðaltími] \dicTerm{staðal··tím|i} \dicIPA{{s}{\textsubring{d}}{a}{\textlengthmark}{ð}{a}{l}{t\smash{\textsuperscript{h}}}{i}{m}{\textsci}} \dicPos{m}[1] \dicFlx{(‑a)}[3] \dicDirectTranslationCS{místní\,/\addthin lokální čas}
\dicEntry[staðaratviksorð] \dicTerm{staðar··at·viks·orð} \dicIPA{{s}{\textsubring{d}}{a}{\textlengthmark}{ð}{a}{r}{a}{\textsubring{d}}{v}{\textsci}{\r{g}}{s}{\textopeno}{r}{\texttheta}} \dicPos{n}[2] \dicFlx{(‑s, ‑)}[5] \dicFieldCat{jaz.} \dicDirectTranslationCS{příslovce místa}
\dicEntry[staðarfall] \dicTerm{staðar··|fall} \dicIPA{{s}{\textsubring{d}}{a}{\textlengthmark}{ð}{a}{\textsubring{r}}{f}{a}{\textsubring{d}}{\textsubring{l}}} \dicPos{n}[2] \dicFlx{(‑falls, ‑föll)}[8] \dicFieldCat{jaz.} \dicDirectTranslationCS{6. pád, lokativ, lokál}
\dicEntry[staðarheiti] \dicTerm{staðar··heiti} \dicIPA{{s}{\textsubring{d}}{a}{\textlengthmark}{ð}{a}{\textsubring{r}}{h}{ei}{\textsubring{d}}{\textsci}} \dicPos{n}[2] \dicFlx{(‑s, ‑)}[14] \dicSynonym{örnefni} \dicDirectTranslationCS{místní název, toponymum}
\dicEntry[staðarnet] \dicTerm{staðar··net} \dicIPA{{s}{\textsubring{d}}{a}{\textlengthmark}{ð}{a}{r}{n}{\textepsilon}{\textsubring{d}}} \dicPos{n}[2] \dicFlx{(‑s, ‑)}[5] \dicFieldCat{poč.} \dicDirectTranslationCS{LAN, místní síť} \dicIndirectTranslationCS{(Local Area Network)};  \dicPhraseIS{þráðlaust staðarnet} \dicFieldCat{poč.} \dicDirectTranslationCS{WLAN, místní bezdrátová síť} \dicIndirectTranslationCS{(Wireless LAN)}
\dicEntry[staðartími] \dicTerm{staðar··tím|i} \dicIPA{{s}{\textsubring{d}}{a}{\textlengthmark}{ð}{a}{\textsubring{r}}{t\smash{\textsuperscript{h}}}{i}{m}{\textsci}} \dicPos{m}[1] \dicFlx{(‑a)}[3] \dicDirectTranslationCS{místní\,/\addthin lokální čas};  \dicPhraseIS{að staðartíma} \dicFlx{adv} \dicDirectTranslationCS{lokálního času}
\dicEntry[staðbundinn] \dicTerm{stað··bundinn} \dicIPA{{s}{\textsubring{d}}{a}{ð}{\textsubring{b}}{\textscy}{n}{\textsubring{d}}{\textsci}{\textsubring{n}}} \dicPos{adj}[6]\dicFlx{}[-2] \dicDirectTranslationCS{místní, lokální} \dicExampleIS{staðbundinn verkur} \dicExampleCS{lokální bolest}
\dicEntry[staðdeyfing] \dicTerm{stað··deyf·ing} \dicIPA{{s}{\textsubring{d}}{a}{ð}{\textsubring{d}}{ei}{v}{i}{\ng}{\r{g}}} \dicPos{f}[4] \dicFlx{(‑ar, ‑ar)}[5] \dicFieldCat{med.} \dicDirectTranslationCS{lokální umrtvení}
\dicEntry[staðfastur] \dicTerm{stað··|fastur} \dicIPA{{s}{\textsubring{d}}{a}{ð}{f}{a}{s}{\textsubring{d}}{\textscy}{\textsubring{r}}} \dicPos{adj}[1] \dicFlx{(f ‑föst)}[11] \dicSynonym{stöðugur} \dicDirectTranslationCS{pev\-ný, rezolutní, neoblomný, neochvějný}
\dicEntry[staðfesta] \dicTerm{stað··fest|a\smash{\textsuperscript{1}}} \dicIPA{{s}{\textsubring{d}}{a}{ð}{f}{\textepsilon}{s}{\textsubring{d}}{a}} \dicPos{f}[1] \dicFlx{(‑u)}[5] \textbf{1.} \dicSynonym{stöðugleiki} \dicDirectTranslationCS{pevnost, neoblomnost, neochvějnost}  \textbf{2.} \dicSynonym{staðfesting} \dicDirectTranslationCS{potvrzení, schválení}
\dicEntry[staðfesta] \dicTerm{stað··fest|a\smash{\textsuperscript{2}}} \dicsymFrequent\  \dicIPA{{s}{\textsubring{d}}{a}{ð}{f}{\textepsilon}{s}{\textsubring{d}}{a}} \dicPos{v}[2] \dicFlx{(‑i, ‑)}[12] \dicFlx{acc} \textbf{1.} \dicSynonym{sannreyna} \dicDirectTranslationCS{potvrdit, potvrzovat, stvrdit, stvrzovat, schválit, schvalovat} \dicExampleIS{staðfesta orð hennar} \dicExampleCS{potvrdit její slova}  \textbf{2.} \dicSynonym{fullgilda} \dicDirectTranslationCS{ratifikovat, schválit} \dicExampleIS{staðfesta lögin með undirskrift} \dicExampleCS{ratifikovat zákon podpisem}  \textbf{3.} \dicPhraseIS{staðfesta ráð sitt} \dicSynonym{kvænast} \dicDirectTranslationCS{oženit se, vdát se};  \dicIdiom{staðfestast}{ \dicPhraseIS{staðfestast}} \dicFlx{refl} \dicDirectTranslationCS{schválit se, ratifikovat se}
\dicEntry[staðfesting] \dicTerm{stað··fest·ing} \dicIPA{{s}{\textsubring{d}}{a}{ð}{f}{\textepsilon}{s}{\textsubring{d}}{i}{\ng}{\r{g}}} \dicPos{f}[4] \dicFlx{(‑ar, ‑ar)}[5] \textbf{1.} \dicSynonym{viðurkenning} \dicDirectTranslationCS{potvrzení, schválení, ratifikování} \dicExampleIS{staðfesting laga} \dicExampleCS{schválení zákona}  \textbf{2.} \dicSynonym{sönnun} \dicDirectTranslationCS{potvrzení, ověření}
\dicEntry[staðfestulaus] \dicTerm{stað·festu··laus} \dicIPA{{s}{\textsubring{d}}{a}{ð}{f}{\textepsilon}{s}{\textsubring{d}}{\textscy}{l}{\oe i}{s}} \dicPos{adj}[5]\dicFlx{}[-1] \textbf{1.} \dicSynonym*{laus í rásinni} \dicDirectTranslationCS{nestálý, přelétavý}  \textbf{2.} \dicSynonym{heimilislaus} \dicDirectTranslationCS{neusazený, (jsoucí) bez domova}
\dicEntry[staðfestur] \dicTerm{stað··festur} \dicIPA{{s}{\textsubring{d}}{a}{ð}{f}{\textepsilon}{s}{\textsubring{d}}{\textscy}{\textsubring{r}}} \dicPos{adj}[1]\dicFlx{}[-13] \dicDirectTranslationCS{potvrzený, schválený, ratifikovaný}
\dicEntry[staðfræði] \dicTerm{stað··fræð|i} \dicIPA{{s}{\textsubring{d}}{a}{ð}{f}{r}{a}{i}{ð}{\textsci}} \dicPos{f}[3] \dicFlx{(‑i)}[3] \dicFieldCat{geog.} \dicDirectTranslationCS{topografie, místopis}
\dicEntry[staðfugl] \dicTerm{stað··fugl} \dicIPA{{s}{\textsubring{d}}{a}{ð}{f}{\textscy}{\r{g}}{\textsubring{l}}} \dicPos{m}[4] \dicFlx{(‑s, ‑ar)}[3] \dicDirectTranslationCS{nestěhovavý pták} \dicAntonym{farfugl}
\dicEntry[staðfæra] \dicTerm{stað··fær|a} \dicIPA{{s}{\textsubring{d}}{a}{ð}{f}{a}{i}{r}{a}} \dicPos{v}[2] \dicFlx{(‑ði, ‑t)}[105] \dicFlx{acc} \dicDirectTranslationCS{lokalizovat} \dicIndirectTranslationCS{(přizpůsobit program ap. zemi, ve které se bude používat)} \dicExampleIS{staðfæra leikrit} \dicExampleCS{lokalizovat divadelní hru}
\dicEntry[staðgengill] \dicTerm{stað··geng|ill} \dicIPA{{s}{\textsubring{d}}{a}{ð}{\r{\textObardotlessj}}{ei}{\textltailn}{\r{\textObardotlessj}}{\textsci}{\textsubring{d}}{\textsubring{l}}} \dicPos{m}[6] \dicFlx{(‑ils, ‑lar)}[35] \dicDirectTranslationCS{zástupce, zástupkyně (lékaře v~jeho nepřítomnosti ap.)} \dicExampleIS{staðgengill læknis} \dicExampleCS{zástupce lékaře}
\dicEntry[staðgóður] \dicTerm{stað··|góður} \dicIPA{{s}{\textsubring{d}}{a}{ð}{\r{g}}{ou}{ð}{\textscy}{\textsubring{r}}} \dicPos{adj}[11] \dicFlx{(comp ‑betri, sup ‑bestur)}[1] \textbf{1.} \dicSynonym{traustur} \dicDirectTranslationCS{solidní, spolehlivý, poctivý} \dicExampleIS{staðgóð þekking} \dicExampleCS{solidní poznatek}  \textbf{2.} \dicSynonym{næringarríkur} \dicDirectTranslationCS{vydatný, výživný} \dicExampleIS{staðgóður matur} \dicExampleCS{výživné jídlo}
\dicEntry[staðgreiða] \dicTerm{stað··grei|ða} \dicIPA{{s}{\textsubring{d}}{a}{ð}{\r{g}}{r}{ei}{ð}{a}} \dicPos{v}[2] \dicFlx{(‑ddi, ‑tt)}[167] \dicFlx{acc} \dicDirectTranslationCS{(za)platit hotově, (za)platit v~hotovosti}
\dicEntry[staðgreiðsla] \dicTerm{stað··greiðsl|a} \dicIPA{{s}{\textsubring{d}}{a}{ð}{\r{g}}{r}{ei}{ð}{s}{\textsubring{d}}{l}{a}} \dicPos{f}[1] \dicFlx{(‑u, ‑ur)}[13] \dicDirectTranslationCS{platba v~hotovosti}
\dicEntry[staðgreiðsluafsláttur] \textls[15]{\dicTerm{stað·greiðslu··af·slátt|ur} \dicIPA{{s}\-{\textsubring{d}}\-{a}\-{ð}\-{\r{g}}\-{r}\-{ei}\-{ð}\-{s}\-{\textsubring{d}}\-{l}\-{\textscy}\-{a}\-{f}\-{s}\-{\textsubring{d}}\-{l}\-{au}\-{h}\-{\textsubring{d}}\-{\textscy}\-{\textsubring{r}}}} \dicPos{m}[12] \dicFlx{(‑ar)}[8] \dicFieldCat{ekon.} \dicDirectTranslationCS{hotovostní sleva}
\dicEntry[staðgreiðslukerfi] \dicTerm{stað·greiðslu··kerfi} \dicIPA{{s}{\textsubring{d}}{a}{ð}{\r{g}}{r}{ei}{ð}{s}{\textsubring{d}}{l}{\textscy}{c\smash{\textsuperscript{h}}}{\textepsilon}{r}{v}{\textsci}} \dicPos{n}[2] \dicFlx{(‑s, ‑)}[14] \dicPhraseIS{staðgreiðslukerfi skatta} \dicFieldCat{ekon.} \dicDirectTranslationCS{srážková daň u~zdroje}
\dicEntry[staðgreiðslukrafa] \textls[15]{\dicTerm{stað·greiðslu··|krafa} \dicIPA{{s}\-{\textsubring{d}}\-{a}\-{ð}\-{\r{g}}\-{r}\-{ei}\-{ð}\-{s}\-{\textsubring{d}}\-{l}\-{\textscy}\-{k\smash{\textsuperscript{h}}}\-{r}\-{a}\-{v}\-{a}\-} \dicPos{f}[1] \dicFlx{(‑kröfu, ‑kröfur)}[8]} \dicFieldCat{ekon.} \dicDirectTranslationCS{dobírka, dobírková služba}
\dicEntry[staðhæfa] \dicTerm{stað··hæf|a} \dicIPA{{s}{\textsubring{d}}{a}{\textlengthmark}{\texttheta}{h}{a}{i}{v}{a}} \dicPos{v}[2] \dicFlx{(‑ði, ‑t)}[110] \dicFlx{acc} \dicSynonym{fullyrða} \dicDirectTranslationCS{tvrdit, prohlásit, prohlašovat} \dicExampleIS{Lögreglan staðhæfir að um morð hafi verið að ræða.} \dicExampleCS{Policie prohlašuje, že se jedná o~vraždu.}
\dicEntry[staðhæfing] \dicTerm{stað··hæf·ing} \dicIPA{{s}{\textsubring{d}}{a}{\textlengthmark}{\texttheta}{h}{a}{i}{v}{i}{\ng}{\r{g}}} \dicPos{f}[4] \dicFlx{(‑ar, ‑ar)}[5] \dicSynonym{fullyrðing} \dicDirectTranslationCS{tvrzení, prohlášení, výrok} \dicExampleIS{sannreyna staðhæfinguna} \dicExampleCS{prověřit tvrzení}
\dicEntry[staðhættir] \dicTerm{stað··hættir} \dicIPA{{s}{\textsubring{d}}{a}{\textlengthmark}{\texttheta}{h}{a}{i}{h}{\textsubring{d}}{\textsci}{\textsubring{r}}} \dicPos{m}[12] \dicFlx{pl}[6] \dicSynonym{aðstæður} \dicDirectTranslationCS{místní podmínky} \dicExampleIS{náttúra og staðhættir} \dicExampleCS{příroda a~místní podmínky}
\dicEntry[staðið] \dicTerm{staðið} \dicIPA{{s}{\textsubring{d}}{a}{\textlengthmark}{ð}{\textsci}{\texttheta}} \dicPos{v} \dicFlx{supin} \dicLink{standa}
\dicEntry[staðla] \dicTerm{staðl|a} \dicIPA{{s}{\textsubring{d}}{a}{ð}{l}{a}} \dicPos{v}[1] \dicFlx{(‑aði)}[13] \dicFlx{acc} \dicDirectTranslationCS{standardizovat, normalizovat} \dicExampleIS{staðla framleiðsluna} \dicExampleCS{standardizovat výrobu}
\dicEntry[staðlaður] \dicTerm{staðl··aður} \dicIPA{{s}{\textsubring{d}}{a}{ð}{l}{a}{ð}{\textscy}{\textsubring{r}}} \dicPos{adj}[3] \dicFlx{(f stöðluð)}[1] \dicDirectTranslationCS{standardizovaný, normalizovaný}
\dicEntry[staðlaus] \dicTerm{stað··laus} \dicIPA{{s}{\textsubring{d}}{a}{ð}{l}{\oe i}{s}} \dicPos{adj}[5]\dicFlx{}[-1] \dicDirectTranslationCS{neopodstatněný, nepodložený, absurdní} \dicExampleIS{staðlausar hugmyndir} \dicExampleCS{nepodložené myšlenky}
\dicEntry[staðna] \dicTerm{staðn|a} \dicIPA{{s}{\textsubring{d}}{a}{ð}{n}{a}} \dicPos{v}[1] \dicFlx{(‑aði)}[13] \dicDirectTranslationCS{stagnovat, ustrnout} \dicExampleIS{Hann er farinn að staðna.} \dicExampleCS{Začíná stagnovat.};  \dicPhraseIS{e‑að staðnar} \dicDirectTranslationCS{(co) stagnuje}
\dicEntry[staðnaður] \dicTerm{staðn··aður} \dicIPA{{s}{\textsubring{d}}{a}{ð}{n}{a}{ð}{\textscy}{\textsubring{r}}} \dicPos{adj}[3] \dicFlx{(f stöðnuð)}[1] \dicDirectTranslationCS{stagnující, ustrnulý, zkostnatělý}
\dicEntry[staðnæmast] \dicTerm{stað··næm|ast} \dicsymFrequent\  \dicIPA{{s}{\textsubring{d}}{a}{ð}{n}{a}{i}{m}{a}{s}{\textsubring{d}}} \dicPos{v}[2] \dicFlx{(‑dist, ‑st)}[205] \dicFlx{refl} \textbf{1.} \dicSynonym{stansa} \dicDirectTranslationCS{zastavit (se) (auto ap.)} \dicIndirectTranslationCS{(ustat v~pohybu)} \dicExampleIS{staðnæmast við húsið} \dicExampleCS{zastavit u~domu}  \textbf{2.} \dicDirectTranslationCS{zastavit (se) (u~určité věty ap.)} \dicIndirectTranslationCS{(přerušit činnost\,/\addthin pohyb)}
\dicEntry[staðráðinn] \dicTerm{stað··ráðinn} \dicIPA{{s}{\textsubring{d}}{a}{ð}{r}{au}{ð}{\textsci}{\textsubring{n}}} \dicPos{adj}[6]\dicFlx{}[-2] \dicDirectTranslationCS{pevně rozhodnutý, odhodlaný};  \dicPhraseIS{vera staðráðinn í e‑u} \dicDirectTranslationCS{být pevně rozhodnutý k~(čemu)} \dicExampleIS{Hún segist vera staðráðin í því að ljúka náminu í vor.} \dicExampleCS{Říká, že je pevně rozhodnutá ukončit studia na jaře.}
\dicEntry[staðreynd] \dicTerm{stað··reynd} \dicsymFrequent\  \dicIPA{{s}{\textsubring{d}}{a}{ð}{r}{ei}{n}{\textsubring{d}}} \dicPos{f}[7] \dicFlx{(‑ar, ‑ir)}[1] \dicDirectTranslationCS{skutečnost, fakt} \dicExampleIS{vísindalegar staðreyndir} \dicExampleCS{vědecká fakta}
\dicEntry[staðset] \dicTerm{stað··set} \dicIPA{{s}{\textsubring{d}}{a}{ð}{s}{\textepsilon}{\textsubring{d}}} \dicPos{v} \dicFlx{ind praes sg 1 pers} \dicLink{staðsetja}
\dicEntry[staðsetja] \dicTerm{stað··|setja} \dicIPA{{s}{\textsubring{d}}{a}{ð}{s}{\textepsilon}{\textsubring{d}}{j}{a}} \dicPos{v}[4] \dicFlx{(‑set, ‑setti, ‑settum, ‑setti, ‑sett)}[26] \dicFlx{acc} \dicDirectTranslationCS{umístit, umísťovat, situovat, lokalizovat} \dicExampleIS{staðsetja húsið} \dicExampleCS{umístit dům}
\dicEntry[staðsetning] \dicTerm{stað··set·ning} \dicsymFrequent\  \dicIPA{{s}{\textsubring{d}}{a}{ð}{s}{\textepsilon}{h}{\textsubring{d}}{n}{i}{\ng}{\r{g}}} \dicPos{f}[4] \dicFlx{(‑ar, ‑ar)}[5] \dicDirectTranslationCS{umístění, situování, lokalizování} \dicExampleIS{staðsetning byggingarinnar} \dicExampleCS{umístění budovy}
\dicEntry[staðsett] \dicTerm{stað··sett} \dicIPA{{s}{\textsubring{d}}{a}{ð}{s}{\textepsilon}{h}{\textsubring{d}}} \dicPos{v} \dicFlx{supin} \dicLink{staðsetja}
\dicEntry[staðsetti] \dicTerm{stað··setti} \dicIPA{{s}{\textsubring{d}}{a}{ð}{s}{\textepsilon}{h}{\textsubring{d}}{\textsci}} \dicPos{v} \dicFlx{ind\,/\addthin con pf sg 1 pers} \dicLink{staðsetja}
\dicEntry[staðsettum] \dicTerm{stað··settum} \dicIPA{{s}{\textsubring{d}}{a}{ð}{s}{\textepsilon}{h}{\textsubring{d}}{\textscy}{\textsubring{m}}} \dicPos{v} \dicFlx{ind pf pl 1 pers} \dicLink{staðsetja}
\dicEntry[staðsettur] \dicTerm{stað··settur} \dicIPA{{s}{\textsubring{d}}{a}{ð}{s}{\textepsilon}{h}{\textsubring{d}}{\textscy}{\textsubring{r}}} \dicPos{adj}[1]\dicFlx{}[-13] \dicDirectTranslationCS{umístěný, situovaný, lokalizovaný}
\dicEntry[staður] \dicTerm{stað|ur\smash{\textsuperscript{1}}} \dicsymFrequent\  \dicIPA{{s}{\textsubring{d}}{a}{\textlengthmark}{ð}{\textscy}{\textsubring{r}}} \dicPos{m}[10] \dicFlx{(‑ar, ‑ir)}[14] \textbf{1.} \dicDirectTranslationCS{místo, prostor}  \textbf{2.} \dicLangCat{zast.} \dicSynonym{bær} \dicDirectTranslationCS{osada, (malé) město};  \dicIdiom{staður}{ \dicPhraseIS{alls staðar, allsstaðar}} \dicFlx{adv} \dicDirectTranslationCS{všude};  \dicPhraseIS{annars staðar, annarsstaðar} \dicFlx{adv} \dicDirectTranslationCS{jinde, někde jinde};  \dicPhraseIS{einhvers staðar} \dicFlx{adv} \dicDirectTranslationCS{někde jinde, na jiném místě};  \dicPhraseIS{e‑að átti sér stað (í fyrra)} \dicDirectTranslationCS{(co) se uskutečnilo (loni)};  \dicPhraseIS{e‑að er til staðar} \dicDirectTranslationCS{(co) je k~dispozici (telefon v~hotelovém pokoji ap.)};  \dicPhraseIS{ganga e‑m í föður\,/\addthin móður stað} \dicDirectTranslationCS{být pro (koho) jako otec\,/\addthin matka};  \dicPhraseIS{gleyma stað og stund} \dicLangCat{přen.} \dicDirectTranslationCS{zapomenout na čas a~místo};  \dicPhraseIS{fara\,/\addthin halda\,/\addthin leggja af stað} \dicSynonym*{hefja ferð} \dicDirectTranslationCS{vydat se, vypravit se, vyrazit, vyjít, vyjet};  \dicPhraseIS{finna e‑u stað} \dicDirectTranslationCS{najít pro (co) zdůvodnění};  \dicPhraseIS{fyrst í stað} \dicFlx{adv} \dicDirectTranslationCS{nejdříve, zpočátku};  \dicPhraseIS{færa e‑ð úr stað} \dicDirectTranslationCS{přesunout (co), přemístit (co)};  \dicPhraseIS{í alla staði} \dicFlx{adv} \dicDirectTranslationCS{v~každém ohledu};  \dicPhraseIS{í annan stað} \dicFlx{adv} \dicSynonym{hins vegar} \dicDirectTranslationCS{za druhé} \dicExampleIS{í fyrsta lagi og í annan stað} \dicExampleCS{za prvé a~za druhé};  \dicPhraseIS{í stað e‑s} \dicFlx{prep} \dicDirectTranslationCS{místo (čeho), namísto (čeho)};  \dicPhraseIS{í staðinn} \dicFlx{adv} \dicDirectTranslationCS{místo\,/\addthin namísto (toho)};  \dicPhraseIS{í staðinn fyrir e‑ð} \dicFlx{prep} \dicDirectTranslationCS{místo (čeho), namísto (čeho)};  \dicPhraseIS{koma e‑u af stað} \dicDirectTranslationCS{začít (co), zahájit (co)};  \dicPhraseIS{neins staðar} \dicFlx{adv} \dicDirectTranslationCS{nikde};  \dicPhraseIS{nema staðar} \dicSynonym{stansa} \dicDirectTranslationCS{zastavit} \dicExampleIS{Hann nam staðar.} \dicExampleCS{Zastavil.};  \dicPhraseIS{standa í stað} \dicLangCat{přen.} \dicDirectTranslationCS{stát na místě, stagnovat};  \dicPhraseIS{sums staðar, sumstaðar} \dicFlx{adv} \dicDirectTranslationCS{sem tam, místy, v~určitých místech};  \dicPhraseIS{þegar í stað} \dicFlx{adv} \dicDirectTranslationCS{okamžitě, hned}
\dicEntry[staður] \dicTerm{staður\smash{\textsuperscript{2}}} \dicIPA{{s}{\textsubring{d}}{a}{\textlengthmark}{ð}{\textscy}{\textsubring{r}}} \dicPos{adj}[2] \dicFlx{(f stöð)}[7] \dicDirectTranslationCS{vzpurný, vzdorovitý} \dicIndirectTranslationCS{(o~koni)}
\dicEntry[staðvindur] \dicTerm{stað··vind|ur} \dicIPA{{s}{\textsubring{d}}{a}{ð}{v}{\textsci}{n}{\textsubring{d}}{\textscy}{\textsubring{r}}} \dicPos{m}[6] \dicFlx{(‑s, ‑ar)}[4] \dicFieldCat{meteo.} \dicDirectTranslationCS{pasát}
\dicEntry[stafa] \dicTerm{staf|a} \dicsymFrequent\  \dicIPA{{s}{\textsubring{d}}{a}{\textlengthmark}{v}{a}} \dicPos{v}[1] \dicFlx{(‑aði)}[46] \dicFlx{acc\,/\addthin dat} \textbf{1.} \dicFlx{acc} \dicSynonym{stafsetja} \dicDirectTranslationCS{slabikovat, hláskovat} \dicExampleIS{stafa orðið} \dicExampleCS{hláskovat slovo}  \textbf{2.} \dicFlx{dat} \dicSynonym*{setja röndum} \dicDirectTranslationCS{vyzařovat, zářit};  \dicPhraseIS{það stafar birtu af e‑u} \dicFlx{impers} \dicDirectTranslationCS{z~(čeho) vychází světlo} \dicExampleIS{Það stafar birtu af lampanum.} \dicExampleCS{Z~lampy vychází světlo.};  \dicPhraseIS{það stafar e‑u frá e‑m} \dicFlx{impers} \dicDirectTranslationCS{(kdo) září (čím) (mladý člověk zdravím ap.)}  \textbf{3.} \dicDirectTranslationCS{způsobit, způsobovat, zapříčinit, zapříčiňovat}  \textbf{4.} \dicPhraseIS{e‑m stafar hætta af e‑u} \dicFlx{impers} \dicDirectTranslationCS{(co) představuje nebezpečí pro (koho)} \dicExampleIS{Það stafar hætta af ungum ökumönnum. } \dicExampleCS{Mladí řidiči představují nebezpečí.}
\dicEntry[stafagerð] \dicTerm{stafa··gerð} \dicIPA{{s}{\textsubring{d}}{a}{\textlengthmark}{v}{a}{\r{\textObardotlessj}}{\textepsilon}{r}{\texttheta}} \dicPos{f}[7] \dicFlx{(‑ar, ‑ir)}[1] \dicDirectTranslationCS{druh\,/\addthin typ písma}
\dicEntry[stafalogn] \dicTerm{stafa··logn} \dicIPA{{s}{\textsubring{d}}{a}{\textlengthmark}{v}{a}{l}{\textopeno}{\r{g}}{\textsubring{n}}} \dicPos{n}[2] \dicFlx{(‑s)}[2] \dicSynonym{blæjalogn} \dicDirectTranslationCS{úplné bezvětří}
\dicEntry[Stafangur] \dicTerm{Staf··ang|ur} \dicIPA{{s}{\textsubring{d}}{a}{\textlengthmark}{v}{au}{\ng}{\r{g}}{\textscy}{\textsubring{r}}} \dicPos{m}[5] \dicFlx{(‑urs)}[6] \dicFieldCat{geog.} \dicDirectTranslationCS{Stavanger} \dicIndirectTranslationCS{(město v~Norsku)}
\dicEntry[stafasett] \dicTerm{stafa··sett} \dicIPA{{s}{\textsubring{d}}{a}{\textlengthmark}{v}{a}{s}{\textepsilon}{h}{\textsubring{d}}} \dicPos{n}[2] \dicFlx{(‑s, ‑)}[5] \dicFieldCat{poč.} \dicDirectTranslationCS{znaková sada}
\dicEntry[stafastrengur] \dicTerm{stafa··streng|ur} \dicIPA{{s}{\textsubring{d}}{a}{\textlengthmark}{v}{a}{s}{\textsubring{d}}{r}{ei}{\ng}{\r{g}}{\textscy}{\textsubring{r}}} \dicPos{m}[9] \dicFlx{(‑s\,/\addthin ‑jar, ‑ir)}[26] \dicFieldCat{poč.} \dicDirectTranslationCS{textový řetězec}
\dicEntry[stafatafla] \dicTerm{stafa··|tafla} \dicIPA{{s}{\textsubring{d}}{a}{\textlengthmark}{v}{a}{t\smash{\textsuperscript{h}}}{a}{\textsubring{b}}{l}{a}} \dicPos{f}[1] \dicFlx{(‑töflu, ‑töflur)}[20] \dicFieldCat{poč.} \dicDirectTranslationCS{tabulka znaků}
\dicEntry[stafavíxl] \dicTerm{stafa··víxl} \dicIPA{{s}{\textsubring{d}}{a}{\textlengthmark}{v}{a}{v}{i}{x}{s}{\textsubring{d}}{\textsubring{l}}} \dicPos{n}[2] \dicFlx{pl}[1] \dicFieldCat{jaz.} \dicDirectTranslationCS{metateze} \dicIndirectTranslationCS{(přesmyknutí hlásek ve slově)}
\dicEntry[stafbil] \dicTerm{staf··bil} \dicIPA{{s}{\textsubring{d}}{a}{v}{\textsubring{b}}{\textsci}{\textsubring{l}}} \dicPos{n}[2] \dicFlx{(‑s, ‑)}[5] \dicFieldCat{poč.} \dicDirectTranslationCS{mezera, prázdný znak}
\dicEntry[stafkrókur] \dicTerm{staf··krók|ur} \dicIPA{{s}{\textsubring{d}}{a}{v}{k\smash{\textsuperscript{h}}}{r}{ou}{\r{g}}{\textscy}{\textsubring{r}}} \dicPos{m}[6] \dicFlx{(‑s, ‑ar)}[22] \textbf{1.} \dicDirectTranslationCS{patka (písmena)}  \textbf{2.} \dicLangCat{přen.} \dicDirectTranslationCS{čárka, zmínka, jedno jediné slovo} \dicIndirectTranslationCS{(nejčastěji v~záporu)} \dicExampleIS{enginn stafkrókur} \dicExampleCS{ani čárka}  \textbf{3.} \dicFieldCat{jaz.} \dicDirectTranslationCS{cedilla, sedila} \dicIndirectTranslationCS{(diakritické znaménko)}
\dicEntry[stafla] \dicTerm{stafl|a} \dicIPA{{s}{\textsubring{d}}{a}{\textsubring{b}}{l}{a}} \dicPos{v}[1] \dicFlx{(‑aði)}[13] \dicFlx{dat} \dicSynonym*{setja í stafla} \dicDirectTranslationCS{(na)skládat na sebe} \dicExampleIS{stafla bókunum} \dicExampleCS{naskládat knížky na sebe}
\dicEntry[stafli] \dicTerm{stafl|i} \dicIPA{{s}{\textsubring{d}}{a}{\textsubring{b}}{l}{\textsci}} \dicPos{m}[1] \dicFlx{(‑a, ‑ar)}[8] \dicDirectTranslationCS{stoh, hromada, pyramida} \dicExampleIS{stafli af blöðum} \dicExampleCS{stoh novin}
\dicEntry[stafmerki] \dicTerm{staf··merki} \dicIPA{{s}{\textsubring{d}}{a}{v}{m}{\textepsilon}{\textsubring{r}}{\r{\textObardotlessj}}{\textsci}} \dicPos{n}[2] \dicFlx{(‑s, ‑)}[16] \dicFieldCat{jaz.} \dicDirectTranslationCS{diakritické znaménko}
\dicEntry[stafn] \dicTerm{stafn} \dicIPA{{s}{\textsubring{d}}{a}{\textsubring{b}}{\textsubring{n}}} \dicPos{m}[4] \dicFlx{(‑s, ‑ar)}[4] \textbf{1.} \dicFieldCat{nám.} \dicSynonym*{framendi skips} \dicDirectTranslationCS{příď} \dicExampleIS{stafninn á skipinu} \dicExampleCS{příď lodi};  \dicPhraseIS{land fyrir stafni} \dicDirectTranslationCS{země na obzoru}  \textbf{2.} \dicSynonym{gafl} \dicDirectTranslationCS{štít (domu ap.)} \dicExampleIS{stafninn á húsinu} \dicExampleCS{štít domu};  \dicPhraseIS{hafa e‑ð fyrir stafni} \dicDirectTranslationCS{mít (co) rozdělaného (práci ap.)}
\dicEntry[stafnlíkan] \dicTerm{stafn··lík|an} \dicIPA{{s}{\textsubring{d}}{a}{\textsubring{b}}{n}{l}{i}{\r{g}}{a}{\textsubring{n}}} \dicPos{n}[2] \dicFlx{(‑ans, ‑ön)}[8] \dicFieldCat{nám.} \dicDirectTranslationCS{galionová figura}
\dicEntry[stafprik] \dicTerm{staf··prik} \dicIPA{{s}{\textsubring{d}}{a}{f}{p\smash{\textsuperscript{h}}}{r}{\textsci}{\r{g}}} \dicPos{n}[2] \dicFlx{(‑s, ‑)}[5] \dicDirectTranslationCS{hůl, špacírka}
\dicEntry[stafréttur] \dicTerm{staf··réttur} \dicIPA{{s}{\textsubring{d}}{a}{v}{r}{j}{\textepsilon}{h}{\textsubring{d}}{\textscy}{\textsubring{r}}} \dicPos{adj}[1]\dicFlx{}[-10] \dicDirectTranslationCS{doslovný}
\dicEntry[stafróf] \dicTerm{staf··róf} \dicIPA{{s}{\textsubring{d}}{a}{v}{r}{ou}{f}} \dicPos{n}[2] \dicFlx{(‑s, ‑)}[5] \dicDirectTranslationCS{abeceda} \dicExampleIS{íslenska stafrófið} \dicExampleCS{islandská abeceda}
\dicEntry[stafrófskver] \dicTerm{staf·rófs··kver} \dicIPA{{s}{\textsubring{d}}{a}{v}{r}{ou}{f}{s}{k\smash{\textsuperscript{h}}}{v}{\textepsilon}{\textsubring{r}}} \dicPos{n}[2] \dicFlx{(‑s, ‑)}[5] \dicDirectTranslationCS{slabikář}
\dicEntry[stafrófsröð] \dicTerm{staf·rófs··|röð} \dicIPA{{s}{\textsubring{d}}{a}{v}{r}{ou}{f}{s}{r}{\oe}{\texttheta}} \dicPos{f}[7] \dicFlx{(‑raðar)}[19] \dicDirectTranslationCS{abecední pořadí};  \dicPhraseIS{raða í stafrófsröð} \dicDirectTranslationCS{seřadit podle abecedy}
\dicEntry[stafrænn] \dicTerm{staf··rænn} \dicIPA{{s}{\textsubring{d}}{a}{v}{r}{a}{i}{\textsubring{d}}{\textsubring{n}}} \dicPos{adj}[7]\dicFlx{}[-1] \dicFieldCat{poč.} \dicDirectTranslationCS{digitální}
\dicEntry[stafset] \dicTerm{staf··set} \dicIPA{{s}{\textsubring{d}}{a}{f}{s}{\textepsilon}{\textsubring{d}}} \dicPos{v} \dicFlx{ind praes sg 1 pers} \dicLink{stafsetja}
\dicEntry[stafsetja] \dicTerm{staf··|setja} \dicIPA{{s}{\textsubring{d}}{a}{f}{s}{\textepsilon}{\textsubring{d}}{j}{a}} \dicPos{v}[4] \dicFlx{(‑set, ‑setti, ‑settum, ‑setti, ‑sett)}[26] \dicFlx{acc} \dicDirectTranslationCS{hláskovat, psát pravopisně}
\dicEntry[stafsetning] \dicTerm{staf··set·ning} \dicIPA{{s}{\textsubring{d}}{a}{f}{s}{\textepsilon}{h}{\textsubring{d}}{n}{i}{\ng}{\r{g}}} \dicPos{f}[4] \dicFlx{(‑ar)}[7] \dicDirectTranslationCS{pravopis, ortografie}
\dicEntry[stafsetningarorðabók] \dicTerm{staf·setningar··orða·|bók} \dicIPA{{s}\-{\textsubring{d}}\-{a}\-{f}\-{s}\-{\textepsilon}\-{h}\-{\textsubring{d}}\-{n}\-{i}\-{\ng}\-{\r{g}}\-{a}\-{r}\-{\textopeno}\-{r}\-{ð}\-{a}\-{\textsubring{b}}\-{ou}\-{\r{g}}\-} \dicPos{f}[8] \dicFlx{(‑bókar, ‑bækur)}[5] \dicDirectTranslationCS{pravopisný slovník, slovník spisovného jazyka}
\dicEntry[stafsetningarregla] \dicTerm{staf·setningar··regl|a} \dicIPA{{s}\-{\textsubring{d}}\-{a}\-{f}\-{s}\-{\textepsilon}\-{h}\-{\textsubring{d}}\-{n}\-{i}\-{\ng}\-{\r{g}}\-{a}\-{r}\-{\textepsilon}\-{\r{g}}\-{l}\-{a}\-} \dicPos{f}[1] \dicFlx{(‑u, ‑ur)}[19] \dicDirectTranslationCS{pravopisné pravidlo}
\dicEntry[stafsett] \dicTerm{staf··sett} \dicIPA{{s}{\textsubring{d}}{a}{f}{s}{\textepsilon}{h}{\textsubring{d}}} \dicPos{v} \dicFlx{supin} \dicLink{stafsetja}
\dicEntry[stafsetti] \dicTerm{staf··setti} \dicIPA{{s}{\textsubring{d}}{a}{f}{s}{\textepsilon}{h}{\textsubring{d}}{\textsci}} \dicPos{v} \dicFlx{ind\,/\addthin con pf sg 1 pers} \dicLink{stafsetja}
\dicEntry[stafsettum] \dicTerm{staf··settum} \dicIPA{{s}{\textsubring{d}}{a}{f}{s}{\textepsilon}{h}{\textsubring{d}}{\textscy}{\textsubring{m}}} \dicPos{v} \dicFlx{ind pf pl 1 pers} \dicLink{stafsetja}
\dicEntry[stafur] \dicTerm{staf|ur} \dicsymFrequent\  \dicIPA{{s}{\textsubring{d}}{a}{\textlengthmark}{v}{\textscy}{\textsubring{r}}} \dicPos{m}[9] \dicFlx{(‑s, ‑ir)}[9] \textbf{1.} \dicSynonym{prik} \dicDirectTranslationCS{hůl, hole}  \textbf{2.} \dicSynonym{bókstafur} \dicDirectTranslationCS{písmeno, litera, grafém} \dicExampleIS{skrifa það með stórum stöfum} \dicExampleCS{napsat to velkými písmeny}  \textbf{3.} \dicSynonym*{dyrastafur} \dicDirectTranslationCS{futro (dveří), dveřní rám}  \textbf{4.} \dicFieldCat{anat.} \dicDirectTranslationCS{tyčinka} \dicIndirectTranslationCS{(zraková buňka sítnice)};  \dicIdiom{stafur}{ \dicPhraseIS{e‑að stendur eins og stafur á bók}} \dicLangCat{přen.} \dicDirectTranslationCS{(co) je pevné jako skála, (co) je spolehlivé}; { \dicPhraseIS{e‑r fellur í stafi}} \dicLangCat{přen.} \dicDirectTranslationCS{(komu) spadne čelist, (kdo) oněmí úžasem};  \dicPhraseIS{lenda á milli stafs og hurðar} \dicLangCat{přen.} \dicDirectTranslationCS{dostat se do úzkých};  \dicPhraseIS{það er enginn stafur fyrir e‑u} \dicFlx{impers} \dicDirectTranslationCS{(co) nemá hlavu ani patu};  \dicPhraseIS{þetta eru staðlausir stafir} \dicLangCat{přen.} \dicDirectTranslationCS{to jsou nesmysly, nemá to hlavu ani patu}
\dicEntry[stafvilla] \dicTerm{staf··vill|a} \dicIPA{{s}{\textsubring{d}}{a}{v}{\textlengthmark}{\textsci}{\textsubring{d}}{l}{a}} \dicPos{f}[1] \dicFlx{(‑u, ‑ur)}[13] \dicDirectTranslationCS{pravopisná chyba}
\dicEntry[stag] \dicTerm{stag} \dicIPA{{s}{\textsubring{d}}{a}{\textlengthmark}{x}} \dicPos{n}[2] \dicFlx{(‑s, stög)}[8] \textbf{1.} \dicDirectTranslationCS{kotevní\,/\addthin kotvicí lano (u~vysílače ap.)}  \textbf{2.} \dicSynonym{þvottasnúra} \dicDirectTranslationCS{prádelní šňůra}  \textbf{3.} \dicFieldCat{nám.} \dicDirectTranslationCS{stěh}
\dicEntry[staga] \dicTerm{stag|a} \dicIPA{{s}{\textsubring{d}}{a}{\textlengthmark}{\textbabygamma}{a}} \dicPos{v}[1] \dicFlx{(‑aði)}[13] \dicFlx{acc} \dicDirectTranslationCS{(za)látat, (vy)spravit (ponožky ap.)} \dicExampleIS{staga sokka} \dicExampleCS{zalátat ponožky};  \dicPhraseIS{staga í e‑u} \dicDirectTranslationCS{zalátat (co)};  \dicIdiom{stagast}[á]{ \dicPhraseIS{stagast á e‑u}} \dicFlx{refl} \dicDirectTranslationCS{omílat (co), přetřásat (co)}
\dicEntry[stagl] \dicTerm{stagl} \dicIPA{{s}{\textsubring{d}}{a}{\r{g}}{\textsubring{l}}} \dicPos{n}[2] \dicFlx{(‑s)}[2] \textbf{1.} \dicSynonym*{stögun} \dicDirectTranslationCS{vyspravení, (za)látání (šatů ap.)}  \textbf{2.} \dicSynonym{klifun} \dicDirectTranslationCS{omílání, přetřásání, opakování (tématu ap.)}  \textbf{3.} \dicSynonym{staut} \dicDirectTranslationCS{zadrhávání se (o~čtení začátečníka ap.)}
\dicEntry[stagla] \dicTerm{stagl|a} \dicIPA{{s}{\textsubring{d}}{a}{\r{g}}{l}{a}} \dicPos{v}[1] \dicFlx{(‑aði)}[13] \textbf{1.} \dicSynonym*{gera við} \dicDirectTranslationCS{(za)látat, vyspravit} \dicExampleIS{stagla í sokka} \dicExampleCS{zalátat ponožky}  \textbf{2.} \dicSynonym{stauta} \dicDirectTranslationCS{zadrhávat se, koktat (při čtení ap.)}  \textbf{3.} \dicPhraseIS{stagla á e‑u} \dicDirectTranslationCS{omílat (co), přetřásat (co), opakovat (co) (do nekonečna)};  \dicIdiom{staglast}[á]{ \dicPhraseIS{staglast á e‑u}} \dicFlx{refl} \dicDirectTranslationCS{omílat (co), přetřásat (co), opakovat (co) do nekonečna}
\dicEntry[staka] \dicTerm{staka} \dicIPA{{s}{\textsubring{d}}{a}{\textlengthmark}{\r{g}}{a}} \dicPos{f}[1] \dicFlx{(stöku, stökur)}[8] \dicFieldCat{lit.} \dicDirectTranslationCS{samostatně stojící kvartet} \dicIndirectTranslationCS{(čtyřveršová strofa)}
\dicEntry[stakk] \dicTerm{stakk} \dicIPA{{s}{\textsubring{d}}{a}{h}{\r{g}}} \dicPos{v} \dicFlx{ind pf sg 1 pers} \dicLink{stinga}
\dicEntry[stakka] \dicTerm{stakk|a} \dicIPA{{s}{\textsubring{d}}{a}{h}{\r{g}}{a}} \dicPos{v}[1] \dicFlx{(‑aði)}[13] \dicFlx{acc} \dicSynonym*{hlaða upp} \dicDirectTranslationCS{(na)skládat, (na)vršit (na sebe)}
\dicEntry[stakkaskipti] \dicTerm{stakka··skipti} \dicIPA{{s}{\textsubring{d}}{a}{h}{\r{g}}{a}{s}{\r{\textObardotlessj}}{\textsci}{f}{\textsubring{d}}{\textsci}} \dicPos{n}[2] \dicFlx{pl}[19] \dicPhraseIS{e‑að tekur stakkaskiptum} \dicDirectTranslationCS{(co) se mění, (co) prochází změnami} \dicExampleIS{Umhverfið hefur tekið miklum stakkaskiptum.} \dicExampleCS{Životní prostředí se znatelně změnilo.}
\dicEntry[stakkur] \dicTerm{stakk|ur} \dicIPA{{s}{\textsubring{d}}{a}{h}{\r{g}}{\textscy}{\textsubring{r}}} \dicPos{m}[6] \dicFlx{(‑s, ‑ar)}[25] \dicSynonym{anórakkur} \dicDirectTranslationCS{anorak, dlouhá bunda s~kapucí};  \dicPhraseIS{sníða sér stakk eftir vexti} \dicLangCat{přen.} \dicDirectTranslationCS{nastříhat si vestu podle svého vzrůstu; nebrat si toho víc, než člověk unese} \dicIndirectTranslationCS{(chovat se v~souladu se svými možnostmi)};  \dicPhraseIS{vera í stakk búinn til að (gera e‑ð)} \dicDirectTranslationCS{být schopen (udělat (co))}
\dicEntry[stakmál] \dicTerm{stak··mál} \dicIPA{{s}{\textsubring{d}}{a}{\textlengthmark}{\r{g}}{m}{au}{\textsubring{l}}} \dicPos{n}[2] \dicFlx{(‑s, ‑)}[5] \dicFieldCat{jaz.} \dicDirectTranslationCS{izolovaný jazyk}
\dicEntry[stakstæður] \dicTerm{stak··stæður} \dicIPA{{s}{\textsubring{d}}{a}{\textlengthmark}{\r{g}}{s}{\textsubring{d}}{a}{i}{ð}{\textscy}{\textsubring{r}}} \dicPos{adj}[2]\dicFlx{}[-6] \dicDirectTranslationCS{osamocený, osamělý, izolovaný}
\dicEntry[stakur] \dicTerm{stakur} \dicIPA{{s}{\textsubring{d}}{a}{\textlengthmark}{\r{g}}{\textscy}{\textsubring{r}}} \dicPos{adj}[1] \dicFlx{(f stök)}[2] \textbf{1.} \dicSynonym*{einn og sér} \dicDirectTranslationCS{osamocený, osamělý, jednotlivý} \dicExampleIS{stakur klettur} \dicExampleCS{osamocená skála}  \textbf{2.} \dicSynonym*{ekki í setti} \dicDirectTranslationCS{lichý} \dicIndirectTranslationCS{(objevující se samostatně bez svého protějšku)} \dicExampleIS{stakur skór} \dicExampleCS{lichá bota}  \textbf{3.} \dicSynonym{frábær} \dicDirectTranslationCS{výborný, jedinečný} \dicExampleIS{með stakri nærgætni} \dicExampleCS{s~jedinečnou ohleduplností};  \dicIdiom{stakur}{ \dicPhraseIS{e‑að er í stakasta lagi}} \dicDirectTranslationCS{(co) je v~nejlepším pořádku}; { \dicPhraseIS{eiga í stökustu vandræðum með að (gera e‑ð)}} \dicLangCat{přen.} \dicDirectTranslationCS{mít obrovské problémy (udělat (co))};  \dicPhraseIS{skilja ekki stakt orð} \dicLangCat{přen.} \dicDirectTranslationCS{nerozumět jednomu jedinému slovu}; { \dicPhraseIS{stöku sinnum}} \dicFlx{adv} \dicDirectTranslationCS{výjimečně};  \dicProverb\  \dicPhraseIS{Sjaldan er ein báran stök.} \dicLangCat{přís.} \dicDirectTranslationCS{Neštěstí nechodí nikdy samo.}
\dicEntry[stal] \dicTerm{stal} \dicIPA{{s}{\textsubring{d}}{a}{\textlengthmark}{\textsubring{l}}} \dicPos{v} \dicFlx{ind pf sg 1 pers} \dicLink{stela}
\dicEntry[staldra] \dicTerm{staldr|a} \dicIPA{{s}{\textsubring{d}}{a}{l}{\textsubring{d}}{r}{a}} \dicPos{v}[1] \dicFlx{(‑aði)}[13] \dicPhraseIS{staldra við} \dicSynonym{stansa} \dicDirectTranslationCS{(po)zastavit se, (po)zdržet se (na chvilku ap.)} \dicExampleIS{staldra við stutta stund} \dicExampleCS{pozastavit se na krátkou chvíli}
\dicEntry[stalínismi] \dicTerm{stalín··ism|i} \dicIPA{{s}{\textsubring{d}}{a}{\textlengthmark}{l}{i}{n}{\textsci}{s}{m}{\textsci}} \dicPos{m}[1] \dicFlx{(‑a)}[3] \dicFieldCat{pol.} \dicDirectTranslationCS{stalinismus}
\dicEntry[stalla] \dicTerm{stalla} \dicIPA{{s}{\textsubring{d}}{a}{\textsubring{d}}{l}{a}} \dicPos{f}[1] \dicFlx{(stöllu, stöllur)}[20] \dicDirectTranslationCS{přítelkyně, kamarádka}
\dicEntry[stallur] \dicTerm{stall|ur} \dicIPA{{s}{\textsubring{d}}{a}{\textsubring{d}}{l}{\textscy}{\textsubring{r}}} \dicPos{m}[6] \dicFlx{(‑s, ‑ar)}[5] \textbf{1.} \dicSynonym{sylla} \dicDirectTranslationCS{římsa, výstupek, odskok}  \textbf{2.} \dicSynonym{undirstaða} \dicDirectTranslationCS{podstavec, sokl};  \dicPhraseIS{hefja e‑n á stall} \dicDirectTranslationCS{vychvalovat (koho)};  \dicPhraseIS{hrinda e‑m af stalli} \dicDirectTranslationCS{svrhnout (koho), zbavit (koho) moci\,/\addthin vlády} \dicExampleIS{velta konunginum af stalli} \dicExampleCS{svrhnout krále}  \textbf{3.} \dicSynonym{jata} \dicDirectTranslationCS{koryto, žlab}
\dicEntry[stam] \dicTerm{stam} \dicIPA{{s}{\textsubring{d}}{a}{\textlengthmark}{\textsubring{m}}} \dicPos{n}[2] \dicFlx{(‑s)}[2] \dicDirectTranslationCS{koktání, zajíkání (se)}
\dicEntry[stama] \dicTerm{stam|a} \dicsymFrequent\  \dicIPA{{s}{\textsubring{d}}{a}{\textlengthmark}{m}{a}} \dicPos{v}[1] \dicFlx{(‑aði)}[13] \dicDirectTranslationCS{koktat, zajíkat se} \dicExampleIS{stama orðunum út úr sér} \dicExampleCS{vykoktat slova ze sebe}
\dicEntry[stampur] \dicTerm{stamp|ur} \dicIPA{{s}{\textsubring{d}}{a}{\textsubring{m}}{\textsubring{b}}{\textscy}{\textsubring{r}}} \dicPos{m}[6] \dicFlx{(‑s, ‑ar)}[5] \dicDirectTranslationCS{káď, sud};  \dicPhraseIS{reka e‑n á stampinn} \dicLangCat{přen.} \dicDirectTranslationCS{porazit (koho) v~debatě}
\dicEntry[stamur] \dicTerm{stamur} \dicIPA{{s}{\textsubring{d}}{a}{\textlengthmark}{m}{\textscy}{\textsubring{r}}} \dicPos{adj}[1] \dicFlx{(f stöm)}[2] \dicDirectTranslationCS{lepkavý, zvlhlý (led ap.)} \dicExampleIS{Gólfið er stamt.} \dicExampleCS{Podlaha lepí.}
\dicEntry[stand] \dicTerm{stand} \dicIPA{{s}{\textsubring{d}}{a}{n}{\textsubring{d}}} \dicPos{n}[2] \dicFlx{(‑s, stönd)}[8] \dicSynonym{ástand} \dicDirectTranslationCS{stav, forma, kondice};  \dicPhraseIS{koma e‑u í stand} \dicDirectTranslationCS{opravit (co) (auto ap.)};  \dicPhraseIS{setja e‑ð í stand} \dicDirectTranslationCS{dát\,/\addthin uvést (co) do pořádku};  \dicPhraseIS{vera (ekki) í standi til að (gera e‑ð)} \dicDirectTranslationCS{(ne)být ve stavu (udělat (co))};  \dicPhraseIS{vera í (góðu) standi} \dicDirectTranslationCS{být v~(dobrém) stavu}
\dicEntry[standa] \dicTerm{standa} \dicsymFrequent\  \dicIPA{{s}{\textsubring{d}}{a}{n}{\textsubring{d}}{a}} \dicPos{v}[6] \dicFlx{(stend, stóð, stóðum, stæði, staðið)}[119] \textbf{1.} \dicDirectTranslationCS{stát (u~zdi ap.)} \dicIndirectTranslationCS{(zaujímat vztyčenou polohu)} \dicExampleIS{standa við vegginn} \dicExampleCS{stát u~zdi}  \textbf{2.} \dicDirectTranslationCS{stát, ležet, nacházet se} \dicExampleIS{Húsið stendur við götuna.} \dicExampleCS{Dům stojí na ulici.}  \textbf{3.} \dicDirectTranslationCS{stát, být napsáno (v~knize ap.)} \dicExampleIS{Það stendur í blaðinu.} \dicExampleCS{Je to napsáno v~novinách.}  \textbf{4.} \dicSynonym*{vera kyrr} \dicDirectTranslationCS{stát} \dicIndirectTranslationCS{(setrvávat na místě)}  \textbf{5.} \dicSynonym{endast} \dicDirectTranslationCS{probíhat, pokračovat, trvat} \dicExampleIS{Samkoman stóð lengi.} \dicExampleCS{Shromáždění trvalo dlouho.}  \textbf{6.} \dicSynonym*{vera ekki í gangi} \dicDirectTranslationCS{stát} \dicIndirectTranslationCS{(nebýt v~činnosti)} \dicExampleIS{Úrið stendur.} \dicExampleCS{Hodinky stojí.}  \textbf{7.} \dicDirectTranslationCS{platit, trvat (nabídka ap.)}  \textbf{8.} \dicPhraseIS{standa sig vel} \dicDirectTranslationCS{vést si dobře}  \textbf{9.} \dicPhraseIS{e‑m stendur} \dicFlx{impers} \dicDirectTranslationCS{(komu) stojí (penis), (kdo) má erekci};  \dicIdiom{standa}[að]{ \dicPhraseIS{standa að e‑u}} \dicDirectTranslationCS{stát za (čím), reprezentovat (co)} \dicIndirectTranslationCS{(být ve vztahu k~něčemu)} \dicExampleIS{standa að fyrirtækinu} \dicExampleCS{stát za společností}; { \dicPhraseIS{standa e‑m\,/\addthin e‑u að baki}} \dicDirectTranslationCS{nevyrovnat se (komu\,/\addthin čemu)}; { \dicPhraseIS{standa e‑n að e‑u}} \dicDirectTranslationCS{přistihnout (koho) při (čem)};  \dicIdiom{standa}[af]{ \dicPhraseIS{e‑m stendur ógn\,/\addthin ótti\,/\addthin beygur af e‑u}} \dicFlx{impers} \dicDirectTranslationCS{(kdo) se (čeho) bojí}; { \dicPhraseIS{standa af sér e‑ð}} \dicDirectTranslationCS{přestát (co), přečkat (co) (déšť ap.)} \dicExampleIS{standa af sér rigninguna} \dicExampleCS{přečkat déšť};  \dicIdiom{standa}[á]{ \dicPhraseIS{e‑m stendur á sama}} \dicFlx{impers} \dicDirectTranslationCS{(komu) je to jedno, (komu) na tom nezáleží}; { \dicPhraseIS{e‑r stendur á e‑m}} \dicDirectTranslationCS{(kdo koho) tlačí, (kdo) na (koho) naléhá}; { \dicPhraseIS{hvernig stendur á e‑u?}} \dicFlx{impers} \dicDirectTranslationCS{jak je (co) možné?}; { \dicPhraseIS{standa á rétti sínum}} \dicDirectTranslationCS{stát si na svých právech};  \dicPhraseIS{standa fast á e‑u} \dicLangCat{přen.} \dicDirectTranslationCS{stát si pevně za (čím)}; { \dicPhraseIS{það stendur á e‑u}} \dicFlx{impers} \dicDirectTranslationCS{(co) má zpoždění (odpověď ap.)}; { \dicPhraseIS{það stendur á e‑u\,/\addthin e‑m}} \dicFlx{impers} \dicDirectTranslationCS{vázne to na (čem\,/\addthin kom), stojí to na (čem\,/\addthin kom)}; { \dicPhraseIS{það stendur ekki á e‑m að (gera e‑ð)}} \dicFlx{impers} \dicDirectTranslationCS{(kdo) je připraven (udělat (co)), (kdo) si pospíší (udělat (co))};  \dicIdiom{standa}[á bak við]{ \dicPhraseIS{standa á bak við e‑ð}} \dicDirectTranslationCS{stát za (čím)} \dicIndirectTranslationCS{(být u~vzniku něčeho)};  \dicIdiom{standa}[á móti]{ \dicPhraseIS{standa á móti e‑u}} \dicDirectTranslationCS{postavit se (proti) (čemu)} \dicExampleIS{standa á móti þessum hugmyndum} \dicExampleCS{postavit se proti těmto nápadům};  \dicIdiom{standa}[fram undan]{ \dicPhraseIS{standa fram undan e‑u}} \dicDirectTranslationCS{vystávat (co), vykukovat (co), čouhat (co) (šátek z~kapsy ap.)};  \dicIdiom{standa}[frammi fyrir]{ \dicPhraseIS{standa frammi fyrir e‑u}} \dicDirectTranslationCS{stát před (čím) (těžkým rozhodnutím ap.)};  \dicIdiom{standa}[fyrir]{ \dicPhraseIS{standa fyrir e‑u}} {\textbf{a.}} \dicSynonym{hindra} \dicDirectTranslationCS{stát v~cestě (čemu), bránit (čemu) (v~chůzi ap.)};  {\textbf{b.}} \dicDirectTranslationCS{uspořádat (co), organizovat (co)} \dicExampleIS{standa fyrir tónleikunum} \dicExampleCS{uspořádat koncert};  \dicIdiom{standa}[gegn]{ \dicPhraseIS{standa gegn e‑u}} \dicDirectTranslationCS{být proti (čemu), postavit se proti (čemu)} \dicExampleIS{standa gegn þessum áformum} \dicExampleCS{postavit se proti těmto plánům};  \dicIdiom{standa}[hjá]{ \dicPhraseIS{standa hjá}} \dicDirectTranslationCS{nečinně stát, přihlížet};  \dicIdiom{standa}[í]{ \dicPhraseIS{e‑að stendur í e‑m}} \dicDirectTranslationCS{(co komu) uvízne v~krku}; { \dicPhraseIS{e‑r stendur í e‑m}} \dicDirectTranslationCS{(kdo) se (komu) staví (čelem)}; { \dicPhraseIS{standa í e‑u}} \dicDirectTranslationCS{být zaměstnaný (čím)};  \dicIdiom{standa}[með]{ \dicPhraseIS{e‑að stendur með blóma}} \dicDirectTranslationCS{(co) vzkvétá, (co) prosperuje (finančně ap.)}; { \dicPhraseIS{e‑r stendur með e‑m}} \dicDirectTranslationCS{(kdo koho) podporuje, (kdo) za (kým) stojí} \dicExampleIS{Ég stend með honum.} \dicExampleCS{Stojím za ním.};  \dicIdiom{standa}[saman]{ \dicPhraseIS{standa saman}} \dicDirectTranslationCS{stát při sobě, držet pohromadě} \dicExampleIS{Þau standa saman á hverju sem gengur.} \dicExampleCS{Stojí při sobě ve všem, co přichází.};  \dicIdiom{standa}[til]{ \dicPhraseIS{e‑að stendur til}} \dicDirectTranslationCS{(co) se chystá, (co) se plánuje};  \dicIdiom{standa}[undir]{ \dicPhraseIS{standa undir e‑u}} \dicDirectTranslationCS{být schopen nést (co) (náklady ap.)} \dicExampleIS{standa undir kostnaðinum} \dicExampleCS{být schopen nést náklady}; { \dicPhraseIS{standa (ekki) undir nafni}} \dicDirectTranslationCS{(ne)dostát svému jménu};  \dicIdiom{standa}[upp]{ \dicPhraseIS{standa upp fyrir e‑m}} \dicDirectTranslationCS{uvolnit (komu) místo};  \dicIdiom{standa}[upp úr]{ \dicPhraseIS{e‑að stendur upp úr e‑u}} \dicDirectTranslationCS{(co) vyčnívá z~(čeho) (tráva ze sněhu ap.)}; { \dicPhraseIS{e‑að stendur upp úr}} \dicDirectTranslationCS{(co) vyčnívá, (co) se vymyká (průměru ap.)};  \dicIdiom{standa}[upp]{ \dicPhraseIS{standa upp}} \dicDirectTranslationCS{vstát, postavit se};  \dicIdiom{standa}[uppi]{ \dicPhraseIS{standa uppi eignalaus}} \dicDirectTranslationCS{přijít o~vše, ztratit vše};  \dicIdiom{standa}[út]{ \dicPhraseIS{standa út}} \dicDirectTranslationCS{vyčnívat, vykukovat};  \dicIdiom{standa}[við]{ \dicPhraseIS{standa við}} \dicSynonym{stansa} \dicDirectTranslationCS{pobýt, pobývat, zdržet se, zdržovat se, setrvat, setrvávat} \dicExampleIS{standa við nokkra stund} \dicExampleCS{zdržet se pár chvil}; { \dicPhraseIS{standa við e‑ð}} \dicDirectTranslationCS{dodržet (co), splnit (co), dostát (čemu)} \dicExampleIS{standa við loforð} \dicExampleCS{dodržet slib};  \dicIdiom{standa}[yfir]{ \dicPhraseIS{e‑að stendur yfir}} \dicDirectTranslationCS{(co) probíhá, (co) se koná} \dicExampleIS{Fundurinn stendur enn yfir.} \dicExampleCS{Zasedání ještě probíhá.}; { \dicPhraseIS{standa yfir e‑m}} \dicDirectTranslationCS{stát nad (kým), dívat se (komu) přes rameno};  \dicIdiom{standast}{ \dicPhraseIS{standast}} \dicFlx{refl acc} \dicSynonym{þola} \dicDirectTranslationCS{vydržet, odolat, odolávat}; { \dicPhraseIS{standast ekki mátið}} \dicFlx{refl} \dicDirectTranslationCS{nemoct si pomoci, nemoct se ubránit (a~něco udělat ap.)};  \dicPhraseIS{standast freistinguna} \dicFlx{refl} \dicDirectTranslationCS{odolat pokušení}; { \dicPhraseIS{standast prófið}} \dicFlx{refl} \dicDirectTranslationCS{složit zkoušku}; { \dicPhraseIS{standast e‑m snúning}} \dicFlx{refl} \dicDirectTranslationCS{vyrovnat se (komu)};  \dicIdiom{standast}[á]{ \dicPhraseIS{standast á}} \dicFlx{refl} \dicDirectTranslationCS{odpovídat, korespondovat, shodovat se};  \dicIdiom{standa}{ \dicPhraseIS{eins og stendur}} \dicFlx{adv} \dicDirectTranslationCS{v~tuto chvíli}; { \dicPhraseIS{sem stendur}} \dicFlx{adv} \dicDirectTranslationCS{v~(tomto) okamžiku, nyní}
\dicEntry[standandi] \dicTerm{stand··andi} \dicIPA{{s}{\textsubring{d}}{a}{n}{\textsubring{d}}{a}{n}{\textsubring{d}}{\textsci}} \dicPos{adj}[13] \dicFlx{indecl}[1] \dicDirectTranslationCS{stojící}
\dicEntry[standard] \dicTerm{standard\smash{\textsuperscript{1}}} \dicIPA{{s}{\textsubring{d}}{a}{n}{\textsubring{d}}{a}{r}{\textsubring{d}}} \dicPos{m}[4] \dicFlx{(‑s, ‑ar)}[23] \dicDirectTranslationCS{standard, úroveň} \dicExampleIS{hækka standardinn á e‑u} \dicExampleCS{zvýšit standard (čeho)}
\dicEntry[standard] \dicTerm{standard-\smash{\textsuperscript{2}}} \dicIPA{{s}{\textsubring{d}}{a}{n}{\textsubring{d}}{a}{r}{\textsubring{d}}} \dicPos{predp} \dicDirectTranslationCS{standardní, normální}
\dicEntry[standari] \dicTerm{stand··ar|i} \dicIPA{{s}{\textsubring{d}}{a}{n}{\textsubring{d}}{a}{r}{\textsci}} \dicPos{m}[1] \dicFlx{(‑a, ‑ar)}[10] \dicDirectTranslationCS{stojan, stojánek} \dicIndirectTranslationCS{(u~jízdního kola ap.)}
\dicEntry[standlampi] \dicTerm{stand··lamp|i} \dicIPA{{s}{\textsubring{d}}{a}{n}{\textsubring{d}}{l}{a}{\textsubring{m}}{\textsubring{b}}{\textsci}} \dicPos{m}[1] \dicFlx{(‑a, ‑ar)}[8] \dicDirectTranslationCS{stojací lampa}
\dicEntry[standpína] \dicTerm{stand··pín|a} \dicIPA{{s}{\textsubring{d}}{a}{n}{\textsubring{d}}{p\smash{\textsuperscript{h}}}{i}{n}{a}} \dicPos{f}[1] \dicFlx{(‑u)}[5] \dicLangCat{hovor.} \dicDirectTranslationCS{trčák, stání} \dicIndirectTranslationCS{(ztopořený penis)}
\dicEntry[standset] \dicTerm{stand··set} \dicIPA{{s}{\textsubring{d}}{a}{n}{\textsubring{d}}{s}{\textepsilon}{\textsubring{d}}} \dicPos{v} \dicFlx{ind praes sg 1 pers} \dicLink{standsetja}
\dicEntry[standsetja] \dicTerm{stand··|setja} \dicIPA{{s}{\textsubring{d}}{a}{n}{\textsubring{d}}{s}{\textepsilon}{\textsubring{d}}{j}{a}} \dicPos{v}[4] \dicFlx{(‑set, ‑setti, ‑settum, ‑setti, ‑sett)}[26] \dicFlx{acc} \dicSynonym{lagfæra} \dicDirectTranslationCS{opravit, opravovat, obnovit, obnovovat, (z)renovovat} \dicExampleIS{standsetja íbúðina} \dicExampleCS{renovovat byt}
\dicEntry[standsetning] \dicTerm{stand··set·ning} \dicIPA{{s}{\textsubring{d}}{a}{n}{\textsubring{d}}{s}{\textepsilon}{\textsubring{d}}{n}{i}{\ng}{\r{g}}} \dicPos{f}[4] \dicFlx{(‑ar)}[7] \dicDirectTranslationCS{oprava, obnovení, renovace}
\dicEntry[standsett] \dicTerm{stand··sett} \dicIPA{{s}{\textsubring{d}}{a}{n}{\textsubring{d}}{s}{\textepsilon}{h}{\textsubring{d}}} \dicPos{v} \dicFlx{supin} \dicLink{standsetja}
\dicEntry[standsetti] \dicTerm{stand··setti} \dicIPA{{s}{\textsubring{d}}{a}{n}{\textsubring{d}}{s}{\textepsilon}{h}{\textsubring{d}}{\textsci}} \dicPos{v} \dicFlx{ind\,/\addthin con pf sg 1 pers} \dicLink{standsetja}
\dicEntry[standsettum] \dicTerm{stand··settum} \dicIPA{{s}{\textsubring{d}}{a}{n}{\textsubring{d}}{s}{\textepsilon}{h}{\textsubring{d}}{\textscy}{\textsubring{m}}} \dicPos{v} \dicFlx{ind pf pl 1 pers} \dicLink{standsetja}
\dicEntry[standur] \dicTerm{stand|ur} \dicIPA{{s}{\textsubring{d}}{a}{n}{\textsubring{d}}{\textscy}{\textsubring{r}}} \dicPos{m}[6] \dicFlx{(‑s, ‑ar)}[5] \dicSynonym{statíf} \dicDirectTranslationCS{stojan, držák};  \dicPhraseIS{hrökkva upp af standinum} \dicLangCat{přen.} \dicSynonym{geggjast} \dicDirectTranslationCS{zbláznit se, pomátnout se}
\dicEntry[stanga] \dicTerm{stang|a} \dicIPA{{s}{\textsubring{d}}{au}{\ng}{\r{g}}{a}} \dicPos{v}[1] \dicFlx{(‑aði)}[13] \dicFlx{acc} \textbf{1.} \dicSynonym*{reka hausinn í} \dicDirectTranslationCS{trknout, (po)trkat} \dicExampleIS{Hrúturinn stangaði mig.} \dicExampleCS{Beran mě potrkal.}  \textbf{2.} \dicSynonym*{sauma í gegn} \dicDirectTranslationCS{prošívat, šít přes více vrstev}  \textbf{3.} \dicPhraseIS{stanga úr tönnunum} \dicDirectTranslationCS{šťourat se v~zubech};  \dicIdiom{stangast}{ \dicPhraseIS{stangast}} \dicFlx{refl} \dicDirectTranslationCS{(po)trkat se};  \dicIdiom{stangast}[á]{ \dicPhraseIS{stangast á við e‑ð}} \dicFlx{refl} \dicDirectTranslationCS{střetat se s~(čím), být v~konfliktu s~(čím)}
\dicEntry[stangar] \dicTerm{stangar} \dicIPA{{s}{\textsubring{d}}{au}{\ng}{\r{g}}{a}{\textsubring{r}}} \dicPos{f} \dicFlx{sg gen} \dicLink{stöng}
\dicEntry[stangarstökk] \dicTerm{stangar··stökk} \dicIPA{{s}{\textsubring{d}}{au}{\ng}{\r{g}}{a}{\textsubring{r}}{s}{\textsubring{d}}{\oe}{h}{\r{g}}} \dicPos{n}[2] \dicFlx{(‑s, ‑)}[5] \dicFieldCat{sport.} \dicDirectTranslationCS{skok o~tyči}
\dicEntry[stangaveiði] \dicTerm{stanga··veið|i} \dicIPA{{s}{\textsubring{d}}{au}{\ng}{\r{g}}{a}{v}{ei}{ð}{\textsci}} \dicPos{f}[2] \dicFlx{(‑i\,/\addthin ‑ar, ‑ar)}[3] \dicDirectTranslationCS{lovení na udici}
\dicEntry[stangaveiðifélag] \dicTerm{stanga·veiði··fé·|lag} \dicIPA{{s}\-{\textsubring{d}}\-{au}\-{\ng}\-{\r{g}}\-{a}\-{v}\-{ei}\-{ð}\-{\textsci}\-{f}\-{j}\-{\textepsilon}\-{l}\-{a}\-{x}\-} \dicPos{n}[2] \dicFlx{(‑lags, ‑lög)}[8] \dicDirectTranslationCS{rybářský klub, rybářská organizace}
\dicEntry[stangl] \dicTerm{stangl} \dicIPA{{s}{\textsubring{d}}{au}{\ng}{\r{g}}{\textsubring{l}}} \dicPos{n}[2] \dicFlx{(‑s, stöngl)}[8] \dicPhraseIS{á stangli} \dicFlx{adv} \dicDirectTranslationCS{tu a~tam, sem tam} \dicExampleIS{Þarna eru nokkur hús á stangli.} \dicExampleCS{Tu a~tam jsou domy.}
\dicEntry[stans] \dicTerm{stans} \dicIPA{{s}{\textsubring{d}}{a}{n}{s}} \dicPos{n}[2] \dicFlx{(‑)}[24] \dicDirectTranslationCS{zastávka, zastavení} \dicExampleIS{eftir hálftíma stans í sæluhúsinu} \dicExampleCS{po půlhodinové zastávce v~chatě}
\dicEntry[stansa] \dicTerm{stans|a} \dicsymFrequent\  \dicIPA{{s}{\textsubring{d}}{a}{n}{s}{a}} \dicPos{v}[1] \dicFlx{(‑aði)}[43] \textbf{1.} \dicSynonym*{nema staðar} \dicDirectTranslationCS{zastavit (se)} \dicIndirectTranslationCS{(ustat v~chůzi\,/\addthin běhu ap.)} \dicExampleIS{Hjólið stansaði við ljósin.} \dicExampleCS{Kolo se zastavilo na semaforech.}  \textbf{2.} \dicFlx{acc} \dicSynonym{stöðva} \dicDirectTranslationCS{zastavit} \dicIndirectTranslationCS{(způsobit, aby někdo\,/\addthin něco ustalo v~chůzi\,/\addthin běhu ap.)} \dicExampleIS{stansa e‑n} \dicExampleCS{zastavit (koho)}  \textbf{3.} \dicDirectTranslationCS{(za)stavit se (u~přátel ap.)} \dicExampleIS{Ég stansaði í klukkutíma hjá ömmu minni.} \dicExampleCS{Zastavil jsem se na hodinu u~své babičky.}
\dicEntry[stanslaus] \dicTerm{stans··laus} \dicIPA{{s}{\textsubring{d}}{a}{n}{s}{l}{\oe i}{s}} \dicPos{adj}[5]\dicFlx{}[-1] \dicSynonym{linnulaus} \dicDirectTranslationCS{nepřetržitý, soustavný} \dicExampleIS{stanslaus vinna} \dicExampleCS{nepřetržitá práce}
\dicEntry[stapi] \dicTerm{stap|i} \dicIPA{{s}{\textsubring{d}}{a}{\textlengthmark}{\textsubring{b}}{\textsci}} \dicPos{m}[1] \dicFlx{(‑a, ‑ar)}[8] \dicDirectTranslationCS{ostroh, mys}
\dicEntry[stapp] \dicTerm{stapp} \dicIPA{{s}{\textsubring{d}}{a}{h}{\textsubring{b}}} \dicPos{n}[2] \dicFlx{(‑s, stöpp)}[8] \textbf{1.} \dicSynonym*{það að stappa} \dicDirectTranslationCS{dupnutí, dupot}  \textbf{2.} \dicSynonym{deila\smash{\textsuperscript{1}}} \dicDirectTranslationCS{spor, konflikt};  \dicPhraseIS{lenda í stappi út af e‑u} \dicDirectTranslationCS{dostat se do sporu kvůli (čemu)}
\dicEntry[stappa] \dicTerm{stappa\smash{\textsuperscript{1}}} \dicIPA{{s}{\textsubring{d}}{a}{h}{\textsubring{b}}{a}} \dicPos{f}[1] \dicFlx{(stöppu, stöppur)}[8] \dicDirectTranslationCS{kaše, míchanina}
\dicEntry[stappa] \dicTerm{stapp|a\smash{\textsuperscript{2}}} \dicIPA{{s}{\textsubring{d}}{a}{h}{\textsubring{b}}{a}} \dicPos{v}[1] \dicFlx{(‑aði)}[13] \dicFlx{acc\,/\addthin dat} \textbf{1.} \dicFlx{acc} \dicSynonym{merja} \dicDirectTranslationCS{rozmačkat (na kaši)} \dicExampleIS{stappa kartöflur} \dicExampleCS{rozmačkat brambory}  \textbf{2.} \dicFlx{dat} \dicDirectTranslationCS{dupnout, dupat} \dicExampleIS{stappa fótunum í gólfið} \dicExampleCS{dupat nohama po podlaze};  \dicPhraseIS{stappa fótunum framan í e‑n} \dicLangCat{přen.} \dicDirectTranslationCS{dupnout si na (koho)};  \dicIdiom{stappa}[í]{ \dicPhraseIS{stappa stálinu í e‑n}} \dicLangCat{přen.} \dicDirectTranslationCS{dodávat (komu) kuráž, dodávat (komu) důvěru v~sebe sama};  \dicIdiom{stappa}[nærri]{ \dicPhraseIS{e‑að stappar nærri e‑u}} \dicDirectTranslationCS{(co) hraničí s~(čím) (šílenstvím ap.)} \dicExampleIS{Það stappar nærri brjálæði.} \dicExampleCS{Hraničí to se šílenstvím.}
\dicEntry[stappaður] \dicTerm{stapp··aður} \dicIPA{{s}{\textsubring{d}}{a}{h}{\textsubring{b}}{a}{ð}{\textscy}{\textsubring{r}}} \dicPos{adj}[3] \dicFlx{(f stöppuð)}[2] \dicPhraseIS{stappaðar kartöflur} \dicFieldCat{kulin.} \dicDirectTranslationCS{šťouchané brambory, bramborová kaše};  \dicPhraseIS{það er stappað e‑s staðar} \dicFlx{impers} \dicDirectTranslationCS{(kde) je namačkáno\,/\addthin našlapáno}
\dicEntry[stara] \dicTerm{star|a} \dicsymFrequent\  \dicIPA{{s}{\textsubring{d}}{a}{\textlengthmark}{r}{a}} \dicPos{v}[2] \dicFlx{(‑ði, ‑að)}[125] \dicDirectTranslationCS{koukat, zírat, dívat se upřeně} \dicExampleIS{Starðu ekki á mig!} \dicExampleCS{Nezírej na mě!};  \dicPhraseIS{stara á e‑ð\,/\addthin e‑n} \dicDirectTranslationCS{koukat na (co\,/\addthin koho)}
\dicEntry[starar] \dicTerm{starar} \dicIPA{{s}{\textsubring{d}}{a}{\textlengthmark}{r}{a}{\textsubring{r}}} \dicPos{f} \dicFlx{sg gen} \dicLink{stör}
\dicEntry[starblinda] \dicTerm{star··blind|a} \dicIPA{{s}{\textsubring{d}}{a}{r}{\textsubring{b}}{l}{\textsci}{n}{\textsubring{d}}{a}} \dicPos{f}[1] \dicFlx{(‑u)}[5] \dicFieldCat{med.} \dicDirectTranslationCS{šedý zákal}
\dicEntry[starf] \dicTerm{starf} \dicsymFrequent\  \dicIPA{{s}{\textsubring{d}}{a}{r}{f}} \dicPos{n}[2] \dicFlx{(‑s, störf)}[8] \textbf{1.} \dicSynonym{staða} \dicDirectTranslationCS{práce, zaměstnání, povolání, profese};  \dicPhraseIS{e‑að tekur til starfa} \dicDirectTranslationCS{(co) zahajuje svou činnost};  \dicPhraseIS{gegna starfi e‑rs} \dicDirectTranslationCS{pracovat jako (kdo)};  \dicPhraseIS{láta af störfum} \dicDirectTranslationCS{odstoupit, ukončit práci};  \dicPhraseIS{ráða e‑n til starfa} \dicDirectTranslationCS{přijmout (koho) do zaměstnání, zaměstnat (koho)};  \dicPhraseIS{vera í fullu starfi} \dicDirectTranslationCS{pracovat na plný úvazek};  \dicPhraseIS{vera í hálfu starfi} \dicDirectTranslationCS{pracovat na poloviční úvazek}  \textbf{2.} \dicSynonym{verkefni} \dicDirectTranslationCS{práce, úkol, zadání}
\dicEntry[starfa] \dicTerm{starf|a} \dicsymFrequent\  \dicIPA{{s}{\textsubring{d}}{a}{r}{v}{a}} \dicPos{v}[1] \dicFlx{(‑aði)}[13] \textbf{1.} \dicSynonym{vinna\smash{\textsuperscript{2}}} \dicDirectTranslationCS{pracovat, dělat} \dicExampleIS{starfa hjá stóru fyrirtæki} \dicExampleCS{pracovat u~velké firmy};  \dicPhraseIS{starfa að e‑u} \dicDirectTranslationCS{pracovat na (čem)}  \textbf{2.} \dicDirectTranslationCS{být aktivní, fungovat, existovat}
\dicEntry[starfhæfur] \textls[10]{\dicTerm{starf··hæfur}\dicTerm{, starfshæfur} \dicIPA{{s}\-{\textsubring{d}}\-{a}\-{r}\-{f}\-{h}\-{a}\-{i}\-{v}\-{\textscy}\-{\textsubring{r}}\-} \dicPos{adj}[1]\dicFlx{}[-1]} \dicDirectTranslationCS{práceschopný, funkční} \dicExampleIS{starfhæfur maður} \dicExampleCS{práceschopný člověk}
\dicEntry[starfrækja] \dicTerm{starf··ræk|ja} \dicIPA{{s}{\textsubring{d}}{a}{r}{v}{r}{a}{i}{\r{\textObardotlessj}}{a}} \dicPos{v}[2] \dicFlx{(‑ti, ‑t)}[25] \dicFlx{acc} \dicSynonym{reka\smash{\textsuperscript{2}}} \dicDirectTranslationCS{provozovat, vést, spravovat, řídit} \dicExampleIS{starfrækja hótel} \dicExampleCS{provozovat hotel}
\dicEntry[starfræksla] \dicTerm{starf··ræksl|a} \dicIPA{{s}{\textsubring{d}}{a}{r}{v}{r}{a}{i}{x}{s}{\textsubring{d}}{l}{a}} \dicPos{f}[1] \dicFlx{(‑u)}[5] \dicDirectTranslationCS{provozování, provoz}
\dicEntry[starfsaldur] \dicTerm{starfs··ald|ur} \dicIPA{{s}{\textsubring{d}}{a}{r}{v}{s}{a}{l}{\textsubring{d}}{\textscy}{\textsubring{r}}} \dicPos{m}[5] \dicFlx{(‑urs, ‑rar)}[3] \textbf{1.} \dicDirectTranslationCS{služební věk, odpracované roky}  \textbf{2.} \dicDirectTranslationCS{aktivní\,/\addthin produktivní věk}
\dicEntry[starfsamur] \dicTerm{starf··|samur} \dicIPA{{s}{\textsubring{d}}{a}{\textsubring{r}}{s}{a}{m}{\textscy}{\textsubring{r}}} \dicPos{adj}[1] \dicFlx{(f ‑söm)}[2] \dicDirectTranslationCS{pracovitý, pilný}
\dicEntry[starfsár] \dicTerm{starfs··ár} \dicIPA{{s}{\textsubring{d}}{a}{r}{v}{s}{au}{\textsubring{r}}} \dicPos{n}[2] \dicFlx{(‑s, ‑)}[5] \dicDirectTranslationCS{pracovní rok}
\dicEntry[starfsáætlun] \dicTerm{starfs··á·ætl|un} \dicIPA{{s}{\textsubring{d}}{a}{r}{v}{s}{au}{a}{i}{h}{\textsubring{d}}{l}{\textscy}{\textsubring{n}}} \dicPos{f}[7] \dicFlx{(‑unar, ‑anir)}[8] \dicDirectTranslationCS{pracovní plán}
\dicEntry[starfsbróðir] \dicTerm{starfs··|bróðir} \dicIPA{{s}{\textsubring{d}}{a}{r}{v}{s}{\textsubring{b}}{r}{ou}{ð}{\textsci}{\textsubring{r}}} \dicPos{m}[13] \dicFlx{(‑bróður, ‑bræður)}[4] \dicDirectTranslationCS{kolega, spolupracovník}
\dicEntry[starfsemi] \dicTerm{starf··sem|i} \dicsymFrequent\  \dicIPA{{s}{\textsubring{d}}{a}{r}{v}{s}{\textepsilon}{m}{\textsci}} \dicPos{f}[2] \dicFlx{(‑i)}[2] \textbf{1.} \dicSynonym*{það að starfa} \dicDirectTranslationCS{činnost, aktivita, fungování} \dicExampleIS{andleg starfsemi} \dicExampleCS{duševní činnost}  \textbf{2.} \dicSynonym{iðni} \dicDirectTranslationCS{pracovitost, píle}
\dicEntry[starfsferill] \dicTerm{starfs··fer|ill} \dicIPA{{s}{\textsubring{d}}{a}{r}{v}{s}{f}{\textepsilon}{r}{\textsci}{\textsubring{d}}{\textsubring{l}}} \dicPos{m}[6] \dicFlx{(‑ils)}[36] \dicDirectTranslationCS{pracovní dráha, kariéra}
\dicEntry[starfsfólk] \dicTerm{starfs··fólk} \dicsymFrequent\  \dicIPA{{s}{\textsubring{d}}{a}{r}{v}{s}{f}{ou}{\textsubring{l}}{\r{g}}} \dicPos{n}[2] \dicFlx{(‑s)}[2] \dicDirectTranslationCS{zaměstnanci, personál, osazenstvo} \dicExampleIS{kröfur starfsfólksins} \dicExampleCS{nároky zaměstnanců}
\dicEntry[starfsfræðsla] \dicTerm{starfs··fræðsl|a} \dicIPA{{s}{\textsubring{d}}{a}{r}{v}{s}{f}{r}{a}{i}{ð}{s}{\textsubring{d}}{l}{a}} \dicPos{f}[1] \dicFlx{(‑u)}[5] \dicDirectTranslationCS{poradenství pro volbu povolání}
\dicEntry[starfsgrein] \dicTerm{starfs··grein} \dicIPA{{s}{\textsubring{d}}{a}{r}{v}{s}{\r{g}}{r}{ei}{\textsubring{n}}} \dicPos{f}[4] \dicFlx{(‑ar, ‑ar)}[1] \dicDirectTranslationCS{pracovní obor, profese}
\dicEntry[starfsheiti] \dicTerm{starfs··heiti} \dicIPA{{s}{\textsubring{d}}{a}{r}{v}{s}{h}{ei}{\textsubring{d}}{\textsci}} \dicPos{n}[2] \dicFlx{(‑s, ‑)}[14] \textbf{1.} \dicDirectTranslationCS{název povolání}  \textbf{2.} \dicDirectTranslationCS{profesní titul}
\dicEntry[starfshópur] \dicTerm{starfs··hóp|ur} \dicIPA{{s}{\textsubring{d}}{a}{r}{v}{s}{h}{ou}{\textsubring{b}}{\textscy}{\textsubring{r}}} \dicPos{m}[6] \dicFlx{(‑s, ‑ar)}[22] \dicDirectTranslationCS{pracovní kolektiv\,/\addthin skupina}
\dicEntry[starfshæfur] \dicTerm{starfs··hæfur} \dicIPA{{s}{\textsubring{d}}{a}{r}{v}{s}{h}{a}{i}{v}{\textscy}{\textsubring{r}}} \dicPos{adj}[1]\dicFlx{}[-1] \dicLink{starfhæfur}
\dicEntry[starfskraftur] \dicTerm{starfs··kraft|ur} \dicIPA{{s}{\textsubring{d}}{a}{r}{v}{s}{k\smash{\textsuperscript{h}}}{r}{a}{f}{\textsubring{d}}{\textscy}{\textsubring{r}}} \dicPos{m}[6] \dicFlx{(‑s\,/\addthin ‑ar, ‑ar)}[64] \dicSynonym{starfsmaður} \dicDirectTranslationCS{pracovní síla, pracovník, pracovnice}
\dicEntry[starfsleyfi] \dicTerm{starfs··leyfi} \dicIPA{{s}{\textsubring{d}}{a}{r}{v}{s}{l}{ei}{v}{\textsci}} \dicPos{n}[2] \dicFlx{(‑s, ‑)}[14] \dicDirectTranslationCS{(pracovní) licence, koncese}
\dicEntry[starfslok] \dicTerm{starfs··lok} \dicIPA{{s}{\textsubring{d}}{a}{r}{v}{s}{l}{\textopeno}{\r{g}}} \dicPos{n}[2] \dicFlx{pl}[1] \dicDirectTranslationCS{odchod (do důchodu ap.), odstoupení}
\dicEntry[starfslýsing] \dicTerm{starfs··lýs·ing} \dicIPA{{s}{\textsubring{d}}{a}{r}{v}{s}{l}{i}{s}{i}{\ng}{\r{g}}} \dicPos{f}[4] \dicFlx{(‑ar, ‑ar)}[5] \dicDirectTranslationCS{popis práce}
\dicEntry[starfsmaður] \textls[10]{\dicTerm{starfs··|maður} \dicsymFrequent\  \dicIPA{{s}{\textsubring{d}}{a}{r}{v}{s}{m}{a}{ð}{\textscy}{\textsubring{r}}} \dicPos{m}[13] \dicFlx{(‑manns, ‑menn)}[2]} \dicDirectTranslationCS{pracovník, pracovnice, zaměstnanec, zaměstnankyně} \dicExampleIS{starfsmenn ríkisins} \dicExampleCS{státní zaměstnanci}
\dicEntry[starfsmannastjóri] \dicTerm{starfs·manna··stjór|i} \dicIPA{{s}\-{\textsubring{d}}\-{a}\-{r}\-{v}\-{s}\-{m}\-{a}\-{n}\-{a}\-{s}\-{\textsubring{d}}\-{j}\-{ou}\-{r}\-{\textsci}\-} \dicPos{m}[1] \dicFlx{(‑a, ‑ar)}[1] \dicDirectTranslationCS{personální ředitel(ka)}
\dicEntry[starfsmat] \dicTerm{starfs··mat} \dicIPA{{s}{\textsubring{d}}{a}{r}{v}{s}{m}{a}{\textsubring{d}}} \dicPos{n}[2] \dicFlx{(‑s)}[2] \dicDirectTranslationCS{(o)hodnocení práce}
\dicEntry[starfsmenntun] \dicTerm{starfs··mennt|un} \dicIPA{{s}{\textsubring{d}}{a}{r}{v}{s}{m}{\textepsilon}{\textsubring{n}}{\textsubring{d}}{\textscy}{\textsubring{n}}} \dicPos{f}[7] \dicFlx{(‑unar)}[9] \dicDirectTranslationCS{profesní vzdělávání\,/\addthin školení}
\dicEntry[starfsnám] \dicTerm{starfs··nám} \dicIPA{{s}{\textsubring{d}}{a}{r}{v}{s}{\textsubring{d}}{n}{au}{\textsubring{m}}} \dicPos{n}[2] \dicFlx{(‑s)}[2] \dicDirectTranslationCS{stáž, odborné vzdělávání}
\dicEntry[starfsorka] \dicTerm{starfs··ork|a} \dicIPA{{s}{\textsubring{d}}{a}{r}{v}{s}{\textopeno}{\textsubring{r}}{\r{g}}{a}} \dicPos{f}[1] \dicFlx{(‑u)}[5] \dicDirectTranslationCS{pracovní schopnost\,/\addthin kapacita}
\dicEntry[starfsregla] \dicTerm{starfs··regl|a} \dicIPA{{s}{\textsubring{d}}{a}{r}{v}{s}{r}{\textepsilon}{\r{g}}{l}{a}} \dicPos{f}[1] \dicFlx{(‑u, ‑ur)}[19] \dicDirectTranslationCS{pravidlo práce}
\dicEntry[starfsreynsla] \dicTerm{starfs··reynsl|a} \dicIPA{{s}{\textsubring{d}}{a}{r}{v}{s}{r}{ei}{n}{s}{\textsubring{d}}{l}{a}} \dicPos{f}[1] \dicFlx{(‑u)}[5] \dicDirectTranslationCS{pracovní zkušenost, praxe}
\dicEntry[starfsskilyrði] \textls[15]{\dicTerm{starfs··skil·yrði} \dicIPA{{s}{\textsubring{d}}{a}{r}{v}{s}{\r{\textObardotlessj}}{\textsci}{l}{\textsci}{r}{ð}{\textsci}} \dicPos{n}[2] \dicFlx{pl}[19] \dicDirectTranslationCS{pracovní podmínky}}
\dicEntry[starfsstjórn] \dicTerm{starfs··stjórn} \dicIPA{{s}{\textsubring{d}}{a}{r}{v}{s}{\textsubring{d}}{j}{ou}{r}{\textsubring{d}}{\textsubring{n}}} \dicPos{f}[7] \dicFlx{(‑ar, ‑ir)}[1] \dicDirectTranslationCS{prozatímní\,/\addthin dočasná vláda}
\dicEntry[starfssvið] \dicTerm{starfs··svið} \dicIPA{{s}{\textsubring{d}}{a}{r}{v}{s}{v}{\textsci}{\texttheta}} \dicPos{n}[2] \dicFlx{(‑s, ‑)}[5] \dicDirectTranslationCS{pracovní oblast}
\dicEntry[starfssystir] \dicTerm{starfs··syst|ir} \dicIPA{{s}{\textsubring{d}}{a}{r}{v}{s}{\textsci}{s}{\textsubring{d}}{\textsci}{\textsubring{r}}} \dicPos{f}[11] \dicFlx{(‑ur, ‑ur)}[3] \dicDirectTranslationCS{kolegyně, spolupracovnice}
\dicEntry[starfsþjálfun] \dicTerm{starfs··þjálf|un} \dicIPA{{s}{\textsubring{d}}{a}{r}{v}{s}{\texttheta}{j}{au}{l}{v}{\textscy}{\textsubring{n}}} \dicPos{f}[7] \dicFlx{(‑unar)}[9] \dicDirectTranslationCS{stáž, odborné vzdělávání}
\dicEntry[stari] \dicTerm{star|i} \dicIPA{{s}{\textsubring{d}}{a}{\textlengthmark}{r}{\textsci}} \dicPos{m}[1] \dicFlx{(‑a, ‑ar)}[8] \dicLink{starri}
\dicEntry[starir] \dicTerm{starir} \dicIPA{{s}{\textsubring{d}}{a}{\textlengthmark}{r}{\textsci}{\textsubring{r}}} \dicPos{f} \dicFlx{pl nom} \dicLink{stör}
\dicEntry[starri] \dicTerm{starr|i}\dicTerm{, stari} \dicIPA{{s}\-{\textsubring{d}}\-{a}\-{r}\-{\textlengthmark}\-{\textsci}\-} \dicPos{m}[1] \dicFlx{(‑a, ‑ar)}[8] \dicFieldCat{zool.} \dicDirectTranslationCS{špaček, špaček obecný} \textit{(l.~{\textLA{Sturnus vulgaris}})}  \dicsymPhoto\ 
\dicFigure{20747.jpg}{Starri}{Starri - Karney, Lee, Biolib, PD}
\dicEntry[start] \dicTerm{start} \dicIPA{{s}{\textsubring{d}}{a}{\textsubring{r}}{\textsubring{d}}} \dicPos{n}[2] \dicFlx{(‑s)}[2] \textbf{1.} \dicSynonym{byrjun} \dicDirectTranslationCS{start, začátek, zahájení}  \textbf{2.} \dicSynonym*{gangsetning} \dicDirectTranslationCS{nastartování, spuštění}
\dicEntry[starta] \dicTerm{start|a} \dicIPA{{s}{\textsubring{d}}{a}{\textsubring{r}}{\textsubring{d}}{a}} \dicPos{v}[1] \dicFlx{(‑aði)}[13] \dicFlx{dat} \textbf{1.} \dicLangCat{hovor.} \dicSynonym{hefja} \dicDirectTranslationCS{(na)startovat, spustit, spouštět, nahodit, nahazovat, uvést\,/\addthin uvádět v~chod} \dicExampleIS{starta bílnum} \dicExampleCS{nastartovat auto}  \textbf{2.} \dicLangCat{hovor.} \dicDirectTranslationCS{odstartovat, začít, zahájit}
\dicEntry[startari] \dicTerm{start··ar|i} \dicIPA{{s}{\textsubring{d}}{a}{\textsubring{r}}{\textsubring{d}}{a}{r}{\textsci}} \dicPos{m}[1] \dicFlx{(‑a, ‑ar)}[10] \dicDirectTranslationCS{startér} \dicIndirectTranslationCS{(spouštěč spalovacího motoru)}
\dicEntry[statisti] \dicTerm{stat··ist|i} \dicIPA{{s}{\textsubring{d}}{a}{\textlengthmark}{\textsubring{d}}{\textsci}{s}{\textsubring{d}}{\textsci}} \dicPos{m}[1] \dicFlx{(‑a, ‑ar)}[1] \dicDirectTranslationCS{statista, statistka}
\dicEntry[statistík] \dicTerm{statistík} \dicIPA{{s}{\textsubring{d}}{a}{\textlengthmark}{\textsubring{d}}{\textsci}{s}{\textsubring{d}}{i}{\r{g}}} \dicPos{f}[10] \dicFlx{(‑ur)}[2] \textbf{1.} \dicDirectTranslationCS{statistika} \dicIndirectTranslationCS{(obor)}  \textbf{2.} \dicDirectTranslationCS{statistika} \dicIndirectTranslationCS{(číselná evidence hromadných jevů)}
\dicEntry[statíf] \dicTerm{statíf} \dicIPA{{s}{\textsubring{d}}{a}{\textlengthmark}{\textsubring{d}}{i}{f}} \dicPos{n}[2] \dicFlx{(‑s, ‑)}[5] \dicLangCat{hovor.} \dicSynonym{standur} \dicDirectTranslationCS{stojan, držák} \dicExampleIS{nótnastatíf} \dicExampleCS{stojan na noty}
\dicEntry[status] \dicTerm{status} \dicIPA{{s}{\textsubring{d}}{a}{\textlengthmark}{\textsubring{d}}{\textscy}{s}} \dicPos{m}[4] \dicFlx{(‑s, ‑ar)}[25] \dicLangCat{hovor.} \dicSynonym{staða} \dicDirectTranslationCS{status, stav, pozice}
\dicEntry[staukur] \dicTerm{stauk|ur} \dicIPA{{s}{\textsubring{d}}{\oe i}{\textlengthmark}{\r{g}}{\textscy}{\textsubring{r}}} \dicPos{m}[6] \dicFlx{(‑s, ‑ar)}[22] \textbf{1.} \dicSynonym*{lítill stampur} \dicDirectTranslationCS{nádobka, lahvička} \dicExampleIS{Hún tók töflur úr mörgum staukum.} \dicExampleCS{Vzala pilulky z~mnoha lahviček.}  \textbf{2.} \dicDirectTranslationCS{sypátko} \dicExampleIS{saltstaukur} \dicExampleCS{slánka}
\dicEntry[staulast] \dicTerm{staul|ast} \dicIPA{{s}{\textsubring{d}}{\oe i}{\textlengthmark}{l}{a}{s}{\textsubring{d}}} \dicPos{v}[1] \dicFlx{(‑aðist)}[95] \dicFlx{refl} \dicDirectTranslationCS{belhat se, vléct se} \dicExampleIS{staulast upp brattann} \dicExampleCS{belhat se do kopce}
\dicEntry[staup] \dicTerm{staup} \dicIPA{{s}{\textsubring{d}}{\oe i}{\textlengthmark}{\textsubring{b}}} \dicPos{n}[2] \dicFlx{(‑s, ‑)}[5] \dicDirectTranslationCS{štamprle, sklenička, kalíšek, stopka};  \dicPhraseIS{fá sér í staupinu} \dicLangCat{přen.} \dicDirectTranslationCS{dát si do nosu};  \dicPhraseIS{gefa e‑m í staupinu} \dicDirectTranslationCS{nalít (komu) alkohol}
\dicEntry[staur] \dicTerm{staur} \dicIPA{{s}{\textsubring{d}}{\oe i}{\textlengthmark}{\textsubring{r}}} \dicPos{m}[4] \dicFlx{(‑s, ‑ar)}[14] \dicDirectTranslationCS{sloup, sloupek} \dicExampleIS{símastaur} \dicExampleCS{telefonní sloup}
\dicEntry[staurblankur] \textls[15]{\dicTerm{staur··|blankur} \dicIPA{{s}{\textsubring{d}}{\oe i}{r}{\textsubring{b}}{l}{au}{\r{\ng}}{\r{g}}{\textscy}{\textsubring{r}}} \dicPos{adj}[1] \dicFlx{(f ‑blönk)}[3] \dicLangCat{hovor.}} \dicSynonym{skít\-blankur} \dicDirectTranslationCS{(jsoucí) úplně bez peněz, (jsoucí) švorc}
\dicEntry[staut] \dicTerm{staut} \dicIPA{{s}{\textsubring{d}}{\oe i}{\textlengthmark}{\textsubring{d}}} \dicPos{n}[2] \dicFlx{(‑s)}[2] \dicDirectTranslationCS{lámavé čtení, koktání (při čtení)}
\dicEntry[stauta] \dicTerm{staut|a} \dicIPA{{s}{\textsubring{d}}{\oe i}{\textlengthmark}{\textsubring{d}}{a}} \dicPos{v}[1] \dicFlx{(‑aði)}[1] \dicDirectTranslationCS{číst s~obtížemi, zakoktávat se (při čtení)};  \dicPhraseIS{stauta sig gegnum e‑ð} \dicLangCat{přen.} \dicDirectTranslationCS{prokousat se (čím) (textem ap.)}
\dicEntry[stál] \dicTerm{stál} \dicIPA{{s}{\textsubring{d}}{au}{\textlengthmark}{\textsubring{l}}} \dicPos{n}[2] \dicFlx{(‑s, ‑)}[5] \dicDirectTranslationCS{ocel};  \dicIdiom{stál}{ \dicPhraseIS{láta sverfa til stáls}} \dicLangCat{přen.} \dicDirectTranslationCS{přitlačit na pilu};  \dicPhraseIS{stappa í e‑n stálinu} \dicLangCat{přen.} \dicDirectTranslationCS{vlít (komu) krev do žil};  \dicPhraseIS{þar mætast stálin stinn} \dicFlx{refl} \dicLangCat{přen.} \dicDirectTranslationCS{padla kosa na kámen}
\dicEntry[stálhraustur] \dicTerm{stál··hraustur} \dicIPA{{s}{\textsubring{d}}{au}{l}{\textsubring{r}}{\oe i}{s}{\textsubring{d}}{\textscy}{\textsubring{r}}} \dicPos{adj}[1]\dicFlx{}[-10] \dicDirectTranslationCS{mající železné zdra\-ví, zdravý jako rybička}
\dicEntry[stállunga] \dicTerm{stál··lung|a} \dicIPA{{s}{\textsubring{d}}{au}{l}{\textlengthmark}{u}{\ng}{\r{g}}{a}} \dicPos{n}[1] \dicFlx{(‑a, ‑u)}[2] \dicFieldCat{med.} \dicDirectTranslationCS{dýchací přístroj, respirátor}
\dicEntry[stálminni] \dicTerm{stál··minni} \dicIPA{{s}{\textsubring{d}}{au}{l}{m}{\textsci}{n}{\textsci}} \dicPos{n}[2] \dicFlx{(‑s, ‑)}[14] \dicDirectTranslationCS{výborná paměť}
\dicEntry[stálnagli] \dicTerm{stál··nagl|i} \dicIPA{{s}{\textsubring{d}}{au}{l}{n}{a}{\r{g}}{l}{\textsci}} \dicPos{m}[1] \dicFlx{(‑a, ‑ar)}[8] \dicDirectTranslationCS{ocelový hřebík}
\dicEntry[stálpaður] \dicTerm{stálp··|aður} \dicIPA{{s}{\textsubring{d}}{au}{\textsubring{l}}{\textsubring{b}}{a}{ð}{\textscy}{\textsubring{r}}} \dicPos{adj}[3] \dicFlx{(f ‑uð)}[3] \dicDirectTranslationCS{dospívající, mladistvý}
\dicEntry[stálpast] \dicTerm{stálp|ast} \dicIPA{{s}{\textsubring{d}}{au}{\textsubring{l}}{\textsubring{b}}{a}{s}{\textsubring{d}}} \dicPos{v}[1] \dicFlx{(‑aðist)}[95] \dicFlx{refl} \dicSynonym*{vaxa upp} \dicDirectTranslationCS{dospívat, dorůstat} \dicExampleIS{Börnin eru farin að stálpast.} \dicExampleCS{Děti začínají dospívat.}
\dicEntry[stálsleginn] \dicTerm{stál··sleginn} \dicIPA{{s}{\textsubring{d}}{au}{l}{s}{\textsubring{d}}{l}{ei}{\textsci}{\textsubring{n}}} \dicPos{adj}[6]\dicFlx{}[-6] \dicDirectTranslationCS{(jsoucí) ve skvělé\,/\addthin vrcholné formě}
\dicEntry[stálull] \dicTerm{stál··ull} \dicIPA{{s}{\textsubring{d}}{au}{\textlengthmark}{l}{\textscy}{\textsubring{d}}{\textsubring{l}}} \dicPos{f}[4] \dicFlx{(‑ar)}[3] \dicDirectTranslationCS{ocelová vlna}
\dicEntry[stálum] \dicTerm{stálum} \dicIPA{{s}{\textsubring{d}}{au}{\textlengthmark}{l}{\textscy}{\textsubring{m}}} \dicPos{v} \dicFlx{ind pf pl 1 pers} \dicLink{stela}
\dicEntry[stáss] \dicTerm{stáss} \dicIPA{{s}{\textsubring{d}}{au}{s}{\textlengthmark}} \dicPos{n}[2] \dicFlx{(‑)}[24] \dicDirectTranslationCS{ozdoba, příkrasa, okrasa}
\dicEntry[státa] \dicTerm{stát|a} \dicIPA{{s}{\textsubring{d}}{au}{\textlengthmark}{\textsubring{d}}{a}} \dicPos{v}[1] \dicFlx{(‑aði)}[1] \dicSynonym{gorta} \dicDirectTranslationCS{pyšnit se, honosit se, holedbat se};  \dicPhraseIS{státa af e‑u} \dicDirectTranslationCS{pyšnit se (čím)} \dicExampleIS{státa af sigrinum} \dicExampleCS{honosit se vítězstvím}
\dicEntry[steðja] \dicTerm{steðj|a} \dicIPA{{s}{\textsubring{d}}{\textepsilon}{ð}{j}{a}} \dicPos{v}[1] \dicFlx{(‑aði)}[1] \dicSynonym*{ganga hratt} \dicDirectTranslationCS{jít rázně\,/\addthin energicky, chvátat, hnát se} \dicExampleIS{steðja áfram stórum skrefum} \dicExampleCS{chvátat dlouhými kroky};  \dicIdiom{steðja}[að]{ \dicPhraseIS{e‑að (illt) steðjar að e‑m}} \dicDirectTranslationCS{(co) (špatného) (koho) postihne, (co) (špatného) se na (koho) žene}
\dicEntry[steðji] \dicTerm{steðj|i} \dicIPA{{s}{\textsubring{d}}{\textepsilon}{ð}{j}{\textsci}} \dicPos{m}[1] \dicFlx{(‑a, ‑ar)}[1] \textbf{1.} \dicDirectTranslationCS{kovadlina}  \textbf{2.} \dicFieldCat{anat.} \dicDirectTranslationCS{kovadlinka}
\dicEntry[stef] \dicTerm{stef} \dicIPA{{s}{\textsubring{d}}{\textepsilon}{\textlengthmark}{f}} \dicPos{n}[2] \dicFlx{(‑s, ‑)}[13] \textbf{1.} \dicFieldCat{lit.} \dicSynonym{viðlag} \dicDirectTranslationCS{refrén}  \textbf{2.} \dicFieldCat{hud.} \dicSynonym{tema} \dicDirectTranslationCS{refrén, téma, motiv}  \textbf{3.} \dicFieldCat{poč.} \dicDirectTranslationCS{standardní program, modul}
\dicEntry[stefna] \dicTerm{stefn|a\smash{\textsuperscript{1}}} \dicsymFrequent\  \dicIPA{{s}{\textsubring{d}}{\textepsilon}{\textsubring{b}}{n}{a}} \dicPos{f}[1] \dicFlx{(‑u, ‑ur)}[7] \textbf{1.} \dicSynonym{átt\smash{\textsuperscript{1}}} \dicDirectTranslationCS{směr, kurz, kurs} \dicExampleIS{ganga í rétta stefnu} \dicExampleCS{jít správným směrem};  \dicPhraseIS{e‑að tekur (nýja) stefnu} \dicLangCat{přen.} \dicDirectTranslationCS{(co) nabírá (nový) směr};  \dicPhraseIS{setja stefnuna á e‑ð} \dicDirectTranslationCS{zacílit na (co), zaměřit se na (co)};  \dicPhraseIS{taka stefnuna e‑s staðar} \dicDirectTranslationCS{zamířit (kam), nabrat kurz (kam)}  \textbf{2.} \dicSynonym*{stefnumið} \dicDirectTranslationCS{politika, program, směr, strategie (vládní ap.)};  \dicPhraseIS{fylgja stefnu e‑s} \dicDirectTranslationCS{následovat program (čeho) (politické strany ap.)}  \textbf{3.} \dicFieldCat{práv.} \dicDirectTranslationCS{předvolání, obsílka};  \dicPhraseIS{birta e‑m stefnu} \dicDirectTranslationCS{předvolat (koho)}
\dicEntry[stefna] \dicTerm{stefn|a\smash{\textsuperscript{2}}} \dicsymFrequent\  \dicIPA{{s}{\textsubring{d}}{\textepsilon}{\textsubring{b}}{n}{a}} \dicPos{v}[2] \dicFlx{(‑di, ‑t)}[157] \dicFlx{dat} \textbf{1.} \dicDirectTranslationCS{mířit, směřovat, mít namířeno} \dicExampleIS{stefna til fjalls} \dicExampleCS{mířit k~hoře}  \textbf{2.} \dicFlx{dat} \dicDirectTranslationCS{svolat, svolávat, pozvat} \dicExampleIS{stefna fjölskyldunni saman heima hjá sér} \dicExampleCS{svolat rodinu k~sobě domů}  \textbf{3.} \dicFlx{dat} \dicFieldCat{práv.} \dicDirectTranslationCS{předvolat} \dicExampleIS{stefna honum fyrir rétt} \dicExampleCS{předvolat ho před soud}  \textbf{4.} \dicDirectTranslationCS{spět, vyvíjet se};  \dicPhraseIS{það stefnir í e‑u} \dicFlx{impers} \dicDirectTranslationCS{spěje to k~(čemu), zadělává se na (co) (rekord ap.)};  \dicIdiom{stefna}[að]{ \dicPhraseIS{stefna að e‑u}} \dicDirectTranslationCS{dospět k~(čemu), směřovat k~(čemu)};  \dicIdiom{stefna}[á]{ \dicPhraseIS{stefna á e‑ð}} \dicDirectTranslationCS{mít v~plánu (co), směřovat k~(čemu)}
\dicEntry[stefnandi] \dicTerm{stefn··|andi} \dicIPA{{s}{\textsubring{d}}{\textepsilon}{\textsubring{b}}{n}{a}{n}{\textsubring{d}}{\textsci}} \dicPos{m}[2] \dicFlx{(‑anda, ‑endur)}[1] \dicFieldCat{práv.} \dicDirectTranslationCS{žalobce, žalobkyně, navrhovatel(ka)}
\dicEntry[stefni] \dicTerm{stefni} \dicIPA{{s}{\textsubring{d}}{\textepsilon}{\textsubring{b}}{n}{\textsci}} \dicPos{n}[2] \dicFlx{(‑s, ‑)}[14] \dicDirectTranslationCS{příď} \dicIndirectTranslationCS{(přední část lodi)}
\dicEntry[stefnubreyting] \dicTerm{stefnu··breyt·ing} \dicIPA{{s}{\textsubring{d}}{\textepsilon}{\textsubring{b}}{n}{\textscy}{\textsubring{b}}{r}{ei}{\textsubring{d}}{i}{\ng}{\r{g}}} \dicPos{f}[4] \dicFlx{(‑ar, ‑ar)}[5] \textbf{1.} \dicDirectTranslationCS{změna směru}  \textbf{2.} \dicDirectTranslationCS{změna směru\,/\addthin politiky (vládní ap.)}
\dicEntry[stefnufastur] \dicTerm{stefnu··|fastur} \dicIPA{{s}{\textsubring{d}}{\textepsilon}{\textsubring{b}}{n}{\textscy}{f}{a}{s}{\textsubring{d}}{\textscy}{\textsubring{r}}} \dicPos{adj}[1] \dicFlx{(f ‑föst)}[11] \dicDirectTranslationCS{cílevědomý, odhodlaný}
\dicEntry[stefnufesta] \dicTerm{stefnu··fest|a} \dicIPA{{s}{\textsubring{d}}{\textepsilon}{\textsubring{b}}{n}{\textscy}{f}{\textepsilon}{s}{\textsubring{d}}{a}} \dicPos{f}[1] \dicFlx{(‑u)}[5] \dicDirectTranslationCS{cílevědomost, odhodlanost}
\dicEntry[stefnulaus] \dicTerm{stefnu··laus} \dicIPA{{s}{\textsubring{d}}{\textepsilon}{\textsubring{b}}{n}{\textscy}{l}{\oe i}{s}} \dicPos{adj}[5]\dicFlx{}[-1] \dicDirectTranslationCS{bezradný, tápavý}
\dicEntry[stefnuljós] \dicTerm{stefnu··ljós} \dicIPA{{s}{\textsubring{d}}{\textepsilon}{\textsubring{b}}{n}{\textscy}{l}{j}{ou}{s}} \dicPos{n}[2] \dicFlx{(‑s, ‑)}[5] \dicDirectTranslationCS{směrovka, blinkr}
\dicEntry[stefnumark] \dicTerm{stefnu··|mark} \dicIPA{{s}{\textsubring{d}}{\textepsilon}{\textsubring{b}}{n}{\textscy}{m}{a}{\textsubring{r}}{\r{g}}} \dicPos{n}[2] \dicFlx{(‑marks, ‑mörk)}[8] \dicSynonym{markmið} \dicDirectTranslationCS{cíl, směřování, záměr}
\dicEntry[stefnumerki] \dicTerm{stefnu··merki} \dicIPA{{s}{\textsubring{d}}{\textepsilon}{\textsubring{b}}{n}{\textscy}{m}{\textepsilon}{\textsubring{r}}{\r{\textObardotlessj}}{\textsci}} \dicPos{n}[2] \dicFlx{(‑s, ‑)}[16] \dicPhraseIS{gefa stefnumerki} \dicDirectTranslationCS{naznačit (směrovkami) směr}
\dicEntry[stefnumót] \dicTerm{stefnu··mót} \dicIPA{{s}{\textsubring{d}}{\textepsilon}{\textsubring{b}}{n}{\textscy}{m}{ou}{\textsubring{d}}} \dicPos{n}[2] \dicFlx{(‑s, ‑)}[5] \textbf{1.} \dicSynonym*{stefnufundur} \dicDirectTranslationCS{schůzka, setkání};  \dicPhraseIS{eiga stefnumót við e‑n} \dicDirectTranslationCS{mít schůzku s~(kým)}  \textbf{2.} \dicSynonym*{ástafundur} \dicDirectTranslationCS{rande, (milostná) schůzka, dostaveníčko};  \dicPhraseIS{fara á stefnumót við e‑n} \dicDirectTranslationCS{jít na rande s~(kým)}
\dicEntry[stefnumótun] \dicTerm{stefnu··mót|un} \dicIPA{{s}{\textsubring{d}}{\textepsilon}{\textsubring{b}}{n}{\textscy}{m}{ou}{\textsubring{d}}{\textscy}{\textsubring{n}}} \dicPos{f}[7] \dicFlx{(‑unar)}[9] \dicDirectTranslationCS{utváření\,/\addthin formování politiky (vládní ap.)}
\dicEntry[stefnuskrá] \dicTerm{stefnu··skrá} \dicsymFrequent\  \dicIPA{{s}{\textsubring{d}}{\textepsilon}{\textsubring{b}}{n}{\textscy}{s}{\r{g}}{r}{au}} \dicPos{f}[4] \dicFlx{(‑r\,/\addthin ‑ar, ‑r)}[21] \dicDirectTranslationCS{(politický) program} \dicExampleIS{Þessi mál eru á stefnuskrá flokksins.} \dicExampleCS{Tyto otázky jsou v~stranickém programu.}
\dicEntry[steggur] \dicTerm{stegg|ur} \dicIPA{{s}{\textsubring{d}}{\textepsilon}{\r{g}}{\textlengthmark}{\textscy}{\textsubring{r}}} \dicPos{m}[9] \dicFlx{(‑s, ‑ir)}[15] \dicSynonym*{karlfugl sundfugla} \dicDirectTranslationCS{houser, kačer}
\dicEntry[steig] \dicTerm{steig} \dicIPA{{s}{\textsubring{d}}{ei}{\textlengthmark}{x}} \dicPos{v} \dicFlx{ind pf sg 1 pers} \dicLink{stíga}
\dicEntry[steigurlátur] \dicTerm{steigur··látur} \dicIPA{{s}{\textsubring{d}}{ei}{\textbabygamma}{\textscy}{r}{l}{au}{\textsubring{d}}{\textscy}{\textsubring{r}}} \dicPos{adj}[1]\dicFlx{}[-1] \dicSynonym{drambsamur} \dicDirectTranslationCS{arogantní, domýšlivý}
\dicEntry[steik] \dicTerm{steik} \dicIPA{{s}{\textsubring{d}}{ei}{\textlengthmark}{\r{g}}} \dicPos{f}[8] \dicFlx{(‑ar, ‑ur)}[1] \dicFieldCat{kulin.} \dicDirectTranslationCS{steak} \dicIndirectTranslationCS{(zprudka opečený plátek masa)};  \dicPhraseIS{e‑að er allt í steik} \dicDirectTranslationCS{(co) je naprostý chaos}
\dicEntry[steikarskúffa] \dicTerm{steikar··skúff|a} \dicIPA{{s}{\textsubring{d}}{ei}{\r{g}}{a}{\textsubring{r}}{s}{\r{g}}{u}{f}{a}} \dicPos{f}[1] \dicFlx{(‑u, ‑ur)}[19] \dicDirectTranslationCS{pekáč}
\dicEntry[steikja] \dicTerm{steik|ja} \dicIPA{{s}{\textsubring{d}}{ei}{\textlengthmark}{\r{\textObardotlessj}}{a}} \dicPos{v}[2] \dicFlx{(‑ti, ‑t)}[25] \dicFlx{acc} \dicDirectTranslationCS{(u)smažit, osmažit, (o)péct} \dicExampleIS{steikja kjötið á pönnu} \dicExampleCS{smažit maso na pánvi};  \dicIdiom{steikjast}{ \dicPhraseIS{steikjast}} \dicFlx{refl} \dicDirectTranslationCS{usmažit se, opéct se}
\dicEntry[steiktur] \dicTerm{steiktur} \dicIPA{{s}{\textsubring{d}}{ei}{x}{\textsubring{d}}{\textscy}{\textsubring{r}}} \dicPos{adj}[1]\dicFlx{}[-10] \dicDirectTranslationCS{smažený, pečený} \dicExampleIS{steiktur laukur} \dicExampleCS{smažená cibule}
\dicEntry[steinafræði] \dicTerm{steina··fræð|i} \dicIPA{{s}{\textsubring{d}}{ei}{\textlengthmark}{n}{a}{f}{r}{a}{i}{ð}{\textsci}} \dicPos{f}[3] \dicFlx{(‑i)}[3] \dicFieldCat{geol.} \dicSynonym{bergfræði} \dicDirectTranslationCS{petrologie}
\dicEntry[steinaldin] \dicTerm{stein··aldin} \dicIPA{{s}{\textsubring{d}}{ei}{\textlengthmark}{n}{a}{l}{\textsubring{d}}{\textsci}{\textsubring{n}}} \dicPos{n}[2] \dicFlx{(‑s, ‑)}[5] \dicFieldCat{bot.} \dicDirectTranslationCS{peckovice}
\dicEntry[steinbítur] \dicTerm{stein··bít|ur} \dicIPA{{s}{\textsubring{d}}{ei}{n}{\textsubring{b}}{i}{\textsubring{d}}{\textscy}{\textsubring{r}}} \dicPos{m}[6] \dicFlx{(‑s, ‑ir\,/\addthin ‑ar)}[70] \dicFieldCat{zool.} \dicDirectTranslationCS{vlkouš, vlkouš obecný} \textit{(l.~{\textLA{Anarhichas lupus}})}
\dicEntry[steinbrjótur] \dicTerm{stein··brjót|ur} \dicIPA{{s}{\textsubring{d}}{ei}{n}{\textsubring{b}}{r}{j}{ou}{\textsubring{d}}{\textscy}{\textsubring{r}}} \dicPos{m}[6] \dicFlx{(‑s, ‑ar)}[10] \dicFieldCat{bot.} \dicDirectTranslationCS{lomikámen} \textit{(l.~{\textLA{Saxifraga}})}  \dicsymPhoto\ 
\dicFigure{ds_image_steinbrjotur_0_2.jpg}{Steinbrjótur}{Steinbrjótur - Bernd Haynold, GFDL}
\dicEntry[steind] \dicTerm{steind} \dicIPA{{s}{\textsubring{d}}{ei}{n}{\textsubring{d}}} \dicPos{f}[7] \dicFlx{(‑ar, ‑ir)}[1] \dicFieldCat{geol.} \dicDirectTranslationCS{minerál, nerost}
\dicEntry[steindafræði] \dicTerm{steinda··fræð|i} \dicIPA{{s}{\textsubring{d}}{ei}{n}{\textsubring{d}}{a}{f}{r}{a}{i}{ð}{\textsci}} \dicPos{f}[3] \dicFlx{(‑i)}[3] \dicFieldCat{geol.} \dicDirectTranslationCS{mineralogie}
\dicEntry[steindauður] \dicTerm{stein··dauður} \dicIPA{{s}{\textsubring{d}}{ei}{n}{\textsubring{d}}{\oe i}{ð}{\textscy}{\textsubring{r}}} \dicPos{adj}[2]\dicFlx{}[-10] \textbf{1.} \dicSynonym{dauður} \dicDirectTranslationCS{dočista\,/\addthin úplně mrtvý}  \textbf{2.} \dicSynonym{líflaus} \dicDirectTranslationCS{mrtvý, (jsoucí) bez života (ulice ap.)}
\dicEntry[steindepill] \dicTerm{stein··dep|ill} \dicIPA{{s}{\textsubring{d}}{ei}{n}{\textsubring{d}}{\textepsilon}{\textsubring{b}}{\textsci}{\textsubring{d}}{\textsubring{l}}} \dicPos{m}[6] \dicFlx{(‑ils, ‑lar)}[35] \dicFieldCat{zool.} \dicDirectTranslationCS{bělořit šedý} \textit{(l.~{\textLA{Oenanthe oenanthe}})}  \dicsymPhoto\ 
\dicFigure{69861.jpg}{Steindepill}{Steindepill - Ševčík Jiří, Biolib, Copyright/PD}
\dicEntry[steindepla] \dicTerm{stein··depl|a} \dicIPA{{s}{\textsubring{d}}{ei}{n}{\textsubring{d}}{\textepsilon}{h}{\textsubring{b}}{l}{a}} \dicPos{f}[1] \dicFlx{(‑u, ‑ur)}[13] \dicFieldCat{bot.} \dicDirectTranslationCS{rozrazil keříčkatý} \textit{(l.~{\textLA{Veronica fruticans}})}  \dicsymPhoto\ 
\dicFigure{73259.jpg}{Steindepla}{Steindepla - Kesl Michael, Biolib, Copyright/CC-BY-NC}
\dicEntry[steinefni] \dicTerm{stein··efni} \dicIPA{{s}{\textsubring{d}}{ei}{\textlengthmark}{n}{\textepsilon}{\textsubring{b}}{n}{\textsci}} \dicPos{n}[2] \dicFlx{(‑s, ‑)}[14] \dicFieldCat{geol.} \dicSynonym{steind} \dicDirectTranslationCS{minerál, nerost}
\dicEntry[steingeit] \dicTerm{stein··geit} \dicIPA{{s}{\textsubring{d}}{ei}{n}{\r{\textObardotlessj}}{ei}{\textsubring{d}}} \dicPos{f}[8] \dicFlx{(‑ar, ‑ur)}[1] \textbf{1.} \dicFieldCat{zool.} \dicDirectTranslationCS{kozorožec, kozorožec horský} \textit{(l.~{\textLA{Capra ibex}})}  \dicsymPhoto\   \textbf{2.} \dicDirectTranslationCS{Kozoroh} \dicIndirectTranslationCS{(znamení zvěrokruhu)} \textit{(l.~{\textLA{Capricornus}})}
\dicFigure{21473.jpg}{Steingeit}{Steingeit - Stolz, Gary M., Biolib, PD}
\dicEntry[steingerður] \dicTerm{stein··gerður} \dicIPA{{s}{\textsubring{d}}{ei}{n}{\r{\textObardotlessj}}{\textepsilon}{r}{ð}{\textscy}{\textsubring{r}}} \dicPos{adj}[2]\dicFlx{}[-4] \dicFieldCat{geol.} \dicSynonym{steinrunninn} \dicDirectTranslationCS{zkamenělý, fosilní, petrifikovaný} \dicExampleIS{steingert bein} \dicExampleCS{zkamenělá kost}
\dicEntry[steingervingafræði] \dicTerm{stein·gervinga··fræð|i} \dicIPA{{s}\-{\textsubring{d}}\-{ei}\-{n}\-{\r{\textObardotlessj}}\-{\textepsilon}\-{r}\-{v}\-{i}\-{\ng}\-{\r{g}}\-{a}\-{f}\-{r}\-{a}\-{i}\-{ð}\-{\textsci}\-} \dicPos{f}[3] \dicFlx{(‑i)}[3] \dicFieldCat{geol.} \dicDirectTranslationCS{paleontologie}

\dicEntry[steingervingafræðingur] \textls[10]{\dicTerm{stein·gervinga·fræð··ing|ur} \dicIPA{{s}\-{\textsubring{d}}\-{ei}\-{n}\-{\r{\textObardotlessj}}\-{\textepsilon}\-{r}\-{v}\-{i}\-{\ng}\-{\r{g}}\-{a}\-{f}\-{r}\-{a}\-{i}\-{ð}\-{i}\-{\ng}\-{\r{g}}\-{\textscy}\-{\textsubring{r}}}} \dicPos{m}[6] \dicFlx{(‑s, ‑ar)}[8] \dicDirectTranslationCS{paleontolog, paleontoložka}

\dicEntry[steingervingur] \dicTerm{stein··gerv·ing|ur} \dicIPA{{s}{\textsubring{d}}{ei}{n}{\r{\textObardotlessj}}{\textepsilon}{r}{v}{i}{\ng}{\r{g}}{\textscy}{\textsubring{r}}} \dicPos{m}[6] \dicFlx{(‑s, ‑ar)}[8] \dicFieldCat{geol.} \dicDirectTranslationCS{zkamenělina, fosilie, petrefakt}
\dicEntry[steingleyma] \dicTerm{stein··gleym|a} \dicIPA{{s}{\textsubring{d}}{ei}{n}{\r{g}}{l}{ei}{m}{a}} \dicPos{v}[2] \dicFlx{(‑di, ‑t)}[141] \dicFlx{dat} \dicDirectTranslationCS{úplně zapomenout}
\dicEntry[steingólf] \dicTerm{stein··gólf} \dicIPA{{s}{\textsubring{d}}{ei}{n}{\r{g}}{ou}{l}{f}} \dicPos{n}[2] \dicFlx{(‑s, ‑)}[5] \dicDirectTranslationCS{kamenná podlaha}
\dicEntry[steinhissa] \dicTerm{stein··hissa} \dicIPA{{s}{\textsubring{d}}{ei}{\textlengthmark}{\textsubring{n}}{h}{\textsci}{s}{a}} \dicPos{adj}[13] \dicFlx{indecl}[1] \dicDirectTranslationCS{ohromený, užaslý, překvapený}
\dicEntry[steinhljóð] \dicTerm{stein··hljóð} \dicIPA{{s}{\textsubring{d}}{ei}{n}{\textsubring{l}}{j}{ou}{\texttheta}} \dicPos{n}[2] \dicFlx{(‑s, ‑)}[5] \dicDirectTranslationCS{hrobové ticho}
\dicEntry[steinhús] \dicTerm{stein··hús} \dicIPA{{s}{\textsubring{d}}{ei}{\textlengthmark}{\textsubring{n}}{h}{u}{s}} \dicPos{n}[2] \dicFlx{(‑s, ‑)}[5] \dicDirectTranslationCS{kamenný dům, dům postavený z~kamene}
\dicEntry[steinhætta] \dicTerm{stein··hætt|a} \dicIPA{{s}{\textsubring{d}}{ei}{\textlengthmark}{\textsubring{n}}{h}{a}{i}{h}{\textsubring{d}}{a}} \dicPos{v}[2] \dicFlx{(‑i, ‑)}[9] \dicFlx{dat} \dicDirectTranslationCS{přestat, zanechat, ustat};  \dicPhraseIS{steinhætta við e‑ð} \dicDirectTranslationCS{zanechat (čeho), upustit od (čeho)}
\dicEntry[steinkol] \dicTerm{stein··kol} \dicIPA{{s}{\textsubring{d}}{ei}{n}{k\smash{\textsuperscript{h}}}{\textopeno}{\textsubring{l}}} \dicPos{n}[2] \dicFlx{pl}[1] \dicFieldCat{geol.} \dicDirectTranslationCS{antracit}
\dicEntry[steinlagður] \dicTerm{stein··|lagður} \dicIPA{{s}{\textsubring{d}}{ei}{n}{l}{a}{\textbabygamma}{ð}{\textscy}{\textsubring{r}}} \dicPos{adj}[2] \dicFlx{(f ‑lögð)}[3] \dicDirectTranslationCS{(vy)dlážděný}
\dicEntry[steinleggja] \dicTerm{stein··|leggja} \dicIPA{{s}{\textsubring{d}}{ei}{n}{l}{\textepsilon}{\r{\textObardotlessj}}{a}} \dicPos{v}[4] \dicFlx{(‑legg, ‑lagði, ‑lögðum, ‑legði, ‑lagt)}[19] \dicFlx{acc} \dicDirectTranslationCS{(vy)dláždit} \dicExampleIS{steinleggja götuna} \dicExampleCS{vydláždit ulici}
\dicEntry[steinlím] \dicTerm{stein··lím} \dicIPA{{s}{\textsubring{d}}{ei}{n}{l}{i}{\textsubring{m}}} \dicPos{n}[2] \dicFlx{(‑s, ‑)}[5] \dicDirectTranslationCS{malta}
\dicEntry[steinn] \dicTerm{stein|n} \dicsymFrequent\  \dicIPA{{s}{\textsubring{d}}{ei}{\textsubring{d}}{\textsubring{n}}} \dicPos{m}[6] \dicFlx{(‑s, ‑ar)}[42] \textbf{1.} \dicDirectTranslationCS{kámen, balvan, šutr} \dicExampleIS{kasta steini} \dicExampleCS{hodit kamenem}  \textbf{2.} \dicLangCat{hovor.} \dicSynonym{fangelsi} \dicDirectTranslationCS{chládek, basa, kriminál} \dicExampleIS{vera í steininum} \dicExampleCS{být v~chládku}  \textbf{3.} \dicSynonym{fræ} \dicDirectTranslationCS{pecka (švestky ap.)}  \textbf{4.} \dicFieldCat{med.} \dicDirectTranslationCS{kámen};  \dicIdiom{steinn}{ \dicPhraseIS{leggja stein í götu e‑s\,/\addthin e‑rs}} \dicLangCat{přen.} \dicDirectTranslationCS{stavět překážky (čemu\,/\addthin komu)};  \dicPhraseIS{setjast í helgan stein} \dicFlx{refl} \dicDirectTranslationCS{stáhnout se do ústraní};  \dicPhraseIS{sofa eins og steinn} \dicLangCat{přen.} \dicDirectTranslationCS{spát jako zabitý};  \dicPhraseIS{vera milli steins og sleggju} \dicLangCat{přen.} \dicDirectTranslationCS{ocitnout se mezi dvěma mlýnskými kameny};  \dicPhraseIS{það liggur fiskur undir steini} \dicLangCat{přen.} \dicDirectTranslationCS{skrývá se za tím nějaký záměr, něco za tím je};  \dicPhraseIS{það stendur ekki steinn yfir steini} \dicLangCat{přen.} \dicDirectTranslationCS{nezůstane kámen na kameni}
\dicEntry[steinolía] \dicTerm{stein··olí|a} \dicIPA{{s}{\textsubring{d}}{ei}{\textlengthmark}{n}{\textopeno}{l}{i}{j}{a}} \dicPos{f}[1] \dicFlx{(‑u)}[5] \dicFieldCat{chem.} \dicDirectTranslationCS{petrolej}
\dicEntry[steinprent] \dicTerm{stein··prent} \dicIPA{{s}{\textsubring{d}}{ei}{n}{p\smash{\textsuperscript{h}}}{r}{\textepsilon}{\textsubring{n}}{\textsubring{d}}} \dicPos{n}[2] \dicFlx{(‑s)}[2] \dicFieldCat{polygr.} \dicDirectTranslationCS{litografie, kamenotisk}
\dicEntry[steinprenta] \dicTerm{stein··prent|a} \dicIPA{{s}{\textsubring{d}}{ei}{n}{p\smash{\textsuperscript{h}}}{r}{\textepsilon}{\textsubring{n}}{\textsubring{d}}{a}} \dicPos{v}[1] \dicFlx{(‑aði)}[1] \dicFlx{acc} \dicDirectTranslationCS{litografovat}
\dicEntry[steinrunninn] \dicTerm{stein··runninn} \dicIPA{{s}{\textsubring{d}}{ei}{n}{r}{\textscy}{n}{\textsci}{\textsubring{n}}} \dicPos{adj}[6]\dicFlx{}[-15] \textbf{1.} \dicSynonym{steingerður} \dicDirectTranslationCS{zkamenělý, fosilní, petrifikovaný} \dicExampleIS{steinrunnar jurtaleifar} \dicExampleCS{zkamenělé pozůstatky rostlin}  \textbf{2.} \dicSynonym{staðnaður} \dicDirectTranslationCS{zkostnatělý, ustrnulý} \dicExampleIS{steinrunnin skoðun} \dicExampleCS{zkostnatělý názor}  \textbf{3.} \dicSynonym{stirðnaður} \dicDirectTranslationCS{zkamenělý, strnulý, ztuhlý, nehybný}
\dicEntry[steinsalt] \dicTerm{stein··|salt} \dicIPA{{s}{\textsubring{d}}{ei}{n}{s}{a}{\textsubring{l}}{\textsubring{d}}} \dicPos{n}[2] \dicFlx{(‑salts, ‑sölt)}[8] \dicFieldCat{geol.} \dicDirectTranslationCS{halit, sůl kamenná}
\dicEntry[steinsef] \dicTerm{stein··sef} \dicIPA{{s}{\textsubring{d}}{ei}{n}{s}{\textepsilon}{f}} \dicPos{v} \dicFlx{ind praes sg 1 pers} \dicLink{steinsofa}
\dicEntry[steinselja] \dicTerm{stein··selj|a} \dicIPA{{s}{\textsubring{d}}{ei}{n}{s}{\textepsilon}{l}{j}{a}} \dicPos{f}[1] \dicFlx{(‑u)}[5] \dicFieldCat{bot.} \dicSynonym{pétursselja} \dicDirectTranslationCS{petržel, petržel zahradní} \textit{(l.~{\textLA{Petroselinum crispum}})}
\dicEntry[steinskrift] \textls[10]{\dicTerm{stein··skrift} \dicIPA{{s}{\textsubring{d}}{ei}{n}{s}{\r{g}}{r}{\textsci}{f}{\textsubring{d}}} \dicPos{f}[7] \dicFlx{(‑ar)}[3] \dicFieldCat{polygr.} \dicDirectTranslationCS{bezpatkové písmo}}
\dicEntry[steinsmiður] \dicTerm{stein··smið|ur} \dicIPA{{s}{\textsubring{d}}{ei}{n}{s}{m}{\textsci}{ð}{\textscy}{\textsubring{r}}} \dicPos{m}[9] \dicFlx{(‑s, ‑ir)}[8] \dicDirectTranslationCS{kameník, kamenice}
\dicEntry[steinsmíði] \dicTerm{stein··smíð|i} \dicIPA{{s}{\textsubring{d}}{ei}{n}{s}{m}{i}{ð}{\textsci}} \dicPos{f}[2] \dicFlx{(‑i, ‑ar)}[1] \dicDirectTranslationCS{kamenictví}
\dicEntry[steinsnar] \dicTerm{stein··snar} \dicIPA{{s}{\textsubring{d}}{ei}{n}{s}{\textsubring{d}}{n}{a}{\textsubring{r}}} \dicPos{n}[2] \dicFlx{(‑s)}[2] \dicDirectTranslationCS{kousek, kousíček} \dicIndirectTranslationCS{(krátká vzdálenost)} \dicExampleIS{Það er ekki steinsnar þangað.} \dicExampleCS{Není to tam zrovna kousek.}
\dicEntry[steinsofa] \dicTerm{stein··|sofa} \dicIPA{{s}{\textsubring{d}}{ei}{n}{s}{\textopeno}{v}{a}} \dicPos{v}[6] \dicFlx{(‑sef, ‑svaf, ‑sváfum, ‑svæfi, ‑sofið)}[3] \dicDirectTranslationCS{tvrdě spát}
\dicEntry[steinsofandi] \dicTerm{stein··sof·andi} \dicIPA{{s}{\textsubring{d}}{ei}{n}{s}{\textopeno}{v}{a}{n}{\textsubring{d}}{\textsci}} \dicPos{adj}[13] \dicFlx{indecl}[1] \dicDirectTranslationCS{tvrdě spící}
\dicEntry[steinsofið] \dicTerm{stein··sofið} \dicIPA{{s}{\textsubring{d}}{ei}{n}{s}{\textopeno}{v}{\textsci}{\texttheta}} \dicPos{v} \dicFlx{supin} \dicLink{steinsofa}
\dicEntry[steinsofna] \dicTerm{stein··sofn|a} \dicIPA{{s}{\textsubring{d}}{ei}{n}{s}{\textopeno}{\textsubring{b}}{n}{a}} \dicPos{v}[1] \dicFlx{(‑aði)}[34] \dicDirectTranslationCS{tvrdě usnout}
\dicEntry[steinsteypa] \dicTerm{stein··steyp|a} \dicIPA{{s}{\textsubring{d}}{ei}{n}{s}{\textsubring{d}}{ei}{\textsubring{b}}{a}} \dicPos{f}[1] \dicFlx{(‑u)}[5] \dicDirectTranslationCS{beton}
\dicEntry[steinsvaf] \dicTerm{stein··svaf} \dicIPA{{s}{\textsubring{d}}{ei}{n}{s}{v}{a}{f}} \dicPos{v} \dicFlx{ind pf sg 1 pers} \dicLink{steinsofa}
\dicEntry[steinsváfum] \dicTerm{stein··sváfum} \dicIPA{{s}{\textsubring{d}}{ei}{n}{s}{v}{au}{v}{\textscy}{\textsubring{m}}} \dicPos{v} \dicFlx{ind pf pl 1 pers} \dicLink{steinsofa}
\dicEntry[steinsvæfi] \dicTerm{stein··svæfi} \dicIPA{{s}{\textsubring{d}}{ei}{n}{s}{v}{a}{i}{v}{\textsci}} \dicPos{v} \dicFlx{con pf sg 1 pers} \dicLink{steinsofa}
\dicEntry[steintau] \dicTerm{stein··tau} \dicIPA{{s}{\textsubring{d}}{ei}{n}{t\smash{\textsuperscript{h}}}{\oe i}} \dicPos{n}[2] \dicFlx{(‑s)}[2] \dicDirectTranslationCS{kamenina, kameninové výrobky}
\dicEntry[steintegund] \dicTerm{stein··tegund} \dicIPA{{s}{\textsubring{d}}{ei}{n}{t\smash{\textsuperscript{h}}}{\textepsilon}{\textbabygamma}{\textscy}{n}{\textsubring{d}}} \dicPos{f}[7] \dicFlx{(‑ar, ‑ir)}[1] \dicFieldCat{geol.} \dicSynonym{steind} \dicDirectTranslationCS{minerál, nerost}
\dicEntry[steinull] \dicTerm{stein··ull} \dicIPA{{s}{\textsubring{d}}{ei}{\textlengthmark}{n}{\textscy}{\textsubring{d}}{\textsubring{l}}} \dicPos{f}[4] \dicFlx{(‑ar)}[3] \dicDirectTranslationCS{minerální vlna\,/\addthin vata}
\dicEntry[steinuppgefinn] \dicTerm{stein··upp·gefinn} \dicIPA{{s}{\textsubring{d}}{ei}{\textlengthmark}{n}{\textscy}{h}{\textsubring{b}}{\r{\textObardotlessj}}{\textepsilon}{v}{\textsci}{\textsubring{n}}} \dicPos{adj}[6]\dicFlx{}[-6] \dicSynonym{dauðþreyttur} \dicDirectTranslationCS{totálně vyčerpaný\,/\addthin zmožený}
\dicEntry[steinveggur] \dicTerm{stein··vegg|ur} \dicIPA{{s}{\textsubring{d}}{ei}{n}{v}{\textepsilon}{\r{g}}{\textscy}{\textsubring{r}}} \dicPos{m}[10] \dicFlx{(‑jar, ‑ir)}[23] \dicDirectTranslationCS{kamenná zeď}
\dicEntry[steinþagað] \dicTerm{stein··þagað} \dicIPA{{s}{\textsubring{d}}{ei}{n}{\texttheta}{a}{\textbabygamma}{a}{\texttheta}} \dicPos{v} \dicFlx{supin} \dicLink{steinþegja}
\dicEntry[steinþagði] \dicTerm{stein··þagði} \dicIPA{{s}{\textsubring{d}}{ei}{n}{\texttheta}{a}{\textbabygamma}{ð}{\textsci}} \dicPos{v} \dicFlx{ind pf sg 1 pers} \dicLink{steinþegja}
\dicEntry[steinþagna] \dicTerm{stein··þagn|a} \dicIPA{{s}{\textsubring{d}}{ei}{n}{\texttheta}{a}{\r{g}}{n}{a}} \dicPos{v}[1] \dicFlx{(‑aði)}[13] \dicDirectTranslationCS{(zcela) utichnout\,/\addthin umlknout}
\dicEntry[steinþegði] \dicTerm{stein··þegði} \dicIPA{{s}{\textsubring{d}}{ei}{n}{\texttheta}{\textepsilon}{\textbabygamma}{ð}{\textsci}} \dicPos{v} \dicFlx{con pf sg 1 pers} \dicLink{steinþegja}
\dicEntry[steinþegi] \dicTerm{stein··þegi} \dicIPA{{s}{\textsubring{d}}{ei}{n}{\texttheta}{ei}{\textsci}} \dicPos{v} \dicFlx{ind praes sg 1 pers} \dicLink{steinþegja}
\dicEntry[steinþegja] \dicTerm{stein··|þegja} \dicIPA{{s}{\textsubring{d}}{ei}{n}{\texttheta}{ei}{j}{a}} \dicPos{v}[4] \dicFlx{(‑þegi, ‑þagði, ‑þögðum, ‑þegði, ‑þagað)}[17] \dicDirectTranslationCS{být naprosto zticha, nedutat, nemukat}
\dicEntry[steinþögðum] \dicTerm{stein··þögðum} \dicIPA{{s}{\textsubring{d}}{ei}{n}{\texttheta}{\oe}{\textbabygamma}{ð}{\textscy}{\textsubring{m}}} \dicPos{v} \dicFlx{ind pf pl 1 pers} \dicLink{steinþegja}
\dicEntry[steinöld] \dicTerm{stein··|öld} \dicIPA{{s}{\textsubring{d}}{ei}{\textlengthmark}{n}{\oe}{l}{\textsubring{d}}} \dicPos{f}[7] \dicFlx{(‑aldar)}[19] \dicFieldCat{hist.} \dicDirectTranslationCS{doba kamenná}
\dicEntry[stekk] \dicTerm{stekk} \dicIPA{{s}{\textsubring{d}}{\textepsilon}{h}{\r{g}}} \dicPos{v} \dicFlx{ind praes sg 1 pers} \dicLink{stökkva\smash{\textsuperscript{2}}}
\dicEntry[Stekkjastaur] \dicTerm{Stekkja··staur} \dicIPA{{s}{\textsubring{d}}{\textepsilon}{h}{\r{\textObardotlessj}}{a}{s}{\textsubring{d}}{\oe i}{\textsubring{r}}} \dicPos{m}[4] \dicFlx{(‑s)}[9] \dicFlx{prop} \dicFieldCat{pov.} \dicDirectTranslationCS{Stekkja\-staur} \dicIndirectTranslationCS{(1. vánoční skřítek, který krade mléko ovcím; přichází 12. prosince)}
\dicEntry[stel] \dicTerm{stel} \dicIPA{{s}{\textsubring{d}}{\textepsilon}{\textlengthmark}{\textsubring{l}}} \dicPos{v} \dicFlx{ind praes sg 1 pers} \dicLink{stela}
\dicEntry[stela] \dicTerm{stela} \dicsymFrequent\  \dicIPA{{s}{\textsubring{d}}{\textepsilon}{\textlengthmark}{l}{a}} \dicPos{v}[6] \dicFlx{(stel, stal, stálum, stæli, stolið)}[14] \dicFlx{dat} \dicDirectTranslationCS{(u)krást, ukradnout, (u)loupit};  \dicPhraseIS{stela e‑u frá e‑m} \dicDirectTranslationCS{ukrást (co komu)} \dicExampleIS{stela peningunum frá henni} \dicExampleCS{ukrást jí peníze};  \dicIdiom{stela}[undan]{ \dicPhraseIS{stela undan skatti}} \dicDirectTranslationCS{obcházet daně};  \dicIdiom{stelast}{ \dicPhraseIS{stelast}} \dicFlx{refl} \dicSynonym*{laumast} \dicDirectTranslationCS{(při)krást se};  \dicIdiom{stelast}[að]{ \dicPhraseIS{stelast að e‑m}} \dicFlx{refl} \dicDirectTranslationCS{přikrást se ke (komu)};  \dicIdiom{stelast}[á burt]{ \dicPhraseIS{stelast á burt}} \dicFlx{refl} \dicDirectTranslationCS{odkrást se, odplížit se}
\dicEntry[stelkur] \dicTerm{stelk|ur} \dicIPA{{s}{\textsubring{d}}{\textepsilon}{\textsubring{l}}{\r{g}}{\textscy}{\textsubring{r}}} \dicPos{m}[6] \dicFlx{(‑s, ‑ar)}[15] \dicFieldCat{zool.} \dicDirectTranslationCS{vodouš, vodouš rudonohý} \textit{(l.~{\textLA{Tringa totanus}})}  \dicsymPhoto\ 
\dicFigure{ds_image_stelkur_0_1.jpg}{Stelkur}{Stelkur - Dorota Nierychlewska-Chejn, GPL}
\dicEntry[stell] \dicTerm{stell} \dicIPA{{s}{\textsubring{d}}{\textepsilon}{l}{\textlengthmark}} \dicPos{n}[2] \dicFlx{(‑s, ‑)}[5] \textbf{1.} \dicSynonym{samstæða} \dicDirectTranslationCS{souprava (jídelní ap.)} \dicExampleIS{bollastell} \dicExampleCS{souprava šálků}  \textbf{2.} \dicSynonym*{stoðgrind} \dicDirectTranslationCS{konstrukce, rám (kola ap.)} \dicExampleIS{hjólastell} \dicExampleCS{rám kola}  \textbf{3.} \dicLangCat{hovor.} \dicSynonym*{falskar tennur} \dicDirectTranslationCS{umělý chrup}
\dicEntry[stelling] \dicTerm{stell··ing} \dicIPA{{s}{\textsubring{d}}{\textepsilon}{\textsubring{d}}{l}{i}{\ng}{\r{g}}} \dicPos{f}[4] \dicFlx{(‑ar, ‑ar)}[5] \dicDirectTranslationCS{postoj, póza, držení těla} \dicExampleIS{vera í góðum\,/\addthin illum stellingum} \dicExampleCS{mít dobré\,/\addthin špatné držení těla}
\dicEntry[stelpa] \dicTerm{stelp|a} \dicsymFrequent\  \dicIPA{{s}{\textsubring{d}}{\textepsilon}{\textsubring{l}}{\textsubring{b}}{a}} \dicPos{f}[1] \dicFlx{(‑u, ‑ur)}[13] \dicSynonym{stúlka} \dicDirectTranslationCS{dívka, holka, děvče} \dicExampleIS{tvær stelpur og fjórir strákar} \dicExampleCS{dvě holky a~čtyři kluci}
\dicEntry[stelpulegur] \dicTerm{stelpu··legur} \dicIPA{{s}{\textsubring{d}}{\textepsilon}{\textsubring{l}}{\textsubring{b}}{\textscy}{l}{\textepsilon}{\textbabygamma}{\textscy}{\textsubring{r}}} \dicPos{adj}[1]\dicFlx{}[-8] \dicDirectTranslationCS{dívčí, holčičí}
\dicEntry[stelpustrákur] \dicTerm{stelpu··strák|ur} \dicIPA{{s}{\textsubring{d}}{\textepsilon}{\textsubring{l}}{\textsubring{b}}{\textscy}{s}{\textsubring{d}}{r}{au}{\r{g}}{\textscy}{\textsubring{r}}} \dicPos{m}[6] \dicFlx{(‑s, ‑ar)}[15] \dicDirectTranslationCS{mazánek, zženštilý kluk}
\dicEntry[stelsjúkur] \dicTerm{stel··sjúkur} \dicIPA{{s}{\textsubring{d}}{\textepsilon}{l}{s}{j}{u}{\r{g}}{\textscy}{\textsubring{r}}} \dicPos{adj}[1]\dicFlx{}[-1] \dicDirectTranslationCS{kleptomanský}
\dicEntry[stelsýki] \dicTerm{stel··sýk|i} \dicIPA{{s}{\textsubring{d}}{\textepsilon}{l}{s}{i}{\r{\textObardotlessj}}{\textsci}} \dicPos{f}[3] \dicFlx{(‑i)}[3] \dicFieldCat{med.} \dicDirectTranslationCS{kleptomanie}
\dicEntry[stelvís] \dicTerm{stel··vís} \dicIPA{{s}{\textsubring{d}}{\textepsilon}{l}{v}{i}{s}} \dicPos{adj}[5]\dicFlx{}[-1] \dicSynonym{þjófóttur} \dicDirectTranslationCS{zlodějský, nenechavý}
\dicEntry[stemma] \dicTerm{stemm|a} \dicIPA{{s}{\textsubring{d}}{\textepsilon}{m}{\textlengthmark}{a}} \dicPos{v}[2] \dicFlx{(‑di, ‑t)}[155] \textbf{1.} \dicDirectTranslationCS{sedět, pasovat, souhlasit (výsledky ap.)} \dicExampleIS{Það stemmir.} \dicExampleCS{To sedí.}  \textbf{2.} \dicPhraseIS{stemma stigu við e‑u} \dicDirectTranslationCS{udělat (čemu) přítrž}
\dicEntry[stemming] \dicTerm{stemm··ing} \dicIPA{{s}{\textsubring{d}}{\textepsilon}{m}{\textlengthmark}{i}{\ng}{\r{g}}} \dicPos{f}[4] \dicFlx{(‑ar)}[7] \dicLink{stemning}
\dicEntry[stemmning] \dicTerm{stemm··ning} \dicIPA{{s}{\textsubring{d}}{\textepsilon}{m}{n}{i}{\ng}{\r{g}}} \dicPos{f}[4] \dicFlx{(‑ar)}[7] \dicLink{stemning}
\dicEntry[stemning] \dicTerm{stem··ning}\dicTerm{, stemmning}\dicTerm{,  stemming} \dicIPA{{s}\-{\textsubring{d}}\-{\textepsilon}\-{m}\-{n}\-{i}\-{\ng}\-{\r{g}}\-} \dicPos{f}[4] \dicFlx{(‑ar)}[7] \dicSynonym{skap} \dicDirectTranslationCS{(příjemné) napětí\,/\addthin vzrušení, (radostné) očekávání, (radostná) nálada, euforie} \dicExampleIS{Hann var í stemningu.} \dicExampleCS{Měl radostnou náladu.}
\dicEntry[stend] \dicTerm{stend} \dicIPA{{s}{\textsubring{d}}{\textepsilon}{n}{\textsubring{d}}} \dicPos{v} \dicFlx{ind praes sg 1 pers} \dicLink{standa}
\dicEntry[stengur] \dicTerm{stengur} \dicIPA{{s}{\textsubring{d}}{ei}{\ng}{\r{g}}{\textscy}{\textsubring{r}}} \dicPos{f} \dicFlx{pl nom} \dicLink{stöng}
\dicEntry[stensill] \dicTerm{stens|ill} \dicIPA{{s}{\textsubring{d}}{\textepsilon}{n}{s}{\textsci}{\textsubring{d}}{\textsubring{l}}} \dicPos{m}[6] \dicFlx{(‑ils, ‑lar)}[35] \dicDirectTranslationCS{šablona} \dicIndirectTranslationCS{(pomůcka s~výřezy pro psaní nebo vytváření opakujících se vzorů)}
\dicEntry[stepp] \dicTerm{stepp} \dicIPA{{s}{\textsubring{d}}{\textepsilon}{h}{\textsubring{b}}} \dicPos{n}[2] \dicFlx{(‑s)}[2] \dicFieldCat{hud.} \dicDirectTranslationCS{step} \dicIndirectTranslationCS{(tanec)}
\dicEntry[steppa] \dicTerm{stepp|a\smash{\textsuperscript{1}}} \dicIPA{{s}{\textsubring{d}}{\textepsilon}{h}{\textsubring{b}}{a}} \dicPos{f}[1] \dicFlx{(‑u, ‑ur)}[19] \dicSynonym{gresja} \dicDirectTranslationCS{step} \dicIndirectTranslationCS{(rovinatá oblast porostlá trávou a~nízkými rostlinami)}
\dicEntry[steppa] \dicTerm{stepp|a\smash{\textsuperscript{2}}} \dicIPA{{s}{\textsubring{d}}{\textepsilon}{h}{\textsubring{b}}{a}} \dicPos{v}[1] \dicFlx{(‑aði)}[1] \dicDirectTranslationCS{stepovat} \dicIndirectTranslationCS{(provádět tanec)}
\dicEntry[steralyf] \dicTerm{stera··lyf} \dicIPA{{s}{\textsubring{d}}{\textepsilon}{\textlengthmark}{r}{a}{l}{\textsci}{f}} \dicPos{n}[2] \dicFlx{(‑s, ‑)}[13] \dicFieldCat{med. } \dicDirectTranslationCS{steroid}
\dicEntry[stereó] \dicTerm{stereó} \dicIPA{{s}{\textsubring{d}}{\textepsilon}{\textlengthmark}{r}{\textepsilon}{ou}} \dicPos{n}[2] \dicFlx{(‑s)}[31] \dicSynonym*{víðómur} \dicDirectTranslationCS{stereo}
\dicEntry[steri] \dicTerm{ster|i} \dicIPA{{s}{\textsubring{d}}{\textepsilon}{\textlengthmark}{r}{\textsci}} \dicPos{m}[1] \dicFlx{(‑a, ‑ar)}[1] \dicFieldCat{med. } \dicDirectTranslationCS{steroid} \dicIndirectTranslationCS{(zvláště v~pl)}
\dicEntry[sterkbyggður] \dicTerm{sterk··byggður} \dicIPA{{s}{\textsubring{d}}{\textepsilon}{\textsubring{r}}{\r{g}}{\textsubring{b}}{\textsci}{\textbabygamma}{ð}{\textscy}{\textsubring{r}}} \dicPos{adj}[2]\dicFlx{}[-1] \textbf{1.} \dicDirectTranslationCS{statný, robustní} \dicExampleIS{sterkbyggður maður} \dicExampleCS{statný muž}  \textbf{2.} \dicDirectTranslationCS{robustní, masivní (stůl ap.)}
\dicEntry[sterkeind] \dicTerm{sterk··eind} \dicIPA{{s}{\textsubring{d}}{\textepsilon}{\textsubring{r}}{\r{g}}{ei}{n}{\textsubring{d}}} \dicPos{f}[7] \dicFlx{(‑ar, ‑ir)}[1] \dicFieldCat{fyz.} \dicDirectTranslationCS{hadron}
\dicEntry[sterkja] \dicTerm{sterkj|a} \dicIPA{{s}{\textsubring{d}}{\textepsilon}{\textsubring{r}}{\r{\textObardotlessj}}{a}} \dicPos{f}[1] \dicFlx{(‑u)}[5] \dicFieldCat{chem.} \dicSynonym{mjölvi} \dicDirectTranslationCS{škrob} \textit{(l.~{\textLA{Amylum}})}
\dicEntry[sterklega] \dicTerm{sterk··lega} \dicIPA{{s}{\textsubring{d}}{\textepsilon}{\textsubring{r}}{\r{g}}{l}{\textepsilon}{\textbabygamma}{a}} \dicPos{adv} \textbf{1.} \dicSynonym*{kraftalega} \dicDirectTranslationCS{silně, pevně, solidně}  \textbf{2.} \dicSynonym{ákveðið\smash{\textsuperscript{2}}} \dicDirectTranslationCS{pevně, silně (přesvědčený ap.)}
\dicEntry[sterklegur] \dicTerm{sterk··legur} \dicIPA{{s}{\textsubring{d}}{\textepsilon}{\textsubring{r}}{\r{g}}{l}{\textepsilon}{\textbabygamma}{\textscy}{\textsubring{r}}} \dicPos{adj}[1]\dicFlx{}[-8] \dicSynonym{þrekinn} \dicDirectTranslationCS{bytelný, mohutný, robustní, masivní}
\dicEntry[sterkur] \dicTerm{sterkur} \dicsymFrequent\  \dicIPA{{s}{\textsubring{d}}{\textepsilon}{\textsubring{r}}{\r{g}}{\textscy}{\textsubring{r}}} \dicPos{adj}[1]\dicFlx{}[-1] \textbf{1.} \dicSynonym{öflugur} \dicDirectTranslationCS{silný};  \dicPhraseIS{vera sterkur í e‑u} \dicDirectTranslationCS{být silný v~(čem) (islandštině ap.)}  \textbf{2.} \dicSynonym*{sem þolir álag} \dicDirectTranslationCS{silný, odolný, pevný} \dicExampleIS{sterkur kaðall} \dicExampleCS{pevný provaz}  \textbf{3.} \dicSynonym*{orkuríkur} \dicDirectTranslationCS{silný, mocný, prudký} \dicExampleIS{sterkur vindur} \dicExampleCS{mocný vítr}  \textbf{4.} \dicSynonym{kraftmikill} \dicDirectTranslationCS{silný, mocný} \dicExampleIS{sterkt kaffi} \dicExampleCS{silná káva};  \dicPhraseIS{sterkt vín} \dicDirectTranslationCS{silný alkohol}  \textbf{5.} \dicSynonym{ákveðinn} \dicDirectTranslationCS{silný, intenzivní (pocit ap.)}  \textbf{6.} \dicFieldCat{jaz.} \dicDirectTranslationCS{silný} \dicIndirectTranslationCS{(náležející do určité skloňovací\,/\addthin časovací třídy)} \dicAntonym{veikur};  \dicPhraseIS{sterk sögn} \dicFieldCat{jaz.} \dicDirectTranslationCS{silné sloveso}
\dicEntry[steypa] \dicTerm{steyp|a\smash{\textsuperscript{1}}} \dicIPA{{s}{\textsubring{d}}{ei}{\textlengthmark}{\textsubring{b}}{a}} \dicPos{f}[1] \dicFlx{(‑u, ‑ur)}[19] \dicSynonym{steinsteypa} \dicDirectTranslationCS{beton};  \dicPhraseIS{þetta er algjör steypa} \dicLangCat{přen.} \dicDirectTranslationCS{to je úplný nesmysl}
\dicEntry[steypa] \dicTerm{steyp|a\smash{\textsuperscript{2}}} \dicsymFrequent\  \dicIPA{{s}{\textsubring{d}}{ei}{\textlengthmark}{\textsubring{b}}{a}} \dicPos{v}[2] \dicFlx{(‑ti, ‑t)}[63] \dicFlx{acc\,/\addthin dat} \textbf{1.} \dicFlx{dat} \dicSynonym*{velta} \dicDirectTranslationCS{svrhnout, shodit, shazovat, svalit, svalovat} \dicExampleIS{steypa e‑u um koll} \dicExampleCS{převrhnout (co)};  \dicPhraseIS{steypa e‑m í glötun} \dicDirectTranslationCS{uvrhnout (koho) do záhuby};  \dicPhraseIS{steypa sér í skuldir} \dicDirectTranslationCS{zadlužit se}  \textbf{2.} \dicPhraseIS{steypa sér} \dicDirectTranslationCS{skočit (po hlavě), vrhnout se} \dicExampleIS{steypa sér út í ána} \dicExampleCS{vrhnout se do řeky};  \dicPhraseIS{steypa sér kollhnís} \dicDirectTranslationCS{udělat kotrmelec\,/\addthin kotoul}  \textbf{3.} \dicFlx{dat} \dicDirectTranslationCS{svrhnout, sesadit} \dicExampleIS{steypa konungnum} \dicExampleCS{svrhnout krále};  \dicPhraseIS{steypa e‑m af stóli} \dicDirectTranslationCS{svrhnout (koho), sesadit (koho)}  \textbf{4.} \dicFlx{acc} \dicDirectTranslationCS{(vy)betonovat, postavit z~betonu} \dicExampleIS{steypa hús} \dicExampleCS{betonovat dům}  \textbf{5.} \dicFlx{acc} \dicDirectTranslationCS{odlít, odlévat (model ap.)} \dicIndirectTranslationCS{(litím zhotovit)} \dicExampleIS{steypa kerti} \dicExampleCS{odlít svíčku};  \dicIdiom{steypa}[saman]{ \dicPhraseIS{steypa e‑u saman}} \dicDirectTranslationCS{míchat (co), mísit (co), směšovat (co)};  \dicIdiom{steypast}{ \dicPhraseIS{steypast}} \dicFlx{refl} \dicDirectTranslationCS{(prudce) (s)padnout\,/\addthin padat, (z)řítit se}
\dicEntry[steypibað] \dicTerm{steypi··|bað} \dicIPA{{s}{\textsubring{d}}{ei}{\textlengthmark}{\textsubring{b}}{\textsci}{\textsubring{b}}{a}{\texttheta}} \dicPos{n}[2] \dicFlx{(‑baðs, ‑böð)}[8] \dicSynonym{sturta\smash{\textsuperscript{1}}} \dicDirectTranslationCS{sprcha}
\dicEntry[steypir] \dicTerm{steyp|ir} \dicIPA{{s}{\textsubring{d}}{ei}{\textlengthmark}{\textsubring{b}}{\textsci}{\textsubring{r}}} \dicPos{m}[7] \dicFlx{(‑is)}[2] \dicPhraseIS{e‑r er komin á steypirinn} \dicDirectTranslationCS{(kdo) (žena ap.) bude brzy rodit}
\dicEntry[steypiregn] \dicTerm{steypi··regn} \dicIPA{{s}{\textsubring{d}}{ei}{\textlengthmark}{\textsubring{b}}{\textsci}{r}{\textepsilon}{\r{g}}{\textsubring{n}}} \dicPos{n}[2] \dicFlx{(‑s)}[2] \dicDirectTranslationCS{průtrž (mračen), liják}
\dicEntry[steypireyður] \dicTerm{steypi··reyð|ur} \dicIPA{{s}{\textsubring{d}}{ei}{\textlengthmark}{\textsubring{b}}{\textsci}{r}{ei}{ð}{\textscy}{\textsubring{r}}} \dicPos{f}[5] \dicFlx{(‑ar, ‑ar)}[1] \dicFieldCat{zool.} \dicDirectTranslationCS{plejtvák obrovský} \textit{(l.~{\textLA{Balaenoptera musculus}})}
\dicEntry[steypubíll] \dicTerm{steypu··bíl|l} \dicIPA{{s}{\textsubring{d}}{ei}{\textlengthmark}{\textsubring{b}}{\textscy}{\textsubring{b}}{i}{\textsubring{d}}{\textsubring{l}}} \dicPos{m}[6] \dicFlx{(‑s, ‑ar)}[48] \dicDirectTranslationCS{míchač betonu} \dicIndirectTranslationCS{(automobil)}
\dicEntry[steypuhrærivél] \dicTerm{steypu··hræri·vél} \dicIPA{{s}{\textsubring{d}}{ei}{\textlengthmark}{\textsubring{b}}{\textscy}{\textsubring{r}}{a}{i}{r}{\textsci}{v}{j}{\textepsilon}{\textsubring{l}}} \dicPos{f}[4] \dicFlx{(‑ar, ‑ar)}[1] \dicDirectTranslationCS{(stavební) míchačka} \dicIndirectTranslationCS{(stroj)}
\dicEntry[steypujárn] \dicTerm{steypu··járn} \dicIPA{{s}{\textsubring{d}}{ei}{\textlengthmark}{\textsubring{b}}{\textscy}{j}{au}{r}{\textsubring{d}}{\textsubring{n}}} \dicPos{n}[2] \dicFlx{(‑s, ‑)}[5] \dicDirectTranslationCS{litina}
\dicEntry[steypumót] \textls[15]{\dicTerm{steypu··mót} \dicIPA{{s}{\textsubring{d}}{ei}{\textlengthmark}{\textsubring{b}}{\textscy}{m}{ou}{\textsubring{d}}} \dicPos{n}[2] \dicFlx{(‑s, ‑)}[5] \dicDirectTranslationCS{odlévací\,/\addthin licí forma}}
\dicEntry[steypuvinna] \dicTerm{steypu··vinn|a} \dicIPA{{s}{\textsubring{d}}{ei}{\textlengthmark}{\textsubring{b}}{\textscy}{v}{\textsci}{n}{a}} \dicPos{f}[1] \dicFlx{(‑u)}[5] \dicDirectTranslationCS{betonování}
\dicEntry[steyta] \dicTerm{steyt|a} \dicIPA{{s}{\textsubring{d}}{ei}{\textlengthmark}{\textsubring{d}}{a}} \dicPos{v}[2] \dicFlx{(‑ti, ‑t)}[54] \dicFlx{acc} \textbf{1.} \dicSynonym{mola} \dicDirectTranslationCS{(roz)drtit} \dicExampleIS{steyta pipar} \dicExampleCS{drtit pepř}  \textbf{2.} \dicPhraseIS{steyta hnefann} \dicDirectTranslationCS{zatnout pěst, sevřít pěst}
\dicEntry[steyttur] \dicTerm{steyttur} \dicIPA{{s}{\textsubring{d}}{ei}{h}{\textsubring{d}}{\textscy}{\textsubring{r}}} \dicPos{adj}[1]\dicFlx{}[-10] \textbf{1.} \dicDirectTranslationCS{(roz)drcený (pepř ap.)}  \textbf{2.} \dicDirectTranslationCS{sevřený, zaťatý (pěst ap.)}
\dicEntry[stél] \dicTerm{stél} \dicIPA{{s}{\textsubring{d}}{j}{\textepsilon}{\textlengthmark}{\textsubring{l}}} \dicPos{n}[2] \dicFlx{(‑s, ‑)}[5] \textbf{1.} \dicDirectTranslationCS{ocas} \dicIndirectTranslationCS{(ptačí)}  \textbf{2.} \dicDirectTranslationCS{ocas} \dicIndirectTranslationCS{(letadla)}
\dicEntry[stétt] \dicTerm{stétt} \dicsymFrequent\  \dicIPA{{s}{\textsubring{d}}{j}{\textepsilon}{h}{\textsubring{d}}} \dicPos{f}[7] \dicFlx{(‑ar, ‑ir)}[1] \textbf{1.} \dicSynonym*{þjóðfélagsstétt} \dicDirectTranslationCS{společenská třída\,/\addthin vrstva} \dicExampleIS{Þarna er fólk úr öllum stéttum.} \dicExampleCS{Jsou tam lidé ze všech společenských vrstev.}  \textbf{2.} \dicSynonym{gangbraut} \dicDirectTranslationCS{chodník} \dicExampleIS{leggja stétt framan við húsið} \dicExampleCS{položit chodník před domem}
\dicEntry[stéttabarátta] \dicTerm{stétta··bar·átt|a} \dicIPA{{s}{\textsubring{d}}{j}{\textepsilon}{h}{\textsubring{d}}{a}{\textsubring{b}}{a}{r}{au}{h}{\textsubring{d}}{a}} \dicPos{f}[1] \dicFlx{(‑u)}[5] \dicFieldCat{pol.} \dicDirectTranslationCS{třídní boj}
\dicEntry[stéttamunur] \dicTerm{stétta··mun|ur} \dicIPA{{s}{\textsubring{d}}{j}{\textepsilon}{h}{\textsubring{d}}{a}{m}{\textscy}{n}{\textscy}{\textsubring{r}}} \dicPos{m}[10] \dicFlx{(‑ar)}[12] \dicFieldCat{pol.} \dicDirectTranslationCS{třídní rozdíl}
\dicEntry[stéttarfélag] \dicTerm{stéttar··fé·|lag} \dicIPA{{s}{\textsubring{d}}{j}{\textepsilon}{h}{\textsubring{d}}{a}{\textsubring{r}}{f}{j}{\textepsilon}{l}{a}{x}} \dicPos{n}[2] \dicFlx{(‑lags, ‑lög)}[8] \dicDirectTranslationCS{odborová organizace, odbory}
\dicEntry[stéttarsamband] \dicTerm{stéttar··sam·|band} \dicIPA{{s}{\textsubring{d}}{j}{\textepsilon}{h}{\textsubring{d}}{a}{\textsubring{r}}{s}{a}{m}{\textsubring{b}}{a}{n}{\textsubring{d}}} \dicPos{n}[2] \dicFlx{(‑bands, ‑bönd)}[8] \dicDirectTranslationCS{odbory, odborový svaz}
\dicEntry[stéttaskipting] \dicTerm{stétta··skipt·ing} \dicIPA{{s}{\textsubring{d}}{j}{\textepsilon}{h}{\textsubring{d}}{a}{s}{\r{\textObardotlessj}}{\textsci}{f}{\textsubring{d}}{i}{\ng}{\r{g}}} \dicPos{f}[4] \dicFlx{(‑ar)}[7] \dicFieldCat{pol.} \dicDirectTranslationCS{třídní rozdělení}
\dicEntry[stéttskiptur] \dicTerm{stétt··skiptur} \dicIPA{{s}{\textsubring{d}}{j}{\textepsilon}{h}{\textsubring{d}}{s}{\r{\textObardotlessj}}{\textsci}{f}{\textsubring{d}}{\textscy}{\textsubring{r}}} \dicPos{adj}[1]\dicFlx{}[-10] \dicDirectTranslationCS{rozdělený podle (sociálních) tříd}
\dicEntry[stéttvís] \dicTerm{stétt··vís} \dicIPA{{s}{\textsubring{d}}{j}{\textepsilon}{h}{\textsubring{d}}{v}{i}{s}} \dicPos{adj}[5]\dicFlx{}[-1] \dicDirectTranslationCS{třídně uvědomělý}
\dicEntry[stéttvísi] \dicTerm{stétt··vís|i} \dicIPA{{s}{\textsubring{d}}{j}{\textepsilon}{h}{\textsubring{d}}{v}{i}{s}{\textsci}} \dicPos{f}[3] \dicFlx{(‑i)}[3] \dicDirectTranslationCS{třídní uvědomělost}
\dicEntry[stig] \dicTerm{stig} \dicsymFrequent\  \dicIPA{{s}{\textsubring{d}}{\textsci}{\textlengthmark}{x}} \dicPos{n}[2] \dicFlx{(‑s, ‑)}[5] \textbf{1.} \dicDirectTranslationCS{stupeň, stadium, fáze} \dicExampleIS{á þessu stigi} \dicExampleCS{v~této fázi}  \textbf{2.} \dicSynonym{ætterni} \dicDirectTranslationCS{původ} \dicExampleIS{vera af háum\,/\addthin lágum stigum} \dicExampleCS{být vysokého\,/\addthin nízkého původu}  \textbf{3.} \dicFieldCat{jaz.} \dicDirectTranslationCS{stupeň} \dicExampleIS{efsta stig} \dicExampleCS{třetí stupeň}  \textbf{4.} \dicFieldCat{geol.} \dicDirectTranslationCS{stupeň (zemětřesení)} \dicExampleIS{Jarðskjálftinn var 5 stig á Richter.} \dicExampleCS{Zemětřesení mělo sílu 5 stupňů Richterovy škály.}  \textbf{5.} \dicFieldCat{fyz.} \dicDirectTranslationCS{stupeň (Celsia ap.)} \dicExampleIS{Það er þriggja stiga frost.} \dicExampleCS{Jsou tři stupně pod nulou.}  \textbf{6.} \dicFieldCat{sport.} \dicDirectTranslationCS{bod (v~soutěži ap.)} \dicExampleIS{vinna með fimm stigum} \dicExampleCS{zvítězit o~pět bodů};  \dicPhraseIS{vinna á stigum} \dicDirectTranslationCS{vyhrát na body}
\dicEntry[stigagangur] \dicTerm{stiga··gang|ur} \dicIPA{{s}{\textsubring{d}}{\textsci}{\textlengthmark}{\textbabygamma}{a}{\r{g}}{au}{\ng}{\r{g}}{\textscy}{\textsubring{r}}} \dicPos{m}[6] \dicFlx{(‑s, ‑ar)}[5] \dicDirectTranslationCS{schodiště, schody}
\dicEntry[stigagat] \dicTerm{stiga··|gat} \dicIPA{{s}{\textsubring{d}}{\textsci}{\textlengthmark}{\textbabygamma}{a}{\r{g}}{a}{\textsubring{d}}} \dicPos{n}[2] \dicFlx{(‑gats, ‑göt)}[8] \dicDirectTranslationCS{(schodišťová) šachta}
\dicEntry[stigamaður] \dicTerm{stiga··|maður} \dicIPA{{s}{\textsubring{d}}{\textsci}{\textlengthmark}{\textbabygamma}{a}{m}{a}{ð}{\textscy}{\textsubring{r}}} \dicPos{m}[13] \dicFlx{(‑manns, ‑menn)}[2] \dicSynonym{ræningi} \dicDirectTranslationCS{loupežník, loupežnice, bandita, banditka}
\dicEntry[stigapallur] \dicTerm{stiga··pall|ur} \dicIPA{{s}{\textsubring{d}}{\textsci}{\textlengthmark}{\textbabygamma}{a}{p\smash{\textsuperscript{h}}}{a}{\textsubring{d}}{l}{\textscy}{\textsubring{r}}} \dicPos{m}[6] \dicFlx{(‑s, ‑ar)}[5] \dicDirectTranslationCS{odpočívadlo, podesta}
\dicEntry[stigbeyging] \dicTerm{stig··beyg·ing} \dicIPA{{s}{\textsubring{d}}{\textsci}{\textbabygamma}{\textsubring{b}}{ei}{j}{i}{\ng}{\r{g}}} \dicPos{f}[4] \dicFlx{(‑ar)}[7] \dicLink{stigbreyting}
\dicEntry[stigbeygja] \dicTerm{stig··beyg|ja} \dicIPA{{s}{\textsubring{d}}{\textsci}{\textbabygamma}{\textsubring{b}}{ei}{j}{a}} \dicPos{v}[2] \dicFlx{(‑ði, ‑t)}[89] \dicFlx{dat} \dicFieldCat{jaz.} \dicDirectTranslationCS{stupňovat};  \dicIdiom{stigbeygjast}{ \dicPhraseIS{stigbeygjast}} \dicFlx{refl} \dicFieldCat{jaz.} \dicDirectTranslationCS{stupňovat se}
\dicEntry[stigbreyta] \dicTerm{stig··breyt|a} \dicIPA{{s}{\textsubring{d}}{\textsci}{\textbabygamma}{\textsubring{b}}{r}{ei}{\textsubring{d}}{a}} \dicPos{v}[2] \dicFlx{(‑ti, ‑t)}[54] \dicFlx{acc} \dicFieldCat{jaz.} \dicSynonym{stigbeygja} \dicDirectTranslationCS{stupňovat};  \dicIdiom{stigbreytast}{ \dicPhraseIS{stigbreytast}} \dicFlx{refl} \dicFieldCat{jaz.} \dicSynonym*{stigbeygjast} \dicDirectTranslationCS{stupňovat se}
\dicEntry[stigbreyting] \dicTerm{stig··breyt·ing}\dicTerm{, stigbeyging} \dicIPA{{s}\-{\textsubring{d}}\-{\textsci}\-{\textbabygamma}\-{\textsubring{b}}\-{r}\-{ei}\-{\textsubring{d}}\-{i}\-{\ng}\-{\r{g}}\-} \dicPos{f}[4] \dicFlx{(‑ar)}[7] \dicFieldCat{jaz.} \dicDirectTranslationCS{stupňování, komparace} \dicIndirectTranslationCS{(tvoření druhého a~třetího stupně u~přídavných jmen a~příslovcí)}
\begin{xtolerant}{}{1pt}
\dicEntry[stighækkandi] \dicTerm{stig··hækk·andi}\addthinS\dicIPA{{s}{\textsubring{d}}{\textsci}{\textlengthmark}{x}{h}{a}{i}{h}{\r{g}}{a}{n}{\textsubring{d}}{\textsci}}\addthinS\dicPos{adj}[13]\addthinS\dicFlx{indecl}[1] \dicDirectTranslationCS{vzestupný, rostoucí, progresivní};  \dicPhraseIS{stighækkandi skattar} \dicFieldCat{ekon.} \dicDirectTranslationCS{progresivní daně}
\end{xtolerant}
\dicEntry[stighækkun] \dicTerm{stig··hækk|un} \dicIPA{{s}{\textsubring{d}}{\textsci}{\textlengthmark}{x}{h}{a}{i}{h}{\r{g}}{\textscy}{\textsubring{n}}} \dicPos{f}[7] \dicFlx{(‑unar, ‑anir)}[8] \dicDirectTranslationCS{(progresivní) nárůst}
\dicEntry[stigi] \dicTerm{stig|i\smash{\textsuperscript{1}}} \dicsymFrequent\  \dicIPA{{s}{\textsubring{d}}{i}{j}{\textlengthmark}{\textsci}} \dicPos{m}[1] \dicFlx{(‑a, ‑ar)}[1] \textbf{1.} \dicSynonym*{tröppur} \dicDirectTranslationCS{schody, schodiště} \dicExampleIS{fara upp stigann} \dicExampleCS{jít nahoru po schodech}  \textbf{2.} \dicDirectTranslationCS{žebřík, štafle} \dicExampleIS{átján rima stigi} \dicExampleCS{žebřík o~osmnácti příčkách}  \textbf{3.} \dicDirectTranslationCS{stupnice, škála, žebříček} \dicExampleIS{tónstigi} \dicExampleCS{hudební stupnice}
\dicEntry[stigi] \dicTerm{stigi\smash{\textsuperscript{2}}} \dicIPA{{s}{\textsubring{d}}{i}{j}{\textlengthmark}{\textsci}} \dicPos{v} \dicFlx{con pf sg 1 pers} \dicLink{stíga}
\dicEntry[stigið] \dicTerm{stigið} \dicIPA{{s}{\textsubring{d}}{i}{j}{\textlengthmark}{\textsci}{\texttheta}} \dicPos{v} \dicFlx{supin} \dicLink{stíga}
\dicEntry[stigmunur] \dicTerm{stig··mun|ur} \dicIPA{{s}{\textsubring{d}}{\textsci}{\textbabygamma}{m}{\textscy}{n}{\textscy}{\textsubring{r}}} \dicPos{m}[10] \dicFlx{(‑ar)}[12] \dicLink{stigsmunur}
\dicEntry[stigmögnun] \dicTerm{stig··mögn|un} \dicIPA{{s}{\textsubring{d}}{\textsci}{\textbabygamma}{m}{\oe}{\r{g}}{n}{\textscy}{\textsubring{n}}} \dicPos{f}[7] \dicFlx{(‑unar)}[12] \dicDirectTranslationCS{(vy)stupňování, eskalace}
\dicEntry[stigsmunur] \dicTerm{stigs··mun|ur}\dicTerm{, stigmunur} \dicIPA{{s}\-{\textsubring{d}}\-{\textsci}\-{x}\-{s}\-{m}\-{\textscy}\-{n}\-{\textscy}\-{\textsubring{r}}\-} \dicPos{m}[10] \dicFlx{(‑ar)}[12] \dicDirectTranslationCS{markantní rozdíl}
\dicEntry[stigum] \dicTerm{stigum} \dicIPA{{s}{\textsubring{d}}{\textsci}{\textlengthmark}{\textbabygamma}{\textscy}{\textsubring{m}}} \dicPos{v} \dicFlx{ind pf pl 1 pers} \dicLink{stíga}
\dicEntry[stigur] \dicTerm{stig|ur} \dicIPA{{s}{\textsubring{d}}{\textsci}{\textlengthmark}{\textbabygamma}{\textscy}{\textsubring{r}}} \dicPos{m}[6] \dicFlx{(‑s, ‑ir\,/\addthin ‑ar)}[71] \dicPhraseIS{kanna ókunna stigu} \dicDirectTranslationCS{poznávat nové (nepoznané) cesty};  \dicPhraseIS{stemma stigu við e‑u} \dicLangCat{přen.} \dicDirectTranslationCS{zabránit (čemu), učinit přítrž (čemu)}
\dicEntry[stigveldi] \dicTerm{stig··veldi} \dicIPA{{s}{\textsubring{d}}{\textsci}{\textbabygamma}{v}{\textepsilon}{l}{\textsubring{d}}{\textsci}} \dicPos{n}[2] \dicFlx{(‑s, ‑)}[14] \dicDirectTranslationCS{hierarchie}
\dicEntry[stika] \dicTerm{stik|a\smash{\textsuperscript{1}}} \dicIPA{{s}{\textsubring{d}}{\textsci}{\textlengthmark}{\r{g}}{a}} \dicPos{f}[1] \dicFlx{(‑u, ‑ur)}[19] \dicSynonym{staur} \dicDirectTranslationCS{kůl, sloupek}
\dicEntry[stika] \dicTerm{stik|a\smash{\textsuperscript{2}}} \dicIPA{{s}{\textsubring{d}}{\textsci}{\textlengthmark}{\r{g}}{a}} \dicPos{v}[1] \dicFlx{(‑aði)}[1] \dicDirectTranslationCS{kráčet, jít (rychle)} \dicExampleIS{stika löngum skrefum} \dicExampleCS{jít dlouhým krokem}
\dicEntry[stikilsber] \dicTerm{stikils··ber} \dicIPA{{s}{\textsubring{d}}{\textsci}{\textlengthmark}{\r{\textObardotlessj}}{\textsci}{l}{s}{\textsubring{b}}{\textepsilon}{\textsubring{r}}} \dicPos{n}[2] \dicFlx{(‑s, ‑)}[13] \dicFieldCat{bot.} \dicDirectTranslationCS{angrešt, meruzalka srstka} \textit{(l.~{\textLA{Ribes uva-crispa}})}  \dicsymPhoto\ 
\dicFigure{ds_image_stikilsber_0_1.jpg}{Stikilsber}{Stikilsber - Kornelia und Hartmut Hafele, CC BY-SA 3.0}
\dicEntry[stikkfrí] \dicTerm{stikk··frí} \dicIPA{{s}{\textsubring{d}}{\textsci}{h}{\r{g}}{f}{r}{i}} \dicPos{adj}[13] \dicFlx{indecl}[1] \dicDirectTranslationCS{nenesoucí odpovědnost, nevinný (o~dítěti ap.)}
\dicEntry[stikkorð] \dicTerm{stikk··orð} \dicIPA{{s}{\textsubring{d}}{\textsci}{h}{\r{g}}{\textopeno}{r}{\texttheta}} \dicPos{n}[2] \dicFlx{(‑s, ‑)}[5] \dicFieldCat{lit.} \dicDirectTranslationCS{narážka} \dicIndirectTranslationCS{(konec textu upozorňující, že v~dialogu pokračuje partner)}
\dicEntry[stikkpilla] \dicTerm{stikk··pill|a} \dicIPA{{s}{\textsubring{d}}{\textsci}{h}{\r{g}}{p\smash{\textsuperscript{h}}}{\textsci}{l}{a}} \dicPos{f}[1] \dicFlx{(‑u, ‑ur)}[19] \dicSynonym{stíll} \dicDirectTranslationCS{čípek} \dicIndirectTranslationCS{(léčivo)}
\dicEntry[stikkprufa] \dicTerm{stikk··pruf|a} \dicIPA{{s}{\textsubring{d}}{\textsci}{h}{\r{g}}{p\smash{\textsuperscript{h}}}{r}{\textscy}{v}{a}} \dicPos{f}[1] \dicFlx{(‑u, ‑ur)}[7] \dicDirectTranslationCS{(náhodný) vzorek};  \dicPhraseIS{taka stikkprufu} \dicDirectTranslationCS{odebrat náhodný vzorek}
\dicEntry[stikla] \dicTerm{stikl|a\smash{\textsuperscript{1}}} \dicIPA{{s}{\textsubring{d}}{\textsci}{h}{\r{g}}{l}{a}} \dicPos{f}[1] \dicFlx{(‑u, ‑ur)}[13] \textbf{1.} \dicDirectTranslationCS{nášlapný kámen} \dicIndirectTranslationCS{(zvláště v~pl)}  \textbf{2.} \dicDirectTranslationCS{chůda}
\dicEntry[stikla] \dicTerm{stikl|a\smash{\textsuperscript{2}}} \dicIPA{{s}{\textsubring{d}}{\textsci}{h}{\r{g}}{l}{a}} \dicPos{v}[1] \dicFlx{(‑aði)}[44] \textbf{1.} \dicDirectTranslationCS{přeskočit\,/\addthin přeskakovat z~kamene na kámen} \dicExampleIS{stikla á steinum yfir lækinn} \dicExampleCS{přeskákat potok po kamenech}  \textbf{2.} \dicPhraseIS{stikla á stóru um e‑ð} \dicDirectTranslationCS{zrekapitulovat (co), shrnout (co)}
\dicEntry[stikna] \dicTerm{stikn|a} \dicIPA{{s}{\textsubring{d}}{\textsci}{h}{\r{g}}{n}{a}} \dicPos{v}[1] \dicFlx{(‑aði)}[1] \dicLangCat{přen.} \dicDirectTranslationCS{upéct se, usmažit se (horkem ap.)} \dicExampleIS{Ég er alveg að stikna.} \dicExampleCS{Asi se upeču (horkem).}
\dicEntry[stilkur] \dicTerm{stilk|ur} \dicIPA{{s}{\textsubring{d}}{\textsci}{\textsubring{l}}{\r{g}}{\textscy}{\textsubring{r}}} \dicPos{m}[6] \dicFlx{(‑s, ‑ar)}[24] \dicFieldCat{bot.} \dicDirectTranslationCS{stopka, stonek}
\dicEntry[stilla] \dicTerm{still|a} \dicsymFrequent\  \dicIPA{{s}{\textsubring{d}}{\textsci}{\textsubring{d}}{l}{a}} \dicPos{v}[2] \dicFlx{(‑ti, ‑t)}[63] \dicFlx{acc} \textbf{1.} \dicSynonym{kyrra} \dicDirectTranslationCS{uklidnit, uklidňovat, zklidnit, zklidňovat, utišit, utišovat} \dicExampleIS{stilla skap sitt} \dicExampleCS{uklidnit se}  \textbf{2.} \dicDirectTranslationCS{nařídit, nařizovat, nastavit, nastavovat (hodiny ap.)} \dicExampleIS{stilla klukkuna rétt} \dicExampleCS{nařídit hodiny na správný čas}  \textbf{3.} \dicDirectTranslationCS{(na)ladit (rádio, hudební nástroj ap.)} \dicExampleIS{stilla útvarpið á aðra rás} \dicExampleCS{naladit rádio na jiný kanál};  \dicIdiom{stilla}[í]{ \dicPhraseIS{stilla e‑u í hóf}} \dicDirectTranslationCS{umírnit (co), omezit (co)};  \dicIdiom{stilla}[til]{ \dicPhraseIS{stilla til friðar}} \dicDirectTranslationCS{zjednat pořádek, nastolit mír}; { \dicPhraseIS{stilla því (þannig) til}} \dicDirectTranslationCS{zařídit si to (tak)};  \dicIdiom{stilla}[um]{ \dicPhraseIS{stilla sig ekki um að (gera e‑ð)}} \dicDirectTranslationCS{neubránit se (udělání (čeho)), neodolat (udělání (čeho))};  \dicIdiom{stilla}[upp]{ \dicPhraseIS{stilla e‑u upp}} \dicSynonym*{raða e‑u upp} \dicDirectTranslationCS{uspořádat (co)} \dicIndirectTranslationCS{(uvést do určitého pořádku)}; { \dicPhraseIS{stilla e‑m upp við vegg}} \dicLangCat{přen.} \dicDirectTranslationCS{postavit (koho) ke zdi};  \dicIdiom{stilla}[út]{ \dicPhraseIS{stilla út e‑u}} \dicDirectTranslationCS{vystavit (co), vyložit (co) (zboží ap.)} \dicExampleIS{stilla út varningnum í glugga} \dicExampleCS{vystavit zboží do okna};  \dicIdiom{stillast}{ \dicPhraseIS{stillast}} \dicFlx{refl} \dicSynonym*{verða rólegri} \dicDirectTranslationCS{uklidnit se, uklidňovat se, zklidnit se, zklidňovat se}
\dicEntry[stillilega] \dicTerm{stilli··lega} \dicIPA{{s}{\textsubring{d}}{\textsci}{\textsubring{d}}{l}{\textsci}{l}{\textepsilon}{\textbabygamma}{a}} \dicPos{adv} \dicDirectTranslationCS{klidně, rozvážně}
\dicEntry[stillilegur] \dicTerm{stilli··legur} \dicIPA{{s}{\textsubring{d}}{\textsci}{\textsubring{d}}{l}{\textsci}{l}{\textepsilon}{\textbabygamma}{\textscy}{\textsubring{r}}} \dicPos{adj}[1]\dicFlx{}[-8] \dicDirectTranslationCS{klidný, rozvážný}
\dicEntry[stilling] \dicTerm{still··ing} \dicIPA{{s}{\textsubring{d}}{\textsci}{\textsubring{d}}{l}{i}{\ng}{\r{g}}} \dicPos{f}[4] \dicFlx{(‑ar, ‑ar)}[5] \textbf{1.} \dicSynonym{rósemi} \dicDirectTranslationCS{klid, klidnost, rozvaha} \dicExampleIS{svara með stillingu} \dicExampleCS{odpovědět s~rozvahou}  \textbf{2.} \dicDirectTranslationCS{nastavení, naladění} \dicExampleIS{vélarstilling} \dicExampleCS{nastavení motoru}  \textbf{3.} \dicFieldCat{poč.} \dicDirectTranslationCS{konfigurace}
\dicEntry[stillir] \dicTerm{still|ir} \dicIPA{{s}{\textsubring{d}}{\textsci}{\textsubring{d}}{l}{\textsci}{\textsubring{r}}} \dicPos{m}[7] \dicFlx{(‑is, ‑ar)}[1] \dicDirectTranslationCS{regulátor} \dicIndirectTranslationCS{(zařízení\,/\addthin látka usměrňující nějakou činnost)} \dicExampleIS{ganghraðastillir} \dicExampleCS{regulátor rychlosti motoru}
\dicEntry[stilltur] \dicTerm{stilltur} \dicsymFrequent\  \dicIPA{{s}{\textsubring{d}}{\textsci}{\textsubring{l}}{\textsubring{d}}{\textscy}{\textsubring{r}}} \dicPos{adj}[1]\dicFlx{}[-10] \textbf{1.} \dicSynonym{rólegur} \dicDirectTranslationCS{klidný, rozvážný} \dicExampleIS{Verið þið stilltir!} \dicExampleCS{Buďte klidní!}  \textbf{2.} \dicDirectTranslationCS{nastavený, naladěný}
\begin{xtolerant}{}{1pt}
\dicEntry[stimamjúkur] \dicTerm{stima··\hspace{-0.06667em}|mjúkur}\addthinS\dicIPA{{s}{\textsubring{d}}{\textsci}{\textlengthmark}{m}{a}{m}{j}{u}{\r{g}}{\textscy}{\textsubring{r}}}\addthinS\dicPos{adj}[10]\addthinS\dicFlx{(comp ‑mýkri, sup ‑mýkst\-ur)}[11] \textbf{1.} \dicSynonym{greiðvikinn} \dicDirectTranslationCS{nápomocný, ochotný pomoci}  \textbf{2.} \dicSynonym*{ofurkurteis} \dicDirectTranslationCS{úslužný, servilní}
\end{xtolerant}
\dicEntry[stimpilgjald] \dicTerm{stimpil··|gjald} \dicIPA{{s}{\textsubring{d}}{\textsci}{\textsubring{m}}{\textsubring{b}}{\textsci}{l}{\r{\textObardotlessj}}{a}{l}{\textsubring{d}}} \dicPos{n}[2] \dicFlx{(‑gjalds, ‑gjöld)}[8] \dicDirectTranslationCS{kolkovné}
\dicEntry[stimpilklukka] \dicTerm{stimpil··klukk|a} \dicIPA{{s}{\textsubring{d}}{\textsci}{\textsubring{m}}{\textsubring{b}}{\textsci}{l}{k\smash{\textsuperscript{h}}}{l}{\textscy}{h}{\r{g}}{a}} \dicPos{f}[1] \dicFlx{(‑u, ‑ur)}[13] \dicDirectTranslationCS{píchací hodiny, píchačky}
\dicEntry[stimpill] \dicTerm{stimp|ill} \dicIPA{{s}{\textsubring{d}}{\textsci}{\textsubring{m}}{\textsubring{b}}{\textsci}{\textsubring{d}}{\textsubring{l}}} \dicPos{m}[6] \dicFlx{(‑ils, ‑lar)}[35] \textbf{1.} \dicDirectTranslationCS{razítko} \dicIndirectTranslationCS{(nástroj k~otiskování textu, značek)}  \textbf{2.} \dicDirectTranslationCS{razítko, štempl} \dicIndirectTranslationCS{(otisk razítka)}  \textbf{3.} \dicFieldCat{techn.} \dicDirectTranslationCS{píst}
\dicEntry[stimpilpúði] \dicTerm{stimpil··púð|i} \dicIPA{{s}{\textsubring{d}}{\textsci}{\textsubring{m}}{\textsubring{b}}{\textsci}{l}{p\smash{\textsuperscript{h}}}{u}{ð}{\textsci}} \dicPos{m}[1] \dicFlx{(‑a, ‑ar)}[1] \dicDirectTranslationCS{razítkovací polštářek}
\dicEntry[stimpla] \dicTerm{stimpl|a} \dicIPA{{s}{\textsubring{d}}{\textsci}{\textsubring{m}}{\textsubring{b}}{l}{a}} \dicPos{v}[1] \dicFlx{(‑aði)}[1] \dicFlx{acc} \textbf{1.} \dicDirectTranslationCS{(o)razítkovat, (o)kolkovat} \dicExampleIS{stimpla bréfið} \dicExampleCS{orazítkovat dopis}  \textbf{2.} \dicDirectTranslationCS{označit, ocejchovat} \dicExampleIS{stimpla e‑n sem þjóf} \dicExampleCS{označit (koho) za zloděje};  \dicIdiom{stimpla}[inn]{ \dicPhraseIS{stimpla sig inn}} \dicDirectTranslationCS{označit svůj příchod (na píchacích hodinách)};  \dicIdiom{stimpla}[út]{ \dicPhraseIS{stimpla sig út}} \dicDirectTranslationCS{označit svůj odchod (na píchacích hodinách)};  \dicIdiom{stimplast}{ \dicPhraseIS{stimplast}} \dicFlx{refl} \dicDirectTranslationCS{být orazítkován\,/\addthin okolkován};  \dicIdiom{stimplast}[inn]{ \dicPhraseIS{e‑að stimplast inn}} \dicFlx{refl} \dicDirectTranslationCS{(co) se vryje, (co) utkvívá (v~paměti národa ap.)}
\dicEntry[sting] \dicTerm{sting} \dicIPA{{s}{\textsubring{d}}{\textsci}{\ng}{\r{g}}} \dicPos{v} \dicFlx{ind praes sg 1 pers} \dicLink{stinga}
\dicEntry[stinga] \dicTerm{stinga} \dicsymFrequent\  \dicIPA{{s}{\textsubring{d}}{i}{\ng}{\r{g}}{a}} \dicPos{v}[6] \dicFlx{(sting, stakk, stungum, styngi, stungið)}[39] \dicFlx{acc\,/\addthin dat} \textbf{1.} \dicFlx{acc} \dicDirectTranslationCS{(za)píchnout, píchat, (za)bodnout, bodat} \dicExampleIS{stinga nál í e‑ð} \dicExampleCS{píchnout jehlou do (čeho)};  \dicPhraseIS{stinga sig} \dicDirectTranslationCS{píchnout se, bodnout se} \dicExampleIS{stinga sig á nagla} \dicExampleCS{píchnout se o~hřebík}  \textbf{2.} \dicFlx{acc} \dicDirectTranslationCS{píchnout, píchat, štípnout, štípat (o~vose ap.)} \dicExampleIS{Geitungur stakk mig.} \dicExampleCS{Píchla mě vosa.}  \textbf{3.} \dicFlx{dat} \dicSynonym{smeygja} \dicDirectTranslationCS{(za)strčit, strkat, vsunout, zasunout} \dicExampleIS{stinga bókinni inn í skáp} \dicExampleCS{zastrčit knížku do skříně}  \textbf{4.} \dicDirectTranslationCS{píchat, kousat (svetr ap.)}  \textbf{5.} \dicPhraseIS{e‑að stingur e‑n} \dicDirectTranslationCS{(co) zaráží (koho)};  \dicIdiom{stinga}[að]{ \dicPhraseIS{stinga e‑u að e‑m}} \dicDirectTranslationCS{podstrčit (co komu), podsunout (co komu)} \dicExampleIS{stinga að henni bréfmiða} \dicExampleCS{podstrčit jí papírek};  \dicIdiom{stinga}[af]{ \dicPhraseIS{stinga af}} \dicDirectTranslationCS{vytratit se, vypařit se}; { \dicPhraseIS{stinga e‑ð af}} \dicDirectTranslationCS{zapudit (co), odstrkávat (co) (špatné myšlenky ap.)};  \dicIdiom{stinga}[á]{ \dicPhraseIS{stinga á e‑u}} \dicDirectTranslationCS{propíchnout (co), probodnout (co)}; { \dicPhraseIS{stinga e‑u á sig}} \dicDirectTranslationCS{strčit si (co) do kapsy};  \dicIdiom{stinga}[inn]{ \dicPhraseIS{stinga e‑m inn}} \dicDirectTranslationCS{strčit (koho) do vězení};  \dicIdiom{stinga}[niður]{ \dicPhraseIS{e‑að stingur sér niður}} \dicDirectTranslationCS{(co) se roznáší\,/\addthin rozšiřuje (epidemie ap.)}; { \dicPhraseIS{stinga niður e‑u}} \dicDirectTranslationCS{zasadit (co), vsadit (co), sázet (co)};  \dicIdiom{stinga}[saman]{ \dicPhraseIS{stinga saman nefjum}} \dicDirectTranslationCS{promluvit si mezi čtyřma očima};  \dicIdiom{stinga}[sér]{ \dicPhraseIS{stinga sér}} \dicDirectTranslationCS{skočit šipku, vrhnout se (po hlavě)} \dicExampleIS{stinga sér í laugina} \dicExampleCS{skočit šipku do bazénu};  \dicIdiom{stinga}[undan]{ \dicPhraseIS{stinga e‑u undan}} \dicDirectTranslationCS{zpronevěřit (co), strkat si (co) do vlastní kapsy} \dicExampleIS{stinga peningunum undan} \dicExampleCS{zpronevěřit peníze};  \dicIdiom{stinga}[upp]{ \dicPhraseIS{stinga e‑ð upp}} \dicDirectTranslationCS{(z)rýt (co), překopat (co)} \dicExampleIS{stinga upp garð} \dicExampleCS{zrýt zahradu};  \dicIdiom{stinga}[upp á]{ \dicPhraseIS{stinga upp á e‑u}} \dicDirectTranslationCS{nadhodit (co), uvést (co) na přetřes (návrh ap.)};  \dicIdiom{stinga}[upp í]{ \dicPhraseIS{stinga e‑ð upp í sig}} \dicDirectTranslationCS{strkat si (co) do pusy};  \dicIdiom{stinga}[við]{ \dicPhraseIS{stinga við}} \dicSynonym{haltra} \dicDirectTranslationCS{kulhat, napadat na nohu};  \dicIdiom{stingast}{ \dicPhraseIS{stingast}} \dicFlx{refl} \dicDirectTranslationCS{zabodnout se, zapíchnout se}
\dicEntry[stingsög] \dicTerm{sting··|sög} \dicIPA{{s}{\textsubring{d}}{i}{\ng}{\r{g}}{s}{\oe}{x}} \dicPos{f}[7] \dicFlx{(‑sagar, ‑sagir)}[16] \dicDirectTranslationCS{přímočará pila, lupenka, lupenková pilka}
\dicEntry[stingur] \dicTerm{sting|ur} \dicIPA{{s}{\textsubring{d}}{i}{\ng}{\r{g}}{\textscy}{\textsubring{r}}} \dicPos{m}[9] \dicFlx{(‑s, ‑ir)}[15] \textbf{1.} \dicSynonym{broddur} \dicDirectTranslationCS{bodec, osten}  \textbf{2.} \dicSynonym{verkur} \dicDirectTranslationCS{bodnutí, píchnutí (v~boku ap.)} \dicExampleIS{fá sting í brjóstið} \dicExampleCS{ucítit píchnutí v~prsou}
\dicEntry[stinnastör] \dicTerm{stinna··|stör} \dicIPA{{s}{\textsubring{d}}{\textsci}{n}{\textlengthmark}{a}{s}{\textsubring{d}}{\oe}{\textsubring{r}}} \dicPos{f}[7] \dicFlx{(‑starar, ‑starir)}[16] \dicFieldCat{bot.} \dicDirectTranslationCS{ostřice Bigelowova} \textit{(l.~{\textLA{Carex bigelowii}})}  \dicsymPhoto\ 
\dicFigure{32487.jpg}{Stinnastör}{Stinnastör - Blanka Brandová, Biolib, Copyright/CC-BY-NC}
\dicEntry[stinning] \dicTerm{stinn··ing} \dicIPA{{s}{\textsubring{d}}{\textsci}{n}{\textlengthmark}{i}{\ng}{\r{g}}} \dicPos{f}[4] \dicFlx{(‑ar, ‑ar)}[5] \dicDirectTranslationCS{ztopoření, erekce}
\dicEntry[stinningarlyf] \dicTerm{stinningar··lyf} \dicIPA{{s}{\textsubring{d}}{\textsci}{n}{\textlengthmark}{i}{\ng}{\r{g}}{a}{r}{l}{\textsci}{f}} \dicPos{n}[2] \dicFlx{(‑s, ‑)}[13] \dicFieldCat{med.} \dicDirectTranslationCS{léky na problémy s~erekcí}
\dicEntry[stinningsgola] \dicTerm{stinnings··gol|a} \dicIPA{{s}{\textsubring{d}}{\textsci}{n}{\textlengthmark}{i}{\ng}{s}{\r{g}}{\textopeno}{l}{a}} \dicPos{f}[1] \dicFlx{(‑u)}[5] \dicSynonym*{allhvöss gola} \dicDirectTranslationCS{dosti čerstvý vítr} \dicIndirectTranslationCS{(Beaufort 4) (5,5--7,9~m/s)}
\dicEntry[stinningskaldi] \dicTerm{stinnings··kald|i} \dicIPA{{s}{\textsubring{d}}{\textsci}{n}{\textlengthmark}{i}{\ng}{s}{k\smash{\textsuperscript{h}}}{a}{l}{\textsubring{d}}{\textsci}} \dicPos{m}[1] \dicFlx{(‑a)}[3] \dicDirectTranslationCS{silný vítr} \dicIndirectTranslationCS{(Beaufort 6) (10,8--13,8~m/s)}
\dicEntry[stinnur] \dicTerm{stinnur} \dicIPA{{s}{\textsubring{d}}{\textsci}{n}{\textlengthmark}{\textscy}{\textsubring{r}}} \dicPos{adj}[1]\dicFlx{}[-1] \dicSynonym{stífur} \dicDirectTranslationCS{tuhý, ztuhlý} \dicExampleIS{stinnur vöðvi} \dicExampleCS{ztuhlý sval}
\dicEntry[stirðbusalegur] \dicTerm{stirð·busa··legur} \dicIPA{{s}{\textsubring{d}}{\textsci}{r}{ð}{\textsubring{b}}{\textscy}{s}{a}{l}{\textepsilon}{\textbabygamma}{\textscy}{\textsubring{r}}} \dicPos{adj}[1]\dicFlx{}[-8] \dicDirectTranslationCS{neohrabaný, nemotorný}
\dicEntry[stirðlega] \dicTerm{stirð··lega} \dicIPA{{s}{\textsubring{d}}{\textsci}{r}{ð}{l}{\textepsilon}{\textbabygamma}{a}} \dicPos{adv} \dicDirectTranslationCS{strnule, těžkopádně}
\dicEntry[stirðlegur] \dicTerm{stirð··legur} \dicIPA{{s}{\textsubring{d}}{\textsci}{r}{ð}{l}{\textepsilon}{\textbabygamma}{\textscy}{\textsubring{r}}} \dicPos{adj}[1]\dicFlx{}[-8] \dicDirectTranslationCS{strnulý, zkostnatělý, těžkopádný, nepružný}
\dicEntry[stirðleiki] \dicTerm{stirð··leik|i} \dicIPA{{s}{\textsubring{d}}{\textsci}{r}{ð}{l}{ei}{\r{\textObardotlessj}}{\textsci}} \dicPos{m}[1] \dicFlx{(‑a)}[3] \textbf{1.} \dicDirectTranslationCS{ztuhlost, strnulost (sva\-lů ap.)}  \textbf{2.} \dicDirectTranslationCS{nepružnost, neohebnost, strnulost, rigidnost}
\dicEntry[stirðna] \dicTerm{stirðn|a} \dicIPA{{s}{\textsubring{d}}{\textsci}{r}{\textsubring{d}}{n}{a}} \dicPos{v}[1] \dicFlx{(‑aði)}[1] \dicSynonym*{verða stífur} \dicDirectTranslationCS{ztuhnout, (u)strnout, stát se nehybným} \dicExampleIS{stirðnað lík} \dicExampleCS{ztuhlé mrtvé tělo};  \dicIdiom{stirðna}[í]{ \dicPhraseIS{stirðna í e‑u}} \dicDirectTranslationCS{ustrnout v~(čem) (ve výuce ap.)}
\dicEntry[stirðnaður] \dicTerm{stirðn··|aður} \dicIPA{{s}{\textsubring{d}}{\textsci}{r}{\textsubring{d}}{n}{a}{ð}{\textscy}{\textsubring{r}}} \dicPos{adj}[3] \dicFlx{(f ‑uð)}[3] \dicDirectTranslationCS{strnulý, ztuhlý (hrůzou ap.)}
\dicEntry[stirður] \dicTerm{stirður} \dicIPA{{s}{\textsubring{d}}{\textsci}{r}{ð}{\textscy}{\textsubring{r}}} \dicPos{adj}[2]\dicFlx{}[-1] \textbf{1.} \dicSynonym{stífur} \dicDirectTranslationCS{ztuhlý, strnulý};  \dicPhraseIS{vera stirður í hreyfingum} \dicDirectTranslationCS{mít strnulé pohyby}  \textbf{2.} \dicSynonym{örðugur} \dicDirectTranslationCS{svízelný, obtížný (vztah ap.)};  \dicPhraseIS{e‑m er stirt um mál} \dicFlx{impers} \dicDirectTranslationCS{(komu) se špatně mluví};  \dicPhraseIS{það er stirt á milli e‑rra} \dicFlx{impers} \dicDirectTranslationCS{mezi (kým) jsou napjaté vztahy}  \textbf{3.} \dicSynonym{önugur} \dicDirectTranslationCS{nevrlý, nerudný, mrzutý};  \dicPhraseIS{vera stirður í skapi} \dicDirectTranslationCS{být nevrlý}  \textbf{4.} \dicDirectTranslationCS{tuhý, těžkopádný (styl ap.)}
\dicEntry[stirfinn] \dicTerm{stirfinn} \dicIPA{{s}{\textsubring{d}}{\textsci}{r}{v}{\textsci}{\textsubring{n}}} \dicPos{adj}[6]\dicFlx{}[-2] \dicSynonym{önugur} \dicDirectTranslationCS{mrzutý, zatrpklý, rozmrzelý} \dicExampleIS{stirfið bros} \dicExampleCS{mrzutý úsměv}
\dicEntry[stirfni] \dicTerm{stirfn|i} \dicIPA{{s}{\textsubring{d}}{\textsci}{r}{\textsubring{d}}{n}{\textsci}} \dicPos{f}[3] \dicFlx{(‑i)}[3] \dicSynonym{þvermóðska} \dicDirectTranslationCS{vzpurnost, svéhlavost}
\dicEntry[stía] \dicTerm{stí|a\smash{\textsuperscript{1}}} \dicIPA{{s}{\textsubring{d}}{i}{\textlengthmark}{j}{a}} \dicPos{f}[1] \dicFlx{(‑u, ‑ur)}[7] \textbf{1.} \dicDirectTranslationCS{chlév, chlívek, kotec, ohrádka} \dicExampleIS{reka féð inn í stíuna} \dicExampleCS{hnát ovce do ohrady}  \textbf{2.} \dicLangCat{přen.} \dicSynonym*{sóðalegur staður} \dicDirectTranslationCS{chlívek} \dicIndirectTranslationCS{(špinavé místo)}
\dicEntry[stía] \dicTerm{stí|a\smash{\textsuperscript{2}}} \dicIPA{{s}{\textsubring{d}}{i}{\textlengthmark}{j}{a}} \dicPos{v}[1] \dicFlx{(‑aði)}[1] \dicFlx{dat} \textbf{1.} \dicDirectTranslationCS{oddělit, oddělovat, separovat (koně, ovce ap.)} \dicExampleIS{stía lömbunum} \dicExampleCS{oddělit jehňata}  \textbf{2.} \dicSynonym{aðskilja} \dicDirectTranslationCS{odloučit, odlučovat, oddělit, oddělovat, separovat, (od)izolovat};  \dicPhraseIS{stía e‑jum í sundur} \dicDirectTranslationCS{separovat (koho) od sebe, odloučit (koho) od sebe} \dicExampleIS{stía fólki í sundur} \dicExampleCS{separovat lidi od sebe}
\dicEntry[stífkrampi] \dicTerm{stíf··kramp|i} \dicIPA{{s}{\textsubring{d}}{i}{v}{k\smash{\textsuperscript{h}}}{r}{a}{\textsubring{m}}{\textsubring{b}}{\textsci}} \dicPos{m}[1] \dicFlx{(‑a)}[3] \dicFieldCat{med.} \dicDirectTranslationCS{tetanus}
\dicEntry[stífla] \dicTerm{stífl|a\smash{\textsuperscript{1}}} \dicIPA{{s}{\textsubring{d}}{i}{\textsubring{b}}{l}{a}} \dicPos{f}[1] \dicFlx{(‑u, ‑ur)}[19] \textbf{1.} \dicSynonym{stíflugarður} \dicDirectTranslationCS{hráz, přehrada, zdrž} \dicExampleIS{hlaða stíflu} \dicExampleCS{stavět hráz}  \textbf{2.} \dicSynonym{teppa\smash{\textsuperscript{1}}} \dicDirectTranslationCS{ucpání, zneprůchodnění, zanesení} \dicExampleIS{stífla í nefi} \dicExampleCS{ucpání nosu}
\dicEntry[stífla] \dicTerm{stífl|a\smash{\textsuperscript{2}}} \dicIPA{{s}{\textsubring{d}}{i}{\textsubring{b}}{l}{a}} \dicPos{v}[1] \dicFlx{(‑aði)}[1] \dicFlx{acc} \textbf{1.} \dicSynonym*{gera stíflu í} \dicDirectTranslationCS{přehradit, přehrazovat, zahradit, zahrazovat} \dicExampleIS{stífla lækinn} \dicExampleCS{přehradit potok}  \textbf{2.} \dicSynonym{teppa\smash{\textsuperscript{2}}} \dicDirectTranslationCS{ucpat, ucpávat, zneprůchodnit, zneprůchodňovat, zanést, zanášet} \dicExampleIS{stífla rör} \dicExampleCS{ucpat trubku};  \dicIdiom{stíflast}{ \dicPhraseIS{stíflast}} \dicFlx{refl} \dicDirectTranslationCS{ucpat se, zneprůchodnit se, zanést se}
\dicEntry[stífleiki] \dicTerm{stíf··leik|i} \dicIPA{{s}{\textsubring{d}}{i}{v}{l}{ei}{\r{\textObardotlessj}}{\textsci}} \dicPos{m}[1] \dicFlx{(‑a)}[3] \dicDirectTranslationCS{tuhost, ztuhlost}
\dicEntry[stíflugarður] \dicTerm{stíflu··garð|ur} \dicIPA{{s}{\textsubring{d}}{i}{\textsubring{b}}{l}{\textscy}{\r{g}}{a}{r}{ð}{\textscy}{\textsubring{r}}} \dicPos{m}[6] \dicFlx{(‑s, ‑ar)}[5] \dicDirectTranslationCS{hráz, přehrada, zdrž}
\dicEntry[stíflyndur] \dicTerm{stíf··lyndur} \dicIPA{{s}{\textsubring{d}}{i}{v}{l}{\textsci}{n}{\textsubring{d}}{\textscy}{\textsubring{r}}} \dicPos{adj}[2]\dicFlx{}[-14] \dicSynonym{stífur} \dicDirectTranslationCS{zatvrzelý, umíněný}
\dicEntry[stífni] \dicTerm{stífn|i} \dicIPA{{s}{\textsubring{d}}{i}{\textsubring{b}}{n}{\textsci}} \dicPos{f}[3] \dicFlx{(‑i)}[3] \dicSynonym{þrjóska} \dicDirectTranslationCS{zatvrzelost, neústupnost, neoblomnost}
\dicEntry[stífur] \dicTerm{stífur} \dicIPA{{s}{\textsubring{d}}{i}{\textlengthmark}{v}{\textscy}{\textsubring{r}}} \dicPos{adj}[1]\dicFlx{}[-1] \textbf{1.} \dicSynonym{stinnur} \dicDirectTranslationCS{ztuhlý, ztvrdlý}  \textbf{2.} \dicSynonym{þrár} \dicDirectTranslationCS{zatvrzelý, tvrdošíjný}  \textbf{3.} \dicSynonym{strangur} \dicDirectTranslationCS{přísný, striktní} \dicExampleIS{stíf regla} \dicExampleCS{přísné pravidlo}  \textbf{4.} \dicSynonym{hvass} \dicDirectTranslationCS{ostrý (vítr ap.)} \dicExampleIS{stífur vindur} \dicExampleCS{ostrý vítr}  \textbf{5.} \dicSynonym{harður} \dicDirectTranslationCS{tuhý, tvrdý} \dicExampleIS{stífur pappír} \dicExampleCS{tuhý papír}
\dicEntry[stífþeyta] \dicTerm{stíf··þeyt|a} \dicIPA{{s}{\textsubring{d}}{i}{v}{\texttheta}{ei}{\textsubring{d}}{a}} \dicPos{v}[2] \dicFlx{(‑ti, ‑t)}[54] \dicFlx{acc} \dicDirectTranslationCS{(na)šlehat (bílek ap.)} \dicExampleIS{stífþeyta eggjahvítur} \dicExampleCS{našlehat bílky}
\dicEntry[stíg] \dicTerm{stíg} \dicIPA{{s}{\textsubring{d}}{i}{\textlengthmark}{x}} \dicPos{v} \dicFlx{ind praes sg 1 pers} \dicLink{stíga}
\dicEntry[stíga] \dicTerm{stíga} \dicsymFrequent\  \dicIPA{{s}{\textsubring{d}}{i}{\textlengthmark}{\textbabygamma}{a}} \dicPos{v}[6] \dicFlx{(stíg, steig, stigum, stigi, stigið)}[73] \dicFlx{acc} \textbf{1.} \dicSynonym*{taka skref} \dicDirectTranslationCS{udělat krok, kráčet} \dicExampleIS{stíga eitt skref áfram} \dicExampleCS{udělat jeden krok dopředu}  \textbf{2.} \dicSynonym{troða} \dicDirectTranslationCS{stoupnout, stoupat, šlápnout, šlapat} \dicExampleIS{stíga á bremsuna} \dicExampleCS{šlápnout na brzdu};  \dicPhraseIS{stíga ofan á e‑ð} \dicDirectTranslationCS{šlápnout na (co), pošlapat (co)} \dicExampleIS{stíga ofan á blómið} \dicExampleCS{šlápnout na květinu};  \dicPhraseIS{stíga dans} \dicSynonym{dansa} \dicDirectTranslationCS{tancovat} \dicExampleIS{stíga dans við hana} \dicExampleCS{tancovat s~ní}  \textbf{3.} \dicDirectTranslationCS{šlapat (na kole ap.)} \dicExampleIS{stíga reiðhjól} \dicExampleCS{šlapat na kole}  \textbf{4.} \dicSynonym{vaxa} \dicDirectTranslationCS{(vy)stoupat, (vy)růst} \dicExampleIS{Hitinn steig.} \dicExampleCS{Teplota stoupla.};  \dicPhraseIS{stíga í verði} \dicDirectTranslationCS{stoupat v~ceně}  \textbf{5.} \dicFieldCat{let.} \dicDirectTranslationCS{(vy)stoupat (letadlo ap.)}  \textbf{6.} \dicDirectTranslationCS{krokovat, počítat kroky};  \dicIdiom{stíga}[af]{ \dicPhraseIS{stíga af baki}} \dicDirectTranslationCS{sesednout z~koně};  \dicIdiom{stíga}[á]{ \dicPhraseIS{stíga á bak}} \dicDirectTranslationCS{nasednout na koně};  \dicIdiom{stíga}[niður]{ \dicPhraseIS{stíga  niður}} \dicDirectTranslationCS{sestoupit, sestupovat};  \dicIdiom{stíga}[til]{ \dicPhraseIS{e‑m stígur e‑að til höfuðs}} \dicFlx{impers} {\textbf{a.}} \dicLangCat{přen.} \dicDirectTranslationCS{(komu co) stoupá do hlavy (alkohol ap.)};  {\textbf{b.}} \dicLangCat{přen.} \dicDirectTranslationCS{(komu co) stoupá do hlavy (moc ap.)};  \dicIdiom{stíga}[upp í]{ \dicPhraseIS{stíga upp í e‑ð}} \dicDirectTranslationCS{nastoupit do (čeho)} \dicExampleIS{stíga upp í bílinn} \dicExampleCS{nastoupit do auta};  \dicIdiom{stíga}[út úr]{ \dicPhraseIS{stíga út úr e‑u}} \dicDirectTranslationCS{vystoupit z~(čeho)} \dicExampleIS{stíga út úr bílnum} \dicExampleCS{vystoupit z~auta};  \dicIdiom{stíga}[við]{ \dicPhraseIS{stíga við e‑n}} \dicDirectTranslationCS{tancovat s~(kým), přešlapovat s~(kým) na místě (při tanci)};  \dicIdiom{stíga}{ \dicPhraseIS{geta ekki stigið í fótinn}} \dicDirectTranslationCS{nemoci se postavit na nohy}
\dicEntry[stígandi] \dicTerm{stíg··|andi\smash{\textsuperscript{1}}} \dicIPA{{s}{\textsubring{d}}{i}{\textlengthmark}{\textbabygamma}{a}{n}{\textsubring{d}}{\textsci}} \dicPos{m}[2] \dicFlx{(‑anda)}[2] \dicDirectTranslationCS{přešlapování}
\dicEntry[stígandi] \dicTerm{stíg··and|i\smash{\textsuperscript{2}}} \dicIPA{{s}{\textsubring{d}}{i}{\textlengthmark}{\textbabygamma}{a}{n}{\textsubring{d}}{\textsci}} \dicPos{f}[3] \dicFlx{(‑i)}[3] \dicDirectTranslationCS{vyvrcholení (v~básni ap.)}
\dicEntry[stígandi] \dicTerm{stíg··andi\smash{\textsuperscript{3}}} \dicIPA{{s}{\textsubring{d}}{i}{\textlengthmark}{\textbabygamma}{a}{n}{\textsubring{d}}{\textsci}} \dicPos{adj}[13] \dicFlx{indecl}[1] \dicDirectTranslationCS{vzrůstající, zvyšující se (napětí ap.)}
\dicEntry[stígur] \dicTerm{stíg|ur} \dicsymFrequent\  \dicIPA{{s}{\textsubring{d}}{i}{\textlengthmark}{\textbabygamma}{\textscy}{\textsubring{r}}} \dicPos{m}[6] \dicFlx{(‑s, ‑ar\,/\addthin ‑ir)}[70] \dicDirectTranslationCS{cesta, stezka, pěšina} \dicExampleIS{ganga stíginn upp brekkuna} \dicExampleCS{jít pěšinou do kopce}
\dicEntry[stígvél] \dicTerm{stíg··vél} \dicIPA{{s}{\textsubring{d}}{i}{\textbabygamma}{v}{j}{\textepsilon}{\textsubring{l}}} \dicPos{n}[2] \dicFlx{(‑s, ‑)}[5] \dicDirectTranslationCS{holínka, gumovka, gumák}
\dicEntry[stíla] \dicTerm{stíl|a} \dicIPA{{s}{\textsubring{d}}{i}{\textlengthmark}{l}{a}} \dicPos{v}[1] \dicFlx{(‑aði)}[1] \dicFlx{acc} \textbf{1.} \dicDirectTranslationCS{adresovat, určit} \dicExampleIS{stíla e‑ð á e‑n} \dicExampleCS{adresovat (co komu)}  \textbf{2.} \dicSynonym{semja} \dicDirectTranslationCS{sepsat, sepisovat, stylizovat, (z)formulovat (text ap.)} \dicExampleIS{stíla bréf} \dicExampleCS{formulovat dopis}  \textbf{3.} \dicPhraseIS{stíla upp á e‑ð} \dicDirectTranslationCS{zaměřit se na (co), orientovat se na (co)}
\dicEntry[stílabók] \dicTerm{stíla··|bók} \dicIPA{{s}{\textsubring{d}}{i}{\textlengthmark}{l}{a}{\textsubring{b}}{ou}{\r{g}}} \dicPos{f}[8] \dicFlx{(‑bókar, ‑bækur)}[5] \dicDirectTranslationCS{zápisník, notes, (školní) sešit} \dicExampleIS{Börnin skrifa í stílabækurnar.} \dicExampleCS{Děti píší do sešitů.}
\begin{xtolerant}{}{6pt}
\dicEntry[stílbragð] \dicTerm{stíl··|bragð} \dicIPA{{s}{\textsubring{d}}{i}{l}{\textsubring{b}}{r}{a}{\textbabygamma}{\texttheta}} \dicPos{n}[2] \dicFlx{(‑bragðs, ‑brögð)}[8] \dicFieldCat{lit.\,/\addthin jaz.} \dicDirectTranslationCS{stylistický prostředek, stylistická\,/\addthinS básnická\,/\addthinS řečnická fi\-gu\-ra}
\end{xtolerant}
\dicEntry[stílfræði] \dicTerm{stíl··fræð|i} \dicIPA{{s}{\textsubring{d}}{i}{l}{f}{r}{a}{i}{ð}{\textsci}} \dicPos{f}[3] \dicFlx{(‑i)}[3] \dicFieldCat{jaz.} \dicDirectTranslationCS{stylistika}
\dicEntry[stílfæra] \dicTerm{stíl··fær|a} \dicIPA{{s}{\textsubring{d}}{i}{l}{f}{a}{i}{r}{a}} \dicPos{v}[2] \dicFlx{(‑ði, ‑t)}[105] \dicFlx{acc} \dicDirectTranslationCS{stylizovat}
\dicEntry[stílfærsla] \dicTerm{stíl··færsl|a} \dicIPA{{s}{\textsubring{d}}{i}{l}{f}{a}{i}{\textsubring{r}}{s}{\textsubring{d}}{l}{a}} \dicPos{f}[1] \dicFlx{(‑u)}[5] \dicDirectTranslationCS{stylizace, stylizování}
\dicEntry[stílhreinn] \dicTerm{stíl··hreinn} \dicIPA{{s}{\textsubring{d}}{i}{l}{\textsubring{r}}{ei}{\textsubring{d}}{\textsubring{n}}} \dicPos{adj}[7]\dicFlx{}[-1] \dicDirectTranslationCS{stylisticky\,/\addthin stylově čistý}
\dicEntry[stílisti] \dicTerm{stíl··ist|i} \dicIPA{{s}{\textsubring{d}}{i}{\textlengthmark}{l}{\textsci}{s}{\textsubring{d}}{\textsci}} \dicPos{m}[1] \dicFlx{(‑a, ‑ar)}[1] \textbf{1.} \dicDirectTranslationCS{stylista, stylistka} \dicIndirectTranslationCS{(člověk schopný obratně sestavovat text)}  \textbf{2.} \dicDirectTranslationCS{stylista, stylistka} \dicIndirectTranslationCS{(člověk profesionálně se zabývající úpravou vzhledu)}
\dicEntry[stíll] \dicTerm{stíl|l} \dicsymFrequent\  \dicIPA{{s}{\textsubring{d}}{i}{\textsubring{d}}{\textsubring{l}}} \dicPos{m}[6] \dicFlx{(‑s, ‑ar)}[48] \textbf{1.} \dicSynonym*{svipmót} \dicDirectTranslationCS{styl, sloh} \dicExampleIS{gotneskur stíll} \dicExampleCS{gotický sloh}  \textbf{2.} \dicSynonym{ritgerð} \dicDirectTranslationCS{cvičení, (školní) kompozice}  \textbf{3.} \dicSynonym*{lyfjastautur} \dicDirectTranslationCS{čípek} \dicIndirectTranslationCS{(léčivo)};  \dicIdiom{stíll}{ \dicPhraseIS{e‑að er í stíl við e‑ð}} \dicDirectTranslationCS{(co) je ve stejném stylu jako (co), (co) se hodí k~(čemu)}; { \dicPhraseIS{í stórum stíl}} \dicFlx{adv} \dicLangCat{přen.} \dicDirectTranslationCS{ve velkém stylu};  \dicPhraseIS{það er stíll yfir e‑n} \dicFlx{impers} \dicLangCat{přen.} \dicDirectTranslationCS{(kdo) má styl}
\dicEntry[stílsnilld] \dicTerm{stíl··snilld} \dicIPA{{s}{\textsubring{d}}{i}{l}{s}{\textsubring{d}}{n}{\textsci}{l}{\textsubring{d}}} \dicPos{f}[7] \dicFlx{(‑ar)}[3] \dicDirectTranslationCS{brilantní styl (psaní ap.)}
\dicEntry[stíma] \dicTerm{stím|a} \dicIPA{{s}{\textsubring{d}}{i}{\textlengthmark}{m}{a}} \dicPos{v}[3] \dicFlx{(‑aði)\,/\addthin (‑di, ‑t)}[8] \textbf{1.} \dicDirectTranslationCS{(vy)plout (rychle), plavit se (rychle)}  \textbf{2.} \dicDirectTranslationCS{namířit, vydat se}
\dicEntry[stímabrak] \dicTerm{stíma··brak} \dicIPA{{s}{\textsubring{d}}{i}{\textlengthmark}{m}{a}{\textsubring{b}}{r}{a}{\r{g}}} \dicPos{n}[2] \dicFlx{(‑s)}[2] \dicDirectTranslationCS{trable, potíže};  \dicPhraseIS{standa í stímabraki til að (gera e‑ð)} \dicDirectTranslationCS{mít nesnáze s~(děláním (čeho))}
\dicEntry[stírur] \dicTerm{stírur} \dicIPA{{s}{\textsubring{d}}{i}{\textlengthmark}{r}{\textscy}{\textsubring{r}}} \dicPos{f}[1] \dicFlx{pl}[12] \dicDirectTranslationCS{ospalky};  \dicPhraseIS{vera með stírur(nar) í augunum} \dicDirectTranslationCS{mít ospalky v~očích}
\dicEntry[stjaka] \dicTerm{stjak|a} \dicIPA{{s}{\textsubring{d}}{j}{a}{\textlengthmark}{\r{g}}{a}} \dicPos{v}[1] \dicFlx{(‑aði)}[13] \dicFlx{dat} \dicSynonym{ýta\smash{\textsuperscript{2}}} \dicDirectTranslationCS{(od)strčit, postrčit, postrkávat, strkat, odstrkovat} \dicExampleIS{Hún stjakaði manninum frá sér.} \dicExampleCS{Odstrčila od sebe toho člověka.};  \dicPhraseIS{stjaka við e‑m} \dicDirectTranslationCS{odstrčit (koho), strkat do (koho)}
\dicEntry[stjaki] \dicTerm{stjak|i} \dicIPA{{s}{\textsubring{d}}{j}{a}{\textlengthmark}{\r{\textObardotlessj}}{\textsci}} \dicPos{m}[1] \dicFlx{(‑a, ‑ar)}[8] \textbf{1.} \dicSynonym{stöng} \dicDirectTranslationCS{bidlo, sochor}  \textbf{2.} \dicSynonym{kertastjaki} \dicDirectTranslationCS{svícen}
\dicEntry[stjana] \dicTerm{stjan|a} \dicIPA{{s}{\textsubring{d}}{j}{a}{\textlengthmark}{n}{a}} \dicPos{v}[1] \dicFlx{(‑aði)}[13] \dicDirectTranslationCS{obskakovat, hýčkat};  \dicPhraseIS{stjana við e‑n} \dicDirectTranslationCS{obskakovat (koho), obsluhovat (koho)}
\dicEntry[stjarfi] \dicTerm{stjarf|i} \dicIPA{{s}{\textsubring{d}}{j}{a}{r}{v}{\textsci}} \dicPos{m}[1] \dicFlx{(‑a)}[3] \textbf{1.} \dicSynonym{stirðleiki} \dicDirectTranslationCS{(u)strnutí, ztuh\-nu\-tí}  \textbf{2.} \dicFieldCat{med.} \dicSynonym{stífkrampi} \dicDirectTranslationCS{tetanus}
\dicEntry[stjarfur] \dicTerm{stjarfur} \dicIPA{{s}{\textsubring{d}}{j}{a}{r}{v}{\textscy}{\textsubring{r}}} \dicPos{adj}[1] \dicFlx{(f stjörf)}[2] \dicSynonym{stífur} \dicDirectTranslationCS{ustrnulý, ochromený, paralyzovaný} \dicExampleIS{stjarfur af skelfingu} \dicExampleCS{ochromený hrůzou}
\dicEntry[stjarna] \dicTerm{stjarna} \dicsymFrequent\  \dicIPA{{s}{\textsubring{d}}{j}{a}{r}{\textsubring{d}}{n}{a}} \dicPos{f}[1] \dicFlx{(stjörnu, stjörnur)}[8] \textbf{1.} \dicFieldCat{astro.} \dicDirectTranslationCS{hvězda} \dicExampleIS{óteljandi stjörnur á himninum} \dicExampleCS{nesčetné množství hvězd na obloze}  \textbf{2.} \dicDirectTranslationCS{hvězda, hvězdice, hvězdička} \dicExampleIS{stjarna á jólatré} \dicExampleCS{hvězdička na vánočním stromku}  \textbf{3.} \dicSynonym*{fræg persóna} \dicDirectTranslationCS{hvězda, celebrita} \dicIndirectTranslationCS{(známá osobnost)};  \dicPhraseIS{vera rísandi stjarna} \dicLangCat{přen.} \dicDirectTranslationCS{být vycházející hvězda}
\dicEntry[stjarneðlisfræði] \dicTerm{stjarn··eðlis·fræð|i} \dicIPA{{s}{\textsubring{d}}{j}{a}{r}{\textsubring{d}}{n}{\textepsilon}{ð}{l}{\textsci}{s}{f}{r}{a}{i}{ð}{\textsci}} \dicPos{f}[3] \dicFlx{(‑i)}[3] \dicFieldCat{astro.} \dicDirectTranslationCS{astrofyzika}
\dicEntry[stjarneðlisfræðingur] \dicTerm{stjarn·eðlis·fræð··ing|ur} \dicIPA{{s}\-{\textsubring{d}}\-{j}\-{a}\-{r}\-{\textsubring{d}}\-{n}\-{\textepsilon}\-{ð}\-{l}\-{\textsci}\-{s}\-{f}\-{r}\-{a}\-{i}\-{ð}\-{i}\-{\ng}\-{\r{g}}\-{\textscy}\-{\textsubring{r}}\-} \dicPos{m}[6] \dicFlx{(‑s, ‑ar)}[8] \dicDirectTranslationCS{astrofyzik, astrofyzička}
\dicEntry[stjarnfræðilegur] \dicTerm{stjarn·fræði··legur} \dicIPA{{s}{\textsubring{d}}{j}{a}{r}{\textsubring{d}}{\textsubring{n}}{f}{r}{a}{i}{ð}{\textsci}{l}{\textepsilon}{\textbabygamma}{\textscy}{\textsubring{r}}} \dicPos{adj}[1]\dicFlx{}[-8] \textbf{1.} \dicDirectTranslationCS{astronomický (dalekohled ap.)}  \textbf{2.} \dicLangCat{přen.} \dicDirectTranslationCS{astronomický, neobyčejně vysoký (cena ap.)}
\dicEntry[stjá] \dicTerm{stjá} \dicIPA{{s}{\textsubring{d}}{j}{au}{\textlengthmark}} \dicPos{n}[2] \dicFlx{(‑s)}[4] \dicPhraseIS{fara á stjá} \dicDirectTranslationCS{vstát, pohnout se z~místa};  \dicPhraseIS{vera á stjái} \dicDirectTranslationCS{být v~pohybu, být na nohou}
\dicEntry[stjóri] \dicTerm{stjór|i} \dicIPA{{s}{\textsubring{d}}{j}{ou}{\textlengthmark}{r}{\textsci}} \dicPos{m}[1] \dicFlx{(‑a, ‑ar)}[1] \textbf{1.} \dicSynonym{akkeri} \dicDirectTranslationCS{(malá) kotva, kámen sloužící jako kotva} \dicExampleIS{Báturinn liggur við stjóra.} \dicExampleCS{Člun je ukotven malou kotvou.}  \textbf{2.} \dicLangCat{hovor.} \dicSynonym{yfirmaður} \dicDirectTranslationCS{šéf(ka), boss}
\dicEntry[stjórn] \dicTerm{stjórn} \dicsymFrequent\  \dicIPA{{s}{\textsubring{d}}{j}{ou}{r}{\textsubring{d}}{\textsubring{n}}} \dicPos{f}[7] \dicFlx{(‑ar, ‑ir)}[1] \textbf{1.} \dicSynonym*{það að stjórna} \dicDirectTranslationCS{vedení, řízení};  \dicPhraseIS{taka við stjórninni} \dicDirectTranslationCS{převzít řízení};  \dicPhraseIS{vera við stjórn} \dicDirectTranslationCS{být ve vedení};  \dicPhraseIS{e‑að er undir stjórn e‑rs} \dicDirectTranslationCS{(co) je pod (čím) vedením}  \textbf{2.} \dicSynonym{hald} \dicDirectTranslationCS{řízení, kontrola};  \dicPhraseIS{hafa stjórn á e‑u} \dicDirectTranslationCS{kontrolovat (co), mít kontrolu nad (čím)};  \dicPhraseIS{missa stjórn á sér} \dicDirectTranslationCS{ztratit nad sebou kontrolu}  \textbf{3.} \dicSynonym*{hópur sem stjórnar} \dicDirectTranslationCS{vláda, vedení, management, představenstvo};  \dicPhraseIS{mynda stjórn} \dicDirectTranslationCS{vytvořit vládu}
\dicEntry[stjórna] \dicTerm{stjórn|a} \dicsymFrequent\  \dicIPA{{s}{\textsubring{d}}{j}{ou}{r}{\textsubring{d}}{n}{a}} \dicPos{v}[1] \dicFlx{(‑aði)}[1] \dicFlx{dat} \textbf{1.} \dicSynonym{ríkja} \dicDirectTranslationCS{vládnout, panovat} \dicExampleIS{stjórna landinu} \dicExampleCS{vládnout zemi}  \textbf{2.} \dicSynonym{stýra} \dicDirectTranslationCS{řídit, vést, ovládat, kontrolovat};  \dicPhraseIS{stjórna bílnum} \dicDirectTranslationCS{řídit auto};  \dicPhraseIS{stjórna fundinum} \dicDirectTranslationCS{vést zasedání};  \dicPhraseIS{stjórna ferðinni} \dicLangCat{přen.} \dicDirectTranslationCS{převzít kontrolu (nad emocemi ap.)};  \dicPhraseIS{stjórna sér} \dicDirectTranslationCS{ovládat se};  \dicIdiom{stjórnast}{ \dicPhraseIS{stjórnast}} \dicFlx{refl} \dicDirectTranslationCS{řídit se, ovládat se};  \dicPhraseIS{láta stjórnast af e‑u} \dicDirectTranslationCS{nechat se ovládnout (čím) (touhou ap.)}
\dicEntry[stjórnandi] \dicTerm{stjórn··|andi} \dicIPA{{s}{\textsubring{d}}{j}{ou}{r}{\textsubring{d}}{n}{a}{n}{\textsubring{d}}{\textsci}} \dicPos{m}[2] \dicFlx{(‑anda, ‑endur)}[1] \textbf{1.} \dicSynonym{foringi} \dicDirectTranslationCS{vládce, vládkyně, lídr(yně), velitel(ka)}  \textbf{2.} \dicSynonym{hljómsveitarstjóri} \dicDirectTranslationCS{dirigent(ka)}
\dicEntry[stjórnarandstaða] \dicTerm{stjórnar··and·|staða} \dicIPA{{s}{\textsubring{d}}{j}{ou}{r}{\textsubring{d}}{n}{a}{r}{a}{n}{s}{\textsubring{d}}{a}{ð}{a}} \dicPos{f}[1] \dicFlx{(‑stöðu)}[2] \dicFieldCat{pol.} \dicDirectTranslationCS{(vládní) opozice}
\dicEntry[stjórnarbylting] \textls[10]{\dicTerm{stjórnar··bylt·ing} \dicIPA{{s}{\textsubring{d}}{j}{ou}{r}{\textsubring{d}}{n}{a}{r}{\textsubring{b}}{\textsci}{\textsubring{l}}{\textsubring{d}}{i}{\ng}{\r{g}}} \dicPos{f}[4] \dicFlx{(‑ar, ‑ar)}[5] \dicFieldCat{pol.}} \dicDirectTranslationCS{vládní převrat, revoluce}
\dicEntry[stjórnarfar] \dicTerm{stjórnar··far} \dicIPA{{s}{\textsubring{d}}{j}{ou}{r}{\textsubring{d}}{n}{a}{\textsubring{r}}{f}{a}{\textsubring{r}}} \dicPos{n}[2] \dicFlx{(‑s)}[2] \dicDirectTranslationCS{způsob vlády}
\dicEntry[stjórnarflokkur] \textls[10]{\dicTerm{stjórnar··flokk|ur} \dicIPA{{s}{\textsubring{d}}{j}{ou}{r}{\textsubring{d}}{n}{a}{\textsubring{r}}{f}{l}{\textopeno}{h}{\r{g}}{\textscy}{\textsubring{r}}} \dicPos{m}[6] \dicFlx{(‑s, ‑ar)}[8] \dicFieldCat{pol.} \dicDirectTranslationCS{vládní strana}}
\dicEntry[stjórnarformaður] \dicTerm{stjórnar··for·|maður} \dicIPA{{s}{\textsubring{d}}{j}{ou}{r}{\textsubring{d}}{n}{a}{\textsubring{r}}{f}{\textopeno}{r}{m}{a}{ð}{\textscy}{\textsubring{r}}} \dicPos{m}[13] \dicFlx{(‑manns, ‑menn)}[2] \dicDirectTranslationCS{předseda\,/\addthin předsedkyně představenstva}
\dicEntry[stjórnarfrumvarp] \textls[15]{\dicTerm{stjórnar··frum·|varp} \dicIPA{{s}{\textsubring{d}}{j}{ou}{r}{\textsubring{d}}{n}{a}{\textsubring{r}}{f}{r}{\textscy}{m}{v}{a}{\textsubring{r}}{\textsubring{b}}} \dicPos{n}[2] \dicFlx{(‑varps, ‑vörp)}[8]} \dicFieldCat{pol.} \dicDirectTranslationCS{vládní návrh (zákona)}
\dicEntry[stjórnarfundur] \dicTerm{stjórnar··fund|ur} \dicIPA{{s}{\textsubring{d}}{j}{ou}{r}{\textsubring{d}}{n}{a}{\textsubring{r}}{f}{\textscy}{n}{\textsubring{d}}{\textscy}{\textsubring{r}}} \dicPos{m}[10] \dicFlx{(‑ar, ‑ir)}[4] \dicDirectTranslationCS{schůze představenstva}
\dicEntry[stjórnarhættir] \dicTerm{stjórnar··hættir} \dicIPA{{s}{\textsubring{d}}{j}{ou}{r}{\textsubring{d}}{n}{a}{\textsubring{r}}{h}{a}{i}{h}{\textsubring{d}}{\textsci}{\textsubring{r}}} \dicPos{m}[12] \dicFlx{pl}[6] \dicSynonym{stjórnarfar} \dicDirectTranslationCS{způsob vládnutí\,/\addthin řízení}
\dicEntry[stjórnarkreppa] \dicTerm{stjórnar··krepp|a} \dicIPA{{s}{\textsubring{d}}{j}{ou}{r}{\textsubring{d}}{n}{a}{\textsubring{r}}{k\smash{\textsuperscript{h}}}{r}{\textepsilon}{h}{\textsubring{b}}{a}} \dicPos{f}[1] \dicFlx{(‑u)}[5] \dicDirectTranslationCS{vládní krize}
\dicEntry[stjórnarmaður] \dicTerm{stjórnar··|maður} \dicIPA{{s}{\textsubring{d}}{j}{ou}{r}{\textsubring{d}}{n}{a}{r}{m}{a}{ð}{\textscy}{\textsubring{r}}} \dicPos{m}[13] \dicFlx{(‑manns, ‑menn)}[2] \dicDirectTranslationCS{člen(ka) představenstva\,/\addthin vlády}
\dicEntry[stjórnarmyndun] \dicTerm{stjórnar··mynd|un} \dicIPA{{s}{\textsubring{d}}{j}{ou}{r}{\textsubring{d}}{n}{a}{r}{m}{\textsci}{n}{\textsubring{d}}{\textscy}{\textsubring{n}}} \dicPos{f}[7] \dicFlx{(‑unar, ‑anir)}[8] \dicFieldCat{pol.} \dicDirectTranslationCS{formování\,/\addthin sestavení vlády}
\dicEntry[stjórnarráð] \dicTerm{stjórnar··ráð} \dicIPA{{s}{\textsubring{d}}{j}{ou}{r}{\textsubring{d}}{n}{a}{r}{au}{\texttheta}} \dicPos{n}[2] \dicFlx{(‑s, ‑)}[5] \dicDirectTranslationCS{vládní úřad\,/\addthin kancelář}
\begin{xtolerant}{}{1pt}
\dicEntry[stjórnarskrá] \dicTerm{stjórnar··skrá}\addthinS\dicIPA{{s}{\textsubring{d}}{j}{ou}{r}{\textsubring{d}}{n}{a}{\textsubring{r}}{s}{\r{g}}{r}{au}}\addthinS\dicPos{f}[4]\addthinS\dicFlx{(‑r\,/\addthin ‑ar, ‑r)}[21]\addthinS\dicFieldCat{práv.}\addthinS\dicDirectTranslationCS{ústava, konstituce}
\end{xtolerant}
\dicEntry[stjórnarstefna] \dicTerm{stjórnar··stefn|a} \dicIPA{{s}{\textsubring{d}}{j}{ou}{r}{\textsubring{d}}{n}{a}{\textsubring{r}}{s}{\textsubring{d}}{\textepsilon}{\textsubring{b}}{n}{a}} \dicPos{f}[1] \dicFlx{(‑u)}[5] \dicFieldCat{pol.} \dicDirectTranslationCS{vládní politika}
\dicEntry[stjórnartíð] \dicTerm{stjórnar··tíð} \dicIPA{{s}{\textsubring{d}}{j}{ou}{r}{\textsubring{d}}{n}{a}{\textsubring{r}}{t\smash{\textsuperscript{h}}}{i}{\texttheta}} \dicPos{f}[7] \dicFlx{(‑ar)}[3] \dicDirectTranslationCS{funkční\,/\addthin vládní období}
\dicEntry[stjórnborð] \dicTerm{stjórn··borð} \dicIPA{{s}{\textsubring{d}}{j}{ou}{r}{\textsubring{d}}{\textsubring{n}}{\textsubring{b}}{\textopeno}{r}{\texttheta}} \dicPos{n}[2] \dicFlx{(‑s)}[2] \dicDirectTranslationCS{řídicí\,/\addthin ovládací\,/\addthin kontrolní panel}
\dicEntry[stjórnborði] \dicTerm{stjórn··borð|i} \dicIPA{{s}{\textsubring{d}}{j}{ou}{r}{\textsubring{d}}{\textsubring{n}}{\textsubring{b}}{\textopeno}{r}{ð}{\textsci}} \dicPos{m}[1] \dicFlx{(‑a, ‑ar)}[1] \dicFieldCat{nám.} \dicDirectTranslationCS{pravobok} \dicExampleIS{á stjórnborða} \dicExampleCS{na pravoboku} \dicAntonym{bakborði}
\dicEntry[stjórnkerfi] \dicTerm{stjórn··kerfi} \dicIPA{{s}{\textsubring{d}}{j}{ou}{r}{\textsubring{d}}{\textsubring{n}}{c\smash{\textsuperscript{h}}}{\textepsilon}{r}{v}{\textsci}} \dicPos{n}[2] \dicFlx{(‑s, ‑)}[14] \textbf{1.} \dicDirectTranslationCS{státní\,/\addthin mocenský aparát}  \textbf{2.} \dicDirectTranslationCS{administrativní systém}
\dicEntry[stjórnklefi] \dicTerm{stjórn··klef|i} \dicIPA{{s}{\textsubring{d}}{j}{ou}{r}{\textsubring{d}}{\textsubring{n}}{k\smash{\textsuperscript{h}}}{l}{\textepsilon}{v}{\textsci}} \dicPos{m}[1] \dicFlx{(‑a, ‑ar)}[1] \dicDirectTranslationCS{kokpit, kabina, kontrolní místnost} \dicExampleIS{stjórnklefi flugvélar} \dicExampleCS{kokpit letadla}
\dicEntry[stjórnlaus] \dicTerm{stjórn··laus} \dicIPA{{s}{\textsubring{d}}{j}{ou}{r}{\textsubring{d}}{n}{l}{\oe i}{s}} \dicPos{adj}[5]\dicFlx{}[-1] \textbf{1.} \dicSynonym*{sem lætur ekki að stjórn} \dicDirectTranslationCS{neovladatelný, nekontrolovatelný}  \textbf{2.} \dicSynonym*{án stjórnar} \dicDirectTranslationCS{(jsoucí) bez vlády, neřízený} \dicExampleIS{stjórnlaust land} \dicExampleCS{země bez vlády}
\dicEntry[stjórnleysi] \dicTerm{stjórn··leysi} \dicIPA{{s}{\textsubring{d}}{j}{ou}{r}{\textsubring{d}}{n}{l}{ei}{s}{\textsci}} \dicPos{n}[2] \dicFlx{(‑s)}[20] \textbf{1.} \dicDirectTranslationCS{anarchie, bezvládí}  \textbf{2.} \dicDirectTranslationCS{nedostatečné\,/\addthin chaotické řízení}
\dicEntry[stjórnleysingi] \dicTerm{stjórn··leys·ing|i} \dicIPA{{s}{\textsubring{d}}{j}{ou}{r}{\textsubring{d}}{n}{l}{ei}{s}{i}{\textltailn}{\r{\textObardotlessj}}{\textsci}} \dicPos{m}[1] \dicFlx{(‑ja, ‑jar)}[14] \dicSynonym{anarkisti} \dicDirectTranslationCS{anarchista, anarchistka}
\dicEntry[stjórnleysisstefna] \dicTerm{stjórn·leysis··stefn|a} \dicIPA{{s}\-{\textsubring{d}}\-{j}\-{ou}\-{r}\-{\textsubring{d}}\-{n}\-{l}\-{ei}\-{s}\-{\textsci}\-{s}\-{\textsubring{d}}\-{\textepsilon}\-{\textsubring{b}}\-{n}\-{a}\-} \dicPos{f}[1] \dicFlx{(‑u)}[5] \dicFieldCat{pol.} \dicSynonym{anarkismi} \dicDirectTranslationCS{anarchismus}
\dicEntry[stjórnmál] \dicTerm{stjórn··mál} \dicsymFrequent\  \dicIPA{{s}{\textsubring{d}}{j}{ou}{r}{\textsubring{d}}{n}{m}{au}{\textsubring{l}}} \dicPos{n}[2] \dicFlx{pl}[1] \dicDirectTranslationCS{politika} \dicExampleIS{þátttaka borgaranna í stjórnmálum} \dicExampleCS{zapojení občanů do politiky}
\dicEntry[stjórnmálaástand] \dicTerm{stjórn·mála··á·stand} \dicIPA{{s}{\textsubring{d}}{j}{ou}{r}{\textsubring{d}}{n}{m}{au}{l}{a}{au}{s}{\textsubring{d}}{a}{n}{\textsubring{d}}} \dicPos{n}[2] \dicFlx{(‑s)}[2] \dicDirectTranslationCS{politická situace}
\dicEntry[stjórnmálaflokkur] \dicTerm{stjórn·mála··flokk|ur} \dicsymFrequent\  \dicIPA{{s}\-{\textsubring{d}}\-{j}\-{ou}\-{r}\-{\textsubring{d}}\-{n}\-{m}\-{au}\-{l}\-{a}\-{f}\-{l}\-{\textopeno}\-{h}\-{\r{g}}\-{\textscy}\-{\textsubring{r}}\-} \dicPos{m}[6] \dicFlx{(‑s, ‑ar)}[8] \dicDirectTranslationCS{politická strana} \dicExampleIS{fulltrúar allra stjórnmálaflokka} \dicExampleCS{zástupci všech politických stran}
\dicEntry[stjórnmálafræði] \dicTerm{stjórn·mála··fræð|i} \dicIPA{{s}{\textsubring{d}}{j}{ou}{r}{\textsubring{d}}{n}{m}{au}{l}{a}{f}{r}{a}{i}{ð}{\textsci}} \dicPos{f}[3] \dicFlx{(‑i)}[3] \dicDirectTranslationCS{politologie}
\begin{xtolerant}{}{1pt}
\dicEntry[stjórnmálafræðingur] \dicTerm{stjórn·mála·fræð··ing|ur} \dicIPA{{s}\-{\textsubring{d}}\-{j}\-{ou}\-{r}\-{\textsubring{d}}\-{n}\-{m}\-{au}\-{l}\-{a}\-{f}\-{r}\-{a}\-{i}\-{ð}\-{i}\-{\ng}\-{\r{g}}\-{\textscy}\-{\textsubring{r}}}\addthinS\dicPos{m}[6] \dicFlx{(‑s, ‑ar)}[8] \dicDirectTranslationCS{politolog, politoložka}
\end{xtolerant}
\dicEntry[stjórnmálahreyfing] \textls[10]{\dicTerm{stjórn·mála··hreyf·ing} \dicIPA{{s}\-{\textsubring{d}}\-{j}\-{ou}\-{r}\-{\textsubring{d}}\-{n}\-{m}\-{au}\-{l}\-{a}\-{\textsubring{r}}\-{ei}\-{v}\-{i}\-{\ng}\-{\r{g}}\-} \dicPos{f}[4] \dicFlx{(‑ar, ‑ar)}[5]} \dicDirectTranslationCS{politické hnutí}
\dicEntry[stjórnmálalegur] \dicTerm{stjórn·mála··legur} \dicIPA{{s}{\textsubring{d}}{j}{ou}{r}{\textsubring{d}}{n}{m}{au}{l}{a}{l}{\textepsilon}{\textbabygamma}{\textscy}{\textsubring{r}}} \dicPos{adj}[1]\dicFlx{}[-8] \dicSynonym{pólitískur} \dicDirectTranslationCS{politický}
\dicEntry[stjórnmálaleiðtogi] \dicTerm{stjórn·mála··leið·tog|i} \dicIPA{{s}\-{\textsubring{d}}\-{j}\-{ou}\-{r}\-{\textsubring{d}}\-{n}\-{m}\-{au}\-{l}\-{a}\-{l}\-{ei}\-{ð}\-{t\smash{\textsuperscript{h}}}\-{\textopeno i}\-{j}\-{\textsci}\-} \dicPos{m}[1] \dicFlx{(‑a, ‑ar)}[1] \dicDirectTranslationCS{politický vůdce\,/\addthin lídr, politická vůdkyně\,/\addthin lídryně}
\dicEntry[stjórnmálamaður] \dicTerm{stjórn·mála··|maður} \dicIPA{{s}{\textsubring{d}}{j}{ou}{r}{\textsubring{d}}{n}{m}{au}{l}{a}{m}{a}{ð}{\textscy}{\textsubring{r}}} \dicPos{m}[13] \dicFlx{(‑manns, ‑menn)}[2] \dicDirectTranslationCS{politik, politička}
\dicEntry[stjórnmálaumræða] \dicTerm{stjórn·mála··um·ræð|a} \dicIPA{{s}\-{\textsubring{d}}\-{j}\-{ou}\-{r}\-{\textsubring{d}}\-{n}\-{m}\-{au}\-{l}\-{a}\-{\textscy}\-{m}\-{r}\-{a}\-{i}\-{ð}\-{a}\-} \dicPos{f}[1] \dicFlx{(‑u, ‑ur)}[19] \dicDirectTranslationCS{politická debata\,/\addthin diskuze}
\dicEntry[stjórnpallur] \dicTerm{stjórn··pall|ur} \dicIPA{{s}{\textsubring{d}}{j}{ou}{r}{\textsubring{d}}{\textsubring{n}}{p\smash{\textsuperscript{h}}}{a}{\textsubring{d}}{l}{\textscy}{\textsubring{r}}} \dicPos{m}[6] \dicFlx{(‑s, ‑ar)}[5] \dicDirectTranslationCS{velitelský\,/\addthin kapitánský můstek}
\dicEntry[stjórnsamur] \dicTerm{stjórn··|samur} \dicIPA{{s}{\textsubring{d}}{j}{ou}{r}{\textsubring{d}}{\textsubring{n}}{s}{a}{m}{\textscy}{\textsubring{r}}} \dicPos{adj}[1] \dicFlx{(f ‑söm)}[2] \dicDirectTranslationCS{panovačný, autoritářský, despotický}
\dicEntry[stjórnskipulag] \dicTerm{stjórn··skipu·lag} \dicIPA{{s}{\textsubring{d}}{j}{ou}{r}{\textsubring{d}}{\textsubring{n}}{s}{\r{\textObardotlessj}}{\textsci}{\textsubring{b}}{\textscy}{l}{a}{x}} \dicPos{n}[2] \dicFlx{(‑s)}[2] \dicSynonym{stjórnskipun} \dicDirectTranslationCS{organizační uspořádání\,/\addthin struktura}
\dicEntry[stjórnskipun] \dicTerm{stjórn··skip|un} \dicIPA{{s}{\textsubring{d}}{j}{ou}{r}{\textsubring{d}}{\textsubring{n}}{s}{\r{\textObardotlessj}}{\textsci}{\textsubring{b}}{\textscy}{\textsubring{n}}} \dicPos{f}[7] \dicFlx{(‑unar, ‑anir)}[8] \dicDirectTranslationCS{ústava, konstituce}
\dicEntry[stjórnskipunarréttur] \dicTerm{stjórn·skipunar··rétt|ur} \dicIPA{{s}\-{\textsubring{d}}\-{j}\-{ou}\-{r}\-{\textsubring{d}}\-{\textsubring{n}}\-{s}\-{\r{\textObardotlessj}}\-{\textsci}\-{\textsubring{b}}\-{\textscy}\-{n}\-{a}\-{r}\-{j}\-{\textepsilon}\-{h}\-{\textsubring{d}}\-{\textscy}\-{\textsubring{r}}\-} \dicPos{m}[10] \dicFlx{(‑ar)}[7] \dicFieldCat{práv.} \dicDirectTranslationCS{ústavní právo}
\dicEntry[stjórnsýsla] \dicTerm{stjórn··sýsl|a} \dicIPA{{s}{\textsubring{d}}{j}{ou}{r}{\textsubring{d}}{\textsubring{n}}{s}{i}{s}{\textsubring{d}}{l}{a}} \dicPos{f}[1] \dicFlx{(‑u)}[5] \dicDirectTranslationCS{státní správa}
\dicEntry[stjórnsýsluréttur] \dicTerm{stjórn·sýslu··rétt|ur} \dicIPA{{s}\-{\textsubring{d}}\-{j}\-{ou}\-{r}\-{\textsubring{d}}\-{\textsubring{n}}\-{s}\-{i}\-{s}\-{\textsubring{d}}\-{l}\-{\textscy}\-{r}\-{j}\-{\textepsilon}\-{h}\-{\textsubring{d}}\-{\textscy}\-{\textsubring{r}}\-} \dicPos{m}[10] \dicFlx{(‑ar)}[7] \dicFieldCat{práv.} \dicDirectTranslationCS{správní právo}
\dicEntry[stjórnturn] \dicTerm{stjórn··turn} \dicIPA{{s}{\textsubring{d}}{j}{ou}{r}{\textsubring{d}}{\textsubring{n}}{t\smash{\textsuperscript{h}}}{\textscy}{r}{\textsubring{d}}{\textsubring{n}}} \dicPos{m}[4] \dicFlx{(‑s, ‑ar)}[3] \dicDirectTranslationCS{kontrolní\,/\addthin řídicí věž} \dicIndirectTranslationCS{(vyvýšené stanoviště letištní dispečerské služby)}
\dicEntry[stjórntæki] \dicTerm{stjórn··tæki} \dicIPA{{s}{\textsubring{d}}{j}{ou}{r}{\textsubring{d}}{\textsubring{n}}{t\smash{\textsuperscript{h}}}{a}{i}{\r{\textObardotlessj}}{\textsci}} \dicPos{n}[2] \dicFlx{(‑s, ‑)}[16] \dicDirectTranslationCS{řídicí\,/\addthin ovládací zařízení, ovládání} \dicExampleIS{stjórntæki vélar} \dicExampleCS{ovládání stroje}
\dicEntry[stjórnun] \dicTerm{stjórn|un} \dicIPA{{s}{\textsubring{d}}{j}{ou}{r}{\textsubring{d}}{n}{\textscy}{\textsubring{n}}} \dicPos{f}[7] \dicFlx{(‑unar)}[9] \textbf{1.} \dicSynonym*{það að stjórna} \dicDirectTranslationCS{vedení, řízení, management} \dicExampleIS{stjórnun fyrirtækja} \dicExampleCS{řízení společností}  \textbf{2.} \dicSynonym{umsjón} \dicDirectTranslationCS{správa, administrativa}
\dicEntry[stjórnvald] \dicTerm{stjórn··|vald} \dicsymFrequent\  \dicIPA{{s}{\textsubring{d}}{j}{ou}{r}{\textsubring{d}}{n}{v}{a}{l}{\textsubring{d}}} \dicPos{n}[2] \dicFlx{(‑valds, ‑völd)}[8] \textbf{1.} \dicFieldCat{pol.} \dicSynonym{ríkisvald} \dicDirectTranslationCS{státní moc\,/\addthin orgán}  \textbf{2.} \dicSynonym*{opinbert yfirvald} \dicDirectTranslationCS{řídicí\,/\addthin výkonný orgán}  \textbf{3.} \dicPhraseIS{stjórnvöld} \dicFlx{pl} \dicFieldCat{pol.} \dicSynonym{ríkisstjórn} \dicDirectTranslationCS{vláda, sbor ministrů} \dicExampleIS{áhrif stjórnvalda á sjávarútveg} \dicExampleCS{vliv vlády na rybářský průmysl}
\dicEntry[stjórnvölur] \dicTerm{stjórn··|völur} \dicIPA{{s}{\textsubring{d}}{j}{ou}{r}{\textsubring{d}}{n}{v}{\oe}{l}{\textscy}{\textsubring{r}}} \dicPos{m}[11] \dicFlx{(‑valar, ‑velir)}[7] \dicDirectTranslationCS{kormidlo};  \dicPhraseIS{sitja\,/\addthin vera við stjórnvölinn} \dicLangCat{přen.} \dicDirectTranslationCS{být u~kormidla, vládnout}
\dicEntry[stjúpa] \dicTerm{stjúp|a} \dicIPA{{s}{\textsubring{d}}{j}{u}{\textlengthmark}{\textsubring{b}}{a}} \dicPos{f}[1] \dicFlx{(‑u, ‑ur)}[19] \textbf{1.} \dicSynonym{stjúpmóðir} \dicDirectTranslationCS{nevlastní matka, macecha}  \textbf{2.} \dicFieldCat{bot.} \dicDirectTranslationCS{maceška, violka zahradní} \textit{(l.~{\textLA{Viola x wittrockiana}})}  \dicsymPhoto\ 
\dicFigure{ds_image_stjupa_0_2.jpg}{Stjúpa}{Stjúpa - JFKCom, GFDL}
\dicEntry[stjúpbarn] \dicTerm{stjúp··|barn} \dicIPA{{s}{\textsubring{d}}{j}{u}{\textlengthmark}{\textsubring{b}}{a}{r}{\textsubring{d}}{\textsubring{n}}} \dicPos{n}[2] \dicFlx{(‑barns, ‑börn)}[8] \dicDirectTranslationCS{nevlastní dítě}
\dicEntry[stjúpbróðir] \textls[10]{\dicTerm{stjúp··|bróðir} \dicIPA{{s}{\textsubring{d}}{j}{u}{\textlengthmark}{\textsubring{b}}{r}{ou}{ð}{\textsci}{\textsubring{r}}} \dicPos{m}[13] \dicFlx{(‑bróður, ‑bræður)}[4] \dicDirectTranslationCS{nevlastní bratr}}
\dicEntry[stjúpdóttir] \dicTerm{stjúp··|dóttir} \dicIPA{{s}{\textsubring{d}}{j}{u}{\textlengthmark}{\textsubring{b}}{\textsubring{d}}{ou}{h}{\textsubring{d}}{\textsci}{\textsubring{r}}} \dicPos{f}[11] \dicFlx{(‑dóttur, ‑dætur)}[2] \dicDirectTranslationCS{nevlastní dcera}
\dicEntry[stjúpfaðir] \dicTerm{stjúp··|faðir} \dicIPA{{s}{\textsubring{d}}{j}{u}{\textlengthmark}{\textsubring{b}}{f}{a}{ð}{\textsci}{\textsubring{r}}} \dicPos{m}[13] \dicFlx{(‑föður, ‑feður)}[3] \dicDirectTranslationCS{nevlastní otec, otčím}
\dicEntry[stjúpi] \dicTerm{stjúp|i} \dicIPA{{s}{\textsubring{d}}{j}{u}{\textlengthmark}{\textsubring{b}}{\textsci}} \dicPos{m}[1] \dicFlx{(‑a, ‑ar)}[1] \dicSynonym{stjúpfaðir} \dicDirectTranslationCS{nevlastní otec, otčím}
\dicEntry[stjúpmóðir] \textls[10]{\dicTerm{stjúp··|móðir} \dicIPA{{s}{\textsubring{d}}{j}{u}{\textlengthmark}{\textsubring{b}}{m}{ou}{ð}{\textsci}{\textsubring{r}}} \dicPos{f}[11] \dicFlx{(‑móður, ‑mæður)}[1] \dicDirectTranslationCS{nevlastní matka, macecha}}
\dicEntry[stjúpsonur] \dicTerm{stjúp··|sonur} \dicIPA{{s}{\textsubring{d}}{j}{u}{\textlengthmark}{\textsubring{b}}{s}{\textopeno}{n}{\textscy}{\textsubring{r}}} \dicPos{m}[12] \dicFlx{(‑sonar, ‑synir)}[9] \dicDirectTranslationCS{nevlastní syn}
\dicEntry[stjúpsystir] \dicTerm{stjúp··syst|ir} \dicIPA{{s}{\textsubring{d}}{j}{u}{\textlengthmark}{\textsubring{b}}{s}{\textsci}{s}{\textsubring{d}}{\textsci}{\textsubring{r}}} \dicPos{f}[11] \dicFlx{(‑ur, ‑ur)}[3] \dicDirectTranslationCS{nevlastní sestra}
\dicEntry[stjúpsystkin] \dicTerm{stjúp··syst·kin} \dicIPA{{s}{\textsubring{d}}{j}{u}{\textlengthmark}{\textsubring{b}}{s}{\textsci}{s}{c\smash{\textsuperscript{h}}}{\textsci}{\textsubring{n}}} \dicPos{n}[2] \dicFlx{pl}[1] \dicLink{stjúpsystkini}
\dicEntry[stjúpsystkini] \dicTerm{stjúp··syst·kini}\dicTerm{, stjúpsystkin} \dicIPA{{s}\-{\textsubring{d}}\-{j}\-{u}\-{\textlengthmark}\-{\textsubring{b}}\-{s}\-{\textsci}\-{s}\-{c\smash{\textsuperscript{h}}}\-{\textsci}\-{n}\-{\textsci}\-} \dicPos{n}[2] \dicFlx{pl}[19] \dicDirectTranslationCS{nevlastní sourozenci}
\dicEntry[stjörnu] \dicTerm{stjörnu} \dicIPA{{s}{\textsubring{d}}{j}{\oe}{r}{\textsubring{d}}{n}{\textscy}} \dicPos{f} \dicFlx{sg gen} \dicLink{stjarna}
\dicEntry[stjörnuár] \dicTerm{stjörnu··ár} \dicIPA{{s}{\textsubring{d}}{j}{\oe}{r}{\textsubring{d}}{n}{\textscy}{au}{\textsubring{r}}} \dicPos{n}[2] \dicFlx{(‑s, ‑)}[5] \dicFieldCat{astro.} \dicDirectTranslationCS{siderický\,/\addthin hvězdný rok}
\begin{xtolerant}{}{1pt}
\dicEntry[stjörnubjartur] \dicTerm{stjörnu··|bjartur} \dicIPA{{s}{\textsubring{d}}{j}{\oe}{r}{\textsubring{d}}{n}{\textscy}{\textsubring{b}}{j}{a}{\textsubring{r}}{\textsubring{d}}{\textscy}{\textsubring{r}}}\addthinS\dicPos{adj}[1]\addthinS\dicFlx{(f ‑björt)}[11] \dicDirectTranslationCS{hvězd\-ný, pokrytý hvězdami (o~noční obloze ap.)}
\end{xtolerant}
\dicEntry[stjörnufræði] \dicTerm{stjörnu··fræð|i} \dicIPA{{s}{\textsubring{d}}{j}{\oe}{r}{\textsubring{d}}{n}{\textscy}{f}{r}{a}{i}{ð}{\textsci}} \dicPos{f}[3] \dicFlx{(‑i)}[3] \dicDirectTranslationCS{astronomie}
\dicEntry[stjörnufræðingur] \dicTerm{stjörnu·fræð··ing|ur} \dicIPA{{s}{\textsubring{d}}{j}{\oe}{r}{\textsubring{d}}{n}{\textscy}{f}{r}{a}{i}{ð}{i}{\ng}{\r{g}}{\textscy}{\textsubring{r}}} \dicPos{m}[6] \dicFlx{(‑s, ‑ar)}[8] \dicDirectTranslationCS{astronom(ka)}
\dicEntry[stjörnuhrap] \dicTerm{stjörnu··|hrap} \dicIPA{{s}{\textsubring{d}}{j}{\oe}{r}{\textsubring{d}}{n}{\textscy}{\textsubring{r}}{a}{\textsubring{b}}} \dicPos{n}[2] \dicFlx{(‑hraps, ‑hröp)}[8] \dicDirectTranslationCS{létavice, meteor}
\dicEntry[stjörnukíkir] \dicTerm{stjörnu··kík|ir} \dicIPA{{s}{\textsubring{d}}{j}{\oe}{r}{\textsubring{d}}{n}{\textscy}{c\smash{\textsuperscript{h}}}{i}{\r{\textObardotlessj}}{\textsci}{\textsubring{r}}} \dicPos{m}[7] \dicFlx{(‑is, ‑jar)}[4] \dicDirectTranslationCS{teleskop, astronomický\,/\addthin hvězdářský dalekohled}
\dicEntry[stjörnulaus] \dicTerm{stjörnu··laus} \dicIPA{{s}{\textsubring{d}}{j}{\oe}{r}{\textsubring{d}}{n}{\textscy}{l}{\oe i}{s}} \dicPos{adj}[5]\dicFlx{}[-4] \dicDirectTranslationCS{bezhvězdný}
\dicEntry[stjörnuljós] \dicTerm{stjörnu··ljós} \dicIPA{{s}{\textsubring{d}}{j}{\oe}{r}{\textsubring{d}}{n}{\textscy}{l}{j}{ou}{s}} \dicPos{n}[2] \dicFlx{(‑s, ‑)}[5] \dicDirectTranslationCS{prskavka}
\dicEntry[stjörnumerki] \dicTerm{stjörnu··merki} \dicIPA{{s}{\textsubring{d}}{j}{\oe}{r}{\textsubring{d}}{n}{\textscy}{m}{\textepsilon}{\textsubring{r}}{\r{\textObardotlessj}}{\textsci}} \dicPos{n}[2] \dicFlx{(‑s, ‑)}[16] \textbf{1.} \dicFieldCat{astro.} \dicDirectTranslationCS{souhvězdí}  \textbf{2.} \dicDirectTranslationCS{znamení zvěrokruhu} \dicExampleIS{Í hvaða stjörnumerki ert þú?} \dicExampleCS{V~jakém jsi znamení?}
\dicEntry[stjörnusjónauki] \dicTerm{stjörnu··sjón·auk|i} \dicIPA{{s}{\textsubring{d}}{j}{\oe}{r}{\textsubring{d}}{n}{\textscy}{s}{j}{ou}{n}{\oe i}{\r{\textObardotlessj}}{\textsci}} \dicPos{m}[1] \dicFlx{(‑a, ‑ar)}[1] \dicDirectTranslationCS{te\-le\-skop, astronomický dalekohled}
\dicEntry[stjörnuskrúfa] \dicTerm{stjörnu··skrúf|a} \dicIPA{{s}{\textsubring{d}}{j}{\oe}{r}{\textsubring{d}}{n}{\textscy}{s}{\r{g}}{r}{u}{v}{a}} \dicPos{f}[1] \dicFlx{(‑u, ‑ur)}[19] \dicDirectTranslationCS{křížový šroub}
\dicEntry[stjörnuskrúfjárn] \dicTerm{stjörnu··skrúf·járn} \dicIPA{{s}{\textsubring{d}}{j}{\oe}{r}{\textsubring{d}}{n}{\textscy}{s}{\r{g}}{r}{u}{v}{j}{au}{r}{\textsubring{d}}{\textsubring{n}}} \dicPos{n}[2] \dicFlx{(‑s, ‑)}[5] \dicDirectTranslationCS{křížový šroubovák}
\dicEntry[stjörnuspá] \dicTerm{stjörnu··spá} \dicIPA{{s}{\textsubring{d}}{j}{\oe}{r}{\textsubring{d}}{n}{\textscy}{s}{\textsubring{b}}{au}} \dicPos{f}[4] \dicFlx{(‑r, ‑r)}[18] \dicDirectTranslationCS{horoskop}
\dicEntry[stjörnuspámaður] \dicTerm{stjörnu··spá·|maður} \dicIPA{{s}{\textsubring{d}}{j}{\oe}{r}{\textsubring{d}}{n}{\textscy}{s}{\textsubring{b}}{au}{m}{a}{ð}{\textscy}{\textsubring{r}}} \dicPos{m}[13] \dicFlx{(‑manns, ‑menn)}[2] \dicDirectTranslationCS{astrolog, astroložka}
\dicEntry[stjörnuspeki] \dicTerm{stjörnu··spek|i} \dicIPA{{s}{\textsubring{d}}{j}{\oe}{r}{\textsubring{d}}{n}{\textscy}{s}{\textsubring{b}}{\textepsilon}{\r{\textObardotlessj}}{\textsci}} \dicPos{f}[3] \dicFlx{(‑i)}[3] \dicDirectTranslationCS{astrologie}
\dicEntry[stjörnusteinbrjótur] \dicTerm{stjörnu··stein·brjót|ur} \dicIPA{{s}\-{\textsubring{d}}\-{j}\-{\oe}\-{r}\-{\textsubring{d}}\-{n}\-{\textscy}\-{s}\-{\textsubring{d}}\-{ei}\-{n}\-{\textsubring{b}}\-{r}\-{j}\-{ou}\-{\textsubring{d}}\-{\textscy}\-{\textsubring{r}}\-} \dicPos{m}[6] \dicFlx{(‑s, ‑ar)}[10] \dicFieldCat{bot.} \dicDirectTranslationCS{lomikámen hvězdovitý} \textit{(l.~{\textLA{Saxifraga stellaris}})}  \rlap{\dicsymPhoto\ }
\dicFigure{45164.jpg}{Stjörnusteinbrjótur}{Stjörnusteinbrjótur - Kesl Michael, Biolib, Copyright/CC-BY-NC}
\dicEntry[stjörnuturn] \dicTerm{stjörnu··turn} \dicIPA{{s}{\textsubring{d}}{j}{\oe}{r}{\textsubring{d}}{n}{\textscy}{t\smash{\textsuperscript{h}}}{\textscy}{r}{\textsubring{d}}{\textsubring{n}}} \dicPos{m}[4] \dicFlx{(‑s, ‑ar)}[3] \dicDirectTranslationCS{hvězdárna, observatoř}
\dicEntry[stjörnuver] \dicTerm{stjörnu··ver} \dicIPA{{s}{\textsubring{d}}{j}{\oe}{r}{\textsubring{d}}{n}{\textscy}{v}{\textepsilon}{\textsubring{r}}} \dicPos{n}[2] \dicFlx{(‑s, ‑)}[5] \dicDirectTranslationCS{planetárium}
\dicEntry[stjörnuþoka] \dicTerm{stjörnu··þok|a} \dicIPA{{s}{\textsubring{d}}{j}{\oe}{r}{\textsubring{d}}{n}{\textscy}{\texttheta}{\textopeno}{\r{g}}{a}} \dicPos{f}[1] \dicFlx{(‑u)}[5] \dicFieldCat{astro.} \dicDirectTranslationCS{mlhovina}
\dicEntry[stjörnuþyrping] \textls[10]{\dicTerm{stjörnu··þyrp·ing} \dicIPA{{s}{\textsubring{d}}{j}{\oe}{r}{\textsubring{d}}{n}{\textscy}{\texttheta}{\textsci}{\textsubring{r}}{\textsubring{b}}{i}{\ng}{\r{g}}} \dicPos{f}[4] \dicFlx{(‑ar, ‑ar)}[5] \dicFieldCat{astro.}} \dicDirectTranslationCS{hvězdokupa}
\dicEntry[stk.] \dicTerm{stk.} \dicPos{zkr} \dicPhraseIS{stykki} \dicDirectTranslationCS{kus}
\dicEntry[stoð] \dicTerm{stoð} \dicIPA{{s}{\textsubring{d}}{\textopeno}{\textlengthmark}{\texttheta}} \dicPos{f}[7] \dicFlx{(‑ar, ‑ir)}[1] \textbf{1.} \dicSynonym{stólpi} \dicDirectTranslationCS{podpěra, vzpěra, opora} \dicExampleIS{stoð í fjósi} \dicExampleCS{podpěra v~kravíně};  \dicPhraseIS{renna stoðum undir e‑ð} \dicDirectTranslationCS{podpořit (argumenty) (co), nalézt důvody pro (co)}  \textbf{2.} \dicSynonym{stuðningur} \dicDirectTranslationCS{podpora, opora, pomoc};  \dicPhraseIS{e‑að á sér enga\,/\addthin ekki stoð} \dicDirectTranslationCS{(co) nemá žádné opodstatnění};  \dicPhraseIS{e‑að kippir stoðunum undan e‑u} \dicLangCat{přen.} \dicDirectTranslationCS{(co) podrývá (samotné) základy pod (čím)};  \dicPhraseIS{vera e‑m stoð í e‑u} \dicDirectTranslationCS{být (komu) oporou v~(čem)} \dicExampleIS{Gunna var ömmu sinni mikil stoð í ellinni.} \dicExampleCS{Gunna byla babičce velkou oporou ve stáří.};  \dicPhraseIS{vera stoð e‑rs og stytta} \dicDirectTranslationCS{být (komu) (životní) oporou}
\dicEntry[stoða] \dicTerm{stoð|a} \dicIPA{{s}{\textsubring{d}}{\textopeno}{\textlengthmark}{ð}{a}} \dicPos{v}[1] \dicFlx{(‑aði)}[79] \dicSynonym{hjálpa} \dicDirectTranslationCS{(vy)pomoct, být užitečný} \dicIndirectTranslationCS{(používané v~otázce nebo záporu)};  \dicPhraseIS{það stoðar ekki að (gera e‑ð)} \dicFlx{impers} \dicDirectTranslationCS{nepomůže (dělat (co)), nemá cenu (dělat (co))}
\dicEntry[stoðkerfi] \dicTerm{stoð··kerfi} \dicIPA{{s}{\textsubring{d}}{\textopeno}{ð}{c\smash{\textsuperscript{h}}}{\textepsilon}{r}{v}{\textsci}} \dicPos{n}[2] \dicFlx{(‑s, ‑)}[14] \textbf{1.} \dicDirectTranslationCS{nosný systém}  \textbf{2.} \dicFieldCat{anat.} \dicDirectTranslationCS{kostra}
\dicEntry[stofa] \dicTerm{stof|a} \dicsymFrequent\  \dicIPA{{s}{\textsubring{d}}{\textopeno}{\textlengthmark}{v}{a}} \dicPos{f}[1] \dicFlx{(‑u, ‑ur)}[7] \textbf{1.} \dicDirectTranslationCS{(hlavní) místnost}  \textbf{2.} \dicSynonym{dagstofa} \dicDirectTranslationCS{obývací pokoj, obývák} \dicExampleIS{sófi í lítilli stofu} \dicExampleCS{gauč v~malém obýváku}  \textbf{3.} \dicSynonym{skrifstofa} \dicDirectTranslationCS{kancelář, úřadovna, pracovna, kabinet, byro} \dicExampleIS{Hann gekk að skrifborði í horni stofunnar.} \dicExampleCS{Šel k~psacímu stolu v~rohu pracovny.}  \textbf{4.} \dicSynonym{læknastofa} \dicDirectTranslationCS{ordinace}  \textbf{5.} \dicSynonym{skólastofa} \dicDirectTranslationCS{třída, učebna} \dicExampleIS{Tölva er í öllum stofum.} \dicExampleCS{Počítač je ve všech učebnách.}  \textbf{6.} \dicSynonym{sjúkrastofa} \dicDirectTranslationCS{ošetřovna}
\dicEntry[stofn] \dicTerm{stofn} \dicsymFrequent\  \dicIPA{{s}{\textsubring{d}}{\textopeno}{\textsubring{b}}{\textsubring{n}}} \dicPos{m}[4] \dicFlx{(‑s, ‑ar)}[3] \textbf{1.} \dicSynonym*{trjástofn} \dicDirectTranslationCS{kmen (stromu ap.)} \dicExampleIS{Sterklegu stofnarnir voru rauðir á litinn.} \dicExampleCS{Mohutné kmeny měly červenou barvu.}  \textbf{2.} \dicSynonym{ætt} \dicDirectTranslationCS{kmen, rod, rasa}  \textbf{3.} \dicFieldCat{jaz.} \dicSynonym{orðstofn} \dicDirectTranslationCS{kmen}  \textbf{4.} \dicSynonym{grundvöllur} \dicDirectTranslationCS{základ, podklad};  \dicPhraseIS{setja á stofn  e‑ð} \dicDirectTranslationCS{založit (co), zřídit (co), vybudovat (co)}
\dicEntry[stofna] \dicTerm{stofn|a} \dicsymFrequent\  \dicIPA{{s}{\textsubring{d}}{\textopeno}{\textsubring{b}}{n}{a}} \dicPos{v}[1] \dicFlx{(‑aði)}[1] \dicFlx{acc\,/\addthin dat} \dicFlx{acc} \dicSynonym*{setja á stofn} \dicDirectTranslationCS{založit, zřídit, vytvořit, vybudovat} \dicExampleIS{stofna félag} \dicExampleCS{založit společnost};  \dicIdiom{stofna}[í]{ \dicPhraseIS{stofna e‑m í e‑ð}} \dicDirectTranslationCS{dostat (koho) do (čeho) (nebezpečí ap.), vystavit (koho čemu)} \dicExampleIS{stofna honum í hættu} \dicExampleCS{dostat ho do nebezpečí};  \dicIdiom{stofna}[til]{ \dicPhraseIS{stofna til e‑s}} \dicSynonym{valda\smash{\textsuperscript{2}}} \dicDirectTranslationCS{přivodit (co), vyvolat (co), zapříčinit (co), zadělat na (co)} \dicExampleIS{stofna til illinda} \dicExampleCS{vyvolat hádku};  \dicPhraseIS{stofna til kynna við e‑n} \dicDirectTranslationCS{seznámit se s~(kým), navázat známost s~(kým)}
\dicEntry[stofnandi] \dicTerm{stofn··|andi} \dicsymFrequent\  \dicIPA{{s}{\textsubring{d}}{\textopeno}{\textsubring{b}}{n}{a}{n}{\textsubring{d}}{\textsci}} \dicPos{m}[2] \dicFlx{(‑anda, ‑endur)}[1] \dicDirectTranslationCS{zakladatel(ka)} \dicExampleIS{stofnendur hljómsveitarinnar} \dicExampleCS{zakladatelé hudební skupiny}
\dicEntry[stofnfé] \dicTerm{stofn··|fé} \dicIPA{{s}{\textsubring{d}}{\textopeno}{\textsubring{b}}{\textsubring{n}}{f}{j}{\textepsilon}} \dicPos{n}[3] \dicFlx{(‑fjár)}[5] \dicFieldCat{ekon.} \dicDirectTranslationCS{počáteční\,/\addthin základní kapitál} \dicExampleIS{stofnfé félags} \dicExampleCS{základní kapitál společnosti}
\dicEntry[stofnfjármagn] \dicTerm{stofn··fjár·magn} \dicIPA{{s}{\textsubring{d}}{\textopeno}{\textsubring{b}}{\textsubring{n}}{f}{j}{au}{r}{m}{a}{\r{g}}{\textsubring{n}}} \dicPos{n}[2] \dicFlx{(‑s)}[2] \dicFieldCat{ekon.} \dicDirectTranslationCS{počáteční\,/\addthin základní kapitál}
\dicEntry[stofnfruma] \dicTerm{stofn··frum|a} \dicIPA{{s}{\textsubring{d}}{\textopeno}{\textsubring{b}}{\textsubring{n}}{f}{r}{\textscy}{m}{a}} \dicPos{f}[1] \dicFlx{(‑u, ‑ur)}[19] \dicFieldCat{biol.} \dicDirectTranslationCS{kmenová buňka}
\dicEntry[stofnfundur] \textls[10]{\dicTerm{stofn··fund|ur} \dicIPA{{s}{\textsubring{d}}{\textopeno}{\textsubring{b}}{\textsubring{n}}{f}{\textscy}{n}{\textsubring{d}}{\textscy}{\textsubring{r}}} \dicPos{m}[10] \dicFlx{(‑ar, ‑ir)}[4] \dicDirectTranslationCS{ustavující schůze}}
\dicEntry[stofnkostnaður] \dicTerm{stofn··kost·nað|ur} \dicIPA{{s}{\textsubring{d}}{\textopeno}{\textsubring{b}}{\textsubring{n}}{k\smash{\textsuperscript{h}}}{\textopeno}{s}{\textsubring{d}}{n}{a}{ð}{\textscy}{\textsubring{r}}} \dicPos{m}[10] \dicFlx{(‑ar)}[9] \dicDirectTranslationCS{zřizovací náklady, náklady na zřízení}
\dicEntry[stofnlán] \dicTerm{stofn··lán} \dicIPA{{s}{\textsubring{d}}{\textopeno}{\textsubring{b}}{n}{l}{au}{\textsubring{n}}} \dicPos{n}[2] \dicFlx{(‑s, ‑)}[5] \dicFieldCat{ekon.} \dicDirectTranslationCS{počáteční půjčka}
\dicEntry[stofnlægur] \dicTerm{stofn··lægur} \dicIPA{{s}{\textsubring{d}}{\textopeno}{\textsubring{b}}{n}{l}{a}{i}{\textbabygamma}{\textscy}{\textsubring{r}}} \dicPos{adj}[1]\dicFlx{}[-1] \dicFieldCat{jaz.} \dicDirectTranslationCS{(jsoucí) součástí kmene}
\dicEntry[stofnsáttmáli] \dicTerm{stofn··sátt·mál|i} \dicIPA{{s}{\textsubring{d}}{\textopeno}{\textsubring{b}}{\textsubring{n}}{s}{au}{h}{\textsubring{d}}{m}{au}{l}{\textsci}} \dicPos{m}[1] \dicFlx{(‑a, ‑ar)}[1] \dicFieldCat{pol.} \dicDirectTranslationCS{pakt}
\dicEntry[stofnset] \dicTerm{stofn··set} \dicIPA{{s}{\textsubring{d}}{\textopeno}{\textsubring{b}}{\textsubring{n}}{s}{\textepsilon}{\textsubring{d}}} \dicPos{v} \dicFlx{ind praes sg 1 pers} \dicLink{stofnsetja}
\dicEntry[stofnsetja] \dicTerm{stofn··|setja} \dicIPA{{s}{\textsubring{d}}{\textopeno}{\textsubring{b}}{\textsubring{n}}{s}{\textepsilon}{\textsubring{d}}{j}{a}} \dicPos{v}[4] \dicFlx{(‑set, ‑setti, ‑settum, ‑setti, ‑sett)}[26] \dicFlx{acc} \dicDirectTranslationCS{založit, zřídit, vytvořit, vybudovat, ustavit} \dicExampleIS{stofnsetja fyrirtæki} \dicExampleCS{zřídit firmu}
\dicEntry[stofnsett] \dicTerm{stofn··sett} \dicIPA{{s}{\textsubring{d}}{\textopeno}{\textsubring{b}}{\textsubring{n}}{s}{\textepsilon}{h}{\textsubring{d}}} \dicPos{v} \dicFlx{supin} \dicLink{stofnsetja}
\dicEntry[stofnsetti] \dicTerm{stofn··setti} \dicIPA{{s}{\textsubring{d}}{\textopeno}{\textsubring{b}}{\textsubring{n}}{s}{\textepsilon}{h}{\textsubring{d}}{\textsci}} \dicPos{v} \dicFlx{ind\,/\addthin con pf sg 1 pers} \dicLink{stofnsetja}
\dicEntry[stofnsettum] \dicTerm{stofn··settum} \dicIPA{{s}{\textsubring{d}}{\textopeno}{\textsubring{b}}{\textsubring{n}}{s}{\textepsilon}{h}{\textsubring{d}}{\textscy}{\textsubring{m}}} \dicPos{v} \dicFlx{ind pf pl 1 pers} \dicLink{stofnsetja}
\dicEntry[stofnsjóður] \dicTerm{stofn··sjóð|ur} \dicIPA{{s}{\textsubring{d}}{\textopeno}{\textsubring{b}}{\textsubring{n}}{s}{j}{ou}{ð}{\textscy}{\textsubring{r}}} \dicPos{m}[9] \dicFlx{(‑s, ‑ir)}[6] \dicFieldCat{ekon.} \dicDirectTranslationCS{základní fond}
\dicEntry[stofnskrá] \dicTerm{stofn··skrá} \dicIPA{{s}{\textsubring{d}}{\textopeno}{\textsubring{b}}{\textsubring{n}}{s}{\r{g}}{r}{au}} \dicPos{f}[4] \dicFlx{(‑r\,/\addthin ‑ar, ‑r)}[21] \dicFieldCat{pol.} \dicDirectTranslationCS{smlouva, pakt, úmluva, dohoda}
\dicEntry[stofnun] \dicTerm{stofn|un} \dicsymFrequent\  \dicIPA{{s}{\textsubring{d}}{\textopeno}{\textsubring{b}}{n}{\textscy}{\textsubring{n}}} \dicPos{f}[7] \dicFlx{(‑unar, ‑anir)}[8] \textbf{1.} \dicSynonym*{það að stofna} \dicDirectTranslationCS{založení, zřízení, vybudování, vytvoření, ustavení}  \textbf{2.} \dicDirectTranslationCS{instituce, institut, ústav};  \dicPhraseIS{opinber stofnun} \dicDirectTranslationCS{veřejná instituce}
\dicEntry[stofnöryggi] \dicTerm{stofn··ör·yggi} \dicIPA{{s}{\textsubring{d}}{\textopeno}{\textsubring{b}}{n}{\oe}{r}{\textsci}{\r{\textObardotlessj}}{\textsci}} \dicPos{n}[2] \dicFlx{(‑s, ‑)}[16] \dicFieldCat{elek.} \dicDirectTranslationCS{hlavní pojistka}
\dicEntry[stofublóm] \dicTerm{stofu··blóm} \dicIPA{{s}{\textsubring{d}}{\textopeno}{\textlengthmark}{v}{\textscy}{\textsubring{b}}{l}{ou}{\textsubring{m}}} \dicPos{n}[2] \dicFlx{(‑s, ‑)}[5] \dicDirectTranslationCS{pokojová rostlina}
\dicEntry[stofufangelsi] \dicTerm{stofu··fang·elsi} \dicIPA{{s}{\textsubring{d}}{\textopeno}{\textlengthmark}{v}{\textscy}{f}{au}{\textltailn}{\r{\textObardotlessj}}{\textepsilon}{l}{s}{\textsci}} \dicPos{n}[2] \dicFlx{(‑s, ‑)}[14] \dicDirectTranslationCS{domácí vězení}
\dicEntry[stofuhiti] \dicTerm{stofu··hit|i} \dicIPA{{s}{\textsubring{d}}{\textopeno}{\textlengthmark}{v}{\textscy}{h}{\textsci}{\textsubring{d}}{\textsci}} \dicPos{m}[1] \dicFlx{(‑a)}[3] \dicDirectTranslationCS{pokojová teplota};  \dicPhraseIS{geyma e‑ð við stofuhita} \dicDirectTranslationCS{uchovávat (co) při pokojové teplotě}
\dicEntry[stofustúlka] \dicTerm{stofu··stúlk|a} \dicIPA{{s}{\textsubring{d}}{\textopeno}{\textlengthmark}{v}{\textscy}{s}{\textsubring{d}}{u}{\textsubring{l}}{\r{g}}{a}} \dicPos{f}[1] \dicFlx{(‑u, ‑ur)}[13] \dicDirectTranslationCS{pokojská}
\dicEntry[stokka] \dicTerm{stokk|a} \dicIPA{{s}{\textsubring{d}}{\textopeno}{h}{\r{g}}{a}} \dicPos{v}[1] \dicFlx{(‑aði)}[1] \dicFlx{acc} \dicSynonym{blanda\smash{\textsuperscript{2}}} \dicDirectTranslationCS{(pro)míchat, zamíchat, smíchat} \dicExampleIS{stokka skjöl} \dicExampleCS{smíchat listiny};  \dicPhraseIS{stokka spil} \dicDirectTranslationCS{(za)míchat karty};  \dicIdiom{stokka}[upp]{ \dicPhraseIS{stokka e‑ð upp}} \dicDirectTranslationCS{zamíchat s~(čím), přeskupit (co) (důkladně)}
\dicEntry[stokkbólginn] \dicTerm{stokk··bólginn} \dicIPA{{s}{\textsubring{d}}{\textopeno}{h}{\r{g}}{\textsubring{b}}{ou}{l}{\r{\textObardotlessj}}{\textsci}{\textsubring{n}}} \dicPos{adj}[6]\dicFlx{}[-6] \dicDirectTranslationCS{velmi oteklý, nateklý, napuchlý}
\dicEntry[stokkfreðinn] \dicTerm{stokk··freðinn} \dicIPA{{s}{\textsubring{d}}{\textopeno}{h}{\r{g}}{f}{r}{\textepsilon}{ð}{\textsci}{\textsubring{n}}} \dicPos{adj}[6]\dicFlx{}[-6] \dicDirectTranslationCS{zmrzlý na kost}
\dicEntry[Stokkhólmur] \dicTerm{Stokk··hólm|ur} \dicIPA{{s}{\textsubring{d}}{\textopeno}{h}{\r{g}}{h}{ou}{l}{m}{\textscy}{\textsubring{r}}} \dicPos{m}[6] \dicFlx{(‑s)}[3] \dicFieldCat{geog.} \dicDirectTranslationCS{Stockholm} \dicIndirectTranslationCS{(hlavní město Švédska)}
\dicEntry[stokkið] \dicTerm{stokkið} \dicIPA{{s}{\textsubring{d}}{\textopeno}{h}{\r{\textObardotlessj}}{\textsci}{\texttheta}} \dicPos{v} \dicFlx{supin} \dicLink{stökkva\smash{\textsuperscript{2}}}
\dicEntry[stokkroðna] \dicTerm{stokk··roðn|a} \dicIPA{{s}{\textsubring{d}}{\textopeno}{h}{\r{g}}{r}{\textopeno}{ð}{n}{a}} \dicPos{v}[1] \dicFlx{(‑aði)}[44] \dicDirectTranslationCS{(velmi) zrudnout, (za)rdít se}
\dicEntry[stokkur] \dicTerm{stokk|ur} \dicIPA{{s}{\textsubring{d}}{\textopeno}{h}{\r{g}}{\textscy}{\textsubring{r}}} \dicPos{m}[6] \dicFlx{(‑s, ‑ar)}[22] \textbf{1.} \dicSynonym{askja} \dicDirectTranslationCS{krabička, škatulka} \dicExampleIS{prjónastokkur} \dicExampleCS{škatulka na špendlíky}  \textbf{2.} \dicSynonym{spilastokkur} \dicDirectTranslationCS{balíček (karet)}  \textbf{3.} \dicSynonym{farvegur} \dicDirectTranslationCS{kanál, průplav, náhon}  \textbf{4.} \dicSynonym{rúmstokkur} \dicDirectTranslationCS{ohrádka (do postýlky)} \dicExampleIS{sofa við stokk} \dicExampleCS{spát s~ohrádkou};  \dicIdiom{stokkur}{ \dicPhraseIS{hleypa e‑u af stokkunum}} \dicDirectTranslationCS{spustit (co), uvést (co) do provozu};  \dicPhraseIS{innan stokks} \dicFlx{adv} \dicDirectTranslationCS{doma, v~domě};  \dicPhraseIS{stíga á stokk og strengja þess heit að (gera e‑ð)} \dicLangCat{přen.} \dicDirectTranslationCS{slavnostně odpřisáhnout (dělání (čeho))}; { \dicPhraseIS{trúa á stokka og steina}} \dicDirectTranslationCS{věřit na nadpřirozené bytosti};  \dicPhraseIS{utan stokks} \dicFlx{adv} \dicDirectTranslationCS{venku, mimo dům}
\dicEntry[stokkönd] \dicTerm{stokk··|önd} \dicIPA{{s}{\textsubring{d}}{\textopeno}{h}{\r{g}}{\oe}{n}{\textsubring{d}}} \dicPos{f}[8] \dicFlx{(‑andar, ‑endur)}[2] \dicFieldCat{zool.} \dicDirectTranslationCS{kachna divoká} \textit{(l.~{\textLA{Anas platyrhynchos}})}  \dicsymPhoto\ 
\dicFigure{98.jpg}{Stokkönd}{Stokkönd - Zicha Ondřej, Biolib, Copyright/CC-BY-NC}
\dicEntry[stolið] \dicTerm{stolið} \dicIPA{{s}{\textsubring{d}}{\textopeno}{\textlengthmark}{l}{\textsci}{\texttheta}} \dicPos{v} \dicFlx{supin} \dicLink{stela}
\dicEntry[stolinn] \dicTerm{stolinn} \dicIPA{{s}{\textsubring{d}}{\textopeno}{\textlengthmark}{l}{\textsci}{\textsubring{n}}} \dicPos{adj}[6]\dicFlx{}[-6] \dicDirectTranslationCS{(u)kradený, zcizený, odcizený}
\dicEntry[stolt] \dicTerm{stolt} \dicsymFrequent\  \dicIPA{{s}{\textsubring{d}}{\textopeno}{\textsubring{l}}{\textsubring{d}}} \dicPos{n}[2] \dicFlx{(‑s)}[2] \dicDirectTranslationCS{hrdost, pýcha} \dicExampleIS{særa stolt e‑rs} \dicExampleCS{zranit (čí) hrdost}
\dicEntry[stoltur] \dicTerm{stoltur} \dicsymFrequent\  \dicIPA{{s}{\textsubring{d}}{\textopeno}{\textsubring{l}}{\textsubring{d}}{\textscy}{\textsubring{r}}} \dicPos{adj}[1]\dicFlx{}[-10] \dicSynonym{rogginn} \dicDirectTranslationCS{hrdý, pyšný};  \dicPhraseIS{vera stoltur af e‑u} \dicDirectTranslationCS{být pyšný na (co)} \dicExampleIS{vera stoltur af afreki sínu} \dicExampleCS{být pyšný na svůj výkon}
\dicEntry[stopp] \dicTerm{stopp\smash{\textsuperscript{1}}} \dicIPA{{s}{\textsubring{d}}{\textopeno}{h}{\textsubring{b}}} \dicPos{n}[2] \dicFlx{(‑s, ‑)}[5] \textbf{1.} \dicSynonym{stans} \dicDirectTranslationCS{zastávka, zastavení}  \textbf{2.} \dicSynonym{tróð\smash{\textsuperscript{1}}} \dicDirectTranslationCS{vycpávka, výplň} \dicExampleIS{troða þurru stoppi í holið} \dicExampleCS{nacpat suchou vycpávku do díry}
\dicEntry[stopp] \dicTerm{stopp\smash{\textsuperscript{2}}} \dicIPA{{s}{\textsubring{d}}{\textopeno}{h}{\textsubring{b}}} \dicPos{adj}[13] \dicFlx{indecl}[2] \dicDirectTranslationCS{zastavený (dopravní ruch ap.)}
\dicEntry[stoppa] \dicTerm{stopp|a} \dicsymFrequent\  \dicIPA{{s}{\textsubring{d}}{\textopeno}{h}{\textsubring{b}}{a}} \dicPos{v}[1] \dicFlx{(‑aði)}[1] \dicFlx{acc} \textbf{1.} \dicSynonym{stansa} \dicDirectTranslationCS{zastavit, zastavovat, stopnout} \dicExampleIS{stoppa þar augnablik} \dicExampleCS{zastavit tam na okamžik};  \dicPhraseIS{stoppa fyrir e‑m} \dicDirectTranslationCS{zastavit (komu)}  \textbf{2.} \dicSynonym{bólstra} \dicDirectTranslationCS{vycpat, vycpávat, (vy)polstrovat, čalounit} \dicExampleIS{stoppa stól} \dicExampleCS{vypolstrovat židli};  \dicPhraseIS{stoppa upp e‑ð} \dicDirectTranslationCS{vycpat (co) (ptáka ap.)}  \textbf{3.} \dicDirectTranslationCS{(za)látat, (za)šít (ponožky ap.)} \dicExampleIS{stoppa í sokkana} \dicExampleCS{zalátat ponožky}
\dicEntry[stoppistöð] \textls[10]{\dicTerm{stoppi··stöð}\dicTerm{, stoppustöð} \dicIPA{{s}\-{\textsubring{d}}\-{\textopeno}\-{h}\-{\textsubring{b}}\-{\textsci}\-{s}\-{\textsubring{d}}\-{\oe}\-{\texttheta}\-} \dicPos{f}[6] \dicFlx{(‑var, ‑var)}[1]} \dicSynonym{biðstöð} \dicDirectTranslationCS{zastávka, stanice}
\dicEntry[stoppunál] \dicTerm{stoppu··nál} \dicIPA{{s}{\textsubring{d}}{\textopeno}{h}{\textsubring{b}}{\textscy}{n}{au}{\textsubring{l}}} \dicPos{f}[4] \dicFlx{(‑ar, ‑ar)}[1] \dicDirectTranslationCS{látací jehla}
\dicEntry[stoppustöð] \dicTerm{stoppu··stöð} \dicIPA{{s}{\textsubring{d}}{\textopeno}{h}{\textsubring{b}}{\textscy}{s}{\textsubring{d}}{\oe}{\texttheta}} \dicPos{f}[6] \dicFlx{(‑var, ‑var)}[1] \dicLink{stoppistöð}
\dicEntry[stopull] \dicTerm{stopull} \dicIPA{{s}{\textsubring{d}}{\textopeno}{\textlengthmark}{\textsubring{b}}{\textscy}{\textsubring{d}}{\textsubring{l}}} \dicPos{adj}[8]\dicFlx{}[-4] \dicDirectTranslationCS{nestálý, nepravidelný} \dicExampleIS{stopul atvinna} \dicExampleCS{nestálé zaměstnání}
\dicEntry[stopult] \dicTerm{stopul|t} \dicIPA{{s}{\textsubring{d}}{\textopeno}{\textlengthmark}{\textsubring{b}}{\textscy}{\textsubring{l}}{\textsubring{d}}} \dicPos{adv} \dicFlx{(comp ‑ar, sup ‑ast)} \dicDirectTranslationCS{nepravidelně, sporadicky}
\dicEntry[storka] \dicTerm{stork|a\smash{\textsuperscript{1}}} \dicIPA{{s}{\textsubring{d}}{\textopeno}{\textsubring{r}}{\r{g}}{a}} \dicPos{f}[1] \dicFlx{(‑u, ‑ur)}[19] \dicDirectTranslationCS{sraženina}
\dicEntry[storka] \dicTerm{stork|a\smash{\textsuperscript{2}}} \dicIPA{{s}{\textsubring{d}}{\textopeno}{\textsubring{r}}{\r{g}}{a}} \dicPos{v}[1] \dicFlx{(‑aði)}[1] \dicFlx{dat} \dicSynonym{erta\smash{\textsuperscript{2}}} \dicDirectTranslationCS{(vy)provokovat, ponoukat} \dicExampleIS{storka e‑m með e‑u} \dicExampleCS{provokovat (koho čím)}
\dicEntry[storkandi] \dicTerm{stork··andi} \dicIPA{{s}{\textsubring{d}}{\textopeno}{\textsubring{r}}{\r{g}}{a}{n}{\textsubring{d}}{\textsci}} \dicPos{adj}[13] \dicFlx{indecl}[1] \dicDirectTranslationCS{provokativní, provokační}
\dicEntry[storkinn] \dicTerm{storkinn} \dicIPA{{s}{\textsubring{d}}{\textopeno}{\textsubring{r}}{\r{\textObardotlessj}}{\textsci}{\textsubring{n}}} \dicPos{adj}[6]\dicFlx{}[-6] \dicDirectTranslationCS{sražený (krev ap.), ztuhlý (láva ap.)} \dicExampleIS{storkið blóð} \dicExampleCS{sražená krev}
\dicEntry[storkna] \dicTerm{storkn|a} \dicIPA{{s}{\textsubring{d}}{\textopeno}{\textsubring{r}}{\r{g}}{n}{a}} \dicPos{v}[1] \dicFlx{(‑aði)}[1] \dicDirectTranslationCS{srazit se, srážet se, (z)tuhnout} \dicExampleIS{Tólgin storknaði á fiskinum.} \dicExampleCS{V~rybě se srazil lůj.}
\dicEntry[storknaður] \dicTerm{storkn··|aður} \dicIPA{{s}{\textsubring{d}}{\textopeno}{\textsubring{r}}{\r{g}}{n}{a}{ð}{\textscy}{\textsubring{r}}} \dicPos{adj}[3] \dicFlx{(f ‑uð)}[3] \dicDirectTranslationCS{sražený, ztuhlý (krev ap.)}
\dicEntry[storknun] \dicTerm{storkn|un} \dicIPA{{s}{\textsubring{d}}{\textopeno}{\textsubring{r}}{\r{g}}{n}{\textscy}{\textsubring{n}}} \dicPos{f}[7] \dicFlx{(‑unar)}[9] \dicDirectTranslationCS{sražení, koagulace}
\dicEntry[storkun] \dicTerm{stork|un} \dicIPA{{s}{\textsubring{d}}{\textopeno}{\textsubring{r}}{\r{g}}{\textscy}{\textsubring{n}}} \dicPos{f}[7] \dicFlx{(‑unar, ‑anir)}[8] \dicSynonym{ertni} \dicDirectTranslationCS{provokace, provokování, ponoukání} \dicExampleIS{segja e‑ð e‑m til storkunar} \dicExampleCS{říct (co komu) k~vyprovokování}
\dicEntry[storkur] \dicTerm{stork|ur} \dicIPA{{s}{\textsubring{d}}{\textopeno}{\textsubring{r}}{\r{g}}{\textscy}{\textsubring{r}}} \dicPos{m}[6] \dicFlx{(‑s, ‑ar)}[10] \dicFieldCat{zool.} \dicDirectTranslationCS{čáp} \textit{(l.~{\textLA{Ciconia}})}
\dicEntry[storma] \dicTerm{storm|a} \dicIPA{{s}{\textsubring{d}}{\textopeno}{r}{m}{a}} \dicPos{v}[1] \dicFlx{(‑aði)}[48] \dicDirectTranslationCS{vtrhnout, vrazit, vpadnout, vřítit se (do místnosti ap.)}
\dicEntry[stormasamur] \dicTerm{storma··|samur} \dicIPA{{s}{\textsubring{d}}{\textopeno}{r}{m}{a}{s}{a}{m}{\textscy}{\textsubring{r}}} \dicPos{adj}[1] \dicFlx{(f ‑söm)}[2] \textbf{1.} \dicDirectTranslationCS{bouřlivý} \dicExampleIS{stormasöm nótt} \dicExampleCS{bouřlivá noc}  \textbf{2.} \dicDirectTranslationCS{bouřlivý, buřičský}
\dicEntry[stormflóð] \dicTerm{storm··flóð} \dicIPA{{s}{\textsubring{d}}{\textopeno}{r}{m}{f}{l}{ou}{\texttheta}} \dicPos{n}[2] \dicFlx{(‑s, ‑)}[5] \dicDirectTranslationCS{přívalová povodeň}
\dicEntry[stormjárn] \dicTerm{storm··járn} \dicIPA{{s}{\textsubring{d}}{\textopeno}{r}{m}{j}{au}{r}{\textsubring{d}}{\textsubring{n}}} \dicPos{n}[2] \dicFlx{(‑s, ‑)}[5] \dicDirectTranslationCS{okenní háček\,/\addthin hák}
\dicEntry[stormmávur] \dicTerm{storm··máv|ur} \dicIPA{{s}{\textsubring{d}}{\textopeno}{r}{m}{\textlengthmark}{au}{v}{\textscy}{\textsubring{r}}} \dicPos{m}[6] \dicFlx{(‑s, ‑ar)}[19] \dicFieldCat{zool.} \dicDirectTranslationCS{racek bouřní} \textit{(l.~{\textLA{Larus canus}})}  \dicsymPhoto\ 
\dicFigure{21095.jpg}{Stormmávur}{Stormmávur - Bowman, Tim, Biolib, PD}
\dicEntry[stormsvala] \dicTerm{storm··|svala} \dicIPA{{s}{\textsubring{d}}{\textopeno}{r}{m}{s}{v}{a}{l}{a}} \dicPos{f}[1] \dicFlx{(‑svölu, ‑svölur)}[20] \dicFieldCat{zool.} \dicDirectTranslationCS{buřňáček malý} \textit{(l.~{\textLA{Hydrobates pelagicus}})}  \dicsymPhoto\ 
\dicFigure{ds_image_stormsvala_0_2.jpg}{Stormsvala}{Stormsvala - J, CC BY-SA 3.0}
\dicEntry[stormur] \dicTerm{storm|ur} \dicsymFrequent\  \dicIPA{{s}{\textsubring{d}}{\textopeno}{r}{m}{\textscy}{\textsubring{r}}} \dicPos{m}[6] \dicFlx{(‑s, ‑ar)}[4] \textbf{1.} \dicSynonym{rok} \dicDirectTranslationCS{bouřka, bouře}  \textbf{2.} \dicDirectTranslationCS{vichřice} \dicIndirectTranslationCS{(Beaufort 9) (20,8--24,4~m/s)}  \textbf{3.} \dicLangCat{přen.} \dicSynonym{ófriður} \dicDirectTranslationCS{bouřlivost};  \dicIdiom{stormur}{ \dicPhraseIS{hafa storminn í fangið}} \dicLangCat{přen.} \dicDirectTranslationCS{narazit na odpor}; { \dicPhraseIS{lognið á undan storminum}} \dicLangCat{přen.} \dicDirectTranslationCS{ticho před bouří};  \dicPhraseIS{þetta er stormur í vatnsglasi} \dicLangCat{přen.} \dicDirectTranslationCS{je to bouře ve sklenici vody} \dicIndirectTranslationCS{(o~povyku nad něčím nepodstatným, nedůležitým)}
\dicEntry[stó] \dicTerm{stó} \dicIPA{{s}{\textsubring{d}}{ou}{\textlengthmark}} \dicPos{f}[4] \dicFlx{(‑r\,/\addthin ‑ar, ‑r\,/\addthin ‑ar)}[25] \dicDirectTranslationCS{krb, krbová kamna}
\dicEntry[stóð] \dicTerm{stóð\smash{\textsuperscript{1}}} \dicIPA{{s}{\textsubring{d}}{ou}{\textlengthmark}{\texttheta}} \dicPos{n}[2] \dicFlx{(‑s, ‑)}[5] \dicDirectTranslationCS{koně, stádo koní} \dicExampleIS{reka stóðið á fjall} \dicExampleCS{hnát koně na horu}
\dicEntry[stóð] \dicTerm{stóð\smash{\textsuperscript{2}}} \dicIPA{{s}{\textsubring{d}}{ou}{\textlengthmark}{\texttheta}} \dicPos{v} \dicFlx{ind pf sg 1 pers} \dicLink{standa}
\dicEntry[stóðhestur] \dicTerm{stóð··hest|ur} \dicIPA{{s}{\textsubring{d}}{ou}{\textlengthmark}{\texttheta}{h}{\textepsilon}{s}{\textsubring{d}}{\textscy}{\textsubring{r}}} \dicPos{m}[6] \dicFlx{(‑s, ‑ar)}[4] \dicDirectTranslationCS{(plemenný) kůň}
\dicEntry[stóðum] \dicTerm{stóðum} \dicIPA{{s}{\textsubring{d}}{ou}{\textlengthmark}{ð}{\textscy}{\textsubring{m}}} \dicPos{v} \dicFlx{ind pf pl 1 pers} \dicLink{standa}
\dicEntry[stóískur] \dicTerm{stóískur} \dicIPA{{s}{\textsubring{d}}{ou}{i}{s}{\r{g}}{\textscy}{\textsubring{r}}} \dicPos{adj}[1]\dicFlx{}[-1] \textbf{1.} \dicFieldCat{filos.} \dicDirectTranslationCS{stoický}  \textbf{2.} \dicSynonym{rólegur} \dicDirectTranslationCS{stoický, klidný} \dicExampleIS{stóísk ró} \dicExampleCS{stoický klid}
\dicEntry[stóla] \dicTerm{stól|a\smash{\textsuperscript{1}}} \dicIPA{{s}{\textsubring{d}}{ou}{\textlengthmark}{l}{a}} \dicPos{f}[1] \dicFlx{(‑u, ‑ur)}[19] \dicDirectTranslationCS{štóla} \dicIndirectTranslationCS{(dlouhý, úzký pruh látky, který nosí katolický kněz při určitých obřadech)}
\dicEntry[stóla] \dicTerm{stól|a\smash{\textsuperscript{2}}} \dicIPA{{s}{\textsubring{d}}{ou}{\textlengthmark}{l}{a}} \dicPos{v}[1] \dicFlx{(‑aði)}[1] \dicPhraseIS{stóla á e‑ð\,/\addthin e‑n} \dicLangCat{hovor.} \dicDirectTranslationCS{spolehnout\,/\addthin spoléhat se na (co\,/\addthin koho)} \dicExampleIS{stóla á okkur sjálfa} \dicExampleCS{spoléhat se sami na sebe}
\dicEntry[stólalyfta] \dicTerm{stóla··lyft|a} \dicIPA{{s}{\textsubring{d}}{ou}{\textlengthmark}{l}{a}{l}{\textsci}{f}{\textsubring{d}}{a}} \dicPos{f}[1] \dicFlx{(‑u, ‑ur)}[7] \dicDirectTranslationCS{sedačkový vlek, sedačková lanovka}
\dicEntry[stólarmur] \dicTerm{stól··arm|ur} \dicIPA{{s}{\textsubring{d}}{ou}{\textlengthmark}{l}{a}{r}{m}{\textscy}{\textsubring{r}}} \dicPos{m}[6] \dicFlx{(‑s, ‑ar)}[5] \dicDirectTranslationCS{područka, loketní opěrka}
\dicEntry[stólbak] \dicTerm{stól··|bak} \dicIPA{{s}{\textsubring{d}}{ou}{l}{\textsubring{b}}{a}{\r{g}}} \dicPos{n}[2] \dicFlx{(‑baks, ‑bök)}[8] \dicDirectTranslationCS{opěradlo\,/\addthin opěrka (židle)}
\dicEntry[stóll] \dicTerm{stól|l} \dicsymFrequent\  \dicIPA{{s}{\textsubring{d}}{ou}{\textsubring{d}}{\textsubring{l}}} \dicPos{m}[6] \dicFlx{(‑s, ‑ar)}[51] \textbf{1.} \dicDirectTranslationCS{židle}  \textbf{2.} \dicDirectTranslationCS{(řečnický) pult\,/\addthin stolek};  \dicPhraseIS{stíga í stól} {\textbf{a.}} \dicDirectTranslationCS{pronést projev};  {\textbf{b.}} \dicDirectTranslationCS{pronést kázání, kázat};  \dicIdiom{stóll}{ \dicPhraseIS{setja e‑m stólinn fyrir dyrnar}} \dicLangCat{přen.} \dicDirectTranslationCS{házet (komu) klacky pod nohy};  \dicPhraseIS{steypa e‑m af stóli} \dicLangCat{přen.} \dicDirectTranslationCS{svrhnout (koho), sesadit (koho)}; { \dicPhraseIS{stinga e‑u undir stól}} \dicLangCat{přen.} \dicDirectTranslationCS{smést (co) pod stůl}
\dicEntry[stólpi] \dicTerm{stólp|i} \dicIPA{{s}{\textsubring{d}}{ou}{\textsubring{l}}{\textsubring{b}}{\textsci}} \dicPos{m}[1] \dicFlx{(‑a, ‑ar)}[1] \textbf{1.} \dicSynonym{undirstaða} \dicDirectTranslationCS{sloup, pilíř} \dicExampleIS{stólparnir undir brúnni} \dicExampleCS{pilíře pod mostem}  \textbf{2.} \dicSynonym{staur} \dicDirectTranslationCS{sloupek, kůl}
\dicEntry[stólpípa] \dicTerm{stól··píp|a} \dicIPA{{s}{\textsubring{d}}{ou}{l}{p\smash{\textsuperscript{h}}}{i}{\textsubring{b}}{a}} \dicPos{f}[1] \dicFlx{(‑u, ‑ur)}[19] \dicFieldCat{med.} \dicDirectTranslationCS{klystýr, výplach}
\dicEntry[stólræða] \dicTerm{stól··ræð|a} \dicIPA{{s}{\textsubring{d}}{ou}{l}{r}{a}{i}{ð}{a}} \dicPos{f}[1] \dicFlx{(‑u, ‑ur)}[19] \dicDirectTranslationCS{kázání} \dicExampleIS{halda stólræðu} \dicExampleCS{kázat}
\dicEntry[stór] \dicTerm{stór} \dicsymFrequent\  \dicIPA{{s}{\textsubring{d}}{ou}{\textlengthmark}{\textsubring{r}}} \dicPos{adj}[10] \dicFlx{(f ‑)}[17] \textbf{1.} \dicSynonym{mikill} \dicDirectTranslationCS{velký, veliký} \dicIndirectTranslationCS{(značný co do počtu)} \dicExampleIS{stórt fyrirtæki} \dicExampleCS{velká firma} \dicAntonym{lítill}  \textbf{2.} \dicDirectTranslationCS{velký, veliký} \dicIndirectTranslationCS{(mající značné rozměry)} \dicExampleIS{stór maður} \dicExampleCS{velký člověk};  \dicPhraseIS{stór stafur} \dicDirectTranslationCS{velké písmeno}  \textbf{3.} \dicSynonym{mikilvægur} \dicDirectTranslationCS{velký, velmi významný} \dicExampleIS{stór fundur} \dicExampleCS{velký nález};  \dicIdiom{stór}{ \dicPhraseIS{stór um sig}} \dicFlx{adj} \dicDirectTranslationCS{značný, rozsáhlý};  \dicPhraseIS{stórum augum} \dicFlx{adv} \dicLangCat{přen.} \dicDirectTranslationCS{velkýma očima (zírat ap.)};  \dicPhraseIS{taka á honum stóra sínum} \dicLangCat{přen.} \dicDirectTranslationCS{dát\,/\addthin pustit se do toho};  \dicPhraseIS{vera stór upp á sig} \dicDirectTranslationCS{být arogantní}
\dicEntry[Stórabeltið] \dicTerm{Stóra··belti|ð} \dicIPA{{s}{\textsubring{d}}{ou}{\textlengthmark}{r}{a}{\textsubring{b}}{\textepsilon}{\textsubring{l}}{\textsubring{d}}{\textsci}{\texttheta}} \dicPos{n}[2] \dicFlx{(‑sins)}[3] \dicFlx{def} \dicDirectTranslationCS{Velký Belt} \dicIndirectTranslationCS{(jedna z~úžin spojujících Baltské moře a~průliv Kattegat)}
\dicEntry[stóratá] \dicTerm{stóra··tá} \dicIPA{{s}{\textsubring{d}}{ou}{\textlengthmark}{r}{a}{t\smash{\textsuperscript{h}}}{au}} \dicPos{f}[9] \dicFlx{(stórutáar, stórutær)}[2] \dicDirectTranslationCS{palec (u~nohy)}
\dicEntry[stórbokki] \dicTerm{stór··bokk|i} \dicIPA{{s}{\textsubring{d}}{ou}{r}{\textsubring{b}}{\textopeno}{h}{\r{\textObardotlessj}}{\textsci}} \dicPos{m}[1] \dicFlx{(‑a, ‑ar)}[1] \dicSynonym{valdsmaður} \dicDirectTranslationCS{magnát(ka), velkopodnikatel(ka)}
\dicEntry[stórborg] \dicTerm{stór··borg} \dicIPA{{s}{\textsubring{d}}{ou}{r}{\textsubring{b}}{\textopeno}{r}{\r{g}}} \dicPos{f}[7] \dicFlx{(‑ar, ‑ir)}[1] \dicDirectTranslationCS{velkoměsto, metropole}
\dicEntry[stórborgarlíf] \dicTerm{stór·borgar··líf} \dicIPA{{s}{\textsubring{d}}{ou}{r}{\textsubring{b}}{\textopeno}{r}{\r{g}}{a}{r}{l}{i}{f}} \dicPos{n}[2] \dicFlx{(‑s)}[2] \dicDirectTranslationCS{velkoměstský život}
\dicEntry[stórbóndi] \dicTerm{stór··|bóndi} \dicIPA{{s}{\textsubring{d}}{ou}{r}{\textsubring{b}}{ou}{n}{\textsubring{d}}{\textsci}} \dicPos{m}[2] \dicFlx{(‑bónda, ‑bændur)}[4] \dicDirectTranslationCS{velkostatkář(ka)}
\dicEntry[stórbrotinn] \dicTerm{stór··brotinn} \dicIPA{{s}{\textsubring{d}}{ou}{r}{\textsubring{b}}{r}{\textopeno}{\textsubring{d}}{\textsci}{\textsubring{n}}} \dicPos{adj}[6]\dicFlx{}[-2] \dicDirectTranslationCS{velkolepý, nádherný, grandiózní} \dicExampleIS{stórbrotið náttúrusmíði} \dicExampleCS{velkolepé dílo přírody}
\dicEntry[stóreflis] \dicTerm{stór··eflis} \dicIPA{{s}{\textsubring{d}}{ou}{\textlengthmark}{r}{\textepsilon}{\textsubring{b}}{l}{\textsci}{s}} \dicPos{adj}[13] \dicFlx{indecl}[1] \dicSynonym{stór} \dicDirectTranslationCS{ohromný, obrovský, gigantický} \dicExampleIS{stóreflis fjall} \dicExampleCS{ohromná hora}

\dicEntry[stóreignamaður] \textls[15]{\dicTerm{stór··eigna·|maður} \dicIPA{{s}{\textsubring{d}}{ou}{\textlengthmark}{r}{ei}{\r{g}}{n}{a}{m}{a}{ð}{\textscy}{\textsubring{r}}} \dicPos{m}[13] \dicFlx{(‑manns, ‑menn)}[2]} \dicDirectTranslationCS{majitel(ka) nemovitostí, magnát(ka)}

\dicEntry[stóreygður] \dicTerm{stór··eygður}\dicTerm{, stóreygur} \dicIPA{{s}\-{\textsubring{d}}\-{ou}\-{\textlengthmark}\-{r}\-{ei}\-{\textbabygamma}\-{ð}\-{\textscy}\-{\textsubring{r}}\-} \dicPos{adj}[2]\dicFlx{}[-1] \dicDirectTranslationCS{(jsoucí)  s~očima dokořán, (jsoucí) s~vyvalenýma očima}
\dicEntry[stóreygur] \dicTerm{stór··eygur} \dicIPA{{s}{\textsubring{d}}{ou}{\textlengthmark}{r}{ei}{\textbabygamma}{\textscy}{\textsubring{r}}} \dicPos{adj}[1]\dicFlx{}[-1] \dicLink{stóreygður}
\dicEntry[stórfelldur] \dicTerm{stór··felldur} \dicIPA{{s}{\textsubring{d}}{ou}{\textsubring{r}}{f}{\textepsilon}{l}{\textsubring{d}}{\textscy}{\textsubring{r}}} \dicPos{adj}[2]\dicFlx{}[-14] \textbf{1.} \dicSynonym{ákafur} \dicDirectTranslationCS{silný, prudký} \dicExampleIS{stórfelld rigning} \dicExampleCS{silný déšť}  \textbf{2.} \dicSynonym{mikill} \dicDirectTranslationCS{velký, rozsáhlý} \dicExampleIS{stórfelldur þjófnaður} \dicExampleCS{rozsáhlá loupež}  \textbf{3.} \dicSynonym{mikilvægur} \dicDirectTranslationCS{velký, významný} \dicExampleIS{stórfelld áhrif} \dicExampleCS{významný vliv}
\dicEntry[stórfengilegur] \dicTerm{stór·fengi··legur} \dicIPA{{s}{\textsubring{d}}{ou}{\textsubring{r}}{f}{ei}{\textltailn}{\r{\textObardotlessj}}{\textsci}{l}{\textepsilon}{\textbabygamma}{\textscy}{\textsubring{r}}} \dicPos{adj}[1]\dicFlx{}[-8] \dicLink{stórfenglegur}
\dicEntry[stórfenglega] \dicTerm{stór··feng·lega} \dicIPA{{s}{\textsubring{d}}{ou}{\textsubring{r}}{f}{ei}{\ng}{\r{g}}{l}{\textepsilon}{\textbabygamma}{a}} \dicPos{adv} \dicDirectTranslationCS{velkolepě, úchvat\-ně}
\dicEntry[stórfenglegur] \dicTerm{stór··feng·legur}\dicTerm{, stórfengilegur} \dicIPA{{s}\-{\textsubring{d}}\-{ou}\-{\textsubring{r}}\-{f}\-{ei}\-{\ng}\-{\r{g}}\-{l}\-{\textepsilon}\-{\textbabygamma}\-{\textscy}\-{\textsubring{r}}\-} \dicPos{adj}[1]\dicFlx{}[-8] \dicSynonym{stórbrotinn} \dicDirectTranslationCS{velkolepý, úchvatný} \dicExampleIS{stórfenglegar hugsjónir} \dicExampleCS{velkolepé ideály}
\dicEntry[stórfengleiki] \dicTerm{stór··feng·leik|i} \dicIPA{{s}{\textsubring{d}}{ou}{\textsubring{r}}{f}{ei}{\ng}{\r{g}}{l}{ei}{\r{\textObardotlessj}}{\textsci}} \dicPos{m}[1] \dicFlx{(‑a)}[3] \dicDirectTranslationCS{velkolepost, úchvatnost}
\dicEntry[stórfé] \dicTerm{stór··fé} \dicIPA{{s}{\textsubring{d}}{ou}{\textsubring{r}}{f}{j}{\textepsilon}} \dicPos{n}[3] \dicFlx{(‑fjár)}[5] \dicDirectTranslationCS{velké peníze, hodně peněz}
\dicEntry[stórfínn] \dicTerm{stór··fínn} \dicIPA{{s}{\textsubring{d}}{ou}{\textsubring{r}}{f}{i}{\textsubring{d}}{\textsubring{n}}} \dicPos{adj}[7]\dicFlx{}[-2] \dicDirectTranslationCS{vynikající, super}
\dicEntry[stórfylki] \dicTerm{stór··fylki} \dicIPA{{s}{\textsubring{d}}{ou}{\textsubring{r}}{f}{\textsci}{\textsubring{l}}{\r{\textObardotlessj}}{\textsci}} \dicPos{n}[2] \dicFlx{(‑s, ‑)}[16] \dicFieldCat{voj.} \dicDirectTranslationCS{brigáda}
\dicEntry[stórfyrirtæki] \dicTerm{stór··fyrir·tæki} \dicIPA{{s}{\textsubring{d}}{ou}{\textsubring{r}}{f}{\textsci}{r}{\textsci}{\textsubring{r}}{t\smash{\textsuperscript{h}}}{a}{i}{\r{\textObardotlessj}}{\textsci}} \dicPos{n}[2] \dicFlx{(‑s, ‑)}[16] \dicDirectTranslationCS{velká firma, korporace}
\dicEntry[stórgrýti] \dicTerm{stór··grýti} \dicIPA{{s}{\textsubring{d}}{ou}{r}{\r{g}}{r}{i}{\textsubring{d}}{\textsci}} \dicPos{n}[2] \dicFlx{(‑s)}[20] \dicSynonym*{klettar} \dicDirectTranslationCS{velké kameny, balvany}
\dicEntry[stórhuga] \dicTerm{stór··huga} \dicIPA{{s}{\textsubring{d}}{ou}{\textlengthmark}{\textsubring{r}}{h}{\textscy}{\textbabygamma}{a}} \dicPos{adj}[13] \dicFlx{indecl}[1] \dicDirectTranslationCS{ambiciózní, odvážný, vizionářský (záměr ap.)}
\dicEntry[stórhætta] \dicTerm{stór··hætt|a} \dicIPA{{s}{\textsubring{d}}{ou}{\textlengthmark}{\textsubring{r}}{h}{a}{i}{h}{\textsubring{d}}{a}} \dicPos{f}[1] \dicFlx{(‑u, ‑ur)}[7] \dicDirectTranslationCS{velké nebezpečí}
\dicEntry[stóriburkni] \dicTerm{stóri··burkni} \dicIPA{{s}{\textsubring{d}}{ou}{\textlengthmark}{r}{\textsci}{\textsubring{b}}{\textscy}{\textsubring{r}}{\r{g}}{n}{\textsci}} \dicPos{m}[1] \dicFlx{(stóraburkna, stóraburknar)}[6] \dicFieldCat{bot.} \dicDirectTranslationCS{kapraď samec} \textit{(l.~{\textLA{Dryopteris filix-mas}})}
\dicEntry[stóriðja] \dicTerm{stór··iðj|a} \dicIPA{{s}{\textsubring{d}}{ou}{\textlengthmark}{r}{\textsci}{ð}{j}{a}} \dicPos{f}[1] \dicFlx{(‑u)}[5] \dicDirectTranslationCS{velkoprůmysl}
\dicEntry[stórkarlalega] \textls[10]{\dicTerm{stór·karla··lega} \dicIPA{{s}{\textsubring{d}}{ou}{\textsubring{r}}{k\smash{\textsuperscript{h}}}{a}{r}{\textsubring{d}}{l}{a}{l}{\textepsilon}{\textbabygamma}{a}} \dicPos{adv}} \dicDirectTranslationCS{hrubě, vulgárně, sprostě}
\dicEntry[stórkarlalegur] \dicTerm{stór··karla·legur} \dicIPA{{s}{\textsubring{d}}{ou}{\textsubring{r}}{k\smash{\textsuperscript{h}}}{a}{r}{\textsubring{d}}{l}{a}{l}{\textepsilon}{\textbabygamma}{\textscy}{\textsubring{r}}} \dicPos{adj}[1]\dicFlx{}[-8] \dicDirectTranslationCS{hrubý, vulgární, sprostý} \dicExampleIS{stórkarlalegt orðbragð} \dicExampleCS{hrubý výraz}
\begin{xtolerant}{}{1pt}
\dicEntry[stórkaupmaður] \dicTerm{stór··kaup·|maður}\addthinSS\dicIPA{{s}{\textsubring{d}}{ou}{\textsubring{r}}{k\smash{\textsuperscript{h}}}{\oe i}{\textsubring{b}}{m}{a}{ð}{\textscy}{\textsubring{r}}}\addthinS\dicPos{m}[1\hspace{-0.06667em}3]\addthinS\dicFlx{(‑manns,\addthinSS‑menn)}[2] \dicDirectTranslationCS{velkoobchodník, velkoobchodnice}
\end{xtolerant}
\dicEntry[stórkostlega] \textls[10]{\dicTerm{stór··kost·lega} \dicIPA{{s}{\textsubring{d}}{ou}{\textsubring{r}}{k\smash{\textsuperscript{h}}}{\textopeno}{s}{\textsubring{d}}{l}{\textepsilon}{\textbabygamma}{a}} \dicPos{adv} \textbf{1.} \dicSynonym{frábærlega} \dicDirectTranslationCS{úžasně, skvěle, fantasticky}  \textbf{2.} \dicSynonym{mjög} \dicDirectTranslationCS{obrovsky, enormně, drasticky}}
\dicEntry[stórkostlegur] \dicTerm{stór··kost·legur} \dicsymFrequent\  \dicIPA{{s}{\textsubring{d}}{ou}{\textsubring{r}}{k\smash{\textsuperscript{h}}}{\textopeno}{s}{\textsubring{d}}{l}{\textepsilon}{\textbabygamma}{\textscy}{\textsubring{r}}} \dicPos{adj}[1]\dicFlx{}[-8] \textbf{1.} \dicSynonym{mikilfenglegur} \dicDirectTranslationCS{velkolepý, úchvatný, fantastický}  \textbf{2.} \dicSynonym{gífurlegur} \dicDirectTranslationCS{ohromný, enormní, kolosální}  \textbf{3.} \dicSynonym{frábær} \dicDirectTranslationCS{výborný, vynikající, skvělý} \dicExampleIS{stórkostleg skemmtun} \dicExampleCS{vynikající zábava}
\dicEntry[stórlax] \dicTerm{stór··lax} \dicIPA{{s}{\textsubring{d}}{ou}{r}{l}{a}{x}{s}} \dicPos{m}[4] \dicFlx{(‑, ‑ar)}[29] \textbf{1.} \dicDirectTranslationCS{velký losos}  \textbf{2.} \dicLangCat{přen.} \dicSynonym{stórbokki} \dicDirectTranslationCS{matador, veličina, legenda, velké zvíře} \dicIndirectTranslationCS{(význačná osoba, dlouho se někde vyskytující)}
\dicEntry[stórlátur] \dicTerm{stór··látur} \dicIPA{{s}{\textsubring{d}}{ou}{r}{l}{au}{\textsubring{d}}{\textscy}{\textsubring{r}}} \dicPos{adj}[1]\dicFlx{}[-1] \dicDirectTranslationCS{pyšný, domýšlivý, nadutý}
\dicEntry[stórlega] \dicTerm{stór··lega} \dicIPA{{s}{\textsubring{d}}{ou}{r}{l}{\textepsilon}{\textbabygamma}{a}} \dicPos{adv} \dicDirectTranslationCS{velmi, silně, značně}
\dicEntry[stórlyndi] \dicTerm{stór··lyndi} \dicIPA{{s}{\textsubring{d}}{ou}{r}{l}{\textsci}{n}{\textsubring{d}}{\textsci}} \dicPos{n}[2] \dicFlx{(‑s)}[20] \textbf{1.} \dicSynonym*{mikið skap} \dicDirectTranslationCS{výbušnost, horkokrevnost, impulzivnost}  \textbf{2.} \dicSynonym{göfuglyndi} \dicDirectTranslationCS{velkodušnost, velkorysost}
\dicEntry[stórlyndur] \dicTerm{stór··lyndur} \dicIPA{{s}{\textsubring{d}}{ou}{r}{l}{\textsci}{n}{\textsubring{d}}{\textscy}{\textsubring{r}}} \dicPos{adj}[2]\dicFlx{}[-14] \textbf{1.} \dicSynonym*{ofsafenginn í skapi} \dicDirectTranslationCS{výbušný, horkokrevný, impulzivní}  \textbf{2.} \dicSynonym*{stórmannlegur í lund} \dicDirectTranslationCS{velkodušný, velkorysý}
\dicEntry[stórmannlega] \dicTerm{stór··mann·lega} \dicIPA{{s}{\textsubring{d}}{ou}{r}{m}{a}{n}{l}{\textepsilon}{\textbabygamma}{a}} \dicPos{adv} \dicDirectTranslationCS{ušlechtile, velkomyslně, důstojně}
\dicEntry[stórmannlegur] \dicTerm{stór··mann·legur} \dicIPA{{s}{\textsubring{d}}{ou}{r}{m}{a}{n}{l}{\textepsilon}{\textbabygamma}{\textscy}{\textsubring{r}}} \dicPos{adj}[1]\dicFlx{}[-8] \dicDirectTranslationCS{ušlechtilý, velkomyslný, důstojný}
\begin{xtolerant}{}{1pt}
\dicEntry[stórmarkaður] \dicTerm{stór··mark·að|ur} \dicIPA{{s}{\textsubring{d}}{ou}{r}{m}{a}{\textsubring{r}}{\r{g}}{a}{ð}{\textscy}{\textsubring{r}}} \dicPos{m}[10] \dicFlx{(‑ar\,/\addthin ‑s, ‑ir)}[31] \dicDirectTranslationCS{supermarket, hypermarket, obchodní dům}
\end{xtolerant}
\dicEntry[stórmál] \dicTerm{stór··mál} \dicIPA{{s}{\textsubring{d}}{ou}{r}{m}{au}{\textsubring{l}}} \dicPos{n}[2] \dicFlx{(‑s, ‑)}[5] \dicDirectTranslationCS{významná otázka\,/\addthin záležitost}
\dicEntry[stórmeistari] \dicTerm{stór··meist·ar|i} \dicIPA{{s}{\textsubring{d}}{ou}{r}{m}{ei}{s}{\textsubring{d}}{a}{r}{\textsci}} \dicPos{m}[1] \dicFlx{(‑a, ‑ar)}[13] \textbf{1.} \dicDirectTranslationCS{velmistr(yně)} \dicIndirectTranslationCS{(nejvyšší hodnostář)}  \textbf{2.} \dicDirectTranslationCS{velmistr(yně)} \dicIndirectTranslationCS{(mezinárodní titul)} \dicExampleIS{stórmeistari í skák} \dicExampleCS{šachový velmistr}
\dicEntry[stórmenni] \dicTerm{stór··menni} \dicIPA{{s}{\textsubring{d}}{ou}{r}{m}{\textepsilon}{n}{\textsci}} \dicPos{n}[2] \dicFlx{(‑s, ‑)}[14] \textbf{1.} \dicSynonym{tignarmaður} \dicDirectTranslationCS{důležitá\,/\addthin velká osobnost, významný člověk}  \textbf{2.} \dicSynonym*{afburðamaður} \dicDirectTranslationCS{průkopník, průkopnice}
\dicEntry[stórmennska] \dicTerm{stór··mennsk|a} \dicIPA{{s}{\textsubring{d}}{ou}{r}{m}{\textepsilon}{n}{s}{\r{g}}{a}} \dicPos{f}[1] \dicFlx{(‑u)}[5] \textbf{1.} \dicSynonym{dramb} \dicDirectTranslationCS{arogance, domýšlivost}  \textbf{2.} \dicSynonym{örlæti} \dicDirectTranslationCS{štědrost, velkorysost}
\dicEntry[stórmennskubrjálæði] \textls[15]{\dicTerm{stór·mennsku··brjál·æði} \dicIPA{{s}\-{\textsubring{d}}\-{ou}\-{r}\-{m}\-{\textepsilon}\-{n}\-{s}\-{\r{g}}\-{\textscy}\-{\textsubring{b}}\-{r}\-{j}\-{au}\-{l}\-{a}\-{i}\-{ð}\-{\textsci}\-} \dicPos{n}[2]} \dicFlx{(‑s)}[20] \dicDirectTranslationCS{velikášství, megalomanie}
\dicEntry[stórmerki] \dicTerm{stór··merki} \dicIPA{{s}{\textsubring{d}}{ou}{r}{m}{\textepsilon}{\textsubring{r}}{\r{\textObardotlessj}}{\textsci}} \dicPos{n}[2] \dicFlx{(‑s, ‑)}[16] \dicSynonym{jarteikn} \dicDirectTranslationCS{znamení, předzvěst}
\dicEntry[stórmerkilegur] \textls[15]{\dicTerm{stór··merki·legur} \dicIPA{{s}{\textsubring{d}}{ou}{r}{m}{\textepsilon}{\textsubring{r}}{\r{\textObardotlessj}}{\textsci}{l}{\textepsilon}{\textbabygamma}{\textscy}{\textsubring{r}}} \dicPos{adj}[1]\dicFlx{}[-6] \dicDirectTranslationCS{velmi významný\,/\addthin důležitý}}
\dicEntry[stórmynd] \dicTerm{stór··mynd} \dicIPA{{s}{\textsubring{d}}{ou}{r}{m}{\textsci}{n}{\textsubring{d}}} \dicPos{f}[7] \dicFlx{(‑ar, ‑ir)}[1] \dicDirectTranslationCS{velkofilm} \dicExampleIS{söguleg stórmynd} \dicExampleCS{historický velkofilm}
\dicEntry[stórorður] \dicTerm{stór··orður} \dicIPA{{s}{\textsubring{d}}{ou}{\textlengthmark}{r}{\textopeno}{r}{ð}{\textscy}{\textsubring{r}}} \dicPos{adj}[2]\dicFlx{}[-1] \dicDirectTranslationCS{používající velká slova}
\dicEntry[stórræði] \dicTerm{stór··ræði} \dicIPA{{s}{\textsubring{d}}{ou}{r}{\textlengthmark}{a}{i}{ð}{\textsci}} \dicPos{n}[2] \dicFlx{(‑s, ‑)}[14] \dicSynonym*{stór fyrirætlun} \dicDirectTranslationCS{velký počin\,/\addthin projekt, velké dílo\,/\addthin předsevzetí} \dicExampleIS{hafa stórræði í hyggju} \dicExampleCS{mít v~plánu velké dílo};  \dicPhraseIS{standa í stórræðum} \dicDirectTranslationCS{účastnit se velkého předsevzetí}
\dicEntry[stórsigla] \dicTerm{stór··sigl|a} \dicIPA{{s}{\textsubring{d}}{ou}{\textsubring{r}}{s}{\textsci}{\r{g}}{l}{a}} \dicPos{f}[1] \dicFlx{(‑u, ‑ur)}[13] \dicFieldCat{nám.} \dicDirectTranslationCS{hlavní stěžeň}
\dicEntry[stórskorinn] \dicTerm{stór··skorinn} \dicIPA{{s}{\textsubring{d}}{ou}{\textsubring{r}}{s}{\r{g}}{\textopeno}{r}{\textsci}{\textsubring{n}}} \dicPos{adj}[6]\dicFlx{}[-2] \textbf{1.} \dicSynonym{grófur} \dicDirectTranslationCS{hrubý, drsný} \dicExampleIS{stórskorið tóbak} \dicExampleCS{hrubý tabák}  \textbf{2.} \dicDirectTranslationCS{velkolepý, úchvatný}
\dicEntry[stórslys] \dicTerm{stór··slys} \dicIPA{{s}{\textsubring{d}}{ou}{\textsubring{r}}{s}{\textsubring{d}}{l}{\textsci}{s}} \dicPos{n}[2] \dicFlx{(‑s, ‑)}[5] \dicDirectTranslationCS{katastrofa, pohroma, tragédie}
\dicEntry[stórstígur] \dicTerm{stór··stígur} \dicIPA{{s}{\textsubring{d}}{ou}{\textsubring{r}}{s}{\textsubring{d}}{i}{\textbabygamma}{\textscy}{\textsubring{r}}} \dicPos{adj}[1]\dicFlx{}[-1] \textbf{1.} \dicDirectTranslationCS{mající dlouhý krok} \dicExampleIS{e‑r er stórstígur} \dicExampleCS{(kdo) má dlouhý krok}  \textbf{2.} \dicDirectTranslationCS{překotný, chvatný, rychlý} \dicExampleIS{stórstígar framfarir} \dicExampleCS{překotný pokrok}
\dicEntry[stórstraumur] \dicTerm{stór··straum|ur} \dicIPA{{s}{\textsubring{d}}{ou}{\textsubring{r}}{s}{\textsubring{d}}{r}{\oe i}{m}{\textscy}{\textsubring{r}}} \dicPos{m}[6] \dicFlx{(‑s, ‑ar)}[22] \dicDirectTranslationCS{skočný příliv}
\dicEntry[stórsvig] \dicTerm{stór··svig} \dicIPA{{s}{\textsubring{d}}{ou}{\textsubring{r}}{s}{v}{\textsci}{x}} \dicPos{n}[2] \dicFlx{(‑s)}[2] \dicFieldCat{sport.} \dicDirectTranslationCS{obří slalom}
\dicEntry[stórtækur] \dicTerm{stór··tækur} \dicIPA{{s}{\textsubring{d}}{ou}{\textsubring{r}}{t\smash{\textsuperscript{h}}}{a}{i}{\r{g}}{\textscy}{\textsubring{r}}} \dicPos{adj}[1]\dicFlx{}[-1] \textbf{1.} \dicDirectTranslationCS{energický, podnikavý}  \textbf{2.} \dicDirectTranslationCS{mocný, výkonný (stroj ap.)}
\dicEntry[stórtölva] \dicTerm{stór··tölv|a} \dicIPA{{s}{\textsubring{d}}{ou}{\textsubring{r}}{t\smash{\textsuperscript{h}}}{\oe}{l}{v}{a}} \dicPos{f}[1] \dicFlx{(‑u, ‑ur)}[7] \dicFieldCat{poč.} \dicDirectTranslationCS{sálový počítač}
\dicEntry[stórutáar] \dicTerm{stóru··táar} \dicIPA{{s}{\textsubring{d}}{ou}{\textlengthmark}{r}{\textscy}{\textsubring{d}}{au}{a}{\textsubring{r}}} \dicPos{f} \dicFlx{sg gen} \dicLink{stóratá}
\dicEntry[stórutær] \dicTerm{stóru··tær} \dicIPA{{s}{\textsubring{d}}{ou}{\textlengthmark}{r}{\textscy}{\textsubring{d}}{a}{i}{\textsubring{r}}} \dicPos{f} \dicFlx{pl nom} \dicLink{stóratá}
\dicEntry[stórvaxinn] \dicTerm{stór··vaxinn} \dicIPA{{s}{\textsubring{d}}{ou}{r}{v}{a}{x}{s}{\textsci}{\textsubring{n}}} \dicPos{adj}[6]\dicFlx{}[-3] \dicDirectTranslationCS{velký, veliký} \dicIndirectTranslationCS{(vzrůstem)} \dicExampleIS{Hesturinn er stórvaxinn.} \dicExampleCS{Kůň je velký.}
\dicEntry[stórveldi] \dicTerm{stór··veldi} \dicIPA{{s}{\textsubring{d}}{ou}{r}{v}{\textepsilon}{l}{\textsubring{d}}{\textsci}} \dicPos{n}[2] \dicFlx{(‑s, ‑)}[14] \dicDirectTranslationCS{mocnost, velmoc}
\dicEntry[stórverslun] \dicTerm{stór··versl|un} \dicIPA{{s}{\textsubring{d}}{ou}{r}{v}{\textepsilon}{\textsubring{r}}{s}{\textsubring{d}}{l}{\textscy}{\textsubring{n}}} \dicPos{f}[7] \dicFlx{(‑unar, ‑anir)}[8] \dicDirectTranslationCS{velkoobchod, supermarket}
\dicEntry[stórvirki] \dicTerm{stór··virki} \dicIPA{{s}{\textsubring{d}}{ou}{r}{v}{\textsci}{\textsubring{r}}{\r{\textObardotlessj}}{\textsci}} \dicPos{n}[2] \dicFlx{(‑s, ‑)}[16] \dicSynonym*{mikið verk} \dicDirectTranslationCS{velký čin\,/\addthin počin} \dicExampleIS{vinna stórvirki} \dicExampleCS{udělat velký čin}
\dicEntry[stórvirkur] \dicTerm{stór··virkur} \dicIPA{{s}{\textsubring{d}}{ou}{r}{v}{\textsci}{\textsubring{r}}{\r{g}}{\textscy}{\textsubring{r}}} \dicPos{adj}[1]\dicFlx{}[-1] \dicDirectTranslationCS{produktivní, výkonný}
\dicEntry[stórvægilega] \textls[15]{\dicTerm{stór··vægi·lega} \dicIPA{{s}{\textsubring{d}}{ou}{r}{v}{a}{i}{j}{\textsci}{l}{\textepsilon}{\textbabygamma}{a}} \dicPos{adv} \dicDirectTranslationCS{významně, zásadně}}
\dicEntry[stórvægilegur] \dicTerm{stór··vægi·legur} \dicIPA{{s}{\textsubring{d}}{ou}{r}{v}{a}{i}{j}{\textsci}{l}{\textepsilon}{\textbabygamma}{\textscy}{\textsubring{r}}} \dicPos{adj}[1]\dicFlx{}[-8] \dicSynonym{mikilvægur} \dicDirectTranslationCS{důležitý, významný, zásadní}
\dicEntry[stóryrði] \dicTerm{stór··yrði} \dicIPA{{s}{\textsubring{d}}{ou}{\textlengthmark}{r}{\textsci}{r}{ð}{\textsci}} \dicPos{n}[2] \dicFlx{(‑s, ‑)}[14] \dicSynonym{skammir} \dicDirectTranslationCS{silný výraz, nadávka} \dicExampleIS{kasta fram stóryrðum} \dicExampleCS{používat silné výrazy}
\dicEntry[stórþjófnaður] \dicTerm{stór··þjóf·nað|ur} \dicIPA{{s}{\textsubring{d}}{ou}{\textsubring{r}}{\texttheta}{j}{ou}{\textsubring{b}}{n}{a}{ð}{\textscy}{\textsubring{r}}} \dicPos{m}[10] \dicFlx{(‑ar, ‑ir)}[8] \dicDirectTranslationCS{velká krádež}
\dicEntry[stóuspeki] \dicTerm{stóu··spek|i} \dicIPA{{s}{\textsubring{d}}{ou}{\textlengthmark}{\textscy}{s}{\textsubring{b}}{\textepsilon}{\r{\textObardotlessj}}{\textsci}} \dicPos{f}[3] \dicFlx{(‑i)}[3] \dicFieldCat{filos.} \dicDirectTranslationCS{stoicismus}
\dicEntry[straff] \dicTerm{straff} \dicIPA{{s}{\textsubring{d}}{r}{a}{f}{\textlengthmark}} \dicPos{n}[2] \dicFlx{(‑s, ströff)}[8] \textbf{1.} \dicLangCat{zast.} \dicDirectTranslationCS{trest}  \textbf{2.} \dicLangCat{hovor.} \dicDirectTranslationCS{karanténa, zákaz vstupu}
\dicEntry[stramma] \dicTerm{stramm|a} \dicIPA{{s}{\textsubring{d}}{r}{a}{m}{\textlengthmark}{a}} \dicPos{v}[1] \dicFlx{(‑aði)}[13] \dicFlx{acc} \dicPhraseIS{stramma sig} \dicDirectTranslationCS{vypínat se, dělat se větším a~silnějším};  \dicIdiom{stramma}[af]{ \dicPhraseIS{stramma sig af}} \dicDirectTranslationCS{dát se dohromady, dodat si odvahy};  \dicIdiom{stramma}[upp]{ \dicPhraseIS{stramma e‑n upp}} \dicDirectTranslationCS{povzbuzovat (koho), dodávat (komu) sílu}
\dicEntry[strand] \dicTerm{strand\smash{\textsuperscript{1}}} \dicIPA{{s}{\textsubring{d}}{r}{a}{n}{\textsubring{d}}} \dicPos{n}[2] \dicFlx{(‑s, strönd)}[8] \textbf{1.} \dicSynonym*{það að stranda} \dicDirectTranslationCS{uváznutí na mělčině};  \dicPhraseIS{e‑að siglir í strand} \dicLangCat{přen.} \dicDirectTranslationCS{(co) (u)vázne, (co) se dostává do mrtvého bodu};  \dicPhraseIS{lenda í strandi með e‑ð} \dicLangCat{přen.} \dicDirectTranslationCS{dostat se s~(čím) do mrtvého bodu, zůstat vězet s~(čím)} \dicExampleIS{lenda í strandi með verkefnið} \dicExampleCS{dostat se s~úkolem do mrtvého bodu}  \textbf{2.} \dicSynonym*{strandað skip} \dicDirectTranslationCS{loď uvázlá na mělčině}
\dicEntry[strand] \dicTerm{strand\smash{\textsuperscript{2}}} \dicIPA{{s}{\textsubring{d}}{r}{a}{n}{\textsubring{d}}} \dicPos{adj}[13] \dicFlx{indecl}[1] \dicPhraseIS{e‑r er strand (með e‑ð)} \dicDirectTranslationCS{(kdo) uvázne (s~(čím)) (prací ap.)}
\dicEntry[stranda] \dicTerm{strand|a} \dicsymFrequent\  \dicIPA{{s}{\textsubring{d}}{r}{a}{n}{\textsubring{d}}{a}} \dicPos{v}[1] \dicFlx{(‑aði)}[43] \textbf{1.} \dicDirectTranslationCS{uváznout na mělčině} \dicExampleIS{Skipið strandaði í fjörunni.} \dicExampleCS{Loď uvázla na mělčině na pobřeží.}  \textbf{2.} \dicPhraseIS{e‑að strandar á e‑u} \dicLangCat{přen.} \dicDirectTranslationCS{(co) (u)vázne na (čem), (co) ztroskotá na (čem) (na nedostatku peněz ap.)} \dicExampleIS{Málið strandaði á peningaskorti.} \dicExampleCS{Záležitost uvázla na nedostatku peněz.}
\dicEntry[strandar] \dicTerm{strandar} \dicIPA{{s}{\textsubring{d}}{r}{a}{n}{\textsubring{d}}{a}{\textsubring{r}}} \dicPos{f} \dicFlx{sg gen} \dicLink{strönd\smash{\textsuperscript{1}}}
\dicEntry[strandgæsla] \dicTerm{strand··gæsl|a} \dicIPA{{s}{\textsubring{d}}{r}{a}{n}{\textsubring{d}}{\r{\textObardotlessj}}{a}{i}{s}{\textsubring{d}}{l}{a}} \dicPos{f}[1] \dicFlx{(‑u)}[5] \dicDirectTranslationCS{pobřežní stráž\,/\addthin hlídka}
\dicEntry[strandlengja] \dicTerm{strand··lengj|a} \dicIPA{{s}{\textsubring{d}}{r}{a}{n}{\textsubring{d}}{l}{ei}{\textltailn}{\r{\textObardotlessj}}{a}} \dicPos{f}[1] \dicFlx{(‑u, ‑ur)}[7] \dicDirectTranslationCS{pobřežní pás, pobřeží}
\dicEntry[strandrauðviður] \textls[5]{\dicTerm{strand··rauð·við|ur}\dicTerm{, strandrisafura} \dicIPA{{s}\-{\textsubring{d}}\-{r}\-{a}\-{n}\-{\textsubring{d}}\-{r}\-{\oe i}\-{ð}\-{v}\-{\textsci}\-{ð}\-{\textscy}\-{\textsubring{r}}}} \textls[15]{\dicPos{m}[10] \dicFlx{(‑ar)}[20] \dicFieldCat{bot.} \dicDirectTranslationCS{sekvoj, sekvoj vždyzelená} \textit{(l.~{\textLA{Sequoia sempervirens}})}}  \dicsymPhoto\ 
\dicFigure{ds_image_strandraudvidur_0_1.jpg}{Strandrauðviður}{Strandrauðviður - Brian Gratwicke, CC BY 2.0}
\dicEntry[strandrisafura] \dicTerm{strand··risa·fur|a} \dicIPA{{s}{\textsubring{d}}{r}{a}{n}{\textsubring{d}}{r}{\textsci}{s}{a}{f}{\textscy}{r}{a}} \dicPos{f}[1] \dicFlx{(‑u, ‑ur)}[7] \dicLink{strandrauðviður}
\dicEntry[strandríki] \dicTerm{strand··ríki} \dicIPA{{s}{\textsubring{d}}{r}{a}{n}{\textsubring{d}}{r}{i}{\r{\textObardotlessj}}{\textsci}} \dicPos{n}[2] \dicFlx{(‑s, ‑)}[16] \dicDirectTranslationCS{pobřežní stát}
\dicEntry[strandtittlingur] \dicTerm{strand··titt·ling|ur} \dicIPA{{s}{\textsubring{d}}{r}{a}{n}{\textsubring{d}}{t\smash{\textsuperscript{h}}}{\textsci}{h}{\textsubring{d}}{l}{i}{\ng}{\r{g}}{\textscy}{\textsubring{r}}} \dicPos{m}[6] \dicFlx{(‑s, ‑ar)}[8] \dicFieldCat{zool.} \dicDirectTranslationCS{linduška vodní} \textit{(l.~{\textLA{Anthus petrosus}})}  \dicsymPhoto\ 
\dicFigure{ds_image_strandtittlingur_0_2.jpg}{Strandtittlingur}{Strandtittlingur - Erik, GFDL}
\dicEntry[strangi] \dicTerm{strang|i} \dicIPA{{s}{\textsubring{d}}{r}{au}{\textltailn}{\r{\textObardotlessj}}{\textsci}} \dicPos{m}[1] \dicFlx{(‑a, ‑ar)}[8] \dicDirectTranslationCS{role, rulička (papíru ap.)}
\dicEntry[strangt] \dicTerm{strangt} \dicIPA{{s}{\textsubring{d}}{r}{au}{\r{\ng}}{\textsubring{d}}} \dicPos{adv} \dicPhraseIS{strangt (til) tekið} \dicFlx{adv} \dicDirectTranslationCS{přísně\,/\addthin striktně vzato}
\dicEntry[strangtrúarmaður] \dicTerm{strang··trúar·|maður} \dicIPA{{s}{\textsubring{d}}{r}{au}{\ng}{\r{g}}{t\smash{\textsuperscript{h}}}{r}{u}{a}{r}{m}{a}{ð}{\textscy}{\textsubring{r}}} \dicPos{m}[13] \dicFlx{(‑manns, ‑menn)}[2] \dicDirectTranslationCS{ortodoxní věřící}
\dicEntry[strangur] \dicTerm{strangur} \dicsymFrequent\  \dicIPA{{s}{\textsubring{d}}{r}{au}{\ng}{\r{g}}{\textscy}{\textsubring{r}}} \dicPos{adj}[1] \dicFlx{(f ströng)}[2] \textbf{1.} \dicSynonym{harður} \dicDirectTranslationCS{přísný, nekompromisní, tvrdý} \dicExampleIS{vera strangur við nemendurna} \dicExampleCS{být přísný na žáky}  \textbf{2.} \dicSynonym{stríður} \dicDirectTranslationCS{prudký, divoký (řeka ap.)} \dicExampleIS{strangur straumur} \dicExampleCS{divoký proud};  \dicPhraseIS{í ströngum\,/\addthin strangasta skilningi} \dicFlx{adv} \dicDirectTranslationCS{v~užším slova smyslu};  \dicPhraseIS{langur og strangur (vinnudagur)} \dicDirectTranslationCS{dlouhý a~těžký (pracovní den)};  \dicPhraseIS{standa í ströngu} \dicLangCat{přen.} \dicDirectTranslationCS{mít plné ruce práce}
\dicEntry[straua] \dicTerm{strau|a} \dicIPA{{s}{\textsubring{d}}{r}{\oe i}{\textlengthmark}{a}} \dicPos{v}[1] \dicFlx{(‑aði)}[1] \dicFlx{acc} \dicLink{strauja}
\dicEntry[straubolti] \dicTerm{strau··bolt|i} \dicIPA{{s}{\textsubring{d}}{r}{\oe i}{\textlengthmark}{\textsubring{b}}{\textopeno}{\textsubring{l}}{\textsubring{d}}{\textsci}} \dicPos{m}[1] \dicFlx{(‑a, ‑ar)}[1] \dicSynonym{straujárn} \dicDirectTranslationCS{žehlička}
\dicEntry[straubretti] \dicTerm{strau··bretti} \dicIPA{{s}{\textsubring{d}}{r}{\oe i}{\textlengthmark}{\textsubring{b}}{r}{\textepsilon}{h}{\textsubring{d}}{\textsci}} \dicPos{n}[2] \dicFlx{(‑s, ‑)}[14] \dicDirectTranslationCS{žehlicí prkno}
\dicEntry[strauja] \dicTerm{strauj|a}\dicTerm{, straua} \dicIPA{{s}\-{\textsubring{d}}\-{r}\-{\oe i}\-{j}\-{\textlengthmark}\-{a}\-} \dicPos{v}[1] \dicFlx{(‑aði)}[1] \dicFlx{acc} \dicDirectTranslationCS{(vy)žehlit} \dicExampleIS{strauja þvottinn} \dicExampleCS{vyžehlit prádlo}
\dicEntry[straujárn] \dicTerm{strau··járn} \dicIPA{{s}{\textsubring{d}}{r}{\oe i}{j}{\textlengthmark}{au}{r}{\textsubring{d}}{\textsubring{n}}} \dicPos{n}[2] \dicFlx{(‑s, ‑)}[5] \dicDirectTranslationCS{žehlička}
\dicEntry[strauk] \dicTerm{strauk} \dicIPA{{s}{\textsubring{d}}{r}{\oe i}{\textlengthmark}{\r{g}}} \dicPos{v} \dicFlx{ind pf sg 1 pers} \dicLink{strjúka}
\dicEntry[straumbreytir] \dicTerm{straum··breyt|ir} \dicIPA{{s}{\textsubring{d}}{r}{\oe i}{m}{\textsubring{b}}{r}{ei}{\textsubring{d}}{\textsci}{\textsubring{r}}} \dicPos{m}[7] \dicFlx{(‑is, ‑ar)}[1] \dicFieldCat{elek.} \dicDirectTranslationCS{transformátor}
\dicEntry[straumfræði] \dicTerm{straum··fræð|i} \dicIPA{{s}{\textsubring{d}}{r}{\oe i}{m}{f}{r}{a}{i}{ð}{\textsci}} \dicPos{f}[3] \dicFlx{(‑i)}[3] \dicDirectTranslationCS{hydrodynamika}
\dicEntry[straumhvörf] \dicTerm{straum··hvörf} \dicIPA{{s}{\textsubring{d}}{r}{\oe i}{m}{k\smash{\textsuperscript{h}}}{v}{\oe}{r}{f}} \dicPos{n}[2] \dicFlx{pl}[9] \dicSynonym*{grundvallarbreyting} \dicDirectTranslationCS{zásadní změna\,/\addthin přeměna} \dicExampleIS{straumhvörf í lífi e‑rs} \dicExampleCS{zásadní změna v~životě (koho)}
\dicEntry[straumiða] \dicTerm{straum··ið|a} \dicIPA{{s}{\textsubring{d}}{r}{\oe i}{\textlengthmark}{m}{\textsci}{ð}{a}} \dicPos{f}[1] \dicFlx{(‑u, ‑ur)}[19] \dicSynonym{hringiða} \dicDirectTranslationCS{vír};  \dicPhraseIS{steypa sér út í straumiðu lífsins} \dicLangCat{přen.} \dicDirectTranslationCS{vrhnout se do víru života}
\dicEntry[straumkast] \dicTerm{straum··|kast} \dicIPA{{s}{\textsubring{d}}{r}{\oe i}{m}{k\smash{\textsuperscript{h}}}{a}{s}{\textsubring{d}}} \dicPos{n}[2] \dicFlx{(‑kasts, ‑köst)}[8] \dicDirectTranslationCS{peřeje, silný proud (řeky ap.)} \dicExampleIS{straumkast árinnar} \dicExampleCS{silný proud řeky}
\dicEntry[straumlína] \dicTerm{straum··lín|a} \dicIPA{{s}{\textsubring{d}}{r}{\oe i}{m}{l}{i}{n}{a}} \dicPos{f}[1] \dicFlx{(‑u, ‑ur)}[7] \dicDirectTranslationCS{proudnice, proudová čára}
\begin{xtolerant}{}{1pt}
\dicEntry[straumlínulagaður] \dicTerm{straum·línu··\hspace{-0.06667em}|lag·aður}\addthinS\dicIPA{{s}\-{\textsubring{d}}\-{r}\-{\oe i}\-{m}\-{l}\-{i}\-{n}\-{\textscy}\-{l}\-{a}\-{\textbabygamma}\-{a}\-{ð}\-{\textscy}\-{\textsubring{r}}}\addthinS\dicPos{adj}[3]\addthinS\dicFlx{(f ‑löguð)}[1] \dicDirectTranslationCS{aerodynamický}
\end{xtolerant}
\dicEntry[straummælir] \dicTerm{straum··mæl|ir} \dicIPA{{s}{\textsubring{d}}{r}{\oe i}{m}{\textlengthmark}{a}{i}{l}{\textsci}{\textsubring{r}}} \dicPos{m}[7] \dicFlx{(‑is, ‑ar)}[1] \dicFieldCat{elek.} \dicDirectTranslationCS{ampérmetr}
\dicEntry[straumrof] \dicTerm{straum··rof} \dicIPA{{s}{\textsubring{d}}{r}{\oe i}{m}{r}{\textopeno}{f}} \dicPos{n}[2] \dicFlx{(‑s, ‑)}[5] \dicFieldCat{elek.} \dicDirectTranslationCS{výpadek\,/\addthin přerušení proudu}
\dicEntry[straumrofi] \dicTerm{straum··rof|i} \dicIPA{{s}{\textsubring{d}}{r}{\oe i}{m}{r}{\textopeno}{v}{\textsci}} \dicPos{m}[1] \dicFlx{(‑a, ‑ar)}[1] \dicFieldCat{elek.} \dicDirectTranslationCS{jistič, proudový vypínač}
\begin{xtolerant}{}{1pt}
\dicEntry[straumstyrkur] \dicTerm{straum··styrk|ur}\addthinS\dicIPA{{s}{\textsubring{d}}{r}{\oe i}{m}{s}{\textsubring{d}}{\textsci}{\textsubring{r}}{\r{g}}{\textscy}{\textsubring{r}}}\addthinS\dicPos{m}[9]\addthinS\dicFlx{(‑s, ‑ir)}[15]\addthinS\dicDirectTranslationCS{síla proudu (elektrického\,/\addthin vodního)}
\end{xtolerant}
\dicEntry[straumur] \dicTerm{straum|ur} \dicsymFrequent\  \dicIPA{{s}{\textsubring{d}}{r}{\oe i}{\textlengthmark}{m}{\textscy}{\textsubring{r}}} \dicPos{m}[6] \dicFlx{(‑s, ‑ar)}[22] \textbf{1.} \dicSynonym*{rennsli vatns} \dicDirectTranslationCS{proud (vody), tok} \dicExampleIS{straumur í á} \dicExampleCS{proud v~řece}  \textbf{2.} \dicSynonym{hneigð} \dicDirectTranslationCS{proud, trend, tendence (v~umění ap.)} \dicIndirectTranslationCS{(zvláště v~pl)}  \textbf{3.} \dicSynonym{fjöldi} \dicDirectTranslationCS{proud, příval (lidí ap.)} \dicExampleIS{fólksstraumur} \dicExampleCS{proud lidí}  \textbf{4.} \dicFieldCat{elek.} \dicSynonym{rafstraumur} \dicDirectTranslationCS{(elektrický) proud};  \dicIdiom{straumur}{ \dicPhraseIS{berast með\,/\addthin fylgja straumnum}} \dicFlx{refl} \dicLangCat{přen.} \dicDirectTranslationCS{jít s~davem\,/\addthin proudem}; { \dicPhraseIS{berjast á móti straumnum}} \dicFlx{refl} \dicLangCat{přen.} \dicDirectTranslationCS{jít proti proudu};  \dicPhraseIS{standa straum af e‑u\,/\addthin e‑m} \dicDirectTranslationCS{pokrýt\,/\addthin nést náklady (čeho\,/\addthin čí)}
\dicEntry[straumönd] \dicTerm{straum··|önd} \dicIPA{{s}{\textsubring{d}}{r}{\oe i}{\textlengthmark}{m}{\oe}{n}{\textsubring{d}}} \dicPos{f}[8] \dicFlx{(‑andar, ‑endur)}[2] \dicFieldCat{zool.} \dicDirectTranslationCS{kačka strakatá} \textit{(l.~{\textLA{Histrionicus histrionicus}})}  \dicsymPhoto\ 
\dicFigure{20685.jpg}{Straumönd}{Straumönd - Whitehouse, Laura L., Biolib, PD}
\dicEntry[strauvél] \dicTerm{strau··vél} \dicIPA{{s}{\textsubring{d}}{r}{\oe i}{\textlengthmark}{v}{j}{\textepsilon}{\textsubring{l}}} \dicPos{f}[4] \dicFlx{(‑ar, ‑ar)}[1] \dicDirectTranslationCS{mandl, žehlicí stroj}
\dicEntry[strax] \dicTerm{strax} \dicsymFrequent\  \dicIPA{{s}{\textsubring{d}}{r}{a}{x}{s}} \dicPos{adv} \dicSynonym{undireins} \dicDirectTranslationCS{okamžitě, (i)hned};  \dicPhraseIS{strax og} \dicFlx{conj} \dicSynonym*{um leið og} \dicDirectTranslationCS{v~okamžiku, kdy; jakmile} \dicExampleIS{Hann fór að sofa strax og hann kom heim.} \dicExampleCS{Jakmile přišel domů, tak si šel lehnout.}
\dicEntry[strá] \dicTerm{strá\smash{\textsuperscript{1}}} \dicIPA{{s}{\textsubring{d}}{r}{au}{\textlengthmark}} \dicPos{n}[2] \dicFlx{(‑s, ‑)}[5] \dicDirectTranslationCS{stéblo} \dicExampleIS{heystrá} \dicExampleCS{stéblo slámy};  \dicPhraseIS{e‑að er á hverju strái} \dicLangCat{přen.} \dicDirectTranslationCS{(co) je na každém rohu};  \dicPhraseIS{vera eins og strá í vindi} \dicLangCat{přen.} \dicDirectTranslationCS{být slabý jako moucha (po nemoci ap.)}
\dicEntry[strá] \dicTerm{strá\smash{\textsuperscript{2}}} \dicIPA{{s}{\textsubring{d}}{r}{au}{\textlengthmark}} \dicPos{v}[2] \dicFlx{(‑ði, ‑ð)}[115] \dicFlx{dat} \dicDirectTranslationCS{(po)sypat, rozsypat} \dicExampleIS{strá salti á göturnar} \dicExampleCS{posypat ulice solí}
\dicEntry[stráhattur] \dicTerm{strá··hatt|ur} \dicIPA{{s}{\textsubring{d}}{r}{au}{\textlengthmark}{h}{a}{h}{\textsubring{d}}{\textscy}{\textsubring{r}}} \dicPos{m}[6] \dicFlx{(‑s, ‑ar)}[5] \dicDirectTranslationCS{slaměný klobouk, slamák}
\dicEntry[stráklingur] \dicTerm{strák··ling|ur} \dicIPA{{s}{\textsubring{d}}{r}{au}{\r{g}}{l}{i}{\ng}{\r{g}}{\textscy}{\textsubring{r}}} \dicPos{m}[6] \dicFlx{(‑s, ‑ar)}[8] \dicDirectTranslationCS{klučík, hošík}
\dicEntry[stráksi] \dicTerm{stráks|i} \dicIPA{{s}{\textsubring{d}}{r}{au}{\r{g}}{s}{\textsci}} \dicPos{m}[1] \dicFlx{(‑a, ‑ar)}[1] \dicDirectTranslationCS{klučík, hošík}
\dicEntry[strákslegur] \dicTerm{stráks··legur} \dicIPA{{s}{\textsubring{d}}{r}{au}{\r{g}}{s}{l}{\textepsilon}{\textbabygamma}{\textscy}{\textsubring{r}}} \dicPos{adj}[1]\dicFlx{}[-8] \dicDirectTranslationCS{chlapecký, klukovský}
\dicEntry[strákur] \dicTerm{strák|ur} \dicsymFrequent\  \dicIPA{{s}{\textsubring{d}}{r}{au}{\textlengthmark}{\r{g}}{\textscy}{\textsubring{r}}} \dicPos{m}[6] \dicFlx{(‑s, ‑ar)}[22] \dicSynonym{piltur} \dicDirectTranslationCS{kluk, chlapec, hoch} \dicExampleIS{kona með tvo stráka} \dicExampleCS{žena se dvěma chlapci};  \dicPhraseIS{sitja á strák sínum} \dicLangCat{přen.} \dicDirectTranslationCS{sedět se založenýma rukama}
\dicEntry[strásykur] \dicTerm{strá··syk|ur} \dicIPA{{s}{\textsubring{d}}{r}{au}{\textlengthmark}{s}{\textsci}{\r{g}}{\textscy}{\textsubring{r}}} \dicPos{m}[5] \dicFlx{(‑urs)}[2] \dicDirectTranslationCS{krupicový\,/\addthin krystalový cukr}
\dicEntry[stráþak] \dicTerm{strá··|þak} \dicIPA{{s}{\textsubring{d}}{r}{au}{\textlengthmark}{\texttheta}{a}{\r{g}}} \dicPos{n}[2] \dicFlx{(‑þaks, ‑þök)}[8] \dicDirectTranslationCS{došková střecha}
\dicEntry[streða] \dicTerm{streð|a} \dicIPA{{s}{\textsubring{d}}{r}{\textepsilon}{\textlengthmark}{ð}{a}} \dicPos{v}[1] \dicFlx{(‑aði)}[1] \dicSynonym{erfiða} \dicDirectTranslationCS{dřít se, pachtit se, mořit se};  \dicPhraseIS{streða við að (gera e‑ð)} \dicDirectTranslationCS{mořit se s~(děláním (čeho))}
\dicEntry[streita] \dicTerm{streit|a} \dicIPA{{s}{\textsubring{d}}{r}{ei}{\textlengthmark}{\textsubring{d}}{a}} \dicPos{f}[1] \dicFlx{(‑u)}[5] \dicSynonym{stress} \dicDirectTranslationCS{stres, (nervové) vypětí} \dicExampleIS{finna fyrir streitu} \dicExampleCS{cítit stres};  \dicPhraseIS{halda e‑u til streitu} \dicDirectTranslationCS{jít za (čím) důrazně, prosazovat (co) důrazně}
\dicEntry[streitast] \dicTerm{strei|tast} \dicIPA{{s}{\textsubring{d}}{r}{ei}{\textlengthmark}{\textsubring{d}}{a}{s}{\textsubring{d}}} \dicPos{v}[2] \dicFlx{(‑ttist, ‑st)}[191] \dicFlx{refl} \dicSynonym{strita} \dicDirectTranslationCS{lopotit se, dřít se};  \dicIdiom{streitast}[á móti]{ \dicPhraseIS{streitast á móti e‑u}} \dicDirectTranslationCS{vzpírat se (čemu), bojovat proti (čemu)};  \dicIdiom{streitast}[við]{ \dicPhraseIS{streitast við e‑ð}} \dicDirectTranslationCS{zápasit s~(čím), lopotit se s~(čím)} \dicExampleIS{streitast við námið} \dicExampleCS{zápasit se studiem}
\dicEntry[strekkingur] \dicTerm{strekk··ing|ur} \dicIPA{{s}{\textsubring{d}}{r}{\textepsilon}{h}{\r{\textObardotlessj}}{i}{\ng}{\r{g}}{\textscy}{\textsubring{r}}} \dicPos{m}[6] \dicFlx{(‑s)}[9] \dicSynonym{stinningskaldi} \dicDirectTranslationCS{silný vítr} \dicIndirectTranslationCS{(Beaufort 6) (10,8--13,8~m/s)}
\dicEntry[strekkja] \dicTerm{strekk|ja} \dicIPA{{s}{\textsubring{d}}{r}{\textepsilon}{h}{\r{\textObardotlessj}}{a}} \dicPos{v}[2] \dicFlx{(‑ti, ‑t)}[25] \dicFlx{acc} \dicSynonym{strengja} \dicDirectTranslationCS{napnout, napínat, natáhnout, natahovat (lano ap.)} \dicExampleIS{strekkja vír} \dicExampleCS{napnout drát};  \dicIdiom{strekkjast}{ \dicPhraseIS{strekkjast}} \dicFlx{refl} \dicDirectTranslationCS{napnout se, natáhnout se}
\dicEntry[strekktur] \dicTerm{strekktur} \dicIPA{{s}{\textsubring{d}}{r}{\textepsilon}{x}{\textsubring{d}}{\textscy}{\textsubring{r}}} \dicPos{adj}[1]\dicFlx{}[-10] \textbf{1.} \dicSynonym{þaninn} \dicDirectTranslationCS{napjatý, napnutý}  \textbf{2.} \dicDirectTranslationCS{napjatý, nervózní}
\dicEntry[strembinn] \dicTerm{strembinn} \dicIPA{{s}{\textsubring{d}}{r}{\textepsilon}{m}{\textsubring{b}}{\textsci}{\textsubring{n}}} \dicPos{adj}[6]\dicFlx{}[-2] \dicSynonym{erfiður} \dicDirectTranslationCS{obtížný, nesnadný, nelehký, těžký}
\dicEntry[strendur] \dicTerm{strendur} \dicIPA{{s}{\textsubring{d}}{r}{\textepsilon}{n}{\textsubring{d}}{\textscy}{\textsubring{r}}} \dicPos{f} \dicFlx{pl nom} \dicLink{strönd\smash{\textsuperscript{1}}}
\dicEntry[strengja] \dicTerm{streng|ja} \dicIPA{{s}{\textsubring{d}}{r}{ei}{\textltailn}{\r{\textObardotlessj}}{a}} \dicPos{v}[2] \dicFlx{(‑di, ‑t)}[129] \dicFlx{(gen +) acc} \textbf{1.} \dicFlx{acc} \dicSynonym*{gera stríðan} \dicDirectTranslationCS{napnout, napínat, natáhnout, natahovat} \dicExampleIS{strengja snúrur á hæla} \dicExampleCS{natáhnout šňůry na kolíky}  \textbf{2.} \dicPhraseIS{strengja þess heit að (gera e‑ð)} \dicFlx{gen + acc} \dicDirectTranslationCS{odpřisáhnout (udělat (co)), slíbit (udělat (co))}
\dicEntry[strengjabrúða] \dicTerm{strengja··brúð|a} \dicIPA{{s}{\textsubring{d}}{r}{ei}{\textltailn}{\r{\textObardotlessj}}{a}{\textsubring{b}}{r}{u}{ð}{a}} \dicPos{f}[1] \dicFlx{(‑u, ‑ur)}[19] \textbf{1.} \dicSynonym{leikbrúða} \dicDirectTranslationCS{marioneta, (závěsná) loutka}  \textbf{2.} \dicDirectTranslationCS{loutka} \dicIndirectTranslationCS{(člověk ovládaný cizí vůlí)}
\dicEntry[strengjahljóðfæri] \dicTerm{strengja··hljóð·færi} \dicIPA{{s}{\textsubring{d}}{r}{ei}{\textltailn}{\r{\textObardotlessj}}{a}{\textsubring{l}}{j}{ou}{\texttheta}{f}{a}{i}{r}{\textsci}} \dicPos{n}[2] \dicFlx{(‑s, ‑)}[14] \dicFieldCat{hud.} \dicDirectTranslationCS{strunný nástroj}
\dicEntry[strengjasteypa] \dicTerm{strengja··steyp|a} \dicIPA{{s}{\textsubring{d}}{r}{ei}{\textltailn}{\r{\textObardotlessj}}{a}{s}{\textsubring{d}}{ei}{\textsubring{b}}{a}} \dicPos{f}[1] \dicFlx{(‑u)}[5] \dicDirectTranslationCS{předpjatý beton}
\dicEntry[strengjasveit] \dicTerm{strengja··sveit} \dicIPA{{s}{\textsubring{d}}{r}{ei}{\textltailn}{\r{\textObardotlessj}}{a}{s}{v}{ei}{\textsubring{d}}} \dicPos{f}[7] \dicFlx{(‑ar, ‑ir)}[1] \dicFieldCat{hud.} \dicDirectTranslationCS{smyčcový orchestr}
\dicEntry[strengur] \dicTerm{streng|ur} \dicsymFrequent\  \dicIPA{{s}{\textsubring{d}}{r}{ei}{\ng}{\r{g}}{\textscy}{\textsubring{r}}} \dicPos{m}[9] \dicFlx{(‑s\,/\addthin ‑jar, ‑ir)}[26] \textbf{1.} \dicSynonym{band} \dicDirectTranslationCS{provaz, lano, šňůra}  \textbf{2.} \dicSynonym{lína} \dicDirectTranslationCS{kabel, drát} \dicExampleIS{leggja streng í jörðu} \dicExampleCS{položit kabel do země}  \textbf{3.} \dicSynonym*{buxnastrengur} \dicDirectTranslationCS{pásek, opasek, řemen}  \textbf{4.} \dicFieldCat{hud.} \dicDirectTranslationCS{struna} \dicExampleIS{strengurinn á fiðlunni} \dicExampleCS{struna houslí}  \textbf{5.} \dicSynonym*{straumfall} \dicDirectTranslationCS{peřeje} \dicExampleIS{Það er strengur í ánni.} \dicExampleCS{Na řece jsou peřeje.}  \textbf{6.} \dicSynonym*{vindstrengur} \dicDirectTranslationCS{závan\,/\addthin poryv (větru)}  \textbf{7.} \dicFieldCat{mat.} \dicDirectTranslationCS{tětiva}  \textbf{8.} \dicFieldCat{poč.} \dicDirectTranslationCS{řetězec};  \dicIdiom{strengur}{ \dicPhraseIS{e‑ir stilla saman strengi sína}} \dicLangCat{přen.} \dicDirectTranslationCS{(kdo) spojí své síly}; { \dicPhraseIS{slá á létta strengi}} \dicLangCat{přen.} \dicDirectTranslationCS{udeřit na lehčí strunu} \dicIndirectTranslationCS{(začít hovořit o~lehčím tématu)}; { \dicPhraseIS{taka í sama streng}} \dicLangCat{přen.} \dicDirectTranslationCS{zahrát na stejnou strunu} \dicIndirectTranslationCS{(být stejného názoru)}; { \dicPhraseIS{vera með strengi}} \dicDirectTranslationCS{mít bolavé svaly}
\dicEntry[stress] \dicTerm{stress} \dicIPA{{s}{\textsubring{d}}{r}{\textepsilon}{s}{\textlengthmark}} \dicPos{n}[2] \dicFlx{(‑)}[24] \dicSynonym{streita} \dicDirectTranslationCS{stres}
\dicEntry[stressa] \dicTerm{stress|a} \dicIPA{{s}{\textsubring{d}}{r}{\textepsilon}{s}{\textlengthmark}{a}} \dicPos{v}[1] \dicFlx{(‑aði)}[1] \dicFlx{acc} \dicDirectTranslationCS{(vy)stresovat} \dicExampleIS{Það stressar mig.} \dicExampleCS{Stresuje mě to.};  \dicIdiom{stressast}{ \dicPhraseIS{stressast}} \dicFlx{refl} \dicDirectTranslationCS{(vy)stresovat se};  \dicIdiom{stressast}[upp]{ \dicPhraseIS{stressast upp}} \dicFlx{refl} \dicDirectTranslationCS{(vy)stresovat se}
\dicEntry[stressaður] \dicTerm{stress··|aður} \dicIPA{{s}{\textsubring{d}}{r}{\textepsilon}{s}{\textlengthmark}{a}{ð}{\textscy}{\textsubring{r}}} \dicPos{adj}[3] \dicFlx{(f ‑uð)}[3] \dicDirectTranslationCS{(vy)stresovaný}
\dicEntry[stressandi] \dicTerm{stress··andi} \dicIPA{{s}{\textsubring{d}}{r}{\textepsilon}{s}{\textlengthmark}{a}{n}{\textsubring{d}}{\textsci}} \dicPos{adj}[13] \dicFlx{indecl}[1] \dicDirectTranslationCS{stresující, stresový};  \dicPhraseIS{það er stressandi að (gera e‑ð)} \dicDirectTranslationCS{je to stresující (dělat (co))}
\dicEntry[stresstaska] \dicTerm{stress··|taska} \dicIPA{{s}{\textsubring{d}}{r}{\textepsilon}{s}{\textlengthmark}{t\smash{\textsuperscript{h}}}{a}{s}{\r{g}}{a}} \dicPos{f}[1] \dicFlx{(‑tösku, ‑töskur)}[8] \dicDirectTranslationCS{diplomat\-ka, kufřík}
\dicEntry[streyma] \dicTerm{streym|a} \dicsymFrequent\  \dicIPA{{s}{\textsubring{d}}{r}{ei}{\textlengthmark}{m}{a}} \dicPos{v}[2] \dicFlx{(‑di, ‑t)}[145] \textbf{1.} \dicSynonym{renna\smash{\textsuperscript{3}}} \dicDirectTranslationCS{proudit, téct, valit se (řeka ap.)} \dicExampleIS{Vatnið streymir úr krananum.} \dicExampleCS{Voda proudí z~kohoutku.}  \textbf{2.} \dicDirectTranslationCS{proudit, valit se (davy ap.)}
\dicEntry[streymi] \dicTerm{streymi} \dicIPA{{s}{\textsubring{d}}{r}{ei}{\textlengthmark}{m}{\textsci}} \dicPos{n}[2] \dicFlx{(‑s, ‑)}[14] \textbf{1.} \dicDirectTranslationCS{proud, tok (řeky ap.)}  \textbf{2.} \dicDirectTranslationCS{proud, příliv (utečenců ap.)}
\dicEntry[strigakjaftur] \dicTerm{striga··kjaft|ur} \dicIPA{{s}{\textsubring{d}}{r}{\textsci}{\textlengthmark}{\textbabygamma}{a}{c\smash{\textsuperscript{h}}}{a}{f}{\textsubring{d}}{\textscy}{\textsubring{r}}} \dicPos{m}[6] \dicFlx{(‑s, ‑ar)}[5] \dicDirectTranslationCS{hulvát, sprosťák, sprosťačka}
\dicEntry[strigaskór] \dicTerm{striga··skó|r} \dicIPA{{s}{\textsubring{d}}{r}{\textsci}{\textlengthmark}{\textbabygamma}{a}{s}{\r{g}}{ou}{\textsubring{r}}} \dicPos{m}[6] \dicFlx{(‑s, ‑r)}[52] \dicDirectTranslationCS{teniska} \dicIndirectTranslationCS{(zvláště v~pl)}
\dicEntry[strigi] \dicTerm{strig|i} \dicIPA{{s}{\textsubring{d}}{r}{i}{j}{\textlengthmark}{\textsci}} \dicPos{m}[1] \dicFlx{(‑a, ‑ar)}[1] \dicDirectTranslationCS{pytlovina} \dicExampleIS{strigapoki} \dicExampleCS{(hrubý) pytel}
\dicEntry[strigsa] \dicTerm{strigs|a}\dicTerm{, strixa}\dicTerm{,  striksa} \dicIPA{{s}\-{\textsubring{d}}\-{r}\-{\textsci}\-{x}\-{s}\-{a}\-} \dicPos{v}[1] \dicFlx{(‑aði)}[44] \dicSynonym{strunsa} \dicDirectTranslationCS{vykračovat si, kráčet}
\dicEntry[strik] \dicTerm{strik}\dicTerm{, stryk} \dicIPA{{s}\-{\textsubring{d}}\-{r}\-{\textsci}\-{\textlengthmark}\-{\r{g}}\-} \dicPos{n}[2] \dicFlx{(‑s, ‑)}[5] \textbf{1.} \dicDirectTranslationCS{čára, čárka, linka, linie} \dicExampleIS{strik á vegi} \dicExampleCS{čára na cestě};  \dicPhraseIS{e‑að setur\,/\addthin gerir strik í reikninginn} \dicLangCat{přen.} \dicDirectTranslationCS{(co) dělá čáru přes rozpočet};  \dicPhraseIS{eins og blátt strik} \dicFlx{adv} \dicDirectTranslationCS{jako blesk, rychlostí blesku};  \dicPhraseIS{draga strik yfir e‑ð} \dicLangCat{přen.} \dicDirectTranslationCS{udělat za (čím) tlustou čáru};  \dicPhraseIS{fara yfir strikið} \dicLangCat{přen.} \dicDirectTranslationCS{překročit mez, zajít příliš daleko};  \dicPhraseIS{komast á strikið með e‑ð} \dicFlx{refl} \dicLangCat{přen.} \dicDirectTranslationCS{dostat se na správnou cestu s~(čím)};  \dicPhraseIS{ná sér á strik} \dicLangCat{přen.} \dicDirectTranslationCS{dostat se zpět do formy}  \textbf{2.} \dicFieldCat{mat.} \dicDirectTranslationCS{úsečka}
\dicEntry[strika] \dicTerm{strik|a}\dicTerm{, stryka} \dicIPA{{s}\-{\textsubring{d}}\-{r}\-{\textsci}\-{\textlengthmark}\-{\r{g}}\-{a}\-} \dicPos{v}[1] \dicFlx{(‑aði)}[1] \dicFlx{acc} \dicSynonym*{gera strik} \dicDirectTranslationCS{(pře)škrtnout, škrtat} \dicIndirectTranslationCS{(dělat čáry na něčem)};  \dicIdiom{strika}[undir]{ \dicPhraseIS{strika undir e‑ð}} {\textbf{a.}} \dicDirectTranslationCS{podškrtnout\,/\addthin podtrhnout (co) (chybu v~sešitě ap.)} \dicExampleIS{strika undir orðið} \dicExampleCS{podtrhnout slovo};  {\textbf{b.}} \dicDirectTranslationCS{podtrhnout (co), zdůraznit (co)};  \dicIdiom{strika}[út]{ \dicPhraseIS{strika e‑ð út}} \dicDirectTranslationCS{vyškrtnout (co)};  \dicIdiom{strika}[yfir]{ \dicPhraseIS{strika yfir e‑ð}} \dicDirectTranslationCS{přeškrtnout (co)}
\dicEntry[strikaður] \dicTerm{strik··|aður} \dicIPA{{s}{\textsubring{d}}{r}{\textsci}{\textlengthmark}{\r{g}}{a}{ð}{\textscy}{\textsubring{r}}} \dicPos{adj}[3] \dicFlx{(f ‑uð)}[3] \dicDirectTranslationCS{linkovaný} \dicExampleIS{strikuð blöð} \dicExampleCS{linkované papíry}
\dicEntry[strikamerki] \dicTerm{strika··merki} \dicIPA{{s}{\textsubring{d}}{r}{\textsci}{\textlengthmark}{\r{g}}{a}{m}{\textepsilon}{\textsubring{r}}{\r{\textObardotlessj}}{\textsci}} \dicPos{n}[2] \dicFlx{(‑s, ‑)}[16] \dicDirectTranslationCS{čárový kód}
\dicEntry[striklota] \dicTerm{strik··lot|a}\dicTerm{, stryklota} \dicIPA{{s}\-{\textsubring{d}}\-{r}\-{\textsci}\-{\textlengthmark}\-{\r{g}}\-{l}\-{\textopeno}\-{\textsubring{d}}\-{a}\-} \dicPos{f}[1] \dicFlx{(‑u, ‑ur)}[19] \dicPhraseIS{í striklotu} \dicFlx{adv} \dicDirectTranslationCS{najednou, naráz, na jeden zátah}
\dicEntry[striksa] \dicTerm{striks|a} \dicIPA{{s}{\textsubring{d}}{r}{\textsci}{\r{g}}{s}{a}} \dicPos{v}[1] \dicFlx{(‑aði)}[44] \dicLink{strigsa}
\dicEntry[strimill] \dicTerm{strim|ill} \dicIPA{{s}{\textsubring{d}}{r}{\textsci}{\textlengthmark}{m}{\textsci}{\textsubring{d}}{\textsubring{l}}} \dicPos{m}[6] \dicFlx{(‑ils, ‑lar)}[35] \textbf{1.} \dicSynonym{ræma} \dicDirectTranslationCS{páska, proužek} \dicExampleIS{rista pappaspjaldið í strimla} \dicExampleCS{nařezat papírový lístek na proužky}  \textbf{2.} \dicSynonym*{kassakvittun} \dicDirectTranslationCS{paragon, pokladní doklad}
\dicEntry[stripl] \dicTerm{stripl} \dicIPA{{s}{\textsubring{d}}{r}{\textsci}{h}{\textsubring{b}}{\textsubring{l}}} \dicPos{n}[2] \dicFlx{(‑s)}[2] \dicDirectTranslationCS{(pohoršující) obnažování, exhibicionismus}
\dicEntry[striplingur] \dicTerm{strip··ling|ur} \dicIPA{{s}{\textsubring{d}}{r}{\textsci}{h}{\textsubring{b}}{l}{i}{\ng}{\r{g}}{\textscy}{\textsubring{r}}} \dicPos{m}[6] \dicFlx{(‑s, ‑ar)}[8] \dicDirectTranslationCS{exhibicionista, exhibicionistka}
\dicEntry[strit] \dicTerm{strit} \dicIPA{{s}{\textsubring{d}}{r}{\textsci}{\textlengthmark}{\textsubring{d}}} \dicPos{n}[2] \dicFlx{(‑s)}[2] \dicSynonym{púl} \dicDirectTranslationCS{dřina, pachtění, robota} \dicExampleIS{endalaust strit} \dicExampleCS{nekonečná robota}
\dicEntry[strita] \dicTerm{strit|a} \dicIPA{{s}{\textsubring{d}}{r}{\textsci}{\textlengthmark}{\textsubring{d}}{a}} \dicPos{v}[1] \dicFlx{(‑aði)}[1] \dicSynonym{erfiða} \dicDirectTranslationCS{dřít (se), pachtit se, makat} \dicExampleIS{strita allan daginn} \dicExampleCS{dřít celý den};  \dicPhraseIS{strita við að (gera e‑ð)} \dicDirectTranslationCS{pachtit se s~(děláním (čeho))};  \dicIdiom{stritast}[við]{ \dicPhraseIS{stritast við að (gera e‑ð)}} \dicFlx{refl} \dicDirectTranslationCS{lopotit\,/\addthin pachtit se s~(děláním (čeho))}
\dicEntry[stritvinna] \dicTerm{strit··vinn|a} \dicIPA{{s}{\textsubring{d}}{r}{\textsci}{\textlengthmark}{\textsubring{d}}{v}{\textsci}{n}{a}} \dicPos{f}[1] \dicFlx{(‑u)}[5] \dicSynonym{púl} \dicDirectTranslationCS{dřina, lopota, otročina} \dicExampleIS{líkamleg stritvinna} \dicExampleCS{fyzická dřina}
\dicEntry[strixa] \dicTerm{strix|a} \dicIPA{{s}{\textsubring{d}}{r}{\textsci}{x}{s}{a}} \dicPos{v}[1] \dicFlx{(‑aði)}[44] \dicLink{strigsa}
\dicEntry[stríð] \dicTerm{stríð} \dicsymFrequent\  \dicIPA{{s}{\textsubring{d}}{r}{i}{\textlengthmark}{\texttheta}} \dicPos{n}[2] \dicFlx{(‑s, ‑)}[5] \dicSynonym{styrjöld} \dicDirectTranslationCS{válka, vojna} \dicExampleIS{heimsstríð} \dicExampleCS{světová válka};  \dicPhraseIS{eiga í stríði við e‑n} \dicDirectTranslationCS{být ve válce s~(kým)};  \dicPhraseIS{heilagt stríð} \dicFieldCat{náb.} \dicDirectTranslationCS{svatá válka};  \dicPhraseIS{heyja stríð} \dicDirectTranslationCS{vést válku};  \dicPhraseIS{kalt stríð} \dicFieldCat{hist.} \dicDirectTranslationCS{studená válka};  \dicPhraseIS{segja e‑u stríð á hendur} \dicLangCat{přen.} \dicDirectTranslationCS{vypovědět (čemu) válku (zločinnosti ap.)}
\dicEntry[stríða] \dicTerm{stríð|a\smash{\textsuperscript{1}}} \dicIPA{{s}{\textsubring{d}}{r}{i}{\textlengthmark}{ð}{a}} \dicPos{f}[1] \dicFlx{(‑u, ‑ur)}[19] \dicSynonym{mótlæti} \dicDirectTranslationCS{protivenství, nepřízeň};  \dicPhraseIS{í blíðu og stríðu} \dicFlx{adv} \dicDirectTranslationCS{v~dobrém i~ve zlém}
\dicEntry[stríða] \dicTerm{strí|ða\smash{\textsuperscript{2}}} \dicsymFrequent\  \dicIPA{{s}{\textsubring{d}}{r}{i}{\textlengthmark}{ð}{a}} \dicPos{v}[2] \dicFlx{(‑ddi, ‑tt)}[177] \dicFlx{dat} \textbf{1.} \dicFlx{dat} \dicSynonym{erta\smash{\textsuperscript{2}}} \dicDirectTranslationCS{zlobit, škádlit, dobírat si, popichovat};  \dicPhraseIS{stríða e‑m} \dicDirectTranslationCS{škádlit (koho)};  \dicPhraseIS{stríða e‑m á\,/\addthin með e‑u} \dicDirectTranslationCS{zlobit (koho) kvůli (čemu)}  \textbf{2.} \dicSynonym*{heyja stríð} \dicDirectTranslationCS{bojovat, válčit};  \dicPhraseIS{stríða í móti e‑m} \dicDirectTranslationCS{bojovat proti (komu)}  \textbf{3.} \dicSynonym{erfiða} \dicDirectTranslationCS{zápasit, bojovat (sám se sebou ap.)};  \dicPhraseIS{eiga við e‑ð að stríða} \dicLangCat{přen.} \dicDirectTranslationCS{muset zápasit s~(čím)}
\dicEntry[stríðinn] \dicTerm{stríðinn} \dicIPA{{s}{\textsubring{d}}{r}{i}{\textlengthmark}{ð}{\textsci}{\textsubring{n}}} \dicPos{adj}[6]\dicFlx{}[-2] \dicDirectTranslationCS{škádlivý, rozverný, rozpustilý, nezbedný}
\dicEntry[stríðni] \dicTerm{stríðn|i} \dicIPA{{s}{\textsubring{d}}{r}{i}{ð}{n}{\textsci}} \dicPos{f}[3] \dicFlx{(‑i)}[3] \dicSynonym{ertni} \dicDirectTranslationCS{zlobení, škádlení, dobírání}
\dicEntry[stríðsár] \dicTerm{stríðs··ár} \dicIPA{{s}{\textsubring{d}}{r}{i}{ð}{s}{au}{\textsubring{r}}} \dicPos{n}[2] \dicFlx{pl}[1] \dicDirectTranslationCS{válečná léta} \dicExampleIS{á stríðsárunum} \dicExampleCS{za války}
\dicEntry[stríðsfangi] \dicTerm{stríðs··fang|i} \dicIPA{{s}{\textsubring{d}}{r}{i}{ð}{s}{f}{au}{\textltailn}{\r{\textObardotlessj}}{\textsci}} \dicPos{m}[1] \dicFlx{(‑a, ‑ar)}[8] \dicDirectTranslationCS{(válečný) zajatec, (válečná) zajatkyně}
\dicEntry[stríðsglæpamaður] \dicTerm{stríðs··glæpa·|maður} \dicIPA{{s}{\textsubring{d}}{r}{i}{ð}{s}{\r{g}}{l}{a}{i}{\textsubring{b}}{a}{m}{a}{ð}{\textscy}{\textsubring{r}}} \dicPos{m}[13] \dicFlx{(‑manns, ‑menn)}[2] \dicDirectTranslationCS{válečný zločinec, válečná zločinkyně}
\dicEntry[stríðsglæpur] \dicTerm{stríðs··glæp|ur} \dicIPA{{s}{\textsubring{d}}{r}{i}{ð}{s}{\r{g}}{l}{a}{i}{\textsubring{b}}{\textscy}{\textsubring{r}}} \dicPos{m}[9] \dicFlx{(‑s, ‑ir)}[8] \dicFieldCat{práv.} \dicDirectTranslationCS{válečný zločin}
\dicEntry[stríðsmaður] \dicTerm{stríðs··|maður} \dicIPA{{s}{\textsubring{d}}{r}{i}{ð}{s}{m}{a}{ð}{\textscy}{\textsubring{r}}} \dicPos{m}[13] \dicFlx{(‑manns, ‑menn)}[2] \dicSynonym{hermaður} \dicDirectTranslationCS{voják(yně), válečník, válečnice}
\dicEntry[stríðsterta] \dicTerm{stríðs··tert|a} \dicIPA{{s}{\textsubring{d}}{r}{i}{ð}{s}{t\smash{\textsuperscript{h}}}{\textepsilon}{\textsubring{r}}{\textsubring{d}}{a}} \dicPos{f}[1] \dicFlx{(‑u, ‑ur)}[19] \dicFieldCat{kulin.} \dicIndirectTranslationCS{zdobený velký koláč\,/\addthin dort}
\dicEntry[stríðstími] \dicTerm{stríðs··tím|i} \dicIPA{{s}{\textsubring{d}}{r}{i}{ð}{s}{t\smash{\textsuperscript{h}}}{i}{m}{\textsci}} \dicPos{m}[1] \dicFlx{(‑a, ‑ar)}[1] \dicPhraseIS{á stríðstímum} \dicFlx{adv} \dicDirectTranslationCS{v~době války, ve válečných letech}
\dicEntry[stríðsöxi] \dicTerm{stríðs··|öxi} \dicIPA{{s}{\textsubring{d}}{r}{i}{ð}{s}{\oe}{x}{s}{\textsci}} \dicPos{f}[7] \dicFlx{(‑axar\,/\addthin ‑öxar, ‑axir)}[23] \dicPhraseIS{grafa stríðsöxina} \dicLangCat{přen.} \dicDirectTranslationCS{zakopat válečnou sekyru}
\dicEntry[stríður] \dicTerm{stríður} \dicsymFrequent\  \dicIPA{{s}{\textsubring{d}}{r}{i}{\textlengthmark}{ð}{\textscy}{\textsubring{r}}} \dicPos{adj}[2]\dicFlx{}[-6] \textbf{1.} \dicSynonym*{flughraður} \dicDirectTranslationCS{dravý, prudký} \dicExampleIS{stríður straumur} \dicExampleCS{prudký proud}  \textbf{2.} \dicSynonym{þaninn} \dicDirectTranslationCS{napnutý, napjatý (lano ap.)}  \textbf{3.} \dicSynonym{óþjáll} \dicDirectTranslationCS{tvrdohlavý, zarputilý} \dicExampleIS{vera stríður í lund} \dicExampleCS{být tvrdohlavý}
\dicEntry[strípa] \dicTerm{stríp|a\smash{\textsuperscript{1}}} \dicIPA{{s}{\textsubring{d}}{r}{i}{\textlengthmark}{\textsubring{b}}{a}} \dicPos{f}[1] \dicFlx{(‑u, ‑ur)}[7] \dicSynonym{rönd} \dicDirectTranslationCS{proužek, pásek}
\dicEntry[strípa] \dicTerm{stríp|a\smash{\textsuperscript{2}}} \dicIPA{{s}{\textsubring{d}}{r}{i}{\textlengthmark}{\textsubring{b}}{a}} \dicPos{v}[1] \dicFlx{(‑aði)}[1] \dicDirectTranslationCS{obnažit, obnažovat, (vy)svléknout};  \dicIdiom{strípast}{ \dicPhraseIS{strípast}} \dicFlx{refl} \dicDirectTranslationCS{obnažit se, obnažovat se}
\dicEntry[strípaður] \dicTerm{stríp··|aður} \dicIPA{{s}{\textsubring{d}}{r}{i}{\textlengthmark}{\textsubring{b}}{a}{ð}{\textscy}{\textsubring{r}}} \dicPos{adj}[3] \dicFlx{(f ‑uð)}[3] \textbf{1.} \dicSynonym{kviknakinn} \dicDirectTranslationCS{(vy)svlečený, obnažený}  \textbf{2.} \dicSynonym{eintómur} \dicDirectTranslationCS{holý, prostý, základní} \dicExampleIS{vera á strípuðum töxtum} \dicExampleCS{být na základní platové úrovni}
\dicEntry[strípalingur] \dicTerm{strípa··ling|ur} \dicIPA{{s}{\textsubring{d}}{r}{i}{\textlengthmark}{\textsubring{b}}{a}{l}{i}{\ng}{\r{g}}{\textscy}{\textsubring{r}}} \dicPos{m}[6] \dicFlx{(‑s, ‑ar)}[8] \dicDirectTranslationCS{exhibicionista, exhibicionistka}
\dicEntry[strípihneigð] \dicTerm{strípi··hneigð} \dicIPA{{s}{\textsubring{d}}{r}{i}{\textlengthmark}{\textsubring{b}}{\textsci}{\textsubring{n}}{ei}{\textbabygamma}{\texttheta}} \dicPos{f}[7] \dicFlx{(‑ar)}[3] \dicFieldCat{psych.} \dicDirectTranslationCS{exhibicionismus}
\dicEntry[strjálbýli] \dicTerm{strjál··býli} \dicIPA{{s}{\textsubring{d}}{r}{j}{au}{l}{\textsubring{b}}{i}{l}{\textsci}} \dicPos{n}[2] \dicFlx{(‑s)}[20] \dicDirectTranslationCS{řídce osídlená oblast, venkov}
\dicEntry[strjálbýll] \dicTerm{strjál··býll} \dicIPA{{s}{\textsubring{d}}{r}{j}{au}{l}{\textsubring{b}}{i}{\textsubring{d}}{\textsubring{l}}} \dicPos{adj}[8]\dicFlx{}[-1] \dicDirectTranslationCS{řídce osídlený} \dicAntonym{þéttbýll}
\dicEntry[strjáll] \dicTerm{strjáll} \dicIPA{{s}{\textsubring{d}}{r}{j}{au}{\textsubring{d}}{\textsubring{l}}} \dicPos{adj}[8]\dicFlx{}[-1] \dicDirectTranslationCS{řídký, sporadický} \dicExampleIS{strjálir hólar} \dicExampleCS{sporadické kopce}
\dicEntry[strjúka] \dicTerm{strjúka} \dicsymFrequent\  \dicIPA{{s}{\textsubring{d}}{r}{j}{u}{\textlengthmark}{\r{g}}{a}} \dicPos{v}[6] \dicFlx{(strýk, strauk, strukum, stryki, strokið)}[49] \dicFlx{dat\,/\addthin acc} \textbf{1.} \dicFlx{dat\,/\addthin acc} \dicDirectTranslationCS{(po)hladit, (po)laskat};  \dicPhraseIS{strjúka e‑m} \dicDirectTranslationCS{pohladit (koho)} \dicExampleIS{strjúka honum ástúðlega um kinnina} \dicExampleCS{pohladit ho láskyplně po tváři };  \dicPhraseIS{strjúka e‑ð með e‑u} \dicDirectTranslationCS{pohladit (co čím) (čelo rukou ap.)}  \textbf{2.} \dicFlx{acc} \dicDirectTranslationCS{utřít, utírat, setřít, stírat, přetřít} \dicExampleIS{strjúka rykið} \dicExampleCS{setřít prach};  \dicPhraseIS{strjúka af e‑u} \dicDirectTranslationCS{setřít (co), přetřít (co)} \dicExampleIS{strjúka af gólfi} \dicExampleCS{přetřít podlahu};  \dicPhraseIS{strjúka yfir e‑ð} \dicDirectTranslationCS{setřít (co)}  \textbf{3.} \dicSynonym*{fara í óleyfi} \dicDirectTranslationCS{utéct, utíkat, uprchnout, prchat} \dicExampleIS{strjúka úr fangelsi} \dicExampleCS{utéct z~vězení};  \dicIdiom{strjúkast}[við]{ \dicPhraseIS{strjúkast við e‑ð}} \dicFlx{refl} \dicDirectTranslationCS{zavadit o~(co), lehce se (čeho) dotknout}
\dicEntry[stroffa] \dicTerm{stroff|a} \dicIPA{{s}{\textsubring{d}}{r}{\textopeno}{f}{\textlengthmark}{a}} \dicPos{f}[1] \dicFlx{(‑u, ‑ur)}[7] \dicDirectTranslationCS{popruh}
\dicEntry[strok] \dicTerm{strok} \dicIPA{{s}{\textsubring{d}}{r}{\textopeno}{\textlengthmark}{\r{g}}} \dicPos{n}[2] \dicFlx{(‑s)}[2] \textbf{1.} \dicSynonym{flótti} \dicDirectTranslationCS{útěk, úprk, uprchnutí, únik} \dicExampleIS{strok úr fangelsi} \dicExampleCS{útěk z~vězení}  \textbf{2.} \dicFieldCat{med. } \dicSynonym*{frumustrok} \dicDirectTranslationCS{výtěr}
\dicEntry[stroka] \dicTerm{strok|a\smash{\textsuperscript{1}}} \dicIPA{{s}{\textsubring{d}}{r}{\textopeno}{\textlengthmark}{\r{g}}{a}} \dicPos{f}[1] \dicFlx{(‑u, ‑ur)}[7] \textbf{1.} \dicSynonym*{það að strjúka e‑u} \dicDirectTranslationCS{(po)hlazení, (po)laskání}  \textbf{2.} \dicSynonym{vindhviða} \dicDirectTranslationCS{poryv\,/\addthin závan větru}  \textbf{3.} \dicSynonym{hýðing} \dicDirectTranslationCS{rána, úder}
\dicEntry[stroka] \dicTerm{strok|a\smash{\textsuperscript{2}}} \dicIPA{{s}{\textsubring{d}}{r}{\textopeno}{\textlengthmark}{\r{g}}{a}} \dicPos{v}[1] \dicFlx{(‑aði)}[1] \dicFlx{acc} \dicPhraseIS{stroka e‑ð út} \dicDirectTranslationCS{vygumovat (co), vymazat (co)}
\dicEntry[strokhljóðfæri] \dicTerm{strok··hljóð·færi} \dicIPA{{s}{\textsubring{d}}{r}{\textopeno}{\textlengthmark}{\r{g}}{\textsubring{l}}{j}{ou}{\texttheta}{f}{a}{i}{r}{\textsci}} \dicPos{n}[2] \dicFlx{(‑s, ‑)}[14] \dicFieldCat{hud.} \dicDirectTranslationCS{smyčcový nástroj}
\dicEntry[strokið] \dicTerm{strokið} \dicIPA{{s}{\textsubring{d}}{r}{\textopeno}{\textlengthmark}{\r{\textObardotlessj}}{\textsci}{\texttheta}} \dicPos{v} \dicFlx{supin} \dicLink{strjúka}
\dicEntry[strokkur] \dicTerm{strokk|ur} \dicIPA{{s}{\textsubring{d}}{r}{\textopeno}{h}{\r{g}}{\textscy}{\textsubring{r}}} \dicPos{m}[6] \dicFlx{(‑s, ‑ar)}[22] \textbf{1.} \dicSynonym*{smjörstrokkur} \dicDirectTranslationCS{máselnice}  \textbf{2.} \dicFieldCat{techn.} \dicDirectTranslationCS{cylindr, válec} \dicIndirectTranslationCS{(motoru)}  \textbf{3.} \dicFieldCat{bot.} \dicDirectTranslationCS{zvonečník} \textit{(l.~{\textLA{Phyteuma}})}  \dicsymPhoto\ 
\dicFigure{66460.jpg}{Strokkur}{Strokkur - Dvořák Václav, Biolib, CC-BY-NC}
\dicEntry[strokkvartett] \dicTerm{strok··kvartett} \dicIPA{{s}{\textsubring{d}}{r}{\textopeno}{\textlengthmark}{\r{g}}{k\smash{\textsuperscript{h}}}{v}{a}{\textsubring{r}}{\textsubring{d}}{\textepsilon}{h}{\textsubring{d}}} \dicPos{m}[4] \dicFlx{(‑s, ‑ar)}[22] \dicFieldCat{hud.} \dicDirectTranslationCS{smyčcový kvartet}
\dicEntry[strokleður] \dicTerm{strok··leður} \dicIPA{{s}{\textsubring{d}}{r}{\textopeno}{\textlengthmark}{\r{g}}{l}{\textepsilon}{ð}{\textscy}{\textsubring{r}}} \dicPos{n}[2] \dicFlx{(‑s, ‑)}[25] \dicDirectTranslationCS{guma} \dicIndirectTranslationCS{(výrobek z~pryže používaný ke gumování)}
\dicEntry[strokufangi] \dicTerm{stroku··fang|i} \dicIPA{{s}{\textsubring{d}}{r}{\textopeno}{\textlengthmark}{\r{g}}{\textscy}{f}{au}{\textltailn}{\r{\textObardotlessj}}{\textsci}} \dicPos{m}[1] \dicFlx{(‑a, ‑ar)}[8] \dicDirectTranslationCS{uprchlý vězeň, uprchlá vězeňkyně, vězeň(kyně) na útěku}
\begin{xtolerant}{}{1pt}
\dicEntry[strokumaður] \dicTerm{stroku··|maður}\addthinS\dicIPA{{s}{\textsubring{d}}{r}{\textopeno}{\textlengthmark}{\r{g}}{\textscy}{m}{a}{ð}{\textscy}{\textsubring{r}}}\addthinS\dicPos{m}[13]\addthinS\dicFlx{(‑manns, ‑menn)}[2] \dicDirectTranslationCS{uprchlík, uprchlice, utečenec, utečenkyně, běženec, běženkyně}
\end{xtolerant}
\dicEntry[strompur] \dicTerm{stromp|ur} \dicIPA{{s}{\textsubring{d}}{r}{\textopeno}{\textsubring{m}}{\textsubring{b}}{\textscy}{\textsubring{r}}} \dicPos{m}[6] \dicFlx{(‑s, ‑ar)}[4] \dicSynonym{reykháfur} \dicDirectTranslationCS{komín} \dicExampleIS{byrgja strompinn} \dicExampleCS{zakrýt komín}
\dicEntry[strókur] \dicTerm{strók|ur} \dicIPA{{s}{\textsubring{d}}{r}{ou}{\textlengthmark}{\r{g}}{\textscy}{\textsubring{r}}} \dicPos{m}[6] \dicFlx{(‑s, ‑ar)}[22] \dicDirectTranslationCS{sloup, sloupec (kouře ap.)} \dicExampleIS{gufustrókur} \dicExampleCS{sloup páry}
\dicEntry[strukum] \dicTerm{strukum} \dicIPA{{s}{\textsubring{d}}{r}{\textscy}{\textlengthmark}{\r{g}}{\textscy}{\textsubring{m}}} \dicPos{v} \dicFlx{ind pf pl 1 pers} \dicLink{strjúka}
\dicEntry[strunsa] \dicTerm{struns|a} \dicIPA{{s}{\textsubring{d}}{r}{\textscy}{n}{s}{a}} \dicPos{v}[1] \dicFlx{(‑aði)}[44] \dicDirectTranslationCS{kráčet, vykračovat si};  \dicIdiom{strunsa}[út]{ \dicPhraseIS{strunsa út}} \dicDirectTranslationCS{odporoučet se, odejít} \dicExampleIS{Mad\-daman strunsaði út.} \dicExampleCS{Madam se odporoučela.}
\dicEntry[strútur] \dicTerm{strút|ur} \dicIPA{{s}{\textsubring{d}}{r}{u}{\textlengthmark}{\textsubring{d}}{\textscy}{\textsubring{r}}} \dicPos{m}[6] \dicFlx{(‑s, ‑ar)}[22] \dicFieldCat{zool.} \dicDirectTranslationCS{pštros, pštros dvouprstý} \textit{(l.~{\textLA{Struthio camelus}})}  \dicsymPhoto\ 
\dicFigure{ds_image_strutur_0_1.jpg}{Strútur}{Strútur - Joachim Huber, CC BY 2.0}
\dicEntry[stryk] \dicTerm{stryk} \dicIPA{{s}{\textsubring{d}}{r}{\textsci}{\textlengthmark}{\r{g}}} \dicPos{n}[2] \dicFlx{(‑s, ‑)}[5] \dicLink{strik}
\dicEntry[stryka] \dicTerm{stryk|a} \dicIPA{{s}{\textsubring{d}}{r}{\textsci}{\textlengthmark}{\r{g}}{a}} \dicPos{v}[1] \dicFlx{(‑aði)}[1] \dicLink{strika}
\dicEntry[stryki] \dicTerm{stryki} \dicIPA{{s}{\textsubring{d}}{r}{\textsci}{\textlengthmark}{\r{\textObardotlessj}}{\textsci}} \dicPos{v} \dicFlx{con pf sg 1 pers} \dicLink{strjúka}
\dicEntry[stryklota] \dicTerm{stryk··lot|a} \dicIPA{{s}{\textsubring{d}}{r}{\textsci}{\textlengthmark}{\r{g}}{l}{\textopeno}{\textsubring{d}}{a}} \dicPos{f}[1] \dicFlx{(‑u, ‑ur)}[19] \dicLink{striklota}
\dicEntry[strý] \dicTerm{strý} \dicIPA{{s}{\textsubring{d}}{r}{i}{\textlengthmark}} \dicPos{n}[2] \dicFlx{(‑s)}[2] \dicDirectTranslationCS{prořídlé\,/\addthin řídké vlasy}
\dicEntry[strýk] \dicTerm{strýk} \dicIPA{{s}{\textsubring{d}}{r}{i}{\textlengthmark}{\r{g}}} \dicPos{v} \dicFlx{ind praes sg 1 pers} \dicLink{strjúka}
\dicEntry[strýta] \dicTerm{strýt|a} \dicIPA{{s}{\textsubring{d}}{r}{i}{\textlengthmark}{\textsubring{d}}{a}} \dicPos{f}[1] \dicFlx{(‑u, ‑ur)}[7] \dicSynonym{keila} \dicDirectTranslationCS{(kuželovitý) vrchol hory}
\dicEntry[strýtulagaður] \dicTerm{strýtu··|lag·aður} \dicIPA{{s}{\textsubring{d}}{r}{i}{\textlengthmark}{\textsubring{d}}{\textscy}{l}{a}{\textbabygamma}{a}{ð}{\textscy}{\textsubring{r}}} \dicPos{adj}[3] \dicFlx{(f ‑löguð)}[2] \dicDirectTranslationCS{kuželovitý, kónický}
\dicEntry[stræka] \dicTerm{stræk|a} \dicIPA{{s}{\textsubring{d}}{r}{a}{i}{\textlengthmark}{\r{g}}{a}} \dicPos{v}[1] \dicFlx{(‑aði)}[1] \dicPhraseIS{stræka á e‑ð} \dicLangCat{hovor.} \dicDirectTranslationCS{odmítnout (co), odmítat (co), odepřít (co), odpírat (co)}
\dicEntry[stræti] \dicTerm{stræti} \dicIPA{{s}{\textsubring{d}}{r}{a}{i}{\textlengthmark}{\textsubring{d}}{\textsci}} \dicPos{n}[2] \dicFlx{(‑s, ‑)}[14] \dicSynonym{gata\smash{\textsuperscript{1}}} \dicDirectTranslationCS{ulice, třída};  \dicPhraseIS{vera í strætinu} \dicLangCat{hovor.} \dicDirectTranslationCS{být na ulici (pobuda ap.)}
\dicEntry[strætisvagn] \dicTerm{strætis··vagn} \dicIPA{{s}{\textsubring{d}}{r}{a}{i}{\textlengthmark}{\textsubring{d}}{\textsci}{s}{v}{a}{\r{g}}{\textsubring{n}}} \dicPos{m}[4] \dicFlx{(‑s, ‑ar)}[4] \dicDirectTranslationCS{autobus} \dicIndirectTranslationCS{(městské hromadné dopravy)}
\dicEntry[strætisvagnastöð] \dicTerm{strætis·vagna·stöð} \dicIPA{{s}{\textsubring{d}}{r}{a}{i}{\textlengthmark}{\textsubring{d}}{\textsci}{s}{v}{a}{\r{g}}{n}{a}{s}{\textsubring{d}}{\oe}{\texttheta}} \dicPos{f}[6] \dicFlx{(‑var, ‑var)}[1] \dicDirectTranslationCS{autobusová zastávka} \dicIndirectTranslationCS{(městské hromadné dopravy)}
\dicEntry[strætó] \dicTerm{strætó} \dicIPA{{s}{\textsubring{d}}{r}{a}{i}{\textlengthmark}{\textsubring{d}}{ou}} \dicPos{m}[4] \dicFlx{(‑s, ‑ar)}[14] \textbf{1.} \dicLangCat{hovor.} \dicSynonym{strætisvagn} \dicDirectTranslationCS{autobus} \dicIndirectTranslationCS{(městské hromadné dopravy)}  \textbf{2.} \dicLangCat{hovor.} \dicDirectTranslationCS{dopravní podnik MHD}
\dicEntry[strætómiði] \dicTerm{strætó··mið|i} \dicIPA{{s}{\textsubring{d}}{r}{a}{i}{\textlengthmark}{\textsubring{d}}{ou}{m}{\textsci}{ð}{\textsci}} \dicPos{m}[1] \dicFlx{(‑a, ‑ar)}[1] \dicDirectTranslationCS{jízdenka\,/\addthin lístek na autobus, autobusový lístek, autobusová jízdenka} \dicIndirectTranslationCS{(městské hromadné dopravy)}
\dicEntry[strætóskýli] \dicTerm{strætó··skýli} \dicIPA{{s}{\textsubring{d}}{r}{a}{i}{\textlengthmark}{\textsubring{d}}{ou}{s}{\r{\textObardotlessj}}{i}{l}{\textsci}} \dicPos{n}[2] \dicFlx{(‑s, ‑)}[14] \dicDirectTranslationCS{přístřešek autobusové zastávky} \dicIndirectTranslationCS{(městské hromadné dopravy)}
\dicEntry[strönd] \dicTerm{strönd\smash{\textsuperscript{1}}} \dicsymFrequent\  \dicIPA{{s}{\textsubring{d}}{r}{\oe}{n}{\textsubring{d}}} \dicPos{f}[8] \dicFlx{(strandar, strendur\,/\addthin strandir)}[9] \textbf{1.} \dicDirectTranslationCS{pobřeží, břeh} \dicExampleIS{verndun hafs og strandar} \dicExampleCS{ochrana moře a~pobřeží}  \textbf{2.} \dicDirectTranslationCS{pláž};  \dicPhraseIS{komast hvorki lönd né strönd} \dicFlx{refl} \dicLangCat{přen.} \dicDirectTranslationCS{zůstat na půli cesty}
\dicEntry[strönd] \dicTerm{strönd\smash{\textsuperscript{2}}} \dicIPA{{s}{\textsubring{d}}{r}{\oe}{n}{\textsubring{d}}} \dicPos{n} \dicFlx{pl nom} \dicLink{strand\smash{\textsuperscript{1}}}
\dicEntry[ströng] \dicTerm{ströng} \dicIPA{{s}{\textsubring{d}}{r}{\oe i}{\ng}{\r{g}}} \dicPos{adj} \dicFlx{f sg nom pos} \dicLink{strangur}
\dicEntry[stubbi] \dicTerm{stubb|i} \dicIPA{{s}{\textsubring{d}}{\textscy}{\textsubring{b}}{\textlengthmark}{\textsci}} \dicPos{m}[1] \dicFlx{(‑a, ‑ar)}[1] \dicLink{stubbur}
\dicEntry[stubbur] \dicTerm{stubb|ur}\dicTerm{, stubbi} \dicIPA{{s}\-{\textsubring{d}}\-{\textscy}\-{\textsubring{b}}\-{\textlengthmark}\-{\textscy}\-{\textsubring{r}}\-} \dicPos{m}[6] \dicFlx{(‑s, ‑ar)}[8] \textbf{1.} \dicSynonym{stúfur\smash{\textsuperscript{1}}} \dicDirectTranslationCS{pahýl, zbytek (tužky ap.)} \dicExampleIS{kertisstubbur} \dicExampleCS{zbytek svíčky}  \textbf{2.} \dicSynonym{sígarettustubbur} \dicDirectTranslationCS{nedopalek, oharek, špaček}  \textbf{3.} \dicSynonym*{drengsnáði} \dicDirectTranslationCS{skrček}
\dicEntry[studdi] \dicTerm{studdi} \dicIPA{{s}{\textsubring{d}}{\textscy}{\textsubring{d}}{\textlengthmark}{\textsci}} \dicPos{v} \dicFlx{ind pf sg 1 pers} \dicLink{styðja}
\dicEntry[studdum] \dicTerm{studdum} \dicIPA{{s}{\textsubring{d}}{\textscy}{\textsubring{d}}{\textlengthmark}{\textscy}{\textsubring{m}}} \dicPos{v} \dicFlx{ind pf pl 1 pers} \dicLink{styðja}
\dicEntry[stuð] \dicTerm{stuð} \dicIPA{{s}{\textsubring{d}}{\textscy}{\textlengthmark}{\texttheta}} \dicPos{n}[2] \dicFlx{(‑s)}[2] \textbf{1.} \dicSynonym*{rafstuð} \dicDirectTranslationCS{rána\,/\addthin pecka (elektrickým proudem)}  \textbf{2.} \dicSynonym*{góð stemning} \dicDirectTranslationCS{ideální rozpoložení, výborná nálada, super pocit};  \dicPhraseIS{vera í stuði} \dicDirectTranslationCS{být v~ideálním rozpoložení}  \textbf{3.} \dicLangCat{hovor.} \dicSynonym{hass} \dicDirectTranslationCS{hašiš, haš}
\dicEntry[stuða] \dicTerm{stuð|a} \dicIPA{{s}{\textsubring{d}}{\textscy}{\textlengthmark}{ð}{a}} \dicPos{v}[1] \dicFlx{(‑aði)}[1] \dicSynonym{særa} \dicDirectTranslationCS{(z)ranit, ublížit, ubližovat, budit rozhořčení} \dicExampleIS{Það stuðar mig hvernig hann talar við konuna sína.} \dicExampleCS{Zraňuje mě, jak mluví se svojí ženou.}
\dicEntry[stuðandi] \dicTerm{stuð··andi} \dicIPA{{s}{\textsubring{d}}{\textscy}{\textlengthmark}{ð}{a}{n}{\textsubring{d}}{\textsci}} \dicPos{adj}[13] \dicFlx{indecl}[1] \dicDirectTranslationCS{hanlivý, urážlivý}
\dicEntry[stuðari] \dicTerm{stuð··ar|i} \dicIPA{{s}{\textsubring{d}}{\textscy}{\textlengthmark}{ð}{a}{r}{\textsci}} \dicPos{m}[1] \dicFlx{(‑a, ‑ar)}[13] \dicDirectTranslationCS{nárazník}
\dicEntry[stuðla] \dicTerm{stuðl|a} \dicsymFrequent\  \dicIPA{{s}{\textsubring{d}}{\textscy}{ð}{l}{a}} \dicPos{v}[1] \dicFlx{(‑aði)}[1] \textbf{1.} \dicFieldCat{lit.} \dicDirectTranslationCS{aliterovat} \dicExampleIS{stuðla við e‑ð} \dicExampleCS{aliterovat s~(čím)}  \textbf{2.} \dicPhraseIS{stuðla að e‑u} \dicDirectTranslationCS{přispět k~(čemu), napomáhat (čemu)}
\dicEntry[stuðlaberg] \dicTerm{stuðla··berg} \dicIPA{{s}{\textsubring{d}}{\textscy}{ð}{l}{a}{\textsubring{b}}{\textepsilon}{r}{\r{g}}} \dicPos{n}[2] \dicFlx{(‑s)}[2] \dicFieldCat{geol.} \dicDirectTranslationCS{sloupcovitý čedič}
\dicEntry[stuðlaður] \dicTerm{stuðl··|aður} \dicIPA{{s}{\textsubring{d}}{\textscy}{ð}{l}{a}{ð}{\textscy}{\textsubring{r}}} \dicPos{adj}[3] \dicFlx{(f ‑uð)}[3] \dicFieldCat{lit.} \dicDirectTranslationCS{aliterační}
\dicEntry[stuðlasetning] \dicTerm{stuðla··set·ning} \dicIPA{{s}{\textsubring{d}}{\textscy}{ð}{l}{a}{s}{\textepsilon}{h}{\textsubring{d}}{n}{i}{\ng}{\r{g}}} \dicPos{f}[4] \dicFlx{(‑ar)}[7] \dicFieldCat{lit.} \dicDirectTranslationCS{aliterace}
\dicEntry[stuðlun] \dicTerm{stuðl|un} \dicIPA{{s}{\textsubring{d}}{\textscy}{ð}{l}{\textscy}{\textsubring{n}}} \dicPos{f}[7] \dicFlx{(‑unar)}[9] \dicFieldCat{lit.} \dicSynonym{stuðlasetning} \dicDirectTranslationCS{aliterace}
\dicEntry[stuðningsmaður] \dicTerm{stuðnings··|maður} \dicIPA{{s}{\textsubring{d}}{\textscy}{ð}{n}{i}{\ng}{s}{m}{a}{ð}{\textscy}{\textsubring{r}}} \dicPos{m}[13] \dicFlx{(‑manns, ‑menn)}[2] \dicDirectTranslationCS{stoupenec, stoupenkyně, fanoušek, fanynka, sympatizant(ka)}
\dicEntry[stuðningur] \dicTerm{stuð··ning|ur} \dicsymFrequent\  \dicIPA{{s}{\textsubring{d}}{\textscy}{ð}{n}{i}{\ng}{\r{g}}{\textscy}{\textsubring{r}}} \dicPos{m}[6] \dicFlx{(‑s)}[9] \dicSynonym{hjálp} \dicDirectTranslationCS{podpora, pomoc} \dicExampleIS{leita eftir stuðningi og skilningi} \dicExampleCS{žádat o~podporu a~porozumění}
\dicEntry[stuðpúði] \dicTerm{stuð··púð|i} \dicIPA{{s}{\textsubring{d}}{\textscy}{ð}{p\smash{\textsuperscript{h}}}{u}{ð}{\textsci}} \dicPos{m}[1] \dicFlx{(‑a, ‑ar)}[1] \dicLangCat{přen.} \dicDirectTranslationCS{tlumič (nárazů), nárazník, tlumicí prvek}
\dicEntry[stuðull] \dicTerm{stuð|ull} \dicIPA{{s}{\textsubring{d}}{\textscy}{\textlengthmark}{ð}{\textscy}{\textsubring{d}}{\textsubring{l}}} \dicPos{m}[6] \dicFlx{(‑uls, ‑lar)}[31] \textbf{1.} \dicFieldCat{lit.} \dicDirectTranslationCS{aliterační hláska}  \textbf{2.} \dicSynonym*{rúmstólpi} \dicDirectTranslationCS{vzpěra, podpěra}  \textbf{3.} \dicFieldCat{mat.} \dicDirectTranslationCS{faktor, koeficient, součinitel, index}  \textbf{4.} \dicFieldCat{geol.} \dicSynonym{stuðlaberg} \dicDirectTranslationCS{sloupcovitý čedič}
\dicEntry[stugga] \dicTerm{stugg|a} \dicIPA{{s}{\textsubring{d}}{\textscy}{\r{g}}{\textlengthmark}{a}} \dicPos{v}[1] \dicFlx{(‑aði)}[3] \dicSynonym*{reka frá} \dicDirectTranslationCS{odehnat, zahnat, zaplašit};  \dicPhraseIS{stugga við e‑u} \dicDirectTranslationCS{odehnat (co), zahnat (co)} \dicExampleIS{stugga við fénu} \dicExampleCS{zahnat ovce};  \dicPhraseIS{e‑n stuggar við e‑u} \dicFlx{impers} \dicDirectTranslationCS{(kdo) má z~(čeho) obavy}
\dicEntry[stuggur] \dicTerm{stugg|ur} \dicIPA{{s}{\textsubring{d}}{\textscy}{\r{g}}{\textlengthmark}{\textscy}{\textsubring{r}}} \dicPos{m}[6] \dicFlx{(‑s)}[17] \dicPhraseIS{e‑m stendur stuggur af e‑u} \dicFlx{impers} \dicDirectTranslationCS{(kdo) má z~(čeho) hrůzu}
\dicEntry[stukkum] \dicTerm{stukkum} \dicIPA{{s}{\textsubring{d}}{\textscy}{h}{\r{g}}{\textscy}{\textsubring{m}}} \dicPos{v} \dicFlx{ind pf pl 1 pers} \dicLink{stökkva\smash{\textsuperscript{2}}}
\dicEntry[stuldur] \dicTerm{stuld|ur} \dicIPA{{s}{\textsubring{d}}{\textscy}{l}{\textsubring{d}}{\textscy}{\textsubring{r}}} \dicPos{m}[10] \dicFlx{(‑ar, ‑ir)}[4] \dicDirectTranslationCS{krádež, rozkrádání, odcizení}
\dicEntry[stumra] \dicTerm{stumr|a} \dicIPA{{s}{\textsubring{d}}{\textscy}{m}{r}{a}} \dicPos{v}[1] \dicFlx{(‑aði)}[1] \dicPhraseIS{stumra yfir e‑m} \dicDirectTranslationCS{poskytnout (komu) první pomoc}
\dicEntry[stuna] \dicTerm{stun|a} \dicIPA{{s}{\textsubring{d}}{\textscy}{\textlengthmark}{n}{a}} \dicPos{f}[1] \dicFlx{(‑u, ‑ur)}[7] \dicSynonym{andvarp} \dicDirectTranslationCS{sten, zasténání} \dicExampleIS{sársaukastuna} \dicExampleCS{zasténání z~bolesti}
\dicEntry[stund] \dicTerm{stund} \dicsymFrequent\  \dicIPA{{s}{\textsubring{d}}{\textscy}{n}{\textsubring{d}}} \dicPos{f}[7] \dicFlx{(‑ar, ‑ir)}[5] \textbf{1.} \dicSynonym*{stutt tími} \dicDirectTranslationCS{chvíle, chvilka, moment, (krátký) čas}  \textbf{2.} \dicSynonym{kennslustund} \dicDirectTranslationCS{(vyučovací) hodina}  \textbf{3.} \dicSynonym{klukkustund} \dicDirectTranslationCS{hodina} \dicIndirectTranslationCS{(šedesát minut)};  \dicIdiom{stund}{ \dicPhraseIS{á elleftu stundu}} \dicFlx{adv} \dicDirectTranslationCS{na poslední chvíli};  \dicPhraseIS{á samri stundu} \dicFlx{adv} \dicSynonym{samtímis} \dicDirectTranslationCS{ve stejnou chvíli};  \dicPhraseIS{á skammri stundu} \dicFlx{adv} \dicDirectTranslationCS{v~krátkém čase};  \dicPhraseIS{á stundinni} \dicFlx{adv} \dicDirectTranslationCS{okamžitě, (i)hned};  \dicPhraseIS{eftir stutta stund} \dicFlx{adv} \dicSynonym{brátt} \dicDirectTranslationCS{za chvilku, zanedlouho};  \dicPhraseIS{góða stund} \dicFlx{adv} \dicDirectTranslationCS{docela dlouho};  \dicPhraseIS{innan stundar} \dicFlx{adv} \dicSynonym*{bráðlega} \dicDirectTranslationCS{brzy, zanedlouho};  \dicPhraseIS{leggja stund á e‑ð} \dicDirectTranslationCS{věnovat se (čemu)};  \dicPhraseIS{lifa fyrir líðandi stund} \dicDirectTranslationCS{žít v~přítomnosti};  \dicPhraseIS{stytta e‑m stundir} \dicDirectTranslationCS{krátit (komu) čas, rozptýlit (koho)};  \dicPhraseIS{stytta sér stundir við e‑ð} \dicDirectTranslationCS{trávit čas (čím), krátit si čas (čím)};  \dicPhraseIS{um stund} \dicFlx{adv} \dicDirectTranslationCS{na chvilku, na (krátkou) chvíli};  \dicPhraseIS{þessa stundina} \dicFlx{adv} \dicDirectTranslationCS{v~tuto chvíli, nyní, teď};  \dicPhraseIS{þegar fram líða stundir} \dicFlx{adv} \dicDirectTranslationCS{(s) časem};  \dicPhraseIS{öllum stundum} \dicFlx{adv} \dicDirectTranslationCS{pořád, celý čas}
\dicEntry[stunda] \dicTerm{stund|a} \dicsymFrequent\  \dicIPA{{s}{\textsubring{d}}{\textscy}{n}{\textsubring{d}}{a}} \dicPos{v}[1] \dicFlx{(‑aði)}[1] \dicFlx{acc} \dicPhraseIS{stunda e‑ð} \dicDirectTranslationCS{pěstovat (co), oddávat se (čemu) (sportu ap.), věnovat se (čemu) (práci ap.), chodit na (co) (kurzy ap.)} \dicExampleIS{stunda nám reglulega} \dicExampleCS{pravidelně chodit na kurzy};  \dicPhraseIS{stunda e‑n} \dicSynonym{gæta} \dicDirectTranslationCS{pečovat o~(koho), docházet ke (komu)} \dicExampleIS{stunda sjúkling} \dicExampleCS{docházet k~pacientovi}
\dicEntry[stundafjöldi] \dicTerm{stunda··fjöld|i} \dicIPA{{s}{\textsubring{d}}{\textscy}{n}{\textsubring{d}}{a}{f}{j}{\oe}{l}{\textsubring{d}}{\textsci}} \dicPos{m}[1] \dicFlx{(‑a)}[3] \dicDirectTranslationCS{počet hodin (vyučovacích ap.)}
\dicEntry[stundaglas] \dicTerm{stunda··|glas} \dicIPA{{s}{\textsubring{d}}{\textscy}{n}{\textsubring{d}}{a}{\r{g}}{l}{a}{s}} \dicPos{n}[2] \dicFlx{(‑glass, ‑glös)}[8] \dicDirectTranslationCS{přesýpací hodiny}
\dicEntry[stundakennari] \dicTerm{stunda··kenn·ar|i} \dicIPA{{s}{\textsubring{d}}{\textscy}{n}{\textsubring{d}}{a}{c\smash{\textsuperscript{h}}}{\textepsilon}{n}{a}{r}{\textsci}} \dicPos{m}[1] \dicFlx{(‑a, ‑ar)}[13] \dicFieldCat{škol.} \dicDirectTranslationCS{externí učitel(ka), externista, externistka}
\dicEntry[stundarfjórðungur] \dicTerm{stundar··fjórð·ung|ur} \dicIPA{{s}\-{\textsubring{d}}\-{\textscy}\-{n}\-{\textsubring{d}}\-{a}\-{\textsubring{r}}\-{f}\-{j}\-{ou}\-{r}\-{ð}\-{u}\-{\ng}\-{\r{g}}\-{\textscy}\-{\textsubring{r}}\-} \dicPos{m}[6] \dicFlx{(‑s, ‑ar)}[8] \dicSynonym{kortér} \dicDirectTranslationCS{čtvrthodina}
\dicEntry[stundarsakir] \dicTerm{stundar··sakir} \dicIPA{{s}{\textsubring{d}}{\textscy}{n}{\textsubring{d}}{a}{\textsubring{r}}{s}{a}{\r{\textObardotlessj}}{\textsci}{\textsubring{r}}} \dicPos{f}[7] \dicFlx{pl}[20] \dicPhraseIS{um stundarsakir} \dicFlx{adv} \dicDirectTranslationCS{dočasně, prozatím} \dicExampleIS{dvelja e‑s staðar um stundarsakir} \dicExampleCS{přebývat (kde) dočasně}
\dicEntry[stundaskrá] \dicTerm{stunda··skrá} \dicIPA{{s}{\textsubring{d}}{\textscy}{n}{\textsubring{d}}{a}{s}{\r{g}}{r}{au}} \dicPos{f}[4] \dicFlx{(‑r\,/\addthin ‑ar, ‑r)}[21] \dicFieldCat{škol.} \dicDirectTranslationCS{rozvrh hodin}
\begin{xtolerant}{}{1pt}
\dicEntry[stundatafla] \dicTerm{stunda··|tafla}\addthinS\dicIPA{{s}{\textsubring{d}}{\textscy}{n}{\textsubring{d}}{a}{t\smash{\textsuperscript{h}}}{a}{\textsubring{b}}{l}{a}}\addthinS\dicPos{f}[1]\addthinS\dicFlx{(‑töflu, ‑töflur)}[20]\addthinS\dicFieldCat{škol.}\addthinS\dicSynonym{stun\-daskrá} \dicDirectTranslationCS{rozvrh hodin}
\end{xtolerant}
\dicEntry[stundi] \dicTerm{stundi} \dicIPA{{s}{\textsubring{d}}{\textscy}{n}{\textsubring{d}}{\textsci}} \dicPos{v} \dicFlx{ind pf sg 1 pers} \dicLink{stynja}
\dicEntry[stundlegur] \dicTerm{stund··legur} \dicIPA{{s}{\textsubring{d}}{\textscy}{n}{\textsubring{d}}{l}{\textepsilon}{\textbabygamma}{\textscy}{\textsubring{r}}} \dicPos{adj}[1]\dicFlx{}[-8] \dicDirectTranslationCS{dočasný, chvilkový, přechodný} \dicAntonym{eilífur};  \dicPhraseIS{hafa aldrei stundlegan frið} \dicLangCat{přen.} \dicDirectTranslationCS{nemít ani chvíli pokoj}
\dicEntry[stundum] \dicTerm{stundum\smash{\textsuperscript{1}}} \dicIPA{{s}{\textsubring{d}}{\textscy}{n}{\textsubring{d}}{\textscy}{\textsubring{m}}} \dicPos{v} \dicFlx{ind pf pl 1 pers} \dicLink{stynja}
\dicEntry[stundum] \dicTerm{stundum\smash{\textsuperscript{2}}} \dicsymFrequent\  \dicIPA{{s}{\textsubring{d}}{\textscy}{n}{\textsubring{d}}{\textscy}{\textsubring{m}}} \dicPos{adv} \dicDirectTranslationCS{občas, někdy, chvílemi} \dicExampleIS{Við hittumst stundum í kaffihúsinu.} \dicExampleCS{Někdy se potkáváme v~kavárně.}
\dicEntry[stundvís] \dicTerm{stund··vís} \dicIPA{{s}{\textsubring{d}}{\textscy}{n}{\textsubring{d}}{v}{i}{s}} \dicPos{adj}[5]\dicFlx{}[-1] \dicDirectTranslationCS{dochvilný, přesný}
\dicEntry[stundvísi] \dicTerm{stund··vís|i} \dicIPA{{s}{\textsubring{d}}{\textscy}{n}{\textsubring{d}}{v}{i}{s}{\textsci}} \dicPos{f}[3] \dicFlx{(‑i)}[3] \dicDirectTranslationCS{dochvilnost, přesnost}
\dicEntry[stunga] \dicTerm{stung|a} \dicIPA{{s}{\textsubring{d}}{u}{\ng}{\r{g}}{a}} \dicPos{f}[1] \dicFlx{(‑u, ‑ur)}[13] \textbf{1.} \dicSynonym*{sár} \dicDirectTranslationCS{bodná rána, bodnutí (zbraní ap.)} \dicExampleIS{sverðstunga} \dicExampleCS{bodnutí mečem}  \textbf{2.} \dicSynonym*{það að stinga} \dicDirectTranslationCS{bodání, píchání}  \textbf{3.} \dicSynonym{saumur} \dicDirectTranslationCS{steh}
\dicEntry[stungið] \dicTerm{stungið} \dicIPA{{s}{\textsubring{d}}{u}{\textltailn}{\r{\textObardotlessj}}{\textsci}{\texttheta}} \dicPos{v} \dicFlx{supin} \dicLink{stinga}
\dicEntry[stungum] \dicTerm{stungum} \dicIPA{{s}{\textsubring{d}}{u}{\ng}{\r{g}}{\textscy}{\textsubring{m}}} \dicPos{v} \dicFlx{ind pf pl 1 pers} \dicLink{stinga}
\dicEntry[stungumý] \dicTerm{stungu··mý} \dicIPA{{s}{\textsubring{d}}{u}{\ng}{\r{g}}{\textscy}{m}{i}} \dicPos{n}[2] \dicFlx{(‑s, ‑)}[13] \dicFieldCat{zool.} \dicDirectTranslationCS{komár} \textit{(l.~{\textLA{Culicidae}})}
\dicEntry[stunið] \dicTerm{stunið} \dicIPA{{s}{\textsubring{d}}{\textscy}{\textlengthmark}{n}{\textsci}{\texttheta}} \dicPos{v} \dicFlx{supin} \dicLink{stynja}
\dicEntry[sturla] \dicTerm{sturl|a} \dicIPA{{s}{\textsubring{d}}{\textscy}{r}{\textsubring{d}}{l}{a}} \dicPos{v}[1] \dicFlx{(‑aði)}[1] \dicSynonym{rugla} \dicDirectTranslationCS{zmást, pomást};  \dicIdiom{sturlast}{ \dicPhraseIS{sturlast}} \dicFlx{refl} \dicDirectTranslationCS{pominout se, zešílet}
\dicEntry[sturlaður] \dicTerm{sturl··|aður} \dicIPA{{s}{\textsubring{d}}{\textscy}{r}{\textsubring{d}}{l}{a}{ð}{\textscy}{\textsubring{r}}} \dicPos{adj}[3] \dicFlx{(f ‑uð)}[3] \dicDirectTranslationCS{pominutý, šílený, nepříčetný}
\dicEntry[sturlun] \dicTerm{sturl|un} \dicIPA{{s}{\textsubring{d}}{\textscy}{r}{\textsubring{d}}{l}{\textscy}{\textsubring{n}}} \dicPos{f}[7] \dicFlx{(‑unar)}[9] \dicDirectTranslationCS{pominutí, šílenství, nepříčetnost};  \dicPhraseIS{e‑m liggur við sturlun} \dicFlx{impers} \dicDirectTranslationCS{(kdo) se pomate, (kdo) zešílí}
\dicEntry[sturta] \dicTerm{sturt|a\smash{\textsuperscript{1}}} \dicIPA{{s}{\textsubring{d}}{\textscy}{\textsubring{r}}{\textsubring{d}}{a}} \dicPos{f}[1] \dicFlx{(‑u, ‑ur)}[19] \textbf{1.} \dicSynonym{steypibað} \dicDirectTranslationCS{sprcha} \dicExampleIS{köld sturta} \dicExampleCS{studená sprcha};  \dicPhraseIS{fara í sturtu} \dicDirectTranslationCS{dát si sprchu, osprchovat se}  \textbf{2.} \dicDirectTranslationCS{sklápěč, vyklápěč} \dicIndirectTranslationCS{(část vozu)}
\dicEntry[sturta] \dicTerm{sturt|a\smash{\textsuperscript{2}}} \dicIPA{{s}{\textsubring{d}}{\textscy}{\textsubring{r}}{\textsubring{d}}{a}} \dicPos{v}[1] \dicFlx{(‑aði)}[1] \dicFlx{dat} \textbf{1.} \dicFlx{dat} \dicDirectTranslationCS{sklopit, sklápět, vyklopit, vyklápět} \dicExampleIS{Vörubíllinn sturtaði sandi.} \dicExampleCS{Nákladní auto vyklopilo písek.}  \textbf{2.} \dicDirectTranslationCS{spláchnout\,/\addthin splachovat (záchod)} \dicExampleIS{Ertu búinn að sturta?} \dicExampleCS{Spláchl jsi?};  \dicPhraseIS{sturta niður} \dicDirectTranslationCS{spláchnout\,/\addthin splachovat (záchod)};  \dicIdiom{sturtast}{ \dicPhraseIS{sturtast}} \dicFlx{refl} \dicDirectTranslationCS{vyklopit se, vypustit se}
\dicEntry[sturtuklefi] \dicTerm{sturtu··klef|i} \dicIPA{{s}{\textsubring{d}}{\textscy}{\textsubring{r}}{\textsubring{d}}{\textscy}{k\smash{\textsuperscript{h}}}{l}{\textepsilon}{v}{\textsci}} \dicPos{m}[1] \dicFlx{(‑a, ‑ar)}[1] \dicDirectTranslationCS{sprchový kout}
\dicEntry[sturtusápa] \dicTerm{sturtu··sáp|a} \dicIPA{{s}{\textsubring{d}}{\textscy}{\textsubring{r}}{\textsubring{d}}{\textscy}{s}{au}{\textsubring{b}}{a}} \dicPos{f}[1] \dicFlx{(‑u, ‑ur)}[19] \dicDirectTranslationCS{sprchový gel, sprchové mýdlo}
\dicEntry[stutt] \dicTerm{stutt\smash{\textsuperscript{1}}} \dicIPA{{s}{\textsubring{d}}{\textscy}{h}{\textsubring{d}}} \dicPos{v} \dicFlx{supin} \dicLink{styðja}
\dicEntry[stutt] \dicTerm{stutt\smash{\textsuperscript{2}}} \dicIPA{{s}{\textsubring{d}}{\textscy}{h}{\textsubring{d}}} \dicPos{adv} \dicFlx{(comp styttra, sup styst)} \textbf{1.} \dicDirectTranslationCS{krátce, krátkodobě}  \textbf{2.} \dicDirectTranslationCS{nedaleko, blízko}
\dicEntry[stuttaralegur] \dicTerm{stuttara··legur} \dicIPA{{s}{\textsubring{d}}{\textscy}{h}{\textsubring{d}}{a}{r}{a}{l}{\textepsilon}{\textbabygamma}{\textscy}{\textsubring{r}}} \dicPos{adj}[1]\dicFlx{}[-8] \textbf{1.} \dicSynonym{stuttorður} \dicDirectTranslationCS{lakonický, úsečný}  \textbf{2.} \dicSynonym*{kaldur í viðmóti} \dicDirectTranslationCS{odměřený, strohý} \dicExampleIS{stuttaralegur koss} \dicExampleCS{odměřený polibek}
\dicEntry[stuttbuxur] \dicTerm{stutt··buxur} \dicIPA{{s}{\textsubring{d}}{\textscy}{h}{\textsubring{d}}{\textsubring{b}}{\textscy}{x}{s}{\textscy}{\textsubring{r}}} \dicPos{f}[12] \dicFlx{pl}[7] \dicDirectTranslationCS{šortky, kraťasy}
\dicEntry[stuttbylgja] \dicTerm{stutt··bylgj|a} \dicIPA{{s}{\textsubring{d}}{\textscy}{h}{\textsubring{d}}{\textsubring{b}}{\textsci}{l}{\r{\textObardotlessj}}{a}} \dicPos{f}[1] \dicFlx{(‑u, ‑ur)}[25] \dicFieldCat{fyz.} \dicDirectTranslationCS{krátká vlna} \dicAntonym{langbylgja}
\dicEntry[stuttbylgjustöð] \dicTerm{stutt·bylgju··stöð} \dicIPA{{s}{\textsubring{d}}{\textscy}{h}{\textsubring{d}}{\textsubring{b}}{\textsci}{l}{\r{\textObardotlessj}}{\textscy}{s}{\textsubring{d}}{\oe}{\texttheta}} \dicPos{f}[6] \dicFlx{(‑var, ‑var)}[1] \dicFieldCat{techn.} \dicDirectTranslationCS{krátkovlnný vysílač}
\dicEntry[stutterma] \dicTerm{stutt··erma} \dicIPA{{s}{\textsubring{d}}{\textscy}{h}{\textsubring{d}}{\textepsilon}{r}{m}{a}} \dicPos{adj}[13] \dicFlx{indecl}[1] \dicDirectTranslationCS{(jsoucí) s~krátkým rukávem}
\dicEntry[stuttermabolur] \dicTerm{stutt·erma··bol|ur} \dicIPA{{s}{\textsubring{d}}{\textscy}{h}{\textsubring{d}}{\textepsilon}{r}{m}{a}{\textsubring{b}}{\textopeno}{l}{\textscy}{\textsubring{r}}} \dicPos{m}[9] \dicFlx{(‑s, ‑ir)}[8] \dicDirectTranslationCS{krátké tričko, tričko s~krátkým rukávem}
\dicEntry[stutthærður] \dicTerm{stutt··hærður} \dicIPA{{s}{t\smash{\textsuperscript{h}}}{\textscy}{h}{\textsubring{d}}{h}{a}{i}{r}{ð}{\textscy}{\textsubring{r}}} \dicPos{adj}[2]\dicFlx{}[-1] \dicDirectTranslationCS{krátkovlasý}
\dicEntry[stuttklipptur] \dicTerm{stutt··klipptur} \dicIPA{{s}{\textsubring{d}}{\textscy}{h}{\textsubring{d}}{k\smash{\textsuperscript{h}}}{l}{\textsci}{f}{\textsubring{d}}{\textscy}{\textsubring{r}}} \dicPos{adj}[1]\dicFlx{}[-10] \dicDirectTranslationCS{nakrátko ostříhaný}
\dicEntry[stuttlega] \dicTerm{stutt··lega} \dicIPA{{s}{\textsubring{d}}{\textscy}{h}{\textsubring{d}}{l}{\textepsilon}{\textbabygamma}{a}} \dicPos{adv} \textbf{1.} \dicSynonym*{í stuttu máli} \dicDirectTranslationCS{krátce, slovem}  \textbf{2.} \dicSynonym*{óvingjarnlega} \dicDirectTranslationCS{úsečně, stroze}
\begin{xtolerant}{}{1pt}
\dicEntry[stuttlegur] \dicTerm{stutt··legur}\addthinS\dicIPA{{s}{\textsubring{d}}{\textscy}{h}{\textsubring{d}}{l}{\textepsilon}{\textbabygamma}{\textscy}{\textsubring{r}}}\addthinS\dicPos{adj}[1]\dicFlx{}[-8]\addthinS\dicSynonym{stuttur} \dicDirectTranslationCS{krátký, struč\-ný}
\end{xtolerant}
\dicEntry[stuttmynd] \dicTerm{stutt··mynd} \dicIPA{{s}{\textsubring{d}}{\textscy}{h}{\textsubring{d}}{m}{\textsci}{n}{\textsubring{d}}} \dicPos{f}[7] \dicFlx{(‑ar, ‑ir)}[1] \dicDirectTranslationCS{krátký film}
\dicEntry[stuttnefja] \dicTerm{stutt··nefj|a} \dicIPA{{s}{\textsubring{d}}{\textscy}{h}{\textsubring{d}}{n}{\textepsilon}{v}{j}{a}} \dicPos{f}[1] \dicFlx{(‑u, ‑ur)}[7] \dicFieldCat{zool.} \dicDirectTranslationCS{alkoun tlustozobý} \textit{(l.~{\textLA{Uria lomvia}})}  \dicsymPhoto\ 
\dicFigure{20873.jpg}{Stuttnefja}{Stuttnefja - Sowls, Art, Biolib, PD}
\dicEntry[stuttorður] \dicTerm{stutt··orður} \dicIPA{{s}{\textsubring{d}}{\textscy}{h}{\textsubring{d}}{\textopeno}{r}{ð}{\textscy}{\textsubring{r}}} \dicPos{adj}[2]\dicFlx{}[-1] \dicDirectTranslationCS{stručný, strohý} \dicExampleIS{stuttorður rithöfundur} \dicExampleCS{stručný autor}
\dicEntry[stuttpils] \dicTerm{stutt··pils} \dicIPA{{s}{\textsubring{d}}{\textscy}{h}{\textsubring{d}}{p\smash{\textsuperscript{h}}}{\textsci}{l}{s}} \dicPos{n}[2] \dicFlx{(‑, ‑)}[22] \dicDirectTranslationCS{krátká sukně, minisukně}
\dicEntry[stuttur] \dicTerm{stuttur} \dicsymFrequent\  \dicIPA{{s}{\textsubring{d}}{\textscy}{h}{\textsubring{d}}{\textscy}{\textsubring{r}}} \dicPos{adj}[10] \dicFlx{(comp styttri, sup stystur)}[9] \textbf{1.} \dicDirectTranslationCS{krátký} \dicIndirectTranslationCS{(vzdálenostně)} \dicExampleIS{stutt vegalengd} \dicExampleCS{krátká vzdálenost};  \dicPhraseIS{það er stutt í e‑u} \dicDirectTranslationCS{je to nedaleko od (čeho) (obchodu ap.)}  \textbf{2.} \dicDirectTranslationCS{krátký} \dicIndirectTranslationCS{(časově)} \dicExampleIS{stutta stund} \dicExampleCS{krátce};  \dicPhraseIS{það er stutt til e‑s} \dicFlx{impers} \dicDirectTranslationCS{(co) se blíží (Velikonoce ap.)};  \dicIdiom{stuttur}{ \dicPhraseIS{fyrir stuttu}} \dicFlx{adv} \dicDirectTranslationCS{nedávno, v~nedávné době};  \dicPhraseIS{í stuttu máli} \dicFlx{adv} \dicDirectTranslationCS{krátce řečeno, zkrátka};  \dicPhraseIS{stuttur í spuna} \dicFlx{adj} \dicDirectTranslationCS{popudlivý}
\dicEntry[stú] \dicTerm{stú} \dicIPA{{s}{\textsubring{d}}{u}{\textlengthmark}} \dicPos{n}[2] \dicFlx{(‑s)}[2] \dicPhraseIS{það er allt á rúi og stúi} \dicLangCat{přen.} \dicDirectTranslationCS{všechno je vzhůru nohama}
\dicEntry[stúdent] \dicTerm{stúdent} \dicsymFrequent\  \dicIPA{{s}{\textsubring{d}}{u}{\textlengthmark}{\textsubring{d}}{\textepsilon}{\textsubring{n}}{\textsubring{d}}} \dicPos{m}[4] \dicFlx{(‑s, ‑ar)}[10] \dicDirectTranslationCS{student(ka)} \dicExampleIS{háskólastúdent} \dicExampleCS{vysokoškolský student}
\dicEntry[stúdentagarður] \dicTerm{stúdenta··garð|ur} \dicIPA{{s}{\textsubring{d}}{u}{\textlengthmark}{\textsubring{d}}{\textepsilon}{\textsubring{n}}{\textsubring{d}}{a}{\r{g}}{a}{r}{ð}{\textscy}{\textsubring{r}}} \dicPos{m}[6] \dicFlx{(‑s, ‑ar)}[5] \dicDirectTranslationCS{(vysokoškolská) kolej, internát, (studentská) ubytovna} \dicExampleIS{búa á stúdentagarði} \dicExampleCS{bydlet na koleji}
\dicEntry[stúdentahreyfing] \dicTerm{stúdenta··hreyf·ing} \dicIPA{{s}{\textsubring{d}}{u}{\textlengthmark}{\textsubring{d}}{\textepsilon}{\textsubring{n}}{\textsubring{d}}{a}{\textsubring{r}}{ei}{v}{i}{\ng}{\r{g}}} \dicPos{f}[4] \dicFlx{(‑ar, ‑ar)}[5] \dicDirectTranslationCS{studentské hnutí}
\dicEntry[stúdentaráð] \dicTerm{stúdenta··ráð} \dicIPA{{s}{\textsubring{d}}{u}{\textlengthmark}{\textsubring{d}}{\textepsilon}{\textsubring{n}}{\textsubring{d}}{a}{r}{au}{\texttheta}} \dicPos{n}[2] \dicFlx{(‑s, ‑)}[5] \dicDirectTranslationCS{studentská rada}
\dicEntry[stúdentaskipti] \dicTerm{stúdenta··skipti} \dicIPA{{s}{\textsubring{d}}{u}{\textlengthmark}{\textsubring{d}}{\textepsilon}{\textsubring{n}}{\textsubring{d}}{a}{s}{\r{\textObardotlessj}}{\textsci}{f}{\textsubring{d}}{\textsci}} \dicPos{n}[2] \dicFlx{pl}[19] \dicDirectTranslationCS{výměna studentů}
\dicEntry[stúdentspróf] \dicTerm{stúdents··próf} \dicIPA{{s}{\textsubring{d}}{u}{\textlengthmark}{\textsubring{d}}{\textepsilon}{\textsubring{n}}{\textsubring{d}}{s}{p\smash{\textsuperscript{h}}}{r}{ou}{f}} \dicPos{n}[2] \dicFlx{(‑s, ‑)}[5] \dicFieldCat{škol.} \dicDirectTranslationCS{maturita}
\dicEntry[stúdera] \dicTerm{stúder|a}\dicTerm{, stúdéra} \dicIPA{{s}\-{\textsubring{d}}\-{u}\-{\textlengthmark}\-{\textsubring{d}}\-{\textepsilon}\-{r}\-{a}\-} \dicPos{v}[1] \dicFlx{(‑aði)}[1] \dicFlx{acc} \textbf{1.} \dicLangCat{hovor.} \dicSynonym*{stunda nám} \dicDirectTranslationCS{studovat, učit se} \dicExampleIS{Hann stúderaði í Frakklandi.} \dicExampleCS{Studoval ve Francii.}  \textbf{2.} \dicLangCat{hovor.} \dicSynonym{íhuga} \dicDirectTranslationCS{(pro)studovat, hloubat} \dicExampleIS{stúdera dulspeki} \dicExampleCS{studovat mysticismus}
\dicEntry[stúdéra] \dicTerm{stúdér|a} \dicIPA{{s}{\textsubring{d}}{u}{\textlengthmark}{\textsubring{d}}{j}{\textepsilon}{r}{a}} \dicPos{v}[1] \dicFlx{(‑aði)}[1] \dicLink{stúdera}
\dicEntry[stúdíó] \dicTerm{stúdíó} \dicIPA{{s}{\textsubring{d}}{u}{\textlengthmark}{\textsubring{d}}{i}{j}{ou}} \dicPos{n}[2] \dicFlx{(‑s, ‑)}[5] \textbf{1.} \dicDirectTranslationCS{(nahrávací) studio}  \textbf{2.} \dicSynonym{vinnustofa} \dicDirectTranslationCS{ateliér, pracovna}  \textbf{3.} \dicLangCat{hovor.} \dicSynonym{stúdíóíbúð} \dicDirectTranslationCS{ateliérový byt}
\dicEntry[stúdíóíbúð] \dicTerm{stúdíó··í·búð} \dicIPA{{s}{\textsubring{d}}{u}{\textlengthmark}{\textsubring{d}}{i}{j}{ou}{i}{\textsubring{b}}{u}{\texttheta}} \dicPos{f}[7] \dicFlx{(‑ar, ‑ir)}[1] \dicDirectTranslationCS{ateliérový byt}
\dicEntry[stúfur] \dicTerm{stúf|ur\smash{\textsuperscript{1}}} \dicIPA{{s}{\textsubring{d}}{u}{\textlengthmark}{v}{\textscy}{\textsubring{r}}} \dicPos{m}[6] \dicFlx{(‑s, ‑ar)}[22] \textbf{1.} \dicSynonym{stubbur} \dicDirectTranslationCS{pahýl, pařez, zbytek (po odříznutí)} \dicExampleIS{fótarstúfur} \dicExampleCS{nožní pahýl}  \textbf{2.} \dicSynonym*{lítill drengur} \dicDirectTranslationCS{hošíček, mrňousek};  \dicIdiom{stúfur}{ \dicPhraseIS{e‑að stingur í stúf við e‑ð}} \dicDirectTranslationCS{(co) si s~(čím) odporuje, (co) s~(čím) kontrastuje};  \dicPhraseIS{fara á stúfana með e‑ð} \dicDirectTranslationCS{zahájit (co), uvést do chodu (co)}
\dicEntry[Stúfur] \dicTerm{Stúf|ur\smash{\textsuperscript{2}}} \dicIPA{{s}{\textsubring{d}}{u}{\textlengthmark}{v}{\textscy}{\textsubring{r}}} \dicPos{m}[6] \dicFlx{(‑s)}[18] \dicFlx{prop} \dicFieldCat{pov.} \dicDirectTranslationCS{Stúfur} \dicIndirectTranslationCS{(3. vánoční skřítek, který vyjídá připálené zbytky jídla z~pánví; přichází 14. prosince)}
\dicEntry[stúka] \dicTerm{stúk|a\smash{\textsuperscript{1}}} \dicIPA{{s}{\textsubring{d}}{u}{\textlengthmark}{\r{g}}{a}} \dicPos{f}[1] \dicFlx{(‑u, ‑ur)}[19] \textbf{1.} \dicSynonym*{sæti á palli í leikhúsi} \dicDirectTranslationCS{lóže (v~divadle ap.)}  \textbf{2.} \dicSynonym*{áhorfendasvæði} \dicDirectTranslationCS{tribuna, hlediště}  \textbf{3.} \dicSynonym{klúbbur} \dicDirectTranslationCS{lóže, klub}  \textbf{4.} \dicFieldCat{anat.} \dicDirectTranslationCS{thalamus, talamus}
\dicEntry[stúka] \dicTerm{stúk|a\smash{\textsuperscript{2}}} \dicIPA{{s}{\textsubring{d}}{u}{\textlengthmark}{\r{g}}{a}} \dicPos{v}[1] \dicFlx{(‑aði)}[1] \dicPhraseIS{stúka e‑ð af} \dicDirectTranslationCS{oddělit (co) přepážkou}
\dicEntry[stúlka] \dicTerm{stúlk|a} \dicsymFrequent\  \dicIPA{{s}{\textsubring{d}}{u}{\textsubring{l}}{\r{g}}{a}} \dicPos{f}[1] \dicFlx{(‑u, ‑ur)}[13] \dicDirectTranslationCS{dívka, holka, děvče} \dicExampleIS{hópur ungra stúlkna} \dicExampleCS{skupina mladých děvčat}
\dicEntry[stúlknaskóli] \dicTerm{stúlkna··skól|i} \dicIPA{{s}{\textsubring{d}}{u}{\textsubring{l}}{\r{g}}{n}{a}{s}{\r{g}}{ou}{l}{\textsci}} \dicPos{m}[1] \dicFlx{(‑a, ‑ar)}[1] \dicDirectTranslationCS{dívčí škola}
\dicEntry[stúlkubarn] \dicTerm{stúlku··|barn} \dicIPA{{s}{\textsubring{d}}{u}{\textsubring{l}}{\r{g}}{\textscy}{\textsubring{b}}{a}{r}{\textsubring{d}}{\textsubring{n}}} \dicPos{n}[2] \dicFlx{(‑barns, ‑börn)}[8] \dicDirectTranslationCS{holčička, děvčátko}
\dicEntry[stúrinn] \dicTerm{stúrinn} \dicIPA{{s}{\textsubring{d}}{u}{\textlengthmark}{r}{\textsci}{\textsubring{n}}} \dicPos{adj}[6]\dicFlx{}[-6] \dicSynonym{sorgbitinn} \dicDirectTranslationCS{zarmoucený, skleslý, sklíčený}
\dicEntry[stúss] \dicTerm{stúss} \dicIPA{{s}{\textsubring{d}}{u}{s}{\textlengthmark}} \dicPos{n}[2] \dicFlx{(‑)}[24] \dicSynonym{umstang} \dicDirectTranslationCS{shon, honička}
\dicEntry[stússa] \dicTerm{stúss|a} \dicIPA{{s}{\textsubring{d}}{u}{s}{\textlengthmark}{a}} \dicPos{v}[1] \dicFlx{(‑aði)}[1] \dicDirectTranslationCS{zaobírat se, činit se};  \dicIdiom{stússa}[í]{ \dicPhraseIS{stússa í að (gera e‑ð)}} \dicDirectTranslationCS{zaobírat se (děláním (čeho))};  \dicIdiom{stússast}{ \dicPhraseIS{stússast}} \dicFlx{refl} \dicDirectTranslationCS{zaobírat se, činit se}
\dicEntry[stúta] \dicTerm{stút|a} \dicIPA{{s}{\textsubring{d}}{u}{\textlengthmark}{\textsubring{d}}{a}} \dicPos{v}[1] \dicFlx{(‑aði)}[1] \dicFlx{dat} \textbf{1.} \dicLangCat{hovor.} \dicSynonym{drepa} \dicDirectTranslationCS{vyřídit, oddělat} \dicIndirectTranslationCS{(zabít)} \dicExampleIS{stúta e‑m} \dicExampleCS{oddělat (koho)}  \textbf{2.} \dicLangCat{hovor.} \dicDirectTranslationCS{roztříštit, rozbít}  \textbf{3.} \dicLangCat{hovor.} \dicDirectTranslationCS{vyžahnout, vyzunknout, vylohnit};  \dicPhraseIS{stúta sig} \dicDirectTranslationCS{přihnout si, loknout si (alkoholu)}
\dicEntry[stútfullur] \dicTerm{stút··fullur} \dicIPA{{s}{\textsubring{d}}{u}{\textlengthmark}{\textsubring{d}}{f}{\textscy}{\textsubring{d}}{l}{\textscy}{\textsubring{r}}} \dicPos{adj}[10]\dicFlx{}[-8] \textbf{1.} \dicSynonym{troðfullur} \dicDirectTranslationCS{naplněný po okraj, napěchovaný}  \textbf{2.} \dicSynonym{pakksaddur} \dicDirectTranslationCS{najedený, nacpaný, přecpaný} \dicIndirectTranslationCS{(nemohoucí sníst více)}
\dicEntry[stútur] \dicTerm{stút|ur} \dicIPA{{s}{\textsubring{d}}{u}{\textlengthmark}{\textsubring{d}}{\textscy}{\textsubring{r}}} \dicPos{m}[6] \dicFlx{(‑s, ‑ar)}[22] \textbf{1.} \dicDirectTranslationCS{hrdlo} \dicIndirectTranslationCS{(láhve)};  \dicPhraseIS{drekka af stút} \dicDirectTranslationCS{pít z~láhve}  \textbf{2.} \dicDirectTranslationCS{hubice, hubička (konvice ap.)};  \dicIdiom{stútur}{ \dicPhraseIS{setja stút á munninn}} \dicDirectTranslationCS{našpulit pusu}
\dicEntry[stybba] \dicTerm{stybb|a} \dicIPA{{s}{\textsubring{d}}{\textsci}{\textsubring{b}}{\textlengthmark}{a}} \dicPos{f}[1] \dicFlx{(‑u)}[5] \dicDirectTranslationCS{výpary, hustý kouř\,/\addthin dým} \dicExampleIS{bensínstybba} \dicExampleCS{benzínové výpary}
\dicEntry[styddi] \dicTerm{styddi} \dicIPA{{s}{\textsubring{d}}{\textsci}{\textsubring{d}}{\textlengthmark}{\textsci}} \dicPos{v} \dicFlx{con pf sg 1 pers} \dicLink{styðja}
\dicEntry[styð] \dicTerm{styð} \dicIPA{{s}{\textsubring{d}}{\textsci}{\textlengthmark}{\texttheta}} \dicPos{v} \dicFlx{ind praes sg 1 pers} \dicLink{styðja}
\dicEntry[styðja] \dicTerm{styðja} \dicsymFrequent\  \dicIPA{{s}{\textsubring{d}}{\textsci}{ð}{j}{a}} \dicPos{v}[4] \dicFlx{(styð, studdi, studdum, styddi, stutt)}[40] \dicFlx{acc\,/\addthin dat} \textbf{1.} \dicFlx{acc} \dicDirectTranslationCS{podepřít, podpírat, přidržet, přidržovat} \dicExampleIS{styðja gömlu konuna út í bílinn} \dicExampleCS{podepřít starší ženu při vystupování z~auta};  \dicPhraseIS{styðja sig við staf} \dicDirectTranslationCS{podpírat se hůlkou}  \textbf{2.} \dicFlx{acc} \dicSynonym{fylgja\smash{\textsuperscript{2}}} \dicDirectTranslationCS{podpořit, podporovat, podepřít, podpírat} \dicExampleIS{styðja stefnu flokksins} \dicExampleCS{podporovat směr politické strany};  \dicPhraseIS{styðja e‑n til e‑s} \dicDirectTranslationCS{podporovat (koho) v~(čem)}  \textbf{3.} \dicFlx{dat} \dicSynonym{þrýsta} \dicDirectTranslationCS{opřít, stlačit} \dicExampleIS{styðja e‑u á e‑ð} \dicExampleCS{opřít (co) o~(co)};  \dicIdiom{styðjast}{ \dicPhraseIS{styðjast}} \dicFlx{refl} \dicDirectTranslationCS{opřít se, opírat se, podepřít se, podepírat se};  \dicIdiom{styðjast}[við]{ \dicPhraseIS{styðjast við e‑ð}} \dicFlx{refl} \dicLangCat{přen.} \dicDirectTranslationCS{opírat se o~(co) (o~důkazy ap.)} \dicExampleIS{styðjast við ákveðna fyrirmynd} \dicExampleCS{opírat se o~určitou předlohu}
\dicEntry[styggð] \dicTerm{styggð} \dicIPA{{s}{\textsubring{d}}{\textsci}{\textbabygamma}{\texttheta}} \dicPos{f}[4] \dicFlx{(‑ar)}[3] \dicSynonym{hræðsla} \dicDirectTranslationCS{vyplašení, úprk ze strachu\,/\addthin z~vylekání};  \dicPhraseIS{það kemur styggð að e‑u} \dicFlx{impers} \dicDirectTranslationCS{(co) se plaší (ovce ap.)}
\dicEntry[styggja] \dicTerm{stygg|ja} \dicIPA{{s}{\textsubring{d}}{\textsci}{\r{\textObardotlessj}}{\textlengthmark}{a}} \dicPos{v}[2] \dicFlx{(‑ði, ‑t)}[89] \dicFlx{acc} \textbf{1.} \dicSynonym{fæla} \dicDirectTranslationCS{(vy)plašit, (vy)lekat} \dicExampleIS{styggja fé} \dicExampleCS{vyplašit ovce}  \textbf{2.} \dicSynonym{móðga} \dicDirectTranslationCS{popudit, popouzet, urazit, urážet, dotknout se} \dicExampleIS{styggja e‑n} \dicExampleCS{popudit (koho)};  \dicIdiom{styggjast}{ \dicPhraseIS{styggjast}} \dicFlx{refl} \dicSynonym*{verða hræddur} \dicDirectTranslationCS{vystrašit se, vyplašit se};  \dicIdiom{styggjast}[við]{ \dicPhraseIS{styggjast við}} \dicFlx{refl} \dicDirectTranslationCS{být dotčen, urazit se}
\dicEntry[stygglyndi] \dicTerm{stygg··lyndi} \dicIPA{{s}{\textsubring{d}}{\textsci}{\r{g}}{l}{\textsci}{n}{\textsubring{d}}{\textsci}} \dicPos{n}[2] \dicFlx{(‑s)}[20] \dicDirectTranslationCS{popudlivost, prchlivost}
\dicEntry[stygglyndur] \dicTerm{stygg··lyndur} \dicIPA{{s}{\textsubring{d}}{\textsci}{\r{g}}{l}{\textsci}{n}{\textsubring{d}}{\textscy}{\textsubring{r}}} \dicPos{adj}[2]\dicFlx{}[-14] \dicDirectTranslationCS{popudlivý, prchlivý}
\dicEntry[styggur] \dicTerm{styggur} \dicIPA{{s}{\textsubring{d}}{\textsci}{\r{g}}{\textlengthmark}{\textscy}{\textsubring{r}}} \dicPos{adj}[1]\dicFlx{}[-1] \textbf{1.} \dicSynonym{skapstyggur} \dicDirectTranslationCS{popuzený, dotčený}  \textbf{2.} \dicSynonym{fælinn} \dicDirectTranslationCS{vyplašený, polekaný}
\dicEntry[stykki] \dicTerm{stykki\smash{\textsuperscript{1}}} \dicIPA{{s}{\textsubring{d}}{\textsci}{h}{\r{\textObardotlessj}}{\textsci}} \dicPos{n}[2] \dicFlx{(‑s, ‑)}[16] \textbf{1.} \dicSynonym{partur} \dicDirectTranslationCS{část, kus, díl} \dicExampleIS{kjötstykki} \dicExampleCS{kus masa}  \textbf{2.} \dicLangCat{hovor.} \dicDirectTranslationCS{kus} \dicIndirectTranslationCS{(umělecké dílo)};  \dicIdiom{stykki}{ \dicPhraseIS{gera öll sín stykki e‑s staðar}} \dicLangCat{přen.} \dicSynonym*{hægja sér} \dicDirectTranslationCS{ulevit si (kde), vykakat se (kde)};  \dicPhraseIS{standa ekki í stykkinu} \dicLangCat{přen.} \dicDirectTranslationCS{neporadit si, nevědět si rady};  \dicPhraseIS{þegar til stykkisins kemur} \dicLangCat{přen.} \dicDirectTranslationCS{když přijde do tuhého}
\dicEntry[stykki] \dicTerm{stykki\smash{\textsuperscript{2}}} \dicIPA{{s}{\textsubring{d}}{\textsci}{h}{\r{\textObardotlessj}}{\textsci}} \dicPos{v} \dicFlx{con pf sg 1 pers} \dicLink{stökkva\smash{\textsuperscript{2}}}
\dicEntry[styn] \dicTerm{styn} \dicIPA{{s}{\textsubring{d}}{\textsci}{\textlengthmark}{\textsubring{n}}} \dicPos{v} \dicFlx{ind praes sg 1 pers} \dicLink{stynja}
\dicEntry[styndi] \dicTerm{styndi} \dicIPA{{s}{\textsubring{d}}{\textsci}{n}{\textsubring{d}}{\textsci}} \dicPos{v} \dicFlx{con pf sg 1 pers} \dicLink{stynja}
\dicEntry[styngi] \dicTerm{styngi} \dicIPA{{s}{\textsubring{d}}{\textsci}{\textltailn}{\r{\textObardotlessj}}{\textsci}} \dicPos{v} \dicFlx{con pf sg 1 pers} \dicLink{stinga}
\dicEntry[stynja] \dicTerm{stynja} \dicsymFrequent\  \dicIPA{{s}{\textsubring{d}}{\textsci}{n}{j}{a}} \dicPos{v}[4] \dicFlx{(styn, stundi, stundum, styndi, stunið)}[33] \dicDirectTranslationCS{(za)sténat, (za)úpět, (za)vzdychat} \dicExampleIS{stynja af sársauka} \dicExampleCS{sténat bolestí};  \dicIdiom{stynja}[upp]{ \dicPhraseIS{stynja e‑u upp}} \dicDirectTranslationCS{vykoktat (co) ze sebe}
\dicEntry[styr] \dicTerm{styr} \dicIPA{{s}{\textsubring{d}}{\textsci}{\textlengthmark}{\textsubring{r}}} \dicPos{m}[10] \dicFlx{(‑jar)}[26] \dicLangCat{básn.} \dicSynonym{stríð} \dicDirectTranslationCS{válka};  \dicPhraseIS{það stendur styr um e‑ð} \dicFlx{impers} \dicDirectTranslationCS{vedou se spory o~(co)}
\dicEntry[styrja] \dicTerm{styrj|a} \dicIPA{{s}{\textsubring{d}}{\textsci}{r}{j}{a}} \dicPos{f}[1] \dicFlx{(‑u, ‑ur)}[7] \dicFieldCat{zool.} \dicDirectTranslationCS{jeseter} \textit{(l.~{\textLA{Acipenser}})}  \dicsymPhoto\ 
\dicFigure{63331.jpg}{Styrja}{Styrja - Bukovský Jiří, Biolib, Copyright/CC-BY-NC}
\dicEntry[styrjuhrogn] \dicTerm{styrju··hrogn} \dicIPA{{s}{\textsubring{d}}{\textsci}{r}{j}{\textscy}{\textsubring{r}}{\textopeno}{\r{g}}{\textsubring{n}}} \dicPos{n}[2] \dicFlx{pl}[1] \dicSynonym{kavíar} \dicDirectTranslationCS{kaviár}
\dicEntry[styrjöld] \dicTerm{styrj··|öld} \dicsymFrequent\  \dicIPA{{s}{\textsubring{d}}{\textsci}{r}{j}{\oe}{l}{\textsubring{d}}} \dicPos{f}[7] \dicFlx{(‑aldar, ‑aldir)}[16] \dicSynonym{stríð} \dicDirectTranslationCS{válka} \dicExampleIS{heimsstyrjöld} \dicExampleCS{světová válka};  \dicPhraseIS{fyrri heimsstyrjöld} \dicDirectTranslationCS{prv\-ní světová válka};  \dicPhraseIS{seinni heimsstyrjöld} \dicDirectTranslationCS{druhá světová válka}
\dicEntry[styrkja] \dicTerm{styrk|ja} \dicsymFrequent\  \dicIPA{{s}{\textsubring{d}}{\textsci}{\textsubring{r}}{\r{\textObardotlessj}}{a}} \dicPos{v}[2] \dicFlx{(‑ti, ‑t)}[25] \dicFlx{acc} \textbf{1.} \dicSynonym*{gera sterkari} \dicDirectTranslationCS{posílit, posilovat, zesílit, zesilovat, posilnit, posilňovat}  \textbf{2.} \dicSynonym*{veita aðstoð} \dicDirectTranslationCS{dotovat, sponzorovat, přispět\,/\addthin přispívat (finančně)} \dicExampleIS{styrkja félagið fjárhagslega} \dicExampleCS{dotovat finančně společnost};  \dicPhraseIS{styrkja e‑n til náms} \dicDirectTranslationCS{financovat (komu) studia};  \dicIdiom{styrkjast}{ \dicPhraseIS{styrkjast}} \dicFlx{refl} \dicDirectTranslationCS{posílit se, zesílit se}
\dicEntry[styrkjandi] \dicTerm{styrkj··andi} \dicIPA{{s}{\textsubring{d}}{\textsci}{\textsubring{r}}{\r{\textObardotlessj}}{a}{n}{\textsubring{d}}{\textsci}} \dicPos{adj}[13] \dicFlx{indecl}[1] \dicDirectTranslationCS{posilující, posilňující, povzbuzující}
\dicEntry[styrkleiki] \dicTerm{styrk··leik|i}\dicTerm{, styrkleikur} \dicIPA{{s}\-{\textsubring{d}}\-{\textsci}\-{\textsubring{r}}\-{\r{g}}\-{l}\-{ei}\-{\r{\textObardotlessj}}\-{\textsci}\-} \dicPos{m}[1] \dicFlx{(‑a)}[3] \textbf{1.} \dicSynonym{styrkur\smash{\textsuperscript{1}}} \dicDirectTranslationCS{síla, silná stránka}  \textbf{2.} \dicSynonym{veigur} \dicDirectTranslationCS{síla, intenzita} \dicExampleIS{styrkleiki vínsins} \dicExampleCS{síla alkoholu}
\dicEntry[styrkleikur] \dicTerm{styrk··leik|ur} \dicIPA{{s}{\textsubring{d}}{\textsci}{\textsubring{r}}{\r{g}}{l}{ei}{\r{g}}{\textscy}{\textsubring{r}}} \dicPos{m}[6] \dicFlx{(‑s)}[17] \dicLink{styrkleiki}
\dicEntry[styrktarfélag] \dicTerm{styrktar··fé·|lag} \dicIPA{{s}{\textsubring{d}}{\textsci}{\textsubring{r}}{\textsubring{d}}{a}{\textsubring{r}}{f}{j}{\textepsilon}{l}{a}{x}} \dicPos{n}[2] \dicFlx{(‑lags, ‑lög)}[8] \dicDirectTranslationCS{příspěvková organizace}
\dicEntry[styrktarmaður] \dicTerm{styrktar··|maður} \dicIPA{{s}{\textsubring{d}}{\textsci}{\textsubring{r}}{\textsubring{d}}{a}{r}{m}{a}{ð}{\textscy}{\textsubring{r}}} \dicPos{m}[13] \dicFlx{(‑manns, ‑menn)}[2] \dicDirectTranslationCS{patron(ka), mecenáš(ka), sponzor(ka)}
\dicEntry[styrktarsjóður] \dicTerm{styrktar··sjóð|ur} \dicIPA{{s}{\textsubring{d}}{\textsci}{\textsubring{r}}{\textsubring{d}}{a}{\textsubring{r}}{s}{j}{ou}{ð}{\textscy}{\textsubring{r}}} \dicPos{m}[9] \dicFlx{(‑s, ‑ir)}[6] \dicDirectTranslationCS{příspěvkový\,/\addthin nadační fond}
\dicEntry[styrkur] \dicTerm{styrk|ur\smash{\textsuperscript{1}}} \dicsymFrequent\  \dicIPA{{s}{\textsubring{d}}{\textsci}{\textsubring{r}}{\r{g}}{\textscy}{\textsubring{r}}} \dicPos{m}[9] \dicFlx{(‑s, ‑ir)}[15] \textbf{1.} \dicSynonym{kraftur} \dicDirectTranslationCS{síla, intenzita, mohutnost}  \textbf{2.} \dicSynonym{fjárstyrkur} \dicDirectTranslationCS{(finanční) příspěvek\,/\addthin podpora, stipendium, subvence} \dicExampleIS{sækja um styrk} \dicExampleCS{požádat o~stipendium}
\dicEntry[styrkur] \dicTerm{styrkur\smash{\textsuperscript{2}}} \dicIPA{{s}{\textsubring{d}}{\textsci}{\textsubring{r}}{\r{g}}{\textscy}{\textsubring{r}}} \dicPos{adj}[1]\dicFlx{}[-1] \textbf{1.} \dicSynonym{sterkur} \dicDirectTranslationCS{silný, mohutný}  \textbf{2.} \dicDirectTranslationCS{silný, pevný (hlas ap.)}
\dicEntry[styrkþegi] \dicTerm{styrk··þeg|i} \dicIPA{{s}{\textsubring{d}}{\textsci}{\textsubring{r}}{\r{g}}{\texttheta}{ei}{\textsci}} \dicPos{m}[1] \dicFlx{(‑a, ‑ar)}[1] \dicDirectTranslationCS{příjemce\,/\addthin příjemkyně stipendia\,/\addthin subvence}
\dicEntry[styst] \dicTerm{styst} \dicIPA{{s}{\textsubring{d}}{\textsci}{s}{\textsubring{d}}} \dicPos{adv} \dicFlx{sup} \dicLink{stutt\smash{\textsuperscript{2}}}
\dicEntry[stystur] \dicTerm{stystur} \dicIPA{{s}{\textsubring{d}}{\textsci}{s}{\textsubring{d}}{\textscy}{\textsubring{r}}} \dicPos{adj}[1]\dicFlx{}[-10] \dicFlx{m sg nom sup} \dicLink{stuttur}
\dicEntry[stytta] \dicTerm{stytt|a\smash{\textsuperscript{1}}} \dicsymFrequent\  \dicIPA{{s}{\textsubring{d}}{\textsci}{h}{\textsubring{d}}{a}} \dicPos{f}[1] \dicFlx{(‑u, ‑ur)}[19] \textbf{1.} \dicSynonym*{myndastytta} \dicDirectTranslationCS{socha} \dicExampleIS{styttan af skáldinu} \dicExampleCS{socha básníka};  \dicPhraseIS{standa eins og stytta} \dicLangCat{přen.} \dicDirectTranslationCS{stát jako socha}  \textbf{2.} \dicSynonym{stuðningur} \dicDirectTranslationCS{opora, podpora};  \dicPhraseIS{vera stoð e‑rs og stytta} \dicDirectTranslationCS{být (komu) (životní) oporou}
\dicEntry[stytta] \dicTerm{stytt|a\smash{\textsuperscript{2}}} \dicsymFrequent\  \dicIPA{{s}{\textsubring{d}}{\textsci}{h}{\textsubring{d}}{a}} \dicPos{v}[2] \dicFlx{(‑i, ‑)}[3] \dicFlx{acc} \dicSynonym*{gera styttri} \dicDirectTranslationCS{zkrátit, zkracovat, zmenšit, zmenšovat} \dicExampleIS{stytta buxurnar} \dicExampleCS{zkrátit kalhoty};  \dicPhraseIS{stytta brotið} \dicFieldCat{mat.} \dicDirectTranslationCS{zkrátit zlomek};  \dicPhraseIS{stytta e‑m stundir} \dicSynonym*{skemmta e‑n} \dicDirectTranslationCS{(za)bavit (koho), (z)krátit (komu) chvíli} \dicExampleIS{stytta henni stundir með því að lesa henni söguna} \dicExampleCS{krátit jí chvíli čtením příběhu};  \dicPhraseIS{stytta sér aldur} \dicDirectTranslationCS{zkrátit si život, spáchat sebevraždu};  \dicPhraseIS{stytta sér leið} \dicDirectTranslationCS{zkrátit si cestu};  \dicIdiom{stytta}[upp]{ \dicPhraseIS{það styttir upp}} \dicFlx{impers} \dicDirectTranslationCS{přestává pršet\,/\addthin sněžit, ustává déšť} \dicExampleIS{Það stytti loks upp.} \dicExampleCS{Konečně přestalo pršet.};  \dicIdiom{styttast}{ \dicPhraseIS{styttast}} \dicFlx{refl} \dicSynonym*{verða styttri} \dicDirectTranslationCS{(z)krátit se} \dicExampleIS{Buxurnar styttust í þvotti.} \dicExampleCS{Kalhoty se při praní zkrátily.};  \dicIdiom{styttast}[í]{ \dicPhraseIS{það styttist í e‑ð}} \dicFlx{refl impers} \dicDirectTranslationCS{(co) se blíží} \dicExampleIS{Það styttist í jólin.} \dicExampleCS{Vánoce se blíží.}
\dicEntry[stytting] \dicTerm{stytt··ing} \dicIPA{{s}{\textsubring{d}}{\textsci}{h}{\textsubring{d}}{i}{\ng}{\r{g}}} \dicPos{f}[4] \dicFlx{(‑ar, ‑ar)}[5] \textbf{1.} \dicSynonym*{það að stytta} \dicDirectTranslationCS{zkrácení, zmenšení} \dicExampleIS{stytting vinnuvikunnar} \dicExampleCS{zkrácení pracovního týdne}  \textbf{2.} \dicFieldCat{jaz.} \dicDirectTranslationCS{zkratka, zkrácené slovo}  \textbf{3.} \dicFieldCat{mat.} \dicDirectTranslationCS{krácení}
\dicEntry[styttingur] \dicTerm{stytt··ing|ur} \dicIPA{{s}{\textsubring{d}}{\textsci}{h}{\textsubring{d}}{i}{\ng}{\r{g}}{\textscy}{\textsubring{r}}} \dicPos{m}[6] \dicFlx{(‑s)}[9] \dicPhraseIS{í styttingi} \dicFlx{adv} \dicDirectTranslationCS{odměřeně, úsečně} \dicExampleIS{kveðjast í styttingi} \dicExampleCS{rozloučit se odměřeně}
\dicEntry[styttra] \dicTerm{styttra} \dicIPA{{s}{\textsubring{d}}{\textsci}{h}{\textsubring{d}}{r}{a}} \dicPos{adv} \dicFlx{comp} \dicLink{stutt\smash{\textsuperscript{2}}}
\dicEntry[styttri] \dicTerm{styttri} \dicIPA{{s}{\textsubring{d}}{\textsci}{h}{\textsubring{d}}{r}{\textsci}} \dicPos{adj} \dicFlx{comp m} \dicLink{stuttur}
\dicEntry[styttur] \dicTerm{styttur} \dicIPA{{s}{\textsubring{d}}{\textsci}{h}{\textsubring{d}}{\textscy}{\textsubring{r}}} \dicPos{adj}[1]\dicFlx{}[-10] \dicDirectTranslationCS{(z)krácený} \dicExampleIS{stytt útgáfa} \dicExampleCS{zkrácená verze}
\dicEntry[stýfa] \dicTerm{stýf|a} \dicIPA{{s}{\textsubring{d}}{i}{\textlengthmark}{v}{a}} \dicPos{v}[2] \dicFlx{(‑ði, ‑t)}[105] \dicFlx{acc} \dicDirectTranslationCS{zkrátit, zkracovat, useknout};  \dicPhraseIS{stýfa e‑ð úr hnefa} \dicDirectTranslationCS{jíst (co) rukama}
\dicEntry[stýfður] \dicTerm{stýfður} \dicIPA{{s}{\textsubring{d}}{i}{v}{ð}{\textscy}{\textsubring{r}}} \dicPos{adj}[2]\dicFlx{}[-1] \dicDirectTranslationCS{(z)krácený, useknutý}
\dicEntry[stýra] \dicTerm{stýr|a} \dicsymFrequent\  \dicIPA{{s}{\textsubring{d}}{i}{\textlengthmark}{r}{a}} \dicPos{v}[2] \dicFlx{(‑ði, ‑t)}[99] \dicFlx{dat} \textbf{1.} \dicSynonym*{vera við stýri} \dicDirectTranslationCS{řídit, kormidlovat, manévrovat} \dicExampleIS{stýra bátnum} \dicExampleCS{kormidlovat loď}  \textbf{2.} \dicSynonym{stjórna} \dicDirectTranslationCS{řídit, vést (firmu ap.)};  \dicIdiom{stýra}[hjá]{ \dicPhraseIS{stýra hjá e‑u}} \dicLangCat{přen.} \dicSynonym*{forðast} \dicDirectTranslationCS{vyhnout se (čemu), minout (co)}
\dicEntry[stýri] \dicTerm{stýri} \dicsymFrequent\  \dicIPA{{s}{\textsubring{d}}{i}{\textlengthmark}{r}{\textsci}} \dicPos{n}[2] \dicFlx{(‑s, ‑)}[14] \textbf{1.} \dicDirectTranslationCS{volant} \dicExampleIS{á bak við stýrið} \dicExampleCS{za volantem}  \textbf{2.} \dicDirectTranslationCS{řídítka}  \textbf{3.} \dicDirectTranslationCS{kormidlo};  \dicPhraseIS{sitja við stýrið} \dicLangCat{přen.} \dicDirectTranslationCS{být u~kormidla}  \textbf{4.} \dicDirectTranslationCS{srdíčko, zlatíčko, pusinka} \dicIndirectTranslationCS{(o~dítěti ap.)} \dicExampleIS{litla stýrið} \dicExampleCS{malé zlatíčko}  \textbf{5.} \dicSynonym*{kattarrófa} \dicDirectTranslationCS{(kočičí) ocas}
\dicEntry[stýribúnaður] \dicTerm{stýri··bú·nað|ur} \dicIPA{{s}{\textsubring{d}}{i}{\textlengthmark}{r}{\textsci}{\textsubring{b}}{u}{n}{a}{ð}{\textscy}{\textsubring{r}}} \dicPos{m}[10] \dicFlx{(‑ar)}[9] \dicDirectTranslationCS{řídicí\,/\addthin ovládací mechanismus, ovládací zařízení}
\dicEntry[stýrifræði] \dicTerm{stýri··fræð|i} \dicIPA{{s}{\textsubring{d}}{i}{\textlengthmark}{r}{\textsci}{f}{r}{a}{i}{ð}{\textsci}} \dicPos{f}[3] \dicFlx{(‑i)}[3] \dicDirectTranslationCS{kybernetika}
\dicEntry[stýrihnappur] \dicTerm{stýri··hnapp|ur} \dicIPA{{s}{\textsubring{d}}{i}{\textlengthmark}{r}{\textsci}{\textsubring{n}}{a}{h}{\textsubring{b}}{\textscy}{\textsubring{r}}} \dicPos{m}[6] \dicFlx{(‑s, ‑ar)}[25] \dicFieldCat{poč.} \dicDirectTranslationCS{klávesa Ctrl}
\dicEntry[stýrikerfi] \dicTerm{stýri··kerfi} \dicIPA{{s}{\textsubring{d}}{i}{\textlengthmark}{r}{\textsci}{c\smash{\textsuperscript{h}}}{\textepsilon}{r}{v}{\textsci}} \dicPos{n}[2] \dicFlx{(‑s, ‑)}[14] \textbf{1.} \dicDirectTranslationCS{řídicí systém}  \textbf{2.} \dicFieldCat{poč.} \dicDirectTranslationCS{operační systém}
\dicEntry[stýrimaður] \dicTerm{stýri··|maður} \dicIPA{{s}{\textsubring{d}}{i}{\textlengthmark}{r}{\textsci}{m}{a}{ð}{\textscy}{\textsubring{r}}} \dicPos{m}[13] \dicFlx{(‑manns, ‑menn)}[2] \dicFieldCat{nám.} \dicDirectTranslationCS{kormidelník, kormidelnice}
\dicEntry[stýring] \dicTerm{stýr··ing} \dicIPA{{s}{\textsubring{d}}{i}{\textlengthmark}{r}{i}{\ng}{\r{g}}} \dicPos{f}[4] \dicFlx{(‑ar)}[7] \textbf{1.} \dicSynonym*{það að stýra} \dicDirectTranslationCS{řízení, ovládání} \dicExampleIS{stýring á e‑u} \dicExampleCS{řízení (čeho)}  \textbf{2.} \dicSynonym{stjórn} \dicDirectTranslationCS{vedení, řízení}
\dicEntry[stýripinni] \dicTerm{stýri··pinn|i} \dicIPA{{s}{\textsubring{d}}{i}{\textlengthmark}{r}{\textsci}{p\smash{\textsuperscript{h}}}{\textsci}{n}{\textsci}} \dicPos{m}[1] \dicFlx{(‑a, ‑ar)}[1] \dicFieldCat{poč.} \dicDirectTranslationCS{joystick, pákový ovladač}
\dicEntry[stýrishús] \dicTerm{stýris··hús} \dicIPA{{s}{\textsubring{d}}{i}{\textlengthmark}{r}{\textsci}{s}{h}{u}{s}} \dicPos{n}[2] \dicFlx{(‑s, ‑)}[5] \dicDirectTranslationCS{kormidelna, řídicí kabina}
\dicEntry[stæði] \dicTerm{stæði\smash{\textsuperscript{1}}} \dicIPA{{s}{\textsubring{d}}{a}{i}{\textlengthmark}{ð}{\textsci}} \dicPos{n}[2] \dicFlx{(‑s, ‑)}[14] \textbf{1.} \dicSynonym{bílastæði} \dicDirectTranslationCS{místo (na parkování), (parkovací) místo}  \textbf{2.} \dicSynonym*{þar sem staðið er} \dicDirectTranslationCS{místo, stanoviště} \dicExampleIS{stæði í íþróttavelli} \dicExampleCS{místo na stadioně}
\dicEntry[stæði] \dicTerm{stæði\smash{\textsuperscript{2}}} \dicIPA{{s}{\textsubring{d}}{a}{i}{\textlengthmark}{ð}{\textsci}} \dicPos{v} \dicFlx{con pf sg 1 pers} \dicLink{standa}
\begin{xtolerant}{}{1pt}
\dicEntry[stæðilegur] \dicTerm{stæði··legur}\addthinS\dicIPA{{s}{\textsubring{d}}{a}{i}{\textlengthmark}{ð}{\textsci}{l}{\textepsilon}{\textbabygamma}{\textscy}{\textsubring{r}}}\addthinS\dicPos{adj}[1]\dicFlx{}[-8]\addthinS\textbf{1.} \dicSynonym{gerðarlegur} \dicDirectTranslationCS{statný, mohutný} \dicExampleIS{stæðilegur maður} \dicExampleCS{mohutný člověk}  \textbf{2.} \dicSynonym{traustlegur} \dicDirectTranslationCS{majestátní, monumentální}
\end{xtolerant}
\dicEntry[stæður] \dicTerm{stæður} \dicIPA{{s}{\textsubring{d}}{a}{i}{\textlengthmark}{ð}{\textscy}{\textsubring{r}}} \dicPos{adj}[2]\dicFlx{}[-6] \dicPhraseIS{vera vel stæður} \dicDirectTranslationCS{prosperovat (hospodářsky), být dobře (finančně) zajištěný};  \dicPhraseIS{vera illa stæður} \dicDirectTranslationCS{upadat (hospodářsky), být špatně (finančně) zajištěný}
\dicEntry[stækja] \dicTerm{stækj|a} \dicIPA{{s}{\textsubring{d}}{a}{i}{\textlengthmark}{\r{\textObardotlessj}}{a}} \dicPos{f}[1] \dicFlx{(‑u)}[5] \textbf{1.} \dicSynonym{ódaunn} \dicDirectTranslationCS{puch, zápach, smrad} \dicExampleIS{stækja af mykjunni} \dicExampleCS{zápach hnoje}  \textbf{2.} \dicSynonym{svæla\smash{\textsuperscript{1}}} \dicDirectTranslationCS{dusivý kouř}
\dicEntry[stækka] \dicTerm{stækk|a} \dicsymFrequent\  \dicIPA{{s}{\textsubring{d}}{a}{i}{h}{\r{g}}{a}} \dicPos{v}[1] \dicFlx{(‑aði)}[1] \dicFlx{acc} \textbf{1.} \dicFlx{acc} \dicSynonym{auka\smash{\textsuperscript{1}}} \dicDirectTranslationCS{zvětšit, zvětšovat, rozšířit, rozšiřovat} \dicExampleIS{stækka herbergið} \dicExampleCS{zvětšit pokoj}  \textbf{2.} \dicSynonym*{aukast} \dicDirectTranslationCS{zvětšit se, zvětšovat se, rozšířit se, rozšiřovat se, expandovat} \dicExampleIS{Þorpið hefur stækkað um helming.} \dicExampleCS{Vesnice vyrostla o~polovinu.}
\dicEntry[stækkanlegur] \dicTerm{stækkan··legur} \dicIPA{{s}{\textsubring{d}}{a}{i}{h}{\r{g}}{a}{n}{l}{\textepsilon}{\textbabygamma}{\textscy}{\textsubring{r}}} \dicPos{adj}[1]\dicFlx{}[-6] \dicDirectTranslationCS{zvětšitelný, rozšiřitelný} \dicExampleIS{stækkanleg mynd} \dicExampleCS{zvětšitelný obrázek}
\dicEntry[stækkari] \dicTerm{stækk··ar|i} \dicIPA{{s}{\textsubring{d}}{a}{i}{h}{\r{g}}{a}{r}{\textsci}} \dicPos{m}[1] \dicFlx{(‑a, ‑ar)}[13] \dicDirectTranslationCS{zvětšovací přístroj, zvětšovák}
\dicEntry[stækkun] \dicTerm{stækk|un} \dicIPA{{s}{\textsubring{d}}{a}{i}{h}{\r{g}}{\textscy}{\textsubring{n}}} \dicPos{f}[7] \dicFlx{(‑unar, ‑anir)}[8] \dicDirectTranslationCS{zvětšování, zvětšení} \dicExampleIS{stækkun á húsinu} \dicExampleCS{zvětšování domu}
\dicEntry[stækkunargler] \dicTerm{stækkunar··gler} \dicIPA{{s}{\textsubring{d}}{a}{i}{h}{\r{g}}{\textscy}{n}{a}{r}{\r{g}}{l}{\textepsilon}{\textsubring{r}}} \dicPos{n}[2] \dicFlx{(‑s, ‑)}[6] \dicDirectTranslationCS{zvětšovací sklo, lupa}
\dicEntry[stækur] \dicTerm{stækur} \dicIPA{{s}{\textsubring{d}}{a}{i}{\textlengthmark}{\r{g}}{\textscy}{\textsubring{r}}} \dicPos{adj}[1]\dicFlx{}[-1] \textbf{1.} \dicSynonym{daunillur} \dicDirectTranslationCS{smradlavý, páchnoucí, zapáchající} \dicExampleIS{stækt kjöt} \dicExampleCS{páchnoucí maso}  \textbf{2.} \dicSynonym{ofstækisfullur} \dicDirectTranslationCS{fanatický, ultraradikální} \dicExampleIS{stækur þjóðernissinni} \dicExampleCS{ultraradikální nacionalista}
\dicEntry[stæla] \dicTerm{stæl|a\smash{\textsuperscript{1}}} \dicIPA{{s}{\textsubring{d}}{a}{i}{\textlengthmark}{l}{a}} \dicPos{f}[1] \dicFlx{(‑u, ‑ur)}[19] \dicSynonym{þræta\smash{\textsuperscript{1}}} \dicDirectTranslationCS{hádka, diskuze, svár} \dicIndirectTranslationCS{(zvláště v~pl)}
\dicEntry[stæla] \dicTerm{stæl|a\smash{\textsuperscript{2}}} \dicIPA{{s}{\textsubring{d}}{a}{i}{\textlengthmark}{l}{a}} \dicPos{v}[2] \dicFlx{(‑ti, ‑t)}[63] \dicFlx{acc} \dicSynonym*{gera sterkan} \dicDirectTranslationCS{utužit, utužovat, zocelit, zocelovat} \dicExampleIS{stæla líkamann} \dicExampleCS{utužit tělo};  \dicIdiom{stælast}{ \dicPhraseIS{stælast}} \dicFlx{refl} \dicDirectTranslationCS{utužit se, zocelit se}
\dicEntry[stæla] \dicTerm{stæl|a\smash{\textsuperscript{3}}} \dicIPA{{s}{\textsubring{d}}{a}{i}{\textlengthmark}{l}{a}} \dicPos{v}[2] \dicFlx{(‑di, ‑t)}[141] \dicFlx{acc} \textbf{1.} \dicFlx{acc} \dicSynonym*{líkja eftir} \dicDirectTranslationCS{napodobovat, imitovat} \dicExampleIS{stæla e‑n í e‑u} \dicExampleCS{napodobovat (koho) v~(čem)}  \textbf{2.} \dicSynonym*{rífast} \dicDirectTranslationCS{hádat se, přít se} \dicExampleIS{stæla um e‑ð} \dicExampleCS{přít se o~(co)}
\dicEntry[stæli] \dicTerm{stæli} \dicIPA{{s}{\textsubring{d}}{a}{i}{\textlengthmark}{l}{\textsci}} \dicPos{v} \dicFlx{con pf sg 1 pers} \dicLink{stela}
\dicEntry[stæling] \dicTerm{stæl··ing} \dicIPA{{s}{\textsubring{d}}{a}{i}{\textlengthmark}{l}{i}{\ng}{\r{g}}} \dicPos{f}[4] \dicFlx{(‑ar, ‑ar)}[5] \textbf{1.} \dicSynonym{herðing} \dicDirectTranslationCS{utužení, zocelení}  \textbf{2.} \dicSynonym{eftirlíking} \dicDirectTranslationCS{napodobenina, imitace}  \textbf{3.} \dicSynonym*{það að stæla} \dicDirectTranslationCS{napodobování, imitování} \dicExampleIS{stæling á rödd hans} \dicExampleCS{napodobování jeho hlasu}
\dicEntry[stæll] \dicTerm{stæl|l} \dicIPA{{s}{\textsubring{d}}{a}{i}{\textsubring{d}}{\textsubring{l}}} \dicPos{m}[6] \dicFlx{(‑s, ‑ar)}[48] \textbf{1.} \dicLangCat{hovor.} \dicDirectTranslationCS{styl, třída, formát};  \dicPhraseIS{með stæl} \dicFlx{adv} \dicDirectTranslationCS{se stylem, stylově}  \textbf{2.} \dicPhraseIS{stælar} \dicFlx{pl} \dicSynonym{ósvífni} \dicDirectTranslationCS{drzost, impertinence} \dicExampleIS{Vertu ekki með þessa stæla!} \dicExampleCS{Nebuď drzý!}
\dicEntry[stæltur] \dicTerm{stæltur} \dicIPA{{s}{\textsubring{d}}{a}{i}{\textsubring{l}}{\textsubring{d}}{\textscy}{\textsubring{r}}} \dicPos{adj}[1]\dicFlx{}[-10] \dicDirectTranslationCS{zocelený, utužený, zpevněný} \dicExampleIS{stæltir vöðvar} \dicExampleCS{zocelené svaly}
\dicEntry[stæra] \dicTerm{stær|a} \dicIPA{{s}{\textsubring{d}}{a}{i}{\textlengthmark}{r}{a}} \dicPos{v}[2] \dicFlx{(‑ði, ‑t)}[100] \dicFlx{acc} \dicDirectTranslationCS{povyšovat se, vyvyšovat se} \dicExampleIS{stæra sig af e‑u} \dicExampleCS{vyvyšovat se kvůli (čemu)}
\dicEntry[stærð] \dicTerm{stærð} \dicsymFrequent\  \dicIPA{{s}{\textsubring{d}}{a}{i}{r}{\texttheta}} \dicPos{f}[7] \dicFlx{(‑ar, ‑ir)}[1] \textbf{1.} \dicDirectTranslationCS{velikost, rozměr} \dicExampleIS{giska á stærðina} \dicExampleCS{hádat velikost};  \dicPhraseIS{af öllum stærðum og gerðum} \dicFlx{adj} \dicDirectTranslationCS{všech tvarů a~velikostí};  \dicPhraseIS{e‑að á stærð við e‑ð} \dicDirectTranslationCS{(co) velikosti (čeho)};  \dicPhraseIS{stærðarinnar (e‑að)} \dicFlx{adj} \dicDirectTranslationCS{obrovský\,/\addthin kolosální (co)} \dicExampleIS{stærðarinnar demanti} \dicExampleCS{obrovský démant}  \textbf{2.} \dicFieldCat{mat.} \dicDirectTranslationCS{hodnota, velikost} \dicExampleIS{föst stærð} \dicExampleCS{fixní hodnota}
\dicEntry[stærðargráða] \dicTerm{stærðar··gráð|a} \dicIPA{{s}{\textsubring{d}}{a}{i}{r}{ð}{a}{r}{\r{g}}{r}{au}{ð}{a}} \dicPos{f}[1] \dicFlx{(‑u, ‑ur)}[19] \dicSynonym{umfang} \dicDirectTranslationCS{řádová velikost\,/\addthin hodnota}
\dicEntry[stærðarmunur] \dicTerm{stærðar··mun|ur} \dicIPA{{s}{\textsubring{d}}{a}{i}{r}{ð}{a}{r}{m}{\textscy}{n}{\textscy}{\textsubring{r}}} \dicPos{m}[10] \dicFlx{(‑ar)}[12] \dicDirectTranslationCS{rozdíl ve velikosti, velikostní rozdíl}
\dicEntry[stærðfræði] \dicTerm{stærð··fræð|i} \dicIPA{{s}{\textsubring{d}}{a}{i}{r}{ð}{f}{r}{a}{i}{ð}{\textsci}} \dicPos{f}[3] \dicFlx{(‑i)}[3] \dicDirectTranslationCS{matematika}
\dicEntry[stærðfræðilegur] \dicTerm{stærð·fræði··legur} \dicIPA{{s}{\textsubring{d}}{a}{i}{r}{ð}{f}{r}{a}{i}{ð}{\textsci}{l}{\textepsilon}{\textbabygamma}{\textscy}{\textsubring{r}}} \dicPos{adj}[1]\dicFlx{}[-8] \dicDirectTranslationCS{matematický}
\dicEntry[stærðfræðingur] \dicTerm{stærð·fræð··ing|ur} \dicIPA{{s}{\textsubring{d}}{a}{i}{r}{ð}{f}{r}{a}{i}{ð}{i}{\ng}{\r{g}}{\textscy}{\textsubring{r}}} \dicPos{m}[6] \dicFlx{(‑s, ‑ar)}[8] \dicDirectTranslationCS{matematik, matematička}
\dicEntry[stærilátur] \dicTerm{stæri··látur} \dicIPA{{s}{\textsubring{d}}{a}{i}{\textlengthmark}{r}{\textsci}{l}{au}{\textsubring{d}}{\textscy}{\textsubring{r}}} \dicPos{adj}[1]\dicFlx{}[-1] \dicSynonym{hrokafullur} \dicDirectTranslationCS{pyšný, samolibý}
\dicEntry[stærilæti] \dicTerm{stæri··læti} \dicIPA{{s}{\textsubring{d}}{a}{i}{\textlengthmark}{r}{\textsci}{l}{a}{i}{\textsubring{d}}{\textsci}} \dicPos{n}[2] \dicFlx{(‑s)}[20] \dicDirectTranslationCS{pýcha, samolibost}
\dicEntry[stærri] \dicTerm{stærri} \dicIPA{{s}{\textsubring{d}}{a}{i}{r}{\textlengthmark}{\textsci}} \dicPos{adj} \dicFlx{comp} \dicLink{stór}
\dicEntry[stærstur] \dicTerm{stærstur} \dicIPA{{s}{\textsubring{d}}{a}{i}{\textsubring{r}}{s}{\textsubring{d}}{\textscy}{\textsubring{r}}} \dicPos{adj}[1]\dicFlx{}[-10] \dicFlx{m sg nom sup} \dicLink{stór}
\dicEntry[stöð] \dicTerm{stöð\smash{\textsuperscript{1}}} \dicsymFrequent\  \dicIPA{{s}{\textsubring{d}}{\oe}{\textlengthmark}{\texttheta}} \dicPos{f}[6] \dicFlx{(‑var, ‑var)}[1] \textbf{1.} \dicSynonym*{miðstöð starfsemi} \dicDirectTranslationCS{stanice, středisko, centrum} \dicExampleIS{gróðrarstöð} \dicExampleCS{zahradnictví}  \textbf{2.} \dicSynonym{stoppistöð} \dicDirectTranslationCS{stanice, stanoviště, zastávka} \dicExampleIS{járnbrautarstöð} \dicExampleCS{železniční stanice}  \textbf{3.} \dicSynonym{sjónvarpsstöð} \dicDirectTranslationCS{(televizní) stanice\,/\addthin kanál, (televizní) program}  \textbf{4.} \dicSynonym{útvarpsstöð} \dicDirectTranslationCS{(rádiová\,/\addthin rozhlasová) stanice, rádio}
\dicEntry[stöð] \dicTerm{stöð\smash{\textsuperscript{2}}} \dicIPA{{s}{\textsubring{d}}{\oe}{\textlengthmark}{\texttheta}} \dicPos{adj} \dicFlx{f sg nom pos} \dicLink{staður\smash{\textsuperscript{2}}}
\dicEntry[stöðlun] \dicTerm{stöðl|un} \dicIPA{{s}{\textsubring{d}}{\oe}{ð}{l}{\textscy}{\textsubring{n}}} \dicPos{f}[7] \dicFlx{(‑unar)}[12] \dicSynonym*{það að staðla} \dicDirectTranslationCS{standardizace}
\dicEntry[stöðnun] \dicTerm{stöðn|un} \dicIPA{{s}{\textsubring{d}}{\oe}{ð}{n}{\textscy}{\textsubring{n}}} \dicPos{f}[7] \dicFlx{(‑unar)}[12] \dicSynonym*{það að staðna} \dicDirectTranslationCS{stagnace} \dicExampleIS{koma í veg fyrir stöðnun} \dicExampleCS{předejít stagnaci}
\dicEntry[stöðu] \dicTerm{stöðu} \dicIPA{{s}{\textsubring{d}}{\oe}{\textlengthmark}{ð}{\textscy}} \dicPos{f} \dicFlx{sg gen} \dicLink{staða}
\dicEntry[stöðufræði] \dicTerm{stöðu··fræð|i} \dicIPA{{s}{\textsubring{d}}{\oe}{\textlengthmark}{ð}{\textscy}{f}{r}{a}{i}{ð}{\textsci}} \dicPos{f}[3] \dicFlx{(‑i)}[3] \dicDirectTranslationCS{statika} \dicIndirectTranslationCS{(obor mechaniky)}
\dicEntry[stöðugleiki] \dicTerm{stöðug··leik|i} \dicIPA{{s}{\textsubring{d}}{\oe}{\textlengthmark}{ð}{\textscy}{\textbabygamma}{l}{ei}{\r{\textObardotlessj}}{\textsci}} \dicPos{m}[1] \dicFlx{(‑a)}[3] \dicDirectTranslationCS{stabilita, stálost, setrvalost} \dicExampleIS{pólitískur stöðugleiki} \dicExampleCS{politická stabilita}
\dicEntry[stöðuglyndur] \textls[15]{\dicTerm{stöðug··lyndur} \dicIPA{{s}{\textsubring{d}}{\oe}{\textlengthmark}{ð}{\textscy}{\textbabygamma}{l}{\textsci}{n}{\textsubring{d}}{\textscy}{\textsubring{r}}} \dicPos{adj}[2]\dicFlx{}[-14] \dicSynonym{skapfastur}} \dicDirectTranslationCS{stálý (povahově), neochvějný, vytrvalý}
\dicEntry[stöðugt] \dicTerm{stöð··ug|t} \dicIPA{{s}{\textsubring{d}}{\oe}{\textlengthmark}{ð}{\textscy}{x}{\textsubring{d}}} \dicPos{adv} \dicFlx{(comp ‑ar, sup ‑ast)} \dicSynonym{sífellt} \dicDirectTranslationCS{(neu)stále, ustavičně, pořád}
\dicEntry[stöðugur] \dicTerm{stöð··ugur} \dicsymFrequent\  \dicIPA{{s}{\textsubring{d}}{\oe}{\textlengthmark}{ð}{\textscy}{\textbabygamma}{\textscy}{\textsubring{r}}} \dicPos{adj}[1]\dicFlx{}[-8] \textbf{1.} \dicSynonym{sífelldur} \dicDirectTranslationCS{(neu)stálý, ustavičný, věčný} \dicExampleIS{stöðugar áhyggjur} \dicExampleCS{neustálé obavy}  \textbf{2.} \dicSynonym*{stabíll} \dicDirectTranslationCS{stabilní}
\dicEntry[stöðuhækkun] \dicTerm{stöðu··hækk|un} \dicIPA{{s}{\textsubring{d}}{\oe}{\textlengthmark}{ð}{\textscy}{h}{a}{i}{h}{\r{g}}{\textscy}{\textsubring{n}}} \dicPos{f}[7] \dicFlx{(‑unar, ‑anir)}[8] \dicDirectTranslationCS{povýšení, postup}
\dicEntry[stöðulína] \dicTerm{stöðu··lín|a} \dicIPA{{s}{\textsubring{d}}{\oe}{\textlengthmark}{ð}{\textscy}{l}{i}{n}{a}} \dicPos{f}[1] \dicFlx{(‑u, ‑ur)}[7] \dicFieldCat{poč.} \dicDirectTranslationCS{stavový řádek}
\dicEntry[stöðuljós] \dicTerm{stöðu··ljós} \dicIPA{{s}{\textsubring{d}}{\oe}{\textlengthmark}{ð}{\textscy}{l}{j}{ou}{s}} \dicPos{n}[2] \dicFlx{(‑s, ‑)}[5] \dicDirectTranslationCS{parkovací světla}
\dicEntry[stöðulækkun] \dicTerm{stöðu··lækk|un} \dicIPA{{s}{\textsubring{d}}{\oe}{\textlengthmark}{ð}{\textscy}{l}{a}{i}{h}{\r{g}}{\textscy}{\textsubring{n}}} \dicPos{f}[7] \dicFlx{(‑unar)}[9] \dicDirectTranslationCS{degradování, sesazení}
\dicEntry[stöðumissir] \dicTerm{stöðu··miss|ir} \dicIPA{{s}{\textsubring{d}}{\oe}{\textlengthmark}{ð}{\textscy}{m}{\textsci}{s}{\textsci}{\textsubring{r}}} \dicPos{m}[7] \dicFlx{(‑is)}[2] \dicDirectTranslationCS{ztráta (pracovního) místa} \dicExampleIS{Það getur kostað hann stöðumissi.} \dicExampleCS{To ho může stát místo.}
\dicEntry[stöðumælasekt] \dicTerm{stöðu·mæla··sekt} \dicIPA{{s}{\textsubring{d}}{\oe}{\textlengthmark}{ð}{\textscy}{m}{a}{i}{l}{a}{s}{\textepsilon}{x}{\textsubring{d}}} \dicPos{f}[7] \dicFlx{(‑ar, ‑ir)}[1] \dicDirectTranslationCS{pokuta za parkování}
\dicEntry[stöðumælir] \dicTerm{stöðu··mæl|ir} \dicIPA{{s}{\textsubring{d}}{\oe}{\textlengthmark}{ð}{\textscy}{m}{a}{i}{l}{\textsci}{\textsubring{r}}} \dicPos{m}[7] \dicFlx{(‑is, ‑ar)}[1] \dicDirectTranslationCS{parkovací hodiny}
\dicEntry[stöðuorka] \dicTerm{stöðu··ork|a} \dicIPA{{s}{\textsubring{d}}{\oe}{\textlengthmark}{ð}{\textscy}{\textopeno}{\textsubring{r}}{\r{g}}{a}} \dicPos{f}[1] \dicFlx{(‑u)}[5] \dicFieldCat{fyz.} \dicDirectTranslationCS{potenciální\,/\addthin polohová energie} \dicAntonym{hreyfiorka}
\dicEntry[stöðurafmagn] \dicTerm{stöðu··raf·magn} \dicIPA{{s}{\textsubring{d}}{\oe}{\textlengthmark}{ð}{\textscy}{r}{a}{m}{a}{\r{g}}{\textsubring{n}}} \dicPos{n}[2] \dicFlx{(‑s)}[2] \dicFieldCat{elek.} \dicDirectTranslationCS{statická elektřina}
\dicEntry[stöðutákn] \dicTerm{stöðu··tákn} \dicIPA{{s}{\textsubring{d}}{\oe}{\textlengthmark}{ð}{\textscy}{t\smash{\textsuperscript{h}}}{au}{h}{\r{g}}{\textsubring{n}}} \dicPos{n}[2] \dicFlx{(‑s, ‑)}[5] \dicDirectTranslationCS{symbol (společenského) postavení}
\dicEntry[stöðuval] \dicTerm{stöðu··val} \dicIPA{{s}{\textsubring{d}}{\oe}{\textlengthmark}{ð}{\textscy}{v}{a}{\textsubring{l}}} \dicPos{n}[2] \dicFlx{(‑s)}[2] \dicDirectTranslationCS{výběr povolání}
\dicEntry[stöðuvatn] \dicTerm{stöðu··|vatn} \dicIPA{{s}{\textsubring{d}}{\oe}{\textlengthmark}{ð}{\textscy}{v}{a}{h}{\textsubring{d}}{\textsubring{n}}} \dicPos{n}[2] \dicFlx{(‑vatns, ‑vötn)}[8] \dicDirectTranslationCS{jezero}
\dicEntry[stöðuveiting] \dicTerm{stöðu··veit·ing} \dicIPA{{s}{\textsubring{d}}{\oe}{\textlengthmark}{ð}{\textscy}{v}{ei}{\textsubring{d}}{i}{\ng}{\r{g}}} \dicPos{f}[4] \dicFlx{(‑ar, ‑ar)}[5] \dicDirectTranslationCS{dosazení do funkce, udělení postu}
\dicEntry[stöðva] \dicTerm{stöðv|a} \dicsymFrequent\  \dicIPA{{s}{\textsubring{d}}{\oe}{ð}{v}{a}} \dicPos{v}[1] \dicFlx{(‑aði)}[1] \dicFlx{acc} \dicSynonym{stansa} \dicDirectTranslationCS{(po)zastavit, (po)zastavovat} \dicExampleIS{stöðva vélina} \dicExampleCS{pozastavit stroj};  \dicIdiom{stöðvast}{ \dicPhraseIS{stöðvast}} \dicFlx{refl} \dicSynonym*{nema staðar} \dicDirectTranslationCS{zastavit se, zastavovat se}
\dicEntry[stöðvarstjóri] \dicTerm{stöðvar··stjór|i} \dicIPA{{s}{\textsubring{d}}{\oe}{ð}{v}{a}{\textsubring{r}}{s}{\textsubring{d}}{j}{ou}{r}{\textsci}} \dicPos{m}[1] \dicFlx{(‑a, ‑ar)}[1] \dicSynonym{yfirmaður} \dicDirectTranslationCS{řídící, mistr(ová), přednosta, přednostka} \dicExampleIS{símstöðvarstjóri} \dicExampleCS{přednosta telefonní ústředny}
\dicEntry[stöðvun] \dicTerm{stöðv|un} \dicIPA{{s}{\textsubring{d}}{\oe}{ð}{v}{\textscy}{\textsubring{n}}} \dicPos{f}[7] \dicFlx{(‑unar, ‑anir)}[8] \dicDirectTranslationCS{(po)zastavení} \dicExampleIS{stöðvun rekstrar} \dicExampleCS{pozastavení podnikání}
\dicEntry[stöðvunarmerki] \dicTerm{stöðvunar··merki} \dicIPA{{s}{\textsubring{d}}{\oe}{ð}{v}{\textscy}{n}{a}{r}{m}{\textepsilon}{\textsubring{r}}{\r{\textObardotlessj}}{\textsci}} \dicPos{n}[2] \dicFlx{(‑s, ‑)}[16] \dicDirectTranslationCS{stopka} \dicIndirectTranslationCS{(dopravní značka přikazující zastavení)}
\dicEntry[stöðvunarvegalengd] \dicTerm{stöðvunar··vega·lengd} \dicIPA{{s}\-{\textsubring{d}}\-{\oe}\-{ð}\-{v}\-{\textscy}\-{n}\-{a}\-{r}\-{v}\-{\textepsilon}\-{\textbabygamma}\-{a}\-{l}\-{\textepsilon}\-{\ng}\-{\textsubring{d}}\-} \dicPos{f}[7] \dicFlx{(‑ar, ‑ir)}[1] \dicDirectTranslationCS{brzdná dráha}
\dicEntry[stög] \dicTerm{stög} \dicIPA{{s}{\textsubring{d}}{\oe}{\textlengthmark}{x}} \dicPos{n} \dicFlx{pl nom} \dicLink{stag}
\dicEntry[stökk] \dicTerm{stökk\smash{\textsuperscript{1}}} \dicsymFrequent\  \dicIPA{{s}{\textsubring{d}}{\oe}{h}{\r{g}}} \dicPos{n}[2] \dicFlx{(‑s, ‑)}[5] \textbf{1.} \dicSynonym*{það að stökkva} \dicDirectTranslationCS{(pře)skok, skočení} \dicExampleIS{stutt stökk} \dicExampleCS{krátký skok}  \textbf{2.} \dicSynonym*{gangtegund hesta} \dicDirectTranslationCS{cval} \dicIndirectTranslationCS{(nejrychlejší chod koně v~třídobém taktu)}
\dicEntry[stökk] \dicTerm{stökk\smash{\textsuperscript{2}}} \dicIPA{{s}{\textsubring{d}}{\oe}{h}{\r{g}}} \dicPos{v} \dicFlx{ind pf sg 1 pers} \dicLink{stökkva\smash{\textsuperscript{2}}}
\dicEntry[stökkbretti] \dicTerm{stökk··bretti} \dicIPA{{s}{\textsubring{d}}{\oe}{h}{\r{g}}{\textsubring{b}}{r}{\textepsilon}{h}{\textsubring{d}}{\textsci}} \dicPos{n}[2] \dicFlx{(‑s, ‑)}[14] \dicFieldCat{sport.} \dicDirectTranslationCS{odrazový můstek};  \dicPhraseIS{nota e‑ð eins og stökkbretti} \dicLangCat{přen.} \dicDirectTranslationCS{použít (co) jako odrazový můstek (k~dalšímu rozvoji ap.)}
\dicEntry[stökkbreyting] \dicTerm{stökk··breyt·ing} \dicIPA{{s}{\textsubring{d}}{\oe}{h}{\r{g}}{\textsubring{b}}{r}{ei}{\textsubring{d}}{i}{\ng}{\r{g}}} \dicPos{f}[4] \dicFlx{(‑ar, ‑ar)}[5] \textbf{1.} \dicFieldCat{biol.} \dicDirectTranslationCS{mutace}  \textbf{2.} \dicSynonym*{mikil breyting} \dicDirectTranslationCS{skoková změna}
\dicEntry[stökkpallur] \dicTerm{stökk··pall|ur} \dicIPA{{s}{\textsubring{d}}{\oe}{h}{\r{g}}{p\smash{\textsuperscript{h}}}{a}{\textsubring{d}}{l}{\textscy}{\textsubring{r}}} \dicPos{m}[6] \dicFlx{(‑s, ‑ar)}[5] \dicFieldCat{sport.} \dicDirectTranslationCS{skokanský můstek}
\dicEntry[stökkur] \dicTerm{stökkur} \dicIPA{{s}{\textsubring{d}}{\oe}{h}{\r{g}}{\textscy}{\textsubring{r}}} \dicPos{adj}[1]\dicFlx{}[-1] \dicDirectTranslationCS{křehký, křupavý} \dicExampleIS{stökk rödd} \dicExampleCS{křehký hlas}
\dicEntry[stökkva] \dicTerm{stökk|va\smash{\textsuperscript{1}}} \dicsymFrequent\  \dicIPA{{s}{\textsubring{d}}{\oe}{h}{\r{g}}{v}{a}} \dicPos{v}[2] \dicFlx{(‑ti, ‑t)}[35] \dicFlx{dat} \textbf{1.} \dicDirectTranslationCS{zahnat, zahánět, odehnat, odhánět, zaplašit} \dicExampleIS{stökkva e‑m burt} \dicExampleCS{odehnat (koho) pryč}  \textbf{2.} \dicDirectTranslationCS{(po)kropit} \dicExampleIS{stökkva vatni á e‑ð} \dicExampleCS{pokropit (co) vodou}
\dicEntry[stökkva] \dicTerm{stökkva\smash{\textsuperscript{2}}} \dicsymFrequent\  \dicIPA{{s}{\textsubring{d}}{\oe}{h}{\r{g}}{v}{a}} \dicPos{v}[6] \dicFlx{(stekk, stökk, stukkum, stykki, stokkið)}[91] \textbf{1.} \dicDirectTranslationCS{skočit, skákat} \dicExampleIS{stökkva tvo metra} \dicExampleCS{skočit dva metry}  \textbf{2.} \dicSynonym{skreppa} \dicDirectTranslationCS{zaskočit, zajít, zaběhnout} \dicExampleIS{stökkva út í búð} \dicExampleCS{zaskočit do obchodu}
\dicEntry[stökkvari] \dicTerm{stökkv··ar|i} \dicIPA{{s}{\textsubring{d}}{\oe}{h}{\r{g}}{v}{a}{r}{\textsci}} \dicPos{m}[1] \dicFlx{(‑a, ‑ar)}[13] \dicDirectTranslationCS{skokan(ka)} \dicExampleIS{hástökkvari} \dicExampleCS{skokan do výšky}
\dicEntry[stöku] \dicTerm{stöku\smash{\textsuperscript{1}}} \dicIPA{{s}{\textsubring{d}}{\oe}{\textlengthmark}{\r{g}}{\textscy}} \dicPos{f} \dicFlx{sg gen} \dicLink{staka}
\dicEntry[stöku] \dicTerm{stöku\smash{\textsuperscript{2}}} \dicsymFrequent\  \dicIPA{{s}{\textsubring{d}}{\oe}{\textlengthmark}{\r{g}}{\textscy}} \dicPos{adj}[13] \dicFlx{indecl}[3] \dicDirectTranslationCS{ojedinělý, osamělý} \dicExampleIS{stöku bændur} \dicExampleCS{ojedinělí farmáři};  \dicPhraseIS{á stöku stað} \dicFlx{adv} \dicDirectTranslationCS{na některých místech, ojediněle} \dicIndirectTranslationCS{(o~pozici)};  \dicPhraseIS{stöku sinnum} \dicFlx{adv} \dicDirectTranslationCS{ojediněle, sporadicky} \dicIndirectTranslationCS{(o~čase)}
\dicEntry[stökur] \dicTerm{stökur} \dicIPA{{s}{\textsubring{d}}{\oe}{\textlengthmark}{\r{g}}{\textscy}{\textsubring{r}}} \dicPos{f} \dicFlx{pl nom} \dicLink{staka}
\dicEntry[stöllu] \dicTerm{stöllu} \dicIPA{{s}{\textsubring{d}}{\oe}{\textsubring{d}}{l}{\textscy}} \dicPos{f} \dicFlx{sg gen} \dicLink{stalla}
\dicEntry[stönd] \dicTerm{stönd} \dicIPA{{s}{\textsubring{d}}{\oe}{n}{\textsubring{d}}} \dicPos{n} \dicFlx{pl nom} \dicLink{stand}
\dicEntry[stöndugur] \dicTerm{stönd··ugur} \dicIPA{{s}{\textsubring{d}}{\oe}{n}{\textsubring{d}}{\textscy}{\textbabygamma}{\textscy}{\textsubring{r}}} \dicPos{adj}[1]\dicFlx{}[-8] \dicSynonym{efnaður} \dicDirectTranslationCS{zámožný, majetný}
\dicEntry[stöng] \dicTerm{stöng} \dicsymFrequent\  \dicIPA{{s}{\textsubring{d}}{\oe i}{\ng}{\r{g}}} \dicPos{f}[8] \dicFlx{(stangar, stengur\,/\addthin stangir)}[9] \textbf{1.} \dicDirectTranslationCS{prut, tyč, tyčka}  \textbf{2.} \dicDirectTranslationCS{stožár, žerď} \dicExampleIS{flaggstöng} \dicExampleCS{vlajkový stožár};  \dicPhraseIS{í hálfa stöng} \dicFlx{adv} \dicDirectTranslationCS{na půl žerdi}  \textbf{3.} \dicSynonym{veiðistöng} \dicDirectTranslationCS{(rybářský) prut}
\dicEntry[stöngl] \dicTerm{stöngl} \dicIPA{{}{s}{\textsubring{d}}{\oe i}{\ng}{\textsubring{l}}} \dicPos{n} \dicFlx{pl nom} \dicLink{stangl}
\dicEntry[stöngull] \dicTerm{stöng|ull} \dicIPA{{s}{\textsubring{d}}{\oe i}{\ng}{\r{g}}{\textscy}{\textsubring{d}}{\textsubring{l}}} \dicPos{m}[6] \dicFlx{(‑uls, ‑lar)}[31] \dicFieldCat{bot.} \dicDirectTranslationCS{stonek, lodyha, stvol} \dicExampleIS{stöngullinn á plöntunni} \dicExampleCS{stonek rostliny}
\dicEntry[stöplarit] \dicTerm{stöpla··rit} \dicIPA{{s}{\textsubring{d}}{\oe}{h}{\textsubring{b}}{l}{a}{r}{\textsci}{\textsubring{d}}} \dicPos{n}[2] \dicFlx{(‑s, ‑)}[5] \dicDirectTranslationCS{histogram, sloupkový diagram}
\dicEntry[stöppu] \dicTerm{stöppu} \dicIPA{{s}{\textsubring{d}}{\oe}{h}{\textsubring{b}}{\textscy}} \dicPos{f} \dicFlx{sg gen} \dicLink{stappa\smash{\textsuperscript{1}}}
\dicEntry[stöpull] \dicTerm{stöp|ull} \dicIPA{{s}{\textsubring{d}}{\oe}{\textlengthmark}{\textsubring{b}}{\textscy}{\textsubring{d}}{\textsubring{l}}} \dicPos{m}[6] \dicFlx{(‑uls, ‑lar)}[31] \dicSynonym{súla} \dicDirectTranslationCS{sloup, pilíř} \dicExampleIS{brúarstöpull} \dicExampleCS{mostní pilíř}
\dicEntry[stör] \dicTerm{stör} \dicIPA{{s}{\textsubring{d}}{\oe}{\textlengthmark}{\textsubring{r}}} \dicPos{f}[7] \dicFlx{(starar, starir)}[16] \dicFieldCat{bot.} \dicDirectTranslationCS{ostřice} \textit{(l.~{\textLA{Carex}})}  \dicsymPhoto\ 
\dicFigure{115960.jpg}{Stör}{Stör - Kubíková Petra, Biolib, Copyright/PD}
\dicEntry[störf] \dicTerm{störf} \dicIPA{{s}{\textsubring{d}}{\oe}{r}{f}} \dicPos{n} \dicFlx{pl nom} \dicLink{starf}
\dicEntry[subba] \dicTerm{subb|a\smash{\textsuperscript{1}}} \dicIPA{{s}{\textscy}{\textsubring{b}}{\textlengthmark}{a}} \dicPos{f}[1] \dicFlx{(‑u, ‑ur)}[7] \dicDirectTranslationCS{(zanedbaná) ženská, cuchta}
\dicEntry[subba] \dicTerm{subb|a\smash{\textsuperscript{2}}} \dicIPA{{s}{\textscy}{\textsubring{b}}{\textlengthmark}{a}} \dicPos{v}[1] \dicFlx{(‑aði)}[1] \dicPhraseIS{subba e‑ð út} \dicDirectTranslationCS{zaneřádit (co), zasvinit (co)}
\dicEntry[subbulegur] \dicTerm{subbu··legur} \dicIPA{{s}{\textscy}{\textsubring{b}}{\textlengthmark}{\textscy}{l}{\textepsilon}{\textbabygamma}{\textscy}{\textsubring{r}}} \dicPos{adj}[1]\dicFlx{}[-8] \dicSynonym{sóðalegur} \dicDirectTranslationCS{zašpině\-ný, ušmudlaný, zanedbaný}
\dicEntry[sudda] \dicTerm{sudd|a} \dicIPA{{s}{\textscy}{\textsubring{d}}{\textlengthmark}{a}} \dicPos{v}[1] \dicFlx{(‑aði)}[14] \dicFlx{impers} \dicPhraseIS{það suddar} \dicDirectTranslationCS{mží, mrholí, poprchává}
\dicEntry[suddi] \dicTerm{sudd|i} \dicIPA{{s}{\textscy}{\textsubring{d}}{\textlengthmark}{\textsci}} \dicPos{m}[1] \dicFlx{(‑a)}[3] \dicSynonym{súld} \dicDirectTranslationCS{mžení, mrholení, poprchávání}
\dicEntry[suð] \dicTerm{suð} \dicIPA{{s}{\textscy}{\textlengthmark}{\texttheta}} \dicPos{n}[2] \dicFlx{(‑s)}[2] \textbf{1.} \dicSynonym*{daufur kliður} \dicDirectTranslationCS{bzučení, bzikání} \dicExampleIS{suð í býflugum} \dicExampleCS{bzučení včel}  \textbf{2.} \dicSynonym{þrábeiðni} \dicDirectTranslationCS{žadonění, škemrání, prošení} \dicExampleIS{Vertu ekki að þessu suði!} \dicExampleCS{Nech toho škemrání!}
\dicEntry[suða] \dicTerm{suð|a\smash{\textsuperscript{1}}} \dicIPA{{s}{\textscy}{\textlengthmark}{ð}{a}} \dicPos{f}[1] \dicFlx{(‑u, ‑ur)}[19] \textbf{1.} \dicSynonym*{það að sjóða} \dicDirectTranslationCS{(po)vaření}  \textbf{2.} \dicSynonym*{málmsuða} \dicDirectTranslationCS{svařování, sváření}
\dicEntry[suða] \dicTerm{suð|a\smash{\textsuperscript{2}}} \dicIPA{{s}{\textscy}{\textlengthmark}{ð}{a}} \dicPos{v}[1] \dicFlx{(‑aði)}[1] \textbf{1.} \dicSynonym{niða} \dicDirectTranslationCS{(za)bzučet, (za)bzikat, příst} \dicExampleIS{Rokkurinn suðaði.} \dicExampleCS{Kolovrátek tiše předl.}  \textbf{2.} \dicSynonym{þrábiðja} \dicDirectTranslationCS{žadonit, škemrat, prosit} \dicExampleIS{suða um e‑ð} \dicExampleCS{žadonit o~(co)}
\dicEntry[suðaustan] \dicTerm{suð··austan} \dicIPA{{s}{\textscy}{\textlengthmark}{ð}{\oe i}{s}{\textsubring{d}}{a}{\textsubring{n}}} \dicPos{adv} \dicDirectTranslationCS{z~jihovýchodu}
\dicEntry[suðaustur] \dicTerm{suð··austur\smash{\textsuperscript{1}}} \dicIPA{{s}{\textscy}{\textlengthmark}{ð}{\oe i}{s}{\textsubring{d}}{\textscy}{\textsubring{r}}} \dicPos{n}[2] \dicFlx{(‑s)}[28] \dicDirectTranslationCS{jihovýchod}
\dicEntry[suðaustur] \dicTerm{suð··aust|ur\smash{\textsuperscript{2}}} \dicIPA{{s}{\textscy}{\textlengthmark}{ð}{\oe i}{s}{\textsubring{d}}{\textscy}{\textsubring{r}}} \dicPos{adv} \dicFlx{(comp ‑ar, sup ‑ast)} \dicDirectTranslationCS{jihovýchodně, na jihovýchod}
\dicEntry[suðlægur] \dicTerm{suð··lægur} \dicIPA{{s}{\textscy}{ð}{l}{a}{i}{\textbabygamma}{\textscy}{\textsubring{r}}} \dicPos{adj}[1]\dicFlx{}[-1] \textbf{1.} \dicDirectTranslationCS{jižní (vítr ap.)}  \textbf{2.} \dicDirectTranslationCS{jižní, jižně položený (místo ap.)} \dicExampleIS{suðlæg lönd} \dicExampleCS{jižní země}
\dicEntry[suðrænn] \dicTerm{suð··rænn} \dicIPA{{s}{\textscy}{ð}{r}{a}{i}{\textsubring{d}}{\textsubring{n}}} \dicPos{adj}[7]\dicFlx{}[-1] \dicDirectTranslationCS{jižní, (jsoucí) z~jižních zemí, (jsoucí) z~teplých krajin} \dicExampleIS{suðræn stemning} \dicExampleCS{atmosféra teplých krajin}
\dicEntry[suðum] \dicTerm{suðum} \dicIPA{{s}{\textscy}{\textlengthmark}{ð}{\textscy}{\textsubring{m}}} \dicPos{v} \dicFlx{ind pf pl 1 pers} \dicLink{sjóða}
\dicEntry[suðumark] \dicTerm{suðu··|mark} \dicIPA{{s}{\textscy}{\textlengthmark}{ð}{\textscy}{m}{a}{\textsubring{r}}{\r{g}}} \dicPos{n}[2] \dicFlx{(‑marks, ‑mörk)}[8] \dicDirectTranslationCS{teplota varu}
\dicEntry[suðupottur] \dicTerm{suðu··pott|ur} \dicIPA{{s}{\textscy}{\textlengthmark}{ð}{\textscy}{p\smash{\textsuperscript{h}}}{\textopeno}{h}{\textsubring{d}}{\textscy}{\textsubring{r}}} \dicPos{m}[6] \dicFlx{(‑s, ‑ar)}[4] \dicDirectTranslationCS{kotel, kotlík}
\dicEntry[suður] \dicTerm{suður\smash{\textsuperscript{1}}} \dicsymFrequent\  \dicIPA{{s}{\textscy}{\textlengthmark}{ð}{\textscy}{\textsubring{r}}} \dicPos{n}[2] \dicFlx{(‑s)}[28] \dicDirectTranslationCS{jih};  \dicPhraseIS{af suðri} \dicFlx{adv} \dicDirectTranslationCS{z~jihu};  \dicPhraseIS{í suðri} \dicFlx{adv} \dicDirectTranslationCS{na jihu (nacházet se ap.)};  \dicPhraseIS{í suður} \dicFlx{adv} \dicDirectTranslationCS{na jih (jet ap.)};  \dicPhraseIS{úr suðri} \dicFlx{adv} \dicDirectTranslationCS{z~jihu}
\dicEntry[suður] \dicTerm{suður\smash{\textsuperscript{2}}} \dicsymFrequent\  \dicIPA{{s}{\textscy}{\textlengthmark}{ð}{\textscy}{\textsubring{r}}} \dicPos{adv} \dicFlx{(comp sunnar, sup syðst)} \textbf{1.} \dicSynonym*{í suðurátt} \dicDirectTranslationCS{jižně, na jih} \dicExampleIS{Hann er farinn suður með spænsku vinkonu sinni.} \dicExampleCS{Odešel na jih se svojí španělskou přítelkyní.}  \textbf{2.} \dicSynonym*{um Suðurland} \dicDirectTranslationCS{na jih, do jižní části země}  \textbf{3.} \dicSynonym*{til Reykjavíkur} \dicDirectTranslationCS{do Reykjavíku, do oblasti hlavního města Islandu};  \dicIdiom{suður}[af]{ \dicPhraseIS{suður af e‑u}} \dicFlx{prep} \dicDirectTranslationCS{na jižní straně (čeho), jižně od (čeho)}; { \dicPhraseIS{suður af}} \dicFlx{adv} \dicDirectTranslationCS{jižně, (směrem) na jih};  \dicIdiom{suður}[eftir]{ \dicPhraseIS{suður eftir}} \dicFlx{adv} \dicDirectTranslationCS{(směrem) k~jihu, (směrem) na jih};  \dicIdiom{suður}[frá]{ \dicPhraseIS{suður frá}} \dicFlx{adv} \dicDirectTranslationCS{na jihu, na jih};  \dicIdiom{suður}[fyrir]{ \dicPhraseIS{suður fyrir e‑ð}} \dicFlx{prep} \dicDirectTranslationCS{na jih od (čeho)};  \dicIdiom{suður}[með]{ \dicPhraseIS{suður með e‑u}} \dicFlx{prep} \dicDirectTranslationCS{na jih podél (čeho)};  \dicIdiom{suður}[um]{ \dicPhraseIS{suður um e‑ð}} \dicFlx{prep} \dicDirectTranslationCS{na jih (čeho)}; { \dicPhraseIS{suður um}} \dicFlx{adv} \dicDirectTranslationCS{na jih, jižně};  \dicIdiom{suður}[undir]{ \dicPhraseIS{suður undir e‑ð}} \dicFlx{prep} \dicDirectTranslationCS{jižně pod (čím)};  \dicIdiom{suður}[úr]{ \dicPhraseIS{suður úr e‑u}} \dicFlx{prep} \dicDirectTranslationCS{jižně z~(čeho)}
\dicEntry[Suður-Afríka] \dicTerm{Suður-Afrík|a} \dicIPA{{s}{\textscy}{\textlengthmark}{ð}{\textscy}{r}{a}{f}{r}{i}{\r{g}}{a}} \dicPos{f}[1] \dicFlx{(‑u)}[6] \dicFieldCat{geog.} \dicDirectTranslationCS{Jihoafrická republika}
\dicEntry[Suður-Afríkumaður] \dicTerm{Suður-Afríku··|maður} \dicIPA{{s}{\textscy}{\textlengthmark}{ð}{\textscy}{r}{a}{f}{r}{i}{\r{g}}{\textscy}{m}{a}{ð}{\textscy}{\textsubring{r}}} \dicPos{m}[13] \dicFlx{(‑manns, ‑menn)}[2] \dicDirectTranslationCS{Jihoafričan(ka)}
\dicEntry[Suður-Kórea] \dicTerm{Suður-Kóre|a} \dicIPA{{s}{\textscy}{\textlengthmark}{ð}{\textscy}{\textsubring{r}}{k\smash{\textsuperscript{h}}}{ou}{r}{\textepsilon}{a}} \dicPos{f}[1] \dicFlx{(‑u)}[6] \dicFieldCat{geog.} \dicDirectTranslationCS{Jižní Korea}
\begin{xtolerant}{}{1pt}
\dicEntry[Suður-Kóreumaður] \dicTerm{Suður-Kóreu··|maður}\addthinM\textls[-5]{\dicIPA{{s}{\textscy}{\textlengthmark}{ð}{\textscy}{\textsubring{r}}{k\smash{\textsuperscript{h}}}{ou}{r}{\textepsilon}{\textscy}{m}{a}{ð}{\textscy}{\textsubring{r}}}} \dicPos{m}[13]\addthinS\dicFlx{(‑manns, ‑menn)}[2] \dicDirectTranslationCS{Jihokorejec, Jihokorejka}
\end{xtolerant}
\dicEntry[Suður-Súdan] \dicTerm{Suður-Súdan} \dicIPA{{s}{\textscy}{\textlengthmark}{ð}{\textscy}{\textsubring{r}}{s}{u}{\textsubring{d}}{a}{\textsubring{n}}} \dicPos{n}[2] \dicFlx{(‑s\,/\addthin ‑)}[35] \dicFieldCat{geog.} \dicDirectTranslationCS{Jižní Súdán}
\dicEntry[suðurafrískur] \dicTerm{suður··afrískur} \dicIPA{{s}{\textscy}{\textlengthmark}{ð}{\textscy}{r}{a}{f}{r}{i}{s}{\r{g}}{\textscy}{\textsubring{r}}} \dicPos{adj}[1]\dicFlx{}[-6] \dicDirectTranslationCS{jihoafrický}
\begin{xtolerant}{}{1pt}
\dicEntry[suðuramerískur] \dicTerm{suður··amerískur}\addthinS\dicIPA{{s}{\textscy}{\textlengthmark}{ð}{\textscy}{r}{a}{m}{\textepsilon}{r}{i}{s}{\r{g}}{\textscy}{\textsubring{r}}}\addthinS\dicPos{adj}[1]\dicFlx{}[-6]\addthinS\dicDirectTranslationCS{jihoamerický}
\end{xtolerant}
\dicEntry[suðurátt] \dicTerm{suður··átt} \dicIPA{{s}{\textscy}{\textlengthmark}{ð}{\textscy}{r}{au}{h}{\textsubring{d}}} \dicPos{f}[7] \dicFlx{(‑ar)}[3] \dicDirectTranslationCS{jižní směr} \dicExampleIS{koma úr suðurátt} \dicExampleCS{přijít z~jihu}
\dicEntry[Suðureyjar] \dicTerm{Suður··eyjar} \dicIPA{{s}{\textscy}{\textlengthmark}{ð}{\textscy}{r}{ei}{j}{a}{\textsubring{r}}} \dicPos{f}[4] \dicFlx{pl}[15] \dicFieldCat{geog.} \dicDirectTranslationCS{Hebridy} \dicIndirectTranslationCS{(skotské souostroví)}
\dicEntry[suðurheimskaut] \dicTerm{suður··heim·skaut} \dicIPA{{s}{\textscy}{\textlengthmark}{ð}{\textscy}{\textsubring{r}}{h}{ei}{m}{s}{\r{g}}{\oe i}{\textsubring{d}}} \dicPos{n}[2] \dicFlx{(‑s)}[2] \dicFieldCat{geog.} \dicDirectTranslationCS{jižní pól}
\dicEntry[suðurheimskautsbaugur] \dicTerm{suður·heim·skauts··baug|ur} \dicIPA{{s}\-{\textscy}\-{\textlengthmark}\-{ð}\-{\textscy}\-{\textsubring{r}}\-{h}\-{ei}\-{m}\-{s}\-{\r{g}}\-{\oe i}\-{\textsubring{d}}\-{s}\-{\textsubring{b}}\-{\oe i}\-{\textbabygamma}\-{\textscy}\-{\textsubring{r}}\-} \dicPos{m}[6] \dicFlx{(‑s, ‑ar)}[24] \dicFieldCat{geog.} \dicDirectTranslationCS{jižní polární kruh}
\dicEntry[suðurhvel] \dicTerm{suður··hvel} \dicIPA{{s}{\textscy}{\textlengthmark}{ð}{\textscy}{\textsubring{r}}{k\smash{\textsuperscript{h}}}{v}{\textepsilon}{\textsubring{l}}} \dicPos{n}[2] \dicFlx{(‑s, ‑)}[5] \dicFieldCat{geog.} \dicDirectTranslationCS{jižní polokoule}
\dicEntry[suðurkóreskur] \dicTerm{suður··kóreskur} \dicIPA{{s}{\textscy}{\textlengthmark}{ð}{\textscy}{\textsubring{r}}{k\smash{\textsuperscript{h}}}{ou}{r}{\textepsilon}{s}{\r{g}}{\textscy}{\textsubring{r}}} \dicPos{adj}[1]\dicFlx{}[-6] \dicDirectTranslationCS{jihokorejský}
\dicEntry[Suðurland] \dicTerm{Suður··land} \dicIPA{{s}{\textscy}{\textlengthmark}{ð}{\textscy}{r}{l}{a}{n}{\textsubring{d}}} \dicPos{n}[2] \dicFlx{(‑s)}[2] \dicFieldCat{geog.} \dicDirectTranslationCS{Jižní Island}
\dicEntry[suðurpóll] \dicTerm{suður··pól|l} \dicIPA{{s}{\textscy}{\textlengthmark}{ð}{\textscy}{\textsubring{r}}{p\smash{\textsuperscript{h}}}{ou}{\textsubring{d}}{\textsubring{l}}} \dicPos{m}[6] \dicFlx{(‑s, ‑ar)}[48] \dicFieldCat{geog.} \dicSynonym{suðurheimskaut} \dicDirectTranslationCS{jižní pól}
\dicEntry[suðurskaut] \dicTerm{suður··skaut} \dicIPA{{s}{\textscy}{\textlengthmark}{ð}{\textscy}{\textsubring{r}}{s}{\r{g}}{\oe i}{\textsubring{d}}} \dicPos{n}[2] \dicFlx{(‑s, ‑)}[5] \dicFieldCat{geog.} \dicDirectTranslationCS{jižní pól}
\dicEntry[Suðurskautslandið] \dicTerm{Suður·skauts··land|ið} \dicIPA{{s}\-{\textscy}\-{\textlengthmark}\-{ð}\-{\textscy}\-{\textsubring{r}}\-{s}\-{\r{g}}\-{\oe i}\-{\textsubring{d}}\-{s}\-{l}\-{a}\-{n}\-{\textsubring{d}}\-{\textsci}\-{\texttheta}\-} \dicPos{n}[2] \dicFlx{(‑sins)}[3] \dicFlx{def} \dicFieldCat{geog.} \dicDirectTranslationCS{Antarktida}
\begin{xtolerant}{}{1pt}
\dicEntry[suðurströnd] \dicTerm{suður··|strönd}\addthinS\dicIPA{{s}{\textscy}{\textlengthmark}{ð}{\textscy}{\textsubring{r}}{s}{\textsubring{d}}{r}{\oe}{n}{\textsubring{d}}}\addthinS\dicPos{f}[8]\addthinS\dicFlx{(‑strandar, ‑strendur\,/\addthin ‑strand\-ir)}[9] \dicDirectTranslationCS{jižní pobřeží}
\end{xtolerant}
\dicEntry[suðutæki] \dicTerm{suðu··tæki} \dicIPA{{s}{\textscy}{\textlengthmark}{ð}{\textscy}{t\smash{\textsuperscript{h}}}{a}{i}{\r{\textObardotlessj}}{\textsci}} \dicPos{n}[2] \dicFlx{(‑s, ‑)}[16] \dicDirectTranslationCS{vařič, sporák}
\dicEntry[suðvestan] \dicTerm{suð··vestan} \dicIPA{{s}{\textscy}{ð}{v}{\textepsilon}{s}{\textsubring{d}}{a}{\textsubring{n}}} \dicPos{adv} \dicDirectTranslationCS{z~jihozápadu}
\dicEntry[suðvestur] \dicTerm{suð··vestur\smash{\textsuperscript{1}}} \dicIPA{{s}{\textscy}{ð}{v}{\textepsilon}{s}{\textsubring{d}}{\textscy}{\textsubring{r}}} \dicPos{n}[2] \dicFlx{(‑s)}[28] \dicDirectTranslationCS{jihozápad}
\dicEntry[suðvestur] \dicTerm{suð··vest|ur\smash{\textsuperscript{2}}} \dicIPA{{s}{\textscy}{ð}{v}{\textepsilon}{s}{\textsubring{d}}{\textscy}{\textsubring{r}}} \dicPos{adv} \dicFlx{(comp ‑ar, sup ‑ast)} \dicDirectTranslationCS{jihozápadně, na jihozápad}
\dicEntry[sugum] \dicTerm{sugum} \dicIPA{{s}{\textscy}{\textlengthmark}{\textbabygamma}{\textscy}{\textsubring{m}}} \dicPos{v} \dicFlx{ind pf pl 1 pers} \dicLink{sjúga}
\dicEntry[sukk] \dicTerm{sukk} \dicIPA{{s}{\textscy}{h}{\r{g}}} \dicPos{n}[2] \dicFlx{(‑s)}[2] \textbf{1.} \dicSynonym{svall\smash{\textsuperscript{1}}} \dicDirectTranslationCS{hýření, flámování}  \textbf{2.} \dicSynonym{eyðslusemi} \dicDirectTranslationCS{hýření, plýtvání, mrhání}
\dicEntry[sukka] \dicTerm{sukk|a} \dicIPA{{s}{\textscy}{h}{\r{g}}{a}} \dicPos{v}[1] \dicFlx{(‑aði)}[1] \textbf{1.} \dicSynonym{svalla} \dicDirectTranslationCS{hýřit, flámovat}  \textbf{2.} \dicDirectTranslationCS{hýřit, plýtvat, mrhat} \dicExampleIS{sukka með peninga} \dicExampleCS{plýtvat penězi}
\dicEntry[sukksamur] \dicTerm{sukk··|samur} \dicIPA{{s}{\textscy}{h}{\r{g}}{s}{a}{m}{\textscy}{\textsubring{r}}} \dicPos{adj}[1] \dicFlx{(f ‑söm)}[2] \textbf{1.} \dicSynonym*{óreglusamur} \dicDirectTranslationCS{hýřivý, zhýralý, prostopášný}  \textbf{2.} \dicSynonym{eyðslusamur} \dicDirectTranslationCS{hýřivý, plýtvající, mrhající}
\dicEntry[sukkum] \dicTerm{sukkum} \dicIPA{{s}{\textscy}{h}{\r{g}}{\textscy}{\textsubring{m}}} \dicPos{v} \dicFlx{ind pf pl 1 pers} \dicLink{sökkva\smash{\textsuperscript{2}}}
\dicEntry[sull] \dicTerm{sull} \dicIPA{{s}{\textscy}{\textsubring{d}}{\textsubring{l}}} \dicPos{n}[2] \dicFlx{(‑s)}[2] \textbf{1.} \dicSynonym{skol} \dicDirectTranslationCS{břečka, brynda} \dicIndirectTranslationCS{(slabý nápoj)} \dicExampleIS{Kaffið er mesta sull.} \dicExampleCS{Ta káva je naprostá břečka.}  \textbf{2.} \dicSynonym{bleyta\smash{\textsuperscript{1}}} \dicDirectTranslationCS{břečka, čvachtanice}
\dicEntry[sulla] \dicTerm{sull|a} \dicIPA{{s}{\textscy}{\textsubring{d}}{l}{a}} \dicPos{v}[1] \dicFlx{(‑aði)}[1] \dicFlx{dat} \textbf{1.} \dicFlx{dat} \dicSynonym{óhreinka} \dicDirectTranslationCS{(po)bryndat, (po)cintat} \dicExampleIS{sulla kaffi á gólfið} \dicExampleCS{pobryndat podlahu kafem}  \textbf{2.} \dicSynonym*{ösla í bleytu} \dicDirectTranslationCS{cákat se, šplouchat se, brouzdat se} \dicExampleIS{sulla í vatninu} \dicExampleCS{cákat se ve vodě};  \dicIdiom{sulla}[í]{ \dicPhraseIS{sulla í e‑u}} \dicLangCat{hovor.} \dicDirectTranslationCS{popíjet (co)};  \dicIdiom{sullast}{ \dicPhraseIS{sullast}} \dicFlx{refl} \dicDirectTranslationCS{(po)cákat se, rozstříknout se}
\dicEntry[sullum] \dicTerm{sullum} \dicIPA{{s}{\textscy}{\textsubring{d}}{l}{\textscy}{\textsubring{m}}} \dicPos{v} \dicFlx{ind pf pl 1 pers} \dicLink{svella\smash{\textsuperscript{2}}}
\dicEntry[sulta] \dicTerm{sult|a\smash{\textsuperscript{1}}} \dicIPA{{s}{\textscy}{\textsubring{l}}{\textsubring{d}}{a}} \dicPos{f}[1] \dicFlx{(‑u, ‑ur)}[19] \textbf{1.} \dicDirectTranslationCS{marmeláda, džem} \dicExampleIS{bláberjasulta} \dicExampleCS{borůvková marmeláda}  \textbf{2.} \dicDirectTranslationCS{sulc, huspenina}
\dicEntry[sulta] \dicTerm{sult|a\smash{\textsuperscript{2}}} \dicIPA{{s}{\textscy}{\textsubring{l}}{\textsubring{d}}{a}} \dicPos{v}[1] \dicFlx{(‑aði)}[1] \dicFlx{acc} \dicDirectTranslationCS{zavařit, konzervovat}
\dicEntry[sultarlaun] \dicTerm{sultar··laun} \dicIPA{{s}{\textscy}{\textsubring{l}}{\textsubring{d}}{a}{r}{l}{\oe i}{\textsubring{n}}} \dicPos{n}[2] \dicFlx{pl}[1] \dicDirectTranslationCS{bídný\,/\addthin mizerný plat}
\dicEntry[sultum] \dicTerm{sultum} \dicIPA{{s}{\textscy}{\textsubring{l}}{\textsubring{d}}{\textscy}{\textsubring{m}}} \dicPos{v} \dicFlx{ind pf pl 1 pers} \dicLink{svelta\smash{\textsuperscript{2}}}
\dicEntry[sultur] \dicTerm{sult|ur} \dicIPA{{s}{\textscy}{\textsubring{l}}{\textsubring{d}}{\textscy}{\textsubring{r}}} \dicPos{m}[10] \dicFlx{(‑ar\,/\addthin ‑s)}[30] \dicDirectTranslationCS{hlad, hladovění} \dicExampleIS{deyja úr sulti} \dicExampleCS{umřít hlady};  \dicPhraseIS{búa\,/\addthin lifa við sult og seyru} \dicDirectTranslationCS{žít o~hladu a~bídě}
\dicEntry[sumar] \dicTerm{sum|ar} \dicsymFrequent\  \dicIPA{{s}{\textscy}{\textlengthmark}{m}{a}{\textsubring{r}}} \dicPos{n}[3] \dicFlx{(‑ars, ‑ur)}[1] \dicDirectTranslationCS{léto};  \dicPhraseIS{á sumrin} \dicFlx{adv} \dicDirectTranslationCS{v~létě, během léta};  \dicPhraseIS{í fyrra sumar} \dicFlx{adv} \dicDirectTranslationCS{minulé léto, loni v~létě};  \dicPhraseIS{í sumar} \dicFlx{adv} \dicDirectTranslationCS{v~létě};  \dicPhraseIS{um sumarið} \dicFlx{adv} \dicDirectTranslationCS{během léta, v~létě}
\dicEntry[sumarauki] \dicTerm{sumar··auk|i} \dicIPA{{s}{\textscy}{\textlengthmark}{m}{a}{r}{\oe i}{\r{\textObardotlessj}}{\textsci}} \dicPos{m}[1] \dicFlx{(‑a)}[3] \dicFieldCat{hist.} \dicIndirectTranslationCS{vložený týden ve starém islandském kalendáři}
\dicEntry[sumaráætlun] \dicTerm{sumar··á·ætl|un} \dicIPA{{s}{\textscy}{\textlengthmark}{m}{a}{r}{au}{a}{i}{h}{\textsubring{d}}{l}{\textscy}{\textsubring{n}}} \dicPos{f}[7] \dicFlx{(‑unar, ‑anir)}[8] \dicDirectTranslationCS{letní jízdní řád}
\dicEntry[sumarbústaður] \dicTerm{sumar··bú·stað|ur} \dicIPA{{s}{\textscy}{\textlengthmark}{m}{a}{r}{\textsubring{b}}{u}{s}{\textsubring{d}}{a}{ð}{\textscy}{\textsubring{r}}} \dicPos{m}[10] \dicFlx{(‑ar, ‑ir)}[14] \dicDirectTranslationCS{letní sídlo, (letní) chata}
\dicEntry[sumardagur] \dicTerm{sumar··dag|ur} \dicIPA{{s}{\textscy}{\textlengthmark}{m}{a}{r}{\textsubring{d}}{a}{\textbabygamma}{\textscy}{\textsubring{r}}} \dicPos{m}[6] \dicFlx{(‑s, ‑ar)}[62] \dicDirectTranslationCS{letní den};  \dicPhraseIS{sumardagurinn fyrsti} \dicDirectTranslationCS{islandský státní svátek prvního letního dne} \dicIndirectTranslationCS{(vždy první čtvrtek po 18. dubnu)}
\dicEntry[sumardekk] \dicTerm{sumar··dekk} \dicIPA{{s}{\textscy}{\textlengthmark}{m}{a}{r}{\textsubring{d}}{\textepsilon}{h}{\r{g}}} \dicPos{n}[2] \dicFlx{(‑s, ‑)}[13] \dicDirectTranslationCS{letní pneumatika} \dicAntonym{vetrardekk}
\dicEntry[sumardvöl] \dicTerm{sumar··|dvöl} \dicIPA{{s}{\textscy}{\textlengthmark}{m}{a}{r}{\textsubring{d}}{v}{\oe}{\textsubring{l}}} \dicPos{f}[7] \dicFlx{(‑dvalar, ‑dvalir)}[16] \dicDirectTranslationCS{letní pobyt}
\dicEntry[sumarfatnaður] \dicTerm{sumar··fat·nað|ur} \dicIPA{{s}{\textscy}{\textlengthmark}{m}{a}{\textsubring{r}}{f}{a}{\textsubring{d}}{n}{a}{ð}{\textscy}{\textsubring{r}}} \dicPos{m}[10] \dicFlx{(‑ar)}[9] \dicDirectTranslationCS{letní oblečení}
\dicEntry[sumarfrí] \dicTerm{sumar··frí} \dicIPA{{s}{\textscy}{\textlengthmark}{m}{a}{\textsubring{r}}{f}{r}{i}} \dicPos{n}[2] \dicFlx{(‑s, ‑)}[5] \dicDirectTranslationCS{letní prázdniny}
\dicEntry[sumarhús] \dicTerm{sumar··hús} \dicIPA{{s}{\textscy}{\textlengthmark}{m}{a}{\textsubring{r}}{h}{u}{s}} \dicPos{n}[2] \dicFlx{(‑s, ‑)}[5] \dicSynonym{sumarbústaður} \dicDirectTranslationCS{chalupa, chata, letní dům}
\dicEntry[sumarlegur] \dicTerm{sumar··legur} \dicIPA{{s}{\textscy}{\textlengthmark}{m}{a}{r}{l}{\textepsilon}{\textbabygamma}{\textscy}{\textsubring{r}}} \dicPos{adj}[1]\dicFlx{}[-8] \dicDirectTranslationCS{letní, připomínající léto}
\dicEntry[sumarleyfi] \dicTerm{sumar··leyfi} \dicIPA{{s}{\textscy}{\textlengthmark}{m}{a}{r}{l}{ei}{v}{\textsci}} \dicPos{n}[2] \dicFlx{(‑s, ‑)}[14] \dicSynonym{sumarfrí} \dicDirectTranslationCS{letní dovolená\,/\addthin prázdniny}
\dicEntry[sumarskóli] \dicTerm{sumar··skól|i} \dicIPA{{s}{\textscy}{\textlengthmark}{m}{a}{\textsubring{r}}{s}{\r{g}}{ou}{l}{\textsci}} \dicPos{m}[1] \dicFlx{(‑a, ‑ar)}[1] \dicFieldCat{škol.} \dicDirectTranslationCS{letní škola}
\dicEntry[sumarsólhvörf] \dicTerm{sumar··sól·hvörf} \dicIPA{{s}{\textscy}{\textlengthmark}{m}{a}{\textsubring{r}}{s}{ou}{l}{k\smash{\textsuperscript{h}}}{v}{\oe}{r}{f}} \dicPos{n}[2] \dicFlx{pl}[9] \dicDirectTranslationCS{letní slunovrat}
\dicEntry[sumarstarf] \dicTerm{sumar··|starf} \dicIPA{{s}{\textscy}{\textlengthmark}{m}{a}{\textsubring{r}}{s}{\textsubring{d}}{a}{r}{f}} \dicPos{n}[2] \dicFlx{(‑starfs, ‑störf)}[8] \dicDirectTranslationCS{letní práce\,/\addthin brigáda}
\dicEntry[sumartími] \dicTerm{sumar··tím|i} \dicIPA{{s}{\textscy}{\textlengthmark}{m}{a}{\textsubring{r}}{t\smash{\textsuperscript{h}}}{i}{m}{\textsci}} \dicPos{m}[1] \dicFlx{(‑a)}[3] \textbf{1.} \dicSynonym{sumar} \dicDirectTranslationCS{léto}  \textbf{2.} \dicDirectTranslationCS{letní čas}
\dicEntry[sumarvinna] \dicTerm{sumar··vinn|a} \dicIPA{{s}{\textscy}{\textlengthmark}{m}{a}{r}{v}{\textsci}{n}{a}} \dicPos{f}[1] \dicFlx{(‑u)}[5] \dicDirectTranslationCS{letní práce\,/\addthin brigáda}
\dicEntry[sumbl] \dicTerm{sumbl} \dicIPA{{s}{\textscy}{m}{\textsubring{b}}{\textsubring{l}}} \dicPos{n}[2] \dicFlx{(‑s)}[2] \dicSynonym*{drykkjuveisla} \dicDirectTranslationCS{pitka, popíjení} \dicExampleIS{sitja að sumbli með e‑m} \dicExampleCS{popíjet s~(kým)}
\dicEntry[sumir] \dicTerm{sumir} \dicIPA{{s}{\textscy}{\textlengthmark}{m}{\textsci}{\textsubring{r}}} \dicPos{pron}[13] \dicFlx{indef pl nom m} \textbf{1.} \dicDirectTranslationCS{některý} \dicIndirectTranslationCS{(užívané predikativně)} \dicExampleIS{Sumir elska að syngja.} \dicExampleCS{Někteří rádi zpívají.}  \textbf{2.} \dicDirectTranslationCS{některý} \dicIndirectTranslationCS{(užívané atributivně)} \dicExampleIS{sumir stafirnir} \dicExampleCS{některá z~písmen}  \textbf{3.} \dicDirectTranslationCS{některý, určitý} \dicIndirectTranslationCS{(užívané atributivně, subs většinou indef)} \dicExampleIS{sum börn} \dicExampleCS{některé děti};  \dicPhraseIS{sumir hverjir} \dicDirectTranslationCS{někteří}
\dicEntry[summa] \dicTerm{summ|a} \dicIPA{{s}{\textscy}{m}{\textlengthmark}{a}} \dicPos{f}[1] \dicFlx{(‑u, ‑ur)}[7] \textbf{1.} \dicFieldCat{mat.} \dicSynonym*{samtala} \dicDirectTranslationCS{suma, součet}  \textbf{2.} \dicSynonym*{fjárupphæð} \dicDirectTranslationCS{suma, obnos, částka}
\dicEntry[sumstaðar] \dicTerm{sum··staðar} \dicIPA{{s}{\textscy}{m}{s}{\textsubring{d}}{a}{ð}{a}{\textsubring{r}}} \dicPos{adv} \dicDirectTranslationCS{sem tam, místy, v~určitých místech}
\dicEntry[sumt] \dicTerm{sumt} \dicIPA{{s}{\textscy}{\textsubring{m}}{\textsubring{d}}} \dicPos{pron}[13] \dicFlx{indef n} \textbf{1.} \dicDirectTranslationCS{něco} \dicIndirectTranslationCS{(užívané predikativně)}  \textbf{2.} \dicDirectTranslationCS{něco} \dicExampleIS{sumt af því} \dicExampleCS{něco z~toho}  \textbf{3.} \dicDirectTranslationCS{některý, určitý} \dicIndirectTranslationCS{(užívané atributivně, subs většinou indef)} \dicExampleIS{sumt fólk} \dicExampleCS{někteří lidé};  \dicPhraseIS{að sumu leyti} \dicFlx{adv} \dicDirectTranslationCS{do určité míry, částečně};  \dicPhraseIS{allt og sumt} \dicFlx{pron} \dicDirectTranslationCS{(úplně) všechno};  \dicPhraseIS{um sumt} \dicFlx{adv} \dicDirectTranslationCS{částečně, zčásti}
\dicEntry[sumur] \dicTerm{sumur\smash{\textsuperscript{1}}} \dicIPA{{s}{\textscy}{\textlengthmark}{m}{\textscy}{\textsubring{r}}} \dicPos{n} \dicFlx{pl nom} \dicLink{sumar}
\dicEntry[sumur] \dicTerm{sumur\smash{\textsuperscript{2}}} \dicsymFrequent\  \dicIPA{{s}{\textscy}{\textlengthmark}{m}{\textscy}{\textsubring{r}}} \dicPos{pron}[13] \dicFlx{indef} \dicDirectTranslationCS{některý, určitý} \dicExampleIS{Sumir búa einir alla tíð.} \dicExampleCS{Někteří žijí celý čas sami.} \dicAntonym{sumir}
\dicEntry[sund] \dicTerm{sund} \dicsymFrequent\  \dicIPA{{s}{\textscy}{n}{\textsubring{d}}} \dicPos{n}[2] \dicFlx{(‑s, ‑)}[5] \textbf{1.} \dicDirectTranslationCS{sund, průliv, úžina} \dicExampleIS{í sundinu milli tveggja eyja} \dicExampleCS{v~úžině mezi dvěma ostrovy}  \textbf{2.} \dicSynonym*{landræma} \dicDirectTranslationCS{úzký pás země} \dicIndirectTranslationCS{(často bažinatý, mezi dvěma kopci nebo řekami)}  \textbf{3.} \dicSynonym*{sundíþrótt} \dicDirectTranslationCS{plavání} \dicExampleIS{heimsmeistarakeppni í sundi} \dicExampleCS{mistrovství světa v~plavání};  \dicIdiom{sund}{ \dicPhraseIS{það eru öll sund lokuð}} \dicFlx{impers} \dicLangCat{přen.} \dicDirectTranslationCS{není cesty ven, všechny možnosti jsou vyčerpány}
\dicEntry[sundbolur] \dicTerm{sund··bol|ur} \dicIPA{{s}{\textscy}{n}{\textsubring{d}}{\textsubring{b}}{\textopeno}{l}{\textscy}{\textsubring{r}}} \dicPos{m}[9] \dicFlx{(‑s, ‑ir)}[8] \dicDirectTranslationCS{(dámské) plavky, koupací úbor} \dicIndirectTranslationCS{(jednodílný)}
\dicEntry[sundfit] \dicTerm{sund··fit} \dicIPA{{s}{\textscy}{n}{\textsubring{d}}{f}{\textsci}{\textsubring{d}}} \dicPos{f}[4] \dicFlx{(‑jar, ‑jar)}[9] \dicFieldCat{biol.} \dicDirectTranslationCS{plovací blána}
\dicEntry[sundfugl] \dicTerm{sund··fugl} \dicIPA{{s}{\textscy}{n}{\textsubring{d}}{f}{\textscy}{\r{g}}{\textsubring{l}}} \dicPos{m}[4] \dicFlx{(‑s, ‑ar)}[3] \dicDirectTranslationCS{vodní pták}
\dicEntry[sundföt] \dicTerm{sund··föt} \dicIPA{{s}{\textscy}{n}{\textsubring{d}}{f}{\oe}{\textsubring{d}}} \dicPos{n}[2] \dicFlx{pl}[9] \dicDirectTranslationCS{plavky, koupací oblečení}
\dicEntry[sundgleraugu] \dicTerm{sund··gler·augu} \dicIPA{{s}{\textscy}{n}{\textsubring{d}}{\r{g}}{l}{\textepsilon}{r}{\oe i}{\textbabygamma}{\textscy}} \dicPos{n}[1] \dicFlx{pl}[3] \dicDirectTranslationCS{plavecké brýle}
\dicEntry[sundhetta] \dicTerm{sund··hett|a} \dicIPA{{s}{\textscy}{n}{\textsubring{d}}{h}{\textepsilon}{h}{\textsubring{d}}{a}} \dicPos{f}[1] \dicFlx{(‑u, ‑ur)}[7] \dicDirectTranslationCS{koupací čepice}
\dicEntry[sundhöll] \dicTerm{sund··|höll} \dicIPA{{s}{\textscy}{n}{\textsubring{d}}{h}{\oe}{\textsubring{d}}{\textsubring{l}}} \dicPos{f}[7] \dicFlx{(‑hallar, ‑hallir)}[16] \dicDirectTranslationCS{(krytý) plavecký bazén}
\dicEntry[sundkennsla] \dicTerm{sund··kennsl|a} \dicIPA{{s}{\textscy}{n}{\textsubring{d}}{c\smash{\textsuperscript{h}}}{\textepsilon}{n}{s}{\textsubring{d}}{l}{a}} \dicPos{f}[1] \dicFlx{(‑u)}[5] \dicDirectTranslationCS{lekce\,/\addthin výuka plavání}
\dicEntry[sundknattleikur] \dicTerm{sund··knatt·leik|ur} \dicIPA{{s}{\textscy}{n}{\textsubring{d}}{k\smash{\textsuperscript{h}}}{n}{a}{h}{\textsubring{d}}{l}{ei}{\r{g}}{\textscy}{\textsubring{r}}} \dicPos{m}[9] \dicFlx{(‑s, ‑ir)}[15] \dicFieldCat{sport.} \dicDirectTranslationCS{vodní pólo}
\dicEntry[sundkona] \dicTerm{sund··kon|a} \dicIPA{{s}{\textscy}{n}{\textsubring{d}}{k\smash{\textsuperscript{h}}}{\textopeno}{n}{a}} \dicPos{f}[1] \dicFlx{(‑u, ‑ur)}[27] \dicDirectTranslationCS{plavkyně}
\dicEntry[sundla] \dicTerm{sundl|a} \dicIPA{{s}{\textscy}{n}{\textsubring{d}}{l}{a}} \dicPos{v}[1] \dicFlx{(‑aði)}[66] \dicFlx{impers} \dicPhraseIS{e‑n sundlar} \dicDirectTranslationCS{(kdo) trpí závratí, (komu) se točí hlava}
\dicEntry[sundlaug] \dicTerm{sund··laug} \dicsymFrequent\  \dicIPA{{s}{\textscy}{n}{\textsubring{d}}{l}{\oe i}{x}} \dicPos{f}[4] \dicFlx{(‑ar, ‑ar)}[1] \dicDirectTranslationCS{(plavecký) bazén, koupaliště, plovárna} \dicExampleIS{bjarga barni frá drukknun í sundlaug} \dicExampleCS{zachránit dítě před utonutím v~bazénu}
\dicEntry[sundmaður] \dicTerm{sund··|maður} \dicIPA{{s}{\textscy}{n}{\textsubring{d}}{m}{a}{ð}{\textscy}{\textsubring{r}}} \dicPos{m}[13] \dicFlx{(‑manns, ‑menn)}[2] \dicDirectTranslationCS{plavec, plavkyně}
\dicEntry[sundmagi] \dicTerm{sund··mag|i} \dicIPA{{s}{\textscy}{n}{\textsubring{d}}{m}{a}{i}{j}{\textsci}} \dicPos{m}[1] \dicFlx{(‑a, ‑ar)}[8] \dicFieldCat{biol.} \dicDirectTranslationCS{plynový\,/\addthin plovací měchýř}
\dicEntry[sundnámskeið] \dicTerm{sund··nám·skeið} \dicIPA{{s}{\textscy}{n}{\textsubring{d}}{n}{au}{m}{s}{\r{\textObardotlessj}}{ei}{\texttheta}} \dicPos{n}[2] \dicFlx{(‑s, ‑)}[5] \dicDirectTranslationCS{plavecký kurz}
\dicEntry[sundra] \dicTerm{sundr|a} \dicIPA{{s}{\textscy}{n}{\textsubring{d}}{r}{a}} \dicPos{v}[1] \dicFlx{(‑aði)}[1] \dicFlx{dat} \textbf{1.} \dicSynonym{rjúfa} \dicDirectTranslationCS{(roz)dělit, rozpoltit, (roz)štěpit} \dicExampleIS{sundra hjónabandi} \dicExampleCS{rozpoltit manželství}  \textbf{2.} \dicFieldCat{chem.} \dicDirectTranslationCS{rozložit, rozkládat};  \dicIdiom{sundrast}{ \dicPhraseIS{sundrast}} \dicFlx{refl} {\textbf{a.}} \dicDirectTranslationCS{rozdělit se, oddělit se};  {\textbf{b.}} \dicFieldCat{chem.} \dicDirectTranslationCS{rozložit se (látka v~těle ap.)}
\dicEntry[sundrun] \dicTerm{sundr|un} \dicIPA{{s}{\textscy}{n}{\textsubring{d}}{r}{\textscy}{\textsubring{n}}} \dicPos{f}[7] \dicFlx{(‑unar)}[9] \textbf{1.} \dicDirectTranslationCS{rozklad, rozpad}  \textbf{2.} \dicFieldCat{chem.} \dicDirectTranslationCS{rozklad, rozložení} \dicExampleIS{sundrun lífrænna næringarefna} \dicExampleCS{rozklad organických výživných látek}
\dicEntry[sundrung] \dicTerm{sundr··ung} \dicIPA{{s}{\textscy}{n}{\textsubring{d}}{r}{u}{\ng}{\r{g}}} \dicPos{f}[4] \dicFlx{(‑ar)}[3] \dicSynonym{aðskilnaður} \dicDirectTranslationCS{rozdělení, oddělení, rozštěpení}
\dicEntry[sundskýla] \dicTerm{sund··skýl|a} \dicIPA{{s}{\textscy}{n}{\textsubring{d}}{s}{\r{\textObardotlessj}}{i}{l}{a}} \dicPos{f}[1] \dicFlx{(‑u, ‑ur)}[19] \dicDirectTranslationCS{(pánské) plavky}
\dicEntry[sundtak] \dicTerm{sund··|tak} \dicIPA{{s}{\textscy}{n}{\textsubring{d}}{t\smash{\textsuperscript{h}}}{a}{\r{g}}} \dicPos{n}[2] \dicFlx{(‑taks, ‑tök)}[8] \dicDirectTranslationCS{plavecké tempo, plavecký záběr}
\dicEntry[sundur] \dicTerm{sundur} \dicsymFrequent\  \dicIPA{{s}{\textscy}{n}{\textsubring{d}}{\textscy}{\textsubring{r}}} \dicPos{adv} \dicPhraseIS{engjast sundur og saman} \dicFlx{refl} \dicDirectTranslationCS{(z)křivit se, (z)kroutit se (bolestí ap.)};  \dicPhraseIS{falla\,/\addthin detta í sundur} \dicDirectTranslationCS{rozpadnout se, rozsypat se};  \dicPhraseIS{í sundur} \dicFlx{adv} \dicDirectTranslationCS{na kusy\,/\addthin části}
\dicEntry[sundurgerð] \dicTerm{sundur··gerð} \dicIPA{{s}{\textscy}{n}{\textsubring{d}}{\textscy}{r}{\r{\textObardotlessj}}{\textepsilon}{r}{\texttheta}} \dicPos{f}[7] \dicFlx{(‑ar)}[3] \dicSynonym{tilgerð} \dicDirectTranslationCS{strojenost, afektovanost}

\dicEntry[sundurgerðarmaður] \textls[35]{\dicTerm{sundur·gerðar··|maður} \dicIPA{{s}\-{\textscy}\-{n}\-{\textsubring{d}}\-{\textscy}\-{r}\-{\r{\textObardotlessj}}\-{\textepsilon}\-{r}\-{ð}\-{a}\-{r}\-{m}\-{a}\-{ð}\-{\textscy}\-{\textsubring{r}}} \dicPos{m}[13]} \dicFlx{(‑manns, ‑menn)}[2] \dicDirectTranslationCS{elegán, švihák, fešák}

\dicEntry[sundurgreina] \dicTerm{sundur··grein|a} \dicIPA{{s}{\textscy}{n}{\textsubring{d}}{\textscy}{r}{\r{g}}{r}{ei}{n}{a}} \dicPos{v}[2] \dicFlx{(‑di, ‑t)}[141] \dicFlx{acc} \dicDirectTranslationCS{analyzovat, rozebrat (analyticky)}
\dicEntry[sundurgreining] \dicTerm{sundur··grein·ing} \dicIPA{{s}{\textscy}{n}{\textsubring{d}}{\textscy}{r}{\r{g}}{r}{ei}{n}{i}{\ng}{\r{g}}} \dicPos{f}[4] \dicFlx{(‑ar)}[7] \dicDirectTranslationCS{analýza, rozbor, rozklad}
\dicEntry[sundurlaus] \dicTerm{sundur··laus} \dicIPA{{s}{\textscy}{n}{\textsubring{d}}{\textscy}{r}{l}{\oe i}{s}} \dicPos{adj}[5]\dicFlx{}[-1] \dicDirectTranslationCS{nesouvislý, nesourodý} \dicExampleIS{sundurlaus ræða} \dicExampleCS{nesouvislá řeč}
\dicEntry[sundurleitni] \dicTerm{sundur··leitn|i} \dicIPA{{s}{\textscy}{n}{\textsubring{d}}{\textscy}{r}{l}{ei}{h}{\textsubring{d}}{n}{\textsci}} \dicPos{f}[3] \dicFlx{(‑i)}[3] \textbf{1.} \dicDirectTranslationCS{různorodost, rozmanitost, heterogenita}  \textbf{2.} \dicDirectTranslationCS{různorodost, nesourodost}
\dicEntry[sundurleitur] \dicTerm{sundur··leitur} \dicIPA{{s}{\textscy}{n}{\textsubring{d}}{\textscy}{r}{l}{ei}{\textsubring{d}}{\textscy}{\textsubring{r}}} \dicPos{adj}[1]\dicFlx{}[-1] \textbf{1.} \dicSynonym{mismunandi\smash{\textsuperscript{1}}} \dicDirectTranslationCS{různorodý, rozmanitý, heterogenní} \dicExampleIS{sundurleitar skoðanir} \dicExampleCS{různorodé názory}  \textbf{2.} \dicSynonym*{sem stingur í stúf} \dicDirectTranslationCS{různorodý, nesourodý} \dicExampleIS{sundurleitir litir} \dicExampleCS{různorodé barvy}
\dicEntry[sundurliða] \dicTerm{sundur··lið|a} \dicIPA{{s}{\textscy}{n}{\textsubring{d}}{\textscy}{r}{l}{\textsci}{ð}{a}} \dicPos{v}[1] \dicFlx{(‑aði)}[1] \dicFlx{acc} \dicDirectTranslationCS{rozložit, rozkládat, rozčlenit, rozčleňovat, rozepsat, rozepisovat (do položek ap.)} \dicExampleIS{sundurliða kostnaðinn} \dicExampleCS{rozepsat náklady}
\dicEntry[sundurliðun] \dicTerm{sundur··lið|un} \dicIPA{{s}{\textscy}{n}{\textsubring{d}}{\textscy}{r}{l}{\textsci}{ð}{\textscy}{\textsubring{n}}} \dicPos{f}[7] \dicFlx{(‑unar)}[9] \dicDirectTranslationCS{rozložení, rozčlenění, rozepsání}
\dicEntry[sundurlyndi] \dicTerm{sundur··lyndi} \dicIPA{{s}{\textscy}{n}{\textsubring{d}}{\textscy}{r}{l}{\textsci}{n}{\textsubring{d}}{\textsci}} \dicPos{n}[2] \dicFlx{(‑s)}[20] \dicSynonym*{ósætti} \dicDirectTranslationCS{neshoda, rozkol, roztržka} \dicExampleIS{sundurlyndi á milli manna} \dicExampleCS{roztržka mezi lidmi}
\dicEntry[sundurorða] \dicTerm{sundur··orða} \dicIPA{{s}{\textscy}{n}{\textsubring{d}}{\textscy}{r}{\textopeno}{r}{ð}{a}} \dicPos{adj}[13] \dicFlx{indecl}[1] \dicPhraseIS{e‑jum verður sundurorða út af e‑u} \dicFlx{impers} \dicDirectTranslationCS{(kdo) se dostane do hádky kvůli (čemu)}
\dicEntry[sundurskilja] \dicTerm{sundur··skil|ja} \dicIPA{{s}{\textscy}{n}{\textsubring{d}}{\textscy}{\textsubring{r}}{s}{\r{\textObardotlessj}}{\textsci}{l}{j}{a}} \dicPos{v}[4] \dicFlx{(‑di, ‑ið)}[30] \dicFlx{acc} \dicSynonym{sundra} \dicDirectTranslationCS{rozdělit, rozdělovat, oddělit, oddělovat}
\dicEntry[sundurþykki] \dicTerm{sundur··þykki} \dicIPA{{s}{\textscy}{n}{\textsubring{d}}{\textscy}{\textsubring{r}}{\texttheta}{\textsci}{h}{\r{\textObardotlessj}}{\textsci}} \dicPos{n}[2] \dicFlx{(‑s)}[20] \dicLink{sundurþykkja}
\dicEntry[sundurþykkja] \dicTerm{sundur··þykkj|a}\dicTerm{, sundurþykki} \dicIPA{{s}\-{\textscy}\-{n}\-{\textsubring{d}}\-{\textscy}\-{\textsubring{r}}\-{\texttheta}\-{\textsci}\-{h}\-{\r{\textObardotlessj}}\-{a}\-} \dicPos{f}[1] \dicFlx{(‑u)}[5] \dicSynonym*{ósæt\-ti} \dicDirectTranslationCS{rozkol, rozpor, neshoda}
\dicEntry[sundurþykkur] \dicTerm{sundur··þykkur} \dicIPA{{s}{\textscy}{n}{\textsubring{d}}{\textscy}{\textsubring{r}}{\texttheta}{\textsci}{h}{\r{g}}{\textscy}{\textsubring{r}}} \dicPos{adj}[1]\dicFlx{}[-6] \dicSynonym{ósammála} \dicDirectTranslationCS{rozporný, rozporuplný, antagonistický}
\dicEntry[sungið] \dicTerm{sungið} \dicIPA{{s}{u}{\textltailn}{\r{\textObardotlessj}}{\textsci}{\texttheta}} \dicPos{v} \dicFlx{supin} \dicLink{syngja}
\dicEntry[sungum] \dicTerm{sungum} \dicIPA{{s}{u}{\ng}{\r{g}}{\textscy}{\textsubring{m}}} \dicPos{v} \dicFlx{ind pf pl 1 pers} \dicLink{syngja}
\dicEntry[sunnan] \dicTerm{sunnan} \dicsymFrequent\  \dicIPA{{s}{\textscy}{n}{\textlengthmark}{a}{\textsubring{n}}} \dicPos{prep\,/\addthin adv} \dicFlx{adv} \dicDirectTranslationCS{z~jihu};  \dicIdiom{að}[sunnan]{ \dicPhraseIS{að sunnan}} \dicFlx{adv} {\textbf{a.}} \dicDirectTranslationCS{z~Reykjavíku};  {\textbf{b.}} \dicDirectTranslationCS{z~jihu (Islandu ap.)};  {\textbf{c.}} \dicDirectTranslationCS{po\,/\addthin na jižní straně};  \dicIdiom{sunnan}[að]{ \dicPhraseIS{sunnan að}} \dicFlx{adv} \dicDirectTranslationCS{z~jihu};  \dicIdiom{sunnan}[á]{ \dicPhraseIS{sunnan á e‑u}} \dicFlx{prep} \dicDirectTranslationCS{z~jižní strany (čeho)};  \dicIdiom{sunnan}[frá]{ \dicPhraseIS{sunnan frá}} \dicFlx{adv} \dicDirectTranslationCS{z~jihu, z~jižní strany};  \dicIdiom{sunnan}[í móti]{ \dicPhraseIS{sunnan í móti}} \dicFlx{adv} \dicDirectTranslationCS{směrem na jih};  \dicIdiom{sunnan}[megin]{ \dicPhraseIS{sunnan megin}} \dicFlx{adv} \dicDirectTranslationCS{na jižní straně, z~jižní části};  \dicIdiom{sunnan}[til]{ \dicPhraseIS{sunnan til}} \dicFlx{adv} \dicDirectTranslationCS{v~jižní části};  \dicIdiom{sunnan}[undir]{ \dicPhraseIS{sunnan undir e‑u}} \dicFlx{prep} \dicDirectTranslationCS{jižně pod (čím), z~jižní strany (čeho)};  \dicIdiom{sunnan}[við]{ \dicPhraseIS{sunnan við e‑ð}} \dicFlx{prep} \dicDirectTranslationCS{na jižní straně (čeho), jižně od (čeho)};  \dicIdiom{sunnan}[yfir]{ \dicPhraseIS{sunnan yfir e‑ð}} \dicFlx{prep} \dicDirectTranslationCS{z~jihu přes (co)}
\dicEntry[sunnanátt] \dicTerm{sunnan··átt} \dicIPA{{s}{\textscy}{n}{\textlengthmark}{a}{n}{au}{h}{\textsubring{d}}} \dicPos{f}[7] \dicFlx{(‑ar, ‑ir)}[1] \dicDirectTranslationCS{jižní vítr}
\dicEntry[sunnanverður] \dicTerm{sunnan··verður} \dicIPA{{s}{\textscy}{n}{\textlengthmark}{a}{n}{v}{\textepsilon}{r}{ð}{\textscy}{\textsubring{r}}} \dicPos{adj}[2]\dicFlx{}[-4] \dicDirectTranslationCS{jižní, jižně položený} \dicExampleIS{við sunnanverðan fjörðinn} \dicExampleCS{na jižní straně fjordu};  \dicPhraseIS{að sunnanverðu} \dicFlx{adv} \dicDirectTranslationCS{po\,/\addthin na jižní straně}
\dicEntry[sunnar] \dicTerm{sunnar} \dicIPA{{s}{\textscy}{n}{\textlengthmark}{a}{\textsubring{r}}} \dicPos{adv} \dicFlx{comp} \dicLink{suður\smash{\textsuperscript{2}}}
\dicEntry[sunnlenskur] \dicTerm{sunn··lenskur} \dicIPA{{s}{\textscy}{n}{l}{\textepsilon}{n}{s}{\r{g}}{\textscy}{\textsubring{r}}} \dicPos{adj}[1]\dicFlx{}[-1] \dicDirectTranslationCS{(jsoucí) z~jihu, jižanský}
\dicEntry[sunnudagaskóli] \dicTerm{sunnu·daga··skól|i} \dicIPA{{s}{\textscy}{n}{\textlengthmark}{\textscy}{\textsubring{d}}{a}{\textbabygamma}{a}{s}{\r{g}}{ou}{l}{\textsci}} \dicPos{m}[1] \dicFlx{(‑a, ‑ar)}[1] \dicDirectTranslationCS{nedělní škola}
\dicEntry[sunnudagur] \dicTerm{sunnu··dag|ur} \dicsymFrequent\  \dicIPA{{s}{\textscy}{n}{\textlengthmark}{\textscy}{\textsubring{d}}{a}{\textbabygamma}{\textscy}{\textsubring{r}}} \dicPos{m}[6] \dicFlx{(‑s, ‑ar)}[62] \dicDirectTranslationCS{neděle};  \dicPhraseIS{á sunnudaginn} \dicFlx{adv} \dicDirectTranslationCS{v~neděli};  \dicPhraseIS{á sunnudaginn kemur} \dicFlx{adv} \dicDirectTranslationCS{příští neděli};  \dicPhraseIS{á sunnudaginn var} \dicFlx{adv} \dicDirectTranslationCS{minulou neděli};  \dicPhraseIS{á sunnudeginum} \dicFlx{adv} \dicDirectTranslationCS{tuto neděli};  \dicPhraseIS{á sunnudögum} \dicFlx{adv} \dicDirectTranslationCS{o~nedělích}
\dicEntry[supum] \dicTerm{supum} \dicIPA{{s}{\textscy}{\textlengthmark}{\textsubring{b}}{\textscy}{\textsubring{m}}} \dicPos{v} \dicFlx{ind pf pl 1 pers} \dicLink{súpa\smash{\textsuperscript{2}}}
\dicEntry[surfum] \dicTerm{surfum} \dicIPA{{s}{\textscy}{r}{v}{\textscy}{\textsubring{m}}} \dicPos{v} \dicFlx{ind pf pl 1 pers} \dicLink{sverfa}
\dicEntry[surtur] \dicTerm{surt|ur} \dicIPA{{s}{\textscy}{\textsubring{r}}{\textsubring{d}}{\textscy}{\textsubring{r}}} \dicPos{m}[6] \dicFlx{(‑s, ‑ar)}[4] \dicLangCat{han.} \dicDirectTranslationCS{negr, negerka, černoch, černoška}
\dicEntry[suss] \dicTerm{suss} \dicIPA{{s}{\textscy}{s}{\textlengthmark}} \dicPos{inter} \dicDirectTranslationCS{pšt, pst}
\dicEntry[sussa] \dicTerm{suss|a} \dicIPA{{s}{\textscy}{s}{\textlengthmark}{a}} \dicPos{v}[1] \dicFlx{(‑aði)}[34] \dicDirectTranslationCS{utišit, umlčet} \dicExampleIS{sussa barn} \dicExampleCS{utišit dítě};  \dicPhraseIS{sussa á e‑n} \dicDirectTranslationCS{syknout na (koho), zasyčet na (koho)} \dicIndirectTranslationCS{(aby ztichl)}
\dicEntry[sussu] \dicTerm{sussu} \dicIPA{{s}{\textscy}{s}{\textlengthmark}{\textscy}} \dicPos{inter} \dicDirectTranslationCS{nono} \dicIndirectTranslationCS{(vyjadřuje přání, aby někdo něčeho nechal (pláče ap.))}
\dicEntry[sú] \dicTerm{sú} \dicIPA{{s}{u}{\textlengthmark}} \dicPos{pron}[5] \dicFlx{dem nom sg f} \dicLink{sá}
\dicEntry[Súdan] \dicTerm{Súdan} \dicIPA{{s}{u}{\textlengthmark}{\textsubring{d}}{a}{\textsubring{n}}} \dicPos{n}[2] \dicFlx{(‑\,/\addthin ‑s)}[35] \dicFieldCat{geog.} \dicDirectTranslationCS{Súdán}
\dicEntry[Súdani] \dicTerm{Súdan|i} \dicIPA{{s}{u}{\textlengthmark}{\textsubring{d}}{a}{n}{\textsci}} \dicPos{m}[3] \dicFlx{(‑a, ‑ar\,/\addthin ‑ir)}[3] \dicDirectTranslationCS{Súdánec, Súdánka}
\dicEntry[súdanskur] \dicTerm{súdanskur} \dicIPA{{s}{u}{\textlengthmark}{\textsubring{d}}{a}{n}{s}{\r{g}}{\textscy}{\textsubring{r}}} \dicPos{adj}[1] \dicFlx{(f súdönsk)}[3] \dicDirectTranslationCS{súdánský}
\dicEntry[súð] \dicTerm{súð} \dicIPA{{s}{u}{\textlengthmark}{\texttheta}} \dicPos{f}[7] \dicFlx{(‑ar, ‑ir)}[1] \textbf{1.} \dicSynonym*{innanvert þak} \dicDirectTranslationCS{šikmý strop};  \dicPhraseIS{undir súð} \dicFlx{adv} \dicDirectTranslationCS{v~podkroví}  \textbf{2.} \dicSynonym*{skipshlið} \dicDirectTranslationCS{lodní bok};  \dicPhraseIS{láta vaða á súðum} \dicLangCat{přen.} \dicDirectTranslationCS{plést páté přes deváté, mluvit z~cesty}
\dicEntry[súgur] \dicTerm{súg|ur} \dicIPA{{s}{u}{\textlengthmark}{\textscy}{\textsubring{r}}} \dicPos{m}[6] \dicFlx{(‑s)}[23] \dicSynonym*{dragsúgur} \dicDirectTranslationCS{průvan} \dicExampleIS{Það er súgur um dyrnar.} \dicExampleCS{Ve dveřích je průvan.};  \dicPhraseIS{fara\,/\addthin ganga í súginn} \dicSynonym*{fara til ónýtis} \dicDirectTranslationCS{přijít vniveč}
\dicEntry[súkkulaði] \dicTerm{súkkulaði} \dicIPA{{s}{u}{h}{\r{g}}{\textscy}{l}{a}{ð}{\textsci}} \dicPos{n}[2] \dicFlx{(‑s)}[20] \textbf{1.} \dicDirectTranslationCS{čokoláda} \dicIndirectTranslationCS{(potravinářský výrobek z~kakaa)}  \textbf{2.} \dicDirectTranslationCS{čokoláda} \dicIndirectTranslationCS{(nápoj)}
\dicEntry[súkkulaðiís] \dicTerm{súkkulaði··ís} \dicIPA{{s}{u}{h}{\r{g}}{\textscy}{l}{a}{ð}{\textsci}{i}{s}} \dicPos{m}[4] \dicFlx{(‑s, ‑ar)}[25] \dicDirectTranslationCS{čokoládová zmrzlina}
\dicEntry[súkkulaðikaka] \dicTerm{súkkulaði··|kaka} \dicIPA{{s}{u}{h}{\r{g}}{\textscy}{l}{a}{ð}{\textsci}{k\smash{\textsuperscript{h}}}{a}{\r{g}}{a}} \dicPos{f}[1] \dicFlx{(‑köku, ‑kökur)}[20] \dicFieldCat{kulin.} \dicDirectTranslationCS{čokoládový dort}
\dicEntry[súkkulaðistykki] \dicTerm{súkkulaði··stykki} \dicIPA{{s}{u}{h}{\r{g}}{\textscy}{l}{a}{ð}{\textsci}{s}{\textsubring{d}}{\textsci}{h}{\r{\textObardotlessj}}{\textsci}} \dicPos{n}[2] \dicFlx{(‑s, ‑)}[16] \dicDirectTranslationCS{tabulka čokolády}
\dicEntry[súla] \dicTerm{súl|a} \dicIPA{{s}{u}{\textlengthmark}{l}{a}} \dicPos{f}[1] \dicFlx{(‑u, ‑ur)}[13] \textbf{1.} \dicSynonym*{burðarsúla} \dicDirectTranslationCS{sloup, pilíř} \dicExampleIS{reisa súlu við vegginn} \dicExampleCS{vztyčit u~zdi pilíř}  \textbf{2.} \dicSynonym{stöng} \dicDirectTranslationCS{sloupek, tyčka (u~stanu ap.)}  \textbf{3.} \dicFieldCat{zool.} \dicDirectTranslationCS{terej bílý} \textit{(l.~{\textLA{Sula bassana}})}  \dicsymPhoto\ ;  \dicPhraseIS{skjóta súlur} \dicLangCat{přen.} \dicDirectTranslationCS{házet žabky (na vodní hladině)}
\dicFigure{21205.jpg}{Súla}{Súla - Hagerty, Ryan, Biolib, PD}
\dicEntry[súld] \dicTerm{súld} \dicIPA{{s}{u}{l}{\textsubring{d}}} \dicPos{f}[7] \dicFlx{(‑ar, ‑ir)}[1] \dicSynonym{úði} \dicDirectTranslationCS{mrholení, mžení}
\dicEntry[súlfíð] \dicTerm{súlfíð} \dicIPA{{s}{u}{l}{f}{i}{\texttheta}} \dicPos{n}[2] \dicFlx{(‑s)}[2] \dicFieldCat{chem.} \dicDirectTranslationCS{sulfid}
\dicEntry[súlnagöng] \dicTerm{súlna··göng} \dicIPA{{s}{u}{l}{n}{a}{\r{g}}{\oe i}{\ng}{\r{g}}} \dicPos{n}[2] \dicFlx{pl}[9] \dicDirectTranslationCS{kolonáda}
\dicEntry[súlurit] \dicTerm{súlu··rit} \dicIPA{{s}{u}{\textlengthmark}{l}{\textscy}{r}{\textsci}{\textsubring{d}}} \dicPos{n}[2] \dicFlx{(‑s, ‑)}[5] \dicDirectTranslationCS{sloupcový graf\,/\addthin diagram, histogram}
\dicEntry[súpa] \dicTerm{súp|a\smash{\textsuperscript{1}}} \dicIPA{{s}{u}{\textlengthmark}{\textsubring{b}}{a}} \dicPos{f}[1] \dicFlx{(‑u, ‑ur)}[19] \dicDirectTranslationCS{polévka, polívka} \dicExampleIS{sveppasúpa} \dicExampleCS{houbová polévka};  \dicPhraseIS{sitja í súpunni} \dicLangCat{přen.} \dicDirectTranslationCS{být v~pěkné bryndě}
\dicEntry[súpa] \dicTerm{súpa\smash{\textsuperscript{2}}} \dicIPA{{s}{u}{\textlengthmark}{\textsubring{b}}{a}} \dicPos{v}[6] \dicFlx{(sýp, saup, supum, sypi, sopið)}[106] \dicSynonym{drekka} \dicDirectTranslationCS{napít se, loknout si} \dicExampleIS{súpa á kaffinu} \dicExampleCS{loknout si kávy}
\dicEntry[súpudiskur] \dicTerm{súpu··disk|ur} \dicIPA{{s}{u}{\textlengthmark}{\textsubring{b}}{\textscy}{\textsubring{d}}{\textsci}{s}{\r{g}}{\textscy}{\textsubring{r}}} \dicPos{m}[6] \dicFlx{(‑s, ‑ar)}[10] \dicDirectTranslationCS{polévkový\,/\addthin hluboký talíř}
\dicEntry[súpukjöt] \dicTerm{súpu··kjöt} \dicIPA{{s}{u}{\textlengthmark}{\textsubring{b}}{\textscy}{c\smash{\textsuperscript{h}}}{\oe}{\textsubring{d}}} \dicPos{n}[2] \dicFlx{(‑s)}[2] \dicDirectTranslationCS{maso na polévku}
\dicEntry[súpuskál] \dicTerm{súpu··skál} \dicIPA{{s}{u}{\textlengthmark}{\textsubring{b}}{\textscy}{s}{\r{g}}{au}{\textsubring{l}}} \dicPos{f}[4] \dicFlx{(‑ar, ‑ar\,/\addthin ‑ir)}[30] \dicDirectTranslationCS{polévková mísa}
\dicEntry[súputeningur] \dicTerm{súpu··ten·ing|ur} \dicIPA{{s}{u}{\textlengthmark}{\textsubring{b}}{\textscy}{t\smash{\textsuperscript{h}}}{\textepsilon}{n}{i}{\ng}{\r{g}}{\textscy}{\textsubring{r}}} \dicPos{m}[6] \dicFlx{(‑s, ‑ar)}[8] \dicDirectTranslationCS{bujón\,/\addthin vývar v~kostce, kostkový bujón}
\dicEntry[súr] \dicTerm{súr\smash{\textsuperscript{1}}} \dicIPA{{s}{u}{\textlengthmark}{\textsubring{r}}} \dicPos{m}[4] \dicFlx{(‑s)}[18] \dicSynonym*{súr mysa} \dicDirectTranslationCS{kyselá syrovátka}
\dicEntry[súr] \dicTerm{súr\smash{\textsuperscript{2}}} \dicIPA{{s}{u}{\textlengthmark}{\textsubring{r}}} \dicPos{adj}[5] \dicFlx{(f ‑)}[8] \textbf{1.} \dicSynonym*{súr á bragðið} \dicDirectTranslationCS{kyselý}  \textbf{2.} \dicSynonym{fúll} \dicDirectTranslationCS{kyselý, mrzutý, zatrpklý};  \dicPhraseIS{vera súr á svipinn} \dicDirectTranslationCS{mít kyselý výraz}  \textbf{3.} \dicFieldCat{chem.} \dicDirectTranslationCS{kyselý};  \dicPhraseIS{súrt regn} \dicDirectTranslationCS{kyselý déšť};  \dicIdiom{súr}{ \dicPhraseIS{súrt og sætt}} \dicFlx{adv} \dicDirectTranslationCS{v~dobrém i~ve zlém}
\dicEntry[súra] \dicTerm{súr|a} \dicIPA{{s}{u}{\textlengthmark}{r}{a}} \dicPos{f}[1] \dicFlx{(‑u, ‑ur)}[7] \dicFieldCat{bot.} \dicDirectTranslationCS{šťovík} \textit{(l.~{\textLA{Rumex}})}  \dicsymPhoto\ 
\dicFigure{ds_image_sura_0_1.jpg}{Súra}{Súra - Sten Porse, CC BY-SA 3.0}
\dicEntry[súrbragð] \dicTerm{súr··bragð} \dicIPA{{s}{u}{r}{\textsubring{b}}{r}{a}{\textbabygamma}{\texttheta}} \dicPos{n}[2] \dicFlx{(‑s)}[2] \dicDirectTranslationCS{kyselá chuť}
\dicEntry[súrdeig] \dicTerm{súr··deig} \dicIPA{{s}{u}{r}{\textsubring{d}}{ei}{x}} \dicPos{n}[2] \dicFlx{(‑s)}[2] \dicDirectTranslationCS{kvas, kvásek}
\dicEntry[súrefni] \dicTerm{súr··efni} \dicIPA{{s}{u}{\textlengthmark}{r}{\textepsilon}{\textsubring{b}}{n}{\textsci}} \dicPos{n}[2] \dicFlx{(‑s)}[20] \dicFieldCat{chem.} \dicDirectTranslationCS{kyslík} \textit{(l.~{\textLA{O, Oxygenium}})}
\dicEntry[súrefnisgríma] \dicTerm{súr·efnis··grím|a} \dicIPA{{s}{u}{\textlengthmark}{r}{\textepsilon}{\textsubring{b}}{n}{\textsci}{s}{\r{g}}{r}{i}{m}{a}} \dicPos{f}[1] \dicFlx{(‑u, ‑ur)}[19] \dicDirectTranslationCS{kyslíková maska}
\dicEntry[súrhey] \dicTerm{súr··hey} \dicIPA{{s}{u}{\textlengthmark}{\textsubring{r}}{h}{ei}} \dicPos{n}[2] \dicFlx{(‑s, ‑)}[13] \dicSynonym{vothey} \dicDirectTranslationCS{siláž}
\dicEntry[Súrínam] \dicTerm{Súrínam} \dicIPA{{s}{u}{\textlengthmark}{r}{i}{n}{a}{\textsubring{m}}} \dicPos{n}[4] \dicFlx{indecl}[2] \dicFieldCat{geog.} \dicDirectTranslationCS{Surinam}
\dicEntry[Súrínami] \dicTerm{Súrínam|i} \dicIPA{{s}{u}{\textlengthmark}{r}{i}{n}{a}{m}{\textsci}} \dicPos{m}[1] \dicFlx{(‑a, ‑ar)}[8] \dicDirectTranslationCS{Surinamec, Surinamka}
\dicEntry[súrínamskur] \dicTerm{súrínamskur} \dicIPA{{s}{u}{\textlengthmark}{r}{i}{n}{a}{m}{s}{\r{g}}{\textscy}{\textsubring{r}}} \dicPos{adj}[1] \dicFlx{(f súrínömsk)}[3] \dicDirectTranslationCS{surinamský}
\dicEntry[súrmatur] \dicTerm{súr··mat|ur}\dicTerm{, súrmeti} \dicIPA{{s}\-{u}\-{r}\-{m}\-{a}\-{\textsubring{d}}\-{\textscy}\-{\textsubring{r}}\-} \dicPos{m}[10] \dicFlx{(‑ar)}[15] \dicFieldCat{kulin.} \dicIndirectTranslationCS{maso nakládané v~kyselé syrovátce}
\dicEntry[súrmeti] \dicTerm{súr··meti} \dicIPA{{s}{u}{r}{m}{\textepsilon}{\textsubring{d}}{\textsci}} \dicPos{n}[2] \dicFlx{(‑s)}[20] \dicLink{súrmatur}
\dicEntry[súrmjólk] \dicTerm{súr··mjólk} \dicIPA{{s}{u}{r}{m}{j}{ou}{\textsubring{l}}{\r{g}}} \dicPos{f}[10] \dicFlx{(‑ur)}[2] \dicDirectTranslationCS{kefír, kysané mléko}
\dicEntry[súrna] \dicTerm{súrn|a} \dicIPA{{s}{u}{r}{\textsubring{d}}{n}{a}} \dicPos{v}[1] \dicFlx{(‑aði)}[61] \textbf{1.} \dicSynonym*{verða súr} \dicDirectTranslationCS{(z)kysnout, zakysnout} \dicExampleIS{Mjólkin súrnaði.} \dicExampleCS{Mléko zkyslo.}  \textbf{2.} \dicPhraseIS{e‑m súrnar í augum} \dicFlx{impers} \dicDirectTranslationCS{(koho) pálí oči} \dicExampleIS{Honum súrnar í augum.} \dicExampleCS{Pálí ho oči.}
\dicEntry[súrnun] \dicTerm{súrn|un} \dicIPA{{s}{u}{r}{\textsubring{d}}{n}{\textscy}{\textsubring{n}}} \dicPos{f}[7] \dicFlx{(‑unar)}[9] \dicDirectTranslationCS{(z)kysnutí}
\dicEntry[súrrealismi] \dicTerm{súrreal··ism|i} \dicIPA{{s}{u}{r}{\textlengthmark}{\textepsilon}{a}{l}{\textsci}{s}{m}{\textsci}} \dicPos{m}[1] \dicFlx{(‑a)}[3] \dicDirectTranslationCS{surrealismus}
\dicEntry[súrrealisti] \dicTerm{súrreal··ist|i} \dicIPA{{s}{u}{r}{\textlengthmark}{\textepsilon}{a}{l}{\textsci}{s}{\textsubring{d}}{\textsci}} \dicPos{m}[1] \dicFlx{(‑a, ‑ar)}[1] \dicDirectTranslationCS{surrealista, surrealistka}
\dicEntry[súrsa] \dicTerm{súrs|a} \dicIPA{{s}{u}{\textsubring{r}}{s}{a}} \dicPos{v}[1] \dicFlx{(‑aði)}[1] \dicFlx{acc} \dicDirectTranslationCS{naložit\,/\addthin nakládat (do kyselého nálevu)}
\dicEntry[sút] \dicTerm{sút} \dicIPA{{s}{u}{\textlengthmark}{\textsubring{d}}} \dicPos{f}[7] \dicFlx{(‑ar, ‑ir)}[1] \textbf{1.} \dicSynonym{kvíði\smash{\textsuperscript{1}}} \dicDirectTranslationCS{znepokojení, úzkost}  \textbf{2.} \dicFieldCat{psych.} \dicSynonym{depurð} \dicDirectTranslationCS{deprese}
\dicEntry[súta] \dicTerm{sút|a\smash{\textsuperscript{1}}} \dicIPA{{s}{u}{\textlengthmark}{\textsubring{d}}{a}} \dicPos{v}[1] \dicFlx{(‑aði)}[1] \dicFlx{acc} \dicDirectTranslationCS{vyčinit, vydělat, vydělávat (kůži ap.)}
\dicEntry[súta] \dicTerm{sút|a\smash{\textsuperscript{2}}} \dicIPA{{s}{u}{\textlengthmark}{\textsubring{d}}{a}} \dicPos{v}[1] \dicFlx{(‑aði)}[1] \dicFlx{acc} \dicSynonym{harma} \dicDirectTranslationCS{truchlit, oplakávat}
\dicEntry[sútaður] \dicTerm{sút··|aður} \dicIPA{{s}{u}{\textlengthmark}{\textsubring{d}}{a}{ð}{\textscy}{\textsubring{r}}} \dicPos{adj}[3] \dicFlx{(f ‑uð)}[4] \dicDirectTranslationCS{vyčiněný, vydělaný (kůže ap.)}
\dicEntry[sútari] \dicTerm{sút··ar|i} \dicIPA{{s}{u}{\textlengthmark}{\textsubring{d}}{a}{r}{\textsci}} \dicPos{m}[1] \dicFlx{(‑a, ‑ar)}[13] \dicDirectTranslationCS{koželuh, koželužka}
\dicEntry[sútun] \dicTerm{sút|un} \dicIPA{{s}{u}{\textlengthmark}{\textsubring{d}}{\textscy}{\textsubring{n}}} \dicPos{f}[7] \dicFlx{(‑unar)}[9] \dicDirectTranslationCS{vyčinění, vydělání, vydělávání (kůže ap.)}
\dicEntry[svað] \dicTerm{svað} \dicIPA{{s}{v}{a}{\textlengthmark}{\texttheta}} \dicPos{n}[2] \dicFlx{(‑s, svöð)}[8] \dicSynonym{fen} \dicDirectTranslationCS{bahnisko, bláto, mokřina} \dicExampleIS{Það fór allt í svað.} \dicExampleCS{Všechno se to proměnilo v~bahnisko.};  \dicPhraseIS{traðka e‑n ofan\,/\addthin niður í svaðið} \dicLangCat{přen.} \dicDirectTranslationCS{zašlapat (koho) do země}
\dicEntry[svaðilför] \dicTerm{svaðil··|för} \dicIPA{{s}{v}{a}{\textlengthmark}{ð}{\textsci}{l}{f}{\oe}{\textsubring{r}}} \dicPos{f}[7] \dicFlx{(‑farar, ‑farir)}[16] \dicDirectTranslationCS{náročná\,/\addthin riskantní cesta (turistická ap.)}
\dicEntry[svaf] \dicTerm{svaf} \dicIPA{{s}{v}{a}{\textlengthmark}{f}} \dicPos{v} \dicFlx{ind pf sg 1 pers} \dicLink{sofa}
\dicEntry[svahílí] \dicTerm{svahílí} \dicIPA{{s}{v}{a}{\textlengthmark}{h}{i}{l}{i}} \dicPos{n}[4] \dicFlx{indecl}[2] \dicDirectTranslationCS{svahilština}
\dicEntry[svaka] \dicTerm{svaka} \dicIPA{{s}{v}{a}{\textlengthmark}{\r{g}}{a}} \dicPos{adv} \dicDirectTranslationCS{fantasticky, super}
\dicEntry[svakalega] \dicTerm{svaka··lega} \dicIPA{{s}{v}{a}{\textlengthmark}{\r{g}}{a}{l}{\textepsilon}{\textbabygamma}{a}} \dicPos{adv} \dicDirectTranslationCS{příšerně, děsně, strašně}
\dicEntry[svakalegur] \dicTerm{svaka··legur} \dicIPA{{s}{v}{a}{\textlengthmark}{\r{g}}{a}{l}{\textepsilon}{\textbabygamma}{\textscy}{\textsubring{r}}} \dicPos{adj}[1]\dicFlx{}[-8] \textbf{1.} \dicDirectTranslationCS{strašný, enormní, obrovský (lhář ap.)} \dicExampleIS{svakalegur glanni} \dicExampleCS{obrovský vtipálek}  \textbf{2.} \dicDirectTranslationCS{příšerný, děsivý, strašný} \dicExampleIS{svakalegt ástand} \dicExampleCS{příšerný stav}
\dicEntry[svala] \dicTerm{svala\smash{\textsuperscript{1}}} \dicIPA{{s}{v}{a}{\textlengthmark}{l}{a}} \dicPos{f}[1] \dicFlx{(svölu, svölur)}[20] \dicFieldCat{zool.} \dicDirectTranslationCS{vlaštovka} \textit{(l.~{\textLA{Hirundo}})}  \dicsymPhoto\ 
\dicFigure{40689.jpg}{Svala}{Svala - Baum Petr, Biolib, CC-BY-SA}
\dicEntry[svala] \dicTerm{sval|a\smash{\textsuperscript{2}}} \dicIPA{{s}{v}{a}{\textlengthmark}{l}{a}} \dicPos{v}[1] \dicFlx{(‑aði)}[50] \dicFlx{dat} \textbf{1.} \dicPhraseIS{svala sér} \dicSynonym*{kæla sig} \dicDirectTranslationCS{osvěžit se, ochladit se};  \dicPhraseIS{e‑að svalar} \dicDirectTranslationCS{(co) chladí, (co) ochlazuje, (co) osvěžuje} \dicExampleIS{Kalt vatn svalar best.} \dicExampleCS{Chladná voda nejlépe osvěžuje.}  \textbf{2.} \dicPhraseIS{svala e‑u} {\textbf{a.}} \dicSynonym*{létta e‑u} \dicDirectTranslationCS{uhasit (co), utišit (co) (žízeň ap.)} \dicExampleIS{Hún svalaði þorsta sínum.} \dicExampleCS{Uhasila svou žízeň.};  {\textbf{b.}} \dicSynonym*{fullnægja e‑u} \dicDirectTranslationCS{uspokojit (co), nasytit (co) (potřebu ap.)} \dicExampleIS{Hann svalaði forvitni sinni.} \dicExampleCS{Uspokojil svou zvědavost.}
\dicEntry[svaladrykkur] \dicTerm{svala··drykk|ur} \dicIPA{{s}{v}{a}{\textlengthmark}{l}{a}{\textsubring{d}}{r}{\textsci}{h}{\r{g}}{\textscy}{\textsubring{r}}} \dicPos{m}[10] \dicFlx{(‑jar, ‑ir)}[23] \dicDirectTranslationCS{osvěžující nápoj}
\dicEntry[svalandi] \dicTerm{sval··andi} \dicIPA{{s}{v}{a}{\textlengthmark}{l}{a}{n}{\textsubring{d}}{\textsci}} \dicPos{adj}[13] \dicFlx{indecl}[1] \dicDirectTranslationCS{osvěžující, ochlazující}
\dicEntry[Svalbarði] \dicTerm{Sval··barð|i} \dicIPA{{s}{v}{a}{l}{\textsubring{b}}{a}{r}{ð}{\textsci}} \dicPos{m}[1] \dicFlx{(‑a)}[5] \dicFieldCat{geog.} \dicDirectTranslationCS{Špicberky}
\dicEntry[svali] \dicTerm{sval|i} \dicIPA{{s}{v}{a}{\textlengthmark}{l}{\textsci}} \dicPos{m}[1] \dicFlx{(‑a)}[3] \dicSynonym{kæla\smash{\textsuperscript{1}}} \dicDirectTranslationCS{chlad, chladno, chladný vzduch} \dicExampleIS{koma út í svalann} \dicExampleCS{vyjít ven na chladný vzduch}
\dicEntry[svalir] \dicTerm{svalir} \dicIPA{{s}{v}{a}{\textlengthmark}{l}{\textsci}{\textsubring{r}}} \dicPos{f}[12] \dicFlx{pl}[9] \dicDirectTranslationCS{balkón}
\dicEntry[svall] \dicTerm{svall\smash{\textsuperscript{1}}} \dicIPA{{s}{v}{a}{\textsubring{d}}{\textsubring{l}}} \dicPos{n}[2] \dicFlx{(‑s)}[2] \textbf{1.} \dicSynonym{slark} \dicDirectTranslationCS{hýření, flámování}  \textbf{2.} \dicSynonym{eyðslusemi} \dicDirectTranslationCS{hýření, plýtvání, mrhání}
\dicEntry[svall] \dicTerm{svall\smash{\textsuperscript{2}}} \dicIPA{{s}{v}{a}{\textsubring{d}}{\textsubring{l}}} \dicPos{v} \dicFlx{ind pf sg 1 pers} \dicLink{svella\smash{\textsuperscript{2}}}
\dicEntry[svalla] \dicTerm{svall|a} \dicIPA{{s}{v}{a}{\textsubring{d}}{l}{a}} \dicPos{v}[1] \dicFlx{(‑aði)}[13] \dicSynonym{slarka} \dicDirectTranslationCS{flámovat, hýřit}
\dicEntry[svallari] \dicTerm{svall··ar|i} \dicIPA{{s}{v}{a}{\textsubring{d}}{l}{a}{r}{\textsci}} \dicPos{m}[1] \dicFlx{(‑a, ‑ar)}[10] \dicDirectTranslationCS{zpustlík, zpustlice, zhýralec, zhýralkyně}
\dicEntry[svalt] \dicTerm{svalt} \dicIPA{{s}{v}{a}{\textsubring{l}}{\textsubring{d}}} \dicPos{v} \dicFlx{ind pf sg 1 pers} \dicLink{svelta\smash{\textsuperscript{2}}}
\dicEntry[svalur] \dicTerm{svalur} \dicsymFrequent\  \dicIPA{{s}{v}{a}{\textlengthmark}{l}{\textscy}{\textsubring{r}}} \dicPos{adj}[1] \dicFlx{(f svöl)}[2] \textbf{1.} \dicSynonym{kaldur} \dicDirectTranslationCS{chladný, studený} \dicExampleIS{svalur vindur} \dicExampleCS{chladný vítr}  \textbf{2.} \dicSynonym{djarfur} \dicDirectTranslationCS{smělý, chladnokrevný} \dicExampleIS{svalur náungi} \dicExampleCS{smělý mladík}
\dicEntry[svampur] \dicTerm{svamp|ur} \dicIPA{{s}{v}{a}{\textsubring{m}}{\textsubring{b}}{\textscy}{\textsubring{r}}} \dicPos{m}[6] \dicFlx{(‑s, ‑ar)}[5] \textbf{1.} \dicFieldCat{zool.} \dicDirectTranslationCS{houbovec, (živočišná) houba} \dicIndirectTranslationCS{(vodní živočich žijící v~trsech)} \textit{(l.~{\textLA{Porifera}})}  \dicsymPhoto\   \textbf{2.} \dicDirectTranslationCS{houba (na mytí nádobí ap.)}
\dicFigure{ds_image_svampur_0_2.jpg}{Svampur}{Svampur - Karakal, GFDL}
\dicEntry[svangur] \dicTerm{svangur} \dicsymFrequent\  \dicIPA{{s}{v}{au}{\ng}{\r{g}}{\textscy}{\textsubring{r}}} \dicPos{adj}[10] \dicFlx{(f svöng,  comp svengri\,/\addthin svangari, sup svengstur\,/\addthin svangastur)}[5] \dicDirectTranslationCS{hladový} \dicExampleIS{Ertu svöng?} \dicExampleCS{Máš hlad?}
\dicEntry[svanur] \dicTerm{svan|ur} \dicIPA{{s}{v}{a}{\textlengthmark}{n}{\textscy}{\textsubring{r}}} \dicPos{m}[9] \dicFlx{(‑s, ‑ir)}[5] \dicFieldCat{zool.} \dicSynonym{álft} \dicDirectTranslationCS{labuť, labuť zpěvná} \textit{(l.~{\textLA{Cygnus cygnus}})}  \dicsymPhoto\ 
\dicFigure{ds_image_svanur_0_1.jpg}{Svanur}{Svanur - Sunnfili, PD}
\dicEntry[svar] \dicTerm{svar} \dicsymFrequent\  \dicIPA{{s}{v}{a}{\textlengthmark}{\textsubring{r}}} \dicPos{n}[2] \dicFlx{(‑s, svör)}[8] \textbf{1.} \dicDirectTranslationCS{odpověď};  \dicPhraseIS{svar við e‑u} \dicDirectTranslationCS{odpověď na (co)} \dicExampleIS{svar við spurningunni} \dicExampleCS{odpověď na otázku}  \textbf{2.} \dicSynonym{verknaður} \dicDirectTranslationCS{odpověď, odezva, reakce};  \dicIdiom{svar}{ \dicPhraseIS{krefja e‑n svara}} \dicDirectTranslationCS{domáhat se (čí) odpovědi}; { \dicPhraseIS{sitja fyrir svörum hjá e‑m}} \dicDirectTranslationCS{zodpovídat (čí) otázky};  \dicPhraseIS{verða fyrir svörum} \dicDirectTranslationCS{odpovědět, odpovídat};  \dicPhraseIS{virða e‑n ekki svars} \dicDirectTranslationCS{ignorovat (koho)}
\dicEntry[svara] \dicTerm{svar|a} \dicsymFrequent\  \dicIPA{{s}{v}{a}{\textlengthmark}{r}{a}} \dicPos{v}[1] \dicFlx{(‑aði)}[13] \dicFlx{(dat +) dat} \textbf{1.} \dicSynonym{ansa} \dicDirectTranslationCS{odpovědět, odpovídat, odvětit};  \dicPhraseIS{svara e‑m e‑u} \dicDirectTranslationCS{odpovědět (komu co)};  \dicPhraseIS{svara í símann} \dicDirectTranslationCS{vzít telefon};  \dicPhraseIS{svara játandi\,/\addthin neitandi} \dicDirectTranslationCS{odpovědět kladně\,/\addthin záporně}  \textbf{2.} \dicSynonym{endurgjalda} \dicDirectTranslationCS{oplatit, oplácet, odplatit, odplácet};  \dicPhraseIS{svara í sömu mynt} \dicLangCat{přen.} \dicDirectTranslationCS{oplatit stejnou mincí}  \textbf{3.} \dicSynonym{borga} \dicDirectTranslationCS{vyvážit, vyrovnat (náklady ap.)};  \dicPhraseIS{e‑að svarar ekki kostnaði} \dicDirectTranslationCS{(co) nezaplatí náklady};  \dicIdiom{svara}[fyrir]{ \dicPhraseIS{svara fyrir e‑ð}} \dicDirectTranslationCS{zodpovídat se za (co)}; { \dicPhraseIS{svara fyrir sig}} \dicDirectTranslationCS{hájit se, obhajovat se};  \dicIdiom{svara}[til]{ \dicPhraseIS{e‑að svarar til e‑s}} \dicDirectTranslationCS{(co) odpovídá (čemu), (co) s~(čím) koresponduje}
\dicEntry[svaramaður] \dicTerm{svara··|maður} \dicIPA{{s}{v}{a}{\textlengthmark}{r}{a}{m}{a}{ð}{\textscy}{\textsubring{r}}} \dicPos{m}[13] \dicFlx{(‑manns, ‑menn)}[2] \dicDirectTranslationCS{svědek\,/\addthin svědkyně (na svatbě)}
\dicEntry[svarf] \dicTerm{svarf\smash{\textsuperscript{1}}} \dicIPA{{s}{v}{a}{r}{f}} \dicPos{n}[2] \dicFlx{(‑s, svörf)}[8] \textbf{1.} \dicSynonym*{sorfið duft} \dicDirectTranslationCS{brusný prach}  \textbf{2.} \dicSynonym{svörfun} \dicDirectTranslationCS{pilování, obrušování}
\dicEntry[svarf] \dicTerm{svarf\smash{\textsuperscript{2}}} \dicIPA{{s}{v}{a}{r}{f}} \dicPos{v} \dicFlx{ind pf sg 1 pers} \dicLink{sverfa}
\dicEntry[svarið] \dicTerm{svarið} \dicIPA{{s}{v}{a}{\textlengthmark}{r}{\textsci}{\texttheta}} \dicPos{v} \dicFlx{supin} \dicLink{sverja}
\dicEntry[svarinn] \dicTerm{svarinn} \dicIPA{{s}{v}{a}{\textlengthmark}{r}{\textsci}{\textsubring{n}}} \dicPos{adj}[6]\dicFlx{}[-3] \dicDirectTranslationCS{zapřisáhlý};  \dicPhraseIS{e‑ir eru svarnir óvinir} \dicDirectTranslationCS{(kdo) jsou zapřisáhlí nepřátelé}
\dicEntry[svarri] \dicTerm{svarr|i} \dicIPA{{s}{v}{a}{r}{\textlengthmark}{\textsci}} \dicPos{m}[1] \dicFlx{(‑a, ‑ar)}[8] \textbf{1.} \dicSynonym{skass} \dicDirectTranslationCS{dračice, čarodějnice}  \textbf{2.} \dicFieldCat{zool.} \dicDirectTranslationCS{ťuhýk} \textit{(l.~{\textLA{Lanius}})}  \dicsymPhoto\ 
\dicFigure{ds_image_svarri_0_1.jpg}{Svarri}{Svarri - Marek Szczepanek, CC BY-SA 3.0}
\dicEntry[svartagaldur] \dicTerm{svarta··gald|ur} \dicIPA{{s}{v}{a}{\textsubring{r}}{\textsubring{d}}{a}{\r{g}}{a}{l}{\textsubring{d}}{\textscy}{\textsubring{r}}} \dicPos{m}[5] \dicFlx{(‑urs)}[2] \dicLink{svartigaldur}
\dicEntry[Svartahaf] \dicTerm{Svarta··haf} \dicIPA{{s}{v}{a}{\textsubring{r}}{\textsubring{d}}{a}{h}{a}{f}} \dicPos{n}[2] \dicFlx{(‑s)}[2] \dicFieldCat{geog.} \dicDirectTranslationCS{Černé moře}
\dicEntry[svartamarkaður] \dicTerm{svarta··mark·að|ur} \dicIPA{{s}{v}{a}{\textsubring{r}}{\textsubring{d}}{a}{m}{a}{\textsubring{r}}{\r{g}}{a}{ð}{\textscy}{\textsubring{r}}} \dicPos{m}[10] \dicFlx{(‑ar\,/\addthin ‑s, ‑ir)}[31] \dicLink{svartimarkaður}
\dicEntry[svartamyrkur] \dicTerm{svarta··myrkur} \dicIPA{{s}{v}{a}{\textsubring{r}}{\textsubring{d}}{a}{m}{\textsci}{\textsubring{r}}{\r{g}}{\textscy}{\textsubring{r}}} \dicPos{n}[2] \dicFlx{(‑s)}[28] \dicDirectTranslationCS{černočerná tma, tma tmoucí}
\dicEntry[svartaþoka] \dicTerm{svarta··þok|a} \dicIPA{{s}{v}{a}{\textsubring{r}}{\textsubring{d}}{a}{\texttheta}{\textopeno}{\r{g}}{a}} \dicPos{f}[1] \dicFlx{(‑u)}[5] \dicDirectTranslationCS{hustá mlha}
\dicEntry[Svartálfheimur] \dicTerm{Svart··álf··heim|ur} \dicIPA{{s}{v}{a}{\textsubring{r}}{\textsubring{d}}{au}{l}{f}{h}{ei}{m}{\textscy}{\textsubring{r}}} \dicPos{m}[6] \dicFlx{(‑s)}[3] \dicFieldCat{myt.} \dicDirectTranslationCS{Svartálfheim} \dicIndirectTranslationCS{(říše temnoty, temných elfů)} \dicAntonym{Ljósálfheimur}
\dicEntry[svartbakur] \dicTerm{svart··bak|ur} \dicIPA{{s}{v}{a}{\textsubring{r}}{\textsubring{d}}{\textsubring{b}}{a}{\r{g}}{\textscy}{\textsubring{r}}} \dicPos{m}[6] \dicFlx{(‑s, ‑ar)}[25] \dicFieldCat{zool.} \dicDirectTranslationCS{racek mořský} \textit{(l.~{\textLA{Larus marinus}})}  \dicsymPhoto\ 
\dicFigure{ds_image_svartbakur_0_1.jpg}{Svartbakur}{Svartbakur - Tim, CC BY 2.0}
\dicEntry[svartbjörn] \dicTerm{svart··|björn} \dicIPA{{s}{v}{a}{\textsubring{r}}{\textsubring{d}}{\textsubring{b}}{j}{\oe}{r}{\textsubring{d}}{\textsubring{n}}} \dicPos{m}[12] \dicFlx{(‑bjarnar, ‑birnir)}[5] \dicFieldCat{zool.} \dicDirectTranslationCS{medvěd baribal, medvěd černý} \textit{(l.~{\textLA{Ursus americanus}})}
\dicEntry[Svartfellingur] \dicTerm{Svart··fell·ing|ur} \dicIPA{{s}{v}{a}{\textsubring{r}}{\textsubring{d}}{f}{\textepsilon}{\textsubring{d}}{l}{i}{\ng}{\r{g}}{\textscy}{\textsubring{r}}} \dicPos{m}[6] \dicFlx{(‑s, ‑ar)}[8] \dicDirectTranslationCS{Čer\-no\-ho\-rec, Černohorka}
\dicEntry[svartfellskur] \dicTerm{svart··fell·skur} \dicIPA{{s}{v}{a}{\textsubring{r}}{\textsubring{d}}{f}{\textepsilon}{\textsubring{d}}{\textsubring{l}}{s}{\r{g}}{\textscy}{\textsubring{r}}} \dicPos{adj}[1]\dicFlx{}[-6] \dicDirectTranslationCS{černohorský}
\dicEntry[Svartfjallaland] \dicTerm{Svart·fjalla··land} \dicIPA{{s}{v}{a}{\textsubring{r}}{\textsubring{d}}{f}{j}{a}{\textsubring{d}}{l}{a}{l}{a}{n}{\textsubring{d}}} \dicPos{n}[2] \dicFlx{(‑s)}[4] \dicFieldCat{geog.} \dicDirectTranslationCS{Černá Hora}
\dicEntry[svartfugl] \dicTerm{svart··fugl} \dicIPA{{s}{v}{a}{\textsubring{r}}{\textsubring{d}}{f}{\textscy}{\r{g}}{\textsubring{l}}} \dicPos{m}[4] \dicFlx{(‑s, ‑ar)}[3] \dicPhraseIS{svartfuglar} \dicFlx{pl} \dicFieldCat{zool.} \dicDirectTranslationCS{alky} \dicIndirectTranslationCS{(ptáci z~podčeledi alek -- alka, alkoun, alkounek, papuchalk)} \textit{(l.~{\textLA{Alcidae}})}
\dicEntry[svarthol] \dicTerm{svart··hol} \dicIPA{{s}{v}{a}{\textsubring{r}}{\textsubring{d}}{h}{\textopeno}{\textsubring{l}}} \dicPos{n}[2] \dicFlx{(‑s, ‑)}[5] \textbf{1.} \dicSynonym{fangelsi} \dicDirectTranslationCS{(vězeňská) kobka, hladomorna}  \textbf{2.} \dicFieldCat{astro.} \dicDirectTranslationCS{černá díra}
\dicEntry[svarthvítur] \dicTerm{svart··hvítur} \dicIPA{{s}{v}{a}{\textsubring{r}}{\textsubring{d}}{k\smash{\textsuperscript{h}}}{v}{i}{\textsubring{d}}{\textscy}{\textsubring{r}}} \dicPos{adj}[1]\dicFlx{}[-6] \dicDirectTranslationCS{černobílý} \dicExampleIS{Myndin er svarthvít.} \dicExampleCS{Fotka je černobílá.}
\dicEntry[svarthærður] \dicTerm{svart··hærður} \dicIPA{{s}{v}{a}{\textsubring{r}}{\textsubring{d}}{h}{a}{i}{r}{ð}{\textscy}{\textsubring{r}}} \dicPos{adj}[2]\dicFlx{}[-1] \dicDirectTranslationCS{černovlasý}
\dicEntry[svartidauði] \dicTerm{svarti··dauði} \dicIPA{{s}{v}{a}{\textsubring{r}}{\textsubring{d}}{\textsci}{\textsubring{d}}{\oe i}{ð}{\textsci}} \dicPos{m}[1] \dicFlx{(svartadauða)}[7] \textbf{1.} \dicDirectTranslationCS{černá smrt} \dicIndirectTranslationCS{(epidemie moru)}  \textbf{2.} \dicLangCat{hovor.} \dicDirectTranslationCS{černá smrt} \dicIndirectTranslationCS{(islandská kořalka)}
\dicEntry[svartigaldur] \dicTerm{svarti··galdur}\dicTerm{, svartagaldur} \dicIPA{{s}\-{v}\-{a}\-{\textsubring{r}}\-{\textsubring{d}}\-{\textsci}\-{\r{g}}\-{a}\-{l}\-{\textsubring{d}}\-{\textscy}\-{\textsubring{r}}\-} \dicPos{m}[5] \dicFlx{(svartagald\-urs)}[4] \dicDirectTranslationCS{černá magie}

\dicEntry[svartimarkaður] \textls[15]{\dicTerm{svarti··mark·aður}\dicTerm{, svartamarkaður} \dicIPA{{s}\-{v}\-{a}\-{\textsubring{r}}\-{\textsubring{d}}\-{\textsci}\-{m}\-{a}\-{\textsubring{r}}\-{\r{g}}\-{a}\-{ð}\-{\textscy}\-{\textsubring{r}}} \dicPos{m}[10]} \dicFlx{(svartamarkaðar\,/\addthin svartamarkaðs, svörtumarkaðir)}[40] \dicDirectTranslationCS{černý trh}

\dicEntry[svartklæddur] \dicTerm{svart··klæddur} \dicIPA{{s}{v}{a}{\textsubring{r}}{\textsubring{d}}{k\smash{\textsuperscript{h}}}{l}{a}{i}{\textsubring{d}}{\textscy}{\textsubring{r}}} \dicPos{adj}[2]\dicFlx{}[-21] \dicDirectTranslationCS{oblečený v~černém}
\dicEntry[svartkráka] \dicTerm{svart··krák|a} \dicIPA{{s}{v}{a}{\textsubring{r}}{\textsubring{d}}{k\smash{\textsuperscript{h}}}{r}{au}{\r{g}}{a}} \dicPos{f}[1] \dicFlx{(‑u, ‑ur)}[19] \dicFieldCat{zool.} \dicDirectTranslationCS{vrána, vrána černá} \textit{(l.~{\textLA{Corvus corone}})}
\dicEntry[svartlist] \dicTerm{svart··list} \dicIPA{{s}{v}{a}{\textsubring{r}}{\textsubring{d}}{l}{\textsci}{s}{\textsubring{d}}} \dicPos{f}[7] \dicFlx{(‑ar)}[3] \dicDirectTranslationCS{černobílá kresba}
\dicEntry[svartsýni] \dicTerm{svart··sýn|i} \dicIPA{{s}{v}{a}{\textsubring{r}}{\textsubring{d}}{s}{i}{n}{\textsci}} \dicPos{f}[3] \dicFlx{(‑i)}[3] \dicDirectTranslationCS{pesimismus, škarohlídství} \dicAntonym{bjartsýni}
\begin{xtolerant}{}{1pt}
\dicEntry[svartsýnismaður] \dicTerm{svart·sýnis··|maður}\addthinS\textls[-5]{\dicIPA{{s}{v}{a}{\textsubring{r}}{\textsubring{d}}{s}{i}{n}{\textsci}{s}{m}{a}{ð}{\textscy}{\textsubring{r}}}}\addthinM\dicPos{m}[13]\addthinS\dicFlx{(‑manns,\addthinSS‑menn)}[2] \dicDirectTranslationCS{pesimista, pesimistka, škarohlíd(ka)}
\end{xtolerant}
\dicEntry[svartsýnn] \dicTerm{svart··sýnn} \dicIPA{{s}{v}{a}{\textsubring{r}}{\textsubring{d}}{s}{i}{\textsubring{d}}{\textsubring{n}}} \dicPos{adj}[7]\dicFlx{}[-1] \dicSynonym{bölsýnn} \dicDirectTranslationCS{pesimistický, škarohlídský} \dicAntonym{bjartsýnn}
\dicEntry[svartur] \dicTerm{svart|ur\smash{\textsuperscript{1}}} \dicIPA{{s}{v}{a}{\textsubring{r}}{\textsubring{d}}{\textscy}{\textsubring{r}}} \dicPos{m}[6] \dicFlx{(‑s)}[3] \dicDirectTranslationCS{černý} \dicIndirectTranslationCS{(hráč mající černé figurky)} \dicAntonym{hvítur\smash{\textsuperscript{1}}};  \dicPhraseIS{svartur á leik} \dicDirectTranslationCS{černý na tahu}
\dicEntry[svartur] \dicTerm{svartur\smash{\textsuperscript{2}}} \dicsymFrequent\  \dicIPA{{s}{v}{a}{\textsubring{r}}{\textsubring{d}}{\textscy}{\textsubring{r}}} \dicPos{adj}[1] \dicFlx{(f svört)}[11] \textbf{1.} \dicDirectTranslationCS{černý} \dicExampleIS{svart hár} \dicExampleCS{černé vlasy}  \textbf{2.} \dicSynonym{dimmur} \dicDirectTranslationCS{černý, chmurný, temný};  \dicPhraseIS{svört hola} \dicFieldCat{astro.} \dicDirectTranslationCS{černá díra};  \dicPhraseIS{svartur kassi} \dicFieldCat{let.} \dicDirectTranslationCS{černá skříňka};  \dicIdiom{svartur}{ \dicPhraseIS{hafa séð það svartara}} \dicLangCat{přen.} \dicDirectTranslationCS{přežít už i~horší věci}; { \dicPhraseIS{setja e‑n á svartan lista}} \dicDirectTranslationCS{umístit (koho) na černou listinu};  \dicPhraseIS{vera svarti sauðurinn} \dicLangCat{přen.} \dicDirectTranslationCS{být černá ovce};  \dicPhraseIS{vinna svart} \dicDirectTranslationCS{pracovat načerno};  \dicPhraseIS{svart á hvítu} \dicFlx{adv} \dicDirectTranslationCS{černé na bílém}
\dicEntry[svartyllir] \dicTerm{svart··yll|ir} \dicIPA{{s}{v}{a}{\textsubring{r}}{\textsubring{d}}{\textsci}{\textsubring{d}}{l}{\textsci}{\textsubring{r}}} \dicPos{m}[7] \dicFlx{(‑is, ‑ar)}[1] \dicFieldCat{bot.} \dicDirectTranslationCS{bez, bez černý} \textit{(l.~{\textLA{Sambucus nigra}})}
\dicEntry[svartþröstur] \dicTerm{svart··|þröstur} \dicIPA{{s}{v}{a}{\textsubring{r}}{\textsubring{d}}{\texttheta}{r}{\oe}{s}{\textsubring{d}}{\textscy}{\textsubring{r}}} \dicPos{m}[11] \dicFlx{(‑þrastar, ‑þrestir)}[5] \dicFieldCat{zool.} \dicDirectTranslationCS{kos, kos černý} \textit{(l.~{\textLA{Turdus merula}})}  \dicsymPhoto\ 
\dicFigure{26666.jpg}{Svartþröstur}{Svartþröstur - Velechovský Marek, Biolib, Copyright/CC-BY-NC}
\dicEntry[Svasíland] \dicTerm{Svasí··land} \dicIPA{{s}{v}{a}{\textlengthmark}{s}{i}{l}{a}{n}{\textsubring{d}}} \dicPos{n}[2] \dicFlx{(‑s)}[4] \dicFieldCat{geog.} \dicDirectTranslationCS{Svazijsko}
\dicEntry[Svaslendingur] \dicTerm{Svas·lend··ing|ur} \dicIPA{{s}{v}{a}{s}{l}{\textepsilon}{n}{\textsubring{d}}{i}{\ng}{\r{g}}{\textscy}{\textsubring{r}}} \dicPos{m}[6] \dicFlx{(‑s, ‑ar)}[8] \dicDirectTranslationCS{Svazijec, Svazijka}
\dicEntry[svaslenskur] \dicTerm{svas··lenskur} \dicIPA{{s}{v}{a}{s}{l}{\textepsilon}{n}{s}{\r{g}}{\textscy}{\textsubring{r}}} \dicPos{adj}[1]\dicFlx{}[-6] \dicDirectTranslationCS{svazijský}
\dicEntry[sváfum] \dicTerm{sváfum} \dicIPA{{s}{v}{au}{\textlengthmark}{v}{\textscy}{\textsubring{m}}} \dicPos{v} \dicFlx{ind pf pl 1 pers} \dicLink{sofa}
\dicEntry[svefn] \dicTerm{svefn} \dicsymFrequent\  \dicIPA{{s}{v}{\textepsilon}{\textsubring{b}}{\textsubring{n}}} \dicPos{m}[4] \dicFlx{(‑s)}[5] \dicDirectTranslationCS{spánek, spaní} \dicExampleIS{í svefni} \dicExampleCS{ve spánku};  \dicPhraseIS{falla í svefn} \dicDirectTranslationCS{usnout};  \dicPhraseIS{leggjast til svefns} \dicFlx{refl} \dicDirectTranslationCS{uložit se ke spánku};  \dicPhraseIS{tala upp úr svefni} \dicDirectTranslationCS{mluvit ze spaní\,/\addthin snu};  \dicPhraseIS{sofna svefninum langa} \dicLangCat{přen.} \dicDirectTranslationCS{ulehnout k~věčnému spánku} \dicIndirectTranslationCS{(zemřít)};  \dicPhraseIS{vera í fasta svefni} \dicDirectTranslationCS{spát tvrdě}
\dicEntry[svefnfriður] \dicTerm{svefn··frið|ur} \dicIPA{{s}{v}{\textepsilon}{\textsubring{b}}{\textsubring{n}}{f}{r}{\textsci}{ð}{\textscy}{\textsubring{r}}} \dicPos{m}[10] \dicFlx{(‑ar)}[7] \dicDirectTranslationCS{noční klid, klid na spaní} \dicExampleIS{hafa ekki svefnfrið fyrir hávaða} \dicExampleCS{nemít kvůli hluku klid na spaní}
\dicEntry[svefnherbergi] \dicTerm{svefn··her·bergi} \dicsymFrequent\  \dicIPA{{s}{v}{\textepsilon}{\textsubring{b}}{\textsubring{n}}{h}{\textepsilon}{r}{\textsubring{b}}{\textepsilon}{r}{\r{\textObardotlessj}}{\textsci}} \dicPos{n}[2] \dicFlx{(‑s, ‑)}[16] \dicDirectTranslationCS{ložnice} \dicExampleIS{í báðum svefnherbergjum} \dicExampleCS{v~obou ložnicích}
\dicEntry[svefnklefi] \dicTerm{svefn··klef|i} \dicIPA{{s}{v}{\textepsilon}{\textsubring{b}}{\textsubring{n}}{k\smash{\textsuperscript{h}}}{l}{\textepsilon}{v}{\textsci}} \dicPos{m}[1] \dicFlx{(‑a, ‑ar)}[1] \dicDirectTranslationCS{spací kupé}
\dicEntry[svefnlaus] \dicTerm{svefn··laus} \dicIPA{{s}{v}{\textepsilon}{\textsubring{b}}{n}{l}{\oe i}{s}} \dicPos{adj}[5]\dicFlx{}[-1] \dicDirectTranslationCS{bezesný} \dicIndirectTranslationCS{(jsoucí bez spánku)}
\dicEntry[svefnleysi] \dicTerm{svefn··leysi} \dicIPA{{s}{v}{\textepsilon}{\textsubring{b}}{n}{l}{ei}{s}{\textsci}} \dicPos{n}[2] \dicFlx{(‑s)}[20] \dicDirectTranslationCS{nespavost, insomnie}
\dicEntry[svefnléttur] \dicTerm{svefn··léttur} \dicIPA{{s}{v}{\textepsilon}{\textsubring{b}}{n}{l}{j}{\textepsilon}{h}{\textsubring{d}}{\textscy}{\textsubring{r}}} \dicPos{adj}[1]\dicFlx{}[-10] \dicDirectTranslationCS{potřebující málo spánku} \dicExampleIS{Hún er svefnlétt.} \dicExampleCS{Nepotřebuje mnoho spánku.}
\dicEntry[svefnlyf] \dicTerm{svefn··lyf} \dicIPA{{s}{v}{\textepsilon}{\textsubring{b}}{n}{l}{\textsci}{f}} \dicPos{n}[2] \dicFlx{(‑s, ‑)}[13] \dicDirectTranslationCS{prášek na spaní}
\dicEntry[svefnpokapláss] \dicTerm{svefn·poka··pláss} \dicIPA{{s}{v}{\textepsilon}{\textsubring{b}}{\textsubring{n}}{p\smash{\textsuperscript{h}}}{\textopeno}{\r{g}}{a}{p\smash{\textsuperscript{h}}}{l}{au}{s}} \dicPos{n}[2] \dicFlx{(‑, ‑)}[22] \dicDirectTranslationCS{ubytování s~vlastními spacáky}
\dicEntry[svefnpoki] \dicTerm{svefn··pok|i} \dicIPA{{s}{v}{\textepsilon}{\textsubring{b}}{\textsubring{n}}{p\smash{\textsuperscript{h}}}{\textopeno}{\r{\textObardotlessj}}{\textsci}} \dicPos{m}[1] \dicFlx{(‑a, ‑ar)}[1] \dicDirectTranslationCS{spacák, spací pytel}
\dicEntry[svefnpurka] \dicTerm{svefn··purk|a} \dicIPA{{s}{v}{\textepsilon}{\textsubring{b}}{\textsubring{n}}{p\smash{\textsuperscript{h}}}{\textscy}{\textsubring{r}}{\r{g}}{a}} \dicPos{f}[1] \dicFlx{(‑u, ‑ur)}[19] \dicDirectTranslationCS{spáč(ka), ospalec, ospalkyně}
\dicEntry[svefnrof] \dicTerm{svefn··rof} \dicIPA{{s}{v}{\textepsilon}{\textsubring{b}}{n}{r}{\textopeno}{f}} \dicPos{n}[2] \dicFlx{pl}[1] \dicDirectTranslationCS{polospánek} \dicExampleIS{í svefnrofunum} \dicExampleCS{v~polospánku}
\dicEntry[svefnsófi] \dicTerm{svefn··sóf|i} \dicIPA{{s}{v}{\textepsilon}{\textsubring{b}}{\textsubring{n}}{s}{ou}{f}{\textsci}} \dicPos{m}[1] \dicFlx{(‑a, ‑ar)}[1] \dicDirectTranslationCS{rozkládací pohov\-ka\,/\addthin gauč}
\dicEntry[svefnsýki] \dicTerm{svefn··sýk|i} \dicIPA{{s}{v}{\textepsilon}{\textsubring{b}}{\textsubring{n}}{s}{i}{\r{\textObardotlessj}}{\textsci}} \dicPos{f}[3] \dicFlx{(‑i)}[3] \dicFieldCat{med. } \dicDirectTranslationCS{spavá nemoc, africká trypanosomiáza}
\dicEntry[svefnvagn] \dicTerm{svefn··vagn} \dicIPA{{s}{v}{\textepsilon}{\textsubring{b}}{n}{v}{a}{\r{g}}{\textsubring{n}}} \dicPos{m}[4] \dicFlx{(‑s, ‑ar)}[4] \dicDirectTranslationCS{spací\,/\addthin lůžkový vůz}
\dicEntry[svefnþungur] \dicTerm{svefn··|þungur} \dicIPA{{s}{v}{\textepsilon}{\textsubring{b}}{\textsubring{n}}{\texttheta}{u}{\ng}{\r{g}}{\textscy}{\textsubring{r}}} \dicPos{adj}[10] \dicFlx{(comp ‑þyngri, sup ‑þyng\-st\-ur)}[7] \dicDirectTranslationCS{tvrdě spící}
\dicEntry[svei] \dicTerm{svei} \dicsymFrequent\  \dicIPA{{s}{v}{ei}{\textlengthmark}} \dicPos{inter} \dicDirectTranslationCS{sakra, kruci};  \dicPhraseIS{svei attan!} \dicDirectTranslationCS{fuj!};  \dicPhraseIS{svei mér þá!} \dicDirectTranslationCS{to mě (tedy) podrž!};  \dicPhraseIS{svei mér} \dicFlx{adv} \dicDirectTranslationCS{fakt, opravdicky} \dicExampleIS{Þetta er svei mér gott!} \dicExampleCS{To je fakt super!};  \dicPhraseIS{svei mér ekki} \dicDirectTranslationCS{fakt\,/\addthin opravdicky ne};  \dicPhraseIS{svei þér!} \dicDirectTranslationCS{fuj!, styď se!}
\dicEntry[sveið] \dicTerm{sveið} \dicIPA{{s}{v}{ei}{\textlengthmark}{\texttheta}} \dicPos{v} \dicFlx{ind pf sg 1 pers} \dicLink{svíða}
\dicEntry[sveif] \dicTerm{sveif\smash{\textsuperscript{1}}} \dicIPA{{s}{v}{ei}{\textlengthmark}{f}} \dicPos{f}[4] \dicFlx{(‑ar, ‑ar)}[1] \textbf{1.} \dicDirectTranslationCS{klika (otočná)}  \textbf{2.} \dicFieldCat{anat.} \dicDirectTranslationCS{vřetenní kost};  \dicIdiom{sveif}{ \dicPhraseIS{hallast\,/\addthin snúast á sveif með e‑m}} \dicFlx{refl} \dicLangCat{přen.} \dicDirectTranslationCS{přiklánět se na (čí) stranu}
\dicEntry[sveif] \dicTerm{sveif\smash{\textsuperscript{2}}} \dicIPA{{s}{v}{ei}{\textlengthmark}{f}} \dicPos{v} \dicFlx{ind pf sg 1 pers} \dicLink{svífa}
\dicEntry[sveifarás] \dicTerm{sveifar··ás} \dicIPA{{s}{v}{ei}{\textlengthmark}{v}{a}{r}{au}{s}} \dicPos{m}[4] \dicFlx{(‑s, ‑ar)}[25] \dicFieldCat{techn.} \dicDirectTranslationCS{kliková hřídel}
\dicEntry[sveifarhús] \dicTerm{sveifar··hús} \dicIPA{{s}{v}{ei}{\textlengthmark}{v}{a}{\textsubring{r}}{h}{u}{s}} \dicPos{n}[2] \dicFlx{(‑s, ‑)}[5] \dicFieldCat{techn.} \dicDirectTranslationCS{kliková skříň}
\dicEntry[sveifla] \dicTerm{sveifl|a\smash{\textsuperscript{1}}} \dicsymFrequent\  \dicIPA{{s}{v}{ei}{\textsubring{b}}{l}{a}} \dicPos{f}[1] \dicFlx{(‑u, ‑ur)}[19] \textbf{1.} \dicDirectTranslationCS{kývání, houpání}  \textbf{2.} \dicSynonym*{flökt} \dicDirectTranslationCS{výkyv, kolísání, fluktuace} \dicExampleIS{dagleg sveifla hlutabréfa á markaði} \dicExampleCS{každodenní kolísání akcií na trhu}  \textbf{3.} \dicFieldCat{fyz.} \dicDirectTranslationCS{kmitání, oscilace}  \textbf{4.} \dicFieldCat{hud.} \dicDirectTranslationCS{improvizace, improvizování} \dicIndirectTranslationCS{(volné rozvíjení určitého tématu)}
\dicEntry[sveifla] \dicTerm{sveifl|a\smash{\textsuperscript{2}}} \dicsymFrequent\  \dicIPA{{s}{v}{ei}{\textsubring{b}}{l}{a}} \dicPos{v}[1] \dicFlx{(‑aði)}[1] \dicFlx{dat} \dicSynonym{veifa\smash{\textsuperscript{1}}} \dicDirectTranslationCS{(roz)kývat, (roz)houpat} \dicExampleIS{sveifla handleggjunum} \dicExampleCS{kývat rukama};  \dicIdiom{sveifla}[yfir]{ \dicPhraseIS{sveifla sér yfir e‑ð}} \dicDirectTranslationCS{přeskočit (co) (zhoupnutím), přehoupnout se přes (co)} \dicExampleIS{sveifla sér yfir girðinguna} \dicExampleCS{přehoupnout se přes plot};  \dicIdiom{sveiflast}{ \dicPhraseIS{sveiflast}} \dicFlx{refl} {\textbf{a.}} \dicSynonym{dingla} \dicDirectTranslationCS{kývat se, houpat se};  {\textbf{b.}} \dicDirectTranslationCS{kolísat, fluktuovat, oscilovat} \dicExampleIS{sveiflast milli vonar og ótta} \dicExampleCS{kolísat mezi nadějí a~strachem}
\dicEntry[sveifluhljóð] \dicTerm{sveiflu··hljóð} \dicIPA{{s}{v}{ei}{\textsubring{b}}{l}{\textscy}{\textsubring{l}}{j}{ou}{\texttheta}} \dicPos{n}[2] \dicFlx{(‑s, ‑)}[5] \dicFieldCat{jaz.} \dicDirectTranslationCS{kmitavá souhláska, vibranta}
\dicEntry[sveiflurit] \dicTerm{sveiflu··rit} \dicIPA{{s}{v}{ei}{\textsubring{b}}{l}{\textscy}{r}{\textsci}{\textsubring{d}}} \dicPos{n}[2] \dicFlx{(‑s, ‑)}[5] \dicFieldCat{fyz.} \dicDirectTranslationCS{oscilogram}
\dicEntry[sveifluriti] \dicTerm{sveiflu··rit|i} \dicIPA{{s}{v}{ei}{\textsubring{b}}{l}{\textscy}{r}{\textsci}{\textsubring{d}}{\textsci}} \dicPos{m}[1] \dicFlx{(‑a, ‑ar)}[1] \dicFieldCat{fyz.} \dicDirectTranslationCS{oscilograf}
\dicEntry[sveiflusjá] \dicTerm{sveiflu··sjá} \dicIPA{{s}{v}{ei}{\textsubring{b}}{l}{\textscy}{s}{j}{au}} \dicPos{f}[4] \dicFlx{(‑r, ‑r)}[18] \dicFieldCat{fyz.} \dicDirectTranslationCS{osciloskop}
\dicEntry[sveiflutíðni] \dicTerm{sveiflu··tíðn|i} \dicIPA{{s}{v}{ei}{\textsubring{b}}{l}{\textscy}{t\smash{\textsuperscript{h}}}{i}{ð}{n}{\textsci}} \dicPos{f}[3] \dicFlx{(‑i)}[3] \dicFieldCat{fyz.} \dicDirectTranslationCS{oscilační frekvence}
\dicEntry[sveigja] \dicTerm{sveigj|a\smash{\textsuperscript{1}}} \dicIPA{{s}{v}{ei}{j}{\textlengthmark}{a}} \dicPos{f}[1] \dicFlx{(‑u, ‑ur)}[7] \textbf{1.} \dicSynonym{bugða} \dicDirectTranslationCS{oblouk, ohyb, (kruhové) zakřivení}  \textbf{2.} \dicSynonym{sveigjanleiki} \dicDirectTranslationCS{ohebnost, elasticita, flexibilita}
\dicEntry[sveigja] \dicTerm{sveig|ja\smash{\textsuperscript{2}}} \dicsymFrequent\  \dicIPA{{s}{v}{ei}{j}{\textlengthmark}{a}} \dicPos{v}[2] \dicFlx{(‑ði, ‑t)}[89] \dicFlx{acc} \textbf{1.} \dicSynonym{beygja\smash{\textsuperscript{2}}} \dicDirectTranslationCS{ohnout, ohýbat, (za)křivit} \dicExampleIS{sveigja líkamann til hliðar} \dicExampleCS{ohnout tělo na stranu}  \textbf{2.} \dicSynonym*{fara til hliðar} \dicDirectTranslationCS{zahnout, zahýbat, zabočit, zabočovat} \dicExampleIS{sveigja  til vinstri\,/\addthin hægri} \dicExampleCS{zabočit doleva\,/\addthin doprava};  \dicPhraseIS{sveigja af leið} \dicDirectTranslationCS{uhnout\,/\addthin odbočit z~cesty};  \dicIdiom{sveigjast}{ \dicPhraseIS{sveigjast}} \dicFlx{refl} {\textbf{a.}} \dicSynonym{bogna} \dicDirectTranslationCS{ohnout se, ohýbat se};  {\textbf{b.}} \dicDirectTranslationCS{stočit se, stáčet se (rozhovor ap.)}
\dicEntry[sveigjanlegur] \dicTerm{sveigjan··legur} \dicIPA{{s}{v}{ei}{j}{\textlengthmark}{a}{n}{l}{\textepsilon}{\textbabygamma}{\textscy}{\textsubring{r}}} \dicPos{adj}[1]\dicFlx{}[-8] \textbf{1.} \dicSynonym{beygjanlegur} \textls[10]{\dicDirectTranslationCS{ohebný, elastický, pružný} \dicExampleIS{sveigjanlegur vinnutími} \dicExampleCS{pružná pracovní doba}  \textbf{2.} \dicDirectTranslationCS{přizpůsobivý, flexibilní, pružný (člověk ap.)}}
\dicEntry[sveigjanleiki] \dicTerm{sveigjan··leik|i} \dicIPA{{s}{v}{ei}{j}{\textlengthmark}{a}{n}{l}{ei}{\r{\textObardotlessj}}{\textsci}} \dicPos{m}[1] \dicFlx{(‑a)}[3] \textbf{1.} \dicDirectTranslationCS{ohebnost, elasticita, pružnost}  \textbf{2.} \dicDirectTranslationCS{přizpůsobivost, flexibilita, pružnost}
\dicEntry[sveigur] \dicTerm{sveig|ur} \dicIPA{{s}{v}{ei}{\textlengthmark}{\textbabygamma}{\textscy}{\textsubring{r}}} \dicPos{m}[6] \dicFlx{(‑s\,/\addthin ‑ar, ‑ar)}[65] \textbf{1.} \dicSynonym{bugða} \dicDirectTranslationCS{oblouk, ohyb, zakřivení};  \dicPhraseIS{taka sveig á leið sinni} \dicDirectTranslationCS{vzít to obloukem, zajet si};  \dicPhraseIS{taka á sig sveig til að (gera e‑ð)} \dicDirectTranslationCS{širokým obloukem se vyhnout (dělání (čeho)), zdaleka se vyhnout (dělání (čeho))}  \textbf{2.} \dicSynonym{krans} \dicDirectTranslationCS{věnec, girlanda} \dicExampleIS{blómsveigur} \dicExampleCS{věnec květin}
\dicEntry[sveik] \dicTerm{sveik} \dicIPA{{s}{v}{ei}{\textlengthmark}{\r{g}}} \dicPos{v} \dicFlx{ind pf sg 1 pers} \dicLink{svíkja}
\dicEntry[sveima] \dicTerm{sveim|a} \dicIPA{{s}{v}{ei}{\textlengthmark}{m}{a}} \dicPos{v}[1] \dicFlx{(‑aði)}[1] \textbf{1.} \dicSynonym{svífa} \dicDirectTranslationCS{vznášet se, poletovat} \dicExampleIS{Fuglahópur sveimaði yfir höfðum okkar.} \dicExampleCS{Hejno ptáků se vznášelo nad našimi hlavami.}  \textbf{2.} \dicSynonym{ráfa} \dicDirectTranslationCS{chodit kolem\,/\addthin okolo, toulat se, potulovat se}
\dicEntry[sveimhugi] \dicTerm{sveim··hug|i} \dicIPA{{s}{v}{ei}{\textlengthmark}{\textsubring{m}}{h}{\textscy}{i}{j}{\textsci}} \dicPos{m}[1] \dicFlx{(‑a, ‑ar)}[1] \dicSynonym{draumóramaður} \dicDirectTranslationCS{snílek, idealista, idealistka}
\dicEntry[sveinbarn] \dicTerm{svein··|barn} \dicIPA{{s}{v}{ei}{n}{\textsubring{b}}{a}{r}{\textsubring{d}}{\textsubring{n}}} \dicPos{n}[2] \dicFlx{(‑barns, ‑börn)}[8] \dicDirectTranslationCS{hošík, klouček, chlapeček} \dicAntonym{meybarn}
\dicEntry[sveinn] \dicTerm{svein|n} \dicIPA{{s}{v}{ei}{\textsubring{d}}{\textsubring{n}}} \dicPos{m}[6] \dicFlx{(‑s, ‑ar)}[42] \textbf{1.} \dicSynonym{drengur} \dicDirectTranslationCS{klouček, hošík}  \textbf{2.} \dicSynonym{piltur} \dicDirectTranslationCS{chlapec, mladík, mládenec}  \textbf{3.} \dicFieldCat{sport.} \dicDirectTranslationCS{junior(ka) (14--16 let)} \dicExampleIS{sveinakeppni} \dicExampleCS{závod juniorů}  \textbf{4.} \dicPhraseIS{hreinn sveinn} \dicDirectTranslationCS{panic}  \textbf{5.} \dicDirectTranslationCS{učeň, učnice}
\dicEntry[sveinspróf] \dicTerm{sveins··próf} \dicIPA{{s}{v}{ei}{n}{s}{p\smash{\textsuperscript{h}}}{r}{ou}{f}} \dicPos{n}[2] \dicFlx{(‑s, ‑)}[5] \dicFieldCat{škol.} \dicDirectTranslationCS{závěrečné učňovské zkoušky}
\dicEntry[sveipa] \dicTerm{sveip|a} \dicIPA{{s}{v}{ei}{\textlengthmark}{\textsubring{b}}{a}} \dicPos{v}[1] \dicFlx{(‑aði)}[1] \dicFlx{acc (+ dat)} \dicDirectTranslationCS{ovinout, ovíjet, zavinout, zavíjet, (za)balit};  \dicPhraseIS{sveipa e‑n e‑u, sveipa e‑n í e‑ð, sveipa e‑u um e‑n} \dicDirectTranslationCS{ovinout (koho čím), zabalit (koho) do (čeho)} \dicExampleIS{sveipa hann hlýju teppi} \dicExampleCS{zavinout ho do teplého přehozu};  \dicIdiom{sveipast}{ \dicPhraseIS{sveipast e‑u}} \dicFlx{refl} \dicDirectTranslationCS{být zahalen (čím) (mlhou ap.)}
\dicEntry[sveipur] \dicTerm{sveip|ur} \dicIPA{{s}{v}{ei}{\textlengthmark}{\textsubring{b}}{\textscy}{\textsubring{r}}} \dicPos{m}[9] \dicFlx{(‑s, ‑ir)}[8] \textbf{1.} \dicSynonym*{liður í hári} \dicDirectTranslationCS{kudrlina, kučera, kadeř}  \textbf{2.} \dicSynonym{gustur} \dicDirectTranslationCS{poryv, závan (větru ap.)}
\dicEntry[sveit] \dicTerm{sveit} \dicsymFrequent\  \dicIPA{{s}{v}{ei}{\textlengthmark}{\textsubring{d}}} \dicPos{f}[7] \dicFlx{(‑ar, ‑ir)}[1] \textbf{1.} \dicSynonym{dreifbýli} \dicDirectTranslationCS{venkov, vesnice} \dicExampleIS{fara í sveit til afa síns} \dicExampleCS{jet na venkov k~dědečkovi}  \textbf{2.} \dicSynonym{hreppur} \dicDirectTranslationCS{okres, region, provincie}  \textbf{3.} \dicSynonym{flokkur} \dicDirectTranslationCS{četa, družstvo, oddíl, skupina}  \textbf{4.} \dicLangCat{hovor.} \dicSynonym{hljómsveit} \dicDirectTranslationCS{(hudební) skupina, kapela}
\dicEntry[sveitabær] \dicTerm{sveita··bæ|r} \dicIPA{{s}{v}{ei}{\textlengthmark}{\textsubring{d}}{a}{\textsubring{b}}{a}{i}{\textsubring{r}}} \dicPos{m}[10] \dicFlx{(‑jar, ‑ir)}[25] \dicDirectTranslationCS{hospodářství, farma}
\dicEntry[sveitafólk] \dicTerm{sveita··fólk} \dicIPA{{s}{v}{ei}{\textlengthmark}{\textsubring{d}}{a}{f}{ou}{\textsubring{l}}{\r{g}}} \dicPos{n}[2] \dicFlx{(‑s)}[2] \dicDirectTranslationCS{venkované, vesničané}
\dicEntry[sveitalegur] \dicTerm{sveita··legur} \dicIPA{{s}{v}{ei}{\textlengthmark}{\textsubring{d}}{a}{l}{\textepsilon}{\textbabygamma}{\textscy}{\textsubring{r}}} \dicPos{adj}[1]\dicFlx{}[-8] \dicDirectTranslationCS{venkovský, vesnický}
\dicEntry[sveitalubbi] \dicTerm{sveita··lubb|i} \dicIPA{{s}{v}{ei}{\textlengthmark}{\textsubring{d}}{a}{l}{\textscy}{\textsubring{b}}{\textsci}} \dicPos{m}[1] \dicFlx{(‑a, ‑ar)}[1] \dicLangCat{han.} \dicDirectTranslationCS{křupan(ka), buran(ka), balík, vidlák}
\dicEntry[sveitamaður] \textls[10]{\dicTerm{sveita··|maður} \dicIPA{{s}{v}{ei}{\textlengthmark}{\textsubring{d}}{a}{m}{a}{ð}{\textscy}{\textsubring{r}}} \dicPos{m}[13] \dicFlx{(‑manns, ‑menn)}[2]} \textbf{1.} \dicSynonym{bóndi} \dicDirectTranslationCS{farmář(ka), hospodář(ka)}  \textbf{2.} \dicSynonym{sveitalubbi} \dicDirectTranslationCS{venkovan(ka), křupan(ka)}
\dicEntry[sveitamenning] \dicTerm{sveita··menn·ing} \dicIPA{{s}{v}{ei}{\textlengthmark}{\textsubring{d}}{a}{m}{\textepsilon}{n}{i}{\ng}{\r{g}}} \dicPos{f}[4] \dicFlx{(‑ar)}[7] \dicDirectTranslationCS{kultura venkova}
\dicEntry[sveitarfélag] \dicTerm{sveitar··fé·|lag} \dicIPA{{s}{v}{ei}{\textsubring{d}}{a}{\textsubring{r}}{f}{j}{\textepsilon}{l}{a}{x}} \dicPos{n}[2] \dicFlx{(‑lags, ‑lög)}[8] \dicSynonym{hreppur} \dicDirectTranslationCS{obec, komuna}
\dicEntry[sveitarstjóri] \dicTerm{sveitar··stjór|i} \dicIPA{{s}{v}{ei}{\textsubring{d}}{a}{\textsubring{r}}{s}{\textsubring{d}}{j}{ou}{r}{\textsci}} \dicPos{m}[1] \dicFlx{(‑a, ‑ar)}[1] \dicDirectTranslationCS{starosta\,/\addthin starostka obce}
\dicEntry[sveitarstjórn] \dicTerm{sveitar··stjórn} \dicIPA{{s}{v}{ei}{\textsubring{d}}{a}{\textsubring{r}}{s}{\textsubring{d}}{j}{ou}{r}{\textsubring{d}}{\textsubring{n}}} \dicPos{f}[7] \dicFlx{(‑ar, ‑ir)}[1] \dicDirectTranslationCS{obecní zastupitelstvo}
\dicEntry[sveitarstjórnarkosningar] \dicTerm{sveitar·stjórnar··kos·ningar} \textls[-10]{\dicIPA{{s}\-{v}\-{ei}\-{\textsubring{d}}\-{a}\-{\textsubring{r}}\-{s}\-{\textsubring{d}}\-{j}\-{ou}\-{\textsubring{r}}\-{\textsubring{d}}\-{n}\-{a}\-{\textsubring{r}}\-{k\smash{\textsuperscript{h}}}\-{\textopeno}\-{s}\-{\textsubring{d}}\-{n}\-{i}\-{\ng}\-{\r{g}}\-{a}\-{\textsubring{r}}}} \dicPos{f}[4] \dicFlx{pl}[6] \dicDirectTranslationCS{volby do obecního zastupitelstva}
\begin{xtolerant}{}{1pt}
\dicEntry[sveitarstjórnarmaður] \dicTerm{sveitar·stjórnar··|maður} \dicIPA{{s}\-{v}\-{ei}\-{\textsubring{d}}\-{a}\-{\textsubring{r}}\-{s}\-{\textsubring{d}}\-{j}\-{ou}\-{r}\-{\textsubring{d}}\-{n}\-{a}\-{r}\-{m}\-{a}\-{ð}\-{\textscy}\-{\textsubring{r}}\-}\addthinS\dicPos{m}[13] \dicFlx{(‑manns, ‑menn)}[2] \dicDirectTranslationCS{člen(ka) obecního zastupitelstva}
\end{xtolerant}
\dicEntry[sveitasæla] \dicTerm{sveita··sæl|a} \dicIPA{{s}{v}{ei}{\textlengthmark}{\textsubring{d}}{a}{s}{a}{i}{l}{a}} \dicPos{f}[1] \dicFlx{(‑u)}[5] \dicDirectTranslationCS{venkovská idyla}
\dicEntry[sveitaþorp] \dicTerm{sveita··þorp} \dicIPA{{s}{v}{ei}{\textlengthmark}{\textsubring{d}}{a}{\texttheta}{\textopeno}{\textsubring{r}}{\textsubring{b}}} \dicPos{n}[2] \dicFlx{(‑s, ‑)}[5] \dicDirectTranslationCS{vesnice, ves}
\dicEntry[sveiti] \dicTerm{sveit|i} \dicIPA{{s}{v}{ei}{\textlengthmark}{\textsubring{d}}{\textsci}} \dicPos{m}[1] \dicFlx{(‑a)}[3] \dicPhraseIS{í sveita síns andlitis} \dicFlx{adv} \dicLangCat{přen.} \dicDirectTranslationCS{v~potu své tváře}
\dicEntry[sveitó] \dicTerm{sveitó} \dicIPA{{s}{v}{ei}{\textlengthmark}{\textsubring{d}}{ou}} \dicPos{adj}[13] \dicFlx{indecl}[1] \dicLangCat{hovor.} \dicSynonym{sveitalegur} \dicDirectTranslationCS{venkovský, vesnický}
\dicEntry[sveittur] \dicTerm{sveittur} \dicIPA{{s}{v}{ei}{h}{\textsubring{d}}{\textscy}{\textsubring{r}}} \dicPos{adj}[1]\dicFlx{}[-10] \dicDirectTranslationCS{zpocený, upocený};  \dicPhraseIS{sitja með sveittan skallann yfir e‑u} \dicLangCat{přen.} \dicDirectTranslationCS{potit se nad (čím)}
\dicEntry[sveitungi] \dicTerm{sveit··ung|i} \dicIPA{{s}{v}{ei}{\textlengthmark}{\textsubring{d}}{u}{\textltailn}{\r{\textObardotlessj}}{\textsci}} \dicPos{m}[1] \dicFlx{(‑a, ‑ar)}[1] \dicDirectTranslationCS{krajan(ka)}
\dicEntry[svekkja] \dicTerm{svekk|ja} \dicIPA{{s}{v}{\textepsilon}{h}{\r{\textObardotlessj}}{a}} \dicPos{v}[2] \dicFlx{(‑ti, ‑t)}[25] \dicFlx{acc} \dicSynonym{stríða\smash{\textsuperscript{2}}} \dicDirectTranslationCS{otrávit, otravovat, (po)zlobit};  \dicIdiom{svekkja}[á]{ \dicPhraseIS{svekkja sig á e‑u}} \dicDirectTranslationCS{být zklamaný z~(čeho)}
\dicEntry[svekktur] \textls[10]{\dicTerm{svekktur} \dicIPA{{s}{v}{\textepsilon}{x}{\textsubring{d}}{\textscy}{\textsubring{r}}} \dicPos{adj}[1]\dicFlx{}[-10] \dicSynonym{mæddur}} \dicDirectTranslationCS{rozčarovaný, zklamaný} \dicExampleIS{svekktur yfir e‑u} \dicExampleCS{rozčarovaný z~(čeho)}
\dicEntry[svelgja] \dicTerm{svelg|ja} \dicIPA{{s}{v}{\textepsilon}{l}{\r{\textObardotlessj}}{a}} \dicPos{v}[2] \dicFlx{(‑di, ‑t)}[133] \dicFlx{acc} \textbf{1.} \dicSynonym*{háma í sig} \dicDirectTranslationCS{zhltnout, hltat} \dicExampleIS{svelgja matinn} \dicExampleCS{hltat jídlo}  \textbf{2.} \dicSynonym*{drekka í stórum teygum} \dicDirectTranslationCS{(vy)chlemtat};  \dicIdiom{svelgjast}[á]{ \dicPhraseIS{e‑m svelgist á e‑u}} \dicFlx{refl impers} \dicDirectTranslationCS{(kdo) se zakucká (čím)} \dicExampleIS{Mér svelgdist á kaffinu.} \dicExampleCS{Zakuckal jsem se kafem.}
\dicEntry[svelgur] \dicTerm{svelg|ur} \dicIPA{{s}{v}{\textepsilon}{l}{\r{g}}{\textscy}{\textsubring{r}}} \dicPos{m}[9] \dicFlx{(‑s, ‑ir)}[15] \textbf{1.} \dicSynonym{hringiða} \dicDirectTranslationCS{vír, víření (vody)}  \textbf{2.} \dicSynonym{átvagl} \dicDirectTranslationCS{žrout(ka), nenasyta, hltoun}
\dicEntry[svell] \dicTerm{svell\smash{\textsuperscript{1}}} \dicIPA{{s}{v}{\textepsilon}{\textsubring{d}}{\textsubring{l}}} \dicPos{n}[2] \dicFlx{(‑s, ‑)}[5] \textbf{1.} \dicSynonym*{íslag} \dicDirectTranslationCS{ledovka, (povrchový) led};  \dicPhraseIS{sterkur\,/\addthin fastur á svellinu} \dicFlx{adj} \dicSynonym{traustur} \dicDirectTranslationCS{solidní, jistý}  \textbf{2.} \dicSynonym{skautasvell} \dicDirectTranslationCS{kluziště}
\dicEntry[svell] \dicTerm{svell\smash{\textsuperscript{2}}} \dicIPA{{s}{v}{\textepsilon}{\textsubring{d}}{\textsubring{l}}} \dicPos{v} \dicFlx{ind praes sg 1 pers} \dicLink{svella\smash{\textsuperscript{2}}}
\dicEntry[svella] \dicTerm{svell|a\smash{\textsuperscript{1}}} \dicIPA{{s}{v}{\textepsilon}{\textsubring{d}}{l}{a}} \dicPos{v}[1] \dicFlx{(‑aði)}[34] \dicDirectTranslationCS{(z)ledovatět, pokrývat se ledem}
\dicEntry[svella] \dicTerm{svella\smash{\textsuperscript{2}}} \dicIPA{{s}{v}{\textepsilon}{\textsubring{d}}{l}{a}} \dicPos{v}[6] \dicFlx{(svell, svall, sullum, sylli, sollið)}[36] \textbf{1.} \dicSynonym{ólga\smash{\textsuperscript{1}}} \dicDirectTranslationCS{vzdouvat se, vřít, kypět} \dicExampleIS{Blóðið svellur í æðum mér.} \dicExampleCS{Krev se mi vzdouvá v~žilách.}  \textbf{2.} \dicPhraseIS{e‑m svellur móður} \dicFlx{impers} \dicDirectTranslationCS{(kdo) je plný entuziasmu}
\dicEntry[svellaður] \dicTerm{svell··|aður} \dicIPA{{s}{v}{\textepsilon}{\textsubring{d}}{l}{a}{ð}{\textscy}{\textsubring{r}}} \dicPos{adj}[3] \dicFlx{(f ‑uð)}[4] \dicDirectTranslationCS{zledovatělý, namrzlý} \dicExampleIS{svellaðar gangstéttir} \dicExampleCS{zledovatělé chodníky}
\dicEntry[svellbunki] \dicTerm{svell··bunk|i} \dicIPA{{s}{v}{\textepsilon}{\textsubring{d}}{\textsubring{l}}{\textsubring{b}}{u}{\r{\textltailn}}{\r{\textObardotlessj}}{\textsci}} \dicPos{m}[1] \dicFlx{(‑a, ‑ar)}[1] \dicDirectTranslationCS{vypouklá a~hladká ledová plocha}
\dicEntry[svelt] \dicTerm{svelt} \dicIPA{{s}{v}{\textepsilon}{\textsubring{l}}{\textsubring{d}}} \dicPos{v} \dicFlx{ind praes sg 1 pers} \dicLink{svelta\smash{\textsuperscript{2}}}
\dicEntry[svelta] \dicTerm{svelt|a\smash{\textsuperscript{1}}} \dicIPA{{s}{v}{\textepsilon}{\textsubring{l}}{\textsubring{d}}{a}} \dicPos{v}[2] \dicFlx{(‑i, ‑)}[16] \dicFlx{acc} \dicDirectTranslationCS{nechat (vy)hladovět, (u)mořit hladem} \dicExampleIS{svelta hann til bana} \dicExampleCS{nechat ho vyhladovět k~smrti}
\dicEntry[svelta] \dicTerm{svelta\smash{\textsuperscript{2}}} \dicIPA{{s}{v}{\textepsilon}{\textsubring{l}}{\textsubring{d}}{a}} \dicPos{v}[6] \dicFlx{(svelt, svalt, sultum, sylti, soltið)}[34] \dicDirectTranslationCS{hladovět, trpět hladem};  \dicPhraseIS{svelta heilu hungri} \dicDirectTranslationCS{hladovět, mít hlad, umírat hlady}
\dicEntry[svelti] \dicTerm{svelti} \dicIPA{{s}{v}{\textepsilon}{\textsubring{l}}{\textsubring{d}}{\textsci}} \dicPos{n}[2] \dicFlx{(‑s, ‑)}[14] \dicDirectTranslationCS{hladovění};  \dicPhraseIS{vera í svelti} \dicLangCat{přen.} \dicDirectTranslationCS{být o~hladu, držet dietu}
\dicEntry[svengja] \dicTerm{sveng|ja} \dicIPA{{s}{v}{ei}{\textltailn}{\r{\textObardotlessj}}{a}} \dicPos{v}[2] \dicFlx{(‑di, ‑t)}[140] \dicPhraseIS{e‑n svengir} \dicFlx{impers} \dicDirectTranslationCS{(kdo) má hlad, (kdo) hladoví, (kdo) trpí hladem}
\dicEntry[sveppafræði] \dicTerm{sveppa··fræð|i} \dicIPA{{s}{v}{\textepsilon}{h}{\textsubring{b}}{a}{f}{r}{a}{i}{ð}{\textsci}} \dicPos{f}[3] \dicFlx{(‑i)}[3] \dicDirectTranslationCS{mykologie}
\begin{xtolerant}{}{1pt}
\dicEntry[sveppasósa] \dicTerm{sveppa··sós|a}\addthinS\dicIPA{{s}{v}{\textepsilon}{h}{\textsubring{b}}{a}{s}{ou}{s}{a}}\addthinS\dicPos{f}[1]\addthinS\dicFlx{(‑u, ‑ur)}[7]\addthinS\dicDirectTranslationCS{houbová omáčka}
\end{xtolerant}
\dicEntry[sveppur] \dicTerm{svepp|ur} \dicIPA{{s}{v}{\textepsilon}{h}{\textsubring{b}}{\textscy}{\textsubring{r}}} \dicPos{m}[6] \dicFlx{(‑s, ‑ir\,/\addthin ‑ar)}[73] \dicFieldCat{bot.} \dicDirectTranslationCS{houba} \textit{(l.~{\textLA{Fungi}})}
\dicEntry[sver] \dicTerm{sver\smash{\textsuperscript{1}}} \dicIPA{{s}{v}{\textepsilon}{\textlengthmark}{\textsubring{r}}} \dicPos{adj}[5] \dicFlx{(f ‑)}[8] \textbf{1.} \dicSynonym{digur} \dicDirectTranslationCS{statný, robustní} \dicExampleIS{sver maður} \dicExampleCS{robustní člověk}  \textbf{2.} \dicDirectTranslationCS{tlustý (kabel ap.)}
\dicEntry[sver] \dicTerm{sver\smash{\textsuperscript{2}}} \dicIPA{{s}{v}{\textepsilon}{\textlengthmark}{\textsubring{r}}} \dicPos{v} \dicFlx{ind praes sg 1 pers} \dicLink{sverja}
\dicEntry[sverð] \dicTerm{sverð} \dicsymFrequent\  \dicIPA{{s}{v}{\textepsilon}{r}{\texttheta}} \dicPos{n}[2] \dicFlx{(‑s, ‑)}[5] \dicDirectTranslationCS{meč} \dicExampleIS{mismunandi hlutar sverðs} \dicExampleCS{různé části meče};  \dicPhraseIS{tvíeggjað sverð} \dicLangCat{přen.} \dicDirectTranslationCS{dvousečný meč}
\dicEntry[sverðfiskur] \dicTerm{sverð··fisk|ur} \dicIPA{{s}{v}{\textepsilon}{r}{ð}{f}{\textsci}{s}{\r{g}}{\textscy}{\textsubring{r}}} \dicPos{m}[6] \dicFlx{(‑s, ‑ar)}[10] \dicFieldCat{zool.} \dicDirectTranslationCS{mečoun, mečoun obecný} \textit{(l.~{\textLA{Xiphias gladius}})}
\dicEntry[sverðköttur] \dicTerm{sverð··|köttur} \dicIPA{{s}{v}{\textepsilon}{r}{ð}{k\smash{\textsuperscript{h}}}{\oe}{h}{\textsubring{d}}{\textscy}{\textsubring{r}}} \dicPos{m}[11] \dicFlx{(‑kattar, ‑kettir)}[5] \dicFieldCat{zool.} \dicDirectTranslationCS{smilodon, tygr šavlozubý} \textit{(l.~{\textLA{Smilodon fatalis}})}
\dicEntry[sverðlilja] \dicTerm{sverð··lilj|a} \dicIPA{{s}{v}{\textepsilon}{r}{ð}{l}{\textsci}{l}{j}{a}} \dicPos{f}[1] \dicFlx{(‑u, ‑ur)}[7] \dicFieldCat{bot.} \dicDirectTranslationCS{kosatec, iris} \textit{(l.~{\textLA{Iris}})}  \dicsymPhoto\ 
\dicFigure{ds_image_sverdlilja_0_1.jpg}{Sverðlilja}{Sverðlilja - A. Barra, CC BY-SA 3.0}
\dicEntry[sverf] \dicTerm{sverf} \dicIPA{{s}{v}{\textepsilon}{r}{f}} \dicPos{v} \dicFlx{ind praes sg 1 pers} \dicLink{sverfa}
\dicEntry[sverfa] \dicTerm{sverfa} \dicIPA{{s}{v}{\textepsilon}{r}{v}{a}} \dicPos{v}[6] \dicFlx{(sverf, svarf, surfum, syrfi, sorfið)}[35] \dicFlx{acc} \dicDirectTranslationCS{pilovat, rašplovat} \dicExampleIS{sverfa neglur sínar} \dicExampleCS{pilovat si nehty};  \dicIdiom{sverfa}[að]{ \dicPhraseIS{e‑að sverfur að e‑m}} \dicDirectTranslationCS{(co) na (koho) doléhá, (co koho) drtí} \dicExampleIS{Hungrið sverfur að honum.} \dicExampleCS{Doléhá na něj hlad.};  \dicIdiom{sverfa}{ \dicPhraseIS{láta sverfa til stáls}} \dicLangCat{přen.} \dicDirectTranslationCS{tlačit na pilu}
\dicEntry[sverja] \dicTerm{sverja} \dicsymFrequent\  \dicIPA{{s}{v}{\textepsilon}{r}{j}{a}} \dicPos{v}[6] \dicFlx{(sver, sór, sórum, sværi, svarið)}[121] \dicFlx{(gen +) acc} \dicDirectTranslationCS{přísahat};  \dicPhraseIS{e‑r sór og sárt við lagði að} \dicDirectTranslationCS{(kdo) se dušoval\,/\addthin zaklínal, že};  \dicPhraseIS{ég get svarið það} \dicLangCat{hovor.} \dicDirectTranslationCS{můžu to odpřisáhnout};  \dicIdiom{sverjast}{ \dicPhraseIS{sverjast}} \dicFlx{refl} \dicDirectTranslationCS{přísahat si (navzájem)}
\dicEntry[sverta] \dicTerm{svert|a\smash{\textsuperscript{1}}} \dicIPA{{s}{v}{\textepsilon}{\textsubring{r}}{\textsubring{d}}{a}} \dicPos{f}[1] \dicFlx{(‑u, ‑ur)}[19] \dicDirectTranslationCS{čerň, černá barva} \dicExampleIS{prentsverta} \dicExampleCS{tiskařská čerň}
\dicEntry[sverta] \dicTerm{svert|a\smash{\textsuperscript{2}}} \dicIPA{{s}{v}{\textepsilon}{\textsubring{r}}{\textsubring{d}}{a}} \dicPos{v}[2] \dicFlx{(‑i, ‑)}[12] \dicFlx{acc} \textbf{1.} \dicSynonym*{gera svartan} \dicDirectTranslationCS{načernit, načerňovat, začernit, začerňovat, (na)malovat načerno}  \textbf{2.} \dicSynonym*{kasta rýrð á} \dicDirectTranslationCS{očernit, očerňovat, pošpinit, pošpiňovat} \dicExampleIS{sverta e‑n í augum almennings} \dicExampleCS{očernit (koho) v~očích veřejnosti}
\dicEntry[svertingi] \dicTerm{svert··ing|i} \dicIPA{{s}{v}{\textepsilon}{\textsubring{r}}{\textsubring{d}}{i}{\textltailn}{\r{\textObardotlessj}}{\textsci}} \dicPos{m}[1] \dicFlx{(‑ja, ‑jar)}[14] \dicSynonym{blökkumaður} \dicDirectTranslationCS{čer\-ný, černoch, černoška}
\dicEntry[sveskja] \dicTerm{sveskj|a} \dicIPA{{s}{v}{\textepsilon}{s}{\r{\textObardotlessj}}{a}} \dicPos{f}[1] \dicFlx{(‑u, ‑ur)}[25] \dicDirectTranslationCS{sušená švestka}
\dicEntry[svið] \dicTerm{svið} \dicsymFrequent\  \dicIPA{{s}{v}{\textsci}{\textlengthmark}{\texttheta}} \dicPos{n}[2] \dicFlx{(‑s, ‑)}[5] \textbf{1.} \dicSynonym{umdæmi} \dicDirectTranslationCS{pole, obor, doména, sféra} \dicExampleIS{á sviði umhverfismála} \dicExampleCS{na poli životního prostředí}  \textbf{2.} \dicSynonym{leiksvið} \dicDirectTranslationCS{jeviště, scéna}  \textbf{3.} \dicSynonym{fiskimið} \dicDirectTranslationCS{loviště ryb}  \textbf{4.} \dicFieldCat{fyz.} \dicDirectTranslationCS{pole}  \textbf{5.} \dicPhraseIS{svið} \dicFlx{pl} \dicFieldCat{kulin.} \dicDirectTranslationCS{vařené ovčí hlavy}
\dicEntry[sviðalykt] \dicTerm{sviða··lykt} \dicIPA{{s}{v}{\textsci}{\textlengthmark}{ð}{a}{l}{\textsci}{x}{\textsubring{d}}} \dicPos{f}[7] \dicFlx{(‑ar)}[3] \dicSynonym*{brunalykt} \dicDirectTranslationCS{zápach dýmu (z~požáru), zápach požáru}
\dicEntry[sviði] \dicTerm{svið|i\smash{\textsuperscript{1}}} \dicIPA{{s}{v}{\textsci}{\textlengthmark}{ð}{\textsci}} \dicPos{m}[1] \dicFlx{(‑a)}[3] \dicDirectTranslationCS{popálení, spálení} \dicExampleIS{Hann finnur til sviða í andlitinu.} \dicExampleCS{Bolí ho popálenina na tváři.}
\dicEntry[sviði] \dicTerm{sviði\smash{\textsuperscript{2}}} \dicIPA{{s}{v}{\textsci}{\textlengthmark}{ð}{\textsci}} \dicPos{v} \dicFlx{con pf sg 1 pers} \dicLink{svíða}
\dicEntry[sviðið] \dicTerm{sviðið} \dicIPA{{s}{v}{\textsci}{\textlengthmark}{ð}{\textsci}{\texttheta}} \dicPos{v} \dicFlx{supin} \dicLink{svíða}
\dicEntry[sviðinn] \dicTerm{sviðinn} \dicIPA{{s}{v}{\textsci}{\textlengthmark}{ð}{\textsci}{\textsubring{n}}} \dicPos{adj}[6]\dicFlx{}[-2] \dicDirectTranslationCS{popálený, připálený, spálený}
\dicEntry[sviðna] \dicTerm{sviðn|a} \dicIPA{{s}{v}{\textsci}{ð}{n}{a}} \dicPos{v}[1] \dicFlx{(‑aði)}[1] \dicSynonym*{verða sviðinn} \dicDirectTranslationCS{popálit se, připálit se, spálit se} \dicExampleIS{Hár hans sviðnaði.} \dicExampleCS{Jeho vlasy se připálily.}
\dicEntry[sviðset] \dicTerm{svið··set} \dicIPA{{s}{v}{\textsci}{ð}{s}{\textepsilon}{\textsubring{d}}} \dicPos{v} \dicFlx{ind praes sg 1 pers} \dicLink{sviðsetja}
\dicEntry[sviðsetja] \dicTerm{svið··|setja} \dicIPA{{s}{v}{\textsci}{ð}{s}{\textepsilon}{\textsubring{d}}{j}{a}} \dicPos{v}[4] \dicFlx{(‑set, ‑setti, ‑settum, ‑setti, ‑sett)}[26] \dicFlx{acc} \textbf{1.} \dicDirectTranslationCS{(z)inscenovat, uvést\,/\addthin uvádět na scéně}  \textbf{2.} \dicDirectTranslationCS{(z)rekonstruovat (událost ap.)}
\dicEntry[sviðsetning] \dicTerm{svið··set·ning} \dicIPA{{s}{v}{\textsci}{ð}{s}{\textepsilon}{h}{\textsubring{d}}{n}{i}{\ng}{\r{g}}} \dicPos{f}[4] \dicFlx{(‑ar)}[7] \textbf{1.} \dicSynonym*{það að sviðsetja} \dicDirectTranslationCS{inscenování, uvedení na scénu}  \textbf{2.} \dicSynonym*{endursköpun atburðar} \dicDirectTranslationCS{rekonstrukce (událostí ap.)} \dicIndirectTranslationCS{(opětovné provedení něčeho)}
\dicEntry[sviðsett] \dicTerm{svið··sett} \dicIPA{{s}{v}{\textsci}{ð}{s}{\textepsilon}{h}{\textsubring{d}}} \dicPos{v} \dicFlx{supin} \dicLink{sviðsetja}
\dicEntry[sviðsetti] \dicTerm{svið··setti} \dicIPA{{s}{v}{\textsci}{ð}{s}{\textepsilon}{h}{\textsubring{d}}{\textsci}} \dicPos{v} \dicFlx{ind\,/\addthin con pf sg 1 pers} \dicLink{sviðsetja}
\dicEntry[sviðsettum] \dicTerm{svið··settum} \dicIPA{{s}{v}{\textsci}{ð}{s}{\textepsilon}{h}{\textsubring{d}}{\textscy}{\textsubring{m}}} \dicPos{v} \dicFlx{ind pf pl 1 pers} \dicLink{sviðsetja}
\dicEntry[sviðsljós] \dicTerm{sviðs··ljós} \dicIPA{{s}{v}{\textsci}{ð}{s}{l}{j}{ou}{s}} \dicPos{n}[2] \dicFlx{(‑s, ‑)}[5] \dicDirectTranslationCS{(divadelní) reflektor\,/\addthin světlomet, jevištní světlo};  \dicPhraseIS{vera í sviðsljósinu} \dicLangCat{přen.} \dicDirectTranslationCS{být v~záři reflektorů}
\dicEntry[sviðsskrekkur] \dicTerm{sviðs··skrekk|ur} \dicIPA{{s}{v}{\textsci}{ð}{s}{\r{g}}{r}{\textepsilon}{h}{\r{g}}{\textscy}{\textsubring{r}}} \dicPos{m}[6] \dicFlx{(‑s)}[17] \dicDirectTranslationCS{tréma}
\dicEntry[sviðum] \dicTerm{sviðum} \dicIPA{{s}{v}{\textsci}{\textlengthmark}{ð}{\textscy}{\textsubring{m}}} \dicPos{v} \dicFlx{ind pf pl 1 pers} \dicLink{svíða}
\dicEntry[svif] \dicTerm{svif} \dicIPA{{s}{v}{\textsci}{\textlengthmark}{f}} \dicPos{n}[2] \dicFlx{(‑s, ‑)}[5] \dicFieldCat{biol.} \dicDirectTranslationCS{plankton};  \dicPhraseIS{í sömu svifum} \dicFlx{adv} \dicDirectTranslationCS{v~tu chvíli, hned}
\dicEntry[svifaseinn] \dicTerm{svifa··seinn} \dicIPA{{s}{v}{\textsci}{\textlengthmark}{v}{a}{s}{ei}{\textsubring{d}}{\textsubring{n}}} \dicPos{adj}[7]\dicFlx{}[-1] \dicSynonym{stirðbusalegur} \dicDirectTranslationCS{zpomalený, zvolněný, neohrabaný}
\dicEntry[svifdreki] \dicTerm{svif··drek|i} \dicIPA{{s}{v}{\textsci}{v}{\textsubring{d}}{r}{\textepsilon}{\r{\textObardotlessj}}{\textsci}} \dicPos{m}[1] \dicFlx{(‑a, ‑ar)}[1] \dicFieldCat{let.} \dicDirectTranslationCS{rogalo}
\dicEntry[svifflug] \dicTerm{svif··flug} \dicIPA{{s}{v}{\textsci}{f}{l}{\textscy}{x}} \dicPos{n}[2] \dicFlx{(‑s)}[2] \dicFieldCat{let.} \dicDirectTranslationCS{plachtění, klouzavý let}
\dicEntry[sviffluga] \dicTerm{svif··flug|a} \dicIPA{{s}{v}{\textsci}{f}{l}{\textscy}{\textbabygamma}{a}} \dicPos{f}[1] \dicFlx{(‑u, ‑ur)}[13] \dicFieldCat{let.} \dicDirectTranslationCS{(závěsný) kluzák, větroň}
\dicEntry[svifi] \dicTerm{svifi} \dicIPA{{s}{v}{\textsci}{\textlengthmark}{v}{\textsci}} \dicPos{v} \dicFlx{con pf sg 1 pers} \dicLink{svífa}
\dicEntry[svifið] \dicTerm{svifið} \dicIPA{{s}{v}{\textsci}{\textlengthmark}{v}{\textsci}{\texttheta}} \dicPos{v} \dicFlx{supin} \dicLink{svífa}
\dicEntry[svifnökkvi] \dicTerm{svif··nökkv|i} \dicIPA{{s}{v}{\textsci}{v}{n}{\oe}{h}{\r{g}}{v}{\textsci}} \dicPos{m}[1] \dicFlx{(‑a, ‑ar)}[1] \dicFieldCat{let.} \dicDirectTranslationCS{vznášedlo}
\dicEntry[svifrá] \dicTerm{svif··rá} \dicIPA{{s}{v}{\textsci}{v}{r}{au}} \dicPos{f}[4] \dicFlx{(‑r, ‑r)}[18] \dicFieldCat{sport.} \dicDirectTranslationCS{hrazda}
\dicEntry[svifum] \dicTerm{svifum} \dicIPA{{s}{v}{\textsci}{\textlengthmark}{v}{\textscy}{\textsubring{m}}} \dicPos{v} \dicFlx{ind pf pl 1 pers} \dicLink{svífa}
\dicEntry[svig] \dicTerm{svig} \dicIPA{{s}{v}{\textsci}{\textlengthmark}{x}} \dicPos{n}[2] \dicFlx{(‑s)}[2] \textbf{1.} \dicSynonym{sveigur} \dicDirectTranslationCS{zatáčka, zákruta};  \dicPhraseIS{fara á svig við e‑ð} \dicDirectTranslationCS{obejít\,/\addthin obcházet (co) (zákon ap.)}  \textbf{2.} \dicFieldCat{sport.} \dicDirectTranslationCS{slalom}
\dicEntry[svigi] \dicTerm{svig|i} \dicIPA{{s}{v}{i}{j}{\textlengthmark}{\textsci}} \dicPos{m}[1] \dicFlx{(‑a, ‑ar)}[1] \dicFieldCat{jaz.} \dicDirectTranslationCS{(kulatá) závorka} \dicIndirectTranslationCS{(interpunkční znaménko)};  \dicPhraseIS{innan sviga} \dicFlx{adv} \dicDirectTranslationCS{v~závorkách}
\dicEntry[svigkeppni] \dicTerm{svig··keppn|i} \dicIPA{{s}{v}{\textsci}{\textbabygamma}{c\smash{\textsuperscript{h}}}{\textepsilon}{h}{\textsubring{b}}{n}{\textsci}} \dicPos{f}[3] \dicFlx{(‑i, ‑ir)}[1] \dicFieldCat{sport.} \dicDirectTranslationCS{slalom}
\dicEntry[svigmaður] \dicTerm{svig··|maður} \dicIPA{{s}{v}{\textsci}{\textbabygamma}{m}{a}{ð}{\textscy}{\textsubring{r}}} \dicPos{m}[13] \dicFlx{(‑manns, ‑menn)}[2] \dicDirectTranslationCS{slalomář(ka)}
\dicEntry[svigna] \dicTerm{svign|a} \dicIPA{{s}{v}{\textsci}{\r{g}}{n}{a}} \dicPos{v}[1] \dicFlx{(‑aði)}[7] \textbf{1.} \dicSynonym{bogna} \dicDirectTranslationCS{ohnout se, ohýbat se, prohnout se, prohýbat se} \dicExampleIS{Greinarnar svignuðu undan snjóþunganum.} \dicExampleCS{Větve se prohnuly pod tíhou sněhu.}  \textbf{2.} \dicPhraseIS{það svignar} \dicFlx{impers} \dicSynonym*{það hlýnar} \dicDirectTranslationCS{otepluje se}
\dicEntry[svigrúm] \dicTerm{svig··rúm} \dicIPA{{s}{v}{\textsci}{\textbabygamma}{r}{u}{\textsubring{m}}} \dicPos{n}[2] \dicFlx{(‑s)}[2] \textbf{1.} \dicSynonym*{tími til að snúast við} \dicDirectTranslationCS{(časová) rezerva, (časový\,/\addthin manévrovací) prostor}  \textbf{2.} \dicSynonym*{rými til að hreyfa sig} \dicDirectTranslationCS{mezera, vůle}  \textbf{3.} \dicFieldCat{fyz.} \dicDirectTranslationCS{orbital}
\dicEntry[svigskíði] \dicTerm{svig··skíði} \dicIPA{{s}{v}{\textsci}{\textbabygamma}{s}{\r{\textObardotlessj}}{i}{ð}{\textsci}} \dicPos{n}[2] \dicFlx{pl}[19] \dicDirectTranslationCS{slalomové lyže}
\dicEntry[svik] \dicTerm{svik} \dicIPA{{s}{v}{\textsci}{\textlengthmark}{\r{g}}} \dicPos{n}[2] \dicFlx{pl}[1] \textbf{1.} \dicSynonym{fals} \dicDirectTranslationCS{podvod, podvádění}  \textbf{2.} \dicSynonym*{það að rjúfa orð sín} \dicDirectTranslationCS{podraz, zrada}
\dicEntry[svikahrappur] \dicTerm{svika··hrapp|ur} \dicIPA{{s}{v}{\textsci}{\textlengthmark}{\r{g}}{a}{\textsubring{r}}{a}{h}{\textsubring{b}}{\textscy}{\textsubring{r}}} \dicPos{m}[6] \dicFlx{(‑s, ‑ar)}[5] \dicSynonym{svikari} \dicDirectTranslationCS{podvodník, podvodnice, podfukář(ka), šejdíř(ka)}
\dicEntry[svikalaust] \dicTerm{svika··laust} \dicIPA{{s}{v}{\textsci}{\textlengthmark}{\r{g}}{a}{l}{\oe i}{s}{\textsubring{d}}} \dicPos{adv} \dicDirectTranslationCS{svědomitě, pečlivě} \dicExampleIS{vinna svikalaust} \dicExampleCS{pracovat svědomitě}
\dicEntry[svikari] \dicTerm{svik··ar|i} \dicIPA{{s}{v}{\textsci}{\textlengthmark}{\r{g}}{a}{r}{\textsci}} \dicPos{m}[1] \dicFlx{(‑a, ‑ar)}[13] \textbf{1.} \dicSynonym{loddari} \dicDirectTranslationCS{podvodník, podvodnice, podfukář(ka)}  \textbf{2.} \dicDirectTranslationCS{podrazák, podrazačka, zrádce, zrádkyně}
\dicEntry[sviki] \dicTerm{sviki} \dicIPA{{s}{v}{\textsci}{\textlengthmark}{\r{\textObardotlessj}}{\textsci}} \dicPos{v} \dicFlx{con pf sg 1 pers} \dicLink{svíkja}
\dicEntry[svikið] \dicTerm{svikið} \dicIPA{{s}{v}{\textsci}{\textlengthmark}{\r{\textObardotlessj}}{\textsci}{\texttheta}} \dicPos{v} \dicFlx{supin} \dicLink{svíkja}
\dicEntry[svikinn] \dicTerm{svikinn} \dicIPA{{s}{v}{\textsci}{\textlengthmark}{\r{\textObardotlessj}}{\textsci}{\textsubring{n}}} \dicPos{adj}[6]\dicFlx{}[-2] \dicSynonym{falsaður} \dicDirectTranslationCS{falešný, padělaný} \dicExampleIS{svikin vara} \dicExampleCS{falešné zboží};  \dicPhraseIS{verða illa svikinn} \dicDirectTranslationCS{být velmi překvapený, velmi se divit}
\dicEntry[svikóttur] \dicTerm{svik··óttur} \dicIPA{{s}{v}{\textsci}{\textlengthmark}{\r{g}}{ou}{h}{\textsubring{d}}{\textscy}{\textsubring{r}}} \dicPos{adj}[1]\dicFlx{}[-10] \dicSynonym*{sviksamur} \dicDirectTranslationCS{proradný, věrolomný}
\dicEntry[svikráð] \dicTerm{svik··ráð} \dicIPA{{s}{v}{\textsci}{\textlengthmark}{\r{g}}{r}{au}{\texttheta}} \dicPos{n}[2] \dicFlx{pl}[1] \dicDirectTranslationCS{úklady, úskoky, nástrahy};  \dicPhraseIS{sitja á svikráðum við e‑n} \dicDirectTranslationCS{strojit úklady proti (komu)}
\dicEntry[sviksamlegur] \dicTerm{svik·sam··legur} \dicIPA{{s}{v}{\textsci}{\textlengthmark}{\r{g}}{s}{a}{m}{l}{\textepsilon}{\textbabygamma}{\textscy}{\textsubring{r}}} \dicPos{adj}[1]\dicFlx{}[-8] \dicDirectTranslationCS{podvodný, proradný, zrádný}
\dicEntry[sviksemi] \dicTerm{svik··sem|i} \dicIPA{{s}{v}{\textsci}{\textlengthmark}{\r{g}}{s}{\textepsilon}{m}{\textsci}} \dicPos{f}[2] \dicFlx{(‑i)}[2] \dicDirectTranslationCS{proradnost, zrádnost}
\dicEntry[svikull] \dicTerm{svikull} \dicIPA{{s}{v}{\textsci}{\textlengthmark}{\r{g}}{\textscy}{\textsubring{d}}{\textsubring{l}}} \dicPos{adj}[8]\dicFlx{}[-4] \textbf{1.} \dicSynonym{ótrúr} \dicDirectTranslationCS{proradný, věrolomný}  \textbf{2.} \dicSynonym{ótraustur} \dicDirectTranslationCS{zrádný, ošidný (led ap.)} \dicExampleIS{svikull ís} \dicExampleCS{zrádný led}
\dicEntry[svikum] \dicTerm{svikum} \dicIPA{{s}{v}{\textsci}{\textlengthmark}{\r{g}}{\textscy}{\textsubring{m}}} \dicPos{v} \dicFlx{ind pf pl 1 pers} \dicLink{svíkja}
\dicEntry[svil] \dicTerm{svil} \dicIPA{{s}{v}{\textsci}{\textlengthmark}{\textsubring{l}}} \dicPos{n}[2] \dicFlx{pl}[7] \dicDirectTranslationCS{mlíčí}
\dicEntry[svili] \dicTerm{svil|i} \dicIPA{{s}{v}{\textsci}{\textlengthmark}{l}{\textsci}} \dicPos{m}[1] \dicFlx{(‑a, ‑ar)}[1] \dicDirectTranslationCS{švagr} \dicIndirectTranslationCS{(manžel sestry nebo sestry druhého z~manželů)} \dicAntonym{mágur}
\dicEntry[svilkona] \dicTerm{svil··kon|a} \dicIPA{{s}{v}{\textsci}{l}{k\smash{\textsuperscript{h}}}{\textopeno}{n}{a}} \dicPos{f}[1] \dicFlx{(‑u, ‑ur)}[27] \dicDirectTranslationCS{švagrová} \dicIndirectTranslationCS{(manželka bratra nebo bratra druhého z~manželů)} \dicAntonym{mágkona}
\dicEntry[svima] \dicTerm{svim|a} \dicIPA{{s}{v}{\textsci}{\textlengthmark}{m}{a}} \dicPos{v}[1] \dicFlx{(‑aði)}[66] \dicFlx{impers} \dicPhraseIS{e‑n svimar} \dicDirectTranslationCS{(kdo) má závratě, (komu) se točí hlava} \dicExampleIS{Mig svimar.} \dicExampleCS{Točí se mi hlava.}
\dicEntry[svimandi] \dicTerm{svim··andi} \dicIPA{{s}{v}{\textsci}{\textlengthmark}{m}{a}{n}{\textsubring{d}}{\textsci}} \dicPos{adj}[13] \dicFlx{indecl}[1] \dicDirectTranslationCS{závratný, nesmírný}
\dicEntry[svimhár] \dicTerm{svim··|hár} \dicIPA{{s}{v}{\textsci}{\textlengthmark}{\textsubring{m}}{h}{au}{\textsubring{r}}} \dicPos{adj}[10] \dicFlx{(comp ‑hærri, sup ‑hæstur)}[14] \dicDirectTranslationCS{závratný, závratně vysoký}
\dicEntry[svimi] \dicTerm{svim|i}\dicTerm{, svími} \dicIPA{{s}\-{v}\-{\textsci}\-{\textlengthmark}\-{m}\-{\textsci}\-} \dicPos{m}[1] \dicFlx{(‑a)}[3] \dicSynonym*{sundl} \dicDirectTranslationCS{závrať, motání hlavy};  \dicPhraseIS{vaða í villu og svima um e‑ð} \dicLangCat{přen.} \dicDirectTranslationCS{tápat v~(čem)}
\dicEntry[svindilbrask] \dicTerm{svindil··brask} \dicIPA{{s}{v}{\textsci}{n}{\textsubring{d}}{\textsci}{l}{\textsubring{b}}{r}{a}{s}{\r{g}}} \dicPos{n}[2] \dicFlx{(‑s)}[2] \dicDirectTranslationCS{podvod, podfuk, švindl}
\dicEntry[svindl] \dicTerm{svindl} \dicIPA{{s}{v}{\textsci}{n}{\textsubring{d}}{\textsubring{l}}} \dicPos{n}[2] \dicFlx{(‑s)}[2] \dicDirectTranslationCS{podvod, podfuk, švindl} \dicExampleIS{vera með svindl} \dicExampleCS{podvádět}
\dicEntry[svindla] \dicTerm{svindl|a} \dicIPA{{s}{v}{\textsci}{n}{\textsubring{d}}{l}{a}} \dicPos{v}[1] \dicFlx{(‑aði)}[1] \dicDirectTranslationCS{švindlovat, podvést, podvádět, (o)šidit} \dicExampleIS{svindla á prófi} \dicExampleCS{podvádět u~zkoušky};  \dicPhraseIS{svindla á e‑n} \dicDirectTranslationCS{podvést (koho), ošidit (koho)}
\dicEntry[svindlari] \dicTerm{svindl··ar|i} \dicIPA{{s}{v}{\textsci}{n}{\textsubring{d}}{l}{a}{r}{\textsci}} \dicPos{m}[1] \dicFlx{(‑a, ‑ar)}[13] \dicSynonym{braskari} \dicDirectTranslationCS{švindlíř(ka), podvodník, podvodnice}
\dicEntry[svipa] \dicTerm{svip|a\smash{\textsuperscript{1}}} \dicIPA{{s}{v}{\textsci}{\textlengthmark}{\textsubring{b}}{a}} \dicPos{f}[1] \dicFlx{(‑u, ‑ur)}[19] \textbf{1.} \dicDirectTranslationCS{bič}  \textbf{2.} \dicFieldCat{biol.} \dicDirectTranslationCS{bičík}
\dicEntry[svipa] \dicTerm{svip|a\smash{\textsuperscript{2}}} \dicsymFrequent\  \dicIPA{{s}{v}{\textsci}{\textlengthmark}{\textsubring{b}}{a}} \dicPos{v}[1] \dicFlx{(‑aði)}[51] \dicSynonym{þjóta} \dicDirectTranslationCS{(pro)svištět};  \dicIdiom{svipa}[til]{ \dicPhraseIS{e‑m svipar til e‑rs}} \dicFlx{impers} \dicDirectTranslationCS{(kdo) se (komu) podobá} \dicExampleIS{Henni svipar til móður sinnar.} \dicExampleCS{Podobá se své matce.};  \dicIdiom{svipast}[um]{ \dicPhraseIS{svipast um}} \dicFlx{refl} \dicDirectTranslationCS{porozhlédnout se, rozhlédnout se kolem} \dicExampleIS{svipast um í borginni} \dicExampleCS{porozhlédnout se po městě}; { \dicPhraseIS{svipast um eftir e‑u}} \dicFlx{refl} \dicDirectTranslationCS{porozhlédnout se po (čem)} \dicIndirectTranslationCS{(hledat)}
\dicEntry[svipaður] \dicTerm{svip··|aður} \dicsymFrequent\  \dicIPA{{s}{v}{\textsci}{\textlengthmark}{\textsubring{b}}{a}{ð}{\textscy}{\textsubring{r}}} \dicPos{adj}[3] \dicFlx{(f ‑uð)}[3] \dicDirectTranslationCS{podobný} \dicExampleIS{svipuð áhrif} \dicExampleCS{podobný vliv}
\dicEntry[svipan] \dicTerm{svipan} \dicIPA{{s}{v}{\textsci}{\textlengthmark}{\textsubring{b}}{a}{\textsubring{n}}} \dicPos{f}[7] \dicFlx{(‑ar, ‑ir)}[1] \dicPhraseIS{í einni svipan} \dicFlx{adv} \dicDirectTranslationCS{naráz, najednou} \dicExampleIS{Þetta breytist ekki í einni svipan.} \dicExampleCS{To se nezmění všechno najednou.};  \dicPhraseIS{í sömu svipan} \dicFlx{adv} \dicDirectTranslationCS{ve stejný okamžik}
\dicEntry[svipbrigðalaus] \textls[10]{\dicTerm{svip·brigða··laus} \dicIPA{{s}{v}{\textsci}{\textlengthmark}{\textsubring{b}}{r}{\textsci}{\textbabygamma}{ð}{a}{l}{\oe i}{s}}} \dicPos{adj}[5]\dicFlx{}[-1] \dicDirectTranslationCS{bezvýrazný, jed\-no\-tvár\-ný, mdlý} \dicExampleIS{svipbrigðalaus leikkona} \dicExampleCS{bezvýrazná herečka}
\dicEntry[svipbrigði] \dicTerm{svip··brigði} \dicIPA{{s}{v}{\textsci}{\textlengthmark}{\textsubring{b}}{r}{\textsci}{\textbabygamma}{ð}{\textsci}} \dicPos{n}[2] \dicFlx{pl}[19] \dicDirectTranslationCS{změna výrazu tváře, změna mimiky}
\dicEntry[svipfar] \dicTerm{svip··far} \dicIPA{{s}{v}{\textsci}{\textlengthmark}{\textsubring{b}}{f}{a}{\textsubring{r}}} \dicPos{n}[2] \dicFlx{(‑s)}[2] \dicLink{svipgerð}
\dicEntry[svipgerð] \dicTerm{svip··gerð}\dicTerm{, svipfar} \dicIPA{{s}\-{v}\-{\textsci}\-{\textlengthmark}\-{\textsubring{b}}\-{\r{\textObardotlessj}}\-{\textepsilon}\-{r}\-{\texttheta}\-} \dicPos{f}[7] \dicFlx{(‑ar, ‑ir)}[1] \dicFieldCat{biol.} \dicDirectTranslationCS{fenotyp}
\dicEntry[sviplaus] \textls[15]{\dicTerm{svip··laus} \dicIPA{{s}{v}{\textsci}{\textlengthmark}{\textsubring{b}}{l}{\oe i}{s}} \dicPos{adj}[5]\dicFlx{}[-1] \dicDirectTranslationCS{bezvýrazný, nevýrazný}} \dicExampleIS{sviplaust landslag} \dicExampleCS{nevýrazná krajina}
\dicEntry[sviplegur] \dicTerm{svip··legur} \dicIPA{{s}{v}{\textsci}{\textlengthmark}{\textsubring{b}}{l}{\textepsilon}{\textbabygamma}{\textscy}{\textsubring{r}}} \dicPos{adj}[1]\dicFlx{}[-8] \textbf{1.} \dicSynonym{skyndilegur} \dicDirectTranslationCS{náhlý, nenadálý} \dicExampleIS{sviplegur dauðdagi} \dicExampleCS{náhlá smrt}  \textbf{2.} \dicSynonym{hörmulegur} \dicDirectTranslationCS{tragický, žalostný} \dicExampleIS{sviplegt slys} \dicExampleCS{tragická nehoda}
\dicEntry[svipmikill] \dicTerm{svip··|mikill} \dicIPA{{s}{v}{\textsci}{\textlengthmark}{\textsubring{b}}{m}{\textsci}{\r{\textObardotlessj}}{\textsci}{\textsubring{d}}{\textsubring{l}}} \dicPos{adj}[11] \dicFlx{(comp ‑meiri, sup ‑mestur)}[12] \textbf{1.} \dicDirectTranslationCS{výrazný, nápadný}  \textbf{2.} \dicSynonym{áhrifamikill} \dicDirectTranslationCS{impozantní, působivý} \dicExampleIS{svipmikil náttúra} \dicExampleCS{impozantní příroda}
\dicEntry[svipmynd] \dicTerm{svip··mynd} \dicIPA{{s}{v}{\textsci}{\textlengthmark}{\textsubring{b}}{m}{\textsci}{n}{\textsubring{d}}} \dicPos{f}[7] \dicFlx{(‑ar, ‑ir)}[1] \dicDirectTranslationCS{momentka}
\dicEntry[svipstund] \dicTerm{svip··stund} \dicIPA{{s}{v}{\textsci}{\textlengthmark}{\textsubring{b}}{s}{\textsubring{d}}{\textscy}{n}{\textsubring{d}}} \dicPos{f}[7] \dicFlx{(‑ar)}[3] \dicSynonym{augnablik} \dicDirectTranslationCS{moment, okamžik};  \dicPhraseIS{á svipstundu} \dicFlx{adv} \dicDirectTranslationCS{během okamžiku}
\dicEntry[svipta] \dicTerm{svipt|a} \dicsymFrequent\  \dicIPA{{s}{v}{\textsci}{f}{\textsubring{d}}{a}} \dicPos{v}[2] \dicFlx{(‑i, ‑)}[12] \dicFlx{(acc +) dat} \textbf{1.} \dicFlx{acc + dat} \dicDirectTranslationCS{zbavit, připravit o~(průkaz ap.)};  \dicPhraseIS{svipta e‑n e‑u} \dicDirectTranslationCS{zbavit (koho čeho), připravit (koho) o~(co)} \dicExampleIS{svipta hana frelsi} \dicExampleCS{připravit ji o~svobodu}  \textbf{2.} \dicFlx{dat} \dicSynonym{rykkja} \dicDirectTranslationCS{strhnout, servat};  \dicIdiom{sviptast}{ \dicPhraseIS{sviptast}} \dicFlx{refl} \dicDirectTranslationCS{být stržen\,/\addthin serván (stan větrem ap.)}
\dicEntry[sviptifall] \dicTerm{svipti··|fall} \dicIPA{{s}{v}{\textsci}{f}{\textsubring{d}}{\textsci}{f}{a}{\textsubring{d}}{\textsubring{l}}} \dicPos{n}[2] \dicFlx{(‑falls, ‑föll)}[8] \dicFieldCat{jaz.} \dicDirectTranslationCS{ablativ}
\dicEntry[svipting] \dicTerm{svipt··ing} \dicIPA{{s}{v}{\textsci}{f}{\textsubring{d}}{i}{\ng}{\r{g}}} \dicPos{f}[4] \dicFlx{(‑ar, ‑ar)}[5] \textbf{1.} \dicSynonym*{það að svipta} \dicDirectTranslationCS{zbavení, odejmutí, odnětí} \dicExampleIS{svipting eigna} \dicExampleCS{odejmutí majetku}  \textbf{2.} \dicPhraseIS{sviptingar} \dicFlx{pl} \dicSynonym{átök} \dicDirectTranslationCS{boj, zápas, střet}
\dicEntry[svipudýr] \dicTerm{svipu··dýr} \dicIPA{{s}{v}{\textsci}{\textlengthmark}{\textsubring{b}}{\textscy}{\textsubring{d}}{i}{\textsubring{r}}} \dicPos{n}[2] \dicFlx{(‑s, ‑)}[5] \dicFieldCat{zool.} \dicDirectTranslationCS{bičíkovec} \textit{(l.~{\textLA{Mastigophora}})}
\dicEntry[svipull] \dicTerm{svipull} \dicIPA{{s}{v}{\textsci}{\textlengthmark}{\textsubring{b}}{\textscy}{\textsubring{d}}{\textsubring{l}}} \dicPos{adj}[8]\dicFlx{}[-4] \dicDirectTranslationCS{proměnlivý, nestálý, kolísavý}
\dicEntry[svipur] \dicTerm{svip|ur} \dicsymFrequent\  \dicIPA{{s}{v}{\textsci}{\textlengthmark}{\textsubring{b}}{\textscy}{\textsubring{r}}} \dicPos{m}[9] \dicFlx{(‑s, ‑ir)}[8] \textbf{1.} \dicSynonym*{andlitssvipur} \dicDirectTranslationCS{výraz, vzezření (tváře ap.)} \dicExampleIS{Það kom undrunarsvipur á hana.} \dicExampleCS{Na její tváři se objevil překvapený výraz.};  \dicPhraseIS{það kemur svipur á e‑n} \dicFlx{impers} \dicDirectTranslationCS{(kdo) se mračí}  \textbf{2.} \dicSynonym*{vofa} \dicDirectTranslationCS{zjevení, přízrak};  \dicIdiom{svipur}{ \dicPhraseIS{e‑að er ekki nema svipur hjá sjón}} \dicLangCat{přen.} \dicDirectTranslationCS{z~(čeho) už je jen stín}; { \dicPhraseIS{í svip}} \dicFlx{adv} \dicDirectTranslationCS{krátce, na krátkou chvíli};  \dicPhraseIS{í svipinn} \dicFlx{adv} \dicSynonym*{í bili} \dicDirectTranslationCS{v~tuto chvíli};  \dicPhraseIS{setja svip á e‑ð} \dicLangCat{přen.} \dicDirectTranslationCS{zanechat stopu na (čem), poznamenat (co), ovlivnit (co)};  \dicPhraseIS{það er svipur með e‑jum} \dicFlx{impers} \dicDirectTranslationCS{(kdo) si je podobný, (kdo) má společné rysy}
\dicEntry[sviss] \dicTerm{sviss\smash{\textsuperscript{1}}} \dicIPA{{s}{v}{\textsci}{s}{\textlengthmark}} \dicPos{m}[4] \dicFlx{(‑, ‑ar)}[35] \dicDirectTranslationCS{vznícení, zážeh}
\dicEntry[Sviss] \dicTerm{Sviss\smash{\textsuperscript{2}}} \dicIPA{{s}{v}{\textsci}{s}{\textlengthmark}} \dicPos{n}[4] \dicFlx{indecl}[2] \dicFieldCat{geog.} \dicDirectTranslationCS{Švýcarsko}
\dicEntry[svissa] \dicTerm{sviss|a} \dicIPA{{s}{v}{\textsci}{s}{\textlengthmark}{a}} \dicPos{v}[1] \dicFlx{(‑aði)}[1] \dicPhraseIS{svissa yfir í e‑u} \dicLangCat{hovor.} \dicDirectTranslationCS{přeskočit na (co), přejít na (co) (na jiný jazyk ap.)}
\dicEntry[Svisslendingur] \dicTerm{Sviss·lend··ing|ur} \dicIPA{{s}{v}{\textsci}{s}{l}{\textepsilon}{n}{\textsubring{d}}{i}{\ng}{\r{g}}{\textscy}{\textsubring{r}}} \dicPos{m}[6] \dicFlx{(‑s, ‑ar)}[8] \dicDirectTranslationCS{Švýcar(ka)}
\dicEntry[svissneskur] \dicTerm{sviss··neskur} \dicIPA{{s}{v}{\textsci}{s}{\textsubring{d}}{n}{\textepsilon}{s}{\r{g}}{\textscy}{\textsubring{r}}} \dicPos{adj}[1]\dicFlx{}[-1] \dicDirectTranslationCS{švýcarský}
\dicEntry[svitabað] \dicTerm{svita··|bað} \dicIPA{{s}{v}{\textsci}{\textlengthmark}{\textsubring{d}}{a}{\textsubring{b}}{a}{\texttheta}} \dicPos{n}[2] \dicFlx{(‑baðs, ‑böð)}[8] \dicPhraseIS{vera í svitabaði} \dicLangCat{přen.} \dicDirectTranslationCS{koupat se v~potu}
\dicEntry[svitahola] \dicTerm{svita··hol|a} \dicIPA{{s}{v}{\textsci}{\textlengthmark}{\textsubring{d}}{a}{h}{\textopeno}{l}{a}} \dicPos{f}[1] \dicFlx{(‑u, ‑ur)}[19] \dicFieldCat{anat.} \dicDirectTranslationCS{potní pór}
\dicEntry[svitakirtill] \dicTerm{svita··kirt|ill} \dicIPA{{s}{v}{\textsci}{\textlengthmark}{\textsubring{d}}{a}{c\smash{\textsuperscript{h}}}{\textsci}{\textsubring{r}}{\textsubring{d}}{\textsci}{\textsubring{d}}{\textsubring{l}}} \dicPos{m}[6] \dicFlx{(‑ils, ‑lar)}[35] \dicFieldCat{anat.} \dicDirectTranslationCS{potní žláza}
\dicEntry[svitakóf] \dicTerm{svita··kóf} \dicIPA{{s}{v}{\textsci}{\textlengthmark}{\textsubring{d}}{a}{k\smash{\textsuperscript{h}}}{ou}{f}} \dicPos{n}[2] \dicFlx{(‑s)}[2] \dicLangCat{přen.} \dicDirectTranslationCS{potní lázeň};  \dicPhraseIS{vera í svitakófi} \dicDirectTranslationCS{potit se jako kůň, být zpocený jako myš}
\dicEntry[svitalykt] \dicTerm{svita··lykt} \dicIPA{{s}{v}{\textsci}{\textlengthmark}{\textsubring{d}}{a}{l}{\textsci}{x}{\textsubring{d}}} \dicPos{f}[7] \dicFlx{(‑ar)}[3] \dicDirectTranslationCS{zápach potu}
\dicEntry[svitalyktareyðir] \dicTerm{svita·lyktar··eyð|ir} \dicIPA{{s}{v}{\textsci}{\textlengthmark}{\textsubring{d}}{a}{l}{\textsci}{x}{\textsubring{d}}{a}{r}{ei}{ð}{\textsci}{\textsubring{r}}} \dicPos{m}[7] \dicFlx{(‑is, ‑ar)}[1] \dicDirectTranslationCS{deodorant, dezodorant}
\dicEntry[sviti] \dicTerm{svit|i} \dicIPA{{s}{v}{\textsci}{\textlengthmark}{\textsubring{d}}{\textsci}} \dicPos{m}[1] \dicFlx{(‑a)}[3] \dicDirectTranslationCS{pot} \dicExampleIS{strjúka svitann af enninu} \dicExampleCS{setřít si pot z~čela}
\dicEntry[svitna] \dicTerm{svitn|a} \dicIPA{{s}{v}{\textsci}{h}{\textsubring{d}}{n}{a}} \dicPos{v}[1] \dicFlx{(‑aði)}[1] \dicDirectTranslationCS{(z)potit se}
\dicEntry[svíð] \dicTerm{svíð} \dicIPA{{s}{v}{i}{\textlengthmark}{\texttheta}} \dicPos{v} \dicFlx{ind praes sg 1 pers} \dicLink{svíða}
\dicEntry[svíða] \dicTerm{svíða} \dicIPA{{s}{v}{i}{\textlengthmark}{ð}{a}} \dicPos{v}[6] \dicFlx{(svíð, sveið, sviðum, sviði, sviðið)}[67] \dicFlx{acc} \textbf{1.} \dicDirectTranslationCS{ožehnout, sežehnout, přismahnout, připálit} \dicExampleIS{svíða á sér hárið} \dicExampleCS{ožehnout si vlasy}  \textbf{2.} \dicPhraseIS{e‑n svíður} \dicFlx{impers} \dicDirectTranslationCS{(koho) pálí, (koho) ostře bolí} \dicExampleIS{Mig svíður í fingurinn.} \dicExampleCS{Pálí mě prst.}  \textbf{3.} \dicPhraseIS{e‑m svíður e‑að} \dicFlx{impers} \dicDirectTranslationCS{(koho co) trápí, (koho co) pálí} \dicExampleIS{Mér svíður framkoma hans.} \dicExampleCS{Trápí mě jeho chování.}
\dicEntry[svíðandi] \dicTerm{svíð··andi} \dicIPA{{s}{v}{i}{\textlengthmark}{ð}{a}{n}{\textsubring{d}}{\textsci}} \dicPos{adj}[13] \dicFlx{indecl}[1] \dicDirectTranslationCS{pálivý, (ostře) bolící} \dicExampleIS{svíðandi sársauki} \dicExampleCS{pálivá bolest}
\dicEntry[svíðingsháttur] \dicTerm{svíðings··hátt|ur} \dicIPA{{s}{v}{i}{\textlengthmark}{ð}{i}{\ng}{s}{h}{au}{h}{\textsubring{d}}{\textscy}{\textsubring{r}}} \dicPos{m}[12] \dicFlx{(‑ar)}[8] \dicDirectTranslationCS{lakomost}
\dicEntry[svíðingur] \dicTerm{svíð··ing|ur} \dicIPA{{s}{v}{i}{\textlengthmark}{ð}{i}{\ng}{\r{g}}{\textscy}{\textsubring{r}}} \dicPos{m}[6] \dicFlx{(‑s, ‑ar)}[8] \dicDirectTranslationCS{lakomec, krkoun}
\dicEntry[svíf] \dicTerm{svíf} \dicIPA{{s}{v}{i}{\textlengthmark}{f}} \dicPos{v} \dicFlx{ind praes sg 1 pers} \dicLink{svífa}
\dicEntry[svífa] \dicTerm{svífa} \dicsymFrequent\  \dicIPA{{s}{v}{i}{\textlengthmark}{v}{a}} \dicPos{v}[6] \dicFlx{(svíf, sveif, svifum, svifi, svifið)}[69] \dicSynonym{fljúga} \dicDirectTranslationCS{vznášet se, poletovat, létat, plout vzduchem} \dicExampleIS{svífa til og frá} \dicExampleCS{poletovat sem a~tam};  \dicIdiom{svífa}[á]{ \dicPhraseIS{e‑að svífur á e‑n}} {\textbf{a.}} \dicDirectTranslationCS{(co) na (koho) padá (spánek ap.)};  {\textbf{b.}} \dicDirectTranslationCS{(co) začíná působit na (koho), (co komu) stoupá do hlavy (alkohol ap.)} \dicExampleIS{Vínið sveif á mig.} \dicExampleCS{Víno mi stouplo do hlavy.}; { \dicPhraseIS{e‑r svífur á e‑n}} \dicLangCat{přen.} \dicDirectTranslationCS{(kdo) se na (koho) vrhá (hostitel na hosty ap.)};  \dicIdiom{svífast}{ \dicPhraseIS{svífast einskis}} \dicFlx{refl} \dicLangCat{přen.} \dicDirectTranslationCS{jít přes mrtvoly}
\dicEntry[svífandi] \dicTerm{svíf··andi} \dicIPA{{s}{v}{i}{\textlengthmark}{v}{a}{n}{\textsubring{d}}{\textsci}} \dicPos{adj}[13] \dicFlx{indecl}[1] \dicDirectTranslationCS{vznášející se, poletující}
\dicEntry[Svíi] \dicTerm{Sví|i} \dicIPA{{s}{v}{i}{j}{\textsci}} \dicPos{m}[1] \dicFlx{(‑a, ‑ar)}[1] \dicDirectTranslationCS{Švéd(ka)}
\dicEntry[svík] \dicTerm{svík} \dicIPA{{s}{v}{i}{\textlengthmark}{\r{g}}} \dicPos{v} \dicFlx{ind praes sg 1 pers} \dicLink{svíkja}
\dicEntry[svíkja] \dicTerm{svíkja} \dicsymFrequent\  \dicIPA{{s}{v}{i}{\textlengthmark}{\r{\textObardotlessj}}{a}} \dicPos{v}[6] \dicFlx{(svík, sveik, svikum, sviki, svikið)}[70] \dicFlx{acc} \textbf{1.} \dicDirectTranslationCS{zradit, podvést, podvádět, (o)klamat, podrazit}  \textbf{2.} \dicDirectTranslationCS{porušit (slib ap.)} \dicExampleIS{svíkja loforðið} \dicExampleCS{porušit slib};  \dicPhraseIS{svíkja lit} \dicLangCat{přen.} \dicDirectTranslationCS{měnit barvu jako chameleon};  \dicIdiom{svíkjast}[um]{ \dicPhraseIS{svíkjast um e‑ð}} \dicFlx{refl} \dicDirectTranslationCS{vyhýbat se (čemu), ulít se z~(čeho)};  \dicIdiom{svíkjast}[undan]{ \dicPhraseIS{svíkjast undan e‑u}} \dicFlx{refl} \dicDirectTranslationCS{vyhýbat se (čemu), unikat (čemu)}; { \dicPhraseIS{svíkjast undan merkjum}} \dicFlx{refl} \dicLangCat{přen.} \dicDirectTranslationCS{zradit svoje ideály, jít proti vlastnímu přesvědčení}
\dicEntry[svími] \dicTerm{svím|i} \dicIPA{{s}{v}{i}{\textlengthmark}{m}{\textsci}} \dicPos{m}[1] \dicFlx{(‑a)}[3] \dicLangCat{zast.} \dicLink{svimi}
\dicEntry[svín] \dicTerm{svín} \dicIPA{{s}{v}{i}{\textlengthmark}{\textsubring{n}}} \dicPos{n}[2] \dicFlx{(‑s, ‑)}[5] \textbf{1.} \dicFieldCat{zool.} \dicDirectTranslationCS{prase, vepř, čuně, svině} \textit{(l.~{\textLA{Suidae}})}  \dicsymPhoto\   \textbf{2.} \dicLangCat{han.} \dicDirectTranslationCS{svině} \dicIndirectTranslationCS{(o~člověku)}
\dicFigure{ds_image_svin_0_1.jpg}{Svín}{Svín - Scott Bauer, PD}
\dicEntry[svína] \dicTerm{svín|a} \dicIPA{{s}{v}{i}{\textlengthmark}{n}{a}} \dicPos{v}[1] \dicFlx{(‑aði)}[1] \dicFlx{acc} \dicSynonym{óhreinka} \dicDirectTranslationCS{zasvinit, zaneřádit} \dicExampleIS{svína e‑ð út} \dicExampleCS{zasvinit (co)};  \dicIdiom{svína}[fyrir]{ \dicPhraseIS{svína fyrir e‑n}} \dicDirectTranslationCS{nacpat se (autem) před (koho)}
\dicEntry[svínabú] \dicTerm{svína··bú} \dicIPA{{s}{v}{i}{\textlengthmark}{n}{a}{\textsubring{b}}{u}} \dicPos{n}[2] \dicFlx{(‑s, ‑)}[5] \dicDirectTranslationCS{vepřín, prasečák}
\dicEntry[svínaflensa] \dicTerm{svína··flens|a}\dicTerm{, svínainflúensa} \dicIPA{{s}\-{v}\-{i}\-{\textlengthmark}\-{n}\-{a}\-{f}\-{l}\-{\textepsilon}\-{n}\-{s}\-{a}\-} \dicPos{f}[1] \dicFlx{(‑u)}[5] \dicFieldCat{med. } \dicDirectTranslationCS{prasečí chřipka}
\dicEntry[svínahakk] \dicTerm{svína··hakk} \dicIPA{{s}{v}{i}{\textlengthmark}{n}{a}{h}{a}{h}{\r{g}}} \dicPos{n}[2] \dicFlx{(‑s)}[2] \dicFieldCat{kulin.} \dicDirectTranslationCS{vepřové mleté maso}
\dicEntry[svínainflúensa] \dicTerm{svína··in·flúens|a} \dicIPA{{s}{v}{i}{\textlengthmark}{n}{a}{\textsci}{n}{f}{l}{u}{\textepsilon}{n}{s}{a}} \dicPos{f}[1] \dicFlx{(‑u)}[5] \dicLink{svínaflensa}
\dicEntry[svínakjöt] \dicTerm{svína··kjöt} \dicIPA{{s}{v}{i}{\textlengthmark}{n}{a}{c\smash{\textsuperscript{h}}}{\oe}{\textsubring{d}}} \dicPos{n}[2] \dicFlx{(‑s)}[2] \dicDirectTranslationCS{vepřové (maso)}
\dicEntry[svínarí] \dicTerm{svínarí} \dicIPA{{s}{v}{i}{\textlengthmark}{n}{a}{r}{i}} \dicPos{n}[2] \dicFlx{(‑s)}[2] \dicDirectTranslationCS{svinstvo, sviňárna, ničemnost}
\dicEntry[svínasteik] \dicTerm{svína··steik} \dicIPA{{s}{v}{i}{\textlengthmark}{n}{a}{s}{\textsubring{d}}{ei}{\r{g}}} \dicPos{f}[8] \dicFlx{(‑ar, ‑ur)}[1] \dicFieldCat{kulin.} \dicDirectTranslationCS{vepřová pečeně, vepřový steak}
\dicEntry[svínastía] \dicTerm{svína··stí|a} \dicIPA{{s}{v}{i}{\textlengthmark}{n}{a}{s}{\textsubring{d}}{i}{j}{a}} \dicPos{f}[1] \dicFlx{(‑u, ‑ur)}[7] \textbf{1.} \dicDirectTranslationCS{prasečí chlívek}  \textbf{2.} \dicLangCat{přen.} \dicDirectTranslationCS{(prasečí) chlívek (dětský pokoj ap.)}
\dicEntry[svínslegur] \dicTerm{svíns··legur} \dicIPA{{s}{v}{i}{n}{s}{l}{\textepsilon}{\textbabygamma}{\textscy}{\textsubring{r}}} \dicPos{adj}[1]\dicFlx{}[-8] \dicDirectTranslationCS{svinský, hanebný}
\dicEntry[svíta] \dicTerm{svít|a} \dicIPA{{s}{v}{i}{\textlengthmark}{\textsubring{d}}{a}} \dicPos{f}[1] \dicFlx{(‑u, ‑ur)}[7] \textbf{1.} \dicFieldCat{hud.} \dicDirectTranslationCS{suita, svita} \dicIndirectTranslationCS{(hudební skladba s~několika volně seřazenými větami)}  \textbf{2.} \dicDirectTranslationCS{(hotelové) apartmá}
\dicEntry[svívirða] \dicTerm{sví··virð|a\smash{\textsuperscript{1}}} \dicIPA{{s}{v}{i}{\textlengthmark}{v}{\textsci}{r}{ð}{a}} \dicPos{f}[1] \dicFlx{(‑u)}[5] \dicSynonym{smán} \dicDirectTranslationCS{hanba, zostuzení, ostuda} \dicExampleIS{gera henni svívirðu} \dicExampleCS{udělat jí ostudu}
\dicEntry[svívirða] \dicTerm{sví··vir|ða\smash{\textsuperscript{2}}} \dicIPA{{s}{v}{i}{\textlengthmark}{v}{\textsci}{r}{ð}{a}} \dicPos{v}[2] \dicFlx{(‑ti, ‑t)}[47] \dicFlx{acc} \dicSynonym{smána} \dicDirectTranslationCS{zostudit, zdiskreditovat, zhanobit, potupit, zneuctít} \dicExampleIS{svívirða minningu hennar} \dicExampleCS{zostudit její památku}
\dicEntry[svívirðilegur] \dicTerm{sví·virði··legur} \dicIPA{{s}{v}{i}{\textlengthmark}{v}{\textsci}{r}{ð}{\textsci}{l}{\textepsilon}{\textbabygamma}{\textscy}{\textsubring{r}}} \dicPos{adj}[1]\dicFlx{}[-8] \dicSynonym{skammarlegur} \dicDirectTranslationCS{ostudný, hanebný} \dicExampleIS{svívirðilegt verk} \dicExampleCS{hanebná činnost}
\dicEntry[svívirðing] \dicTerm{sví··virð·ing} \dicIPA{{s}{v}{i}{\textlengthmark}{v}{\textsci}{r}{ð}{i}{\ng}{\r{g}}} \dicPos{f}[4] \dicFlx{(‑ar, ‑ar)}[5] \textbf{1.} \dicSynonym{smán} \dicDirectTranslationCS{zostuzení, zdiskreditování, zhanobení, potupení, zneuctění}  \textbf{2.} \dicPhraseIS{svívirðingar} \dicFlx{pl} \dicSynonym*{ónotaorð} \dicDirectTranslationCS{urážky, nadávky} \dicExampleIS{vera með svívirðingar} \dicExampleCS{nadávat}
\dicEntry[Svíþjóð] \dicTerm{Sví··þjóð} \dicsymFrequent\  \dicIPA{{s}{v}{i}{\textlengthmark}{\texttheta}{j}{ou}{\texttheta}} \dicPos{f}[7] \dicFlx{(‑ar)}[4] \dicFieldCat{geog.} \dicDirectTranslationCS{Švédsko} \dicExampleIS{við suðausturströnd Svíþjóðar} \dicExampleCS{u~jihovýchodního pobřeží Švédska}
\dicEntry[svo] \dicTerm{svo\smash{\textsuperscript{1}}} \dicsymFrequent\  \dicIPA{{s}{v}{\textopeno}{\textlengthmark}} \dicPos{adv} \dicDirectTranslationCS{tak} \dicIndirectTranslationCS{(užívané s~adj)} \dicExampleIS{Hún er svo falleg.} \dicExampleCS{Je tak krásná.};  \dicPhraseIS{er það svo?} \dicDirectTranslationCS{není to tak?, nemám pravdu?}
\dicEntry[svo] \dicTerm{svo\smash{\textsuperscript{2}}} \dicIPA{{s}{v}{\textopeno}{\textlengthmark}} \dicPos{conj} \textbf{1.} \dicDirectTranslationCS{tak, takže} \dicIndirectTranslationCS{(označuje výsledek nebo závěr)}  \textbf{2.} \dicDirectTranslationCS{tak, aby; aby} \dicIndirectTranslationCS{(označuje účel)};  \dicIdiom{svo}[að]{ \dicPhraseIS{svo að}} \dicFlx{conj} {\textbf{a.}} \dicDirectTranslationCS{tak, že} \dicIndirectTranslationCS{(o~následcích)};  {\textbf{b.}} \dicDirectTranslationCS{tak, aby} \dicIndirectTranslationCS{(o~účelu)};  \dicIdiom{svo}[og]{ \dicPhraseIS{svo og}} \dicFlx{conj} \dicDirectTranslationCS{také, rovněž};  \dicIdiom{svo}[sem]{ \dicPhraseIS{svo gott sem}} \dicFlx{adv} \dicDirectTranslationCS{téměř, skoro}; { \dicPhraseIS{svo framarlega\,/\addthin fremi sem}} \dicFlx{conj} \dicDirectTranslationCS{pokud}; { \dicPhraseIS{svo sem}} \dicFlx{conj} \dicDirectTranslationCS{jak} \dicExampleIS{svo sem vitað er} \dicExampleCS{jak je známo}; { \dicPhraseIS{svo sem}} \dicFlx{adv} \dicDirectTranslationCS{vlastně, ve skutečnosti};  \dicIdiom{svo}[til]{ \dicPhraseIS{svo til}} \dicFlx{adv} \dicDirectTranslationCS{téměř, skoro, takřka};  \dicIdiom{svo}{ \dicPhraseIS{svo ekki verður um villst}} \dicFlx{adv} \dicDirectTranslationCS{jednoznačně, nesporně}
\dicEntry[svoddan] \dicTerm{svoddan} \dicIPA{{s}{v}{\textopeno}{\textsubring{d}}{\textlengthmark}{a}{\textsubring{n}}} \dicPos{adj}[13] \dicFlx{indecl}[3] \dicLink{soddan}
\dicEntry[svofelldur] \dicTerm{svo··felldur} \dicIPA{{s}{v}{\textopeno}{\textlengthmark}{f}{\textepsilon}{l}{\textsubring{d}}{\textscy}{\textsubring{r}}} \dicPos{adj}[2]\dicFlx{}[-17] \dicSynonym{slíkur} \dicDirectTranslationCS{takovýto, následovný} \dicExampleIS{á svofelldan hátt} \dicExampleCS{následovným způsobem}
\dicEntry[svohljóðandi] \dicTerm{svo··hljóð·andi} \dicIPA{{s}{v}{\textopeno}{\textlengthmark}{\textsubring{l}}{j}{ou}{ð}{a}{n}{\textsubring{d}}{\textsci}} \dicPos{adj}[13] \dicFlx{indecl}[1] \dicSynonym{slíkur} \dicDirectTranslationCS{takový, následující, takto znějící} \dicExampleIS{svohljóðandi tillaga} \dicExampleCS{takový návrh}
\dicEntry[svokallaður] \dicTerm{svo··|kall·aður} \dicsymFrequent\  \dicIPA{{s}{v}{\textopeno}{\textlengthmark}{k\smash{\textsuperscript{h}}}{a}{\textsubring{d}}{l}{a}{ð}{\textscy}{\textsubring{r}}} \dicPos{adj}[3] \dicFlx{(f ‑kölluð)}[2] \dicDirectTranslationCS{takzvaný} \dicExampleIS{svokölluð borgaraleg blöð} \dicExampleCS{takzvané městské noviny}
\dicEntry[svoleiðis] \dicTerm{svo··leiðis\smash{\textsuperscript{1}}} \dicIPA{{s}{v}{\textopeno}{\textlengthmark}{l}{ei}{ð}{\textsci}{s}} \dicPos{adj}[13] \dicFlx{indecl}[1] \dicDirectTranslationCS{takový}
\dicEntry[svoleiðis] \dicTerm{svo··leiðis\smash{\textsuperscript{2}}} \dicsymFrequent\  \dicIPA{{s}{v}{\textopeno}{\textlengthmark}{l}{ei}{ð}{\textsci}{s}} \dicPos{adv} \dicSynonym{þannig\smash{\textsuperscript{2}}} \dicDirectTranslationCS{tak, takto, takhle} \dicExampleIS{Það er ekki svoleiðis.} \dicExampleCS{Tak to není.}
\dicEntry[svolgra] \dicTerm{svolgr|a} \dicIPA{{s}{v}{\textopeno}{l}{\r{g}}{r}{a}} \dicPos{v}[1] \dicFlx{(‑aði)}[1] \dicFlx{acc} \dicSynonym{svelgja} \dicDirectTranslationCS{lokat, (vy)chlastat, (vy)chlemtat};  \dicPhraseIS{svolgra í sig e‑ð} \dicDirectTranslationCS{lokat (co) (kafe ap.)}
\dicEntry[svoli] \dicTerm{svol|i} \dicIPA{{s}{v}{\textopeno}{\textlengthmark}{l}{\textsci}} \dicPos{m}[1] \dicFlx{(‑a, ‑ar)}[1] \dicDirectTranslationCS{klacek, výrostek}
\dicEntry[svolítið] \dicTerm{svo··lítið} \dicsymFrequent\  \dicIPA{{s}{v}{\textopeno}{\textlengthmark}{l}{i}{\textsubring{d}}{\textsci}{\texttheta}} \dicPos{adv} \dicSynonym{aðeins} \dicDirectTranslationCS{(tak) trochu, nepatrně, poněkud} \dicExampleIS{Það kom mér svolítið á óvart.} \dicExampleCS{Trochu mě to překvapilo.}
\dicEntry[svolítill] \dicTerm{svo··lítill} \dicsymFrequent\  \dicIPA{{s}{v}{\textopeno}{\textlengthmark}{l}{i}{\textsubring{d}}{\textsci}{\textsubring{d}}{\textsubring{l}}} \dicPos{adj}[11]\dicFlx{}[-11] \dicDirectTranslationCS{maličký, nepatrný} \dicExampleIS{vekja svolitla athygli} \dicExampleCS{vzbudit nepatrnou pozornost}
\dicEntry[svona] \dicTerm{svona\smash{\textsuperscript{1}}} \dicIPA{{s}{v}{\textopeno}{\textlengthmark}{n}{a}} \dicPos{adj}[13] \dicFlx{indecl}[1] \dicDirectTranslationCS{takový} \dicExampleIS{svona fólk} \dicExampleCS{takoví lidé}
\dicEntry[svona] \dicTerm{svona\smash{\textsuperscript{2}}} \dicsymFrequent\  \dicIPA{{s}{v}{\textopeno}{\textlengthmark}{n}{a}} \dicPos{adv} \textbf{1.} \dicSynonym*{á þennan hátt} \dicDirectTranslationCS{tak, takovým způsobem} \dicExampleIS{Svona fór það.} \dicExampleCS{Tak to šlo.}  \textbf{2.} \dicDirectTranslationCS{tak} \dicIndirectTranslationCS{(k~zdůraznění)} \dicExampleIS{svona klár} \dicExampleCS{tak chytrý}  \textbf{3.} \dicSynonym{sirka} \dicDirectTranslationCS{tak asi\,/\addthin přibližně}  \textbf{4.} \dicSynonym{eiginlega} \dicDirectTranslationCS{tak nějak}  \textbf{5.} \dicPhraseIS{svona svona} \dicDirectTranslationCS{nono, no tak (přestaň plakat ap.)}
\dicEntry[svonefndur] \dicTerm{svo··nefndur} \dicsymFrequent\  \dicIPA{{s}{v}{\textopeno}{\textlengthmark}{n}{\textepsilon}{m}{\textsubring{d}}{\textscy}{\textsubring{r}}} \dicPos{adj}[2]\dicFlx{}[-17] \dicDirectTranslationCS{takzvaný} \dicExampleIS{svonefnt gregoríanskt tímatal} \dicExampleCS{takzvaný gregoriánský kalendář}
\dicEntry[svunta] \dicTerm{svunt|a} \dicIPA{{s}{v}{\textscy}{\textsubring{n}}{\textsubring{d}}{a}} \dicPos{f}[1] \dicFlx{(‑u, ‑ur)}[19] \textbf{1.} \dicDirectTranslationCS{zástěra} \dicIndirectTranslationCS{(kus pracovního oblečení)}  \textbf{2.} \dicFieldCat{ekon.} \dicDirectTranslationCS{kontrolní útržek\,/\addthin ústřižek}
\dicEntry[svæðamót] \dicTerm{svæða··mót} \dicIPA{{s}{v}{a}{i}{\textlengthmark}{ð}{a}{m}{ou}{\textsubring{d}}} \dicPos{n}[2] \dicFlx{(‑s, ‑)}[5] \dicDirectTranslationCS{oblastní\,/\addthin regionální turnaj}
\dicEntry[svæði] \dicTerm{svæði} \dicsymFrequent\  \dicIPA{{s}{v}{a}{i}{\textlengthmark}{ð}{\textsci}} \dicPos{n}[2] \dicFlx{(‑s, ‑)}[14] \dicDirectTranslationCS{oblast, území, prostor, prostranství} \dicExampleIS{á öllum svæðum landsins} \dicExampleCS{ve všech oblastech země}
\dicEntry[svæðisnúmer] \dicTerm{svæðis··númer} \dicIPA{{s}{v}{a}{i}{\textlengthmark}{ð}{\textsci}{s}{n}{u}{m}{\textepsilon}{\textsubring{r}}} \dicPos{n}[2] \dicFlx{(‑s, ‑)}[5] \dicDirectTranslationCS{(meziměstská) předvolba, (telefonní) směrové číslo}
\dicEntry[svæfa] \dicTerm{svæf|a} \dicIPA{{s}{v}{a}{i}{\textlengthmark}{v}{a}} \dicPos{v}[2] \dicFlx{(‑ði, ‑t)}[110] \dicFlx{acc} \textbf{1.} \dicSynonym*{láta sofna} \dicDirectTranslationCS{uspat, uspávat (pohádkou ap.)} \dicExampleIS{svæfa barnið} \dicExampleCS{uspat dítě}  \textbf{2.} \dicDirectTranslationCS{uspat (anestetiky), narkotizovat} \dicExampleIS{svæfa sjúklinginn} \dicExampleCS{uspat pacienta}  \textbf{3.} \dicSynonym*{þagga niður} \dicDirectTranslationCS{ututlat, utajit} \dicExampleIS{svæfa málið} \dicExampleCS{ututlat záležitost}  \textbf{4.} \dicSynonym*{aflífa með lyfjagjöf} \dicDirectTranslationCS{uspat, utratit (domácí zvíře ap.)}
\dicEntry[svæfandi] \dicTerm{svæf··andi} \dicIPA{{s}{v}{a}{i}{\textlengthmark}{v}{a}{n}{\textsubring{d}}{\textsci}} \dicPos{adj}[13] \dicFlx{indecl}[1] \dicDirectTranslationCS{uspávající, uspávací} \dicExampleIS{svæfandi hljómlist} \dicExampleCS{uspávající hudba}
\dicEntry[svæfi] \dicTerm{svæfi} \dicIPA{{s}{v}{a}{i}{\textlengthmark}{v}{\textsci}} \dicPos{v} \dicFlx{con pf sg 1 pers} \dicLink{sofa}
\dicEntry[svæfill] \dicTerm{svæf|ill} \dicIPA{{s}{v}{a}{i}{\textlengthmark}{v}{\textsci}{\textsubring{d}}{\textsubring{l}}} \dicPos{m}[6] \dicFlx{(‑ils, ‑lar)}[35] \dicDirectTranslationCS{polštářek (pod hlavu), (malý) podhlavník, podhlavníček}
\dicEntry[svæfilyf] \dicTerm{svæfi··lyf} \dicIPA{{s}{v}{a}{i}{\textlengthmark}{v}{\textsci}{l}{\textsci}{f}} \dicPos{n}[2] \dicFlx{(‑s, ‑)}[13] \dicLink{svæfingarlyf}
\dicEntry[svæfing] \dicTerm{svæf··ing} \dicIPA{{s}{v}{a}{i}{\textlengthmark}{v}{i}{\ng}{\r{g}}} \dicPos{f}[4] \dicFlx{(‑ar, ‑ar)}[5] \textbf{1.} \dicSynonym*{það að svæfa} \dicDirectTranslationCS{uspání}  \textbf{2.} \dicFieldCat{med.} \dicDirectTranslationCS{narkóza}  \textbf{3.} \dicSynonym*{sefun} \dicDirectTranslationCS{utišení, zklidnění} \dicExampleIS{svæfing harma} \dicExampleCS{utišení zármutku}  \textbf{4.} \dicSynonym*{það að aflífa skepnu} \dicDirectTranslationCS{porážka (dobytka ap.)}
\dicEntry[svæfingarlyf] \dicTerm{svæfingar··lyf}\dicTerm{, svæfilyf} \dicIPA{{s}\-{v}\-{a}\-{i}\-{\textlengthmark}\-{v}\-{i}\-{\ng}\-{\r{g}}\-{a}\-{r}\-{l}\-{\textsci}\-{f}\-} \dicPos{n}[2] \dicFlx{(‑s, ‑)}[13] \dicFieldCat{med.} \dicDirectTranslationCS{anestetikum}
\dicEntry[svæfingarlæknir] \dicTerm{svæfingar··lækn|ir} \dicIPA{{s}{v}{a}{i}{\textlengthmark}{v}{i}{\ng}{\r{g}}{a}{r}{l}{a}{i}{h}{\r{g}}{n}{\textsci}{\textsubring{r}}} \dicPos{m}[7] \dicFlx{(‑is, ‑ar)}[1] \dicFieldCat{med.} \dicDirectTranslationCS{anesteziolog, anestezioložka}
\dicEntry[svækja] \dicTerm{svækj|a} \dicIPA{{s}{v}{a}{i}{\textlengthmark}{\r{\textObardotlessj}}{a}} \dicPos{f}[1] \dicFlx{(‑u, ‑ur)}[17] \dicDirectTranslationCS{dusno, dusné počasí} \dicExampleIS{svækja og loftleysi} \dicExampleCS{dusno a~nedostatek vzduchu}
\dicEntry[svækjuloft] \dicTerm{svækju··loft} \dicIPA{{s}{v}{a}{i}{\textlengthmark}{\r{\textObardotlessj}}{\textscy}{l}{\textopeno}{f}{\textsubring{d}}} \dicPos{n}[2] \dicFlx{(‑s, ‑)}[5] \dicDirectTranslationCS{dusný vzduch}
\dicEntry[svæla] \dicTerm{svæl|a\smash{\textsuperscript{1}}} \dicIPA{{s}{v}{a}{i}{\textlengthmark}{l}{a}} \dicPos{f}[1] \dicFlx{(‑u)}[5] \dicDirectTranslationCS{hustý dým\,/\addthin kouř} \dicExampleIS{reykjarsvæla} \dicExampleCS{hustý cigaretový kouř}
\dicEntry[svæla] \dicTerm{svæl|a\smash{\textsuperscript{2}}} \dicIPA{{s}{v}{a}{i}{\textlengthmark}{l}{a}} \dicPos{v}[2] \dicFlx{(‑di, ‑t)}[141] \dicFlx{acc} \textbf{1.} \dicSynonym*{kæfa í reyk} \dicDirectTranslationCS{vykouřit (zvíře z~nory ap.)} \dicExampleIS{svæla tófu í greni} \dicExampleCS{vykouřit lišku z~brlohu}  \textbf{2.} \dicDirectTranslationCS{čoudit, čudit, bafat} \dicExampleIS{svæla vindil} \dicExampleCS{bafat doutník};  \dicIdiom{svæla}[í]{ \dicPhraseIS{svæla e‑u í sig}} \dicDirectTranslationCS{donutit se sníst (co)} \dicIndirectTranslationCS{(o~něčem nechutném)};  \dicIdiom{svæla}[út]{ \dicPhraseIS{svæla e‑ð út úr e‑m}} \dicDirectTranslationCS{vymámit (co) z~(koho) (dotěrností ap.)}
\dicEntry[sværi] \dicTerm{sværi} \dicIPA{{s}{v}{a}{i}{\textlengthmark}{r}{\textsci}} \dicPos{v} \dicFlx{con pf sg 1 pers} \dicLink{sverja}
\dicEntry[svæsinn] \dicTerm{svæsinn} \dicIPA{{s}{v}{a}{i}{\textlengthmark}{s}{\textsci}{\textsubring{n}}} \dicPos{adj}[6]\dicFlx{}[-2] \textbf{1.} \dicSynonym{ofsafenginn} \dicDirectTranslationCS{zarytý, drsný} \dicExampleIS{svæsinn andstæðingur} \dicExampleCS{zarytý protivník}  \textbf{2.} \dicSynonym{ákafur} \dicDirectTranslationCS{prudký, úporný} \dicExampleIS{svæsinn vindur} \dicExampleCS{prudký vítr}
\dicEntry[svöð] \dicTerm{svöð} \dicIPA{{s}{v}{\oe}{\textlengthmark}{\texttheta}} \dicPos{n} \dicFlx{pl nom} \dicLink{svað}
\dicEntry[svölu] \dicTerm{svölu} \dicIPA{{s}{v}{\oe}{\textlengthmark}{l}{\textscy}} \dicPos{f} \dicFlx{sg gen} \dicLink{svala\smash{\textsuperscript{1}}}
\dicEntry[svölun] \dicTerm{svöl|un} \dicIPA{{s}{v}{\oe}{\textlengthmark}{l}{\textscy}{\textsubring{n}}} \dicPos{f}[7] \dicFlx{(‑unar)}[12] \textbf{1.} \dicSynonym*{það að svala} \dicDirectTranslationCS{ochlazení, osvěžení} \dicExampleIS{svölun við þorsta} \dicExampleCS{osvěžení v~žízni}  \textbf{2.} \dicSynonym{fullnæging} \dicDirectTranslationCS{uspokojení, ukojení} \dicExampleIS{svölun girnda} \dicExampleCS{ukojení chtíče}
\dicEntry[svölur] \dicTerm{svölur} \dicIPA{{s}{v}{\oe}{\textlengthmark}{l}{\textscy}{\textsubring{r}}} \dicPos{f} \dicFlx{pl nom} \dicLink{svala\smash{\textsuperscript{1}}}
\dicEntry[svör] \dicTerm{svör} \dicIPA{{s}{v}{\oe}{\textlengthmark}{\textsubring{r}}} \dicPos{n} \dicFlx{pl nom} \dicLink{svar}
\dicEntry[svörfun] \dicTerm{svörf|un} \dicIPA{{s}{v}{\oe}{r}{v}{\textscy}{\textsubring{n}}} \dicPos{f}[7] \dicFlx{(‑unar)}[12] \dicFieldCat{geol.} \dicDirectTranslationCS{abraze}
\dicEntry[svörun] \dicTerm{svör|un} \dicIPA{{s}{v}{\oe}{\textlengthmark}{r}{\textscy}{\textsubring{n}}} \dicPos{f}[7] \dicFlx{(‑unar)}[12] \textbf{1.} \dicSynonym{viðbrögð} \dicDirectTranslationCS{reakce, odezva}  \textbf{2.} \dicFieldCat{biol.} \dicDirectTranslationCS{reakce (imunitní ap.)}
\dicEntry[sybbinn] \dicTerm{sybbinn} \dicIPA{{s}{\textsci}{\textsubring{b}}{\textlengthmark}{\textsci}{\textsubring{n}}} \dicPos{adj}[6]\dicFlx{}[-2] \dicSynonym{syfjaður} \dicDirectTranslationCS{ospalý} \dicIndirectTranslationCS{(zvláště o~dětech)}
\dicEntry[syði] \dicTerm{syði} \dicIPA{{s}{\textsci}{\textlengthmark}{ð}{\textsci}} \dicPos{v} \dicFlx{con pf sg 1 pers} \dicLink{sjóða}
\dicEntry[syðra] \dicTerm{syðra} \dicIPA{{s}{\textsci}{ð}{r}{a}} \dicPos{adv} \dicDirectTranslationCS{na jihu} \dicIndirectTranslationCS{(myšleno zvláště o~Reykjavíku)}
\dicEntry[syðri] \dicTerm{syðri} \dicIPA{{s}{\textsci}{ð}{r}{\textsci}} \dicPos{adj}[12] \dicFlx{comp (sup syðstur)}[6] \dicDirectTranslationCS{jižnější}
\dicEntry[syðst] \dicTerm{syðst} \dicIPA{{s}{\textsci}{ð}{s}{\textsubring{d}}} \dicPos{adv} \dicFlx{sup} \dicLink{suður\smash{\textsuperscript{2}}}
\dicEntry[syðstur] \dicTerm{syðstur} \dicIPA{{s}{\textsci}{ð}{s}{\textsubring{d}}{\textscy}{\textsubring{r}}} \dicPos{adj}[1]\dicFlx{}[-10] \dicFlx{m sg nom sup} \dicLink{syðri}
\dicEntry[syfja] \dicTerm{syfj|a\smash{\textsuperscript{1}}} \dicIPA{{s}{\textsci}{v}{j}{a}} \dicPos{f}[1] \dicFlx{(‑u)}[5] \dicDirectTranslationCS{ospalost}
\dicEntry[syfja] \dicTerm{syfj|a\smash{\textsuperscript{2}}} \dicIPA{{s}{\textsci}{v}{j}{a}} \dicPos{v}[1] \dicFlx{(‑aði)}[59] \dicFlx{impers} \dicPhraseIS{e‑n syfjar} \dicDirectTranslationCS{(kdo) je ospalý, (kdo) usíná}
\dicEntry[syfjaður] \dicTerm{syfj··|aður} \dicIPA{{s}{\textsci}{v}{j}{a}{ð}{\textscy}{\textsubring{r}}} \dicPos{adj}[3] \dicFlx{(f ‑uð)}[3] \dicDirectTranslationCS{ospalý}
\dicEntry[syfjulega] \dicTerm{syfju··lega} \dicIPA{{s}{\textsci}{v}{j}{\textscy}{l}{\textepsilon}{\textbabygamma}{a}} \dicPos{adv} \dicDirectTranslationCS{ospale, netečně}
\dicEntry[syfjulegur] \dicTerm{syfju··legur} \dicIPA{{s}{\textsci}{v}{j}{\textscy}{l}{\textepsilon}{\textbabygamma}{\textscy}{\textsubring{r}}} \dicPos{adj}[1]\dicFlx{}[-8] \dicDirectTranslationCS{ospalý, netečný}
\dicEntry[sygi] \dicTerm{sygi} \dicIPA{{s}{\textsci}{j}{\textlengthmark}{\textsci}} \dicPos{v} \dicFlx{con pf sg 1 pers} \dicLink{sjúga}
\dicEntry[sykki] \dicTerm{sykki} \dicIPA{{s}{\textsci}{h}{\r{g}}{\textsci}} \dicPos{v} \dicFlx{con pf sg 1 pers} \dicLink{sökkva\smash{\textsuperscript{2}}}
\dicEntry[sykra] \dicTerm{sykr|a\smash{\textsuperscript{1}}} \dicIPA{{s}{\textsci}{\textlengthmark}{\r{g}}{r}{a}} \dicPos{f}[1] \dicFlx{(‑u, ‑ur)}[7] \dicFieldCat{chem.} \dicDirectTranslationCS{sacharid}
\dicEntry[sykra] \dicTerm{sykr|a\smash{\textsuperscript{2}}} \dicIPA{{s}{\textsci}{\textlengthmark}{\r{g}}{r}{a}} \dicPos{v}[1] \dicFlx{(‑aði)}[34] \dicFlx{acc} \dicSynonym*{gera sætt} \dicDirectTranslationCS{(po)cukrovat, (o)sladit}
\dicEntry[sykraður] \dicTerm{sykr··|aður} \dicIPA{{s}{\textsci}{\textlengthmark}{\r{g}}{r}{a}{ð}{\textscy}{\textsubring{r}}} \dicPos{adj}[3] \dicFlx{(f ‑uð)}[3] \dicDirectTranslationCS{(po)cukrovaný, (o)slazený}
\dicEntry[sykur] \dicTerm{syk|ur} \dicIPA{{s}{\textsci}{\textlengthmark}{\r{g}}{\textscy}{\textsubring{r}}} \dicPos{m}[5] \dicFlx{(‑urs)}[2] \dicDirectTranslationCS{cukr}
\dicEntry[sykurhúð] \dicTerm{sykur··húð} \dicIPA{{s}{\textsci}{\textlengthmark}{\r{g}}{\textscy}{\textsubring{r}}{h}{u}{\texttheta}} \dicPos{f}[7] \dicFlx{(‑ar)}[3] \dicFieldCat{kulin.} \dicDirectTranslationCS{glazura, cukrová poleva}
\dicEntry[sykurkar] \dicTerm{sykur··|kar}\dicTerm{, sykurker} \dicIPA{{s}\-{\textsci}\-{\textlengthmark}\-{\r{g}}\-{\textscy}\-{\textsubring{r}}\-{k\smash{\textsuperscript{h}}}\-{a}\-{\textsubring{r}}\-} \dicPos{n}[2] \dicFlx{(‑kars, ‑kör)}[8] \dicDirectTranslationCS{cukřenka}
\dicEntry[sykurker] \dicTerm{sykur··ker} \dicIPA{{s}{\textsci}{\textlengthmark}{\r{g}}{\textscy}{\textsubring{r}}{c\smash{\textsuperscript{h}}}{\textepsilon}{\textsubring{r}}} \dicPos{n}[2] \dicFlx{(‑s, ‑)}[6] \dicLink{sykurkar}
\dicEntry[sykurlaus] \dicTerm{sykur··laus} \dicIPA{{s}{\textsci}{\textlengthmark}{\r{g}}{\textscy}{r}{l}{\oe i}{s}} \dicPos{adj}[5]\dicFlx{}[-1] \dicDirectTranslationCS{neslazený, neoslazený}
\dicEntry[sykurmoli] \dicTerm{sykur··mol|i} \dicIPA{{s}{\textsci}{\textlengthmark}{\r{g}}{\textscy}{r}{m}{\textopeno}{l}{\textsci}} \dicPos{m}[1] \dicFlx{(‑a, ‑ar)}[1] \dicDirectTranslationCS{kostka cukru}
\dicEntry[sykurreyr] \dicTerm{sykur··reyr} \dicIPA{{s}{\textsci}{\textlengthmark}{\r{g}}{\textscy}{r}{\textlengthmark}{ei}{\textsubring{r}}} \dicPos{m}[4] \dicFlx{(‑s)}[18] \dicFieldCat{bot.} \dicDirectTranslationCS{cukrová třtina} \textit{(l.~{\textLA{Saccharum officinarum}})}  \dicsymPhoto\ 
\dicFigure{50999.jpg}{Sykurreyr}{Sykurreyr - Forest \& Kim [[p:2684;Starr]], Biolib, CC-BY}
\dicEntry[sykurrófa] \dicTerm{sykur··róf|a} \dicIPA{{s}{\textsci}{\textlengthmark}{\r{g}}{\textscy}{r}{\textlengthmark}{ou}{v}{a}} \dicPos{f}[1] \dicFlx{(‑u, ‑ur)}[19] \dicFieldCat{bot.} \dicDirectTranslationCS{řepa obecná} \textit{(l.~{\textLA{Beta vulgaris}})}
\dicEntry[sykursjúkur] \dicTerm{sykur··sjúkur} \dicIPA{{s}{\textsci}{\textlengthmark}{\r{g}}{\textscy}{\textsubring{r}}{s}{j}{u}{\r{g}}{\textscy}{\textsubring{r}}} \dicPos{adj}[1]\dicFlx{}[-6] \dicFieldCat{med.} \dicDirectTranslationCS{nemocný cukrovkou, diabetický}
\dicEntry[sykursýki] \dicTerm{sykur··sýk|i} \dicIPA{{s}{\textsci}{\textlengthmark}{\r{g}}{\textscy}{\textsubring{r}}{s}{i}{\r{\textObardotlessj}}{\textsci}} \dicPos{f}[3] \dicFlx{(‑i)}[3] \dicFieldCat{med.} \dicDirectTranslationCS{cukrovka, diabetes}
\dicEntry[sylgja] \dicTerm{sylgj|a} \dicIPA{{s}{\textsci}{l}{\r{\textObardotlessj}}{a}} \dicPos{f}[1] \dicFlx{(‑u, ‑ur)}[25] \dicSynonym{spenna\smash{\textsuperscript{1}}} \dicDirectTranslationCS{přezka, spona (na opasku)} \dicExampleIS{Á beltinu er sylgja til að loka því.} \dicExampleCS{Na opasku je přezka k~zapnutí.}
\dicEntry[sylla] \dicTerm{syll|a} \dicIPA{{s}{\textsci}{\textsubring{d}}{l}{a}} \dicPos{f}[1] \dicFlx{(‑u, ‑ur)}[19] \textbf{1.} \dicSynonym{stallur} \dicDirectTranslationCS{římsa, výstupek, výklenek}  \textbf{2.} \dicFieldCat{stav.} \dicDirectTranslationCS{pozednice}
\dicEntry[sylli] \dicTerm{sylli} \dicIPA{{s}{\textsci}{\textsubring{d}}{l}{\textsci}} \dicPos{v} \dicFlx{con pf sg 1 pers} \dicLink{svella\smash{\textsuperscript{2}}}
\dicEntry[sylti] \dicTerm{sylti} \dicIPA{{s}{\textsci}{\textsubring{l}}{\textsubring{d}}{\textsci}} \dicPos{v} \dicFlx{con pf sg 1 pers} \dicLink{svelta\smash{\textsuperscript{2}}}
\dicEntry[symból] \dicTerm{symból} \dicIPA{{s}{\textsci}{m}{\textsubring{b}}{ou}{\textsubring{l}}} \dicPos{n}[2] \dicFlx{(‑s, ‑)}[5] \dicSynonym{tákn} \dicDirectTranslationCS{symbol}
\dicEntry[symbólismi] \dicTerm{symból··ism|i} \dicIPA{{s}{\textsci}{m}{\textsubring{b}}{ou}{l}{\textsci}{s}{m}{\textsci}} \dicPos{m}[1] \dicFlx{(‑a)}[3] \dicDirectTranslationCS{symbolismus}
\dicEntry[symbólisti] \dicTerm{symból··ist|i} \dicIPA{{s}{\textsci}{m}{\textsubring{b}}{ou}{l}{\textsci}{s}{\textsubring{d}}{\textsci}} \dicPos{m}[1] \dicFlx{(‑a, ‑ar)}[1] \dicDirectTranslationCS{symbolista, symbolistka}
\dicEntry[synd] \dicTerm{synd} \dicIPA{{s}{\textsci}{n}{\textsubring{d}}} \dicPos{f}[7] \dicFlx{(‑ar, ‑ir)}[1] \textbf{1.} \dicSynonym{yfirsjón} \dicDirectTranslationCS{provinění, prohřešek}  \textbf{2.} \dicDirectTranslationCS{hřích};  \dicPhraseIS{segja e‑m til syndanna} \dicDirectTranslationCS{vynadat (komu), vyhubovat (komu)};  \dicPhraseIS{það er synd að} \dicLangCat{přen.} \dicDirectTranslationCS{je to hřích, že}
\dicEntry[synda] \dicTerm{syn|da} \dicsymFrequent\  \dicIPA{{s}{\textsci}{n}{\textsubring{d}}{a}} \dicPos{v}[2] \dicFlx{(‑ti, ‑t)}[43] \dicDirectTranslationCS{plavat, zaplavat si} \dicExampleIS{synda í sjónum} \dicExampleCS{plavat v~moři};  \dicPhraseIS{synda e‑ð} \dicDirectTranslationCS{plavat (čím)} \dicExampleIS{synda skriðsund} \dicExampleCS{plavat kraulem}
\dicEntry[syndaaflát] \dicTerm{synda··af·lát} \dicIPA{{s}{\textsci}{n}{\textsubring{d}}{a}{a}{v}{l}{au}{\textsubring{d}}} \dicPos{n}[2] \dicFlx{(‑s, ‑)}[5] \dicFieldCat{náb.} \dicDirectTranslationCS{odpustek}
\dicEntry[syndaflóð] \dicTerm{synda··flóð} \dicIPA{{s}{\textsci}{n}{\textsubring{d}}{a}{f}{l}{ou}{\texttheta}} \dicPos{n}[2] \dicFlx{(‑s)}[2] \dicFieldCat{náb.} \dicDirectTranslationCS{potopa světa}
\dicEntry[syndari] \dicTerm{synd··ar|i} \dicIPA{{s}{\textsci}{n}{\textsubring{d}}{a}{r}{\textsci}} \dicPos{m}[1] \dicFlx{(‑a, ‑ar)}[13] \dicDirectTranslationCS{hříšník, hříšnice}
\dicEntry[syndaselur] \dicTerm{synda··sel|ur} \dicIPA{{s}{\textsci}{n}{\textsubring{d}}{a}{s}{\textepsilon}{l}{\textscy}{\textsubring{r}}} \dicPos{m}[9] \dicFlx{(‑s, ‑ir)}[8] \dicSynonym{þrjótur} \dicDirectTranslationCS{darebák, darebačka, ničema}
\dicEntry[syndga] \dicTerm{syndg|a} \dicIPA{{s}{\textsci}{\ng}{\r{g}}{a}} \dicPos{v}[1] \dicFlx{(‑aði)}[1] \dicDirectTranslationCS{(z)hřešit}
\dicEntry[syndlaus] \dicTerm{synd··laus} \dicIPA{{s}{\textsci}{n}{\textsubring{d}}{l}{\oe i}{s}} \dicPos{adj}[5]\dicFlx{}[-1] \dicDirectTranslationCS{(jsoucí) bez hříchu, nehřešící}
\dicEntry[syndsamlegur] \dicTerm{synd·sam··legur} \dicIPA{{s}{\textsci}{n}{\textsubring{d}}{s}{a}{m}{l}{\textepsilon}{\textbabygamma}{\textscy}{\textsubring{r}}} \dicPos{adj}[1]\dicFlx{}[-8] \dicDirectTranslationCS{hříšný, hřešící} \dicExampleIS{syndsamlegt athæfi} \dicExampleCS{hříšné chování}
\dicEntry[syndugur] \dicTerm{synd··ugur} \dicIPA{{s}{\textsci}{n}{\textsubring{d}}{\textscy}{\textbabygamma}{\textscy}{\textsubring{r}}} \dicPos{adj}[1]\dicFlx{}[-8] \dicDirectTranslationCS{hříšný}
\dicEntry[syndur] \dicTerm{syndur} \dicIPA{{s}{\textsci}{n}{\textsubring{d}}{\textscy}{\textsubring{r}}} \dicPos{adj}[2]\dicFlx{}[-14] \dicDirectTranslationCS{umějící plavat} \dicExampleIS{vera vel syndur} \dicExampleCS{umět dobře plavat}
\dicEntry[syng] \dicTerm{syng} \dicIPA{{s}{i}{\ng}{\r{g}}} \dicPos{v} \dicFlx{ind praes sg 1 pers} \dicLink{syngja}
\dicEntry[syngi] \dicTerm{syngi} \dicIPA{{s}{i}{\textltailn}{\r{\textObardotlessj}}{\textsci}} \dicPos{v} \dicFlx{con pf sg 1 pers} \dicLink{syngja}
\dicEntry[syngja] \dicTerm{syngja} \dicsymFrequent\  \dicIPA{{s}{i}{\textltailn}{\r{\textObardotlessj}}{a}} \dicPos{v}[6] \dicFlx{(syng, söng, sungum, syngi, sungið)}[92] \dicFlx{acc} \textbf{1.} \dicDirectTranslationCS{(za)zpívat, (za)pět} \dicIndirectTranslationCS{(o~člověku)} \dicExampleIS{syngja í kór} \dicExampleCS{zpívat ve sboru}  \textbf{2.} \dicDirectTranslationCS{(za)zpívat, (za)pět} \dicIndirectTranslationCS{(o~ptácích)};  \dicIdiom{syngja}[eftir]{ \dicPhraseIS{syngja eftir nótum}} {\textbf{a.}} \dicDirectTranslationCS{zpívat podle not};  {\textbf{b.}} \dicLangCat{přen.} \dicDirectTranslationCS{zpívat podle not} \dicIndirectTranslationCS{(dodržovat stanovená pravidla)};  \dicIdiom{syngja}[fyrir]{ \dicPhraseIS{syngja fyrir}} \dicDirectTranslationCS{předzpívávat};  \dicIdiom{syngja}{ \dicPhraseIS{syngja e‑n lof}} \dicDirectTranslationCS{pět (komu) chválu}; { \dicPhraseIS{hvað syngur í þér?}} \dicDirectTranslationCS{jak se daří?}
\dicEntry[syngjandi] \dicTerm{syngj··andi} \dicIPA{{s}{i}{\textltailn}{\r{\textObardotlessj}}{a}{n}{\textsubring{d}}{\textsci}} \dicPos{adj}[13] \dicFlx{indecl}[1] \dicDirectTranslationCS{zpívající, pějící}
\dicEntry[synir] \dicTerm{synir} \dicIPA{{s}{\textsci}{\textlengthmark}{n}{\textsci}{\textsubring{r}}} \dicPos{m} \dicFlx{pl nom} \dicLink{sonur}
\dicEntry[synja] \dicTerm{synj|a} \dicIPA{{s}{\textsci}{n}{j}{a}} \dicPos{v}[1] \dicFlx{(‑aði)}[1] \dicFlx{dat (+ gen)} \dicSynonym{neita} \dicDirectTranslationCS{odmítnout, odmítat, zamítnout, zamítat, odříct, odříkat};  \dicPhraseIS{synja e‑u} \dicDirectTranslationCS{odmítnout (co), zamítnout (co)};  \dicPhraseIS{synja e‑m e‑s} \dicDirectTranslationCS{odmítnout (komu co)} \dicExampleIS{synja honum leyfis} \dicExampleCS{zamítnout mu povolení};  \dicPhraseIS{synja e‑m um e‑ð} \dicDirectTranslationCS{odmítnout (komu co), odříct (komu co)};  \dicPhraseIS{synja fyrir e‑ð} \dicSynonym{afneita} \dicDirectTranslationCS{odmítnout (co), popřít (co)} \dicExampleIS{synja fyrir barn} \dicExampleCS{odmítnout dítě}
\dicEntry[synjun] \dicTerm{synj|un} \dicIPA{{s}{\textsci}{n}{j}{\textscy}{\textsubring{n}}} \dicPos{f}[7] \dicFlx{(‑unar, ‑anir)}[8] \dicSynonym{neitun} \dicDirectTranslationCS{odmítnutí, zamítnutí, odřeknutí} \dicExampleIS{synjun við beiðninni} \dicExampleCS{odmítnutí žádosti}
\dicEntry[sypi] \dicTerm{sypi} \dicIPA{{s}{\textsci}{\textlengthmark}{\textsubring{b}}{\textsci}} \dicPos{v} \dicFlx{con pf sg 1 pers} \dicLink{súpa\smash{\textsuperscript{2}}}
\dicEntry[syrfi] \dicTerm{syrfi} \dicIPA{{s}{\textsci}{r}{v}{\textsci}} \dicPos{v} \dicFlx{con pf sg 1 pers} \dicLink{sverfa}
\dicEntry[syrgja] \dicTerm{syrg|ja} \dicIPA{{s}{\textsci}{r}{\r{\textObardotlessj}}{a}} \dicPos{v}[2] \dicFlx{(‑ði, ‑t)}[94] \dicFlx{acc} \dicSynonym{harma} \dicDirectTranslationCS{oplakávat, truchlit} \dicExampleIS{syrgja ástvin sinn} \dicExampleCS{truchlit nad svým přítelem}
\dicEntry[syrgjandi] \dicTerm{syrgj··|andi\smash{\textsuperscript{1}}} \dicIPA{{s}{\textsci}{r}{\r{\textObardotlessj}}{a}{n}{\textsubring{d}}{\textsci}} \dicPos{m}[2] \dicFlx{(‑anda, ‑endur)}[1] \dicDirectTranslationCS{truchlící, smuteční host} \dicExampleIS{Dómkirkjan var full af syrgjendum.} \dicExampleCS{Katedrála byla plná truchlících.}
\dicEntry[syrjgandi] \dicTerm{syrgj··andi\smash{\textsuperscript{2}}} \dicIPA{{s}{\textsci}{r}{\r{\textObardotlessj}}{a}{n}{\textsubring{d}}{\textsci}} \dicPos{adj}[13] \dicFlx{indecl}[1] \dicDirectTranslationCS{truchlící, oplakávající}
\dicEntry[syrpa] \dicTerm{syrp|a} \dicIPA{{s}{\textsci}{\textsubring{r}}{\textsubring{b}}{a}} \dicPos{f}[1] \dicFlx{(‑u, ‑ur)}[19] \textbf{1.} \dicSynonym*{röð af e‑u} \dicDirectTranslationCS{výběr, směsice, směs, mix} \dicExampleIS{lagasyrpa} \dicExampleCS{směsice písniček}  \textbf{2.} \dicSynonym{þáttaröð} \dicDirectTranslationCS{televizní série\,/\addthin seriál}
\dicEntry[systir] \dicTerm{syst|ir} \dicsymFrequent\  \dicIPA{{s}{\textsci}{s}{\textsubring{d}}{\textsci}{\textsubring{r}}} \dicPos{f}[11] \dicFlx{(‑ur, ‑ur)}[3] \dicDirectTranslationCS{sestra} \dicExampleIS{Hún var elst fjögurra systra.} \dicExampleCS{Byla ze čtyř sester nejstarší.}
\dicEntry[systkin] \dicTerm{syst··kin} \dicsymFrequent\  \dicIPA{{s}{\textsci}{s}{c\smash{\textsuperscript{h}}}{\textsci}{\textsubring{n}}} \dicPos{n}[2] \dicFlx{pl}[1] \dicLink{systkini}
\dicEntry[systkinabörn] \dicTerm{syst·kina··börn} \dicIPA{{s}{\textsci}{s}{c\smash{\textsuperscript{h}}}{\textsci}{n}{a}{\textsubring{b}}{\oe}{r}{\textsubring{d}}{\textsubring{n}}} \dicPos{n}[2] \dicFlx{pl}[9] \dicDirectTranslationCS{synovci, neteře} \dicIndirectTranslationCS{(děti sourozence\,/\addthin sourozenců)}
\dicEntry[systkini] \dicTerm{syst··kini}\dicTerm{, systkin} \dicIPA{{s}\-{\textsci}\-{s}\-{c\smash{\textsuperscript{h}}}\-{\textsci}\-{n}\-{\textsci}\-} \dicPos{n}[2] \dicFlx{pl}[19] \dicDirectTranslationCS{sourozenci}
\dicEntry[systrabörn] \dicTerm{systra··börn} \dicIPA{{s}{\textsci}{s}{\textsubring{d}}{r}{a}{\textsubring{b}}{\oe}{r}{\textsubring{d}}{\textsubring{n}}} \dicPos{n}[2] \dicFlx{pl}[9] \dicDirectTranslationCS{synovci, neteře} \dicIndirectTranslationCS{(děti sestry)}
\dicEntry[systur] \dicTerm{systur} \dicIPA{{s}{\textsci}{s}{\textsubring{d}}{\textscy}{\textsubring{r}}} \dicPos{f} \dicFlx{sg gen} \dicLink{systir}
\dicEntry[systurdóttir] \dicTerm{systur··|dóttir} \dicIPA{{s}{\textsci}{s}{\textsubring{d}}{\textscy}{r}{\textsubring{d}}{ou}{h}{\textsubring{d}}{\textsci}{\textsubring{r}}} \dicPos{f}[11] \dicFlx{(‑dóttur, ‑dætur)}[2] \dicDirectTranslationCS{neteř} \dicIndirectTranslationCS{(dcera sestry)} \dicAntonym{bróðurdóttir}
\dicEntry[systurfélag] \dicTerm{systur··fé·|lag} \dicIPA{{s}{\textsci}{s}{\textsubring{d}}{\textscy}{\textsubring{r}}{f}{j}{\textepsilon}{l}{a}{x}} \dicPos{n}[2] \dicFlx{(‑lags, ‑lög)}[8] \dicDirectTranslationCS{sesterská společnost\,/\addthin organizace}
\dicEntry[systursonur] \dicTerm{systur··|sonur} \dicIPA{{s}{\textsci}{s}{\textsubring{d}}{\textscy}{\textsubring{r}}{s}{\textopeno}{n}{\textscy}{\textsubring{r}}} \dicPos{m}[12] \dicFlx{(‑sonar, ‑synir)}[9] \dicDirectTranslationCS{synovec} \dicIndirectTranslationCS{(syn sestry)} \dicAntonym{bróðursonur}
\dicEntry[sýð] \dicTerm{sýð} \dicIPA{{s}{i}{\textlengthmark}{\texttheta}} \dicPos{v} \dicFlx{ind praes sg 1 pers} \dicLink{sjóða}
\dicEntry[sýfilis] \dicTerm{sýfilis} \dicIPA{{s}{i}{\textlengthmark}{f}{\textsci}{l}{\textsci}{s}} \dicPos{m}[4] \dicFlx{(‑s)}[18] \dicFieldCat{med.} \dicSynonym{sárasótt} \dicDirectTranslationCS{příjice, syfilis}
\dicEntry[sýg] \dicTerm{sýg} \dicIPA{{s}{i}{\textlengthmark}{x}} \dicPos{v} \dicFlx{ind praes sg 1 pers} \dicLink{sjúga}
\dicEntry[sýki] \dicTerm{sýk|i} \dicIPA{{s}{i}{\textlengthmark}{\r{\textObardotlessj}}{\textsci}} \dicPos{f}[3] \dicFlx{(‑i, ‑ir)}[2] \dicSynonym{pest} \dicDirectTranslationCS{choroba, nákaza} \dicExampleIS{Það kom upp sýki í fénu.} \dicExampleCS{Mezi ovcemi se objevila nákaza.}
\dicEntry[sýkill] \dicTerm{sýk|ill} \dicIPA{{s}{i}{\textlengthmark}{\r{\textObardotlessj}}{\textsci}{\textsubring{d}}{\textsubring{l}}} \dicPos{m}[6] \dicFlx{(‑ils, ‑lar)}[35] \dicFieldCat{med.} \dicDirectTranslationCS{patogen, patogenní agens, choroboplodný zárodek}
\dicEntry[sýking] \dicTerm{sýk··ing} \dicIPA{{s}{i}{\textlengthmark}{\r{\textObardotlessj}}{i}{\ng}{\r{g}}} \dicPos{f}[4] \dicFlx{(‑ar, ‑ar)}[5] \textbf{1.} \dicFieldCat{med.} \dicDirectTranslationCS{nákaza, infekce} \dicExampleIS{fá sýkingu í sárið} \dicExampleCS{dostat do rány infekci}  \textbf{2.} \dicSynonym*{það að smitast} \dicDirectTranslationCS{nakažení (se), infikování (se)}
\dicEntry[sýkja] \dicTerm{sýk|ja} \dicIPA{{s}{i}{\textlengthmark}{\r{\textObardotlessj}}{a}} \dicPos{v}[2] \dicFlx{(‑ti, ‑t)}[25] \dicFlx{acc} \dicSynonym{smita} \dicDirectTranslationCS{nakazit, infikovat} \dicExampleIS{sýkja e‑n} \dicExampleCS{nakazit (koho)};  \dicIdiom{sýkjast}{ \dicPhraseIS{sýkjast}} \dicFlx{refl} \dicSynonym*{smitast} \dicDirectTranslationCS{nakazit se, infikovat se} \dicExampleIS{sýkjast af berklum} \dicExampleCS{nakazit se tuberkulózou}
\dicEntry[sýkla] \dicTerm{sýkl|a} \dicIPA{{s}{i}{h}{\r{g}}{l}{a}} \dicPos{f}[1] \dicFlx{(‑u)}[5] \dicFieldCat{biol.} \dicDirectTranslationCS{plak}
\dicEntry[sýklahernaður] \dicTerm{sýkla··her·nað|ur} \dicIPA{{s}{i}{h}{\r{g}}{l}{a}{h}{\textepsilon}{r}{\textsubring{d}}{n}{a}{ð}{\textscy}{\textsubring{r}}} \dicPos{m}[10] \dicFlx{(‑ar)}[9] \dicDirectTranslationCS{bio\-lo\-gic\-ká válka}
\dicEntry[sýklalyf] \dicTerm{sýkla··lyf} \dicIPA{{s}{i}{h}{\r{g}}{l}{a}{l}{\textsci}{f}} \dicPos{n}[2] \dicFlx{(‑s, ‑)}[13] \dicFieldCat{med.} \dicDirectTranslationCS{antibiotikum}
\dicEntry[sýklavopn] \dicTerm{sýkla··vopn} \dicIPA{{s}{i}{h}{\r{g}}{l}{a}{v}{\textopeno}{h}{\textsubring{b}}{\textsubring{n}}} \dicPos{n}[2] \dicFlx{(‑s, ‑)}[5] \dicDirectTranslationCS{biologická zbraň}
\dicEntry[sýkn] \dicTerm{sýkn} \dicIPA{{s}{i}{h}{\r{g}}{\textsubring{n}}} \dicPos{adj}[5]\dicFlx{}[-5] \dicSynonym{saklaus} \dicDirectTranslationCS{nevinný, zproštěný viny};  \dicPhraseIS{dæma e‑n sýkna saka} \dicDirectTranslationCS{zprostit (koho) viny};  \dicPhraseIS{sýknt og heilagt} \dicFlx{adv} \dicDirectTranslationCS{celý čas, pořád, v~jednom kuse};  \dicPhraseIS{vera sýkn saka} \dicDirectTranslationCS{být zproštěn viny}
\dicEntry[sýkna] \dicTerm{sýkn|a\smash{\textsuperscript{1}}} \dicIPA{{s}{i}{h}{\r{g}}{n}{a}} \dicPos{f}[1] \dicFlx{(‑u)}[5] \dicFieldCat{práv.} \dicDirectTranslationCS{nevina};  \dicPhraseIS{sanna sýknu sína} \dicDirectTranslationCS{dokázat svoji nevinu}
\dicEntry[sýkna] \dicTerm{sýkn|a\smash{\textsuperscript{2}}} \dicIPA{{s}{i}{h}{\r{g}}{n}{a}} \dicPos{v}[1] \dicFlx{(‑aði)}[1] \dicFlx{acc} \dicFieldCat{práv.} \dicDirectTranslationCS{zprostit viny, ospravedlnit} \dicExampleIS{sýkna e‑n af e‑u} \dicExampleCS{ospravedlnit (koho) z~(čeho)}
\dicEntry[sýknudómur] \dicTerm{sýknu··dóm|ur} \dicIPA{{s}{i}{h}{\r{g}}{n}{\textscy}{\textsubring{d}}{ou}{m}{\textscy}{\textsubring{r}}} \dicPos{m}[6] \dicFlx{(‑s, ‑ar)}[10] \dicFieldCat{práv.} \dicDirectTranslationCS{zproštění viny}
\dicEntry[sýknun] \dicTerm{sýkn|un} \dicIPA{{s}{i}{h}{\r{g}}{n}{\textscy}{\textsubring{n}}} \dicPos{f}[7] \dicFlx{(‑unar, ‑anir)}[8] \dicDirectTranslationCS{zproštění viny, ospravedlnění}
\dicEntry[sýn] \dicTerm{sýn} \dicsymFrequent\  \dicIPA{{s}{i}{\textlengthmark}{\textsubring{n}}} \dicPos{f}[7] \dicFlx{(‑ar, ‑ir)}[1] \textbf{1.} \dicSynonym{sjón} \dicDirectTranslationCS{vidění, zrak}  \textbf{2.} \dicSynonym*{það sem sést} \dicDirectTranslationCS{pohled, podívaná, výhled} \dicExampleIS{undarleg sýn} \dicExampleCS{kouzelný pohled};  \dicPhraseIS{e‑að byrgir e‑m sýn} \dicDirectTranslationCS{(co komu) brání ve výhledu};  \dicPhraseIS{vera fríður sýnum} \dicDirectTranslationCS{být krásný (na pohled)};  \dicPhraseIS{við fyrstu sýn} \dicFlx{adv} \dicDirectTranslationCS{na první pohled}  \textbf{3.} \dicSynonym{skilningur} \dicDirectTranslationCS{pohled, hledisko, úhel pohledu};  \dicPhraseIS{fá nýja sýn á e‑ð} \dicDirectTranslationCS{získat nový pohled na (co)}  \textbf{4.} \dicSynonym{vitrun} \dicDirectTranslationCS{vidění, vidina};  \dicPhraseIS{sjá sýnir} \dicDirectTranslationCS{mít vidiny}
\dicEntry[sýna] \dicTerm{sýn|a} \dicsymFrequent\  \dicIPA{{s}{i}{\textlengthmark}{n}{a}} \dicPos{v}[2] \dicFlx{(‑di, ‑t)}[144] \dicFlx{(dat +) acc} \textbf{1.} \dicFlx{acc} \dicSynonym*{sýna verk} \dicDirectTranslationCS{vystavit, vystavovat (na výstavě ap.)} \dicExampleIS{sýna málverk eftir sig á safninu} \dicExampleCS{vystavovat svůj obraz v~muzeu}  \textbf{2.} \dicFlx{dat + acc} \dicDirectTranslationCS{ukázat, ukazovat, předvést, předvádět, demonstrovat} \dicExampleIS{sýna henni myndina} \dicExampleCS{ukázat jí obraz}  \textbf{3.} \dicFlx{dat + acc} \dicDirectTranslationCS{projevit, projevovat, ukázat, ukazovat, dát\,/\addthin dávat najevo, vykazovat} \dicExampleIS{sýna henni kurteisi} \dicExampleCS{projevit jí zdvořilost}  \textbf{4.} \dicFlx{acc} \dicDirectTranslationCS{vysílat, promítat (film ap.)};  \dicIdiom{sýna}[fram]{ \dicPhraseIS{sýna fram á e‑ð}} \dicSynonym{sanna} \dicDirectTranslationCS{dokázat (co), prokázat (co)};  \dicIdiom{sýna}[í]{ \dicPhraseIS{sýna e‑m í tvo heimana}} \dicLangCat{přen.} \dicDirectTranslationCS{dát (komu) nůž na krk} \dicIndirectTranslationCS{(dát někomu drastické ultimátum)};  \dicIdiom{sýnast}{ \dicPhraseIS{e‑r\,/\addthin e‑að sýnist}} \dicFlx{refl} \dicDirectTranslationCS{(kdo\,/\addthin co) vypadá, (kdo\,/\addthin co) se jeví}; { \dicPhraseIS{e‑m sýnist}} \dicFlx{refl impers} {\textbf{a.}} \dicDirectTranslationCS{(komu) se zdá, (komu) připadá} \dicExampleIS{mér sýnist að} \dicExampleCS{zdá se mi, že};  {\textbf{b.}} \dicDirectTranslationCS{(komu) se líbí, (komu) se zachce (udělat něco ap.)} \dicExampleIS{Ég geri það sem mér sýnist.} \dicExampleCS{Dělám to, co se mi líbí.}; { \dicPhraseIS{ekki er allt sem sýnist}} \dicFlx{refl} \dicDirectTranslationCS{není všechno tak, jak to vypadá}; { \dicPhraseIS{það sýnist sitt hverjum}} \dicFlx{refl impers} \dicLangCat{přen.} \dicDirectTranslationCS{tisíc lidí, tisíc názorů; sto lidí, sto chutí};  \dicPhraseIS{það sýnist} \dicFlx{refl impers} \dicDirectTranslationCS{zdá se} \dicExampleIS{Það er ekki eins auðvelt og það sýnist.} \dicExampleCS{Není to tak jednoduché, jak se zdá.}
\dicEntry[sýnagóga] \dicTerm{sýnagóg|a} \dicIPA{{s}{i}{\textlengthmark}{n}{a}{\r{g}}{ou}{\r{g}}{a}} \dicPos{f}[1] \dicFlx{(‑u, ‑ur)}[7] \dicSynonym{samkunduhús} \dicDirectTranslationCS{synagoga}
\dicEntry[sýndarmennska] \dicTerm{sýndar··mennsk|a} \dicIPA{{s}{i}{n}{\textsubring{d}}{a}{r}{m}{\textepsilon}{n}{s}{\r{g}}{a}} \dicPos{f}[1] \dicFlx{(‑u)}[5] \dicSynonym{yfirborðsmennska} \dicDirectTranslationCS{přetvářka, předstírání}
\dicEntry[sýndarveruleiki] \dicTerm{sýndar··veru·leik|i} \dicIPA{{s}{i}{n}{\textsubring{d}}{a}{r}{v}{\textepsilon}{r}{\textscy}{l}{ei}{\r{\textObardotlessj}}{\textsci}} \dicPos{m}[1] \dicFlx{(‑a)}[3] \dicFieldCat{poč.} \dicDirectTranslationCS{virtuální realita}
\dicEntry[sýni] \dicTerm{sýni} \dicsymFrequent\  \dicIPA{{s}{i}{\textlengthmark}{n}{\textsci}} \dicPos{n}[2] \dicFlx{(‑s, ‑)}[14] \textbf{1.} \dicSynonym{sýnishorn} \dicDirectTranslationCS{vzorek, ukázka} \dicExampleIS{taka sýni af blóði} \dicExampleCS{vzít vzorek krve}  \textbf{2.} \dicSynonym{sýning} \dicDirectTranslationCS{prezentace, představení} \dicExampleIS{hafa e‑ð til sýnis} \dicExampleCS{mít (co) k~představení}
\dicEntry[sýnilega] \dicTerm{sýni··lega} \dicIPA{{s}{i}{\textlengthmark}{n}{\textsci}{l}{\textepsilon}{\textbabygamma}{a}} \dicPos{adv} \dicDirectTranslationCS{viditelně, očividně, zjevně}
\dicEntry[sýnilegur] \dicTerm{sýni··legur} \dicsymFrequent\  \dicIPA{{s}{i}{\textlengthmark}{n}{\textsci}{l}{\textepsilon}{\textbabygamma}{\textscy}{\textsubring{r}}} \dicPos{adj}[1]\dicFlx{}[-8] \dicSynonym{augljós} \dicDirectTranslationCS{viditelný, očividný, zjevný} \dicExampleIS{sýnilegir fordómar} \dicExampleCS{zjevné předsudky}
\dicEntry[sýning] \dicTerm{sýn··ing} \dicsymFrequent\  \dicIPA{{s}{i}{\textlengthmark}{n}{i}{\ng}{\r{g}}} \dicPos{f}[4] \dicFlx{(‑ar, ‑ar)}[5] \textbf{1.} \dicDirectTranslationCS{ukázání, předvedení, demonstrace}  \textbf{2.} \dicSynonym{leiksýning} \dicDirectTranslationCS{představení, vystoupení (divadelní ap.)} \dicExampleIS{Sú sýning var bönnuð stuttu síðar.} \dicExampleCS{To představení bylo krátce na to zakázané.}  \textbf{3.} \dicSynonym*{málverkasýning} \dicDirectTranslationCS{výstava, expozice}  \textbf{4.} \dicSynonym{útsending} \dicDirectTranslationCS{vysílání, promítání, projekce}
\dicEntry[sýningargluggi] \dicTerm{sýningar··glugg|i} \dicIPA{{s}{i}{\textlengthmark}{n}{i}{\ng}{\r{g}}{a}{r}{\r{g}}{l}{\textscy}{\r{\textObardotlessj}}{\textsci}} \dicPos{m}[1] \dicFlx{(‑a, ‑ar)}[1] \dicDirectTranslationCS{výloha, výklad, vitrína}
\dicEntry[sýningarskrá] \dicTerm{sýningar··skrá} \dicIPA{{s}{i}{\textlengthmark}{n}{i}{\ng}{\r{g}}{a}{\textsubring{r}}{s}{\r{g}}{r}{au}} \dicPos{f}[4] \dicFlx{(‑r\,/\addthin ‑ar, ‑r)}[21] \dicDirectTranslationCS{katalog výstavy}
\dicEntry[sýningarstjóri] \dicTerm{sýningar··stjór|i} \dicIPA{{s}{i}{\textlengthmark}{n}{i}{\ng}{\r{g}}{a}{\textsubring{r}}{s}{\textsubring{d}}{j}{ou}{r}{\textsci}} \dicPos{m}[1] \dicFlx{(‑a, ‑ar)}[1] \textbf{1.} \dicDirectTranslationCS{vedoucí produkce (v~divadle)}  \textbf{2.} \dicDirectTranslationCS{kurátor(ka), ředitel(ka) výstavy}
\dicEntry[sýningarsvæði] \dicTerm{sýningar··svæði} \dicIPA{{s}{i}{\textlengthmark}{n}{i}{\ng}{\r{g}}{a}{\textsubring{r}}{s}{v}{a}{i}{ð}{\textsci}} \dicPos{n}[2] \dicFlx{(‑s, ‑)}[14] \dicDirectTranslationCS{výstaviště, výstavní areál\,/\addthin prostor}
\dicEntry[sýningarvél] \dicTerm{sýningar··vél} \dicIPA{{s}{i}{\textlengthmark}{n}{i}{\ng}{\r{g}}{a}{r}{v}{j}{\textepsilon}{\textsubring{l}}} \dicPos{f}[4] \dicFlx{(‑ar, ‑ar)}[1] \dicDirectTranslationCS{projektor, promítačka, promítací přístroj}
\dicEntry[sýningarvíxill] \dicTerm{sýningar··víx|ill} \dicIPA{{s}{i}{\textlengthmark}{n}{i}{\ng}{\r{g}}{a}{r}{v}{i}{x}{s}{\textsci}{\textsubring{d}}{\textsubring{l}}} \dicPos{m}[6] \dicFlx{(‑ils, ‑lar)}[35] \dicFieldCat{ekon.} \dicDirectTranslationCS{vista směnka, směnka na viděnou}
\dicEntry[sýnisbók] \dicTerm{sýnis··|bók} \dicIPA{{s}{i}{\textlengthmark}{n}{\textsci}{s}{\textsubring{b}}{ou}{\r{g}}} \dicPos{f}[8] \dicFlx{(‑bókar, ‑bækur)}[5] \dicDirectTranslationCS{antologie, výbor, výběr}
\dicEntry[sýnishorn] \dicTerm{sýnis··horn} \dicIPA{{s}{i}{\textlengthmark}{n}{\textsci}{s}{h}{\textopeno}{r}{\textsubring{d}}{\textsubring{n}}} \dicPos{n}[2] \dicFlx{(‑s, ‑)}[5] \dicDirectTranslationCS{ukázka, vzorek (výrobku ap.)} \dicExampleIS{sýnishorn af e‑u} \dicExampleCS{ukázka (čeho)}
\dicEntry[sýnn] \dicTerm{sýnn} \dicIPA{{s}{i}{\textlengthmark}{\textsubring{d}}{\textsubring{n}}} \dicPos{adj}[7]\dicFlx{}[-1] \textbf{1.} \dicSynonym{augljós} \dicDirectTranslationCS{viditelný, zjevný, jasný} \dicExampleIS{Það er sýnt.} \dicExampleCS{Je to jasné.}  \textbf{2.} \dicSynonym{vænlegur} \dicDirectTranslationCS{příhodný, slibný};  \dicPhraseIS{e‑m er sýnt um e‑ð} \dicFlx{impers} \dicDirectTranslationCS{(kdo) je v~(čem) dobrý, (komu) se (co) daří} \dicExampleIS{Honum er sýnt um að fá  fólk til að vinna saman.} \dicExampleCS{Daří se mu přimět lidi ke spolupráci.}  \textbf{3.} \dicSynonym*{sjáandi} \dicDirectTranslationCS{vidící, schopný vidět} \dicExampleIS{sýnn á báðum augum} \dicExampleCS{vidící na obě oči}
\dicEntry[sýp] \dicTerm{sýp} \dicIPA{{s}{i}{\textlengthmark}{\textsubring{b}}} \dicPos{v} \dicFlx{ind praes sg 1 pers} \dicLink{súpa\smash{\textsuperscript{2}}}
\dicEntry[sýprus] \dicTerm{sýprus} \dicIPA{{s}{i}{\textlengthmark}{\textsubring{b}}{r}{\textscy}{s}} \dicPos{m}[4] \dicFlx{(‑s, ‑ar)}[25] \dicLink{kýprusviður}
\dicEntry[sýra] \dicTerm{sýr|a\smash{\textsuperscript{1}}} \dicIPA{{s}{i}{\textlengthmark}{r}{a}} \dicPos{f}[1] \dicFlx{(‑u, ‑ur)}[7] \textbf{1.} \dicFieldCat{chem.} \dicDirectTranslationCS{kyselina} \dicAntonym{basi}  \textbf{2.} \dicFieldCat{kulin.} \dicSynonym*{súr mysa} \dicDirectTranslationCS{syrovátka} \dicExampleIS{súrsa mat í sýru} \dicExampleCS{naložit potraviny do syrovátky}  \textbf{3.} \dicLangCat{slang.} \dicDirectTranslationCS{LSD}
\dicEntry[sýra] \dicTerm{sýr|a\smash{\textsuperscript{2}}} \dicIPA{{s}{i}{\textlengthmark}{r}{a}} \dicPos{v}[2] \dicFlx{(‑ði, ‑t)}[99] \dicDirectTranslationCS{okyselit, okyselovat}
\dicEntry[sýrður] \dicTerm{sýrður} \dicIPA{{s}{i}{r}{ð}{\textscy}{\textsubring{r}}} \dicPos{adj}[2]\dicFlx{}[-1] \dicDirectTranslationCS{okyselený, (za)kysaný}
\dicEntry[sýrena] \dicTerm{sýren|a} \dicIPA{{s}{i}{\textlengthmark}{r}{\textepsilon}{n}{a}} \dicPos{f}[1] \dicFlx{(‑u, ‑ur)}[7] \dicFieldCat{bot.} \dicDirectTranslationCS{šeřík} \textit{(l.~{\textLA{Syringa}})}  \dicsymPhoto\ 
\dicFigure{ds_image_syrena_0_1.jpg}{Sýrena}{Sýrena - Marisa DeMeglio, CC BY 2.0}
\dicEntry[Sýrland] \dicTerm{Sýr··land} \dicIPA{{s}{i}{r}{l}{a}{n}{\textsubring{d}}} \dicPos{n}[2] \dicFlx{(‑s)}[4] \dicFieldCat{geog.} \dicDirectTranslationCS{Sýrie}
\dicEntry[Sýrlendingur] \dicTerm{Sýr·lend··ing|ur} \dicIPA{{s}{i}{r}{l}{\textepsilon}{n}{\textsubring{d}}{i}{\ng}{\r{g}}{\textscy}{\textsubring{r}}} \dicPos{m}[6] \dicFlx{(‑s, ‑ar)}[8] \dicDirectTranslationCS{Syřan(ka)}
\dicEntry[sýrlenska] \dicTerm{sýr··lensk|a} \dicIPA{{s}{i}{r}{l}{\textepsilon}{n}{s}{\r{g}}{a}} \dicPos{f}[1] \dicFlx{(‑u)}[5] \dicDirectTranslationCS{syrština}
\dicEntry[sýrlenskur] \dicTerm{sýr··lenskur} \dicIPA{{s}{i}{r}{l}{\textepsilon}{n}{s}{\r{g}}{\textscy}{\textsubring{r}}} \dicPos{adj}[1]\dicFlx{}[-1] \dicDirectTranslationCS{syrský}
\dicEntry[sýsla] \dicTerm{sýsl|a\smash{\textsuperscript{1}}} \dicIPA{{s}{i}{s}{\textsubring{d}}{l}{a}} \dicPos{f}[1] \dicFlx{(‑u, ‑ur)}[13] \textbf{1.} \dicFieldCat{hist.} \dicDirectTranslationCS{správní oblast, okres}  \textbf{2.} \dicDirectTranslationCS{(geografická) oblast, region}
\dicEntry[sýsla] \dicTerm{sýsl|a\smash{\textsuperscript{2}}} \dicIPA{{s}{i}{s}{\textsubring{d}}{l}{a}} \dicPos{v}[1] \dicFlx{(‑aði)}[1] \dicDirectTranslationCS{zabývat se, zaobírat se, věnovat se} \dicExampleIS{sýsla við e‑ð} \dicExampleCS{zaobírat se (čím)}
\dicEntry[sýslumaður] \dicTerm{sýslu··|maður} \dicsymFrequent\  \dicIPA{{s}{i}{s}{\textsubring{d}}{l}{\textscy}{m}{a}{ð}{\textscy}{\textsubring{r}}} \dicPos{m}[13] \dicFlx{(‑manns, ‑menn)}[2] \textbf{1.} \dicFieldCat{práv.} \dicDirectTranslationCS{okresní soudce\,/\addthin soudkyně} \dicExampleIS{bera mál undir sýslumann} \dicExampleCS{předložit případ okresnímu soudci}  \textbf{2.} \dicFieldCat{hist.} \dicDirectTranslationCS{okresní správce\,/\addthin správkyně}
\dicEntry[sæ.] \dicTerm{sæ.} \dicPos{zkr} \dicPhraseIS{sænska} \dicDirectTranslationCS{švédština}
\dicEntry[sæbjúga] \dicTerm{sæ··bjúg|a} \dicIPA{{s}{a}{i}{\textlengthmark}{\textsubring{b}}{j}{u}{a}} \dicPos{f}[1] \dicFlx{(‑u, ‑ur)}[13] \dicFieldCat{zool.} \dicDirectTranslationCS{sumýš} \textit{(l.~{\textLA{Holothuroidea}})}
\dicEntry[sædýr] \dicTerm{sæ··dýr} \dicIPA{{s}{a}{i}{\textlengthmark}{\textsubring{d}}{i}{\textsubring{r}}} \dicPos{n}[2] \dicFlx{(‑s, ‑)}[5] \dicDirectTranslationCS{mořský živočich}
\dicEntry[sædýrasafn] \dicTerm{sæ·dýra··|safn} \dicIPA{{s}{a}{i}{\textlengthmark}{\textsubring{d}}{i}{r}{a}{s}{a}{\textsubring{b}}{\textsubring{n}}} \dicPos{n}[2] \dicFlx{(‑safns, ‑söfn)}[8] \dicDirectTranslationCS{oceánografické\,/\addthin podmořské muzeum}
\dicEntry[sæða] \dicTerm{sæ|ða} \dicIPA{{s}{a}{i}{\textlengthmark}{ð}{a}} \dicPos{v}[2] \dicFlx{(‑ddi, ‑tt)}[167] \dicFlx{acc} \dicDirectTranslationCS{inseminovat} \dicExampleIS{sæða kúna} \dicExampleCS{inseminovat krávu}
\dicEntry[sæði] \dicTerm{sæði} \dicIPA{{s}{a}{i}{\textlengthmark}{ð}{\textsci}} \dicPos{n}[2] \dicFlx{(‑s, ‑)}[14] \dicFieldCat{biol.} \dicDirectTranslationCS{sperma, semeno}
\dicEntry[sæðing] \dicTerm{sæð··ing} \dicIPA{{s}{a}{i}{\textlengthmark}{ð}{i}{\ng}{\r{g}}} \dicPos{f}[4] \dicFlx{(‑ar)}[7] \dicDirectTranslationCS{inseminace} \dicExampleIS{sæðing búfjár} \dicExampleCS{inseminace dobytka}
\dicEntry[sæðisfruma] \dicTerm{sæðis··frum|a} \dicIPA{{s}{a}{i}{\textlengthmark}{ð}{\textsci}{s}{f}{r}{\textscy}{m}{a}} \dicPos{f}[1] \dicFlx{(‑u, ‑ur)}[19] \dicFieldCat{biol.} \dicSynonym{sáðfruma} \dicDirectTranslationCS{spermie, spermatická buňka}
\dicEntry[sæfari] \dicTerm{sæ··far|i} \dicIPA{{s}{a}{i}{\textlengthmark}{f}{a}{r}{\textsci}} \dicPos{m}[1] \dicFlx{(‑a, ‑ar)}[8] \dicLangCat{kniž.} \dicDirectTranslationCS{námořník, námořnice}
\dicEntry[sæfíll] \dicTerm{sæ··fíl|l} \dicIPA{{s}{a}{i}{\textlengthmark}{f}{i}{\textsubring{d}}{\textsubring{l}}} \dicPos{m}[6] \dicFlx{(‑s, ‑ar)}[48] \dicFieldCat{zool.} \dicDirectTranslationCS{rypouš severní} \textit{(l.~{\textLA{Mirounga angustirostris}})}  \dicsymPhoto\ 
\dicFigure{ds_image_saefill_0_1.jpg}{Sæfíll}{Sæfíll - Frank Schulenburg, CC BY-SA 3.0}
\dicEntry[sægur] \dicTerm{sæg|ur} \dicIPA{{s}{a}{i}{\textlengthmark}{\textbabygamma}{\textscy}{\textsubring{r}}} \dicPos{m}[9] \dicFlx{(‑s, ‑ir)}[8] \dicSynonym{mergð} \dicDirectTranslationCS{spousta, množství} \dicExampleIS{sægur af fólki} \dicExampleCS{spousta lidí}
\dicEntry[sæhestur] \dicTerm{sæ··hest|ur} \dicIPA{{s}{a}{i}{\textlengthmark}{h}{\textepsilon}{s}{\textsubring{d}}{\textscy}{\textsubring{r}}} \dicPos{m}[6] \dicFlx{(‑s, ‑ar)}[4] \dicFieldCat{zool.} \dicDirectTranslationCS{koníček, koníček\,/\addthin koník mořský} \textit{(l.~{\textLA{Hippocampus}})}  \dicsymPhoto\ 
\dicFigure{ds_image_saehestur_0_1.jpg}{Sæhestur}{Sæhestur - Florin DUMITRESCU, CC BY-SA 3.0}
\dicEntry[sæi] \dicTerm{sæi} \dicIPA{{s}{a}{i}{j}{\textsci}} \dicPos{v} \dicFlx{con pf sg 1 pers} \dicLink{sjá}
\dicEntry[sæki] \dicTerm{sæki} \dicIPA{{s}{a}{i}{\textlengthmark}{\r{\textObardotlessj}}{\textsci}} \dicPos{v} \dicFlx{ind praes sg 1 pers} \dicLink{sækja}
\dicEntry[sækja] \dicTerm{sækja} \dicsymFrequent\  \dicIPA{{s}{a}{i}{\textlengthmark}{\r{\textObardotlessj}}{a}} \dicPos{v}[5] \dicFlx{(sæki, sótti, sóttum, sækti, sótt)}[6] \dicFlx{acc} \textbf{1.} \dicSynonym*{ná í} \dicDirectTranslationCS{vyzvednout (si), vyzvedat (si),  dojít (si) pro, přijít (si) pro} \dicExampleIS{sækja póstinn} \dicExampleCS{vyzvednout si poštu}  \textbf{2.} \dicSynonym*{berjast við} \dicDirectTranslationCS{(za)útočit, napadnout} \dicExampleIS{sækja hann með vopnum} \dicExampleCS{napadnout ho zbraní}  \textbf{3.} \dicSynonym*{koma á e‑n stað} \dicDirectTranslationCS{navštívit, navštěvovat, docházet (na přednášky ap.)};  \dicPhraseIS{sækja vinnu} \dicDirectTranslationCS{pracovat}  \textbf{4.} \dicFieldCat{poč.} \dicSynonym*{hlaða niður} \dicDirectTranslationCS{stáhnout (si), stahovat (si)};  \dicIdiom{sækja}[á]{ \dicPhraseIS{e‑að sækir á e‑n}} \dicDirectTranslationCS{(co koho) přemáhá, (co koho) drtí (těžké myšlenky ap.)};  \dicIdiom{sækja}[að]{ \dicPhraseIS{sækja að e‑m}} \dicDirectTranslationCS{zaútočit na (koho)} \dicExampleIS{sækja að óvinunum} \dicExampleCS{zaútočit na nepřátele};  \dicIdiom{sækja}[fram]{ \dicPhraseIS{sækja fram}} \dicDirectTranslationCS{postupovat, jít kupředu (vojsko ap.)}; { \dicPhraseIS{þegar fram í sækir}} \dicFlx{adv} \dicDirectTranslationCS{(s) časem, výhledově};  \dicIdiom{sækja}[heim]{ \dicPhraseIS{sækja e‑ð heim}} \dicDirectTranslationCS{navštívit (co)};  \dicIdiom{sækja}[í]{ \dicPhraseIS{sækja í e‑ð}} \dicDirectTranslationCS{vyhledávat (co) (dobrou společnost ap.)} \dicExampleIS{Hann sækir í vondan félagsskap.} \dicExampleCS{Vyhledává špatnou společnost.}; { \dicPhraseIS{sækja í sig veðrið}} \dicLangCat{přen.} \dicDirectTranslationCS{nabrat druhý dech};  \dicIdiom{sækja}[móti]{ \dicPhraseIS{sækja móti e‑m\,/\addthin e‑u}} \dicDirectTranslationCS{bojovat proti (komu\,/\addthin čemu)}; { \dicPhraseIS{sækja móti vindi}} \dicDirectTranslationCS{zápasit s~větrem};  \dicIdiom{sækja}[um]{ \dicPhraseIS{sækja um e‑ð}} \dicDirectTranslationCS{(po)žádat\,/\addthin zažádat o~(co)} \dicExampleIS{sækja um lán} \dicExampleCS{požádat o~půjčku};  \dicIdiom{sækjast}{ \dicPhraseIS{e‑m sækist e‑að (vel)}} \dicFlx{refl impers} \dicDirectTranslationCS{(komu) se (co) daří (dobře)}; { \dicPhraseIS{e‑að sækist vel\,/\addthin illa}} \dicFlx{refl} \dicDirectTranslationCS{(co) se daří\,/\addthin nedaří} \dicExampleIS{Ferðin sækist vel.} \dicExampleCS{Cesta se daří.};  \dicIdiom{sækjast}[eftir]{ \dicPhraseIS{sækjast eftir e‑u}} \dicFlx{refl} \dicDirectTranslationCS{jít za (čím), usilovat o~(co)}
\dicEntry[sækjandi] \dicTerm{sækj··|andi} \dicIPA{{s}{a}{i}{\textlengthmark}{\r{\textObardotlessj}}{a}{n}{\textsubring{d}}{\textsci}} \dicPos{m}[2] \dicFlx{(‑anda, ‑endur)}[1] \dicFieldCat{práv.} \dicDirectTranslationCS{žalobce, žalobkyně}
\dicEntry[sækti] \dicTerm{sækti} \dicIPA{{s}{a}{i}{x}{\textsubring{d}}{\textsci}} \dicPos{v} \dicFlx{con pf sg 1 pers} \dicLink{sækja}
\dicEntry[sæla] \dicTerm{sæl|a} \dicIPA{{s}{a}{i}{\textlengthmark}{l}{a}} \dicPos{f}[1] \dicFlx{(‑u, ‑ur)}[19] \dicSynonym{velsæld} \dicDirectTranslationCS{blaho, blaženost, štěstí} \dicExampleIS{finna til óumræðilegrar sælu} \dicExampleCS{cítit nevyslovitelné blaho}
\dicEntry[sæld] \dicTerm{sæld} \dicIPA{{s}{a}{i}{l}{\textsubring{d}}} \dicPos{f}[7] \dicFlx{(‑ar)}[3] \dicSynonym{sæla} \dicDirectTranslationCS{blaho, blaženost, štěstí};  \dicPhraseIS{e‑að er ekki tekið út með (sitjandi) sældinni} \dicLangCat{přen.} \dicDirectTranslationCS{(co) není procházka růžovou zahradou}
\dicEntry[sælda] \dicTerm{sæld|a} \dicIPA{{s}{a}{i}{l}{\textsubring{d}}{a}} \dicPos{v}[1] \dicFlx{(‑aði)}[44] \dicPhraseIS{hafa (lítið) saman við e‑n að sælda} \dicDirectTranslationCS{mít s~(kým) (málo) do činění}
\dicEntry[sælgæti] \dicTerm{sæl··gæti} \dicIPA{{s}{a}{i}{l}{\r{\textObardotlessj}}{a}{i}{\textsubring{d}}{\textsci}} \dicPos{n}[2] \dicFlx{(‑s)}[20] \textbf{1.} \dicSynonym{sætindi} \dicDirectTranslationCS{sladkost, cukrovinka}  \textbf{2.} \dicSynonym*{bragðgóður matur} \dicDirectTranslationCS{lahůdka, delikatesa}
\dicEntry[sælína] \dicTerm{sæ··lín|a} \dicIPA{{s}{a}{i}{\textlengthmark}{l}{i}{n}{a}} \dicPos{f}[1] \dicFlx{(‑u, ‑ur)}[7] \dicDirectTranslationCS{podmořský kabel}
\dicEntry[sæljón] \dicTerm{sæ··ljón} \dicIPA{{s}{a}{i}{\textlengthmark}{l}{j}{ou}{\textsubring{n}}} \dicPos{n}[2] \dicFlx{(‑s, ‑)}[5] \dicFieldCat{zool.} \dicDirectTranslationCS{lachtan} \textit{(l.~{\textLA{Otaria}})}  \dicsymPhoto\ 
\dicFigure{ds_image_saeljon_0_1.jpg}{Sæljón}{Sæljón - Nestor Galina, CC BY 2.0}
\dicEntry[sælkerabúð] \dicTerm{sæl·kera··búð} \dicIPA{{s}{a}{i}{l}{c\smash{\textsuperscript{h}}}{\textepsilon}{r}{a}{\textsubring{b}}{u}{\texttheta}} \dicPos{f}[7] \dicFlx{(‑ar, ‑ir)}[1] \dicDirectTranslationCS{lahůdky, lahůdkářství} \dicIndirectTranslationCS{(obchod s~lahůdkami)}
\dicEntry[sælkeri] \dicTerm{sæl··ker|i} \dicIPA{{s}{a}{i}{l}{c\smash{\textsuperscript{h}}}{\textepsilon}{r}{\textsci}} \dicPos{m}[1] \dicFlx{(‑a, ‑ar)}[1] \dicDirectTranslationCS{labužník, labužnice, gurmán(ka)}
\dicEntry[sæll] \dicTerm{sæll} \dicsymFrequent\  \dicIPA{{s}{a}{i}{\textsubring{d}}{\textsubring{l}}} \dicPos{adj}[8]\dicFlx{}[-1] \textbf{1.} \dicSynonym{hamingjusamur} \dicDirectTranslationCS{šťastný, spokojený, blažený}  \textbf{2.} \dicSynonym{kær} \dicDirectTranslationCS{drahý, milý} \dicExampleIS{sonur sæll} \dicExampleCS{drahý synu}  \textbf{3.} \dicSynonym{látinn} \dicDirectTranslationCS{zesnulý} \dicExampleIS{hann faðir minn sæll} \dicExampleCS{můj zesnulý otec}  \textbf{4.} \dicSynonym{blessaður} \dicDirectTranslationCS{požehnaný, svatý}  \textbf{5.} \dicIndirectTranslationCS{(ve vítáních či loučeních)};  \dicPhraseIS{komdu sæll!} \dicDirectTranslationCS{dobrý den!, ahoj!, nazdar!} \dicIndirectTranslationCS{(pozdrav určený muži)};  \dicPhraseIS{komdu sæl!} \dicDirectTranslationCS{dobrý den!, ahoj!, nazdar!} \dicIndirectTranslationCS{(pozdrav určený ženě)};  \dicPhraseIS{sæll!} \dicDirectTranslationCS{ahoj!, čau!} \dicIndirectTranslationCS{(pozdrav určený muži)};  \dicPhraseIS{sæl!} \dicDirectTranslationCS{ahoj!, čau!} \dicIndirectTranslationCS{(pozdrav určený ženě)};  \dicPhraseIS{sælt veri fólkið!} \dicDirectTranslationCS{dobrý den všem!};  \dicPhraseIS{vertu sæll!} \dicDirectTranslationCS{na shledanou!, sbohem!} \dicIndirectTranslationCS{(pozdrav určený muži)};  \dicPhraseIS{vertu sæl!} \dicDirectTranslationCS{na shledanou!, sbohem!} \dicIndirectTranslationCS{(pozdrav určený ženě)};  \dicIdiom{sæll}{ \dicPhraseIS{þakka sínum sæla fyrir e‑ð}} \dicDirectTranslationCS{děkovat Bohu za (co)}
\dicEntry[sællegur] \dicTerm{sæl··legur} \dicIPA{{s}{a}{i}{l}{\textlengthmark}{\textepsilon}{\textbabygamma}{\textscy}{\textsubring{r}}} \dicPos{adj}[1]\dicFlx{}[-8] \dicDirectTranslationCS{spokojený, dobře živený}
\dicEntry[sællífi] \dicTerm{sæl··lífi} \dicIPA{{s}{a}{i}{l}{\textlengthmark}{i}{v}{\textsci}} \dicPos{n}[2] \dicFlx{(‑s)}[20] \dicSynonym{óhóf} \dicDirectTranslationCS{přepych, luxus, nadbytek}
\dicEntry[sæluhús] \dicTerm{sælu··hús} \dicIPA{{s}{a}{i}{\textlengthmark}{l}{\textscy}{h}{u}{s}} \dicPos{n}[2] \dicFlx{(‑s, ‑)}[5] \dicDirectTranslationCS{(horská záchranná) chata, (nouzové) přístřeší}
\dicEntry[sæluríki] \dicTerm{sælu··ríki} \dicIPA{{s}{a}{i}{\textlengthmark}{l}{\textscy}{r}{i}{\r{\textObardotlessj}}{\textsci}} \dicPos{n}[2] \dicFlx{(‑s, ‑)}[16] \dicDirectTranslationCS{utopie, utopický stát}
\dicEntry[sæluvíma] \dicTerm{sælu··vím|a} \dicIPA{{s}{a}{i}{\textlengthmark}{l}{\textscy}{v}{i}{m}{a}} \dicPos{f}[1] \dicFlx{(‑u)}[5] \dicDirectTranslationCS{příval štěstí, extáze, vytržení}
\dicEntry[sæma] \dicTerm{sæm|a} \dicIPA{{s}{a}{i}{\textlengthmark}{m}{a}} \dicPos{v}[2] \dicFlx{(‑di, ‑t)}[158] \dicFlx{acc + dat} \dicSynonym{heiðra} \dicDirectTranslationCS{honorovat, vyznamenat, dekorovat} \dicExampleIS{sæma hann riddarakrossi} \dicExampleCS{vyznamenat ho rytířským křížem};  \dicPhraseIS{e‑m sæmir e‑að} \dicFlx{impers} \dicDirectTranslationCS{(co komu) přísluší\,/\addthin náleží (po právu ap.)} \dicExampleIS{Mér sæmir ekki að níðast á þér.} \dicExampleCS{Nepřísluší mi, abych tě tak využíval.}
\dicEntry[sæmd] \dicTerm{sæmd} \dicIPA{{s}{a}{i}{m}{\textsubring{d}}} \dicPos{f}[7] \dicFlx{(‑ar, ‑ir)}[1] \dicSynonym{sómi} \dicDirectTranslationCS{(po)čest, pocta};  \dicPhraseIS{vera e‑m til sæmdar} \dicDirectTranslationCS{být (komu) ke cti}
\dicEntry[sæmdur] \dicTerm{sæmdur} \dicIPA{{s}{a}{i}{m}{\textsubring{d}}{\textscy}{\textsubring{r}}} \dicPos{adj}[2]\dicFlx{}[-17] \dicDirectTranslationCS{poctěný, oceněný}
\dicEntry[sæmilega] \dicTerm{sæmi··lega} \dicsymFrequent\  \dicIPA{{s}{a}{i}{\textlengthmark}{m}{\textsci}{l}{\textepsilon}{\textbabygamma}{a}} \dicPos{adv} \dicDirectTranslationCS{obstojně, solidně, slušně} \dicExampleIS{Hann talar sæmilega íslensku.} \dicExampleCS{Mluví obstojně islandsky.}
\dicEntry[sæmilegur] \dicTerm{sæmi··legur} \dicsymFrequent\  \dicIPA{{s}{a}{i}{\textlengthmark}{m}{\textsci}{l}{\textepsilon}{\textbabygamma}{\textscy}{\textsubring{r}}} \dicPos{adj}[1]\dicFlx{}[-8] \dicSynonym{bærilegur} \dicDirectTranslationCS{obstojný, slušný, ucházející};  \dicPhraseIS{vera sæmilegur í e‑u} \dicDirectTranslationCS{být celkem dobrý v~(čem)}
\dicEntry[sæng] \dicTerm{sæng} \dicsymFrequent\  \dicIPA{{s}{a}{i}{\ng}{\r{g}}} \dicPos{f}[10] \dicFlx{(‑ur, ‑ur)}[1] \dicDirectTranslationCS{peřina, přikrývka, deka} \dicExampleIS{Það eru tvær sængur í hjónarúminu.} \dicExampleCS{V~manželské posteli jsou dvě peřiny.};  \dicPhraseIS{ganga í eina sæng með e‑m} {\textbf{a.}} \dicDirectTranslationCS{oženit se s~(kým)};  {\textbf{b.}} \dicDirectTranslationCS{zahájit spolupráci s~(kým)};  \dicPhraseIS{sjá sína sæng uppreidda\,/\addthin útbreidda} \dicLangCat{přen.} \dicDirectTranslationCS{uvědomit si, že už je člověk mimo hru}
\dicEntry[sænga] \dicTerm{sæng|a} \dicIPA{{s}{a}{i}{\ng}{\r{g}}{a}} \dicPos{v}[1] \dicFlx{(‑aði)}[1] \dicPhraseIS{e‑ir sænga saman} \dicDirectTranslationCS{(kdo) spolu spí (pod jednou peřinou)};  \dicPhraseIS{sænga hjá\,/\addthin með e‑m} \dicDirectTranslationCS{spát s~(kým) (pod jednou peřinou)}
\dicEntry[sængurfatnaður] \dicTerm{sængur··fat·nað|ur} \dicIPA{{s}{a}{i}{\ng}{\r{g}}{\textscy}{\textsubring{r}}{f}{a}{h}{\textsubring{d}}{n}{a}{ð}{\textscy}{\textsubring{r}}} \dicPos{m}[10] \dicFlx{(‑ar)}[9] \dicDirectTranslationCS{ložní prádlo, lůžkoviny}
\dicEntry[sængurhiminn] \dicTerm{sængur··him|inn} \dicIPA{{s}{a}{i}{\ng}{\r{g}}{\textscy}{\textsubring{r}}{h}{\textsci}{m}{\textsci}{\textsubring{n}}} \dicPos{m}[6] \dicFlx{(‑ins, ‑nar)}[39] \dicDirectTranslationCS{baldachýn} \dicIndirectTranslationCS{(nad lůžkem)}
\dicEntry[sængurkona] \dicTerm{sængur··kon|a} \dicIPA{{s}{a}{i}{\ng}{\r{g}}{\textscy}{\textsubring{r}}{k\smash{\textsuperscript{h}}}{\textopeno}{n}{a}} \dicPos{f}[1] \dicFlx{(‑u, ‑ur)}[27] \dicDirectTranslationCS{šestinedělka, novopečená matka}
\dicEntry[sængurkvennadeild] \dicTerm{sængur·kvenna··deild} \dicIPA{{s}{a}{i}{\ng}{\r{g}}{\textscy}{\textsubring{r}}{k\smash{\textsuperscript{h}}}{v}{\textepsilon}{n}{a}{\textsubring{d}}{ei}{l}{\textsubring{d}}} \dicPos{f}[7] \dicFlx{(‑ar, ‑ir)}[1] \dicDirectTranslationCS{porodnické oddělení}
\dicEntry[sængurver] \dicTerm{sængur··ver} \dicIPA{{s}{a}{i}{\ng}{\r{g}}{\textscy}{r}{v}{\textepsilon}{\textsubring{r}}} \dicPos{n}[2] \dicFlx{(‑s, ‑)}[5] \dicDirectTranslationCS{cícha, povlak na peřinu}
\dicEntry[sænska] \dicTerm{sænsk|a} \dicIPA{{s}{a}{i}{n}{s}{\r{g}}{a}} \dicPos{f}[1] \dicFlx{(‑u)}[5] \dicDirectTranslationCS{švédština}
\dicEntry[sænskur] \dicTerm{sænskur} \dicsymFrequent\  \dicIPA{{s}{a}{i}{n}{s}{\r{g}}{\textscy}{\textsubring{r}}} \dicPos{adj}[1]\dicFlx{}[-1] \dicDirectTranslationCS{švédský} \dicExampleIS{á sænskri vefsíðu} \dicExampleCS{na švédské webové stránce}
\dicEntry[sæotur] \dicTerm{sæ··ot|ur} \dicIPA{{s}{a}{i}{\textlengthmark}{\textopeno}{\textsubring{d}}{\textscy}{\textsubring{r}}} \dicPos{m}[5] \dicFlx{(‑urs, ‑rar)}[1] \dicFieldCat{zool.} \dicDirectTranslationCS{vydra mořská, kalan} \textit{(l.~{\textLA{Enhydra lutris}})}  \dicsymPhoto\ 
\dicFigure{ds_image_saeotur_0_1.jpg}{Sæotur}{Sæotur - Michael L. Baird, CC BY 2.0}
\dicEntry[sær] \dicTerm{sæ|r\smash{\textsuperscript{1}}} \dicIPA{{s}{a}{i}{\textlengthmark}{\textsubring{r}}} \dicPos{m}[8] \dicFlx{(‑var)}[4] \dicLangCat{kniž.} \dicSynonym{sjór} \dicDirectTranslationCS{moře}
\dicEntry[sær] \dicTerm{sær\smash{\textsuperscript{2}}} \dicIPA{{s}{a}{i}{\textlengthmark}{\textsubring{r}}} \dicPos{adj}[5] \dicFlx{(f ‑)}[10] \dicPhraseIS{e‑m er eiður sær} \dicFlx{impers} \dicDirectTranslationCS{(kdo) může zaručit} \dicExampleIS{Mér er eiður sær að það var þannig.} \dicExampleCS{Můžu zaručit, že to tak bylo.}
\dicEntry[særa] \dicTerm{sær|a} \dicsymFrequent\  \dicIPA{{s}{a}{i}{\textlengthmark}{r}{a}} \dicPos{v}[2] \dicFlx{(‑ði, ‑t)}[99] \dicFlx{acc} \textbf{1.} \dicDirectTranslationCS{zranit, zraňovat, poranit, poraňovat, způsobit\,/\addthin způsobovat zranění} \dicExampleIS{særa hana djúpu sári} \dicExampleCS{způsobit jí hlubokou ránu}  \textbf{2.} \dicDirectTranslationCS{ranit, zasáhnout, ublížit, ubližovat} \dicIndirectTranslationCS{(psychicky)} \dicExampleIS{særa tilfinningar hans} \dicExampleCS{ranit jeho city}  \textbf{3.} \dicDirectTranslationCS{zastřihnout, přistřihnout (vlasy ap.)} \dicExampleIS{særa hárið} \dicExampleCS{zastřihnout vlasy}  \textbf{4.} \dicSynonym*{fara með særingar} \dicDirectTranslationCS{zaklínat, zaříkat};  \dicIdiom{særa}[fram]{ \dicPhraseIS{særa fram e‑ð}} \dicDirectTranslationCS{vyvolat (co) (ducha ap.)};  \dicIdiom{særast}{ \dicPhraseIS{særast}} \dicFlx{refl} \dicDirectTranslationCS{být zraněn, zranit se}
\dicEntry[særandi] \dicTerm{sær··andi} \dicIPA{{s}{a}{i}{\textlengthmark}{r}{a}{n}{\textsubring{d}}{\textsci}} \dicPos{adj}[13] \dicFlx{indecl}[1] \dicDirectTranslationCS{zraňující, bolestivý}
\dicEntry[særður] \dicTerm{særður} \dicIPA{{s}{a}{i}{r}{ð}{\textscy}{\textsubring{r}}} \dicPos{adj}[2]\dicFlx{}[-1] \textbf{1.} \dicSynonym{meiddur} \dicDirectTranslationCS{(z)raněný, poraněný}  \textbf{2.} \dicSynonym*{miður sín} \dicDirectTranslationCS{raněný, dotčený} \dicIndirectTranslationCS{(psychicky)} \dicExampleIS{vera djúpt særður eftir rifrildið} \dicExampleCS{být hluboce raněný po té hádce}
\dicEntry[særindi] \dicTerm{sær··indi} \dicIPA{{s}{a}{i}{\textlengthmark}{r}{\textsci}{n}{\textsubring{d}}{\textsci}} \dicPos{n}[2] \dicFlx{pl}[19] \dicSynonym{sárindi} \dicDirectTranslationCS{bolest, bolení} \dicExampleIS{særindi í hálsi} \dicExampleCS{bolest v~krku}
\dicEntry[særing] \dicTerm{sær··ing} \dicIPA{{s}{a}{i}{\textlengthmark}{r}{i}{\ng}{\r{g}}} \dicPos{f}[4] \dicFlx{(‑ar, ‑ar)}[5] \dicDirectTranslationCS{zaklínání, zaříkání} \dicIndirectTranslationCS{(zvláště v~pl)} \dicExampleIS{fara með særingar} \dicExampleCS{zaklínat}
\dicEntry[særok] \dicTerm{sæ··rok} \dicIPA{{s}{a}{i}{\textlengthmark}{r}{\textopeno}{\r{g}}} \dicPos{n}[2] \dicFlx{(‑s)}[2] \dicSynonym*{rok af hafi} \dicDirectTranslationCS{mořská bouře, bouře na moři}
\dicEntry[sæskjaldbaka] \dicTerm{sæ··skjald·|baka} \dicIPA{{s}{a}{i}{\textlengthmark}{s}{\r{\textObardotlessj}}{a}{l}{\textsubring{d}}{\textsubring{b}}{a}{\r{g}}{a}} \dicPos{f}[1] \dicFlx{(‑böku, ‑bökur)}[20] \dicFieldCat{zool.} \dicDirectTranslationCS{kareta, mořská želva} \textit{(l.~{\textLA{Cheloniidae}})}  \dicsymPhoto\ 
\dicFigure{21397.jpg}{Sæskjaldbaka}{Sæskjaldbaka - Hagerty, Ryan, Biolib, PD}
\dicEntry[sæstrengur] \dicTerm{sæ··streng|ur} \dicIPA{{s}{a}{i}{\textlengthmark}{s}{\textsubring{d}}{r}{ei}{\ng}{\r{g}}{\textscy}{\textsubring{r}}} \dicPos{m}[9] \dicFlx{(‑s\,/\addthin ‑jar, ‑ir)}[26] \dicDirectTranslationCS{podmořský kabel}
\dicEntry[sæta] \dicTerm{sæt|a\smash{\textsuperscript{1}}} \dicIPA{{s}{a}{i}{\textlengthmark}{\textsubring{d}}{a}} \dicPos{f}[1] \dicFlx{(‑u, ‑ur)}[19] \dicSynonym{sætleiki} \dicDirectTranslationCS{sladkost} \dicIndirectTranslationCS{(vlastnost sladké věci)}
\dicEntry[sæta] \dicTerm{sæt|a\smash{\textsuperscript{2}}} \dicIPA{{s}{a}{i}{\textlengthmark}{\textsubring{d}}{a}} \dicPos{v}[2] \dicFlx{(‑ti, ‑t)}[62] \dicFlx{dat} \textbf{1.} \dicPhraseIS{sæta færi að (gera e‑ð)} \dicDirectTranslationCS{chopit se příležitosti (udělat (co))};  \dicPhraseIS{sæta lagi að (gera e‑ð)} \dicDirectTranslationCS{využít příležitosti k~(udělání (čeho))}  \textbf{2.} \dicPhraseIS{hverju sætir það?} \dicDirectTranslationCS{z~jakého důvodu?, co je příčinou?} \dicExampleIS{Hverju sætir það að þetta hefur gengið svona hægt?} \dicExampleCS{Jak je možné, že to šlo tak pomalu?};  \dicPhraseIS{það sætir furðu} \dicFlx{impers} \dicDirectTranslationCS{je (po)divné}  \textbf{3.} \dicSynonym{bera} \dicDirectTranslationCS{nést, podstoupit};  \dicPhraseIS{sæta ábyrgð fyrir e‑ð} \dicDirectTranslationCS{nést za (co) odpovědnost};  \dicPhraseIS{sæta refsingu fyrir e‑ð} \dicDirectTranslationCS{podstoupit trest za (co)}
\dicEntry[sæta] \dicTerm{sæt|a\smash{\textsuperscript{3}}} \dicIPA{{s}{a}{i}{\textlengthmark}{\textsubring{d}}{a}} \dicPos{v}[2] \dicFlx{(‑ti, ‑t)}[62] \dicFlx{acc} \dicSynonym{sykra\smash{\textsuperscript{2}}} \dicDirectTranslationCS{(o)sladit} \dicExampleIS{sæta e‑ð} \dicExampleCS{osladit (co)}
\dicEntry[sætabrauð] \dicTerm{sæta··brauð} \dicIPA{{s}{a}{i}{\textlengthmark}{\textsubring{d}}{a}{\textsubring{b}}{r}{\oe i}{\texttheta}} \dicPos{n}[2] \dicFlx{(‑s, ‑)}[5] \dicFieldCat{kulin.} \dicDirectTranslationCS{moučník, závin}
\dicEntry[sætaröð] \dicTerm{sæta··|röð} \dicIPA{{s}{a}{i}{\textlengthmark}{\textsubring{d}}{a}{r}{\oe}{\texttheta}} \dicPos{f}[7] \dicFlx{(‑raðar, ‑raðir)}[16] \dicDirectTranslationCS{řada (sedadel)} \dicExampleIS{Við sátum í annarri sætaröð fyrir aftan nunnurnar.} \dicExampleCS{Seděli jsme v~druhé řadě za jeptiškami.}
\dicEntry[sætaskipti] \dicTerm{sæta··skipti} \dicIPA{{s}{a}{i}{\textlengthmark}{\textsubring{d}}{a}{s}{\r{\textObardotlessj}}{\textsci}{f}{\textsubring{d}}{\textsci}} \dicPos{n}[2] \dicFlx{pl}[19] \dicDirectTranslationCS{změna\,/\addthin výměna místa};  \dicPhraseIS{hafa sætaskipti við e‑n} \dicDirectTranslationCS{vyměnit si s~(kým) místo}
\dicEntry[sætavísa] \dicTerm{sæta··vís|a} \dicIPA{{s}{a}{i}{\textlengthmark}{\textsubring{d}}{a}{v}{i}{s}{a}} \dicPos{f}[1] \dicFlx{(‑u, ‑ur)}[13] \dicDirectTranslationCS{uvaděč(ka) (v~kině ap.)}
\dicEntry[sæti] \dicTerm{sæti\smash{\textsuperscript{1}}} \dicsymFrequent\  \dicIPA{{s}{a}{i}{\textlengthmark}{\textsubring{d}}{\textsci}} \dicPos{n}[2] \dicFlx{(‑s, ‑)}[14] \textbf{1.} \dicSynonym{set} \dicDirectTranslationCS{místo (na sezení)} \dicExampleIS{fá sér sæti á stólnum} \dicExampleCS{usednout na židli}  \textbf{2.} \dicSynonym*{staður í röð} \dicDirectTranslationCS{místo, pozice, umístění} \dicExampleIS{fyrsta sæti} \dicExampleCS{první místo}  \textbf{3.} \dicSynonym{aðild} \dicDirectTranslationCS{místo, účast (v~radě ap.)};  \dicPhraseIS{eiga sæti í e‑u} \dicDirectTranslationCS{sedět\,/\addthin zasedat v~(čem) (radě ap.)};  \dicPhraseIS{taka sæti í e‑u} \dicDirectTranslationCS{stát se členem (čeho) (vlády ap.)}  \textbf{4.} \dicSynonym*{heysæti} \dicDirectTranslationCS{kupka (sena), stoh};  \dicPhraseIS{setja hey í sæti} \dicDirectTranslationCS{rovnat seno do kupek}
\dicEntry[sæti] \dicTerm{sæti\smash{\textsuperscript{2}}} \dicIPA{{s}{a}{i}{\textlengthmark}{\textsubring{d}}{\textsci}} \dicPos{v} \dicFlx{con pf sg 1 pers} \dicLink{sitja}
\dicEntry[sætindaverslun] \dicTerm{sætinda··versl|un} \dicIPA{{s}{a}{i}{\textlengthmark}{\textsubring{d}}{\textsci}{n}{\textsubring{d}}{a}{v}{\textepsilon}{\textsubring{r}}{s}{\textsubring{d}}{l}{\textscy}{\textsubring{n}}} \dicPos{f}[7] \dicFlx{(‑unar, ‑anir)}[8] \dicDirectTranslationCS{cukrárna}
\dicEntry[sætindi] \dicTerm{sæt··indi} \dicIPA{{s}{a}{i}{\textlengthmark}{\textsubring{d}}{\textsci}{n}{\textsubring{d}}{\textsci}} \dicPos{n}[2] \dicFlx{pl}[19] \dicSynonym{sælgæti} \dicDirectTranslationCS{sladkosti, cukrovinky}
\dicEntry[sætisbak] \dicTerm{sætis··|bak} \dicIPA{{s}{a}{i}{\textlengthmark}{\textsubring{d}}{\textsci}{s}{\textsubring{b}}{a}{\r{g}}} \dicPos{n}[2] \dicFlx{(‑baks, ‑bök)}[8] \dicDirectTranslationCS{opěradlo (sedadla)}
\dicEntry[sætisbelti] \dicTerm{sætis··belti} \dicIPA{{s}{a}{i}{\textlengthmark}{\textsubring{d}}{\textsci}{s}{\textsubring{b}}{\textepsilon}{\textsubring{l}}{\textsubring{d}}{\textsci}} \dicPos{n}[2] \dicFlx{(‑s, ‑)}[14] \dicDirectTranslationCS{(sedadlový\,/\addthin bezpečnostní) pás}
\dicEntry[sætistala] \dicTerm{sætis··|tala} \dicIPA{{s}{a}{i}{\textlengthmark}{\textsubring{d}}{\textsci}{s}{t\smash{\textsuperscript{h}}}{a}{l}{a}} \dicPos{f}[1] \dicFlx{(‑tölu, ‑tölur)}[14] \dicFieldCat{fyz.} \dicDirectTranslationCS{protonové\,/\addthin atomové číslo}
\dicEntry[sætleiki] \dicTerm{sæt··leik|i}\dicTerm{, sætleikur} \dicIPA{{s}\-{a}\-{i}\-{\textlengthmark}\-{\textsubring{d}}\-{l}\-{ei}\-{\r{\textObardotlessj}}\-{\textsci}\-} \dicPos{m}[1] \dicFlx{(‑a)}[3] \dicDirectTranslationCS{sladkost} \dicIndirectTranslationCS{(vlastnost sladké věci)}
\dicEntry[sætleikur] \dicTerm{sæt··leik|ur} \dicIPA{{s}{a}{i}{\textlengthmark}{\textsubring{d}}{l}{ei}{\r{g}}{\textscy}{\textsubring{r}}} \dicPos{m}[6] \dicFlx{(‑s)}[17] \dicLink{sætleiki}
\dicEntry[sætt] \dicTerm{sætt} \dicIPA{{s}{a}{i}{h}{\textsubring{d}}} \dicPos{f}[7] \dicFlx{(‑ar, ‑ir)}[1] \textbf{1.} \dicSynonym{sátt} \dicDirectTranslationCS{(u)smíření, smír} \dicExampleIS{koma á sættum} \dicExampleCS{usmířit se}  \textbf{2.} \dicSynonym{sáttargerð} \dicDirectTranslationCS{dohoda, ujednání, usnesení}
\dicEntry[sætta] \dicTerm{sætt|a} \dicsymFrequent\  \dicIPA{{s}{a}{i}{h}{\textsubring{d}}{a}} \dicPos{v}[2] \dicFlx{(‑i, ‑)}[1] \dicFlx{acc} \dicDirectTranslationCS{(u)smířit, (u)smiřovat, udobřit, udobřovat} \dicExampleIS{sætta þær} \dicExampleCS{usmířit je};  \dicPhraseIS{sætta sig við e‑ð} \dicDirectTranslationCS{smířit se s~(čím)} \dicExampleIS{sætta sig við niðurstöðuna} \dicExampleCS{smířit se s~rozhodnutím};  \dicIdiom{sættast}{ \dicPhraseIS{sættast}} \dicFlx{refl} \dicDirectTranslationCS{(u)smířit se, udobřit se} \dicExampleIS{sættast við hana} \dicExampleCS{usmířit se s~ní};  \dicIdiom{sættast}[á]{ \dicPhraseIS{sættast á e‑ð}} \dicFlx{refl} \dicDirectTranslationCS{shodnout se na (čem)}
\dicEntry[sættargerð] \dicTerm{sættar··gerð}\dicTerm{, sættargjörð} \dicIPA{{s}\-{a}\-{i}\-{h}\-{\textsubring{d}}\-{a}\-{r}\-{\r{\textObardotlessj}}\-{\textepsilon}\-{r}\-{\texttheta}\-} \dicPos{f}[7] \dicFlx{(‑ar, ‑ir)}[1] \dicSynonym{sáttargerð} \dicDirectTranslationCS{(u)smíření, dohoda}
\dicEntry[sættargjörð] \dicTerm{sættar··gjörð} \dicIPA{{s}{a}{i}{h}{\textsubring{d}}{a}{r}{\r{\textObardotlessj}}{\oe}{r}{\texttheta}} \dicPos{f}[7] \dicFlx{(‑ar, ‑ir)}[1] \dicLink{sættargerð}
\dicEntry[sættir] \dicTerm{sættir} \dicIPA{{s}{a}{i}{h}{\textsubring{d}}{\textsci}{\textsubring{r}}} \dicPos{f} \dicFlx{pl nom} \dicLink{sátt}
\dicEntry[sætuefni] \dicTerm{sætu··efni} \dicIPA{{s}{a}{i}{\textlengthmark}{\textsubring{d}}{\textscy}{\textepsilon}{\textsubring{b}}{n}{\textsci}} \dicPos{n}[2] \dicFlx{(‑s, ‑)}[14] \dicDirectTranslationCS{sladidlo};  \dicPhraseIS{tilbúið sætuefni} \dicDirectTranslationCS{umělé sladidlo}
\dicEntry[sætur] \dicTerm{sætur} \dicsymFrequent\  \dicIPA{{s}{a}{i}{\textlengthmark}{\textsubring{d}}{\textscy}{\textsubring{r}}} \dicPos{adj}[1]\dicFlx{}[-1] \textbf{1.} \dicSynonym{sykraður} \dicDirectTranslationCS{sladký} \dicIndirectTranslationCS{(chuťově)}  \textbf{2.} \dicSynonym{elskulegur} \dicDirectTranslationCS{sladký, roztomilý, líbezný, milý} \dicExampleIS{sætir kettlingar} \dicExampleCS{roztomilá koťátka}  \textbf{3.} \dicSynonym{snotur} \dicDirectTranslationCS{sladký, půvabný, rozkošný}
\dicEntry[södd] \dicTerm{södd} \dicIPA{{s}{\oe}{\textsubring{d}}{\textlengthmark}} \dicPos{adj} \dicFlx{f sg nom pos} \dicLink{saddur}
\dicEntry[söddum] \dicTerm{söddum} \dicIPA{{s}{\oe}{\textsubring{d}}{\textlengthmark}{\textscy}{\textsubring{m}}} \dicPos{v} \dicFlx{ind pf pl 1 pers} \dicLink{seðja}
\dicEntry[söðla] \dicTerm{söðl|a} \dicIPA{{s}{\oe}{ð}{l}{a}} \dicPos{v}[1] \dicFlx{(‑aði)}[1] \dicFlx{acc} \dicDirectTranslationCS{(o)sedlat} \dicExampleIS{söðla hest} \dicExampleCS{osedlat koně};  \dicIdiom{söðla}[um]{ \dicPhraseIS{söðla um}} \dicLangCat{přen.} \dicDirectTranslationCS{přesedlat} \dicIndirectTranslationCS{(změnit dosavadní zaměstnání, zájmy ap.)}
\dicEntry[söðull] \dicTerm{söð|ull} \dicIPA{{s}{\oe}{\textlengthmark}{ð}{\textscy}{\textsubring{d}}{\textsubring{l}}} \dicPos{m}[6] \dicFlx{(‑uls, ‑lar)}[31] \dicDirectTranslationCS{dámské (jezdecké) sedlo} \dicExampleIS{leggja söðul á hestinn} \dicExampleCS{osedlat koně dámským sedlem};  \dicPhraseIS{hafa ekki úr háum söðli að detta} \dicLangCat{přen.} \dicDirectTranslationCS{nemít co ztratit}
\dicEntry[söfn] \dicTerm{söfn} \dicIPA{{s}{\oe}{\textsubring{b}}{\textsubring{n}}} \dicPos{n} \dicFlx{pl nom} \dicLink{safn}
\dicEntry[söfnuður] \dicTerm{söfn··uður} \dicIPA{{s}{\oe}{\textsubring{b}}{n}{\textscy}{ð}{\textscy}{\textsubring{r}}} \dicPos{m}[10] \dicFlx{(söfnuðar\,/\addthin safnaðar, söfnuðir)}[44] \textbf{1.} \dicSynonym{fjöldi} \dicDirectTranslationCS{shromáždění (lidí ap.)}  \textbf{2.} \dicFieldCat{náb.} \dicSynonym*{trúarsöfnuður} \dicDirectTranslationCS{(církevní) shromáždění, kongregace} \dicExampleIS{kristinn söfnuður} \dicExampleCS{shromáždění křesťanů}
\dicEntry[söfnun] \dicTerm{söfnun} \dicIPA{{s}{\oe}{\textsubring{b}}{n}{\textscy}{\textsubring{n}}} \dicPos{f}[7] \dicFlx{(söfnunar, safnanir)}[11] \textbf{1.} \dicSynonym*{það að safna} \dicDirectTranslationCS{sběr, sbírání, kolekce}  \textbf{2.} \dicSynonym{samskot} \dicDirectTranslationCS{sbírka (dobročinná ap.)}
\dicEntry[sög] \dicTerm{sög} \dicIPA{{s}{\oe}{\textlengthmark}{x}} \dicPos{f}[7] \dicFlx{(sagar, sagir)}[16] \dicDirectTranslationCS{pila}
\dicEntry[sögðum] \dicTerm{sögðum} \dicIPA{{s}{\oe}{\textbabygamma}{ð}{\textscy}{\textsubring{m}}} \dicPos{v} \dicFlx{ind pf pl 1 pers} \dicLink{segja}
\dicEntry[sögn] \dicTerm{sögn} \dicsymFrequent\  \dicIPA{{s}{\oe}{\r{g}}{\textsubring{n}}} \dicPos{f}[7] \dicFlx{(sagnar, sagnir)}[16] \textbf{1.} \dicSynonym*{saga} \dicDirectTranslationCS{vyprávění, příběh, historka, povídka} \dicExampleIS{það er sögn manna að} \dicExampleCS{povídá se, že}  \textbf{2.} \dicSynonym{frásögn} \dicDirectTranslationCS{tvrzení, prohlášení, oznámení};  \dicPhraseIS{að sögn e‑rs} \dicFlx{prep} \dicDirectTranslationCS{podle (tvrzení) (koho)}  \textbf{3.} \dicFieldCat{jaz.} \dicSynonym{sagnorð} \dicDirectTranslationCS{sloveso} \textit{(l.~{\textLA{verbum}})} \dicExampleIS{persónuleg sögn} \dicExampleCS{osobní sloveso}  \textbf{4.} \dicDirectTranslationCS{(trumfová) hláška}
\dicEntry[sögu] \dicTerm{sögu} \dicIPA{{s}{\oe}{\textlengthmark}{\textbabygamma}{\textscy}} \dicPos{f} \dicFlx{sg gen} \dicLink{saga\smash{\textsuperscript{1}}}
\dicEntry[sögubók] \dicTerm{sögu··|bók} \dicIPA{{s}{\oe}{\textlengthmark}{\textbabygamma}{\textscy}{\textsubring{b}}{ou}{\r{g}}} \dicPos{f}[8] \dicFlx{(‑bókar, ‑bækur)}[5] \textbf{1.} \dicFieldCat{škol.} \dicDirectTranslationCS{učebnice dějepisu}  \textbf{2.} \dicFieldCat{lit.} \dicDirectTranslationCS{sbírka povídek}
\dicEntry[söguburður] \dicTerm{sögu··burð|ur} \dicIPA{{s}{\oe}{\textlengthmark}{\textbabygamma}{\textscy}{\textsubring{b}}{\textscy}{r}{ð}{\textscy}{\textsubring{r}}} \dicPos{m}[10] \dicFlx{(‑ar)}[7] \dicDirectTranslationCS{drb, klep, pomluva}
\dicEntry[sögufrægur] \dicTerm{sögu··frægur} \dicIPA{{s}{\oe}{\textlengthmark}{\textbabygamma}{\textscy}{f}{r}{a}{i}{\textbabygamma}{\textscy}{\textsubring{r}}} \dicPos{adj}[1]\dicFlx{}[-1] \dicDirectTranslationCS{legendární, historický, slavný}
\dicEntry[söguhetja] \dicTerm{sögu··hetj|a} \dicIPA{{s}{\oe}{\textlengthmark}{\textbabygamma}{\textscy}{h}{\textepsilon}{\textsubring{d}}{j}{a}} \dicPos{f}[1] \dicFlx{(‑u, ‑ur)}[7] \dicFieldCat{lit.} \dicDirectTranslationCS{protagonista, protagonistka, románový hrdina, románová hrdinka}
\dicEntry[sögukennari] \dicTerm{sögu··kenn·ar|i} \dicIPA{{s}{\oe}{\textlengthmark}{\textbabygamma}{\textscy}{c\smash{\textsuperscript{h}}}{\textepsilon}{n}{a}{r}{\textsci}} \dicPos{m}[1] \dicFlx{(‑a, ‑ar)}[13] \dicDirectTranslationCS{učitel(ka) dějepisu\,/\addthin historie}
\dicEntry[sögulegur] \dicTerm{sögu··legur} \dicsymFrequent\  \dicIPA{{s}{\oe}{\textlengthmark}{\textbabygamma}{\textscy}{l}{\textepsilon}{\textbabygamma}{\textscy}{\textsubring{r}}} \dicPos{adj}[1]\dicFlx{}[-8] \textbf{1.} \dicSynonym*{sagnfræðilegur} \dicDirectTranslationCS{historický, dějinný, dějepisný}  \textbf{2.} \dicSynonym*{frásagnaverður} \dicDirectTranslationCS{historický, významný (vítězství ap.)} \dicExampleIS{sögulegur sigur} \dicExampleCS{historické vítězství}
\dicEntry[söguljóð] \dicTerm{sögu··ljóð} \dicIPA{{s}{\oe}{\textlengthmark}{\textbabygamma}{\textscy}{l}{j}{ou}{\texttheta}} \dicPos{n}[2] \dicFlx{(‑s, ‑)}[5] \dicFieldCat{lit.} \dicDirectTranslationCS{epos, epická báseň}
\dicEntry[sögulok] \dicTerm{sögu··lok} \dicIPA{{s}{\oe}{\textlengthmark}{\textbabygamma}{\textscy}{l}{\textopeno}{\r{g}}} \dicPos{n}[2] \dicFlx{pl}[1] \dicDirectTranslationCS{konec příběhu\,/\addthin vyprávění}
\begin{xtolerant}{}{1pt}
\dicEntry[sögumaður] \dicTerm{sögu··|maður}\addthinS\dicIPA{{s}{\oe}{\textlengthmark}{\textbabygamma}{\textscy}{m}{a}{ð}{\textscy}{\textsubring{r}}}\addthinS\dicPos{m}[13]\addthinS\dicFlx{(‑manns, ‑menn)}[2]\addthinS\textbf{1.}\addthinS\dicSynonym{sagnfræðingur} \dicDirectTranslationCS{historik, historička}  \textbf{2.} \dicFieldCat{lit.} \dicDirectTranslationCS{vypravěč(ka) (v~románu ap.)}
\end{xtolerant}
\dicEntry[sögun] \dicTerm{sög|un} \dicIPA{{s}{\oe}{\textlengthmark}{\textbabygamma}{\textscy}{\textsubring{n}}} \dicPos{f}[7] \dicFlx{(‑unar)}[12] \dicDirectTranslationCS{řezání (pilou)}
\dicEntry[sögunarverksmiðja] \dicTerm{sögunar··verk·smiðj|a} \dicIPA{{s}\-{\oe}\-{\textbabygamma}\-{\textscy}\-{n}\-{a}\-{r}\-{v}\-{\textepsilon}\-{\textsubring{r}}\-{\r{g}}\-{s}\-{m}\-{\textsci}\-{ð}\-{j}\-{a}\-} \dicPos{f}[1] \dicFlx{(‑u, ‑ur)}[7] \dicDirectTranslationCS{pila} \dicIndirectTranslationCS{(závod na zpracování dřeva)}
\dicEntry[sögupersóna] \dicTerm{sögu··persón|a} \dicIPA{{s}{\oe}{\textlengthmark}{\textbabygamma}{\textscy}{p\smash{\textsuperscript{h}}}{\textepsilon}{\textsubring{r}}{s}{ou}{n}{a}} \dicPos{f}[1] \dicFlx{(‑u, ‑ur)}[7] \dicFieldCat{lit.} \dicDirectTranslationCS{románová postava}
\dicEntry[söguskoðun] \dicTerm{sögu··skoð|un} \dicIPA{{s}{\oe}{\textlengthmark}{\textbabygamma}{\textscy}{s}{\r{g}}{\textopeno}{ð}{\textscy}{\textsubring{n}}} \dicPos{f}[7] \dicFlx{(‑unar)}[9] \dicDirectTranslationCS{pohled na historii, pojetí historie\,/\addthin dějin}
\dicEntry[sögustaður] \dicTerm{sögu··stað|ur} \dicIPA{{s}{\oe}{\textlengthmark}{\textbabygamma}{\textscy}{s}{\textsubring{d}}{a}{ð}{\textscy}{\textsubring{r}}} \dicPos{m}[10] \dicFlx{(‑ar, ‑ir)}[14] \dicDirectTranslationCS{historické místo}
\dicEntry[sögusvið] \dicTerm{sögu··svið} \dicIPA{{s}{\oe}{\textlengthmark}{\textbabygamma}{\textscy}{s}{v}{\textsci}{\texttheta}} \dicPos{n}[2] \dicFlx{(‑s, ‑)}[5] \dicFieldCat{lit.} \dicDirectTranslationCS{dějiště (románu ap.)}
\dicEntry[sögusögn] \dicTerm{sögu··|sögn} \dicIPA{{s}{\oe}{\textlengthmark}{\textbabygamma}{\textscy}{s}{\oe}{\r{g}}{\textsubring{n}}} \dicPos{f}[7] \dicFlx{(‑sagnar, ‑sagnir)}[16] \dicDirectTranslationCS{mýtus, úst\-ně předávaný příběh, historka}
\dicEntry[söguþráður] \dicTerm{sögu··|þráður} \dicIPA{{s}{\oe}{\textlengthmark}{\textbabygamma}{\textscy}{\texttheta}{r}{au}{ð}{\textscy}{\textsubring{r}}} \dicPos{m}[12] \dicFlx{(‑þráðar, ‑þræðir)}[7] \dicFieldCat{lit.} \dicDirectTranslationCS{zápletka, děj}
\dicEntry[söguöld] \dicTerm{sögu··|öld} \dicIPA{{s}{\oe}{\textlengthmark}{\textbabygamma}{\textscy}{\oe}{l}{\textsubring{d}}} \dicPos{f}[7] \dicFlx{(‑aldar)}[19] \dicFieldCat{hist.} \dicDirectTranslationCS{období ság} \dicIndirectTranslationCS{(v~islandské historii, v~letech 930--1030)}
\dicEntry[sök] \dicTerm{sök} \dicsymFrequent\  \dicIPA{{s}{\oe}{\textlengthmark}{\r{g}}} \dicPos{f}[7] \dicFlx{(sakar, sakir)}[16] \textbf{1.} \dicSynonym*{kæra} \dicDirectTranslationCS{obvinění};  \dicPhraseIS{bera af sér sakir} \dicDirectTranslationCS{vyvrátit obvinění};  \dicPhraseIS{bera sakir á e‑n} \dicDirectTranslationCS{(ob)vinit (koho)};  \dicPhraseIS{gefa e‑m að sök að (gera e‑ð)} \dicDirectTranslationCS{dávat (komu) za vinu, že (dělá (co))}  \textbf{2.} \dicSynonym{sekt} \dicDirectTranslationCS{vina, provinění};  \dicPhraseIS{eiga sök á e‑u} \dicDirectTranslationCS{mít vinu na (čem)};  \dicPhraseIS{gefa e‑m upp sakir} \dicDirectTranslationCS{omilostnit (koho), odpustit (komu) provinění};  \dicPhraseIS{gera upp sakirnar við e‑n} \dicDirectTranslationCS{volat (koho) k~odpovědnosti};  \dicPhraseIS{hafa e‑n fyrir rangri sök} \dicDirectTranslationCS{obvinit (koho) mylně\,/\addthin neprávem};  \dicPhraseIS{taka á sig sökina} \dicDirectTranslationCS{vzít na sebe vinu};  \dicPhraseIS{vera sýkn saka} \dicDirectTranslationCS{být zproštěn viny};  \dicPhraseIS{verða sannur að sök} \dicDirectTranslationCS{být shledán vinným}  \textbf{3.} \dicSynonym{orsök} \dicDirectTranslationCS{osud, sběh událostí};  \dicPhraseIS{af þeim sökum} \dicFlx{adv} \dicDirectTranslationCS{z~těchto důvodů};  \dicPhraseIS{fyrir hvaða sök} \dicDirectTranslationCS{z~jakého důvodu};  \dicIdiom{sök}{ \dicPhraseIS{e‑að kemur ekki að sök}} \dicDirectTranslationCS{(co) nic nedělá, (co) nijak neuškodí};  \dicPhraseIS{eins og sakir standa} \dicFlx{adv} \dicDirectTranslationCS{v~tuto chvíli, v~tento moment}; { \dicPhraseIS{fara hægt í sakirnar}} \dicLangCat{přen.} \dicDirectTranslationCS{jít pomalu na věc};  \dicPhraseIS{sakir\,/\addthin sökum þess að} \dicFlx{conj} \dicDirectTranslationCS{protože; z~toho důvodu, že; kvůli tomu, že}
\dicEntry[sökk] \dicTerm{sökk\smash{\textsuperscript{1}}} \dicIPA{{s}{\oe}{h}{\r{g}}} \dicPos{n}[2] \dicFlx{(‑s, ‑)}[5] \dicFieldCat{med.} \dicSynonym*{blóðsökk} \dicDirectTranslationCS{krevní sedimentace}
\dicEntry[sökk] \dicTerm{sökk\smash{\textsuperscript{2}}} \dicIPA{{s}{\oe}{h}{\r{g}}} \dicPos{v} \dicFlx{ind pf sg 1 pers} \dicLink{sökkva\smash{\textsuperscript{2}}}
\dicEntry[sökku] \dicTerm{sökku} \dicIPA{{s}{\oe}{h}{\r{g}}{\textscy}} \dicPos{f} \dicFlx{sg gen} \dicLink{sakka}
\dicEntry[sökkull] \dicTerm{sökk|ull} \dicIPA{{s}{\oe}{h}{\r{g}}{\textscy}{\textsubring{d}}{\textsubring{l}}} \dicPos{m}[6] \dicFlx{(‑uls, ‑lar)}[31] \textbf{1.} \dicDirectTranslationCS{(betonový) základ}  \textbf{2.} \dicDirectTranslationCS{sokl, podstavec}
\dicEntry[sökkva] \dicTerm{sökk|va\smash{\textsuperscript{1}}} \dicsymFrequent\  \dicIPA{{s}{\oe}{h}{\r{g}}{v}{a}} \dicPos{v}[2] \dicFlx{(‑ti, ‑t)}[36] \dicFlx{dat} \dicSynonym*{láta sökkva} \dicDirectTranslationCS{ponořit, potopit} \dicExampleIS{sökkva bátnum} \dicExampleCS{potopit loď};  \dicIdiom{sökkva}[í]{ \dicPhraseIS{sökkva sér í skuldir}} \dicLangCat{přen.} \dicDirectTranslationCS{topit se v~dluzích};  \dicIdiom{sökkva}[niður]{ \dicPhraseIS{sökkva sér niður\,/\addthin ofan í e‑ð}} \dicDirectTranslationCS{ponořit se do (čeho), zabrat se do (čeho)} \dicExampleIS{sökkva sér niður í bók} \dicExampleCS{ponořit se do knížky};  \dicIdiom{sökkvast}{ \dicPhraseIS{sökkvast}} \dicFlx{refl} \dicDirectTranslationCS{ponořit se}
\dicEntry[sökkva] \dicTerm{sökkva\smash{\textsuperscript{2}}} \dicsymFrequent\  \dicIPA{{s}{\oe}{h}{\r{g}}{v}{a}} \dicPos{v}[6] \dicFlx{(sekk, sökk, sukkum, sykki, sokkið)}[91] \dicDirectTranslationCS{ponořit se, potopit se} \dicExampleIS{Fætur hans sukku í gljúpan jarðveginn.} \dicExampleCS{Nohy se mu zabořily do porézní půdy.}
\dicEntry[söknuður] \dicTerm{sökn··uður} \dicIPA{{s}{\oe}{h}{\r{g}}{n}{\textscy}{ð}{\textscy}{\textsubring{r}}} \dicPos{m}[10] \dicFlx{(saknaðar\,/\addthin söknuðar)}[43] \dicSynonym{eftirsjá} \dicDirectTranslationCS{lítost, zármutek, stesk} \dicExampleIS{finna til saknaðar} \dicExampleCS{cítit lítost}
\dicEntry[sökudólgur] \dicTerm{söku··dólg|ur} \dicIPA{{s}{\oe}{\textlengthmark}{\r{g}}{\textscy}{\textsubring{d}}{ou}{l}{\r{g}}{\textscy}{\textsubring{r}}} \dicPos{m}[6] \dicFlx{(‑s, ‑ar)}[24] \dicDirectTranslationCS{viník, vinice, provinilec, provinilkyně}
\dicEntry[sökum] \dicTerm{sökum} \dicsymFrequent\  \dicIPA{{s}{\oe}{\textlengthmark}{\r{g}}{\textscy}{\textsubring{m}}} \dicPos{prep} \dicFlx{gen} \dicSynonym{vegna\smash{\textsuperscript{2}}} \dicDirectTranslationCS{kvůli, zásluhou} \dicExampleIS{sökum óvissu} \dicExampleCS{kvůli nejistotě};  \dicPhraseIS{sökum þess að} \dicFlx{conj} \dicDirectTranslationCS{díky tomu, že; kvůli tomu, že}
\dicEntry[söl] \dicTerm{söl} \dicIPA{{s}{\oe}{\textlengthmark}{\textsubring{l}}} \dicPos{n}[2] \dicFlx{pl}[9] \dicFieldCat{bot.} \dicIndirectTranslationCS{druh mořské řasy sloužící k~obživě} \textit{(l.~{\textLA{Palmaria palmata}})}  \dicsymPhoto\ 
\dicFigure{ds_image_sol_0_1.jpg}{Söl}{Söl - Cwmhiraeth, CC BY-SA 3.0}
\dicEntry[sölna] \dicTerm{söln|a} \dicIPA{{s}{\oe}{l}{n}{a}} \dicPos{v}[1] \dicFlx{(‑aði)}[1] \dicSynonym{fölna} \dicDirectTranslationCS{(z)vadnout, uvadat, uschnout, usychat} \dicExampleIS{Grasið sölnar.} \dicExampleCS{Tráva vadne.}
\dicEntry[sölnaður] \dicTerm{söln··|aður} \dicIPA{{s}{\oe}{l}{n}{a}{ð}{\textscy}{\textsubring{r}}} \dicPos{adj}[3] \dicFlx{(f ‑uð)}[3] \dicDirectTranslationCS{zvadlý, uvadlý, uschlý}
\dicEntry[sölsa] \dicTerm{söls|a} \dicIPA{{s}{\oe}{l}{s}{a}} \dicPos{v}[1] \dicFlx{(‑aði)}[1] \dicFlx{acc} \dicPhraseIS{sölsa e‑ð undir sig} \dicDirectTranslationCS{uchvátit (co) pro sebe, zmocnit se (čeho)}
\dicEntry[sölt] \dicTerm{sölt\smash{\textsuperscript{1}}} \dicIPA{{s}{\oe}{\textsubring{l}}{\textsubring{d}}} \dicPos{n} \dicFlx{pl nom} \dicLink{salt}
\dicEntry[sölt] \dicTerm{sölt\smash{\textsuperscript{2}}} \dicIPA{{s}{\oe}{\textsubring{l}}{\textsubring{d}}} \dicPos{adj} \dicFlx{f sg nom pos} \dicLink{saltur}
\dicEntry[söltun] \dicTerm{sölt|un} \dicIPA{{s}{\oe}{\textsubring{l}}{\textsubring{d}}{\textscy}{\textsubring{n}}} \dicPos{f}[7] \dicFlx{(‑unar)}[12] \dicSynonym*{það að salta} \dicDirectTranslationCS{(o)solení, posolení}
\dicEntry[sölu] \dicTerm{sölu} \dicIPA{{s}{\oe}{\textlengthmark}{l}{\textscy}} \dicPos{f} \dicFlx{sg gen} \dicLink{sala}
\dicEntry[söluaðferð] \dicTerm{sölu··að·ferð} \dicIPA{{s}{\oe}{\textlengthmark}{l}{\textscy}{a}{ð}{f}{\textepsilon}{r}{\texttheta}} \dicPos{f}[7] \dicFlx{(‑ar, ‑ir)}[1] \dicDirectTranslationCS{prodejní strategie, marketing}
\dicEntry[söluaðili] \dicTerm{sölu··aðil|i} \dicIPA{{s}{\oe}{\textlengthmark}{l}{\textscy}{a}{ð}{\textsci}{l}{\textsci}} \dicPos{m}[1] \dicFlx{(‑a\,/\addthin ‑ja, ‑ar\,/\addthin ‑jar)}[15] \dicDirectTranslationCS{prodejce, prodejkyně, dealer(ka)}
\dicEntry[söluaukning] \dicTerm{sölu··auk·ning} \dicIPA{{s}{\oe}{\textlengthmark}{l}{\textscy}{\oe i}{h}{\r{g}}{n}{i}{\ng}{\r{g}}} \dicPos{f}[4] \dicFlx{(‑ar)}[7] \dicDirectTranslationCS{nárůst\,/\addthin zvýšení prodeje}
\dicEntry[sölubúð] \dicTerm{sölu··búð} \dicIPA{{s}{\oe}{\textlengthmark}{l}{\textscy}{\textsubring{b}}{u}{\texttheta}} \dicPos{f}[7] \dicFlx{(‑ar, ‑ir)}[1] \dicSynonym{búð} \dicDirectTranslationCS{prodejna, obchod}
\dicEntry[sölugengi] \dicTerm{sölu··gengi} \dicIPA{{s}{\oe}{\textlengthmark}{l}{\textscy}{\r{\textObardotlessj}}{ei}{\textltailn}{\r{\textObardotlessj}}{\textsci}} \dicPos{n}[2] \dicFlx{(‑s)}[20] \dicFieldCat{ekon.} \dicDirectTranslationCS{prodejní kurz}
\dicEntry[söluhagnaður] \dicTerm{sölu··hag·nað|ur} \dicIPA{{s}{\oe}{\textlengthmark}{l}{\textscy}{h}{a}{\r{g}}{n}{a}{ð}{\textscy}{\textsubring{r}}} \dicPos{m}[10] \dicFlx{(‑ar)}[9] \dicFieldCat{ekon.} \dicDirectTranslationCS{zisk\,/\addthin výnos z~prodeje}
\dicEntry[söluhæfur] \dicTerm{sölu··hæfur} \dicIPA{{s}{\oe}{\textlengthmark}{l}{\textscy}{h}{a}{i}{v}{\textscy}{\textsubring{r}}} \dicPos{adj}[1]\dicFlx{}[-1] \dicDirectTranslationCS{prodejný} \dicIndirectTranslationCS{(který je na prodej nebo jde na odbyt)} \dicExampleIS{söluhæf vara} \dicExampleCS{prodejné zboží}
\dicEntry[sölukostnaður] \dicTerm{sölu··kost·nað|ur} \dicIPA{{s}{\oe}{\textlengthmark}{l}{\textscy}{k\smash{\textsuperscript{h}}}{\textopeno}{s}{\textsubring{d}}{n}{a}{ð}{\textscy}{\textsubring{r}}} \dicPos{m}[10] \dicFlx{(‑ar)}[9] \dicDirectTranslationCS{prodejní náklady}
\dicEntry[sölumaður] \dicTerm{sölu··|maður} \dicIPA{{s}{\oe}{\textlengthmark}{l}{\textscy}{m}{a}{ð}{\textscy}{\textsubring{r}}} \dicPos{m}[13] \dicFlx{(‑manns, ‑menn)}[2] \dicSynonym{seljandi} \dicDirectTranslationCS{prodavač(ka), prodejce, prodejkyně, obchodník, obchodnice}
\dicEntry[sölumennska] \textls[15]{\dicTerm{sölu··mennsk|a} \dicIPA{{s}{\oe}{\textlengthmark}{l}{\textscy}{m}{\textepsilon}{n}{s}{\r{g}}{a}}} \dicPos{f}[1] \dicFlx{(‑u)}[5] \dicDirectTranslationCS{obchodování} \dicExampleIS{Hann hefur mikla reynslu af sölumennsku.} \dicExampleCS{Má velké zkušenosti s~obchodováním.}
\dicEntry[sölumiðstöð] \dicTerm{sölu··mið·stöð} \dicIPA{{s}{\oe}{\textlengthmark}{l}{\textscy}{m}{\textsci}{ð}{s}{\textsubring{d}}{\oe}{\texttheta}} \dicPos{f}[6] \dicFlx{(‑var, ‑var)}[1] \dicDirectTranslationCS{obchodní centrum}
\dicEntry[sölusamband] \dicTerm{sölu··sam·|band} \dicIPA{{s}{\oe}{\textlengthmark}{l}{\textscy}{s}{a}{m}{\textsubring{b}}{a}{n}{\textsubring{d}}} \dicPos{n}[2] \dicFlx{(‑bands, ‑bönd)}[8] \dicDirectTranslationCS{obchodní organizace}
\dicEntry[sölusamningur] \dicTerm{sölu··samn·ing|ur} \dicIPA{{s}{\oe}{\textlengthmark}{l}{\textscy}{s}{a}{m}{n}{i}{\ng}{\r{g}}{\textscy}{\textsubring{r}}} \dicPos{m}[6] \dicFlx{(‑s, ‑ar)}[8] \dicDirectTranslationCS{kupní smlouva, smlouva o~prodeji}
\dicEntry[söluskattur] \dicTerm{sölu··skatt|ur} \dicIPA{{s}{\oe}{\textlengthmark}{l}{\textscy}{s}{\r{g}}{a}{h}{\textsubring{d}}{\textscy}{\textsubring{r}}} \dicPos{m}[6] \dicFlx{(‑s)}[7] \dicFieldCat{ekon.} \dicDirectTranslationCS{daň z~prodeje}
\dicEntry[söluskrifstofa] \dicTerm{sölu··skrif·stof|a} \dicIPA{{s}{\oe}{\textlengthmark}{l}{\textscy}{s}{\r{g}}{r}{\textsci}{f}{s}{\textsubring{d}}{\textopeno}{v}{a}} \dicPos{f}[1] \dicFlx{(‑u, ‑ur)}[7] \dicDirectTranslationCS{prodejní kancelář\,/\addthin agentura}
\dicEntry[sölustaður] \dicTerm{sölu··stað|ur} \dicIPA{{s}{\oe}{\textlengthmark}{l}{\textscy}{s}{\textsubring{d}}{a}{ð}{\textscy}{\textsubring{r}}} \dicPos{m}[10] \dicFlx{(‑ar, ‑ir)}[14] \dicDirectTranslationCS{prodejní místo, prodejna}
\dicEntry[sölustjóri] \dicTerm{sölu··stjór|i} \dicIPA{{s}{\oe}{\textlengthmark}{l}{\textscy}{s}{\textsubring{d}}{j}{ou}{r}{\textsci}} \dicPos{m}[1] \dicFlx{(‑a, ‑ar)}[1] \dicDirectTranslationCS{ředitel(ka) prodeje}
\dicEntry[sölusvæði] \dicTerm{sölu··svæði} \dicIPA{{s}{\oe}{\textlengthmark}{l}{\textscy}{s}{v}{a}{i}{ð}{\textsci}} \dicPos{n}[2] \dicFlx{(‑s, ‑)}[14] \dicDirectTranslationCS{prodejní areál}
\dicEntry[sölusýning] \dicTerm{sölu··sýn·ing} \dicIPA{{s}{\oe}{\textlengthmark}{l}{\textscy}{s}{i}{n}{i}{\ng}{\r{g}}} \dicPos{f}[4] \dicFlx{(‑ar, ‑ar)}[5] \dicDirectTranslationCS{prodejní výstava}
\dicEntry[sölutorg] \dicTerm{sölu··torg} \dicIPA{{s}{\oe}{\textlengthmark}{l}{\textscy}{t\smash{\textsuperscript{h}}}{\textopeno}{r}{\r{g}}} \dicPos{n}[2] \dicFlx{(‑s, ‑)}[5] \dicDirectTranslationCS{tržiště, tržnice, trh}
\dicEntry[söluturn] \dicTerm{sölu··turn} \dicIPA{{s}{\oe}{\textlengthmark}{l}{\textscy}{t\smash{\textsuperscript{h}}}{\textscy}{r}{\textsubring{d}}{\textsubring{n}}} \dicPos{m}[4] \dicFlx{(‑s, ‑ar)}[3] \dicDirectTranslationCS{stánek, kiosk}
\dicEntry[söluvara] \dicTerm{sölu··|vara} \dicIPA{{s}{\oe}{\textlengthmark}{l}{\textscy}{v}{a}{r}{a}} \dicPos{f}[1] \dicFlx{(‑vöru, ‑vörur)}[8] \dicDirectTranslationCS{zboží, artikl, tovar} \dicIndirectTranslationCS{(k~prodeji)}
\dicEntry[söluverð] \dicTerm{sölu··verð} \dicIPA{{s}{\oe}{\textlengthmark}{l}{\textscy}{v}{\textepsilon}{r}{\texttheta}} \dicPos{n}[2] \dicFlx{(‑s)}[2] \dicDirectTranslationCS{prodejní cena}
\dicEntry[söm] \dicTerm{söm} \dicIPA{{s}{\oe}{\textlengthmark}{\textsubring{m}}} \dicPos{adj} \dicFlx{f sg nom pos} \dicLink{samur}
\dicEntry[sömdum] \dicTerm{sömdum} \dicIPA{{s}{\oe}{m}{\textsubring{d}}{\textscy}{\textsubring{m}}} \dicPos{v} \dicFlx{ind pf pl 1 pers} \dicLink{semja}
\dicEntry[sömuleiðis] \dicTerm{sömu··leiðis} \dicIPA{{s}{\oe}{\textlengthmark}{m}{\textscy}{l}{ei}{ð}{\textsci}{s}} \dicPos{adv} \textbf{1.} \dicDirectTranslationCS{obdobně, srovnatelně}  \textbf{2.} \dicDirectTranslationCS{nápodobně, také tak} \dicExampleIS{Takk fyrir síðast. -- Sömuleiðis.} \dicExampleCS{Díky za posledně. -- Nápodobně.}
\dicEntry[söng] \dicTerm{söng\smash{\textsuperscript{1}}} \dicIPA{{s}{\oe i}{\ng}{\r{g}}} \dicPos{m} \dicFlx{sg dat\,/\addthin acc} \dicLink{söngur}
\dicEntry[söng] \dicTerm{söng\smash{\textsuperscript{2}}} \dicIPA{{s}{\oe i}{\ng}{\r{g}}} \dicPos{adj} \dicFlx{f sg nom pos} \dicLink{sangur}
\dicEntry[söng] \dicTerm{söng\smash{\textsuperscript{3}}} \dicIPA{{s}{\oe i}{\ng}{\r{g}}} \dicPos{v} \dicFlx{ind pf sg 1 pers} \dicLink{syngja}
\dicEntry[söngbók] \dicTerm{söng··|bók} \dicIPA{{s}{\oe i}{\ng}{\r{g}}{\textsubring{b}}{ou}{\r{g}}} \dicPos{f}[8] \dicFlx{(‑bókar, ‑bækur)}[5] \dicDirectTranslationCS{zpěvník}
\dicEntry[söngfugl] \dicTerm{söng··fugl} \dicIPA{{s}{\oe i}{\ng}{\r{g}}{f}{\textscy}{\r{g}}{\textsubring{l}}} \dicPos{m}[4] \dicFlx{(‑s, ‑ar)}[3] \textbf{1.} \dicDirectTranslationCS{zpěvný pták}  \textbf{2.} \dicDirectTranslationCS{dobrý pěvec, dobrá pěvkyně} \dicIndirectTranslationCS{(o~člověku)}
\dicEntry[sönghneigður] \dicTerm{söng··hneigður} \dicIPA{{s}{\oe i}{\ng}{\r{g}}{\textsubring{n}}{ei}{\textbabygamma}{ð}{\textscy}{\textsubring{r}}} \dicPos{adj}[2]\dicFlx{}[-1] \dicDirectTranslationCS{muzikální, pěvecky nadaný}
\dicEntry[söngkona] \dicTerm{söng··kon|a} \dicIPA{{s}{\oe i}{\ng}{\r{g}}{k\smash{\textsuperscript{h}}}{\textopeno}{n}{a}} \dicPos{f}[1] \dicFlx{(‑u, ‑ur)}[27] \dicDirectTranslationCS{zpěvačka, pěvkyně}
\dicEntry[söngl] \dicTerm{söngl} \dicIPA{{s}{\oe i}{\ng}{\textsubring{l}}} \dicPos{n}[2] \dicFlx{(‑s)}[2] \dicDirectTranslationCS{broukání, pobrukování}
\dicEntry[söngla] \dicTerm{söngl|a} \dicIPA{{s}{\oe i}{\ng}{\r{g}}{l}{a}} \dicPos{v}[1] \dicFlx{(‑aði)}[1] \dicFlx{acc} \dicSynonym{raula} \dicDirectTranslationCS{broukat si, pobrukovat si, zpívat si tiše} \dicExampleIS{söngla lag} \dicExampleCS{broukat si písničku}
\dicEntry[sönglag] \dicTerm{söng··|lag} \dicIPA{{s}{\oe i}{\ng}{\r{g}}{l}{a}{x}} \dicPos{n}[2] \dicFlx{(‑lags, ‑lög)}[8] \dicDirectTranslationCS{melodie, popěvek, píseň}
\dicEntry[söngleikur] \dicTerm{söng··leik|ur} \dicIPA{{s}{\oe i}{\ng}{\r{g}}{l}{ei}{\r{g}}{\textscy}{\textsubring{r}}} \dicPos{m}[9] \dicFlx{(‑s, ‑ir)}[15] \dicFieldCat{hud.} \dicDirectTranslationCS{muzikál}
\dicEntry[söngmaður] \dicTerm{söng··|maður} \dicIPA{{s}{\oe i}{\ng}{\r{g}}{m}{a}{ð}{\textscy}{\textsubring{r}}} \dicPos{m}[13] \dicFlx{(‑manns, ‑menn)}[2] \dicDirectTranslationCS{zpěvák, pěvec}
\dicEntry[söngrödd] \dicTerm{söng··|rödd} \dicIPA{{s}{\oe i}{\ng}{\r{g}}{r}{\oe}{\textsubring{d}}} \dicPos{f}[7] \dicFlx{(‑raddar, ‑raddir)}[16] \dicDirectTranslationCS{pěvecký\,/\addthin zpěvný hlas};  \dicPhraseIS{hafa góða\,/\addthin vonda söngrödd} \dicDirectTranslationCS{mít dobrý\,/\addthin špatný (pěvecký) hlas}
\dicEntry[söngskrá] \dicTerm{söng··skrá} \dicIPA{{s}{\oe i}{\ng}{\r{g}}{s}{\r{g}}{r}{au}} \dicPos{f}[4] \dicFlx{(‑r\,/\addthin ‑ar, ‑r)}[21] \dicDirectTranslationCS{koncertní program} \dicIndirectTranslationCS{(seznam skladeb představených během koncertu)}
\dicEntry[söngstjóri] \dicTerm{söng··stjór|i} \dicIPA{{s}{\oe i}{\ng}{\r{g}}{s}{\textsubring{d}}{j}{ou}{r}{\textsci}} \dicPos{m}[1] \dicFlx{(‑a, ‑ar)}[1] \dicDirectTranslationCS{sbormistr(yně)}
\dicEntry[söngtexti] \dicTerm{söng··text|i} \dicIPA{{s}{\oe i}{\ng}{\r{g}}{t\smash{\textsuperscript{h}}}{\textepsilon}{x}{s}{\textsubring{d}}{\textsci}} \dicPos{m}[1] \dicFlx{(‑a, ‑ar)}[1] \dicDirectTranslationCS{text písně}
\dicEntry[söngtifa] \dicTerm{söng··tif|a} \dicIPA{{s}{\oe i}{\ng}{\r{g}}{t\smash{\textsuperscript{h}}}{\textsci}{v}{a}} \dicPos{f}[1] \dicFlx{(‑u, ‑ur)}[7] \dicFieldCat{zool.} \dicDirectTranslationCS{cikáda} \textit{(l.~{\textLA{Cicadidae}})}  \dicsymPhoto\ 
\dicFigure{1085.jpg}{Söngtifa}{Söngtifa - Zicha Ondřej, Biolib, Copyright/CC-BY-NC}
\dicEntry[söngtími] \dicTerm{söng··tím|i} \dicIPA{{s}{\oe i}{\ng}{\r{g}}{t\smash{\textsuperscript{h}}}{i}{m}{\textsci}} \dicPos{m}[1] \dicFlx{(‑a, ‑ar)}[1] \dicDirectTranslationCS{hodina zpěvu, lekce zpívání}
\dicEntry[söngur] \dicTerm{söng|ur} \dicsymFrequent\  \dicIPA{{s}{\oe i}{\ng}{\r{g}}{\textscy}{\textsubring{r}}} \dicPos{m}[8] \dicFlx{(‑s, ‑var)}[2] \textbf{1.} \dicSynonym*{það að syngja} \dicDirectTranslationCS{zpěv, zpívání, prozpěvování} \dicExampleIS{læra börnin söngva} \dicExampleCS{učit děti zpěvu}  \textbf{2.} \dicSynonym{nauð\smash{\textsuperscript{2}}} \dicDirectTranslationCS{škemrání, žadonění, kňourání}
\dicEntry[söngvakeppni] \dicTerm{söngva··keppn|i} \dicIPA{{s}{\oe i}{\ng}{\r{g}}{v}{a}{c\smash{\textsuperscript{h}}}{\textepsilon}{h}{\textsubring{b}}{n}{\textsci}} \dicPos{f}[3] \dicFlx{(‑i, ‑ir)}[1] \dicDirectTranslationCS{zpěvácká\,/\addthin pěvecká soutěž}
\dicEntry[söngvari] \dicTerm{söngv··ar|i} \dicIPA{{s}{\oe i}{\ng}{\r{g}}{v}{a}{r}{\textsci}} \dicPos{m}[1] \dicFlx{(‑a, ‑ar)}[13] \dicDirectTranslationCS{zpěvák, pěvec}
\dicEntry[söngvaskáld] \dicTerm{söngva··skáld} \dicIPA{{s}{\oe i}{\ng}{\r{g}}{v}{a}{s}{\r{g}}{au}{l}{\textsubring{d}}} \dicPos{n}[2] \dicFlx{(‑s, ‑)}[5] \dicDirectTranslationCS{písničkář(ka), textař(ka)}
\dicEntry[söngvinn] \dicTerm{söngvinn} \dicIPA{{s}{\oe i}{\ng}{\r{g}}{v}{\textsci}{\textsubring{n}}} \dicPos{adj}[7]\dicFlx{}[-7] \dicSynonym{sönghneigður} \dicDirectTranslationCS{muzikální, pěvecky nadaný}
\dicEntry[söngþröstur] \dicTerm{söng··|þröstur} \dicIPA{{s}{\oe i}{\ng}{\r{g}}{\texttheta}{r}{\oe}{s}{\textsubring{d}}{\textscy}{\textsubring{r}}} \dicPos{m}[11] \dicFlx{(‑þrastar, ‑þrestir)}[5] \dicFieldCat{zool.} \dicDirectTranslationCS{drozd zpěvný} \textit{(l.~{\textLA{Turdus philomelos}})}  \dicsymPhoto\ 
\dicFigure{ds_image_songthrostur_0_1.jpg}{Söngþröstur}{Söngþröstur - Andreas Trepte, CC BY-SA 2.5}
\dicEntry[sönn] \dicTerm{sönn} \dicIPA{{s}{\oe}{n}{\textlengthmark}} \dicPos{adj} \dicFlx{f sg nom pos} \dicLink{sannur\smash{\textsuperscript{2}}}
\dicEntry[sönnun] \dicTerm{sönnun} \dicsymFrequent\  \dicIPA{{s}{\oe}{n}{\textlengthmark}{\textscy}{\textsubring{n}}} \dicPos{f}[7] \dicFlx{(sönnunar, sannanir)}[11] \dicDirectTranslationCS{důkaz, dokázání} \dicExampleIS{nægar sannanir fyrir sekt hans} \dicExampleCS{dostatečné důkazy pro prokázání jeho viny}
\dicEntry[sönnunarbyrði] \dicTerm{sönnunar··byrð|i} \dicIPA{{s}{\oe}{n}{\textlengthmark}{\textscy}{n}{a}{r}{\textsubring{b}}{\textsci}{r}{ð}{\textsci}} \dicPos{f}[2] \dicFlx{(‑i\,/\addthin ‑ar)}[4] \dicFieldCat{práv.} \dicDirectTranslationCS{důkazní břemeno}
\dicEntry[sönnunargagn] \dicTerm{sönnunar··|gagn} \dicIPA{{s}{\oe}{n}{\textlengthmark}{\textscy}{n}{a}{r}{\r{g}}{a}{\r{g}}{\textsubring{n}}} \dicPos{n}[2] \dicFlx{(‑gagns, ‑gögn)}[8] \dicFieldCat{práv.} \dicDirectTranslationCS{důkaz, předmět doličný}
\dicEntry[sönnunarskylda] \dicTerm{sönnunar··skyld|a} \dicIPA{{s}{\oe}{n}{\textlengthmark}{\textscy}{n}{a}{\textsubring{r}}{s}{\r{\textObardotlessj}}{\textsci}{l}{\textsubring{d}}{a}} \dicPos{f}[1] \dicFlx{(‑u)}[5] \dicDirectTranslationCS{důkazní břemeno}
\dicEntry[sönnur] \dicTerm{sönnur} \dicIPA{{s}{\oe}{n}{\textlengthmark}{\textscy}{\textsubring{r}}} \dicPos{f}[1] \dicFlx{pl}[12] \dicPhraseIS{færa sönnur á e‑ð} \dicDirectTranslationCS{dokázat (co), prokázat (co)}
\dicEntry[sötra] \dicTerm{sötr|a} \dicIPA{{s}{\oe}{\textlengthmark}{\textsubring{d}}{r}{a}} \dicPos{v}[1] \dicFlx{(‑aði)}[1] \dicFlx{acc} \dicDirectTranslationCS{srkat, usrkávat} \dicExampleIS{sötra kaffið} \dicExampleCS{srkat kávu}
\dicEntry[söx] \dicTerm{söx} \dicIPA{{s}{\oe}{x}{s}} \dicPos{n} \dicFlx{pl nom} \dicLink{sax}

\else
\foreach \l in \alphabetF{
  \dicLetter{l}{letter15}
\dicEntry[l] \dicTerm{l} \dicPos{zkr} \dicPhraseIS{lítri} \dicDirectTranslationCS{litr}
\dicEntry[l.] \dicTerm{l.} \dicPos{zkr} \dicPhraseIS{lýsingarorð} \dicFieldCat{jaz.} \dicDirectTranslationCS{přídavné jméno}
\dicEntry[labb] \dicTerm{labb} \dicIPA{{l}{a}{\textsubring{b}}{\textlengthmark}} \dicPos{n}[2] \dicFlx{(‑s)}[2] \dicDirectTranslationCS{procházka, špacír}
\dicEntry[labba] \dicTerm{labb|a} \dicsymFrequent\  \dicIPA{{l}{a}{\textsubring{b}}{\textlengthmark}{a}} \dicPos{v}[1] \dicFlx{(‑aði)}[13] \dicDirectTranslationCS{projít se, procházet se, jít na procházku, špacírovat} \dicExampleIS{Ég labbaði í bæinn.} \dicExampleCS{Prošel jsem se do města.}
\dicEntry[labbakútur] \dicTerm{labba··kút|ur} \dicIPA{{l}{a}{\textsubring{b}}{\textlengthmark}{a}{k\smash{\textsuperscript{h}}}{u}{\textsubring{d}}{\textscy}{\textsubring{r}}} \dicPos{m}[6] \dicFlx{(‑s, ‑ar)}[22] \textbf{1.} \dicSynonym{barn} \dicDirectTranslationCS{klouček, prcek}  \textbf{2.} \dicSynonym*{mannleysa} \dicDirectTranslationCS{trdlo, janek}
\dicEntry[laða] \dicTerm{lað|a} \dicIPA{{l}{a}{\textlengthmark}{ð}{a}} \dicPos{v}[1] \dicFlx{(‑aði)}[13] \dicFlx{acc} \dicDirectTranslationCS{přitahovat, (při)lákat, (při)vábit} \dicExampleIS{laða e‑n að sér} \dicExampleCS{přilákat (koho) k~sobě};  \dicIdiom{laða}[fram]{ \dicPhraseIS{laða e‑u fram}} \dicDirectTranslationCS{vyvolat (co), přivolat (co), přivodit (co)};  \dicIdiom{laðast}{ \dicPhraseIS{laðast}} \dicFlx{refl} \dicDirectTranslationCS{být přitahován} \dicExampleIS{laðast að e‑u} \dicExampleCS{být přitahován (čím)}
\dicEntry[lafa] \dicTerm{laf|a} \dicIPA{{l}{a}{\textlengthmark}{v}{a}} \dicPos{v}[2] \dicFlx{(‑ði, ‑að)}[125] \textbf{1.} \dicSynonym{hanga} \dicDirectTranslationCS{viset} \dicExampleIS{láta fæturna lafa} \dicExampleCS{nechat nohy viset}  \textbf{2.} \dicDirectTranslationCS{prodlévat, zdržovat se} \dicExampleIS{lafa í þorpinu} \dicExampleCS{zdržovat se na vsi}
\dicEntry[lag] \dicTerm{lag} \dicsymFrequent\  \dicIPA{{l}{a}{\textlengthmark}{x}} \dicPos{n}[2] \dicFlx{(‑s, lög)}[8] \textbf{1.} \dicSynonym{samfella} \dicDirectTranslationCS{vrstva} \dicExampleIS{lag af kolum} \dicExampleCS{vrstva uhlí}  \textbf{2.} \dicDirectTranslationCS{(pro)bodnutí (zbraní ap.)} \dicExampleIS{Lagið nam í hjartastað.} \dicExampleCS{Bodnutí zasáhlo srdce.}  \textbf{3.} \dicSynonym{aðferð} \dicDirectTranslationCS{způsob, metoda};  \dicPhraseIS{finna lagið á e‑u} \dicDirectTranslationCS{přijít na to, jak se (co) dělá};  \dicPhraseIS{hafa lag á e‑u} \dicDirectTranslationCS{mít talent na (co)};  \dicPhraseIS{koma e‑m (upp) á lagið með e‑ð} \dicDirectTranslationCS{naučit (koho), jak se (co) dělá};  \dicPhraseIS{með þessu lagi} \dicFlx{adv} \dicDirectTranslationCS{takto, tímto způsobem}  \textbf{4.} \dicSynonym{magn} \dicDirectTranslationCS{množství};  \dicPhraseIS{í einu lagi} \dicFlx{adv} \dicSynonym*{í heild} \dicDirectTranslationCS{najednou, naráz};  \dicPhraseIS{í fyrsta lagi} \dicFlx{adv} {\textbf{a.}} \dicDirectTranslationCS{ne dříve než} \dicExampleIS{Ég kem í fyrsta lagi kl. þrjú.} \dicExampleCS{Nepřijdu dřív než ve tři.};  {\textbf{b.}} \dicDirectTranslationCS{za prvé, v~první řadě};  \dicPhraseIS{í hæsta lagi} \dicFlx{adv} \dicDirectTranslationCS{nanejvýš, maximálně, nejvýše} \dicExampleIS{Þetta tekur í hæsta lagi fimm mínútur.} \dicExampleCS{Zabere to nejvýše pět minut.};  \dicPhraseIS{í meira lagi} \dicFlx{adv} \dicDirectTranslationCS{dosti, značně};  \dicPhraseIS{í mesta lagi} \dicFlx{adv} \dicDirectTranslationCS{ne více než};  \dicPhraseIS{í öðru lagi} \dicFlx{adv} \dicDirectTranslationCS{za druhé}  \textbf{5.} \dicSynonym{regla} \dicDirectTranslationCS{řád, pořádek};  \dicPhraseIS{allt í lagi!} \dicDirectTranslationCS{ok!, všechno v~pořádku!};  \dicPhraseIS{e‑að er fjarri lagi} \dicDirectTranslationCS{(co) není zdaleka v~pořádku};  \dicPhraseIS{koma e‑u í lag} \dicSynonym{lagfæra} \dicDirectTranslationCS{dát (co) do pořádku};  \dicPhraseIS{nærri lagi} \dicFlx{adv} \dicDirectTranslationCS{téměř v~pořádku}  \textbf{6.} \dicSynonym{tónverk} \dicDirectTranslationCS{píseň, písnička, melodie} \dicExampleIS{lag eftir e‑n} \dicExampleCS{(čí) písnička};  \dicPhraseIS{fara út af laginu} \dicDirectTranslationCS{vypadnout z~rytmu (v~písni), zazpívat falešný tón};  \dicPhraseIS{taka lagið} \dicDirectTranslationCS{(začít) zpívat}  \textbf{7.} \dicSynonym{verðlag} \dicDirectTranslationCS{cena, hodnota};  \dicPhraseIS{leggja lag á e‑ð} \dicDirectTranslationCS{ocenit (co), ohodnotit (co), určit cenu (čeho)}  \textbf{8.} \dicSynonym{samneyti} \dicDirectTranslationCS{společnost, přítomnost} \dicIndirectTranslationCS{(společenský styk)};  \dicPhraseIS{leggja lag sitt við e‑n} \dicDirectTranslationCS{dělat (komu) společnost}  \textbf{9.} \dicPhraseIS{lög} \dicFlx{pl} \dicLink{lög};  \dicIdiom{lag}{ \dicPhraseIS{slá e‑n út af laginu}} \dicDirectTranslationCS{vyvést (koho) z~míry, vytočit (koho)}; { \dicPhraseIS{sæta lagi}} \dicDirectTranslationCS{počkat si (na příležitost)}
\dicEntry[laga] \dicTerm{lag|a} \dicsymFrequent\  \dicIPA{{l}{a}{\textlengthmark}{\textbabygamma}{a}} \dicPos{v}[1] \dicFlx{(‑aði)}[13] \dicFlx{acc} \textbf{1.} \dicSynonym{prýða} \dicDirectTranslationCS{upravit, upravovat, urovnat, urovnávat} \dicExampleIS{laga á sér hárið} \dicExampleCS{upravit si vlasy}  \textbf{2.} \dicSynonym*{koma e‑u í lag} \dicDirectTranslationCS{opravit, opravovat, spravit, spravovat, dát\,/\addthin dávat do pořádku} \dicExampleIS{laga vélina} \dicExampleCS{opravit stroj}  \textbf{3.} \dicDirectTranslationCS{připravit, připravovat (kávu ap.)} \dicExampleIS{laga kaffi og te} \dicExampleCS{připravit kávu a~čaj}  \textbf{4.} \dicSynonym{streyma} \dicDirectTranslationCS{proudit, prýštit (krev z~rány ap.)};  \dicIdiom{laga}[að]{ \dicPhraseIS{laga sig að e‑u}} \dicSynonym{aðhæfa} \dicDirectTranslationCS{přizpůsobit se (čemu)};  \dicIdiom{laga}[til]{ \dicPhraseIS{laga til}} \dicDirectTranslationCS{uklízet, poklízet, uklidit} \dicExampleIS{laga til í íbúðinni} \dicExampleCS{uklidit v~bytě}; { \dicPhraseIS{laga sig til}} \dicDirectTranslationCS{poupravit se};  \dicIdiom{lagast}{ \dicPhraseIS{lagast}} \dicFlx{refl} {\textbf{a.}} \dicDirectTranslationCS{(z)lepšit se, vylepšit se} \dicExampleIS{Ástandið lagast af sjálfu sér.} \dicExampleCS{Situace se zlepšuje sama od sebe.};  {\textbf{b.}} \dicDirectTranslationCS{ustoupit, povolit (bolest ap.)}
\dicEntry[lagaákvæði] \dicTerm{laga··á·kvæði} \dicIPA{{l}{a}{\textlengthmark}{\textbabygamma}{a}{au}{k\smash{\textsuperscript{h}}}{v}{a}{i}{ð}{\textsci}} \dicPos{n}[2] \dicFlx{(‑s, ‑)}[14] \dicFieldCat{práv.} \dicDirectTranslationCS{právní předpis\,/\addthin ustanovení}
\dicEntry[lagabálkur] \dicTerm{laga··bálk|ur} \dicIPA{{l}{a}{\textlengthmark}{\textbabygamma}{a}{\textsubring{b}}{au}{\textsubring{l}}{\r{g}}{\textscy}{\textsubring{r}}} \dicPos{m}[6] \dicFlx{(‑s, ‑ar)}[24] \dicFieldCat{práv.} \dicDirectTranslationCS{kodex, zákoník}
\dicEntry[lagabreyting] \dicTerm{laga··breyt·ing} \dicIPA{{l}{a}{\textlengthmark}{\textbabygamma}{a}{\textsubring{b}}{r}{ei}{\textsubring{d}}{i}{\ng}{\r{g}}} \dicPos{f}[4] \dicFlx{(‑ar, ‑ar)}[5] \dicFieldCat{práv.} \dicDirectTranslationCS{novela zákona}
\dicEntry[lagabrot] \dicTerm{laga··brot} \dicIPA{{l}{a}{\textlengthmark}{\textbabygamma}{a}{\textsubring{b}}{r}{\textopeno}{\textsubring{d}}} \dicPos{n}[2] \dicFlx{(‑s, ‑)}[5] \dicDirectTranslationCS{porušení zákona}
\dicEntry[lagadeild] \dicTerm{laga··deild} \dicIPA{{l}{a}{\textlengthmark}{\textbabygamma}{a}{\textsubring{d}}{ei}{l}{\textsubring{d}}} \dicPos{f}[7] \dicFlx{(‑ar, ‑ir)}[1] \dicDirectTranslationCS{právnická fakulta}
\dicEntry[lagafrumvarp] \dicTerm{laga··frum·|varp} \dicIPA{{l}{a}{\textlengthmark}{\textbabygamma}{a}{f}{r}{\textscy}{m}{v}{a}{\textsubring{r}}{\textsubring{b}}} \dicPos{n}[2] \dicFlx{(‑varps, ‑vörp)}[8] \dicFieldCat{práv.} \dicDirectTranslationCS{návrh zákona}
\dicEntry[lagagrein] \dicTerm{laga··grein} \dicIPA{{l}{a}{\textlengthmark}{\textbabygamma}{a}{\r{g}}{r}{ei}{\textsubring{n}}} \dicPos{f}[4] \dicFlx{(‑ar, ‑ar)}[1] \dicFieldCat{práv.} \dicSynonym{grein} \dicDirectTranslationCS{odstavec zákona, paragraf}
\dicEntry[lagaheimild] \dicTerm{laga··heimild} \dicIPA{{l}{a}{\textlengthmark}{\textbabygamma}{a}{h}{ei}{m}{\textsci}{l}{\textsubring{d}}} \dicPos{f}[7] \dicFlx{(‑ar, ‑ir)}[1] \dicDirectTranslationCS{moc zákona, zákonnost}
\dicEntry[lagahöfundur] \dicTerm{laga··höf·und|ur} \dicIPA{{l}{a}{\textlengthmark}{\textbabygamma}{a}{h}{\oe}{v}{\textscy}{n}{\textsubring{d}}{\textscy}{\textsubring{r}}} \dicPos{m}[6] \dicFlx{(‑ar, ‑ar)}[56] \dicDirectTranslationCS{písničkář(ka), skladatel(ka)}
\dicEntry[lagalegur] \dicTerm{laga··legur} \dicIPA{{l}{a}{\textlengthmark}{\textbabygamma}{a}{l}{\textepsilon}{\textbabygamma}{\textscy}{\textsubring{r}}} \dicPos{adj}[1]\dicFlx{}[-8] \dicDirectTranslationCS{zákonný, právní}
\dicEntry[laganir] \dicTerm{laganir} \dicIPA{{l}{a}{\textlengthmark}{\textbabygamma}{a}{n}{\textsci}{\textsubring{r}}} \dicPos{f} \dicFlx{pl nom} \dicLink{lögun}
\dicEntry[lagar] \dicTerm{lagar} \dicIPA{{l}{a}{\textlengthmark}{\textbabygamma}{a}{\textsubring{r}}} \dicPos{m} \dicFlx{sg gen} \dicLink{lögur}
\dicEntry[lagaregla] \dicTerm{laga··regl|a} \dicIPA{{l}{a}{\textlengthmark}{\textbabygamma}{a}{r}{\textepsilon}{\r{g}}{l}{a}} \dicPos{f}[1] \dicFlx{(‑u, ‑ur)}[19] \dicFieldCat{práv.} \dicDirectTranslationCS{právní předpis}
\dicEntry[lagasafn] \dicTerm{laga··|safn} \dicIPA{{l}{a}{\textlengthmark}{\textbabygamma}{a}{s}{a}{\textsubring{b}}{\textsubring{n}}} \dicPos{n}[2] \dicFlx{(‑safns, ‑söfn)}[8] \dicFieldCat{práv.} \dicDirectTranslationCS{sbírka zákonů}
\dicEntry[lager] \dicTerm{lager} \dicIPA{{l}{a}{\textlengthmark}{\textbabygamma}{\textepsilon}{\textsubring{r}}} \dicPos{m}[4] \dicFlx{(‑s, ‑ar)}[14] \textbf{1.} \dicSynonym{vörubirgðir} \dicDirectTranslationCS{sklad, zásoba} \dicExampleIS{Vörur eru sendar samdægurs séu þær til á lager.} \dicExampleCS{Zboží je odesíláno ve stejný den, pokud je na skladě.};  \dicPhraseIS{hafa e‑ð á lager} \dicDirectTranslationCS{mít (co) na skladě}  \textbf{2.} \dicSynonym{vörugeymsla} \dicDirectTranslationCS{sklad, skladiště} \dicExampleIS{vinna á lager} \dicExampleCS{pracovat ve skladu}
\dicEntry[lagermaður] \dicTerm{lager··|maður} \dicIPA{{l}{a}{\textlengthmark}{\textbabygamma}{\textepsilon}{r}{m}{a}{ð}{\textscy}{\textsubring{r}}} \dicPos{m}[13] \dicFlx{(‑manns, ‑menn)}[2] \dicDirectTranslationCS{skladník, skladnice}
\dicEntry[lagfæra] \dicTerm{lag··fær|a} \dicIPA{{l}{a}{x}{f}{a}{i}{r}{a}} \dicPos{v}[2] \dicFlx{(‑ði, ‑t)}[105] \dicFlx{acc} \textbf{1.} \dicDirectTranslationCS{dát\,/\addthin uvést do pořádku, opravit} \dicExampleIS{lagfæra skemmdir á húsi} \dicExampleCS{dát do pořádku poškození na domě}  \textbf{2.} \dicSynonym{endurbæta} \dicDirectTranslationCS{zlepšit, vylepšit, uzpůsobit}
\dicEntry[lagfæring] \dicTerm{lag··fær·ing} \dicIPA{{l}{a}{x}{f}{a}{i}{r}{i}{\ng}{\r{g}}} \dicPos{f}[4] \dicFlx{(‑ar, ‑ar)}[5] \dicSynonym{endurbót} \dicDirectTranslationCS{oprava, náprava, korekce}
\dicEntry[laggar] \dicTerm{laggar} \dicIPA{{l}{a}{\r{g}}{\textlengthmark}{a}{\textsubring{r}}} \dicPos{f} \dicFlx{sg gen} \dicLink{lögg}
\dicEntry[laggir] \dicTerm{laggir} \dicIPA{{l}{a}{\r{\textObardotlessj}}{\textlengthmark}{\textsci}{\textsubring{r}}} \dicPos{f} \dicFlx{pl nom} \dicLink{lögg}
\dicEntry[laggóður] \textls[15]{\dicTerm{lag··|góður} \dicIPA{{l}{a}{\textbabygamma}{\r{g}}{ou}{ð}{\textscy}{\textsubring{r}}} \dicPos{adj}[11] \dicFlx{(comp ‑betri, sup ‑bestur)}[1] \dicSynonym{hnyttinn} \dicDirectTranslationCS{výstižný, trefný} \dicExampleIS{stutt og laggott svar} \dicExampleCS{krátká a~výstižná odpověď}}
\dicEntry[laghendur] \dicTerm{lag··hendur}\dicTerm{, laghentur} \dicIPA{{l}\-{a}\-{\textlengthmark}\-{x}\-{h}\-{\textepsilon}\-{n}\-{\textsubring{d}}\-{\textscy}\-{\textsubring{r}}\-} \dicPos{adj}[2]\dicFlx{}[-14] \dicSynonym{handlaginn} \dicDirectTranslationCS{obratný, zručný}
\dicEntry[laghentur] \dicTerm{lag··hentur} \dicIPA{{l}{a}{\textlengthmark}{x}{h}{\textepsilon}{\textsubring{n}}{\textsubring{d}}{\textscy}{\textsubring{r}}} \dicPos{adj}[1]\dicFlx{}[-10] \dicLink{laghendur}
\dicEntry[laginn] \dicTerm{laginn} \dicIPA{{l}{a}{i}{j}{\textlengthmark}{\textsci}{\textsubring{n}}} \dicPos{adj}[6]\dicFlx{}[-3] \dicDirectTranslationCS{zručný, šikovný};  \dicPhraseIS{e‑r er laginn í höndunum} \dicDirectTranslationCS{(kdo) má šikovné ruce};  \dicPhraseIS{e‑r er laginn við e‑n} \dicDirectTranslationCS{(kdo) má dobrý přístup ke (komu) (zvířatům ap.)}
\dicEntry[lagkaka] \dicTerm{lag··|kaka} \dicIPA{{l}{a}{\textlengthmark}{\textbabygamma}{k\smash{\textsuperscript{h}}}{a}{\r{g}}{a}} \dicPos{f}[1] \dicFlx{(‑köku, ‑kökur)}[20] \dicFieldCat{kulin.} \dicDirectTranslationCS{řezy (čokoládové ap.), vrstvený koláč, dortík, zákusek}
\dicEntry[laglaus] \dicTerm{lag··laus} \dicIPA{{l}{a}{\textbabygamma}{l}{\oe i}{s}} \dicPos{adj}[5]\dicFlx{}[-1] \dicDirectTranslationCS{(jsoucí) bez hudebního sluchu, hluchý (k~hudbě ap.)}
\dicEntry[laglegur] \dicTerm{lag··legur} \dicsymFrequent\  \dicIPA{{l}{a}{\textbabygamma}{l}{\textepsilon}{\textbabygamma}{\textscy}{\textsubring{r}}} \dicPos{adj}[1]\dicFlx{}[-8] \textbf{1.} \dicSynonym{fríður} \dicDirectTranslationCS{pohledný, půvabný} \dicExampleIS{lagleg stúlka} \dicExampleCS{pohledná dívka}  \textbf{2.} \dicSynonym{snotur} \dicDirectTranslationCS{krásný, hezký}  \textbf{3.} \dicDirectTranslationCS{pěkný, pořádný} \dicIndirectTranslationCS{(k~zdůraznění)}
\dicEntry[laglína] \dicTerm{lag··lín|a} \dicIPA{{l}{a}{\textbabygamma}{l}{i}{n}{a}} \dicPos{f}[1] \dicFlx{(‑u, ‑ur)}[7] \dicDirectTranslationCS{melodie}
\dicEntry[lagnar] \dicTerm{lagnar} \dicIPA{{l}{a}{\r{g}}{n}{a}{\textsubring{r}}} \dicPos{f} \dicFlx{sg gen} \dicLink{lögn}
\dicEntry[lagni] \dicTerm{lagn|i} \dicIPA{{l}{a}{\r{g}}{n}{\textsci}} \dicPos{f}[3] \dicFlx{(‑i)}[3] \dicSynonym{leikni} \dicDirectTranslationCS{šikovnost, zručnost, obratnost}
\dicEntry[lagning] \dicTerm{lag··ning} \dicIPA{{l}{a}{\r{g}}{n}{i}{\ng}{\r{g}}} \dicPos{f}[4] \dicFlx{(‑ar, ‑ar)}[5] \textbf{1.} \dicDirectTranslationCS{kladení, pokládání, instalace} \dicExampleIS{lagning vegar} \dicExampleCS{kladení cesty}  \textbf{2.} \dicSynonym*{hárlagning} \dicDirectTranslationCS{úprava\,/\addthin upravení vlasů}
\dicEntry[lagnir] \dicTerm{lagnir} \dicIPA{{l}{a}{\r{g}}{n}{\textsci}{\textsubring{r}}} \dicPos{f} \dicFlx{pl nom} \dicLink{lögn}
\dicEntry[lagskipting] \dicTerm{lag··skipt·ing} \dicIPA{{l}{a}{x}{s}{\r{\textObardotlessj}}{\textsci}{f}{\textsubring{d}}{i}{\ng}{\r{g}}} \dicPos{f}[4] \dicFlx{(‑ar)}[7] \dicDirectTranslationCS{(roz)vrstvení}
\dicEntry[lagsmaður] \dicTerm{lags··|maður} \dicIPA{{l}{a}{x}{s}{m}{a}{ð}{\textscy}{\textsubring{r}}} \dicPos{m}[13] \dicFlx{(‑manns, ‑menn)}[2] \dicSynonym{félagi} \dicDirectTranslationCS{kamarád(ka), přítel(kyně), společník, společnice}
\dicEntry[lagviss] \dicTerm{lag··viss} \dicIPA{{l}{a}{\textbabygamma}{v}{\textsci}{s}} \dicPos{adj}[5]\dicFlx{}[-1] \dicDirectTranslationCS{zpívající čistě, muzikální}
\dicEntry[lak] \dicTerm{lak\smash{\textsuperscript{1}}} \dicIPA{{l}{a}{\textlengthmark}{\r{g}}} \dicPos{n}[2] \dicFlx{(‑s, lök)}[8] \dicDirectTranslationCS{prostěradlo}
\dicEntry[lak] \dicTerm{lak\smash{\textsuperscript{2}}} \dicIPA{{l}{a}{\textlengthmark}{\r{g}}} \dicPos{v} \dicFlx{ind pf sg 1 pers} \dicLink{leka}
\dicEntry[lakk] \dicTerm{lakk} \dicIPA{{l}{a}{h}{\r{g}}} \dicPos{n}[2] \dicFlx{(‑s, lökk)}[8] \textbf{1.} \dicSynonym*{lakkmálning} \dicDirectTranslationCS{lak, nátěr}  \textbf{2.} \dicSynonym{innsigli} \dicDirectTranslationCS{pečeť}  \textbf{3.} \dicSynonym*{innsiglislakk} \dicDirectTranslationCS{pečetní vosk}
\dicEntry[lakka] \dicTerm{lakk|a} \dicIPA{{l}{a}{h}{\r{g}}{a}} \dicPos{v}[1] \dicFlx{(‑aði)}[13] \dicFlx{acc} \textbf{1.} \dicDirectTranslationCS{(na)lakovat} \dicExampleIS{lakka gluggana} \dicExampleCS{nalakovat okna}  \textbf{2.} \dicDirectTranslationCS{zapečetit} \dicExampleIS{lakka bréf} \dicExampleCS{zapečetit dopis}
\dicEntry[lakkrís] \dicTerm{lakk··rís} \dicIPA{{l}{a}{h}{\r{g}}{r}{i}{s}} \dicPos{m}[4] \dicFlx{(‑s)}[18] \dicDirectTranslationCS{lékořice} \dicExampleIS{fá einn poka af lakkrís} \dicExampleCS{dostat jeden pytlík lékořice}
\dicEntry[lakkrísrót] \dicTerm{lakk·rís··|rót} \dicIPA{{l}{a}{h}{\r{g}}{r}{i}{s}{r}{ou}{\textsubring{d}}} \dicPos{f}[8] \dicFlx{(‑rótar, ‑rætur)}[5] \dicFieldCat{bot.} \dicDirectTranslationCS{kořen lékořice lysé} \textit{(l.~{\textLA{Glycyrrhiza glabra}})}
\dicEntry[laktósi] \dicTerm{laktós|i} \dicIPA{{l}{a}{x}{\textsubring{d}}{ou}{s}{\textsci}} \dicPos{m}[1] \dicFlx{(‑a)}[3] \dicFieldCat{chem.} \dicSynonym*{mjólkursykur} \dicDirectTranslationCS{laktóza, mléčný cukr}
\dicEntry[lakur] \dicTerm{lakur} \dicIPA{{l}{a}{\textlengthmark}{\r{g}}{\textscy}{\textsubring{r}}} \dicPos{adj}[1] \dicFlx{(f lök)}[2] \textbf{1.} \dicSynonym{lélegur} \dicDirectTranslationCS{chabý, mizerný (výsledek ap.)} \dicExampleIS{lakur árangur} \dicExampleCS{mizerný výsledek}  \textbf{2.} \dicSynonym{lasinn} \dicDirectTranslationCS{chorý, neduživý}  \textbf{3.} \dicSynonym{tæpur} \dicDirectTranslationCS{necelý, pouhý}
\dicEntry[lama] \dicTerm{lam|a} \dicIPA{{l}{a}{\textlengthmark}{m}{a}} \dicPos{v}[1] \dicFlx{(‑aði)}[13] \dicDirectTranslationCS{ochromit, paralyzovat} \dicExampleIS{Sjúkdómur lamar líkamann.} \dicExampleCS{Nemoc ochromuje tělo.};  \dicIdiom{lamast}{ \dicPhraseIS{lamast}} \dicFlx{refl} {\textbf{a.}} \dicDirectTranslationCS{ochrnout (o~člověku ap.)} \dicExampleIS{lamast við slysið} \dicExampleCS{ochrnout při nehodě};  {\textbf{b.}} \dicDirectTranslationCS{být paralyzován\,/\addthin ochromen (dopravní ruch ap.)}
\dicEntry[lamadýr] \dicTerm{lama··dýr} \dicIPA{{l}{a}{\textlengthmark}{m}{a}{\textsubring{d}}{i}{\textsubring{r}}} \dicPos{n}[2] \dicFlx{(‑s, ‑)}[5] \dicFieldCat{zool.} \dicDirectTranslationCS{lama} \textit{(l.~{\textLA{Lama}})}  \dicsymPhoto\ 
\dicFigure{ds_image_lamadyr_0_1.jpg}{Lamadýr}{Lamadýr - Luca Galuzzi, CC BY-SA 2.5}
\dicEntry[lamaður] \dicTerm{lam··aður} \dicIPA{{l}{a}{\textlengthmark}{m}{a}{ð}{\textscy}{\textsubring{r}}} \dicPos{adj}[3] \dicFlx{(f lömuð)}[1] \dicDirectTranslationCS{chromý, ochrnutý}
\dicEntry[lamasess] \dicTerm{lama··sess} \dicIPA{{l}{a}{\textlengthmark}{m}{a}{s}{\textepsilon}{s}} \dicPos{m}[4] \dicFlx{(‑)}[34] \dicPhraseIS{liggja\,/\addthin vera í lamasessi} \dicDirectTranslationCS{být mimo provoz, nefungovat (o~věcech)}
\dicEntry[lamb] \dicTerm{lamb} \dicsymFrequent\  \dicIPA{{l}{a}{m}{\textsubring{b}}} \dicPos{n}[2] \dicFlx{(‑s, lömb)}[8] \textbf{1.} \dicDirectTranslationCS{jehně} \dicExampleIS{nýfædd lömb} \dicExampleCS{novorozená jehňata}  \textbf{2.} \dicDirectTranslationCS{drahoušek, miláček};  \dicPhraseIS{launa e‑m lambið gráa} \dicLangCat{přen.} \dicDirectTranslationCS{oplatit (komu) stejnou mincí}
\dicEntry[lambagras] \dicTerm{lamba··|gras} \dicIPA{{l}{a}{m}{\textsubring{b}}{a}{\r{g}}{r}{a}{s}} \dicPos{n}[2] \dicFlx{(‑grass, ‑grös)}[8] \dicFieldCat{bot.} \dicDirectTranslationCS{silenka bezlodyžná} \textit{(l.~{\textLA{Silene acaulis}})}  \dicsymPhoto\ 
\dicFigure{9260.jpg}{Lambagras}{Lambagras - Daněk Pavel, Biolib, Copyright/CC-BY-SA}
\dicEntry[lambakjöt] \dicTerm{lamba··kjöt} \dicIPA{{l}{a}{m}{\textsubring{b}}{a}{c\smash{\textsuperscript{h}}}{\oe}{\textsubring{d}}} \dicPos{n}[2] \dicFlx{(‑s)}[2] \dicDirectTranslationCS{jehněčí (maso)}
\dicEntry[lambalæri] \dicTerm{lamba··læri} \dicIPA{{l}{a}{m}{\textsubring{b}}{a}{l}{a}{i}{r}{\textsci}} \dicPos{n}[2] \dicFlx{(‑s, ‑)}[14] \dicDirectTranslationCS{jehněčí kýta}
\dicEntry[lamdi] \dicTerm{lamdi} \dicIPA{{l}{a}{m}{\textsubring{d}}{\textsci}} \dicPos{v} \dicFlx{ind pf sg 1 pers} \dicLink{lemja}
\dicEntry[lamið] \dicTerm{lamið} \dicIPA{{l}{a}{\textlengthmark}{m}{\textsci}{\texttheta}} \dicPos{v} \dicFlx{supin} \dicLink{lemja}
\dicEntry[lampi] \dicTerm{lamp|i} \dicsymFrequent\  \dicIPA{{l}{a}{\textsubring{m}}{\textsubring{b}}{\textsci}} \dicPos{m}[1] \dicFlx{(‑a, ‑ar)}[8] \textbf{1.} \dicDirectTranslationCS{lampa, svítilna} \dicExampleIS{kveikja á lampa} \dicExampleCS{zapálit lampu}  \textbf{2.} \dicFieldCat{elek.} \dicDirectTranslationCS{výbojka}
\dicEntry[land] \dicTerm{land} \dicsymFrequent\  \dicIPA{{l}{a}{n}{\textsubring{d}}} \dicPos{n}[2] \dicFlx{(‑s, lönd)}[8] \textbf{1.} \dicSynonym{þurrlendi} \dicDirectTranslationCS{souš, pevnina, země} \dicExampleIS{á sjó og landi} \dicExampleCS{na moři a~na pevnině}  \textbf{2.} \dicSynonym{árbakki} \dicDirectTranslationCS{břeh} \dicExampleIS{á öðru landinu} \dicExampleCS{na druhém břehu}  \textbf{3.} \dicSynonym{ríki} \dicDirectTranslationCS{země, stát, vlast} \dicExampleIS{höfuðborg landsins} \dicExampleCS{hlavní město země}  \textbf{4.} \dicSynonym{landareign} \dicDirectTranslationCS{země, půda, pozemek};  \dicIdiom{land}{ \dicPhraseIS{e‑að á langt í land}} \dicLangCat{přen.} \dicDirectTranslationCS{(co) je ještě daleko (cíl ap.)}; { \dicPhraseIS{eiga langt í land}} \dicLangCat{přen.} \dicDirectTranslationCS{mít ještě hodně před sebou (větší část práce ap.)}; { \dicPhraseIS{draga í land}} \dicLangCat{přen.} \dicDirectTranslationCS{ustoupit (ze zásad ap.), rezignovat, upustit (od plánů ap.)}; { \dicPhraseIS{draga e‑n að landi}} \dicDirectTranslationCS{dojíst po (kom)}; { \dicPhraseIS{fara með löndum}} \dicLangCat{přen.} \dicDirectTranslationCS{postupovat\,/\addthin jít opatrně}; { \dicPhraseIS{láta e‑ð lönd og leið}} \dicLangCat{přen.} \dicDirectTranslationCS{nestarat se o~(co), nechat (co) být}; { \dicPhraseIS{leggja land undir fót}} \dicDirectTranslationCS{jít pěšky, pochodovat, kráčet}
\dicEntry[landa] \dicTerm{land|a} \dicIPA{{l}{a}{n}{\textsubring{d}}{a}} \dicPos{v}[1] \dicFlx{(‑aði)}[13] \dicFlx{dat} \dicDirectTranslationCS{vyložit, vylodit} \dicExampleIS{landa aflanum} \dicExampleCS{vyložit úlovek}
\dicEntry[landabréf] \dicTerm{landa··bréf} \dicIPA{{l}{a}{n}{\textsubring{d}}{a}{\textsubring{b}}{r}{j}{\textepsilon}{f}} \dicPos{n}[2] \dicFlx{(‑s, ‑)}[5] \dicSynonym{landakort} \dicDirectTranslationCS{mapa (zeměpisná ap.)}
\dicEntry[landafræði] \dicTerm{landa··fræð|i} \dicIPA{{l}{a}{n}{\textsubring{d}}{a}{f}{r}{a}{i}{ð}{\textsci}} \dicPos{f}[3] \dicFlx{(‑i)}[3] \dicDirectTranslationCS{zeměpis, geografie}
\dicEntry[landakort] \dicTerm{landa··kort} \dicIPA{{l}{a}{n}{\textsubring{d}}{a}{k\smash{\textsuperscript{h}}}{\textopeno}{\textsubring{r}}{\textsubring{d}}} \dicPos{n}[2] \dicFlx{(‑s, ‑)}[5] \dicDirectTranslationCS{mapa}
\dicEntry[landakortabók] \dicTerm{landa·korta··|bók} \dicIPA{{l}{a}{n}{\textsubring{d}}{a}{k\smash{\textsuperscript{h}}}{\textopeno}{\textsubring{r}}{\textsubring{d}}{a}{\textsubring{b}}{ou}{\r{g}}} \dicPos{f}[8] \dicFlx{(‑bókar, ‑bækur)}[5] \dicSynonym{atlas} \dicDirectTranslationCS{atlas}
\dicEntry[landalda] \dicTerm{land··|alda} \dicIPA{{l}{a}{n}{\textsubring{d}}{a}{l}{\textsubring{d}}{a}} \dicPos{f}[1] \dicFlx{(‑öldu, ‑öldur)}[20] \dicDirectTranslationCS{pobřežní vlna}
\dicEntry[landamerki] \dicTerm{landa··merki} \dicIPA{{l}{a}{n}{\textsubring{d}}{a}{m}{\textepsilon}{\textsubring{r}}{\r{\textObardotlessj}}{\textsci}} \dicPos{n}[2] \dicFlx{(‑s, ‑)}[16] \dicSynonym{landamæri} \dicDirectTranslationCS{pomezí, hranice (státní ap.)}
\dicEntry[landamæri] \dicTerm{landa··mæri} \dicIPA{{l}{a}{n}{\textsubring{d}}{a}{m}{a}{i}{r}{\textsci}} \dicPos{n}[2] \dicFlx{pl}[19] \dicDirectTranslationCS{(státní) hranice} \dicExampleIS{vera kominn yfir landamærin} \dicExampleCS{překročit státní hranici}
\dicEntry[landareign] \dicTerm{landar··eign} \dicIPA{{l}{a}{n}{\textsubring{d}}{a}{r}{ei}{\r{g}}{\textsubring{n}}} \dicPos{f}[7] \dicFlx{(‑ar, ‑ir)}[1] \dicDirectTranslationCS{pozemek, půda, pozemkové vlastnictví}
\dicEntry[landbúnaðarframleiðsla] \dicTerm{land·búnaðar··fram·leiðsl|a} \dicIPA{{l}\-{a}\-{n}\-{\textsubring{d}}\-{\textsubring{b}}\-{u}\-{n}\-{a}\-{ð}\-{a}\-{\textsubring{r}}\-{f}\-{r}\-{a}\-{m}\-{l}\-{ei}\-{ð}\-{s}\-{\textsubring{d}}\-{l}\-{a}\-} \dicPos{f}[1] \dicFlx{(‑u)}[5] \dicDirectTranslationCS{zemědělská produkce\,/\addthin výroba}
\dicEntry[landbúnaðarráðherra] \dicTerm{land·búnaðar··ráð·herr|a} \dicIPA{{l}\-{a}\-{n}\-{\textsubring{d}}\-{\textsubring{b}}\-{u}\-{n}\-{a}\-{ð}\-{a}\-{r}\-{au}\-{\texttheta}\-{h}\-{\textepsilon}\-{r}\-{a}\-} \dicPos{m}[1] \dicFlx{(‑a, ‑ar)}[17] \dicDirectTranslationCS{ministr(yně) zemědělství}
\dicEntry[landbúnaðarráðuneyti] \textls[15]{\dicTerm{land·búnaðar··ráðu·neyti} \dicIPA{{l}\-{a}\-{n}\-{\textsubring{d}}\-{\textsubring{b}}\-{u}\-{n}\-{a}\-{ð}\-{a}\-{r}\-{au}\-{ð}\-{\textscy}\-{n}\-{ei}\-{\textsubring{d}}\-{\textsci}}} \dicPos{n}[2] \dicFlx{(‑s, ‑)}[14] \dicDirectTranslationCS{ministerstvo zemědělství}
\dicEntry[landbúnaður] \dicTerm{land··bú·nað|ur} \dicIPA{{l}{a}{n}{\textsubring{d}}{\textsubring{b}}{u}{n}{a}{ð}{\textscy}{\textsubring{r}}} \dicPos{m}[10] \dicFlx{(‑ar)}[9] \dicDirectTranslationCS{zemědělství}
\dicEntry[landeigandi] \dicTerm{land··eig·|andi} \dicIPA{{l}{a}{n}{\textsubring{d}}{ei}{\textbabygamma}{a}{n}{\textsubring{d}}{\textsci}} \dicPos{m}[2] \dicFlx{(‑anda, ‑endur)}[1] \dicDirectTranslationCS{majitel(ka) půdy\,/\addthin pozemku}
\dicEntry[landfastur] \dicTerm{land··|fastur} \dicIPA{{l}{a}{n}{\textsubring{d}}{f}{a}{s}{\textsubring{d}}{\textscy}{\textsubring{r}}} \dicPos{adj}[1] \dicFlx{(f ‑föst)}[11] \dicDirectTranslationCS{(jsoucí) při pobřeží (led ap.)}
\dicEntry[landflótta] \dicTerm{land··flótta} \dicIPA{{l}{a}{n}{\textsubring{d}}{f}{l}{ou}{h}{\textsubring{d}}{a}} \dicPos{adj}[13] \dicFlx{indecl}[1] \dicDirectTranslationCS{(jsoucí) v~exilu, (jsoucí) ve vyhnanství}
\dicEntry[landfræðilegur] \dicTerm{land·fræði··legur} \dicIPA{{l}{a}{n}{\textsubring{d}}{f}{r}{a}{i}{ð}{\textsci}{l}{\textepsilon}{\textbabygamma}{\textscy}{\textsubring{r}}} \dicPos{adj}[1]\dicFlx{}[-8] \dicDirectTranslationCS{geografický}
\dicEntry[landfræðingur] \dicTerm{land·fræð··ing|ur} \dicIPA{{l}{a}{n}{\textsubring{d}}{f}{r}{a}{i}{ð}{i}{\ng}{\r{g}}{\textscy}{\textsubring{r}}} \dicPos{m}[6] \dicFlx{(‑s, ‑ar)}[8] \dicDirectTranslationCS{geograf(ka)}
\dicEntry[landgrunn] \dicTerm{land··grunn} \dicIPA{{l}{a}{n}{\textsubring{d}}{\r{g}}{r}{\textscy}{\textsubring{n}}} \dicPos{n}[2] \dicFlx{(‑s, ‑)}[5] \dicFieldCat{geol.} \dicDirectTranslationCS{kontinentální\,/\addthin pevninský šelf, pevninský práh}
\dicEntry[landgræðsla] \dicTerm{land··græðsl|a} \dicIPA{{l}{a}{n}{\textsubring{d}}{\r{g}}{r}{a}{i}{ð}{s}{\textsubring{d}}{l}{a}} \dicPos{f}[1] \dicFlx{(‑u)}[5] \dicDirectTranslationCS{meliorace}
\dicEntry[landgöngubrú] \dicTerm{land··göngu·|brú} \dicIPA{{l}{a}{n}{\textsubring{d}}{\r{g}}{\oe i}{\ng}{\r{g}}{\textscy}{\textsubring{b}}{r}{u}} \dicPos{f}[9] \dicFlx{(‑brúar, ‑brýr)}[4] \dicDirectTranslationCS{lodní lávka}
\dicEntry[landhelgi] \dicTerm{land··helg|i} \dicsymFrequent\  \dicIPA{{l}{a}{n}{\textsubring{d}}{h}{\textepsilon}{l}{\r{\textObardotlessj}}{\textsci}} \dicPos{f}[3] \dicFlx{(‑i)}[3] \dicDirectTranslationCS{výsostné\,/\addthin teritoriální vody} \dicExampleIS{í landhelgi Íslands} \dicExampleCS{v~islandských výsostných vodách}
\dicEntry[landhelgisgæsla] \dicTerm{land·helgis··gæsl|a} \dicIPA{{l}{a}{n}{\textsubring{d}}{h}{\textepsilon}{l}{\r{\textObardotlessj}}{\textsci}{s}{\r{\textObardotlessj}}{a}{i}{s}{\textsubring{d}}{l}{a}} \dicPos{f}[1] \dicFlx{(‑u)}[5] \dicDirectTranslationCS{pobřežní stráž\,/\addthin hlídka}
\dicEntry[landher] \dicTerm{land··her} \dicIPA{{l}{a}{n}{\textsubring{d}}{h}{\textepsilon}{\textsubring{r}}} \dicPos{m}[9] \dicFlx{(‑s, ‑ir)}[21] \dicDirectTranslationCS{pozemní vojsko}
\dicEntry[landhreinsun] \dicTerm{land··hreins|un} \dicIPA{{l}{a}{n}{\textsubring{d}}{\textsubring{r}}{ei}{n}{s}{\textscy}{\textsubring{n}}} \dicPos{f}[7] \dicFlx{(‑unar)}[9] \dicPhraseIS{það er landhreinsun að e‑m} \dicLangCat{přen.} \dicDirectTranslationCS{no sláva, že už je (kdo) pryč}
\dicEntry[landi] \dicTerm{land|i} \dicIPA{{l}{a}{n}{\textsubring{d}}{\textsci}} \dicPos{m}[1] \dicFlx{(‑a, ‑ar)}[8] \textbf{1.} \dicSynonym{samlandi} \dicDirectTranslationCS{krajan(ka)}  \textbf{2.} \dicSynonym{brugg} \dicDirectTranslationCS{(domácí) pálenka}
\dicEntry[landkrabbi] \dicTerm{land··krabb|i} \dicIPA{{l}{a}{n}{\textsubring{d}}{k\smash{\textsuperscript{h}}}{r}{a}{\textsubring{b}}{\textsci}} \dicPos{m}[1] \dicFlx{(‑a, ‑ar)}[8] \textbf{1.} \dicFieldCat{zool.} \dicDirectTranslationCS{krab} \textit{(l.~{\textLA{Gecarcinus}})}  \dicsymPhoto\   \textbf{2.} \dicDirectTranslationCS{suchozemská krysa} \dicIndirectTranslationCS{(o~člověku, který ví málo o~práci na moři)}
\dicFigure{ds_image_landkrabbi_0_1.jpg}{Landkrabbi}{Landkrabbi - Dozenist, CC BY-SA 3.0}
\dicEntry[landkynning] \dicTerm{land··kynn·ing} \dicIPA{{l}{a}{n}{\textsubring{d}}{c\smash{\textsuperscript{h}}}{\textsci}{n}{i}{\ng}{\r{g}}} \dicPos{f}[4] \dicFlx{(‑ar, ‑ar)}[5] \dicDirectTranslationCS{propagace země}
\dicEntry[landkönnuður] \dicTerm{land··könn·uð|ur} \dicIPA{{l}{a}{n}{\textsubring{d}}{k\smash{\textsuperscript{h}}}{\oe}{n}{\textscy}{ð}{\textscy}{\textsubring{r}}} \dicPos{m}[10] \dicFlx{(‑ar, ‑ir)}[4] \dicDirectTranslationCS{cestovatel(ka), průzkumník, průzkumnice, objevitel(ka) (neznámých končin ap.)}
\dicEntry[landluktur] \dicTerm{land··luktur} \dicIPA{{l}{a}{n}{\textsubring{d}}{l}{\textscy}{x}{\textsubring{d}}{\textscy}{\textsubring{r}}} \dicPos{adj}[1]\dicFlx{}[-10] \dicDirectTranslationCS{vnitrozemský}
\dicEntry[landlægur] \dicTerm{land··lægur} \dicIPA{{l}{a}{n}{\textsubring{d}}{l}{a}{i}{\textbabygamma}{\textscy}{\textsubring{r}}} \dicPos{adj}[1]\dicFlx{}[-1] \textbf{1.} \dicFieldCat{med.} \dicDirectTranslationCS{endemický} \dicExampleIS{landlægur sjúkdómur} \dicExampleCS{endemické onemocnění}  \textbf{2.} \dicSynonym{fastur} \dicDirectTranslationCS{zakořeněný} \dicExampleIS{landlægur siður} \dicExampleCS{zakořeněný zvyk}
\dicEntry[landlæknir] \dicTerm{land··lækn|ir} \dicIPA{{l}{a}{n}{\textsubring{d}}{l}{a}{i}{h}{\r{g}}{n}{\textsci}{\textsubring{r}}} \dicPos{m}[7] \dicFlx{(‑is, ‑ar)}[1] \dicDirectTranslationCS{ředitel(ka) zdravotnictví}
\dicEntry[Landmannalaugar] \dicTerm{Land·manna··laugar} \dicIPA{{l}{a}{n}{\textsubring{d}}{m}{a}{n}{a}{l}{\oe i}{\textbabygamma}{a}{\textsubring{r}}} \dicPos{f}[4] \dicFlx{pl}[2] \dicFieldCat{geog.} \dicDirectTranslationCS{Landmannalaugar} \dicIndirectTranslationCS{(oblast s~horkými prameny a~barevnými skalami na jihu Islandu, oblíbený turistický cíl)}
\dicEntry[landmótunarfræði] \dicTerm{land·mótunar··fræð|i} \dicIPA{{l}{a}{n}{\textsubring{d}}{m}{ou}{\textsubring{d}}{\textscy}{n}{a}{\textsubring{r}}{f}{r}{a}{i}{ð}{\textsci}} \dicPos{f}[3] \dicFlx{(‑i)}[3] \dicDirectTranslationCS{geomorfologie}
\dicEntry[landmælingafræði] \dicTerm{land·mælinga··fræð|i} \dicIPA{{l}{a}{n}{\textsubring{d}}{m}{a}{i}{l}{i}{\ng}{\r{g}}{a}{f}{r}{a}{i}{ð}{\textsci}} \dicPos{f}[3] \dicFlx{(‑i)}[3] \dicDirectTranslationCS{geodézie, zeměměřičství}
\dicEntry[landmælingar] \dicTerm{land··mæl·ingar} \dicIPA{{l}{a}{n}{\textsubring{d}}{m}{a}{i}{l}{i}{\ng}{\r{g}}{a}{\textsubring{r}}} \dicPos{f}[4] \dicFlx{pl}[6] \dicDirectTranslationCS{geodetické měření}
\dicEntry[landnám] \dicTerm{land··nám} \dicIPA{{l}{a}{n}{\textsubring{d}}{n}{au}{\textsubring{m}}} \dicPos{n}[2] \dicFlx{(‑s)}[2] \dicDirectTranslationCS{osídlení, osídlování, kolonizace, kolonizování}

\dicEntry[landnámsmaður] \textls[19]{\dicTerm{land·náms··|maður} \dicIPA{{l}{a}{n}{\textsubring{d}}{n}{au}{m}{s}{m}{a}{ð}{\textscy}{\textsubring{r}}} \dicPos{m}[13] \dicFlx{(‑manns, ‑menn)}[2]} \dicDirectTranslationCS{osadník, osadnice, kolonista, kolonistka}

\dicEntry[landnámsöld] \dicTerm{land·náms··|öld} \dicIPA{{l}{a}{n}{\textsubring{d}}{n}{au}{m}{s}{\oe}{l}{\textsubring{d}}} \dicPos{f}[7] \dicFlx{(‑aldar)}[19] \dicFieldCat{hist.} \dicDirectTranslationCS{doba osídlování}
\dicEntry[landnemi] \dicTerm{land··nem|i} \dicIPA{{l}{a}{n}{\textsubring{d}}{n}{\textepsilon}{m}{\textsci}} \dicPos{m}[1] \dicFlx{(‑a, ‑ar)}[1] \dicDirectTranslationCS{osadník, osadnice, kolonista, kolonistka} \dicExampleIS{íslensku landnemarnir í Vesturheimi} \dicExampleCS{islandští osadníci v~Severní Americe}
\dicEntry[landráð] \dicTerm{land··ráð} \dicIPA{{l}{a}{n}{\textsubring{d}}{r}{au}{\texttheta}} \dicPos{n}[2] \dicFlx{pl}[1] \dicFieldCat{práv.} \dicDirectTranslationCS{velezrada, vlastizrada}
\dicEntry[landrek] \dicTerm{land··rek} \dicIPA{{l}{a}{n}{\textsubring{d}}{r}{\textepsilon}{\r{g}}} \dicPos{n}[2] \dicFlx{(‑s)}[2] \dicFieldCat{geol.} \dicDirectTranslationCS{kontinentální drift}
\dicEntry[landsbanki] \dicTerm{lands··bank|i} \dicIPA{{l}{a}{n}{\textsubring{d}}{s}{\textsubring{b}}{au}{\r{\textltailn}}{\r{\textObardotlessj}}{\textsci}} \dicPos{m}[1] \dicFlx{(‑a, ‑ar)}[8] \dicDirectTranslationCS{národní banka}
\dicEntry[landsbyggð] \dicTerm{lands··byggð} \dicIPA{{l}{a}{n}{\textsubring{d}}{s}{\textsubring{b}}{\textsci}{\textbabygamma}{\texttheta}} \dicPos{f}[7] \dicFlx{(‑ar)}[3] \dicDirectTranslationCS{venkov, venkovský kraj} \dicExampleIS{úti á landsbyggðinni} \dicExampleCS{na venkově}
\dicEntry[landsfrægur] \dicTerm{lands··frægur} \dicIPA{{l}{a}{n}{\textsubring{d}}{s}{f}{r}{a}{i}{\textbabygamma}{\textscy}{\textsubring{r}}} \dicPos{adj}[1]\dicFlx{}[-1] \dicDirectTranslationCS{známý\,/\addthin proslavený po celé zemi}
\dicEntry[landshluti] \dicTerm{lands··hlut|i} \dicIPA{{l}{a}{n}{\textsubring{d}}{s}{\textsubring{l}}{\textscy}{\textsubring{d}}{\textsci}} \dicPos{m}[1] \dicFlx{(‑a, ‑ar)}[1] \dicDirectTranslationCS{část země, region}
\dicEntry[landsími] \dicTerm{land··sím|i} \dicIPA{{l}{a}{n}{\textsubring{d}}{s}{i}{m}{\textsci}} \dicPos{m}[1] \dicFlx{(‑a, ‑ar)}[1] \dicFieldCat{hist.} \dicDirectTranslationCS{národní telefonní společnost}
\dicEntry[landskjálfti] \dicTerm{land··skjálft|i} \dicIPA{{l}{a}{n}{\textsubring{d}}{s}{\r{\textObardotlessj}}{au}{\textsubring{l}}{\textsubring{d}}{\textsci}} \dicPos{m}[1] \dicFlx{(‑a, ‑ar)}[1] \dicFieldCat{geol.} \dicLangCat{zast.} \dicSynonym{jarðskjálfti} \dicDirectTranslationCS{zemětřesení}
\dicEntry[landslag] \dicTerm{lands··lag} \dicsymFrequent\  \dicIPA{{l}{a}{n}{\textsubring{d}}{s}{l}{a}{x}} \dicPos{n}[2] \dicFlx{(‑s)}[2] \dicDirectTranslationCS{krajina, terén} \dicExampleIS{skoða landslag í firðinum} \dicExampleCS{zkoumat krajinu ve fjordu}
\dicEntry[landslagsarkitekt] \dicTerm{lands·lags··arkitekt} \dicIPA{{l}{a}{n}{\textsubring{d}}{s}{l}{a}{x}{s}{a}{\textsubring{r}}{\r{\textObardotlessj}}{\textsci}{\textsubring{d}}{\textepsilon}{x}{\textsubring{d}}} \dicPos{m}[4] \dicFlx{(‑s, ‑ar)}[10] \dicDirectTranslationCS{krajinný architekt, krajinná architektka}
\dicEntry[landslagsmynd] \dicTerm{lands·lags··mynd} \dicIPA{{l}{a}{n}{\textsubring{d}}{s}{l}{a}{x}{s}{m}{\textsci}{n}{\textsubring{d}}} \dicPos{f}[7] \dicFlx{(‑ar, ‑ir)}[1] \dicDirectTranslationCS{krajina (obraz), krajinomalba}
\dicEntry[landsleikur] \dicTerm{lands··leik|ur} \dicIPA{{l}{a}{n}{\textsubring{d}}{s}{l}{ei}{\r{g}}{\textscy}{\textsubring{r}}} \dicPos{m}[9] \dicFlx{(‑s, ‑ir)}[15] \dicDirectTranslationCS{mezistátní\,/\addthin mezinárodní zápas\,/\addthin utkání (ve fotbale ap.)}
\dicEntry[landslið] \dicTerm{lands··lið} \dicIPA{{l}{a}{n}{\textsubring{d}}{s}{l}{\textsci}{\texttheta}} \dicPos{n}[2] \dicFlx{(‑s, ‑)}[5] \dicDirectTranslationCS{národní tým\,/\addthin mužstvo}
\dicEntry[landsliðsþjálfari] \dicTerm{lands·liðs··þjálf·ar|i} \dicIPA{{l}\-{a}\-{n}\-{\textsubring{d}}\-{s}\-{l}\-{\textsci}\-{ð}\-{s}\-{\texttheta}\-{j}\-{au}\-{l}\-{v}\-{a}\-{r}\-{\textsci}\-} \dicPos{m}[1] \dicFlx{(‑a, ‑ar)}[13] \dicDirectTranslationCS{trenér(ka) národního týmu}
\dicEntry[landslýður] \dicTerm{lands··lýð|ur} \dicIPA{{l}{a}{n}{\textsubring{d}}{s}{l}{i}{ð}{\textscy}{\textsubring{r}}} \dicPos{m}[9] \dicFlx{(‑s, ‑ir)}[8] \dicDirectTranslationCS{obyvatelstvo}
\dicEntry[landsmaður] \dicTerm{lands··|maður} \dicsymFrequent\  \dicIPA{{l}{a}{n}{\textsubring{d}}{s}{m}{a}{ð}{\textscy}{\textsubring{r}}} \dicPos{m}[13] \dicFlx{(‑manns, ‑menn)}[2] \dicDirectTranslationCS{obyvatel(ka)} \dicExampleIS{allir landsmenn undir fimmtugu} \dicExampleCS{všichni obyvatelé pod padesát let}
\dicEntry[landsmót] \dicTerm{lands··mót} \dicIPA{{l}{a}{n}{\textsubring{d}}{s}{m}{ou}{\textsubring{d}}} \dicPos{n}[2] \dicFlx{(‑s, ‑)}[5] \textbf{1.} \dicDirectTranslationCS{celorepublikový\,/\addthin národní sjezd\,/\addthin setkání}  \textbf{2.} \dicDirectTranslationCS{celorepublikový\,/\addthin národní turnaj (v~házené ap.)}
\dicEntry[landspítali] \dicTerm{land··spítal|i} \dicIPA{{l}{a}{n}{\textsubring{d}}{s}{\textsubring{b}}{i}{t\smash{\textsuperscript{h}}}{a}{l}{\textsci}} \dicPos{m}[1] \dicFlx{(‑a, ‑ar)}[8] \dicDirectTranslationCS{národní univerzitní nemocnice}
\begin{xtolerant}{}{1pt}
\dicEntry[landssamband] \dicTerm{lands··sam·|band} \dicIPA{{l}{a}{n}{\textsubring{d}}{s}{a}{m}{\textsubring{b}}{a}{n}{\textsubring{d}}} \dicPos{n}[2] \dicFlx{(‑bands, ‑bönd)}[8] \dicDirectTranslationCS{národní\,/\addthin celorepublikový svaz}
\end{xtolerant}
\dicEntry[landstjóri] \dicTerm{land··stjór|i} \dicIPA{{l}{a}{n}{\textsubring{d}}{s}{\textsubring{d}}{j}{ou}{r}{\textsci}} \dicPos{m}[1] \dicFlx{(‑a, ‑ar)}[1] \dicDirectTranslationCS{místodržitel(ka), guvernér(ka), zemský správce, zemská správkyně}
\dicEntry[landsvala] \dicTerm{land··|svala} \dicIPA{{l}{a}{n}{\textsubring{d}}{s}{v}{a}{l}{a}} \dicPos{f}[1] \dicFlx{(‑svölu, ‑svölur)}[20] \dicFieldCat{zool.} \dicDirectTranslationCS{vlaštovka, vlaštovka obecná} \textit{(l.~{\textLA{Hirundo rustica}})}  \dicsymPhoto\ 
\dicFigure{20731.jpg}{Landsvala}{Landsvala - Menke, Dave, Biolib, PD}
\dicEntry[Landsvirkjun] \dicTerm{Lands··virkj|un} \dicIPA{{l}{a}{n}{\textsubring{d}}{s}{v}{\textsci}{\textsubring{r}}{\r{\textObardotlessj}}{\textscy}{\textsubring{n}}} \dicPos{f}[7] \dicFlx{(‑unar)}[10] \dicIndirectTranslationCS{islandská společnost provozující elektrárny}
\dicEntry[landsvæði] \dicTerm{land··svæði} \dicIPA{{l}{a}{n}{\textsubring{d}}{s}{v}{a}{i}{ð}{\textsci}} \dicPos{n}[2] \dicFlx{(‑s, ‑)}[14] \dicDirectTranslationCS{území, oblast, teritorium}
\dicEntry[langa] \dicTerm{langa\smash{\textsuperscript{1}}} \dicIPA{{l}{au}{\ng}{\r{g}}{a}} \dicPos{f}[1] \dicFlx{(löngu, löngur)}[20] \dicFieldCat{zool.} \dicDirectTranslationCS{mník mořský} \textit{(l.~{\textLA{Molva molva}})}
\dicEntry[langa] \dicTerm{lang|a\smash{\textsuperscript{2}}} \dicsymFrequent\  \dicIPA{{l}{au}{\ng}{\r{g}}{a}} \dicPos{v}[1] \dicFlx{(‑aði)}[79] \dicFlx{impers} \dicPhraseIS{e‑n langar} \dicDirectTranslationCS{(kdo) touží, (kdo) chce, (kdo) má chuť} \dicExampleIS{Mig langar að fara.} \dicExampleCS{Chce se mi jít.};  \dicPhraseIS{e‑n langar í e‑ð} \dicDirectTranslationCS{(kdo) touží po (čem), (kdo) chce (co), (kdo) má chuť na (co)} \dicExampleIS{Mig langar í sælgæti.} \dicExampleCS{Mám chuť na něco sladkého.};  \dicPhraseIS{e‑n langar til e‑s} \dicDirectTranslationCS{(kdo) touží po (čem), (kdo) chce (co)} \dicExampleIS{Mig langar til að lesa bókina.} \dicExampleCS{Mám chuť si číst knížku.};  \dicPhraseIS{e‑n langar eftir e‑u} \dicDirectTranslationCS{(kdo) touží po (čem), (kdo) prahne po (čem)} \dicExampleIS{Hana langaði eftir friði.} \dicExampleCS{Toužila po klidu.}
\dicEntry[langafasta] \dicTerm{langa··fasta} \dicIPA{{l}{au}{\ng}{\r{g}}{a}{f}{a}{s}{\textsubring{d}}{a}} \dicPos{f}[1] \dicFlx{(lönguföstu, lönguföstur)}[21] \dicFieldCat{náb.} \dicDirectTranslationCS{půst} \dicIndirectTranslationCS{(trvající 7 týdnů od Popeleční středy do Velikonoc)}
\dicEntry[langafi] \dicTerm{lang··af|i} \dicIPA{{l}{au}{\ng}{\r{g}}{a}{v}{\textsci}} \dicPos{m}[1] \dicFlx{(‑a, ‑ar)}[8] \dicDirectTranslationCS{pradědeček} \dicIndirectTranslationCS{(otec dědečka nebo babičky)}
\dicEntry[langamma] \dicTerm{lang··|amma} \dicIPA{{l}{au}{\ng}{\r{g}}{a}{m}{a}} \dicPos{f}[1] \dicFlx{(‑ömmu, ‑ömmur)}[8] \dicDirectTranslationCS{prababička} \dicIndirectTranslationCS{(matka dědečka nebo babičky)}
\dicEntry[langanir] \dicTerm{langanir} \dicIPA{{l}{au}{\ng}{\r{g}}{a}{n}{\textsci}{\textsubring{r}}} \dicPos{f} \dicFlx{pl nom} \dicLink{löngun}
\dicEntry[langatöng] \dicTerm{langa··töng}\dicTerm{, löngutöng} \dicIPA{{l}\-{au}\-{\ng}\-{\r{g}}\-{a}\-{t\smash{\textsuperscript{h}}}\-{\oe i}\-{\ng}\-{\r{g}}\-} \dicPos{f}[8] \dicFlx{(löngutangar, löngutengur\,/\addthin löngutangir)}[10] \dicDirectTranslationCS{prostředník (prst), prostředníček}
\dicEntry[langbestur] \dicTerm{lang··bestur} \dicIPA{{l}{au}{\ng}{\r{g}}{\textsubring{b}}{\textepsilon}{s}{\textsubring{d}}{\textscy}{\textsubring{r}}} \dicPos{adj}[11]\dicFlx{}[-2] \dicFlx{m sg nom sup} \dicDirectTranslationCS{úplně\,/\addthin zdaleka nejlepší}
\dicEntry[langbylgja] \dicTerm{lang··bylgj|a} \dicIPA{{l}{au}{\ng}{\r{g}}{\textsubring{b}}{\textsci}{l}{\r{\textObardotlessj}}{a}} \dicPos{f}[1] \dicFlx{(‑u, ‑ur)}[25] \dicFieldCat{fyz.} \dicDirectTranslationCS{dlouhá vlna} \dicAntonym{stuttbylgja}
\dicEntry[langdreginn] \dicTerm{lang··dreginn} \dicIPA{{l}{au}{\ng}{\r{g}}{\textsubring{d}}{r}{ei}{\textsci}{\textsubring{n}}} \dicPos{adj}[6]\dicFlx{}[-2] \textbf{1.} \dicDirectTranslationCS{zdlouhavý, rozvláč\-ný} \dicExampleIS{langdreginn ræða} \dicExampleCS{zdlouhavý proslov}  \textbf{2.} \dicDirectTranslationCS{dlouhotrvající} \dicExampleIS{langdreginn sjúkdómur} \dicExampleCS{dlouhotrvající onemocnění}
\dicEntry[langdvöl] \dicTerm{lang··|dvöl} \dicIPA{{l}{au}{\ng}{\r{g}}{\textsubring{d}}{v}{\oe}{\textsubring{l}}} \dicPos{f}[7] \dicFlx{(‑dvalar, ‑dvalir)}[16] \dicDirectTranslationCS{dlouhotrvající pobyt};  \dicPhraseIS{dvelja\,/\addthin vera e‑s staðar langdvölum} \dicDirectTranslationCS{pobývat (kde) dlouho}
\dicEntry[langeygur] \dicTerm{lang··eygur} \dicIPA{{l}{au}{\ng}{\r{g}}{ei}{\textbabygamma}{\textscy}{\textsubring{r}}} \dicPos{adj}[1]\dicFlx{}[-1] \dicPhraseIS{vera langeygur eftir e‑u} \dicDirectTranslationCS{být nedočkavý (čeho), nemoct se dočkat (čeho)}
\dicEntry[langframi] \dicTerm{lang··fram|i} \dicIPA{{l}{au}{\ng}{\r{g}}{f}{r}{a}{m}{\textsci}} \dicPos{m}[1] \dicFlx{(‑a)}[3] \dicPhraseIS{til langframa} \dicFlx{adv} \dicDirectTranslationCS{dlouhodobě}
\dicEntry[langhlaup] \dicTerm{lang··hlaup} \dicIPA{{l}{au}{\ng}{\r{g}}{\textsubring{l}}{\oe i}{\textsubring{b}}} \dicPos{n}[2] \dicFlx{(‑s, ‑)}[5] \dicFieldCat{sport.} \dicDirectTranslationCS{dálkový běh}
\dicEntry[langlífi] \dicTerm{lang··lífi} \dicIPA{{l}{au}{\ng}{\r{g}}{l}{i}{v}{\textsci}} \dicPos{n}[2] \dicFlx{(‑s)}[20] \dicDirectTranslationCS{dlouhověkost}
\dicEntry[langlífur] \dicTerm{lang··lífur} \dicIPA{{l}{au}{\ng}{\r{g}}{l}{i}{v}{\textscy}{\textsubring{r}}} \dicPos{adj}[1]\dicFlx{}[-1] \dicDirectTranslationCS{dlouhověký}
\dicEntry[langloka] \dicTerm{lang··lok|a} \dicIPA{{l}{au}{\ng}{\r{g}}{l}{\textopeno}{\r{g}}{a}} \dicPos{f}[1] \dicFlx{(‑u, ‑ur)}[19] \textbf{1.} \dicSynonym{romsa\smash{\textsuperscript{1}}} \dicDirectTranslationCS{litanie, nudná a~zdlouhavá řeč, nudný a~zdlouhavý text}  \textbf{2.} \dicSynonym{samloka} \dicDirectTranslationCS{bageta}
\dicEntry[langorður] \dicTerm{lang··orður} \dicIPA{{l}{au}{\ng}{\r{g}}{\textopeno}{r}{ð}{\textscy}{\textsubring{r}}} \dicPos{adj}[2]\dicFlx{}[-1] \dicDirectTranslationCS{rozvláčný, upovídaný}
\dicEntry[langreyður] \dicTerm{lang··reyð|ur} \dicIPA{{l}{au}{\ng}{\r{g}}{r}{ei}{ð}{\textscy}{\textsubring{r}}} \dicPos{f}[5] \dicFlx{(‑ar, ‑ar)}[1] \dicFieldCat{zool.} \dicDirectTranslationCS{plejtvák myšok} \textit{(l.~{\textLA{Balaenoptera physalus}})}
\dicEntry[langrækinn] \dicTerm{lang··rækinn} \dicIPA{{l}{au}{\ng}{\r{g}}{r}{a}{i}{\r{\textObardotlessj}}{\textsci}{\textsubring{n}}} \dicPos{adj}[6]\dicFlx{}[-2] \dicDirectTranslationCS{neodpouštějící, nesmiřitelný}
\dicEntry[langskip] \dicTerm{lang··skip} \dicIPA{{l}{au}{\ng}{\r{g}}{s}{\r{\textObardotlessj}}{\textsci}{\textsubring{b}}} \dicPos{n}[2] \dicFlx{(‑s, ‑)}[5] \dicFieldCat{nám.} \dicDirectTranslationCS{vikinská válečná loď}
\dicEntry[langsnið] \dicTerm{lang··snið} \dicIPA{{l}{au}{\ng}{\r{g}}{s}{\textsubring{d}}{n}{\textsci}{\texttheta}} \dicPos{n}[2] \dicFlx{(‑s, ‑)}[5] \dicFieldCat{poč.} \dicDirectTranslationCS{formát na délku}
\dicEntry[langsóttur] \dicTerm{lang··sóttur} \dicIPA{{l}{au}{\ng}{\r{g}}{s}{ou}{h}{\textsubring{d}}{\textscy}{\textsubring{r}}} \dicPos{adj}[1]\dicFlx{}[-10] \dicDirectTranslationCS{nevěrohodný, přitažený za vlasy} \dicExampleIS{langsótt sönnun} \dicExampleCS{nevěrohodný důkaz}
\dicEntry[langspil] \dicTerm{lang··spil} \dicIPA{{l}{au}{\ng}{\r{g}}{s}{\textsubring{b}}{\textsci}{\textsubring{l}}} \dicPos{n}[2] \dicFlx{(‑s, ‑)}[5] \dicFieldCat{hud.} \dicIndirectTranslationCS{islandský strunný hudební nástroj}
\dicEntry[langstökk] \dicTerm{lang··stökk} \dicIPA{{l}{au}{\ng}{\r{g}}{s}{\textsubring{d}}{\oe}{h}{\r{g}}} \dicPos{n}[2] \dicFlx{(‑s, ‑)}[5] \dicFieldCat{sport.} \dicDirectTranslationCS{skok do dálky, skok daleký}
\dicEntry[langt] \dicTerm{langt} \dicsymFrequent\  \dicIPA{{l}{au}{\r{\ng}}{\textsubring{d}}} \dicPos{adv} \dicFlx{(comp lengra, sup lengst)} \dicDirectTranslationCS{daleko} \dicExampleIS{Þetta liggur langt inni í fortíðinni.} \dicExampleCS{Má to původ v~daleké minulosti.};  \dicPhraseIS{langt í burtu} \dicFlx{adv} \dicDirectTranslationCS{daleko};  \dicPhraseIS{langt síðan} \dicFlx{adv} \dicDirectTranslationCS{dávno}
\dicEntry[langtíma] \dicTerm{lang··tíma-} \dicIPA{{l}{au}{\ng}{\r{g}}{t\smash{\textsuperscript{h}}}{i}{m}{a}} \dicPos{predp} \dicDirectTranslationCS{dlouhodobý}
\dicEntry[langtímalán] \dicTerm{lang·tíma··lán} \dicIPA{{l}{au}{\ng}{\r{g}}{t\smash{\textsuperscript{h}}}{i}{m}{a}{l}{au}{\textsubring{n}}} \dicPos{n}[2] \dicFlx{(‑s, ‑)}[5] \dicFieldCat{ekon.} \dicDirectTranslationCS{dlouhodobá půjčka}
\dicEntry[langtímasamningur] \dicTerm{lang·tíma··samn·ing|ur} \dicIPA{{l}\-{au}\-{\ng}\-{\r{g}}\-{t\smash{\textsuperscript{h}}}\-{i}\-{m}\-{a}\-{s}\-{a}\-{m}\-{n}\-{i}\-{\ng}\-{\r{g}}\-{\textscy}\-{\textsubring{r}}\-} \dicPos{m}[6] \dicFlx{(‑s, ‑ar)}[8] \dicDirectTranslationCS{dlouhodobá smlouva}
\dicEntry[langur] \dicTerm{lang|ur\smash{\textsuperscript{1}}} \dicIPA{{l}{au}{\ng}{\r{g}}{\textscy}{\textsubring{r}}} \dicPos{m}[6] \dicFlx{(‑s, ‑ar)}[16] \dicPhraseIS{draga e‑ð á langinn} \dicDirectTranslationCS{odložit (co), odkládat (co)}
\dicEntry[langur] \dicTerm{langur\smash{\textsuperscript{2}}} \dicsymFrequent\  \dicIPA{{l}{au}{\ng}{\r{g}}{\textscy}{\textsubring{r}}} \dicPos{adj}[10] \dicFlx{(f löng, comp lengri, sup lengstur)}[2] \mbox{\textbf{1.} \dicDirectTranslationCS{dlouhý}} \dicIndirectTranslationCS{(o~vzdálenosti)} \dicExampleIS{löng leið} \dicExampleCS{dlouhá cesta} \dicAntonym{stuttur}  \textbf{2.} \dicDirectTranslationCS{vysoký} \dicExampleIS{langur maður} \dicExampleCS{vysoký člověk}  \textbf{3.} \dicDirectTranslationCS{dlouhý} \dicIndirectTranslationCS{(o~čase)} \dicExampleIS{langur tími} \dicExampleCS{dlouhý čas};  \dicPhraseIS{fyrir löngu} \dicFlx{adv} \dicDirectTranslationCS{dávno, před dlouhou dobou};  \dicPhraseIS{löngum stundum} \dicFlx{adv} \dicDirectTranslationCS{dlouho, dlouhé hodiny};  \dicPhraseIS{það er langt til e‑s} \dicDirectTranslationCS{je daleko do (čeho) (Vánoc ap.)}
\dicEntry[langvarandi] \dicTerm{lang··var·andi} \dicsymFrequent\  \dicIPA{{l}{au}{\ng}{\r{g}}{v}{a}{r}{a}{n}{\textsubring{d}}{\textsci}} \dicPos{adj}[13] \dicFlx{indecl}[1] \dicDirectTranslationCS{dlouhotrvající, prodlužovaný} \dicExampleIS{langvarandi erfiðleikar} \dicExampleCS{dlouhotrvající potíže}
\dicEntry[langvinnur] \textls[15]{\dicTerm{lang··vinnur} \dicIPA{{l}{au}{\ng}{\r{g}}{v}{\textsci}{n}{\textscy}{\textsubring{r}}} \dicPos{adj}[1]\dicFlx{}[-1] \dicDirectTranslationCS{přetrvávající, táhlý, dlouhotrvající} \dicExampleIS{langvinnur sjúkdómur} \dicExampleCS{přetrvávající onemocnění}}
\dicEntry[langvía] \dicTerm{lang··ví|a} \dicIPA{{l}{au}{\ng}{\r{g}}{v}{i}{j}{a}} \dicPos{f}[1] \dicFlx{(‑u, ‑ur)}[7] \dicFieldCat{zool.} \dicDirectTranslationCS{alkoun úzkozobý} \textit{(l.~{\textLA{Uria aalge}})}  \dicsymPhoto\ 
\dicFigure{20877.jpg}{Langvía}{Langvía - Sowls, Art, Biolib, PD}
\dicEntry[langþráður] \dicTerm{lang··þráður} \dicIPA{{l}{au}{\ng}{\r{g}}{\texttheta}{r}{au}{ð}{\textscy}{\textsubring{r}}} \dicPos{adj}[2]\dicFlx{}[-12] \dicDirectTranslationCS{vytoužený, (dlouho) očekávaný}
\dicEntry[langþreyttur] \dicTerm{lang··þreyttur} \dicIPA{{l}{au}{\ng}{\r{g}}{\texttheta}{r}{ei}{h}{\textsubring{d}}{\textscy}{\textsubring{r}}} \dicPos{adj}[1]\dicFlx{}[-13] \dicDirectTranslationCS{znavený, (velmi) unavený}
\dicEntry[lantan] \dicTerm{lantan} \dicIPA{{l}{a}{\textsubring{n}}{\textsubring{d}}{a}{\textsubring{n}}} \dicPos{n}[2] \dicFlx{(‑s)}[2] \dicFieldCat{chem.} \dicDirectTranslationCS{lanthan} \textit{(l.~{\textLA{La, Lanthanum}})}
\dicEntry[Laos] \dicTerm{Laos} \dicIPA{{l}{a}{\textlengthmark}{\textopeno}{s}} \dicPos{n}[4] \dicFlx{indecl}[2] \dicFieldCat{geog.} \dicDirectTranslationCS{Laos}
\dicEntry[Laosi] \dicTerm{Laos|i} \dicIPA{{l}{a}{\textlengthmark}{\textopeno}{s}{\textsci}} \dicPos{m}[1] \dicFlx{(‑a, ‑ar)}[1] \dicDirectTranslationCS{Laosan(ka)}
\dicEntry[laoska] \dicTerm{laosk|a} \dicIPA{{l}{a}{\textlengthmark}{\textopeno}{s}{\r{g}}{a}} \dicPos{f}[1] \dicFlx{(‑u)}[5] \dicDirectTranslationCS{laoština}
\dicEntry[laoskur] \dicTerm{laoskur} \dicIPA{{l}{a}{\textlengthmark}{\textopeno}{s}{\r{g}}{\textscy}{\textsubring{r}}} \dicPos{adj}[1]\dicFlx{}[-6] \dicDirectTranslationCS{laoský}
\dicEntry[lapið] \dicTerm{lapið} \dicIPA{{l}{a}{\textlengthmark}{\textsubring{b}}{\textsci}{\texttheta}} \dicPos{v} \dicFlx{supin} \dicLink{lepja}
\dicEntry[lappa] \dicTerm{lapp|a} \dicIPA{{l}{a}{h}{\textsubring{b}}{a}} \dicPos{v}[1] \dicFlx{(‑aði)}[13] \dicDirectTranslationCS{ochodit, prochodit (boty ap.)} \dicExampleIS{lappa skóna} \dicExampleCS{prochodit boty};  \dicPhraseIS{e‑að lappar e‑n} \dicDirectTranslationCS{(co) je (komu) volné (boty ap.)} \dicExampleIS{Skórinn lappar mig.} \dicExampleCS{Boty mi jsou volné.};  \dicIdiom{lappa}[upp á]{ \dicPhraseIS{lappa upp á e‑ð}} \dicDirectTranslationCS{osvěžit (co), vyspravit (co)}
\dicEntry[lappajaðrakan] \dicTerm{lappa··jað·rakan} \dicIPA{{l}{a}{h}{\textsubring{b}}{a}{j}{a}{ð}{r}{a}{\r{g}}{a}{\textsubring{n}}} \dicPos{m}[4] \dicFlx{(‑s, ‑ar)}[12] \dicFieldCat{zool.} \dicDirectTranslationCS{břehouš rudý} \textit{(l.~{\textLA{Limosa lapponica}})}  \dicsymPhoto\ 
\dicFigure{ds_image_lappajadrakan_0_1.jpg}{Lappajaðrakan}{Lappajaðrakan - Andreas Trepte, CC BY-SA 2.5}
\dicEntry[lappalaus] \dicTerm{lappa··laus} \dicIPA{{l}{a}{h}{\textsubring{b}}{a}{l}{\oe i}{s}} \dicPos{adj}[5]\dicFlx{}[-1] \dicPhraseIS{vera orðinn lappalaus} \dicLangCat{přen.} \dicDirectTranslationCS{být k~smrti unavený}
\dicEntry[lappar] \dicTerm{lappar} \dicIPA{{l}{a}{h}{\textsubring{b}}{a}{\textsubring{r}}} \dicPos{f} \dicFlx{sg gen} \dicLink{löpp}
\dicEntry[Lappi] \dicTerm{Lapp|i} \dicIPA{{l}{a}{h}{\textsubring{b}}{\textsci}} \dicPos{m}[1] \dicFlx{(‑a, ‑ar)}[8] \dicDirectTranslationCS{Laponec, Laponka}
\dicEntry[lappir] \dicTerm{lappir} \dicIPA{{l}{a}{h}{\textsubring{b}}{\textsci}{\textsubring{r}}} \dicPos{f} \dicFlx{pl nom} \dicLink{löpp}
\dicEntry[Lappland] \dicTerm{Lapp··land} \dicIPA{{l}{a}{h}{\textsubring{b}}{l}{a}{n}{\textsubring{d}}} \dicPos{n}[2] \dicFlx{(‑s)}[4] \dicFieldCat{geog.} \dicSynonym{Samaland} \dicDirectTranslationCS{Laponsko}
\dicEntry[lappneska] \dicTerm{lapp··nesk|a} \dicIPA{{l}{a}{h}{\textsubring{b}}{n}{\textepsilon}{s}{\r{g}}{a}} \dicPos{f}[1] \dicFlx{(‑u)}[5] \dicDirectTranslationCS{laponština}
\dicEntry[lappneskur] \dicTerm{lapp··neskur} \dicIPA{{l}{a}{h}{\textsubring{b}}{n}{\textepsilon}{s}{\r{g}}{\textscy}{\textsubring{r}}} \dicPos{adj}[1]\dicFlx{}[-1] \dicDirectTranslationCS{laponský}
\dicEntry[lapti] \dicTerm{lapti} \dicIPA{{l}{a}{f}{\textsubring{d}}{\textsci}} \dicPos{v} \dicFlx{ind pf sg 1 pers} \dicLink{lepja}
\dicEntry[las] \dicTerm{las} \dicIPA{{l}{a}{\textlengthmark}{s}} \dicPos{v} \dicFlx{ind pf sg 1 pers} \dicLink{lesa}
\dicEntry[lasanja] \dicTerm{lasanja} \dicIPA{{l}{a}{\textlengthmark}{s}{a}{n}{j}{a}} \dicPos{n}[4] \dicFlx{indecl}[1] \dicFieldCat{kulin.} \dicDirectTranslationCS{lasagne}
\dicEntry[lasburða] \dicTerm{las··burða} \dicIPA{{l}{a}{\textlengthmark}{s}{\textsubring{b}}{\textscy}{r}{ð}{a}} \dicPos{adj}[13] \dicFlx{indecl}[1] \dicDirectTranslationCS{churavý, neduživý}
\dicEntry[lasinn] \dicTerm{lasinn} \dicIPA{{l}{a}{\textlengthmark}{s}{\textsci}{\textsubring{n}}} \dicPos{adj}[6]\dicFlx{}[-3] \textbf{1.} \dicSynonym{veikur} \dicDirectTranslationCS{nemocný, indisponovaný}  \textbf{2.} \dicSynonym{lélegur} \dicDirectTranslationCS{opotřebovaný}
\dicEntry[laska] \dicTerm{lask|a} \dicIPA{{l}{a}{s}{\r{g}}{a}} \dicPos{v}[1] \dicFlx{(‑aði)}[13] \dicFlx{acc} \dicSynonym{skadda} \dicDirectTranslationCS{poškodit, způsobit škodu}
\dicEntry[lasleiki] \dicTerm{las··leik|i} \dicIPA{{l}{a}{\textlengthmark}{s}{l}{ei}{\r{\textObardotlessj}}{\textsci}} \dicPos{m}[1] \dicFlx{(‑a)}[3] \dicDirectTranslationCS{indispozice, churavost}
\dicEntry[last] \dicTerm{last} \dicIPA{{l}{a}{s}{\textsubring{d}}} \dicPos{n}[2] \dicFlx{(‑s)}[2] \dicSynonym*{baktal} \dicDirectTranslationCS{pomluva, osočení};  \dicPhraseIS{leggja e‑m e‑ð til lasts} \dicDirectTranslationCS{dávat (co komu) za vinu}
\dicEntry[lasta] \dicTerm{last|a} \dicIPA{{l}{a}{s}{\textsubring{d}}{a}} \dicPos{v}[1] \dicFlx{(‑aði)}[13] \dicFlx{acc} \dicSynonym{gagnrýna} \dicDirectTranslationCS{očernit, očerňovat, (po)hanit, pohanět} \dicExampleIS{lasta þýðingu á bók} \dicExampleCS{očerňovat překlad knížky}
\dicEntry[lastar] \dicTerm{lastar} \dicIPA{{l}{a}{s}{\textsubring{d}}{a}{\textsubring{r}}} \dicPos{m} \dicFlx{sg gen} \dicLink{löstur}
\dicEntry[lastyrði] \dicTerm{last··yrði} \dicIPA{{l}{a}{s}{\textsubring{d}}{\textsci}{r}{ð}{\textsci}} \dicPos{n}[2] \dicFlx{(‑s, ‑)}[14] \dicSynonym{last} \dicDirectTranslationCS{hana, hanobení}
\dicEntry[latína] \dicTerm{latín|a} \dicIPA{{l}{a}{\textlengthmark}{t\smash{\textsuperscript{h}}}{i}{n}{a}} \dicPos{f}[1] \dicFlx{(‑u)}[5] \dicDirectTranslationCS{latina}
\dicEntry[latínusegl] \dicTerm{latínu··segl} \dicIPA{{l}{a}{\textlengthmark}{t\smash{\textsuperscript{h}}}{i}{n}{\textscy}{s}{\textepsilon}{\r{g}}{\textsubring{l}}} \dicPos{n}[2] \dicFlx{(‑s, ‑)}[5] \dicFieldCat{nám.} \dicDirectTranslationCS{latinská plachta}
\dicEntry[latneskur] \dicTerm{latn··eskur} \dicIPA{{l}{a}{h}{\textsubring{d}}{n}{\textepsilon}{s}{\r{g}}{\textscy}{\textsubring{r}}} \dicPos{adj}[1]\dicFlx{}[-1] \dicDirectTranslationCS{latinský}
\dicEntry[latt] \dicTerm{latt} \dicIPA{{l}{a}{h}{\textsubring{d}}} \dicPos{v} \dicFlx{supin} \dicLink{letja}
\dicEntry[latti] \dicTerm{latti} \dicIPA{{l}{a}{h}{\textsubring{d}}{\textsci}} \dicPos{v} \dicFlx{ind pf sg 1 pers} \dicLink{letja}
\dicEntry[latur] \dicTerm{latur} \dicIPA{{l}{a}{\textlengthmark}{\textsubring{d}}{\textscy}{\textsubring{r}}} \dicPos{adj}[1] \dicFlx{(f löt)}[2] \dicDirectTranslationCS{líný, lenivý}
\dicEntry[lauf] \dicTerm{lauf} \dicsymFrequent\  \dicIPA{{l}{\oe i}{\textlengthmark}{f}} \dicPos{n}[2] \dicFlx{(‑s, ‑)}[5] \textbf{1.} \dicSynonym{laufblað} \dicDirectTranslationCS{list (stromu ap.)} \dicExampleIS{gulnuð lauf að hausti} \dicExampleCS{zežloutlé listy na podzim}  \textbf{2.} \dicSynonym*{litur í spilum} \dicDirectTranslationCS{trefy, kříže} \dicIndirectTranslationCS{(barva v~kartách)}
\dicEntry[laufabrauð] \dicTerm{laufa··brauð} \dicIPA{{l}{\oe i}{\textlengthmark}{v}{a}{\textsubring{b}}{r}{\oe i}{\texttheta}} \dicPos{n}[2] \dicFlx{(‑s, ‑)}[5] \dicFieldCat{kulin.} \dicIndirectTranslationCS{islandský vánoční chléb ručně zdobený různými motivy} \dicsymPhoto\ 
\dicFigure{ds_image_laufabraud_0_2.jpg}{Laufabrauð}{Laufabrauð - Fany Larota Catunta, PD}
\dicEntry[laufblað] \dicTerm{lauf··|blað} \dicIPA{{l}{\oe i}{v}{\textsubring{b}}{l}{a}{\texttheta}} \dicPos{n}[2] \dicFlx{(‑blaðs, ‑blöð)}[8] \dicSynonym{lauf} \dicDirectTranslationCS{list (stromu ap.)}
\dicEntry[laufgast] \dicTerm{laufg|ast} \dicIPA{{l}{\oe i}{v}{\r{g}}{a}{s}{\textsubring{d}}} \dicPos{v}[1] \dicFlx{(‑aðist)}[95] \dicFlx{refl} \dicDirectTranslationCS{nasadit listy, zazelenat se}
\dicEntry[laufglói] \dicTerm{lauf··gló|i} \dicIPA{{l}{\oe i}{v}{\r{g}}{l}{ou}{\textsci}} \dicPos{m}[1] \dicFlx{(‑a, ‑ar)}[1] \dicFieldCat{zool.} \dicSynonym{gullþröstur} \dicDirectTranslationCS{žluva, žluva hajní} \textit{(l.~{\textLA{Oriolus oriolus}})}  \dicsymPhoto\ 
\dicFigure{ds_image_laufgloi_0_1.jpg}{Laufglói}{Laufglói - Dixi, GFDL}
\dicEntry[laufgræna] \dicTerm{lauf··græn|a} \dicIPA{{l}{\oe i}{v}{\r{g}}{r}{a}{i}{n}{a}} \dicPos{f}[1] \dicFlx{(‑u)}[5] \dicFieldCat{bot.} \dicSynonym{blaðgræna} \dicDirectTranslationCS{chlorofyl}
\dicEntry[laufskógur] \dicTerm{lauf··skóg|ur} \dicIPA{{l}{\oe i}{f}{s}{\r{g}}{ou}{\textscy}{\textsubring{r}}} \dicPos{m}[6] \dicFlx{(‑ar, ‑ar)}[56] \dicDirectTranslationCS{listnatý les}
\dicEntry[laufsöngvari] \dicTerm{lauf··söngv·ar|i} \dicIPA{{l}{\oe i}{f}{s}{\oe i}{\ng}{\r{g}}{v}{a}{r}{\textsci}} \dicPos{m}[1] \dicFlx{(‑a, ‑ar)}[13] \dicFieldCat{zool.} \dicDirectTranslationCS{budníček větší} \textit{(l.~{\textLA{Phylloscopus trochilus}})}  \dicsymPhoto\ 
\dicFigure{ds_image_laufsongvari_0_1.jpg}{Laufsöngvari}{Laufsöngvari - Vogelartinfo, CC BY-SA 3.0}
\dicEntry[lauftré] \dicTerm{lauf··tré} \dicIPA{{l}{\oe i}{v}{t\smash{\textsuperscript{h}}}{r}{j}{\textepsilon}} \dicPos{n}[2] \dicFlx{(‑s, ‑)}[36] \dicDirectTranslationCS{listnatý strom, listnáč}
\dicEntry[laug] \dicTerm{laug\smash{\textsuperscript{1}}} \dicIPA{{l}{\oe i}{\textlengthmark}{x}} \dicPos{f}[4] \dicFlx{(‑ar, ‑ar)}[1] \textbf{1.} \dicSynonym{sundlaug} \dicDirectTranslationCS{(plavecký) bazén, koupaliště, plovárna}  \textbf{2.} \dicSynonym*{heit lind} \dicDirectTranslationCS{vřídlo, horký pramen}
\dicEntry[laug] \dicTerm{laug\smash{\textsuperscript{2}}} \dicIPA{{l}{\oe i}{\textlengthmark}{x}} \dicPos{v} \dicFlx{ind pf sg 1 pers} \dicLink{ljúga}
\dicEntry[lauga] \dicTerm{laug|a} \dicIPA{{l}{\oe i}{\textlengthmark}{\textbabygamma}{a}} \dicPos{v}[1] \dicFlx{(‑aði)}[1] \dicFlx{acc} \dicLangCat{zast.} \dicSynonym{þvo} \dicDirectTranslationCS{(vy)koupat} \dicExampleIS{lauga barn} \dicExampleCS{umýt dítě}
\dicEntry[laugardagskvöld] \dicTerm{laugar·dags··kvöld} \dicIPA{{l}{\oe i}{\textlengthmark}{\textbabygamma}{a}{r}{\textsubring{d}}{a}{x}{s}{k\smash{\textsuperscript{h}}}{v}{\oe}{l}{\textsubring{d}}} \dicPos{n}[2] \dicFlx{(‑s, ‑)}[5] \dicDirectTranslationCS{sobotní večer, sobota večer}
\dicEntry[laugardagur] \dicTerm{laugar··dag|ur} \dicIPA{{l}{\oe i}{\textlengthmark}{\textbabygamma}{a}{r}{\textsubring{d}}{a}{\textbabygamma}{\textscy}{\textsubring{r}}} \dicPos{m}[6] \dicFlx{(‑s, ‑ar)}[62] \dicDirectTranslationCS{sobota}
\dicEntry[Laugavegur] \dicTerm{Lauga··veg|ur} \dicIPA{{l}{\oe i}{\textlengthmark}{\textbabygamma}{a}{v}{\textepsilon}{\textbabygamma}{\textscy}{\textsubring{r}}} \dicPos{m}[10] \dicFlx{(‑ar\,/\addthin ‑s)}[28] \dicFlx{prop} \dicDirectTranslationCS{Laugavegur} \dicIndirectTranslationCS{(hlavní nákupní třída v~Reykjavíku)}
\dicEntry[lauk] \dicTerm{lauk} \dicIPA{{l}{\oe i}{\textlengthmark}{\r{g}}} \dicPos{v} \dicFlx{ind pf sg 1 pers} \dicLink{ljúka}
\dicEntry[laukur] \dicTerm{lauk|ur} \dicIPA{{l}{\oe i}{\textlengthmark}{\r{g}}{\textscy}{\textsubring{r}}} \dicPos{m}[6] \dicFlx{(‑s, ‑ar)}[22] \dicFieldCat{bot.} \dicDirectTranslationCS{cibule, cibule kuchyňská} \textit{(l.~{\textLA{Allium cepa}})};  \dicPhraseIS{laukur ættar} \dicLangCat{přen.} \dicDirectTranslationCS{pýcha rodiny}
\dicEntry[laum] \dicTerm{laum} \dicIPA{{l}{\oe i}{\textlengthmark}{\textsubring{m}}} \dicPos{n}[2] \dicFlx{(‑s)}[2] \dicPhraseIS{í laumi} \dicFlx{adv} \dicDirectTranslationCS{potají, v~tajnosti, tajně}
\dicEntry[lauma] \dicTerm{laum|a} \dicsymFrequent\  \dicIPA{{l}{\oe i}{\textlengthmark}{m}{a}} \dicPos{v}[1] \dicFlx{(‑aði)}[1] \dicFlx{dat} \dicDirectTranslationCS{(nenápadně) dát, vsunout, podstrčit, propašovat} \dicExampleIS{lauma bréfi að e‑m} \dicExampleCS{nenápadně (komu) podstrčit dopis};  \dicIdiom{laumast}{ \dicPhraseIS{laumast}} \dicFlx{refl} \dicSynonym*{læðast} \dicDirectTranslationCS{(v)plížit se, (v)krást se} \dicExampleIS{laumast úr bænum} \dicExampleCS{vykrást se z~města}
\dicEntry[laumufarþegi] \dicTerm{laumu··far·þeg|i} \dicIPA{{l}{\oe i}{\textlengthmark}{m}{\textscy}{f}{a}{\textsubring{r}}{\texttheta}{ei}{\textsci}} \dicPos{m}[1] \dicFlx{(‑a, ‑ar)}[1] \dicDirectTranslationCS{černý pasažér, černá pasažérka}
\dicEntry[laumulegur] \dicTerm{laumu··legur} \dicIPA{{l}{\oe i}{\textlengthmark}{m}{\textscy}{l}{\textepsilon}{\textbabygamma}{\textscy}{\textsubring{r}}} \dicPos{adj}[1]\dicFlx{}[-8] \dicDirectTranslationCS{tajný, skrytý, kradmý}
\dicEntry[laumuspil] \dicTerm{laumu··spil} \dicIPA{{l}{\oe i}{\textlengthmark}{m}{\textscy}{s}{\textsubring{b}}{\textsci}{\textsubring{l}}} \dicPos{n}[2] \dicFlx{(‑s, ‑)}[5] \dicDirectTranslationCS{tajnosti, tajnůstkaření}
\dicEntry[laun] \dicTerm{laun\smash{\textsuperscript{1}}} \dicIPA{{l}{\oe i}{\textlengthmark}{\textsubring{n}}} \dicPos{f}[4] \dicFlx{(‑ar)}[3] \dicSynonym{leynd} \dicDirectTranslationCS{tajemství};  \dicPhraseIS{á laun} \dicFlx{adv} \dicDirectTranslationCS{potají, v~tajnosti}
\dicEntry[laun] \dicTerm{laun\smash{\textsuperscript{2}}} \dicsymFrequent\  \dicIPA{{l}{\oe i}{\textlengthmark}{\textsubring{n}}} \dicPos{n}[2] \dicFlx{pl}[1] \textbf{1.} \dicSynonym{kaup} \dicDirectTranslationCS{plat, výplata, mzda} \dicExampleIS{há laun} \dicExampleCS{vysoký plat}  \textbf{2.} \dicSynonym{endurgjald} \dicDirectTranslationCS{vděk, odměna}  \textbf{3.} \dicSynonym{borgun} \dicDirectTranslationCS{(peněžitá) odměna, honorář}
\dicEntry[launa] \dicTerm{laun|a} \dicIPA{{l}{\oe i}{\textlengthmark}{n}{a}} \dicPos{v}[1] \dicFlx{(‑aði)}[1] \dicFlx{dat + acc} \textbf{1.} \dicSynonym{endurgjalda} \dicDirectTranslationCS{odplatit, oplatit, odvděčit se};  \dicPhraseIS{launa e‑m e‑ð} \dicDirectTranslationCS{odvděčit se (komu čím), oplatit (komu čím)}  \textbf{2.} \dicDirectTranslationCS{ohodnotit, honorovat, platit} \dicExampleIS{illa launað starf} \dicExampleCS{špatně honorovaná práce}
\dicEntry[launaauki] \dicTerm{launa··auk|i} \dicIPA{{l}{\oe i}{\textlengthmark}{n}{a}{\oe i}{\r{\textObardotlessj}}{\textsci}} \dicPos{m}[1] \dicFlx{(‑a)}[3] \dicDirectTranslationCS{příplatek\,/\addthin přídavek k~platu, prémie, bonus}
\dicEntry[launadeila] \dicTerm{launa··deil|a} \dicIPA{{l}{\oe i}{\textlengthmark}{n}{a}{\textsubring{d}}{ei}{l}{a}} \dicPos{f}[1] \dicFlx{(‑u, ‑ur)}[13] \dicDirectTranslationCS{mzdový spor}
\dicEntry[launaflokkur] \dicTerm{launa··flokk|ur} \dicIPA{{l}{\oe i}{\textlengthmark}{n}{a}{f}{l}{\textopeno}{h}{\r{g}}{\textscy}{\textsubring{r}}} \dicPos{m}[6] \dicFlx{(‑s, ‑ar)}[8] \dicDirectTranslationCS{platové zařazení, platová skupina\,/\addthin třída}
\dicEntry[launagreiðsla] \dicTerm{launa··greiðsl|a} \dicIPA{{l}{\oe i}{\textlengthmark}{n}{a}{\r{g}}{r}{ei}{ð}{s}{\textsubring{d}}{l}{a}} \dicPos{f}[1] \dicFlx{(‑u, ‑ur)}[13] \dicDirectTranslationCS{výplata mzdy}
\dicEntry[launahækkun] \dicTerm{launa··hækk|un} \dicIPA{{l}{\oe i}{\textlengthmark}{n}{a}{h}{a}{i}{h}{\r{g}}{\textscy}{\textsubring{n}}} \dicPos{f}[7] \dicFlx{(‑unar, ‑anir)}[8] \dicDirectTranslationCS{zvýšení platu, mzdový růst}
\dicEntry[launakjör] \dicTerm{launa··kjör} \dicIPA{{l}{\oe i}{\textlengthmark}{n}{a}{c\smash{\textsuperscript{h}}}{\oe}{\textsubring{r}}} \dicPos{n}[2] \dicFlx{pl}[9] \dicDirectTranslationCS{platové\,/\addthin mzdové podmínky}
\begin{xtolerant}{}{1pt}
\dicEntry[launakostnaður] \dicTerm{launa··kost·nað|ur}\addthinS\dicIPA{{l}{\oe i}{\textlengthmark}{n}{a}{k\smash{\textsuperscript{h}}}{\textopeno}{s}{\textsubring{d}}{n}{a}{ð}{\textscy}{\textsubring{r}}}\addthinS\dicPos{m}[10]\addthinS\dicFlx{(‑ar)}[9] \dicDirectTranslationCS{mzdové náklady}
\end{xtolerant}
\dicEntry[launalækkun] \dicTerm{launa··lækk|un} \dicIPA{{l}{\oe i}{\textlengthmark}{n}{a}{l}{a}{i}{h}{\r{g}}{\textscy}{\textsubring{n}}} \dicPos{f}[7] \dicFlx{(‑unar)}[9] \dicDirectTranslationCS{snížení platu, mzdový pokles}
\dicEntry[launamaður] \dicTerm{launa··|maður} \dicIPA{{l}{\oe i}{\textlengthmark}{n}{a}{m}{a}{ð}{\textscy}{\textsubring{r}}} \dicPos{m}[13] \dicFlx{(‑manns, ‑menn)}[2] \dicDirectTranslationCS{zaměstnanec, zaměstnankyně}
\dicEntry[launaskrá] \dicTerm{launa··skrá} \dicIPA{{l}{\oe i}{\textlengthmark}{n}{a}{s}{\r{g}}{r}{au}} \dicPos{f}[4] \dicFlx{(‑r\,/\addthin ‑ar, ‑r)}[21] \dicDirectTranslationCS{výplatní listina}
\dicEntry[launaskrið] \dicTerm{launa··skrið} \dicIPA{{l}{\oe i}{\textlengthmark}{n}{a}{s}{\r{g}}{r}{\textsci}{\texttheta}} \dicPos{n}[2] \dicFlx{(‑s)}[2] \dicFieldCat{ekon.} \dicDirectTranslationCS{stoupání mezd}
\dicEntry[launastigi] \dicTerm{launa··stig|i} \dicIPA{{l}{\oe i}{\textlengthmark}{n}{a}{s}{\textsubring{d}}{i}{j}{\textsci}} \dicPos{m}[1] \dicFlx{(‑a)}[3] \dicDirectTranslationCS{mzdový tarif}
\dicEntry[launatengdur] \dicTerm{launa··tengdur} \dicIPA{{l}{\oe i}{\textlengthmark}{n}{a}{t\smash{\textsuperscript{h}}}{\textepsilon}{\ng}{\textsubring{d}}{\textscy}{\textsubring{r}}} \dicPos{adj}[2]\dicFlx{}[-17] \dicDirectTranslationCS{mzdový, spojený se mzdou (náklady ap.)};  \dicPhraseIS{launatengd gjöld} \dicDirectTranslationCS{příspěvky zaměstnavatele}
\dicEntry[launauppbót] \dicTerm{launa··upp·|bót} \dicIPA{{l}{\oe i}{\textlengthmark}{n}{a}{\textscy}{h}{\textsubring{b}}{ou}{\textsubring{d}}} \dicPos{f}[8] \dicFlx{(‑bótar, ‑bætur)}[5] \dicDirectTranslationCS{platový bonus}
\dicEntry[launráð] \dicTerm{laun··ráð} \dicIPA{{l}{\oe i}{n}{r}{au}{\texttheta}} \dicPos{n}[2] \dicFlx{pl}[1] \dicSynonym{ráðabrugg} \dicDirectTranslationCS{pikle, intriky, komplot}
\dicEntry[launsátur] \dicTerm{laun··sátur} \dicIPA{{l}{\oe i}{n}{s}{au}{\textsubring{d}}{\textscy}{\textsubring{r}}} \dicPos{n}[2] \dicFlx{(‑s, ‑)}[25] \dicSynonym{fyrirsát} \dicDirectTranslationCS{léčka, nástraha}
\dicEntry[launung] \dicTerm{laun··ung} \dicIPA{{l}{\oe i}{\textlengthmark}{n}{u}{\ng}{\r{g}}} \dicPos{f}[4] \dicFlx{(‑ar)}[3] \dicSynonym{leynd} \dicDirectTranslationCS{tajemství, tajnost};  \dicPhraseIS{með launung} \dicFlx{adv} \dicSynonym*{leynilega} \dicDirectTranslationCS{tajně};  \dicPhraseIS{e‑m er engin launung á e‑u} \dicFlx{impers} \dicDirectTranslationCS{(kdo) se netají (čím)}
\dicEntry[launþegi] \dicTerm{laun··þeg|i} \dicIPA{{l}{\oe i}{n}{\texttheta}{ei}{\textsci}} \dicPos{m}[1] \dicFlx{(‑a, ‑ar)}[1] \dicDirectTranslationCS{zaměstnanec, zaměstnankyně, pracovník, pracovnice}
\dicEntry[laus] \dicTerm{laus} \dicsymFrequent\  \dicIPA{{l}{\oe i}{\textlengthmark}{s}} \dicPos{adj}[5] \dicFlx{(f ‑)}[1] \textbf{1.} \dicSynonym{óbundinn} \dicDirectTranslationCS{volný, uvolněný, povolený} \dicExampleIS{Steinninn er orðinn laus frá veggnum.} \dicExampleCS{Kámen se už uvolnil ze zdi.}  \textbf{2.} \dicSynonym{frjáls} \dicDirectTranslationCS{volný, svobodný} \dicExampleIS{Ertu laus í kvöld?} \dicExampleCS{Máš večer čas?};  \dicPhraseIS{ganga laus} \dicDirectTranslationCS{jít volně (pes ap.)};  \dicPhraseIS{láta e‑n lausan} \dicDirectTranslationCS{pustit (koho) na svobodu};  \dicPhraseIS{vera laus við e‑ð} \dicDirectTranslationCS{zbavit se (čeho), oprostit se od (čeho)}  \textbf{3.} \dicSynonym*{auður} \dicDirectTranslationCS{volný, prázdný, neobsazený} \dicExampleIS{Er þetta sæti laust?} \dicExampleCS{Je toto místo volné?} \dicAntonym{upptekinn}  \textbf{4.} \dicSynonym{léttur} \dicDirectTranslationCS{lehký} \dicExampleIS{laust högg} \dicExampleCS{lehká rána}  \textbf{5.} \dicSynonym{óáreiðanlegur} \dicDirectTranslationCS{nespolehlivý, nestabilní};  \dicPhraseIS{vera á lausu} \dicDirectTranslationCS{být nezadaný};  \dicPhraseIS{vera laus á kostunum} \dicSynonym{lauslátur} \dicDirectTranslationCS{mít volné mravy} \dicIndirectTranslationCS{(zvláště o~ženě)}  \textbf{6.} \dicSynonym{örlátur} \dicDirectTranslationCS{štědrý};  \dicPhraseIS{vera laus á fé} \dicDirectTranslationCS{být štědrý}
\dicEntry[lausafé] \dicTerm{lausa··|fé} \dicIPA{{l}{\oe i}{\textlengthmark}{s}{a}{f}{j}{\textepsilon}} \dicPos{n}[3] \dicFlx{(‑fjár)}[5] \dicDirectTranslationCS{hotovost}
\dicEntry[lausafjárstaða] \dicTerm{lausa··fjár·|staða} \dicIPA{{l}{\oe i}{\textlengthmark}{s}{a}{f}{j}{au}{\textsubring{r}}{s}{\textsubring{d}}{a}{ð}{a}} \dicPos{f}[1] \dicFlx{(‑stöðu)}[2] \dicFieldCat{ekon.} \dicDirectTranslationCS{likvidita}
\dicEntry[lausaleikur] \dicTerm{lausa··leik|ur} \dicIPA{{l}{\oe i}{\textlengthmark}{s}{a}{l}{ei}{\r{g}}{\textscy}{\textsubring{r}}} \dicPos{m}[6] \dicFlx{(‑s)}[17] \dicPhraseIS{eiga\,/\addthin eignast barn í lausaleik} \dicDirectTranslationCS{mít nemanželské dítě}
\dicEntry[lausamennska] \dicTerm{lausa··mennsk|a} \dicIPA{{l}{\oe i}{\textlengthmark}{s}{a}{m}{\textepsilon}{n}{s}{\r{g}}{a}} \dicPos{f}[1] \dicFlx{(‑u)}[5] \dicDirectTranslationCS{práce na volné noze}
\dicEntry[lausamjöll] \dicTerm{lausa··|mjöll} \dicIPA{{l}{\oe i}{\textlengthmark}{s}{a}{m}{j}{\oe}{\textsubring{d}}{\textsubring{l}}} \dicPos{f}[7] \dicFlx{(‑mjallar)}[19] \dicDirectTranslationCS{čerstvě napadlý sníh, prašan}
\dicEntry[lausasala] \dicTerm{lausa··|sala} \dicIPA{{l}{\oe i}{\textlengthmark}{s}{a}{s}{a}{l}{a}} \dicPos{f}[1] \dicFlx{(‑sölu, ‑sölur)}[20] \dicDirectTranslationCS{volný prodej}
\dicEntry[lausavísa] \dicTerm{lausa··vís|a} \dicIPA{{l}{\oe i}{\textlengthmark}{s}{a}{v}{i}{s}{a}} \dicPos{f}[1] \dicFlx{(‑u, ‑ur)}[13] \dicFieldCat{lit.} \dicDirectTranslationCS{epigram} \dicIndirectTranslationCS{(často od neznámého autora, velmi populární na Islandu)}
\dicEntry[lausblaðabók] \dicTerm{laus·blaða··|bók} \dicIPA{{l}{\oe i}{s}{\textsubring{b}}{l}{a}{ð}{a}{\textsubring{b}}{ou}{\r{g}}} \dicPos{f}[8] \dicFlx{(‑bókar, ‑bækur)}[5] \dicDirectTranslationCS{zápisník\,/\addthin sešit s~volnými listy}

\dicEntry[lausblaðamappa] \textls[15]{\dicTerm{laus·blaða··|mappa} \dicIPA{{l}{\oe i}{s}{\textsubring{b}}{l}{a}{ð}{a}{m}{a}{h}{\textsubring{b}}{a}} \dicPos{f}[1]} \dicFlx{(‑möppu, ‑möpp\-ur)}[8] \dicDirectTranslationCS{pořadač, šanon}

\dicEntry[lauslátur] \dicTerm{laus··látur} \dicIPA{{l}{\oe i}{s}{l}{au}{\textsubring{d}}{\textscy}{\textsubring{r}}} \dicPos{adj}[1]\dicFlx{}[-1] \dicDirectTranslationCS{nezávazný (flirt ap.), promiskuitní}
\dicEntry[lauslegur] \dicTerm{laus··legur} \dicIPA{{l}{\oe i}{s}{l}{\textepsilon}{\textbabygamma}{\textscy}{\textsubring{r}}} \dicPos{adj}[1]\dicFlx{}[-8] \textbf{1.} \dicSynonym{yfirborðslegur} \dicDirectTranslationCS{zběžný, letmý, hrubý (odhad ap.)} \dicExampleIS{lauslegur uppdráttur} \dicExampleCS{zběžný nákres}  \textbf{2.} \dicSynonym{óstöðugur} \dicDirectTranslationCS{volný, nepřipevněný}
\dicEntry[lauslæti] \dicTerm{laus··læti} \dicIPA{{l}{\oe i}{s}{l}{a}{i}{\textsubring{d}}{\textsci}} \dicPos{n}[2] \dicFlx{(‑s)}[20] \dicDirectTranslationCS{nezávaznost, promiskuita}
\dicEntry[lausmáll] \dicTerm{laus··máll} \dicIPA{{l}{\oe i}{s}{m}{au}{\textsubring{d}}{\textsubring{l}}} \dicPos{adj}[8]\dicFlx{}[-1] \textbf{1.} \dicSynonym{málgefinn} \dicDirectTranslationCS{upovídaný, užvaněný}  \textbf{2.} \dicSynonym*{óorðheldinn} \dicDirectTranslationCS{indiskrétní, netaktní}
\dicEntry[lausn] \dicTerm{lausn} \dicsymFrequent\  \dicIPA{{l}{\oe i}{s}{\textsubring{d}}{\textsubring{n}}} \dicPos{f}[7] \dicFlx{(‑ar, ‑ir)}[1] \textbf{1.} \dicSynonym{frelsun} \dicDirectTranslationCS{osvobození, vysvobození} \dicExampleIS{lausn úr þrældómi} \dicExampleCS{vysvobození z~otroctví}  \textbf{2.} \dicDirectTranslationCS{uvolnění, propuštění (z~práce ap.)} \dicExampleIS{lausn frá embætti} \dicExampleCS{uvolnění z~úřadu};  \dicPhraseIS{biðjast lausnar} \dicFlx{refl} \dicDirectTranslationCS{žádat o~uvolnění (z~pozice ap.)}  \textbf{3.} \dicSynonym{úrslit} \dicDirectTranslationCS{(vy)řešení, rozřešení, rozuzlení} \dicExampleIS{lausn á vandanum} \dicExampleCS{vyřešení problému}  \textbf{4.} \dicSynonym{endurlausn} \dicDirectTranslationCS{vykoupení, spasení}  \textbf{5.} \dicFieldCat{chem.} \dicDirectTranslationCS{roztok};  \dicPhraseIS{mettuð lausn} \dicFieldCat{chem.} \dicDirectTranslationCS{nasycený roztok}
\dicEntry[lausnarbeiðni] \dicTerm{lausnar··beiðn|i} \dicIPA{{l}{\oe i}{s}{\textsubring{d}}{n}{a}{r}{\textsubring{b}}{ei}{ð}{n}{\textsci}} \dicPos{f}[3] \dicFlx{(‑i, ‑ir)}[1] \dicDirectTranslationCS{výpověď, žádost o~uvolnění (z~práce ap.)};  \dicPhraseIS{senda inn lausnarbeiðni} \dicDirectTranslationCS{podat výpověď (z~práce ap.)}
\dicEntry[lausnargjald] \dicTerm{lausnar··|gjald} \dicIPA{{l}{\oe i}{s}{\textsubring{d}}{n}{a}{r}{\r{\textObardotlessj}}{a}{l}{\textsubring{d}}} \dicPos{n}[2] \dicFlx{(‑gjalds, ‑gjöld)}[8] \dicDirectTranslationCS{výkupné}
\dicEntry[lausnarhnappur] \dicTerm{lausnar··hnapp|ur} \dicIPA{{l}{\oe i}{s}{\textsubring{d}}{n}{a}{\textsubring{r}}{\textsubring{n}}{a}{h}{\textsubring{b}}{\textscy}{\textsubring{r}}} \dicPos{m}[6] \dicFlx{(‑s, ‑ar)}[25] \dicFieldCat{poč.} \dicDirectTranslationCS{klávesa escape}
\dicEntry[lausnari] \dicTerm{lausn··ar|i} \dicIPA{{l}{\oe i}{s}{\textsubring{d}}{n}{a}{r}{\textsci}} \dicPos{m}[1] \dicFlx{(‑a, ‑ar)}[13] \dicFieldCat{náb.} \dicSynonym{frelsari} \dicDirectTranslationCS{spasitel, vykupitel}
\dicEntry[laust] \dicTerm{laust\smash{\textsuperscript{1}}} \dicIPA{{l}{\oe i}{s}{\textsubring{d}}} \dicPos{v} \dicFlx{ind pf sg 1 pers} \dicLink{ljósta}
\dicEntry[laust] \dicTerm{laus|t\smash{\textsuperscript{2}}} \dicIPA{{l}{\oe i}{s}{\textsubring{d}}} \dicPos{adv} \dicFlx{(comp ‑ar, sup ‑ast)} \textbf{1.} \dicDirectTranslationCS{lehce, slabě, mírně} \dicExampleIS{berja laust að dyrum} \dicExampleCS{zaklepat lehce na dveře}  \textbf{2.} \dicDirectTranslationCS{krátce, těsně} \dicExampleIS{laust fyrir klukkan fimm} \dicExampleCS{krátce před pátou hodinou}
\dicEntry[lausung] \dicTerm{laus··ung} \dicIPA{{l}{\oe i}{\textlengthmark}{s}{u}{\ng}{\r{g}}} \dicPos{f}[4] \dicFlx{(‑ar)}[3] \textbf{1.} \dicSynonym*{staðfestuleysi} \dicDirectTranslationCS{nestabilita, nestálost}  \textbf{2.} \dicSynonym{lauslæti} \dicDirectTranslationCS{promiskuita};  \dicIdiom{lausung}{ \dicPhraseIS{gjalda lausung við lygi}} \dicLangCat{přen.} \dicDirectTranslationCS{oplatit faleš lží}
\dicEntry[laut] \dicTerm{laut\smash{\textsuperscript{1}}} \dicIPA{{l}{\oe i}{\textlengthmark}{\textsubring{d}}} \dicPos{f}[7] \dicFlx{(‑ar, ‑ir)}[6] \dicSynonym{lægð} \dicDirectTranslationCS{prohlubeň, dolík}
\dicEntry[laut] \dicTerm{laut\smash{\textsuperscript{2}}} \dicIPA{{l}{\oe i}{\textlengthmark}{\textsubring{d}}} \dicPos{v} \dicFlx{ind pf sg 1 pers} \dicLink{lúta}
\dicEntry[lautinant] \dicTerm{lautinant} \dicIPA{{l}{\oe i}{\textlengthmark}{\textsubring{d}}{\textsci}{n}{a}{\textsubring{n}}{\textsubring{d}}} \dicPos{m}[4] \dicFlx{(‑s, ‑ar)}[11] \dicFieldCat{voj.} \dicDirectTranslationCS{poručík, poručice}
\dicEntry[lax] \dicTerm{lax}\dicTerm{, laxfiskur} \dicIPA{{l}\-{a}\-{x}\-{s}\-} \dicPos{m}[4] \dicFlx{(‑, ‑ar)}[29] \dicFieldCat{zool.} \dicDirectTranslationCS{losos, losos obecný} \textit{(l.~{\textLA{Salmo salar}})}
\dicEntry[laxá] \dicTerm{lax··á} \dicIPA{{l}{a}{x}{s}{au}} \dicPos{f}[4] \dicFlx{(‑r, ‑r)}[18] \dicDirectTranslationCS{lososí řeka}
\dicEntry[laxerolía] \dicTerm{laxer··olí|a} \dicIPA{{l}{a}{x}{s}{\textepsilon}{r}{\textopeno}{l}{i}{j}{a}} \dicPos{f}[1] \dicFlx{(‑u)}[5] \dicDirectTranslationCS{ricinový olej}
\dicEntry[laxfiskur] \dicTerm{lax··fisk|ur} \dicIPA{{l}{a}{x}{s}{f}{\textsci}{s}{\r{g}}{\textscy}{\textsubring{r}}} \dicPos{m}[6] \dicFlx{(‑s, ‑ar)}[10] \dicLink{lax}
\dicEntry[laxveiði] \dicTerm{lax··veið|i} \dicIPA{{l}{a}{x}{s}{v}{ei}{ð}{\textsci}} \dicPos{f}[2] \dicFlx{(‑i\,/\addthin ‑ar, ‑ar)}[3] \textbf{1.} \dicDirectTranslationCS{lososí úlovek}  \textbf{2.} \dicDirectTranslationCS{chytání\,/\addthin lov lososů}
\dicEntry[lá] \dicTerm{lá\smash{\textsuperscript{1}}} \dicIPA{{l}{au}{\textlengthmark}} \dicPos{v}[2] \dicFlx{(‑ði, ‑ð)}[117] \dicFlx{dat} \dicSynonym{ásaka} \dicDirectTranslationCS{vyčíst, vyčítat, dávat za vinu} \dicIndirectTranslationCS{(používané v~otázce nebo záporu)} \dicExampleIS{Ég lái þér það ekki.} \dicExampleCS{Já ti to nevyčítám.};  \dicIdiom{lást}{ \dicPhraseIS{e‑m láist að (gera e‑ð)}} \dicFlx{refl impers} \dicDirectTranslationCS{(kdo) zapomíná (udělat (co)), (komu) vypadá z~paměti (udělat (co))} \dicExampleIS{Mér láðist að senda bréfið.} \dicExampleCS{Vypadlo mi (z~paměti), že mám poslat dopis.}
\dicEntry[lá] \dicTerm{lá\smash{\textsuperscript{2}}} \dicIPA{{l}{au}{\textlengthmark}} \dicPos{v} \dicFlx{ind pf sg 1 pers} \dicLink{liggja}
\dicEntry[láð] \dicTerm{láð} \dicIPA{{l}{au}{\textlengthmark}{\texttheta}} \dicPos{n}[2] \dicFlx{(‑s, ‑)}[5] \dicSynonym{land} \dicDirectTranslationCS{souš, země} \dicExampleIS{um láð og lög} \dicExampleCS{po souši a~po vodě}
\dicEntry[lág] \dicTerm{lág} \dicIPA{{l}{au}{\textlengthmark}} \dicPos{f}[4] \dicFlx{(‑ar, ‑ar)}[1] \dicSynonym{dæld} \dicDirectTranslationCS{dolina, kotlina};  \dicPhraseIS{láta e‑ð liggja í láginni} \dicLangCat{přen.} \dicDirectTranslationCS{nechat (co) ležet u~ledu}
\dicEntry[lágfiðla] \dicTerm{lág··fiðl|a} \dicIPA{{l}{au}{\textlengthmark}{f}{\textsci}{ð}{l}{a}} \dicPos{f}[1] \dicFlx{(‑u, ‑ur)}[7] \dicFieldCat{hud.} \dicDirectTranslationCS{viola}
\dicEntry[lágkúra] \dicTerm{lág··kúr|a} \dicIPA{{l}{au}{\textlengthmark}{k\smash{\textsuperscript{h}}}{u}{r}{a}} \dicPos{f}[1] \dicFlx{(‑u)}[5] \dicSynonym*{svipleysa} \dicDirectTranslationCS{banalita, přízemnost}
\dicEntry[lágkúrulegur] \textls[15]{\dicTerm{lág·kúru··legur} \dicIPA{{l}{au}{\textlengthmark}{k\smash{\textsuperscript{h}}}{u}{r}{\textscy}{l}{\textepsilon}{\textbabygamma}{\textscy}{\textsubring{r}}} \dicPos{adj}[1]\dicFlx{}[-8] \textbf{1.} \dicSynonym*{flatneskjulegur} \dicDirectTranslationCS{banální, všední, nezajímavý}} \dicExampleIS{lágkúrulegt hús} \dicExampleCS{všední dům}  \textbf{2.} \dicSynonym{lítilmótlegur} \dicDirectTranslationCS{přízemní, nízký, mrzký} \dicExampleIS{lágkúrulegur maður} \dicExampleCS{přízemní člověk}
\dicEntry[lágkveða] \dicTerm{lág··kveð|a} \dicIPA{{l}{au}{\textlengthmark}{k\smash{\textsuperscript{h}}}{v}{\textepsilon}{ð}{a}} \dicPos{f}[1] \dicFlx{(‑u, ‑ur)}[19] \dicFieldCat{lit.} \dicDirectTranslationCS{sudá stopa se slabším přízvukem}
\dicEntry[láglaunafólk] \dicTerm{lág·launa··fólk} \dicIPA{{l}{au}{\textlengthmark}{l}{\oe i}{n}{a}{f}{ou}{\textsubring{l}}{\r{g}}} \dicPos{n}[2] \dicFlx{(‑s)}[2] \dicDirectTranslationCS{osoby s~nízkým příjmem}
\dicEntry[láglendi] \dicTerm{lág··lendi} \dicIPA{{l}{au}{\textlengthmark}{l}{\textepsilon}{n}{\textsubring{d}}{\textsci}} \dicPos{n}[2] \dicFlx{(‑s)}[20] \dicDirectTranslationCS{nížina} \dicAntonym{hálendi}
\dicEntry[lágmark] \dicTerm{lág··|mark} \dicIPA{{l}{au}{\textlengthmark}{m}{a}{\textsubring{r}}{\r{g}}} \dicPos{n}[2] \dicFlx{(‑marks, ‑mörk)}[8] \dicDirectTranslationCS{minimum} \dicExampleIS{Kuldinn er í algjöru lágmarki.} \dicExampleCS{Zima je na úplném minimu.} \dicAntonym{hámark}
\dicEntry[lágmarksframfærsla] \dicTerm{lág·marks··fram·færsl|a} \dicIPA{{l}\-{au}\-{\textlengthmark}\-{m}\-{a}\-{\textsubring{r}}\-{x}\-{s}\-{f}\-{r}\-{a}\-{m}\-{f}\-{a}\-{i}\-{\textsubring{r}}\-{s}\-{\textsubring{d}}\-{l}\-{a}\-} \dicPos{f}[1] \dicFlx{(‑u)}[5] \dicFieldCat{ekon.} \dicDirectTranslationCS{životní minimum}
\dicEntry[lágmarksgjald] \dicTerm{lág·marks··|gjald} \dicIPA{{l}{au}{\textlengthmark}{m}{a}{\textsubring{r}}{x}{s}{\r{\textObardotlessj}}{a}{l}{\textsubring{d}}} \dicPos{n}[2] \dicFlx{(‑gjalds, ‑gjöld)}[8] \dicDirectTranslationCS{minimální sazba}
\begin{xtolerant}{}{1pt}
\dicEntry[lágmarkslaun] \dicTerm{lág·marks··laun}\addthinS\dicIPA{{l}{au}{\textlengthmark}{m}{a}{\textsubring{r}}{x}{s}{l}{\oe i}{\textsubring{n}}}\addthinS\dicPos{n}[2]\addthinS\dicFlx{pl}[1]\addthinS\dicDirectTranslationCS{minimální mzda\hspace{-0.06667em}\,/\addthin plat}
\end{xtolerant}
\dicEntry[lágmarkstímasókn] \dicTerm{lág·marks··tíma·sókn} \dicIPA{{l}{au}{\textlengthmark}{m}{a}{\textsubring{r}}{x}{s}{t\smash{\textsuperscript{h}}}{i}{m}{a}{s}{ou}{h}{\r{g}}{\textsubring{n}}} \dicPos{f}[7] \dicFlx{(‑ar)}[3] \dicDirectTranslationCS{minimální počet hodin školní docházky} \dicIndirectTranslationCS{(pro ukončení kurzu ap.)}
\dicEntry[lágmarksverð] \dicTerm{lág·marks··verð} \dicIPA{{l}{au}{\textlengthmark}{m}{a}{\textsubring{r}}{x}{s}{v}{\textepsilon}{r}{\texttheta}} \dicPos{n}[2] \dicFlx{(‑s)}[2] \dicDirectTranslationCS{minimální cena}
\dicEntry[lágmynd] \dicTerm{lág··mynd} \dicIPA{{l}{au}{\textlengthmark}{m}{\textsci}{n}{\textsubring{d}}} \dicPos{f}[7] \dicFlx{(‑ar, ‑ir)}[1] \dicDirectTranslationCS{reliéf}
\dicEntry[lágnætti] \dicTerm{lág··nætti} \dicIPA{{l}{au}{\textlengthmark}{n}{a}{i}{h}{\textsubring{d}}{\textsci}} \dicPos{n}[2] \dicFlx{(‑s, ‑)}[14] \dicDirectTranslationCS{půlnoc} \dicIndirectTranslationCS{(zvláště v~létě)}
\dicEntry[lágspenna] \dicTerm{lág··spenn|a} \dicIPA{{l}{au}{\textlengthmark}{s}{\textsubring{b}}{\textepsilon}{n}{a}} \dicPos{f}[1] \dicFlx{(‑u)}[5] \dicFieldCat{elek.} \dicDirectTranslationCS{nízké napětí} \dicAntonym{háspenna}
\dicEntry[lágstétt] \dicTerm{lág··stétt} \dicIPA{{l}{au}{\textlengthmark}{s}{\textsubring{d}}{j}{\textepsilon}{h}{\textsubring{d}}} \dicPos{f}[7] \dicFlx{(‑ar, ‑ir)}[1] \dicDirectTranslationCS{nižší (společenská) třída, proletariát} \dicAntonym{hástétt}
\dicEntry[lágt] \dicTerm{lágt} \dicsymFrequent\  \dicIPA{{l}{au}{x}{\textsubring{d}}} \dicPos{adv} \dicFlx{(comp lægra, sup lægst)} \textbf{1.} \dicSynonym*{lítill hljóðstyrkur} \dicDirectTranslationCS{tiše, potichu, tlumeně} \dicExampleIS{tala lágt} \dicExampleCS{mluvit potichu}  \textbf{2.} \dicSynonym{niðri} \dicDirectTranslationCS{nízko, dole}
\dicEntry[lágum] \dicTerm{lágum} \dicIPA{{l}{au}{\textlengthmark}{\textbabygamma}{\textscy}{\textsubring{m}}} \dicPos{v} \dicFlx{ind pf pl 1 pers} \dicLink{liggja}
\dicEntry[lágur] \dicTerm{lágur} \dicsymFrequent\  \dicIPA{{l}{au}{\textlengthmark}{\textscy}{\textsubring{r}}} \dicPos{adj}[10] \dicFlx{(comp lægri, sup lægstur)}[16] \textbf{1.} \dicSynonym{stuttur} \dicDirectTranslationCS{nízký, malý} \dicIndirectTranslationCS{(o~výšce)};  \dicPhraseIS{vera lágur vexti} \dicDirectTranslationCS{být malého vzrůstu}  \textbf{2.} \dicDirectTranslationCS{nízký} \dicIndirectTranslationCS{(o~čísle, ceně ap.)}  \textbf{3.} \dicDirectTranslationCS{tichý, tlumený} \dicExampleIS{lágur söngur} \dicExampleCS{tichý zpěv};  \dicPhraseIS{í lágum hljóðum} \dicFlx{adv} \dicDirectTranslationCS{tlumeně, tiše}  \textbf{4.} \dicDirectTranslationCS{nízký (tón ap.)} \dicExampleIS{lágur tónn} \dicExampleCS{nízký tón};  \dicIdiom{lágur}{ \dicPhraseIS{bera\,/\addthin bíða lægri hlut}} \dicLangCat{přen.} \dicDirectTranslationCS{tahat za kratší konec provazu}
\dicEntry[lágvaxinn] \dicTerm{lág··vaxinn} \dicIPA{{l}{au}{\textlengthmark}{v}{a}{x}{s}{\textsci}{\textsubring{n}}} \dicPos{adj}[6]\dicFlx{}[-3] \dicDirectTranslationCS{malý (vzrůstem), (jsoucí) malého vzrůstu} \dicAntonym{hávaxinn}
\dicEntry[lágþrýstingur] \dicTerm{lág··þrýst·ing|ur} \dicIPA{{l}{au}{\textlengthmark}{\texttheta}{r}{i}{s}{\textsubring{d}}{i}{\ng}{\r{g}}{\textscy}{\textsubring{r}}} \dicPos{m}[6] \dicFlx{(‑s)}[9] \textbf{1.} \dicFieldCat{med.} \dicDirectTranslationCS{nízký tlak}  \textbf{2.} \dicFieldCat{meteo.} \dicDirectTranslationCS{nízký tlak}
\dicEntry[lágþrýstisvæði] \dicTerm{lág·þrýsti··svæði} \dicIPA{{l}{au}{\textlengthmark}{\texttheta}{r}{i}{s}{\textsubring{d}}{\textsci}{s}{v}{a}{i}{ð}{\textsci}} \dicPos{n}[2] \dicFlx{(‑s, ‑)}[14] \dicFieldCat{meteo.} \dicDirectTranslationCS{oblast nízkého tlaku} \dicAntonym{háþrýstisvæði}
\dicEntry[lágþýska] \dicTerm{lág··þýsk|a} \dicIPA{{l}{au}{\textlengthmark}{\texttheta}{i}{s}{\r{g}}{a}} \dicPos{f}[1] \dicFlx{(‑u)}[5] \dicDirectTranslationCS{dolní němčina} \dicIndirectTranslationCS{(dialekt němčiny používaný v~severních částech Německa)} \dicAntonym{háþýska}
\dicEntry[lákum] \dicTerm{lákum} \dicIPA{{l}{au}{\textlengthmark}{\r{g}}{\textscy}{\textsubring{m}}} \dicPos{v} \dicFlx{ind pf pl 1 pers} \dicLink{leka}
\dicEntry[lán] \dicTerm{lán} \dicsymFrequent\  \dicIPA{{l}{au}{\textlengthmark}{\textsubring{n}}} \dicPos{n}[2] \dicFlx{(‑s, ‑)}[5] \textbf{1.} \dicSynonym{lánsfé} \dicDirectTranslationCS{(vý)půjčka, úvěr} \dicExampleIS{borga lán} \dicExampleCS{splatit půjčku};  \dicPhraseIS{taka lán} \dicDirectTranslationCS{vzít si půjčku\,/\addthin úvěr}  \textbf{2.} \dicSynonym*{það að lána} \dicDirectTranslationCS{(vy)půjčení};  \dicPhraseIS{fá e‑ð að láni} \dicDirectTranslationCS{(vy)půjčit si (co)};  \dicPhraseIS{hafa e‑ð að láni, vera með e‑ð í láni} \dicDirectTranslationCS{(vy)půjčit si (co), mít (co) půjčené}  \textbf{3.} \dicSynonym{gæfa} \dicDirectTranslationCS{štěstí, štěstěna};  \dicPhraseIS{lánið leikur við e‑n} \dicLangCat{přen.} \dicDirectTranslationCS{(kdo) má štěstí, na (koho) se usmívá štěstí};  \dicPhraseIS{það er lán í óláni að} \dicLangCat{přen.} \dicDirectTranslationCS{je to štěstí v~neštěstí, že}
\dicEntry[lána] \dicTerm{lán|a} \dicsymFrequent\  \dicIPA{{l}{au}{\textlengthmark}{n}{a}} \dicPos{v}[1] \dicFlx{(‑aði)}[1] \dicFlx{dat + acc} \dicDirectTranslationCS{(vy)půjčit, zapůjčit};  \dicPhraseIS{lána e‑m e‑ð} \dicDirectTranslationCS{půjčit (komu co)} \dicExampleIS{lána e‑m bókina} \dicExampleCS{půjčit (komu) knížku};  \dicPhraseIS{fá e‑ð lánað} \dicDirectTranslationCS{(vy)půjčit si (co)};  \dicIdiom{lána}[út]{ \dicPhraseIS{lána út e‑ð}} \dicDirectTranslationCS{(za)půjčovat (co)}
\dicEntry[lánardrottinn] \dicTerm{lánar··drott|inn} \dicIPA{{l}{au}{\textlengthmark}{n}{a}{r}{\textsubring{d}}{r}{\textopeno}{h}{\textsubring{d}}{\textsci}{\textsubring{n}}} \dicPos{m}[6] \dicFlx{(‑ins, ‑nar)}[39] \dicDirectTranslationCS{věřitel(ka)}
\dicEntry[lánasjóður] \dicTerm{lána··sjóð|ur} \dicIPA{{l}{au}{\textlengthmark}{n}{a}{s}{j}{ou}{ð}{\textscy}{\textsubring{r}}} \dicPos{m}[9] \dicFlx{(‑s, ‑ir)}[6] \dicFieldCat{ekon.} \dicDirectTranslationCS{úvěrový fond}
\dicEntry[lánastofnun] \dicTerm{lána··stofn|un} \dicIPA{{l}{au}{\textlengthmark}{n}{a}{s}{\textsubring{d}}{\textopeno}{\textsubring{b}}{n}{\textscy}{\textsubring{n}}} \dicPos{f}[7] \dicFlx{(‑unar, ‑anir)}[8] \dicFieldCat{ekon.} \dicDirectTranslationCS{úvěrová společnost}
\dicEntry[lánaveitandi] \dicTerm{lána··veit·|andi} \dicIPA{{l}{au}{\textlengthmark}{n}{a}{v}{ei}{\textsubring{d}}{a}{n}{\textsubring{d}}{\textsci}} \dicPos{m}[2] \dicFlx{(‑anda, ‑endur)}[1] \dicDirectTranslationCS{věřitel(ka)}
\dicEntry[lánlaus] \dicTerm{lán··laus} \dicIPA{{l}{au}{n}{l}{\oe i}{s}} \dicPos{adj}[5]\dicFlx{}[-1] \dicDirectTranslationCS{smolařský, nemající štěstí}
\dicEntry[lánsamur] \dicTerm{lán··|samur} \dicIPA{{l}{au}{n}{s}{a}{m}{\textscy}{\textsubring{r}}} \dicPos{adj}[1] \dicFlx{(f ‑söm)}[2] \dicSynonym{gæfusamur} \dicDirectTranslationCS{šťastný, mající štěstí}
\dicEntry[lánsfé] \dicTerm{láns··|fé} \dicIPA{{l}{au}{n}{s}{f}{j}{\textepsilon}} \dicPos{n}[3] \dicFlx{(‑fjár)}[5] \dicDirectTranslationCS{půjčka, půjčené peníze}
\dicEntry[lánsfjármögnun] \dicTerm{láns·fjár··mögn|un} \dicIPA{{l}{au}{n}{s}{f}{j}{au}{r}{m}{\oe}{\r{g}}{n}{\textscy}{\textsubring{n}}} \dicPos{f}[7] \dicFlx{(‑unar)}[12] \dicFieldCat{ekon.} \dicDirectTranslationCS{financování dluhu}
\begin{xtolerant}{}{1pt}
\dicEntry[lánskjaravísitala] \dicTerm{láns·kjara··vísi·|tala}\addthinS\dicIPA{{l}\-{au}\-{n}\-{s}\-{c\smash{\textsuperscript{h}}}\-{a}\-{r}\-{a}\-{v}\-{i}\-{s}\-{\textsci}\-{t\smash{\textsuperscript{h}}}\-{a}\-{l}\-{a}\-}\addthinS\dicPos{f}[1]\addthinS\dicFlx{(‑tölu, ‑tölur)}[14] \dicFieldCat{ekon.} \dicDirectTranslationCS{index úvěrových podmínek}
\end{xtolerant}
\dicEntry[lánskjör] \dicTerm{láns··kjör} \dicIPA{{l}{au}{n}{s}{c\smash{\textsuperscript{h}}}{\oe}{\textsubring{r}}} \dicPos{n}[2] \dicFlx{pl}[9] \dicFieldCat{ekon.} \dicDirectTranslationCS{úvěrové podmínky}
\dicEntry[lánstími] \dicTerm{láns··tím|i} \dicIPA{{l}{au}{n}{s}{t\smash{\textsuperscript{h}}}{i}{m}{\textsci}} \dicPos{m}[1] \dicFlx{(‑a)}[3] \textbf{1.} \dicFieldCat{ekon.} \dicDirectTranslationCS{doba splatnosti úvěru}  \textbf{2.} \dicDirectTranslationCS{výpůjční doba (v~knihovně ap.)}
\dicEntry[lánstraust] \dicTerm{láns··traust} \dicIPA{{l}{au}{n}{s}{t\smash{\textsuperscript{h}}}{r}{\oe i}{s}{\textsubring{d}}} \dicPos{n}[2] \dicFlx{(‑s)}[2] \dicFieldCat{ekon.} \dicDirectTranslationCS{úvěruschopnost}
\dicEntry[lántakandi] \dicTerm{lán··tak·|andi} \dicIPA{{l}{au}{n}{t\smash{\textsuperscript{h}}}{a}{\r{g}}{a}{n}{\textsubring{d}}{\textsci}} \dicPos{m}[2] \dicFlx{(‑anda, ‑endur)}[1] \dicSynonym{lánþegi} \dicDirectTranslationCS{vypůjčovatel(ka), příjemce\,/\addthin příjemkyně půjčky}
\dicEntry[lántaki] \dicTerm{lán··tak|i} \dicIPA{{l}{au}{n}{t\smash{\textsuperscript{h}}}{a}{\r{\textObardotlessj}}{\textsci}} \dicPos{m}[1] \dicFlx{(‑a, ‑ar)}[8] \dicSynonym{lánþegi} \dicDirectTranslationCS{vypůjčovatel(ka), příjemce\,/\addthin příjemkyně půjčky}
\dicEntry[lánveiting] \dicTerm{lán··veit·ing} \dicIPA{{l}{au}{n}{v}{ei}{\textsubring{d}}{i}{\ng}{\r{g}}} \dicPos{f}[4] \dicFlx{(‑ar, ‑ar)}[5] \dicFieldCat{ekon.} \dicDirectTranslationCS{poskytnutí\,/\addthin poskytování půjčky\,/\addthin úvěru}
\dicEntry[lánþegi] \dicTerm{lán··þeg|i} \dicIPA{{l}{au}{n}{\texttheta}{ei}{\textsci}} \dicPos{m}[1] \dicFlx{(‑a, ‑ar)}[1] \dicDirectTranslationCS{vypůjčovatel(ka), příjemce\,/\addthin příjemkyně půjčky}
\dicEntry[láréttur] \dicTerm{lá··réttur} \dicIPA{{l}{au}{\textlengthmark}{r}{j}{\textepsilon}{h}{\textsubring{d}}{\textscy}{\textsubring{r}}} \dicPos{adj}[1]\dicFlx{}[-10] \dicDirectTranslationCS{vodorovný, horizontální}
\dicEntry[lárpera] \dicTerm{lár··per|a} \dicIPA{{l}{au}{r}{p\smash{\textsuperscript{h}}}{\textepsilon}{r}{a}} \dicPos{f}[1] \dicFlx{(‑u, ‑ur)}[7] \dicDirectTranslationCS{avokádo} \dicIndirectTranslationCS{(plod)}
\dicEntry[lárperutré] \dicTerm{lár·peru··tré} \dicIPA{{l}{au}{r}{p\smash{\textsuperscript{h}}}{\textepsilon}{r}{\textscy}{t\smash{\textsuperscript{h}}}{r}{j}{\textepsilon}} \dicPos{n}[2] \dicFlx{(‑s, ‑)}[36] \dicFieldCat{bot.} \dicDirectTranslationCS{hruškovec přelahodný} \textit{(l.~{\textLA{Persea americana}})}  \dicsymPhoto\ 
\dicFigure{54128.jpg}{Lárperutré}{Lárperutré - Forest \& Kim [[p:2684;Starr]], Biolib, CC-BY}
\dicEntry[lárviðarlauf] \dicTerm{lár·viðar··lauf} \dicIPA{{l}{au}{r}{v}{\textsci}{ð}{a}{r}{l}{\oe i}{f}} \dicPos{n}[2] \dicFlx{(‑s, ‑)}[5] \dicDirectTranslationCS{bobkový list}
\dicEntry[lárviðarrós] \dicTerm{lár·viðar··rós} \dicIPA{{l}{au}{r}{v}{\textsci}{ð}{a}{r}{ou}{s}} \dicPos{f}[7] \dicFlx{(‑ar, ‑ir)}[1] \dicFieldCat{bot.} \dicSynonym{nería} \dicDirectTranslationCS{oleandr, oleandr obecný} \textit{(l.~{\textLA{Nerium oleander}})}  \dicsymPhoto\ 
\dicFigure{ds_image_larvidarros_0_1.jpg}{Lárviðarrós}{Lárviðarrós - Alvesgaspar, CC BY-SA 3.0}
\dicEntry[lárviður] \dicTerm{lár··við|ur} \dicIPA{{l}{au}{r}{v}{\textsci}{ð}{\textscy}{\textsubring{r}}} \dicPos{m}[10] \dicFlx{(‑ar, ‑ir)}[19] \dicFieldCat{bot.} \dicDirectTranslationCS{vavřín, vavřín vznešený} \textit{(l.~{\textLA{Laurus nobilis}})}  \dicsymPhoto\ 
\dicFigure{55905.jpg}{Lárviður}{Lárviður - Forest \& Kim [[p:2684;Starr]], Biolib, CC-BY}
\dicEntry[lás] \dicTerm{lás} \dicIPA{{l}{au}{\textlengthmark}{s}} \dicPos{m}[4] \dicFlx{(‑s, ‑ar)}[25] \textbf{1.} \dicDirectTranslationCS{zámek (u~dveří ap.)} \dicExampleIS{lásinn á hjólinu} \dicExampleCS{zámek na kolo};  \dicPhraseIS{skella e‑u í lás} \dicDirectTranslationCS{zabouchnout (si) (co) (dveře ap.)}  \textbf{2.} \dicDirectTranslationCS{pojistka} \dicIndirectTranslationCS{(součást střelné zbraně)} \dicExampleIS{lás á byssu} \dicExampleCS{pojistka pistole}
\dicEntry[lásasmiður] \dicTerm{lása··smið|ur} \dicIPA{{l}{au}{\textlengthmark}{s}{a}{s}{m}{\textsci}{ð}{\textscy}{\textsubring{r}}} \dicPos{m}[9] \dicFlx{(‑s, ‑ir)}[8] \dicDirectTranslationCS{zámečník, zámečnice}
\dicEntry[lásbogi] \dicTerm{lás··bog|i} \dicIPA{{l}{au}{\textlengthmark}{s}{\textsubring{b}}{\textopeno i}{j}{\textsci}} \dicPos{m}[1] \dicFlx{(‑a, ‑ar)}[1] \dicDirectTranslationCS{kuše, samostříl}
\dicEntry[lásum] \dicTerm{lásum} \dicIPA{{l}{au}{\textlengthmark}{s}{\textscy}{\textsubring{m}}} \dicPos{v} \dicFlx{ind pf pl 1 pers} \dicLink{lesa}
\dicEntry[lát] \dicTerm{lát} \dicIPA{{l}{au}{\textlengthmark}{\textsubring{d}}} \dicPos{n}[2] \dicFlx{(‑s, ‑)}[5] \textbf{1.} \dicSynonym{andlát} \dicDirectTranslationCS{skon, úmrtí} \dicExampleIS{auglýsa lát e‑rs} \dicExampleCS{oznámit (čí) úmrtí}  \textbf{2.} \dicSynonym{hlé} \dicDirectTranslationCS{přestávka, chvíle oddechu} \dicExampleIS{Það var ekkert lát á storminum alla nóttina.} \dicExampleCS{Bouře zuřila bez přestávky celou noc.}  \textbf{3.} \dicSynonym*{undanlát} \dicDirectTranslationCS{poddajnost} \dicExampleIS{Það er ekkert lát á honum.} \dicExampleCS{Není v~něm kousek poddajnosti.}
\dicEntry[láta] \dicTerm{láta} \dicsymFrequent\  \dicIPA{{l}{au}{\textlengthmark}{\textsubring{d}}{a}} \dicPos{v}[6] \dicFlx{(læt, lét, létum, léti, látið)}[84] \dicFlx{acc} \textbf{1.} \dicSynonym*{leyfa að e‑að gerist} \dicDirectTranslationCS{(po)nechat, (po)nechávat} \dicExampleIS{láta e‑n lesa} \dicExampleCS{nechat (koho) číst};  \dicPhraseIS{láta e‑n í friði} \dicDirectTranslationCS{nechat (koho) na pokoji};  \dicPhraseIS{láta e‑n vera} \dicDirectTranslationCS{nechat (koho) být}  \textbf{2.} \dicDirectTranslationCS{(po)nechat, (po)nechávat, zanechat, zanechávat} \dicExampleIS{láta bókina þar} \dicExampleCS{nechat tam knížku};  \dicPhraseIS{láta e‑ð af hendi} \dicDirectTranslationCS{dát (co) z~ruky};  \dicPhraseIS{láta e‑m e‑ð í té} \dicDirectTranslationCS{poskytnout (komu co), předat (komu co)}  \textbf{3.} \dicSynonym*{hegða sér} \dicDirectTranslationCS{chovat se, jednat} \dicExampleIS{láta eins og fífl} \dicExampleCS{chovat se jako hlupák};  \dicPhraseIS{láta sem ekkert sé} \dicDirectTranslationCS{dělat jakoby nic; dělat, jako by se nic nestalo};  \dicPhraseIS{láta öllum illum látum} \dicDirectTranslationCS{dělat problémy\,/\addthin potíže};  \dicPhraseIS{láttu ekki svona} \dicFlx{imper} \dicDirectTranslationCS{nech toho}  \textbf{4.} \dicPhraseIS{láta lífið} \dicDirectTranslationCS{přijít o~život}  \textbf{5.} \dicPhraseIS{láta í haf} \dicDirectTranslationCS{vyplout na moře};  \dicIdiom{láta}[að]{ \dicPhraseIS{láta vel að e‑n}} \dicDirectTranslationCS{být milý\,/\addthin laskavý ke (komu)};  \dicIdiom{láta}[af]{ \dicPhraseIS{láta af e‑u}} \dicDirectTranslationCS{(za)nechat (čeho) (práce ap.)};  \dicIdiom{láta}[aftur]{ \dicPhraseIS{láta aftur e‑ð}} \dicDirectTranslationCS{zavřít\,/\addthin zavírat (co)} \dicExampleIS{láta aftur hurðina\,/\addthin dyrnar} \dicExampleCS{zavřít dveře};  \dicIdiom{láta}[eftir]{ \dicPhraseIS{láta allt eftir sér}} \dicDirectTranslationCS{dopřávat si, nic si neodepřít}; { \dicPhraseIS{láta e‑ð eftir e‑m}} \dicDirectTranslationCS{nechat (co) na (kom), neodepřít (komu co)}; { \dicPhraseIS{láta eftir}} \dicDirectTranslationCS{podřídit se, podvolit se} \dicExampleIS{Hann varð að láta eftir.} \dicExampleCS{Musel se podřídit.};  \dicIdiom{láta}[nærri]{ \dicPhraseIS{það lætur nærri}} \dicFlx{impers} \dicDirectTranslationCS{moc tomu nechybí, téměř};  \dicIdiom{láta}[sér]{ \dicPhraseIS{láta sér e‑ð vel líka}} \dicDirectTranslationCS{být s~(čím) spokojený}; { \dicPhraseIS{láta sér ekki bregða}} \dicDirectTranslationCS{nenechat se překvapit\,/\addthin zaskočit}; { \dicPhraseIS{láta sér fátt um finnast}} \dicDirectTranslationCS{nedat najevo příliš velký zájem};  \dicIdiom{láta}[sig]{ \dicPhraseIS{láta sig ekki}} \dicDirectTranslationCS{nevzdat se, nevzdávat se};  \dicIdiom{láta}[undan]{ \dicPhraseIS{láta undan e‑m}} \dicSynonym{beygja\smash{\textsuperscript{2}}} \dicDirectTranslationCS{podvolit se (komu), podrobit se (komu)}; { \dicPhraseIS{láta undan}} \dicDirectTranslationCS{povolit, spadnout (střecha ap.)};  \dicIdiom{láta}[uppi]{ \dicPhraseIS{láta e‑ð uppi}} \dicDirectTranslationCS{vyjevit (co), dát (co) najevo};  \dicIdiom{láta}[út úr]{ \dicPhraseIS{láta e‑ð út úr sér}} \dicDirectTranslationCS{(dovolit si) říct (co)};  \dicIdiom{láta}[yfir]{ \dicPhraseIS{láta illa yfir e‑u}} \dicDirectTranslationCS{neschvalovat (co), odsuzovat (co)} \dicExampleIS{láta illa yfir ástandinu} \dicExampleCS{neschvalovat takovýto stav};  \dicIdiom{látast}{ \dicPhraseIS{látast}} \dicFlx{refl} {\textbf{a.}} \dicDirectTranslationCS{předstírat} \dicExampleIS{látast vera reiður} \dicExampleCS{předstírat rozčilení};  {\textbf{b.}} \dicDirectTranslationCS{skonat, zemřít} \dicExampleIS{látast úr elli} \dicExampleCS{zemřít v~důsledku vysokého věku};  \dicIdiom{láta}{ \dicPhraseIS{láta e‑ð eiga sig}} \dicDirectTranslationCS{odpustit si (co)}; { \dicPhraseIS{láta e‑ð ekki á sig fá}} \dicDirectTranslationCS{nepoddat se (čemu)}; { \dicPhraseIS{láta e‑ð í ljós}} \dicDirectTranslationCS{dát (co) najevo, vyjádřit (co)} \dicExampleIS{láta skoðun sína í ljós} \dicExampleCS{dát najevo své názory}; { \dicPhraseIS{láta fara vel um sig}} \dicDirectTranslationCS{udělat si pohodlí} \dicExampleIS{Afi lætur fara vel um sig í hægindastólnum.} \dicExampleCS{Děda si dělá pohodlí v~křesle.}; { \dicPhraseIS{láta ekki hugfallast}} \dicDirectTranslationCS{neztratit kuráž}; { \dicPhraseIS{láta ekki e‑ð hjá líða}} \dicDirectTranslationCS{nenechat si ujít (co)}; { \dicPhraseIS{láta til skarar skríða}} \dicLangCat{přen.} \dicDirectTranslationCS{jít na věc}
\dicEntry[látalæti] \dicTerm{láta··læti} \dicIPA{{l}{au}{\textlengthmark}{\textsubring{d}}{a}{l}{a}{i}{\textsubring{d}}{\textsci}} \dicPos{n}[2] \dicFlx{pl}[19] \dicSynonym{uppgerð} \dicDirectTranslationCS{předstírání, přetvářka, pokrytectví}
\dicEntry[látbragð] \dicTerm{lát··bragð} \dicIPA{{l}{au}{\textlengthmark}{\textsubring{d}}{\textsubring{b}}{r}{a}{\textbabygamma}{\texttheta}} \dicPos{n}[2] \dicFlx{(‑s)}[2] \textbf{1.} \dicDirectTranslationCS{posunek, gesto}  \textbf{2.} \dicDirectTranslationCS{pantomima}
\dicEntry[látbragðsleikur] \dicTerm{lát·bragðs··leik|ur} \dicIPA{{l}{au}{\textlengthmark}{\textsubring{d}}{\textsubring{b}}{r}{a}{\textbabygamma}{ð}{s}{l}{ei}{\r{g}}{\textscy}{\textsubring{r}}} \dicPos{m}[6] \dicFlx{(‑s)}[17] \dicDirectTranslationCS{němohra, pantomima}
\dicEntry[látið] \dicTerm{látið} \dicIPA{{l}{au}{\textlengthmark}{\textsubring{d}}{\textsci}{\texttheta}} \dicPos{v} \dicFlx{supin} \dicLink{láta}
\dicEntry[látinn] \dicTerm{látinn} \dicsymFrequent\  \dicIPA{{l}{au}{\textlengthmark}{\textsubring{d}}{\textsci}{\textsubring{n}}} \dicPos{adj}[6]\dicFlx{}[-6] \dicSynonym{dáinn} \dicDirectTranslationCS{zesnulý, zemřelý} \dicExampleIS{níu látnir og fjölmargir slasaðir} \dicExampleCS{devět mrtvých a~mnoho zraněných};  \dicPhraseIS{vera vant við látinn} \dicDirectTranslationCS{mít napilno, být zaneprázdněný};  \dicPhraseIS{vera vel látinn} \dicDirectTranslationCS{být populární\,/\addthin uznávaný}
\dicEntry[látlaus] \dicTerm{lát··laus} \dicIPA{{l}{au}{\textlengthmark}{\textsubring{d}}{l}{\oe i}{s}} \dicPos{adj}[5]\dicFlx{}[-1] \textbf{1.} \dicSynonym{tilgerðarlaus} \dicDirectTranslationCS{prostý, skromný, nenáročný}  \textbf{2.} \dicSynonym{linnulaus} \dicDirectTranslationCS{vytrvalý, soustavný, nepřetržitý} \dicExampleIS{látlaus rigning} \dicExampleCS{soustavný déšť}
\dicEntry[látleysi] \dicTerm{lát··leysi} \dicIPA{{l}{au}{\textlengthmark}{\textsubring{d}}{l}{ei}{s}{\textsci}} \dicPos{n}[2] \dicFlx{(‑s)}[20] \dicSynonym*{tilgerðarleysi} \dicDirectTranslationCS{prostota, přirozenost, jednoduchost}
\dicEntry[látún] \dicTerm{látún} \dicIPA{{l}{au}{\textlengthmark}{t\smash{\textsuperscript{h}}}{u}{\textsubring{n}}} \dicPos{n}[2] \dicFlx{(‑s)}[2] \dicDirectTranslationCS{mosaz}
\dicEntry[látæði] \dicTerm{lát··æði} \dicIPA{{l}{au}{\textlengthmark}{\textsubring{d}}{a}{i}{ð}{\textsci}} \dicPos{n}[2] \dicFlx{(‑s)}[20] \textbf{1.} \dicSynonym{hegðun} \dicDirectTranslationCS{chování, jednání}  \textbf{2.} \dicSynonym{látbragð} \dicDirectTranslationCS{posunek, gesto}
\dicEntry[lávarður] \dicTerm{lávarð|ur} \dicIPA{{l}{au}{\textlengthmark}{v}{a}{r}{ð}{\textscy}{\textsubring{r}}} \dicPos{m}[6] \dicFlx{(‑s\,/\addthin ‑ar, ‑ar)}[64] \dicDirectTranslationCS{lord}
\dicEntry[Lbs.] \dicTerm{Lbs.} \dicPos{zkr} \dicPhraseIS{Landsbókasafn} \dicDirectTranslationCS{Národní knihovna}
\dicEntry[ld.] \dicTerm{ld.} \dicPos{zkr} \dicPhraseIS{laugardagur} \dicDirectTranslationCS{sobota}
\dicEntry[leðja] \dicTerm{leðj|a} \dicIPA{{l}{\textepsilon}{ð}{j}{a}} \dicPos{f}[1] \dicFlx{(‑u)}[5] \dicSynonym{for\smash{\textsuperscript{1}}} \dicDirectTranslationCS{bláto, bahno} \dicExampleIS{sökkva ofan í leðjuna} \dicExampleCS{zapadnout do bláta}
\dicEntry[leður] \dicTerm{leður} \dicsymFrequent\  \dicIPA{{l}{\textepsilon}{\textlengthmark}{ð}{\textscy}{\textsubring{r}}} \dicPos{n}[2] \dicFlx{(‑s, ‑)}[25] \dicDirectTranslationCS{kůže} \dicIndirectTranslationCS{(materiál)} \dicExampleIS{vörur úr leðri} \dicExampleCS{výrobky z~kůže}
\dicEntry[leðurblaka] \dicTerm{leður··|blaka} \dicIPA{{l}{\textepsilon}{\textlengthmark}{ð}{\textscy}{r}{\textsubring{b}}{l}{a}{\r{g}}{a}} \dicPos{f}[1] \dicFlx{(‑blöku, ‑blökur)}[20] \textbf{1.} \dicPhraseIS{leðurblökur} \dicFlx{pl} \dicFieldCat{zool.} \dicDirectTranslationCS{letouni} \textit{(l.~{\textLA{Chiroptera}})}  \dicsymPhoto\   \textbf{2.} \dicFieldCat{zool.} \dicDirectTranslationCS{netopýr} \textit{(l.~{\textLA{Microchiroptera}})}
\dicFigure{62444.jpg}{Leðurblaka}{Leðurblaka - Krejčík Stanislav, Biolib, Copyright/CC-BY}
\dicEntry[leðurjakki] \dicTerm{leður··jakk|i} \dicIPA{{l}{\textepsilon}{\textlengthmark}{ð}{\textscy}{r}{j}{a}{h}{\r{\textObardotlessj}}{\textsci}} \dicPos{m}[1] \dicFlx{(‑a, ‑ar)}[8] \dicDirectTranslationCS{kožená bunda}
\dicEntry[leðurlíki] \dicTerm{leður··líki} \dicIPA{{l}{\textepsilon}{\textlengthmark}{ð}{\textscy}{r}{l}{i}{\r{\textObardotlessj}}{\textsci}} \dicPos{n}[2] \dicFlx{(‑s, ‑)}[16] \dicDirectTranslationCS{koženka}
\dicEntry[leg] \dicTerm{leg} \dicIPA{{l}{\textepsilon}{\textlengthmark}{x}} \dicPos{n}[2] \dicFlx{(‑s, ‑)}[5] \textbf{1.} \dicDirectTranslationCS{umístění, položení}  \textbf{2.} \dicFieldCat{anat.} \dicDirectTranslationCS{děloha}
\dicEntry[lega] \dicTerm{leg|a} \dicsymFrequent\  \dicIPA{{l}{\textepsilon}{\textlengthmark}{\textbabygamma}{a}} \dicPos{f}[1] \dicFlx{(‑u, ‑ur)}[19] \textbf{1.} \dicSynonym{staðsetning} \dicDirectTranslationCS{uložení, pozice} \dicExampleIS{landfræðileg lega Íslands} \dicExampleCS{geografická pozice Islandu}  \textbf{2.} \dicSynonym{rúmlega\smash{\textsuperscript{1}}} \dicDirectTranslationCS{ležení (na lůžku)}  \textbf{3.} \dicFieldCat{techn.} \dicDirectTranslationCS{ložisko}
\dicEntry[leggja] \dicTerm{leggja} \dicsymFrequent\  \dicIPA{{l}{\textepsilon}{\r{\textObardotlessj}}{\textlengthmark}{a}} \dicPos{v}[4] \dicFlx{(legg, lagði, lögðum, legði, lagt)}[20] \dicFlx{acc\,/\addthin dat} \textbf{1.} \dicFlx{acc} \dicDirectTranslationCS{položit, pokládat, umístit, umísťovat, uložit, ukládat} \dicExampleIS{leggja hlutinn á borðið} \dicExampleCS{položit věc na stůl};  \dicPhraseIS{leggja sig} \dicDirectTranslationCS{položit se, lehnout si, natáhnout se}  \textbf{2.} \dicFlx{acc} \dicDirectTranslationCS{(na)instalovat (potrubí ap.), (na)táhnout (elek. vedení ap.), položit, pokládat (dlaždice ap.)} \dicExampleIS{leggja nýjan veg} \dicExampleCS{položit novou silnici}  \textbf{3.} \dicDirectTranslationCS{vyjít, vyjet, vyrazit};  \dicPhraseIS{leggja af\,/\addthin á stað} \dicDirectTranslationCS{vyrazit (na cestu ap.), vydat se} \dicExampleIS{Hún leggur af stað til Akureyrar.} \dicExampleCS{Vyráží do Akureyri.}  \textbf{4.} \dicFlx{dat} \dicDirectTranslationCS{(za)parkovat} \dicExampleIS{leggja bílnum hjá búðinni} \dicExampleCS{zaparkovat auto u~obchodu}  \textbf{5.} \dicFlx{dat} \dicDirectTranslationCS{vyřadit, vyřazovat} \dicExampleIS{leggja gömlu kaffivélinni} \dicExampleCS{vyřadit starý kávovar}  \textbf{6.} \dicFlx{acc} \dicDirectTranslationCS{porazit, porážet (protivníka ap.)}  \textbf{7.} \dicFlx{acc} \dicDirectTranslationCS{(za)útočit, napadnout, napadat}  \textbf{8.} \dicPhraseIS{e‑ð leggur} \dicFlx{impers} {\textbf{a.}} \dicDirectTranslationCS{(co) se šíří, (co) se rozšiřuje (vůně ap.)};  {\textbf{b.}} \dicDirectTranslationCS{(co) zamrzá};  \dicIdiom{leggja}[að]{ \dicPhraseIS{leggja hart að sér}} \dicDirectTranslationCS{tvrdě pracovat, vyvinout veškeré úsilí}; { \dicPhraseIS{leggja fast að e‑m}} \dicDirectTranslationCS{vyvíjet tlak na (koho)}; { \dicPhraseIS{leggja e‑ð að veði}} \dicDirectTranslationCS{dát (co) v~sázku, vsadit (co)};  \dicIdiom{leggja}[af]{ \dicPhraseIS{leggja af}} \dicDirectTranslationCS{(z)hubnout} \dicExampleIS{Hún lagði af vegna þess að hún var veik.} \dicExampleCS{Zhubla, protože byla nemocná.};  \dicIdiom{leggja}[aftur]{ \dicPhraseIS{leggja aftur e‑ð}} \dicSynonym{loka\smash{\textsuperscript{2}}} \dicDirectTranslationCS{zavřít (co), zamknout (co)} \dicExampleIS{leggja aftur augun} \dicExampleCS{zavřít oči};  \dicIdiom{leggja}[á]{ \dicPhraseIS{leggja á borð}} \dicDirectTranslationCS{prostřít (stůl)}; { \dicPhraseIS{leggja ást á e‑n}} \dicDirectTranslationCS{zamilovat se do (koho), milovat (koho)}; { \dicPhraseIS{leggja áherslu á e‑ð}} \dicDirectTranslationCS{zdůraznit\,/\addthin zdůrazňovat (co), klást na (co) důraz}; { \dicPhraseIS{leggja mikið á sig}} \dicSynonym*{beita sig hörku} \dicDirectTranslationCS{opřít se do toho, vynasnažit se}; { \dicPhraseIS{leggja á}} \dicDirectTranslationCS{zavěsit, položit (telefon)} \dicExampleIS{Hann lagði símann á.} \dicExampleCS{Položil telefon.}; { \dicPhraseIS{leggja e‑ð á minnið, leggja e‑ð í minni}} \dicDirectTranslationCS{uložit (co) do paměti, zapamatovat si (co)}; { \dicPhraseIS{leggja skatt á e‑ð}} \dicDirectTranslationCS{(z)danit (co), uvalit daň na (co)}; { \dicPhraseIS{leggja hatur á e‑n}} \dicDirectTranslationCS{nenávidět (koho)}; { \dicPhraseIS{leggja dul á e‑ð}} \dicDirectTranslationCS{(u)tajit (co)}; { \dicPhraseIS{leggja stund á e‑ð}} \dicDirectTranslationCS{věnovat se (čemu)};  \dicIdiom{leggja}[fram]{ \dicPhraseIS{leggja sig allan fram}} \dicDirectTranslationCS{vynasnažit se, vyvinout úsilí}; { \dicPhraseIS{leggja fram e‑ð}} \dicDirectTranslationCS{předložit\,/\addthin předkládat (co) (k~diskuzi ap.)};  \dicIdiom{leggja}[fyrir]{ \dicPhraseIS{leggja e‑ð fyrir e‑n}} \dicDirectTranslationCS{předložit\,/\addthin předkládat (komu co) (otázku ap.)}; { \dicPhraseIS{leggja fyrir sig e‑ð}} \dicDirectTranslationCS{zabývat se (čím), být zapojen do (čeho)}; { \dicPhraseIS{leggja fyrir}} \dicDirectTranslationCS{odkládat, šetřit};  \dicIdiom{leggja}[inn]{ \dicPhraseIS{leggja e‑ð inn}} \dicDirectTranslationCS{uložit\,/\addthin ukládat (co)} \dicExampleIS{Ég legg peninga inn í bankann.} \dicExampleCS{Ukládám si peníze do banky.};  \dicIdiom{leggja}[í]{ \dicPhraseIS{leggja í e‑ð}} \dicDirectTranslationCS{odvážit se na (co)}; { \dicPhraseIS{leggja í hann}} \dicDirectTranslationCS{vydat se na cestu, vyrazit};  \dicIdiom{leggja}[niður]{ \dicPhraseIS{leggja niður e‑ð}} \dicDirectTranslationCS{skončit s~(čím), skoncovat s~(čím)} \dicExampleIS{leggja niður starf} \dicExampleCS{skončit s~prací}; { \dicPhraseIS{leggja niður vopn}} \dicDirectTranslationCS{složit zbraně}; \textls[15]{{ \dicPhraseIS{leggja sig niður við e‑ð}} \dicDirectTranslationCS{sní\nobreak žit se k~(čemu)};}  \dicIdiom{leggja}[saman]{ \dicPhraseIS{leggja saman e‑ð}} \dicFieldCat{mat.} \dicDirectTranslationCS{sčítat (co), přičítat (co)} \dicExampleIS{Barnið leggur saman tvo og þrjá og fær fimm.} \dicExampleCS{Dítě sečte dva plus tři a~vyjde mu pět.}; { \dicPhraseIS{leggja saman tvo og tvo}} \dicLangCat{přen.} \dicDirectTranslationCS{dát si dvě a~dvě dohromady};  \dicIdiom{leggja}[til]{ \dicPhraseIS{leggja e‑m til e‑ð}} \dicDirectTranslationCS{poskytnout\,/\addthin poskytovat (komu co)}; { \dicPhraseIS{leggja e‑ð til}} \dicDirectTranslationCS{navrhnout\,/\addthin navrhovat (co), předložit návrh o~(čem)} \dicExampleIS{Ég legg til að þessi plaköt verði lagfærð.} \dicExampleCS{Navrhuji, aby tyto plakáty byly upraveny.};  \dicPhraseIS{leggja sér e‑ð til munns} \dicDirectTranslationCS{jíst (co), pojídat (co)};  \dicIdiom{leggja}[undir]{ \dicPhraseIS{leggja undir sig e‑ð}} \dicDirectTranslationCS{podmanit\,/\addthin podmaňovat si (co), podrobit\,/\addthin podrobovat si (co)};  \dicIdiom{leggja}[upp]{ \dicPhraseIS{leggja upp}} \dicDirectTranslationCS{vyrazit, vydat se (na cestu ap.)} \dicExampleIS{leggja upp í ferðina} \dicExampleCS{vydat se na cestu};  \dicIdiom{leggja}[upp úr]{ \dicPhraseIS{leggja mikið\,/\addthin lítið upp úr e‑u}} \dicDirectTranslationCS{pokládat (co) za hodně\,/\addthin málo důležité};  \dicIdiom{leggja}[út]{ \dicPhraseIS{leggja út e‑ð}} {\textbf{a.}} \dicSynonym{borga} \dicDirectTranslationCS{vyplatit (co), přispět (čím) (částkou ap.)};  {\textbf{b.}} \dicSynonym{túlka} \dicDirectTranslationCS{přeložit (co), tlumočit (co)};  \dicIdiom{leggja}[við]{ \dicPhraseIS{leggja rækt við e‑ð}} \dicDirectTranslationCS{kultivovat (co), rozvíjet (co)}; { \dicPhraseIS{leggja e‑ð við e‑ð}} \dicDirectTranslationCS{přidat (co) k~(čemu)} \dicExampleIS{leggja tvo við þrjá} \dicExampleCS{přidat dvě ke třem};  \dicIdiom{leggjast}{ \dicPhraseIS{leggjast}} \dicFlx{refl} \dicDirectTranslationCS{lehnout si, natáhnout se};  \dicIdiom{leggjast}[að]{ \dicPhraseIS{leggjast að bryggju}} \dicFlx{refl} \dicDirectTranslationCS{přistát u~mola};  \dicIdiom{leggjast}[af]{ \dicPhraseIS{leggjast af}} \dicFlx{refl} \dicDirectTranslationCS{vymizet, zaniknout};  \dicIdiom{leggjast}[á]{ \dicPhraseIS{leggjast á eitt}} \dicFlx{refl} \dicDirectTranslationCS{spojit síly, táhnout za jeden provaz}; { \dicPhraseIS{e‑að leggst þungt á e‑n}} \dicFlx{refl} \dicLangCat{přen.} \dicDirectTranslationCS{(co koho) tíží (obavy ap.)};  \dicIdiom{leggjast}[gegn]{ \dicPhraseIS{leggjast gegn e‑u}} \dicFlx{refl} \dicDirectTranslationCS{postavit\,/\addthin stavět se proti (čemu)};  \dicIdiom{leggjast}[í]{ \dicPhraseIS{leggjast í drykkju}} \dicFlx{refl} \dicDirectTranslationCS{oddávat se pití (alkoholu)}; { \dicPhraseIS{leggjast í þunglyndi}} \dicFlx{refl} \dicDirectTranslationCS{upadnout do deprese};  \dicIdiom{leggjast}[með]{ \dicPhraseIS{leggjast með e‑m}} \dicFlx{refl} \dicDirectTranslationCS{vyspat se s~(kým), spát s~(kým)};  \dicIdiom{leggjast}[niður]{ \dicPhraseIS{leggjast niður}} \dicFlx{refl} \dicDirectTranslationCS{být přerušen, přestat};  \dicIdiom{leggjast}[yfir]{ \dicPhraseIS{e‑að leggst yfir e‑ð}} \dicFlx{refl} \dicDirectTranslationCS{(co) se usazuje nad (čím) (mlha nad městem ap.)}; { \dicPhraseIS{leggjast yfir e‑ð}} \dicFlx{refl} \dicDirectTranslationCS{pročíst (co), (důkladně) přečíst (co)}
\dicEntry[leggur] \dicTerm{legg|ur} \dicsymFrequent\  \dicIPA{{l}{\textepsilon}{\r{g}}{\textlengthmark}{\textscy}{\textsubring{r}}} \dicPos{m}[9] \dicFlx{(‑jar\,/\addthin ‑s, ‑ir)}[26] \textbf{1.} \dicSynonym{bein} \dicDirectTranslationCS{kost}  \textbf{2.} \dicSynonym{fótleggur} \dicDirectTranslationCS{noha, dolní končetina} \dicExampleIS{Ég er tilfinningalaus í leggnum.} \dicExampleCS{Nemám v~noze žádný cit.}  \textbf{3.} \dicFieldCat{bot.} \dicSynonym{stilkur} \dicDirectTranslationCS{stopka, stonek};  \dicIdiom{leggur}{ \dicPhraseIS{komast á legg}} \dicFlx{refl} \dicLangCat{přen.} \dicDirectTranslationCS{postavit se na vlastní nohy}
\dicEntry[leggöng] \dicTerm{leg··göng} \dicIPA{{l}{\textepsilon}{\textbabygamma}{\r{g}}{\oe i}{\ng}{\r{g}}} \dicPos{n}[2] \dicFlx{pl}[9] \dicFieldCat{anat.} \dicDirectTranslationCS{vagina, pochva}
\dicEntry[legið] \dicTerm{legið} \dicIPA{{l}{ei}{\textsci}{\texttheta}} \dicPos{v} \dicFlx{supin} \dicLink{liggja}
\dicEntry[legir] \dicTerm{legir} \dicIPA{{l}{ei}{\textsci}{\textsubring{r}}} \dicPos{m} \dicFlx{pl nom} \dicLink{lögur}
\dicEntry[legkaka] \dicTerm{leg··|kaka} \dicIPA{{l}{\textepsilon}{\textbabygamma}{k\smash{\textsuperscript{h}}}{a}{\r{g}}{a}} \dicPos{f}[1] \dicFlx{(‑köku, ‑kökur)}[20] \dicFieldCat{anat.} \dicDirectTranslationCS{placenta, plodový obal}
\dicEntry[legstaður] \dicTerm{leg··stað|ur} \dicIPA{{l}{\textepsilon}{x}{s}{\textsubring{d}}{a}{ð}{\textscy}{\textsubring{r}}} \dicPos{m}[10] \dicFlx{(‑ar, ‑ir)}[14] \dicSynonym{gröf} \dicDirectTranslationCS{hrob(ka)}
\dicEntry[legsteinn] \dicTerm{leg··stein|n} \dicIPA{{l}{\textepsilon}{x}{s}{\textsubring{d}}{ei}{\textsubring{d}}{\textsubring{n}}} \dicPos{m}[6] \dicFlx{(‑s, ‑ar)}[42] \dicDirectTranslationCS{náhrobní kámen}
\dicEntry[legubekkur] \dicTerm{legu··bekk|ur} \dicIPA{{l}{\textepsilon}{\textlengthmark}{\textbabygamma}{\textscy}{\textsubring{b}}{\textepsilon}{h}{\r{g}}{\textscy}{\textsubring{r}}} \dicPos{m}[9] \dicFlx{(‑s\,/\addthin ‑jar, ‑ir)}[26] \dicDirectTranslationCS{pohovka, lehátko}
\dicEntry[legudagur] \dicTerm{legu··dag|ur} \dicIPA{{l}{\textepsilon}{\textlengthmark}{\textbabygamma}{\textscy}{\textsubring{d}}{a}{\textbabygamma}{\textscy}{\textsubring{r}}} \dicPos{m}[6] \dicFlx{(‑s, ‑ar)}[62] \textbf{1.} \dicDirectTranslationCS{den na lůžku (v~nemocnici ap.)}  \textbf{2.} \dicDirectTranslationCS{den (lodi ap.) v~přístavu}
\dicEntry[legugjald] \dicTerm{legu··|gjald} \dicIPA{{l}{\textepsilon}{\textlengthmark}{\textbabygamma}{\textscy}{\r{\textObardotlessj}}{a}{l}{\textsubring{d}}} \dicPos{n}[2] \dicFlx{(‑gjalds, ‑gjöld)}[8] \dicFieldCat{nám.} \dicDirectTranslationCS{kotevné} \dicIndirectTranslationCS{(poplatek za kotvení)}
\dicEntry[legusár] \dicTerm{legu··sár} \dicIPA{{l}{\textepsilon}{\textlengthmark}{\textbabygamma}{\textscy}{s}{au}{\textsubring{r}}} \dicPos{n}[2] \dicFlx{(‑s, ‑)}[5] \dicFieldCat{med.} \dicDirectTranslationCS{proleženina}
\dicEntry[leið] \dicTerm{leið\smash{\textsuperscript{1}}} \dicsymFrequent\  \dicIPA{{l}{ei}{\textlengthmark}{\texttheta}} \dicPos{f}[7] \dicFlx{(‑ar, ‑ir)}[1] \textbf{1.} \dicSynonym{vegur} \dicDirectTranslationCS{cesta, stezka} \dicExampleIS{bein leið} \dicExampleCS{přímá cesta};  \dicPhraseIS{eiga leið hjá} \dicDirectTranslationCS{mít cestu kolem};  \dicPhraseIS{vera á leið(inni) (þangað)} \dicDirectTranslationCS{být na cestě (tam)}  \textbf{2.} \dicSynonym{háttur\smash{\textsuperscript{1}}} \dicDirectTranslationCS{způsob, postup};  \dicPhraseIS{það er engin leið að \dots{}} \dicDirectTranslationCS{není žádný způsob, jak\dots{}}  \textbf{3.} \dicSynonym{vegalengd} \dicDirectTranslationCS{cesta, vzdálenost} \dicExampleIS{dagleið} \dicExampleCS{jednodenní cesta}  \textbf{4.} \dicSynonym{áætlunarleið} \dicDirectTranslationCS{(dopravní) linka};  \dicIdiom{leið}{ \dicPhraseIS{á miðri leið}} \dicFlx{adv} \dicDirectTranslationCS{v~půli cesty}; { \dicPhraseIS{heim á leið}} \dicFlx{adv} \dicDirectTranslationCS{domů}; { \dicPhraseIS{leggja leið sína, leggja leiðir sínar}} \dicDirectTranslationCS{vydat se (do cizí země ap.)} \dicExampleIS{leggja leið sína til Íslands} \dicExampleCS{vydat se na Island};  \dicPhraseIS{um leið} \dicFlx{adv} \dicDirectTranslationCS{zároveň, současně, naráz}; { \dicPhraseIS{um leið og}} \dicFlx{conj} \dicSynonym*{strax og} \dicDirectTranslationCS{jakmile, hned jak} \dicExampleIS{Hann fór að sofa um leið og hann kom heim.} \dicExampleCS{Jakmile přišel domů, šel si lehnout.}; { \dicPhraseIS{vera komin (fimm) mánuði á leið}} \dicDirectTranslationCS{být v~(pátém) měsíci (těhotenství)}
\dicEntry[leið] \dicTerm{leið\smash{\textsuperscript{2}}} \dicIPA{{l}{ei}{\textlengthmark}{\texttheta}} \dicPos{v} \dicFlx{ind pf sg 1 pers} \dicLink{líða}
\dicEntry[leiða] \dicTerm{lei|ða} \dicsymFrequent\  \dicIPA{{l}{ei}{\textlengthmark}{ð}{a}} \dicPos{v}[2] \dicFlx{(‑ddi, ‑tt)}[169] \dicFlx{acc} \textbf{1.} \dicDirectTranslationCS{(při)vést, přivádět, vodit} \dicExampleIS{leiða e‑n þangað} \dicExampleCS{přivést (koho) tam}  \textbf{2.} \dicDirectTranslationCS{(při)vést, přivádět (elektřinu ap.)} \dicExampleIS{leiða rafmagn í hús} \dicExampleCS{přivádět elektřinu do domu};  \dicIdiom{leiða}[að]{ \dicPhraseIS{leiða hugann að e‑u}} \dicDirectTranslationCS{zavést myšlenku na (co)}; { \dicPhraseIS{leiða tal að e‑u}} \dicDirectTranslationCS{zavést řeč na (co), nadhodit (co)};  \dicIdiom{leiða}[af]{ \dicPhraseIS{leiða e‑ð af sér}} \dicDirectTranslationCS{mít (co) za následek, vést k~(čemu), skončit se (čím)};  \dicIdiom{leiða}[fyrir]{ \dicPhraseIS{leiða e‑ð e‑m fyrir sjónir}} \dicDirectTranslationCS{ukázat (komu co), vysvětlit (komu co)};  \dicIdiom{leiða}[hjá]{ \dicPhraseIS{leiða e‑ð hjá sér}} \dicDirectTranslationCS{nevšímat si (čeho), nebrat na vědomí (co)} \dicExampleIS{leiða orð hans hjá sér} \dicExampleCS{nevšímat si jeho slov};  \dicIdiom{leiða}[í]{ \dicPhraseIS{leiða e‑u í ljós}} \dicDirectTranslationCS{(po)odhalit (co)};  \dicIdiom{leiða}[til]{ \dicPhraseIS{leiða til e‑s}} \dicDirectTranslationCS{vést k~(čemu), vyústit v~(co)};  \dicIdiom{leiðast}{ \dicPhraseIS{leiðast}} \dicFlx{refl} \dicDirectTranslationCS{vést se za ruce} \dicExampleIS{Þau leiddust niður götuna.} \dicExampleCS{Šli po ulici a~vedli se za ruce.}; { \dicPhraseIS{leiðast út í drykkju}} \dicFlx{refl} \dicDirectTranslationCS{začít pít (alkohol)}; { \dicPhraseIS{láta til leiðast}} \dicDirectTranslationCS{nechat se přesvědčit};  \dicPhraseIS{e‑m leiðist} \dicFlx{refl impers} \dicDirectTranslationCS{(kdo) se nudí, (komu) je dlouhá chvíle} \dicExampleIS{Mér leiðist.} \dicExampleCS{Nudím se.}; { \dicPhraseIS{e‑m leiðist e‑að}} \dicFlx{refl impers} \dicDirectTranslationCS{(koho) nudí (co), (komu) přijde nudné (co)}
\dicEntry[leiðabók] \dicTerm{leiða··|bók} \dicIPA{{l}{ei}{\textlengthmark}{ð}{a}{\textsubring{b}}{ou}{\r{g}}} \dicPos{f}[8] \dicFlx{(‑bókar, ‑bækur)}[5] \dicDirectTranslationCS{jízdní řád}
\dicEntry[leiðangur] \dicTerm{leið··ang|ur} \dicsymFrequent\  \dicIPA{{l}{ei}{\textlengthmark}{ð}{au}{\ng}{\r{g}}{\textscy}{\textsubring{r}}} \dicPos{m}[5] \dicFlx{(‑urs, ‑rar)}[3] \dicSynonym{könnunarferð} \dicDirectTranslationCS{výprava, expedice, mise} \dicExampleIS{Leiðangurinn tekur tvær vikur.} \dicExampleCS{Expedice potrvá dva týdny.}
\dicEntry[leiðarendi] \dicTerm{leiðar··end|i} \dicIPA{{l}{ei}{\textlengthmark}{ð}{a}{r}{\textepsilon}{n}{\textsubring{d}}{\textsci}} \dicPos{m}[1] \dicFlx{(‑a, ‑ar)}[1] \dicSynonym{ákvörðunarstaður} \dicDirectTranslationCS{cíl cesty, destinace} \dicExampleIS{komast á leiðarenda} \dicExampleCS{dorazit do cíle cesty}
\dicEntry[leiðari] \dicTerm{leið··ar|i} \dicIPA{{l}{ei}{\textlengthmark}{ð}{a}{r}{\textsci}} \dicPos{m}[1] \dicFlx{(‑a, ‑ar)}[13] \textbf{1.} \dicSynonym{grein} \dicDirectTranslationCS{úvodník} \dicIndirectTranslationCS{(článek umístěný na první stránce novin)}  \textbf{2.} \dicFieldCat{fyz.} \dicDirectTranslationCS{vodič (elektřiny ap.)}
\dicEntry[leiðarljós] \dicTerm{leiðar··ljós} \dicIPA{{l}{ei}{\textlengthmark}{ð}{a}{r}{l}{j}{ou}{s}} \dicPos{n}[2] \dicFlx{(‑s, ‑)}[5] \dicLangCat{přen.} \dicDirectTranslationCS{záchytný bod, vodítko};  \dicPhraseIS{hafa e‑ð að leiðarljósi} \dicLangCat{přen.} \dicDirectTranslationCS{mít (co) jako záchytný bod}
\dicEntry[leiðarlok] \dicTerm{leiðar··lok} \dicIPA{{l}{ei}{\textlengthmark}{ð}{a}{r}{l}{\textopeno}{\r{g}}} \dicPos{n}[2] \dicFlx{pl}[1] \dicDirectTranslationCS{závěr\,/\addthin konec cesty (silnice i~putování)}
\dicEntry[leiðarlýsing] \dicTerm{leiðar··lýs·ing} \dicIPA{{l}{ei}{\textlengthmark}{ð}{a}{r}{l}{i}{s}{i}{\ng}{\r{g}}} \dicPos{f}[4] \dicFlx{(‑ar, ‑ar)}[5] \dicDirectTranslationCS{itinerář, cestovní plán}
\dicEntry[leiðarminni] \dicTerm{leiðar··minni} \dicIPA{{l}{ei}{\textlengthmark}{ð}{a}{r}{m}{\textsci}{n}{\textsci}} \dicPos{n}[2] \dicFlx{(‑s, ‑)}[14] \dicDirectTranslationCS{příznačný motiv, leitmotiv}
\dicEntry[leiðarvísir] \dicTerm{leiðar··vís|ir} \dicIPA{{l}{ei}{\textlengthmark}{ð}{a}{r}{v}{i}{s}{\textsci}{\textsubring{r}}} \dicPos{m}[7] \dicFlx{(‑is, ‑ar)}[1] \dicDirectTranslationCS{příručka, manuál, návod} \dicExampleIS{leiðarvísir um vél} \dicExampleCS{manuál k~zařízení}
\dicEntry[leiðbeina] \dicTerm{leið··bein|a} \dicIPA{{l}{ei}{ð}{\textsubring{b}}{ei}{n}{a}} \dicPos{v}[2] \dicFlx{(‑di, ‑t)}[152] \dicFlx{dat} \dicSynonym{kenna} \dicDirectTranslationCS{učit, vyučovat, cvičit, školit} \dicExampleIS{leiðbeina e‑m í e‑u} \dicExampleCS{školit (koho) v~(čem)}
\dicEntry[leiðbeinandi] \dicTerm{leið··bein·|andi} \dicIPA{{l}{ei}{ð}{\textsubring{b}}{ei}{n}{a}{n}{\textsubring{d}}{\textsci}} \dicPos{m}[2] \dicFlx{(‑anda, ‑endur)}[1] \textbf{1.} \dicDirectTranslationCS{školitel(ka), instruktor(ka), cvičitel(ka)}  \textbf{2.} \dicFieldCat{škol.} \dicDirectTranslationCS{nekvalifikovaný učitel, nekvalifikovaná učitelka}
\dicEntry[leiðbeining] \dicTerm{leið··bein·ing} \dicIPA{{l}{ei}{ð}{\textsubring{b}}{ei}{n}{i}{\ng}{\r{g}}} \dicPos{f}[4] \dicFlx{(‑ar, ‑ar)}[5] \textbf{1.} \dicSynonym{leiðsögn} \dicDirectTranslationCS{pokyn, instrukce}  \textbf{2.} \dicPhraseIS{leiðbeiningar} \dicFlx{pl} \dicSynonym{leiðarvísir} \dicDirectTranslationCS{pokyny, návod, manuál, instrukce} \dicExampleIS{leiðbeiningar um notkun} \dicExampleCS{návod k~použití}
\dicEntry[leiði] \dicTerm{leið|i\smash{\textsuperscript{1}}} \dicIPA{{l}{ei}{\textlengthmark}{ð}{\textsci}} \dicPos{m}[1] \dicFlx{(‑a)}[3] \dicSynonym{leiðindi} \dicDirectTranslationCS{nuda, dlouhá chvíle} \dicExampleIS{fá leiða á leikjum} \dicExampleCS{nudit se u~hry}
\dicEntry[leiði] \dicTerm{leiði\smash{\textsuperscript{2}}} \dicIPA{{l}{ei}{\textlengthmark}{ð}{\textsci}} \dicPos{n}[2] \dicFlx{(‑s, ‑)}[14] \textbf{1.} \dicSynonym{gröf} \dicDirectTranslationCS{náhrobek, hrobka} \dicExampleIS{setja kross á leiðið} \dicExampleCS{umístit kříž na náhrobek}  \textbf{2.} \dicSynonym{meðvindur} \dicDirectTranslationCS{příznivý vítr}  \textbf{3.} \dicSynonym{færð} \dicDirectTranslationCS{sjízdnost}
\dicEntry[leiðindamál] \dicTerm{leiðinda··mál} \dicIPA{{l}{ei}{\textlengthmark}{ð}{\textsci}{n}{\textsubring{d}}{a}{m}{au}{\textsubring{l}}} \dicPos{n}[2] \dicFlx{(‑s, ‑)}[5] \dicDirectTranslationCS{nepříjemnost, nepříjemná\,/\addthin otravná věc}
\dicEntry[leiðindi] \dicTerm{leið··indi} \dicIPA{{l}{ei}{\textlengthmark}{ð}{\textsci}{n}{\textsubring{d}}{\textsci}} \dicPos{n}[2] \dicFlx{pl}[19] \textbf{1.} \dicSynonym{leiði\smash{\textsuperscript{1}}} \dicDirectTranslationCS{nuda, nudnost, dlouhá chvíle}  \textbf{2.} \dicSynonym{óánægja} \dicDirectTranslationCS{protivnost, špatná nálada}
\dicEntry[leiðinlegur] \dicTerm{leiðin··legur} \dicsymFrequent\  \dicIPA{{l}{ei}{\textlengthmark}{ð}{\textsci}{n}{l}{\textepsilon}{\textbabygamma}{\textscy}{\textsubring{r}}} \dicPos{adj}[1]\dicFlx{}[-8] \textbf{1.} \dicDirectTranslationCS{nudný, fádní} \dicExampleIS{leiðinlegur kennari} \dicExampleCS{nudný učitel}  \textbf{2.} \dicSynonym{dapurlegur} \dicDirectTranslationCS{smutný, tristní}  \textbf{3.} \dicDirectTranslationCS{protivný, nepříjemný (počasí ap.)}
\dicEntry[leiðitamur] \textls[15]{\dicTerm{leiði··|tamur} \dicIPA{{l}{ei}{\textlengthmark}{ð}{\textsci}{t\smash{\textsuperscript{h}}}{a}{m}{\textscy}{\textsubring{r}}} \dicPos{adj}[1] \dicFlx{(f ‑töm)}[2] \dicSynonym{eftirlátur}} \dicDirectTranslationCS{svolný, povolný}
\dicEntry[leiðni] \dicTerm{leiðn|i} \dicIPA{{l}{ei}{ð}{n}{\textsci}} \dicPos{f}[3] \dicFlx{(‑i)}[3] \dicFieldCat{fyz.} \dicDirectTranslationCS{vodivost}
\dicEntry[leiðrétta] \dicTerm{leið··rétt|a} \dicIPA{{l}{ei}{ð}{r}{j}{\textepsilon}{h}{\textsubring{d}}{a}} \dicPos{v}[2] \dicFlx{(‑i, ‑)}[1] \dicFlx{acc} \dicDirectTranslationCS{opravit, opravovat, (z)korigovat (chyby ap.)} \dicExampleIS{leiðrétta verkefnið} \dicExampleCS{opravit cvičení}
\dicEntry[leiðrétting] \dicTerm{leið··rétt·ing} \dicIPA{{l}{ei}{ð}{r}{j}{\textepsilon}{h}{\textsubring{d}}{i}{\ng}{\r{g}}} \dicPos{f}[4] \dicFlx{(‑ar, ‑ar)}[5] \dicDirectTranslationCS{oprava, opravení, korekce, korektura} \dicExampleIS{leiðrétting  á mistökum} \dicExampleCS{oprava chyb}
\dicEntry[leiðsla] \dicTerm{leiðsl|a} \dicIPA{{l}{ei}{ð}{s}{\textsubring{d}}{l}{a}} \dicPos{f}[1] \dicFlx{(‑u, ‑ur)}[19] \textbf{1.} \dicSynonym{leiðsögn} \dicDirectTranslationCS{pokyn, instrukce}  \textbf{2.} \dicSynonym{rör} \dicDirectTranslationCS{roura, trubka, potrubí}  \textbf{3.} \dicSynonym{rafleiðsla} \dicDirectTranslationCS{(elektrické) vedení}  \textbf{4.} \dicSynonym*{ofurhrifning} \dicDirectTranslationCS{vytržení, trans} \dicExampleIS{vera í leiðslu} \dicExampleCS{být v~transu}
\dicEntry[leiðsögn] \dicTerm{leið··|sögn} \dicIPA{{l}{ei}{ð}{s}{\oe}{\r{g}}{\textsubring{n}}} \dicPos{f}[7] \dicFlx{(‑sagnar)}[19] \dicDirectTranslationCS{vedení, navádění, provádění}
\dicEntry[leiðsögumaður] \dicTerm{leið·sögu··|maður} \dicIPA{{l}{ei}{ð}{s}{\oe}{\textbabygamma}{\textscy}{m}{a}{ð}{\textscy}{\textsubring{r}}} \dicPos{m}[13] \dicFlx{(‑manns, ‑menn)}[2] \dicDirectTranslationCS{průvodce, průvodkyně} \dicExampleIS{leiðsögumaður ferðamanna} \dicExampleCS{průvodce turistů}
\dicEntry[leiðtogi] \dicTerm{leið··tog|i} \dicIPA{{l}{ei}{ð}{t\smash{\textsuperscript{h}}}{\textopeno i}{j}{\textsci}} \dicPos{m}[1] \dicFlx{(‑a, ‑ar)}[1] \dicSynonym{foringi} \dicDirectTranslationCS{vedoucí představitel(ka), lídr(yně)} \dicExampleIS{Flokkurinn hefur kosið sér leiðtoga.} \dicExampleCS{Strana si zvolila vedoucího představitele.}
\dicEntry[leiður] \dicTerm{leiður} \dicsymFrequent\  \dicIPA{{l}{ei}{\textlengthmark}{ð}{\textscy}{\textsubring{r}}} \dicPos{adj}[2]\dicFlx{}[-6] \textbf{1.} \dicSynonym{sorgmæddur} \dicDirectTranslationCS{zarmoucený, cítící lítost} \dicExampleIS{leiður út af e‑u} \dicExampleCS{zarmoucený kvůli čemu}  \textbf{2.} \dicSynonym{hvimleiður} \dicDirectTranslationCS{znuděný, unavený (z~poslouchání nářků ap.)} \dicExampleIS{vera orðinn leiður á verkinu} \dicExampleCS{být už unavený z~práce}
\dicEntry[leif] \dicTerm{leif} \dicsymFrequent\  \dicIPA{{l}{ei}{\textlengthmark}{f}} \dicPos{f}[4] \dicFlx{(‑ar, ‑ar)}[1] \textbf{1.} \dicLangCat{básn.} \dicSynonym{arfur} \dicDirectTranslationCS{dědictví, pozůstalost}  \textbf{2.} \dicPhraseIS{leifar} \dicFlx{pl} \dicDirectTranslationCS{zbytky, ostatky} \dicExampleIS{matarleifar} \dicExampleCS{zbytky jídla}  \textbf{3.} \dicFieldCat{mat.} \dicSynonym{afgangur} \dicDirectTranslationCS{zbytek}
\dicEntry[leifa] \dicTerm{leif|a} \dicIPA{{l}{ei}{\textlengthmark}{v}{a}} \dicPos{v}[2] \dicFlx{(‑ði, ‑t)}[99] \dicFlx{dat} \dicDirectTranslationCS{(za)nechat (jídlo nedojedené ap.)} \dicExampleIS{leifa matnum} \dicExampleCS{nechat jídlo}
\dicEntry[leiftra] \dicTerm{leiftr|a} \dicIPA{{l}{ei}{f}{\textsubring{d}}{r}{a}} \dicPos{v}[1] \dicFlx{(‑aði)}[1] \dicSynonym{glitra} \dicDirectTranslationCS{blýskat se, blyštět se, třpytit se} \dicExampleIS{Stjörnurnar leiftra á himninum.} \dicExampleCS{Hvězdy se blyští na obloze.}
\dicEntry[leiftur] \dicTerm{leift|ur\smash{\textsuperscript{1}}} \dicIPA{{l}{ei}{f}{\textsubring{d}}{\textscy}{\textsubring{r}}} \dicPos{m}[5] \dicFlx{(‑urs, ‑rar)}[1] \dicFieldCat{zool.} \dicDirectTranslationCS{plískavice běloboká} \textit{(l.~{\textLA{Lagenorhynchus acutus}})}  \dicsymPhoto\ 
\dicFigure{ds_image_leiftur_1_1.jpg}{Leiftur}{Leiftur - Twp, CC BY-SA 3.0}
\dicEntry[leiftur] \dicTerm{leiftur\smash{\textsuperscript{2}}} \dicIPA{{l}{ei}{f}{\textsubring{d}}{\textscy}{\textsubring{r}}} \dicPos{n}[2] \dicFlx{(‑s, ‑)}[25] \dicSynonym{blossi} \dicDirectTranslationCS{(za)blýsknutí, zá\-blesk}
\dicEntry[leifturhraði] \dicTerm{leiftur··hrað|i} \dicIPA{{l}{ei}{f}{\textsubring{d}}{\textscy}{\textsubring{r}}{a}{ð}{\textsci}} \dicPos{m}[1] \dicFlx{(‑a)}[3] \dicLangCat{přen.} \dicDirectTranslationCS{rychlost blesku, blesková rychlost}
\dicEntry[leiftursókn] \dicTerm{leiftur··sókn} \dicIPA{{l}{ei}{f}{\textsubring{d}}{\textscy}{\textsubring{r}}{s}{ou}{h}{\r{g}}{\textsubring{n}}} \dicPos{f}[7] \dicFlx{(‑ar, ‑ir)}[1] \dicDirectTranslationCS{bleskový útok}
\dicEntry[leiga] \dicTerm{leig|a} \dicsymFrequent\  \dicIPA{{l}{ei}{\textlengthmark}{\textbabygamma}{a}} \dicPos{f}[1] \dicFlx{(‑u, ‑ur)}[19] \textbf{1.} \dicSynonym*{það að leigja} \dicDirectTranslationCS{(pro)nájem} \dicExampleIS{leiga á íbúð} \dicExampleCS{pronájem bytu}  \textbf{2.} \dicSynonym*{leigugjald} \dicDirectTranslationCS{nájem, nájemné} \dicExampleIS{borga leiguna} \dicExampleCS{zaplatit nájem}
\dicEntry[leigja] \dicTerm{leig|ja} \dicsymFrequent\  \dicIPA{{l}{ei}{j}{\textlengthmark}{a}} \dicPos{v}[2] \dicFlx{(‑ði, ‑t)}[89] \dicFlx{acc} \dicDirectTranslationCS{(pro)najmout, (pro)najímat} \dicExampleIS{leigja íbúð af e‑m} \dicExampleCS{pronajmout si od (koho) byt}
\dicEntry[leigjandi] \dicTerm{leigj··|andi} \dicIPA{{l}{ei}{j}{\textlengthmark}{a}{n}{\textsubring{d}}{\textsci}} \dicPos{m}[2] \dicFlx{(‑anda, ‑endur)}[1] \dicDirectTranslationCS{(pod)nájemník, (pod)nájemnice}
\dicEntry[leigubifreið] \dicTerm{leigu··bif·reið} \dicIPA{{l}{ei}{\textlengthmark}{\textbabygamma}{\textscy}{\textsubring{b}}{\textsci}{v}{r}{ei}{\texttheta}} \dicPos{f}[4] \dicFlx{(‑ar, ‑ar\,/\addthin ‑ir)}[30] \dicLangCat{form.} \dicDirectTranslationCS{taxi, taxislužba}
\dicEntry[leigubílastöð] \textls[10]{\dicTerm{leigu·bíla··stöð} \dicIPA{{l}{ei}{\textlengthmark}{\textbabygamma}{\textscy}{\textsubring{b}}{i}{l}{a}{s}{\textsubring{d}}{\oe}{\texttheta}} \dicPos{f}[6] \dicFlx{(‑var, ‑var)}[1] \textbf{1.} \dicDirectTranslationCS{taxi\-služba}} \dicIndirectTranslationCS{(firma)}  \textbf{2.} \dicDirectTranslationCS{stanoviště taxi}
\dicEntry[leigubíll] \dicTerm{leigu··bíl|l} \dicIPA{{l}{ei}{\textlengthmark}{\textbabygamma}{\textscy}{\textsubring{b}}{i}{\textsubring{d}}{\textsubring{l}}} \dicPos{m}[6] \dicFlx{(‑s, ‑ar)}[48] \dicDirectTranslationCS{taxi, taxík}
\dicEntry[leigubílstjóri] \dicTerm{leigu··bíl·stjór|i} \dicIPA{{l}{ei}{\textlengthmark}{\textbabygamma}{\textscy}{\textsubring{b}}{i}{l}{s}{\textsubring{d}}{j}{ou}{r}{\textsci}} \dicPos{m}[1] \dicFlx{(‑a, ‑ar)}[1] \dicDirectTranslationCS{taxikář(ka), řidič(ka) taxi}
\dicEntry[leiguflug] \dicTerm{leigu··flug} \dicIPA{{l}{ei}{\textlengthmark}{\textbabygamma}{\textscy}{f}{l}{\textscy}{x}} \dicPos{n}[2] \dicFlx{(‑s, ‑)}[5] \dicDirectTranslationCS{charterový let}
\dicEntry[leiguíbúð] \dicTerm{leigu··í·búð} \dicIPA{{l}{ei}{\textlengthmark}{\textbabygamma}{\textscy}{i}{\textsubring{b}}{u}{\texttheta}} \dicPos{f}[7] \dicFlx{(‑ar, ‑ir)}[1] \dicDirectTranslationCS{pronajatý byt}
\dicEntry[leiguliði] \dicTerm{leigu··lið|i} \dicIPA{{l}{ei}{\textlengthmark}{\textbabygamma}{\textscy}{l}{\textsci}{ð}{\textsci}} \dicPos{m}[1] \dicFlx{(‑a, ‑ar)}[1] \dicDirectTranslationCS{nájemce, nájemkyně}
\dicEntry[leigumiðlari] \dicTerm{leigu··miðl·ar|i} \dicIPA{{l}{ei}{\textlengthmark}{\textbabygamma}{\textscy}{m}{\textsci}{ð}{l}{a}{r}{\textsci}} \dicPos{m}[1] \dicFlx{(‑a, ‑ar)}[13] \dicDirectTranslationCS{zprostředkovatel(ka) pronájmu}
\dicEntry[leigumorð] \dicTerm{leigu··morð} \dicIPA{{l}{ei}{\textlengthmark}{\textbabygamma}{\textscy}{m}{\textopeno}{r}{\texttheta}} \dicPos{n}[2] \dicFlx{(‑s, ‑)}[5] \dicDirectTranslationCS{nájemná vražda}
\dicEntry[leigumorðingi] \dicTerm{leigu··morð·ing|i} \dicIPA{{l}{ei}{\textlengthmark}{\textbabygamma}{\textscy}{m}{\textopeno}{r}{ð}{i}{\textltailn}{\r{\textObardotlessj}}{\textsci}} \dicPos{m}[1] \dicFlx{(‑ja, ‑jar)}[14] \dicDirectTranslationCS{nájemný vrah, nájemná vražedkyně}
\dicEntry[leigusali] \dicTerm{leigu··sal|i} \dicIPA{{l}{ei}{\textlengthmark}{\textbabygamma}{\textscy}{s}{a}{l}{\textsci}} \dicPos{m}[1] \dicFlx{(‑a, ‑ar)}[8] \dicDirectTranslationCS{pronajímatel(ka)}
\dicEntry[leigusamningur] \dicTerm{leigu··samn·ing|ur} \dicIPA{{l}{ei}{\textlengthmark}{\textbabygamma}{\textscy}{s}{a}{m}{n}{i}{\ng}{\r{g}}{\textscy}{\textsubring{r}}} \dicPos{m}[6] \dicFlx{(‑s, ‑ar)}[8] \dicDirectTranslationCS{nájem\-ní smlouva}
\dicEntry[leigutaki] \dicTerm{leigu··tak|i} \dicIPA{{l}{ei}{\textlengthmark}{\textbabygamma}{\textscy}{t\smash{\textsuperscript{h}}}{a}{\r{\textObardotlessj}}{\textsci}} \dicPos{m}[1] \dicFlx{(‑a, ‑ar)}[8] \dicDirectTranslationCS{nájemce, nájemkyně}
\dicEntry[leik] \dicTerm{leik} \dicIPA{{l}{ei}{\textlengthmark}{\r{g}}} \dicPos{v} \dicFlx{ind praes sg 1 pers} \dicLink{leika}
\dicEntry[leika] \dicTerm{leika} \dicsymFrequent\  \dicIPA{{l}{ei}{\textlengthmark}{\r{g}}{a}} \dicPos{v}[6] \dicFlx{(leik, lék, lékum, léki, leikið)}[63] \dicFlx{acc} \textbf{1.} \dicDirectTranslationCS{(za)hrát} \dicIndirectTranslationCS{(bavit se pro zábavu)};  \dicPhraseIS{leika sér} \dicDirectTranslationCS{hrát si} \dicExampleIS{Börnin leika sér saman.} \dicExampleCS{Děti si spolu hrají.};  \dicPhraseIS{leika sér að e‑u} \dicDirectTranslationCS{hrát si s~(čím)} \dicExampleIS{leika sér að dúkkum} \dicExampleCS{hrát si s~panenkami}  \textbf{2.} \dicDirectTranslationCS{(za)hrát (si)} \dicIndirectTranslationCS{(provozovat jako sport)} \dicExampleIS{leika knattspyrnu} \dicExampleCS{hrát fotbal}  \textbf{3.} \dicDirectTranslationCS{(za)hrát} \dicIndirectTranslationCS{(provozovat hudbu)};  \dicPhraseIS{leika á hljóðfæri} \dicDirectTranslationCS{hrát na hudební nástroj} \dicExampleIS{Hún lék á fiðlu sína.} \dicExampleCS{Hrála na housle.}  \textbf{4.} \dicDirectTranslationCS{(za)hrát} \dicIndirectTranslationCS{(předvádět divadelní hru nebo osobu v~ní)} \dicExampleIS{leika Hamlet} \dicExampleCS{hrát Hamleta}  \textbf{5.} \dicPhraseIS{leika e‑n grátt\,/\addthin hart\,/\addthin illa} \dicDirectTranslationCS{zatočit s~(kým), nemazlit se s~(kým)};  \dicIdiom{leika}[af]{ \dicPhraseIS{leika af sér}} \dicDirectTranslationCS{udělat špatný tah (v~šachu ap.)};  \dicIdiom{leika}[á]{ \dicPhraseIS{leika á e‑n}} \dicDirectTranslationCS{vyzrát na (koho), přelstít (koho)}; { \dicPhraseIS{leika á als oddi, leika á alsoddi}} \dicLangCat{přen.} \dicDirectTranslationCS{mít výbornou náladu, být v~dobrém rozmaru}; { \dicPhraseIS{það leikur ekki á tveim tungum}} \dicFlx{impers} \dicLangCat{přen.} \dicDirectTranslationCS{o~tom není pochyb};  \dicIdiom{leika}[eftir]{ \dicPhraseIS{leika e‑ð eftir e‑m}} \dicDirectTranslationCS{zopakovat (co) po (kom)};  \dicIdiom{leika}[í]{ \dicPhraseIS{leika í e‑u}} \dicDirectTranslationCS{hrát v~(čem) (ve filmu ap.)};  \dicIdiom{leika}[undir]{ \dicPhraseIS{leika undir (e‑ð)}} \dicFieldCat{hud.} \dicDirectTranslationCS{doprovodit\,/\addthin doprovázet ((co))};  \dicIdiom{leika}[við]{ \dicPhraseIS{leika við e‑n}} \dicDirectTranslationCS{hrát si s~(kým)} \dicExampleIS{Lilja leikur við mömmu sína.} \dicExampleCS{Lilja si hraje s~maminkou.}; { \dicPhraseIS{leika við hvern sinn fingur}} \dicLangCat{přen.} \dicDirectTranslationCS{být dobře naladěn, mít dobrou náladu}
\dicEntry[leikandi] \dicTerm{leik··|andi\smash{\textsuperscript{1}}} \dicIPA{{l}{ei}{\textlengthmark}{\r{g}}{a}{n}{\textsubring{d}}{\textsci}} \dicPos{m}[2] \dicFlx{(‑anda, ‑endur)}[1] \textbf{1.} \dicDirectTranslationCS{hráč(ka) (ve fotbale ap.)}  \textbf{2.} \dicDirectTranslationCS{herec, herečka}
\dicEntry[leikandi] \dicTerm{leik··andi\smash{\textsuperscript{2}}} \dicIPA{{l}{ei}{\textlengthmark}{\r{g}}{a}{n}{\textsubring{d}}{\textsci}} \dicPos{adv} \dicDirectTranslationCS{hravě, s~lehkostí} \dicExampleIS{stökkva leikandi yfir skurðinn} \dicExampleCS{přeskočit hravě (přes) příkop}
\dicEntry[leikaraskapur] \dicTerm{leikara··skap|ur} \dicIPA{{l}{ei}{\textlengthmark}{\r{g}}{a}{r}{a}{s}{\r{g}}{a}{\textsubring{b}}{\textscy}{\textsubring{r}}} \dicPos{m}[10] \dicFlx{(‑ar)}[15] \textbf{1.} \dicSynonym{fíflalæti} \dicDirectTranslationCS{šaškoviny, komediantství}  \textbf{2.} \dicSynonym{látalæti} \dicDirectTranslationCS{předstírání, hraní}
\dicEntry[leikari] \dicTerm{leik··ar|i} \dicsymFrequent\  \dicIPA{{l}{ei}{\textlengthmark}{\r{g}}{a}{r}{\textsci}} \dicPos{m}[1] \dicFlx{(‑a, ‑ar)}[13] \dicDirectTranslationCS{herec, herečka} \dicExampleIS{Leikarinn hlaut verðlaun.} \dicExampleCS{Herec obdržel ocenění.}
\dicEntry[leikbrúða] \dicTerm{leik··brúð|a} \dicIPA{{l}{ei}{\textlengthmark}{\r{g}}{\textsubring{b}}{r}{u}{ð}{a}} \dicPos{f}[1] \dicFlx{(‑u, ‑ur)}[19] \textbf{1.} \dicSynonym{brúða} \dicDirectTranslationCS{loutka} \dicIndirectTranslationCS{(v~loutkovém divadle)}  \textbf{2.} \dicSynonym{handbendi} \dicDirectTranslationCS{loutka} \dicIndirectTranslationCS{(zmanipulovaný člověk)}
\dicEntry[leikdómur] \dicTerm{leik··dóm|ur} \dicIPA{{l}{ei}{\textlengthmark}{\r{g}}{\textsubring{d}}{ou}{m}{\textscy}{\textsubring{r}}} \dicPos{m}[6] \dicFlx{(‑s, ‑ar)}[10] \dicDirectTranslationCS{divadelní kritika}
\dicEntry[leikfang] \dicTerm{leik··|fang} \dicIPA{{l}{ei}{\textlengthmark}{\r{g}}{f}{au}{\ng}{\r{g}}} \dicPos{n}[2] \dicFlx{(‑fangs, ‑föng)}[8] \dicDirectTranslationCS{hračka (na hraní)}
\dicEntry[leikfangabúð] \dicTerm{leik·fanga··búð} \dicIPA{{l}{ei}{\textlengthmark}{\r{g}}{f}{au}{\ng}{\r{g}}{a}{\textsubring{b}}{u}{\texttheta}} \dicPos{f}[7] \dicFlx{(‑ar, ‑ir)}[1] \dicDirectTranslationCS{hračkářství}
\dicEntry[leikfélag] \dicTerm{leik··fé·|lag} \dicIPA{{l}{ei}{\textlengthmark}{\r{g}}{f}{j}{\textepsilon}{l}{a}{x}} \dicPos{n}[2] \dicFlx{(‑lags, ‑lög)}[8] \dicDirectTranslationCS{divadelní společnost}
\dicEntry[leikfélagi] \dicTerm{leik··fé·lag|i} \dicIPA{{l}{ei}{\textlengthmark}{\r{g}}{f}{j}{\textepsilon}{l}{a}{i}{j}{\textsci}} \dicPos{m}[1] \dicFlx{(‑a, ‑ar)}[8] \dicDirectTranslationCS{kamarád(ka) z~dětství}
\dicEntry[leikfimi] \dicTerm{leik··fim|i} \dicIPA{{l}{ei}{\textlengthmark}{\r{g}}{f}{\textsci}{m}{\textsci}} \dicPos{f}[3] \dicFlx{(‑i)}[3] \dicDirectTranslationCS{tělocvik, (tělesné) cvičení}
\dicEntry[leikfimisalur] \dicTerm{leik·fimi··sal|ur} \dicIPA{{l}{ei}{\textlengthmark}{\r{g}}{f}{\textsci}{m}{\textsci}{s}{a}{l}{\textscy}{\textsubring{r}}} \dicPos{m}[10] \dicFlx{(‑ar, ‑ir)}[14] \dicDirectTranslationCS{tělocvična}
\dicEntry[leikfimiskennari] \dicTerm{leik·fimis··kenn·ar|i} \dicIPA{{l}\-{ei}\-{\textlengthmark}\-{\r{g}}\-{f}\-{\textsci}\-{m}\-{\textsci}\-{s}\-{c\smash{\textsuperscript{h}}}\-{\textepsilon}\-{n}\-{a}\-{r}\-{\textsci}\-} \dicPos{m}[1] \dicFlx{(‑a, ‑ar)}[13] \dicDirectTranslationCS{učitel(ka) tělesné výchovy, tělocvikář(ka)}
\dicEntry[leikgagnrýnandi] \dicTerm{leik··gagn·rýn·|andi} \dicIPA{{l}{ei}{\r{g}}{\textlengthmark}{a}{\r{g}}{n}{r}{i}{n}{a}{n}{\textsubring{d}}{\textsci}} \dicPos{m}[2] \dicFlx{(‑anda, ‑endur)}[1] \dicDirectTranslationCS{divadelní kritik\,/\addthin kritička}
\dicEntry[leikgerð] \dicTerm{leik··gerð} \dicIPA{{l}{ei}{\textlengthmark}{\r{g}}{\r{\textObardotlessj}}{\textepsilon}{r}{\texttheta}} \dicPos{f}[7] \dicFlx{(‑ar, ‑ir)}[1] \dicFieldCat{lit.} \dicDirectTranslationCS{dramatizace} \dicExampleIS{leikgerð Sölku Völku} \dicExampleCS{dramatizace Salky Valky}
\dicEntry[leikgrind] \dicTerm{leik··grind} \dicIPA{{l}{ei}{\r{g}}{\textlengthmark}{r}{\textsci}{n}{\textsubring{d}}} \dicPos{f}[8] \dicFlx{(‑ar, ‑ur)}[1] \dicDirectTranslationCS{(dětská) ohrádka}
\dicEntry[leikhús] \dicTerm{leik··hús} \dicIPA{{l}{ei}{\textlengthmark}{\r{g}}{h}{u}{s}} \dicPos{n}[2] \dicFlx{(‑s, ‑)}[5] \dicDirectTranslationCS{divadlo} \dicExampleIS{sýning í leikhúsinu} \dicExampleCS{představení v~divadle}
\dicEntry[leikhússtjóri] \dicTerm{leik·hús··stjór|i} \dicIPA{{l}{ei}{\textlengthmark}{\r{g}}{h}{u}{s}{\textsubring{d}}{j}{ou}{r}{\textsci}} \dicPos{m}[1] \dicFlx{(‑a, ‑ar)}[1] \dicDirectTranslationCS{divadelní ředitel(ka), ředitel(ka) divadla}
\dicEntry[leikið] \dicTerm{leikið} \dicIPA{{l}{ei}{\textlengthmark}{\r{\textObardotlessj}}{\textsci}{\texttheta}} \dicPos{v} \dicFlx{supin} \dicLink{leika}
\dicEntry[leikinn] \dicTerm{leikinn} \dicIPA{{l}{ei}{\textlengthmark}{\r{\textObardotlessj}}{\textsci}{\textsubring{n}}} \dicPos{adj}[6]\dicFlx{}[-2] \dicSynonym{snjall} \dicDirectTranslationCS{obratný, šikovný, dovedný} \dicExampleIS{vera leikinn í e‑u} \dicExampleCS{být šikovný v~(čem)};  \dicPhraseIS{vera grátt\,/\addthin illa leikinn} \dicDirectTranslationCS{nést známky špatného zacházení}
\dicEntry[leikjanámskeið] \dicTerm{leikja··nám·skeið} \dicIPA{{l}{ei}{\textlengthmark}{\r{\textObardotlessj}}{a}{n}{au}{m}{s}{\r{\textObardotlessj}}{ei}{\texttheta}} \dicPos{n}[2] \dicFlx{(‑s, ‑)}[5] \dicDirectTranslationCS{tábor, (dětský) prázdninový kurz}
\dicEntry[leikjatölva] \dicTerm{leikja··tölv|a} \dicIPA{{l}{ei}{\textlengthmark}{\r{\textObardotlessj}}{a}{t\smash{\textsuperscript{h}}}{\oe}{l}{v}{a}} \dicPos{f}[1] \dicFlx{(‑u, ‑ur)}[7] \dicFieldCat{poč.} \dicDirectTranslationCS{(hrací) konzole}
\dicEntry[leikkerahjól] \dicTerm{leik·kera··hjól} \dicIPA{{l}{ei}{\textlengthmark}{\r{g}}{c\smash{\textsuperscript{h}}}{\textepsilon}{r}{a}{\c{c}}{ou}{\textsubring{l}}} \dicPos{n}[2] \dicFlx{(‑s, ‑)}[5] \dicDirectTranslationCS{hrnčířské kolo}
\dicEntry[leikkerasmiður] \dicTerm{leik·kera··smið|ur} \dicIPA{{l}{ei}{\textlengthmark}{\r{g}}{c\smash{\textsuperscript{h}}}{\textepsilon}{r}{a}{s}{m}{\textsci}{ð}{\textscy}{\textsubring{r}}} \dicPos{m}[9] \dicFlx{(‑s, ‑ir)}[11] \dicDirectTranslationCS{hrnčíř(ka), keramik, keramička}
\dicEntry[leikkona] \dicTerm{leik··kon|a} \dicIPA{{l}{ei}{\textlengthmark}{\r{g}}{k\smash{\textsuperscript{h}}}{\textopeno}{n}{a}} \dicPos{f}[1] \dicFlx{(‑u, ‑ur)}[27] \dicDirectTranslationCS{herečka}
\dicEntry[leiklist] \dicTerm{leik··list} \dicIPA{{l}{ei}{\textlengthmark}{\r{g}}{l}{\textsci}{s}{\textsubring{d}}} \dicPos{f}[7] \dicFlx{(‑ar)}[3] \dicDirectTranslationCS{divadlo, divadelní umění}
\dicEntry[leiklistarskóli] \dicTerm{leik·listar··skól|i} \dicIPA{{l}{ei}{\textlengthmark}{\r{g}}{l}{\textsci}{s}{\textsubring{d}}{a}{\textsubring{r}}{s}{\r{g}}{ou}{l}{\textsci}} \dicPos{m}[1] \dicFlx{(‑a, ‑ar)}[1] \dicFieldCat{škol.} \dicDirectTranslationCS{divadelní škola}
\dicEntry[leikmaður] \dicTerm{leik··|maður} \dicIPA{{l}{ei}{\textlengthmark}{\r{g}}{m}{a}{ð}{\textscy}{\textsubring{r}}} \dicPos{m}[13] \dicFlx{(‑manns, ‑menn)}[2] \textbf{1.} \dicDirectTranslationCS{laik, laička, neodborník, neodbornice}  \textbf{2.} \dicDirectTranslationCS{hráč(ka)}  \textbf{3.} \dicSynonym{leikari} \dicDirectTranslationCS{herec, herečka}
\dicEntry[leikmunur] \dicTerm{leik··mun|ur} \dicIPA{{l}{ei}{\textlengthmark}{\r{g}}{m}{\textscy}{n}{\textscy}{\textsubring{r}}} \dicPos{m}[10] \dicFlx{(‑ar, ‑ir)}[11] \dicDirectTranslationCS{rekvizita}
\dicEntry[leikmynd] \dicTerm{leik··mynd} \dicIPA{{l}{ei}{\textlengthmark}{\r{g}}{m}{\textsci}{n}{\textsubring{d}}} \dicPos{f}[7] \dicFlx{(‑ar, ‑ir)}[1] \dicDirectTranslationCS{scénografie, scéna}
\dicEntry[leikmyndahönnuður] \dicTerm{leik·mynda··hönn·uð|ur} \dicIPA{{l}\-{ei}\-{\textlengthmark}\-{\r{g}}\-{m}\-{\textsci}\-{n}\-{\textsubring{d}}\-{a}\-{h}\-{\oe}\-{n}\-{\textscy}\-{ð}\-{\textscy}\-{\textsubring{r}}\-} \dicPos{m}[10] \dicFlx{(‑ar, ‑ir)}[4] \dicDirectTranslationCS{jevištní výtvarník\,/\addthin výtvarnice, scénograf(ka)}
\dicEntry[leikni] \dicTerm{leikn|i} \dicIPA{{l}{ei}{h}{\r{g}}{n}{\textsci}} \dicPos{f}[3] \dicFlx{(‑i)}[3] \dicDirectTranslationCS{šikovnost, obratnost, zručnost, dovednost} \dicExampleIS{leikni í íþróttum} \dicExampleCS{dovednost ve sportu}
\dicEntry[leikregla] \dicTerm{leik··regl|a} \dicIPA{{l}{ei}{\textlengthmark}{\r{g}}{r}{\textepsilon}{\r{g}}{l}{a}} \dicPos{f}[1] \dicFlx{(‑u, ‑ur)}[19] \dicDirectTranslationCS{pravidlo hry}
\dicEntry[leikrit] \dicTerm{leik··rit} \dicIPA{{l}{ei}{\textlengthmark}{\r{g}}{r}{\textsci}{\textsubring{d}}} \dicPos{n}[2] \dicFlx{(‑s, ‑)}[5] \dicDirectTranslationCS{divadelní hra}
\dicEntry[leikritahöfundur] \dicTerm{leik·rita··höf·und|ur} \dicIPA{{l}\-{ei}\-{\textlengthmark}\-{\r{g}}\-{r}\-{\textsci}\-{\textsubring{d}}\-{a}\-{h}\-{\oe}\-{v}\-{\textscy}\-{n}\-{\textsubring{d}}\-{\textscy}\-{\textsubring{r}}\-} \dicPos{m}[6] \dicFlx{(‑ar, ‑ar)}[56] \dicDirectTranslationCS{(divadelní) dramatik, dramatička} \dicExampleIS{frægur tékkneskur leikritahöfundur} \dicExampleCS{slavný český dramatik}
\dicEntry[leikrænn] \dicTerm{leik··rænn} \dicIPA{{l}{ei}{\textlengthmark}{\r{g}}{r}{a}{i}{\textsubring{d}}{\textsubring{n}}} \dicPos{adj}[7]\dicFlx{}[-1] \dicDirectTranslationCS{divadelní}
\dicEntry[leikskáld] \dicTerm{leik··skáld} \dicIPA{{l}{ei}{\textlengthmark}{\r{g}}{s}{\r{g}}{au}{l}{\textsubring{d}}} \dicPos{n}[2] \dicFlx{(‑s, ‑)}[5] \dicDirectTranslationCS{dramatik, dramatička, autor(ka) divadelních her}
\dicEntry[leikskóladeild] \dicTerm{leik·skóla··deild} \dicIPA{{l}{ei}{\textlengthmark}{\r{g}}{s}{\r{g}}{ou}{l}{a}{\textsubring{d}}{ei}{l}{\textsubring{d}}} \dicPos{f}[7] \dicFlx{(‑ar, ‑ir)}[1] \dicDirectTranslationCS{oddělení mateřské školky} \dicIndirectTranslationCS{(v~rámci školy ap.)}
\dicEntry[leikskóladvöl] \dicTerm{leik·skóla··|dvöl} \dicIPA{{l}{ei}{\textlengthmark}{\r{g}}{s}{\r{g}}{ou}{l}{a}{\textsubring{d}}{v}{\oe}{\textsubring{l}}} \dicPos{f}[7] \dicFlx{(‑dvalar, ‑dvalir)}[16] \dicDirectTranslationCS{pobyt ve školce}
\dicEntry[leikskólakennari] \dicTerm{leik·skóla··kenn·ar|i} \dicIPA{{l}\-{ei}\-{\textlengthmark}\-{\r{g}}\-{s}\-{\r{g}}\-{ou}\-{l}\-{a}\-{c\smash{\textsuperscript{h}}}\-{\textepsilon}\-{n}\-{a}\-{r}\-{\textsci}\-} \dicPos{m}[1] \dicFlx{(‑a, ‑ar)}[13] \dicDirectTranslationCS{učitel(ka) v~mateřské školce}
\dicEntry[leikskóli] \dicTerm{leik··skól|i} \dicIPA{{l}{ei}{\textlengthmark}{\r{g}}{s}{\r{g}}{ou}{l}{\textsci}} \dicPos{m}[1] \dicFlx{(‑a, ‑ar)}[1] \dicDirectTranslationCS{(mateřská) školka}
\dicEntry[leikskrá] \dicTerm{leik··skrá} \dicIPA{{l}{ei}{\textlengthmark}{\r{g}}{s}{\r{g}}{r}{au}} \dicPos{f}[4] \dicFlx{(‑r\,/\addthin ‑ar, ‑r)}[21] \dicDirectTranslationCS{divadelní program}
\dicEntry[leikslok] \dicTerm{leiks··lok} \dicIPA{{l}{ei}{\textlengthmark}{\r{g}}{s}{l}{\textopeno}{\r{g}}} \dicPos{n}[2] \dicFlx{pl}[1] \textbf{1.} \dicSynonym{niðurstaða} \dicDirectTranslationCS{závěr, rozhodnutí}  \textbf{2.} \dicDirectTranslationCS{výsledek (zápasu, utkání ap.)}
\dicEntry[leiksoppur] \dicTerm{leik··sopp|ur} \dicIPA{{l}{ei}{\textlengthmark}{\r{g}}{s}{\textopeno}{h}{\textsubring{b}}{\textscy}{\textsubring{r}}} \dicPos{m}[6] \dicFlx{(‑s, ‑ar)}[22] \dicDirectTranslationCS{loutka} \dicIndirectTranslationCS{(člověk vnitřně prázdný)}
\dicEntry[leikstjóri] \dicTerm{leik··stjór|i} \dicIPA{{l}{ei}{\textlengthmark}{\r{g}}{s}{\textsubring{d}}{j}{ou}{r}{\textsci}} \dicPos{m}[1] \dicFlx{(‑a, ‑ar)}[1] \textbf{1.} \dicDirectTranslationCS{režisér(ka)}  \textbf{2.} \dicFieldCat{sport.} \dicDirectTranslationCS{rozhodčí}
\dicEntry[leikstjórn] \dicTerm{leik··stjórn} \dicIPA{{l}{ei}{\textlengthmark}{\r{g}}{s}{\textsubring{d}}{j}{ou}{r}{\textsubring{d}}{\textsubring{n}}} \dicPos{f}[7] \dicFlx{(‑ar)}[3] \dicDirectTranslationCS{režie, režírování}
\dicEntry[leikstýra] \dicTerm{leik··stýr|a} \dicIPA{{l}{ei}{\textlengthmark}{\r{g}}{s}{\textsubring{d}}{i}{r}{a}} \dicPos{v}[2] \dicFlx{(‑ði, ‑t)}[105] \dicFlx{dat} \dicDirectTranslationCS{režírovat}
\dicEntry[leiksvið] \dicTerm{leik··svið} \dicIPA{{l}{ei}{\textlengthmark}{\r{g}}{s}{v}{\textsci}{\texttheta}} \dicPos{n}[2] \dicFlx{(‑s, ‑)}[5] \dicDirectTranslationCS{jeviště, scéna}
\dicEntry[leiksýning] \dicTerm{leik··sýn·ing} \dicIPA{{l}{ei}{\textlengthmark}{\r{g}}{s}{i}{n}{i}{\ng}{\r{g}}} \dicPos{f}[4] \dicFlx{(‑ar, ‑ar)}[5] \dicDirectTranslationCS{(divadelní) představení}
\dicEntry[leiktíð] \dicTerm{leik··tíð} \dicIPA{{l}{ei}{\textlengthmark}{\r{g}}{t\smash{\textsuperscript{h}}}{i}{\texttheta}} \dicPos{f}[7] \dicFlx{(‑ar, ‑ir)}[1] \dicFieldCat{sport.} \dicDirectTranslationCS{sezóna (fotbalová ap.)}
\dicEntry[leikur] \dicTerm{leik|ur} \dicsymFrequent\  \dicIPA{{l}{ei}{\textlengthmark}{\r{g}}{\textscy}{\textsubring{r}}} \dicPos{m}[6] \dicFlx{(‑s, ‑ir\,/\addthin ‑ar)}[67] \textbf{1.} \dicFlx{pl (-ir)} \dicSynonym{barnaleikur} \dicDirectTranslationCS{hra, hraní} \dicExampleIS{leikur barna} \dicExampleCS{hra dětí};  \dicPhraseIS{bregða á leik} \dicDirectTranslationCS{hrát si, skotačit}  \textbf{2.} \dicFlx{pl (-ar)} \dicSynonym{viðureign} \dicDirectTranslationCS{boj, střetnutí, utkání}  \textbf{3.} \dicFlx{pl (-ar)} \dicSynonym*{íþróttakeppni} \dicDirectTranslationCS{zápas, utkání, hra} \dicExampleIS{Ólympíuleikar} \dicExampleCS{olympijské hry};  \dicPhraseIS{fara á leikinn} \dicDirectTranslationCS{jít na zápas}  \textbf{4.} \dicSynonym{leikrit} \dicDirectTranslationCS{hra (divadelní ap.), představení}  \textbf{5.} \dicDirectTranslationCS{tah} \dicIndirectTranslationCS{(posunutí šachové figury ap.)};  \dicPhraseIS{eiga leik} \dicDirectTranslationCS{být na tahu};  \dicIdiom{leikur}{ \dicPhraseIS{á nýjan leik}} \dicFlx{adv} \dicDirectTranslationCS{znovu, opět}; { \dicPhraseIS{e‑að er leikur einn}} \dicLangCat{přen.} \dicDirectTranslationCS{(co) je hračka}; { \dicPhraseIS{með leik}} \dicFlx{adv} \dicDirectTranslationCS{snadno, lehce}; { \dicPhraseIS{vera úr leik}} \dicLangCat{přen.} \dicDirectTranslationCS{být mimo hru}; { \dicPhraseIS{við illan leik}} \dicFlx{adv} \dicDirectTranslationCS{horko těžko}
\dicEntry[leikvangur] \dicTerm{leik··vang|ur} \dicIPA{{l}{ei}{\textlengthmark}{\r{g}}{v}{au}{\ng}{\r{g}}{\textscy}{\textsubring{r}}} \dicPos{m}[6] \dicFlx{(‑s, ‑ar)}[25] \dicSynonym{íþróttavöllur} \dicDirectTranslationCS{stadión}
\dicEntry[leikvöllur] \dicTerm{leik··|völlur} \dicIPA{{l}{ei}{\textlengthmark}{\r{g}}{v}{\oe}{\textsubring{d}}{l}{\textscy}{\textsubring{r}}} \dicPos{m}[11] \dicFlx{(‑vallar, ‑vellir)}[5] \textbf{1.} \dicFieldCat{sport.} \dicSynonym{í\-þrótta\-völlur} \dicDirectTranslationCS{(sportovní) hřiště}  \textbf{2.} \dicSynonym*{barnaleikvöllur} \dicDirectTranslationCS{dětské hřiště}
\dicEntry[leir] \dicTerm{leir} \dicsymFrequent\  \dicIPA{{l}{ei}{\textlengthmark}{\textsubring{r}}} \dicPos{m}[4] \dicFlx{(‑s)}[18] \dicDirectTranslationCS{jíl, hlína} \dicExampleIS{hvít englabörn úr leir} \dicExampleCS{bílí andělíčci z~hlíny}
\dicEntry[leira] \dicTerm{leir|a\smash{\textsuperscript{1}}} \dicIPA{{l}{ei}{\textlengthmark}{r}{a}} \dicPos{f}[1] \dicFlx{(‑u, ‑ur)}[19] \dicDirectTranslationCS{jílovitá půda}
\dicEntry[Leira] \dicTerm{Leir|a\smash{\textsuperscript{2}}} \dicIPA{{l}{ei}{\textlengthmark}{r}{a}} \dicPos{f}[1] \dicFlx{(‑u)}[6] \dicFieldCat{geog.} \dicDirectTranslationCS{Loira} \dicIndirectTranslationCS{(řeka ve Francii)}
\dicEntry[leirburður] \dicTerm{leir··burð|ur} \dicIPA{{l}{ei}{r}{\textsubring{b}}{\textscy}{r}{ð}{\textscy}{\textsubring{r}}} \dicPos{m}[10] \dicFlx{(‑ar)}[7] \dicDirectTranslationCS{špatná poezie}
\dicEntry[leirker] \dicTerm{leir··ker} \dicIPA{{l}{ei}{\textsubring{r}}{c\smash{\textsuperscript{h}}}{\textepsilon}{\textsubring{r}}} \dicPos{n}[2] \dicFlx{(‑s, ‑)}[6] \dicDirectTranslationCS{keramika} \dicIndirectTranslationCS{(především z~archeologických nálezů)}
\dicEntry[leirlist] \dicTerm{leir··list} \dicIPA{{l}{ei}{r}{l}{\textsci}{s}{\textsubring{d}}} \dicPos{f}[7] \dicFlx{(‑ar)}[3] \dicDirectTranslationCS{keramika, keramické umění}
\dicEntry[leirmunur] \dicTerm{leir··mun|ur} \dicIPA{{l}{ei}{r}{m}{\textscy}{n}{\textscy}{\textsubring{r}}} \dicPos{m}[10] \dicFlx{(‑ar, ‑ir)}[11] \dicDirectTranslationCS{keramika, keramický výrobek}
\dicEntry[leit] \dicTerm{leit\smash{\textsuperscript{1}}} \dicsymFrequent\  \dicIPA{{l}{ei}{\textlengthmark}{\textsubring{d}}} \dicPos{f}[7] \dicFlx{(‑ar, ‑ir)}[1] \textbf{1.} \dicDirectTranslationCS{hledání, pátrání, vyhledávání} \dicExampleIS{Leitin er mjög umfangsmikil.} \dicExampleCS{Pátrání je velmi rozsáhlé.}  \textbf{2.} \dicPhraseIS{leitir} \dicFlx{pl} \dicSynonym{smölun} \dicDirectTranslationCS{shánění ovcí};  \dicPhraseIS{fara í leitir} \dicDirectTranslationCS{jet shánět ovce}
\dicEntry[leit] \dicTerm{leit\smash{\textsuperscript{2}}} \dicIPA{{l}{ei}{\textlengthmark}{\textsubring{d}}} \dicPos{v} \dicFlx{ind pf sg 1 pers} \dicLink{líta}
\dicEntry[leita] \dicTerm{leit|a} \dicsymFrequent\  \dicIPA{{l}{ei}{\textlengthmark}{\textsubring{d}}{a}} \dicPos{v}[1] \dicFlx{(‑aði)}[1] \dicFlx{gen} \textbf{1.} \dicDirectTranslationCS{(vy)hledat, (vy)pátrat} \dicIndirectTranslationCS{(snažit se nalézt)};  \dicPhraseIS{leita e‑s} \dicDirectTranslationCS{hledat (co)} \dicExampleIS{leita bátsins} \dicExampleCS{hledat člun};  \dicPhraseIS{leita að e‑u\,/\addthin e‑m} \dicDirectTranslationCS{hledat (co\,/\addthin koho)} \dicExampleIS{leita að lyklunum} \dicExampleCS{hledat klíče}  \textbf{2.} \dicDirectTranslationCS{vyhledat, navštívit (lékaře ap.)};  \dicPhraseIS{leita læknis} \dicDirectTranslationCS{vyhledat doktora};  \dicIdiom{leita}[á]{ \dicPhraseIS{leita á e‑m}} \dicDirectTranslationCS{prohledat (koho), prošacovat (koho)};  \dicIdiom{leita}[eftir]{ \dicPhraseIS{leita eftir e‑u við e‑n}} \dicDirectTranslationCS{požádat (koho) o~(co)};  \dicIdiom{leita}[fyrir]{ \dicPhraseIS{leita fyrir sér um e‑ð}} \dicDirectTranslationCS{shánět se po (čem), shánět (co) (práci ap.)};  \dicIdiom{leita}[til]{ \dicPhraseIS{leita til e‑rs}} \dicDirectTranslationCS{obrátit se na (koho)} \dicExampleIS{leita til hennar með vandamál sín} \dicExampleCS{obrátit se na ni se svými problémy};  \dicIdiom{leita}[uppi]{ \dicPhraseIS{leita e‑n uppi}} \dicDirectTranslationCS{vyhledat (koho), vypátrat (koho)};  \dicIdiom{leitast}[við]{ \dicPhraseIS{leitast við}} \dicFlx{refl} \dicDirectTranslationCS{pokoušet se, snažit se}
\dicEntry[leitandi] \dicTerm{leit··andi} \dicIPA{{l}{ei}{\textlengthmark}{\textsubring{d}}{a}{n}{\textsubring{d}}{\textsci}} \dicPos{adj}[13] \dicFlx{indecl}[1] \dicDirectTranslationCS{hledající, pátrající}
\dicEntry[leitarflokkur] \dicTerm{leitar··flokk|ur} \dicIPA{{l}{ei}{\textlengthmark}{\textsubring{d}}{a}{\textsubring{r}}{f}{l}{\textopeno}{h}{\r{g}}{\textscy}{\textsubring{r}}} \dicPos{m}[6] \dicFlx{(‑s, ‑ar)}[8] \dicDirectTranslationCS{pátrací skupina\,/\addthin četa}
\dicEntry[leitarvél] \dicTerm{leitar··vél} \dicIPA{{l}{ei}{\textlengthmark}{\textsubring{d}}{a}{r}{v}{j}{\textepsilon}{\textsubring{l}}} \dicPos{f}[4] \dicFlx{(‑ar, ‑ar)}[1] \dicFieldCat{poč.} \dicDirectTranslationCS{(internetový) vyhledávač}
\dicEntry[leiti] \dicTerm{leiti} \dicIPA{{l}{ei}{\textlengthmark}{\textsubring{d}}{\textsci}} \dicPos{n}[2] \dicFlx{(‑s, ‑)}[14] \dicSynonym{hæð} \dicDirectTranslationCS{vrch, vyvýšenina, kopec};  \dicPhraseIS{á næsta leiti} \dicFlx{adv} \dicDirectTranslationCS{blízko, nedaleko (časově), za dveřmi\,/\addthin rohem} \dicExampleIS{Hátíðin er á næsta leiti.} \dicExampleCS{Svátky jsou za dveřmi.}
\dicEntry[leitun] \dicTerm{leit|un} \dicIPA{{l}{ei}{\textlengthmark}{\textsubring{d}}{\textscy}{\textsubring{n}}} \dicPos{f}[7] \dicFlx{(‑unar)}[9] \dicDirectTranslationCS{hledání, vyhledávání, pátrání};  \dicPhraseIS{leitun að e‑u} \dicDirectTranslationCS{hledání (čeho)};  \dicPhraseIS{það er leitun á\,/\addthin að e‑u} \dicDirectTranslationCS{je vzácné\,/\addthin těžké nalézt (co)}
\dicEntry[lek] \dicTerm{lek} \dicIPA{{l}{\textepsilon}{\textlengthmark}{\r{g}}} \dicPos{v} \dicFlx{ind praes sg 1 pers} \dicLink{leka}
\dicEntry[leka] \dicTerm{leka} \dicsymFrequent\  \dicIPA{{l}{\textepsilon}{\textlengthmark}{\r{g}}{a}} \dicPos{v}[6] \dicFlx{(lek, lak, lákum, læki, lekið)}[17] \textbf{1.} \dicSynonym*{falla í dropum} \dicDirectTranslationCS{téct, stékat, kapat} \dicExampleIS{Vatnið lekur niður á gólfið.} \dicExampleCS{Voda stéká dolů na podlahu.}  \textbf{2.} \dicSynonym*{vera lekur} \dicDirectTranslationCS{téct, prosakovat} \dicIndirectTranslationCS{(propouštět kapalinu)}  \textbf{3.} \dicLangCat{přen.} \dicDirectTranslationCS{uniknout (na veřejnost), prosáknout (na veřejnost) (zpráva ap.)}
\dicEntry[lekandi] \dicTerm{lek··|andi} \dicIPA{{l}{\textepsilon}{\textlengthmark}{\r{g}}{a}{n}{\textsubring{d}}{\textsci}} \dicPos{m}[2] \dicFlx{(‑anda)}[2] \dicFieldCat{med.} \dicDirectTranslationCS{kapavka}
\dicEntry[leki] \dicTerm{lek|i} \dicIPA{{l}{\textepsilon}{\textlengthmark}{\r{\textObardotlessj}}{\textsci}} \dicPos{m}[1] \dicFlx{(‑a, ‑ar)}[1] \textbf{1.} \dicDirectTranslationCS{kapání, odkapávání} \dicExampleIS{leki úr þaki} \dicExampleCS{kapání ze střechy}  \textbf{2.} \dicDirectTranslationCS{prosakování, zatékání} \dicExampleIS{Það kom leki að bátnum.} \dicExampleCS{Loď začala zatékat.}  \textbf{3.} \dicDirectTranslationCS{únik, unikání (informací ap.)}
\dicEntry[lekið] \dicTerm{lekið} \dicIPA{{l}{\textepsilon}{\textlengthmark}{\r{\textObardotlessj}}{\textsci}{\texttheta}} \dicPos{v} \dicFlx{supin} \dicLink{leka}
\dicEntry[lektor] \dicTerm{lektor} \dicIPA{{l}{\textepsilon}{x}{\textsubring{d}}{\textopeno}{\textsubring{r}}} \dicPos{m}[4] \dicFlx{(‑s, ‑ar)}[14] \dicFieldCat{škol.} \dicDirectTranslationCS{odborný asistent, odborná asistentka}
\dicEntry[lekur] \dicTerm{lekur} \dicIPA{{l}{\textepsilon}{\textlengthmark}{\r{g}}{\textscy}{\textsubring{r}}} \dicPos{adj}[1]\dicFlx{}[-1] \dicDirectTranslationCS{zatékající, prosakující} \dicExampleIS{lekur bátur} \dicExampleCS{zatékající člun}
\dicEntry[lem] \dicTerm{lem} \dicIPA{{l}{\textepsilon}{\textlengthmark}{\textsubring{m}}} \dicPos{v} \dicFlx{ind praes sg 1 pers} \dicLink{lemja}
\dicEntry[lemdi] \dicTerm{lemdi} \dicIPA{{l}{\textepsilon}{m}{\textsubring{d}}{\textsci}} \dicPos{v} \dicFlx{con pf sg 1 pers} \dicLink{lemja}
\dicEntry[lemja] \dicTerm{lemja} \dicsymFrequent\  \dicIPA{{l}{\textepsilon}{m}{j}{a}} \dicPos{v}[4] \dicFlx{(lem, lamdi, lömdum, lemdi, lamið)}[3] \dicFlx{acc} \dicSynonym{slá\smash{\textsuperscript{3}}} \dicDirectTranslationCS{uhodit, udeřit, tlouct, mlátit} \dicExampleIS{lemja e‑n með barefli} \dicExampleCS{uhodit (koho) pálkou}
\dicEntry[lemstra] \dicTerm{lemstr|a} \dicIPA{{l}{\textepsilon}{m}{s}{\textsubring{d}}{r}{a}} \dicPos{v}[1] \dicFlx{(‑aði)}[1] \dicFlx{acc} \dicSynonym{limlesta} \dicDirectTranslationCS{zmrzačit}
\dicEntry[lend] \dicTerm{lend} \dicIPA{{l}{\textepsilon}{n}{\textsubring{d}}} \dicPos{f}[4] \dicFlx{(‑ar, ‑ar)}[1] \textbf{1.} \dicDirectTranslationCS{zadek (koně ap.)}  \textbf{2.} \dicPhraseIS{lendar} \dicFlx{pl} \dicSynonym{mjaðmir} \dicDirectTranslationCS{boky, bedra}
\dicEntry[lenda] \dicTerm{len|da} \dicsymFrequent\  \dicIPA{{l}{\textepsilon}{n}{\textsubring{d}}{a}} \dicPos{v}[2] \dicFlx{(‑ti, ‑t)}[41] \dicFlx{dat} \textbf{1.} \dicSynonym*{koma að landi} \dicDirectTranslationCS{přistát, přistávat, dosednout, dosedat} \dicExampleIS{lenda flugvélinni} \dicExampleCS{přistát s~letadlem}  \textbf{2.} \dicDirectTranslationCS{přistát, dopadnout (na zadek ap.)};  \dicIdiom{lenda}[í]{ \dicPhraseIS{lenda í e‑u}} \dicSynonym*{rata í e‑ð} \dicDirectTranslationCS{spadnout do (čeho), dostat se do (čeho)} \dicExampleIS{lenda í vandræðum} \dicExampleCS{dostat se do potíží};  \dicIdiom{lenda}[saman]{ \dicPhraseIS{e‑jum lendir saman}} \dicFlx{impers} \dicLangCat{přen.} \dicDirectTranslationCS{(kdo) se dostává do křížku} \dicExampleIS{Þeim lenti saman.} \dicExampleCS{Dostali se do křížku.}
\dicEntry[lending] \dicTerm{lend··ing} \dicIPA{{l}{\textepsilon}{n}{\textsubring{d}}{i}{\ng}{\r{g}}} \dicPos{f}[4] \dicFlx{(‑ar, ‑ar)}[5] \dicDirectTranslationCS{přistání, přistávání, dosednutí} \dicExampleIS{mjúk lending} \dicExampleCS{měkké přistání}
\dicEntry[lengd] \dicTerm{lengd} \dicsymFrequent\  \dicIPA{{l}{\textepsilon}{\ng}{\textsubring{d}}} \dicPos{f}[7] \dicFlx{(‑ar, ‑ir)}[1] \textbf{1.} \dicDirectTranslationCS{délka} \dicExampleIS{í fullri lengd} \dicExampleCS{v~plné délce}  \textbf{2.} \dicFieldCat{geog.} \dicDirectTranslationCS{zeměpisná délka};  \dicIdiom{lengd}{ \dicPhraseIS{til lengdar}} \dicFlx{adv} \dicDirectTranslationCS{dlouhodobě, v~dlouhém časovém období}
\dicEntry[lengdarbaugur] \dicTerm{lengdar··baug|ur} \dicIPA{{l}{\textepsilon}{\ng}{\textsubring{d}}{a}{r}{\textsubring{b}}{\oe i}{\textbabygamma}{\textscy}{\textsubring{r}}} \dicPos{m}[6] \dicFlx{(‑s, ‑ar)}[24] \dicFieldCat{geog.} \dicDirectTranslationCS{poledník}
\dicEntry[lengi] \dicTerm{lengi} \dicsymFrequent\  \dicIPA{{l}{ei}{\textltailn}{\r{\textObardotlessj}}{\textsci}} \dicPos{adv} \dicFlx{(comp lengur, sup lengst)} \dicSynonym*{langan tíma} \dicDirectTranslationCS{dlouho, dlouze} \dicExampleIS{lengi vetrar} \dicExampleCS{po většinu zimy};  \dicPhraseIS{lengi vel} \dicFlx{adv} \dicDirectTranslationCS{dlouho, dlouhou dobu}
\dicEntry[lenging] \dicTerm{leng··ing} \dicIPA{{l}{ei}{\textltailn}{\r{\textObardotlessj}}{i}{\ng}{\r{g}}} \dicPos{f}[4] \dicFlx{(‑ar, ‑ar)}[5] \textbf{1.} \dicSynonym*{það að lengja e‑ð} \dicDirectTranslationCS{prodloužení, prodlužování} \dicExampleIS{lenging á jólafríinu} \dicExampleCS{prodloužení vánočních prázdnin}  \textbf{2.} \dicSynonym{viðbót} \dicDirectTranslationCS{vylepšení, vylepšování} \dicExampleIS{lenging á bátnum} \dicExampleCS{vylepšení člunu}
\dicEntry[lengja] \dicTerm{lengj|a\smash{\textsuperscript{1}}} \dicIPA{{l}{ei}{\textltailn}{\r{\textObardotlessj}}{a}} \dicPos{f}[1] \dicFlx{(‑u, ‑ur)}[7] \textbf{1.} \dicSynonym{ræma} \dicDirectTranslationCS{pruh, proužek}  \textbf{2.} \dicDirectTranslationCS{karton (cigaret)}
\dicEntry[lengja] \dicTerm{leng|ja\smash{\textsuperscript{2}}} \dicIPA{{l}{ei}{\textltailn}{\r{\textObardotlessj}}{a}} \dicPos{v}[2] \dicFlx{(‑di, ‑t)}[130] \dicFlx{acc} \textbf{1.} \dicDirectTranslationCS{prodloužit, prodlužovat} \dicExampleIS{lengja vinnutímann} \dicExampleCS{prodloužit pracovní dobu}  \textbf{2.} \dicPhraseIS{e‑ð lengir} \dicFlx{impers} \dicDirectTranslationCS{(co) se prodlužuje} \dicExampleIS{Daginn lengir.} \dicExampleCS{Den se prodlužuje.};  \dicIdiom{lengja}[eftir]{ \dicPhraseIS{e‑n lengir eftir e‑m\,/\addthin e‑u}} \dicFlx{impers} \dicDirectTranslationCS{(kdo) touží po (kom\,/\addthin čem)} \dicExampleIS{Mig lengir eftir henni.} \dicExampleCS{Toužím po ní.};  \dicIdiom{lengjast}{ \dicPhraseIS{lengjast}} \dicFlx{refl} \dicDirectTranslationCS{stávat se delší, prodlužovat se} \dicExampleIS{Dagurinn lengist.} \dicExampleCS{Den se prodlužuje.}
\dicEntry[lengra] \dicTerm{lengra} \dicIPA{{l}{\textepsilon}{\ng}{\r{g}}{r}{a}} \dicPos{adv} \dicFlx{comp} \dicLink{langt}
\dicEntry[lengri] \dicTerm{lengri} \dicIPA{{l}{\textepsilon}{\ng}{\r{g}}{r}{\textsci}} \dicPos{adj} \dicFlx{comp m} \dicLink{langur\smash{\textsuperscript{2}}}
\dicEntry[lengst] \dicTerm{lengst} \dicIPA{{l}{\textepsilon}{\ng}{s}{\textsubring{d}}} \dicPos{adv} \dicFlx{sup} \textbf{1.} \dicLink{langt}  \textbf{2.} \dicLink{lengi}
\dicEntry[lengstur] \dicTerm{lengstur} \dicIPA{{l}{\textepsilon}{\ng}{s}{\textsubring{d}}{\textscy}{\textsubring{r}}} \dicPos{adj} \dicFlx{m sg nom sup} \dicLink{langur\smash{\textsuperscript{2}}}
\dicEntry[lengur] \dicTerm{lengur} \dicIPA{{l}{ei}{\ng}{\r{g}}{\textscy}{\textsubring{r}}} \dicPos{adv} \dicFlx{comp} \dicLink{lengi}
\dicEntry[lens] \dicTerm{lens} \dicIPA{{l}{\textepsilon}{n}{s}} \dicPos{n}[2] \dicFlx{(‑)}[24] \dicFieldCat{nám.} \dicSynonym{byr} \dicDirectTranslationCS{zadní vítr}
\dicEntry[lenska] \dicTerm{lensk|a} \dicIPA{{l}{\textepsilon}{n}{s}{\r{g}}{a}} \dicPos{f}[1] \dicFlx{(‑u)}[5] \dicSynonym*{landsvenja} \dicDirectTranslationCS{zvyk, obyčej, zvyklost}
\dicEntry[lenskur] \dicTerm{lenskur} \dicIPA{{l}{\textepsilon}{n}{s}{\r{g}}{\textscy}{\textsubring{r}}} \dicPos{adj}[1]\dicFlx{}[-1] \dicPhraseIS{hvers lenskur} \dicDirectTranslationCS{jaké národnosti} \dicExampleIS{Ég vissi ekki hvers lenskur hann var.} \dicExampleCS{Nevěděl jsem, jaké je národnosti.}
\dicEntry[lep] \dicTerm{lep} \dicIPA{{l}{\textepsilon}{\textlengthmark}{\textsubring{b}}} \dicPos{v} \dicFlx{ind praes sg 1 pers} \dicLink{lepja}
\dicEntry[lepja] \dicTerm{lepja} \dicIPA{{l}{\textepsilon}{\textlengthmark}{\textsubring{b}}{j}{a}} \dicPos{v}[4] \dicFlx{(lep, lapti, löptum, lepti, lapið)}[47] \dicFlx{acc} \textbf{1.} \dicDirectTranslationCS{(s)lízat, vylízat} \dicIndirectTranslationCS{(o~zvířatech)} \dicExampleIS{lepja úr skálinni} \dicExampleCS{lízat z~misky}  \textbf{2.} \dicDirectTranslationCS{papouškovat, opakovat (bez porozumění ap.)}
\dicEntry[leppa] \dicTerm{lepp|a} \dicIPA{{l}{\textepsilon}{h}{\textsubring{b}}{a}} \dicPos{v}[1] \dicFlx{(‑aði)}[1] \dicFlx{acc} \textbf{1.} \dicPhraseIS{leppa sig} \dicSynonym*{klæða sig} \dicDirectTranslationCS{obléct se, oblékat se}  \textbf{2.} \dicSynonym{bæta} \dicDirectTranslationCS{opravit, vyspravit} \dicExampleIS{leppa upp á e‑ð} \dicExampleCS{vyspravit (co)}  \textbf{3.} \dicDirectTranslationCS{vázat} \dicIndirectTranslationCS{(v~šachách)}
\dicEntry[Leppalúði] \dicTerm{Leppa··lúð|i} \dicIPA{{l}{\textepsilon}{h}{\textsubring{b}}{a}{l}{u}{ð}{\textsci}} \dicPos{m}[1] \dicFlx{(‑a)}[5] \dicFlx{prop} \dicFieldCat{pov.} \dicDirectTranslationCS{Leppalúði} \dicIndirectTranslationCS{(nesmírně líný obr, otec 13 vánočních skřítků)}
\dicEntry[leppríki] \dicTerm{lepp··ríki} \dicIPA{{l}{\textepsilon}{h}{\textsubring{b}}{r}{i}{\r{\textObardotlessj}}{\textsci}} \dicPos{n}[2] \dicFlx{(‑s, ‑)}[16] \dicFieldCat{pol.} \dicDirectTranslationCS{loutkový\,/\addthin satelitní stát}
\dicEntry[leppstjórn] \dicTerm{lepp··stjórn} \dicIPA{{l}{\textepsilon}{h}{\textsubring{b}}{s}{\textsubring{d}}{j}{ou}{r}{\textsubring{d}}{\textsubring{n}}} \dicPos{f}[7] \dicFlx{(‑ar, ‑ir)}[1] \dicDirectTranslationCS{loutková vláda}
\dicEntry[leppur] \dicTerm{lepp|ur} \dicIPA{{l}{\textepsilon}{h}{\textsubring{b}}{\textscy}{\textsubring{r}}} \dicPos{m}[6] \dicFlx{(‑s, ‑ar)}[22] \textbf{1.} \dicSynonym{bót} \dicDirectTranslationCS{záplata}  \textbf{2.} \dicDirectTranslationCS{vložka do bot}  \textbf{3.} \dicPhraseIS{leppar} \dicFlx{pl} \dicSynonym*{tötrar} \dicDirectTranslationCS{hadry} \dicIndirectTranslationCS{(o~šatech)}  \textbf{4.} \dicSynonym{handbendi} \dicDirectTranslationCS{loutka, nohsled}  \textbf{5.} \dicDirectTranslationCS{klapka (na oku ap.)}  \textbf{6.} \dicDirectTranslationCS{vazba} \dicIndirectTranslationCS{(v~šachách)}
\dicEntry[lepti] \dicTerm{lepti} \dicIPA{{l}{\textepsilon}{f}{\textsubring{d}}{\textsci}} \dicPos{v} \dicFlx{con pf sg 1 pers} \dicLink{lepja}
\dicEntry[lerki] \dicTerm{lerki} \dicIPA{{l}{\textepsilon}{\textsubring{r}}{\r{\textObardotlessj}}{\textsci}} \dicPos{n}[2] \dicFlx{(‑s)}[20] \dicFieldCat{bot.} \dicDirectTranslationCS{modřín} \textit{(l.~{\textLA{Larix}})}  \dicsymPhoto\ 
\dicFigure{286.jpg}{Lerki}{Lerki - Zicha Ondřej, Biolib, Copyright/CC-BY-NC}
\dicEntry[les] \dicTerm{les} \dicIPA{{l}{\textepsilon}{\textlengthmark}{s}} \dicPos{v} \dicFlx{ind praes sg 1 pers} \dicLink{lesa}
\dicEntry[lesa] \dicTerm{lesa} \dicsymFrequent\  \dicIPA{{l}{\textepsilon}{\textlengthmark}{s}{a}} \dicPos{v}[6] \dicFlx{(les, las, lásum, læsi, lesið)}[16] \dicFlx{acc} \textbf{1.} \dicDirectTranslationCS{(pře)číst} \dicExampleIS{lesa söguna upphátt} \dicExampleCS{číst příběh nahlas}  \textbf{2.} \dicSynonym{læra} \dicDirectTranslationCS{studovat, učit se} \dicExampleIS{Hann les sagnfræði í háskóla.} \dicExampleCS{Studuje historii na vysoké škole.}  \textbf{3.} \dicLangCat{zast.} \dicDirectTranslationCS{sbírat (květy ap.)} \dicExampleIS{lesa aldin af trjánum} \dicExampleCS{sbírat plody ze stromů};  \dicIdiom{lesa}[fyrir]{ \dicPhraseIS{lesa fyrir}} \dicDirectTranslationCS{diktovat}; { \dicPhraseIS{lesa fyrir e‑m e‑ð}} \dicDirectTranslationCS{(na)diktovat (komu co)} \dicExampleIS{lesa henni bréfið fyrir} \dicExampleCS{diktovat jí dopis};  \dicIdiom{lesa}[í]{ \dicPhraseIS{lesa í lófa}} \dicLangCat{přen.} \dicDirectTranslationCS{číst z~dlaně};  \dicIdiom{lesa}[saman]{ \dicPhraseIS{lesa saman e‑ð og e‑ð}} \dicDirectTranslationCS{srovnávat\,/\addthin porovnávat (co) s~(čím)} \dicExampleIS{lesa frumritið og afritið saman} \dicExampleCS{srovnávat původní text s~opisem};  \dicIdiom{lesa}[til]{ \dicPhraseIS{lesa sér til e‑ð}} \dicDirectTranslationCS{seznámit se s~(čím), dozvědět se (co) nového (čtením)};  \dicIdiom{lesa}[undir]{ \dicPhraseIS{lesa undir próf}} \dicDirectTranslationCS{učit se na zkoušku, studovat na zkoušku} \dicExampleIS{lesa undir skóla} \dicExampleCS{učit se do školy};  \dicIdiom{lesa}[upp]{ \dicPhraseIS{lesa upp e‑ð}} \dicDirectTranslationCS{recitovat (co), přednášet (co)} \dicExampleIS{lesa upp ljóð} \dicExampleCS{recitovat báseň}; { \dicPhraseIS{lesa sig upp}} \dicSynonym{klifra} \dicDirectTranslationCS{šplhat, vyšplhat (se)};  \dicIdiom{lesa}[yfir]{ \dicPhraseIS{lesa e‑ð yfir}} \dicDirectTranslationCS{pročíst si (co)}
\dicEntry[lesandi] \dicTerm{les··|andi} \dicsymFrequent\  \dicIPA{{l}{\textepsilon}{\textlengthmark}{s}{a}{n}{\textsubring{d}}{\textsci}} \dicPos{m}[2] \dicFlx{(‑anda, ‑endur)}[1] \dicDirectTranslationCS{čtenář(ka)} \dicExampleIS{Lesendurnir eru ánægðir með bókina.} \dicExampleCS{Čtenářům se knížka líbí.}
\dicEntry[lesbía] \dicTerm{lesbí|a} \dicIPA{{l}{\textepsilon}{s}{\textsubring{b}}{i}{j}{a}} \dicPos{f}[1] \dicFlx{(‑u, ‑ur)}[7] \dicDirectTranslationCS{lesba, lesbička}
\dicEntry[lesblinda] \dicTerm{les··blind|a} \dicIPA{{l}{\textepsilon}{\textlengthmark}{s}{\textsubring{b}}{l}{\textsci}{n}{\textsubring{d}}{a}} \dicPos{f}[1] \dicFlx{(‑u)}[5] \dicFieldCat{med.} \dicDirectTranslationCS{dyslexie}
\dicEntry[lesbók] \dicTerm{les··|bók} \dicIPA{{l}{\textepsilon}{\textlengthmark}{s}{\textsubring{b}}{ou}{\r{g}}} \dicPos{f}[8] \dicFlx{(‑bókar, ‑bækur)}[5] \dicDirectTranslationCS{čítanka}
\dicEntry[lesefni] \dicTerm{les··efni} \dicIPA{{l}{\textepsilon}{\textlengthmark}{s}{\textepsilon}{\textsubring{b}}{n}{\textsci}} \dicPos{n}[2] \dicFlx{(‑s, ‑)}[14] \dicDirectTranslationCS{čtení, četba, čtivo}
\dicEntry[lesgleraugu] \dicTerm{les··gler·augu} \dicIPA{{l}{\textepsilon}{\textlengthmark}{s}{\r{g}}{l}{\textepsilon}{r}{\oe i}{\textbabygamma}{\textscy}} \dicPos{n}[1] \dicFlx{pl}[3] \dicDirectTranslationCS{brýle na čtení}
\dicEntry[lesið] \dicTerm{lesið} \dicIPA{{l}{\textepsilon}{\textlengthmark}{s}{\textsci}{\texttheta}} \dicPos{v} \dicFlx{supin} \dicLink{lesa}
\dicEntry[lesinn] \dicTerm{lesinn} \dicIPA{{l}{\textepsilon}{\textlengthmark}{s}{\textsci}{\textsubring{n}}} \dicPos{adj}[6]\dicFlx{}[-2] \dicDirectTranslationCS{sečtělý}
\dicEntry[lesminni] \dicTerm{les··minni} \dicIPA{{l}{\textepsilon}{\textlengthmark}{s}{m}{\textsci}{n}{\textsci}} \dicPos{n}[2] \dicFlx{(‑s, ‑)}[14] \dicFieldCat{poč.} \dicDirectTranslationCS{ROM} \dicIndirectTranslationCS{(typ elektronické paměti používaný pro uložení firmware v~elektronických přístrojích)} \dicAntonym{vinnsluminni}
\dicEntry[lesning] \dicTerm{les··ning} \dicIPA{{l}{\textepsilon}{s}{\textsubring{d}}{n}{i}{\ng}{\r{g}}} \dicPos{f}[4] \dicFlx{(‑ar, ‑ar)}[5] \dicDirectTranslationCS{četba, čtení}
\dicEntry[Lesótó] \dicTerm{Lesótó} \dicIPA{{l}{\textepsilon}{\textlengthmark}{s}{ou}{\textsubring{d}}{ou}} \dicPos{n}[4] \dicFlx{indecl}[2] \dicFieldCat{geog.} \dicDirectTranslationCS{Lesotho}
\dicEntry[Lesótói] \dicTerm{Lesótó|i} \dicIPA{{l}{\textepsilon}{\textlengthmark}{s}{ou}{\textsubring{d}}{ou}{\textsci}} \dicPos{m}[1] \dicFlx{(‑a, ‑ar)}[1] \dicLink{Lesótómaður}
\dicEntry[Lesótómaður] \dicTerm{Lesótó··|maður}\dicTerm{, Lesótói} \dicIPA{{l}\-{\textepsilon}\-{\textlengthmark}\-{s}\-{ou}\-{\textsubring{d}}\-{ou}\-{m}\-{a}\-{ð}\-{\textscy}\-{\textsubring{r}}\-} \dicPos{m}[13] \dicFlx{(‑manns, ‑menn)}[2] \dicDirectTranslationCS{Lesothan(ka)}
\dicEntry[lesótóskur] \dicTerm{lesótóskur} \dicIPA{{l}{\textepsilon}{\textlengthmark}{s}{ou}{\textsubring{d}}{ou}{s}{\r{g}}{\textscy}{\textsubring{r}}} \dicPos{adj}[1]\dicFlx{}[-6] \dicDirectTranslationCS{lesothský}
\dicEntry[lesskilningur] \dicTerm{les··skiln·ing|ur} \dicIPA{{l}{\textepsilon}{\textlengthmark}{s}{\r{\textObardotlessj}}{\textsci}{l}{n}{i}{\ng}{\r{g}}{\textscy}{\textsubring{r}}} \dicPos{m}[6] \dicFlx{(‑s)}[9] \dicDirectTranslationCS{porozumění (psanému) textu}
\dicEntry[lest] \dicTerm{lest} \dicsymFrequent\  \dicIPA{{l}{\textepsilon}{s}{\textsubring{d}}} \dicPos{f}[7] \dicFlx{(‑ar, ‑ir)}[1] \textbf{1.} \dicDirectTranslationCS{vlak, vláček} \dicExampleIS{fara með lest} \dicExampleCS{jet vlakem};  \dicPhraseIS{missa af lestinni} {\textbf{a.}} \dicDirectTranslationCS{zmeškat vlak};  {\textbf{b.}} \dicLangCat{přen.} \dicDirectTranslationCS{nechat si ujet vlak} \dicIndirectTranslationCS{(promeškat příležitost)}  \textbf{2.} \dicFieldCat{nám.} \dicSynonym{farmur} \dicDirectTranslationCS{náklad}  \textbf{3.} \dicSynonym{tonn} \dicDirectTranslationCS{(lodní) tuna}  \textbf{4.} \dicFieldCat{nám.} \dicSynonym*{lestarrúm} \dicDirectTranslationCS{nákladní prostor}  \textbf{5.} \dicSynonym{röð} \dicDirectTranslationCS{řada, karavana} \dicExampleIS{löng lest af bílum} \dicExampleCS{dlouhá řada aut}
\dicEntry[lesta] \dicTerm{lest|a} \dicIPA{{l}{\textepsilon}{s}{\textsubring{d}}{a}} \dicPos{v}[1] \dicFlx{(‑aði)}[1] \dicSynonym{hlaða\smash{\textsuperscript{2}}} \dicDirectTranslationCS{naložit, nakládat} \dicExampleIS{lesta skipið} \dicExampleCS{nakládat loď}
\dicEntry[lestarstjóri] \dicTerm{lestar··stjór|i} \dicIPA{{l}{\textepsilon}{s}{\textsubring{d}}{a}{\textsubring{r}}{s}{\textsubring{d}}{j}{ou}{r}{\textsci}} \dicPos{m}[1] \dicFlx{(‑a, ‑ar)}[1] \dicDirectTranslationCS{strojvedoucí, strojvůdce, strojvůdkyně}
\dicEntry[lestir] \dicTerm{lestir} \dicIPA{{l}{\textepsilon}{s}{\textsubring{d}}{\textsci}{\textsubring{r}}} \dicPos{m} \dicFlx{pl nom} \dicLink{löstur}
\dicEntry[lestrarbók] \dicTerm{lestrar··|bók} \dicIPA{{l}{\textepsilon}{s}{\textsubring{d}}{r}{a}{r}{\textsubring{b}}{ou}{\r{g}}} \dicPos{f}[8] \dicFlx{(‑bókar, ‑bækur)}[5] \dicDirectTranslationCS{čítanka}
\dicEntry[lestrarsalur] \dicTerm{lestrar··sal|ur} \dicIPA{{l}{\textepsilon}{s}{\textsubring{d}}{r}{a}{\textsubring{r}}{s}{a}{l}{\textscy}{\textsubring{r}}} \dicPos{m}[10] \dicFlx{(‑ar, ‑ir)}[14] \dicDirectTranslationCS{čítárna}
\dicEntry[lestur] \dicTerm{lest|ur} \dicsymFrequent\  \dicIPA{{l}{\textepsilon}{s}{\textsubring{d}}{\textscy}{\textsubring{r}}} \dicPos{m}[5] \dicFlx{(‑rar\,/\addthin ‑urs, ‑rar)}[9] \textbf{1.} \dicSynonym*{það að lesa} \dicDirectTranslationCS{čtení, četba} \dicExampleIS{skemmtilegur lestur} \dicExampleCS{zábavné čtení}  \textbf{2.} \dicSynonym{nám} \dicDirectTranslationCS{studium, učení se}  \textbf{3.} \dicSynonym{klifur} \dicDirectTranslationCS{šplhání}
\dicEntry[let] \dicTerm{let} \dicIPA{{l}{\textepsilon}{\textlengthmark}{\textsubring{d}}} \dicPos{v} \dicFlx{ind praes sg 1 pers} \dicLink{letja}
\dicEntry[leti] \dicTerm{let|i} \dicIPA{{l}{\textepsilon}{\textlengthmark}{\textsubring{d}}{\textsci}} \dicPos{f}[3] \dicFlx{(‑i)}[3] \dicDirectTranslationCS{lenost, lenivost, zahálčivost}
\dicEntry[letidýr] \dicTerm{leti··dýr} \dicIPA{{l}{\textepsilon}{\textlengthmark}{\textsubring{d}}{\textsci}{\textsubring{d}}{i}{\textsubring{r}}} \dicPos{n}[2] \dicFlx{(‑s, ‑)}[5] \dicFieldCat{zool.} \dicDirectTranslationCS{lenochod} \textit{(l.~{\textLA{Bradypus}})}  \dicsymPhoto\ 
\dicFigure{ds_image_letidyr_0_1.jpg}{Letidýr}{Letidýr - Stefan Laube, PD}
\dicEntry[letihaugur] \dicTerm{leti··haug|ur} \dicIPA{{l}{\textepsilon}{\textlengthmark}{\textsubring{d}}{\textsci}{h}{\oe i}{\textbabygamma}{\textscy}{\textsubring{r}}} \dicPos{m}[6] \dicFlx{(‑s, ‑ar)}[22] \dicDirectTranslationCS{lenoch, lenoška, povaleč(ka)}
\dicEntry[letilegur] \dicTerm{leti··legur} \dicIPA{{l}{\textepsilon}{\textlengthmark}{\textsubring{d}}{\textsci}{l}{\textepsilon}{\textbabygamma}{\textscy}{\textsubring{r}}} \dicPos{adj}[1]\dicFlx{}[-8] \dicDirectTranslationCS{líný, lenivý}
\dicEntry[letingi] \dicTerm{let··ing|i} \dicIPA{{l}{\textepsilon}{\textlengthmark}{\textsubring{d}}{i}{\textltailn}{\r{\textObardotlessj}}{\textsci}} \dicPos{m}[1] \dicFlx{(‑ja, ‑jar)}[14] \dicDirectTranslationCS{lenoch, lenoška, povaleč(ka)}
\dicEntry[letja] \dicTerm{letja} \dicIPA{{l}{\textepsilon}{\textlengthmark}{\textsubring{d}}{j}{a}} \dicPos{v}[4] \dicFlx{(let, latti, löttum, letti, latt)}[48] \dicFlx{acc} \dicSynonym*{telja úr} \dicDirectTranslationCS{odradit, odrazovat};  \dicPhraseIS{letja e‑n e‑s} \dicDirectTranslationCS{odrazovat (koho) od (čeho)} \dicExampleIS{letja e‑n fararinnar} \dicExampleCS{odrazovat (koho) od cesty}
\dicEntry[letra] \dicTerm{letr|a} \dicIPA{{l}{\textepsilon}{\textlengthmark}{\textsubring{d}}{r}{a}} \dicPos{v}[1] \dicFlx{(‑aði)}[1] \dicFlx{acc} \textbf{1.} \dicSynonym{skrifa} \dicDirectTranslationCS{psát, napsat}  \textbf{2.} \dicSynonym*{höggva stafi í stein} \dicDirectTranslationCS{vysázet\,/\addthin vyrýt písmem}
\dicEntry[Letti] \dicTerm{Lett|i\smash{\textsuperscript{1}}} \dicIPA{{l}{\textepsilon}{h}{\textsubring{d}}{\textsci}} \dicPos{m}[1] \dicFlx{(‑a, ‑ar)}[1] \dicDirectTranslationCS{Lotyš(ka)}
\dicEntry[letti] \dicTerm{letti\smash{\textsuperscript{2}}} \dicIPA{{l}{\textepsilon}{h}{\textsubring{d}}{\textsci}} \dicPos{v} \dicFlx{con pf sg 1 pers} \dicLink{letja}
\dicEntry[Lettland] \dicTerm{Lett··land} \dicIPA{{l}{\textepsilon}{h}{\textsubring{d}}{l}{a}{n}{\textsubring{d}}} \dicPos{n}[2] \dicFlx{(‑s)}[4] \dicFieldCat{geog.} \dicDirectTranslationCS{Lotyšsko}
\dicEntry[lettneska] \dicTerm{lett··nesk|a} \dicIPA{{l}{\textepsilon}{h}{\textsubring{d}}{n}{\textepsilon}{s}{\r{g}}{a}} \dicPos{f}[1] \dicFlx{(‑u)}[5] \dicDirectTranslationCS{lotyština}
\dicEntry[lettneskur] \dicTerm{lett··neskur} \dicIPA{{l}{\textepsilon}{h}{\textsubring{d}}{n}{\textepsilon}{s}{\r{g}}{\textscy}{\textsubring{r}}} \dicPos{adj}[1]\dicFlx{}[-6] \dicDirectTranslationCS{lotyšský}
\dicEntry[letur] \dicTerm{letur} \dicIPA{{l}{\textepsilon}{\textlengthmark}{\textsubring{d}}{\textscy}{\textsubring{r}}} \dicPos{n}[2] \dicFlx{(‑s, ‑)}[25] \textbf{1.} \dicSynonym*{stafategund} \dicDirectTranslationCS{písmo, druh\,/\addthin typ písma, font}  \textbf{2.} \dicSynonym{áletrun} \dicDirectTranslationCS{nápis, epigraf} \dicExampleIS{letur á steini} \dicExampleCS{nápis na kameni}
\dicEntry[leturbreyting] \dicTerm{letur··breyt·ing} \dicIPA{{l}{\textepsilon}{\textlengthmark}{\textsubring{d}}{\textscy}{r}{\textsubring{b}}{r}{ei}{\textsubring{d}}{i}{\ng}{\r{g}}} \dicPos{f}[4] \dicFlx{(‑ar, ‑ar)}[5] \dicDirectTranslationCS{změna písma} \dicIndirectTranslationCS{(z~tučného písma na normální ap.)}
\dicEntry[leturgerð] \dicTerm{letur··gerð} \dicIPA{{l}{\textepsilon}{\textlengthmark}{\textsubring{d}}{\textscy}{r}{\r{\textObardotlessj}}{\textepsilon}{r}{\texttheta}} \dicPos{f}[7] \dicFlx{(‑ar, ‑ir)}[1] \dicSynonym*{leturtegund} \dicDirectTranslationCS{font, znaková sada}
\dicEntry[leturgrafari] \dicTerm{letur··graf·ar|i} \dicIPA{{l}{\textepsilon}{\textlengthmark}{\textsubring{d}}{\textscy}{r}{\r{g}}{r}{a}{v}{a}{r}{\textsci}} \dicPos{m}[1] \dicFlx{(‑a, ‑ar)}[10] \dicDirectTranslationCS{rytec\,/\addthin rytečka (písma)}
\dicEntry[leturstærð] \dicTerm{letur··stærð} \dicIPA{{l}{\textepsilon}{\textlengthmark}{\textsubring{d}}{\textscy}{\textsubring{r}}{s}{\textsubring{d}}{a}{i}{r}{\texttheta}} \dicPos{f}[7] \dicFlx{(‑ar, ‑ir)}[1] \dicDirectTranslationCS{velikost písma}
\dicEntry[lexía] \dicTerm{lexí|a} \dicIPA{{l}{\textepsilon}{x}{s}{i}{j}{a}} \dicPos{f}[1] \dicFlx{(‑u, ‑ur)}[7] \textbf{1.} \dicSynonym{námsefni} \dicDirectTranslationCS{lekce} \dicIndirectTranslationCS{(úsek učební látky)};  \dicPhraseIS{lesa lexíuna sína} \dicDirectTranslationCS{naučit se, udělat domácí úkoly}  \textbf{2.} \dicSynonym{áminning} \dicDirectTranslationCS{lekce, (nepříjemné) poučení}
\dicEntry[leyfa] \dicTerm{leyf|a} \dicsymFrequent\  \dicIPA{{l}{ei}{\textlengthmark}{v}{a}} \dicPos{v}[2] \dicFlx{(‑ði, ‑t)}[106] \dicFlx{dat + acc} \dicDirectTranslationCS{dovolit, dovolovat, povolit, povolovat, nechat, nechávat};  \dicPhraseIS{leyfa e‑m e‑ð} \dicDirectTranslationCS{dovolit (co komu)} \dicExampleIS{leyfa henni að halda á hestinum} \dicExampleCS{dovolit jí vést koně};  \dicIdiom{leyfast}{ \dicPhraseIS{e‑m leyfist að (gera e‑ð)}} \dicFlx{refl impers} \dicDirectTranslationCS{(komu) je dovoleno (udělat (co))} \dicExampleIS{ef mér leyfist að spyrja} \dicExampleCS{pokud se můžu zeptat}
\dicEntry[leyfður] \dicTerm{leyfður} \dicIPA{{l}{ei}{v}{ð}{\textscy}{\textsubring{r}}} \dicPos{adj}[2]\dicFlx{}[-1] \dicDirectTranslationCS{dovolený, povolený} \dicExampleIS{leyfður hraði} \dicExampleCS{povolená rychlost}
\dicEntry[leyfi] \dicTerm{leyfi} \dicsymFrequent\  \dicIPA{{l}{ei}{\textlengthmark}{v}{\textsci}} \dicPos{n}[2] \dicFlx{(‑s, ‑)}[14] \textbf{1.} \dicSynonym{heimild} \dicDirectTranslationCS{licence, oprávnění}  \textbf{2.} \dicDirectTranslationCS{povolení, dovolení, souhlas, svolení} \dicExampleIS{án leyfis} \dicExampleCS{bez povolení};  \dicPhraseIS{með leyfi} \dicFlx{adv} \dicDirectTranslationCS{s~dovolením}  \textbf{3.} \dicSynonym{frí} \dicDirectTranslationCS{dovolená, volno}
\dicEntry[leyfilegur] \dicTerm{leyfi··legur} \dicIPA{{l}{ei}{\textlengthmark}{v}{\textsci}{l}{\textepsilon}{\textbabygamma}{\textscy}{\textsubring{r}}} \dicPos{adj}[1]\dicFlx{}[-8] \dicSynonym{heimill} \dicDirectTranslationCS{povolený, dovolený, přípustný}
\dicEntry[leyfisleysi] \dicTerm{leyfis··leysi} \dicIPA{{l}{ei}{\textlengthmark}{v}{\textsci}{s}{l}{ei}{s}{\textsci}} \dicPos{n}[2] \dicFlx{(‑s)}[20] \dicPhraseIS{gera e‑ð í leyfisleysi} \dicDirectTranslationCS{udělat (co) bez povolení\,/\addthin dovolení\,/\addthin svolení}
\dicEntry[leyfisveiting] \dicTerm{leyfis··veit·ing} \dicIPA{{l}{ei}{\textlengthmark}{v}{\textsci}{s}{v}{ei}{\textsubring{d}}{i}{\ng}{\r{g}}} \dicPos{f}[4] \dicFlx{(‑ar, ‑ar)}[5] \dicDirectTranslationCS{udělení\,/\addthin vydání oprávnění\,/\addthin povolení}
\dicEntry[leyna] \dicTerm{leyn|a} \dicsymFrequent\  \dicIPA{{l}{ei}{\textlengthmark}{n}{a}} \dicPos{v}[2] \dicFlx{(‑di, ‑t)}[141] \dicFlx{dat} \dicDirectTranslationCS{ukrýt, ukrývat, skrýt, skrývat, (u)tajit, zatajovat};  \dicPhraseIS{leyna e‑n e‑u, leyna e‑u fyrir e‑m} \dicDirectTranslationCS{skrýt (co) před (kým), zatajit (co) před (kým), zatajit (komu co)} \dicExampleIS{leyna hana vitneskju um málið} \dicExampleCS{zatajit jí informace o~té záležitosti};  \dicIdiom{leyna}[á]{ \dicPhraseIS{leyna á sér}} \dicDirectTranslationCS{nezdát se} \dicIndirectTranslationCS{(neukazovat pravou skutečnost)};  \dicIdiom{leynast}{ \dicPhraseIS{leynast}} \dicFlx{refl} \dicDirectTranslationCS{skrýt se, skrývat se}
\dicEntry[leynd] \dicTerm{leynd} \dicIPA{{l}{ei}{n}{\textsubring{d}}} \dicPos{f}[7] \dicFlx{(‑ar, ‑ir)}[1] \textbf{1.} \dicSynonym{launung} \dicDirectTranslationCS{tajnost};  \dicPhraseIS{með leynd} \dicFlx{adv} \dicDirectTranslationCS{v~tajnosti, potají}  \textbf{2.} \dicSynonym*{leynistaður} \dicDirectTranslationCS{(místo) utajení}
\begin{xtolerant}{}{1pt}
\dicEntry[leyndardómsfullur] \dicTerm{leyndar·dóms··|fullur} \dicsymFrequent\addthinS\textls[-10]{\dicIPA{{l}\-{ei}\-{n}\-{\textsubring{d}}\-{a}\-{r}\-{\textsubring{d}}\-{ou}\-{m}\-{s}\-{f}\-{\textscy}\-{\textsubring{d}}\-{l}\-{\textscy}\-{\textsubring{r}}}} \dicPos{adj}[10]\addthinS\dicFlx{(comp ‑fyllri, sup ‑fyllstur)}[7] \dicDirectTranslationCS{tajuplný, tajemný, záhadný} \dicExampleIS{leyndardómsfullur á svip} \dicExampleCS{s~tajuplným výrazem ve tváři}
\end{xtolerant}
\dicEntry[leyndardómur] \dicTerm{leyndar··dóm|ur} \dicIPA{{l}{ei}{n}{\textsubring{d}}{a}{r}{\textsubring{d}}{ou}{m}{\textscy}{\textsubring{r}}} \dicPos{m}[6] \dicFlx{(‑s, ‑ar)}[10] \dicDirectTranslationCS{tajemství, tajemnost, záhada} \dicExampleIS{leyndardómur lífsins} \dicExampleCS{tajemství života}
\dicEntry[leyndarmál] \dicTerm{leyndar··mál} \dicsymFrequent\  \dicIPA{{l}{ei}{n}{\textsubring{d}}{a}{r}{m}{au}{\textsubring{l}}} \dicPos{n}[2] \dicFlx{(‑s, ‑)}[5] \dicDirectTranslationCS{tajemství} \dicExampleIS{opinbert leyndarmál} \dicExampleCS{veřejné tajemství}
\dicEntry[leyndur] \dicTerm{leyndur} \dicsymFrequent\  \dicIPA{{l}{ei}{n}{\textsubring{d}}{\textscy}{\textsubring{r}}} \dicPos{adj}[2]\dicFlx{}[-14] \textbf{1.} \dicDirectTranslationCS{skrytý, schovaný} \dicExampleIS{leyndir hæfileikar} \dicExampleCS{skryté schopnosti}  \textbf{2.} \dicDirectTranslationCS{tajný, utajený};  \dicPhraseIS{fara leynt með e‑ð} \dicDirectTranslationCS{tajit (co), držet (co) v~tajnosti};  \dicPhraseIS{halda e‑ð leyndum} \dicDirectTranslationCS{tajit (co), držet (co) v~tajnosti}
\dicEntry[leyni] \dicTerm{leyni} \dicIPA{{l}{ei}{\textlengthmark}{n}{\textsci}} \dicPos{n}[2] \dicFlx{(‑s, ‑)}[14] \dicDirectTranslationCS{skrýš, úkryt};  \dicPhraseIS{í leyni} \dicFlx{adv} \dicDirectTranslationCS{tajně, potají};  \dicPhraseIS{liggja í leyni} \dicDirectTranslationCS{skrývat se, ukrývat se}
\dicEntry[leynifélag] \dicTerm{leyni··fé·|lag} \dicIPA{{l}{ei}{\textlengthmark}{n}{\textsci}{f}{j}{\textepsilon}{l}{a}{x}} \dicPos{n}[2] \dicFlx{(‑lags, ‑lög)}[8] \dicDirectTranslationCS{tajné společenství, tajná společnost}
\dicEntry[leynilegur] \dicTerm{leyni··legur} \dicIPA{{l}{ei}{\textlengthmark}{n}{\textsci}{l}{\textepsilon}{\textbabygamma}{\textscy}{\textsubring{r}}} \dicPos{adj}[1]\dicFlx{}[-8] \dicSynonym{dulinn} \dicDirectTranslationCS{tajný, utajený}
\dicEntry[leynilögregla] \dicTerm{leyni··lög·regl|a} \dicIPA{{l}{ei}{\textlengthmark}{n}{\textsci}{l}{\oe}{\textbabygamma}{r}{\textepsilon}{\r{g}}{l}{a}} \dicPos{f}[1] \dicFlx{(‑u)}[5] \dicDirectTranslationCS{tajná policie}

\dicEntry[leynilögreglumaður] \textls[23]{\dicTerm{leyni··lög·reglu·|maður} \dicIPA{{l}\-{ei}\-{\textlengthmark}\-{n}\-{\textsci}\-{l}\-{\oe}\-{\textbabygamma}\-{r}\-{\textepsilon}\-{\r{g}}\-{l}\-{\textscy}\-{m}\-{a}\-{ð}\-{\textscy}\-{\textsubring{r}}} \dicPos{m}[13] \dicFlx{(‑manns, ‑menn)}[2]} \dicDirectTranslationCS{detektiv, tajný policista, tajná policistka}

\dicEntry[leynilögreglusaga] \dicTerm{leyni·lög·reglu··|saga} \dicIPA{{l}\-{ei}\-{\textlengthmark}\-{n}\-{\textsci}\-{l}\-{\oe}\-{\textbabygamma}\-{r}\-{\textepsilon}\-{\r{g}}\-{l}\-{\textscy}\-{s}\-{a}\-{\textbabygamma}\-{a}\-} \dicPos{f}[1] \dicFlx{(‑sögu, ‑sög\-ur)}[14] \dicDirectTranslationCS{detektivka, detektivní příběh}
\dicEntry[leyniskytta] \dicTerm{leyni··skytt|a} \dicIPA{{l}{ei}{\textlengthmark}{n}{\textsci}{s}{\r{\textObardotlessj}}{\textsci}{h}{\textsubring{d}}{a}} \dicPos{f}[1] \dicFlx{(‑u, ‑ur)}[19] \dicDirectTranslationCS{odstřelovač(ka), snajper(ka)}
\dicEntry[leyniþjónusta] \textls[15]{\dicTerm{leyni··þjón·ust|a} \dicIPA{{l}{ei}{\textlengthmark}{n}{\textsci}{\texttheta}{j}{ou}{n}{\textscy}{s}{\textsubring{d}}{a}} \dicPos{f}[1] \dicFlx{(‑u)}[5] \dicDirectTranslationCS{tajná\,/\addthin vý\-zvěd\-ná služba, rozvědka}}
\dicEntry[leysa] \dicTerm{leys|a} \dicsymFrequent\  \dicIPA{{l}{ei}{\textlengthmark}{s}{a}} \dicPos{v}[2] \dicFlx{(‑ti, ‑t)}[64] \dicFlx{acc} \textbf{1.} \dicSynonym{losa} \dicDirectTranslationCS{rozvázat, rozvazovat, rozmotat, rozmotávat, uvolnit, uvolňovat} \dicExampleIS{leysa bandið} \dicExampleCS{rozvázat provázek}  \textbf{2.} \dicSynonym*{finna lausn} \dicDirectTranslationCS{vyřešit, rozřešit, rozuzlit}  \textbf{3.} \dicPhraseIS{e‑ð leysir} \dicFlx{impers} \dicDirectTranslationCS{(co) taje, (co) roztává} \dicExampleIS{Ísa leysir.} \dicExampleCS{Ledy tají.};  \dicIdiom{leysa}[af]{ \dicPhraseIS{leysa e‑n af}} \dicDirectTranslationCS{zastoupit (koho)}; { \dicPhraseIS{leysa e‑ð af hólmi}} \dicDirectTranslationCS{nahradit (co), zastoupit (co)};  \dicIdiom{leysa}[niður]{ \dicPhraseIS{leysa niður um sig}} \dicDirectTranslationCS{stáhnout si kalhoty};  \dicIdiom{leysa}[undan]{ \dicPhraseIS{leysa e‑n undan e‑u}} \dicDirectTranslationCS{oprostit (koho) od (čeho), osvobodit (koho) od (čeho)};  \dicIdiom{leysa}[upp]{ \dicPhraseIS{leysa e‑ð upp}} \dicDirectTranslationCS{rozpustit\,/\addthin rozpouštět (co) (sůl ve vodě ap.)};  \dicIdiom{leysa}[úr]{ \dicPhraseIS{leysa úr e‑u}} \dicDirectTranslationCS{vyřešit (co), nalézt řešení (čeho)} \dicExampleIS{leysa úr vanda} \dicExampleCS{vyřešit problém};  \dicIdiom{leysa}[út]{ \dicPhraseIS{leysa e‑ð út}} \dicDirectTranslationCS{zaplatit za (co) (za objednané zboží)};  \dicPhraseIS{leysa e‑n út með e‑ð} \dicDirectTranslationCS{obdarovat (koho čím) na odchodu};  \dicIdiom{leysast}{ \dicPhraseIS{leysast}} \dicFlx{refl} {\textbf{a.}} \dicDirectTranslationCS{rozvázat se, rozmotat se};  {\textbf{b.}} \dicDirectTranslationCS{vyřešit se, rozuzlit se} \dicExampleIS{Deilan leystist með samkomulagi.} \dicExampleCS{Spor se vyřešil dohodou.}
\dicEntry[leysanlegur] \dicTerm{leysan··legur} \dicIPA{{l}{ei}{\textlengthmark}{s}{a}{n}{l}{\textepsilon}{\textbabygamma}{\textscy}{\textsubring{r}}} \dicPos{adj}[1]\dicFlx{}[-8] \textbf{1.} \dicFieldCat{chem.} \dicDirectTranslationCS{rozpustný}  \textbf{2.} \dicDirectTranslationCS{(vy)řešitelný}
\dicEntry[leysiefni] \dicTerm{leysi··efni} \dicIPA{{l}{ei}{\textlengthmark}{s}{\textsci}{\textepsilon}{\textsubring{b}}{n}{\textsci}} \dicPos{n}[2] \dicFlx{(‑s, ‑)}[14] \dicFieldCat{chem.} \dicDirectTranslationCS{rozpouštědlo}
\dicEntry[leysigeisli] \dicTerm{leysi··geisl|i} \dicIPA{{l}{ei}{\textlengthmark}{s}{\textsci}{\r{\textObardotlessj}}{ei}{s}{\textsubring{d}}{l}{\textsci}} \dicPos{m}[1] \dicFlx{(‑a, ‑ar)}[1] \dicFieldCat{fyz.} \dicDirectTranslationCS{laserový paprsek}
\dicEntry[leysing] \dicTerm{leys··ing} \dicIPA{{l}{ei}{\textlengthmark}{s}{i}{\ng}{\r{g}}} \dicPos{f}[4] \dicFlx{(‑ar, ‑ar)}[5] \textbf{1.} \dicSynonym{lausn} \dicDirectTranslationCS{(vy)řešení, rozuzlení}  \textbf{2.} \dicSynonym*{snjóbráð} \dicDirectTranslationCS{tání, obleva}  \textbf{3.} \dicSynonym{upplausn} \dicDirectTranslationCS{rozpuštění, rozpouštění (látky ve vodě ap.)} \dicExampleIS{leysing sykurs í vatni} \dicExampleCS{rozpuštění cukru ve vodě}
\dicEntry[leysingi] \dicTerm{leys··ing|i} \dicIPA{{l}{ei}{\textlengthmark}{s}{i}{\textltailn}{\r{\textObardotlessj}}{\textsci}} \dicPos{m}[1] \dicFlx{(‑ja, ‑jar)}[14] \dicFieldCat{hist.} \dicDirectTranslationCS{osvobozený otrok, osvobozená otrokyně}
\dicEntry[leysinn] \dicTerm{leysinn} \dicIPA{{l}{ei}{\textlengthmark}{s}{\textsci}{\textsubring{n}}} \dicPos{adj}[6]\dicFlx{}[-6] \dicDirectTranslationCS{rozpustný}
\dicEntry[leysiprentari] \dicTerm{leysi··prent·ar|i} \dicIPA{{l}{ei}{\textlengthmark}{s}{\textsci}{p\smash{\textsuperscript{h}}}{r}{\textepsilon}{\textsubring{n}}{\textsubring{d}}{a}{r}{\textsci}} \dicPos{m}[1] \dicFlx{(‑a, ‑ar)}[13] \dicFieldCat{poč.} \dicDirectTranslationCS{laserová tiskárna}
\dicEntry[leysir] \dicTerm{leys|ir} \dicIPA{{l}{ei}{\textlengthmark}{s}{\textsci}{\textsubring{r}}} \dicPos{m}[7] \dicFlx{(‑is, ‑ar)}[1] \dicFieldCat{fyz.} \dicDirectTranslationCS{laser}
\dicEntry[leysni] \dicTerm{leysn|i} \dicIPA{{l}{ei}{s}{\textsubring{d}}{n}{\textsci}} \dicPos{f}[3] \dicFlx{(‑i)}[3] \dicFieldCat{chem.} \dicDirectTranslationCS{rozpustnost (materiálu ap.)}
\dicEntry[leyti] \dicTerm{leyti} \dicsymFrequent\  \dicIPA{{l}{ei}{\textlengthmark}{\textsubring{d}}{\textsci}} \dicPos{n}[2] \dicFlx{(‑s)}[20] \textbf{1.} \dicSynonym{tillit} \dicDirectTranslationCS{ohled, zřetel, míra};  \dicPhraseIS{að mestu leyti} \dicFlx{adv} \dicDirectTranslationCS{převážně, většinou};  \dicPhraseIS{að miklu leyti} \dicFlx{adv} \dicDirectTranslationCS{do velké míry, z~velké části};  \dicPhraseIS{að mörgu leyti} \dicFlx{adv} \dicDirectTranslationCS{do značné míry};  \dicPhraseIS{að nokkru leyti} \dicFlx{adv} \dicDirectTranslationCS{do určité míry, částečně};  \dicPhraseIS{að vissu leyti} \dicFlx{adv} \dicDirectTranslationCS{do jisté míry};  \dicPhraseIS{að öðru leyti} \dicFlx{adv} \dicDirectTranslationCS{jinak};  \dicPhraseIS{að þessu leyti} \dicFlx{adv} \dicDirectTranslationCS{v~tomto ohledu\,/\addthin směru}  \textbf{2.} \dicSynonym{tími} \dicDirectTranslationCS{doba, čas} \dicExampleIS{um sjöleytið} \dicExampleCS{kolem sedmé hodiny};  \dicPhraseIS{um hvaða leyti} \dicDirectTranslationCS{v~kolik (hodin)};  \dicPhraseIS{um sama leyti} \dicFlx{adv} \dicDirectTranslationCS{(přibližně) ve stejný čas};  \dicPhraseIS{um það leyti} \dicFlx{adv} \dicDirectTranslationCS{tehdy, tenkrát}
\dicEntry[léð] \dicTerm{léð} \dicIPA{{l}{j}{\textepsilon}{\textlengthmark}{\texttheta}} \dicPos{v} \dicFlx{supin} \dicLink{ljá}
\dicEntry[léði] \dicTerm{léði} \dicIPA{{l}{j}{\textepsilon}{\textlengthmark}{ð}{\textsci}} \dicPos{v} \dicFlx{ind\,/\addthin con pf sg 1 pers} \dicLink{ljá}
\dicEntry[léðum] \dicTerm{léðum} \dicIPA{{l}{j}{\textepsilon}{\textlengthmark}{ð}{\textscy}{\textsubring{m}}} \dicPos{v} \dicFlx{ind pf pl 1 pers} \dicLink{ljá}
\dicEntry[lék] \dicTerm{lék} \dicIPA{{l}{j}{\textepsilon}{\textlengthmark}{\r{g}}} \dicPos{v} \dicFlx{ind pf sg 1 pers} \dicLink{leika}
\dicEntry[léki] \dicTerm{léki} \dicIPA{{l}{j}{\textepsilon}{\textlengthmark}{\r{\textObardotlessj}}{\textsci}} \dicPos{v} \dicFlx{con pf sg 1 pers} \dicLink{leika}
\dicEntry[lékum] \dicTerm{lékum} \dicIPA{{l}{j}{\textepsilon}{\textlengthmark}{\r{g}}{\textscy}{\textsubring{m}}} \dicPos{v} \dicFlx{ind pf pl 1 pers} \dicLink{leika}
\dicEntry[lélega] \dicTerm{lé··lega} \dicIPA{{l}{j}{\textepsilon}{\textlengthmark}{l}{\textepsilon}{\textbabygamma}{a}} \dicPos{adv} \dicDirectTranslationCS{špatně, bídně, mizerně}
\dicEntry[lélegur] \dicTerm{lé··legur} \dicsymFrequent\  \dicIPA{{l}{j}{\textepsilon}{\textlengthmark}{l}{\textepsilon}{\textbabygamma}{\textscy}{\textsubring{r}}} \dicPos{adj}[1]\dicFlx{}[-8] \dicSynonym{slæmur} \dicDirectTranslationCS{špatný, bídný, mizerný, chabý} \dicExampleIS{benda á lélegan árangur} \dicExampleCS{upozornit na špatný výsledek}
\dicEntry[lén] \dicTerm{lén} \dicIPA{{l}{j}{\textepsilon}{\textlengthmark}{\textsubring{n}}} \dicPos{n}[2] \dicFlx{(‑s, ‑)}[5] \textbf{1.} \dicSynonym{ríki} \dicDirectTranslationCS{léno}  \textbf{2.} \dicFieldCat{poč.} \dicSynonym{umdæmi} \dicDirectTranslationCS{doména}  \textbf{3.} \dicFieldCat{biol.} \dicDirectTranslationCS{doména}
\dicEntry[lénskur] \dicTerm{lénskur} \dicIPA{{l}{j}{\textepsilon}{n}{s}{\r{g}}{\textscy}{\textsubring{r}}} \dicPos{adj}[1]\dicFlx{}[-1] \dicDirectTranslationCS{feudální}
\dicEntry[lénsmaður] \dicTerm{léns··|maður} \dicIPA{{l}{j}{\textepsilon}{n}{s}{m}{a}{ð}{\textscy}{\textsubring{r}}} \dicPos{m}[13] \dicFlx{(‑manns, ‑menn)}[2] \dicDirectTranslationCS{leník, vazal(ka)}
\dicEntry[léreft] \dicTerm{léreft} \dicIPA{{l}{j}{\textepsilon}{\textlengthmark}{r}{\textepsilon}{f}{\textsubring{d}}} \dicPos{n}[2] \dicFlx{(‑s, ‑)}[5] \dicDirectTranslationCS{(lněné) plátno}
\dicEntry[lét] \dicTerm{lét} \dicIPA{{l}{j}{\textepsilon}{\textlengthmark}{\textsubring{d}}} \dicPos{v} \dicFlx{ind pf sg 1 pers} \dicLink{láta}
\dicEntry[léti] \dicTerm{léti} \dicIPA{{l}{j}{\textepsilon}{\textlengthmark}{\textsubring{d}}{\textsci}} \dicPos{v} \dicFlx{con pf sg 1 pers} \dicLink{láta}
\dicEntry[létt] \dicTerm{létt} \dicsymFrequent\  \dicIPA{{l}{j}{\textepsilon}{h}{\textsubring{d}}} \dicPos{adv} \dicFlx{(comp ‑ar, sup ‑ast)} \dicDirectTranslationCS{lehce, mírně, jemně, zlehka} \dicExampleIS{Hann klappaði létt á bakið á mér.} \dicExampleCS{Lehce mě poklepal po zádech.}
\dicEntry[létta] \dicTerm{létt|a} \dicsymFrequent\  \dicIPA{{l}{j}{\textepsilon}{h}{\textsubring{d}}{a}} \dicPos{v}[2] \dicFlx{(‑i, ‑)}[2] \dicFlx{acc\,/\addthin dat} \textbf{1.} \dicDirectTranslationCS{odlehčit, odlehčovat, ulehčit, ulehčovat, (u)činit lehčím} \dicExampleIS{létta byrðina} \dicExampleCS{odlehčit břemeno}  \textbf{2.} \dicPhraseIS{e‑m léttir} \dicFlx{impers} \dicDirectTranslationCS{(komu) se ulehčuje, (komu) se ulevuje}  \textbf{3.} \dicPhraseIS{e‑u léttir} \dicFlx{impers} \dicDirectTranslationCS{(čeho) ubývá (mlhy ap.), (co) mizí};  \dicIdiom{létta}[af]{ \dicPhraseIS{létta af e‑m e‑u}} \dicDirectTranslationCS{odlehčit (komu) od (čeho)}; { \dicPhraseIS{það léttir af e‑m e‑u}} \dicFlx{impers} \dicDirectTranslationCS{(komu) se ulevuje od (čeho)};  \dicIdiom{létta}[á]{ \dicPhraseIS{létta á sér}} \dicDirectTranslationCS{ulevit si} \dicIndirectTranslationCS{(vykonat tělesnou potřebu)};  \dicIdiom{létta}[sér]{ \dicPhraseIS{létta sér lífið}} \dicDirectTranslationCS{pobavit se, povyrazit si};  \dicIdiom{létta}[til]{ \dicPhraseIS{létta til}} \dicDirectTranslationCS{vyčasit se, vyjasnit se, vyjasňovat se} \dicExampleIS{Það er að létta til.} \dicExampleCS{Začíná se vyjasňovat.};  \dicIdiom{létta}[undir]{ \dicPhraseIS{létta undir með e‑m}} \dicDirectTranslationCS{ulehčit (komu), učinit to (komu) lehčím};  \dicIdiom{léttast}{ \dicPhraseIS{léttast}} \dicFlx{refl} \dicDirectTranslationCS{(z)hubnout};  \dicIdiom{létta}{ \dicPhraseIS{létta akkerum}} \dicFieldCat{nám.} \dicDirectTranslationCS{zvednout kotvu}
\dicEntry[léttasótt] \dicTerm{létta··sótt} \dicIPA{{l}{j}{\textepsilon}{h}{\textsubring{d}}{a}{s}{ou}{h}{\textsubring{d}}} \dicPos{f}[7] \dicFlx{(‑ar)}[3] \dicDirectTranslationCS{porodní bolesti, kontrakce} \dicExampleIS{taka léttasóttina} \dicExampleCS{mít porodní bolesti}
\dicEntry[létteind] \dicTerm{létt··eind} \dicIPA{{l}{j}{\textepsilon}{h}{\textsubring{d}}{ei}{n}{\textsubring{d}}} \dicPos{f}[7] \dicFlx{(‑ar, ‑ir)}[1] \dicFieldCat{fyz.} \dicDirectTranslationCS{lepton}
\dicEntry[létti] \dicTerm{létt|i} \dicIPA{{l}{j}{\textepsilon}{h}{\textsubring{d}}{\textsci}} \dicPos{m}[1] \dicFlx{(‑a, ‑ar)}[1] \textbf{1.} \dicSynonym{léttir} \dicDirectTranslationCS{ulehčení, úleva}  \textbf{2.} \dicSynonym{axlaband} \dicDirectTranslationCS{popruh, šle}  \textbf{3.} \dicSynonym*{fúsleiki} \dicDirectTranslationCS{ochota};  \dicPhraseIS{segja allt af létta} \dicDirectTranslationCS{ochotně všechno povědět}
\dicEntry[léttir] \dicTerm{létt|ir} \dicIPA{{l}{j}{\textepsilon}{h}{\textsubring{d}}{\textsci}{\textsubring{r}}} \dicPos{m}[7] \dicFlx{(‑is, ‑ar)}[1] \textbf{1.} \dicSynonym{huggun} \dicDirectTranslationCS{útěcha, úleva}  \textbf{2.} \dicSynonym{hjálp} \dicDirectTranslationCS{pomoc};  \dicPhraseIS{vera e‑m til léttis} \dicDirectTranslationCS{pomoci (komu), ulehčit (komu)}  \textbf{3.} \dicFieldCat{zool.} \dicDirectTranslationCS{delfín obecný} \textit{(l.~{\textLA{Delphinus delphis}})}  \dicsymPhoto\ 
\dicFigure{26462.jpg}{Léttir}{Léttir - Shane Anderson, Biolib, PD}
\dicEntry[léttklæddur] \dicTerm{létt··klæddur} \dicIPA{{l}{j}{\textepsilon}{h}{\textsubring{d}}{k\smash{\textsuperscript{h}}}{l}{a}{i}{\textsubring{d}}{\textscy}{\textsubring{r}}} \dicPos{adj}[2]\dicFlx{}[-18] \dicDirectTranslationCS{lehce oblečený}
\dicEntry[léttleiki] \dicTerm{létt··leik|i} \dicIPA{{l}{j}{\textepsilon}{h}{\textsubring{d}}{l}{ei}{\r{\textObardotlessj}}{\textsci}} \dicPos{m}[1] \dicFlx{(‑a)}[3] \dicDirectTranslationCS{lehkost} \dicExampleIS{óbærilegur léttleiki tilverunnar} \dicExampleCS{nesnesitelná lehkost bytí}
\dicEntry[léttlyndi] \dicTerm{létt··lyndi} \dicIPA{{l}{j}{\textepsilon}{h}{\textsubring{d}}{l}{\textsci}{n}{\textsubring{d}}{\textsci}} \dicPos{n}[2] \dicFlx{(‑s)}[20] \dicDirectTranslationCS{dobrá nálada, uvolněnost, bezstarostnost}
\dicEntry[léttlyndur] \dicTerm{létt··lyndur} \dicIPA{{l}{j}{\textepsilon}{h}{\textsubring{d}}{l}{\textsci}{n}{\textsubring{d}}{\textscy}{\textsubring{r}}} \dicPos{adj}[2]\dicFlx{}[-14] \dicDirectTranslationCS{uvolněný, bezstarostný}
\dicEntry[léttmjólk] \dicTerm{létt··mjólk} \dicIPA{{l}{j}{\textepsilon}{h}{\textsubring{d}}{m}{j}{ou}{\textsubring{l}}{\r{g}}} \dicPos{f}[10] \dicFlx{(‑ur)}[2] \dicDirectTranslationCS{nízkotučné mléko}
\dicEntry[léttskýjaður] \dicTerm{létt··skýj·|aður} \dicIPA{{l}{j}{\textepsilon}{h}{\textsubring{d}}{s}{\r{\textObardotlessj}}{i}{j}{a}{ð}{\textscy}{\textsubring{r}}} \dicPos{adj}[3] \dicFlx{(f ‑uð)}[3] \dicPhraseIS{það er léttskýjað} \dicDirectTranslationCS{je polojasno}
\dicEntry[léttur] \dicTerm{léttur} \dicsymFrequent\  \dicIPA{{l}{j}{\textepsilon}{h}{\textsubring{d}}{\textscy}{\textsubring{r}}} \dicPos{adj}[1]\dicFlx{}[-10] \textbf{1.} \dicDirectTranslationCS{lehký, snadný} \dicExampleIS{létt verk} \dicExampleCS{lehká práce} \dicAntonym{erfiður}  \textbf{2.} \dicDirectTranslationCS{lehký} \dicIndirectTranslationCS{(vážící málo)} \dicAntonym{þungur}  \textbf{3.} \dicSynonym{glaður} \dicDirectTranslationCS{uvolněný, klidný, bezstarostný} \dicExampleIS{Hún er létt í lund.} \dicExampleCS{Je bezstarostná.}  \textbf{4.} \dicDirectTranslationCS{dietní, lehký (strava ap.)} \dicExampleIS{léttar veitingar} \dicExampleCS{lehké občerstvení}  \textbf{5.} \dicDirectTranslationCS{lehký, slabý (alkohol ap.)} \dicExampleIS{létt vín} \dicExampleCS{slabý alkohol}
\dicEntry[léttúð] \dicTerm{létt··úð} \dicIPA{{l}{j}{\textepsilon}{h}{\textsubring{d}}{u}{\texttheta}} \dicPos{f}[4] \dicFlx{(‑ar)}[3] \textbf{1.} \dicSynonym{alvöruleysi} \dicDirectTranslationCS{lehkovážnost}  \textbf{2.} \dicSynonym{lauslæti} \dicDirectTranslationCS{rozpustilost, frivolnost}
\dicEntry[léttúðugur] \dicTerm{létt·úð··ugur} \dicIPA{{l}{j}{\textepsilon}{h}{\textsubring{d}}{u}{ð}{\textscy}{\textbabygamma}{\textscy}{\textsubring{r}}} \dicPos{adj}[1]\dicFlx{}[-8] \dicSynonym{lauslátur} \dicDirectTranslationCS{frivolní, prostopášný}
\dicEntry[léttvín] \dicTerm{létt··vín} \dicIPA{{l}{j}{\textepsilon}{h}{\textsubring{d}}{v}{i}{\textsubring{n}}} \dicPos{n}[2] \dicFlx{(‑s, ‑)}[5] \dicDirectTranslationCS{víno}
\dicEntry[léttvægur] \dicTerm{létt··vægur} \dicIPA{{l}{j}{\textepsilon}{h}{\textsubring{d}}{v}{a}{i}{\textbabygamma}{\textscy}{\textsubring{r}}} \dicPos{adj}[1]\dicFlx{}[-1] \dicDirectTranslationCS{nedůležitý, nepodstat\-ný, nevýznamný}
\dicEntry[létum] \dicTerm{létum} \dicIPA{{l}{j}{\textepsilon}{\textlengthmark}{\textsubring{d}}{\textscy}{\textsubring{m}}} \dicPos{v} \dicFlx{ind pf pl 1 pers} \dicLink{láta}
\dicEntry[lh. nt.] \dicTerm{lh. nt.} \dicPos{zkr} \dicPhraseIS{lýsingarháttur nútíðar} \dicFieldCat{jaz.} \dicDirectTranslationCS{příčestí přítomné}
\dicEntry[lh. þt.] \dicTerm{lh. þt.} \dicPos{zkr} \dicPhraseIS{lýsingarháttur þátíðar} \dicFieldCat{jaz.} \dicDirectTranslationCS{příčestí minulé}
\dicEntry[lið] \dicTerm{lið} \dicsymFrequent\  \dicIPA{{l}{\textsci}{\textlengthmark}{\texttheta}} \dicPos{n}[2] \dicFlx{(‑s, ‑)}[5] \textbf{1.} \dicSynonym{hópur} \dicDirectTranslationCS{družstvo, oddíl, četa};  \dicPhraseIS{ganga í lið við e‑n} \dicDirectTranslationCS{spojit síly s~(kým)};  \dicPhraseIS{safna liði} \dicDirectTranslationCS{shromažďovat podporu}  \textbf{2.} \dicFieldCat{sport.} \dicDirectTranslationCS{družina, tým, mužstvo}  \textbf{3.} \dicSynonym{aðstoð} \dicDirectTranslationCS{výpomoc, pomocná ruka};  \dicPhraseIS{koma til liðs við e‑n} \dicDirectTranslationCS{vypomoct (komu), přijít (komu) na pomoc};  \dicPhraseIS{verða e‑m að liði} \dicDirectTranslationCS{vypomoct (komu), asistovat (komu)}  \textbf{4.} \dicLangCat{hovor.} \dicSynonym{fólk} \dicDirectTranslationCS{lidi}
\dicEntry[liða] \dicTerm{lið|a} \dicIPA{{l}{\textsci}{\textlengthmark}{ð}{a}} \dicPos{v}[1] \dicFlx{(‑aði)}[1] \dicFlx{acc} \textbf{1.} \dicSynonym{skipta} \dicDirectTranslationCS{(roz)dělit (do skupin)}  \textbf{2.} \dicSynonym*{gera hrokkið} \dicDirectTranslationCS{natočit, natáčet, (na)kadeřit (vlasy ap.)} \dicExampleIS{liða hár} \dicExampleCS{natáčet vlasy};  \dicIdiom{liðast}{ \dicPhraseIS{liðast}} \dicFlx{refl} \dicDirectTranslationCS{vlnit se, kroutit se, kadeřit se}
\dicEntry[liðaður] \dicTerm{lið··|aður} \dicIPA{{l}{\textsci}{\textlengthmark}{ð}{a}{ð}{\textscy}{\textsubring{r}}} \dicPos{adj}[3] \dicFlx{(f ‑uð)}[3] \dicSynonym*{bylgjaður} \dicDirectTranslationCS{natočený, nakadeřený, navlněný}
\dicEntry[liðagigt] \dicTerm{liða··gigt} \dicIPA{{l}{\textsci}{\textlengthmark}{ð}{a}{\r{\textObardotlessj}}{\textsci}{x}{\textsubring{d}}} \dicPos{f}[4] \dicFlx{(‑ar)}[3] \dicFieldCat{med.} \dicDirectTranslationCS{revmatoidní artritida}
\dicEntry[liðamót] \dicTerm{liða··mót} \dicIPA{{l}{\textsci}{\textlengthmark}{ð}{a}{m}{ou}{\textsubring{d}}} \dicPos{n}[2] \dicFlx{pl}[1] \dicFieldCat{anat.} \dicDirectTranslationCS{kloub}
\dicEntry[liðband] \dicTerm{lið··|band} \dicIPA{{l}{\textsci}{ð}{\textsubring{b}}{a}{n}{\textsubring{d}}} \dicPos{n}[2] \dicFlx{(‑bands, ‑bönd)}[8] \dicFieldCat{anat.} \dicDirectTranslationCS{vazivo}
\dicEntry[liðdýr] \dicTerm{lið··dýr} \dicIPA{{l}{\textsci}{ð}{\textsubring{d}}{i}{\textsubring{r}}} \dicPos{n}[2] \dicFlx{(‑s, ‑)}[5] \dicFieldCat{zool.} \dicDirectTranslationCS{členovec} \textit{(l.~{\textLA{Arthropoda}})}
\dicEntry[liðfall] \dicTerm{lið··|fall} \dicIPA{{l}{\textsci}{ð}{f}{a}{\textsubring{d}}{\textsubring{l}}} \dicPos{n}[2] \dicFlx{(‑falls, ‑föll)}[8] \dicFieldCat{jaz.} \dicSynonym{brottfall} \dicDirectTranslationCS{elipsa, výpustka}
\dicEntry[liðfelldur] \dicTerm{lið··felldur} \dicIPA{{l}{\textsci}{ð}{f}{\textepsilon}{l}{\textsubring{d}}{\textscy}{\textsubring{r}}} \dicPos{adj}[2]\dicFlx{}[-14] \dicFieldCat{jaz.} \dicDirectTranslationCS{eliptický} \dicExampleIS{liðfellt orðtak} \dicExampleCS{eliptický výraz}
\dicEntry[liðhlaupi] \dicTerm{lið··hlaup|i} \dicIPA{{l}{\textsci}{ð}{\textsubring{l}}{\oe i}{\textsubring{b}}{\textsci}} \dicPos{m}[1] \dicFlx{(‑a, ‑ar)}[1] \dicFieldCat{voj.} \dicDirectTranslationCS{zběh, dezertér(ka)}
\dicEntry[liði] \dicTerm{lið|i\smash{\textsuperscript{1}}} \dicIPA{{l}{\textsci}{\textlengthmark}{ð}{\textsci}} \dicPos{m}[1] \dicFlx{(‑a, ‑ar)}[1] \dicDirectTranslationCS{člen(ka) týmu (vojenského, zdravotního ap.)} \dicIndirectTranslationCS{(vyskytuje se jako druhá část složeného slova)} \dicExampleIS{skæruliði} \dicExampleCS{partyzán}
\dicEntry[liði] \dicTerm{liði\smash{\textsuperscript{2}}} \dicIPA{{l}{\textsci}{\textlengthmark}{ð}{\textsci}} \dicPos{v} \dicFlx{con pf sg 1 pers} \dicLink{líða}
\dicEntry[liðið] \dicTerm{liðið} \dicIPA{{l}{\textsci}{\textlengthmark}{ð}{\textsci}{\texttheta}} \dicPos{v} \dicFlx{supin} \dicLink{líða}
\dicEntry[liðinn] \dicTerm{liðinn} \dicsymFrequent\  \dicIPA{{l}{\textsci}{\textlengthmark}{ð}{\textsci}{\textsubring{n}}} \dicPos{adj}[6]\dicFlx{}[-6] \textbf{1.} \dicDirectTranslationCS{minulý, uplynulý} \dicExampleIS{í liðnum mánuði} \dicExampleCS{v~minulém měsíci}  \textbf{2.} \dicDirectTranslationCS{zesnulý};  \dicPhraseIS{vera vel liðinn} \dicDirectTranslationCS{být populární\,/\addthin oblíbený}
\dicEntry[liðka] \dicTerm{liðk|a} \dicIPA{{l}{\textsci}{\texttheta}{\r{g}}{a}} \dicPos{v}[1] \dicFlx{(‑aði)}[1] \dicFlx{acc} \dicDirectTranslationCS{rozhýbat, rozpohybovat} \dicExampleIS{liðka handlegginn} \dicExampleCS{rozhýbat ruku};  \dicPhraseIS{liðka sig} \dicDirectTranslationCS{rozhýbat se};  \dicIdiom{liðka}[til]{ \dicPhraseIS{liðka e‑ð til}} \dicDirectTranslationCS{usnadnit (co)};  \dicIdiom{liðkast}{ \dicPhraseIS{liðkast}} \dicFlx{refl} \dicDirectTranslationCS{rozhýbat se}
\dicEntry[liðlangur] \dicTerm{lið··|langur} \dicIPA{{l}{\textsci}{ð}{l}{au}{\ng}{\r{g}}{\textscy}{\textsubring{r}}} \dicPos{adj}[10] \dicFlx{(f ‑löng, comp ‑lengri, sup ‑leng\-st\-ur)}[2] \dicPhraseIS{allan liðlangan daginn} \dicFlx{adv} \dicDirectTranslationCS{(po) celý celičký den}
\dicEntry[liðlegur] \dicTerm{lið··legur} \dicIPA{{l}{\textsci}{ð}{l}{\textepsilon}{\textbabygamma}{\textscy}{\textsubring{r}}} \dicPos{adj}[1]\dicFlx{}[-8] \textbf{1.} \dicSynonym{fimur} \dicDirectTranslationCS{obratný}  \textbf{2.} \dicSynonym{hjálpsamur} \dicDirectTranslationCS{vstřícný, úslužný, ochotný pomoci}
\dicEntry[liðleiki] \dicTerm{lið··leik|i} \dicIPA{{l}{\textsci}{ð}{l}{ei}{\r{\textObardotlessj}}{\textsci}} \dicPos{m}[1] \dicFlx{(‑a)}[3] \textbf{1.} \dicDirectTranslationCS{pohyblivost, obratnost}  \textbf{2.} \dicDirectTranslationCS{vstřícnost, ochota pomoci}
\dicEntry[liðsauki] \dicTerm{liðs··auk|i} \dicIPA{{l}{\textsci}{ð}{s}{\oe i}{\r{\textObardotlessj}}{\textsci}} \dicPos{m}[1] \dicFlx{(‑a)}[3] \dicDirectTranslationCS{posila, posílení\,/\addthin rozšíření řad}
\dicEntry[liðsemd] \textls[30]{\dicTerm{lið··semd} \dicIPA{{l}{\textsci}{ð}{s}{\textepsilon}{m}{\textsubring{d}}} \dicPos{f}[7] \dicFlx{(‑ar)}[3] \dicSynonym{aðstoð} \dicDirectTranslationCS{výpomoc, asis\nobreak tence}}
\dicEntry[liðsforingi] \dicTerm{liðs··for·ing|i} \dicIPA{{l}{\textsci}{ð}{s}{f}{\textopeno}{r}{i}{\textltailn}{\r{\textObardotlessj}}{\textsci}} \dicPos{m}[1] \dicFlx{(‑ja, ‑jar)}[14] \dicFieldCat{voj.} \dicDirectTranslationCS{poručík, poručice}
\dicEntry[liðsinna] \dicTerm{lið··sinn|a} \dicIPA{{l}{\textsci}{ð}{s}{\textsci}{n}{a}} \dicPos{v}[2] \dicFlx{(‑ti, ‑t)}[78] \dicFlx{dat} \dicSynonym{hjálpa} \dicDirectTranslationCS{asistovat, (vy)pomáhat} \dicExampleIS{liðsinna e‑m við e‑ð} \dicExampleCS{asistovat (komu) při (čem)}
\dicEntry[liðsinni] \dicTerm{lið··sinni} \dicIPA{{l}{\textsci}{ð}{s}{\textsci}{n}{\textsci}} \dicPos{n}[2] \dicFlx{(‑s)}[20] \dicSynonym{aðstoð} \dicDirectTranslationCS{asistence, (vý)pomoc}
\dicEntry[liðsmaður] \dicTerm{liðs··|maður} \dicIPA{{l}{\textsci}{ð}{s}{m}{a}{ð}{\textscy}{\textsubring{r}}} \dicPos{m}[13] \dicFlx{(‑manns, ‑menn)}[2] \textbf{1.} \dicSynonym{hermaður} \dicDirectTranslationCS{voják(yně)}  \textbf{2.} \dicSynonym{stuðningsmaður} \dicDirectTranslationCS{přívrženec, přívrženkyně, stoupenec, stoupenkyně}  \textbf{3.} \dicDirectTranslationCS{(spolu)hráč(ka), člen(ka) družiny\,/\addthin týmu}
\dicEntry[liðtækur] \dicTerm{lið··tækur} \dicIPA{{l}{\textsci}{ð}{t\smash{\textsuperscript{h}}}{a}{i}{\r{g}}{\textscy}{\textsubring{r}}} \dicPos{adj}[1]\dicFlx{}[-6] \dicSynonym{hæfur} \dicDirectTranslationCS{kompetentní, způsobilý}
\dicEntry[liðugur] \dicTerm{lið··ugur} \dicIPA{{l}{\textsci}{\textlengthmark}{ð}{\textscy}{\textbabygamma}{\textscy}{\textsubring{r}}} \dicPos{adj}[1]\dicFlx{}[-8] \dicSynonym{fimur} \dicDirectTranslationCS{hbitý, mrštný} \dicExampleIS{vera liðugur í hreyfingum} \dicExampleCS{pohybovat se hbitě};  \dicPhraseIS{vera laus og liðugur} \dicDirectTranslationCS{být volný, užívat si svobody}
\dicEntry[liðum] \dicTerm{liðum} \dicIPA{{l}{\textsci}{\textlengthmark}{ð}{\textscy}{\textsubring{m}}} \dicPos{v} \dicFlx{ind pf pl 1 pers} \dicLink{líða}
\dicEntry[liður] \dicTerm{lið|ur} \dicsymFrequent\  \dicIPA{{l}{\textsci}{\textlengthmark}{ð}{\textscy}{\textsubring{r}}} \dicPos{m}[10] \dicFlx{(‑ar\,/\addthin ‑s, ‑ir)}[37] \textbf{1.} \dicFieldCat{anat.} \dicSynonym{liðamót} \dicDirectTranslationCS{kloub} \dicExampleIS{axlarliður} \dicExampleCS{ramenní kloub};  \dicPhraseIS{geta hvorki hrært legg né lið} \dicLangCat{přen.} \dicDirectTranslationCS{nemoci se vůbec pohnout}  \textbf{2.} \dicSynonym{bylgja} \dicDirectTranslationCS{kudrna, vlnka (ve vlasech)} \dicExampleIS{liður í hári} \dicExampleCS{kudrna ve vlasech}  \textbf{3.} \dicSynonym{atriði} \dicDirectTranslationCS{položka, bod} \dicExampleIS{Þessi liður er rangur.} \dicExampleCS{Tato položka je špatně.}  \textbf{4.} \dicSynonym{ættliður} \dicDirectTranslationCS{koleno, pokolení, generace}  \textbf{5.} \dicFieldCat{mat.} \dicDirectTranslationCS{člen}
\dicEntry[liðveisla] \dicTerm{lið··veisl|a} \dicIPA{{l}{\textsci}{ð}{v}{ei}{s}{\textsubring{d}}{l}{a}} \dicPos{f}[1] \dicFlx{(‑u, ‑ur)}[13] \dicDirectTranslationCS{asistence, (vý)pomoc}
\dicEntry[liðþjálfi] \dicTerm{lið··þjálf|i} \dicIPA{{l}{\textsci}{ð}{\texttheta}{j}{au}{l}{v}{\textsci}} \dicPos{m}[1] \dicFlx{(‑a, ‑ar)}[1] \dicFieldCat{voj.} \dicDirectTranslationCS{seržant(ka)}
\dicEntry[Liechtenstein] \dicTerm{Liechtenstein} \dicIPA{{l}{i}{x}{\textsubring{d}}{\textepsilon}{n}{s}{\textsubring{d}}{a}{i}{\textsubring{n}}} \dicPos{n}[4] \dicFlx{indecl}[2] \dicFieldCat{geog.} \dicDirectTranslationCS{Lichtenštejnsko}
\dicEntry[Liechtensteini] \dicTerm{Liechtenstein|i} \dicIPA{{l}{i}{x}{\textsubring{d}}{\textepsilon}{n}{s}{\textsubring{d}}{a}{i}{n}{\textsci}} \dicPos{m}[1] \dicFlx{(‑a, ‑ar)}[1] \dicDirectTranslationCS{Lichtenštejňan(ka), Lichtenštejnec, Lichtenštejnka}
\dicEntry[liechtensteinskur] \dicTerm{liechtensteinskur} \dicIPA{{l}{i}{x}{\textsubring{d}}{\textepsilon}{n}{s}{\textsubring{d}}{a}{i}{n}{s}{\r{g}}{\textscy}{\textsubring{r}}} \dicPos{adj}[1]\dicFlx{}[-6] \dicDirectTranslationCS{lichtenštejnský}
\dicEntry[lifa] \dicTerm{lif|a} \dicsymFrequent\  \dicIPA{{l}{\textsci}{\textlengthmark}{v}{a}} \dicPos{v}[2] \dicFlx{(‑ði, ‑að)}[124] \dicFlx{acc} \textbf{1.} \dicDirectTranslationCS{žít} \dicExampleIS{lifa langa ævi} \dicExampleCS{žít dlouhý život}  \textbf{2.} \dicDirectTranslationCS{existovat}  \textbf{3.} \dicDirectTranslationCS{zažít, prožít} \dicExampleIS{lifa þessa atburði} \dicExampleCS{zažít tyto události}  \textbf{4.} \dicSynonym*{vera logandi} \dicDirectTranslationCS{planout, svítit};  \dicIdiom{lifa}[af]{ \dicPhraseIS{lifa af}} \dicSynonym{endast} \dicDirectTranslationCS{přežít}; { \dicPhraseIS{lifa af e‑u}} \dicDirectTranslationCS{vyžít z~(čeho)};  \dicIdiom{lifa}[á]{ \dicPhraseIS{lifa á e‑u}} \dicDirectTranslationCS{živit se (čím)};  \dicIdiom{lifa}[eftir]{ \dicPhraseIS{lifa eftir e‑u}} \dicDirectTranslationCS{žít podle (čeho), žít v~souladu s~(čím)};  \dicIdiom{lifa}[fyrir]{ \dicPhraseIS{lifa fyrir e‑n}} \dicDirectTranslationCS{žít pro (koho)};  \dicIdiom{lifa}[inn í]{ \dicPhraseIS{lifa sig inn í e‑ð}} \dicDirectTranslationCS{vžít se do (čeho)} \dicExampleIS{lifa sig inn í starfið} \dicExampleCS{vžít se do práce};  \dicIdiom{lifa}[við]{ \dicPhraseIS{lifa við e‑ð}} \dicDirectTranslationCS{žít s~(čím) (pocitem viny ap.)}
\dicEntry[lifandi] \dicTerm{lif··|andi\smash{\textsuperscript{1}}} \dicIPA{{l}{\textsci}{\textlengthmark}{v}{a}{n}{\textsubring{d}}{\textsci}} \dicPos{m}[2] \dicFlx{(‑anda, ‑endur)}[1] \dicDirectTranslationCS{živý člověk, živá bytost} \dicIndirectTranslationCS{(zvláště v~pl)}
\dicEntry[lifandi] \dicTerm{lif··andi\smash{\textsuperscript{2}}} \dicsymFrequent\  \dicIPA{{l}{\textsci}{\textlengthmark}{v}{a}{n}{\textsubring{d}}{\textsci}} \dicPos{adj}[13] \dicFlx{indecl}[1] \textbf{1.} \dicSynonym*{á lífi} \dicDirectTranslationCS{živý, (jsoucí) naživu} \dicExampleIS{ná birninum lifandi} \dicExampleCS{chytit medvěda živého}  \textbf{2.} \dicSynonym{líflegur} \dicDirectTranslationCS{živý, čilý};  \dicIdiom{lifandi}{ \dicPhraseIS{ekki nokkur\,/\addthin enginn lifandi maður}} \dicLangCat{přen.} \dicDirectTranslationCS{ani živáček};  \dicPhraseIS{lifandi tónlist} \dicDirectTranslationCS{živá hudba}
\dicEntry[lifna] \dicTerm{lifn|a} \dicsymFrequent\  \dicIPA{{l}{\textsci}{\textsubring{b}}{n}{a}} \dicPos{v}[1] \dicFlx{(‑aði)}[1] \textbf{1.} \dicDirectTranslationCS{oživit se, rozhořet se}  \textbf{2.} \dicPhraseIS{lifna við} \dicSynonym*{vakna til lífs} \dicDirectTranslationCS{oživnout, ožít} \dicExampleIS{Fjöllin lifnuðu við eftir fimmtíu ár.} \dicExampleCS{Hory ožily po padesáti letech.}  \textbf{3.} \dicPhraseIS{það lifnar yfir e‑m} \dicFlx{impers} \dicLangCat{přen.} \dicDirectTranslationCS{(kdo) ožívá, (kdo) se stává živějším}
\dicEntry[lifnaður] \dicTerm{lif··nað|ur} \dicIPA{{l}{\textsci}{\textsubring{b}}{n}{a}{ð}{\textscy}{\textsubring{r}}} \dicPos{m}[10] \dicFlx{(‑ar)}[7] \dicSynonym{líferni} \dicDirectTranslationCS{způsob života, životní styl}
\dicEntry[lifrarkæfa] \dicTerm{lifrar··kæf|a} \dicIPA{{l}{\textsci}{v}{r}{a}{\textsubring{r}}{c\smash{\textsuperscript{h}}}{a}{i}{v}{a}} \dicPos{f}[1] \dicFlx{(‑u)}[5] \dicFieldCat{kulin.} \dicDirectTranslationCS{játrová paštika, játrovka}
\dicEntry[lifrarpylsa] \dicTerm{lifrar··pyls|a} \dicIPA{{l}{\textsci}{v}{r}{a}{\textsubring{r}}{p\smash{\textsuperscript{h}}}{\textsci}{l}{s}{a}} \dicPos{f}[1] \dicFlx{(‑u, ‑ur)}[13] \dicFieldCat{kulin.} \dicDirectTranslationCS{jitrnice}
\dicEntry[lifur] \dicTerm{lif|ur} \dicIPA{{l}{\textsci}{\textlengthmark}{v}{\textscy}{\textsubring{r}}} \dicPos{f}[6] \dicFlx{(‑rar, ‑rar)}[4] \dicFieldCat{anat.} \dicDirectTranslationCS{játra}
\dicEntry[ligg] \dicTerm{ligg} \dicIPA{{l}{\textsci}{\r{g}}{\textlengthmark}} \dicPos{v} \dicFlx{ind praes sg 1 pers} \dicLink{liggja}
\dicEntry[liggja] \dicTerm{liggja} \dicsymFrequent\  \dicIPA{{l}{\textsci}{\r{\textObardotlessj}}{\textlengthmark}{a}} \dicPos{v}[7] \dicFlx{(ligg, lá, lágum, lægi, legið)}[9] \textbf{1.} \dicDirectTranslationCS{ležet} \dicIndirectTranslationCS{(spočívat na něčem podélně)} \dicExampleIS{Blaðið liggur á gólfinu.} \dicExampleCS{Noviny leží na podlaze.}  \textbf{2.} \dicDirectTranslationCS{ležet} \dicIndirectTranslationCS{(být uložen ke spánku)}  \textbf{3.} \dicDirectTranslationCS{ležet, rozkládat se, prostírat se} \dicExampleIS{Vegurinn liggur meðfram ánni.} \dicExampleCS{Cesta vede podél řeky.};  \dicIdiom{liggja}[að]{ \dicPhraseIS{liggja að e‑u}} \dicDirectTranslationCS{sousedit s~(čím)};  \dicIdiom{liggja}[á]{ \dicPhraseIS{ekkert liggur á}} \dicDirectTranslationCS{všechno popořadě, vše ve správný čas}; { \dicPhraseIS{e‑m liggur á}} \dicFlx{impers} \dicDirectTranslationCS{(kdo) má naspěch, (kdo) pospíchá, (kdo) chvátá} \dicExampleIS{Honum lá á.} \dicExampleCS{Měl naspěch.}; { \dicPhraseIS{e‑m liggur á e‑u}} \dicFlx{impers} \dicDirectTranslationCS{(kdo) má naspěch s~(čím), (komu) záleží na (čem)} \dicExampleIS{Honum liggur mikið á sendingunni.} \dicExampleCS{Na té zásilce mu moc záleží.}; { \dicPhraseIS{liggja á eggjum}} \dicDirectTranslationCS{sedět na vejcích}; { \dicPhraseIS{það liggur vel á e‑m}} \dicFlx{impers} \dicDirectTranslationCS{(kdo) se dobře cítí, (kdo) je v~dobrém rozpoložení};  \dicIdiom{liggja}[fyrir]{ \dicPhraseIS{e‑að á fyrir e‑m að liggja}} \dicDirectTranslationCS{(co) na (koho) čeká, (co) má (kdo) před sebou}; { \dicPhraseIS{e‑að liggur fyrir e‑m}} \dicDirectTranslationCS{(co) na (koho) čeká, (co) leží před (kým) (rozhodnutí ap.)}; { \dicPhraseIS{liggja fyrir}} \dicDirectTranslationCS{ležet v~posteli};  \dicIdiom{liggja}[í]{ \dicPhraseIS{liggja í e‑u}} \dicDirectTranslationCS{ležet s~(čím) (chřipkou ap.)} \dicExampleIS{liggja í flensu} \dicExampleCS{ležet s~chřipkou}; { \dicPhraseIS{e‑að liggur í loftinu}} \dicLangCat{přen.} \dicDirectTranslationCS{(co) visí ve vzduchu, (co) se vznáší ve vzduchu (napětí ap.)}; { \dicPhraseIS{liggja í því}} \dicDirectTranslationCS{prohrát, dostat nakládačku};  \dicIdiom{liggja}[með]{ \dicPhraseIS{liggja með e‑m}} \dicDirectTranslationCS{spát s~(kým)};  \dicIdiom{liggja}[nærri]{ \dicPhraseIS{það lá nærri e‑u}} \dicFlx{impers} \dicDirectTranslationCS{schylovalo se k~(čemu), (co) bylo na spadnutí};  \dicIdiom{liggja}[undir]{ \dicPhraseIS{liggja undir grun}} \dicDirectTranslationCS{čelit podezření, být v~podezření};  \dicIdiom{liggja}[úti]{ \dicPhraseIS{liggja úti}} \dicDirectTranslationCS{spát venku, spát pod širým nebem};  \dicIdiom{liggja}[við]{ \dicPhraseIS{e‑m liggur við e‑u}} \dicFlx{impers} \dicDirectTranslationCS{(kdo) málem (co) (upadl ap.)}; { \dicPhraseIS{e‑að liggur beint við}} \dicDirectTranslationCS{(co) je přirozené}; { \dicPhraseIS{ef\,/\addthin þegar mikið liggur við}} \dicDirectTranslationCS{pokud to je nutné}; { \dicPhraseIS{það lá við e‑u}} \dicFlx{impers} \dicDirectTranslationCS{schylovalo se k~(čemu)};  \dicIdiom{liggja}[yfir]{ \dicPhraseIS{liggja yfir e‑u}} \dicDirectTranslationCS{ležet nad (čím), ležet v~(čem) (v~učení ap.)};  \dicIdiom{liggja}{ \dicPhraseIS{liggja í augum uppi}} \dicDirectTranslationCS{být jasný\,/\addthin zřejmý\,/\addthin očividný, bít do očí}
\dicEntry[lilja] \dicTerm{lilj|a} \dicIPA{{l}{\textsci}{l}{j}{a}} \dicPos{f}[1] \dicFlx{(‑u, ‑ur)}[7] \dicFieldCat{bot.} \dicDirectTranslationCS{lilie} \textit{(l.~{\textLA{Lilium}})}  \dicsymPhoto\ 
\dicFigure{1249.jpg}{Lilja}{Lilja - Zicha Ondřej, Biolib, Copyright/CC-BY-NC}
\dicEntry[lim] \dicTerm{lim} \dicIPA{{l}{\textsci}{\textlengthmark}{\textsubring{m}}} \dicPos{f}[4] \dicFlx{(‑ar, ‑ar)}[1] \dicDirectTranslationCS{listoví}
\dicEntry[limaburður] \dicTerm{lima··burð|ur} \dicIPA{{l}{\textsci}{\textlengthmark}{m}{a}{\textsubring{b}}{\textscy}{r}{ð}{\textscy}{\textsubring{r}}} \dicPos{m}[10] \dicFlx{(‑ar)}[7] \dicDirectTranslationCS{postoj, póza} \dicIndirectTranslationCS{(držení těla při stání)}
\dicEntry[limgerði] \dicTerm{lim··gerði} \dicIPA{{l}{\textsci}{m}{\r{\textObardotlessj}}{\textepsilon}{r}{ð}{\textsci}} \dicPos{n}[2] \dicFlx{(‑s, ‑)}[14] \dicDirectTranslationCS{živý plot}
\dicEntry[limlesta] \dicTerm{lim··lest|a} \dicIPA{{l}{\textsci}{m}{l}{\textepsilon}{s}{\textsubring{d}}{a}} \dicPos{v}[2] \dicFlx{(‑i, ‑)}[12] \dicFlx{acc} \dicDirectTranslationCS{zmrzačit, polámat kosti}
\dicEntry[limra] \dicTerm{limr|a} \dicIPA{{l}{\textsci}{m}{r}{a}} \dicPos{f}[1] \dicFlx{(‑u, ‑ur)}[7] \dicFieldCat{lit.} \dicDirectTranslationCS{limerik} \dicIndirectTranslationCS{(zábavný poetický útvar)}
\dicEntry[limur] \dicTerm{lim|ur} \dicIPA{{l}{\textsci}{\textlengthmark}{m}{\textscy}{\textsubring{r}}} \dicPos{m}[9] \dicFlx{(‑s, ‑ir)}[8] \textbf{1.} \dicSynonym{útlimur} \dicDirectTranslationCS{úd, končetina} \dicExampleIS{limir líkamans} \dicExampleCS{údy těla}  \textbf{2.} \dicSynonym{getnaðarlimur} \dicDirectTranslationCS{pyj}
\dicEntry[lina] \dicTerm{lin|a} \dicIPA{{l}{\textsci}{\textlengthmark}{n}{a}} \dicPos{v}[1] \dicFlx{(‑aði)}[1] \dicFlx{acc} \dicDirectTranslationCS{(z)mírnit, (z)tlumit, (z)měkčit} \dicExampleIS{Kærleikurinn linar sorg hennar.} \dicExampleCS{Láska mírní její žal.};  \dicIdiom{linast}{ \dicPhraseIS{linast}} \dicFlx{refl} {\textbf{a.}} \dicSynonym*{mýkjast} \dicDirectTranslationCS{(z)měknout} \dicIndirectTranslationCS{(stát\,/\addthin stávat se měkkým)};  {\textbf{b.}} \dicSynonym*{slá af fyrri skoðun} \dicDirectTranslationCS{obměkčit se, ustoupit};  {\textbf{c.}} \dicSynonym*{missa krafta} \dicDirectTranslationCS{(ze)slábnout, ztratit\,/\addthin ztrácet sílu}
\dicEntry[lind] \dicTerm{lind} \dicsymFrequent\  \dicIPA{{l}{\textsci}{n}{\textsubring{d}}} \dicPos{f}[7] \dicFlx{(‑ar, ‑ir)}[1] \textbf{1.} \dicSynonym{uppspretta} \dicDirectTranslationCS{zřídlo, pramen} \dicExampleIS{Hún vísaði þyrstum hermönnum á lind.} \dicExampleCS{Ukázala žíznivým vojákům cestu k~pramenu.}  \textbf{2.} \dicLangCat{přen.} \dicSynonym{upptök} \dicDirectTranslationCS{zřídlo, původ}  \textbf{3.} \dicFieldCat{bot.} \dicSynonym{linditré} \dicDirectTranslationCS{lípa} \textit{(l.~{\textLA{Tilia}})}  \dicsymPhoto\ 
\dicFigure{16367.jpg}{Lind}{Lind - Hanzlík Václav, Biolib, Copyright/CC-BY-NC}
\dicEntry[lindá] \dicTerm{lind··á} \dicIPA{{l}{\textsci}{n}{\textsubring{d}}{au}} \dicPos{f}[4] \dicFlx{(‑r, ‑r)}[18] \dicFieldCat{geol.} \dicDirectTranslationCS{řeka napájená zřídly}
\dicEntry[lindi] \dicTerm{lind|i} \dicIPA{{l}{\textsci}{n}{\textsubring{d}}{\textsci}} \dicPos{m}[1] \dicFlx{(‑a, ‑ar)}[1] \dicSynonym{band} \dicDirectTranslationCS{opasek, pás}
\dicEntry[linditré] \dicTerm{lindi··tré} \dicIPA{{l}{\textsci}{n}{\textsubring{d}}{\textsci}{t\smash{\textsuperscript{h}}}{r}{j}{\textepsilon}} \dicPos{n}[2] \dicFlx{(‑s, ‑)}[36] \dicFieldCat{bot.} \dicSynonym{lind} \dicDirectTranslationCS{lípa} \textit{(l.~{\textLA{Tilia}})}
\dicEntry[lindýr] \dicTerm{lin··dýr} \dicIPA{{l}{\textsci}{n}{\textsubring{d}}{i}{\textsubring{r}}} \dicPos{n}[2] \dicFlx{(‑s, ‑)}[5] \dicFieldCat{zool.} \dicDirectTranslationCS{měkkýš} \textit{(l.~{\textLA{Mollusca}})}
\dicEntry[lingerður] \dicTerm{lin··gerður} \dicIPA{{l}{\textsci}{n}{\r{\textObardotlessj}}{\textepsilon}{r}{ð}{\textscy}{\textsubring{r}}} \dicPos{adj}[2]\dicFlx{}[-1] \textbf{1.} \dicSynonym*{kveifarlegur} \dicDirectTranslationCS{změkčilý}  \textbf{2.} \dicSynonym{veikburða} \dicDirectTranslationCS{zesláblý, neduživý}
\dicEntry[linka] \dicTerm{link|a} \dicIPA{{l}{i}{\r{\ng}}{\r{g}}{a}} \dicPos{f}[1] \dicFlx{(‑u)}[5] \textbf{1.} \dicSynonym{lasleiki} \dicDirectTranslationCS{churavost, neduživost}  \textbf{2.} \dicSynonym{dugleysi} \dicDirectTranslationCS{neschopnost, nezpůsobilost}
\dicEntry[linmæli] \dicTerm{lin··mæli} \dicIPA{{l}{\textsci}{n}{m}{a}{i}{l}{\textsci}} \dicPos{n}[2] \dicFlx{(‑s)}[20] \dicFieldCat{jaz.} \dicIndirectTranslationCS{druh výslovnosti -- používání neaspirované plozivy ([\textsubring{b}], [\textsubring{d}], [\r{g}], [\r{\textObardotlessj}]) uprostřed slova po dlouhé samohlásce (standardní výslovnost)} \dicAntonym{harðmæli}
\dicEntry[linna] \dicTerm{linn|a} \dicIPA{{l}{\textsci}{n}{\textlengthmark}{a}} \dicPos{v}[2] \dicFlx{(‑ti, ‑t)}[80] \dicPhraseIS{e‑u linnir} \dicFlx{impers} \dicDirectTranslationCS{(co) přestává\,/\addthin ustává\,/\addthin končí} \dicExampleIS{Óveðrinu er tekið að linna.} \dicExampleCS{Bouře pomalu ustává.};  \dicPhraseIS{linna ekki látum} \dicDirectTranslationCS{nepřestat, neustat (s~prosbami ap.)}
\dicEntry[linnulaus] \dicTerm{linnu··laus} \dicIPA{{l}{\textsci}{n}{\textlengthmark}{\textscy}{l}{\oe i}{s}} \dicPos{adj}[5]\dicFlx{}[-1] \dicSynonym{sífelldur} \dicDirectTranslationCS{nepřetržitý, neustálý, trvalý} \dicExampleIS{linnulausar árásir} \dicExampleCS{nepřetržité útoky}
\dicEntry[linsa] \dicTerm{lins|a} \dicIPA{{l}{\textsci}{n}{s}{a}} \dicPos{f}[1] \dicFlx{(‑u, ‑ur)}[7] \textbf{1.} \dicFieldCat{fyz.} \dicDirectTranslationCS{čočka}  \textbf{2.} \dicSynonym*{augnlinsa} \dicDirectTranslationCS{(kontaktní) čočka}
\dicEntry[linsoðinn] \dicTerm{lin··soðinn} \dicIPA{{l}{\textsci}{n}{s}{\textopeno}{ð}{\textsci}{\textsubring{n}}} \dicPos{adj}[6]\dicFlx{}[-2] \dicDirectTranslationCS{(jsoucí) naměkko} \dicExampleIS{linsoðið egg} \dicExampleCS{vejce uvařené naměkko}
\dicEntry[linsubaun] \dicTerm{linsu··baun} \dicIPA{{l}{\textsci}{n}{s}{\textscy}{\textsubring{b}}{\oe i}{\textsubring{n}}} \dicPos{f}[7] \dicFlx{(‑ar, ‑ir)}[1] \dicFieldCat{bot.} \dicDirectTranslationCS{čočka} \textit{(l.~{\textLA{Lens culinaris}})}
\dicEntry[linun] \dicTerm{lin|un} \dicIPA{{l}{\textsci}{\textlengthmark}{n}{\textscy}{\textsubring{n}}} \dicPos{f}[7] \dicFlx{(‑unar)}[9] \textbf{1.} \dicSynonym{fró} \dicDirectTranslationCS{uklidnění, úleva}  \textbf{2.} \dicSynonym{mildun} \dicDirectTranslationCS{zmírnění} \dicExampleIS{linun refsingar} \dicExampleCS{zmírnění trestu}
\dicEntry[linur] \dicTerm{linur} \dicIPA{{l}{\textsci}{\textlengthmark}{n}{\textscy}{\textsubring{r}}} \dicPos{adj}[1]\dicFlx{}[-1] \textbf{1.} \dicSynonym{mjúkur} \dicDirectTranslationCS{hebký, měkký, jemný}  \textbf{2.} \dicSynonym{slakur} \dicDirectTranslationCS{slabý, chabý} \dicExampleIS{linur í fótbolta} \dicExampleCS{slabý ve fotbale}  \textbf{3.} \dicSynonym{slappur} \dicDirectTranslationCS{indisponovaný, slabý}
\dicEntry[lipur] \dicTerm{lipur} \dicIPA{{l}{\textsci}{\textlengthmark}{\textsubring{b}}{\textscy}{\textsubring{r}}} \dicPos{adj}[9] \dicFlx{(f ‑)}[1] \textbf{1.} \dicSynonym*{flinkur} \dicDirectTranslationCS{zručný, šikovný} \dicExampleIS{Hún er lipur í höndunum.} \dicExampleCS{Má šikovné ruce.}  \textbf{2.} \dicSynonym{þægilegur} \dicDirectTranslationCS{vstřícný, přívětivý}  \textbf{3.} \dicSynonym{meðfærilegur} \dicDirectTranslationCS{(uživatelsky) přívětivý}
\dicEntry[lipurð] \dicTerm{lipurð} \dicIPA{{l}{\textsci}{\textlengthmark}{\textsubring{b}}{\textscy}{r}{\texttheta}} \dicPos{f}[4] \dicFlx{(‑ar)}[3] \textbf{1.} \dicSynonym*{fimleiki} \dicDirectTranslationCS{obratnost, zručnost}  \textbf{2.} \dicSynonym*{lipur í viðmóti} \dicDirectTranslationCS{vstřícnost, přívětivost}
\dicEntry[lirfa] \dicTerm{lirf|a} \dicIPA{{l}{\textsci}{r}{v}{a}} \dicPos{f}[1] \dicFlx{(‑u, ‑ur)}[19] \dicDirectTranslationCS{larva}
\dicEntry[list] \dicTerm{list} \dicsymFrequent\  \dicIPA{{l}{\textsci}{s}{\textsubring{d}}} \dicPos{f}[7] \dicFlx{(‑ar, ‑ir)}[1] \textbf{1.} \dicDirectTranslationCS{umění} \dicIndirectTranslationCS{(odvětví lidské tvorby)} \dicExampleIS{unna listum} \dicExampleCS{zbožňovat umění}  \textbf{2.} \dicSynonym{færni} \dicDirectTranslationCS{umění, um, dovednost}
\dicEntry[listaháskóli] \dicTerm{lista··há·skól|i} \dicIPA{{l}{\textsci}{s}{\textsubring{d}}{a}{h}{au}{s}{\r{g}}{ou}{l}{\textsci}} \dicPos{m}[1] \dicFlx{(‑a, ‑ar)}[1] \dicFieldCat{škol.} \dicDirectTranslationCS{akademie umění}
\dicEntry[listahátíð] \dicTerm{lista··há·tíð} \dicIPA{{l}{\textsci}{s}{\textsubring{d}}{a}{h}{au}{t\smash{\textsuperscript{h}}}{i}{\texttheta}} \dicPos{f}[7] \dicFlx{(‑ar, ‑ir)}[1] \dicDirectTranslationCS{umělecký festival}
\dicEntry[listakona] \dicTerm{lista··kon|a} \dicIPA{{l}{\textsci}{s}{\textsubring{d}}{a}{k\smash{\textsuperscript{h}}}{\textopeno}{n}{a}} \dicPos{f}[1] \dicFlx{(‑u, ‑ur)}[27] \dicDirectTranslationCS{umělkyně}
\dicEntry[listamaður] \dicTerm{lista··|maður} \dicsymFrequent\  \dicIPA{{l}{\textsci}{s}{\textsubring{d}}{a}{m}{a}{ð}{\textscy}{\textsubring{r}}} \dicPos{m}[13] \dicFlx{(‑manns, ‑menn)}[2] \dicDirectTranslationCS{umělec, umělkyně} \dicExampleIS{Það er eitt dáðasta verk listamannsins.} \dicExampleCS{Je to jedno z~nejvíce obdivovaných děl umělce.}
\dicEntry[listasafn] \dicTerm{lista··|safn} \dicIPA{{l}{\textsci}{s}{\textsubring{d}}{a}{s}{a}{\textsubring{b}}{\textsubring{n}}} \dicPos{n}[2] \dicFlx{(‑safns, ‑söfn)}[8] \dicDirectTranslationCS{galerie (umění)}
\dicEntry[listasaga] \dicTerm{lista··|saga} \dicIPA{{l}{\textsci}{s}{\textsubring{d}}{a}{s}{a}{\textbabygamma}{a}} \dicPos{f}[1] \dicFlx{(‑sögu)}[2] \dicDirectTranslationCS{dějiny umění}
\dicEntry[listaskóli] \dicTerm{lista··skól|i} \dicIPA{{l}{\textsci}{s}{\textsubring{d}}{a}{s}{\r{g}}{ou}{l}{\textsci}} \dicPos{m}[1] \dicFlx{(‑a, ‑ar)}[1] \dicFieldCat{škol.} \dicDirectTranslationCS{umělecká škola}
\dicEntry[listaverk] \dicTerm{lista··verk} \dicIPA{{l}{\textsci}{s}{\textsubring{d}}{a}{v}{\textepsilon}{\textsubring{r}}{\r{g}}} \dicPos{n}[2] \dicFlx{(‑s, ‑)}[5] \dicDirectTranslationCS{umělecké dílo}
\dicEntry[listdans] \dicTerm{list··dans} \dicIPA{{l}{\textsci}{s}{\textsubring{d}}{a}{n}{s}} \dicPos{m}[4] \dicFlx{(‑, ‑ar)}[29] \dicSynonym{ballett} \dicDirectTranslationCS{umělecký tanec, balet}
\dicEntry[listfengur] \dicTerm{list··fengur} \dicIPA{{l}{\textsci}{s}{\textsubring{d}}{f}{ei}{\ng}{\r{g}}{\textscy}{\textsubring{r}}} \dicPos{adj}[1]\dicFlx{}[-1] \dicDirectTranslationCS{umělecky nadaný}
\dicEntry[listfræðingur] \dicTerm{list·fræð··ing|ur} \dicIPA{{l}{\textsci}{s}{\textsubring{d}}{f}{r}{a}{i}{ð}{i}{\ng}{\r{g}}{\textscy}{\textsubring{r}}} \dicPos{m}[6] \dicFlx{(‑s, ‑ar)}[8] \dicDirectTranslationCS{historik\,/\addthin historička umění, kunsthistorik, kunsthistorička}
\dicEntry[listgáfa] \dicTerm{list··gáf|a} \dicIPA{{l}{\textsci}{s}{\textsubring{d}}{\r{g}}{au}{v}{a}} \dicPos{f}[1] \dicFlx{(‑u, ‑ur)}[13] \dicDirectTranslationCS{umělecký talent}
\dicEntry[listi] \dicTerm{list|i} \dicsymFrequent\  \dicIPA{{l}{\textsci}{s}{\textsubring{d}}{\textsci}} \dicPos{m}[1] \dicFlx{(‑a, ‑ar)}[1] \textbf{1.} \dicSynonym{skrá\smash{\textsuperscript{1}}} \dicDirectTranslationCS{seznam, soupis, listina, rejstřík, žebříček} \dicExampleIS{óskalisti} \dicExampleCS{seznam přání};  \dicPhraseIS{listi yfir e‑n\,/\addthin e‑ð} \dicDirectTranslationCS{seznam (koho\,/\addthin čeho)}  \textbf{2.} \dicSynonym{ræma} \dicDirectTranslationCS{lišta, laťka}
\dicEntry[listiðn] \dicTerm{list··iðn} \dicIPA{{l}{\textsci}{s}{\textsubring{d}}{\textsci}{ð}{\textsubring{n}}} \dicPos{f}[7] \dicFlx{(‑ar, ‑ir)}[1] \dicDirectTranslationCS{umělecké řemeslo}
\dicEntry[listiðnaður] \dicTerm{list··iðn·að|ur} \dicIPA{{l}{\textsci}{s}{\textsubring{d}}{\textsci}{ð}{n}{a}{ð}{\textscy}{\textsubring{r}}} \dicPos{m}[10] \dicFlx{(‑ar)}[9] \dicDirectTranslationCS{užité umění}
\dicEntry[listmálari] \dicTerm{list··mál·ar|i} \dicIPA{{l}{\textsci}{s}{\textsubring{d}}{m}{au}{l}{a}{r}{\textsci}} \dicPos{m}[1] \dicFlx{(‑a, ‑ar)}[13] \dicDirectTranslationCS{malíř(ka)} \dicIndirectTranslationCS{(umělec malující obrazy)}
\dicEntry[listnám] \dicTerm{list··nám} \dicIPA{{l}{\textsci}{s}{\textsubring{d}}{n}{au}{\textsubring{m}}} \dicPos{n}[2] \dicFlx{(‑s)}[2] \dicDirectTranslationCS{studium umění}
\dicEntry[listrænn] \dicTerm{list··rænn} \dicIPA{{l}{\textsci}{s}{\textsubring{d}}{r}{a}{i}{\textsubring{d}}{\textsubring{n}}} \dicPos{adj}[7]\dicFlx{}[-1] \dicDirectTranslationCS{umělecký}
\dicEntry[listskautar] \dicTerm{list··skautar} \dicIPA{{l}{\textsci}{s}{\textsubring{d}}{s}{\r{g}}{\oe i}{\textsubring{d}}{a}{\textsubring{r}}} \dicPos{m}[1] \dicFlx{pl}[2] \dicFieldCat{sport.} \dicDirectTranslationCS{krasobruslení}
\dicEntry[listunnandi] \dicTerm{list··unn·|andi} \dicIPA{{l}{\textsci}{s}{\textsubring{d}}{\textscy}{n}{a}{n}{\textsubring{d}}{\textsci}} \dicPos{m}[2] \dicFlx{(‑anda, ‑endur)}[1] \dicDirectTranslationCS{milovník\,/\addthin milovnice umění}
\dicEntry[lita] \dicTerm{lit|a} \dicsymFrequent\  \dicIPA{{l}{\textsci}{\textlengthmark}{\textsubring{d}}{a}} \dicPos{v}[1] \dicFlx{(‑aði)}[1] \dicFlx{acc} \textbf{1.} \dicSynonym{mála} \dicDirectTranslationCS{(na)barvit, nabarvovat, obarvit, obarvovat} \dicExampleIS{lita fötin} \dicExampleCS{nabarvit oblečení}  \textbf{2.} \dicDirectTranslationCS{zabarvit, zkreslit (skutečnost ap.)} \dicExampleIS{Skáldskapurinn litar frá sér.} \dicExampleCS{Literární dílo zkresluje.};  \dicIdiom{litast}{ \dicPhraseIS{litast}} \dicFlx{refl} \dicDirectTranslationCS{zbarvit se, zabarvit se, obarvit se};  \dicIdiom{litast}[um]{ \dicPhraseIS{litast um}} \dicFlx{refl} \dicDirectTranslationCS{(po)rozhlédnout se};  \dicIdiom{litast}[upp]{ \dicPhraseIS{litast upp}} \dicFlx{refl} \dicSynonym*{tapa lit} \dicDirectTranslationCS{(vy)blednout, ztratit\,/\addthin ztrácet barvu}
\dicEntry[litaraft] \dicTerm{litar··aft} \dicIPA{{l}{\textsci}{\textlengthmark}{\textsubring{d}}{a}{r}{a}{f}{\textsubring{d}}} \dicPos{n}[2] \dicFlx{(‑s)}[2] \dicDirectTranslationCS{barva\,/\addthin zbarvení obličeje}
\dicEntry[litarháttur] \dicTerm{litar··hátt|ur} \dicIPA{{l}{\textsci}{\textlengthmark}{\textsubring{d}}{a}{\textsubring{r}}{h}{au}{h}{\textsubring{d}}{\textscy}{\textsubring{r}}} \dicPos{m}[12] \dicFlx{(‑ar)}[8] \dicDirectTranslationCS{barva pleti} \dicExampleIS{fordómar vegna litarháttar} \dicExampleCS{předsudky k~barvě pleti}
\dicEntry[litasamsetning] \dicTerm{lita··sam·set·ning} \dicIPA{{l}{\textsci}{\textlengthmark}{\textsubring{d}}{a}{s}{a}{m}{s}{\textepsilon}{h}{\textsubring{d}}{n}{i}{\ng}{\r{g}}} \dicPos{f}[4] \dicFlx{(‑ar, ‑ar)}[5] \dicDirectTranslationCS{barevná kombinace, barevné schéma}
\dicEntry[litaspjald] \dicTerm{lita··|spjald} \dicIPA{{l}{\textsci}{\textlengthmark}{\textsubring{d}}{a}{s}{\textsubring{b}}{j}{a}{l}{\textsubring{d}}} \dicPos{n}[2] \dicFlx{(‑spjalds, ‑spjöld)}[8] \textbf{1.} \dicDirectTranslationCS{paleta} \dicIndirectTranslationCS{(malířská deska k~míchání a~roztírání barev)}  \textbf{2.} \dicFieldCat{poč.} \dicDirectTranslationCS{paleta} \dicIndirectTranslationCS{(na výběr barev)}
\dicEntry[litblinda] \dicTerm{lit··blind|a} \dicIPA{{l}{\textsci}{\textlengthmark}{\textsubring{d}}{\textsubring{b}}{l}{\textsci}{n}{\textsubring{d}}{a}} \dicPos{f}[1] \dicFlx{(‑u)}[5] \dicDirectTranslationCS{barvoslepost}
\dicEntry[litblindur] \dicTerm{lit··blindur} \dicIPA{{l}{\textsci}{\textlengthmark}{\textsubring{d}}{\textsubring{b}}{l}{\textsci}{n}{\textsubring{d}}{\textscy}{\textsubring{r}}} \dicPos{adj}[2]\dicFlx{}[-14] \dicFieldCat{med.} \dicDirectTranslationCS{barvoslepý}
\dicEntry[litbrigði] \dicTerm{lit··brigði} \dicIPA{{l}{\textsci}{\textlengthmark}{\textsubring{d}}{\textsubring{b}}{r}{\textsci}{\textbabygamma}{ð}{\textsci}} \dicPos{n}[2] \dicFlx{(‑s, ‑)}[14] \dicDirectTranslationCS{odstín (barvy)}
\dicEntry[litgreining] \dicTerm{lit··grein·ing} \dicIPA{{l}{\textsci}{\textlengthmark}{\textsubring{d}}{\r{g}}{r}{ei}{n}{i}{\ng}{\r{g}}} \dicPos{f}[4] \dicFlx{(‑ar)}[7] \textbf{1.} \dicFieldCat{polygr.} \dicDirectTranslationCS{separace barev}  \textbf{2.} \dicDirectTranslationCS{barevná analýza} \dicIndirectTranslationCS{(osobní kombinace barev)}
\dicEntry[Litháen] \dicTerm{Litháen} \dicIPA{{l}{\textsci}{\textlengthmark}{t\smash{\textsuperscript{h}}}{au}{\textepsilon}{\textsubring{n}}} \dicPos{n}[2] \dicFlx{(‑s)}[32] \dicFieldCat{geog.} \dicDirectTranslationCS{Litva}
\dicEntry[Lithái] \dicTerm{Lithá|i} \dicIPA{{l}{\textsci}{\textlengthmark}{t\smash{\textsuperscript{h}}}{au}{\textsci}} \dicPos{m}[1] \dicFlx{(‑a, ‑ar)}[1] \dicDirectTranslationCS{Litevec, Litevka}
\dicEntry[litháíska] \dicTerm{litháísk|a} \dicIPA{{l}{\textsci}{\textlengthmark}{t\smash{\textsuperscript{h}}}{au}{i}{s}{\r{g}}{a}} \dicPos{f}[1] \dicFlx{(‑u)}[5] \dicLink{litháska}
\dicEntry[litháískur] \dicTerm{litháískur} \dicIPA{{l}{\textsci}{\textlengthmark}{t\smash{\textsuperscript{h}}}{au}{i}{s}{\r{g}}{\textscy}{\textsubring{r}}} \dicPos{adj}[1]\dicFlx{}[-6] \dicLink{litháskur}
\dicEntry[litháska] \dicTerm{lithásk|a}\dicTerm{, litháíska} \dicIPA{{l}\-{\textsci}\-{\textlengthmark}\-{t\smash{\textsuperscript{h}}}\-{au}\-{s}\-{\r{g}}\-{a}\-} \dicPos{f}[1] \dicFlx{(‑u)}[5] \dicDirectTranslationCS{litevština}
\dicEntry[litháskur] \dicTerm{litháskur}\dicTerm{, litháískur} \dicIPA{{l}\-{\textsci}\-{\textlengthmark}\-{t\smash{\textsuperscript{h}}}\-{au}\-{s}\-{\r{g}}\-{\textscy}\-{\textsubring{r}}\-} \dicPos{adj}[1]\dicFlx{}[-6] \dicDirectTranslationCS{litevský}
\dicEntry[lithimna] \dicTerm{lit··himn|a} \dicIPA{{l}{\textsci}{\textlengthmark}{\textsubring{d}}{h}{\textsci}{m}{n}{a}} \dicPos{f}[1] \dicFlx{(‑u, ‑ur)}[7] \dicFieldCat{anat.} \dicDirectTranslationCS{duhovka}
\dicEntry[liti] \dicTerm{liti} \dicIPA{{l}{\textsci}{\textlengthmark}{\textsubring{d}}{\textsci}} \dicPos{v} \dicFlx{con pf sg 1 pers} \dicLink{líta}
\dicEntry[litið] \dicTerm{litið} \dicIPA{{l}{\textsci}{\textlengthmark}{\textsubring{d}}{\textsci}{\texttheta}} \dicPos{v} \dicFlx{supin} \dicLink{líta}
\dicEntry[litlafingurs] \dicTerm{litla··fingurs} \dicIPA{{l}{\textsci}{h}{\textsubring{d}}{l}{a}{f}{\textsci}{\ng}{\r{g}}{\textscy}{\textsubring{r}}{s}} \dicPos{m} \dicFlx{sg gen} \dicLink{litlifingur}
\dicEntry[litlatá] \dicTerm{litla··tá} \dicIPA{{l}{\textsci}{h}{\textsubring{d}}{l}{a}{t\smash{\textsuperscript{h}}}{au}} \dicPos{f}[9] \dicFlx{(litlutáar, litlutær)}[2] \dicDirectTranslationCS{malík\,/\addthin malíček (u~nohy)}
\dicEntry[litlaus] \dicTerm{lit··laus} \dicIPA{{l}{\textsci}{\textlengthmark}{\textsubring{d}}{l}{\oe i}{s}} \dicPos{adj}[5]\dicFlx{}[-1] \textbf{1.} \dicSynonym*{ekki í lit} \dicDirectTranslationCS{bezbarvý}  \textbf{2.} \dicSynonym{dauflegur} \dicDirectTranslationCS{mdlý, fádní, bezbarvý} \dicExampleIS{litlaus frásögn} \dicExampleCS{bezbarvé vyprávění}
\dicEntry[litlifingur] \dicTerm{litli··fingur} \dicIPA{{l}{\textsci}{h}{\textsubring{d}}{l}{\textsci}{f}{i}{\ng}{\r{g}}{\textscy}{\textsubring{r}}} \dicPos{m}[5] \dicFlx{(litlafingurs, litlufingur)}[14] \dicDirectTranslationCS{malík\,/\addthin malíček (u~ruky)}
\dicEntry[litlufingur] \dicTerm{litlu··fingur} \dicIPA{{l}{\textsci}{h}{\textsubring{d}}{l}{\textscy}{f}{\textsci}{\ng}{\r{g}}{\textscy}{\textsubring{r}}} \dicPos{m} \dicFlx{pl nom} \dicLink{litlifingur}
\dicEntry[litmynd] \dicTerm{lit··mynd} \dicIPA{{l}{\textsci}{\textlengthmark}{\textsubring{d}}{m}{\textsci}{n}{\textsubring{d}}} \dicPos{f}[7] \dicFlx{(‑ar, ‑ir)}[1] \dicDirectTranslationCS{barevný obraz, barevná fotografie}
\dicEntry[litningur] \dicTerm{litn··ing|ur} \dicIPA{{l}{\textsci}{h}{\textsubring{d}}{n}{i}{\ng}{\r{g}}{\textscy}{\textsubring{r}}} \dicPos{m}[6] \dicFlx{(‑s, ‑ar)}[8] \dicFieldCat{biol.} \dicDirectTranslationCS{chromozom}
\dicEntry[litprentun] \dicTerm{lit··prent|un} \dicIPA{{l}{\textsci}{\textlengthmark}{\textsubring{d}}{p\smash{\textsuperscript{h}}}{r}{\textepsilon}{\textsubring{n}}{\textsubring{d}}{\textscy}{\textsubring{n}}} \dicPos{f}[7] \dicFlx{(‑unar)}[9] \dicFieldCat{polygr.} \dicDirectTranslationCS{barevný tisk}
\dicEntry[litríkur] \dicTerm{lit··ríkur} \dicIPA{{l}{\textsci}{\textlengthmark}{\textsubring{d}}{r}{i}{\r{g}}{\textscy}{\textsubring{r}}} \dicPos{adj}[1]\dicFlx{}[-1] \textbf{1.} \dicDirectTranslationCS{pestrobarevný, pestrý} \dicExampleIS{litrík föt} \dicExampleCS{pestrobarevné oblečení}  \textbf{2.} \dicDirectTranslationCS{pestrý, rozmanitý} \dicExampleIS{Söngkonan á litríkan feril að baki.} \dicExampleCS{Zpěvačka má za sebou pestrou kariéru.}
\dicEntry[litróf] \dicTerm{lit··róf} \dicIPA{{l}{\textsci}{\textlengthmark}{\textsubring{d}}{r}{ou}{f}} \dicPos{n}[2] \dicFlx{(‑s, ‑)}[5] \textbf{1.} \dicFieldCat{fyz.} \dicDirectTranslationCS{barevné spektrum}  \textbf{2.} \dicLangCat{přen.} \dicDirectTranslationCS{spektrum (názorové ap.)}
\dicEntry[litskrúð] \dicTerm{lit··skrúð} \dicIPA{{l}{\textsci}{\textlengthmark}{\textsubring{d}}{s}{\r{g}}{r}{u}{\texttheta}} \dicPos{n}[2] \dicFlx{(‑s)}[2] \dicDirectTranslationCS{pestrobarevnost}
\dicEntry[litskyggna] \dicTerm{lit··skyggn|a} \dicIPA{{l}{\textsci}{\textlengthmark}{\textsubring{d}}{s}{\r{\textObardotlessj}}{\textsci}{\r{g}}{n}{a}} \dicPos{f}[1] \dicFlx{(‑u, ‑ur)}[7] \dicDirectTranslationCS{barevný diapozitiv}
\dicEntry[litum] \dicTerm{litum} \dicIPA{{l}{\textsci}{\textlengthmark}{\textsubring{d}}{\textscy}{\textsubring{m}}} \dicPos{v} \dicFlx{ind pf pl 1 pers} \dicLink{líta}
\dicEntry[litun] \dicTerm{lit|un} \dicIPA{{l}{\textsci}{\textlengthmark}{\textsubring{d}}{\textscy}{\textsubring{n}}} \dicPos{f}[7] \dicFlx{(‑unar)}[9] \dicDirectTranslationCS{(na)barvení}
\dicEntry[litur] \dicTerm{lit|ur} \dicsymFrequent\  \dicIPA{{l}{\textsci}{\textlengthmark}{\textsubring{d}}{\textscy}{\textsubring{r}}} \dicPos{m}[10] \dicFlx{(‑ar, ‑ir)}[13] \textbf{1.} \dicDirectTranslationCS{barva, barevný odstín} \dicExampleIS{Hvernig er það á litinn?} \dicExampleCS{Jakou to má barvu?}  \textbf{2.} \dicSynonym*{litarefni} \dicDirectTranslationCS{barva, nátěr}  \textbf{3.} \dicSynonym*{tegund spila} \dicDirectTranslationCS{barva} \dicIndirectTranslationCS{(v~kartách)}
\dicEntry[Líbani] \dicTerm{Líban|i} \dicIPA{{l}{i}{\textlengthmark}{\textsubring{b}}{a}{n}{\textsci}} \dicPos{m}[3] \dicFlx{(‑a, ‑ar\,/\addthin ‑ir)}[3] \dicDirectTranslationCS{Libanonec, Libanonka}
\dicEntry[Líbanon] \dicTerm{Líbanon} \dicIPA{{l}{i}{\textlengthmark}{\textsubring{b}}{a}{n}{\textopeno}{\textsubring{n}}} \dicPos{n}[2] \dicFlx{(‑s)}[32] \dicFieldCat{geog.} \dicDirectTranslationCS{Libanon}
\dicEntry[líbanskur] \dicTerm{líbanskur} \dicIPA{{l}{i}{\textlengthmark}{\textsubring{b}}{a}{n}{s}{\r{g}}{\textscy}{\textsubring{r}}} \dicPos{adj}[1] \dicFlx{(f líbönsk)}[3] \dicDirectTranslationCS{libanonský}
\dicEntry[Líbería] \dicTerm{Líberí|a} \dicIPA{{l}{i}{\textlengthmark}{\textsubring{b}}{\textepsilon}{r}{i}{j}{a}} \dicPos{f}[1] \dicFlx{(‑u)}[6] \dicFieldCat{geog.} \dicDirectTranslationCS{Libérie}
\dicEntry[líberískur] \dicTerm{líberískur} \dicIPA{{l}{i}{\textlengthmark}{\textsubring{b}}{\textepsilon}{r}{i}{s}{\r{g}}{\textscy}{\textsubring{r}}} \dicPos{adj}[1]\dicFlx{}[-6] \dicDirectTranslationCS{liberijský}
\dicEntry[Líberíumaður] \dicTerm{Líberíu··|maður} \dicIPA{{l}{i}{\textlengthmark}{\textsubring{b}}{\textepsilon}{r}{i}{j}{\textscy}{m}{a}{ð}{\textscy}{\textsubring{r}}} \dicPos{m}[13] \dicFlx{(‑manns, ‑menn)}[2] \dicDirectTranslationCS{Liberijec, Liberijka}
\dicEntry[Líbía] \dicTerm{Líbí|a}\dicTerm{, Líbýa} \dicIPA{{l}\-{i}\-{\textlengthmark}\-{\textsubring{b}}\-{i}\-{j}\-{a}\-} \dicPos{f}[1] \dicFlx{(‑u)}[6] \dicFieldCat{geog.} \dicDirectTranslationCS{Libye}
\dicEntry[líbískur] \dicTerm{líbískur}\dicTerm{, líbýskur} \dicIPA{{l}\-{i}\-{\textlengthmark}\-{\textsubring{b}}\-{i}\-{s}\-{\r{g}}\-{\textscy}\-{\textsubring{r}}\-} \dicPos{adj}[1]\dicFlx{}[-6] \dicDirectTranslationCS{libyjský}
\dicEntry[Líbíumaður] \dicTerm{Líbíu··|maður} \dicIPA{{l}{i}{\textlengthmark}{\textsubring{b}}{i}{j}{\textscy}{m}{a}{ð}{\textscy}{\textsubring{r}}} \dicPos{m}[13] \dicFlx{(‑manns, ‑menn)}[2] \dicDirectTranslationCS{Libyjec, Libyjka}
\dicEntry[Líbýa] \dicTerm{Líbý|a} \dicIPA{{l}{i}{\textlengthmark}{\textsubring{b}}{i}{j}{a}} \dicPos{f}[1] \dicFlx{(‑u)}[6] \dicLink{Líbía}
\dicEntry[líbýskur] \dicTerm{líbýskur} \dicIPA{{l}{i}{\textlengthmark}{\textsubring{b}}{i}{s}{\r{g}}{\textscy}{\textsubring{r}}} \dicPos{adj}[1]\dicFlx{}[-6] \dicLink{líbískur}
\dicEntry[líð] \dicTerm{líð} \dicIPA{{l}{i}{\textlengthmark}{\texttheta}} \dicPos{v} \dicFlx{ind praes sg 1 pers} \dicLink{líða}
\dicEntry[líða] \dicTerm{líða} \dicsymFrequent\  \dicIPA{{l}{i}{\textlengthmark}{ð}{a}} \dicPos{v}[6] \dicFlx{(líð, leið, liðum, liði, liðið)}[66] \dicFlx{acc\,/\addthin dat} \textbf{1.} \dicSynonym*{um tíma} \dicDirectTranslationCS{plynout, ubíhat, utíkat} \dicIndirectTranslationCS{(o~čase)} \dicExampleIS{Árið líður hratt.} \dicExampleCS{Rok rychle ubíhá.}  \textbf{2.} \dicFlx{dat} \dicDirectTranslationCS{pokračovat, postupovat} \dicExampleIS{Hvernig líður verkinu?} \dicExampleCS{Jak pokračuje práce?}  \textbf{3.} \dicFlx{acc} \dicSynonym{þola} \dicDirectTranslationCS{snést, snášet, tolerovat} \dicIndirectTranslationCS{(často se záporem)} \dicExampleIS{líða þetta ekki} \dicExampleCS{nesnést to}  \textbf{4.} \dicPhraseIS{e‑m líður} \dicFlx{impers} \dicDirectTranslationCS{(kdo) se cítí (dobře ap.)} \dicExampleIS{Mér líður vel.} \dicExampleCS{Cítím se dobře.\,/\addthin Je mi dobře.};  \dicPhraseIS{hvernig líður þér?} \dicDirectTranslationCS{jak se máš?\,/\addthin jak se cítíš?};  \dicIdiom{líða}[að]{ \dicPhraseIS{það líður að e‑u}} \dicFlx{impers} \dicDirectTranslationCS{(co) se blíží (svátky ap.)};  \dicIdiom{líða}[hjá]{ \dicPhraseIS{e‑að líður hjá}} \dicDirectTranslationCS{(co) míjí, (co) pomine};  \dicIdiom{líða}[í]{ \dicPhraseIS{líða í ómegin}} \dicDirectTranslationCS{omdlít, upadnout do mdlob};  \dicIdiom{líða}[undir]{ \dicPhraseIS{e‑að líður undir lok}} \dicDirectTranslationCS{(co) se schyluje\,/\addthin chýlí ke konci};  \dicIdiom{líða}[út af]{ \dicPhraseIS{líða út af}} {\textbf{a.}} \dicDirectTranslationCS{omdlít};  {\textbf{b.}} \dicDirectTranslationCS{usnout} \dicExampleIS{líða út af í hægindastólnum} \dicExampleCS{usnout v~křesle};  \dicIdiom{líða}[yfir]{ \dicPhraseIS{það líður yfir e‑n}} \dicFlx{impers} \dicDirectTranslationCS{(kdo) omdlívá, (kdo) ztrácí vědomí} \dicExampleIS{Það leið yfir mig.} \dicExampleCS{Omdlel jsem.};  \dicIdiom{líðast}{ \dicPhraseIS{e‑m líðst e‑að}} \dicFlx{refl impers} \dicDirectTranslationCS{(co) se (komu) toleruje} \dicIndirectTranslationCS{(často se záporem)};  \dicIdiom{líða}{ \dicPhraseIS{hvað sem öðru líður}} \dicDirectTranslationCS{beztak, když se to vezme kolem a~kolem}
\dicEntry[líðan] \dicTerm{líðan} \dicIPA{{l}{i}{\textlengthmark}{ð}{a}{\textsubring{n}}} \dicPos{f}[4] \dicFlx{(‑ar)}[3] \dicDirectTranslationCS{(zdravotní) stav} \dicExampleIS{Líðan sjúklingsins fer batnandi.} \dicExampleCS{Zdravotní stav pacienta se zlepšuje.}
\dicEntry[líðandi] \dicTerm{líð··andi} \dicIPA{{l}{i}{\textlengthmark}{ð}{a}{n}{\textsubring{d}}{\textsci}} \dicPos{adj}[13] \dicFlx{indecl}[1] \dicPhraseIS{á líðandi stund} \dicFlx{adv} \dicSynonym{núna} \dicDirectTranslationCS{teď, v~současné chvíli}
\dicEntry[líf] \dicTerm{líf} \dicsymFrequent\  \dicIPA{{l}{i}{\textlengthmark}{f}} \dicPos{n}[2] \dicFlx{(‑s, ‑)}[5] \textbf{1.} \dicSynonym*{það að lifa} \dicDirectTranslationCS{život, žití, bytí};  \dicPhraseIS{vera á lífi} \dicDirectTranslationCS{být naživu}  \textbf{2.} \dicSynonym*{lífverur} \dicDirectTranslationCS{život, existence (organismu)} \dicExampleIS{lífið á jörðinni} \dicExampleCS{život na Zemi}  \textbf{3.} \dicSynonym*{dagleg tilvera} \dicDirectTranslationCS{život (všední ap.)}  \textbf{4.} \dicSynonym{fjör} \dicDirectTranslationCS{živo, rušno};  \dicIdiom{líf}{ \dicPhraseIS{af lífi og sál}} \dicFlx{adv} \dicDirectTranslationCS{z~celých sil}; { \dicPhraseIS{blása lífið í e‑ð}} \dicLangCat{přen.} \dicDirectTranslationCS{vdechnout život (čemu)}; { \dicPhraseIS{elska e‑n út af lífinu}} \dicDirectTranslationCS{milovat (koho) z~celé duše};  \dicPhraseIS{láta lífið} \dicDirectTranslationCS{přijít o~život}; { \dicPhraseIS{svipta sig lífi}} \dicDirectTranslationCS{vzít si život}; { \dicPhraseIS{taka e‑n af lífi}} \dicDirectTranslationCS{vzít (komu) život}; { \dicPhraseIS{vekja til lífsins}} \dicLangCat{přen.} \dicDirectTranslationCS{probudit k~životu}
\dicEntry[lífbátur] \dicTerm{líf··bát|ur} \dicIPA{{l}{i}{v}{\textsubring{b}}{au}{\textsubring{d}}{\textscy}{\textsubring{r}}} \dicPos{m}[6] \dicFlx{(‑s, ‑ar)}[22] \dicDirectTranslationCS{záchranný člun}
\dicEntry[lífbein] \dicTerm{líf··bein} \dicIPA{{l}{i}{v}{\textsubring{b}}{ei}{\textsubring{n}}} \dicPos{n}[2] \dicFlx{(‑s, ‑)}[5] \dicFieldCat{anat.} \dicDirectTranslationCS{stydká kost}
\dicEntry[lífeðlisfræði] \dicTerm{líf··eðlis·fræð|i} \dicIPA{{l}{i}{\textlengthmark}{v}{\textepsilon}{ð}{l}{\textsci}{s}{f}{r}{a}{i}{ð}{\textsci}} \dicPos{f}[3] \dicFlx{(‑i)}[3] \dicDirectTranslationCS{fyziologie}
\dicEntry[lífeðlisfræðingur] \dicTerm{líf·eðlis·fræð··ing|ur} \dicIPA{{l}\-{i}\-{\textlengthmark}\-{v}\-{\textepsilon}\-{ð}\-{l}\-{\textsci}\-{s}\-{f}\-{r}\-{a}\-{i}\-{ð}\-{i}\-{\ng}\-{\r{g}}\-{\textscy}\-{\textsubring{r}}\-} \dicPos{m}[6] \dicFlx{(‑s, ‑ar)}[8] \dicDirectTranslationCS{fyziolog, fyzioložka}
\dicEntry[lífefnafræði] \dicTerm{líf··efna·fræð|i} \dicIPA{{l}{i}{\textlengthmark}{v}{\textepsilon}{\textsubring{b}}{n}{a}{f}{r}{a}{i}{ð}{\textsci}} \dicPos{f}[3] \dicFlx{(‑i)}[3] \dicDirectTranslationCS{biochemie}
\dicEntry[lífefnafræðingur] \dicTerm{líf·efna·fræð··ing|ur} \dicIPA{{l}\-{i}\-{\textlengthmark}\-{v}\-{\textepsilon}\-{\textsubring{b}}\-{n}\-{a}\-{f}\-{r}\-{a}\-{i}\-{ð}\-{i}\-{\ng}\-{\r{g}}\-{\textscy}\-{\textsubring{r}}\-} \dicPos{m}[6] \dicFlx{(‑s, ‑ar)}[8] \dicDirectTranslationCS{biochemik, biochemička}
\dicEntry[líferni] \dicTerm{líf··erni} \dicIPA{{l}{i}{\textlengthmark}{v}{\textepsilon}{r}{\textsubring{d}}{n}{\textsci}} \dicPos{n}[2] \dicFlx{(‑s)}[20] \dicDirectTranslationCS{životospráva, životní styl}
\dicEntry[lífeyrir] \dicTerm{líf··eyr|ir} \dicIPA{{l}{i}{\textlengthmark}{v}{ei}{r}{\textsci}{\textsubring{r}}} \dicPos{m}[7] \dicFlx{(‑is)}[2] \dicDirectTranslationCS{důchod, penze} \dicIndirectTranslationCS{(pravidelný příjem osob plynoucí z~pojištění)} \dicExampleIS{taka lífeyris} \dicExampleCS{příjem důchodu}
\dicEntry[lífeyrisgreiðsla] \dicTerm{líf·eyris··greiðsl|a} \dicIPA{{l}{i}{\textlengthmark}{v}{ei}{r}{\textsci}{s}{\r{g}}{r}{ei}{ð}{s}{\textsubring{d}}{l}{a}} \dicPos{f}[1] \dicFlx{(‑u, ‑ur)}[13] \dicDirectTranslationCS{výplata důchodu}
\dicEntry[lífeyrisréttindi] \dicTerm{líf·eyris··rétt·indi} \dicIPA{{l}{i}{\textlengthmark}{v}{ei}{r}{\textsci}{s}{r}{j}{\textepsilon}{h}{\textsubring{d}}{\textsci}{n}{\textsubring{d}}{\textsci}} \dicPos{n}[2] \dicFlx{pl}[19] \dicDirectTranslationCS{nárok na důchod}
\dicEntry[lífeyrissjóðstekjur] \dicTerm{líf·eyris·sjóðs··tekjur} \dicIPA{{l}\-{i}\-{\textlengthmark}\-{v}\-{ei}\-{r}\-{\textsci}\-{s}\-{j}\-{ou}\-{ð}\-{s}\-{t\smash{\textsuperscript{h}}}\-{\textepsilon}\-{\r{\textObardotlessj}}\-{\textscy}\-{\textsubring{r}}\-} \dicPos{f}[1] \dicFlx{pl}[18] \dicDirectTranslationCS{příjmy z~penzijního fondu}
\dicEntry[lífeyrissjóður] \dicTerm{líf·eyris··sjóð|ur} \dicIPA{{l}{i}{\textlengthmark}{v}{ei}{r}{\textsci}{s}{j}{ou}{ð}{\textscy}{\textsubring{r}}} \dicPos{m}[9] \dicFlx{(‑s, ‑ir)}[6] \dicFieldCat{ekon.} \dicDirectTranslationCS{penzijní fond}
\dicEntry[lífeyrissparnaður] \dicTerm{líf·eyris··spar·nað|ur} \dicIPA{{l}\-{i}\-{\textlengthmark}\-{v}\-{ei}\-{r}\-{\textsci}\-{s}\-{\textsubring{b}}\-{a}\-{r}\-{\textsubring{d}}\-{n}\-{a}\-{ð}\-{\textscy}\-{\textsubring{r}}\-} \dicPos{m}[10] \dicFlx{(‑ar)}[9] \dicDirectTranslationCS{důchodové spoření}
\dicEntry[lífeyristryggingar] \dicTerm{líf··eyris··trygg·ingar} \dicIPA{{l}\-{i}\-{\textlengthmark}\-{v}\-{ei}\-{r}\-{\textsci}\-{s}\-{t\smash{\textsuperscript{h}}}\-{r}\-{\textsci}\-{\r{\textObardotlessj}}\-{i}\-{\ng}\-{\r{g}}\-{a}\-{\textsubring{r}}\-} \dicPos{f}[4] \dicFlx{pl}[6] \dicDirectTranslationCS{důchodové pojištění}
\dicEntry[lífeyrisþegi] \dicTerm{líf·eyris··þeg|i} \dicIPA{{l}{i}{\textlengthmark}{v}{ei}{r}{\textsci}{s}{\texttheta}{ei}{\textsci}} \dicPos{m}[1] \dicFlx{(‑a, ‑ar)}[1] \dicDirectTranslationCS{důchodce, důchodkyně, penzista, penzistka}
\dicEntry[líffræði] \dicTerm{líf··fræð|i} \dicIPA{{l}{i}{f}{\textlengthmark}{r}{a}{i}{ð}{\textsci}} \dicPos{f}[3] \dicFlx{(‑i)}[3] \dicDirectTranslationCS{biologie}
\dicEntry[líffræðingur] \dicTerm{líf·fræð··ing|ur} \dicIPA{{l}{i}{f}{\textlengthmark}{r}{a}{i}{ð}{i}{\ng}{\r{g}}{\textscy}{\textsubring{r}}} \dicPos{m}[6] \dicFlx{(‑s, ‑ar)}[8] \dicDirectTranslationCS{biolog, bioložka}
\dicEntry[líffærafræði] \dicTerm{líf·færa··fræð|i} \dicIPA{{l}{i}{f}{\textlengthmark}{a}{i}{r}{a}{f}{r}{a}{i}{ð}{\textsci}} \dicPos{f}[3] \dicFlx{(‑i)}[3] \dicDirectTranslationCS{anatomie}
\dicEntry[líffæri] \dicTerm{líf··færi} \dicIPA{{l}{i}{f}{\textlengthmark}{a}{i}{r}{\textsci}} \dicPos{n}[2] \dicFlx{(‑s, ‑)}[14] \dicFieldCat{anat.} \dicDirectTranslationCS{orgán, ústrojí}
\dicEntry[lífga] \dicTerm{lífg|a} \dicIPA{{l}{i}{v}{\r{g}}{a}} \dicPos{v}[1] \dicFlx{(‑aði)}[1] \dicFlx{acc} \textbf{1.} \dicSynonym{fjörga} \dicDirectTranslationCS{oživit, oživovat, vzkřísit} \dicIndirectTranslationCS{(činit živějším)} \dicExampleIS{lífga eldinn} \dicExampleCS{oživit oheň};  \dicPhraseIS{lífga upp á e‑ð} \dicDirectTranslationCS{oživit\,/\addthin osvěžit (co) (pokoj barvami ap.)}  \textbf{2.} \dicSynonym*{vekja til lífs} \dicDirectTranslationCS{oživit, oživovat, vzkřísit, resuscitovat} \dicIndirectTranslationCS{(probudit člověka k~životu)};  \dicPhraseIS{lífga e‑n úr dauðadái} \dicDirectTranslationCS{probudit (koho) z~kómatu};  \dicPhraseIS{lífga e‑n við} \dicDirectTranslationCS{oživit (koho), probudit (koho) k~životu}
\dicEntry[lífgjafi] \dicTerm{líf··gjaf|i} \dicIPA{{l}{i}{v}{\r{\textObardotlessj}}{a}{v}{\textsci}} \dicPos{m}[1] \dicFlx{(‑a, ‑ar)}[8] \dicDirectTranslationCS{zachránce\,/\addthin zachránkyně (života)}
\dicEntry[lífgjöf] \dicTerm{líf··|gjöf} \dicIPA{{l}{i}{v}{\r{\textObardotlessj}}{\oe}{f}} \dicPos{f}[7] \dicFlx{(‑gjafar)}[19] \dicDirectTranslationCS{záchrana života}
\dicEntry[lífgun] \dicTerm{lífg|un} \dicIPA{{l}{i}{v}{\r{g}}{\textscy}{\textsubring{n}}} \dicPos{f}[7] \dicFlx{(‑unar, ‑anir)}[8] \dicDirectTranslationCS{resuscitace, reanimace, oživování}
\dicEntry[lífhimna] \dicTerm{líf··himn|a} \dicIPA{{l}{i}{\textlengthmark}{f}{h}{\textsci}{m}{n}{a}} \dicPos{f}[1] \dicFlx{(‑u, ‑ur)}[7] \dicFieldCat{anat.} \dicDirectTranslationCS{pobřišnice}
\dicEntry[lífklukka] \dicTerm{líf··klukk|a} \dicIPA{{l}{i}{v}{k\smash{\textsuperscript{h}}}{l}{\textscy}{h}{\r{g}}{a}} \dicPos{f}[1] \dicFlx{(‑u, ‑ur)}[13] \dicDirectTranslationCS{biologické hodiny}
\dicEntry[líflaus] \dicTerm{líf··laus} \dicIPA{{l}{i}{v}{l}{\oe i}{s}} \dicPos{adj}[5]\dicFlx{}[-1] \textbf{1.} \dicDirectTranslationCS{(jsoucí) bez života, mrtvý}  \textbf{2.} \dicFieldCat{biol.} \dicLangCat{zast.} \dicSynonym{ólífrænn} \dicDirectTranslationCS{neživý}
\dicEntry[líflát] \dicTerm{líf··lát} \dicIPA{{l}{i}{v}{l}{au}{\textsubring{d}}} \dicPos{n}[2] \dicFlx{(‑s, ‑)}[5] \dicSynonym{aftaka\smash{\textsuperscript{1}}} \dicDirectTranslationCS{poprava} \dicExampleIS{hóta e‑m  líflát} \dicExampleCS{hrozit (komu) popravou}
\dicEntry[lífláta] \dicTerm{líf··|láta} \dicIPA{{l}{i}{v}{l}{au}{\textsubring{d}}{a}} \dicPos{v}[6] \dicFlx{(‑læt, ‑lét, ‑létum, ‑léti, ‑látið)}[87] \dicFlx{acc} \dicDirectTranslationCS{popravit} \dicExampleIS{lífláta e‑n} \dicExampleCS{popravit (koho)}
\dicEntry[líflátið] \dicTerm{líf··látið} \dicIPA{{l}{i}{v}{l}{au}{\textsubring{d}}{\textsci}{\texttheta}} \dicPos{v} \dicFlx{supin} \dicLink{lífláta}
\dicEntry[líflegur] \dicTerm{líf··legur} \dicIPA{{l}{i}{v}{l}{\textepsilon}{\textbabygamma}{\textscy}{\textsubring{r}}} \dicPos{adj}[1]\dicFlx{}[-8] \dicSynonym{fjörugur} \dicDirectTranslationCS{živý, svěží, plný života}
\dicEntry[líflét] \dicTerm{líf··lét} \dicIPA{{l}{i}{v}{l}{j}{\textepsilon}{\textsubring{d}}} \dicPos{v} \dicFlx{ind pf sg 1 pers} \dicLink{lífláta}
\dicEntry[lífléti] \dicTerm{líf··léti} \dicIPA{{l}{i}{v}{l}{j}{\textepsilon}{\textsubring{d}}{\textsci}} \dicPos{v} \dicFlx{con pf sg 1 pers} \dicLink{lífláta}
\dicEntry[líflétum] \dicTerm{líf··létum} \dicIPA{{l}{i}{v}{l}{j}{\textepsilon}{\textsubring{d}}{\textscy}{\textsubring{m}}} \dicPos{v} \dicFlx{ind pf pl 1 pers} \dicLink{lífláta}
\dicEntry[líflæt] \dicTerm{líf··læt} \dicIPA{{l}{i}{v}{l}{a}{i}{\textsubring{d}}} \dicPos{v} \dicFlx{ind praes sg 1 pers} \dicLink{lífláta}
\dicEntry[lífríki] \dicTerm{líf··ríki} \dicIPA{{l}{i}{v}{r}{i}{\r{\textObardotlessj}}{\textsci}} \dicPos{n}[2] \dicFlx{(‑s, ‑)}[16] \dicDirectTranslationCS{biosystém, biologický systém}
\dicEntry[lífrænn] \dicTerm{líf··rænn} \dicIPA{{l}{i}{v}{r}{a}{i}{\textsubring{d}}{\textsubring{n}}} \dicPos{adj}[7]\dicFlx{}[-1] \textbf{1.} \dicDirectTranslationCS{živý, živoucí, rušný} \dicExampleIS{lífrænt starf} \dicExampleCS{živá práce}  \textbf{2.} \dicSynonym{náttúrulegur} \dicDirectTranslationCS{organický, ekologický} \dicExampleIS{lífrænn búskapur} \dicExampleCS{organické zemědělství} \dicAntonym{ólífrænn}
\dicEntry[lífsbarátta] \dicTerm{lífs··bar·átt|a} \dicIPA{{l}{i}{f}{s}{\textsubring{b}}{a}{r}{au}{h}{\textsubring{d}}{a}} \dicPos{f}[1] \dicFlx{(‑u)}[5] \dicDirectTranslationCS{existenční\,/\addthin životní boj}
\dicEntry[lífseigur] \dicTerm{líf··seigur} \dicIPA{{l}{i}{f}{s}{ei}{\textbabygamma}{\textscy}{\textsubring{r}}} \dicPos{adj}[1]\dicFlx{}[-1] \dicDirectTranslationCS{houževnatý, (jsoucí) s~tuhým kořínkem}
\dicEntry[lífsflótti] \dicTerm{lífs··flótt|i} \dicIPA{{l}{i}{f}{s}{f}{l}{ou}{h}{\textsubring{d}}{\textsci}} \dicPos{m}[1] \dicFlx{(‑a)}[3] \dicDirectTranslationCS{útěk před realitou, eskapismus}
\dicEntry[lífsglaður] \dicTerm{lífs··|glaður} \dicIPA{{l}{i}{f}{s}{\r{g}}{l}{a}{ð}{\textscy}{\textsubring{r}}} \dicPos{adj}[2] \dicFlx{(f ‑glöð)}[7] \dicSynonym{glaðlyndur} \dicDirectTranslationCS{radostný, šťastný, radující se ze života}
\dicEntry[lífsgæði] \dicTerm{lífs··gæði} \dicIPA{{l}{i}{f}{s}{\r{\textObardotlessj}}{a}{i}{ð}{\textsci}} \dicPos{n}[2] \dicFlx{pl}[19] \dicDirectTranslationCS{životní úroveň, kvalita života}
\dicEntry[lífsháski] \dicTerm{lífs··hásk|i} \dicIPA{{l}{i}{f}{s}{h}{au}{s}{\r{\textObardotlessj}}{\textsci}} \dicPos{m}[1] \dicFlx{(‑a)}[3] \dicDirectTranslationCS{ohrožení života}
\dicEntry[lífshætta] \dicTerm{lífs··hætt|a} \dicIPA{{l}{i}{f}{s}{h}{a}{i}{h}{\textsubring{d}}{a}} \dicPos{f}[1] \dicFlx{(‑u)}[5] \dicDirectTranslationCS{nebezpečí života}
\dicEntry[lífshættulegur] \dicTerm{lífs··hættu·legur} \dicIPA{{l}{i}{f}{s}{h}{a}{i}{h}{\textsubring{d}}{\textscy}{l}{\textepsilon}{\textbabygamma}{\textscy}{\textsubring{r}}} \dicPos{adj}[1]\dicFlx{}[-8] \dicDirectTranslationCS{životu nebezpečný}
\dicEntry[lífskjör] \dicTerm{lífs··kjör} \dicIPA{{l}{i}{f}{s}{c\smash{\textsuperscript{h}}}{\oe}{\textsubring{r}}} \dicPos{n}[2] \dicFlx{pl}[9] \dicDirectTranslationCS{životní podmínky}
\dicEntry[lífsleiði] \dicTerm{lífs··leið|i} \dicIPA{{l}{i}{f}{s}{l}{ei}{ð}{\textsci}} \dicPos{m}[1] \dicFlx{(‑a)}[3] \dicDirectTranslationCS{životní nuda, omrzelost životem}
\dicEntry[lífsleikni] \dicTerm{lífs··leikn|i} \dicIPA{{l}{i}{f}{s}{l}{ei}{h}{\r{g}}{n}{\textsci}} \dicPos{f}[3] \dicFlx{(‑i)}[3] \dicFieldCat{škol.} \dicDirectTranslationCS{praktická výchova} \dicIndirectTranslationCS{(předmět ve škole)}
\dicEntry[lífsnauðsyn] \dicTerm{lífs··nauð·syn} \dicIPA{{l}{i}{f}{s}{n}{\oe i}{ð}{s}{\textsci}{\textsubring{n}}} \dicPos{f}[4] \dicFlx{(‑jar, ‑jar)}[9] \dicDirectTranslationCS{životní nutnost}
\dicEntry[lífsorka] \dicTerm{lífs··ork|a} \dicIPA{{l}{i}{f}{s}{\textopeno}{\textsubring{r}}{\r{g}}{a}} \dicPos{f}[1] \dicFlx{(‑u)}[5] \dicDirectTranslationCS{vitalita, životní energie}
\dicEntry[lífsregla] \dicTerm{lífs··regl|a} \dicIPA{{l}{i}{f}{s}{r}{\textepsilon}{\r{g}}{l}{a}} \dicPos{f}[1] \dicFlx{(‑u, ‑ur)}[19] \dicDirectTranslationCS{životní zásada}
\dicEntry[lífsreynsla] \dicTerm{lífs··reynsl|a} \dicIPA{{l}{i}{f}{s}{r}{ei}{n}{s}{\textsubring{d}}{l}{a}} \dicPos{f}[1] \dicFlx{(‑u)}[5] \dicDirectTranslationCS{životní zkušenost}
\dicEntry[lífsskilyrði] \dicTerm{lífs··skil·yrði} \dicIPA{{l}{i}{f}{s}{\r{\textObardotlessj}}{\textsci}{l}{\textsci}{r}{ð}{\textsci}} \dicPos{n}[2] \dicFlx{pl}[19] \dicDirectTranslationCS{životní podmínky}
\dicEntry[lífsskoðun] \dicTerm{lífs··skoð|un} \dicIPA{{l}{i}{f}{s}{\r{g}}{\textopeno}{ð}{\textscy}{\textsubring{n}}} \dicPos{f}[7] \dicFlx{(‑unar, ‑anir)}[8] \dicDirectTranslationCS{životní postoj\,/\addthin názor}
\dicEntry[lífsspeki] \dicTerm{lífs··spek|i} \dicIPA{{l}{i}{f}{s}{\textsubring{b}}{\textepsilon}{\r{\textObardotlessj}}{\textsci}} \dicPos{f}[3] \dicFlx{(‑i)}[3] \dicDirectTranslationCS{životní moudrost}
\dicEntry[lífsstarf] \dicTerm{lífs··|starf} \dicIPA{{l}{i}{f}{s}{\textsubring{d}}{a}{r}{f}} \dicPos{n}[2] \dicFlx{(‑starfs, ‑störf)}[8] \dicDirectTranslationCS{životní poslání}
\dicEntry[lífsstíll] \dicTerm{lífs··stíl|l} \dicIPA{{l}{i}{f}{s}{\textsubring{d}}{i}{\textsubring{d}}{\textsubring{l}}} \dicPos{m}[6] \dicFlx{(‑s)}[50] \dicDirectTranslationCS{životní styl}
\dicEntry[lífstykki] \dicTerm{líf··stykki} \dicIPA{{l}{i}{f}{s}{\textsubring{d}}{\textsci}{h}{\r{\textObardotlessj}}{\textsci}} \dicPos{n}[2] \dicFlx{(‑s, ‑)}[16] \dicDirectTranslationCS{korzet, šněrovačka}
\dicEntry[lífsveig] \dicTerm{lífs··veig} \dicIPA{{l}{i}{f}{s}{v}{ei}{x}} \dicPos{f}[4] \dicFlx{(‑ar, ‑ar)}[1] \dicDirectTranslationCS{elixír života}
\dicEntry[lífsviðurværi] \dicTerm{lífs··viður·væri} \dicIPA{{l}{i}{f}{s}{v}{\textsci}{ð}{\textscy}{r}{v}{a}{i}{r}{\textsci}} \dicPos{n}[2] \dicFlx{(‑s)}[20] \dicDirectTranslationCS{obživa (jídlo ap.)}
\dicEntry[lífsþægindi] \dicTerm{lífs··þæg·indi} \dicIPA{{l}{i}{f}{s}{\texttheta}{a}{i}{j}{\textsci}{n}{\textsubring{d}}{\textsci}} \dicPos{n}[2] \dicFlx{pl}[19] \dicDirectTranslationCS{životní komfort\,/\addthin pohodlí}
\dicEntry[líftóra] \dicTerm{líf··tór|a} \dicIPA{{l}{i}{f}{t\smash{\textsuperscript{h}}}{ou}{r}{a}} \dicPos{f}[1] \dicFlx{(‑u)}[5] \dicDirectTranslationCS{jiskra života};  \dicPhraseIS{hræða úr e‑m líftóruna} \dicLangCat{přen.} \dicDirectTranslationCS{vyděsit (koho) k~smrti}
\dicEntry[líftrygging] \dicTerm{líf··trygg·ing} \dicIPA{{l}{i}{f}{t\smash{\textsuperscript{h}}}{r}{\textsci}{\r{\textObardotlessj}}{i}{\ng}{\r{g}}} \dicPos{f}[4] \dicFlx{(‑ar, ‑ar)}[5] \dicDirectTranslationCS{životní pojištění}
\dicEntry[líftryggja] \textls[15]{\dicTerm{líf··trygg|ja} \dicIPA{{l}{i}{f}{t\smash{\textsuperscript{h}}}{r}{\textsci}{\r{\textObardotlessj}}{a}} \dicPos{v}[2] \dicFlx{(‑ði, ‑t)}[92] \dicFlx{acc} \dicDirectTranslationCS{uzavřít životní pojiš\nobreak tění}}
\dicEntry[lífvera] \dicTerm{líf··ver|a} \dicIPA{{l}{i}{v}{\textlengthmark}{\textepsilon}{r}{a}} \dicPos{f}[1] \dicFlx{(‑u, ‑ur)}[7] \dicDirectTranslationCS{živá bytost, živý organismus}
\dicEntry[lífvænlegur] \dicTerm{líf··væn·legur} \dicIPA{{l}{i}{v}{\textlengthmark}{a}{i}{n}{l}{\textepsilon}{\textbabygamma}{\textscy}{\textsubring{r}}} \dicPos{adj}[1]\dicFlx{}[-8] \textbf{1.} \dicDirectTranslationCS{rentabilní} \dicExampleIS{lífvænleg atvinna} \dicExampleCS{rentabilní zaměstnání}  \textbf{2.} \dicDirectTranslationCS{příznivý pro život} \dicExampleIS{lífvænlegur staður} \dicExampleCS{místo příznivé pro život}
\begin{xtolerant}{}{1pt}
\dicEntry[lífvörður] \dicTerm{líf··|vörður}\addthinS\dicIPA{{l}{i}{v}{\textlengthmark}{\oe}{r}{ð}{\textscy}{\textsubring{r}}}\addthinS\dicPos{m}[11]\addthinS\dicFlx{(‑varðar, ‑verðir)}[5] \dicDirectTranslationCS{osobní strážce\,/\addthin strážkyně, tělesná stráž, bodyguard}
\end{xtolerant}
\dicEntry[lík] \dicTerm{lík} \dicsymFrequent\  \dicIPA{{l}{i}{\textlengthmark}{\r{g}}} \dicPos{n}[2] \dicFlx{(‑s, ‑)}[5] \textbf{1.} \dicSynonym*{dauður líkami} \dicDirectTranslationCS{mrtvola} \dicExampleIS{líkið af e‑m} \dicExampleCS{(čí) mrtvola}  \textbf{2.} \dicLangCat{hovor.} \dicSynonym*{tóm vínflaska} \dicDirectTranslationCS{prázdná láhev (po alkoholu)}
\dicEntry[líka] \dicTerm{lík|a\smash{\textsuperscript{1}}} \dicsymFrequent\  \dicIPA{{l}{i}{\textlengthmark}{\r{g}}{a}} \dicPos{v}[1] \dicFlx{(‑aði)}[82] \dicFlx{impers} \dicPhraseIS{e‑m líkar e‑að} \dicDirectTranslationCS{(co) se (komu) líbí, (co) má (kdo) rád} \dicExampleIS{Mér líkar það ekki.} \dicExampleCS{Mně se to nelíbí.};  \dicPhraseIS{e‑m líkar vel við e‑n} \dicDirectTranslationCS{(kdo) má (koho) rád}
\dicEntry[líka] \dicTerm{líka\smash{\textsuperscript{2}}} \dicsymFrequent\  \dicIPA{{l}{i}{\textlengthmark}{\r{g}}{a}} \dicPos{adv} \dicSynonym{einnig} \dicDirectTranslationCS{také, taky, rovněž} \dicExampleIS{Dagblaðið er líka á netinu.} \dicExampleCS{Deník je také na internetu.}
\dicEntry[líkami] \dicTerm{líkam|i} \dicsymFrequent\  \dicIPA{{l}{i}{\textlengthmark}{\r{g}}{a}{m}{\textsci}} \dicPos{m}[1] \dicFlx{(‑a, ‑ar)}[8] \dicDirectTranslationCS{tělo} \dicExampleIS{æfingar fyrir líkama} \dicExampleCS{cvičení pro tělo}
\dicEntry[líkamlegur] \dicTerm{líkam··legur} \dicsymFrequent\  \dicIPA{{l}{i}{\r{g}}{a}{m}{l}{\textepsilon}{\textbabygamma}{\textscy}{\textsubring{r}}} \dicPos{adj}[1]\dicFlx{}[-8] \dicDirectTranslationCS{tělesný, fyzický} \dicExampleIS{líkamleg atvinna} \dicExampleCS{fyzická práce}
\dicEntry[líkamsárás] \dicTerm{líkams··á·rás} \dicIPA{{l}{i}{\textlengthmark}{\r{g}}{a}{m}{s}{au}{r}{au}{s}} \dicPos{f}[7] \dicFlx{(‑ar, ‑ir)}[1] \dicDirectTranslationCS{tělesný útok, fyzické napadení}
\dicEntry[líkamsástand] \dicTerm{líkams··á·stand} \dicIPA{{l}{i}{\textlengthmark}{\r{g}}{a}{m}{s}{au}{s}{\textsubring{d}}{a}{n}{\textsubring{d}}} \dicPos{n}[2] \dicFlx{(‑s)}[2] \dicDirectTranslationCS{tělesný stav}
\dicEntry[líkamsburðir] \dicTerm{líkams··burðir} \dicIPA{{l}{i}{\textlengthmark}{\r{g}}{a}{m}{s}{\textsubring{b}}{\textscy}{r}{ð}{\textsci}{\textsubring{r}}} \dicPos{m}[10] \dicFlx{pl}[1] \dicDirectTranslationCS{tělesná\,/\addthin fyzická síla}
\dicEntry[líkamshiti] \dicTerm{líkams··hit|i} \dicIPA{{l}{i}{\textlengthmark}{\r{g}}{a}{m}{s}{h}{\textsci}{\textsubring{d}}{\textsci}} \dicPos{m}[1] \dicFlx{(‑a)}[3] \dicDirectTranslationCS{tělesná teplota}
\dicEntry[líkamsrækt] \dicTerm{líkams··rækt} \dicIPA{{l}{i}{\textlengthmark}{\r{g}}{a}{m}{s}{r}{a}{i}{x}{\textsubring{d}}} \dicPos{f}[4] \dicFlx{(‑ar)}[3] \dicDirectTranslationCS{tělesné cvičení, posilování, fitness}
\dicEntry[líkamsræktarstöð] \dicTerm{líkams·ræktar··stöð} \dicIPA{{l}{i}{\textlengthmark}{\r{g}}{a}{m}{s}{r}{a}{i}{x}{\textsubring{d}}{a}{\textsubring{r}}{s}{\textsubring{d}}{\oe}{\texttheta}} \dicPos{f}[6] \dicFlx{(‑var, ‑var)}[1] \dicDirectTranslationCS{posilovna, fitness centrum}
\dicEntry[líkamsstærð] \dicTerm{líkams··stærð} \dicIPA{{l}{i}{\textlengthmark}{\r{g}}{a}{m}{s}{\textsubring{d}}{a}{i}{r}{\texttheta}} \dicPos{f}[7] \dicFlx{(‑ar)}[3] \dicDirectTranslationCS{životní velikost} \dicExampleIS{í fullri líkamsstærð} \dicExampleCS{v~životní velikosti}
\dicEntry[líkan] \dicTerm{lík|an} \dicIPA{{l}{i}{\textlengthmark}{\r{g}}{a}{\textsubring{n}}} \dicPos{n}[2] \dicFlx{(‑ans, ‑ön)}[8] \dicDirectTranslationCS{model, maketa} \dicIndirectTranslationCS{(zmenšené provedení předmětu)} \dicExampleIS{líkan af húsi} \dicExampleCS{model domu}
\dicEntry[líkbrennsla] \dicTerm{lík··brennsl|a} \dicIPA{{l}{i}{\textlengthmark}{\r{g}}{\textsubring{b}}{r}{\textepsilon}{n}{s}{\textsubring{d}}{l}{a}} \dicPos{f}[1] \dicFlx{(‑u, ‑ur)}[19] \dicSynonym{bálför} \dicDirectTranslationCS{kremace}
\dicEntry[líkbörur] \dicTerm{lík··börur} \dicIPA{{l}{i}{\textlengthmark}{\r{g}}{\textsubring{b}}{\oe}{r}{\textscy}{\textsubring{r}}} \dicPos{f}[12] \dicFlx{pl}[4] \dicDirectTranslationCS{máry}
\dicEntry[líkfylgd] \dicTerm{lík··fylgd} \dicIPA{{l}{i}{\textlengthmark}{\r{g}}{f}{\textsci}{l}{\textsubring{d}}} \dicPos{f}[7] \dicFlx{(‑ar, ‑ir)}[1] \dicDirectTranslationCS{pohřební průvod}
\dicEntry[líkhús] \dicTerm{lík··hús} \dicIPA{{l}{i}{\textlengthmark}{\r{g}}{h}{u}{s}} \dicPos{n}[2] \dicFlx{(‑s, ‑)}[5] \dicDirectTranslationCS{márnice}
\dicEntry[líki] \dicTerm{lík|i\smash{\textsuperscript{1}}} \dicIPA{{l}{i}{\textlengthmark}{\r{\textObardotlessj}}{\textsci}} \dicPos{m}[1] \dicFlx{(‑a, ‑ar)}[1] \dicSynonym{jafningi} \dicDirectTranslationCS{sobě rovný, podobný (člověk)} \dicExampleIS{hann og hans líkar} \dicExampleCS{on a~jemu podobní};  \dicPhraseIS{eiga engan sinn líka} \dicDirectTranslationCS{nemít sobě rovného}
\dicEntry[líki] \dicTerm{líki\smash{\textsuperscript{2}}} \dicIPA{{l}{i}{\textlengthmark}{\r{\textObardotlessj}}{\textsci}} \dicPos{n}[2] \dicFlx{(‑s, ‑)}[16] \dicSynonym{sköpulag} \dicDirectTranslationCS{podoba, forma};  \dicPhraseIS{í heilu líki} \dicFlx{adj} \dicDirectTranslationCS{celý, netknutý}
\dicEntry[líkindi] \dicTerm{lík··indi} \dicsymFrequent\  \dicIPA{{l}{i}{\textlengthmark}{\r{\textObardotlessj}}{\textsci}{n}{\textsubring{d}}{\textsci}} \dicPos{n}[2] \dicFlx{pl}[19] \dicDirectTranslationCS{pravděpodobnost} \dicExampleIS{Það eru mikil líkindi til þess að \dots{}} \dicExampleCS{Je velká pravděpodobnost, že\dots{}};  \dicPhraseIS{að öllum líkindum} \dicFlx{adv} \dicDirectTranslationCS{s~největší pravděpodobností}
\dicEntry[líking] \dicTerm{lík··ing} \dicIPA{{l}{i}{\textlengthmark}{\r{\textObardotlessj}}{i}{\ng}{\r{g}}} \dicPos{f}[4] \dicFlx{(‑ar, ‑ar)}[5] \textbf{1.} \dicSynonym{líki\smash{\textsuperscript{2}}} \dicDirectTranslationCS{podoba, podobnost, obdoba, analogie}  \textbf{2.} \dicDirectTranslationCS{přirovnání, podobenství} \dicExampleIS{tala í líkingum} \dicExampleCS{mluvit v~přirovnáních}  \textbf{3.} \dicFieldCat{lit.} \dicSynonym{samlíking} \dicDirectTranslationCS{metafora} \dicExampleIS{líkingar í skáldskap} \dicExampleCS{metafora v~literatuře}  \textbf{4.} \dicFieldCat{mat.} \dicSynonym{jafna\smash{\textsuperscript{1}}} \dicDirectTranslationCS{rovnost}
\dicEntry[líkja] \dicTerm{lík|ja} \dicsymFrequent\  \dicIPA{{l}{i}{\textlengthmark}{\r{\textObardotlessj}}{a}} \dicPos{v}[2] \dicFlx{(‑ti, ‑t)}[25] \dicFlx{dat} \textbf{1.} \dicDirectTranslationCS{přirovnat, přirovnávat};  \dicPhraseIS{líkja e‑u saman} \dicDirectTranslationCS{srovnávat (co)};  \dicPhraseIS{líkja e‑u við e‑ð} \dicDirectTranslationCS{přirovnávat (co) k~(čemu)} \dicExampleIS{líkja ævi mannsins við blóm} \dicExampleCS{přirovnávat život člověka ke květině}  \textbf{2.} \dicDirectTranslationCS{napodobit, napodobovat};  \dicPhraseIS{líkja eftir e‑u} \dicDirectTranslationCS{napodobit (co), imitovat (co)} \dicExampleIS{líkja eftir hljóði fuglsins} \dicExampleCS{napodobovat zvuk ptáčka};  \dicIdiom{líkjast}{ \dicPhraseIS{líkjast e‑m}} \dicFlx{refl} \dicDirectTranslationCS{připomínat (koho), podobat se (komu)} \dicExampleIS{Hverjum líkist ég?} \dicExampleCS{Komu se podobám?}
\dicEntry[líkjör] \dicTerm{líkjör} \dicIPA{{l}{i}{\textlengthmark}{\r{\textObardotlessj}}{\oe}{\textsubring{r}}} \dicPos{m}[4] \dicFlx{(‑s, ‑ar)}[14] \dicDirectTranslationCS{likér}
\dicEntry[líkkista] \dicTerm{lík··kist|a} \dicIPA{{l}{i}{\textlengthmark}{\r{g}}{c\smash{\textsuperscript{h}}}{\textsci}{s}{\textsubring{d}}{a}} \dicPos{f}[1] \dicFlx{(‑u, ‑ur)}[19] \dicDirectTranslationCS{rakev}
\dicEntry[líkklæði] \dicTerm{lík··klæði} \dicIPA{{l}{i}{\textlengthmark}{\r{g}}{k\smash{\textsuperscript{h}}}{l}{a}{i}{ð}{\textsci}} \dicPos{n}[2] \dicFlx{pl}[19] \dicDirectTranslationCS{pohřební roucho}
\dicEntry[líkl.] \dicTerm{líkl.} \dicPos{zkr} \dicPhraseIS{líklega} \dicFlx{adv} \dicDirectTranslationCS{pravděpodobně}
\dicEntry[líklega] \dicTerm{lík··leg|a} \dicsymFrequent\  \dicIPA{{l}{i}{\textlengthmark}{\r{g}}{l}{\textepsilon}{\textbabygamma}{a}} \dicPos{adv} \dicFlx{(comp ‑ar, sup ‑ast)} \dicSynonym{sennilega} \dicDirectTranslationCS{pravděpodobně, nejspíš, asi} \dicExampleIS{Þau segja líklega nei.} \dicExampleCS{Pravděpodobně řeknou ne.}
\dicEntry[líklegur] \dicTerm{lík··legur} \dicsymFrequent\  \dicIPA{{l}{i}{\textlengthmark}{\r{g}}{l}{\textepsilon}{\textbabygamma}{\textscy}{\textsubring{r}}} \dicPos{adj}[1]\dicFlx{}[-8] \textbf{1.} \dicSynonym{efnilegur} \dicDirectTranslationCS{slibný, nadějný}  \textbf{2.} \dicSynonym{sennilegur} \dicDirectTranslationCS{pravděpodobný, předpokládaný} \dicExampleIS{Þeir eru líklegir til að feta í sömu fótspor.} \dicExampleCS{Pravděpodobně budou kráčet ve stejných stopách.}
\dicEntry[líkn] \dicTerm{líkn} \dicIPA{{l}{i}{h}{\r{g}}{\textsubring{n}}} \dicPos{f}[7] \dicFlx{(‑ar, ‑ir)}[1] \textbf{1.} \dicSynonym{miskunn} \dicDirectTranslationCS{milost, smilování}  \textbf{2.} \dicSynonym{hjúkrun} \dicDirectTranslationCS{opatrování, péče}
\dicEntry[líkna] \dicTerm{líkn|a} \dicIPA{{l}{i}{h}{\r{g}}{n}{a}} \dicPos{v}[1] \dicFlx{(‑aði)}[1] \dicFlx{dat} \dicSynonym{hjúkra} \dicDirectTranslationCS{pečovat, poskytnout péči, opatrovat} \dicExampleIS{líkna e‑m} \dicExampleCS{pečovat o~(koho)}
\dicEntry[líknarbelgur] \dicTerm{líknar··belg|ur} \dicIPA{{l}{i}{h}{\r{g}}{n}{a}{r}{\textsubring{b}}{\textepsilon}{l}{\r{g}}{\textscy}{\textsubring{r}}} \dicPos{m}[9] \dicFlx{(‑s\,/\addthin ‑jar, ‑ir)}[26] \textbf{1.} \dicFieldCat{anat.} \dicSynonym*{fósturfylgja} \dicDirectTranslationCS{amnion, vnitřní plodový obal}  \textbf{2.} \dicSynonym*{loftpúði í bíl} \dicDirectTranslationCS{airbag} \dicIndirectTranslationCS{(zařízení pasivní bezpečnosti v~autě)}
\dicEntry[líknardeild] \dicTerm{líknar··deild} \dicIPA{{l}{i}{h}{\r{g}}{n}{a}{r}{\textsubring{d}}{ei}{l}{\textsubring{d}}} \dicPos{f}[7] \dicFlx{(‑ar, ‑ir)}[1] \dicDirectTranslationCS{hospic}
\dicEntry[líknardráp] \dicTerm{líknar··dráp} \dicIPA{{l}{i}{h}{\r{g}}{n}{a}{r}{\textsubring{d}}{r}{au}{\textsubring{b}}} \dicPos{n}[2] \dicFlx{(‑s, ‑)}[5] \dicDirectTranslationCS{eutanazie}
\dicEntry[líknarmál] \dicTerm{líknar··mál} \dicIPA{{l}{i}{h}{\r{g}}{n}{a}{r}{m}{au}{\textsubring{l}}} \dicPos{n}[2] \dicFlx{(‑s, ‑)}[5] \dicDirectTranslationCS{charita, dobročinnost}
\dicEntry[líknarstofnun] \dicTerm{líknar··stofn|un} \dicIPA{{l}{i}{h}{\r{g}}{n}{a}{\textsubring{r}}{s}{\textsubring{d}}{\textopeno}{\textsubring{b}}{n}{\textscy}{\textsubring{n}}} \dicPos{f}[7] \dicFlx{(‑unar, ‑anir)}[8] \dicSynonym{góðgerðastofnun} \dicDirectTranslationCS{dobročinná\,/\addthin charitativní organizace}
\dicEntry[líkneski] \dicTerm{lík··neski} \dicIPA{{l}{i}{h}{\r{g}}{n}{\textepsilon}{s}{\r{\textObardotlessj}}{\textsci}} \dicPos{n}[2] \dicFlx{(‑s, ‑)}[16] \dicSynonym{höggmynd} \dicDirectTranslationCS{podobizna}
\dicEntry[líknsamur] \dicTerm{líkn··|samur} \dicIPA{{l}{i}{h}{\r{g}}{\textsubring{n}}{s}{a}{m}{\textscy}{\textsubring{r}}} \dicPos{adj}[1] \dicFlx{(f ‑söm)}[2] \dicDirectTranslationCS{milosrdný, milostivý, soucitný}
\dicEntry[líkræða] \dicTerm{lík··ræð|a} \dicIPA{{l}{i}{\textlengthmark}{\r{g}}{r}{a}{i}{ð}{a}} \dicPos{f}[1] \dicFlx{(‑u, ‑ur)}[19] \dicDirectTranslationCS{smuteční proslov}
\dicEntry[líkskoðun] \dicTerm{lík··skoð|un} \dicIPA{{l}{i}{\textlengthmark}{\r{g}}{s}{\r{g}}{\textopeno}{ð}{\textscy}{\textsubring{n}}} \dicPos{f}[7] \dicFlx{(‑unar)}[9] \dicDirectTranslationCS{ohledání mrtvoly}
\dicEntry[líkskurður] \dicTerm{lík··skurð|ur} \dicIPA{{l}{i}{\textlengthmark}{\r{g}}{s}{\r{g}}{\textscy}{r}{ð}{\textscy}{\textsubring{r}}} \dicPos{m}[10] \dicFlx{(‑ar, ‑ir)}[4] \dicDirectTranslationCS{pitva}
\dicEntry[líkur] \dicTerm{líkur\smash{\textsuperscript{1}}} \dicsymFrequent\  \dicIPA{{l}{i}{\textlengthmark}{\r{g}}{\textscy}{\textsubring{r}}} \dicPos{f}[12] \dicFlx{pl}[5] \dicSynonym{líkindi} \dicDirectTranslationCS{šance, pravděpodobnost} \dicExampleIS{líkur á e‑u\,/\addthin líkur til e‑s} \dicExampleCS{šance na (co)};  \dicPhraseIS{ef að líkum lætur} \dicFlx{adv} \dicDirectTranslationCS{pravděpodobně, patrně}
\dicEntry[líkur] \dicTerm{líkur\smash{\textsuperscript{2}}} \dicsymFrequent\  \dicIPA{{l}{i}{\textlengthmark}{\r{g}}{\textscy}{\textsubring{r}}} \dicPos{adj}[1]\dicFlx{}[-1] \textbf{1.} \dicDirectTranslationCS{podobný} \dicExampleIS{lík vandamál} \dicExampleCS{podobné problémy};  \dicPhraseIS{vera líkur e‑m} \dicDirectTranslationCS{být (komu) podobný, podobat se (komu)}  \textbf{2.} \dicDirectTranslationCS{obdobný, analogický};  \dicIdiom{líkur}{ \dicPhraseIS{líkt og}} \dicFlx{conj} \dicDirectTranslationCS{jak, jako} \dicExampleIS{líta út líkt og Íslendingur} \dicExampleCS{vypadat jako Islanďan}
\dicEntry[líkþorn] \dicTerm{lík··þorn} \dicIPA{{l}{i}{\textlengthmark}{\r{g}}{\texttheta}{\textopeno}{r}{\textsubring{d}}{\textsubring{n}}} \dicPos{n}[2] \dicFlx{(‑s, ‑)}[5] \dicDirectTranslationCS{kuří oko}
\dicEntry[líkþrá] \dicTerm{lík··þrá} \dicIPA{{l}{i}{\textlengthmark}{\r{g}}{\texttheta}{r}{au}} \dicPos{f}[4] \dicFlx{(‑r)}[19] \dicFieldCat{med.} \dicDirectTranslationCS{lepra, malomocenství}
\dicEntry[líkþrár] \dicTerm{lík··þrár} \dicIPA{{l}{i}{\textlengthmark}{\r{g}}{\texttheta}{r}{au}{\textsubring{r}}} \dicPos{adj}[4]\dicFlx{}[-3] \dicDirectTranslationCS{malomocný}
\dicEntry[lím] \dicTerm{lím} \dicIPA{{l}{i}{\textlengthmark}{\textsubring{m}}} \dicPos{n}[2] \dicFlx{(‑s, ‑)}[5] \dicDirectTranslationCS{lepidlo, lep, klih}
\dicEntry[líma] \dicTerm{lím|a} \dicsymFrequent\  \dicIPA{{l}{i}{\textlengthmark}{m}{a}} \dicPos{v}[2] \dicFlx{(‑di, ‑t)}[141] \dicFlx{acc} \dicDirectTranslationCS{(na)lepit, nalepovat, přilepit, přilepovat} \dicExampleIS{líma mynd á blað} \dicExampleCS{nalepit obrázek na papír};  \dicIdiom{límast}{ \dicPhraseIS{límast}} \dicFlx{refl} \dicDirectTranslationCS{(při)lepit se}
\dicEntry[límband] \dicTerm{lím··|band} \dicIPA{{l}{i}{m}{\textsubring{b}}{a}{n}{\textsubring{d}}} \dicPos{n}[2] \dicFlx{(‑bands, ‑bönd)}[8] \dicDirectTranslationCS{lepicí páska, izolepa}
\dicEntry[límmiði] \dicTerm{lím··mið|i} \dicIPA{{l}{i}{m}{\textlengthmark}{\textsci}{ð}{\textsci}} \dicPos{m}[1] \dicFlx{(‑a, ‑ar)}[1] \dicDirectTranslationCS{nálepka, štítek, etiketa}
\dicEntry[lín] \dicTerm{lín} \dicIPA{{l}{i}{\textlengthmark}{\textsubring{n}}} \dicPos{n}[2] \dicFlx{(‑s)}[2] \textbf{1.} \dicFieldCat{bot.} \dicSynonym{hör} \dicDirectTranslationCS{len setý} \textit{(l.~{\textLA{Linum usitatissimum}})}  \dicsymPhoto\   \textbf{2.} \dicSynonym{léreft} \dicDirectTranslationCS{len, lněné plátno} \dicExampleIS{spinna lín} \dicExampleCS{tkát len}
\dicFigure{ds_image_lin_0_2.jpg}{Lín}{Lín - Franz Eugen Kohler, PD}
\dicEntry[lína] \dicTerm{lín|a} \dicsymFrequent\  \dicIPA{{l}{i}{\textlengthmark}{n}{a}} \dicPos{f}[1] \dicFlx{(‑u, ‑ur)}[7] \textbf{1.} \dicSynonym{strik} \dicDirectTranslationCS{čára, linka, linie}  \textbf{2.} \dicSynonym*{textalína} \dicDirectTranslationCS{řádek, řádka} \dicExampleIS{nokkrar línur úr textanum} \dicExampleCS{několik řádek z~textu}  \textbf{3.} \dicSynonym{band} \dicDirectTranslationCS{lano, šňůra, provázek}  \textbf{4.} \dicSynonym{lóð\smash{\textsuperscript{2}}} \dicDirectTranslationCS{lanová síť}  \textbf{5.} \dicSynonym{stefna\smash{\textsuperscript{1}}} \dicDirectTranslationCS{politický směr, politická linie}  \textbf{6.} \dicFieldCat{mat.} \dicDirectTranslationCS{přímka};  \dicIdiom{lína}{ \dicPhraseIS{lesa milli línanna}} \dicLangCat{přen.} \dicDirectTranslationCS{číst mezi řádky} \dicIndirectTranslationCS{(porozumět více, než je napsáno)}; { \dicPhraseIS{passa upp á línurnar}} \dicLangCat{přen.} \dicDirectTranslationCS{dávat pozor na linii} \dicIndirectTranslationCS{(snažit se neztloustnout)}
\dicEntry[línolía] \dicTerm{lín··olí|a} \dicIPA{{l}{i}{\textlengthmark}{n}{\textopeno}{l}{i}{j}{a}} \dicPos{f}[1] \dicFlx{(‑u)}[5] \dicDirectTranslationCS{lněný olej}
\dicEntry[línubil] \dicTerm{línu··bil} \dicIPA{{l}{i}{\textlengthmark}{n}{\textscy}{\textsubring{b}}{\textsci}{\textsubring{l}}} \dicPos{n}[2] \dicFlx{(‑s, ‑)}[5] \dicFieldCat{poč.} \dicDirectTranslationCS{řádkování}
\dicEntry[línumaður] \dicTerm{línu··|maður} \dicIPA{{l}{i}{\textlengthmark}{n}{\textscy}{m}{a}{ð}{\textscy}{\textsubring{r}}} \dicPos{m}[13] \dicFlx{(‑manns, ‑menn)}[2] \dicFieldCat{sport.} \dicDirectTranslationCS{krajní útočník\,/\addthin útočnice, křídlo}
\dicEntry[línurit] \dicTerm{línu··rit} \dicIPA{{l}{i}{\textlengthmark}{n}{\textscy}{r}{\textsci}{\textsubring{d}}} \dicPos{n}[2] \dicFlx{(‑s, ‑)}[5] \dicDirectTranslationCS{graf, diagram} \dicExampleIS{töflur og línurit} \dicExampleCS{tabulky a~grafy}
\dicEntry[línuskauti] \dicTerm{línu··skaut|i} \dicIPA{{l}{i}{\textlengthmark}{n}{\textscy}{s}{\r{g}}{\oe i}{\textsubring{d}}{\textsci}} \dicPos{m}[1] \dicFlx{(‑a, ‑ar)}[1] \dicPhraseIS{línuskautar} \dicFlx{pl} \dicFieldCat{sport.} \dicDirectTranslationCS{kolečkové brusle}
\dicEntry[línuskil] \dicTerm{línu··skil} \dicIPA{{l}{i}{\textlengthmark}{n}{\textscy}{s}{\r{\textObardotlessj}}{\textsci}{\textsubring{l}}} \dicPos{n}[2] \dicFlx{pl}[1] \dicFieldCat{poč.} \dicDirectTranslationCS{zalomení řádku}
\dicEntry[líra] \dicTerm{lír|a} \dicIPA{{l}{i}{\textlengthmark}{r}{a}} \dicPos{f}[1] \dicFlx{(‑u, ‑ur)}[7] \textbf{1.} \dicDirectTranslationCS{lira} \dicIndirectTranslationCS{(bývalá italská měna)}  \textbf{2.} \dicFieldCat{hud.} \dicDirectTranslationCS{lyra}
\dicEntry[lít] \dicTerm{lít} \dicIPA{{l}{i}{\textlengthmark}{\textsubring{d}}} \dicPos{v} \dicFlx{ind praes sg 1 pers} \dicLink{líta}
\dicEntry[líta] \dicTerm{líta} \dicsymFrequent\  \dicIPA{{l}{i}{\textlengthmark}{\textsubring{d}}{a}} \dicPos{v}[6] \dicFlx{(lít, leit, litum, liti, litið)}[76] \dicFlx{acc} \dicDirectTranslationCS{(po)dívat se, hledět, pohlédnout, pohlížet} \dicExampleIS{líta á klukkuna} \dicExampleCS{podívat se na hodiny};  \dicPhraseIS{líta e‑n} \dicDirectTranslationCS{pohlédnout na (koho)};  \dicPhraseIS{líta á e‑n} \dicDirectTranslationCS{podívat se na (koho)} \dicExampleIS{líta spyrjandi á hana} \dicExampleCS{podívat se na ni tázavě};  \dicIdiom{líta}[af]{ \dicPhraseIS{líta af e‑u\,/\addthin e‑m}} \dicDirectTranslationCS{moct spustit oči z~(čeho\,/\addthin koho)} \dicIndirectTranslationCS{(většinou v~záporu)};  \dicIdiom{líta}[á]{ \dicPhraseIS{líta á e‑ð}} \dicDirectTranslationCS{podívat se na (co), kouknout na (co), prozkoumat (co)} \dicExampleIS{líta á vélina} \dicExampleCS{kouknout na motor};  \dicPhraseIS{líta stórt á sig} \dicDirectTranslationCS{mít o~sobě vysoké mínění};  \dicIdiom{líta}[eftir]{ \dicPhraseIS{líta eftir e‑m}} \dicDirectTranslationCS{dohlédnout\,/\addthin dohlížet na (koho)} \dicExampleIS{líta eftir barninu} \dicExampleCS{dohlížet na dítě};  \dicIdiom{líta}[fram hjá]{ \dicPhraseIS{líta fram hjá e‑u}} \dicDirectTranslationCS{přehlédnout\,/\addthin přehlížet (co)} \dicExampleIS{líta fram hjá mikilvægum staðreyndum} \dicExampleCS{přehlížet důležitá fakta};  \dicIdiom{líta}[inn]{ \dicPhraseIS{líta inn hjá e‑m}} \dicDirectTranslationCS{(na)kouknout ke (komu), (za)stavit se u~(koho) (na návštěvu ap.)};  \dicIdiom{líta}[í]{ \dicPhraseIS{líta í e‑ð}} \dicDirectTranslationCS{podívat se do (čeho), nahlédnout do (čeho) (novin ap.)};  \dicIdiom{líta}[niður]{ \dicPhraseIS{líta niður á e‑n}} \dicDirectTranslationCS{dívat se na (koho) spatra, shlížet na (koho) povýšeně};  \dicIdiom{líta}[til]{ \dicPhraseIS{líta til e‑rs}} \dicDirectTranslationCS{zastavit se u~(koho), podívat se ke (komu)};  \dicIdiom{líta}[undan]{ \dicPhraseIS{líta undan}} \dicDirectTranslationCS{sklopit zrak, podívat se dolů};  \dicIdiom{líta}[upp]{ \dicPhraseIS{líta upp til e‑rs}} \dicDirectTranslationCS{obdivovat (koho), vzhlížet ke (komu) (s~úctou ap.)};  \dicIdiom{líta}[út]{ \dicPhraseIS{líta út}} \dicDirectTranslationCS{vypadat, vyhlížet} \dicExampleIS{líta vel út} \dicExampleCS{vypadat dobře}; { \dicPhraseIS{líta út fyrir að (vera ríkur)}} \dicDirectTranslationCS{vypadat na (boháče), vypadat jako (bohatý)};  \dicIdiom{líta}[við]{ \dicPhraseIS{líta ekki við e‑u}} \dicDirectTranslationCS{netknout se (čeho), nejevit zájem o~(co)}; { \dicPhraseIS{líta við (þar)}} \dicDirectTranslationCS{stavit se (tam) na skok, zaskočit (tam)};  \dicIdiom{líta}[yfir]{ \dicPhraseIS{líta yfir e‑ð}} \dicDirectTranslationCS{projít (co), prohlédnout (co)} \dicExampleIS{líta yfir textann} \dicExampleCS{projít text};  \dicIdiom{lítast}[á]{ \dicPhraseIS{e‑m líst vel á e‑ð}} \dicFlx{refl impers} \dicDirectTranslationCS{(co) se (komu) líbí}
\dicEntry[lítið] \dicTerm{lítið} \dicsymFrequent\  \dicIPA{{l}{i}{\textlengthmark}{\textsubring{d}}{\textsci}{\texttheta}} \dicPos{adv} \dicFlx{(comp minna, sup minnst)} \dicDirectTranslationCS{málo, trochu} \dicExampleIS{vera lítið áberandi} \dicExampleCS{být trochu nápadný}
\dicEntry[lítilfjörlegur] \dicTerm{lítil··fjör·legur} \dicIPA{{l}{i}{\textlengthmark}{\textsubring{d}}{\textsci}{l}{f}{j}{\oe}{r}{l}{\textepsilon}{\textbabygamma}{\textscy}{\textsubring{r}}} \dicPos{adj}[1]\dicFlx{}[-8] \textbf{1.} \dicSynonym{ómerkilegur} \dicDirectTranslationCS{nezajímavý, nezáživný, všední}  \textbf{2.} \dicSynonym{lasinn} \dicDirectTranslationCS{indisponovaný, slabý}
\dicEntry[lítill] \dicTerm{lítill} \dicsymFrequent\  \dicIPA{{l}{i}{\textlengthmark}{\textsubring{d}}{\textsci}{\textsubring{d}}{\textsubring{l}}} \dicPos{adj}[11] \dicFlx{(comp minni, sup minnstur)}[10] \textbf{1.} \dicSynonym{smár} \dicDirectTranslationCS{malý, drobný} \dicExampleIS{lítill maður} \dicExampleCS{malý člověk} \dicAntonym{stór}  \textbf{2.} \dicSynonym*{þann yngri til aðgreiningar frá þeim eldri} \dicDirectTranslationCS{malý, mladší} \dicExampleIS{Sölvi litli} \dicExampleCS{malý Sölvi}  \textbf{3.} \dicSynonym{óverulegur} \dicDirectTranslationCS{malý, zanedbatelný, nepatrný}  \textbf{4.} \dicSynonym{óduglegur} \dicDirectTranslationCS{zahálčivý, málo snaživý} \dicExampleIS{lítill íþróttamaður} \dicExampleCS{malý sportovec};  \dicIdiom{lítill}{ \dicPhraseIS{vera lítill í sér}} \dicDirectTranslationCS{být zbabělý}
\dicEntry[lítillátur] \dicTerm{lítil··látur} \dicIPA{{l}{i}{\textlengthmark}{\textsubring{d}}{\textsci}{l}{au}{\textsubring{d}}{\textscy}{\textsubring{r}}} \dicPos{adj}[1]\dicFlx{}[-1] \textbf{1.} \dicSynonym*{hrokalaus} \dicDirectTranslationCS{pokorný}  \textbf{2.} \dicSynonym{auðmjúkur} \dicDirectTranslationCS{nenáročný, skromný}
\dicEntry[lítillækka] \dicTerm{lítil··lækk|a} \dicIPA{{l}{i}{\textlengthmark}{\textsubring{d}}{\textsci}{l}{a}{i}{h}{\r{g}}{a}} \dicPos{v}[1] \dicFlx{(‑aði)}[1] \dicFlx{acc} \dicDirectTranslationCS{pokořit, pokořovat, ponížit, ponižovat} \dicExampleIS{lítillækka sig} \dicExampleCS{ponížit se}
\dicEntry[lítillækkun] \dicTerm{lítil··lækk|un} \dicIPA{{l}{i}{\textlengthmark}{\textsubring{d}}{\textsci}{l}{a}{i}{h}{\r{g}}{\textscy}{\textsubring{n}}} \dicPos{f}[7] \dicFlx{(‑unar)}[9] \dicDirectTranslationCS{pokoření, ponížení}
\dicEntry[lítillæti] \dicTerm{lítil··læti} \dicIPA{{l}{i}{\textlengthmark}{\textsubring{d}}{\textsci}{l}{a}{i}{\textsubring{d}}{\textsci}} \dicPos{n}[2] \dicFlx{(‑s)}[20] \textbf{1.} \dicSynonym{yfirlætisleysi} \dicDirectTranslationCS{pokora}  \textbf{2.} \dicSynonym{auðmýkt} \dicDirectTranslationCS{skromnost}
\dicEntry[lítilmagni] \dicTerm{lítil··magn|i} \dicIPA{{l}{i}{\textsubring{d}}{\textsci}{l}{m}{a}{\r{g}}{n}{\textsci}} \dicPos{m}[1] \dicFlx{(‑a, ‑ar)}[8] \dicDirectTranslationCS{padavka, outsider(ka)}
\dicEntry[lítilmenni] \dicTerm{lítil··menni} \dicIPA{{l}{i}{\textlengthmark}{\textsubring{d}}{\textsci}{l}{m}{\textepsilon}{n}{\textsci}} \dicPos{n}[2] \dicFlx{(‑s, ‑)}[14] \textbf{1.} \dicDirectTranslationCS{malý člověk (vzrůstem)}  \textbf{2.} \dicDirectTranslationCS{malicherný\,/\addthin nanicovatý člověk}
\dicEntry[lítilmótlegur] \dicTerm{lítil··mót·legur} \dicIPA{{l}{i}{\textlengthmark}{\textsubring{d}}{\textsci}{l}{m}{ou}{\textsubring{d}}{l}{\textepsilon}{\textbabygamma}{\textscy}{\textsubring{r}}} \dicPos{adj}[1]\dicFlx{}[-8] \dicSynonym{ómerkilegur} \dicDirectTranslationCS{ubohý, mizerný, mrzký}
\dicEntry[lítilræði] \dicTerm{lítil··ræði} \dicIPA{{l}{i}{\textlengthmark}{\textsubring{d}}{\textsci}{l}{r}{a}{i}{ð}{\textsci}} \dicPos{n}[2] \dicFlx{(‑s)}[20] \dicDirectTranslationCS{maličkost, maličko}
\dicEntry[lítilsháttar] \dicTerm{lítils··háttar\smash{\textsuperscript{1}}} \dicIPA{{l}{i}{\textlengthmark}{\textsubring{d}}{\textsci}{l}{s}{h}{au}{h}{\textsubring{d}}{a}{\textsubring{r}}} \dicPos{adj}[13] \dicFlx{indecl}[1] \dicDirectTranslationCS{malý, drobný}
\dicEntry[lítilsháttar] \dicTerm{lítils··háttar\smash{\textsuperscript{2}}} \dicIPA{{l}{i}{\textlengthmark}{\textsubring{d}}{\textsci}{l}{s}{h}{au}{h}{\textsubring{d}}{a}{\textsubring{r}}} \dicPos{adv} \dicDirectTranslationCS{málo, trochu}
\dicEntry[lítilsigldur] \dicTerm{lítil··sigldur} \dicIPA{{l}{i}{\textlengthmark}{\textsubring{d}}{\textsci}{l}{s}{\textsci}{l}{\textsubring{d}}{\textscy}{\textsubring{r}}} \dicPos{adj}[2]\dicFlx{}[-14] \dicSynonym{lítilfjörlegur} \dicDirectTranslationCS{nevýznamný, ubohý, nicotný}
\dicEntry[lítilsvirða] \dicTerm{lítils··vir|ða} \dicIPA{{l}{i}{\textlengthmark}{\textsubring{d}}{\textsci}{l}{s}{v}{\textsci}{r}{ð}{a}} \dicPos{v}[2] \dicFlx{(‑ti, ‑t)}[52] \dicFlx{acc} \dicDirectTranslationCS{pohrdat, opovrhovat}
\dicEntry[lítilsvirðing] \dicTerm{lítils··virð·ing} \dicIPA{{l}{i}{\textlengthmark}{\textsubring{d}}{\textsci}{l}{s}{v}{\textsci}{r}{ð}{i}{\ng}{\r{g}}} \dicPos{f}[4] \dicFlx{(‑ar)}[7] \dicSynonym{óvirðing} \dicDirectTranslationCS{opovržení, pohrdání}
\dicEntry[lítilvægur] \dicTerm{lítil··vægur} \dicIPA{{l}{i}{\textlengthmark}{\textsubring{d}}{\textsci}{l}{v}{a}{i}{\textbabygamma}{\textscy}{\textsubring{r}}} \dicPos{adj}[1]\dicFlx{}[-1] \dicSynonym*{þýðingarlítill} \dicDirectTranslationCS{nedůležitý, nepodstatný, bezvýznamný}
\dicEntry[lítilþægur] \dicTerm{lítil··þægur} \dicIPA{{l}{i}{\textlengthmark}{\textsubring{d}}{\textsci}{l}{\texttheta}{a}{i}{\textbabygamma}{\textscy}{\textsubring{r}}} \dicPos{adj}[1]\dicFlx{}[-1] \dicDirectTranslationCS{skromný, nenáročný} \dicIndirectTranslationCS{(často používané v~záporu)}
\dicEntry[lítri] \dicTerm{lítr|i} \dicIPA{{l}{i}{\textlengthmark}{\textsubring{d}}{r}{\textsci}} \dicPos{m}[1] \dicFlx{(‑a, ‑ar)}[1] \dicDirectTranslationCS{litr}
\dicEntry[lítt] \dicTerm{lítt} \dicsymFrequent\  \dicIPA{{l}{i}{h}{\textsubring{d}}} \dicPos{adv} \dicFlx{(comp miður, sup minnst)} \textbf{1.} \dicSynonym{lítið} \dicDirectTranslationCS{trochu, málo} \dicExampleIS{lítt þekktur ráðherra} \dicExampleCS{málo známý ministr}  \textbf{2.} \dicSynonym{illa} \dicDirectTranslationCS{špatně}
\dicEntry[ljá] \dicTerm{ljá} \dicIPA{{l}{j}{au}{\textlengthmark}} \dicPos{v}[5] \dicFlx{(ljæ, léði, léðum, léði, léð)}[15] \dicFlx{dat + acc} \textbf{1.} \dicLangCat{zast.} \dicDirectTranslationCS{(vy)půjčit} \dicExampleIS{ljá e‑m e‑ð} \dicExampleCS{půjčit (komu co)}  \textbf{2.} \dicPhraseIS{ljá e‑m eyra} \dicLangCat{přen.} \dicDirectTranslationCS{vyslyšet (koho), dopřát (komu) sluchu}  \textbf{3.} \dicDirectTranslationCS{vnést, vnášet, přidat, přidávat} \dicIndirectTranslationCS{(způsobit uplatnění, začlenění něčeho)}  \textbf{4.} \dicPhraseIS{ljá máls á e‑u} \dicDirectTranslationCS{zavést téma na (co)}
\dicEntry[ljár] \dicTerm{ljá|r} \dicIPA{{l}{j}{au}{\textlengthmark}{\textsubring{r}}} \dicPos{m}[9] \dicFlx{(‑s, ‑ir)}[12] \dicDirectTranslationCS{kosa} \dicIndirectTranslationCS{(ruční nástroj)};  \dicPhraseIS{slá með orfi og ljá} \dicDirectTranslationCS{sekat kosou};  \dicPhraseIS{vera e‑m óþægur ljár í þúfu} \dicLangCat{přen.} \dicDirectTranslationCS{být (komu) trnem v~oku}
\dicEntry[ljóð] \dicTerm{ljóð} \dicsymFrequent\  \dicIPA{{l}{j}{ou}{\textlengthmark}{\texttheta}} \dicPos{n}[2] \dicFlx{(‑s, ‑)}[5] \dicDirectTranslationCS{báseň} \dicExampleIS{ljóð eftir e‑n} \dicExampleCS{(čí) báseň}
\dicEntry[ljóðabók] \dicTerm{ljóða··|bók} \dicIPA{{l}{j}{ou}{\textlengthmark}{ð}{a}{\textsubring{b}}{ou}{\r{g}}} \dicPos{f}[8] \dicFlx{(‑bókar, ‑bækur)}[5] \dicDirectTranslationCS{sbírka\,/\addthin kniha poezie}
\dicEntry[ljóðaháttur] \dicTerm{ljóða··|háttur} \dicIPA{{l}{j}{ou}{\textlengthmark}{ð}{a}{h}{au}{h}{\textsubring{d}}{\textscy}{\textsubring{r}}} \dicPos{m}[12] \dicFlx{(‑háttar, ‑hættir)}[7] \dicFieldCat{lit.} \dicDirectTranslationCS{ed\nobreak dické metrum (šestiverší)}
\dicEntry[ljóðlist] \dicTerm{ljóð··list} \dicIPA{{l}{j}{ou}{ð}{l}{\textsci}{s}{\textsubring{d}}} \dicPos{f}[7] \dicFlx{(‑ar)}[3] \dicDirectTranslationCS{básnictví, básnické umění}
\dicEntry[ljóðmæli] \dicTerm{ljóð··mæli} \dicIPA{{l}{j}{ou}{ð}{m}{a}{i}{l}{\textsci}} \dicPos{n}[2] \dicFlx{pl}[19] \dicDirectTranslationCS{sbírka poezie}
\dicEntry[ljóðræna] \dicTerm{ljóð··ræn|a} \dicIPA{{l}{j}{ou}{ð}{r}{a}{i}{n}{a}} \dicPos{f}[1] \dicFlx{(‑u)}[5] \dicFieldCat{lit.} \dicDirectTranslationCS{lyrika}
\dicEntry[ljóðrænn] \dicTerm{ljóð··rænn} \dicIPA{{l}{j}{ou}{ð}{r}{a}{i}{\textsubring{d}}{\textsubring{n}}} \dicPos{adj}[7]\dicFlx{}[-1] \dicDirectTranslationCS{poetický, lyrický}
\dicEntry[ljóðskáld] \dicTerm{ljóð··skáld} \dicIPA{{l}{j}{ou}{ð}{s}{\r{g}}{au}{l}{\textsubring{d}}} \dicPos{n}[2] \dicFlx{(‑s, ‑)}[5] \dicDirectTranslationCS{básník, básnířka, poeta, poetka}
\dicEntry[ljóður] \dicTerm{ljóð|ur} \dicIPA{{l}{j}{ou}{\textlengthmark}{ð}{\textscy}{\textsubring{r}}} \dicPos{m}[6] \dicFlx{(‑s)}[26] \dicSynonym{lýti} \dicDirectTranslationCS{kaz, vada}
\dicEntry[ljóma] \dicTerm{ljóm|a} \dicIPA{{l}{j}{ou}{\textlengthmark}{m}{a}} \dicPos{v}[1] \dicFlx{(‑aði)}[1] \dicDirectTranslationCS{zářit, vyzařovat} \dicExampleIS{Sólin ljómar á himninum.} \dicExampleCS{Slunce září na obloze.};  \dicPhraseIS{ljóma af gleði} \dicLangCat{přen.} \dicDirectTranslationCS{zářit štěstím}
\dicEntry[ljómandi] \dicTerm{ljóm··andi\smash{\textsuperscript{1}}} \dicIPA{{l}{j}{ou}{\textlengthmark}{m}{a}{n}{\textsubring{d}}{\textsci}} \dicPos{adj}[13] \dicFlx{indecl}[1] \dicDirectTranslationCS{zářivý, zářící} \dicExampleIS{ljómandi sól} \dicExampleCS{zářící slunce}
\dicEntry[ljómandi] \dicTerm{ljóm··andi\smash{\textsuperscript{2}}} \dicIPA{{l}{j}{ou}{\textlengthmark}{m}{a}{n}{\textsubring{d}}{\textsci}} \dicPos{adv} \dicDirectTranslationCS{úžasně, fantasticky} \dicExampleIS{ljómandi vel} \dicExampleCS{fantasticky dobře}
\dicEntry[ljómi] \dicTerm{ljóm|i} \dicIPA{{l}{j}{ou}{\textlengthmark}{m}{\textsci}} \dicPos{m}[1] \dicFlx{(‑a)}[3] \textbf{1.} \dicSynonym{birta\smash{\textsuperscript{1}}} \dicDirectTranslationCS{zář(e), svit}  \textbf{2.} \dicFieldCat{fyz.} \dicDirectTranslationCS{jas}
\dicEntry[ljón] \dicTerm{ljón} \dicsymFrequent\  \dicIPA{{l}{j}{ou}{\textlengthmark}{\textsubring{n}}} \dicPos{n}[2] \dicFlx{(‑s, ‑)}[5] \textbf{1.} \dicFieldCat{zool.} \dicDirectTranslationCS{lev} \textit{(l.~{\textLA{Panthera leo}})}  \dicsymPhoto\  \dicExampleIS{Ljón hleypur í átt að manni.} \dicExampleCS{Směrem k~člověku běží lev.}  \textbf{2.} \dicDirectTranslationCS{Lev} \dicIndirectTranslationCS{(znamení zvěrokruhu)} \textit{(l.~{\textLA{Leo}})}
\dicFigure{20786.jpg}{Ljón}{Ljón - Stansell, Ken, Biolib, PD}
\dicEntry[ljónatemjari] \dicTerm{ljóna··temj·ar|i} \dicIPA{{l}{j}{ou}{\textlengthmark}{n}{a}{t\smash{\textsuperscript{h}}}{\textepsilon}{m}{j}{a}{r}{\textsci}} \dicPos{m}[1] \dicFlx{(‑a, ‑ar)}[13] \dicDirectTranslationCS{krotitel(ka) lvů\,/\addthin šelem\,/\addthin dravé zvěře}
\dicEntry[ljónslappi] \dicTerm{ljóns··lapp|i} \dicIPA{{l}{j}{ou}{n}{s}{l}{a}{h}{\textsubring{b}}{\textsci}} \dicPos{m}[1] \dicFlx{(‑a)}[3] \dicFieldCat{bot.} \dicDirectTranslationCS{kontryhel alpinský} \textit{(l.~{\textLA{Alchemilla alpina}})}  \dicsymPhoto\ 
\dicFigure{88632.jpg}{Ljónslappi}{Ljónslappi - Pleva František, Biolib, PD}
\dicEntry[ljónynja] \dicTerm{ljón··ynj|a} \dicIPA{{l}{j}{ou}{\textlengthmark}{n}{\textsci}{n}{j}{a}} \dicPos{f}[1] \dicFlx{(‑u, ‑ur)}[7] \dicDirectTranslationCS{lvice}
\dicEntry[ljóri] \dicTerm{ljór|i} \dicIPA{{l}{j}{ou}{\textlengthmark}{r}{\textsci}} \dicPos{m}[1] \dicFlx{(‑a, ‑ar)}[1] \dicDirectTranslationCS{světlík}
\dicEntry[ljós] \dicTerm{ljós\smash{\textsuperscript{1}}} \dicsymFrequent\  \dicIPA{{l}{j}{ou}{\textlengthmark}{s}} \dicPos{n}[2] \dicFlx{(‑s, ‑)}[5] \textbf{1.} \dicSynonym{birta\smash{\textsuperscript{1}}} \dicDirectTranslationCS{světlo, záře, jas, svit} \dicExampleIS{ljósið frá lampanum} \dicExampleCS{světlo lampy};  \dicPhraseIS{kveikja ljósið} \dicDirectTranslationCS{rozsvítit, zapnout\,/\addthin zapálit světlo};  \dicPhraseIS{slökkva ljósið} \dicDirectTranslationCS{zhasnout, vypnout světlo}  \textbf{2.} \dicSynonym*{þægt barn} \dicDirectTranslationCS{zlatíčko, andílek} \dicIndirectTranslationCS{(o~poslušném dítěti)};  \dicIdiom{ljós}{ \dicPhraseIS{í nýju ljósi}} \dicFlx{adv} \dicLangCat{přen.} \dicDirectTranslationCS{v~novém světle}; { \dicPhraseIS{það kom í ljós að}} \dicDirectTranslationCS{ukázalo se, že; vyšlo najevo, že}
\dicEntry[ljós] \dicTerm{ljós\smash{\textsuperscript{2}}} \dicsymFrequent\  \dicIPA{{l}{j}{ou}{\textlengthmark}{s}} \dicPos{adj}[5]\dicFlx{}[-1] \dicDirectTranslationCS{jasný, světlý} \dicExampleIS{vera ljós á húð} \dicExampleCS{mít světlou pokožku};  \dicPhraseIS{gera sér e‑ð ljóst} \dicDirectTranslationCS{uvědomit si (co)};  \dicPhraseIS{e‑m er ljóst e‑að} \dicFlx{impers} \dicDirectTranslationCS{(co) je (komu) jasné};  \dicPhraseIS{e‑að er deginum ljósara} \dicLangCat{přen.} \dicDirectTranslationCS{(co) je nad slunce jasnější, (co) je jasné jak facka}
\dicEntry[ljósakróna] \dicTerm{ljósa··krón|a} \dicIPA{{l}{j}{ou}{\textlengthmark}{s}{a}{k\smash{\textsuperscript{h}}}{r}{ou}{n}{a}} \dicPos{f}[1] \dicFlx{(‑u, ‑ur)}[7] \dicDirectTranslationCS{(mnohoramenný) lustr}
\dicEntry[ljósapera] \dicTerm{ljósa··per|a} \dicIPA{{l}{j}{ou}{\textlengthmark}{s}{a}{p\smash{\textsuperscript{h}}}{\textepsilon}{r}{a}} \dicPos{f}[1] \dicFlx{(‑u, ‑ur)}[7] \dicFieldCat{elek.} \dicDirectTranslationCS{žárovka}
\dicEntry[ljósaskilti] \dicTerm{ljósa··skilti} \dicIPA{{l}{j}{ou}{\textlengthmark}{s}{a}{s}{\r{\textObardotlessj}}{\textsci}{\textsubring{l}}{\textsubring{d}}{\textsci}} \dicPos{n}[2] \dicFlx{(‑s, ‑)}[14] \dicDirectTranslationCS{neonová reklama}
\dicEntry[ljósaskipti] \dicTerm{ljósa··skipti} \dicIPA{{l}{j}{ou}{\textlengthmark}{s}{a}{s}{\r{\textObardotlessj}}{\textsci}{f}{\textsubring{d}}{\textsci}} \dicPos{n}[2] \dicFlx{pl}[19] \dicSynonym{rökkur} \dicDirectTranslationCS{soumrak}
\dicEntry[ljósastaur] \dicTerm{ljósa··staur} \dicIPA{{l}{j}{ou}{\textlengthmark}{s}{a}{s}{\textsubring{d}}{\oe i}{\textsubring{r}}} \dicPos{m}[4] \dicFlx{(‑s, ‑ar)}[14] \dicDirectTranslationCS{kandelábr, stojan veřejného osvětlení}
\dicEntry[ljósastæði] \dicTerm{ljósa··stæði} \dicIPA{{l}{j}{ou}{\textlengthmark}{s}{a}{s}{\textsubring{d}}{a}{i}{ð}{\textsci}} \dicPos{n}[2] \dicFlx{(‑s, ‑)}[14] \dicDirectTranslationCS{přípojka na světlo}
\dicEntry[Ljósálfheimur] \dicTerm{Ljós··álf·heim|ur}\dicTerm{, Álfheimur} \dicIPA{{l}\-{j}\-{ou}\-{\textlengthmark}\-{s}\-{au}\-{l}\-{f}\-{h}\-{ei}\-{m}\-{\textscy}\-{\textsubring{r}}\-} \dicPos{m}[6] \dicFlx{(‑s)}[3] \dicFieldCat{myt.} \dicDirectTranslationCS{Álfheim} \dicIndirectTranslationCS{(říše světla, světlých elfů)} \dicAntonym{Svartálfheimur}
\dicEntry[ljósár] \dicTerm{ljós··ár} \dicIPA{{l}{j}{ou}{\textlengthmark}{s}{au}{\textsubring{r}}} \dicPos{n}[2] \dicFlx{(‑s, ‑)}[5] \dicFieldCat{astro.} \dicDirectTranslationCS{světelný rok}
\dicEntry[ljósberi] \dicTerm{ljós··ber|i} \dicIPA{{l}{j}{ou}{s}{\textsubring{b}}{\textepsilon}{r}{\textsci}} \dicPos{m}[1] \dicFlx{(‑a, ‑ar)}[1] \dicFieldCat{bot.} \dicDirectTranslationCS{kohoutek alpský} \textit{(l.~{\textLA{Lychnis alpina}})}  \dicsymPhoto\ 
\dicFigure{72890.jpg}{Ljósberi}{Ljósberi - Kesl Michael, Biolib, Copyright/CC-BY-NC}
\dicEntry[ljóseind] \dicTerm{ljós··eind} \dicIPA{{l}{j}{ou}{\textlengthmark}{s}{ei}{n}{\textsubring{d}}} \dicPos{f}[7] \dicFlx{(‑ar, ‑ir)}[1] \dicFieldCat{fyz.} \dicDirectTranslationCS{foton}
\dicEntry[ljósfræði] \dicTerm{ljós··fræð|i} \dicIPA{{l}{j}{ou}{s}{f}{r}{a}{i}{ð}{\textsci}} \dicPos{f}[3] \dicFlx{(‑i)}[3] \dicFieldCat{fyz.} \dicDirectTranslationCS{optika}
\dicEntry[ljósgeisli] \dicTerm{ljós··geisl|i} \dicIPA{{l}{j}{ou}{s}{\r{\textObardotlessj}}{ei}{s}{\textsubring{d}}{l}{\textsci}} \dicPos{m}[1] \dicFlx{(‑a, ‑ar)}[1] \dicDirectTranslationCS{světelný paprsek, paprsek světla}
\dicEntry[ljósgulur] \dicTerm{ljós··gulur} \dicIPA{{l}{j}{ou}{s}{\r{g}}{\textscy}{l}{\textscy}{\textsubring{r}}} \dicPos{adj}[1]\dicFlx{}[-1] \dicDirectTranslationCS{světle žlutý, nažloutlý}
\dicEntry[ljóshraði] \dicTerm{ljós··hrað|i} \dicIPA{{l}{j}{ou}{\textlengthmark}{s}{\textsubring{r}}{a}{ð}{\textsci}} \dicPos{m}[1] \dicFlx{(‑a)}[3] \dicFieldCat{fyz.} \dicDirectTranslationCS{rychlost světla}
\dicEntry[ljóshærður] \dicTerm{ljós··hærður} \dicsymFrequent\  \dicIPA{{l}{j}{ou}{\textlengthmark}{s}{h}{a}{i}{r}{ð}{\textscy}{\textsubring{r}}} \dicPos{adj}[2]\dicFlx{}[-1] \dicDirectTranslationCS{blonďatý, světlovlasý} \dicExampleIS{Hann er ljóshærður.} \dicExampleCS{Má světlé vlasy.}
\dicEntry[ljóshöfðaönd] \dicTerm{ljós·höfða··|önd} \dicIPA{{l}{j}{ou}{\textlengthmark}{s}{h}{\oe}{v}{ð}{a}{\oe}{n}{\textsubring{d}}} \dicPos{f}[8] \dicFlx{(‑andar, ‑endur)}[2] \dicFieldCat{zool.} \dicDirectTranslationCS{hvízdák americký} \textit{(l.~{\textLA{Anas americana}})}  \dicsymPhoto\ 
\dicFigure{ds_image_ljoshofdaond_0_2.jpg}{Ljóshöfðaönd}{Ljóshöfðaönd - Donna Dewhurst, USFWS, PD}
\dicEntry[ljóska] \dicTerm{ljósk|a} \dicIPA{{l}{j}{ou}{s}{\r{g}}{a}} \dicPos{f}[1] \dicFlx{(‑u, ‑ur)}[19] \dicDirectTranslationCS{blondýna, blondýnka}
\dicEntry[ljóskastari] \dicTerm{ljós··kast·ar|i} \dicIPA{{l}{j}{ou}{s}{k\smash{\textsuperscript{h}}}{a}{s}{\textsubring{d}}{a}{r}{\textsci}} \dicPos{m}[1] \dicFlx{(‑a, ‑ar)}[10] \dicDirectTranslationCS{světlomet, reflektor}
\dicEntry[ljósker] \dicTerm{ljós··ker} \dicIPA{{l}{j}{ou}{s}{c\smash{\textsuperscript{h}}}{\textepsilon}{\textsubring{r}}} \dicPos{n}[2] \dicFlx{(‑s, ‑)}[6] \dicDirectTranslationCS{lampa, svítilna, lucerna}
\dicEntry[ljóslifandi] \dicTerm{ljós··lif·andi} \dicIPA{{l}{j}{ou}{s}{l}{\textsci}{v}{a}{n}{\textsubring{d}}{\textsci}} \dicPos{adj}[13] \dicFlx{indecl}[1] \dicDirectTranslationCS{živý, barvitý} \dicExampleIS{e‑að stendur e‑m ljóslifandi fyrir hugskotssjónum} \dicExampleCS{(co) má (kdo) v~živé paměti}
\dicEntry[ljósmerki] \dicTerm{ljós··merki} \dicIPA{{l}{j}{ou}{s}{m}{\textepsilon}{\textsubring{r}}{\r{\textObardotlessj}}{\textsci}} \dicPos{n}[2] \dicFlx{(‑s, ‑)}[16] \dicDirectTranslationCS{světelný signál, světelné znamení}
\dicEntry[ljósmóðir] \dicTerm{ljós··|móðir} \dicIPA{{l}{j}{ou}{s}{m}{ou}{ð}{\textsci}{\textsubring{r}}} \dicPos{f}[11] \dicFlx{(‑móður, ‑mæður)}[1] \dicDirectTranslationCS{porodní asistentka}
\dicEntry[ljósmynd] \dicTerm{ljós··mynd} \dicsymFrequent\  \dicIPA{{l}{j}{ou}{s}{m}{\textsci}{n}{\textsubring{d}}} \dicPos{f}[7] \dicFlx{(‑ar, ‑ir)}[1] \dicDirectTranslationCS{fotografie, fotka} \dicExampleIS{svarthvít ljósmynd af ungri konu} \dicExampleCS{černobílá fotografie mladé ženy}
\dicEntry[ljósmynda] \dicTerm{ljós··mynd|a} \dicIPA{{l}{j}{ou}{s}{m}{\textsci}{n}{\textsubring{d}}{a}} \dicPos{v}[1] \dicFlx{(‑aði)}[1] \dicDirectTranslationCS{(vy)fotit, (vy)fotografovat, (u)dělat fotografii}
\dicEntry[ljósmyndari] \dicTerm{ljós··mynd·ar|i} \dicIPA{{l}{j}{ou}{s}{m}{\textsci}{n}{\textsubring{d}}{a}{r}{\textsci}} \dicPos{m}[1] \dicFlx{(‑a, ‑ar)}[13] \dicDirectTranslationCS{fotograf(ka)}
\dicEntry[ljósmyndavél] \dicTerm{ljós··mynda·vél} \dicIPA{{l}{j}{ou}{s}{m}{\textsci}{n}{\textsubring{d}}{a}{v}{j}{\textepsilon}{\textsubring{l}}} \dicPos{f}[4] \dicFlx{(‑ar, ‑ar)}[1] \dicDirectTranslationCS{fotoaparát}
\dicEntry[ljósmyndun] \dicTerm{ljós··mynd|un} \dicIPA{{l}{j}{ou}{s}{m}{\textsci}{n}{\textsubring{d}}{\textscy}{\textsubring{n}}} \dicPos{f}[7] \dicFlx{(‑unar, ‑anir)}[8] \textbf{1.} \dicDirectTranslationCS{fotografování}  \textbf{2.} \dicDirectTranslationCS{fotografie} \dicIndirectTranslationCS{(umělecký obor)}
\dicEntry[ljósmælir] \dicTerm{ljós··mæl|ir} \dicIPA{{l}{j}{ou}{s}{m}{a}{i}{l}{\textsci}{\textsubring{r}}} \dicPos{m}[7] \dicFlx{(‑is, ‑ar)}[1] \dicFieldCat{fyz.} \dicDirectTranslationCS{expozimetr}
\dicEntry[ljósnemi] \dicTerm{ljós··nem|i} \dicIPA{{l}{j}{ou}{s}{n}{\textepsilon}{m}{\textsci}} \dicPos{m}[1] \dicFlx{(‑a, ‑ar)}[1] \dicFieldCat{fyz.} \dicDirectTranslationCS{fotočlánek}
\dicEntry[ljósnæmur] \dicTerm{ljós··næmur} \dicIPA{{l}{j}{ou}{s}{n}{a}{i}{m}{\textscy}{\textsubring{r}}} \dicPos{adj}[1]\dicFlx{}[-1] \dicDirectTranslationCS{citlivý na světlo, fotosenzitivní}
\dicEntry[ljósrit] \dicTerm{ljós··rit} \dicIPA{{l}{j}{ou}{s}{r}{\textsci}{\textsubring{d}}} \dicPos{n}[2] \dicFlx{(‑s, ‑)}[5] \dicDirectTranslationCS{fotokopie}
\dicEntry[ljósrita] \dicTerm{ljós··rit|a} \dicIPA{{l}{j}{ou}{s}{r}{\textsci}{\textsubring{d}}{a}} \dicPos{v}[1] \dicFlx{(‑aði)}[1] \dicFlx{acc} \dicDirectTranslationCS{(o)kopírovat (pomocí fotokopírky)}
\dicEntry[ljósritunarvél] \dicTerm{ljós·ritunar··vél} \dicIPA{{l}{j}{ou}{s}{r}{\textsci}{\textsubring{d}}{\textscy}{n}{a}{r}{v}{j}{\textepsilon}{\textsubring{l}}} \dicPos{f}[4] \dicFlx{(‑ar, ‑ar)}[1] \dicDirectTranslationCS{(foto)kopírka}
\dicEntry[ljósta] \dicTerm{ljósta} \dicsymFrequent\  \dicIPA{{l}{j}{ou}{s}{\textsubring{d}}{a}} \dicPos{v}[6] \dicFlx{(lýst, laust, lustum, lysti, lostið)}[97] \dicFlx{acc} \dicLangCat{zast.} \dicDirectTranslationCS{udeřit, uhodit} \dicExampleIS{ljósta e‑n í andlitið með e‑u} \dicExampleCS{uhodit (koho) do tváře (čím)};  \dicIdiom{ljósta}[niður]{ \dicPhraseIS{e‑u lýstur niður}} \dicFlx{impers} \dicDirectTranslationCS{(co) udeří, (co) zasáhne (blesk ap.)} \dicExampleIS{Eldingu laust niður í húsið.} \dicExampleCS{Blesk udeřil do domu.};  \dicIdiom{ljósta}[saman]{ \dicPhraseIS{e‑m lýstur saman}} \dicFlx{impers} \dicDirectTranslationCS{(kdo) se střetne} \dicExampleIS{Herjunum laust saman.} \dicExampleCS{Armády se střetly.};  \dicIdiom{ljósta}[upp]{ \dicPhraseIS{ljósta upp ópi}} \dicDirectTranslationCS{vykřiknout}
\dicEntry[ljóstillífun] \dicTerm{ljós··til·líf|un} \dicIPA{{l}{j}{ou}{s}{t\smash{\textsuperscript{h}}}{\textsci}{l}{i}{v}{\textscy}{\textsubring{n}}} \dicPos{f}[7] \dicFlx{(‑unar)}[9] \dicFieldCat{biol.} \dicDirectTranslationCS{fotosyntéza}
\dicEntry[ljóstra] \dicTerm{ljóstr|a} \dicIPA{{l}{j}{ou}{s}{\textsubring{d}}{r}{a}} \dicPos{v}[1] \dicFlx{(‑aði)}[1] \dicFlx{dat} \dicPhraseIS{ljóstra e‑u upp} \dicDirectTranslationCS{odhalit\,/\addthin odhalovat (co), odkrýt\,/\addthin odkrývat (co)} \dicExampleIS{ljóstra upp leyndarmáli} \dicExampleCS{odkrýt tajemství}
\dicEntry[ljósvaki] \dicTerm{ljós··vak|i} \dicIPA{{l}{j}{ou}{\textlengthmark}{s}{v}{a}{\r{\textObardotlessj}}{\textsci}} \dicPos{m}[1] \dicFlx{(‑a, ‑ar)}[8] \dicDirectTranslationCS{éter} \dicIndirectTranslationCS{(podle starých představ látka zaplňující vesmírný prostor)};  \dicPhraseIS{á öldum ljósvakans} \dicFlx{adv} \dicDirectTranslationCS{na vlnách éteru}
\dicEntry[ljótleiki] \dicTerm{ljót··leik|i} \dicIPA{{l}{j}{ou}{\textlengthmark}{\textsubring{d}}{l}{ei}{\r{\textObardotlessj}}{\textsci}} \dicPos{m}[1] \dicFlx{(‑a)}[3] \dicSynonym{lýti} \dicDirectTranslationCS{šerednost, ohyzdnost (vzhled ap.)}
\dicEntry[ljótur] \dicTerm{ljótur} \dicsymFrequent\  \dicIPA{{l}{j}{ou}{\textlengthmark}{\textsubring{d}}{\textscy}{\textsubring{r}}} \dicPos{adj}[1]\dicFlx{}[-1] \dicDirectTranslationCS{šeredný, ošklivý, škaredý} \dicExampleIS{ljótar athugasemdir} \dicExampleCS{ošklivé poznámky}
\dicEntry[Ljubljana] \dicTerm{Ljubljana} \dicIPA{{l}{j}{\textscy}{\textsubring{b}}{l}{j}{a}{n}{a}} \dicPos{subs} \dicFlx{indecl} \dicFieldCat{geog.} \dicDirectTranslationCS{Lublaň} \dicIndirectTranslationCS{(hlavní město Slovinska)}
\dicEntry[ljúffengur] \dicTerm{ljúf··fengur} \dicIPA{{l}{j}{u}{f}{\textlengthmark}{ei}{\ng}{\r{g}}{\textscy}{\textsubring{r}}} \dicPos{adj}[1]\dicFlx{}[-1] \dicDirectTranslationCS{chutný, lahodný, delikátní} \dicIndirectTranslationCS{(o~jídle)}
\dicEntry[ljúflingur] \dicTerm{ljúf··ling|ur} \dicIPA{{l}{j}{u}{v}{l}{i}{\ng}{\r{g}}{\textscy}{\textsubring{r}}} \dicPos{m}[6] \dicFlx{(‑s, ‑ar)}[8] \textbf{1.} \dicSynonym{huldumaður} \dicDirectTranslationCS{elf}  \textbf{2.} \dicSynonym{eftirlæti} \dicDirectTranslationCS{miláček, zlato}
\dicEntry[ljúflyndi] \dicTerm{ljúf··lyndi} \dicIPA{{l}{j}{u}{v}{l}{\textsci}{n}{\textsubring{d}}{\textsci}} \dicPos{n}[2] \dicFlx{(‑s)}[20] \dicDirectTranslationCS{přívětivost, laskavost}
\dicEntry[ljúfmenni] \dicTerm{ljúf··menni} \dicIPA{{l}{j}{u}{v}{m}{\textepsilon}{n}{\textsci}} \dicPos{n}[2] \dicFlx{(‑s, ‑)}[14] \dicSynonym{góðmenni} \dicDirectTranslationCS{milý\,/\addthin přívětivý člověk}
\dicEntry[ljúfmennska] \dicTerm{ljúf··mennsk|a} \dicIPA{{l}{j}{u}{v}{m}{\textepsilon}{n}{s}{\r{g}}{a}} \dicPos{f}[1] \dicFlx{(‑u)}[5] \dicDirectTranslationCS{dobrota, přívětivost, vlídnost}
\dicEntry[ljúfur] \dicTerm{ljúf|ur\smash{\textsuperscript{1}}} \dicIPA{{l}{j}{u}{\textlengthmark}{v}{\textscy}{\textsubring{r}}} \dicPos{m}[6] \dicFlx{(‑s, ‑ar)}[24] \dicDirectTranslationCS{poklad, drahoušek, miláček} \dicExampleIS{ljúfur(inn) (minn)} \dicExampleCS{ty můj poklade}
\dicEntry[ljúfur] \dicTerm{ljúfur\smash{\textsuperscript{2}}} \dicsymFrequent\  \dicIPA{{l}{j}{u}{\textlengthmark}{v}{\textscy}{\textsubring{r}}} \dicPos{adj}[1]\dicFlx{}[-1] \textbf{1.} \dicSynonym{blíður} \dicDirectTranslationCS{laskavý, přívětivý, vlídný}  \textbf{2.} \dicSynonym{þægilegur} \dicDirectTranslationCS{příjemný, milý} \dicExampleIS{ljúfar minningar} \dicExampleCS{příjemné vzpomínky};  \dicPhraseIS{e‑m er ljúft að (gera e‑ð)} \dicFlx{impers} \dicDirectTranslationCS{(komu) je příjemné (udělat (co))}
\dicEntry[ljúga] \dicTerm{ljúga} \dicsymFrequent\  \dicIPA{{l}{j}{u}{\textlengthmark}{a}} \dicPos{v}[6] \dicFlx{(lýg, laug, lugum, lygi, logið)}[52] \dicFlx{dat} \dicDirectTranslationCS{lhát, zalhat} \dicExampleIS{ljúga e‑u} \dicExampleCS{zalhat o~(čem)};  \dicPhraseIS{ljúga að e‑m} \dicDirectTranslationCS{lhát (komu)};  \dicPhraseIS{ljúga til um e‑ð} \dicDirectTranslationCS{lhát o~(čem)} \dicExampleIS{ljúga til um aldur} \dicExampleCS{lhát o~(svém) věku}
\dicEntry[ljúgvitni] \dicTerm{ljúg··vitni} \dicIPA{{l}{j}{u}{\textlengthmark}{v}{\textsci}{h}{\textsubring{d}}{n}{\textsci}} \dicPos{n}[2] \dicFlx{(‑s, ‑)}[14] \dicDirectTranslationCS{křivá přísaha, křivé doznání};  \dicPhraseIS{bera ljúgvitni gegn e‑m} \dicDirectTranslationCS{křivě svědčit proti (komu)}
\dicEntry[ljúka] \dicTerm{ljúka} \dicsymFrequent\  \dicIPA{{l}{j}{u}{\textlengthmark}{\r{g}}{a}} \dicPos{v}[6] \dicFlx{(lýk, lauk, lukum, lyki, lokið)}[47] \dicFlx{dat} \textbf{1.} \dicSynonym{klára} \dicDirectTranslationCS{(u)končit, ukončovat, zakončit, zakončovat, dokončit, dokončovat};  \dicPhraseIS{ljúka e‑u} \dicDirectTranslationCS{ukončit (co)};  \dicPhraseIS{ljúka við e‑ð} \dicDirectTranslationCS{ukončit (co), (s)končit s~(čím)}  \textbf{2.} \dicPhraseIS{e‑u lýkur} \dicFlx{impers} \dicDirectTranslationCS{(co) končí} \dicExampleIS{Skólanum lýkur í júní.} \dicExampleCS{Škola končí v~červnu.};  \dicIdiom{ljúka}[af]{ \dicPhraseIS{ljúka e‑u af}} \dicDirectTranslationCS{ukončit (co), skončit (co)};  \dicIdiom{ljúka}[upp]{ \dicPhraseIS{ljúka upp e‑u}} \dicLangCat{zast.} \dicSynonym{opna\smash{\textsuperscript{2}}} \dicDirectTranslationCS{otevřít (co)} \dicExampleIS{ljúka dyrunum upp} \dicExampleCS{otevřít dveře};  \dicIdiom{ljúka}[yfir]{ \dicPhraseIS{þar til yfir lýkur}} \dicFlx{adv} \dicDirectTranslationCS{až do samého konce}; { \dicPhraseIS{áður en yfir lýkur}} \dicFlx{adv} \dicDirectTranslationCS{dříve než bude po všem};  \dicIdiom{ljúkast}[saman]{ \dicPhraseIS{ljúkast saman}} \dicFlx{refl} \dicLangCat{zast.} \dicDirectTranslationCS{zavřít se};  \dicIdiom{ljúkast}[upp]{ \dicPhraseIS{ljúkast upp}} \dicFlx{refl} \dicLangCat{zast.} \dicDirectTranslationCS{otevřít se, rozevřít se}; { \dicPhraseIS{e‑að lýkst upp fyrir e‑m}} \dicFlx{refl} \dicLangCat{přen.} \dicDirectTranslationCS{(co) se otevírá před (kým) (nová myšlenka ap.)}
\dicEntry[ljæ] \dicTerm{ljæ} \dicIPA{{l}{j}{a}{i}{\textlengthmark}} \dicPos{v} \dicFlx{ind praes sg 1 pers} \dicLink{ljá}
\dicEntry[lo.] \dicTerm{lo.} \dicPos{zkr} \dicPhraseIS{lýsingarorð} \dicFieldCat{jaz.} \dicDirectTranslationCS{přídavné jméno}
\dicEntry[lobbíisti] \dicTerm{lobbí··ist|i} \dicIPA{{l}{\textopeno}{\textsubring{b}}{\textlengthmark}{i}{j}{\textsci}{s}{\textsubring{d}}{\textsci}} \dicPos{m}[1] \dicFlx{(‑a, ‑ar)}[1] \dicDirectTranslationCS{lobbista, lobbistka}
\dicEntry[loddari] \dicTerm{lodd··ar|i} \dicIPA{{l}{\textopeno}{\textsubring{d}}{\textlengthmark}{a}{r}{\textsci}} \dicPos{m}[1] \dicFlx{(‑a, ‑ar)}[13] \textbf{1.} \dicSynonym{svikari} \dicDirectTranslationCS{šarlatán(ka)}  \textbf{2.} \dicSynonym{trúður} \dicDirectTranslationCS{šašek, klaun(ka)}
\dicEntry[loða] \dicTerm{lo|ða} \dicIPA{{l}{\textopeno}{\textlengthmark}{ð}{a}} \dicPos{v}[2] \dicFlx{(‑ddi, ‑ðað)}[182] \dicSynonym{tolla\smash{\textsuperscript{2}}} \dicDirectTranslationCS{přiléhat, lnout, přimknout se, přimykat};  \dicIdiom{loða}[saman]{ \dicPhraseIS{loða saman}} \dicDirectTranslationCS{přilnout k~sobě};  \dicIdiom{loða}[við]{ \dicPhraseIS{loða við e‑ð}} \dicDirectTranslationCS{přilnout k~(čemu)} \dicExampleIS{Málningin loðir við vegginn.} \dicExampleCS{Nátěr přilne ke zdi.}
\dicEntry[loðdýr] \dicTerm{loð··dýr} \dicIPA{{l}{\textopeno}{ð}{\textsubring{d}}{i}{\textsubring{r}}} \dicPos{n}[2] \dicFlx{(‑s, ‑)}[5] \dicDirectTranslationCS{kožešinové zvíře}
\dicEntry[loðfeldur] \dicTerm{loð··feld|ur} \dicIPA{{l}{\textopeno}{ð}{f}{\textepsilon}{l}{\textsubring{d}}{\textscy}{\textsubring{r}}} \dicPos{m}[10] \dicFlx{(‑ar, ‑ir)}[4] \dicSynonym{loðskinn} \dicDirectTranslationCS{kožešina}
\dicEntry[loðfíll] \dicTerm{loð··fíl|l} \dicIPA{{l}{\textopeno}{ð}{f}{i}{\textsubring{d}}{\textsubring{l}}} \dicPos{m}[6] \dicFlx{(‑s, ‑ar)}[48] \dicFieldCat{zool.} \dicDirectTranslationCS{mamut} \textit{(l.~{\textLA{Mammuthus primigenius}})}
\dicEntry[loðhúfa] \dicTerm{loð··húf|a} \dicIPA{{l}{\textopeno}{\textlengthmark}{\texttheta}{h}{u}{v}{a}} \dicPos{f}[1] \dicFlx{(‑u, ‑ur)}[19] \dicDirectTranslationCS{kožešinová čepice, beranice}
\dicEntry[loðinn] \dicTerm{loðinn} \dicIPA{{l}{\textopeno}{\textlengthmark}{ð}{\textsci}{\textsubring{n}}} \dicPos{adj}[6]\dicFlx{}[-2] \textbf{1.} \dicDirectTranslationCS{chlupatý, srstnatý, huňatý}  \textbf{2.} \dicSynonym{óhreinn} \dicDirectTranslationCS{vyhýbavý, mlhavý (odpověď ap.)}
\dicEntry[loðkápa] \dicTerm{loð··káp|a} \dicIPA{{l}{\textopeno}{ð}{k\smash{\textsuperscript{h}}}{au}{\textsubring{b}}{a}} \dicPos{f}[1] \dicFlx{(‑u, ‑ur)}[19] \dicDirectTranslationCS{kožich}
\dicEntry[loðna] \dicTerm{loðn|a} \dicIPA{{l}{\textopeno}{ð}{n}{a}} \dicPos{f}[1] \dicFlx{(‑u, ‑ur)}[7] \textbf{1.} \dicDirectTranslationCS{ochlupení}  \textbf{2.} \dicFieldCat{zool.} \dicDirectTranslationCS{huňáček severní} \textit{(l.~{\textLA{Mallotus villosus}})}
\dicEntry[loðskinn] \dicTerm{loð··skinn} \dicIPA{{l}{\textopeno}{ð}{s}{\r{\textObardotlessj}}{\textsci}{\textsubring{n}}} \dicPos{n}[2] \dicFlx{(‑s, ‑)}[5] \dicDirectTranslationCS{kožešina, kožka}
\dicEntry[loðvíðir] \dicTerm{loð··víð|ir} \dicIPA{{l}{\textopeno}{ð}{v}{i}{ð}{\textsci}{\textsubring{r}}} \dicPos{m}[7] \dicFlx{(‑is)}[2] \dicFieldCat{bot.} \dicDirectTranslationCS{vrba vlnatá} \textit{(l.~{\textLA{Salix lanata}})}  \dicsymPhoto\ 
\dicFigure{66869.jpg}{Loðvíðir}{Loðvíðir - Hroneš Michal, Biolib, Copyright/CC-BY-NC}
\dicEntry[lof] \dicTerm{lof} \dicIPA{{l}{\textopeno}{\textlengthmark}{f}} \dicPos{n}[2] \dicFlx{(‑s)}[2] \dicSynonym{hrós} \dicDirectTranslationCS{chvála, pochvala};  \dicPhraseIS{bera lof á e‑n} \dicDirectTranslationCS{chválit (koho)};  \dicPhraseIS{guði sé lof} \dicDirectTranslationCS{díky bohu};  \dicPhraseIS{klappa lof í lófa} \dicDirectTranslationCS{tleskat, aplaudovat}
\dicEntry[lofa] \dicTerm{lof|a} \dicsymFrequent\  \dicIPA{{l}{\textopeno}{\textlengthmark}{v}{a}} \dicPos{v}[1] \dicFlx{(‑aði)}[1] \dicFlx{acc\,/\addthin dat} \textbf{1.} \dicFlx{acc} \dicSynonym{hrósa} \dicDirectTranslationCS{(po)chválit, vychvalovat};  \dicPhraseIS{lofa e‑n\,/\addthin e‑ð} \dicDirectTranslationCS{chválit (koho\,/\addthin co)} \dicExampleIS{lofa listaverkið} \dicExampleCS{pochválit umělecké dílo}  \textbf{2.} \dicFlx{dat + dat} \dicSynonym{heita} \dicDirectTranslationCS{(při)slíbit, slibovat};  \dicPhraseIS{lofa e‑m e‑u} \dicDirectTranslationCS{slíbit (komu co)}  \textbf{3.} \mbox{\dicFlx{dat + acc}} \dicSynonym{leyfa} \dicDirectTranslationCS{dovolit, dovolovat, povolit, povolovat};  \dicPhraseIS{lofa e‑m e‑ð} \dicDirectTranslationCS{dovolit (komu co)} \dicExampleIS{Ég lofaði henni að dansa.} \dicExampleCS{Dovolil jsem jí tančit.};  \dicIdiom{lofast}{ \dicPhraseIS{lofast e‑m}} \dicFlx{refl} \dicLangCat{zast.} \dicDirectTranslationCS{zaslíbit se (komu)}
\dicEntry[lofgerð] \dicTerm{lof··gerð} \dicIPA{{l}{\textopeno}{v}{\r{\textObardotlessj}}{\textepsilon}{r}{\texttheta}} \dicPos{f}[7] \dicFlx{(‑ar, ‑ir)}[1] \dicSynonym{lof} \dicDirectTranslationCS{chvála, pochvala}
\dicEntry[loforð] \dicTerm{lof··orð} \dicIPA{{l}{\textopeno}{\textlengthmark}{v}{\textopeno}{r}{\texttheta}} \dicPos{n}[2] \dicFlx{(‑s, ‑)}[5] \textbf{1.} \dicDirectTranslationCS{slib, slovo};  \dicPhraseIS{standa við loforð} \dicDirectTranslationCS{držet slovo}  \textbf{2.} \dicDirectTranslationCS{přísaha}
\dicEntry[lofsamlegur] \dicTerm{lof·sam··legur} \dicIPA{{l}{\textopeno}{f}{s}{a}{m}{l}{\textepsilon}{\textbabygamma}{\textscy}{\textsubring{r}}} \dicPos{adj}[1]\dicFlx{}[-8] \dicDirectTranslationCS{pochvalný}
\dicEntry[lofsorð] \dicTerm{lofs··orð} \dicIPA{{l}{\textopeno}{f}{s}{\textopeno}{r}{\texttheta}} \dicPos{n}[2] \dicFlx{(‑s, ‑)}[5] \dicPhraseIS{ljúka lofsorði á e‑ð} \dicDirectTranslationCS{chválit (co), vyjádřit se o~(čem) pochvalně}
\dicEntry[lofsverður] \dicTerm{lofs··verður} \dicIPA{{l}{\textopeno}{f}{s}{v}{\textepsilon}{r}{ð}{\textscy}{\textsubring{r}}} \dicPos{adj}[2]\dicFlx{}[-1] \dicDirectTranslationCS{chvályhodný, záslužný} \dicExampleIS{lofsverð viðleitni} \dicExampleCS{chvályhodné úsilí}
\dicEntry[lofsöngur] \dicTerm{lof··söng|ur} \dicIPA{{l}{\textopeno}{f}{s}{\oe i}{\ng}{\r{g}}{\textscy}{\textsubring{r}}} \dicPos{m}[8] \dicFlx{(‑s, ‑var)}[2] \dicDirectTranslationCS{chvalozpěv}
\dicEntry[loft] \dicTerm{loft} \dicsymFrequent\  \dicIPA{{l}{\textopeno}{f}{\textsubring{d}}} \dicPos{n}[2] \dicFlx{(‑s, ‑)}[5] \textbf{1.} \dicSynonym{andrúmsloft} \dicDirectTranslationCS{vzduch, povětří, ovzduší} \dicExampleIS{ferskt loft} \dicExampleCS{čerstvý vzduch}  \textbf{2.} \dicSynonym{himinn} \dicDirectTranslationCS{vzduch, obloha} \dicExampleIS{kasta e‑u upp í loftið} \dicExampleCS{vyhodit (co) do vzduchu}  \textbf{3.} \dicDirectTranslationCS{poschodí, podlaží};  \dicPhraseIS{uppi á lofti} \dicFlx{adv} \dicDirectTranslationCS{v~horním patře}  \textbf{4.} \dicSynonym{háaloft} \dicDirectTranslationCS{podkroví};  \dicIdiom{loft}{ \dicPhraseIS{það er allt upp í loft}} \dicLangCat{přen.} \dicDirectTranslationCS{všechno je vzhůru nohama};  \dicPhraseIS{e‑að er úr lausu lofti gripið} \dicLangCat{přen.} \dicDirectTranslationCS{(co) je neopodstatněné}
\dicEntry[loftárás] \dicTerm{loft··á·rás} \dicIPA{{l}{\textopeno}{f}{\textsubring{d}}{au}{r}{au}{s}} \dicPos{f}[7] \dicFlx{(‑ar, ‑ir)}[1] \dicDirectTranslationCS{nálet, letecký útok}
\dicEntry[loftbelgur] \dicTerm{loft··belg|ur} \dicIPA{{l}{\textopeno}{f}{\textsubring{d}}{\textsubring{b}}{\textepsilon}{l}{\r{g}}{\textscy}{\textsubring{r}}} \dicPos{m}[9] \dicFlx{(‑s\,/\addthin ‑jar, ‑ir)}[26] \dicFieldCat{let.} \dicDirectTranslationCS{(horkovzdušný) balón}
\dicEntry[loftbor] \dicTerm{loft··bor} \dicIPA{{l}{\textopeno}{f}{\textsubring{d}}{\textsubring{b}}{\textopeno}{\textsubring{r}}} \dicPos{m}[4] \dicFlx{(‑s, ‑ar)}[14] \dicDirectTranslationCS{sbíječka, sbíjecí kladivo}
\dicEntry[loftbréf] \dicTerm{loft··bréf} \dicIPA{{l}{\textopeno}{f}{\textsubring{d}}{\textsubring{b}}{r}{j}{\textepsilon}{f}} \dicPos{n}[2] \dicFlx{(‑s, ‑)}[5] \dicDirectTranslationCS{letecké psaní}
\dicEntry[loftbrú] \dicTerm{loft··|brú} \dicIPA{{l}{\textopeno}{f}{\textsubring{d}}{\textsubring{b}}{r}{u}} \dicPos{f}[9] \dicFlx{(‑brúar, ‑brýr)}[4] \dicDirectTranslationCS{letecký most}
\dicEntry[loftfar] \dicTerm{loft··|far} \dicIPA{{l}{\textopeno}{f}{\textsubring{d}}{f}{a}{\textsubring{r}}} \dicPos{n}[2] \dicFlx{(‑fars, ‑för)}[8] \dicDirectTranslationCS{létající prostředek (letadlo, balón, vzducholoď ap.)}
\dicEntry[loftfimleikar] \dicTerm{loft··fim·leikar} \dicIPA{{l}{\textopeno}{f}{\textsubring{d}}{f}{\textsci}{m}{l}{ei}{\r{g}}{a}{\textsubring{r}}} \dicPos{m}[6] \dicFlx{pl}[2] \dicDirectTranslationCS{letecká akrobacie}
\dicEntry[lofthræddur] \dicTerm{loft··hræddur} \dicIPA{{l}{\textopeno}{f}{\textsubring{d}}{\textsubring{r}}{a}{i}{\textsubring{d}}{\textscy}{\textsubring{r}}} \dicPos{adj}[2]\dicFlx{}[-18] \dicDirectTranslationCS{mající strach z~výšek}
\dicEntry[lofthræðsla] \dicTerm{loft··hræðsl|a} \dicIPA{{l}{\textopeno}{f}{\textsubring{d}}{\textsubring{r}}{a}{i}{ð}{s}{\textsubring{d}}{l}{a}} \dicPos{f}[1] \dicFlx{(‑u)}[5] \dicDirectTranslationCS{strach z~výšek}
\dicEntry[loftkastali] \dicTerm{loft··kastal|i} \dicIPA{{l}{\textopeno}{f}{\textsubring{d}}{k\smash{\textsuperscript{h}}}{a}{s}{\textsubring{d}}{a}{l}{\textsci}} \dicPos{m}[1] \dicFlx{(‑a, ‑ar)}[12] \dicDirectTranslationCS{vzdušný zámek};  \dicPhraseIS{byggja loftkastala} \dicLangCat{přen.} \dicDirectTranslationCS{stavět vzdušné zámky}
\dicEntry[loftkenndur] \dicTerm{loft··kenndur} \dicIPA{{l}{\textopeno}{f}{\textsubring{d}}{c\smash{\textsuperscript{h}}}{\textepsilon}{n}{\textsubring{d}}{\textscy}{\textsubring{r}}} \dicPos{adj}[2]\dicFlx{}[-14] \textbf{1.} \dicDirectTranslationCS{vzdušný}  \textbf{2.} \dicDirectTranslationCS{plynný} \dicExampleIS{loftkennd gosefni} \dicExampleCS{plynné magma}
\dicEntry[loftkældur] \dicTerm{loft··kældur} \dicIPA{{l}{\textopeno}{f}{\textsubring{d}}{c\smash{\textsuperscript{h}}}{a}{i}{l}{\textsubring{d}}{\textscy}{\textsubring{r}}} \dicPos{adj}[2]\dicFlx{}[-17] \dicDirectTranslationCS{klimatizovaný} \dicExampleIS{loftkæld skrifstofa} \dicExampleCS{klimatizovaná kancelář}
\dicEntry[loftkæling] \dicTerm{loft··kæl·ing} \dicIPA{{l}{\textopeno}{f}{\textsubring{d}}{c\smash{\textsuperscript{h}}}{a}{i}{l}{i}{\ng}{\r{g}}} \dicPos{f}[4] \dicFlx{(‑ar, ‑ar)}[5] \dicDirectTranslationCS{klimatizace}
\begin{xtolerant}{}{1pt}
\dicEntry[loftmengun] \dicTerm{loft··meng|un}\addthinS\dicIPA{{l}{\textopeno}{f}{\textsubring{d}}{m}{ei}{\ng}{\r{g}}{\textscy}{\textsubring{n}}}\addthinS\dicPos{f}[7]\addthinS\dicFlx{(‑unar)}[9] \dicDirectTranslationCS{znečištění ovzduší}
\end{xtolerant}
\dicEntry[loftnet] \dicTerm{loft··net} \dicIPA{{l}{\textopeno}{f}{\textsubring{d}}{n}{\textepsilon}{\textsubring{d}}} \dicPos{n}[2] \dicFlx{(‑s, ‑)}[5] \dicDirectTranslationCS{anténa}
\dicEntry[loftræsta] \dicTerm{loft··ræst|a} \dicIPA{{l}{\textopeno}{f}{\textsubring{d}}{r}{a}{i}{s}{\textsubring{d}}{a}} \dicPos{v}[2] \dicFlx{(‑i, ‑)}[16] \dicFlx{acc} \dicDirectTranslationCS{(vy)větrat, (vy)ventilovat} \dicExampleIS{loftræsta sal} \dicExampleCS{vyvětrat sál}
\dicEntry[loftræsting] \dicTerm{loft··ræst·ing} \dicIPA{{l}{\textopeno}{f}{\textsubring{d}}{r}{a}{i}{s}{\textsubring{d}}{i}{\ng}{\r{g}}} \dicPos{f}[4] \dicFlx{(‑ar, ‑ar)}[5] \dicDirectTranslationCS{větrání, ventilace}
\begin{xtolerant}{}{1pt}
\dicEntry[loftskeytamaður] \dicTerm{loft·skeyta··|maður}\addthinSS\dicIPA{{l}{\textopeno}{f}{\textsubring{d}}{s}{c\smash{\textsuperscript{h}}}{ei}{\textsubring{d}}{a}{m}{a}{ð}{\textscy}{\textsubring{r}}}\addthinSS\dicPos{m}[1\hspace{-0.10667em}3]\addthinS\dicFlx{(‑manns,\addthinS‑menn)}[2] \dicDirectTranslationCS{radista, radistka, rádiový operátor, rádiová operátorka}
\end{xtolerant}
\dicEntry[loftskip] \dicTerm{loft··skip} \dicIPA{{l}{\textopeno}{f}{\textsubring{d}}{s}{\r{\textObardotlessj}}{\textsci}{\textsubring{b}}} \dicPos{n}[2] \dicFlx{(‑s, ‑)}[5] \dicFieldCat{let.} \dicDirectTranslationCS{vzducholoď}
\dicEntry[loftslag] \dicTerm{lofts··lag} \dicIPA{{l}{\textopeno}{f}{s}{l}{a}{x}} \dicPos{n}[2] \dicFlx{(‑s)}[2] \dicSynonym{veðurfar} \dicDirectTranslationCS{podnebí, klima}
\dicEntry[loftslagsbelti] \dicTerm{lofts·lags··belti} \dicIPA{{l}{\textopeno}{f}{s}{l}{a}{x}{s}{\textsubring{b}}{\textepsilon}{\textsubring{l}}{\textsubring{d}}{\textsci}} \dicPos{n}[2] \dicFlx{(‑s, ‑)}[14] \dicFieldCat{meteo.} \dicDirectTranslationCS{podnebný pás}
\dicEntry[loftslagsfræði] \dicTerm{lofts·lags··fræð|i} \dicIPA{{l}{\textopeno}{f}{s}{l}{a}{x}{s}{f}{r}{a}{i}{ð}{\textsci}} \dicPos{f}[3] \dicFlx{(‑i)}[3] \dicSynonym{veðurfarsfræði} \dicDirectTranslationCS{klimatologie}
\dicEntry[loftsteinn] \dicTerm{loft··stein|n} \dicIPA{{l}{\textopeno}{f}{\textsubring{d}}{s}{\textsubring{d}}{ei}{\textsubring{d}}{\textsubring{n}}} \dicPos{m}[6] \dicFlx{(‑s, ‑ar)}[42] \dicFieldCat{astro.} \dicDirectTranslationCS{meteorit}
\dicEntry[lofttegund] \dicTerm{loft··tegund} \dicIPA{{l}{\textopeno}{f}{t\smash{\textsuperscript{h}}}{\textepsilon}{\textbabygamma}{\textscy}{n}{\textsubring{d}}} \dicPos{f}[7] \dicFlx{(‑ar, ‑ir)}[1] \dicFieldCat{chem.} \dicDirectTranslationCS{plyn}
\dicEntry[lofttóm] \textls[15]{\dicTerm{loft··tóm} \dicIPA{{l}{\textopeno}{f}{t\smash{\textsuperscript{h}}}{ou}{\textsubring{m}}} \dicPos{n}[2] \dicFlx{(‑s)}[2] \dicFieldCat{fyz.} \dicDirectTranslationCS{vakuum, vzduchoprázdno}}
\dicEntry[loftvarnir] \dicTerm{loft··varnir} \dicIPA{{l}{\textopeno}{f}{\textsubring{d}}{v}{a}{r}{\textsubring{d}}{n}{\textsci}{\textsubring{r}}} \dicPos{f}[7] \dicFlx{pl}[18] \dicFieldCat{voj.} \dicDirectTranslationCS{protiletecká\,/\addthin protiletadlová obrana}
\dicEntry[loftveiki] \dicTerm{loft··veik|i} \dicIPA{{l}{\textopeno}{f}{\textsubring{d}}{v}{ei}{\r{\textObardotlessj}}{\textsci}} \dicPos{f}[3] \dicFlx{(‑i)}[3] \dicFieldCat{med.} \dicDirectTranslationCS{výšková\,/\addthin horská nemoc}
\dicEntry[loftvog] \dicTerm{loft··vog} \dicIPA{{l}{\textopeno}{f}{\textsubring{d}}{v}{\textopeno}{x}} \dicPos{f}[7] \dicFlx{(‑ar, ‑ir)}[1] \dicDirectTranslationCS{barometr}
\dicEntry[loftþéttur] \dicTerm{loft··þéttur} \dicIPA{{l}{\textopeno}{f}{\textsubring{d}}{\texttheta}{j}{\textepsilon}{h}{\textsubring{d}}{\textscy}{\textsubring{r}}} \dicPos{adj}[1]\dicFlx{}[-10] \dicDirectTranslationCS{vzduchotěsný, hermetický}
\dicEntry[loftþrýstingur] \dicTerm{loft··þrýst·ing|ur} \dicIPA{{l}{\textopeno}{f}{\textsubring{d}}{\texttheta}{r}{i}{s}{\textsubring{d}}{i}{\ng}{\r{g}}{\textscy}{\textsubring{r}}} \dicPos{m}[6] \dicFlx{(‑s)}[9] \dicDirectTranslationCS{tlak vzduchu, atmosférický tlak}
\dicEntry[loga] \dicTerm{log|a} \dicsymFrequent\  \dicIPA{{l}{\textopeno}{\textlengthmark}{\textbabygamma}{a}} \dicPos{v}[1] \dicFlx{(‑aði)}[44] \textbf{1.} \dicDirectTranslationCS{hořet, planout};  \dicPhraseIS{það logar á e‑u} \dicFlx{impers} \dicDirectTranslationCS{(co) hoří} \dicExampleIS{Það logar vel á kertinu.} \dicExampleCS{Svíčka dobře hoří.}  \textbf{2.} \dicDirectTranslationCS{hořet, zuřit (město ap.)} \dicIndirectTranslationCS{(o~nepokojích nebo válce)};  \dicPhraseIS{e‑að logar í e‑u} \dicDirectTranslationCS{v~(čem) zuří (co) (ve městě nepokoje ap.)}
\dicEntry[logandi] \dicTerm{log··andi\smash{\textsuperscript{1}}} \dicIPA{{l}{\textopeno}{\textlengthmark}{\textbabygamma}{a}{n}{\textsubring{d}}{\textsci}} \dicPos{adj}[13] \dicFlx{indecl}[1] \dicDirectTranslationCS{hořící, planoucí} \dicExampleIS{logandi kerti} \dicExampleCS{hořící svíce}
\dicEntry[logandi] \dicTerm{log··andi\smash{\textsuperscript{2}}} \dicIPA{{l}{\textopeno}{\textlengthmark}{\textbabygamma}{a}{n}{\textsubring{d}}{\textsci}} \dicPos{adv} \dicDirectTranslationCS{hrozně, strašně} \dicIndirectTranslationCS{(k~zdůraznění)} \dicExampleIS{Hann var logandi hræddur.} \dicExampleCS{Strašně se bál.};  \dicPhraseIS{hvert þó í logandi!} \dicDirectTranslationCS{panebože!, svatá dobroto!}
\dicEntry[logga] \dicTerm{logg|a} \dicIPA{{l}{\textopeno}{\r{g}}{\textlengthmark}{a}} \dicPos{v}[1] \dicFlx{(‑aði)}[1] \dicFieldCat{poč.} \dicDirectTranslationCS{logovat, zaznamenat do souboru};  \dicIdiom{logga}[inn]{ \dicPhraseIS{logga sig inn}} \dicFieldCat{poč.} \dicDirectTranslationCS{(za)logovat se, nalogovat se, přihlásit se};  \dicIdiom{logga}[út]{ \dicPhraseIS{logga sig út}} \dicFieldCat{poč.} \dicDirectTranslationCS{vylogovat se, odlogovat se, odhlásit se}
\dicEntry[logi] \dicTerm{log|i} \dicsymFrequent\  \dicIPA{{l}{\textopeno i}{j}{\textlengthmark}{\textsci}} \dicPos{m}[1] \dicFlx{(‑a, ‑ar)}[1] \dicDirectTranslationCS{plamen, planutí} \dicExampleIS{Loginn hefur dofnað.} \dicExampleCS{Plamen zeslábl.};  \dicPhraseIS{fara eins og logi yfir akur} \dicLangCat{přen.} \dicDirectTranslationCS{přehnat se jako bouře}
\dicEntry[logið] \dicTerm{logið} \dicIPA{{l}{\textopeno i}{j}{\textlengthmark}{\textsci}{\texttheta}} \dicPos{v} \dicFlx{supin} \dicLink{ljúga}
\dicEntry[logn] \dicTerm{logn} \dicsymFrequent\  \dicIPA{{l}{\textopeno}{\r{g}}{\textsubring{n}}} \dicPos{n}[2] \dicFlx{(‑s, ‑)}[5] \dicSynonym*{vindleysa} \dicDirectTranslationCS{bezvětří} \dicIndirectTranslationCS{(Beaufort 0) \mbox{(0--0,2~m/s)}} \dicExampleIS{Það er logn í dag.} \dicExampleCS{Dnes je bezvětří.}
\dicEntry[lognmolla] \dicTerm{logn··moll|a} \dicIPA{{l}{\textopeno}{\r{g}}{n}{m}{\textopeno}{\textsubring{d}}{l}{a}} \dicPos{f}[1] \dicFlx{(‑u)}[5] \textbf{1.} \dicSynonym{andleysi} \dicDirectTranslationCS{nečinnost, neaktivita}  \textbf{2.} \dicDirectTranslationCS{poklidné počasí}
\dicEntry[logsauð] \dicTerm{log··sauð} \dicIPA{{l}{\textopeno}{x}{s}{\oe i}{\texttheta}} \dicPos{v} \dicFlx{ind pf sg 1 pers} \dicLink{logsjóða}
\dicEntry[logsjóða] \dicTerm{log··|sjóða} \dicIPA{{l}{\textopeno}{x}{s}{j}{ou}{ð}{a}} \dicPos{v}[6] \dicFlx{(‑sýð, ‑sauð, ‑suðum, ‑syði, ‑soðið)}[104] \dicFlx{acc} \dicFieldCat{techn.} \dicDirectTranslationCS{svářet, svařovat}
\dicEntry[logsoðið] \dicTerm{log··soðið} \dicIPA{{l}{\textopeno}{x}{s}{\textopeno}{ð}{\textsci}{\texttheta}} \dicPos{v} \dicFlx{supin} \dicLink{logsjóða}
\dicEntry[logsuða] \dicTerm{log··suð|a} \dicIPA{{l}{\textopeno}{x}{s}{\textscy}{ð}{a}} \dicPos{f}[1] \dicFlx{(‑u)}[5] \dicSynonym*{það að logsjóða} \dicDirectTranslationCS{sváření, svařování}
\dicEntry[logsuðum] \dicTerm{log··suðum} \dicIPA{{l}{\textopeno}{x}{s}{\textscy}{ð}{\textscy}{\textsubring{m}}} \dicPos{v} \dicFlx{ind pf pl 1 pers} \dicLink{logsjóða}
\dicEntry[logsyði] \dicTerm{log··syði} \dicIPA{{l}{\textopeno}{x}{s}{\textsci}{ð}{\textsci}} \dicPos{v} \dicFlx{con pf sg 1 pers} \dicLink{logsjóða}
\dicEntry[logsýð] \dicTerm{log··sýð} \dicIPA{{l}{\textopeno}{x}{s}{i}{\texttheta}} \dicPos{v} \dicFlx{ind praes sg 1 pers} \dicLink{logsjóða}
\dicEntry[lok] \dicTerm{lok} \dicsymFrequent\  \dicIPA{{l}{\textopeno}{\textlengthmark}{\r{g}}} \dicPos{n}[2] \dicFlx{(‑s, ‑)}[5] \textbf{1.} \dicDirectTranslationCS{uzávěr, víčko, víko, poklička, poklice} \dicExampleIS{lokið á dósinni} \dicExampleCS{víčko plechovky}  \textbf{2.} \dicPhraseIS{lok} \dicFlx{pl} \dicSynonym{endir} \dicDirectTranslationCS{konec, závěr};  \dicPhraseIS{að lokum} \dicFlx{adv} \dicDirectTranslationCS{nakonec, na závěr};  \dicPhraseIS{e‑að líður undir lok} \dicDirectTranslationCS{(co) se chýlí k~závěru};  \dicPhraseIS{í lokin} \dicFlx{adv} \dicDirectTranslationCS{nakonec}
\dicEntry[loka] \dicTerm{lok|a\smash{\textsuperscript{1}}} \dicIPA{{l}{\textopeno}{\textlengthmark}{\r{g}}{a}} \dicPos{f}[1] \dicFlx{(‑u, ‑ur)}[19] \textbf{1.} \dicSynonym{hespa\smash{\textsuperscript{1}}} \dicDirectTranslationCS{západka, petlice, závora}  \textbf{2.} \dicFieldCat{anat.} \dicDirectTranslationCS{chlopeň}
\dicEntry[loka] \dicTerm{lok|a\smash{\textsuperscript{2}}} \dicsymFrequent\  \dicIPA{{l}{\textopeno}{\textlengthmark}{\r{g}}{a}} \dicPos{v}[1] \dicFlx{(‑aði)}[1] \dicFlx{dat} \dicDirectTranslationCS{(u)zavřít, zavírat} \dicExampleIS{loka glugganum} \dicExampleCS{zavřít okno};  \dicIdiom{loka}[að]{ \dicPhraseIS{loka að sér}} \dicDirectTranslationCS{zavřít za sebou (dveře)};  \dicIdiom{loka}[fyrir]{ \dicPhraseIS{loka fyrir e‑ð}} \dicDirectTranslationCS{zavřít (co), vypnout (co) (elektřinu ap.)};  \dicIdiom{loka}[inni]{ \dicPhraseIS{loka e‑n inni}} \dicDirectTranslationCS{zavřít (koho) (do blázince ap.)};  \dicIdiom{loka}[úti]{ \dicPhraseIS{loka e‑n úti}} \dicDirectTranslationCS{vyloučit (koho) (z~party ap.), nevpustit (koho) dovnitř};  \dicIdiom{lokast}{ \dicPhraseIS{lokast}} \dicFlx{refl} \dicDirectTranslationCS{zavřít se, zavírat se}
\dicEntry[lokaáfangi] \dicTerm{loka··á·fang|i} \dicIPA{{l}{\textopeno}{\textlengthmark}{\r{g}}{a}{au}{f}{au}{\textltailn}{\r{\textObardotlessj}}{\textsci}} \dicPos{m}[1] \dicFlx{(‑a, ‑ar)}[8] \dicDirectTranslationCS{poslední etapa}
\dicEntry[lokaður] \dicTerm{lok··|aður} \dicsymFrequent\  \dicIPA{{l}{\textopeno}{\textlengthmark}{\r{g}}{a}{ð}{\textscy}{\textsubring{r}}} \dicPos{adj}[3] \dicFlx{(f ‑uð)}[3] \textbf{1.} \dicSynonym*{ekki opinn} \dicDirectTranslationCS{zavřený, uzavřený} \dicExampleIS{lokaður vegur} \dicExampleCS{uzavřená cesta}  \textbf{2.} \dicLangCat{přen.} \dicDirectTranslationCS{uzavřený, odměřený, zdrženlivý, rezervovaný}
\dicEntry[lokafrestur] \dicTerm{loka··frest|ur} \dicIPA{{l}{\textopeno}{\textlengthmark}{\r{g}}{a}{f}{r}{\textepsilon}{s}{\textsubring{d}}{\textscy}{\textsubring{r}}} \dicPos{m}[6] \dicFlx{(‑s)}[7] \dicDirectTranslationCS{uzávěrka, (konečný) termín, deadline}
\dicEntry[lokaorð] \dicTerm{loka··orð} \dicIPA{{l}{\textopeno}{\textlengthmark}{\r{g}}{a}{\textopeno}{r}{\texttheta}} \dicPos{n}[2] \dicFlx{(‑s, ‑)}[5] \dicDirectTranslationCS{konečné\,/\addthin závěrečné slovo, doslov};  \dicPhraseIS{eiga lokaorðið} \dicLangCat{přen.} \dicDirectTranslationCS{mít poslední slovo}
\dicEntry[lokapróf] \dicTerm{loka··próf} \dicIPA{{l}{\textopeno}{\textlengthmark}{\r{g}}{a}{p\smash{\textsuperscript{h}}}{r}{ou}{f}} \dicPos{n}[2] \dicFlx{(‑s, ‑)}[5] \dicFieldCat{škol.} \dicDirectTranslationCS{závěrečná zkouška}
\dicEntry[lokaritgerð] \dicTerm{loka··rit·gerð} \dicIPA{{l}{\textopeno}{\textlengthmark}{\r{g}}{a}{r}{\textsci}{\textsubring{d}}{\r{\textObardotlessj}}{\textepsilon}{r}{\texttheta}} \dicPos{f}[7] \dicFlx{(‑ar, ‑ir)}[1] \dicFieldCat{škol.} \dicDirectTranslationCS{závěrečná práce}
\dicEntry[lokasjóður] \dicTerm{loka··sjóð|ur} \dicIPA{{l}{\textopeno}{\textlengthmark}{\r{g}}{a}{s}{j}{ou}{ð}{\textscy}{\textsubring{r}}} \dicPos{m}[9] \dicFlx{(‑s, ‑ir)}[6] \dicFieldCat{bot.} \dicDirectTranslationCS{kokrhel menší} \textit{(l.~{\textLA{Rhinanthus minor}})}  \dicsymPhoto\ 
\dicFigure{71849.jpg}{Lokasjóður}{Lokasjóður - ŠARŽÍK František, Biolib, Copyright/PD}
\dicEntry[lokaæfing] \dicTerm{loka··æf·ing} \dicIPA{{l}{\textopeno}{\textlengthmark}{\r{g}}{a}{a}{i}{v}{i}{\ng}{\r{g}}} \dicPos{f}[4] \dicFlx{(‑ar, ‑ar)}[5] \dicDirectTranslationCS{generální zkouška, generálka}
\dicEntry[lokhljóð] \dicTerm{lok··hljóð} \dicIPA{{l}{\textopeno}{\textlengthmark}{\r{g}}{\textsubring{l}}{j}{ou}{\texttheta}} \dicPos{n}[2] \dicFlx{(‑s, ‑)}[5] \dicFieldCat{jaz.} \dicDirectTranslationCS{ražená souhláska, ploziva}
\dicEntry[loki] \dicTerm{lok|i} \dicIPA{{l}{\textopeno}{\textlengthmark}{\r{\textObardotlessj}}{\textsci}} \dicPos{m}[1] \dicFlx{(‑a, ‑ar)}[1] \textbf{1.} \dicSynonym{hnútur} \dicDirectTranslationCS{uzel, zauzlení}  \textbf{2.} \dicFieldCat{techn.} \dicSynonym*{öryggisloki} \dicDirectTranslationCS{ventil, ventilek}
\dicEntry[lokið] \dicTerm{lokið} \dicIPA{{l}{\textopeno}{\textlengthmark}{\r{\textObardotlessj}}{\textsci}{\texttheta}} \dicPos{v} \dicFlx{supin} \dicLink{ljúka}
\dicEntry[lokinn] \dicTerm{lokinn} \dicsymFrequent\  \dicIPA{{l}{\textopeno}{\textlengthmark}{\r{\textObardotlessj}}{\textsci}{\textsubring{n}}} \dicPos{adj}[6]\dicFlx{}[-6] \dicSynonym{búinn} \dicDirectTranslationCS{skončený, uzavřený} \dicExampleIS{Rannsókn er lokið.} \dicExampleCS{Vyšetřování je uzavřené.};  \dicPhraseIS{e‑m er öllum lokið} \dicFlx{impers} \dicLangCat{přen.} \dicDirectTranslationCS{(kdo) toho má dost, (kdo) je hotový}
\dicEntry[lokka] \dicTerm{lokk|a} \dicIPA{{l}{\textopeno}{h}{\r{g}}{a}} \dicPos{v}[1] \dicFlx{(‑aði)}[1] \dicFlx{acc} \textbf{1.} \dicSynonym{ginna} \dicDirectTranslationCS{(na)lákat, zlákat} \dicExampleIS{lokka e‑n til sín} \dicExampleCS{nalákat (koho) k~sobě}  \textbf{2.} \dicDirectTranslationCS{vylákat, vymámit, vyloudit} \dicExampleIS{lokka e‑ð út úr e‑m} \dicExampleCS{vylákat (co) od (koho)}
\dicEntry[lokkur] \dicTerm{lokk|ur} \dicIPA{{l}{\textopeno}{h}{\r{g}}{\textscy}{\textsubring{r}}} \dicPos{m}[6] \dicFlx{(‑s, ‑ar)}[22] \dicDirectTranslationCS{lokna, kudrna}
\dicEntry[lokræsi] \dicTerm{lok··ræsi} \dicIPA{{l}{\textopeno}{\textlengthmark}{\r{g}}{r}{a}{i}{s}{\textsci}} \dicPos{n}[2] \dicFlx{(‑s, ‑)}[14] \dicLangCat{zast.} \dicDirectTranslationCS{stoka, odpad}
\dicEntry[loks] \dicTerm{loks} \dicsymFrequent\  \dicIPA{{l}{\textopeno}{\r{g}}{s}} \dicPos{adv} \dicSynonym*{að lokum} \dicDirectTranslationCS{nakonec, konečně} \dicExampleIS{Hurðin var loks opnuð.} \dicExampleCS{Dveře se nakonec otevřely.}
\dicEntry[loksins] \dicTerm{loksins} \dicsymFrequent\  \dicIPA{{l}{\textopeno}{\r{g}}{s}{\textsci}{n}{s}} \dicPos{adv} \dicSynonym*{á endanum} \dicDirectTranslationCS{konečně} \dicExampleIS{Fundur er loksins að hefjast.} \dicExampleCS{Schůze konečně začíná.}
\dicEntry[lokun] \dicTerm{lok|un} \dicIPA{{l}{\textopeno}{\textlengthmark}{\r{g}}{\textscy}{\textsubring{n}}} \dicPos{f}[7] \dicFlx{(‑unar, ‑anir)}[8] \textbf{1.} \dicSynonym*{það að loka} \dicDirectTranslationCS{zavírání, (u)zavření} \dicExampleIS{lokun búða} \dicExampleCS{zavření obchodu}  \textbf{2.} \dicSynonym*{það að setja hindrun} \dicDirectTranslationCS{uzavírka (cesty ap.)} \dicExampleIS{lokun vegar} \dicExampleCS{uzavírka cesty}
\dicEntry[lokunartími] \textls[15]{\dicTerm{lokunar··tím|i} \dicIPA{{l}{\textopeno}{\r{g}}{\textscy}{n}{a}{\textsubring{r}}{t\smash{\textsuperscript{h}}}{i}{m}{\textsci}} \dicPos{m}[1] \dicFlx{(‑a)}[3] \dicDirectTranslationCS{zavírací doba}} \dicAntonym{opnunartími}
\dicEntry[lopapeysa] \dicTerm{lopa··peys|a} \dicIPA{{l}{\textopeno}{\textlengthmark}{\textsubring{b}}{a}{p\smash{\textsuperscript{h}}}{ei}{s}{a}} \dicPos{f}[1] \dicFlx{(‑u, ‑ur)}[7] \dicIndirectTranslationCS{pletený vlněný svetr s~typickými islandskými vzory}
\dicEntry[lopi] \dicTerm{lop|i} \dicIPA{{l}{\textopeno}{\textlengthmark}{\textsubring{b}}{\textsci}} \dicPos{m}[1] \dicFlx{(‑a, ‑ar)}[1] \textbf{1.} \dicSynonym*{ullarþráður} \dicDirectTranslationCS{mykaná vlna}  \textbf{2.} \dicFieldCat{med.} \dicSynonym{bjúgur} \dicDirectTranslationCS{edém, otok};  \dicIdiom{lopi}{ \dicPhraseIS{spinna\,/\addthin teygja lopann}} \dicLangCat{přen.} \dicDirectTranslationCS{pustit si pusu na špacír}
\dicEntry[loppa] \dicTerm{lopp|a} \dicIPA{{l}{\textopeno}{h}{\textsubring{b}}{a}} \dicPos{f}[1] \dicFlx{(‑u, ‑ur)}[19] \dicDirectTranslationCS{tlapa, tlapka}
\dicEntry[los] \dicTerm{los} \dicIPA{{l}{\textopeno}{\textlengthmark}{s}} \dicPos{n}[2] \dicFlx{(‑s, ‑)}[5] \dicDirectTranslationCS{nestabilita, nestabilnost, nesoudržnost} \dicExampleIS{það er los á e‑u} \dicExampleCS{(co) je nestabilní}
\dicEntry[losa] \dicTerm{los|a} \dicsymFrequent\  \dicIPA{{l}{\textopeno}{\textlengthmark}{s}{a}} \dicPos{v}[1] \dicFlx{(‑aði)}[1] \dicFlx{acc} \textbf{1.} \dicDirectTranslationCS{uvolnit, uvolňovat, povolit, povolovat (šroubek ap.)} \dicExampleIS{Hann losar hjólið af bílnum.} \dicExampleCS{Uvolňuje kolo u~auta.}  \textbf{2.} \dicDirectTranslationCS{uvolnit, uvolňovat (pokoj ap.)} \dicExampleIS{losa hótelherbergi} \dicExampleCS{uvolnit hotelový pokoj}  \textbf{3.} \dicDirectTranslationCS{vyprázdnit, vyprazdňovat} \dicExampleIS{losa flöskuna} \dicExampleCS{vyprázdnit láhev}  \textbf{4.} \dicPhraseIS{e‑að losar (20 tonn)} \dicDirectTranslationCS{(co) má něco málo přes (20 tun), (co) přesahuje (20 tun)};  \dicIdiom{losa}[um]{ \dicPhraseIS{losa um e‑ð}} \dicDirectTranslationCS{uvolnit (co), zmírnit (co) (o\nobreak mezení ap.)};  \dicIdiom{losa}[úr]{ \dicPhraseIS{losa e‑n úr e‑u}} \dicDirectTranslationCS{uvolnit (koho) z~(čeho), vymanit (koho) z~(čeho)};  \dicIdiom{losa}[við]{ \dicPhraseIS{losa e‑n við e‑ð}} \dicDirectTranslationCS{pomoct (komu) s~(čím)}; { \dicPhraseIS{losa sig við e‑ð}} \dicDirectTranslationCS{zbavit se (čeho)}
\dicEntry[losaralegur] \dicTerm{losara··legur} \dicIPA{{l}{\textopeno}{\textlengthmark}{s}{a}{r}{a}{l}{\textepsilon}{\textbabygamma}{\textscy}{\textsubring{r}}} \dicPos{adj}[1]\dicFlx{}[-8] \dicSynonym{sundurlaus} \dicDirectTranslationCS{náhod\-ný, nahodilý, neplánovaný}
\dicEntry[losna] \dicTerm{losn|a} \dicsymFrequent\  \dicIPA{{l}{\textopeno}{s}{\textsubring{d}}{n}{a}} \dicPos{v}[1] \dicFlx{(‑aði)}[1] \textbf{1.} \dicDirectTranslationCS{uvolnit se, uprázdnit se (místo ap.)} \dicExampleIS{Starfið losnar bráðlega.} \dicExampleCS{Zaměstnání se brzy uvolní.}  \textbf{2.} \dicDirectTranslationCS{uvolnit\,/\addthin uvolňovat se, povolit, povolovat (šroubek ap.)} \dicExampleIS{Það hefur losnað hjól af bílnum.} \dicExampleCS{U~auta se uvolnilo kolo.}  \textbf{3.} \dicDirectTranslationCS{uvolnit\,/\addthin uvolňovat se (z~práce ap.)} \dicExampleIS{losna úr vinnunni} \dicExampleCS{uvolnit se z~práce};  \dicIdiom{losna}[undan]{ \dicPhraseIS{losna undan e‑u}} \dicDirectTranslationCS{vyhnout\,/\addthin vyhýbat se (čemu), uniknout\,/\addthin unikat (čemu)} \dicExampleIS{losna undan mörgum skyldum} \dicExampleCS{vyhýbat se mnoha povinnostem};  \dicIdiom{losna}[við]{ \dicPhraseIS{losna við e‑ð}} \dicSynonym*{koma e‑u af sér} \dicDirectTranslationCS{zbavit se (čeho)}
\dicEntry[lost] \dicTerm{lost} \dicIPA{{l}{\textopeno}{s}{\textsubring{d}}} \dicPos{n}[2] \dicFlx{(‑s, ‑)}[5] \textbf{1.} \dicFieldCat{med.} \dicDirectTranslationCS{šok}  \textbf{2.} \dicLangCat{hovor.} \dicDirectTranslationCS{šok, překvapení} \dicExampleIS{Ég var alveg í losti á eftir.} \dicExampleCS{Byl jsem po tom v~šoku.}
\dicEntry[lostafullur] \dicTerm{losta··|fullur} \dicIPA{{l}{\textopeno}{s}{\textsubring{d}}{a}{f}{\textscy}{\textsubring{d}}{l}{\textscy}{\textsubring{r}}} \dicPos{adj}[10] \dicFlx{(comp ‑fyllri, sup ‑fyllstur)}[7] \dicSynonym*{munaðargjarn} \dicDirectTranslationCS{chlípný, žádostivý, chtivý, vilný}
\dicEntry[losti] \dicTerm{lost|i} \dicIPA{{l}{\textopeno}{s}{\textsubring{d}}{\textsci}} \dicPos{m}[1] \dicFlx{(‑a)}[3] \dicSynonym{girnd} \dicDirectTranslationCS{chlípnost, žádostivost, chtíč, vilnost}
\dicEntry[lostið] \dicTerm{lostið} \dicIPA{{l}{\textopeno}{s}{\textsubring{d}}{\textsci}{\texttheta}} \dicPos{v} \dicFlx{supin} \dicLink{ljósta}
\dicEntry[lostinn] \dicTerm{lostinn} \dicIPA{{l}{\textopeno}{s}{\textsubring{d}}{\textsci}{\textsubring{n}}} \dicPos{adj}[6]\dicFlx{}[-6] \dicPhraseIS{skelfingu lostinn} \dicFlx{adj} \dicDirectTranslationCS{vyděšený, vystrašený};  \dicPhraseIS{vera furðu lostinn yfir e‑u} \dicDirectTranslationCS{oněmět úžasem nad (čím)};  \dicPhraseIS{vera þrumu lostinn yfir e‑u} \dicDirectTranslationCS{být šokován (čím)}
\dicEntry[lostæti] \dicTerm{lost··æti} \dicIPA{{l}{\textopeno}{s}{\textsubring{d}}{a}{i}{\textsubring{d}}{\textsci}} \dicPos{n}[2] \dicFlx{(‑s, ‑)}[14] \dicSynonym{góðgæti} \dicDirectTranslationCS{pochoutka, lahůdka, delikatesa}
\dicEntry[losun] \dicTerm{los|un} \dicIPA{{l}{\textopeno}{\textlengthmark}{s}{\textscy}{\textsubring{n}}} \dicPos{f}[7] \dicFlx{(‑unar)}[9] \textbf{1.} \dicSynonym*{losun sorps o.fl.} \dicDirectTranslationCS{vypouštění, vyprazdňování (odpadu, odpadků ap.)}  \textbf{2.} \dicSynonym*{afferming} \dicDirectTranslationCS{vykládka, vyložení} \dicExampleIS{losun á farminum} \dicExampleCS{vyložení nákladu}
\dicEntry[lota] \dicTerm{lot|a} \dicIPA{{l}{\textopeno}{\textlengthmark}{\textsubring{d}}{a}} \dicPos{f}[1] \dicFlx{(‑u, ‑ur)}[19] \dicDirectTranslationCS{období, cyklus, perioda, (časový) úsek} \dicExampleIS{í einni lotu} \dicExampleCS{na jeden zátah}
\dicEntry[lotið] \dicTerm{lotið} \dicIPA{{l}{\textopeno}{\textlengthmark}{\textsubring{d}}{\textsci}{\texttheta}} \dicPos{v} \dicFlx{supin} \dicLink{lúta}
\dicEntry[lotinn] \dicTerm{lotinn} \dicIPA{{l}{\textopeno}{\textlengthmark}{\textsubring{d}}{\textsci}{\textsubring{n}}} \dicPos{adj}[6]\dicFlx{}[-2] \dicSynonym{boginn} \dicDirectTranslationCS{shrbený} \dicExampleIS{lotinn í herðum} \dicExampleCS{shrbený v~zádech}
\dicEntry[lotning] \dicTerm{lot··ning} \dicIPA{{l}{\textopeno}{h}{\textsubring{d}}{n}{i}{\ng}{\r{g}}} \dicPos{f}[4] \dicFlx{(‑ar)}[7] \dicSynonym{virðing} \dicDirectTranslationCS{úcta, uctívání} \dicExampleIS{sýna e‑m lotningu} \dicExampleCS{projevovat ke (komu) úctu}
\begin{xtolerant}{}{1pt}
\dicEntry[lotningarhreimur] \dicTerm{lotningar··hreim|ur}\addthinS\dicIPA{{l}{\textopeno}{h}{\textsubring{d}}{n}{i}{\ng}{\r{g}}{a}{\textsubring{r}}{ei}{m}{\textscy}{\textsubring{r}}}\addthinS\dicPos{m}[6]\addthinS\dicFlx{(‑s, ‑ar)}[22] \dicDirectTranslationCS{pochlebovačný tón}
\end{xtolerant}
\dicEntry[lottó] \dicTerm{lottó} \dicIPA{{l}{\textopeno}{h}{\textsubring{d}}{ou}} \dicPos{n}[2] \dicFlx{(‑s, ‑)}[34] \dicDirectTranslationCS{loterie, sportka}
\dicEntry[lotugræðgi] \dicTerm{lotu··græðg|i} \dicIPA{{l}{\textopeno}{\textlengthmark}{\textsubring{d}}{\textscy}{\r{g}}{r}{a}{i}{ð}{\r{g}}{\textsci}} \dicPos{f}[3] \dicFlx{(‑i)}[3] \dicFieldCat{med.} \dicDirectTranslationCS{bulimie}
\dicEntry[lotukerfi] \dicTerm{lotu··kerfi} \dicIPA{{l}{\textopeno}{\textlengthmark}{\textsubring{d}}{\textscy}{c\smash{\textsuperscript{h}}}{\textepsilon}{r}{v}{\textsci}} \dicPos{n}[2] \dicFlx{(‑s)}[20] \dicFieldCat{chem.} \dicDirectTranslationCS{periodická tabulka\,/\addthin soustava prvků}
\dicEntry[lotutúlkun] \dicTerm{lotu··túlk|un} \dicIPA{{l}{\textopeno}{\textlengthmark}{\textsubring{d}}{\textscy}{t\smash{\textsuperscript{h}}}{u}{\textsubring{l}}{\r{g}}{\textscy}{\textsubring{n}}} \dicPos{f}[7] \dicFlx{(‑unar, ‑anir)}[8] \dicDirectTranslationCS{konsekutivní tlumočení}
\dicEntry[lóa] \dicTerm{ló|a} \dicIPA{{l}{ou}{\textlengthmark}{a}} \dicPos{f}[1] \dicFlx{(‑u, ‑ur)}[7] \dicFieldCat{zool.} \dicSynonym{heiðlóa} \dicDirectTranslationCS{kulík, kulík zlatý} \textit{(l.~{\textLA{Pluvialis apricaria}})}  \dicsymPhoto\ 
\dicFigure{21284.jpg}{Lóa}{Lóa - Karney, Lee, Biolib, PD}
\dicEntry[lóð] \dicTerm{lóð\smash{\textsuperscript{1}}} \dicIPA{{l}{ou}{\textlengthmark}{\texttheta}} \dicPos{f}[7] \dicFlx{(‑ar, ‑ir)}[1] \dicSynonym*{landspilda} \dicDirectTranslationCS{parcela, pozemek} \dicExampleIS{Húsið stendur á lóðinni.} \dicExampleCS{Dům stojí na pozemku.}
\dicEntry[lóð] \dicTerm{lóð\smash{\textsuperscript{2}}} \dicIPA{{l}{ou}{\textlengthmark}{\texttheta}} \dicPos{f}[7] \dicFlx{(‑ar, ‑ir)}[1] \dicFieldCat{nám.} \dicSynonym{fiskilína} \dicDirectTranslationCS{síť s~mnoha háčky};  \dicPhraseIS{leggja lóðirnar} {\textbf{a.}} \dicDirectTranslationCS{položit síť} \dicIndirectTranslationCS{(síť s~mnoha háčky)};  {\textbf{b.}} \dicSynonym*{selja upp af sjóveiki} \dicDirectTranslationCS{zvracet} \dicIndirectTranslationCS{(kvůli mořské nemoci)}
\dicEntry[lóð] \dicTerm{lóð\smash{\textsuperscript{3}}} \dicIPA{{l}{ou}{\textlengthmark}{\texttheta}} \dicPos{n}[2] \dicFlx{(‑s, ‑)}[5] \dicDirectTranslationCS{závaží} \dicExampleIS{lóð á vog} \dicExampleCS{závaží na váze};  \dicPhraseIS{það er lóðið} \dicLangCat{přen.} \dicDirectTranslationCS{to je jádro věci}
\dicEntry[lóða] \dicTerm{lóð|a} \dicIPA{{l}{ou}{\textlengthmark}{ð}{a}} \dicPos{v}[1] \dicFlx{(‑aði)}[1] \dicFlx{acc} \textbf{1.} \dicFieldCat{techn.} \dicDirectTranslationCS{letovat, pájet} \dicExampleIS{lóða e‑ð saman} \dicExampleCS{sletovat (co)}  \textbf{2.} \dicSynonym*{mæla dýpt} \dicDirectTranslationCS{sondovat (dno nebo moře pro výskyt ryb)};  \dicPhraseIS{lóða á fisk(i)} \dicFieldCat{nám.} \dicDirectTranslationCS{hledat ryby sonarem}
\dicEntry[lóðbolti] \dicTerm{lóð··bolt|i} \dicIPA{{l}{ou}{ð}{\textsubring{b}}{\textopeno}{\textsubring{l}}{\textsubring{d}}{\textsci}} \dicPos{m}[1] \dicFlx{(‑a, ‑ar)}[1] \dicFieldCat{techn.} \dicDirectTranslationCS{páječka, pájedlo}
\dicEntry[lóðréttur] \dicTerm{lóð··réttur} \dicIPA{{l}{ou}{ð}{r}{j}{\textepsilon}{h}{\textsubring{d}}{\textscy}{\textsubring{r}}} \dicPos{adj}[1]\dicFlx{}[-10] \dicDirectTranslationCS{svislý, vertikální, kolmý}
\dicEntry[lóðun] \dicTerm{lóð|un} \dicIPA{{l}{ou}{\textlengthmark}{ð}{\textscy}{\textsubring{n}}} \dicPos{f}[7] \dicFlx{(‑unar)}[9] \dicFieldCat{techn.} \dicSynonym*{það að lóða} \dicDirectTranslationCS{letování, pájení}
\dicEntry[lófaklapp] \dicTerm{lófa··klapp} \dicIPA{{l}{ou}{\textlengthmark}{v}{a}{k\smash{\textsuperscript{h}}}{l}{a}{h}{\textsubring{b}}} \dicPos{n}[2] \dicFlx{(‑s)}[2] \dicDirectTranslationCS{potlesk, aplaus}
\dicEntry[lófatak] \dicTerm{lófa··|tak} \dicIPA{{l}{ou}{\textlengthmark}{v}{a}{t\smash{\textsuperscript{h}}}{a}{\r{g}}} \dicPos{n}[2] \dicFlx{(‑taks, ‑tök)}[8] \dicSynonym{lófaklapp} \dicDirectTranslationCS{potlesk, aplaus}
\dicEntry[lófatölva] \dicTerm{lófa··tölv|a} \dicIPA{{l}{ou}{\textlengthmark}{v}{a}{t\smash{\textsuperscript{h}}}{\oe}{l}{v}{a}} \dicPos{f}[1] \dicFlx{(‑u, ‑ur)}[7] \dicFieldCat{poč.} \dicDirectTranslationCS{palmtop, PDA, kapesní počítač}
\dicEntry[lófi] \dicTerm{lóf|i} \dicsymFrequent\  \dicIPA{{l}{ou}{\textlengthmark}{v}{\textsci}} \dicPos{m}[1] \dicFlx{(‑a, ‑ar)}[1] \dicDirectTranslationCS{dlaň} \dicExampleIS{geyma hring í lófa sínum} \dicExampleCS{držet prsten ve své dlani}
\dicEntry[lófótur] \textls[15]{\dicTerm{ló··|fótur} \dicIPA{{l}{ou}{\textlengthmark}{f}{ou}{\textsubring{d}}{\textscy}{\textsubring{r}}} \dicPos{m}[13] \dicFlx{(‑fótar, ‑fætur)}[5] \dicFieldCat{bot.} \dicDirectTranslationCS{prustka obecná}} \textit{(l.~{\textLA{Hippuris vulgaris}})}  \dicsymPhoto\ 
\dicFigure{3847.jpg}{Lófótur}{Lófótur - Zicha Ondřej, Biolib, Copyright/CC-BY-NC}
\dicEntry[lóga] \dicTerm{lóg|a} \dicIPA{{l}{ou}{\textlengthmark}{a}} \dicPos{v}[1] \dicFlx{(‑aði)}[34] \dicFlx{dat} \textbf{1.} \dicSynonym{slátra} \dicDirectTranslationCS{porazit, porážet (ovci ap.)} \dicExampleIS{lóga fé} \dicExampleCS{porážet ovce}  \textbf{2.} \dicSynonym*{aflífa} \dicDirectTranslationCS{utratit, uspat (psa ap.)} \dicExampleIS{lóga hundi} \dicExampleCS{uspat psa}
\dicEntry[lógaritmi] \dicTerm{lógaritm|i} \dicIPA{{l}{ou}{\textlengthmark}{\r{g}}{a}{r}{\textsci}{h}{\textsubring{d}}{m}{\textsci}} \dicPos{m}[1] \dicFlx{(‑a, ‑ar)}[1] \dicFieldCat{mat.} \dicDirectTranslationCS{logaritmus}
\dicEntry[lómur] \dicTerm{lóm|ur} \dicIPA{{l}{ou}{\textlengthmark}{m}{\textscy}{\textsubring{r}}} \dicPos{m}[6] \dicFlx{(‑s, ‑ar)}[24] \textbf{1.} \dicFieldCat{zool.} \dicDirectTranslationCS{potáplice malá} \textit{(l.~{\textLA{Gavia stellata}})}  \dicsymPhoto\   \textbf{2.} \dicSynonym{kveinstafir} \dicDirectTranslationCS{nářky, stížnosti};  \dicPhraseIS{berja\,/\addthin lemja lóminn} \dicDirectTranslationCS{stěžovat si, naříkat}
\dicFigure{21071.jpg}{Lómur}{Lómur - Bowman, Tim, Biolib, PD}
\dicEntry[lón] \dicTerm{lón} \dicIPA{{l}{ou}{\textlengthmark}{\textsubring{n}}} \dicPos{n}[2] \dicFlx{(‑s, ‑)}[5] \dicDirectTranslationCS{laguna};  \dicPhraseIS{Bláa lónið} \dicFieldCat{geog.} \dicDirectTranslationCS{Modrá laguna} \dicIndirectTranslationCS{(laguna v~blízkosti Keflavíku)}
\dicEntry[lóuþræll] \dicTerm{lóu··þræl|l} \dicIPA{{l}{ou}{\textlengthmark}{\textscy}{\texttheta}{r}{a}{i}{\textsubring{d}}{\textsubring{l}}} \dicPos{m}[6] \dicFlx{(‑s, ‑ar)}[51] \dicFieldCat{zool.} \dicDirectTranslationCS{jespák, jespák obecný} \textit{(l.~{\textLA{Calidris alpina}})}  \dicsymPhoto\ 
\dicFigure{21101.jpg}{Lóuþræll}{Lóuþræll - Bowman, Tim, Biolib, PD}
\dicEntry[luggarasegl] \dicTerm{luggara··segl} \dicIPA{{l}{\textscy}{\r{g}}{\textlengthmark}{a}{r}{a}{s}{\textepsilon}{\r{g}}{\textsubring{l}}} \dicPos{n}[2] \dicFlx{(‑s, ‑)}[5] \dicFieldCat{nám.} \dicDirectTranslationCS{lugrová plachta}
\dicEntry[lugt] \dicTerm{lugt} \dicIPA{{l}{\textscy}{x}{\textsubring{d}}} \dicPos{f}[7] \dicFlx{(‑ar, ‑ir)}[1] \dicLink{lukt}
\dicEntry[lugum] \dicTerm{lugum} \dicIPA{{l}{\textscy}{\textlengthmark}{\textbabygamma}{\textscy}{\textsubring{m}}} \dicPos{v} \dicFlx{ind pf pl 1 pers} \dicLink{ljúga}
\dicEntry[lukið] \dicTerm{lukið} \dicIPA{{l}{\textscy}{\textlengthmark}{\r{\textObardotlessj}}{\textsci}{\texttheta}} \dicPos{v} \dicFlx{supin} \dicLink{lykja}
\dicEntry[lukka] \dicTerm{lukk|a} \dicIPA{{l}{\textscy}{h}{\r{g}}{a}} \dicPos{f}[1] \dicFlx{(‑u)}[5] \dicSynonym{hamingja} \dicDirectTranslationCS{štěstí, zdar} \dicExampleIS{óska e‑m til lukku með afmælið} \dicExampleCS{popřát (komu) štěstí k~narozeninám}
\dicEntry[lukkast] \dicTerm{lukk|ast} \dicIPA{{l}{\textscy}{h}{\r{g}}{a}{s}{\textsubring{d}}} \dicPos{v}[1] \dicFlx{(‑aðist)}[97] \dicFlx{refl} \dicSynonym{heppnast} \dicDirectTranslationCS{(po)dařit se, zdařit se} \dicExampleIS{Hjónabandið lukkast vel.} \dicExampleCS{Manželství se daří.}
\dicEntry[lukkulegur] \dicTerm{lukku··legur} \dicIPA{{l}{\textscy}{h}{\r{g}}{\textscy}{l}{\textepsilon}{\textbabygamma}{\textscy}{\textsubring{r}}} \dicPos{adj}[1]\dicFlx{}[-8] \dicSynonym{ánægður} \dicDirectTranslationCS{šťastný, spokojený, potěšený}
\dicEntry[lukt] \dicTerm{lukt}\dicTerm{, lugt} \dicIPA{{l}\-{\textscy}\-{x}\-{\textsubring{d}}\-} \dicPos{f}[7] \dicFlx{(‑ar, ‑ir)}[1] \textbf{1.} \dicDirectTranslationCS{lucerna, svítilna} \dicExampleIS{vasalukt} \dicExampleCS{kapesní svítilna}  \textbf{2.} \dicDirectTranslationCS{přední světlo (u~auta, kola ap.)}
\dicEntry[lukti] \dicTerm{lukti} \dicIPA{{l}{\textscy}{x}{\textsubring{d}}{\textsci}} \dicPos{v} \dicFlx{ind pf sg 1 pers} \dicLink{lykja}
\dicEntry[luktum] \dicTerm{luktum} \dicIPA{{l}{\textscy}{x}{\textsubring{d}}{\textscy}{\textsubring{m}}} \dicPos{v} \dicFlx{ind pf pl 1 pers} \dicLink{lykja}
\dicEntry[luktur] \dicTerm{luktur} \dicIPA{{l}{\textscy}{x}{\textsubring{d}}{\textscy}{\textsubring{r}}} \dicPos{adj}[1]\dicFlx{}[-10] \dicDirectTranslationCS{(u)zavřený, obklopený}
\dicEntry[lukum] \dicTerm{lukum} \dicIPA{{l}{\textscy}{\textlengthmark}{\r{g}}{\textscy}{\textsubring{m}}} \dicPos{v} \dicFlx{ind pf pl 1 pers} \dicLink{ljúka}
\dicEntry[luma] \dicTerm{lum|a} \dicIPA{{l}{\textscy}{\textlengthmark}{m}{a}} \dicPos{v}[1] \dicFlx{(‑aði)}[1] \dicPhraseIS{luma á e‑u} \dicDirectTranslationCS{ukrýt si (co), schovat si (co), ulít si (co)} \dicExampleIS{Hann lumar á peningum.} \dicExampleCS{Ukrývá si peníze.}
\dicEntry[lumbra] \dicTerm{lumbr|a} \dicIPA{{l}{\textscy}{m}{\textsubring{b}}{r}{a}} \dicPos{v}[1] \dicFlx{(‑aði)}[1] \dicDirectTranslationCS{zašantročit, zastrčit} \dicExampleIS{lumbra e‑u e‑s staðar} \dicExampleCS{zašantročit (co kde)};  \dicIdiom{lumbra}[á]{ \dicPhraseIS{lumbra á e‑m}} \dicSynonym{berja} \dicDirectTranslationCS{(z)bít (koho), (z)mlátit (koho)}
\dicEntry[lumma] \dicTerm{lumm|a} \dicIPA{{l}{\textscy}{m}{\textlengthmark}{a}} \dicPos{f}[1] \dicFlx{(‑u, ‑ur)}[7] \dicFieldCat{kulin.} \dicDirectTranslationCS{lívanec, lívaneček};  \dicPhraseIS{e‑að selst eins og heitar lummur} \dicFlx{refl} \dicLangCat{přen.} \dicDirectTranslationCS{(co) jde na dračku};  \dicPhraseIS{það er gömul lumma} \dicLangCat{přen.} \dicDirectTranslationCS{to je ohraná písnička}
\dicEntry[lummulegur] \dicTerm{lummu··legur} \dicIPA{{l}{\textscy}{m}{\textlengthmark}{\textscy}{l}{\textepsilon}{\textbabygamma}{\textscy}{\textsubring{r}}} \dicPos{adj}[1]\dicFlx{}[-8] \dicLangCat{hovor.} \dicSynonym{tilkomulítill} \dicDirectTranslationCS{nijaký, nezáživný, nudný}
\dicEntry[lund] \dicTerm{lund} \dicsymFrequent\  \dicIPA{{l}{\textscy}{n}{\textsubring{d}}} \dicPos{f}[7] \dicFlx{(‑ar, ‑ir)}[1] \textbf{1.} \dicSynonym{geð} \dicDirectTranslationCS{povaha, temperament} \dicExampleIS{vera léttur í lund} \dicExampleCS{být šťastný}  \textbf{2.} \dicSynonym{háttur\smash{\textsuperscript{1}}} \dicDirectTranslationCS{způsob}
\dicEntry[Lundey] \dicTerm{Lund··ey} \dicIPA{{l}{\textscy}{n}{\textsubring{d}}{ei}} \dicPos{f}[4] \dicFlx{(‑jar)}[17] \dicFieldCat{geog.} \dicDirectTranslationCS{Lundey} \dicIndirectTranslationCS{(ostrov severně od Islandu)}
\dicEntry[lundi] \dicTerm{lund|i} \dicIPA{{l}{\textscy}{n}{\textsubring{d}}{\textsci}} \dicPos{m}[1] \dicFlx{(‑a, ‑ar)}[1] \dicFieldCat{zool.} \dicDirectTranslationCS{papuchálek, papuchalk, papuchalk ploskozobý\,/\addthin bělobradý} \textit{(l.~{\textLA{Fratercula arctica}})}  \dicsymPhoto\ 
\dicFigure{ds_image_lundi_0_1.jpg}{Lundi}{Lundi - Tgo2002, CC BY-SA 2.5}
\dicEntry[lundir] \dicTerm{lundir} \dicIPA{{l}{\textscy}{n}{\textsubring{d}}{\textsci}{\textsubring{r}}} \dicPos{f}[7] \dicFlx{pl}[2] \dicDirectTranslationCS{jemné filety, svíčková}
\dicEntry[lundur] \dicTerm{lund|ur} \dicIPA{{l}{\textscy}{n}{\textsubring{d}}{\textscy}{\textsubring{r}}} \dicPos{m}[10] \dicFlx{(‑ar, ‑ir)}[4] \dicDirectTranslationCS{háj, lesík}
\dicEntry[Lundúnir] \dicTerm{Lundúnir} \dicIPA{{l}{\textscy}{n}{\textsubring{d}}{u}{n}{\textsci}{\textsubring{r}}} \dicPos{f}[12] \dicFlx{pl}[6] \dicFieldCat{geog.} \dicDirectTranslationCS{Londýn} \dicIndirectTranslationCS{(hlavní město Spojeného království)}
\dicEntry[lunga] \dicTerm{lung|a} \dicIPA{{l}{u}{\ng}{\r{g}}{a}} \dicPos{n}[1] \dicFlx{(‑a, ‑u)}[2] \dicFieldCat{anat.} \dicDirectTranslationCS{plíce}
\dicEntry[lungnabólga] \dicTerm{lungna··bólg|a} \dicIPA{{l}{u}{\ng}{n}{a}{\textsubring{b}}{ou}{l}{\r{g}}{a}} \dicPos{f}[1] \dicFlx{(‑u)}[5] \dicFieldCat{med.} \dicDirectTranslationCS{zápal plic, pneumonie}
\dicEntry[lungnakvef] \dicTerm{lungna··kvef} \dicIPA{{l}{u}{\ng}{n}{a}{k\smash{\textsuperscript{h}}}{v}{\textepsilon}{f}} \dicPos{n}[2] \dicFlx{(‑s)}[2] \dicFieldCat{med.} \dicDirectTranslationCS{zánět průdušek, bronchitida}
\dicEntry[lungnapípa] \dicTerm{lungna··píp|a} \dicIPA{{l}{u}{\ng}{n}{a}{p\smash{\textsuperscript{h}}}{i}{\textsubring{b}}{a}} \dicPos{f}[1] \dicFlx{(‑u, ‑ur)}[19] \dicFieldCat{anat.} \dicDirectTranslationCS{průduška}
\dicEntry[luntalegur] \dicTerm{lunta··legur} \dicIPA{{l}{\textscy}{\textsubring{n}}{\textsubring{d}}{a}{l}{\textepsilon}{\textbabygamma}{\textscy}{\textsubring{r}}} \dicPos{adj}[1]\dicFlx{}[-8] \dicSynonym{ólundarlegur} \dicDirectTranslationCS{mrzutý, nevrlý}
\dicEntry[luralegur] \dicTerm{lura··legur} \dicIPA{{l}{\textscy}{\textlengthmark}{r}{a}{l}{\textepsilon}{\textbabygamma}{\textscy}{\textsubring{r}}} \dicPos{adj}[1]\dicFlx{}[-8] \dicSynonym*{þyngslalegur} \dicDirectTranslationCS{nemotorný, neohrabaný}
\dicEntry[lustum] \dicTerm{lustum} \dicIPA{{l}{\textscy}{s}{\textsubring{d}}{\textscy}{\textsubring{m}}} \dicPos{v} \dicFlx{ind pf pl 1 pers} \dicLink{ljósta}
\dicEntry[lutum] \dicTerm{lutum} \dicIPA{{l}{\textscy}{\textlengthmark}{\textsubring{d}}{\textscy}{\textsubring{m}}} \dicPos{v} \dicFlx{ind pf pl 1 pers} \dicLink{lúta}
\dicEntry[lúalegur] \dicTerm{lúa··legur} \dicIPA{{l}{u}{\textlengthmark}{a}{l}{\textepsilon}{\textbabygamma}{\textscy}{\textsubring{r}}} \dicPos{adj}[1]\dicFlx{}[-8] \textbf{1.} \dicSynonym{þreytulegur} \dicDirectTranslationCS{vyčerpaný, unavený}  \textbf{2.} \dicSynonym{ódrengilegur} \dicDirectTranslationCS{hanebný, sprostý} \dicExampleIS{lúalegar getsakir} \dicExampleCS{hanebné narážky}
\dicEntry[lúð] \dicTerm{lúð} \dicIPA{{l}{u}{\textlengthmark}{\texttheta}} \dicPos{v} \dicFlx{supin} \dicLink{lýja}
\dicEntry[lúða] \dicTerm{lúð|a} \dicIPA{{l}{u}{\textlengthmark}{ð}{a}} \dicPos{f}[1] \dicFlx{(‑u, ‑ur)}[19] \dicFieldCat{zool.} \dicSynonym{flyðra} \dicDirectTranslationCS{platýs, platýs obecný} \textit{(l.~{\textLA{Hippoglossus hippoglossus}})}
\dicEntry[lúðalegur] \dicTerm{lúða··legur} \dicIPA{{l}{u}{\textlengthmark}{ð}{a}{l}{\textepsilon}{\textbabygamma}{\textscy}{\textsubring{r}}} \dicPos{adj}[1]\dicFlx{}[-8] \dicDirectTranslationCS{slizký, odporný}
\dicEntry[lúði] \dicTerm{lúð|i\smash{\textsuperscript{1}}} \dicIPA{{l}{u}{\textlengthmark}{ð}{\textsci}} \dicPos{m}[1] \dicFlx{(‑a, ‑ar)}[1] \dicDirectTranslationCS{křupan, balík}
\dicEntry[lúði] \dicTerm{lúði\smash{\textsuperscript{2}}} \dicIPA{{l}{u}{\textlengthmark}{ð}{\textsci}} \dicPos{v} \dicFlx{ind pf sg 1 pers} \dicLink{lýja}
\dicEntry[lúðrasveit] \dicTerm{lúðra··sveit} \dicIPA{{l}{u}{ð}{r}{a}{s}{v}{ei}{\textsubring{d}}} \dicPos{f}[7] \dicFlx{(‑ar, ‑ir)}[1] \dicFieldCat{hud.} \dicDirectTranslationCS{dechová kapela, dechovka}
\dicEntry[lúðum] \dicTerm{lúðum} \dicIPA{{l}{u}{\textlengthmark}{ð}{\textscy}{\textsubring{m}}} \dicPos{v} \dicFlx{ind pf pl 1 pers} \dicLink{lýja}
\dicEntry[lúður] \dicTerm{lúð|ur} \dicIPA{{l}{u}{\textlengthmark}{ð}{\textscy}{\textsubring{r}}} \dicPos{m}[5] \dicFlx{(‑urs, ‑rar)}[1] \dicFieldCat{hud.} \dicDirectTranslationCS{trumpeta} \dicExampleIS{blása í lúður} \dicExampleCS{foukat na trumpetu}
\dicEntry[lúðurþeytari] \dicTerm{lúður··þeyt·ar|i} \dicIPA{{l}{u}{\textlengthmark}{ð}{\textscy}{\textsubring{r}}{\texttheta}{ei}{\textsubring{d}}{a}{r}{\textsci}} \dicPos{m}[1] \dicFlx{(‑a, ‑ar)}[13] \dicDirectTranslationCS{trubač(ka)}
\dicEntry[lúga] \dicTerm{lúg|a} \dicIPA{{l}{u}{\textlengthmark}{a}} \dicPos{f}[1] \dicFlx{(‑u, ‑ur)}[7] \dicDirectTranslationCS{poklop, příklop}
\dicEntry[lúi] \dicTerm{lú|i} \dicIPA{{l}{u}{\textlengthmark}{\textsci}} \dicPos{m}[1] \dicFlx{(‑a)}[3] \dicSynonym{þreyta\smash{\textsuperscript{1}}} \dicDirectTranslationCS{znavení, vyčerpanost, vysílení} \dicExampleIS{finna fyrir lúa eftir ferðina} \dicExampleCS{cítit vysílení po cestě}
\dicEntry[lúið] \dicTerm{lúið} \dicIPA{{l}{u}{\textlengthmark}{\textsci}{\texttheta}} \dicPos{v} \dicFlx{supin} \dicLink{lýja}
\dicEntry[lúinn] \dicTerm{lúinn} \dicIPA{{l}{u}{\textlengthmark}{\textsci}{\textsubring{n}}} \dicPos{adj}[6]\dicFlx{}[-2] \textbf{1.} \dicSynonym{þreyttur} \dicDirectTranslationCS{znavený, vyčerpaný, vysílený}  \textbf{2.} \dicSynonym{slitinn} \dicDirectTranslationCS{opotřebovaný, obnošený}
\dicEntry[lúka] \dicTerm{lúk|a} \dicIPA{{l}{u}{\textlengthmark}{\r{g}}{a}} \dicPos{f}[1] \dicFlx{(‑u, ‑ur)}[19] \textbf{1.} \dicSynonym{lófi} \dicDirectTranslationCS{dlaň};  \dicPhraseIS{vera með lífið í lúkunum} \dicLangCat{přen.} \dicDirectTranslationCS{mít srdce v~kalhotách}  \textbf{2.} \dicSynonym{handfylli} \dicDirectTranslationCS{hrst} \dicExampleIS{lúka af salti} \dicExampleCS{hrst soli}
\dicEntry[lúkar] \dicTerm{lúkar} \dicIPA{{l}{u}{\textlengthmark}{\r{g}}{a}{\textsubring{r}}} \dicPos{m}[4] \dicFlx{(‑s, ‑ar)}[15] \dicFieldCat{nám.} \dicSynonym{káeta} \dicDirectTranslationCS{důstojnická kabina\,/\addthin kajuta}
\dicEntry[lúlla] \dicTerm{lúll|a} \dicIPA{{l}{u}{l}{\textlengthmark}{a}} \dicPos{v}[1] \dicFlx{(‑aði)}[44] \dicLangCat{dět.} \dicSynonym{sofa} \dicDirectTranslationCS{spinkat, hajat} \dicExampleIS{Farðu nú að lúlla.} \dicExampleCS{Jdi spinkat.}
\dicEntry[lúmskur] \dicTerm{lúmskur} \dicIPA{{l}{u}{m}{s}{\r{g}}{\textscy}{\textsubring{r}}} \dicPos{adj}[1]\dicFlx{}[-1] \dicSynonym{undirförull} \dicDirectTranslationCS{podlý, zákeřný};  \dicPhraseIS{hafa lúmskt gaman af að (gera e‑ð)} \dicDirectTranslationCS{mít škodolibou radost z~(dělání (čeho))}
\dicEntry[lúpína] \dicTerm{lúpín|a} \dicIPA{{l}{u}{\textlengthmark}{p\smash{\textsuperscript{h}}}{i}{n}{a}} \dicPos{f}[1] \dicFlx{(‑u, ‑ur)}[7] \dicFieldCat{bot.} \dicDirectTranslationCS{lupina nutkajská, vlčí bob} \textit{(l.~{\textLA{Lupinus nootkatensis}})}  \dicsymPhoto\ 
\dicFigure{ds_image_lupina_0_1.jpg}{Lúpína}{Lúpína - Jutta234, GFDL}
\dicEntry[lúpulegur] \dicTerm{lúpu··legur} \dicIPA{{l}{u}{\textlengthmark}{\textsubring{b}}{\textscy}{l}{\textepsilon}{\textbabygamma}{\textscy}{\textsubring{r}}}\addthinS\dicPos{adj}[1]\dicFlx{}[-8] \dicSynonym{skömmustulegur} \dicDirectTranslationCS{zahanbený, (jsoucí) v~rozpacích}
\dicEntry[lúr] \dicTerm{lúr} \dicIPA{{l}{u}{\textlengthmark}{\textsubring{r}}} \dicPos{m}[4] \dicFlx{(‑s, ‑ar)}[14] \dicSynonym{blundur} \dicDirectTranslationCS{zdřímnutí, šlofík} \dicExampleIS{fá sér svolítinn lúr} \dicExampleCS{trochu si zdřímnout}
\dicEntry[lúra] \dicTerm{lúr|a} \dicIPA{{l}{u}{\textlengthmark}{r}{a}} \dicPos{v}[2] \dicFlx{(‑ði, ‑að\,/\addthin ‑t)}[184] \dicSynonym{sofa} \dicDirectTranslationCS{zdřímnout si, dřímat, podřimovat};  \dicIdiom{lúra}[á]{ \dicPhraseIS{lúra á e‑u}} \dicDirectTranslationCS{schovávat (co), ukrýt (si) (co)} \dicExampleIS{lúra á góðgæti} \dicExampleCS{schovávat dobroty}
\dicEntry[lús] \dicTerm{lús} \dicIPA{{l}{u}{\textlengthmark}{s}} \dicPos{f}[9] \dicFlx{(‑ar, lýs)}[6] \dicFieldCat{zool.} \dicDirectTranslationCS{veš} \textit{(l.~{\textLA{Anoplura}})};  \dicPhraseIS{læðast eins og lús með saum(i)} \dicFlx{refl} \dicLangCat{přen.} \dicDirectTranslationCS{loudat se jako hlemýžď};  \dicPhraseIS{lús kviknaði milli þeirra} \dicLangCat{přen.} \dicDirectTranslationCS{dostali se do křížku};  \dicPhraseIS{nú detta mér allar dauðar lýs úr höfði} \dicLangCat{přen.} \dicDirectTranslationCS{to jsem z~toho jelen}
\dicEntry[lúsarlegur] \dicTerm{lúsar··legur} \dicIPA{{l}{u}{\textlengthmark}{s}{a}{r}{l}{\textepsilon}{\textbabygamma}{\textscy}{\textsubring{r}}} \dicPos{adj}[1]\dicFlx{}[-8] \dicSynonym{smásálarlegur} \dicDirectTranslationCS{ubohý, tristní, mizerný (plat ap.)}
\dicEntry[lúta] \dicTerm{lúta} \dicsymFrequent\  \dicIPA{{l}{u}{\textlengthmark}{\textsubring{d}}{a}} \dicPos{v}[6] \dicFlx{(lýt, laut, lutum, lyti, lotið)}[106] \dicFlx{dat} \textbf{1.} \dicSynonym{beygja\smash{\textsuperscript{2}}} \dicDirectTranslationCS{sklonit, sklánět, sehnout, shýbat};  \dicPhraseIS{lúta höfði} \dicDirectTranslationCS{sklonit hlavu}  \textbf{2.} \dicDirectTranslationCS{sklonit se, podrobit se};  \dicPhraseIS{lúta e‑m} \dicDirectTranslationCS{sklonit se před (kým), podrobit se (komu)};  \dicIdiom{lúta}[að]{ \dicPhraseIS{lúta að e‑u}} \dicDirectTranslationCS{vztahovat se na (co), týkat se (čeho)};  \dicIdiom{lúta}{ \dicPhraseIS{lúta (svo) lágt}} \dicLangCat{přen.} \dicDirectTranslationCS{snížit se, klesnout} \dicIndirectTranslationCS{(propůjčit se k~něčemu špatnému)}; { \dicPhraseIS{lúta í lægra haldi fyrir e‑m}} \dicLangCat{přen.} \dicDirectTranslationCS{sklonit se před (kým)} \dicIndirectTranslationCS{(přijmout prohru)}
\dicEntry[lúterskur] \dicTerm{lúterskur} \dicIPA{{l}{u}{\textlengthmark}{t\smash{\textsuperscript{h}}}{e}{\textsubring{r}}{s}{\r{g}}{\textscy}{\textsubring{r}}} \dicPos{adj}[1]\dicFlx{}[-1] \dicFieldCat{náb.} \dicDirectTranslationCS{luteránský}
\dicEntry[lúterstrú] \dicTerm{lúters··trú} \dicIPA{{l}{u}{\textlengthmark}{t\smash{\textsuperscript{h}}}{\textepsilon}{\textsubring{r}}{s}{t\smash{\textsuperscript{h}}}{r}{u}} \dicPos{f}[4] \dicFlx{(‑ar)}[27] \dicFieldCat{náb.} \dicDirectTranslationCS{luteránství}
\dicEntry[lútur] \dicTerm{lút|ur} \dicIPA{{l}{u}{\textlengthmark}{\textsubring{d}}{\textscy}{\textsubring{r}}} \dicPos{m}[6] \dicFlx{(‑s\,/\addthin ‑ar, ‑ar)}[65] \dicFieldCat{chem.} \dicDirectTranslationCS{louh}
\dicEntry[Lúxemborg] \dicTerm{Lúxem··borg} \dicIPA{{l}{u}{x}{s}{\textepsilon}{m}{\textsubring{b}}{\textopeno}{r}{\r{g}}} \dicPos{f}[7] \dicFlx{(‑ar)}[4] \textbf{1.} \dicFieldCat{geog.} \dicDirectTranslationCS{Lucembursko}  \textbf{2.} \dicFieldCat{geog.} \dicDirectTranslationCS{Lucemburk} \dicIndirectTranslationCS{(hlavní město Lucemburska)}
\dicEntry[Lúxemborgari] \dicTerm{Lúxem··borg·ar|i} \dicIPA{{l}{u}{x}{s}{\textepsilon}{m}{\textsubring{b}}{\textopeno}{r}{\r{g}}{a}{r}{\textsci}} \dicPos{m}[1] \dicFlx{(‑a, ‑ar)}[13] \dicDirectTranslationCS{Lucemburčan(ka)}
\dicEntry[lúxemborgíska] \dicTerm{lúxem··borgísk|a} \dicIPA{{l}{u}{x}{s}{\textepsilon}{m}{\textsubring{b}}{\textopeno}{r}{\r{\textObardotlessj}}{i}{s}{\r{g}}{a}} \dicPos{f}[1] \dicFlx{(‑u)}[5] \dicDirectTranslationCS{lucemburština}
\dicEntry[lúxemborgskur] \dicTerm{lúxem··borgskur} \dicIPA{{l}{u}{x}{s}{\textepsilon}{m}{\textsubring{b}}{\textopeno}{r}{\r{g}}{s}{\r{g}}{\textscy}{\textsubring{r}}} \dicPos{adj}[1]\dicFlx{}[-6] \dicDirectTranslationCS{lucemburský}
\dicEntry[lúxus] \dicTerm{lúxus} \dicIPA{{l}{u}{x}{s}{\textscy}{s}} \dicPos{m}[4] \dicFlx{(‑s)}[21] \dicDirectTranslationCS{luxus, přepych}
\dicEntry[lyf] \dicTerm{lyf} \dicsymFrequent\  \dicIPA{{l}{\textsci}{\textlengthmark}{f}} \dicPos{n}[2] \dicFlx{(‑s, ‑)}[13] \dicDirectTranslationCS{lék} \dicExampleIS{lyf gegn lyfseðli} \dicExampleCS{lék na předpis};  \dicPhraseIS{lyf gegn e‑u} \dicDirectTranslationCS{lék proti (čemu), lék na (co)} \dicExampleIS{lyf gegn kvefi} \dicExampleCS{lék na rýmu};  \dicPhraseIS{lyf við e‑ð} \dicDirectTranslationCS{lék na (co)};  \dicPhraseIS{taka inn lyf} \dicDirectTranslationCS{brát léky};  \dicPhraseIS{verja á lyfjum} \dicDirectTranslationCS{být na prášcích}
\dicEntry[lyfjabúð] \dicTerm{lyfja··búð} \dicIPA{{l}{\textsci}{v}{j}{a}{\textsubring{b}}{u}{\texttheta}} \dicPos{f}[7] \dicFlx{(‑ar, ‑ir)}[1] \dicDirectTranslationCS{lékárna}
\dicEntry[lyfjafræði] \dicTerm{lyfja··fræð|i} \dicIPA{{l}{\textsci}{v}{j}{a}{f}{r}{a}{i}{ð}{\textsci}} \dicPos{f}[3] \dicFlx{(‑i)}[3] \textbf{1.} \dicDirectTranslationCS{lékárnictví, farmacie}  \textbf{2.} \dicDirectTranslationCS{farmakologie}
\dicEntry[lyfjafræðingur] \dicTerm{lyfja·fræð··ing|ur} \dicIPA{{l}{\textsci}{v}{j}{a}{f}{r}{a}{i}{ð}{i}{\ng}{\r{g}}{\textscy}{\textsubring{r}}} \dicPos{m}[6] \dicFlx{(‑s, ‑ar)}[8] \textbf{1.} \dicDirectTranslationCS{lékárník, lékárnice, farmaceut(ka)}  \textbf{2.} \dicDirectTranslationCS{farmakolog, farmakoložka}
\dicEntry[lyfjagras] \dicTerm{lyfja··|gras} \dicIPA{{l}{\textsci}{v}{j}{a}{\r{g}}{r}{a}{s}} \dicPos{n}[2] \dicFlx{(‑grass, ‑grös)}[8] \dicFieldCat{bot.} \dicDirectTranslationCS{tučnice obecná} \textit{(l.~{\textLA{Pinguicula vulgaris}})}  \dicsymPhoto\ 
\dicFigure{31254.jpg}{Lyfjagras}{Lyfjagras - Dvořák Václav, Biolib, Copyright/CC-BY-NC}
\dicEntry[lyfjakort] \dicTerm{lyfja··kort} \dicIPA{{l}{\textsci}{v}{j}{a}{k\smash{\textsuperscript{h}}}{\textopeno}{\textsubring{r}}{\textsubring{d}}} \dicPos{n}[2] \dicFlx{(‑s, ‑)}[5] \dicIndirectTranslationCS{karta pacienta s~informacemi o~užívaných lécích}
\dicEntry[lyfjakostnaður] \dicTerm{lyfja··kost·nað|ur} \dicIPA{{l}{\textsci}{v}{j}{a}{k\smash{\textsuperscript{h}}}{\textopeno}{s}{\textsubring{d}}{n}{a}{ð}{\textscy}{\textsubring{r}}} \dicPos{m}[10] \dicFlx{(‑ar)}[9] \dicDirectTranslationCS{náklady na léky}
\dicEntry[lyfjaskírteini] \dicTerm{lyfja··skír·teini} \dicIPA{{l}{\textsci}{v}{j}{a}{s}{\r{\textObardotlessj}}{i}{\textsubring{r}}{t\smash{\textsuperscript{h}}}{ei}{n}{\textsci}} \dicPos{n}[2] \dicFlx{(‑s, ‑)}[14] \dicIndirectTranslationCS{průkaz opravňující ke slevě na léky}
\dicEntry[lyfjaverð] \dicTerm{lyfja··verð} \dicIPA{{l}{\textsci}{v}{j}{a}{v}{\textepsilon}{r}{\texttheta}} \dicPos{n}[2] \dicFlx{(‑s)}[2] \dicDirectTranslationCS{cena léku}
\dicEntry[lyfjaverslun] \dicTerm{lyfja··versl|un} \dicIPA{{l}{\textsci}{v}{j}{a}{v}{\textepsilon}{\textsubring{r}}{s}{\textsubring{d}}{l}{\textscy}{\textsubring{n}}} \dicPos{f}[7] \dicFlx{(‑unar, ‑anir)}[8] \dicSynonym{lyfjabúð} \dicDirectTranslationCS{lékárna}
\dicEntry[lyflækningar] \dicTerm{lyf··lækn·ingar} \dicIPA{{l}{\textsci}{v}{l}{a}{i}{h}{\r{g}}{n}{i}{\ng}{\r{g}}{a}{\textsubring{r}}} \dicPos{f}[4] \dicFlx{pl}[6] \dicFieldCat{med.} \dicDirectTranslationCS{klinická medicína}
\dicEntry[lyfsali] \dicTerm{lyf··sal|i} \dicIPA{{l}{\textsci}{f}{s}{a}{l}{\textsci}} \dicPos{m}[1] \dicFlx{(‑a, ‑ar)}[8] \dicDirectTranslationCS{lékárník, lékárnice}
\dicEntry[lyfseðill] \dicTerm{lyf··seð|ill} \dicIPA{{l}{\textsci}{f}{s}{\textepsilon}{ð}{\textsci}{\textsubring{d}}{\textsubring{l}}} \dicPos{m}[6] \dicFlx{(‑ils, ‑lar)}[35] \dicDirectTranslationCS{(lékařský) předpis}
\dicEntry[lyfta] \dicTerm{lyft|a\smash{\textsuperscript{1}}} \dicIPA{{l}{\textsci}{f}{\textsubring{d}}{a}} \dicPos{f}[1] \dicFlx{(‑u, ‑ur)}[7] \textbf{1.} \dicDirectTranslationCS{výtah}  \textbf{2.} \dicDirectTranslationCS{zdviž}  \textbf{3.} \dicSynonym{skíðalyfta} \dicDirectTranslationCS{vlek (lyžařský ap.)}
\dicEntry[lyfta] \dicTerm{lyft|a\smash{\textsuperscript{2}}} \dicsymFrequent\  \dicIPA{{l}{\textsci}{f}{\textsubring{d}}{a}} \dicPos{v}[2] \dicFlx{(‑i, ‑)}[14] \dicFlx{dat} \dicSynonym{hefja} \dicDirectTranslationCS{zvednout, zvedat, zdvihat, nadzvednout, nadzdvihnout} \dicExampleIS{lyfta steini} \dicExampleCS{zvednout kámen};  \dicPhraseIS{lyfta brúnum} \dicDirectTranslationCS{zvednout obočí};  \dicIdiom{lyfta}[sér]{ \dicPhraseIS{lyfta sér}} \dicDirectTranslationCS{(na)kynout (těsto ap.)};  \dicIdiom{lyfta}[upp]{ \dicPhraseIS{lyfta sér upp}} \dicSynonym*{skemmta sér} \dicDirectTranslationCS{povyrazit si, odvázat se};  \dicIdiom{lyftast}{ \dicPhraseIS{lyftast}} \dicFlx{refl} \dicDirectTranslationCS{zvednout\,/\addthin zvedat se (letadlo ap.)}
\dicEntry[lyftari] \dicTerm{lyft··ar|i} \dicIPA{{l}{\textsci}{f}{\textsubring{d}}{a}{r}{\textsci}} \dicPos{m}[1] \dicFlx{(‑a, ‑ar)}[13] \dicFieldCat{techn.} \dicDirectTranslationCS{vysokozdvižný vozík}
\dicEntry[lyftiduft] \dicTerm{lyfti··duft} \dicIPA{{l}{\textsci}{f}{\textsubring{d}}{\textsci}{\textsubring{d}}{\textscy}{f}{\textsubring{d}}} \dicPos{n}[2] \dicFlx{(‑s)}[2] \dicDirectTranslationCS{prášek do pečiva}
\dicEntry[lyfting] \dicTerm{lyft··ing} \dicIPA{{l}{\textsci}{f}{\textsubring{d}}{i}{\ng}{\r{g}}} \dicPos{f}[4] \dicFlx{(‑ar, ‑ar)}[5] \textbf{1.} \dicSynonym*{það að lyfta} \dicDirectTranslationCS{zvednutí, zvedání, (vy)zdvižení, nadzvednutí}  \textbf{2.} \dicSynonym*{fjörgun} \dicDirectTranslationCS{pozvednutí, povznesení (na duchu ap.)} \dicExampleIS{lyfting andans} \dicExampleCS{povznesení na duchu}  \textbf{3.} \dicPhraseIS{lyftingar} \dicFlx{pl} \dicFieldCat{sport.} \dicDirectTranslationCS{vzpírání}
\begin{xtolerant}{}{1pt}
\dicEntry[lyftingamaður] \dicTerm{lyftinga··|maður}\addthinS\dicIPA{{l}{\textsci}{f}{\textsubring{d}}{i}{\ng}{\r{g}}{a}{m}{a}{ð}{\textscy}{\textsubring{r}}}\addthinS\dicPos{m}[13]\addthinS\dicFlx{(‑manns, ‑menn)}[2] \dicDirectTranslationCS{vzpě\-rač(ka)}
\end{xtolerant}
\dicEntry[lygamælir] \dicTerm{lyga··mæl|ir} \dicIPA{{l}{\textsci}{\textlengthmark}{\textbabygamma}{a}{m}{a}{i}{l}{\textsci}{\textsubring{r}}} \dicPos{m}[7] \dicFlx{(‑is, ‑ar)}[1] \dicDirectTranslationCS{detektor lži}
\dicEntry[lygar] \dicTerm{lygar} \dicIPA{{l}{\textsci}{\textlengthmark}{\textbabygamma}{a}{\textsubring{r}}} \dicPos{f} \dicFlx{pl nom} \dicLink{lygi\smash{\textsuperscript{1}}}
\dicEntry[lygari] \dicTerm{lyg··ar|i} \dicIPA{{l}{\textsci}{\textlengthmark}{\textbabygamma}{a}{r}{\textsci}} \dicPos{m}[1] \dicFlx{(‑a, ‑ar)}[13] \dicDirectTranslationCS{lhář(ka)}
\dicEntry[lygasaga] \dicTerm{lyga··|saga} \dicIPA{{l}{\textsci}{\textlengthmark}{\textbabygamma}{a}{s}{a}{\textbabygamma}{a}} \dicPos{f}[1] \dicFlx{(‑sögu, ‑sögur)}[14] \dicDirectTranslationCS{báchorka, pohádka}
\dicEntry[lygi] \dicTerm{lyg|i\smash{\textsuperscript{1}}} \dicsymFrequent\  \dicIPA{{l}{i}{j}{\textlengthmark}{\textsci}} \dicPos{f}[2] \dicFlx{(‑i, ‑ar)}[1] \dicDirectTranslationCS{lež, nepravda} \dicExampleIS{Bæjarstjórinn segir lygar.} \dicExampleCS{Starosta lže.};  \dicPhraseIS{fara með lygar} \dicDirectTranslationCS{lhát}
\dicEntry[lygi] \dicTerm{lygi\smash{\textsuperscript{2}}} \dicIPA{{l}{i}{j}{\textlengthmark}{\textsci}} \dicPos{v} \dicFlx{con pf sg 1 pers} \dicLink{ljúga}
\dicEntry[lygilegur] \textls[15]{\dicTerm{lygi··legur} \dicIPA{{l}{i}{j}{\textlengthmark}{\textsci}{l}{\textepsilon}{\textbabygamma}{\textscy}{\textsubring{r}}} \dicPos{adj}[1]\dicFlx{}[-8] \textbf{1.} \dicSynonym{ótrúlegur} \dicDirectTranslationCS{nevěrohodný, nepravděpodobný}}  \textbf{2.} \dicSynonym{ósannur} \dicDirectTranslationCS{nepravdivý, lživý}
\dicEntry[lyginn] \dicTerm{lyginn} \dicIPA{{l}{i}{j}{\textlengthmark}{\textsci}{\textsubring{n}}} \dicPos{adj}[6]\dicFlx{}[-2] \dicDirectTranslationCS{ulhaný, prolhaný}
\dicEntry[lygn] \dicTerm{lygn} \dicIPA{{l}{\textsci}{\r{g}}{\textsubring{n}}} \dicPos{adj}[5]\dicFlx{}[-5] \textbf{1.} \dicSynonym*{vindlaus} \dicDirectTranslationCS{bezvětrný, klidný} \dicIndirectTranslationCS{(o~počasí)}  \textbf{2.} \dicSynonym*{straumlaus} \dicDirectTranslationCS{nehybný, nezčeřený, klidný} \dicIndirectTranslationCS{(o~vodě)}
\dicEntry[lygna] \dicTerm{lygn|a\smash{\textsuperscript{1}}} \dicIPA{{l}{\textsci}{\r{g}}{n}{a}} \dicPos{f}[1] \dicFlx{(‑u, ‑ur)}[7] \textbf{1.} \dicSynonym*{kyrrt vatn} \dicDirectTranslationCS{klidná voda}  \textbf{2.} \dicSynonym*{kyrrt loft} \dicDirectTranslationCS{bezvětří}
\dicEntry[lygna] \dicTerm{lygn|a\smash{\textsuperscript{2}}} \dicIPA{{l}{\textsci}{\r{g}}{n}{a}} \dicPos{v}[2] \dicFlx{(‑di, ‑t)}[150] \dicFlx{dat} \textbf{1.} \dicPhraseIS{það lygnir} \dicFlx{impers} \dicSynonym*{kyrrast} \dicDirectTranslationCS{utišuje se (vítr), přestává foukat}  \textbf{2.} \dicPhraseIS{lygna aftur augunum} \dicDirectTranslationCS{zavřít\,/\addthin přivřít oči}
\dicEntry[lyk] \dicTerm{lyk} \dicIPA{{l}{\textsci}{\textlengthmark}{\r{g}}} \dicPos{v} \dicFlx{ind praes sg 1 pers} \dicLink{lykja}
\dicEntry[lyki] \dicTerm{lyki} \dicIPA{{l}{\textsci}{\textlengthmark}{\r{\textObardotlessj}}{\textsci}} \dicPos{v} \dicFlx{con pf sg 1 pers} \dicLink{ljúka}
\dicEntry[lykilaðstaða] \dicTerm{lykil··að·|staða} \dicIPA{{l}{\textsci}{\textlengthmark}{\r{\textObardotlessj}}{\textsci}{l}{a}{ð}{s}{\textsubring{d}}{a}{ð}{a}} \dicPos{f}[1] \dicFlx{(‑stöðu)}[2] \dicDirectTranslationCS{klíčové postavení}
\dicEntry[lykilatriði] \dicTerm{lykil··at·riði} \dicIPA{{l}{\textsci}{\textlengthmark}{\r{\textObardotlessj}}{\textsci}{l}{a}{\textsubring{d}}{r}{\textsci}{ð}{\textsci}} \dicPos{n}[2] \dicFlx{(‑s, ‑)}[14] \dicDirectTranslationCS{klíčový bod, klíčová otázka}
\dicEntry[lykill] \dicTerm{lyk|ill} \dicsymFrequent\  \dicIPA{{l}{\textsci}{\textlengthmark}{\r{\textObardotlessj}}{\textsci}{\textsubring{d}}{\textsubring{l}}} \dicPos{m}[6] \dicFlx{(‑ils, ‑lar)}[35] \textbf{1.} \dicDirectTranslationCS{klíč (od dveří ap.)};  \dicPhraseIS{lykill að e‑u} \dicDirectTranslationCS{klíč od (čeho)} \dicExampleIS{lykill að húsinu} \dicExampleCS{klíč od domu}  \textbf{2.} \dicFieldCat{techn.} \dicDirectTranslationCS{klíč} \dicIndirectTranslationCS{(ruční nástroj)}  \textbf{3.} \dicFieldCat{hud.} \dicDirectTranslationCS{klíč} \dicIndirectTranslationCS{(hudební značka)} \dicExampleIS{nótnalykill} \dicExampleCS{hudební klíč}  \textbf{4.} \dicSynonym{skýring} \dicDirectTranslationCS{klíč, nápověda} \dicIndirectTranslationCS{(pomůcka podstatná pro vysvětlení něčeho)} \dicExampleIS{orðalykill} \dicExampleCS{konkordance}
\dicEntry[lykilmaður] \dicTerm{lykil··|maður} \dicIPA{{l}{\textsci}{\textlengthmark}{\r{\textObardotlessj}}{\textsci}{l}{m}{a}{ð}{\textscy}{\textsubring{r}}} \dicPos{m}[13] \dicFlx{(‑manns, ‑menn)}[2] \dicDirectTranslationCS{klíčová postava}
\dicEntry[lykilorð] \dicTerm{lykil··orð} \dicIPA{{l}{\textsci}{\textlengthmark}{\r{\textObardotlessj}}{\textsci}{l}{\textopeno}{r}{\texttheta}} \dicPos{n}[2] \dicFlx{(‑s, ‑)}[5] \textbf{1.} \dicSynonym{aðalatriði} \dicDirectTranslationCS{klíčové slovo}  \textbf{2.} \dicFieldCat{poč.} \dicSynonym*{aðgangsorð} \dicDirectTranslationCS{(přístupové) heslo}
\dicEntry[lykja] \dicTerm{lykja} \dicIPA{{l}{\textsci}{\textlengthmark}{\r{\textObardotlessj}}{a}} \dicPos{v}[4] \dicFlx{(lyk, lukti, luktum, lykti, lukið)}[61] \dicDirectTranslationCS{obklopit, obklopovat, obstoupit, obstupovat} \dicExampleIS{Vatn lukti um hana.} \dicExampleCS{Voda ji obklopila.};  \dicIdiom{lykjast}{ \dicPhraseIS{lykjast}} \dicFlx{refl} \dicDirectTranslationCS{obklopit, obklopovat, obstoupit, obstupovat}
\dicEntry[lykkja] \dicTerm{lykkj|a} \dicIPA{{l}{\textsci}{h}{\r{\textObardotlessj}}{a}} \dicPos{f}[1] \dicFlx{(‑u, ‑ur)}[25] \textbf{1.} \dicSynonym{hanki} \dicDirectTranslationCS{smyčka, oko, klička (na provázku ap.)} \dicExampleIS{lykkja á bandi} \dicExampleCS{smyčka na provázku}  \textbf{2.} \dicDirectTranslationCS{oko} \dicIndirectTranslationCS{(v~pletení)}  \textbf{3.} \dicSynonym{getnaðarvörn} \dicDirectTranslationCS{spirála} \dicIndirectTranslationCS{(druh antikoncepce)}
\dicEntry[lykkjufall] \dicTerm{lykkju··|fall} \dicIPA{{l}{\textsci}{h}{\r{\textObardotlessj}}{\textscy}{f}{a}{\textsubring{d}}{\textsubring{l}}} \dicPos{n}[2] \dicFlx{(‑falls, ‑föll)}[8] \dicDirectTranslationCS{puštěné oko} \dicIndirectTranslationCS{(v~pletení)} \dicExampleIS{lykkjufall á sokkabuxurnar} \dicExampleCS{oko na punčoše}
\dicEntry[lyklaborð] \dicTerm{lykla··borð} \dicsymFrequent\  \dicIPA{{l}{\textsci}{h}{\r{g}}{l}{a}{\textsubring{b}}{\textopeno}{r}{\texttheta}} \dicPos{n}[2] \dicFlx{(‑s, ‑)}[5] \dicFieldCat{poč.} \dicDirectTranslationCS{klávesnice} \dicExampleIS{íslenskt lyklaborð} \dicExampleCS{islandská klávesnice}
\dicEntry[lyklakippa] \dicTerm{lykla··kipp|a} \dicIPA{{l}{\textsci}{h}{\r{g}}{l}{a}{c\smash{\textsuperscript{h}}}{\textsci}{h}{\textsubring{b}}{a}} \dicPos{f}[1] \dicFlx{(‑u, ‑ur)}[19] \dicDirectTranslationCS{svazek klíčů}
\dicEntry[lykt] \dicTerm{lykt\smash{\textsuperscript{1}}} \dicIPA{{l}{\textsci}{x}{\textsubring{d}}} \dicPos{f}[7] \dicFlx{(‑ar, ‑ir)}[1] \dicSynonym{endir} \dicDirectTranslationCS{konec, závěr};  \dicPhraseIS{að lyktum} \dicFlx{adv} \dicSynonym*{að lokum} \dicDirectTranslationCS{nakonec, na závěr};  \dicPhraseIS{leiða e‑ð til lykta} \dicSynonym*{ljúka e‑u} \dicDirectTranslationCS{dovést (co) do konce, dokončit (co)};  \dicPhraseIS{gera lyktir á e‑u} \dicSynonym*{ljúka við e‑ð} \dicDirectTranslationCS{dokončit\,/\addthin dokončovat (co)}
\dicEntry[lykt] \dicTerm{lykt\smash{\textsuperscript{2}}} \dicsymFrequent\  \dicIPA{{l}{\textsci}{x}{\textsubring{d}}} \dicPos{f}[7] \dicFlx{(‑ar, ‑ir)}[1] \dicSynonym{ilmur} \dicDirectTranslationCS{(zá)pach, vůně, aroma} \dicExampleIS{lykt af blómunum} \dicExampleCS{vůně květin}
\dicEntry[lykta] \dicTerm{lykt|a\smash{\textsuperscript{1}}} \dicIPA{{l}{\textsci}{x}{\textsubring{d}}{a}} \dicPos{v}[1] \dicFlx{(‑aði)}[5] \dicPhraseIS{e‑u lyktar} \dicFlx{impers} \dicDirectTranslationCS{(co) končí} \dicExampleIS{Leiknum lyktaði með jafntefli.} \dicExampleCS{Zápas skončil remízou.}
\dicEntry[lykta] \dicTerm{lykt|a\smash{\textsuperscript{2}}} \dicIPA{{l}{\textsci}{x}{\textsubring{d}}{a}} \dicPos{v}[1] \dicFlx{(‑aði)}[34] \textbf{1.} \dicSynonym{ilma} \dicDirectTranslationCS{páchnout, být cítit, vonět, vydávat vůni} \dicExampleIS{lykta vel} \dicExampleCS{vonět dobře};  \dicPhraseIS{lykta af e‑u} \dicDirectTranslationCS{páchnout (čím), být cítit (čím), vonět (čím)} \dicExampleIS{lykta af ilmvatni} \dicExampleCS{vonět parfémem}  \textbf{2.} \dicSynonym{þefa} \dicDirectTranslationCS{přivonět (si), přičichnout (si)};  \dicPhraseIS{lykta af e‑u} \dicSynonym{þefa} \dicDirectTranslationCS{přivonět (si) k~(čemu)} \dicExampleIS{lykta af blóminu} \dicExampleCS{přivonět si ke květině}
\dicEntry[lyktandi] \dicTerm{lykt··andi} \dicIPA{{l}{\textsci}{x}{\textsubring{d}}{a}{n}{\textsubring{d}}{\textsci}} \dicPos{adj}[13] \dicFlx{indecl}[1] \dicDirectTranslationCS{vonící, šířící vůni};  \dicPhraseIS{e‑að er illa lyktandi} \dicDirectTranslationCS{(co) smrdí\,/\addthin páchne}
\dicEntry[lykti] \dicTerm{lykti} \dicIPA{{l}{\textsci}{x}{\textsubring{d}}{\textsci}} \dicPos{v} \dicFlx{con pf sg 1 pers} \dicLink{lykja}
\dicEntry[lyktnæmur] \dicTerm{lykt··næmur} \dicIPA{{l}{\textsci}{x}{\textsubring{d}}{n}{a}{i}{m}{\textscy}{\textsubring{r}}} \dicPos{adj}[1]\dicFlx{}[-1] \dicDirectTranslationCS{mající dobrý čich}
\dicEntry[lymska] \dicTerm{lymsk|a} \dicIPA{{l}{\textsci}{m}{s}{\r{g}}{a}} \dicPos{f}[1] \dicFlx{(‑u)}[5] \dicSynonym*{lævísi} \dicDirectTranslationCS{záludnost, zákeřnost, proradnost}
\dicEntry[lymskur] \dicTerm{lymskur} \dicIPA{{l}{\textsci}{m}{s}{\r{g}}{\textscy}{\textsubring{r}}} \dicPos{adj}[1]\dicFlx{}[-1] \dicSynonym{lævís} \dicDirectTranslationCS{záludný, zákeřný, proradný}
\dicEntry[lynda] \dicTerm{lyn|da} \dicIPA{{l}{\textsci}{n}{\textsubring{d}}{a}} \dicPos{v}[2] \dicFlx{(‑ti, ‑t)}[44] \dicPhraseIS{e‑m lyndir við e‑n} \dicFlx{impers} \dicDirectTranslationCS{(kdo) si rozumí s~(kým), (kdo) má pochopení pro (koho)} \dicExampleIS{Mér lyndir ekki við hana.} \dicExampleCS{Nerozumím si s~ní.};  \dicPhraseIS{e‑jum lyndir} \dicFlx{impers} \dicDirectTranslationCS{(kdo) si rozumí, (kdo) má pochopení jeden pro druhého};  \dicPhraseIS{láta sér e‑ð lynda} \dicDirectTranslationCS{spokojit se s~(čím)}
\dicEntry[lyndi] \dicTerm{lyndi} \dicIPA{{l}{\textsci}{n}{\textsubring{d}}{\textsci}} \dicPos{n}[2] \dicFlx{(‑s)}[20] \dicSynonym{skap} \dicDirectTranslationCS{nálada, rozpoložení};  \dicPhraseIS{e‑að leikur í lyndi} \dicDirectTranslationCS{(co) klape, (co) hraje do karet};  \dicPhraseIS{e‑m leikur allt í lyndi} \dicFlx{impers} \dicDirectTranslationCS{(komu) vše vychází, (komu) vše hraje do noty}
\dicEntry[lyndiseinkunn] \dicTerm{lyndis··ein·kunn} \dicIPA{{l}{\textsci}{n}{\textsubring{d}}{\textsci}{s}{ei}{\r{\ng}}{\r{g}}{\textscy}{\textsubring{n}}} \dicPos{f}[7] \dicFlx{(‑ar, ‑ir)}[1] \dicSynonym*{skaphöfn} \dicDirectTranslationCS{charakter}
\dicEntry[lyng] \dicTerm{lyng} \dicIPA{{l}{i}{\ng}{\r{g}}} \dicPos{n}[2] \dicFlx{(‑s)}[2] \dicFieldCat{bot.} \dicDirectTranslationCS{vřesovec} \textit{(l.~{\textLA{Erica}})}  \dicsymPhoto\ 
\dicFigure{74636.jpg}{Lyng}{Lyng - Kesl Michael, Biolib, Copyright/CC-BY-NC}
\dicEntry[lyst] \dicTerm{lyst} \dicIPA{{l}{\textsci}{s}{\textsubring{d}}} \dicPos{f}[7] \dicFlx{(‑ar, ‑ir)}[1] \textbf{1.} \dicSynonym{matarlyst} \dicDirectTranslationCS{chuť (k~jídlu), apetit};  \dicPhraseIS{hafa lyst á e‑u} \dicDirectTranslationCS{mít chuť na (co)} \dicExampleIS{hafa góða lyst á matnum} \dicExampleCS{mít dobrý apetit}  \textbf{2.} \dicSynonym{löngun} \dicDirectTranslationCS{chuť, přání};  \dicPhraseIS{af hjartans lyst} \dicFlx{adv} \dicDirectTranslationCS{z~celého srdce}
\dicEntry[lystarlaus] \dicTerm{lystar··laus} \dicIPA{{l}{\textsci}{s}{\textsubring{d}}{a}{r}{l}{\oe i}{s}} \dicPos{adj}[5]\dicFlx{}[-1] \dicDirectTranslationCS{(jsoucí) bez apetitu, nemající chuť (na jídlo)}
\dicEntry[lystarleysi] \dicTerm{lystar··leysi} \dicIPA{{l}{\textsci}{s}{\textsubring{d}}{a}{r}{l}{ei}{s}{\textsci}} \dicPos{n}[2] \dicFlx{(‑s)}[20] \dicDirectTranslationCS{nechutenství}
\dicEntry[lystarstol] \dicTerm{lystar··stol} \dicIPA{{l}{\textsci}{s}{\textsubring{d}}{a}{\textsubring{r}}{s}{\textsubring{d}}{\textopeno}{\textsubring{l}}} \dicPos{n}[2] \dicFlx{(‑s)}[2] \dicFieldCat{med.} \dicDirectTranslationCS{(mentální) anorexie}
\dicEntry[lysti] \dicTerm{lysti} \dicIPA{{l}{\textsci}{s}{\textsubring{d}}{\textsci}} \dicPos{v} \dicFlx{con pf sg 1 pers} \dicLink{ljósta}
\dicEntry[lystibátur] \dicTerm{lysti··bát|ur} \dicIPA{{l}{\textsci}{s}{\textsubring{d}}{\textsci}{\textsubring{b}}{au}{\textsubring{d}}{\textscy}{\textsubring{r}}} \dicPos{m}[6] \dicFlx{(‑s, ‑ar)}[22] \dicFieldCat{nám.} \dicDirectTranslationCS{jachta}
\dicEntry[lystigarður] \dicTerm{lysti··garð|ur} \dicIPA{{l}{\textsci}{s}{\textsubring{d}}{\textsci}{\r{g}}{a}{r}{ð}{\textscy}{\textsubring{r}}} \dicPos{m}[6] \dicFlx{(‑s, ‑ar)}[5] \dicSynonym{skemmtigarður} \dicDirectTranslationCS{(zábavní) park}
\dicEntry[lystiskip] \dicTerm{lysti··skip} \dicIPA{{l}{\textsci}{s}{\textsubring{d}}{\textsci}{s}{\r{\textObardotlessj}}{\textsci}{\textsubring{b}}} \dicPos{n}[2] \dicFlx{(‑s, ‑)}[5] \dicDirectTranslationCS{jachta, výletní loď}
\dicEntry[lystugur] \dicTerm{lyst··ugur} \dicIPA{{l}{\textsci}{s}{\textsubring{d}}{\textscy}{\textbabygamma}{\textscy}{\textsubring{r}}} \dicPos{adj}[1]\dicFlx{}[-8] \textbf{1.} \dicSynonym{sólginn} \dicDirectTranslationCS{mající dobrý apetit} \dicExampleIS{vera lystugur á matinn} \dicExampleCS{mít dobrý apetit}  \textbf{2.} \dicSynonym*{sem vekur lyst} \dicDirectTranslationCS{chutný, vábný} \dicIndirectTranslationCS{(o~jídle)}
\dicEntry[lyti] \dicTerm{lyti} \dicIPA{{l}{\textsci}{\textlengthmark}{\textsubring{d}}{\textsci}} \dicPos{v} \dicFlx{con pf sg 1 pers} \dicLink{lúta}
\dicEntry[lýðfrjáls] \dicTerm{lýð··frjáls} \dicIPA{{l}{i}{ð}{f}{r}{j}{au}{l}{s}} \dicPos{adj}[5]\dicFlx{}[-1] \dicDirectTranslationCS{demokratický, svobodný} \dicExampleIS{lýðfrjálst land} \dicExampleCS{demokratická země}
\dicEntry[lýðháskóli] \dicTerm{lýð··há·skól|i} \dicIPA{{l}{i}{\textlengthmark}{\texttheta}{h}{au}{s}{\r{g}}{ou}{l}{\textsci}} \dicPos{m}[1] \dicFlx{(‑a, ‑ar)}[1] \dicFieldCat{škol.} \dicDirectTranslationCS{veřejná vysoká škola}
\dicEntry[lýðhylli] \dicTerm{lýð··hyll|i} \dicIPA{{l}{i}{\textlengthmark}{\texttheta}{h}{\textsci}{\textsubring{d}}{l}{\textsci}} \dicPos{f}[3] \dicFlx{(‑i)}[3] \dicDirectTranslationCS{přízeň lidu, popularita}
\dicEntry[lýði] \dicTerm{lýði} \dicIPA{{l}{i}{\textlengthmark}{ð}{\textsci}} \dicPos{v} \dicFlx{con pf sg 1 pers} \dicLink{lýja}
\dicEntry[lýðréttindi] \dicTerm{lýð··rétt·indi} \dicIPA{{l}{i}{ð}{r}{j}{\textepsilon}{h}{\textsubring{d}}{\textsci}{n}{\textsubring{d}}{\textsci}} \dicPos{n}[2] \dicFlx{pl}[19] \dicDirectTranslationCS{občanská práva}
\dicEntry[lýðræði] \dicTerm{lýð··ræði} \dicIPA{{l}{i}{ð}{r}{a}{i}{ð}{\textsci}} \dicPos{n}[2] \dicFlx{(‑s)}[20] \dicDirectTranslationCS{demokracie}
\dicEntry[lýðræðislegur] \dicTerm{lýð·ræðis··legur} \dicIPA{{l}{i}{ð}{r}{a}{i}{ð}{\textsci}{s}{l}{\textepsilon}{\textbabygamma}{\textscy}{\textsubring{r}}} \dicPos{adj}[1]\dicFlx{}[-8] \dicDirectTranslationCS{demokratický}
\dicEntry[lýðræðissinnaður] \dicTerm{lýð·ræðis··sinn·|aður} \dicIPA{{l}\-{i}\-{ð}\-{r}\-{a}\-{i}\-{ð}\-{\textsci}\-{s}\-{\textsci}\-{n}\-{a}\-{ð}\-{\textscy}\-{\textsubring{r}}\-} \dicPos{adj}[3] \dicFlx{(f ‑uð)}[3] \dicDirectTranslationCS{(pro)demokratický}
\dicEntry[lýðræðissinni] \dicTerm{lýð·ræðis··sinn|i} \dicIPA{{l}{i}{ð}{r}{a}{i}{ð}{\textsci}{s}{\textsci}{n}{\textsci}} \dicPos{m}[1] \dicFlx{(‑a, ‑ar)}[1] \dicDirectTranslationCS{demokrat(ka)}
\dicEntry[lýðskrum] \dicTerm{lýð··skrum} \dicIPA{{l}{i}{ð}{s}{\r{g}}{r}{\textscy}{\textsubring{m}}} \dicPos{n}[2] \dicFlx{(‑s)}[2] \dicDirectTranslationCS{demagogie, populismus}
\dicEntry[lýðskrumari] \dicTerm{lýð··skrum·ar|i} \dicIPA{{l}{i}{ð}{s}{\r{g}}{r}{\textscy}{m}{a}{r}{\textsci}} \dicPos{m}[1] \dicFlx{(‑a, ‑ar)}[13] \dicDirectTranslationCS{demagog, demagožka, populista, populistka}
\dicEntry[lýður] \dicTerm{lýð|ur} \dicIPA{{l}{i}{\textlengthmark}{ð}{\textscy}{\textsubring{r}}} \dicPos{m}[9] \dicFlx{(‑s, ‑ir)}[8] \textbf{1.} \dicSynonym{almenningur} \dicDirectTranslationCS{lid(é), občané}  \textbf{2.} \dicSynonym{skríll} \dicDirectTranslationCS{lůza, holota};  \dicPhraseIS{e‑að er enn við lýði} \dicDirectTranslationCS{(co) stále existuje, (co) se stále vyskytuje (zvyk ap.)}
\dicEntry[lýðveldi] \dicTerm{lýð··veldi} \dicIPA{{l}{i}{ð}{v}{\textepsilon}{l}{\textsubring{d}}{\textsci}} \dicPos{n}[2] \dicFlx{(‑s, ‑)}[14] \dicDirectTranslationCS{republika}
\dicEntry[lýg] \dicTerm{lýg} \dicIPA{{l}{i}{\textlengthmark}{x}} \dicPos{v} \dicFlx{ind praes sg 1 pers} \dicLink{ljúga}
\dicEntry[lýi] \dicTerm{lýi} \dicIPA{{l}{i}{j}{\textsci}} \dicPos{v} \dicFlx{ind praes sg 1 pers} \dicLink{lýja}
\dicEntry[lýja] \dicTerm{lýja} \dicIPA{{l}{i}{j}{\textlengthmark}{a}} \dicPos{v}[4] \dicFlx{(lýi, lúði, lúðum, lýði, lúið\,/\addthin lúð)}[54] \dicFlx{acc} \textbf{1.} \dicDirectTranslationCS{kovat, kout} \dicExampleIS{lýja járn} \dicExampleCS{kovat železo}  \textbf{2.} \dicSynonym{þreyta\smash{\textsuperscript{2}}} \dicDirectTranslationCS{unavit, utahat} \dicExampleIS{Þetta lýir mann.} \dicExampleCS{To člověka unaví.};  \dicIdiom{lýjast}{ \dicPhraseIS{lýjast}} \dicFlx{refl} \dicSynonym*{þreytast} \dicDirectTranslationCS{vyčerpat se, unavit se} \dicExampleIS{Strákurinn er farinn að lýjast.} \dicExampleCS{Kluk se začíná unavovat.}
\dicEntry[lýjandi] \dicTerm{lýj··andi} \dicIPA{{l}{i}{j}{\textlengthmark}{a}{n}{\textsubring{d}}{\textsci}} \dicPos{adj}[13] \dicFlx{indecl}[1] \dicDirectTranslationCS{vyčerpávající, vysilující} \dicExampleIS{lýjandi starf} \dicExampleCS{vyčerpávající práce}
\dicEntry[lýk] \dicTerm{lýk} \dicIPA{{l}{i}{\textlengthmark}{\r{g}}} \dicPos{v} \dicFlx{ind praes sg 1 pers} \dicLink{ljúka}
\dicEntry[lýrik] \dicTerm{lýrik} \dicIPA{{l}{i}{\textlengthmark}{r}{\textsci}{\r{g}}} \dicPos{f}[10] \dicFlx{(‑ur)}[2] \dicFieldCat{lit.} \dicSynonym{ljóðræna} \dicDirectTranslationCS{lyrika}
\dicEntry[lýrískur] \dicTerm{lýrískur} \dicIPA{{l}{i}{\textlengthmark}{r}{i}{s}{\r{g}}{\textscy}{\textsubring{r}}} \dicPos{adj}[1]\dicFlx{}[-1] \dicFieldCat{lit.} \dicSynonym{ljóðrænn} \dicDirectTranslationCS{lyrický}
\dicEntry[lýs] \dicTerm{lýs} \dicIPA{{l}{i}{\textlengthmark}{s}} \dicPos{f} \dicFlx{pl nom} \dicLink{lús}
\dicEntry[lýsa] \dicTerm{lýs|a\smash{\textsuperscript{1}}} \dicIPA{{l}{i}{\textlengthmark}{s}{a}} \dicPos{f}[1] \dicFlx{(‑u, ‑ur)}[7] \textbf{1.} \dicSynonym{ljómi} \dicDirectTranslationCS{světlo, paprsek} \dicExampleIS{fyrsta lýsa morgunsins} \dicExampleCS{první ranní paprsek}  \textbf{2.} \dicFieldCat{zool.} \dicDirectTranslationCS{treska bezvousá} \textit{(l.~{\textLA{Merlangius merlangus}})}  \dicsymPhoto\ 
\dicFigure{ds_image_lysa_1_2.jpg}{Lýsa}{Lýsa - Georges Jansoone (JoJan), CC BY 3.0}
\dicEntry[lýsa] \dicTerm{lýs|a\smash{\textsuperscript{2}}} \dicsymFrequent\  \dicIPA{{l}{i}{\textlengthmark}{s}{a}} \dicPos{v}[2] \dicFlx{(‑ti, ‑t)}[67] \dicFlx{dat} \textbf{1.} \dicSynonym*{bera birtu} \dicDirectTranslationCS{osvětlit, osvětlovat, osvítit, ozářit, ozařovat} \dicExampleIS{Tunglið lýsir mér.} \dicExampleCS{Měsíc mě osvětluje.};  \dicPhraseIS{það lýsir (af degi)} \dicFlx{impers} \dicSynonym*{það birtir} \dicDirectTranslationCS{rozednívá se}  \textbf{2.} \dicSynonym*{lýsa mynd} \dicDirectTranslationCS{vyvolat, vyvolávat (film ap.)} \dicExampleIS{lýsa myndina} \dicExampleCS{vyvolat film}  \textbf{3.} \dicSynonym{sýna} \dicDirectTranslationCS{projevit, projevovat, vyjádřit, vyjadřovat}  \textbf{4.} \dicSynonym*{gefa lýsingu} \dicDirectTranslationCS{popsat, popisovat, (vy)líčit, vykreslit, vykreslovat} \dicExampleIS{lýsa reynslu sinni} \dicExampleCS{popsat svou zkušenost}  \textbf{5.} \dicSynonym{tilkynna} \dicDirectTranslationCS{prohlásit, prohlašovat, deklarovat};  \dicIdiom{lýsa}[eftir]{ \dicPhraseIS{lýsa eftir e‑m}} \dicDirectTranslationCS{vyhlásit pátrání po (kom)};  \dicIdiom{lýsa}[upp]{ \dicPhraseIS{lýsa upp e‑ð}} \dicDirectTranslationCS{osvětlit (co) (místnost ap.)};  \dicIdiom{lýsa}[yfir]{ \dicPhraseIS{lýsa yfir e‑u}} \dicDirectTranslationCS{ohlásit (co), oznámit (co)}
\dicEntry[lýsi] \dicTerm{lýsi} \dicIPA{{l}{i}{\textlengthmark}{s}{\textsci}} \dicPos{n}[2] \dicFlx{(‑s, ‑)}[14] \dicDirectTranslationCS{rybí tuk} \dicIndirectTranslationCS{(medicína)}
\dicEntry[lýsing] \dicTerm{lýs··ing} \dicsymFrequent\  \dicIPA{{l}{i}{\textlengthmark}{s}{i}{\ng}{\r{g}}} \dicPos{f}[4] \dicFlx{(‑ar, ‑ar)}[5] \textbf{1.} \dicDirectTranslationCS{osvětlení, nasvícení, iluminace} \dicExampleIS{lýsing í eldhúsinu} \dicExampleCS{osvětlení v~kuchyni}  \textbf{2.} \dicSynonym{frásögn} \dicDirectTranslationCS{popis, vylíčení, vykreslení} \dicExampleIS{lýsing á ferðinni} \dicExampleCS{vylíčení cesty}  \textbf{3.} \dicDirectTranslationCS{osvit, expozice}  \textbf{4.} \dicSynonym{skreyting} \dicDirectTranslationCS{ilustrace, ilustrování (rukopisu, knihy ap.)}
\dicEntry[lýsingarháttur] \dicTerm{lýsingar··|háttur} \dicIPA{{l}{i}{\textlengthmark}{s}{i}{\ng}{\r{g}}{a}{\textsubring{r}}{h}{au}{h}{\textsubring{d}}{\textscy}{\textsubring{r}}} \dicPos{m}[12] \dicFlx{(‑háttar, ‑hættir)}[7] \dicFieldCat{jaz.} \dicDirectTranslationCS{příčestí};  \dicPhraseIS{lýsingarháttur nútíðar} \dicFieldCat{jaz.} \dicDirectTranslationCS{příčestí přítomné};  \dicPhraseIS{lýsingarháttur þátíðar} \dicFieldCat{jaz.} \dicDirectTranslationCS{příčestí minulé}
\dicEntry[lýsingarorð] \dicTerm{lýsingar··orð} \dicIPA{{l}{i}{\textlengthmark}{s}{i}{\ng}{\r{g}}{a}{r}{\textopeno}{r}{\texttheta}} \dicPos{n}[2] \dicFlx{(‑s, ‑)}[5] \dicFieldCat{jaz.} \dicDirectTranslationCS{přídavné jméno, adjektivum} \textit{(l.~{\textLA{nomen adiectivum}})}
\dicEntry[lýst] \dicTerm{lýst} \dicIPA{{l}{i}{s}{\textsubring{d}}} \dicPos{v} \dicFlx{ind praes sg 1 pers} \dicLink{ljósta}
\dicEntry[lýt] \dicTerm{lýt} \dicIPA{{l}{i}{\textlengthmark}{\textsubring{d}}} \dicPos{v} \dicFlx{ind praes sg 1 pers} \dicLink{lúta}
\dicEntry[lýtalaus] \textls[15]{\dicTerm{lýta··laus} \dicIPA{{l}{i}{\textlengthmark}{\textsubring{d}}{a}{l}{\oe i}{s}} \dicPos{adj}[5]\dicFlx{}[-1] \dicSynonym{gallalaus} \dicDirectTranslationCS{bezchybný, (jsoucí) bez poskvrnky}}
\dicEntry[lýtalækningar] \dicTerm{lýta··lækn·ingar} \dicIPA{{l}{i}{\textlengthmark}{\textsubring{d}}{a}{l}{a}{i}{h}{\r{g}}{n}{i}{\ng}{\r{g}}{a}{\textsubring{r}}} \dicPos{f}[4] \dicFlx{pl}[6] \dicFieldCat{med.} \dicDirectTranslationCS{plastická chirurgie}
\dicEntry[lýtalæknir] \dicTerm{lýta··lækn|ir} \dicIPA{{l}{i}{\textlengthmark}{\textsubring{d}}{a}{l}{a}{i}{h}{\r{g}}{n}{\textsci}{\textsubring{r}}} \dicPos{m}[7] \dicFlx{(‑is, ‑ar)}[1] \dicDirectTranslationCS{plastický chirurg, plastická chirurgyně\,/\addthin chiruržka}
\dicEntry[lýti] \dicTerm{lýti} \dicIPA{{l}{i}{\textlengthmark}{\textsubring{d}}{\textsci}} \dicPos{n}[2] \dicFlx{(‑s, ‑)}[14] \textbf{1.} \dicSynonym{ljóður} \dicDirectTranslationCS{poskvrna, kaz, vada} \dicExampleIS{lýti á bókinni} \dicExampleCS{vada na knížce}  \textbf{2.} \dicSynonym{ljótleiki} \dicDirectTranslationCS{šerednost, ošklivost (vzhled člověka ap.)}
\dicEntry[læða] \dicTerm{læð|a\smash{\textsuperscript{1}}} \dicIPA{{l}{a}{i}{\textlengthmark}{ð}{a}} \dicPos{f}[1] \dicFlx{(‑u, ‑ur)}[19] \dicDirectTranslationCS{kočka} \dicIndirectTranslationCS{(samice)}
\dicEntry[læða] \dicTerm{læ|ða\smash{\textsuperscript{2}}} \dicsymFrequent\  \dicIPA{{l}{a}{i}{\textlengthmark}{ð}{a}} \dicPos{v}[2] \dicFlx{(‑ddi, ‑tt)}[172] \dicFlx{dat} \dicPhraseIS{læða e‑u að e‑m} \dicDirectTranslationCS{propašovat (komu co), vstrčit (komu co) skrytě} \dicExampleIS{læða þessari frétt að henni} \dicExampleCS{nenápadně jí sdělit zprávu};  \dicIdiom{læðast}{ \dicPhraseIS{læðast}} \dicFlx{refl} \dicDirectTranslationCS{(v)plížit se, plazit se, krást se} \dicExampleIS{læðast burt} \dicExampleCS{odplížit se pryč}
\dicEntry[læðupokast] \dicTerm{læðu··pok|ast} \dicIPA{{l}{a}{i}{\textlengthmark}{ð}{\textscy}{p\smash{\textsuperscript{h}}}{\textopeno}{\r{g}}{a}{s}{\textsubring{d}}} \dicPos{v}[1] \dicFlx{(‑aðist)}[91] \dicFlx{refl} \dicDirectTranslationCS{(při)krást se, plížit se}
\dicEntry[lægð] \dicTerm{lægð} \dicIPA{{l}{a}{i}{\textbabygamma}{\texttheta}} \dicPos{f}[7] \dicFlx{(‑ar, ‑ir)}[1] \textbf{1.} \dicSynonym{dæld} \dicDirectTranslationCS{prohlubeň, proláklina} \dicExampleIS{Það myndast djúp lægð í götuna.} \dicExampleCS{Na ulici se tvoří hluboká prohlubeň.}  \textbf{2.} \dicFieldCat{meteo.} \dicSynonym{lágþrýstisvæði} \dicDirectTranslationCS{tlaková níže} \dicAntonym{hægð}  \textbf{3.} \dicSynonym{deyfð} \dicDirectTranslationCS{útlum};  \dicPhraseIS{vera í lægð} \dicDirectTranslationCS{být v~útlumu}
\dicEntry[lægi] \dicTerm{lægi\smash{\textsuperscript{1}}} \dicIPA{{l}{a}{i}{j}{\textlengthmark}{\textsci}} \dicPos{n}[2] \dicFlx{(‑s, ‑)}[16] \dicSynonym{skipalægi} \dicDirectTranslationCS{kotviště}
\dicEntry[lægi] \dicTerm{lægi\smash{\textsuperscript{2}}} \dicIPA{{l}{a}{i}{j}{\textlengthmark}{\textsci}} \dicPos{v} \dicFlx{con pf sg 1 pers} \dicLink{liggja}
\dicEntry[lægja] \dicTerm{læg|ja} \dicIPA{{l}{a}{i}{j}{\textlengthmark}{a}} \dicPos{v}[2] \dicFlx{(‑ði, ‑t)}[90] \dicFlx{acc} \dicSynonym{auðmýkja} \dicDirectTranslationCS{zklidnit, zklidňovat, zmírnit, zmírňovat} \dicExampleIS{lægja æsing} \dicExampleCS{zklidnit vztek};  \dicPhraseIS{lægja rostann í e‑m} \dicDirectTranslationCS{zklidnit (koho)};  \dicPhraseIS{e‑ð lægir} \dicFlx{impers} {\textbf{a.}} \dicDirectTranslationCS{(co) klesá, (co) zapadá (slunce ap.)} \dicExampleIS{Sólina lægir.} \dicExampleCS{Slunce klesá.};  {\textbf{b.}} \dicDirectTranslationCS{(co) ustává (o~větru ap.), (co) se uklidňuje} \dicExampleIS{Veðrið lægir.} \dicExampleCS{Počasí se uklidňuje.}
\dicEntry[lægra] \dicTerm{lægra} \dicIPA{{l}{a}{i}{\textbabygamma}{r}{a}} \dicPos{adv} \dicFlx{comp} \dicLink{lágt}
\dicEntry[lægri] \dicTerm{lægri} \dicIPA{{l}{a}{i}{\textbabygamma}{r}{\textsci}} \dicPos{adj} \dicFlx{comp m} \dicLink{lágur}
\dicEntry[lægst] \dicTerm{lægst} \dicIPA{{l}{a}{i}{x}{s}{\textsubring{d}}} \dicPos{adv} \dicFlx{sup} \dicLink{lágt}
\dicEntry[lægstur] \dicTerm{lægstur} \dicIPA{{l}{a}{i}{x}{s}{\textsubring{d}}{\textscy}{\textsubring{r}}} \dicPos{adj} \dicFlx{m sg nom sup} \dicLink{lágur}
\dicEntry[læki] \dicTerm{læki} \dicIPA{{l}{a}{i}{\textlengthmark}{\r{\textObardotlessj}}{\textsci}} \dicPos{v} \dicFlx{con pf sg 1 pers} \dicLink{leka}
\dicEntry[lækka] \dicTerm{lækk|a} \dicsymFrequent\  \dicIPA{{l}{a}{i}{h}{\r{g}}{a}} \dicPos{v}[1] \dicFlx{(‑aði)}[48] \dicFlx{acc} \textbf{1.} \dicSynonym{minnka} \dicDirectTranslationCS{snížit, snižovat, (z)redukovat} \dicExampleIS{lækka skatta} \dicExampleCS{snížit daně};  \dicPhraseIS{lækka í e‑u} \dicDirectTranslationCS{zeslabit (co), ztišit (co), ztlumit (co), zmírnit (co)} \dicExampleIS{lækka í útvarpinu} \dicExampleCS{ztišit rádio};  \dicPhraseIS{lækka röddina\,/\addthin róminn} \dicDirectTranslationCS{snížit hlas}  \textbf{2.} \dicDirectTranslationCS{klesnout, klesat, sestoupit, sestupovat, snášet se} \dicExampleIS{lækka flugið} \dicExampleCS{klesat s~letadlem};  \dicPhraseIS{e‑að lækkar} \dicDirectTranslationCS{(co) klesá, (co) ustupuje (voda ap.)};  \dicPhraseIS{það lækkar í e‑u} \dicFlx{impers} \dicDirectTranslationCS{(co) klesá, (co) opadává} \dicExampleIS{Það lækkar í ánni.} \dicExampleCS{Voda v~řece opadává.}
\dicEntry[lækkun] \dicTerm{lækk|un} \dicIPA{{l}{a}{i}{h}{\r{g}}{\textscy}{\textsubring{n}}} \dicPos{f}[7] \dicFlx{(‑unar)}[9] \textbf{1.} \dicSynonym*{það að lækka} \dicDirectTranslationCS{snížení, snižování, redukce, redukování} \dicExampleIS{lækkun skatta} \dicExampleCS{snížení daní}  \textbf{2.} \dicSynonym*{það að e‑að lækkar} \dicDirectTranslationCS{klesání, opadání, ubývání} \dicExampleIS{lækkun á yfirborði sjávar} \dicExampleCS{klesání hladiny moře}
\dicEntry[lækna] \dicTerm{lækn|a} \dicsymFrequent\  \dicIPA{{l}{a}{i}{h}{\r{g}}{n}{a}} \dicPos{v}[1] \dicFlx{(‑aði)}[1] \dicFlx{acc} \dicDirectTranslationCS{(vy)léčit};  \dicPhraseIS{lækna e‑n af e‑u} \dicDirectTranslationCS{léčit (koho) z~(čeho)} \dicExampleIS{lækna e‑n af berklum} \dicExampleCS{léčit (koho) z~tuberkulózy}
\dicEntry[læknadeild] \dicTerm{lækna··deild} \dicIPA{{l}{a}{i}{h}{\r{g}}{n}{a}{\textsubring{d}}{ei}{l}{\textsubring{d}}} \dicPos{f}[7] \dicFlx{(‑ar, ‑ir)}[1] \dicFieldCat{škol.} \dicDirectTranslationCS{lékařská fakulta}
\dicEntry[læknandi] \dicTerm{lækn··andi} \dicIPA{{l}{a}{i}{h}{\r{g}}{n}{a}{n}{\textsubring{d}}{\textsci}} \dicPos{adj}[13] \dicFlx{indecl}[1] \dicDirectTranslationCS{léčivý, léčebný}
\dicEntry[læknanemi] \dicTerm{lækna··nem|i} \dicIPA{{l}{a}{i}{h}{\r{g}}{n}{a}{n}{\textepsilon}{m}{\textsci}} \dicPos{m}[1] \dicFlx{(‑a, ‑ar)}[1] \dicDirectTranslationCS{student(ka) medicíny\,/\addthin lékařství}
\dicEntry[læknanlegur] \dicTerm{læknan··legur} \dicIPA{{l}{a}{i}{h}{\r{g}}{n}{a}{n}{l}{\textepsilon}{\textbabygamma}{\textscy}{\textsubring{r}}} \dicPos{adj}[1]\dicFlx{}[-8] \dicDirectTranslationCS{(vy)léčitelný}
\dicEntry[læknastofa] \dicTerm{lækna··stof|a} \dicIPA{{l}{a}{i}{h}{\r{g}}{n}{a}{s}{\textsubring{d}}{\textopeno}{v}{a}} \dicPos{f}[1] \dicFlx{(‑u, ‑ur)}[7] \dicDirectTranslationCS{ordinace}
\dicEntry[lækning] \dicTerm{lækn··ing} \dicsymFrequent\  \dicIPA{{l}{a}{i}{h}{\r{g}}{n}{i}{\ng}{\r{g}}} \dicPos{f}[4] \dicFlx{(‑ar, ‑ar)}[5] \dicDirectTranslationCS{léčení, léčba};  \dicPhraseIS{lækning við e‑u} \dicDirectTranslationCS{léčba (čeho)} \dicExampleIS{lækning við sjúkdómi} \dicExampleCS{léčba nemoci}
\dicEntry[lækningastofa] \dicTerm{lækninga··stof|a} \dicIPA{{l}{a}{i}{h}{\r{g}}{n}{i}{\ng}{\r{g}}{a}{s}{\textsubring{d}}{\textopeno}{v}{a}} \dicPos{f}[1] \dicFlx{(‑u, ‑ur)}[7] \dicDirectTranslationCS{ordinace}
\dicEntry[læknir] \dicTerm{lækn|ir} \dicsymFrequent\  \dicIPA{{l}{a}{i}{h}{\r{g}}{n}{\textsci}{\textsubring{r}}} \dicPos{m}[7] \dicFlx{(‑is, ‑ar)}[1] \dicDirectTranslationCS{doktor(ka), lékař(ka)} \dicExampleIS{panta tíma hjá lækni} \dicExampleCS{objednat se k~lékaři}
\dicEntry[læknisfræði] \dicTerm{læknis··fræð|i} \dicIPA{{l}{a}{i}{h}{\r{g}}{n}{\textsci}{s}{f}{r}{a}{i}{ð}{\textsci}} \dicPos{f}[3] \dicFlx{(‑i)}[3] \dicDirectTranslationCS{lékařství, medicína} \dicIndirectTranslationCS{(nauka)}
\dicEntry[læknisskoðun] \dicTerm{læknis··skoð|un} \dicIPA{{l}{a}{i}{h}{\r{g}}{n}{\textsci}{s}{\r{g}}{\textopeno}{ð}{\textscy}{\textsubring{n}}} \dicPos{f}[7] \dicFlx{(‑unar)}[9] \dicDirectTranslationCS{zdravotní\,/\addthin lékařská prohlídka}
\dicEntry[læknisvottorð] \dicTerm{læknis··vott·orð} \dicIPA{{l}{a}{i}{h}{\r{g}}{n}{\textsci}{s}{v}{\textopeno}{h}{\textsubring{d}}{\textopeno}{r}{\texttheta}} \dicPos{n}[2] \dicFlx{(‑s, ‑)}[5] \dicDirectTranslationCS{lékařská zpráva}
\dicEntry[læknisþjónusta] \dicTerm{læknis··þjón·ust|a} \dicIPA{{l}{a}{i}{h}{\r{g}}{n}{\textsci}{s}{\texttheta}{j}{ou}{n}{\textscy}{s}{\textsubring{d}}{a}} \dicPos{f}[1] \dicFlx{(‑u)}[5] \dicDirectTranslationCS{zdravot\-ní\,/\addthin lékařská služba}
\dicEntry[lækur] \dicTerm{læk|ur} \dicsymFrequent\  \dicIPA{{l}{a}{i}{\textlengthmark}{\r{g}}{\textscy}{\textsubring{r}}} \dicPos{m}[10] \dicFlx{(‑jar, ‑ir)}[23] \textbf{1.} \dicDirectTranslationCS{potok, říčka} \dicExampleIS{ár og lækir} \dicExampleCS{řeky a~potoky}  \textbf{2.} \dicFieldCat{anat.} \dicSynonym{naflastrengur} \dicDirectTranslationCS{pupeční šňůra}
\dicEntry[læmingi] \dicTerm{læm··ing|i} \dicIPA{{l}{a}{i}{\textlengthmark}{m}{i}{\textltailn}{\r{\textObardotlessj}}{\textsci}} \dicPos{m}[1] \dicFlx{(‑ja, ‑jar)}[14] \dicFieldCat{zool.} \dicDirectTranslationCS{lumík norský} \textit{(l.~{\textLA{Lemmus lemmus}})}  \dicsymPhoto\ 
\dicFigure{ds_image_laemingi_0_1.jpg}{Læmingi}{Læmingi - Argus fin, PD}
\dicEntry[læra] \dicTerm{lær|a} \dicsymFrequent\  \dicIPA{{l}{a}{i}{\textlengthmark}{r}{a}} \dicPos{v}[2] \dicFlx{(‑ði, ‑t)}[103] \dicFlx{acc} \dicDirectTranslationCS{(na)učit se, studovat} \dicExampleIS{læra íslensku} \dicExampleCS{učit se islandštinu};  \dicPhraseIS{læra e‑ð utanbókar} \dicDirectTranslationCS{učit se (co) zpaměti, memorovat (co)} \dicExampleIS{læra utanbókar ljóð} \dicExampleCS{učit se básničku nazpaměť};  \dicIdiom{læra}[af]{ \dicPhraseIS{læra af e‑m}} \dicDirectTranslationCS{učit se od (koho)};  \dicIdiom{læra}[á]{ \dicPhraseIS{læra á e‑ð}} \dicDirectTranslationCS{učit se na (co) (na hudební nástroj ap.)} \dicExampleIS{Ég lærði á píanó.} \dicExampleCS{Učil jsem se na piano.};  \dicIdiom{læra}[fyrir]{ \dicPhraseIS{læra fyrir e‑ð}} \dicDirectTranslationCS{učit se na (co)} \dicExampleIS{læra fyrir próf} \dicExampleCS{učit se na zkoušku};  \dicIdiom{læra}[undir]{ \dicPhraseIS{læra undir e‑ð}} \dicDirectTranslationCS{učit se na (co) (na zkoušku ap.)} \dicExampleIS{læra undir skóla} \dicExampleCS{učit se do školy};  \dicIdiom{læra}[utan að]{ \dicPhraseIS{læra utan að}} \dicDirectTranslationCS{učit se zpaměti, memorovat};  \dicIdiom{lærast}{ \dicPhraseIS{e‑að lærist vel}} \dicFlx{refl} \dicDirectTranslationCS{(co) se učí dobře}; { \dicPhraseIS{e‑m lærist að (gera e‑ð)}} \dicFlx{refl impers} \dicDirectTranslationCS{(kdo) se učí (dělat (co)), (kdo) si osvojuje (dělat (co))};  \dicProverb\  \dicPhraseIS{Svo lengi lærir sem lifir.} \dicLangCat{přís.} \dicDirectTranslationCS{Člověk se učí po celý život.}
\dicEntry[lærbein] \dicTerm{lær··bein} \dicIPA{{l}{a}{i}{r}{\textsubring{b}}{ei}{\textsubring{n}}} \dicPos{n}[2] \dicFlx{(‑s, ‑)}[5] \dicFieldCat{anat.} \dicDirectTranslationCS{stehenní kost}
\dicEntry[lærdómur] \dicTerm{lær··dóm|ur} \dicIPA{{l}{a}{i}{r}{\textsubring{d}}{ou}{m}{\textscy}{\textsubring{r}}} \dicPos{m}[6] \dicFlx{(‑s, ‑ar)}[10] \dicSynonym{þekking} \dicDirectTranslationCS{učení, učenost, vzdělání};  \dicPhraseIS{draga lærdóma af e‑u} \dicDirectTranslationCS{poučit se z~(čeho)}
\dicEntry[lærður] \dicTerm{lærður} \dicsymFrequent\  \dicIPA{{l}{a}{i}{r}{ð}{\textscy}{\textsubring{r}}} \dicPos{adj}[2]\dicFlx{}[-1] \dicDirectTranslationCS{učený, vzdělaný} \dicExampleIS{mjög lærður maður} \dicExampleCS{velmi vzdělaný člověk}
\dicEntry[læri] \dicTerm{læri} \dicsymFrequent\  \dicIPA{{l}{a}{i}{\textlengthmark}{r}{\textsci}} \dicPos{n}[2] \dicFlx{(‑s, ‑)}[14] \textbf{1.} \dicFieldCat{anat.} \dicDirectTranslationCS{stehno}  \textbf{2.} \dicDirectTranslationCS{kýta} \dicExampleIS{læri af lambi} \dicExampleCS{jehněčí kýta}  \textbf{3.} \dicSynonym{nám} \dicDirectTranslationCS{učení} \dicIndirectTranslationCS{(učňovská léta)};  \dicPhraseIS{vera í læri hjá e‑m} \dicDirectTranslationCS{být u~(koho) v~učení}
\dicEntry[lærifaðir] \dicTerm{læri··|faðir} \dicIPA{{l}{a}{i}{\textlengthmark}{r}{\textsci}{f}{a}{ð}{\textsci}{\textsubring{r}}} \dicPos{m}[13] \dicFlx{(‑föður, ‑feður)}[3] \dicSynonym{kennari} \dicDirectTranslationCS{duchovní otec, mistr, guru}
\dicEntry[lærisveinn] \dicTerm{læri··svein|n} \dicIPA{{l}{a}{i}{\textlengthmark}{r}{\textsci}{s}{v}{ei}{\textsubring{d}}{\textsubring{n}}} \dicPos{m}[6] \dicFlx{(‑s, ‑ar)}[42] \dicSynonym{nemandi} \dicDirectTranslationCS{učeň(ka), učnice, učedník, učednice, žák(yně), žačka} \dicExampleIS{lærisveinar Jesú} \dicExampleCS{učedníci Ježíše}
\dicEntry[lærleggur] \dicTerm{lær··legg|ur} \dicIPA{{l}{a}{i}{r}{l}{\textepsilon}{\r{g}}{\textscy}{\textsubring{r}}} \dicPos{m}[9] \dicFlx{(‑jar\,/\addthin ‑s, ‑ir)}[26] \dicFieldCat{anat.} \dicDirectTranslationCS{stehenní kost}
\dicEntry[lærlingur] \dicTerm{lær··ling|ur} \dicIPA{{l}{a}{i}{r}{l}{i}{\ng}{\r{g}}{\textscy}{\textsubring{r}}} \dicPos{m}[6] \dicFlx{(‑s, ‑ar)}[8] \dicSynonym{nemandi} \dicDirectTranslationCS{učedník, učednice, učeň(ka), učnice}
\dicEntry[læs] \dicTerm{læs} \dicIPA{{l}{a}{i}{\textlengthmark}{s}} \dicPos{adj}[5]\dicFlx{}[-1] \dicDirectTranslationCS{gramotný, umějící číst}
\dicEntry[læsa] \dicTerm{læs|a} \dicsymFrequent\  \dicIPA{{l}{a}{i}{\textlengthmark}{s}{a}} \dicPos{v}[2] \dicFlx{(‑ti, ‑t)}[63] \dicFlx{dat\,/\addthin acc} \textbf{1.} \dicFlx{dat} \dicSynonym*{loka með lás} \dicDirectTranslationCS{(u)zamknout, (u)zamykat, (u)zavřít, (u)zavírat} \dicExampleIS{Dyrunum var læst.} \dicExampleCS{Dveře byly zamčeny.};  \dicPhraseIS{læsa að sér} \dicDirectTranslationCS{zamknout (za sebou) dveře};  \dicPhraseIS{læsa á eftir sér} \dicDirectTranslationCS{zavřít za sebou (dveře)};  \dicPhraseIS{læsa e‑n inni} \dicDirectTranslationCS{zavřít (koho), dát (koho) pod zámek (do vězení ap.)}  \textbf{2.} \dicFlx{dat} \dicSynonym{innsigla} \dicDirectTranslationCS{zalepit, zalepovat, zapečetit, zapečeťovat (obálku ap.)}  \textbf{3.} \dicFlx{acc} \dicSynonym*{breiðast út} \dicDirectTranslationCS{(roz)šířit se, rozšiřovat se};  \dicPhraseIS{læsa sig um e‑ð} \dicDirectTranslationCS{rozšířit se po (čem)}  \textbf{4.} \dicFlx{dat} \dicDirectTranslationCS{sevřít, svírat (v~drápech\,/\addthin zubech)};  \dicIdiom{læsast}{ \dicPhraseIS{læsast}} \dicFlx{refl} {\textbf{a.}} \dicDirectTranslationCS{zavřít se (klíče v~domě ap.)};  {\textbf{b.}} \dicDirectTranslationCS{sevřít, chopit se}
\dicEntry[læsi] \dicTerm{læsi\smash{\textsuperscript{1}}} \dicIPA{{l}{a}{i}{\textlengthmark}{s}{\textsci}} \dicPos{n}[2] \dicFlx{(‑s)}[20] \dicDirectTranslationCS{gramotnost, dovednost čtení}
\dicEntry[læsi] \dicTerm{læsi\smash{\textsuperscript{2}}} \dicIPA{{l}{a}{i}{\textlengthmark}{s}{\textsci}} \dicPos{v} \dicFlx{con pf sg 1 pers} \dicLink{lesa}
\dicEntry[læsilegur] \dicTerm{læsi··legur} \dicIPA{{l}{a}{i}{\textlengthmark}{s}{\textsci}{l}{\textepsilon}{\textbabygamma}{\textscy}{\textsubring{r}}} \dicPos{adj}[1]\dicFlx{}[-8] \textbf{1.} \dicDirectTranslationCS{čitelný} \dicExampleIS{læsileg skrift} \dicExampleCS{čitelné písmo}  \textbf{2.} \dicDirectTranslationCS{čtivý} \dicExampleIS{læsileg bók} \dicExampleCS{čtivá knížka}
\dicEntry[læsing] \dicTerm{læs··ing} \dicIPA{{l}{a}{i}{\textlengthmark}{s}{i}{\ng}{\r{g}}} \dicPos{f}[4] \dicFlx{(‑ar, ‑ar)}[5] \textbf{1.} \dicSynonym{lás} \dicDirectTranslationCS{zámek, zamykání (u~dveří ap.)} \dicExampleIS{læsing á dyrunum} \dicExampleCS{zámek u~dveří}  \textbf{2.} \dicSynonym{lokun} \dicDirectTranslationCS{(u)zamčení, (u)zamknutí}
\dicEntry[læstur] \dicTerm{læstur} \dicsymFrequent\  \dicIPA{{l}{a}{i}{s}{\textsubring{d}}{\textscy}{\textsubring{r}}} \dicPos{adj}[1]\dicFlx{}[-10] \dicSynonym*{aflokaður} \dicDirectTranslationCS{(u)zamčený, (u)zamknutý} \dicExampleIS{læstar dyr} \dicExampleCS{zamknuté dveře}
\dicEntry[læt] \dicTerm{læt} \dicIPA{{l}{a}{i}{\textlengthmark}{\textsubring{d}}} \dicPos{v} \dicFlx{ind praes sg 1 pers} \dicLink{láta}
\dicEntry[læti] \dicTerm{læti} \dicsymFrequent\  \dicIPA{{l}{a}{i}{\textlengthmark}{\textsubring{d}}{\textsci}} \dicPos{n}[2] \dicFlx{pl}[19] \dicDirectTranslationCS{rozruch, povyk} \dicExampleIS{vera með læti} \dicExampleCS{dělat povyk};  \dicPhraseIS{kunna sér ekki læti} \dicLangCat{přen.} \dicDirectTranslationCS{být štěstím bez sebe}
\dicEntry[lævirki] \dicTerm{læ··virk|i} \dicIPA{{l}{a}{i}{\textlengthmark}{v}{\textsci}{\textsubring{r}}{\r{\textObardotlessj}}{\textsci}} \dicPos{m}[1] \dicFlx{(‑ja, ‑jar)}[14] \dicFieldCat{zool.} \dicDirectTranslationCS{skřivan} \textit{(l.~{\textLA{Alauda}})}  \dicsymPhoto\ 
\dicFigure{ds_image_laevirki_0_2.jpg}{Lævirki}{Lævirki - Daniel Pettersson, CC BY-SA 2.5}
\dicEntry[lævís] \dicTerm{læ··vís} \dicIPA{{l}{a}{i}{\textlengthmark}{v}{i}{s}} \dicPos{adj}[5]\dicFlx{}[-1] \dicSynonym{slóttugur} \dicDirectTranslationCS{záludný, zákeřný}
\dicEntry[löðrungur] \dicTerm{löðr··ung|ur} \dicIPA{{l}{\oe}{ð}{r}{u}{\ng}{\r{g}}{\textscy}{\textsubring{r}}} \dicPos{m}[6] \dicFlx{(‑s, ‑ar)}[8] \dicSynonym{kinnhestur} \dicDirectTranslationCS{facka, políček} \dicExampleIS{gefa e‑m löðrungur} \dicExampleCS{dát (komu) facku}
\dicEntry[löður] \dicTerm{löður} \dicIPA{{l}{\oe}{\textlengthmark}{ð}{\textscy}{\textsubring{r}}} \dicPos{n}[2] \dicFlx{(‑s)}[28] \dicSynonym{froða} \dicDirectTranslationCS{(vodní) pěna}
\dicEntry[lög] \dicTerm{lög} \dicIPA{{l}{\oe}{\textlengthmark}{x}} \dicPos{n}[2] \dicFlx{pl}[9] \dicFieldCat{práv.} \dicDirectTranslationCS{zákon};  \dicPhraseIS{leiða e‑ð í lög} \dicDirectTranslationCS{uzákonit (co)};  \dicPhraseIS{setja lög} \dicDirectTranslationCS{vydat zákon}
\dicEntry[lögaldur] \dicTerm{lög··ald|ur} \dicIPA{{l}{\oe}{\textlengthmark}{\textbabygamma}{a}{l}{\textsubring{d}}{\textscy}{\textsubring{r}}} \dicPos{m}[5] \dicFlx{(‑urs, ‑rar)}[3] \dicFieldCat{práv.} \dicDirectTranslationCS{plnoletost, zletilost, zákonný věk};  \dicPhraseIS{ná lögaldri} \dicDirectTranslationCS{dosáhnout plnoletosti}
\dicEntry[lögbann] \dicTerm{lög··bann} \dicIPA{{l}{\oe}{\textbabygamma}{\textsubring{b}}{a}{\textsubring{n}}} \dicPos{n}[2] \dicFlx{(‑s)}[2] \dicFieldCat{práv.} \dicDirectTranslationCS{soudní zákaz}
\dicEntry[lögbrot] \dicTerm{lög··brot} \dicIPA{{l}{\oe}{\textbabygamma}{\textsubring{b}}{r}{\textopeno}{\textsubring{d}}} \dicPos{n}[2] \dicFlx{(‑s, ‑)}[5] \dicFieldCat{práv.} \dicDirectTranslationCS{porušení zákona, delikt}
\dicEntry[lögfesta] \dicTerm{lög··fest|a} \dicIPA{{l}{\oe}{x}{f}{\textepsilon}{s}{\textsubring{d}}{a}} \dicPos{v}[2] \dicFlx{(‑i, ‑)}[12] \dicFlx{acc} \dicFieldCat{práv.} \dicDirectTranslationCS{uzákonit, legalizovat, kodifikovat}
\dicEntry[lögformlegur] \dicTerm{lög··form·legur} \dicIPA{{l}{\oe}{x}{f}{\textopeno}{r}{m}{l}{\textepsilon}{\textbabygamma}{\textscy}{\textsubring{r}}} \dicPos{adj}[1]\dicFlx{}[-8] \dicFieldCat{práv.} \dicDirectTranslationCS{zákonný}
\dicEntry[lögfræði] \dicTerm{lög··fræð|i} \dicIPA{{l}{\oe}{x}{f}{r}{a}{i}{ð}{\textsci}} \dicPos{f}[3] \dicFlx{(‑i)}[3] \dicDirectTranslationCS{právní věda, práva}
\dicEntry[lögfræðideild] \dicTerm{lög·fræði··deild} \dicIPA{{l}{\oe}{x}{f}{r}{a}{i}{ð}{\textsci}{\textsubring{d}}{ei}{l}{\textsubring{d}}} \dicPos{f}[7] \dicFlx{(‑ar, ‑ir)}[1] \dicFieldCat{škol.} \dicDirectTranslationCS{právnická fakulta}
\dicEntry[lögfræðilegur] \dicTerm{lög·fræði··legur} \dicIPA{{l}{\oe}{x}{f}{r}{a}{i}{ð}{\textsci}{l}{\textepsilon}{\textbabygamma}{\textscy}{\textsubring{r}}} \dicPos{adj}[1]\dicFlx{}[-8] \dicDirectTranslationCS{právní} \dicExampleIS{lögfræðilegur ráðgjafi} \dicExampleCS{právní poradce}
\dicEntry[lögfræðingur] \dicTerm{lög·fræð··ing|ur} \dicsymFrequent\  \dicIPA{{l}{\oe}{x}{f}{r}{a}{i}{ð}{i}{\ng}{\r{g}}{\textscy}{\textsubring{r}}} \dicPos{m}[6] \dicFlx{(‑s, ‑ar)}[8] \dicDirectTranslationCS{právník, právnička} \dicExampleIS{aðstoð lögfræðings} \dicExampleCS{pomoc právníka}
\dicEntry[lögg] \dicTerm{lögg} \dicIPA{{l}{\oe}{\r{g}}{\textlengthmark}} \dicPos{f}[7] \dicFlx{(laggar\,/\addthin ‑var, laggir\,/\addthin ‑ir\,/\addthin ‑var)}[30] \dicDirectTranslationCS{kapka, doušek (v~láhvi ap.)} \dicExampleIS{kaffilögg} \dicExampleCS{doušek kávy};  \dicPhraseIS{setja e‑ð á laggirnar} \dicDirectTranslationCS{založit (co), zřídit (co)} \dicExampleIS{setja fyrirtæki á laggirnar} \dicExampleCS{zřídit společnost};  \dicPhraseIS{koma e‑u á laggirnar} \dicDirectTranslationCS{založit (co), zřídit (co)}
\dicEntry[lögga] \dicTerm{lögg|a} \dicsymFrequent\  \dicIPA{{l}{\oe}{\r{g}}{\textlengthmark}{a}} \dicPos{f}[1] \dicFlx{(‑u, ‑ur)}[7] \dicLangCat{hovor.} \dicSynonym{lögreglumaður} \dicDirectTranslationCS{polda, policajt(ka)} \dicExampleIS{löggur í kvikmyndum} \dicExampleCS{poldové ve filmech}
\dicEntry[löggilda] \dicTerm{lög··gil|da} \dicIPA{{l}{\oe}{\textbabygamma}{\r{\textObardotlessj}}{\textsci}{l}{\textsubring{d}}{a}} \dicPos{v}[2] \dicFlx{(‑ti, ‑t)}[37] \dicFlx{acc} \dicFieldCat{práv.} \dicDirectTranslationCS{oprávnit, zmocnit, autorizovat}
\dicEntry[löggilding] \dicTerm{lög··gild·ing} \dicIPA{{l}{\oe}{\textbabygamma}{\r{\textObardotlessj}}{\textsci}{l}{\textsubring{d}}{i}{\ng}{\r{g}}} \dicPos{f}[4] \dicFlx{(‑ar, ‑ar)}[5] \dicFieldCat{práv.} \dicDirectTranslationCS{oprávnění, zmocnění, autorizace}
\dicEntry[löggiltur] \dicTerm{lög··giltur} \dicIPA{{l}{\oe}{\textbabygamma}{\r{\textObardotlessj}}{\textsci}{\textsubring{l}}{\textsubring{d}}{\textscy}{\textsubring{r}}} \dicPos{adj}[1]\dicFlx{}[-13] \dicFieldCat{práv.} \dicDirectTranslationCS{oprávněný, autorizovaný};  \dicPhraseIS{löggiltur skjalaþýðandi} \dicDirectTranslationCS{úřední překladatel(ka)}
\dicEntry[löggjafarvald] \dicTerm{lög·gjafar··vald} \dicIPA{{l}{\oe}{\textbabygamma}{\r{\textObardotlessj}}{a}{v}{a}{r}{v}{a}{l}{\textsubring{d}}} \dicPos{n}[2] \dicFlx{(‑s)}[2] \dicDirectTranslationCS{zákonodárná moc}
\dicEntry[löggjafarþing] \dicTerm{lög·gjafar··þing} \dicIPA{{l}{\oe}{\textbabygamma}{\r{\textObardotlessj}}{a}{f}{a}{\textsubring{r}}{\texttheta}{i}{\ng}{\r{g}}} \dicPos{n}[2] \dicFlx{(‑s, ‑)}[5] \dicFieldCat{práv.} \dicDirectTranslationCS{zákonodárný sbor}
\dicEntry[löggjafi] \dicTerm{lög··gjaf|i} \dicIPA{{l}{\oe}{\textbabygamma}{\r{\textObardotlessj}}{a}{v}{\textsci}} \dicPos{m}[1] \dicFlx{(‑a, ‑ar)}[8] \dicFieldCat{pol.} \dicDirectTranslationCS{zákonodárce, zákonodárkyně}
\dicEntry[löggjöf] \dicTerm{lög··|gjöf} \dicIPA{{l}{\oe}{\textbabygamma}{\r{\textObardotlessj}}{\oe}{f}} \dicPos{f}[7] \dicFlx{(‑gjafar)}[19] \dicFieldCat{práv.} \dicDirectTranslationCS{legislativa, zákonodárství}
\dicEntry[löggæsla] \dicTerm{lög··gæsl|a} \dicIPA{{l}{\oe}{\textbabygamma}{\r{\textObardotlessj}}{a}{i}{s}{\textsubring{d}}{l}{a}} \dicPos{f}[1] \dicFlx{(‑u)}[5] \dicDirectTranslationCS{prosazování práva}
\dicEntry[löghald] \dicTerm{lög··hald} \dicIPA{{l}{\oe}{\textlengthmark}{x}{h}{a}{l}{\textsubring{d}}} \dicPos{n}[2] \dicFlx{(‑s)}[2] \dicFieldCat{práv.} \dicDirectTranslationCS{(úřední) zabavení, konfiskace};  \dicPhraseIS{leggja löghald á e‑u} \dicDirectTranslationCS{zkonfiskovat (co)}
\dicEntry[lögheimili] \dicTerm{lög··heimili} \dicIPA{{l}{\oe}{\textlengthmark}{x}{h}{ei}{m}{\textsci}{l}{\textsci}} \dicPos{n}[2] \dicFlx{(‑s, ‑)}[14] \dicFieldCat{práv.} \dicDirectTranslationCS{trvalé bydliště}
\dicEntry[löghlýðinn] \dicTerm{lög··hlýðinn} \dicIPA{{l}{\oe}{x}{\textsubring{l}}{i}{ð}{\textsci}{\textsubring{n}}} \dicPos{adj}[6]\dicFlx{}[-2] \dicDirectTranslationCS{dodržující zákony}
\dicEntry[löghyggja] \dicTerm{lög··hyggj|a} \dicIPA{{l}{\oe}{\textlengthmark}{x}{h}{\textsci}{\r{\textObardotlessj}}{a}} \dicPos{f}[1] \dicFlx{(‑u)}[5] \dicFieldCat{filos.} \dicDirectTranslationCS{determinismus}
\dicEntry[löglegur] \dicTerm{lög··legur} \dicIPA{{l}{\oe}{\textbabygamma}{l}{\textepsilon}{\textbabygamma}{\textscy}{\textsubring{r}}} \dicPos{adj}[1]\dicFlx{}[-8] \dicDirectTranslationCS{legální, právoplatný}
\dicEntry[lögmaður] \dicTerm{lög··|maður} \dicsymFrequent\  \dicIPA{{l}{\oe}{\textbabygamma}{m}{a}{ð}{\textscy}{\textsubring{r}}} \dicPos{m}[13] \dicFlx{(‑manns, ‑menn)}[2] \dicDirectTranslationCS{advokát(ka)} \dicExampleIS{Vantar þig lögmann?} \dicExampleCS{Potřebuješ advokáta?}
\dicEntry[lögmál] \dicTerm{lög··mál} \dicIPA{{l}{\oe}{\textbabygamma}{m}{au}{\textsubring{l}}} \dicPos{n}[2] \dicFlx{(‑s, ‑)}[5] \textbf{1.} \dicSynonym{lög} \dicDirectTranslationCS{zákon}  \textbf{2.} \dicDirectTranslationCS{zákon, řád (přírody ap.)} \dicExampleIS{lögmál lífsins} \dicExampleCS{zákon života}  \textbf{3.} \dicFieldCat{filos.} \dicSynonym{grundvallarregla} \dicDirectTranslationCS{princip, zákonitost}
\dicEntry[lögmæti] \dicTerm{lög··mæti} \dicIPA{{l}{\oe}{\textbabygamma}{m}{a}{i}{\textsubring{d}}{\textsci}} \dicPos{n}[2] \dicFlx{(‑s)}[20] \dicDirectTranslationCS{legálnost, zákonnost}
\dicEntry[lögmætur] \dicTerm{lög··mætur} \dicIPA{{l}{\oe}{\textbabygamma}{m}{a}{i}{\textsubring{d}}{\textscy}{\textsubring{r}}} \dicPos{adj}[1]\dicFlx{}[-1] \dicSynonym{gildur} \dicDirectTranslationCS{legální, zákonný}
\dicEntry[lögn] \dicTerm{lögn} \dicIPA{{l}{\oe}{\r{g}}{\textsubring{n}}} \dicPos{f}[7] \dicFlx{(lagnar, lagnir)}[16] \textbf{1.} \dicSynonym{leiðsla} \dicDirectTranslationCS{vedení, instalace} \dicExampleIS{raflögn} \dicExampleCS{elektrické vedení}  \textbf{2.} \dicSynonym*{það að leggja net} \dicDirectTranslationCS{pokládání sítě}
\dicEntry[lögregla] \dicTerm{lög··regl|a} \dicsymFrequent\  \dicIPA{{l}{\oe}{\textbabygamma}{r}{\textepsilon}{\r{g}}{l}{a}} \dicPos{f}[1] \dicFlx{(‑u, ‑ur)}[19] \textbf{1.} \dicSynonym*{stofnun til löggæslu} \dicDirectTranslationCS{policie} \dicExampleIS{Lögreglan leitar að vitnum að árekstri.} \dicExampleCS{Policie hledá svědky srážky vozidel.}  \textbf{2.} \dicSynonym{lögreglumaður} \dicDirectTranslationCS{policista, policistka}
\dicEntry[lögreglubíll] \dicTerm{lög·reglu··bíl|l} \dicsymFrequent\  \dicIPA{{l}{\oe}{\textbabygamma}{r}{\textepsilon}{\r{g}}{l}{\textscy}{\textsubring{b}}{i}{\textsubring{d}}{\textsubring{l}}} \dicPos{m}[6] \dicFlx{(‑s, ‑ar)}[48] \dicDirectTranslationCS{policejní auto\,/\addthin vůz} \dicExampleIS{Lögreglubílar lentu í árekstri.} \dicExampleCS{Policejní auta se srazila.}
\dicEntry[lögreglumaður] \dicTerm{lög·reglu··|maður} \dicsymFrequent\  \dicIPA{{l}{\oe}{\textbabygamma}{r}{\textepsilon}{\r{g}}{l}{\textscy}{m}{a}{ð}{\textscy}{\textsubring{r}}} \dicPos{m}[13] \dicFlx{(‑manns, ‑menn)}[2] \dicDirectTranslationCS{policista, policistka} \dicExampleIS{Lögreglumenn réðust inn í húsið.} \dicExampleCS{Policisté vtrhli do domu.}
\dicEntry[lögreglusamþykkt] \dicTerm{lög·reglu··sam·þykkt} \dicIPA{{l}{\oe}{\textbabygamma}{r}{\textepsilon}{\r{g}}{l}{\textscy}{s}{a}{m}{\texttheta}{\textsci}{x}{\textsubring{d}}} \dicPos{f}[7] \dicFlx{(‑ar, ‑ir)}[1] \dicDirectTranslationCS{policejní nařízení}
\dicEntry[lögreglustjóri] \dicTerm{lög·reglu··stjór|i} \dicIPA{{l}{\oe}{\textbabygamma}{r}{\textepsilon}{\r{g}}{l}{\textscy}{s}{\textsubring{d}}{j}{ou}{r}{\textsci}} \dicPos{m}[1] \dicFlx{(‑a, ‑ar)}[1] \dicDirectTranslationCS{policejní prezident(ka)}
\dicEntry[lögreglustöð] \dicTerm{lög·reglu··stöð} \dicIPA{{l}{\oe}{\textbabygamma}{r}{\textepsilon}{\r{g}}{l}{\textscy}{s}{\textsubring{d}}{\oe}{\texttheta}} \dicPos{f}[6] \dicFlx{(‑var, ‑var)}[1] \dicDirectTranslationCS{policejní stanice}
\dicEntry[lögregluvarðstjóri] \dicTerm{lög·reglu··varð·stjór|i} \dicIPA{{l}\-{\oe}\-{\textbabygamma}\-{r}\-{\textepsilon}\-{\r{g}}\-{l}\-{\textscy}\-{v}\-{a}\-{r}\-{ð}\-{s}\-{\textsubring{d}}\-{j}\-{ou}\-{r}\-{\textsci}\-} \dicPos{m}[1] \dicFlx{(‑a, ‑ar)}[1] \dicDirectTranslationCS{policejní komisař(ka)}
\dicEntry[lögregluvarðstofa] \dicTerm{lög·reglu··varð·stof|a} \dicIPA{{l}\-{\oe}\-{\textbabygamma}\-{r}\-{\textepsilon}\-{\r{g}}\-{l}\-{\textscy}\-{v}\-{a}\-{r}\-{ð}\-{s}\-{\textsubring{d}}\-{\textopeno}\-{v}\-{a}\-} \dicPos{f}[1] \dicFlx{(‑u, ‑ur)}[7] \dicDirectTranslationCS{okrsková služebna, policejní stanice}
\dicEntry[lögregluþjónn] \dicTerm{lög·reglu··þjón|n} \dicsymFrequent\  \dicIPA{{l}{\oe}{\textbabygamma}{r}{\textepsilon}{\r{g}}{l}{\textscy}{\texttheta}{j}{ou}{\textsubring{d}}{\textsubring{n}}} \dicPos{m}[6] \dicFlx{(‑s, ‑ar)}[42] \dicDirectTranslationCS{policista, policistka} \dicExampleIS{lögregluþjónar á næturvakt} \dicExampleCS{policisté na noční službě}
\dicEntry[Lögrétta] \dicTerm{Lög··rétt|a} \dicIPA{{l}{\oe}{\textbabygamma}{r}{j}{\textepsilon}{h}{\textsubring{d}}{a}} \dicPos{f}[1] \dicFlx{(‑u)}[6] \dicFieldCat{geog., hist.} \dicDirectTranslationCS{Lögrétta} \dicIndirectTranslationCS{(zákonodárný a~soudní sněm nacházející se v~Altingu)}
\dicEntry[lögræði] \dicTerm{lög··ræði} \dicIPA{{l}{\oe}{\textbabygamma}{r}{a}{i}{ð}{\textsci}} \dicPos{n}[2] \dicFlx{(‑s)}[20] \dicFieldCat{práv.} \dicDirectTranslationCS{svéprávnost};  \dicPhraseIS{svipta e‑n lögræði} \dicDirectTranslationCS{zbavit (koho) svéprávnosti}
\dicEntry[lögsótt] \dicTerm{lög··sótt} \dicIPA{{l}{\oe}{x}{s}{ou}{h}{\textsubring{d}}} \dicPos{v} \dicFlx{supin} \dicLink{lögsækja}
\dicEntry[lögsótti] \dicTerm{lög··sótti} \dicIPA{{l}{\oe}{x}{s}{ou}{h}{\textsubring{d}}{\textsci}} \dicPos{v} \dicFlx{ind pf sg 1 pers} \dicLink{lögsækja}
\dicEntry[lögsóttum] \dicTerm{lög··sóttum} \dicIPA{{l}{\oe}{x}{s}{ou}{h}{\textsubring{d}}{\textscy}{\textsubring{m}}} \dicPos{v} \dicFlx{ind pf pl 1 pers} \dicLink{lögsækja}
\dicEntry[lögsæki] \dicTerm{lög··sæki} \dicIPA{{l}{\oe}{x}{s}{a}{i}{\r{\textObardotlessj}}{\textsci}} \dicPos{v} \dicFlx{ind praes sg 1 pers} \dicLink{lögsækja}
\dicEntry[lögsækja] \dicTerm{lög··|sækja} \dicIPA{{l}{\oe}{x}{s}{a}{i}{\r{\textObardotlessj}}{a}} \dicPos{v}[5] \dicFlx{(‑sæki, ‑sótti, ‑sóttum, ‑sækti, ‑sótt)}[7] \dicFlx{acc} \dicFieldCat{práv.} \dicDirectTranslationCS{žalovat, podat\,/\addthin podávat žalobu} \dicExampleIS{lögsækja e‑n fyrir e‑ð} \dicExampleCS{žalovat (koho) za (co)}
\dicEntry[lögsækti] \dicTerm{lög··sækti} \dicIPA{{l}{\oe}{x}{s}{a}{i}{x}{\textsubring{d}}{\textsci}} \dicPos{v} \dicFlx{con pf sg 1 pers} \dicLink{lögsækja}
\dicEntry[lögtak] \dicTerm{lög··|tak} \dicIPA{{l}{\oe}{x}{t\smash{\textsuperscript{h}}}{a}{\r{g}}} \dicPos{n}[2] \dicFlx{(‑taks, ‑tök)}[8] \dicFieldCat{práv.} \dicDirectTranslationCS{exekuce}
\dicEntry[lögun] \dicTerm{lögun} \dicsymFrequent\  \dicIPA{{l}{\oe}{\textlengthmark}{\textbabygamma}{\textscy}{\textsubring{n}}} \dicPos{f}[7] \dicFlx{(‑ar, laganir)}[16] \textbf{1.} \dicSynonym{form} \dicDirectTranslationCS{tvar, podoba} \dicExampleIS{breyta löguninni} \dicExampleCS{změnit podobu}  \textbf{2.} \dicSynonym*{kaffilögun o.þ.h.} \dicDirectTranslationCS{příprava (kávy ap.)} \dicExampleIS{lögun á kaffi} \dicExampleCS{příprava kávy}
\dicEntry[lögur] \dicTerm{lögur} \dicIPA{{l}{\oe}{\textlengthmark}{\textbabygamma}{\textscy}{\textsubring{r}}} \dicPos{m}[11] \dicFlx{(lagar, legir)}[5] \textbf{1.} \dicDirectTranslationCS{kapalina, tekutina}  \textbf{2.} \dicSynonym{vatn} \dicDirectTranslationCS{voda, vodní plocha} \dicExampleIS{á láði og legi} \dicExampleCS{na souši a~na vodě}
\dicEntry[lök] \dicTerm{lök\smash{\textsuperscript{1}}} \dicIPA{{l}{\oe}{\textlengthmark}{\r{g}}} \dicPos{n} \dicFlx{pl nom} \dicLink{lak\smash{\textsuperscript{1}}}
\dicEntry[lök] \dicTerm{lök\smash{\textsuperscript{2}}} \dicIPA{{l}{\oe}{\textlengthmark}{\r{g}}} \dicPos{adj} \dicFlx{f sg nom pos} \dicLink{lakur}
\dicEntry[lökk] \dicTerm{lökk} \dicIPA{{l}{\oe}{h}{\r{g}}} \dicPos{n} \dicFlx{pl nom} \dicLink{lakk}
\dicEntry[lömb] \dicTerm{lömb} \dicIPA{{l}{\oe}{m}{\textsubring{b}}} \dicPos{n} \dicFlx{pl nom} \dicLink{lamb}
\dicEntry[lömdum] \dicTerm{lömdum} \dicIPA{{l}{\oe}{m}{\textsubring{d}}{\textscy}{\textsubring{m}}} \dicPos{v} \dicFlx{ind pf pl 1 pers} \dicLink{lemja}
\dicEntry[lömuð] \dicTerm{lömuð} \dicIPA{{l}{\oe}{\textlengthmark}{m}{\textscy}{\texttheta}} \dicPos{adj} \dicFlx{f sg nom pos} \dicLink{lamaður}
\dicEntry[lömun] \dicTerm{löm|un} \dicIPA{{l}{\oe}{\textlengthmark}{m}{\textscy}{\textsubring{n}}} \dicPos{f}[7] \dicFlx{(‑unar)}[12] \dicFieldCat{med.} \dicDirectTranslationCS{ochrnutí, paralýza} \dicExampleIS{fá lömun í fótinn} \dicExampleCS{ochrnout na nohu}
\dicEntry[lömunarveiki] \dicTerm{lömunar··veik|i} \dicIPA{{l}{\oe}{\textlengthmark}{m}{\textscy}{n}{a}{r}{v}{ei}{\r{\textObardotlessj}}{\textsci}} \dicPos{f}[3] \dicFlx{(‑i)}[3] \dicFieldCat{med.} \dicDirectTranslationCS{obrna}
\dicEntry[lönd] \dicTerm{lönd} \dicIPA{{l}{\oe}{n}{\textsubring{d}}} \dicPos{n} \dicFlx{pl nom} \dicLink{land}
\dicEntry[löndun] \dicTerm{löndun} \dicIPA{{l}{\oe}{n}{\textsubring{d}}{\textscy}{\textsubring{n}}} \dicPos{f}[7] \dicFlx{(‑ar, landanir)}[16] \dicSynonym*{uppskipun} \dicDirectTranslationCS{vykládka, vylodění}
\dicEntry[löng] \dicTerm{löng} \dicIPA{{l}{\oe i}{\ng}{\r{g}}} \dicPos{adj} \dicFlx{f sg nom pos} \dicLink{langur\smash{\textsuperscript{2}}}
\dicEntry[löngu] \dicTerm{löngu\smash{\textsuperscript{1}}} \dicIPA{{l}{\oe i}{\ng}{\r{g}}{\textscy}} \dicPos{f} \dicFlx{sg gen} \dicLink{langa\smash{\textsuperscript{1}}}
\dicEntry[löngu] \dicTerm{löngu\smash{\textsuperscript{2}}} \dicIPA{{l}{\oe i}{\ng}{\r{g}}{\textscy}} \dicPos{adv} \dicDirectTranslationCS{dlouho, dávno};  \dicPhraseIS{endur fyrir löngu} \dicFlx{adv} \dicDirectTranslationCS{dávno, před dávnými dobami};  \dicPhraseIS{fyrir (langa\,/\addthin margt) löngu} \dicFlx{adv} \dicDirectTranslationCS{dávno, před dávnou dobou};  \dicPhraseIS{fyrir lifandi löngu} \dicFlx{adv} \dicDirectTranslationCS{už\,/\addthin již (velmi) dávno};  \dicPhraseIS{löngu áður} \dicFlx{adv} \dicDirectTranslationCS{dlouho předtím};  \dicPhraseIS{löngu fyrr} \dicFlx{adv} \dicDirectTranslationCS{mnohem dříve, dlouho předtím};  \dicPhraseIS{löngu síðar\,/\addthin seinna} \dicFlx{adv} \dicDirectTranslationCS{dlouho poté, o~dost později}
\dicEntry[löngum] \dicTerm{löngum} \dicsymFrequent\  \dicIPA{{l}{\oe i}{\ng}{\r{g}}{\textscy}{\textsubring{m}}} \dicPos{adv} \dicDirectTranslationCS{dlouho} \dicExampleIS{Það hefur löngum verið vitað að sólböð eru hættuleg.} \dicExampleCS{Už se dlouho ví, že opalování je nebezpečné.}
\dicEntry[löngun] \dicTerm{löngun} \dicsymFrequent\  \dicIPA{{l}{\oe i}{\ng}{\r{g}}{\textscy}{\textsubring{n}}} \dicPos{f}[7] \dicFlx{(‑ar, langanir)}[16] \dicDirectTranslationCS{touha, toužení, tužba} \dicExampleIS{löngun í nýjan bíl} \dicExampleCS{touha po novém autě}
\dicEntry[löngutangar] \dicTerm{löngu··tangar} \dicIPA{{l}{\oe i}{\ng}{\r{g}}{\textscy}{t\smash{\textsuperscript{h}}}{au}{\ng}{\r{g}}{a}{\textsubring{r}}} \dicPos{f} \dicFlx{sg gen} \dicLink{langatöng}
\dicEntry[löngutengur] \dicTerm{löngu··tengur} \dicIPA{{l}{\oe i}{\ng}{\r{g}}{\textscy}{t\smash{\textsuperscript{h}}}{\textepsilon}{i}{\ng}{\r{g}}{\textscy}{\textsubring{r}}} \dicPos{f} \dicFlx{pl nom} \dicLink{langatöng}
\dicEntry[löngutöng] \dicTerm{löngu··|töng} \dicIPA{{l}{\oe i}{\ng}{\r{g}}{\textscy}{t\smash{\textsuperscript{h}}}{\oe i}{\ng}{\r{g}}} \dicPos{f}[8] \dicFlx{(‑tangar, ‑tengur\,/\addthin ‑tangir)}[9] \dicLink{langatöng}
\dicEntry[löpp] \dicTerm{löpp} \dicsymFrequent\  \dicIPA{{l}{\oe}{h}{\textsubring{b}}} \dicPos{f}[7] \dicFlx{(lappar, lappir)}[16] \dicDirectTranslationCS{noha, chodidlo} \dicExampleIS{vöðvar í löppum} \dicExampleCS{svaly na nohou};  \dicPhraseIS{koma sér á lappir} \dicDirectTranslationCS{vstát z~postele}
\dicEntry[löptum] \dicTerm{löptum} \dicIPA{{l}{\oe}{f}{\textsubring{d}}{\textscy}{\textsubring{m}}} \dicPos{v} \dicFlx{ind pf pl 1 pers} \dicLink{lepja}
\dicEntry[löstur] \dicTerm{löstur} \dicIPA{{l}{\oe}{s}{\textsubring{d}}{\textscy}{\textsubring{r}}} \dicPos{m}[11] \dicFlx{(lastar, lestir)}[5] \dicDirectTranslationCS{zlozvyk, neřest} \dicExampleIS{drykkjuskapur og aðrir lestir} \dicExampleCS{opilství a~jiné neřesti}
\dicEntry[löt] \dicTerm{löt} \dicIPA{{l}{\oe}{\textlengthmark}{\textsubring{d}}} \dicPos{adj} \dicFlx{f sg nom pos} \dicLink{latur}
\dicEntry[lötra] \dicTerm{lötr|a} \dicIPA{{l}{\oe}{\textlengthmark}{\textsubring{d}}{r}{a}} \dicPos{v}[1] \dicFlx{(‑aði)}[44] \dicDirectTranslationCS{loudat se, šinout se, courat se} \dicExampleIS{Hann lötraði á eftir hópnum.} \dicExampleCS{Coural se za skupinou.}
\dicEntry[löttum] \dicTerm{löttum} \dicIPA{{l}{\oe}{h}{\textsubring{d}}{\textscy}{\textsubring{m}}} \dicPos{v} \dicFlx{ind pf pl 1 pers} \dicLink{letja}
}
\fi
\restoregeometry
\pagestyle{plain}

\fi

% ================================================================ PHONETICS =

\ifinputphon

\cleardoublepage

\chapter{Seznam islandských fonémů}              \label{sec:phon_phonems}

\begin{table}[h]
\begin{tabular}{lllll} \toprule

\textbf{foném} & \textbf{vysvětlivka} & \textbf{příklad} & \textbf{fonetický přepis} & \textbf{překlad} \\
\midrule
{\textipa{[{a}]}} & otevřená přední nezaokrouhlená samohl. & rata & {\textipa{[{r}{a}{\textlengthmark}{\textsubring{d}}{a}]}} & trefit \\ 
{\textipa{[{\textepsilon}]}} & polootevřená přední nezaokrouhlená samohl. & meta & {\textipa{[{m}{\textepsilon}{\textlengthmark}{\textsubring{d}}{a}]}} & ocenit \\ 
{\textipa{[{e}]}} & polozavřená přední nezaokrouhlená samohl. & heiður & {\textipa{[{h}{ei}{\textlengthmark}{ð}{\textscy}{\textsubring{r}}]}} & čest \\
{\textipa{[{\textsci}]}} & téměř zavřená téměř přední nezaokrouhlená samohl. & vinna & {\textipa{[{v}{\textsci}{n}{\textlengthmark}{a}]}} & pracovat \\ 
{\textipa{[{i}]}} & zavřená přední nezaokrouhlená samohl. & fínn & {\textipa{[{f}{i}{\textsubring{d}}{\textsubring{n}}]}} & jemný \\ 
{\textipa{[{\textopeno}]}} & polootevřená zadní zaokrouhlená samohl. & lofa & {\textipa{[{l}{\textopeno}{\textlengthmark}{v}{a}]}} & slibovat \\ 
{\textipa{[{o}]}} & polozavřená zadní zaokrouhlená samohl. & bóndi & {\textipa{[{\textsubring{b}}{ou}{n}{\textsubring{d}}{\textsci}]}} & farmář \\ 
{\textipa{[{\textscy}]}} & téměř zavřená téměř přední zaokrouhlená samohl. & hundur & {\textipa{[{h}{\textscy}{n}{\textsubring{d}}{\textscy}{\textsubring{r}}]}} & pes \\ 
{\textipa{[{u}]}} & zavřená zadní zaokrouhlená samohl. & súkkulaði & {\textipa{[{s}{u}{h}{\r{g}}{\textscy}{l}{a}{ð}{\textsci}]}} & čokoláda \\ 
{\textipa{[{\oe}]}} & téměř otevřená přední nezaokrouhlená samohl. & lönd & {\textipa{[{l}{\oe}{n}{\textsubring{d}}]}} & země \\ 
{\textipa{[{m}]}} & bilabiální nazála & mýri & {\textipa{[{m}{i}{\textlengthmark}{r}{\textsci}]}} & mokřina \\ 
{\textipa{[{\textsubring{m}}]}} & neznělá bilabiální nazála & þyrmt & {\textipa{[{\texttheta}{\textsci}{\textsubring{r}}{\textsubring{m}}{\textsubring{d}}]}} & {\textit{$\shortrightarrow$ þyrma}} \\  
{\textipa{[{n}]}} & alveolární nazála & nemandi & {\textipa{[{n}{\textepsilon}{\textlengthmark}{m}{a}{n}{\textsubring{d}}{\textsci}]}} & žák \\ 
{\textipa{[{\textsubring{n}}]}} & neznělá alveolární nazála & mennt & {\textipa{[{m}{\textepsilon}{\textsubring{n}}{\textsubring{d}}]}} & vzdělání \\ 
{\textipa{[{\textltailn}]}} & palatální nazála & angi & {\textipa{[{au}{\textltailn}{\r{\textObardotlessj}}{\textsci}]}} & drobek \\ 
{\textipa{[{\r{\textltailn}}]}} & neznělá palatální nazála & banki & {\textipa{[{\textsubring{b}}{au}{\r{\textltailn}}{\r{\textObardotlessj}}{\textsci}]}} & banka \\
{\textipa{[{\ng}]}} & velární nazála & skyggndur & {\textipa{[{s}{\r{\textObardotlessj}}{\textsci}{\ng}{\textsubring{d}}{\textscy}{\textsubring{r}}]}} & {\textit{$\shortrightarrow$ skyggna}} \\ 
{\textipa{[{\r{\ng}}]}} & neznělá velární nazála & skyggnt & {\textipa{[{s}{\r{\textObardotlessj}}{\textsci}{\r{\ng}}{\textsubring{d}}]}} & {\textit{$\shortrightarrow$ skyggna}} \\ 
{\textipa{[{\textsubring{b}}]}} & neznělá bilabiální ploziva & bara & {\textipa{[{\textsubring{b}}{a}{\textlengthmark}{r}{a}]}} & jenom \\ 
{\textipa{[{p\textsuperscript{h}}]}} & neznělá bilabiální ploziva s přídechem & prestur & {\textipa{[{p\textsuperscript{h}}{r}{\textepsilon}{s}{\textsubring{d}}{\textscy}{\textsubring{r}}]}} & kněz \\ 
{\textipa{[{\textsubring{d}}]}} & neznělá alveolární ploziva & dæmi & {\textipa{[{\textsubring{d}}{a}{i}{\textlengthmark}{m}{\textsci}]}} & příklad \\ 
{\textipa{[{t\textsuperscript{h}}]}} & neznělá alveolární ploziva s přídechem & tómur & {\textipa{[{t\textsuperscript{h}}{ou}{\textlengthmark}{m}{\textscy}{\textsubring{r}}]}} & prázdný \\ 
{\textipa{[{\r{\textObardotlessj}}]}} & neznělá palatální ploziva & gefa & {\textipa{[{\r{\textObardotlessj}}{\textepsilon}{\textlengthmark}{v}{a}]}} & dát \\ 
{\textipa{[{c\textsuperscript{h}}]}} & neznělá palatální ploziva s přídechem & kenna & {\textipa{[{c\textsuperscript{h}}{\textepsilon}{n}{\textlengthmark}{a}]}} & učit \\
{\textipa{[{\r{g}}]}} & neznělá velární ploziva & gata & {\textipa{[{\r{g}}{a}{\textlengthmark}{\textsubring{d}}{a}]}} & ulice \\ 
{\textipa{[{k\textsuperscript{h}}]}} & neznělá velární ploziva s přídechem & koma & {\textipa{[{k\textsuperscript{h}}{\textopeno}{\textlengthmark}{m}{a}]}} & přijít \\
{\textipa{[{v}]}} & znělá labiodentální frikativa & hafa & {\textipa{[{h}{a}{\textlengthmark}{v}{a}]}} & mít \\ 
{\textipa{[{f}]}} & neznělá labiodentální frikativa & finna & {\textipa{[{f}{\textsci}{n}{\textlengthmark}{a}]}} & najít \\ 
{\textipa{[{ð}]}} & znělá dentální frikativa & borða & {\textipa{[{\textsubring{b}}{\textopeno}{r}{ð}{a}]}} & jíst \\ 
{\textipa{[{\texttheta}]}} & neznělá dentální frikativa & þola & {\textipa{[{\texttheta}{\textopeno}{\textlengthmark}{l}{a}]}} & snést \\
{\textipa{[{s}]}} & neznělá alveolární frikativa & systir & {\textipa{[{s}{\textsci}{s}{\textsubring{d}}{\textsci}{\textsubring{r}}]}} & sestra \\ 
{\textipa{[{\c{c}}]}} & neznělá palatální frikativa & hjóla & {\textipa{[{\c{c}}{ou}{\textlengthmark}{l}{a}]}} & jet na kole \\ 
{\textipa{[{\textbabygamma}]}} & znělá velární frikativa & sagði & {\textipa{[{s}{a}{\textbabygamma}{ð}{\textsci}]}} & řekl \\ 
{\textipa{[{x}]}} & neznělá velární frikativa & lag & {\textipa{[{l}{a}{\textlengthmark}{x}]}} & píseň \\
{\textipa{[{h}]}} & neznělá glotální frikativa & heima & {\textipa{[{h}{ei}{\textlengthmark}{m}{a}]}} & doma \\ 
{\textipa{[{l}]}} & alveolární laterální aproximanta & loforð & {\textipa{[{l}{\textopeno}{\textlengthmark}{v}{\textopeno}{r}{\texttheta}]}} & slib \\ 
{\textipa{[{\textsubring{l}}]}} & neznělá alveolární laterální aproximanta & hlutur & {\textipa{[{\textsubring{l}}{\textscy}{\textlengthmark}{\textsubring{d}}{\textscy}{\textsubring{r}}]}} & věc \\  
{\textipa{[{j}]}} & palatální aproximanta & jól & {\textipa{[{j}{ou}{\textlengthmark}{\textsubring{l}}]}} & Vánoce \\ 
{\textipa{[{r}]}} & alveolární vibranta & reynsla & {\textipa{[{r}{ei}{n}{s}{\textsubring{d}}{l}{a}]}} & zkušenost \\ 
{\textipa{[{\textsubring{r}}]}} & neznělá alveolární vibranta & hraun & {\textipa{[{\textsubring{r}}{\oe i}{\textlengthmark}{\textsubring{n}}]}} & láva \\ 
\bottomrule
\end{tabular}
\end{table}

\clearpage

%\section{Výslovnost islandských samohlásek}      \label{sec:phon_vowels}
%\LTXtable{\columnwidth}{phon/vowels}
%\clearpage

%\section{Výslovnost islandských souhlásek}       \label{sec:phon_consonants}
%\LTXtable{\columnwidth}{phon/consonants}
%\clearpage
\fi

% =============================================================== MORPHOLOGY =

\ifinputmorpho

\cleardoublepage

\chapter{Morfologický klíč}                      \label{sec:morpho}

\section{Podstatné jméno, rod mužský}            \label{sec:morpho_m}
{\small\LTXfw{morpho/m}}
\clearpage

\section{Podstatné jméno, rod ženský}            \label{sec:morpho_f}
{\small\LTXfw{morpho/f}}
\clearpage

\section{Podstatné jméno, rod střední}           \label{sec:morpho_n}
{\small\LTXfw{morpho/n}}
\clearpage

\section{Přídavné jméno}                         \label{sec:morpho_adj}
{\small\LTXfw{morpho/adj}}
\clearpage

\section{Zájmeno}                                \label{sec:morpho_pron}
{\small\LTXfw{morpho/pron}}
 \begin{center}
\begin{minipage}[t]{.45\textwidth}
\begin{tabular}{l>{\footnotesize\itshape}lll}
 {\textbf{\textit{třída}}} & {\textit{pád}}   & \textit{sg} & \textit{pl}  \\ 
 \hline
\multirow{3}{*}{{{\textbf{pron} \Large{\textbf{17}}}}}  &  nom & ég & við   \\*
 & acc &  mig  & okkur  \\*
 & dat & mér & okkur   \\*
 & gen & mín  & okkar  \\
\hline

\multirow{3}{*}{{{\textbf{pron} \Large{\textbf{17}}}}}  &  nom & hann & þeir   \\*
 & acc &  hann  & þá  \\*
 & dat & honum & þeim   \\*
 & gen & hans  & þeirra  \\
\hline

\multirow{3}{*}{{{\textbf{pron} \Large{\textbf{17}}}}}  &  nom & hún & þær   \\*
 & acc &  hana  & þær  \\*
 & dat & henni & þeim   \\*
 & gen & hennar  & þeirra  \\
\hline

\multirow{3}{*}{{{\textbf{pron} \Large{\textbf{17}}}}}  &  nom &  & vér   \\*
 & acc &    & oss  \\*
 & dat &  & oss   \\*
 & gen &   & vor  \\
\hline

\multirow{3}{*}{{{\textbf{pron} \Large{\textbf{17}}}}}  &  nom &  & þér   \\*
 & acc &    & yður  \\*
 & dat &  & yður   \\*
 & gen &   & yðar  \\
\hline

\multirow{3}{*}{{{\textbf{pron} \Large{\textbf{17}}}}}  &  nom & þú & þið   \\*
 & acc &  þig  & ykkur  \\*
 & dat & þér & ykkur   \\*
 & gen & þín  & ykkar  \\
\hline

\multirow{3}{*}{{{\textbf{pron} \Large{\textbf{17}}}}}  &  nom & það & þau   \\*
 & acc &  það  & þau  \\*
 & dat & því & þeim   \\*
 & gen & þess  & þeirra  \\
\hline
\end{tabular}
\end{minipage}\hfil
\begin{minipage}[t]{.45\textwidth}
\begin{tabular}{l>{\footnotesize\itshape}lll}
\hline
\multirow{3}{*}{{{\textbf{pron} \Large{\textbf{18}}}}}  &  nom &  &    \\*
 & acc &  sig  &   \\*
 & dat & sér &    \\*
 & gen & sín  &   \\
 \hline
\end{tabular}
\end{minipage}
\end{center}

\clearpage

\section{Číslovka}                               \label{sec:morpho_num}
 {\small\LTXfw{morpho/num}}
\clearpage

\begin{landscape}
\section{Sloveso}                                \label{sec:morpho_v}
{\small\LTXfw{morpho/v}}
\end{landscape}

\clearpage

\fi

% ============================================================== BACK MATTER =

\ifmakebackmatter

% ------------------------------------------------ Photographs authors index -

\cleardoublepage
%\pagestyle{empty}
%\pagestyle{indexstyle}
%\phantomsection
%\addcontentsline{toc}{chapter}
%  {Seznam autorů fotografií a ilustrací}        
\label{sec:photo}
\printindex[figures]

% ------------------------------------------------------------- Bibliography -
\cleardoublepage
\chapter{Bibliografie}
%\twocolumn
%\phantomsection
%\addcontentsline{toc}{chapter}{Bibliografie}
\begin{multicols}{2}
\nocite{*}\printbibliography[heading=none]
\end{multicols}

% ------------------------------------------------------------------ Apendix -
\cleardoublepage
\onecolumn
\ifPDF
\else
\chapter{Apendix: Islandská gramatika}
\clearpage
\covergeometry  
\makeatletter
\includepdf[
  pages={1,2,3,4,5,6,7,8,9,10,11,12,13,14,15,16,17,18,19,20,21,22,23,24,25,26,27,28,29,30,31,32,33,34,35,36,37,38,39,40,41,42,43,44,45,46,47,48,49,50,51,52,53,54,55,56,57,58,59,60,61,62,63,64,65,66,67,68,69,70,71,72,73},
  turn=false,
  scale=0.79,
  keepaspectratio,
  offset={\dimexpr(\paperwidth-\csname Gm@layoutwidth\endcsname)/2-\csname Gm@layouthoffset\endcsname\relax}
  {\dimexpr(\paperheight-\csname  Gm@layoutheight\endcsname)/2-\csname Gm@layoutvoffset\endcsname\relax}
   ]{gram.pdf}
\makeatother
\restoregeometry  
\fi

% ------------------------------------------------------------------ Barcode -
\pagenumbering{gobble}
\ifPDF
\else
\cleardoublepage\null\clearpage % clear to verso
\vspace*\fill
\EANisbn[SC3]
\fi
\fi

% --------------------------------------------------------------- Back cover -

\ifmakecovers
  \cleardoublepage\null\clearpage % inner side of the cover
\covergeometry  
  \makecoverwith\backcoverimages
\restoregeometry  
\fi

\end{document}

